\documentclass[lang=cn,11pt,scheme=chinese,math=cm,thmcnt=section,usesamecnt]{elegantbook}
\usepackage{unicode-math}
\setmathfont{NewCMMath-Regular.otf}
\renewcommand{\Box}{\mdlgwhtsquare}
\renewcommand{\dashrightarrow}{\mathrel{\rightdasharrow}}
\usepackage{longtable}
\usepackage{float}
\usepackage{tikz-cd}
\usepackage{imakeidx}
\makeindex[intoc,title=索引]
\renewcommand{\bibname}{参考文献}
\elegantnewtheorem{physicsconjecture}{物理学猜想}{thmstyle}{physconj}
\elegantnewtheorem{physicsdefinition}{物理学定义}{defstyle}{physdef}
\elegantnewtheorem{conjecture}{猜想}{thmstyle}{conj}
% 原书的例题、练习和问题按所在小节编号；保留 ElegantBook 原生环境。
\counterwithin{exam}{section}
\counterwithin{exer}{section}
\counterwithin{prob}{section}
\renewcommand{\theHexam}{\thesection.\arabic{exam}}
\renewcommand{\theHexer}{\thesection.\arabic{exer}}
\renewcommand{\theHprob}{\thesection.\arabic{prob}}

% 保留原书例题采用的节或小节前缀，仍使用原生 example 环境。
\newcommand{\SourceExampleNumber}[2]{%
  \renewcommand{\theexam}{#1.\arabic{exam}}%
  \renewcommand{\theHexam}{#1.\arabic{exam}}%
  \setcounter{exam}{\numexpr#2-1\relax}}
\title{镜像对称与代数几何}
\subtitle{Mirror Symmetry and Algebraic Geometry}
\author{David A. Cox \quad Sheldon Katz}
\date{}
\cover{cover.jpg}
\hypersetup{pdftitle={镜像对称与代数几何},pdfauthor={David A. Cox; Sheldon Katz},bookmarksdepth=1}
\setcounter{tocdepth}{1}
\setlength{\emergencystretch}{2em}

% ElegantBook 先加载 hyperref 再加载 titlesec；显式启用 hyperref 的
% titlesec 兼容处理，把目标随标题安放，避免分页时落在上一页。
\makeatletter
\def\ttl@Hy@steplink#1{%
  \Hy@MakeCurrentHrefAuto{#1*}%
  \edef\ttl@Hy@saveanchor{%
    \noexpand\Hy@raisedlink{%
      \noexpand\hyper@anchorstart{\@currentHref}%
      \noexpand\hyper@anchorend
      \def\noexpand\ttl@Hy@SavedCurrentHref{\@currentHref}%
      \noexpand\ttl@Hy@PatchSaveWrite}}}
\def\ttl@Hy@PatchSaveWrite{%
  \begingroup
    \toks@\expandafter{\ttl@savewrite}%
    \edef\x{\endgroup
      \def\noexpand\ttl@savewrite{%
        \let\noexpand\@currentHref\noexpand\ttl@Hy@SavedCurrentHref
        \the\toks@}}%
    \x}
\def\ttl@Hy@refstepcounter#1{%
  \let\ttl@b\Hy@raisedlink
  \def\Hy@raisedlink##1{%
    \def\ttl@Hy@saveanchor{\Hy@raisedlink{##1}}}%
  \refstepcounter{#1}%
  \let\Hy@raisedlink\ttl@b}
\makeatother

% 源页锚点仅用于校对定位；原书 PDF 物理页码从 1 起算。
\makeatletter
\newcommand{\SourcePage}[1]{\hypertarget{source:#1}{}%
  \begingroup
  \def\@currentlabel{}\def\@currentlabelname{}\def\@currentHref{source:#1}%
  \let\@bsphack\relax\let\@esphack\relax
  \label{orig:#1}\endgroup}
\newcommand{\IndexLocation}[1]{\hypertarget{#1}{}%
  \begingroup\def\@currentlabel{}\def\@currentlabelname{}\def\@currentHref{#1}%
  \let\@bsphack\relax\let\@esphack\relax\label{#1}\endgroup}
\newcommand{\IndexSubSubSubItem}{\@idxitem\hspace*{40\p@}}
\makeatother
% 原索引页码由源页标签换算为译本页码；同页去重，相邻页合并为范围。
\ExplSyntaxOn
\seq_new:N \l__mirror_index_pages_seq
\int_new:N \l__mirror_index_start_int
\int_new:N \l__mirror_index_end_int
\bool_new:N \l__mirror_index_first_bool
\cs_new_protected:Npn \__mirror_index_output_range:
 {
  \bool_if:NF \l__mirror_index_first_bool {,\allowbreak\space}
  \bool_set_false:N \l__mirror_index_first_bool
  \hyperlink{page.\int_use:N\l__mirror_index_start_int}
   {\int_use:N\l__mirror_index_start_int}
  \int_compare:nNnF {\l__mirror_index_start_int}={\l__mirror_index_end_int}
   {--\hyperlink{page.\int_use:N\l__mirror_index_end_int}
      {\int_use:N\l__mirror_index_end_int}}
 }
\NewDocumentCommand{\IndexPageList}{o m}
 {
  \group_begin:
  \seq_clear:N \l__mirror_index_pages_seq
  \clist_map_inline:nn {#2}
   {
    \int_set:Nn \l_tmpa_int {\getpagerefnumber{orig:##1}}
    \IfValueT{#1}
     {\cs_if_exist:cT {r@indexloc:#1:##1}
       {\int_set:Nn \l_tmpa_int {\getpagerefnumber{indexloc:#1:##1}}}}
    \int_compare:nNnT {\l_tmpa_int}>{0}
     {\seq_put_right:Nx \l__mirror_index_pages_seq {\int_use:N\l_tmpa_int}}
   }
  \seq_remove_duplicates:N \l__mirror_index_pages_seq
  \seq_sort:Nn \l__mirror_index_pages_seq
   {\int_compare:nNnTF {##1}>{##2}{\sort_return_swapped:}{\sort_return_same:}}
  \int_zero:N \l__mirror_index_start_int
  \int_zero:N \l__mirror_index_end_int
  \bool_set_true:N \l__mirror_index_first_bool
  \seq_map_inline:Nn \l__mirror_index_pages_seq
   {
    \int_compare:nNnTF {\l__mirror_index_start_int}={0}
     {\int_set:Nn\l__mirror_index_start_int{##1}\int_set:Nn\l__mirror_index_end_int{##1}}
     {
      \int_compare:nNnTF {##1}={\l__mirror_index_end_int+1}
       {\int_set:Nn\l__mirror_index_end_int{##1}}
       {\__mirror_index_output_range:\int_set:Nn\l__mirror_index_start_int{##1}\int_set:Nn\l__mirror_index_end_int{##1}}
     }
   }
  \int_compare:nNnTF {\l__mirror_index_start_int}>{0}
   {\__mirror_index_output_range:}{\textbf{??}}
  \group_end:
 }
\ExplSyntaxOff
% 原书的字母式书目标签逐项链接到完整参考文献条目。
\ExplSyntaxOn
\bool_new:N \l__mirror_cite_first_bool
\NewDocumentCommand{\SourceCite}{o m}
 {
  \group_begin:
  \upshape [
  \bool_set_true:N \l__mirror_cite_first_bool
  \clist_map_inline:nn {#2}
   {
    \bool_if:NF \l__mirror_cite_first_bool {,\allowbreak\space}
    \bool_set_false:N \l__mirror_cite_first_bool
    \hyperlink{bib:##1}
     {\str_case:nnF {##1}{{Huebsch}{Hübsch}{GoP}{GöP}{Tjotta}{Tjøtta}{BjoD}{BD}{KaS}{KS}{Gottsche}{Göttsche}{Kollar}{Kollár}}{##1}}
   }
  \IfValueT{#1}{,\nobreakspace #1}
  ]\group_end:
 }
\ExplSyntaxOff

\begin{document}
\maketitle
\hypersetup{pageanchor=true}
\frontmatter
\tableofcontents
\chapter*{前言}
\markboth{前言}{前言}
\phantomsection\addcontentsline{toc}{chapter}{前言}
\SourcePage{14}

近年来，镜像对称作为一个研究领域迅速进入数学界的视野。这是量子场论与数学诸多分支之间联系日益紧密的一个表现。\index{镜像对称|hyperpage}

有人说过，量子场论把二十世纪的物理学与二十一世纪的数学结合在了一起。在尚未得到数学严格论证的情形下进行数学运算，物理学家已经积累了丰富的经验。他们之所以能这样做，部分原因在于实验可以帮助他们区分哪些运算行得通，哪些显然是错的。凡是经受住所有已知检验的运算，在出现相反证据以前，就被认为是有效的。

基于这些证据，物理学家确信镜像对称是成立的。他们运用得极其娴熟的一种工具是 Feynman 路径积分：它要在无限维空间上，借助复测度进行积分，例如在从 Riemann 曲面到 Calabi--Yau 三维簇的所有 $C^\infty$ 映射构成的空间上积分。这并不是严格的数学；然而，正是这些方法促成了 Candelas、de la Ossa、Green 和 Parkes 于 1991 年发表的论文 \SourceCite{CdGP}。该论文对五次三维簇上的有理曲线作出了一些惊人的预测，其内容远远超出了当时代数几何所能证明的范围。\index{五次三维簇|hyperpage}\index{Feynman路径积分@Feynman 路径积分|hyperpage}

数学家面临的挑战是弄清其中究竟发生了什么，更重要的是证明物理学家作出的一些预测。本书将表明，代数几何学家已经取得了实质性的进展，尽管前面仍有很长的路要走。为镜像对称的某些方面建立数学基础的过程，推动了代数几何中新领域的发展。例如，量子上同调、Kontsevich 对稳定映射的定义、Calabi--Yau 三维簇的复化 Kähler 模空间、Batyrev 关于某些环面簇的对偶性，以及 Givental 的量子微分方程概念，都是由此而来的研究方向。镜像对称还推动了形变理论的发展，进而产生了虚拟基本类理论；它也揭示了代数形变理论与辛形变理论之间一种此前未知的联系。虽然我们至今仍不知道镜像对称究竟“是什么”，但它对 Gromov--Witten 不变量作出的预测，如今在许多情形下已经能够得到数学证明。\index{量子上同调|hyperpage}\index{稳定映射|hyperpage}\index{虚拟基本类|hyperpage}\index{Gromov--Witten不变量@Gromov--Witten 不变量|hyperpage}

\section*{本书的目标}
\phantomsection\addcontentsline{toc}{section}{本书的目标}
对于想学习镜像对称的数学家来说，最大的障碍或许是无从着手。目前，许多参考资料分散在各类期刊中，而不少数学思想只见于物理学文献；这些文献对数学家来说很难读懂。本书的首要目标，是介绍镜像对称的代数几何方面。\SourcePage{15}我们会把主要思想和定义阐述清楚，提供充分的细节；篇幅不容许展开更多细节之处，则明确指出可供查阅的文献。我们既解释严格的数学，也说明那些借自物理学、尚未成为定理的直观想法。这样安排，是因为我们主要面向两类读者：希望学习镜像对称的数学家，以及已经了解镜像对称、想进一步认识其数学方面的物理学家。

镜像对称与数学的几个分支都有联系；此外，不同维数的物理理论与数学中许多领域之间，还存在范围更广的联系。本书着重讨论代数几何与镜像对称之间的联系，但也会涉及辛几何等关系密切的领域。我们希望，通过这样的取舍，能够对这一主题作出相当自足的介绍。

本书\hyperref[chap:introduction]{第 1 章}先从总体上介绍镜像对称的思想。随后，\hyperref[chap:quintic]{第 2 章}讨论五次三维簇，并说明镜像对称如何导出 \SourceCite{CdGP} 中的枚举预测。\hyperref[chap:toric]{第 3 章}回顾环面几何，\hyperref[chap:mirror-constructions]{第 4 章}介绍 Batyrev、Batyrev--Borisov 和 Voisin--Borcea 的镜像构造。接下来的四章，即第 \hyperref[chap:hodge-yukawa]{\ref*{chap:hodge-yukawa}}、\hyperref[chap:moduli]{\ref*{chap:moduli}}、\hyperref[chap:gromov-witten]{\ref*{chap:gromov-witten}}、\hyperref[chap:quantum-cohomology]{\ref*{chap:quantum-cohomology}} 章，将建立精确表述镜像对称所需的数学基础，内容涵盖最大幂幺单值性、Yukawa 耦合、复模与 Kähler 模、镜像映射、Gromov--Witten 不变量，以及量子上同调。在此基础上，我们将在\hyperref[chap:quantum-cohomology]{第~\ref*{chap:quantum-cohomology} 章}末提出若干镜像猜想。\index{环面几何|hyperpage}\index{镜像映射|hyperpage}

之后的三章，即第 \hyperref[chap:localization]{\ref*{chap:localization}}、\hyperref[chap:quantum-differential-equations]{\ref*{chap:quantum-differential-equations}}、\hyperref[chap:mirror-theorem]{\ref*{chap:mirror-theorem}} 章，专门证明镜像对称的若干具体情形。等变上同调和局部化在这些证明中起着关键作用，因此\hyperref[chap:localization]{第~\ref*{chap:localization} 章}将回顾相关内容。这些方法还提供了计算 Gromov--Witten 不变量的有力工具。为了解释 Givental 证明镜像定理的途径，我们需要\hyperref[chap:quantum-differential-equations]{第~\ref*{chap:quantum-differential-equations} 章}讨论的引力关联数与量子微分方程。最后，\hyperref[chap:mirror-theorem]{第~\ref*{chap:mirror-theorem} 章}介绍 Lian、Liu 和 Yau \SourceCite{LLY} 以及 Givental \SourceCite{Givental2,Givental4} 关于镜像定理的工作。

第 \hyperref[chap:introduction]{\ref*{chap:introduction}}--\hyperref[chap:mirror-theorem]{\ref*{chap:mirror-theorem}} 章所讨论的数学十分精彩，但也相当艰深。稍后，我们会就如何阅读这些章节提出一些建议。

全书以\hyperref[chap:conclusion]{第~\ref*{chap:conclusion} 章}收尾。这一章汇集前文提到的所有开放问题，并讨论书中未涉及的镜像对称的若干其他方面。最后还有关于奇异簇和物理理论的附录。

我们曾努力使参考文献尽可能完备，但关于镜像对称及相关课题的高质量论文层出不穷，要及时跟上实在困难。对于参考文献中未能列入的许多近期论文，我们向同行致歉。

\section*{与物理学的关系}
\phantomsection\addcontentsline{toc}{section}{与物理学的关系}
对数学家而言，研究镜像对称的一项挫折，是很难真正理解物理学家的直观想法。这些直观想法的力量毋庸置疑，因为镜像现象正是由此被发现的。但是，要理解它们，就必须深入学习量子场论。只要浏览一下讨论相关物理理论的\hyperref[app:physical-theories]{附录 B}，就能看出这项任务有多艰巨。

\SourcePage{16}
要读懂关于镜像对称的物理学文献，需要广泛的背景知识，而本书所提供的知识远远不够。因此，\hyperref[app:physical-theories]{附录 B} 的目标要谦逊得多：仅向读者介绍物理学文献中与镜像对称有关的若干主题。

虽然本书着眼于镜像对称的数学，我们也希望说明这些数学如何体现物理学的精神。基于这一考虑，\hyperref[chap:introduction]{第 1 章}将从促成镜像对称的物理学谈起。我们会经常使用物理学术语，但不假定读者学过量子场论。我们的用意，是让读者感受到：从某些共形场论的观点看，镜像对称完全是自然的。本书中这一章的“物理味”最浓。后面的章节将把注意力集中在数学上，不过我们也会不时停下来，谈谈数学与物理学之间的关系。

物理学所起作用的一个重要方面是：有深厚数学修养的物理学家参与发现了镜像对称的数学基础。代数几何学家有理由为自己创立的、能够解释镜像对称某些方面的优美理论感到自豪；同时，我们也应当承认，物理学家所提供的远不止预测。他们往往还指出了应研究哪些对象，在有些情形下甚至给出了严格的数学描述。查阅书中的参考文献就能明白这一点：即使在本书纯数学的部分，所引文献中也有多得令人意外的物理学论文。我们对物理学同行的贡献负有毋庸置疑的学术之债。

\section*{如何阅读本书}
\phantomsection\addcontentsline{toc}{section}{如何阅读本书}
镜像对称是一个精彩的故事，但讲述它必须涉及代数几何诸多领域的大量细节。读者很容易迷失其间，尤其是试图从头到尾逐页通读本书时。好在这并非唯一的读法。

我们的基本建议是先读第 \hyperref[chap:introduction]{\ref*{chap:introduction}}、\hyperref[chap:quintic]{\ref*{chap:quintic}} 章。如前所述，\hyperref[chap:introduction]{第 1 章}解释一些物理背景，还引入两个关键思想：Calabi--Yau 流形 $V$ 的 A 模型，它编码了我们所关注的枚举信息；以及镜像 $V^\circ$ 的 B 模型，它可以用 Hodge 理论来计算。然后，\hyperref[chap:quintic]{第 2 章}以五次三维簇 $V\subset\mathbb P^4$ 为例，展示这些思想的具体面貌，并在此过程中推导出 \SourceCite{CdGP} 中的枚举预测。这一章的结尾还将预览\hyperref[chap:mirror-theorem]{第~\ref*{chap:mirror-theorem} 章}给出的镜像对称证明。

读完前两章以后，读者可以根据自己的数学兴趣和专长，选择不同的路线。为了帮助读者作出选择，下面介绍余下各章的一些重点：

\noindent\textbf{\hyperref[chap:toric]{第 3 章}。} 熟悉环面几何的读者可以跳过本章大部分内容。\hyperref[sec:fano-reflexive]{第 3.5 节}介绍自反多胞形；Batyrev 的镜像构造将用到它们。

\noindent\textbf{\hyperref[chap:mirror-constructions]{第 4 章}。} \hyperref[sec:batyrev-mirror]{第 4.1 节}介绍 Batyrev 的镜像构造，并给出支持镜像关系的一些证据。\hyperref[sec:quintic-revisited]{第 4.2 节}说明这一构造如何用于五次三维簇。\hyperref[sec:voisin-borcea]{第 4.4 节}则利用 K3 曲面构造若干有趣的 Calabi--Yau 三维簇镜像对。

\SourcePage{17}
\noindent\textbf{\hyperref[chap:hodge-yukawa]{第 5 章}。} 最大幂幺单值性是镜像对称的重要组成部分，其定义见\hyperref[sec:maximal-unipotent]{第 5.2 节}。对射影超曲面、环面超曲面及超几何方程的计算方法感兴趣的读者，应阅读\hyperref[sec:griffiths-dwork]{第 5.3}、\hyperref[sec:hodge-examples]{5.4}、\hyperref[sec:hypergeometric]{5.5} 节；关注 Calabi--Yau 三维簇的 Hodge 理论的读者，则应仔细阅读\hyperref[sec:yukawa-couplings]{第~\ref*{sec:yukawa-couplings} 节}。

\noindent\textbf{\hyperref[chap:moduli]{第~\ref*{chap:moduli} 章}。} \hyperref[sec:complex-moduli]{第~\ref*{sec:complex-moduli} 节}考察复模，\hyperref[sec:kahler-moduli]{第~\ref*{sec:kahler-moduli} 节}考察 Kähler 模。由于镜像对称预言二者之间存在联系，这两部分的讨论相互交织。我们详细分析的主要例子涉及环面超曲面，因此读者需要了解\hyperref[chap:mirror-constructions]{第 4 章}的 Batyrev 镜像构造。对 Calabi--Yau 流形的模空间、Kähler 锥及镜像对称的整体性质感兴趣的读者，应细读这两节。\hyperref[sec:mirror-map]{第~\ref*{sec:mirror-map} 节}讨论镜像映射，并进一步介绍超几何方程；在环面情形下，正是用这些方程来构造镜像映射。

\noindent\textbf{\hyperref[chap:gromov-witten]{第~\ref*{chap:gromov-witten} 章}。} 除了少数例子，本章独立于前面的章节。其主要研究对象是 Gromov--Witten 不变量。\hyperref[subsec:stable-maps]{第 7.1.1}、\hyperref[subsec:moduli-to-gw]{7.1.2}、\hyperref[subsec:gw-axioms]{7.3.1} 小节是必读内容。此外：
\begin{itemize}
\item 第 \hyperref[subsec:normal-cone]{\ref*{subsec:normal-cone}}--\hyperref[subsec:gw-algebraic-classes]{\ref*{subsec:gw-algebraic-classes}} 小节对虚拟基本类的讨论较为技术性，初读时可以略过。唯一的例外是例题~\ref{ex:quintic-gw-euler-integral}，它给出了五次三维簇某些 Gromov--Witten 不变量的重要公式。第 \hyperref[chap:localization]{\ref*{chap:localization}}、\hyperref[chap:quantum-differential-equations]{\ref*{chap:quantum-differential-equations}}、\hyperref[chap:mirror-theorem]{\ref*{chap:mirror-theorem}} 章中有多处会用到虚拟基本类。
\item 对辛几何感兴趣的读者，应仔细阅读\hyperref[sec:gw-symplectic-definition]{第~\ref*{sec:gw-symplectic-definition} 节}和\hyperref[subsec:instanton-numbers]{第~\ref*{subsec:instanton-numbers} 小节}。
\item 对枚举几何感兴趣的读者，可以重点阅读\hyperref[sec:computing-gw-i]{第~\ref*{sec:computing-gw-i} 节}。其中的一些例子将在\hyperref[chap:quantum-cohomology]{第~\ref*{chap:quantum-cohomology} 章}重新讨论。
\end{itemize}
\hyperref[subsec:instanton-numbers]{第~\ref*{subsec:instanton-numbers} 小节}中一个出人意料之处，是瞬子数 $n_{10}$ 与五次三维簇上次数为 $10$ 的有理曲线数目之间的微妙关系。

\noindent\textbf{\hyperref[chap:quantum-cohomology]{第~\ref*{chap:quantum-cohomology} 章}。} 本章利用\hyperref[chap:gromov-witten]{第~\ref*{chap:gromov-witten} 章}的 Gromov--Witten 不变量，定义量子上同调的两种形式：小量子上同调和大量子上同调。每位读者都应阅读\hyperref[subsec:small-quantum-definition]{第~\ref*{subsec:small-quantum-definition} 小节}关于小量子乘积的内容，以及\hyperref[subsec:a-model-correlation]{第~\ref*{subsec:a-model-correlation} 小节}中的若干例子。了解 Gromov--Witten 势也有帮助；\hyperref[subsec:gw-potential]{第 8.2.2}、\hyperref[subsec:gw-potential-structure]{8.3.1}、\hyperref[subsec:quintic-gw-potential]{8.3.3} 小节讨论了它。然后：
\begin{itemize}
\item 对枚举几何感兴趣的读者，应仔细阅读\hyperref[sec:small-quantum-cohomology]{第 8.1}、\hyperref[sec:big-quantum-cohomology]{8.2}、\hyperref[sec:computing-gw-ii]{8.3} 节。
\item 对 Hodge 理论感兴趣的读者，可以阅读\hyperref[sec:a-model-vhs]{第~\ref*{sec:a-model-vhs} 节}；该节利用量子上同调，在 Calabi--Yau 流形的上同调上构造 A 型 Hodge 结构变分。
\end{itemize}
本章的一项重点是\hyperref[sec:mirror-symmetry-revisited]{第~\ref*{sec:mirror-symmetry-revisited} 节}，其中表述了镜像猜想的若干 Hodge 理论版本。

\noindent\textbf{\hyperref[chap:localization]{第~\ref*{chap:localization} 章}。} 对任何想理解\hyperref[chap:mirror-theorem]{第~\ref*{chap:mirror-theorem} 章}中镜像对称证明的读者，\hyperref[sec:localization-theorem]{第~\ref*{sec:localization-theorem} 节}和\hyperref[subsec:kontsevich-localization]{第~\ref*{subsec:kontsevich-localization} 小节}都是必读内容。尤其需要看例题~\ref{ex:quintic-quartics-localization}：该例采用例题~\ref{ex:quintic-gw-euler-integral} 所给公式的等变版本，计算 Gromov--Witten 不变量 $\langle I_{0,0,d}\rangle$。\hyperref[subsec:multiple-cover-localization]{第 9.2.2}、\hyperref[subsec:nodal-quintic-double-covers]{9.2.3} 小节证明\hyperref[subsec:instanton-numbers]{第~\ref*{subsec:instanton-numbers} 小节}中关于 Calabi--Yau 三维簇的 Gromov--Witten 不变量的若干断言；阅读它们需要深入了解虚拟基本类。

\SourcePage{18}
\noindent\textbf{\hyperref[chap:quantum-differential-equations]{第~\ref*{chap:quantum-differential-equations} 章}。} 只对 \SourceCite{LLY} 证明镜像定理的方法感兴趣的读者，可以跳过本章。但若要理解 Givental 的方法，就需要阅读引力关联数（第 \hyperref[subsec:gravitational-correlator-definition]{\ref*{subsec:gravitational-correlator-definition}}--\hyperref[subsec:gravitational-correlator-examples]{\ref*{subsec:gravitational-correlator-examples}} 小节）、Givental 联络的平坦截面（\hyperref[subsec:giventals-flat-sections]{第~\ref*{subsec:giventals-flat-sections} 小节}），以及 $J$ 函数和量子微分方程（\hyperref[sec:quantum-cohomology-relations]{第~\ref*{sec:quantum-cohomology-relations} 节}）。对 Hodge 理论感兴趣的读者，可以在第 \hyperref[subsec:a-model-gravitational-flat-sections]{\ref*{subsec:a-model-gravitational-flat-sections}}、\hyperref[subsec:cy-threefold-j-function]{\ref*{subsec:cy-threefold-j-function}} 小节读到 A 模型的 Hodge 结构变分。这引出 Picard--Fuchs 算子与量子上同调中的关系式之间一个很好的联系。

\noindent\textbf{\hyperref[chap:mirror-theorem]{第~\ref*{chap:mirror-theorem} 章}。} 这一章讨论当时新近给出的镜像定理证明。需要考察两条路线：
\begin{itemize}
\item \hyperref[sec:quintic-mirror-theorem]{第~\ref*{sec:quintic-mirror-theorem} 节}介绍 \SourceCite{LLY} 对五次三维簇的镜像定理所采用的方法。读者需要掌握前述第 \hyperref[chap:gromov-witten]{\ref*{chap:gromov-witten}}、\hyperref[chap:quantum-cohomology]{\ref*{chap:quantum-cohomology}}、\hyperref[chap:localization]{\ref*{chap:localization}} 章的核心小节。
\item 对于 Givental 版本的镜像定理，读者还需要阅读上面指出的\hyperref[chap:quantum-differential-equations]{第~\ref*{chap:quantum-differential-equations} 章}相关小节。在第 \hyperref[subsec:givental-nef-projective-mirror]{\ref*{subsec:givental-nef-projective-mirror}}、\hyperref[subsec:givental-projective-space-proof]{\ref*{subsec:givental-projective-space-proof}}、\hyperref[subsec:givental-general-bundle-proof]{\ref*{subsec:givental-general-bundle-proof}} 小节，我们讨论 $\mathbb P^n$ 中的数值有效完全交的镜像定理；随后，\hyperref[subsec:givental-toric-mirror]{第~\ref*{subsec:givental-toric-mirror} 小节}考察环境空间为光滑环面簇的情形。
\end{itemize}
特别地，我们会解释这两种方法如何证明\hyperref[chap:quintic]{第 2 章}对五次三维簇作出的全部预测。另一个很有意思的情形见例题~\ref{ex:octic-toric-mirror-theorem}：其中研究的 Calabi--Yau 三维簇，是加权射影空间 $\mathbb P(1,1,2,2,2)$ 中八次超曲面的极小消解。这个环面超曲面例子在第 \hyperref[chap:toric]{\ref*{chap:toric}}、\hyperref[chap:hodge-yukawa]{\ref*{chap:hodge-yukawa}}、\hyperref[chap:moduli]{\ref*{chap:moduli}}、\hyperref[chap:quantum-cohomology]{\ref*{chap:quantum-cohomology}} 章中屡次出现；读者需要回顾那些较早的例子，才能充分理解例题~\ref{ex:octic-toric-mirror-theorem} 的内容。

还须说明，在许多情形下，我们给出的证明并不完整，因为论证的某些细节往往需要读者查阅文献。背景材料也是如此：有些主题，如\hyperref[chap:localization]{第~\ref*{chap:localization} 章}的等变上同调，我们会回顾基本事实；另一些主题，如\hyperref[chap:gromov-witten]{第~\ref*{chap:gromov-witten} 章}的代数叠，则只给出文献指引。希望这种详略程度上的不一致，不至于让读者感到太不适应。

\begingroup\linespread{1.22}\selectfont
\section*{致谢}
\phantomsection\addcontentsline{toc}{section}{致谢}
许多人以不同方式帮助过我们，我们很高兴在此致谢。尤其要感谢 P. Aspinwall、P. Berglund、C. Borcea、P. Candelas、R. Dijkgraaf、B. Greene、T. Hübsch、A. Klemm、P. Mayr、G. Moore、D. R. Morrison、M. R. Plesser、D. van Straten 和 C. Vafa。过去几年中，我们与他们的多次交谈，帮助我们形成了对镜像对称的看法，也加深了我们对相关数学和物理学的理解。

\looseness=-1 1997 年春，Institut Mittag-Leffler 热情接待了我们，谨致谢意。那年春天举办的一场镜像对称研讨班讨论了本书的部分内容，极大地帮助我们改进了论述。我们很高兴感谢所有参加者。虽然这样列举难免无意中漏掉某些人，我们仍特别感谢 P. Aluffi、V. Batyrev、K. Behrend、P. Belorousski、I. Ciocan-Fontanine、A. Elezi、G. Ellingsrud、C. Faber、B. Fantechi、W. Fulton、E. Getzler、L. Göttsche、T. Graber、B. Kim、A. Kresch、D. Laksov、J. Li、A. Libgober、
\SourcePage{19}
R. Pandharipande、K. Ranestad、E. Rødland、D. E. Sommervoll、S. A. Strømme、M. Thaddeus 和 E. Tjøtta 的帮助。

\looseness=-1 在写作本书期间，我们曾将各章草稿寄给多位同行。除上面已经提到的一些人外，我们还要感谢 E. Cattani、E. Ionel、B. Lian、K. Liu、A. Mavlyutov、T. Parker、G. Pearlstein、Y. Ruan、C. Taubes、G. Tian、P. Vermeire 和 S.-T. Yau 提出的有益意见。我们也要感谢 B. Dwork 和 J. Kollár 与我们进行的有益交谈，以及本书早期草稿的审稿人给予的反馈。A. Greenspoon 仔细阅读了最终稿，我们同样深表感谢。当然，书中若仍有错误，责任在我们自己。

我们还感谢 Oklahoma State University 研讨班和 Valley Geometry Seminar 的成员，他们聆听了许多关于镜像对称的报告。他们的反馈与支持令我们十分感激。

本书写作期间，D.A.C. 的研究部分得到 NSF 资助，S.K. 的研究得到 NSF 和 NSA 资助。我们感谢这些机构的支持。

最后，感谢 AMS 的 Sergei Gelfand 和 Sarah Donnelly 在编辑工作上给予指导，也感谢 AMS 技术支持人员 Barbara Beeton 在 AMS-\LaTeX{} 方面提供帮助。

\section*{我们的期望}
\phantomsection\addcontentsline{toc}{section}{我们的期望}
镜像对称是代数几何中一个活跃的研究领域，在未来许多年里，数学家仍有大量工作可做。本书应被视为对这一课题当时状况的一份初步报告；关于镜像对称的定本尚未写成。尽管如此，我们仍希望读者会觉得，本书探讨的种种联系构成了一个激动人心、并且仍在延续的故事。

\noindent\begin{tabular*}{\textwidth}{@{}l@{\extracolsep{\fill}}r@{}}1998 年 11 月 & David A. Cox\\ & Sheldon Katz\end{tabular*}

\noindent\emph{PS}：本书的排印错误清单可在网站 \url{https://math.illinois.edu/~katz/msag/} 查到。
\par\endgroup

\chapter*{符号说明}
\markboth{符号说明}{符号说明}
\phantomsection\addcontentsline{toc}{chapter}{符号说明}
\SourcePage{20}
下面列出本书将用到的部分符号。

\begingroup\linespread{1.18}\selectfont
\setlength{\LTpre}{4pt}\setlength{\LTpost}{4pt}\renewcommand{\arraystretch}{0.96}
\begin{longtable}{@{}p{0.38\textwidth}p{0.57\textwidth}@{}}
\toprule
\multicolumn{2}{@{}l}{\textbf{一般符号}}\\
\midrule
$M$ & 复流形或辛流形\\
$X$ 与 $V$ & 代数簇与 Calabi--Yau 簇\\
$\Theta_X$ 与 $T_X$ & $X$ 的切层与切丛\\
$\omega_X$ 与 $K_X$ & $X$ 的典范丛与典范除子\\
$\widehat\Omega_X^p$ & 轨形 $X$ 上的 Zariski $p$-形式\\
$M(X)$ & $X$ 上有效一维圈的 Mori 锥\\
$A_k(X)$ & $X$ 的第 $k$ 个 Chow 群\\
$\langle a,b\rangle$ 或 $g(a,b)$ & 上同调上的杯积配对\\
$\mathbb P^n$ 与 $\mathbb P(q_0,\ldots,q_n)$ & 射影空间与加权射影空间\\
\midrule
\multicolumn{2}{@{}l}{\textbf{环面簇与多胞形（\hyperref[chap:toric]{第 3 章}）}}\\
\midrule
$\Sigma$ 与 $\Sigma(1)$ & 扇及其一维锥\\
$\Delta$ 与 $\Delta^\circ$ & 多胞形及其极多胞形（或对偶多胞形）\\
$L(\Delta)$ & 指数属于 $\Delta$ 的 Laurent 多项式\\
$X_\Sigma$ 与 $\mathbb P_\Delta$ & 分别由扇 $\Sigma$ 和多胞形 $\Delta$ 确定的环面簇\\
$\operatorname{cpl}(\Sigma)$ & $\Sigma$ 上凸支撑函数构成的锥\\
\midrule
\multicolumn{2}{@{}l}{\textbf{Hodge 理论与 Yukawa 耦合（\hyperref[chap:hodge-yukawa]{第 5 章}）}}\\
\midrule
$\nabla=\nabla^{\mathrm{GM}}$ & Gauss--Manin 联络\\
$\mathcal F^\bullet$ 与 $\mathcal W_\bullet$ & Hodge 滤过与权滤过\\
$\mathcal H$ 与 $\mathcal H_{\mathbb Z}$ & Hodge 丛 $\mathcal F^0$ 及其中的整数子层\\
$T_j$ 与 $N_j$ & 单值性变换及其对数\\
$\left\langle\frac{\partial}{\partial t_i},\frac{\partial}{\partial t_j},\frac{\partial}{\partial t_k}\right\rangle=Y_{ijk}$ & 归一化 Yukawa 耦合，或 B 模型关联函数\\
\midrule
\multicolumn{2}{@{}l}{\SourcePage{21}\textbf{复模与 Kähler 模（\hyperref[chap:moduli]{第~\ref*{chap:moduli} 章}）}}\\
\midrule
$\mathcal M$ 与 $\overline{\mathcal M}$ & 复模空间及其紧化\\
$\mathcal M_{\mathrm{poly}}$ 与 $\overline{\mathcal M}_{\mathrm{poly}}$ & 多项式模空间及其紧化\\
$\mathcal M_{\mathrm{simp}}$ 与 $\overline{\mathcal M}_{\mathrm{simp}}$ & 简化模空间及其紧化\\
$K(V)$ 与 $\overline{K(V)}_{\mathbb C}$ & Kähler 锥与复化 Kähler 空间\\
$\mathcal{KM}$ 与 $\overline{\mathcal{KM}}$ & Kähler 模空间及其紧化\\
$\mathcal{KM}_{\mathrm{toric}}$ 与 $\overline{\mathcal{KM}}_{\mathrm{toric}}$ & 环面 Kähler 模空间及其紧化\\
\midrule
\multicolumn{2}{@{}l}{\textbf{Gromov--Witten 不变量与量子上同调（第 \hyperref[chap:gromov-witten]{\ref*{chap:gromov-witten}}、\hyperref[chap:quantum-cohomology]{\ref*{chap:quantum-cohomology}} 章）}}\\
\midrule
$\overline M_{g,n}(X,\beta)$ 与 $\overline{\mathcal M}_{g,n}(X,\beta)$ & 类为 $\beta$、亏格为 $g$、带 $n$ 个标记点的稳定映射的粗模空间与精细模叠\\
$[\overline{\mathcal M}_{g,n}(X,\beta)]^{\mathrm{virt}}$ & 虚拟基本类\\
$I_{g,n,\beta}(\alpha_1,\ldots,\alpha_n)$ 与 $\langle I_{g,n,\beta}\rangle(\alpha_1,\ldots,\alpha_n)$ & Gromov--Witten 类与 Gromov--Witten 不变量\\
$n_\beta$ 与 $n_d$ & 一般情形及五次三维簇的瞬子数\\
$N_\beta=\langle I_{0,0,\beta}\rangle$ 与 $N_d=\langle I_{0,0,d}\rangle$ & 一般情形及五次三维簇的零标记点、亏格零 Gromov--Witten 不变量\\
$*_{\mathrm{small}}$ 与 $*$ & 小量子乘积与大量子乘积（小量子乘积有时也记为 $*$）\\
$\langle a,b,c\rangle$ & 三点函数，或 A 模型关联函数\\
$\Phi(\gamma)$ & 亏格零 Gromov--Witten 势\\
$\nabla=\nabla^{\mathrm{GW}}$ & A 模型联络\\
\midrule
\multicolumn{2}{@{}l}{\textbf{Chern 类与等变上同调（\hyperref[chap:localization]{第~\ref*{chap:localization} 章}）}}\\
\midrule
$c_i(\mathcal E)$ 与 $c_i^T(\mathcal E)$ & $\mathcal E$ 的 Chern 类与等变 Chern 类\\
$\operatorname{Euler}(\mathcal E)$ 与 $\operatorname{Euler}_T(\mathcal E)$ & $\mathcal E$ 的 Euler 类与等变 Euler 类\\
$\int_{X_T}$ & 等变积分\\
$\lambda_0,\ldots,\lambda_n$ 与 $\hbar$ & $H^*(BT)$ 与 $H^*(B\mathbb C^*)$ 的环生成元；当 $G=\mathbb C^*\times T$、$T=(\mathbb C^*)^{n+1}$ 时，它们生成 $H^*(BG)$\\
$H$ 与 $p$ & 超平面类与等变超平面类\\
$q_0,\ldots,q_n$ & $\mathbb P^n$ 上标准 $T$ 作用的不动点\\
\midrule
\multicolumn{2}{@{}l}{\SourcePage{22}\textbf{引力关联数与 $J$ 函数（\hyperref[chap:quantum-differential-equations]{第~\ref*{chap:quantum-differential-equations} 章}）}}\\
\midrule
$\langle\tau_{d_1}\gamma_1,\ldots,\tau_{d_n}\gamma_n\rangle_{g,\beta}$ & 引力关联数\\
$\langle\gamma_1,\ldots,\gamma_n\rangle_{g,\beta}$ & $\langle I_{g,n,\beta}\rangle(\gamma_1,\ldots,\gamma_n)$ 的另一种记法\\
$\langle\!\langle\tau_{d_1}\gamma_1,\ldots,\tau_{d_n}\gamma_n\rangle\!\rangle_g$ & 亏格为 $g$ 的引力耦合\\
$\Phi_g^{\mathrm{grav}}(\gamma)$ & 亏格为 $g$ 的引力 Gromov--Witten 势\\
$*_{\mathrm{grav}}$ & 引力量子乘积\\
$\nabla^g$ 与 $\widetilde\nabla^g$ & Givental 联络及其对偶联络\\
$s_j$ & Givental 联络的平坦截面\\
$\left\langle\!\left\langle\dfrac{T_a}{\hbar-c},T_j\right\rangle\!\right\rangle_0$ & $\displaystyle\sum_{n=0}^{\infty}\hbar^{-(n+1)}\langle\!\langle\tau_n T_a,T_j\rangle\!\rangle_0$ 的符号记法\\
$J=J_X$ & 光滑簇 $X$ 的 Givental $J$ 函数\\
\midrule
\multicolumn{2}{@{}l}{\textbf{镜像定理（\hyperref[chap:mirror-theorem]{第~\ref*{chap:mirror-theorem} 章}）}}\\
\midrule
$\mathcal V=\bigoplus_{i=1}^{\ell}\mathcal O_{\mathbb P^n}(a_i)$ & 镜像定理中使用的向量丛\\
$\mathcal V_d$ 与 $\mathcal V_{d,k}$ & 由 $\mathcal V$ 诱导的 $\overline M_{0,0}(\mathbb P^n,d)$ 与 $\overline M_{0,k}(\mathbb P^n,d)$ 上的丛\\
$M_d$ 与 $N_d$ & $\overline M_{0,0}(\mathbb P^1\times\mathbb P^n,(1,d))$ 与 $\mathbb P(H^0(\mathbb P^1,\mathcal O_{\mathbb P^1}(d))^{n+1})$ 的简记\\
$p_{i,r}$ & $N_d$ 上 $G=\mathbb C^*\times T$ 作用的不动点\\
$\kappa$ & $H_G^*(N_d)$ 中的等变超平面类\\
$\mathcal R_T$ 与 $\mathcal R_G$ & $H^*(BT)$ 与 $H^*(BG)$ 的分式域\\
$\widehat P$ 与 $\widehat Q$ & 重要的 Euler 数据\\
$HG[B](t)$ & 集合 $B=\{B_d\}$ 所确定的上同调值函数\\
$\Psi(t)$ & 镜像映射\\
$I_{\mathcal V}$ 与 $J_{\mathcal V}$ & Givental 的上同调值函数\\
$\mathcal V'_{d,k,1}$ & 定义 $J_{\mathcal V}$ 所用的 $\mathcal V_{d,k}$ 的某个子丛\\
$*_X$ & 由完全交 $X\subset\mathbb P^n$ 决定的 $H^*(\mathbb P^n)$ 上的修正量子乘积\\
$\overline I_{\mathcal V}$ & 量子超平面截面原理中使用的上同调值函数\\
$I_T$ 与 $J_T$ & $I_{\mathcal V}$ 与 $J_{\mathcal V}$ 的等变版本\\
$Z_{i,\mathcal V}$ & 对 $0\le i\le n$ 定义的一组函数，它们唯一确定 $J_T$\\
\bottomrule
\end{longtable}
\endgroup

参考文献的引用约定，在\hyperlink{source:460}{参考文献表}的开头说明。

\mainmatter
\chapter{引言}\label{chap:introduction}
\SourcePage{24}

镜像对称在代数几何中作出了一些令人惊讶的预测，所涉及的对象从五次三维簇上的有理曲线数目，到某些模空间的结构。这些都是极有价值的研究问题，而且如前言所述，已经催生了许多非常有趣的数学。然而，若要弄清这些预测从何而来，代数几何学家就必须深入物理学的语言；这种语言既陌生，有时又因欠缺严格性而令人沮丧。因此，在开始考察镜像对称的代数几何之前，我们先讨论促成镜像对称的动机，以及镜像对称在物理学中意味着什么。由于会涉及许多读者可能不熟悉的物理学术语，我们的介绍难免不够完整。尽管如此，我们仍希望传达这一非凡现象背后的一些直观想法。

接着，我们将讨论三点函数——它们对于镜像对称的枚举预测至关重要——以及 Calabi--Yau 流形出现在这一理论中的物理原因。最后，在本章结尾，我们会回到读者更熟悉的数学世界，预览本书其余章节将要探索的代数几何。

\section{镜像对称的物理学}\label{sec:physics}

本节旨在让读者体会镜像对称为什么会出现，以及它应当蕴含什么。从物理学的观点看，镜像对称自然地产生于超对称弦理论的标准构造；因此，我们会先对弦和超对称作一些初步说明。请读者放心：我们不预设任何物理学背景！我们的目的，是传达这些物理理论的风貌，并借此增强读者对相应数学的直观认识。读者不必详尽掌握其中的物理学，虽然后面的章节会引用本章介绍的物理内容。关于弦理论的一般参考资料，可参阅 \SourceCite{GSW,Polchinski2}。\index{超对称弦理论|hyperpage}

在弦理论中，物理过程由弦在时空中的传播来描述。传播中的弦扫出一个曲面，称为弦的\emph{世界面} $\Sigma$。经典场可以描述为世界面上的函数、丛的截面等；将它们量子化，便得到一个二维量子场论。这个理论具有一个广义的态 Hilbert 空间，也有一组可观测量；后者成为态空间上的自伴算子。理论的另外两个关键组成部分，是把 Lagrange 量在世界面 $\Sigma$ 上积分得到的作用量 $S$，以及\emph{关联函数}
\begin{equation}\label{eq:correlation}
 \bigl\langle\phi(x_1),\ldots,\phi(x_n)\bigr\rangle
 =\int[D\phi]\,\phi(x_1)\cdots\phi(x_n)e^{iS(\phi)}.
\end{equation}
\SourcePage{25}
这里 $\phi$ 是一个可观测量。这个 Feynman 积分遍及所有可能的世界面；到目前为止，它在数学上尚无定义。\hyperref[app:physical-theories]{附录 B} 会更详细地讨论这类理论。现在，我们先关注它们的一般特征。

弦理论有一个优点：它消除了粒子分裂为两个粒子时出现的某些问题。下图表明，如果把粒子表示为点，就会出现一个奇点；用弦来表示时，得到的则是一个光滑的带边二维流形：
\begin{figure}[H]
 \centering\includegraphics[width=0.79219\textwidth]{1-1-string-splitting.png}
\end{figure}

不过，弦理论仍有一些不理想的特征，其中包括许多需要重整化处理的无穷大。近来一个非凡的发现是，超对称能够消除其中许多困难。超对称把玻色子——自旋为整数、波函数对称的粒子——变为费米子——自旋为半整数、波函数反对称的粒子——反之亦然。虽然迄今尚无实验证实超对称，但由于超对称理论具有良好的性质，它们已在理论物理学中变得十分重要。\index{超对称|hyperpage}

我们还需要另一个要素：世界面 $\Sigma$ 带有共形结构，我们的超对称弦理论必须在共形等价下保持等价。因此，该理论是一种\emph{超共形场论}，简称 SCFT。这类理论的对称群的 Lie 代数称为超共形代数。它以共形代数——世界面共形变换群的 Lie 代数——为子代数，同时还包含超对称变换。\index{超共形场论|hyperpage}

\IndexLocation{indexloc:1306:25}\IndexLocation{indexloc:1307:25}\IndexLocation{indexloc:1308:25}超弦理论有四种基本类型：I 型、IIA 型、IIB 型，以及我们最感兴趣的杂化弦理论。杂化弦理论是一种 $N=2$ 的 SCFT，因为费米子解耦为左行解和右行解；这意味着实际上有四种超对称，其中两种左行、两种右行。因此，更准确地说，杂化弦理论是一种 $(2,2)$ 理论：其左行部分与右行部分各有两种独立的超对称。

$N=2$ 超共形代数包含通常的超共形代数的两个副本，因此有两个 $u(1)$ 子代数。一个位于理论的右行部分，对两种超对称施加无穷小旋转；另一个位于左行部分，作用方式相似。给两种超对称变换排定顺序，便可为每个 $u(1)$ 选取一个生成元，但这种选取并非典范的。如果颠倒两种超对称的顺序，$u(1)$ 生成元的符号就会改变。这两个子代数各自的生成元记为
\SourcePage{26}
$(Q,\overline Q)$，它们除符号外唯一确定。我们把 $(Q,\overline Q)$ 看作作用于态 Hilbert 空间的算子；在 $u(1)\times u(1)$ 的作用下，该空间分解为特征子空间。正如即将看到的，这些子空间可能非常有趣。

到目前为止，我们只是抽象地讨论杂化弦理论。下一步是实际构造这样的理论；代数几何也由此登场。构造方法有好几种，对我们最重要的是由 Calabi--Yau 三维簇\footnote{眼下，我们把 Calabi--Yau 三维簇定义为第一 Betti 数为零、典范类平凡的光滑、紧致、连通三维簇。稍后我们会允许某些奇点。} $V$ 以及 $V$ 上的\emph{复化 Kähler 类} $\omega=B+iJ$ 所确定的\emph{非线性 $\sigma$ 模型}，简称 $\sigma$ 模型。这里 $B,J\in H^2(V,\mathbb R)$，且 $J$ 是 Kähler 类。\index{非线性sigma模型@非线性 $\sigma$ 模型|hyperpage}

从输入数据 $(V,\omega)$ 出发，有一种几何构造可以得到 $N=2$ 的 SCFT，并且明确赋予各种超对称不同的作用。因此，在这一背景下，$(Q,\overline Q)$ 有一个典范选取。于是，我们的 Hilbert 空间上的 $u(1)\times u(1)$ 表示也得到明确选取。可以计算出：对 $p,q\geq0$，$(Q,\overline Q)$ 的特征子空间为
\begin{equation}\label{eq:eigenspaces}
\begin{aligned}
(p,q)\text{ 特征子空间}&\simeq H^q(V,\wedge^pT_V),\\
(-p,q)\text{ 特征子空间}&\simeq H^q(V,\Omega_V^p).
\end{aligned}
\end{equation}
\hyperref[sec:nonlinear-sigma-models]{附录 B.2} 更完整地解释什么是非线性 $\sigma$ 模型；第 \ref{sec:calabi-yau-physics} 节则说明 Calabi--Yau 条件如何从物理学中产生。

杂化弦理论中最重要的场与 $H^1(V,T_V)$ 和 $H^1(V,\Omega_V^1)$ 中的元素相联系；它们分别对应特征值 $(1,1)$ 和 $(-1,1)$。与这些空间中的元素相应的算子，称为\emph{边缘算子}。它们之所以重要，部分原因在于它们与来自 $(V,\omega)$ 的 $\sigma$ 模型的模密切相关。直观地说，同时改变 $V$ 的复结构和复化 Kähler 类 $\omega=B+iJ$，便得到 SCFT 的模；不过，SCFT 模空间还存在一些额外的离散等同，暂且不予考虑。读者或许熟悉复模，但“复化 Kähler 模”这一概念可能比较陌生。我们将在第 \ref{sec:math-preview} 节详细介绍。\index{边缘算子|hyperpage}\index{复化Kähler模@复化 Kähler 模|hyperpage}

由此可知，SCFT 模空间有两个叶状结构：其中一个的叶可描述为“$V$ 不变”，另一个则是“$\omega$ 不变”。这就得到下图：
\begin{figure}[H]
 \centering\refstepcounter{equation}\label{eq:scft-foliations}
 \makebox[\textwidth]{\makebox[0pt][l]{(\theequation)}\hfill
 \raisebox{-.5\height}{\includegraphics[width=0.48536\textwidth]{1-2-scft-moduli.png}}\hfill\mbox{}}
\end{figure}

\SourcePage{27}
尽管有这幅图，我们仍要强调：SCFT 模空间\emph{并非}复结构模空间与 Kähler 结构模空间的乘积，即使在局部也不是。事实上，$\omega$ 的 Kähler 模空间可能依赖于 $V$ 的复结构 \SourceCite{Wilson2}。一般而言，只有先固定 $V$ 的复结构，它才得到良好定义。另一方面，\SourceCite{Wilson2} 表明，对足够一般的 Calabi--Yau 三维簇 $V$，$\omega$ 的 Kähler 模与 $V$ 的复结构无关。\hyperref[chap:moduli]{第~\ref*{chap:moduli} 章}将更详细地讨论这些问题。因此，上图虽有助于概念上的理解，却没有反映 SCFT 模空间的微妙之处。

现在，我们对 Calabi--Yau 三维簇和复化 Kähler 类如何给出有趣的物理理论已有了较清楚的认识，可以解释镜像对称从何而来了。基本出发点是 $(Q,\overline Q)$ 的符号不确定性。前面说过，$(Q,\overline Q)$ 除符号外唯一确定，但由 $(V,\omega)$ 得到的 $\sigma$ 模型却作出了一个十分具体的选取。如果将 $Q$ 改为 $-Q$，而保持 $\overline Q$ 不变，那么 $(p,q)$ 与 $(-p,q)$ 特征子空间就会互换；依照式~\eqref{eq:eigenspaces}，这会交换 $H^q(V,\wedge^pT_V)$ 与 $H^q(V,\Omega_V^p)$。一般来说，这两个向量空间的维数不同，所以对同一个模型不可能如此交换。然而，从物理学的观点看，这种符号改变是合理的。这种不对称性暗示：符号改变或许对应于由另一个不同的配对 $(V^\circ,\omega^\circ)$ 所产生的 $\sigma$ 模型。如果这样的配对存在，就称 $(V,\omega)$ 与 $(V^\circ,\omega^\circ)$ 构成一个\emph{镜像对}。更正式地说，我们有如下来自物理学的定义。\index{镜像对|hyperpage}

\begin{physicsdefinition}\label{def:physics-mirror-pair}
若 $(V,\omega)$ 与 $(V^\circ,\omega^\circ)$ 的 $\sigma$ 模型所诱导的超共形场论同构，而且这些理论的 $N=2$ 超共形表示只相差上述符号改变，就称它们构成一个镜像对。
\end{physicsdefinition}

注意，这不是一个数学定义，因为与 $(V,\omega)$ 相联系的 SCFT 尚未得到严格定义。不过，\hyperref[chap:quantum-cohomology]{第~\ref*{chap:quantum-cohomology} 章}将给出\emph{数学镜像对}的严谨定义，其中会包含镜像对称所预言的许多性质。

下面就来考察这些性质。如果 $(V,\omega)$ 与 $(V^\circ,\omega^\circ)$ 构成镜像对，我们便得到同构的 SCFT。但这在数学上意味着什么？理解镜像对称的数学后果，正是本书的主要目标之一。

要看镜像对称对 $V$ 与 $V^\circ$ 有何启示，先把式~\eqref{eq:eigenspaces} 与特征值的变换 $(p,q)\leftrightarrow(-p,q)$ 结合起来，就得到同构
\begin{equation}\label{eq:mirror-cohomology-tangent}
\begin{aligned}
H^q(V,\wedge^pT_V)&\simeq H^q(V^\circ,\Omega^p_{V^\circ}),\\
H^q(V,\Omega^p_V)&\simeq H^q(V^\circ,\wedge^pT_{V^\circ}).
\end{aligned}
\end{equation}
由于 $V$ 是 Calabi--Yau 簇，它有一个处处非零的全纯三形式 $\Omega$；与 $\Omega$ 作杯积，便给出一个非典范同构 $H^q(V,\wedge^pT_V)\simeq H^q(V,\Omega_V^{3-p})$。$V^\circ$ 也一样。因此，式~\eqref{eq:mirror-cohomology-tangent} 可改写为
\begin{equation}\label{eq:mirror-cohomology-forms}
\begin{aligned}
H^q(V,\Omega_V^{3-p})&\simeq H^q(V^\circ,\Omega_{V^\circ}^p),\\
H^q(V,\Omega_V^p)&\simeq H^q(V^\circ,\Omega_{V^\circ}^{3-p}).
\end{aligned}
\end{equation}
这些同构可以借助 Hodge 菱形得到一个很好的解释。由于 $V$ 是光滑三维簇，它的 Hodge 数 $h^{p,q}(V)=\dim H^q(V,\Omega_V^p)$ 满足对称关系
\[
h^{p,q}(V)=h^{q,p}(V)=h^{3-p,3-q}(V)=h^{3-q,3-p}(V).
\]
又因为 $V$ 是 Calabi--Yau 簇，所以 $b_1(V)=0$ 且 $\Omega_V^3\simeq\mathcal O_V$。这推出 $h^{1,0}(V)=0$、$h^{3,0}(V)=1$。此外，$h^{0,2}(V)=\dim H^2(V,\mathcal O_V)=\dim H^1(V,\mathcal O_V)=$
\SourcePage{28}
$h^{0,1}(V)=0$；其中第二个等号来自 Serre 对偶和 $\Omega_V^3\simeq\mathcal O_V$。因此，$V$ 的 Hodge 菱形为
\[
\begin{array}{ccccccc}
&&&1&&&\\
&&0&&0&&\\
&0&&h^{1,1}(V)&&0&\\
1&&h^{2,1}(V)&&h^{2,1}(V)&&1\\
&0&&h^{1,1}(V)&&0&\\
&&0&&0&&\\
&&&1&&&
\end{array}
\]
现在比较镜像对 $V$ 与 $V^\circ$ 的 Hodge 菱形。式~\eqref{eq:mirror-cohomology-forms} 蕴含 $h^{p,q}(V)=h^{3-p,q}(V^\circ)$，表明 $V^\circ$ 的 Hodge 菱形是 $V$ 的 Hodge 菱形关于一条 $45^\circ$ 直线的反射，即镜像。“镜像对称”这个名称正是由此而来。\index{Hodge菱形@Hodge 菱形|hyperpage}

式~\eqref{eq:mirror-cohomology-tangent} 和~\eqref{eq:mirror-cohomology-forms} 其实只是一系列愈发引人瞩目的镜像对称后果中的头两项。构成镜像对的下一个有趣后果涉及模空间。要理解模空间如何出现，注意式~\eqref{eq:mirror-cohomology-tangent} 给出 $H^1(V,T_V)\simeq H^1(V^\circ,\Omega^1_{V^\circ})$ 和 $H^1(V,\Omega_V^1)\simeq H^1(V^\circ,T_{V^\circ})$。这自然地把 $V$ 的复模空间的切空间与 $\omega^\circ$ 的 Kähler 模空间的切空间等同起来（定义见第 \ref{sec:math-preview} 节）；同样，它把 $\omega$ 的 Kähler 模空间的切空间与 $V^\circ$ 的复模空间的切空间等同起来。因此，$V$ 的复模空间局部同构于 $\omega^\circ$ 的 Kähler 模空间；同样，$V^\circ$ 的复模空间局部同构于 $\omega$ 的 Kähler 模空间。这些局部同构合称为\emph{镜像映射}。\hyperref[chap:moduli]{第~\ref*{chap:moduli} 章}将更详细地研究复模和 Kähler 模，并给出镜像映射的严谨定义。

回想我们那个略有失真的示意图~\eqref{eq:scft-foliations}：$(V,\omega)$ 的 SCFT 模空间有两个叶状结构，其中一个的叶满足“$V$ 不变”，另一个满足“$\omega$ 不变”。因此，若 $(V,\omega)$ 与 $(V^\circ,\omega^\circ)$ 构成镜像对，就得到下列 SCFT 模空间的局部同构：
\begin{figure}[H]
 \centering\refstepcounter{equation}\label{eq:mirror-foliations}
 \makebox[\textwidth]{\makebox[0pt][l]{(\theequation)}\hfill
 \raisebox{-.5\height}{\includegraphics[width=0.86471\textwidth]{1-3-mirror-moduli.png}}\hfill\mbox{}}
\end{figure}
这幅图也说明，谈论 Calabi--Yau 流形 $V$ 的“那个”镜像流形并无意义。给定 $V$ 后，我们可以在 SCFT 模空间的一个叶内任意改变 $\omega$；在镜像一侧，这将使 $\omega^\circ$ 保持不变，而 $V^\circ$ 自由变化。因此，真要说“$V$ 的镜像”，它应当是类 $\omega^\circ$，
\SourcePage{29}
连同 $V^\circ$ 的那些形变所构成的模空间；在这些形变上，$\omega^\circ$ 仍可被理解为一个复化 Kähler 类。

除了上面所讨论的，镜像对的存在还有更多后果，其中并非所有后果都已得到理解。基本思想是：原则上，凡是能够用 SCFT 定义的量，都可以利用 SCFT 的两种不同构造来计算。在最理想的情形下，还可以用 $V$ 与 $V^\circ$ 的几何来计算这些量。注意，$Q$ 的符号发生了改变，而 $H^q(V,\Omega_V^p)$ 又对应于 $H^q(V^\circ,\wedge^pT_{V^\circ})$；因此，预期两个模型中的几何计算会有所不同。最好的例子是式~\eqref{eq:correlation} 中已经提到的 SCFT 关联函数。第 \ref{sec:three-point} 节将更详细地讨论它们；它们也是镜像对称作出枚举预测的关键。我们将看到，在镜像族上进行计算，可以得出关于原来 Calabi--Yau 流形的惊人结论。

在物理学文献中，镜像对称是一个内容丰富的现象。除了前面讨论的非线性 $\sigma$ 模型的镜像对称，一些非几何型的 SCFT 中也观察到了镜像对称，其中包括极小模型和 Landau--Ginzburg 轨形。要解释镜像对称在这些情形下如何起作用，需要对 SCFT 按一个有限群取“轨形”。首先取由不变场组成的子理论；但所得子理论在时间演化下并不稳定，也就是说，它不是幺正的。因此需要增添一些场，得到再次成为 SCFT 的幺正理论。实际的物理过程更微妙，细节见\hyperref[sec:landau-ginzburg-models]{附录 B.4}。如果按经过精心选取的群作用取商，就会得到与出发时相同的物理理论，只是 $Q$ 特征值的符号发生了改变。这种形式的“镜像对称”甚至早于非线性 $\sigma$ 模型中镜像对称的发现。

Calabi--Yau 三维簇存在镜像对称的早期证据，来自加权射影空间中的 Calabi--Yau 超曲面，或它们按有限群取商后所得对象的列表。这些超曲面的 Hodge 数展现出醒目但远非完美的对称性。对于其中一些超曲面，\SourceCite{GPl} 先证明某些 Landau--Ginzburg 理论的镜像对称，再把这些理论与超曲面的 $\sigma$ 模型联系起来，从而证明了相应的镜像对称。正如\hyperref[chap:mirror-constructions]{第 4 章}将说明的，\SourceCite{CdK} 指出，这些加权射影超曲面都属于 Batyrev 的自反多胞形构造 \SourceCite{Batyrev4} 所产生的对象。人们猜想——而且普遍相信——由自反多胞形得到的 Calabi--Yau 三维簇具有镜像对称；更一般地，范围更大的环面完全交类 \SourceCite{Borisov1} 也具有镜像对称。过去几年已有大量证据支持这些看法 \SourceCite{CdGP,Font,Morrison1,CdFKM,HKTY1,CFKM,HKTY2,BK1,AGM1,ES1,BKK,Kontsevich2,MP1,Givental2,Givental4,LLY}。\hyperref[chap:mirror-theorem]{第~\ref*{chap:mirror-theorem} 章}将概述证明镜像定理的两种相关方法；该定理确立 Calabi--Yau 环面完全交与其猜想镜像的某些关联函数相等。\index{自反多胞形|hyperpage}

\section{三点函数}\label{sec:three-point}
式~\eqref{eq:correlation} 定义的关联函数，在 SCFT 中本身就是值得研究的对象。在物理学里，研究粒子的连续世代时，它们会自然出现。最常见的关联函数是\emph{三点函数}，
\SourcePage{30}
它描述来自三个世代——不一定互不相同——的粒子之间的相互作用。在基本粒子物理的标准模型中，粒子的一个\emph{世代}，是指在电磁力、弱核力与强核力作用下具有特定相互作用类型的一组粒子。实验表明，粒子共有三个世代。其中一个世代包含最熟悉的粒子，即电子及与它相伴的中微子，以及构成质子和中子的上夸克与下夸克。已知的另外两个世代则包含更奇特的夸克和轻子。

我们将考察由 Calabi--Yau 三维簇 $V$ 和复化 Kähler 类 $\omega$ 所产生的非线性 $\sigma$ 模型的三点函数。最有趣的三点函数是 Yukawa 耦合，它们来自前一节讨论的边缘算子。边缘算子对应于 $H^1(V,\Omega_V^1)$ 和 $H^1(V,T_V)$，因此需要考虑两类 Yukawa 耦合。\index{Yukawa耦合@Yukawa 耦合|hyperpage}

先看由 $H^1(V,\Omega_V^1)$ 产生的 Yukawa 耦合。对于这个空间中的每个元素，$\sigma$ 模型都对应一个由场组成的 $27$ 维向量空间；这些场构成 $E_6$ 的一个不可约表示。为了把弦理论与现实的物理世界联系起来，人们认为这个向量空间包含一代基本粒子。令 $\omega_i\in H^1(V,\Omega_V^1)$，$i=1,2,3$，分别对应三个世代。那么这些世代之间的 Yukawa 耦合，是从每一代各取一个粒子而形成的、在物理学上重要的三粒子耦合。

在此情形下，用 Feynman 路径积分方法算得 Yukawa 耦合为
\begin{equation}\label{eq:a-yukawa}
\begin{aligned}
\langle\omega_1,\omega_2,\omega_3\rangle
={}&\int_V\omega_1\wedge\omega_2\wedge\omega_3\\
&+\sum_{\substack{\beta\ne0\\\beta\text{ 有效}}}n_\beta
 \left(\int_\beta\omega_1\right)
 \left(\int_\beta\omega_2\right)
 \left(\int_\beta\omega_3\right)
 \frac{e^{2\pi i\int_\beta\omega}}{1-e^{2\pi i\int_\beta\omega}}.
\end{aligned}
\end{equation}
求和遍及非零有效同调类 $\beta\in H_2(V,\mathbb Z)$；粗略地说，$n_\beta$ 是同调类为 $\beta$ 的有理曲线的数目。要严格定义 $n_\beta$，必须理解 Gromov--Witten 不变量；\hyperref[chap:gromov-witten]{第~\ref*{chap:gromov-witten} 章}将对此作详细讨论。理论预言，当 $V$ 发生形变时，$n_\beta$ 不变。因此，这个 Yukawa 耦合依赖于 $\omega$，却不依赖于 $V$。\hyperref[chap:quantum-cohomology]{第~\ref*{chap:quantum-cohomology} 章}将说明，式~\eqref{eq:a-yukawa} 中的耦合与 $V$ 的量子上同调环密切相关。

刚才描述的 Yukawa 耦合有时称为\emph{A 模型关联函数}。这个名称来自\hyperref[sec:nonlinear-sigma-models]{附录 B.2} 所介绍的 A 模型。A 模型是 SCFT 的一个“扭曲”版本：某些场局部看来与 $\sigma$ 模型中的场相同，但从整体来看，它们成为原来所在丛的某种扭曲的截面。这些耦合与 A 模型中相应的三点函数等同。

再简要说明式~\eqref{eq:a-yukawa} 从何而来。注意，其第一项只是杯积：
\begin{equation}\label{eq:a-cup}
\int_V\omega_1\wedge\omega_2\wedge\omega_3.
\end{equation}
从物理学的观点看，由式~\eqref{eq:a-cup} 走向 A 模型关联函数~\eqref{eq:a-yukawa} 时，发生了两件事。首先有一个\emph{非重整化定理}：从微扰的观点看，不需要任何量子修正。正如上一节所说，这是超对称的一个良好后果。不过，还须考虑世界面上的非微扰修正；
\SourcePage{31}
在这里，这些修正来自\emph{全纯瞬子}。全纯瞬子是从紧致 Riemann 曲面 $\Sigma$ 到 $V$ 的非常值全纯映射。第 \hyperref[chap:gromov-witten]{\ref*{chap:gromov-witten}}、\hyperref[chap:quantum-cohomology]{\ref*{chap:quantum-cohomology}} 章更仔细地处理 A 模型关联函数时，我们会看到，$\Sigma$ 可以有结点奇点，也可以有不止一个分支。在 A 模型关联函数中，只须考虑 $\Sigma$ 的亏格为零的瞬子。粗略地说，式~\eqref{eq:a-yukawa} 的 $n_\beta$ 计数的正是它们。按照\hyperref[chap:gromov-witten]{第~\ref*{chap:gromov-witten} 章}的术语，我们称 $n_\beta$ 为\emph{瞬子数}。\index{瞬子数|hyperpage}

第二类需要考虑的 Yukawa 耦合来自 $H^1(V,T_V)$。此时，该空间中的元素对应于 $E_6$ 的共轭 $27$ 维表示。取 $\theta_i\in H^1(V,T_V)$，$i=1,2,3$，分别对应三个世代，便得到相应的 Yukawa 耦合：
\begin{equation}\label{eq:b-yukawa}
\langle\theta_1,\theta_2,\theta_3\rangle
=\int_V\Omega\wedge(\theta_1\cdot\theta_2\cdot\theta_3\cdot\Omega),
\end{equation}
这里 $\Omega$ 是 $V$ 上的一个全纯三形式。表达式 $\theta_1\cdot\theta_2\cdot\theta_3\cdot\Omega$ 由以下复合映射定义：
\[
S^3H^1(V,T_V)\otimes H^0(V,\Omega_V^3)
\longmapsto H^3(V,\wedge^3T_V\otimes\Omega_V^3)
\simeq H^3(V,\mathcal O_V)\simeq H^{0,3}(V).
\]
也可以把式~\eqref{eq:b-yukawa} 理解为
\[
\int_V\Omega\wedge
\bigl(\nabla_{\theta_1}\nabla_{\theta_2}\nabla_{\theta_3}\Omega\bigr),
\]
其中 $\nabla$ 是 Gauss--Manin 联络。\footnote{\SourceCite{CdGP} 中的 Yukawa 耦合就是以这种形式出现的。正如命题~\ref{prop:normalized-yukawa-pf} 将说明的，Hodge 理论促使我们在 Yukawa 耦合的定义中再加一个负号。\SourceCite{Cd} 还讨论了加负号的物理理由。由于 $\Omega$ 可以乘以任意常数，眼下符号并不重要；但在\hyperref[subsec:normalized-yukawa-mirror]{第~\ref*{subsec:normalized-yukawa-mirror} 小节}，我们会固定 $\Omega$ 的一种选取，那时就必须留意符号。我们对归一化 Yukawa 耦合的定义将包含这个负号。}注意，$\langle\theta_1,\theta_2,\theta_3\rangle$ 此时尚未得到良好定义，因为 $\Omega$ 可以乘以任意常数。不过，存在一种自然的归一化方法，后文将予介绍。显然，这个 Yukawa 耦合不依赖于复化 Kähler 类 $\omega$，却依赖于 $V$ 的复结构，因为 $\Omega$ 是全纯三形式。

式~\eqref{eq:b-yukawa} 中的 Yukawa 耦合有时称为\emph{B 模型关联函数}，因为它与另一种扭曲理论——\hyperref[sec:nonlinear-sigma-models]{附录 B.2} 中介绍的 B 模型——里相应的三点函数相同。

要解释式~\eqref{eq:b-yukawa} 的来历，我们仿照 A 模型的情形，从“杯积”
\begin{equation}\label{eq:b-cup}
\int_V\Omega\wedge(\theta_1\cdot\theta_2\cdot\theta_3\cdot\Omega)
\end{equation}
出发。这里同样适用\IndexLocation{indexloc:0944:31}非重整化定理，也有与前面相同的全纯瞬子。关键在于，瞬子通过杯积发挥作用。在 A 模型中，$\omega_i$ 经由 $\int_\beta\omega_i$ 出现在式~\eqref{eq:a-yukawa} 中；这可看成杯积 $g_\beta\smile\omega_i\in H^6(V,\mathbb C)\simeq\mathbb C$，其中 $g_\beta\in H^4(V,\mathbb C)=H^{2,2}(V)$ 是 $\beta$ 的 Poincaré 对偶。而在 B 模型中，$\theta_i\in H^1(V,T_V)\simeq H^1(V,\Omega_V^2)=H^{2,1}(V)$，且 $g_\beta\in H^{2,2}(V)$；因此它们的杯积 $g_\beta\smile\theta_i$ 属于 $H^{4,3}(V)=0$。于是，瞬子不会与式~\eqref{eq:b-cup} 相互作用。这就是式~\eqref{eq:b-cup} 等于 B 模型关联函数~\eqref{eq:b-yukawa} 的一个粗略理由。

\SourcePage{32}
式~\eqref{eq:b-yukawa} 中没有瞬子修正，这一点极为重要。它意味着我们能够用代数几何的标准方法，精确计算 B 模型的 Yukawa 耦合。\hyperref[chap:hodge-yukawa]{第 5 章}将解释具体做法。

现在设 $(V,\omega)$ 与 $(V^\circ,\omega^\circ)$ 构成镜像对，也就是说，与 $(V,\omega)$ 相联系的 $\sigma$ 模型，在适当改变符号后，给出的 SCFT 与 $(V^\circ,\omega^\circ)$ 的 $\sigma$ 模型相同。于是自然得到同构 $H^1(V,\Omega_V^1)\simeq H^1(V^\circ,T_{V^\circ})$。又因为三点函数是 SCFT 的内在对象，应当可以推出：来自 $(V,\omega)$ 上 $\sigma$ 模型的 A 模型关联函数，与来自 $(V^\circ,\omega^\circ)$ 上 $\sigma$ 模型的 B 模型关联函数，在适当归一化后相等。这是镜像对称的一个主要数学后果。由于镜像映射也参与其中，实际的等同细节要更复杂一些；但 A 模型与 B 模型关联函数相等，显然是合理的预期。

关联函数的性质也与镜像对称相吻合。我们已经看到，与 $(V,\omega)$ 相联系的 A 模型关联函数依赖于 $\omega$，却不依赖于 $V$。根据镜像对称以及上面讨论的模空间局部对应，与 $(V^\circ,\omega^\circ)$ 相联系的 B 模型关联函数应当依赖于 $V^\circ$，却不依赖于 $\omega^\circ$。如前所述，B 模型确实具有这一性质。

在五次三维簇的例子中我们将看到，\SourceCite{CdGP} 正是利用 A 模型的三点函数与其镜像的 B 模型三点函数相等，才对五次三维簇上的有理曲线数目作出了一些非凡预测。下一章将详细讨论这个例子。一般而言，镜像对称的一些最重要的后果，来自下面两个事实的结合：
\begin{itemize}
\item A 模型与 B 模型的关联函数相等；
\item B 模型的函数可以精确计算。
\end{itemize}
二者合起来，就能计算 Calabi--Yau 三维簇上的 Gromov--Witten 不变量，并由此得到丰富的枚举信息。当然，这一切的前提是我们能够用数学证明：这些物理后果确实成立。

刚才介绍的途径可以看作研究镜像对称后果的“经典”方法，其中 A 模型与 B 模型的关联函数是主要对象。不过，在镜像对称发现后的那些年里，研究重点略有转移。特别是，人们对 B 模型关联函数与 Gauss--Manin 联络之间的关系理解得更深了；在 A 模型一侧，这推动了量子上同调及 A 型 Hodge 结构变分的发展。Givental 以及 Lian、Liu、Yau 关于镜像定理的工作，又引入了其他值得研究的新对象。我们将看到，等变上同调与局部化在镜像定理的证明中起着重要作用。后面的章节会讨论所有这些思想。

最后，再谈谈式~\eqref{eq:a-yukawa} 与~\eqref{eq:b-yukawa}。二者看似不对称，因为前者远比后者复杂。好在镜像对称很容易解释这种差异。上面说过，A 模型关联函数应随复化 Kähler 类 $\omega$ 变化，而 B 模型关联函数应随 $V$ 的复结构变化。式~\eqref{eq:a-cup} 中的杯积是纯拓扑不变量；所以，若要得到对 $\omega$ 的非平凡依赖，显然还需要其他项。另一方面，
\SourcePage{33}
式~\eqref{eq:b-cup} 中的杯积已经依赖于复结构，因为 $\Omega$ 是 $V$ 上的全纯三形式。因此，不再需要额外的项也是可能的。

\begingroup\linespread{1.22}\selectfont
\section{为什么是 Calabi--Yau 流形？}\label{sec:calabi-yau-physics}

接下来的任务，是解释镜像对称为什么只适用于 Calabi--Yau 流形，尤其是 Calabi--Yau 三维簇。要完整说明非线性 $\sigma$ 模型为什么需要 Calabi--Yau 流形，我们必须大幅转入物理学。\hyperref[sec:nonlinear-sigma-models]{附录 B.2} 和 \hyperref[sec:conformal-field-theories]{B.3} 将对此作部分介绍。这里先概述其中一些想法。出发点是假设时空应当是一个带有号型为 $(9,1)$ 的半 Riemann 度量的十维流形。这个流形局部应为乘积 $M_{3,1}\times V$，其中 $M_{3,1}$ 是狭义相对论中的通常时空，$V$ 是一个紧致的六维 Riemann 流形。基本直观是：$V$ 太小，以致在宏观尺度上看不见，但它对理论的量子方面却不可或缺。若时空维数不是十维，本节前面讨论的共形对称性就无法在量子化过程中保持下来。

如往常一样，其中的物理学并不简单，还涉及一些陌生术语。思路是先考察从物理学观点看更为基本的其他理论，再由此接近非线性 $\sigma$ 模型。具体来说，从低能极限下的一个 $N=1$ 超引力理论出发，先去掉费米场；随后，超对称和其他相容性要求对 Riemann 流形 $V$ 上的度量施加强约束。特别地，可得到该度量的和乐群等于 $SU(3)$。这有一些很好的后果：
\begin{itemize}
\item \textbf{微分几何方面：}$V$ 上的度量是 Ricci 平坦的，即其 Ricci 曲率张量处处为零。这蕴含 $b_1(V)=0$。
\item \textbf{代数几何方面：}$V$ 具有一个复结构，使得 $c_1(V)=0$，而且该度量对于这一复结构是 Kähler 的。
\end{itemize}
由此可见，$V$ 正是我们所需要的 Calabi--Yau 流形。

到目前为止，在时空中我们只有 $N=1$ 超对称。经过进一步处理，这个理论可以重新解释为世界面上的 $N=2$ SCFT；不过，为了保持超共形不变性，还须形变上述 Ricci 平坦度量。我们尚不知道这种形变的精确描述——这仍是物理学中的一个开放问题——但可以证明新度量与旧度量属于同一个 Kähler 类。因此，得到的仍然是一个 Calabi--Yau 三维簇。

如果读者想更完整地了解从物理学观点看 Ricci 曲率如何出现，并查阅相关文献，可参阅 \SourceCite{Huebsch} 的第 0 章。

\section{镜像对称的数学}\label{sec:math-preview}

从数学的观点看，第 \ref{sec:physics} 节给出的镜像对称表述存在严重问题。例如，$N=2$ SCFT 的定义涉及对所有映射 $\Sigma\to V$ 构成的空间积分，而这类积分尚未得到严格定义。因此，镜像对称的数学证明似乎要求在两个数学上尚未得到良好定义的对象之间——即 $(V,\omega)$ 与 $(V^\circ,\omega^\circ)$ 的 $\sigma$ 模型之间——建立同构。甚至第 \ref{sec:physics} 节示意图里的 $N=2$ SCFT 模空间也没有得到良好定义！

\begingroup\emergencystretch=1em
面对这种情况，数学家该怎么办？一种办法是完全避开 SCFT，转而严谨定义 Kähler 模空间、\SourcePage{34}Gromov--Witten 不变量等对象，再证明它们的行为符合镜像对称的预测。另一种办法是接纳物理学及其直观，以更深入地洞察代数几何，特别是镜像对称对 Calabi--Yau 环面完全交的 Gromov--Witten 不变量所作的预测。第三种办法，是表述一个纯数学版本的镜像对称。例如，\SourceCite{Kontsevich3} 提议：复流形 $V$ 的镜像是某个辛流形 $V^\circ$；镜像对称或可表述为导出范畴的等价，将 $V$ 上的凝聚层与由 $V^\circ$ 的 Lagrange 子流形构成的范畴联系起来。还有一种同样迷人、但目前仍带有一定猜测性质的办法：尝试把镜像流形几何地\emph{构造}为特殊 Lagrange 子流形的模空间 \SourceCite{SYZ}。\footnote{一般而言，辛几何在镜像对称中扮演重要角色，不过本书将更多地关注代数几何方面。}实际上，所有这些办法都得到了应用，这正是镜像对称如此令人兴奋的原因。\par\endgroup
本书将主要采用第一种办法，同时不时对相关物理学加以评论。因此，我们的目标，是讨论理解第 \ref{sec:physics}、\ref{sec:three-point} 节的数学方面所涉及的代数几何。与物理理论不同，我们将研究任意维数的 Calabi--Yau 流形，而非仅限于三维。然而，到\hyperref[chap:quantum-cohomology]{第~\ref*{chap:quantum-cohomology} 章}尝试给镜像对一个数学定义时，我们在 Calabi--Yau 三维簇的情形下将取得最好的结果。

先从严格定义 Calabi--Yau 开始。由于\hyperref[chap:mirror-constructions]{第 4 章}将研究的镜像对称构造有时会产生奇异簇，我们需要一个允许某些类型奇点的定义。

\begin{definition}\label{def:calabi-yau-variety}
一个 $d$ 维正规紧致簇 $V$ 称为\emph{Calabi--Yau 簇}，如果它满足：
\begin{enumerate}[(i)]
\item $V$ 至多具有 Gorenstein 典范奇点；
\item $V$ 的对偶化层平凡，即 $\widehat\Omega_V^d\simeq\mathcal O_V$；
\item $H^1(V,\mathcal O_V)=\cdots=H^{d-1}(V,\mathcal O_V)=\{0\}$。
\end{enumerate}
如果 $V$ 还至多具有 Gorenstein、$\mathbb Q$ 析因的终端奇点，则称 $V$ 为\emph{极小 Calabi--Yau 簇}。
\end{definition}
\hyperref[app:singular-varieties]{附录 A} 将回顾对偶化层 $\widehat\Omega_V^d$，以及 Gorenstein 奇点、典范奇点和终端奇点的定义。第 \ref{sec:physics} 节考虑的 Calabi--Yau 三维簇显然符合这里的定义。

第 \ref{sec:physics} 节用到的其他一些术语也与各种模有关。给定流形 $V$ 上所有复结构构成的空间，是代数几何中熟知的对象；Kähler 模的概念则可能不那么熟悉。回想第 \ref{sec:physics} 节：我们有一个 Calabi--Yau 三维簇 $V$，以及一个复化 Kähler 类 $\omega=B+iJ\in H^2(V,\mathbb C)$，其中 $J$ 是 Kähler 类。不过，如果用一个整上同调类改变 $\omega$，物理理论不会改变，因为非线性 $\sigma$ 模型的定义对映射 $\phi:\Sigma\to V$ 只用到 $\exp\bigl(2\pi i\int_\Sigma\phi^*(\omega)\bigr)$。\footnote{从 A 模型关联函数~\eqref{eq:a-yukawa} 中也可以看出这一点。完整细节见\hyperref[sec:nonlinear-sigma-models]{附录 B.2}。}当 $\omega$ 改变一个 $H^2(V,\mathbb Z)$ 中的元素时，这个量保持不变。因此，在定义 Kähler 模时，应当对 $H^2(V,\mathbb Z)$ 的像取商。由此得到下面的定义；和前面一样，它也允许某些奇点。

\SourcePage{35}
\begin{definition}\label{def:kahler-moduli}
设 $V$ 是一个射影轨形，且 $h^{2,0}(V)=0$。那么：
\begin{enumerate}[(i)]
\item \emph{Kähler 锥} $K(V)\subset H^2(V,\mathbb R)=H^{1,1}(V,\mathbb R)$，是所有 Kähler 类构成的子集；
\item \emph{复化 Kähler 空间}是商空间
\[
K_{\mathbb C}(V)=\{\omega\in H^2(V,\mathbb C):\operatorname{Im}(\omega)\in K(V)\}/\operatorname{im}H^2(V,\mathbb Z),
\]
其中 $\operatorname{im}H^2(V,\mathbb Z)$ 是自然映射 $H^2(V,\mathbb Z)\to H^2(V,\mathbb C)$ 的像；
\item \emph{复化 Kähler 模空间}是商空间 $K_{\mathbb C}(V)/\operatorname{Aut}(V)$。
\end{enumerate}
\end{definition}
事实上，在 SCFT 中出现的 Kähler 模空间\footnote{讨论 Kähler 模空间时，只要上下文足以表明含义，我们常省去“复化”这个限定语。}与这里的定义略有不同。我们定义的 Kähler 模空间会受到量子修正，其性质因此会发生细微变化。特别是，理论可能并非对所有 $\omega$ 都收敛，所以我们不得不把注意力限制在 $\omega=B+iJ$ 中 $J$ 充分正的复化 Kähler 类上。尽管如此，这里定义的较大空间在数学上仍然有趣，也是我们的主要研究对象之一。\hyperref[chap:moduli]{第~\ref*{chap:moduli} 章}将讨论其中一些微妙之处。

\hyperref[app:singular-varieties]{附录 A} 将回顾轨形及轨形上 Kähler 类的定义。因为 $h^{2,0}=0$，Kähler 锥 $K(V)$ 是 $H^2(V,\mathbb R)$ 中的开凸锥。这说明复化 Kähler 空间 $K_{\mathbb C}(V)$ 是一个性质良好的对象。但是，若要确定 Kähler 模空间 $K_{\mathbb C}(V)/\operatorname{Aut}(V)$ 的结构，还须知道\IndexLocation{indexloc:0088:35}自同构群 $\operatorname{Aut}(V)$ 如何作用于 Kähler 锥。\hyperref[chap:moduli]{第~\ref*{chap:moduli} 章}将回到这一主题。

注意，定义~\ref{def:calabi-yau-variety} 与定义~\ref{def:kahler-moduli} 似乎涉及略微不同的奇点类型。幸运的是，\SourceCite{Reid1} 的一个结果表明，Gorenstein 轨形具有\IndexLocation{indexloc:1259:35}典范奇点。这意味着，凡满足定义~\ref{def:calabi-yau-variety} 第 (ii)、(iii) 项的轨形，都自动是 Calabi--Yau 簇。因此，对于镜像对称，Calabi--Yau 轨形是一类自然的研究对象。

\section{后续内容}\label{sec:whats-next}

下面介绍接下来的十章将如何把我们从这里引向镜像定理的证明。\hyperref[chap:quintic]{第 2 章}首先仔细描述光滑五次三维簇 $V\subset\mathbb P^4$ 的镜像。沿着第 \ref{sec:three-point} 节概述的策略，我们将对 $V$ 上给定次数的有理曲线数目作出一些明确预测。在后续章节中，随着数学背景逐步建立，我们还会多次回到这个例子，并最终论证\hyperref[chap:quintic]{第 2 章}中出现的全部计算。

人们很快意识到，要推广五次三维簇的例子，环面几何有重要作用。因此，\hyperref[chap:toric]{第 3 章}将探讨描述环面簇的各种方式，\hyperref[chap:mirror-constructions]{第 4 章}再介绍已知的镜像对称构造，其中许多构造都用到环面几何。这会为我们提供大量预期具有镜像对称的例子，成为检验数学理论的良好场所。还应提到，有些物理理论也明确使用环面簇，例如\hyperref[sec:gauged-linear-sigma-models]{附录 B.5} 将介绍的 \SourceCite{Witten5} 中的\IndexLocation{indexloc:0354:35}规范线性 $\sigma$ 模型。\SourceCite{MP1} 提供了在这种模型背景下支持镜像对称的证据。

下一个任务是描述并计算 B 模型的 Yukawa 耦合，这将在\hyperref[chap:hodge-yukawa]{第 5 章}完成。事实表明，为了找到计算这一耦合所需的正确坐标，
\SourcePage{36}
计算这个 Yukawa 耦合，需要充分理解模空间在某些边界点附近的情形。因此，\hyperref[chap:moduli]{第~\ref*{chap:moduli} 章}将研究这些模空间的结构及其紧化。正如前面提到的，我们也会考察 Kähler 模。事实上，Kähler 模的若干基本事实，会帮助我们理解通常的模空间的紧化应当是什么样子。\hyperref[chap:gromov-witten]{第~\ref*{chap:gromov-witten} 章}将讨论 Gromov--Witten 不变量。辛几何范畴和代数范畴中都已有相应定义；任一种定义都可用于从数学上定义 A 模型的 Yukawa 耦合。

到\hyperref[chap:gromov-witten]{第~\ref*{chap:gromov-witten} 章}末，我们将具备精确表述“经典”镜像对称所需的一切：它断言某些关联函数在镜像映射下彼此相容。不过，从\hyperref[chap:quantum-cohomology]{第~\ref*{chap:quantum-cohomology} 章}起，我们将探索对这个主题更深层的理解。新引入的两个要素是量子上同调和 A 型 Hodge 结构变分。前者可视为对 Gromov--Witten 不变量的代数与枚举后果的系统展开；后者则自然地产生于量子上同调。基本想法是：镜像对称真正涉及两个 Hodge 结构变分在镜像映射下的同构——一个来自 B 模型，即由复模产生的通常 Hodge 结构变分；另一个来自 A 模型，即 A 型 Hodge 结构变分。\hyperref[chap:quantum-cohomology]{第~\ref*{chap:quantum-cohomology} 章}末将以这种形式表述镜像定理。所要求的关联函数相等，将直接从 Hodge 结构变分的这一同构推出。

用 Hodge 结构变分表述的镜像定理仍是一个开放问题；不过，Givental \SourceCite{Givental2,Givental4} 以及 Lian、Liu、Yau \SourceCite{LLY} 的近期工作，已朝证明这种形式的定理迈出实质性的一步。\hyperref[chap:mirror-theorem]{第~\ref*{chap:mirror-theorem} 章}将讨论这些论文的思想。我们将看到，这需要引入新的技术和新的研究对象，尤其是等变上同调与局部化，它们在证明中扮演重要角色。\hyperref[chap:localization]{第~\ref*{chap:localization} 章}回顾若干基本定义和定理，并运用这些方法证明一些关于 Gromov--Witten 不变量的有趣结果。随后，\hyperref[chap:quantum-differential-equations]{第~\ref*{chap:quantum-differential-equations} 章}讨论 Gromov--Witten 不变量的一种推广，即引力关联数。这些不变量使我们能够描述 A 模型联络的平坦截面。我们还将定义 Givental 的 $J$ 函数，并解释它与量子微分方程的关系。最后，\hyperref[chap:mirror-theorem]{第~\ref*{chap:mirror-theorem} 章}讨论 \SourceCite{Givental2,Givental4} 与 \SourceCite{LLY} 所陈述的镜像定理。\hyperref[chap:quintic]{第 2 章}结尾会简短预览其中涉及的内容，\hyperref[chap:mirror-theorem]{第~\ref*{chap:mirror-theorem} 章}再给出五次三维簇情形的完整细节。

这些章节涉及大量数学内容，从头到尾逐章读下去并非轻松之事。我们建议读者先从\hyperref[chap:quintic]{第 2 章}开始。此后，可依个人兴趣选择几条不同的路线。前言给出了阅读指引，浏览目录也会有所帮助。我们希望，一旦读者开始阅读，这个主题本身的吸引力便会引领读者继续前行。镜像对称的故事引人入胜，而看清各个部分如何拼合起来，也是一件有趣的事。

\par\endgroup

\chapter{五次三维簇}\label{chap:quintic}
\SourcePage{38}

五次三维簇是人们首次运用镜像对称作出枚举预测的例子 \SourceCite{CdGP}。本章将回顾这一例子，基本沿用 \SourceCite{CdGP} 的方法，同时根据 \SourceCite{Morrison2} 更偏重数学的论述作一些调整。第 \ref{sec:quintic-other-methods} 节将依据 \SourceCite{Givental2} 与 \SourceCite{LLY}，介绍计算同一组枚举数据的另外两种方法。在后面的章节，我们还会多次回到五次三维簇，用它来说明研究镜像对称的不同方法。

\section{五次三维簇的 A 模型关联函数}\label{sec:quintic-a-model}

众所周知，光滑五次超曲面 $V\subset\mathbb P^4$ 是 Calabi--Yau 簇；\hyperref[chap:mirror-constructions]{第 4 章}将建立的理论也会推出这一点。我们考虑与 $(V,\omega)$ 相联系的非线性 $\sigma$ 模型，其中 $\omega=B+iJ$ 是 $V$ 上的复化 Kähler 类。若以 $H$ 表示超平面类，则\hyperref[chap:introduction]{第 1 章}中的 A 模型关联函数~\eqref{eq:a-yukawa} 简化为
\begin{equation}\label{eq:quintic-a}
\langle H,H,H\rangle=5+\sum_{d=1}^{\infty}n_d d^3\frac{q^d}{1-q^d}.
\end{equation}
粗略地说，$n_d$ 是 $V$ 中次数为 $d$ 的有理曲线的数目，而 $q=e^{2\pi i\int_\ell\omega}$，$\ell$ 是 $V$ 中的一条直线。若写 $\omega=tH$，则 $q$ 也可表示为 $e^{2\pi it}$。下一节将看到，$q$ 实际上是紧化 Kähler 模空间上的一个局部坐标。另外，正如\hyperref[chap:introduction]{第 1 章}所解释的，由于这是 A 模型关联函数，式~\eqref{eq:quintic-a} 依赖于 $\omega$，却不依赖于 $V$ 的复结构。

代数几何学家初次见到式~\eqref{eq:quintic-a} 时，觉得它有些奇怪。其中一个原因是，我们甚至还不知道一般五次三维簇对给定次数 $d$ 是否只含有限多条有理曲线。按照 Clemens 猜想，答案应当是肯定的；但迄今只在 $d\leq9$ 的情形下知道该猜想成立 \SourceCite{JK1}。式~\eqref{eq:quintic-a} 却似乎预设它对所有 $d$ 都成立。更令人不安的是，镜像对称竟声称能给出这些曲线的数目 $n_d$ 的显式公式，而且对所有 $d\geq1$ 都适用。\index{Clemens猜想@Clemens 猜想|hyperpage}

\hyperref[chap:gromov-witten]{第~\ref*{chap:gromov-witten} 章}将为式~\eqref{eq:quintic-a} 建立牢固的基础。关键对象是 Gromov--Witten 不变量 $\langle I_{0,3,d}\rangle(H,H,H)$。这个不变量有多种定义方式。从辛几何的观点出发，可以考虑满足适当形变后的 Cauchy--Riemann 方程的 $C^\infty$ 映射 $f:\mathbb P^1\to V$。对 $\mathbb P^1$ 上三个点的像施加合适的条件后，可以证明，在每个同调类中这样的曲线都只有有限多条。于是，$\langle I_{0,3,d}\rangle(H,H,H)$ 被定义为这些曲线按定向所赋符号计数后的总数；这里已不再有来自复结构的优选定向。

\SourcePage{39}
事实上，要把这一想法严格化，有几种不同的做法；\hyperref[sec:gw-symplectic-definition]{第~\ref*{sec:gw-symplectic-definition} 节}将逐一回顾。

\SourceCite{Kontsevich2} 提出了定义 $n_d$ 的另一条途径：使用稳定映射 $f:C\to V$ 的模空间，其中 $C$ 的亏格为零，但不必不可约。由此得到的 Gromov--Witten 不变量定义相当精巧，涉及代数叠和虚拟基本类。\hyperref[sec:gw-algebraic-definition]{第~\ref*{sec:gw-algebraic-definition} 节}将讨论这些主题。

一旦 Gromov--Witten 不变量得到严格定义，我们便可把 A 模型关联函数定义为形式和
\[
\langle H,H,H\rangle=\sum_{d=0}^{\infty}\langle I_{0,3,d}\rangle(H,H,H)q^d.
\]
为弄清它与式~\eqref{eq:quintic-a} 的关系，固定 $d$，设 $k\mid d$。若 $C\subset V$ 是次数为 $k$、法丛为 $\mathcal O_{\mathbb P^1}(-1)\oplus\mathcal O_{\mathbb P^1}(-1)$ 的光滑有理曲线，那么次数为 $d/k$ 的映射 $f:\mathbb P^1\to C$ 就给出被 $\langle I_{0,3,d}\rangle(H,H,H)$ 计入的一类稳定映射。\hyperref[chap:gromov-witten]{第~\ref*{chap:gromov-witten} 章}将看到，这些 $f$ 所组成的族对该不变量贡献一个 $k^3$ 因子。由于每一条次数为 $k\mid d$ 的曲线 $C$ 都如此，因而有理由期待
\begin{equation}\label{eq:multi-cover}
\langle I_{0,3,d}\rangle(H,H,H)=\sum_{k\mid d}n_k k^3,
\end{equation}
其中 $d\geq1$。由这个等式和上面对 $\langle H,H,H\rangle$ 的定义，经过简单的幂级数运算就能得到式~\eqref{eq:quintic-a}。

遗憾的是，这套论证远不完整，而且似乎预设了 Clemens 猜想。为使一切严格化，我们反过来用式~\eqref{eq:multi-cover} \emph{定义} $n_d$。这样，问题就变成了：$n_d$ 与 $V$ 中的有理曲线究竟有何关系？\hyperref[subsec:instanton-numbers]{第~\ref*{subsec:instanton-numbers} 小节}将证明，当 $d\leq9$ 时，$n_d$ 恰好是五次三维簇上次数为 $d$ 的有理曲线数目。这一点并不像表面看来那样显然。虽然 Clemens 猜想在 $d\leq9$ 时成立，但 $V$ 中有些有理曲线是奇异的。特别地，$V$ 含有 $17{,}601{,}000$ 条次数为 $5$、具有六个结点的有理曲线 \SourceCite{Vainsencher}；要弄清它们如何计入 $n_5$，还须仔细分析。

长期以来，人们以为一旦 Clemens 猜想得到证明，就能立即推出式~\eqref{eq:multi-cover} 中的 $n_d$ 是 $V$ 上次数为 $d$ 的有理曲线数目。然而，R. Pandharipande 最近注意到，即使 Clemens 猜想对 $d=10$ 成立，式~\eqref{eq:multi-cover} 定义的 $n_{10}$ \emph{也不是} $V$ 上十次有理曲线的数目。原因在于，上述六结点五次有理曲线的二重覆盖，所贡献的量超过预期的 $5^3$ 因子。第 \hyperref[chap:gromov-witten]{\ref*{chap:gromov-witten}}、\hyperref[chap:mirror-theorem]{\ref*{chap:mirror-theorem}} 章将对此作严格证明。因此，我们不再把 $n_d$ 称为“次数为 $d$ 的有理曲线的数目”，而采用\hyperref[chap:gromov-witten]{第~\ref*{chap:gromov-witten} 章}中较为中性的名称：\emph{瞬子数}。

对 $n_{10}$ 的讨论体现了\hyperref[chap:gromov-witten]{第~\ref*{chap:gromov-witten} 章}的一个主题：Gromov--Witten 不变量与\IndexLocation{indexloc:0241:39}枚举几何之间存在微妙的相互关系。尽管会遇到这样的困难，\hyperref[chap:gromov-witten]{第~\ref*{chap:gromov-witten} 章}仍将表明，Gromov--Witten 不变量包含大量枚举信息，而镜像对称也推动了枚举代数几何中的一股研究热潮。

本章将依照\hyperref[chap:introduction]{第 1 章}概述的策略，用镜像对称计算五次三维簇 $V$ 的 A 模型关联函数~\eqref{eq:quintic-a}。接下来的几节有一些十分精彩的计算，最终会给出对 $n_d$ 的明确预测。由于我们是用尚不严格的物理理论，来计算枚举意义尚不明朗的数值，这样的讨论难免令人感到不够满意。
\SourcePage{40}
不过，后面的章节将说明，计算中的每一步都可以严格化。特别地，\hyperref[chap:mirror-theorem]{第~\ref*{chap:mirror-theorem} 章}将证明：若用式~\eqref{eq:multi-cover} 以及 $\langle I_{0,3,d}\rangle(H,H,H)$ 的辛几何定义或代数定义来确定 $n_d$，那么镜像对称对 $n_d$ 的预测，便是镜像定理的直接推论。

\section{五次三维簇的镜像}\label{sec:quintic-mirror}

五次三维簇 $V$ 满足 $h^{1,1}(V)=1$，因此它的镜像应满足 $\dim H^1(V^\circ,T_{V^\circ})=1$。所以，五次三维簇的镜像是一个单参数的 Calabi--Yau 流形族。Greene 和 Plesser \SourceCite{GPl} 把它描述为 $\mathbb P^4/G$ 中某个超曲面族的解奇异化。这里，$G$ 是群
\begin{equation}\label{eq:greene-plesser-group}
G=\bigl\{(a_1,\ldots,a_5)\in\mathbb Z_5^5:\textstyle\sum_i a_i\equiv0\pmod5\bigr\}/\mathbb Z_5,
\end{equation}
其中分母中的 $\mathbb Z_5$ 以对角方式嵌入。元素 $g=(a_1,\ldots,a_5)\in G$ 在 $\mathbb P^4$ 上的作用为
\[
g\cdot(x_1,\ldots,x_5)=(\mu^{a_1}x_1,\ldots,\mu^{a_5}x_5),
\]
其中 $\mu=e^{2\pi i/5}$ 是一个本原五次单位根。

以 $\psi$ 为参数，这个超曲面族由方程
\begin{equation}\label{eq:quintic-pencil}
x_1^5+x_2^5+x_3^5+x_4^5+x_5^5+\psi x_1x_2x_3x_4x_5=0
\end{equation}
定义。该方程在 $G$ 的作用下不变，因而在 $\mathbb P^4/G$ 中定义了一个超曲面族。（\SourceCite{CdGP} 用 $-5\psi$ 代替这里的 $\psi$。）商空间 $\mathbb P^4/G$ 是奇异的，式~\eqref{eq:quintic-pencil} 给出的超曲面也继承了这些奇点。不过，只要 $\psi\ne-5\mu^i$，$0\leq i\leq4$，该超曲面就是 $\mathbb P^4/G$ 中的 Calabi--Yau 轨形。更好的是，我们可以同时对所有 $\psi$ 解奇异化，而不破坏 Calabi--Yau 性质；这一点可在局部直接验证。因此，我们把由式~\eqref{eq:quintic-pencil} 得到的解奇异化超曲面族定义为 $V^\circ$。为了强调这是一个单参数簇族，有时记作 $V_\psi^\circ$。

$V_\psi^\circ$ 的这一构造归功于 Greene 和 Plesser，它涉及\hyperref[chap:introduction]{第 1 章}讨论的超共形场论。不过，\hyperref[chap:mirror-constructions]{第 4 章}将说明，Batyrev 对 $V_\psi^\circ$ 给出了纯数学的描述，并将其推广到环面簇中的 Calabi--Yau 超曲面。

如上所述，式~\eqref{eq:quintic-pencil} 定义了奇异簇。由于我们研究的是三维簇，对它的解奇异化并无典范选择。\SourceCite{Morrison2} 给出了一个明确的选择，但还有许多其他选择。幸运的是，不同选择的复模相同，B \IndexLocation{indexloc:1167:40}模型关联函数也相同。我们将用 $V_\psi^\circ$ 指代其中任一选择。（\hyperref[chap:moduli]{第~\ref*{chap:moduli} 章}将看到，这些不同选择反映了镜像的 Kähler 模；在物理学中，它们对应同一个物理理论的不同“相”。）

接下来的任务，是充实第 \ref{sec:three-point} 节讨论的策略，用镜像 $V_\psi^\circ$ 来计算式~\eqref{eq:quintic-a}。第一步是理解有关模空间。对于镜像的复模，注意映射 $(x_1,\ldots,x_5)\mapsto(\mu^{-1}x_1,x_2,\ldots,x_5)$ 诱导同构 $V_\psi^\circ\simeq V_{\mu\psi}^\circ$。\footnote{这一点对于 $\mathbb P^4/G$ 中的超曲面是显然的；还可以证明该同构延伸到解奇异化，见 \SourceCite{Morrison2}。}因此，$\psi^5$ 在 $V_\psi^\circ$ 的复结构模空间上得到良好定义，
\SourcePage{41}
于是
\[
x=\psi^{-5}
\]
是复模的一个局部坐标；选择 $x$ 的原因很快就会明白。此外，可验证 $V_\psi^\circ$ 在 $\psi=-5\mu^i$（$0\leq i\leq4$）以及 $\psi=\infty$ 时发生奇异。用模坐标 $x$ 表示，这意味着复模空间具有边界点 $x=-5^{-5}$ 和 $x=0$。

根据镜像对称，$V^\circ$ 的复模应对应五次三维簇 $V$ 的 Kähler 模。由于 $H$ 生成 $H^2(V,\mathbb Z)\simeq\mathbb Z$，我们可以写 $\omega=tH$。又因为 $H$ 是 Kähler 的，当且仅当 $t$ 位于上半平面时，$\omega=tH$ 是复化 Kähler 类。依照定义~\ref{def:kahler-moduli}，Kähler 空间 $K_{\mathbb C}(V)$ 是复化 Kähler 类对 $H^2(V,\mathbb Z)$ 取商。因此，$q=e^{2\pi it}$ 诱导一个同构
\[
q:K_{\mathbb C}(V)\simeq\Delta^*.
\]
所以，$q$ 是 Kähler 模空间的局部参数，$q=0$ 是一个边界点。一个重要而容易验证的事实是：若 $\ell$ 是 $V$ 中的一条直线，则 $q$ 可写成
\begin{equation}\label{eq:q-kahler}
q=e^{2\pi it}=e^{2\pi i\int_\ell\omega}.
\end{equation}
因此，A 模型关联函数~\eqref{eq:quintic-a} 自然是 Kähler 模局部参数 $q$ 的函数。

假定镜像对称成立，$V$ 的 Kähler 模与 $V^\circ$ 的复模之间便有一个局部同构。参照图式~\eqref{eq:mirror-foliations}，“$V$ 不变、$\omega$ 变化”的截片，局部对应于镜像一侧“$\omega^\circ$ 不变、$V^\circ$ 变化”的截片。然而，关键在于“局部”二字：我们需要给两个模空间各选合适的局部坐标。对 $V$ 的 Kähler 模而言，由于三点函数~\eqref{eq:quintic-a} 是 $q$ 的幂级数，所需坐标显然是 $q$。既然 $q$ 是一个边界点处的局部坐标，镜像一侧也应如此。

在 $V^\circ$ 的复模空间中，可以从两个边界点 $x=0$ 与 $x=-5^{-5}$ 里选择。我们断言 $x=0$ 对应 $q=0$，所以 $x$ 就是需要的局部坐标。理由如下。回到五次三维簇 $V$：与 $H$ 作杯积给出 $\bigoplus_{p=0}^3H^{p,p}(V)$ 上的自同态 $\smile H$，它是最大幂零的，即 $(\smile H)^3\ne0$ 而 $(\smile H)^4=0$。第 \hyperref[chap:hodge-yukawa]{\ref*{chap:hodge-yukawa}}、\hyperref[chap:moduli]{\ref*{chap:moduli}} 章将说明，在镜像对称下，我们预期 $\smile H$ 对应于 $V^\circ$ 的复模空间上、与 $q=0$ 相应的点周围的单值性之对数。因此，镜像对称告诉我们要寻找\emph{最大幺幂单值性}。在上述边界点中，只有 $x=0$ 具有这一性质。\hyperref[chap:hodge-yukawa]{第 5 章}将严格定义最大幺幂单值性。

为了计算 A 模型关联函数~\eqref{eq:quintic-a}，还剩下两项任务：
\begin{itemize}
\item 局部坐标 $q$ 和 $x$ 描述 Kähler 模空间与复模空间中相应的边界点，但镜像对称给出的局部同构\emph{不是} $q=x$。从物理学的观点看，还要考虑量子修正。因此，需要把 $x$ 表示为 $q$ 的函数，反过来也一样。完成后，便得到 $V^\circ$ 的复模空间与 $V$ 的 Kähler 模空间之间的显式局部映射；这通常称为\emph{镜像映射}。

\SourcePage{42}
\item 我们还必须计算 B 模型关联函数并正确地将它归一化，然后用镜像映射把它写成 $q$ 的函数。这样就会得到式~\eqref{eq:quintic-a} 的显式公式。
\end{itemize}
现在先着手第一项任务。

\section{镜像映射}\label{sec:quintic-mirror-map}

理解镜像映射的一种方式是：$q$ 是 $V$ 的 Kähler 模空间上典范选出的坐标。因此，我们希望在 $V^\circ$ 的模空间上找到一个“典范”坐标。\hyperref[chap:moduli]{第~\ref*{chap:moduli} 章}将解释，之所以能这样做，是因为 $x=0$ 处的单值性是\IndexLocation{indexloc:0772:42}最大幺幂的。按照 \SourceCite{Morrison2}，在 $x=0$ 附近存在一个极小整消失圈 $\gamma_0$，它在单值性作用下保持不变。此外，还有一个极小整圈 $\gamma_1$，绕 $x=0$ 的单值性将其变为 $\gamma_1+m\gamma_0$，其中 $m\in\mathbb Z$。对\IndexLocation{indexloc:1181:42}\IndexLocation{indexloc:1195:42}五次三维簇，$m=1$ \SourceCite{Morrison2}。因此，若 $\Omega$ 是 $V^\circ$ 上的全纯三形式，$\int_{\gamma_1}\Omega/\int_{\gamma_0}\Omega$ 在单值性作用下变为它自身加 $1$，从而
\[
\exp\!\left(2\pi i\frac{\int_{\gamma_1}\Omega}{\int_{\gamma_0}\Omega}\right)
\]
是 $x=0$ 邻域中一个得到良好定义的函数。此外，$\gamma_0$ 除符号外唯一；$\gamma_1$ 则在取相同符号以及加上 $\gamma_0$ 的整数倍的意义下唯一。因此，上式在 $V^\circ$ 的模空间上典范地确定。由于 $V^\circ$ 是 Calabi--Yau 簇，它也与 $\Omega$ 的选取无关。于是自然可以断言：镜像映射由 $t=\int_{\gamma_1}\Omega/\int_{\gamma_0}\Omega$ 以及
\begin{equation}\label{eq:mirror-map-periods}
q=\exp\!\left(2\pi i\frac{\int_{\gamma_1}\Omega}{\int_{\gamma_0}\Omega}\right)
\end{equation}
给出。\SourceCite{BCOV2} 给出了支持这一断言的物理学论证；我们还可以从数学上证明，这个镜像映射确实正确，见定理~\ref{thm:lly-quintic-mirror} 或定理~\ref{thm:givental-projective-mirror}。

为了把 $\int_{\gamma_1}\Omega$ 与 $\int_{\gamma_0}\Omega$ 明确表示为 $x$ 的函数，我们沿用 \SourceCite{Morrison1} 的方法。对于 $V^\circ$ 上任意三维圈 $\gamma$，积分 $y=\int_\gamma\Omega$ 称为该族的一个\emph{周期}。众所周知，周期满足 Picard--Fuchs 方程。对 $h^3(V^\circ)=4$ 的三维簇，该方程的形式为
\[
y^{(4)}+f_1y^{(3)}+f_2y''+f_3y'+f_4y=0,
\]
其中 $f_i$ 依赖于模坐标 $x$，导数也对 $x$ 取。一般地，Picard--Fuchs 方程可通过多种方法计算：直接计算周期 \SourceCite{CdGP}；用 Griffiths--Dwork 方法降低极点阶数 \SourceCite{Morrison1,Font,KT1}；利用广义超几何方程 \SourceCite{BvS,HKTY1}；或者，按照 \SourceCite{Givental4} 的一个断言，用环面完全交的显式公式。后三种方法将在\hyperref[chap:hodge-yukawa]{第 5 章}讨论。\hyperref[chap:mirror-theorem]{第~\ref*{chap:mirror-theorem} 章}还会给出整理这些数据的其他方式，以及不同版本镜像定理的证明。第 \ref{sec:quintic-other-methods} 节将先作简要预览。

若把 $\Omega$ 选为三形式
\begin{equation}\label{eq:quintic-residue}
\begin{aligned}
\Omega=\operatorname{Res}\Bigg(\frac{\psi}{x_1^5+\cdots+x_5^5+\psi x_1\cdots x_5}
\Big(&x_1\,dx_2\wedge dx_3\wedge dx_4\wedge dx_5\\
&-x_2\,dx_1\wedge dx_3\wedge dx_4\wedge dx_5\\
&+x_3\,dx_1\wedge dx_2\wedge dx_4\wedge dx_5\\
&-x_4\,dx_1\wedge dx_2\wedge dx_3\wedge dx_5\\
&+x_5\,dx_1\wedge dx_2\wedge dx_3\wedge dx_4\Big)\Bigg),
\end{aligned}
\end{equation}
\SourcePage{43}
那么，\hyperref[chap:hodge-yukawa]{第 5 章}的方法便给出 Picard--Fuchs 方程
\begin{equation}\label{eq:quintic-pf}
\begin{aligned}
0={}&\left(x\frac{d}{dx}\right)^4y
+\frac{2\cdot5^5x}{1+5^5x}\left(x\frac{d}{dx}\right)^3y
+\frac{7\cdot5^4x}{1+5^5x}\left(x\frac{d}{dx}\right)^2y\\
&+\frac{2\cdot5^4x}{1+5^5x}\left(x\frac{d}{dx}\right)y
+\frac{24\cdot5x}{1+5^5x}y.
\end{aligned}
\end{equation}
该方程在 $x=0$ 处有一个正则奇点。还要注意，在另一个边界点 $x=-5^{-5}$ 处也有奇点。

周期 $y_0=\int_{\gamma_0}\Omega$ 与 $y_1=\int_{\gamma_1}\Omega$ 都满足这个方程。$\gamma_0$ 与 $\gamma_1$ 的单值性性质说明，$y_0$ 和 $2\pi i y_1-(\log x)y_0$ 在 $x=0$ 处都是单值的。再用求解式~\eqref{eq:quintic-pf} 的标准方法，可知 $y_0,y_1$ 在下述变换意义下唯一：
\[
y_0\longmapsto b_1y_0,\qquad y_1\longmapsto b_1y_1+b_2y_0,
\]
其中 $b_1,b_2\in\mathbb C$ 且 $b_1\ne0$。这与 $\gamma_0,\gamma_1$ 的唯一性类似，只是这里丢失了整数格。因此，我们可以确定 $q=e^{2\pi i y_1/y_0}$，但仍相差一个常数因子 $c_1=\exp(2\pi i b_2/b_1)$。所得镜像映射为
\begin{equation}\label{eq:q-series-c1}
q=c_1(x-770x^2+\cdots),
\end{equation}
其逆映射为
\begin{equation}\label{eq:x-series-c1}
x=\frac q{c_1}+770\left(\frac q{c_1}\right)^2+\cdots.
\end{equation}
第 \ref{sec:quintic-assemble} 节将利用\IndexLocation{indexloc:1197:43}镜像对称证明 $c_1=-1$。也有其他方法可得到这种归一化；例如，在求解\IndexLocation{indexloc:0618:43}广义\IndexLocation{indexloc:0617:43}超几何函数时采用 Frobenius 方法，这种归一化就是一种自然的选择。\SourceCite{BvS} 提到了这个想法。\hyperref[chap:moduli]{第~\ref*{chap:moduli} 章}一般性地讨论镜像映射时，我们将再回到这一点。

利用广义超几何方程的求解方法，许多公式都能写成闭式。例如，在常数倍意义下，$y_0$ 为
\begin{equation}\label{eq:y0-hypergeometric}
y_0=\sum_{n=0}^{\infty}\frac{(5n)!}{(n!)^5}(-1)^n x^n=1-120x+\cdots;
\end{equation}
若假定 $c_1=-1$，则镜像映射可写为
\begin{equation}\label{eq:q-hypergeometric}
q=-x\exp\!\left(\frac{5}{y_0(x)}\sum_{n=1}^{\infty}\frac{(5n)!}{(n!)^5}
\left[\sum_{j=n+1}^{5n}\frac1j\right](-1)^n x^n\right).
\end{equation}
推导细节见\hyperref[sec:mirror-map]{第~\ref*{sec:mirror-map} 节}；第 \ref{sec:quintic-other-methods} 节还会介绍另一种方法。结合式~\eqref{eq:q-hypergeometric} 与 Dwork 的 $p$ 进超几何级数理论，可以证明 $q$ 是 $x$ 的一个整系数幂级数 \SourceCite{LY2}。

\SourcePage{44}
\section{B 模型关联函数}\label{sec:quintic-b-model}

现在转向计算 $V^\circ$ 上的 Yukawa 耦合。一个关键但初等的观察是：Yukawa 耦合本身也满足一个线性微分方程，而且可以很容易地从 Picard--Fuchs 方程推导出来。依照第 \ref{sec:three-point} 节的式~\eqref{eq:b-yukawa}，有
\begin{equation}\label{eq:b-x-derivatives}
\left\langle x\frac d{dx},x\frac d{dx},x\frac d{dx}\right\rangle
=\int_{V^\circ}\Omega\wedge\Omega''',
\end{equation}
其中撇号表示算子 $x\frac d{dx}$。这里假设 $\Omega$ 由式~\eqref{eq:quintic-residue} 给出。我们也可以把 $\Omega$ 乘以 $x$ 的任意处处非零函数；究竟选哪个倍数，就是\emph{归一化问题}，留待本节稍后处理。正如第 \ref{sec:three-point} 节脚注所说，归一化还会在式~\eqref{eq:b-x-derivatives} 右边引入一个负号；\hyperref[subsec:normalized-yukawa-mirror]{第~\ref*{subsec:normalized-yukawa-mirror} 小节}将对此作解释。

我们断言，Yukawa 耦合 $Y=\int_{V^\circ}\Omega\wedge\Omega'''$ 满足微分方程
\begin{equation}\label{eq:yukawa-ode}
\left(x\frac d{dx}\right)Y=-\frac{5^5x}{1+5^5x}Y.
\end{equation}
为验证这一点，先注意：由 Griffiths 横截性，$\Omega'$ 和 $\Omega''$ 都没有 $(0,3)$ 分量。因此
\begin{equation}\label{eq:omega-zero}
\int_{V^\circ}\Omega\wedge\Omega
=\int_{V^\circ}\Omega\wedge\Omega'
=\int_{V^\circ}\Omega\wedge\Omega''=0.
\end{equation}
将最后一个等式求导两次，得到
\begin{equation}\label{eq:omega-identity}
\int_{V^\circ}\Omega\wedge\Omega^{(4)}
+2\int_{V^\circ}\Omega'\wedge\Omega'''=0.
\end{equation}
另一方面，由 $Y=\int_{V^\circ}\Omega\wedge\Omega'''$ 得
$Y'=\int_{V^\circ}\Omega'\wedge\Omega'''+\int_{V^\circ}\Omega\wedge\Omega^{(4)}$。利用式~\eqref{eq:omega-identity}，可知 $Y'=\frac12\int_{V^\circ}\Omega\wedge\Omega^{(4)}$。而 Picard--Fuchs 方程~\eqref{eq:quintic-pf} 蕴含上同调中的等式
\begin{equation}\label{eq:omega-pf}
0=\Omega^{(4)}+\frac{2\cdot5^5x}{1+5^5x}\Omega'''
+\frac{7\cdot5^4x}{1+5^5x}\Omega''
+\frac{2\cdot5^4x}{1+5^5x}\Omega'
+\frac{24\cdot5x}{1+5^5x}\Omega.
\end{equation}
把 $\Omega$ 与式~\eqref{eq:omega-pf} 右边各项作楔积，再用式~\eqref{eq:omega-zero}，便得到式~\eqref{eq:yukawa-ode}。

式~\eqref{eq:yukawa-ode} 的解为
\begin{equation}\label{eq:yukawa-unnormalized}
Y=\frac{c_2}{1+5^5x},
\end{equation}
其中常数 $c_2$ 将在下一节确定。因此，我们已确定 Yukawa 耦合，但仍相差一个常数因子。可以看到，它在边界点 $x=-5^{-5}$ 处是奇异的，这也符合预期。

但工作尚未完成。式~\eqref{eq:quintic-residue} 给出的 $\Omega$ 可乘以 $x$ 的任意处处非零函数，所以还须归一化 Yukawa 耦合。一种自然的做法，是要求归一化后 $\int_{\gamma_0}\Omega=1$。沿用式~\eqref{eq:quintic-residue} 的 $\Omega$，取周期 $y_0(x)=\int_{\gamma_0}\Omega$，再用 $\Omega/y_0(x)$ 代替 $\Omega$，就得到所需归一化。由式~\eqref{eq:omega-zero} 很容易证明，这种归一化对 Yukawa 耦合
\SourcePage{45}
$Y$ 的作用，是把 $Y$ 换成 $Y/y_0(x)^2$。因此，在这一 $\Omega$ 的选取下，得到\emph{归一化 Yukawa 耦合}
\begin{equation}\label{eq:yukawa-normalized}
\left\langle x\frac d{dx},x\frac d{dx},x\frac d{dx}\right\rangle
=\frac{c_2}{(1+5^5x)y_0(x)^2}.
\end{equation}
注意，上面提到的负号实质上不会影响这个表达式，因为它可以吸收到 $c_2$ 的选取中。据我们所知，这或许也是该负号此前未出现在数学文献中的原因。\hyperref[chap:hodge-yukawa]{第 5 章}将给出计算归一化 Yukawa 耦合的一般方法，并讨论它与 Gauss--Manin 联络之间出乎意料地紧密的联系。

\section{综合起来}\label{sec:quintic-assemble}

现在可以利用前两节的公式，计算五次三维簇的 A 模型关联函数。镜像对称给出同构
\[
H^1(V,\Omega_V^1)\simeq H^1(V^\circ,T_{V^\circ}),
\]
并使 A 模型与 B 模型的关联函数相容。于是，若 $\theta\in H^1(V^\circ,T_{V^\circ})$ 在镜像对称下对应于 $H\in H^1(V,\Omega_V^1)$，依照\hyperref[chap:introduction]{第 1 章}的解释，就应有 $\langle H,H,H\rangle=\langle\theta,\theta,\theta\rangle$。

下面确定 $\theta$。把 $H^1(V,\Omega_V^1)$ 看作 Kähler 模空间的切空间，$H$ 就是 $V$ 的 Kähler 模空间上的一个向量场。若把一般的复化 Kähler 类写为 $tH$，该向量场便对应于 $d/dt$。用 Kähler 模的局部坐标 $q=e^{2\pi it}$ 表示，有 $\frac d{dt}=2\pi iq\frac d{dq}$。在镜像映射 $q=q(x)$ 下，$\frac d{dq}$ 映为 $\frac{dx}{dq}\frac d{dx}$。于是
\[
H=2\pi iq\frac d{dq}\quad\longmapsto\quad
\theta=2\pi iq\frac{dx}{dq}\frac d{dx}.
\]
由 $\langle H,H,H\rangle=\langle\theta,\theta,\theta\rangle$，得到
\begin{equation}\label{eq:quintic-mirror-yukawa}
\begin{aligned}
\langle H,H,H\rangle
&=\left(2\pi iq\frac{dx}{dq}\right)^3
  \left\langle\frac d{dx},\frac d{dx},\frac d{dx}\right\rangle\\
&=\left(2\pi i\frac qx\frac{dx}{dq}\right)^3
  \left\langle x\frac d{dx},x\frac d{dx},x\frac d{dx}\right\rangle.
\end{aligned}
\end{equation}
这一断言就是 A 模型与 B 模型关联函数通过镜像映射相等，对五次三维簇而言，这正是镜像对称。

为使式~\eqref{eq:quintic-mirror-yukawa} 更明确，注意式~\eqref{eq:yukawa-normalized} 已计算出右边的归一化 Yukawa 耦合。把它与镜像映射~\eqref{eq:q-series-c1}、\eqref{eq:x-series-c1} 以及 $y_0(x)$ 的公式~\eqref{eq:y0-hypergeometric} 结合，便有
\[
\begin{aligned}
\langle H,H,H\rangle
&=\frac{c_2}{(1+5^5x)y_0(x)^2}
  \left(2\pi i\frac qx\frac{dx}{dq}\right)^3\\
&=\frac{c_2}{(1+5^5x)(1-120x+\cdots)^2}
 \left(2\pi i(1-770x+\cdots)
 \left(1+1540\frac q{c_1}+\cdots\right)\right)^3.
\end{aligned}
\]
再用式~\eqref{eq:x-series-c1} 全部换成 $q$ 坐标，化简为
\begin{equation}\label{eq:quintic-c-series}
\langle H,H,H\rangle
=\frac{(2\pi i)^3c_2}{
 (1+5^5(q/c_1)+\cdots)(1-240(q/c_1)+\cdots)}
 \left(1+770\frac q{c_1}+\cdots\right)^3.
\end{equation}

\SourcePage{46}
现在确定常数 $c_1$ 和 $c_2$。基本思路是让上述幂级数与我们已知的、式~(2.1) 中给出的公式
\begin{equation}\label{eq:quintic-a-expand}
\langle H,H,H\rangle=5+\sum_{d=1}^{\infty}n_d d^3\frac{q^d}{1-q^d}
=5+2875q+\cdots
\end{equation}
相吻合。首先，由于这个级数的常数项为 $5$，与式~\eqref{eq:quintic-c-series} 比较便得到 $c_2=5/(2\pi i)^3$。其次，我们断言 $c_1=-1$。事实上，采用上述 $c_2$ 后，式~\eqref{eq:quintic-c-series} 的幂级数展开为
\[
\langle H,H,H\rangle=5-2875\frac q{c_1}+\cdots.
\]
经典结果表明，一般五次三维簇包含 $2875$ 条直线，所以式~\eqref{eq:quintic-a-expand} 中 $n_1=2875$。因此，为使两个级数相等，必须有 $c_1=-1$。这并非一厢情愿；\SourceCite{CdGP} 的原始计算确实给出了所需的值 $c_1=-1$。如前所说，\hyperref[chap:moduli]{第~\ref*{chap:moduli} 章}研究镜像映射时，还会给出系统确定 $c_1$ 的方法。

确定 $c_1,c_2$ 后，可以把五次三维簇 $V$ 的镜像对称概括如下。首先，根据式~\eqref{eq:y0-hypergeometric}，
\begin{equation}\label{eq:quintic-y0-final}
y_0(x)=\sum_{n=0}^{\infty}\frac{(5n)!}{(n!)^5}(-1)^nx^n
\end{equation}
是镜像 $V^\circ$ 的 Picard--Fuchs 方程~\eqref{eq:quintic-pf} 唯一的全纯解，唯一性只差一个常数倍。还有一个具有对数奇点的解
\begin{equation}\label{eq:quintic-y1-final}
y_1(x)=y_0(x)\log(-x)
+5\sum_{n=1}^{\infty}\frac{(5n)!}{(n!)^5}
\left[\sum_{j=n+1}^{5n}\frac1j\right](-1)^nx^n.
\end{equation}
依照式~\eqref{eq:q-hypergeometric}，镜像映射为
\begin{equation}\label{eq:quintic-q-final}
q=\exp\!\left(\frac{y_1}{y_0}\right)
=-x\exp\!\left(\frac5{y_0(x)}\sum_{n=1}^{\infty}\frac{(5n)!}{(n!)^5}
\left[\sum_{j=n+1}^{5n}\frac1j\right](-1)^nx^n\right).
\end{equation}
那么，式~\eqref{eq:quintic-mirror-yukawa} 所述五次三维簇的镜像对称断言
\begin{equation}\label{eq:quintic-mirror-identity}
5+\sum_{d=1}^{\infty}n_dd^3\frac{q^d}{1-q^d}
=\frac5{(1+5^5x)y_0(x)^2}\left(\frac qx\frac{dx}{dq}\right)^3.
\end{equation}
请记住：左边是 $V$ 的 A 模型关联函数；右边是 $V^\circ$ 的 B 模型关联函数，即把归一化 Yukawa 耦合视为 $q$ 的函数。

展开式~\eqref{eq:quintic-mirror-identity} 的右边，可得
\begin{equation}\label{eq:quintic-first-coefficients}
5+\sum_{d=1}^{\infty}n_dd^3\frac{q^d}{1-q^d}
=5+2875\frac q{1-q}+609250\cdot2^3\frac{q^2}{1-q^2}+\cdots.
\end{equation}
所以，镜像对称蕴含 $n_2=609250$。继续计算这个级数的其他项，便得到五次三维簇瞬子数 $n_d$ 的下列预测：

\SourcePage{47}
\begin{center}
\begin{tabular}{r|r}
$d$ & $n_d$\\\hline
1 & 2875\\
2 & 609250\\
3 & 317206375\\
4 & 242467530000\\
5 & 229305888887625
\end{tabular}
\end{center}
前四项与已知的枚举结果吻合：$n_1=2875$ 是经典结果（如前所述）；$n_2=609250$ 由 \SourceCite{Katz2} 证明；$n_3=317206375$ 由 \SourceCite{ES1} 证明；$n_4=242467530000$ 由 \SourceCite{Kontsevich2} 证明。\SourceCite{CdGP} 计算了所有 $d\leq10$ 的这些数。

对于一般的 $d$，式~\eqref{eq:quintic-mirror-identity} 断言，式~\eqref{eq:quintic-first-coefficients} 右边级数中出现的 $2875$、$609250$ 等数，正是由五次三维簇的 Gromov--Witten 不变量通过式~\eqref{eq:multi-cover} 定义的瞬子数 $n_d$。\hyperref[subsec:instanton-numbers]{第~\ref*{subsec:instanton-numbers} 小节}将证明：当 $d\leq9$ 时，$n_d$ 的确等于一般五次三维簇上次数为 $d$ 的有理曲线数目。对 $d\geq10$，瞬子数 $n_d$ 仍然存在，但正如已经指出的，它与 $d$ 次有理曲线之间的关系更微妙。

上表中的数暗示 $n_d$ 具有一些非凡的\IndexLocation{indexloc:0630:47}整除性质。容易看出 $n_d$ 是有理数，但尚不知道它们是否总是整数，尽管已计算出的数都是整数。如果假定 $n_d$ 都是整数，那么当 $d$ 不被 $5$ 整除时，可以证明 $5^3\mid n_d$ \SourceCite{LY1}。人们猜想，对\emph{所有} $d$ 都有 $5^3\mid n_d$。

到目前为止，我们还没有谈如何证明式~\eqref{eq:quintic-mirror-identity}。这恰恰是镜像定理的核心内容，也就是下面将要讨论的主题。

\begingroup\linespread{1.22}\selectfont
\section{镜像定理}\label{sec:quintic-other-methods}

前几节的式~\eqref{eq:quintic-mirror-identity} 基于镜像对称的“经典”表述，其首要目标是证明某些关联函数可通过镜像映射相互等同。为了论证这些计算，必须证明一个镜像定理。本节将介绍镜像定理的三个不同版本，它们都蕴含五次三维簇的“经典”镜像对称。我们暂不一般地陈述这些定理，而是考察它们对于五次三维簇 $V$ 及其镜像 $V^\circ$ 的具体形式。完整陈述见第 \hyperref[chap:quantum-cohomology]{\ref*{chap:quantum-cohomology}}、\hyperref[chap:mirror-theorem]{\ref*{chap:mirror-theorem}} 章。

\subsection{Hodge 理论}\label{subsec:quintic-hodge-theory}
第 \ref{sec:quintic-mirror-map} 节构造的镜像映射，把五次三维簇 $V$ 的 Kähler 参数 $q$ 明确表示为镜像 $V^\circ$ 的模参数 $x$ 的函数。因此，我们可以用 $q$ 作为 $V^\circ$ 的模空间的新局部参数。第 \ref{sec:quintic-b-model} 节又确定了归一化 Yukawa 耦合~\eqref{eq:yukawa-normalized}；\hyperref[chap:hodge-yukawa]{第 5 章}将把这一耦合记为 $Y$。前一节的枚举预测，是通过把 $Y$ 写成 $q$ 的函数，再令它等于 A 模型关联函数~\eqref{eq:quintic-a-expand} 得到的。

不过，\hyperref[chap:hodge-yukawa]{第 5 章}将告诉我们，一旦用 $q$ 表示 $Y$，也就完全确定了 $H^3(V^\circ,\mathbb C)$ 上的 Hodge 结构变分。更准确地说，在命题~\ref{prop:normalized-yukawa-pf} 的证明中，我们将看到，在适当的基下，
\SourcePage{48}
$V^\circ$ 的 Gauss--Manin 联络 $\nabla$ 的联络矩阵为
\begin{equation}\label{eq:connection-matrix}
\begin{pmatrix}
0&0&0&0\\
1&0&0&0\\
0&Y&0&0\\
0&0&1&0
\end{pmatrix},
\end{equation}
这里使用 $q$ 作为模参数。

这表明，归一化 B 模型关联函数 $Y$ 并不是物理学所提示的一个奇特函数，而是镜像族的 Hodge 结构不可缺少的一部分。此外，$Y$ 应与 A 模型关联函数对应，这暗示 A 模型一侧也应存在一个 Hodge 结构变分。这是我们第一次看到，镜像对称或可用 Hodge 结构变分来表述。\hyperref[chap:quantum-cohomology]{第~\ref*{chap:quantum-cohomology} 章}定义\emph{A 型 Hodge 结构变分}时，将更充分地发展这一思想。与之相伴的联络称为\emph{A 模型联络}，由量子上同调构造。特别地，对于五次三维簇 $V$，例题~\ref{ex:a-model-yukawa-connection} 将说明：只要把 $Y$ 取为 A 模型关联函数~\eqref{eq:quintic-a-expand}，A 模型联络的联络矩阵也正是式~\eqref{eq:connection-matrix}。

\hyperref[chap:quantum-cohomology]{第~\ref*{chap:quantum-cohomology} 章}将严格定义\emph{数学镜像对}。对五次三维簇及其镜像而言，这意味着 $V$ 的 A 型 Hodge 结构变分，应通过镜像映射同构于镜像族 $V^\circ$ 所产生的 Hodge 结构变分。\hyperref[chap:quantum-cohomology]{第~\ref*{chap:quantum-cohomology} 章}还将证明，这一版本的镜像对称，与式~\eqref{eq:quintic-mirror-identity} 给出的 A 模型、B 模型关联函数相等是等价的。

无论证明镜像对称的经典形式，还是刚讨论的更精细的 Hodge 理论形式，都不是易事。特别是，证明需要引入一些新的研究对象。下面将说明，这些新对象又引出了镜像定理的新版本。

\subsection{Givental 的 $I$ 函数与 $J$ 函数}\label{subsec:givental-ij}

在 Givental 证明镜像定理的方法中，两个上同调值形式函数 $I_{\mathcal V}$ 和 $J_{\mathcal V}$ 起关键作用。先讨论 $J$。这个函数由引力关联数定义；引力关联数可视为 Gromov--Witten 不变量的推广。事实上，我们要先从 Givental 的 $J$ 函数谈起。\hyperref[subsec:cy-threefold-j-function]{第~\ref*{subsec:cy-threefold-j-function} 小节}将看到，对五次三维簇有
\begin{equation}\label{eq:givental-j}
J=e^{(t_0+t_1H)/\hbar}
\left(1+\hbar^{-2}\sum_{d>0}N_dq^d\,d\ell
-2\hbar^{-3}\sum_{d>0}N_dq^d\,\mathrm{pt}\right),
\end{equation}
其中 $q=e^{t_1}$，$\ell$ 是 $V$ 中一条直线的类，$\mathrm{pt}$ 是一个点的类，而且
\begin{equation}\label{eq:nd-gw}
N_d=\langle I_{0,0,d}\rangle=\sum_{k\mid d}n_{d/k}k^{-3}.
\end{equation}
这里的 $n_{d/k}$ 是我们想计算的瞬子数。本章前面提到了 Gromov--Witten 不变量 $\langle I_{0,3,d}\rangle(H,H,H)$。\hyperref[chap:gromov-witten]{第~\ref*{chap:gromov-witten} 章}的 Gromov--Witten 不变量性质蕴含 $\langle I_{0,3,d}\rangle(H,H,H)=d^3\langle I_{0,0,d}\rangle$，所以 $N_d$ 的上述公式与式~\eqref{eq:multi-cover} 等价。

$J$ 函数在\hyperref[subsec:giventals-j-function]{第~\ref*{subsec:giventals-j-function} 小节} Givental 的量子微分算子理论中居于核心地位。对于眼下的目的，式~\eqref{eq:givental-j} 最关键的性质，是 $J$ 包含我们所关注的全部枚举信息。

\SourcePage{49}
函数 $J$ 取值于 $V$ 的上同调环。但若要把它与下面将介绍的 $I_{\mathcal V}$ 比较，我们需要一个取值于 $H^*(\mathbb P^4)$ 的对象。设 $i:V\to\mathbb P^4$ 是包含映射，$i_!$ 是 Gysin 映射，定义 $J_{\mathcal V}=i_!(J)$ 即可。由式~\eqref{eq:givental-j} 容易得到
\begin{equation}\label{eq:givental-jv}
J_{\mathcal V}=e^{(t_0+t_1H)/\hbar}5H
\left(1+\hbar^{-2}\sum_{d>0}N_dq^d\,\frac d5H^2
-2\hbar^{-3}\sum_{d>0}N_dq^d\,\frac15H^3\right).
\end{equation}
（\hyperref[chap:mirror-theorem]{第~\ref*{chap:mirror-theorem} 章}将以略有不同的方式定义 $J_{\mathcal V}$。）这里重要的是：$J_{\mathcal V}$ 唯一确定各个 $n_d$。

Givental 的第二个函数是 $I_{\mathcal V}$，它也取值于上同调。对五次三维簇 $V$，例题~\ref{ex:givental-quintic-mirror-potential} 给出的公式是
\[
I_{\mathcal V}=e^{(t_0+t_1H)/\hbar}5H
\sum_{d=0}^{\infty}e^{dt_1}
\frac{(5H+\hbar)\cdots(5H+5d\hbar)}{
\bigl((H+\hbar)\cdots(H+d\hbar)\bigr)^5}.
\]
虽然 $I_{\mathcal V}$ 取值于 $\mathbb P^4$ 的上同调，但 Givental 的一个观察是：$I_{\mathcal V}$ 由镜像 $V^\circ$ 的周期确定。因此，$I_{\mathcal V}$ 包含镜像的模信息。Givental 版本的镜像定理断言，$I_{\mathcal V}$ 与 $J_{\mathcal V}$ 之间有一种特别好的关系。我们对证明的解释将使用\emph{量子超平面截面原理}；该原理涉及稳定映射模空间上等变上同调的局部化公式。\hyperref[chap:mirror-theorem]{第~\ref*{chap:mirror-theorem} 章}将讨论所有这些内容。

为看清它与前几节的计算有什么关系，令 $t_0=0$、$\hbar=1$，把 $I_{\mathcal V}$ 的上述公式简化为
\begin{equation}\label{eq:iv-simple}
I_{\mathcal V}=e^{t_1H}5H\sum_{d=0}^{\infty}e^{dt_1}
\frac{\prod_{m=1}^{5d}(5H+m)}{\prod_{m=1}^{d}(H+m)^5}.
\end{equation}
可把 $I_{\mathcal V}$ 展开为
\begin{equation}\label{eq:iv-period-expansion}
I_{\mathcal V}=5H\bigl(y_0(t_1)+y_1(t_1)H+y_2(t_1)H^2+y_3(t_1)H^3\bigr).
\end{equation}
作变量替换 $x=-e^{t_1}$ 后，\hyperref[chap:hodge-yukawa]{第 5 章}将证明 $I_{\mathcal V}$ 满足镜像 $V^\circ$ 的 Picard--Fuchs 方程~\eqref{eq:quintic-pf}。实际上，利用 Frobenius 方法，\hyperref[chap:moduli]{第~\ref*{chap:moduli} 章}将证明 $y_0,y_1,y_2,y_3$ 构成解的一组基。而且，$y_0$ 就是式~\eqref{eq:quintic-y0-final} 的级数，$y_1$ 则是式~\eqref{eq:quintic-y1-final}。因此，在 Givental 的表述中，镜像映射~\eqref{eq:quintic-q-final} 自然以 $q=\exp(y_1/y_0)$ 的形式出现。

Givental 版本的镜像定理等价于如下断言：通过镜像映射变换坐标之后，$J_{\mathcal V}=I_{\mathcal V}/y_0$，见定理~\ref{thm:givental-projective-mirror}。由于式~\eqref{eq:nd-gw}、\eqref{eq:givental-jv} 说明 $J_{\mathcal V}$ 包含 $n_d$，而式~\eqref{eq:iv-simple} 明确给出了 $I_{\mathcal V}$，我们可以据此计算任意多项瞬子数 $n_d$。

不过，为使这一组思想首尾相接，还须说明它与关联函数的关系。先注意：第 \ref{sec:quintic-assemble} 节 $V^\circ$ 的归一化 Yukawa 耦合 $Y$，会自然地从 $I_{\mathcal V}$ 产生。为说明原因，考虑 $I_{\mathcal V}/y_0$。只展开到含 $e^{t_1}$ 乘以 $t_1$ 的多项式的项，得到
\[
\begin{aligned}
I_{\mathcal V}/y_0=5H\Big(&1+(t_1+770e^{t_1}+\cdots)H\\
&+\bigl(\tfrac{t_1^2}{2}+770t_1e^{t_1}+575e^{t_1}+\cdots\bigr)H^2\\
&+\bigl(\tfrac{t_1^3}{6}+385t_1^2e^{t_1}+575t_1e^{t_1}-1150e^{t_1}+\cdots\bigr)H^3\Big).
\end{aligned}
\]

\SourcePage{50}
若改用 $q$ 表示，则
\[
I_{\mathcal V}/y_0=5H\Big(1+sH+\bigl(\tfrac{s^2}{2}+575q+\cdots\bigr)H^2
+\bigl(\tfrac{s^3}{6}+575sq-1150q+\cdots\bigr)H^3\Big),
\]
其中 $s=\log q$。于是，若撇号表示 $d/ds=q\,d/dq$，就可算得
\begin{equation}\label{eq:iv-second-derivative}
\begin{aligned}
(I_{\mathcal V}/y_0)''
&=5(1+575q+\cdots)H^3+5(s+575sq+\cdots)H^4\\
&=(5+2875q+\cdots)(H^3+sH^4).
\end{aligned}
\end{equation}
最后一行的第一个因子以 $5+2875q+\cdots$ 开头；至少到 $q$ 的一次项为止，它与归一化 Yukawa 耦合 $Y$ 相符。实际上，这个因子正是 $Y$。这就是 Givental 关于 $I_{\mathcal V}$ 由镜像周期确定的观察。要证明这一点，可用联络矩阵~\eqref{eq:connection-matrix}。它说明，归一化 Picard--Fuchs 方程可写为
\[
\frac{d^2}{ds^2}\left(\frac{d^2y/ds^2}{Y}\right)=0.
\]
容易看出，此方程有解 $f_0,f_1,f_2,f_3$，使得 $f_0=1$、$f_1=s$、$f_2''=Y$ 且 $f_3''=Ys$。式~\eqref{eq:iv-period-expansion} 给出了未归一化的 Picard--Fuchs 方程~\eqref{eq:quintic-pf} 的解 $y_0,y_1,y_2,y_3$。归一化就是除以 $y_0$；所得解为 $y_0/y_0=1$、$y_1/y_0=s$、$y_2/y_0$ 和 $y_3/y_0$。精彩之处在于，$y_2/y_0$ 与 $y_3/y_0$ 基本上就是 $f_2$ 与 $f_3$。严格地说，有恒等式
\begin{equation}\label{eq:period-derivative-identities}
5\frac{d^2}{ds^2}\frac{y_2}{y_0}=Y,
\qquad
5\frac{d^2}{ds^2}\frac{y_3}{y_0}=Ys.
\end{equation}
要证明第一个等式，注意对某些常数 $a,b,c,d$，有 $y_2/y_0=a+bs+cf_2+df_3$，所以 $(y_2/y_0)''=cY+dYs$。由式~\eqref{eq:iv-second-derivative} 可知 $(y_2/y_0)''=1+575q+\cdots$；又因为 $Y=5+2875q+\cdots$，于是 $c=1/5$、$d=0$。第一个等式得证；第二个等式的论证类似。

利用式~\eqref{eq:period-derivative-identities} 可得 $(I_{\mathcal V}/y_0)''=Y(H^3+sH^4)$。命题~\ref{prop:cy-threefold-j-potential} 将类似地证明 $J_{\mathcal V}''=\langle H,H,H\rangle(H^3+sH^4)$。因此，按上述方式理解的定理~\ref{thm:givental-projective-mirror} 蕴含 $\langle H,H,H\rangle=Y$，从而证明式~\eqref{eq:quintic-mirror-identity} 所要求的 A 模型与 B 模型关联函数相等。这就建立了五次三维簇的“经典”镜像对称。

\subsection{Euler 数据与 Gromov--Witten 势}\label{subsec:euler-data}

\SourceCite{LLY} 给出了证明镜像定理的另一种方法。这篇论文包含五次三维簇镜像对称的完整证明，\hyperref[chap:mirror-theorem]{第~\ref*{chap:mirror-theorem} 章}将予讨论。这里先对其中的内容作一个简短预览。

\SourceCite{LLY} 引入了几个值得我们考察的对象。第一个是\emph{Euler 数据}：它是某些射影空间的等变上同调中相互相容的一组元素，能够把稳定映射的模空间“线性化”。有两组特殊的 Euler 数据，记为 $\widehat Q$ 和 $\widehat P$；粗略地说，它们分别编码 A 模型与 B 模型关联函数的信息。

由 Euler 数据可以构造一些上同调值函数，不过这里的值属于等变上同调。例如，五次三维簇的 Euler 数据 $\widehat P$ 给出函数

\SourcePage{51}
\[
HG[\mathcal I(\widehat P)](t_1)
=e^{-pt_1/\hbar}\sum_{d\geq0}
\frac{\prod_{m=0}^{5d}(5p-m\hbar)}{
\prod_{m=1}^{d}\prod_{k=0}^{4}(p-\lambda_k-m\hbar)}e^{dt_1}.
\]
参见定义~\ref{def:euler-data-hypergeometric-function} 后的讨论。这里 $p$ 是等变超平面类，而 $\lambda_k$ 是某些等变上同调类。如果令 $\lambda_k\to0$——\hyperref[chap:localization]{第~\ref*{chap:localization} 章}将称之为“\IndexLocation{indexloc:0945:51}非等变极限”——就得到普通上同调。此时 $p$ 变成 $H$；再令 $\hbar=-1$，便得到式~\eqref{eq:iv-simple} 中的 $I_{\mathcal V}$。特别地，式~\eqref{eq:iv-period-expansion} 变为
\begin{equation}\label{eq:euler-p-limit}
\lim_{\lambda\to0}HG[\mathcal I(\widehat P)](t_1)
=I_{\mathcal V}=5H\bigl(y_0(t_1)+y_1(t_1)H+y_2(t_1)H^2+y_3(t_1)H^3\bigr).
\end{equation}
和前面一样，这使我们能够把镜像映射写为 $q=\exp(y_1/y_0)$。我们用 $\Psi(t_1)$ 表示商 $y_1/y_0$，于是镜像映射为 $q=\exp(\Psi(t_1))$。若记 $q=e^s$，便有 $s=\Psi(t_1)=y_1(t_1)/y_0(t_1)$。

迄今处理的是描述 $V^\circ$ 复模的 Euler 数据 $\widehat P$。还有由 $V$ 的 Gromov--Witten 不变量决定的 Euler 数据 $\widehat Q$。理解它最容易的途径，是借助\emph{Gromov--Witten 势}，其定义为
\begin{equation}\label{eq:gw-potential-quintic}
\Phi(s)=\frac56s^3+\sum_{d=1}^{\infty}N_dq^d
=\frac56s^3+\sum_{d=1}^{\infty}\langle I_{0,0,d}\rangle q^d,
\end{equation}
其中 $N_d$ 由式~\eqref{eq:nd-gw} 给出，且 $q=e^s$。函数 $\Phi(s)$ 与 A 模型关联函数关系很近；事实上，很容易看出式~\eqref{eq:quintic-a-expand} 就是 $d^3\Phi(s)/ds^3$。\hyperref[chap:quantum-cohomology]{第~\ref*{chap:quantum-cohomology} 章}将详细研究 Gromov--Witten 势。

现在考察 Euler 数据 $\widehat Q$。它给出等变上同调值函数 $HG[\mathcal I(\widehat Q)](s)$；\hyperref[sec:quintic-mirror-theorem]{第~\ref*{sec:quintic-mirror-theorem} 节}将证明
\begin{equation}\label{eq:euler-q-limit}
\lim_{\lambda\to0}HG[\mathcal I(\widehat Q)](s)
=5H\left(1+sH+\frac{\Phi'}5H^2+\frac{s\Phi'-2\Phi}5H^3\right),
\end{equation}
其中撇号表示 $d/ds$。（\hyperref[chap:mirror-theorem]{第~\ref*{chap:mirror-theorem} 章}给出的这个公式还包含一个 $\hbar$；与上面一样，这里已令它等于 $-1$。）此式体现了前面所说的：$\widehat Q$ 由 Gromov--Witten 不变量确定。

证明中的关键一步，是说明作镜像映射给出的代换 $s=\Psi(t_1)$ 后，式~\eqref{eq:euler-q-limit} 的极限与式~\eqref{eq:euler-p-limit} 的极限相同，只差一个 $y_0$ 因子。更准确地说，推论~\ref{cor:quintic-hypergeometric-mirror-equality} 将证明
\[
\lim_{\lambda\to0}HG[\mathcal I(\widehat Q)](\Psi(t_1))
=\frac1{y_0}\lim_{\lambda\to0}HG[\mathcal I(\widehat P)](t_1).
\]
把它与式~\eqref{eq:euler-q-limit}、\eqref{eq:euler-p-limit} 结合并比较系数，很容易得到恒等式
\begin{equation}\label{eq:lly-potential}
\Phi(\Psi(t_1))=\frac52\left(\frac{y_1y_2}{y_0y_0}-\frac{y_3}{y_0}\right).
\end{equation}
这就是 \SourceCite{LLY} 所给出的五次三维簇\IndexLocation{indexloc:0864:51}\IndexLocation{indexloc:0867:51}镜像定理。

由于 $\Phi(s)$ 是式~\eqref{eq:gw-potential-quintic} 的 Gromov--Witten 势，公式~\eqref{eq:lly-potential} 使我们能计算任意 $d$ 的 $n_d$。但它与三点函数有什么关系？我们已指出，$d^3\Phi/ds^3$ 就是 A 模型关联函数。五次三维簇的镜像对称断言，经由镜像映射，这个函数等于
\SourcePage{52}
B 模型关联函数 $Y$。因此，只要证明
\begin{equation}\label{eq:lly-yukawa}
\frac{d^3}{ds^3}\frac52\left(\frac{y_1y_2}{y_0y_0}-\frac{y_3}{y_0}\right)=Y,
\end{equation}
式~\eqref{eq:lly-potential} 就会推出式~\eqref{eq:quintic-mirror-identity} 所述的镜像对称表述。这里，与前面一样，$y_i$ 经由镜像映射被看作 $s$ 的函数。

为证明式~\eqref{eq:lly-yukawa}，先用式~\eqref{eq:period-derivative-identities} 得到
\[
\frac{d^2}{ds^2}\frac52\left(\frac{y_1y_2}{y_0y_0}-\frac{y_3}{y_0}\right)
=\frac{d^2}{ds^2}\frac52\left(s\frac{y_2}{y_0}-\frac{y_3}{y_0}\right)
=5\frac d{ds}\frac{y_2}{y_0}.
\]
这个直接的证明留给读者。然后再求一次导数，并再次使用式~\eqref{eq:period-derivative-identities}，式~\eqref{eq:lly-yukawa} 立即得证。至此，五次三维簇的镜像对称证明便完成了；当然，从这里到\hyperref[chap:mirror-theorem]{第~\ref*{chap:mirror-theorem} 章}之间，仍有大量实质性的工作要做。

请注意，在这条证明镜像定理的路线中，归一化 Yukawa 耦合 $Y$ 直到计算的最后才出现；此前所有步骤都涉及新的研究对象。Givental 的方法也是如此，只是中间对象略有不同。因此，在第 \hyperref[chap:toric]{\ref*{chap:toric}}--\hyperref[chap:quantum-cohomology]{\ref*{chap:quantum-cohomology}} 章建立表述镜像对称所需的数学时，读者应记住：要在\hyperref[chap:mirror-theorem]{第~\ref*{chap:mirror-theorem} 章}证明镜像对称，还需要一些根本性的新思想。这会花费大量功夫，但最终结果之奇妙，足以回报这些努力。

\par\endgroup

\chapter{环面几何}\label{chap:toric}
\SourcePage{54}

本章将研究复数域上的环面簇理论。我们假定读者对用锥与扇描述环面簇的经典方法已有一定了解，参见 \SourceCite{Danilov,Fulton3,Oda}；不过，我们也会介绍较近出现的、利用多胞形与齐次坐标构造环面簇的方法。我们还将讨论环面几何中与镜像对称关系最密切的内容，包括 Kähler 锥、辛约化、GKZ 分解、Fano 环面簇与自反多胞形，以及自同构群。本章最后给出一些例子。

\section{锥与扇}\label{sec:cones-fans}

首先概述环面簇，主要目的是确定记号和术语。有关证明可参见 \SourceCite{Danilov,Fulton3,Oda}。

设 $M\simeq\mathbb Z^n$ 是秩为 $n$ 的自由 Abel 群，$N=\operatorname{Hom}(M,\mathbb Z)$ 是它的对偶群。$m\in M$ 与 $v\in N$ 的配对记为 $\langle m,v\rangle\in\mathbb Z$。$N_{\mathbb R}=N\otimes\mathbb R$ 中的\emph{有理多面锥} $\sigma$ 是形如
\[
\sigma=\left\{\sum_{i=1}^{s}\lambda_i u_i:\lambda_i\geq0\right\}
\]
的子集，其中 $u_1,\ldots,u_s\in N$。若 $\sigma\cap(-\sigma)=\{0\}$，则称 $\sigma$\emph{强凸}。$\sigma$ 的维数是它在 $N_{\mathbb R}$ 中张成的子空间的维数。每个锥 $\sigma$ 都有\emph{对偶锥}
\[
\check\sigma=\{m\in M_{\mathbb R}:\langle m,v\rangle\geq0\text{ 对所有 }v\in\sigma\}.
\]
标准理论表明，$\check\sigma$ 也是有理多面锥。给定 $m\in M\cap\check\sigma$，子集
\[
\tau=\{v\in\sigma:\langle m,v\rangle=0\}\subset\sigma
\]
称为 $\sigma$ 的一个\emph{面}。每个面仍是有理多面锥；面的面仍是原锥的面。

$N_{\mathbb R}$ 中的\emph{扇} $\Sigma$ 是一个由有限多个强凸有理多面锥组成的集合，满足：
\begin{itemize}
\item 若 $\sigma\in\Sigma$，则 $\sigma$ 的每个面也属于 $\Sigma$；
\item 若 $\sigma,\tau\in\Sigma$，则 $\sigma\cap\tau$ 是二者各自的一个面。
\end{itemize}
集合 $|\Sigma|=\bigcup_{\sigma\in\Sigma}\sigma\subset N_{\mathbb R}$ 称为 $\Sigma$ 的\emph{支集}。对每个 $d$，用 $\Sigma(d)$ 表示 $\Sigma$ 中所有 $d$ 维锥的集合。

一维锥有时称为扇的\emph{一维骨架}，尤为重要。我们约定用 $\rho$ 表示 $\Sigma(1)$ 中的元素。对每个 $\rho$，令 $v_\rho$ 为半群 $\rho\cap N$ 的唯一生成元。借助这些生成元，$\Sigma$ 中的锥 $\sigma$ 可写为
\[
\sigma=\left\{\sum_{\rho\subset\sigma}\lambda_\rho v_\rho:\lambda_\rho\geq0\right\}.
\]

\SourcePage{55}
属于 $\sigma$ 的各 $v_\rho$ 称为 $\sigma$ 的生成元。按照惯常的记号简化，我们常把 $\rho$ 与 $v_\rho$ 等同起来，于是也写 $v_\rho\in\Sigma(1)$。若 $r=|\Sigma(1)|$ 是一维锥的数目，有时还把这些 $v_\rho$ 记为 $v_1,\ldots,v_r$。\index{扇|hyperpage}\index{有理多面锥|hyperpage}\index{环面簇|hyperpage}

下面回顾环面簇的经典定义。给定扇 $\Sigma$，每个锥 $\sigma\in\Sigma$ 都给出仿射环面簇
\[
X_\sigma=\operatorname{Spec}\bigl(\mathbb C[M\cap\check\sigma]\bigr).
\]
这里，$\mathbb C[M\cap\check\sigma]$ 是这样一个 $\mathbb C$ 代数：对每个 $m\in M\cap\check\sigma$ 有生成元 $\chi^m$，并满足关系 $\chi^m\chi^{m'}=\chi^{m+m'}$。沿着 $X_{\sigma\cap\tau}$ 将各 $X_\sigma$ 与 $X_\tau$ 粘合，便得到环面簇 $X_\Sigma$。

若 $\Sigma$ 是 $N_{\mathbb R}\simeq\mathbb R^n$ 中的扇，则 $X_\Sigma$ 是维数为 $n$ 的 Cohen--Macaulay 代数簇。由于平凡锥 $\{0\}$ 对应的仿射环面簇是代数环面 $T_N=N\otimes\mathbb C^*=\operatorname{Spec}(\mathbb C[M])$，故 $T_N$ 是 $X_\Sigma$ 的仿射开子集。环面在自身上的作用可延拓为 $T_N$ 在 $X_\Sigma$ 上的代数作用，“环面簇”由此得名。在不会混淆时，我们简写 $X$，不写 $X_\Sigma$。

环面在 $X$ 上的作用只有有限多个轨道，轨道闭包与 $\Sigma$ 中的锥之间存在反转包含关系的一一对应。特别地，每个 $\rho\in\Sigma(1)$ 对应一个不可约的 $T_N$ 不变除子 $D_\rho\subset X$。此外，$m\in M$ 给出一个特征标 $\chi^m:T_N\to\mathbb C^*$。把 $\chi^m$ 看作 $X$ 上的有理函数，就有
\begin{equation}\label{eq:toric-principal-divisor}
\operatorname{div}(\chi^m)=\sum_\rho\langle m,v_\rho\rangle D_\rho .
\end{equation}

扇 $\Sigma$ 的性质对环面簇 $X$ 的几何性质有深刻影响。例如：
\begin{itemize}
\item $X$ 是完备的，即紧的，当且仅当 $|\Sigma|=N_{\mathbb R}$；这样的扇称为\emph{完备扇}。
\item $X$ 光滑，即无奇点，当且仅当 $\Sigma$ 中每个锥的生成元都是 $N$ 的某组 $\mathbb Z$ 基的一部分；这样的扇称为\emph{光滑扇}。
\item $X$ 是 轨形，即 $V$ 流形，当且仅当 $\Sigma$ 中每个锥的生成元在 $\mathbb R$ 上线性无关；此时称 $\Sigma$ 和 $X$ 为\emph{单纯}的。对环面簇 $X$，“$X$ 单纯 $\Rightarrow$ $X$ 是轨形”是初等的；反之，“$X$ 是轨形 $\Rightarrow$ $X$ 单纯”则深刻得多，其证明见 M.~Brion, \emph{Rational smoothness and fixed points of torus actions}, Transformation Groups 4 (1999), 127--156。
\end{itemize}
一般的环面簇只有较温和的奇点，因为它是 Cohen--Macaulay 的；不过，我们有时仍会要求 $X$ 单纯。\hyperref[app:singular-varieties]{附录 A} 将说明，\IndexLocation{indexloc:1375:55}单纯环面簇在许多方面具有流形的性质。

由 Weil 除子模去线性等价所得的 Chow 群 $A_{n-1}(X)$，可直接通过扇计算。我们始终假定一维锥张成 $N_{\mathbb R}$。由式~\eqref{eq:toric-principal-divisor} 得到正合列
\begin{equation}\label{eq:toric-chow-sequence}
0\longrightarrow M\longrightarrow\mathbb Z^{\Sigma(1)}
\longrightarrow A_{n-1}(X)\longrightarrow0 .
\end{equation}
其中 $m\in M$ 映到 $(\langle m,v_\rho\rangle)\in\mathbb Z^{\Sigma(1)}$，而 $(a_\rho)\in\mathbb Z^{\Sigma(1)}$ 映到除子 $\sum_\rho a_\rho D_\rho$ 的类。因此 $A_{n-1}(X)$ 的秩为 $r-n$。

$A_{n-1}(X)$ 中包含 Picard 群 $\operatorname{Pic}(X)$，它由 Cartier 除子模去线性等价构成。Weil 除子 $D=\sum_\rho a_\rho D_\rho$ 是 Cartier 除子，当且仅当对每个 $\sigma\in\Sigma$，都有 $m_\sigma\in M$，使得只要 $\rho\subset\sigma$ 就有 $\langle m_\sigma,v_\rho\rangle=-a_\rho$。函数 $\phi_D:|\Sigma|\to\mathbb R$，在 $v\in\sigma$ 时定义为 $\phi_D(v)=\langle m_\sigma,v\rangle$，称为 $D$ 的\emph{\IndexLocation{indexloc:1314:55}支撑函数}。

若 $X$ 光滑，则 $\operatorname{Pic}(X)=A_{n-1}(X)$。在单纯情形，Picard 群是 Chow 群的有限指数子群。$X$ 完备时 Picard 群总是无挠的；而 $A_{n-1}(X)$ 即使在 $X$ 单纯时也可能有挠。

\SourcePage{56}
若 $X$ 完备且除子 $D=\sum_\rho a_\rho D_\rho$ 是 Cartier 除子，可按以下条件判定 $D$ 是否丰沛，或是否由整体截面生成：
\begin{itemize}
\item $D$ 由整体截面生成，当且仅当对任意 $\rho\not\subset\sigma$，均有 $\langle m_\sigma,v_\rho\rangle\geq-a_\rho$。
\item $D$ 丰沛，当且仅当对任意 $n$ 维锥 $\sigma$ 及 $\rho\not\subset\sigma$，均有 $\langle m_\sigma,v_\rho\rangle>-a_\rho$。
\end{itemize}
第一种情形称 $\phi_D$\emph{上凸}，第二种情形称它\emph{严格上凸}。下一节将用某些多胞形解释这些条件。

\section{多胞形与齐次坐标}\label{sec:polytopes-coordinates}

上一节给出的环面簇经典定义，需要将仿射环面簇粘合起来。近来又发现了其他构造方法。本节研究其中两种：一种来自多胞形，另一种推广了 $\mathbb P^n$ 的通常齐次坐标。涉及辛约化的第三种构造将在\hyperref[subsec:toric-symplectic]{第 3.3.3 小节}讨论。

\subsection{多胞形}\label{subsec:toric-polytopes}

先确定术语。$M_{\mathbb R}$ 中的\emph{多胞形} $\Delta$ 是有限个点的凸包。$\Delta$ 的维数，是所有差 $m_1-m_2$（$m_1,m_2\in\Delta$）所张成子空间的维数。若 $\Delta$ 的顶点都属于 $M$，即 $\Delta$ 是 $M$ 中有限子集的凸包，则称其为\emph{整多胞形}。\emph{刻面}是余维数为一的面。

给定 $M_{\mathbb R}$ 中的多胞形 $\Delta_1,\ldots,\Delta_k$，可以用几种方法构造新多胞形。例如，它们并集的凸包记为
\[
\operatorname{Conv}(\Delta_1,\ldots,\Delta_k).
\]
还可定义\emph{Minkowski 和}：
\[
\Delta_1+\cdots+\Delta_k
=\{m_1+\cdots+m_k:m_i\in\Delta_i\}.
\]
若 $k\Delta=\Delta+\cdots+\Delta$（共 $k$ 项），则由凸性可知 $k\Delta=\{km:m\in\Delta\}$。\hyperref[sec:toric-complete-intersections]{第 4.3 节}将讨论的镜像对称构造，关键地用到了凸包与 Minkowski 和。

研究环面簇时，多胞形自然出现。例如，若 $X$ 完备，$D=\sum_\rho a_\rho D_\rho$ 是 Cartier 除子，则
\begin{equation}\label{eq:divisor-polytope}
\begin{aligned}
\Delta_D
&=\{m\in M_{\mathbb R}:\langle m,v\rangle\geq\phi_D(v)
\text{ 对所有 }v\in N_{\mathbb R}\}\\
&=\{m\in M_{\mathbb R}:\langle m,v_\rho\rangle\geq-a_\rho
\text{ 对所有 }\rho\}
\end{aligned}
\end{equation}
是一个多胞形。\footnote{即使 $D$ 不是 Cartier 除子，也可用第二个等式定义 $\Delta_D$。}
此外，$\Delta_{kD}=k\Delta_D$，$\Delta_{D+\operatorname{div}(\chi^m)}=\Delta_D-m$，并且：
\begin{itemize}
\item $D$ 由整体截面生成，当且仅当 $\Delta_D$ 是集合 $\{m_\sigma:\sigma\in\Sigma(n)\}$ 的凸包；
\item $D$ 丰沛，当且仅当对 $\Sigma(n)$ 中不同的 $\sigma,\tau$ 有 $m_\sigma\ne m_\tau$，且 $\Delta_D$ 是以 $\{m_\sigma:\sigma\in\Sigma(n)\}$ 为顶点集合的 $n$ 维多胞形。
\end{itemize}
在这两种情形，$\Delta_D$ 都是整多胞形。此外，当 $D$ 丰沛时，$\Delta_D$ 的非空面与 $\Sigma$ 中的锥一一对应。特别地，$\Delta_D$ 的刻面对应于 $\Sigma(1)$ 中的元素。这些断言的证明可参见 \SourceCite{Danilov,Fulton3,Oda}。\index{多胞形|hyperpage}\index{Minkowski和@Minkowski 和|hyperpage}

\SourcePage{57}
另一个有用的事实是，$\Delta_D\cap M$ 与 $H^0(X,\mathcal O(D))$ 中的 $T_N$ 特征向量自然对应。换言之，有一个 $T_N$ 等变同构
\begin{equation}\label{eq:toric-sections}
H^0(X,\mathcal O(D))\simeq
\bigoplus_{m\in\Delta_D\cap M}\mathbb C\chi^m .
\end{equation}
要理解这个等式，可将 $\mathcal O(D)$ 的截面看作 $X$ 上满足 $\operatorname{div}(f)+D\geq0$ 的有理函数 $f$。由式~\eqref{eq:toric-principal-divisor}，$\operatorname{div}(\chi^m)+D\geq0$ 等价于对所有 $\rho$ 均有 $\langle m,v_\rho\rangle\geq-a_\rho$。上述同构由此很容易推出，详情参见 \SourceCite[第 3.4 节]{Fulton3}。

\subsection{由多胞形构造环面簇}\label{subsec:toric-from-polytopes}

下面说明如何从一个 $n$ 维整多胞形出发构造环面簇。我们采用 \SourceCite{Batyrev4} 的方法。给定 $\Delta$，考虑 $m\in k\Delta\cap M$ 时的“单项式” $t_0^k\chi^m$。它们按规则
$t_0^k\chi^m\cdot t_0^\ell\chi^{m'}
=t_0^{k+\ell}\chi^{m+m'}$
相乘；因为 $m\in k\Delta\cap M$、$m'\in\ell\Delta\cap M$ 意味着 $m+m'\in(k+\ell)\Delta\cap M$。由这些 $t_0^k\chi^m$ 生成的 $\mathbb C$ 代数记为 $S_\Delta$，并规定
\[
\deg(t_0^k\chi^m)=k
\]
来赋予它分次。我们称 $S_\Delta$ 为 $\Delta$ 的\emph{多胞形环}。令
\[
\mathbb P_\Delta=\operatorname{Proj}(S_\Delta)
\]
为相应的射影簇。在语境清楚时，用 $\mathbb P$ 简记 $\mathbb P_\Delta$。

要把 $\mathbb P=\mathbb P_\Delta$ 描述为环面簇，须给出它在 $N_{\mathbb R}$ 中的扇。给定非空面 $F\subset\Delta$，考虑锥
\[
\check\sigma_F
=\{\lambda(m-m'):m\in\Delta,\ m'\in F,\ \lambda\geq0\}
\subset M_{\mathbb R}.
\]
它的对偶是 $N_{\mathbb R}$ 中的锥 $\sigma_F$。所有这些锥组成扇
$\{\sigma_F:F\text{ 是 }\Delta\text{ 的非空面}\}$，称为 $\Delta$ 的\emph{法扇}。这是 $N_{\mathbb R}$ 中的完备扇，而且可以证明 $\mathbb P$ 正是由该法扇确定的环面簇 \SourceCite{Batyrev4}。\index{法扇|hyperpage}

当原点是 $\Delta$ 的内点时，法扇特别容易形象化。先定义\emph{\IndexLocation{indexloc:1077:57}极多胞形}
\[
\Delta^\circ
=\{v\in N_{\mathbb R}:\langle m,v\rangle\geq-1
\text{ 对所有 }m\in\Delta\}\subset N_{\mathbb R}.
\]
$\Delta^\circ$ 也以原点为内点，且 $(\Delta^\circ)^\circ=\Delta$。此外，$\Delta$ 的每个 $i$ 维面 $F$ 对应于 $\Delta^\circ$ 的一个 $(n-1-i)$ 维面 $F^\circ$；其特征是，当 $m'\in F$、$v\in F^\circ$ 时有 $\langle m',v\rangle=-1$。利用 $\Delta^\circ$，可如下描述 $\mathbb P$。

\begin{lemma}\label{lem:normal-polar}
在 $N_{\mathbb R}$ 中，由 $\Delta^\circ$ 的各真面上的锥组成的扇，就是 $\Delta$ 的法扇，因此给出环面簇 $\mathbb P=\mathbb P_\Delta$。
\end{lemma}
\begin{proof}
只须证明：对 $\Delta$ 的每个真面 $F$，
$\sigma_F=\{\lambda v:v\in F^\circ,\ \lambda\geq0\}$。
首先，$v\in\sigma_F$ 当且仅当对所有 $m\in\Delta$ 和 $m'\in F$，都有 $\langle m,v\rangle\geq\langle m',v\rangle$。还要注意，若 $v\ne0$，则 $\langle m',v\rangle<0$；否则就会对所有 $m\in\Delta$ 有 $\langle m,v\rangle\geq0$，这与原点是 $\Delta$ 的内点矛盾。

若 $v\in F^\circ$，则 $\langle m,v\rangle\geq\langle m',v\rangle=-1$，所以容易得到 $\lambda v\in\sigma_F$。反过来，设对所有 $m\in\Delta$、$m'\in F$ 均有 $\langle m,v\rangle\geq\langle m',v\rangle$。可设 $v\ne0$。固定 $F$ 相对内部的一个 $m'$；由上面的观察，$\langle m',v\rangle<0$。因此可写 $v=\lambda v'$，其中 $\lambda>0$，并且对所有 $m\in\Delta$ 有 $\langle m,v'\rangle\geq\langle m',v'\rangle=-1$。于是 $v'\in\Delta^\circ$；又因 $m'$ 在 $F$ 的相对内部，$\langle m',v'\rangle=-1$ 进一步表明 $v'\in F^\circ$。
\end{proof}

\SourcePage{58}
由引理~\ref{lem:normal-polar} 可知，在 $\mathbb P_\Delta$ 中，$T_N$ 轨道闭包与 $\Delta^\circ$ 的真面一一对应；再由极对偶性，它们也与 $\Delta$ 的非空面 $F$ 一一对应。实际上，从多胞形环出发也能直接看出这种对应。

设 $F\subset\Delta$ 是一个面。多胞形环之间有典范满射 $\pi_F:S_\Delta\to S_F$，定义为
\[
\pi_F(t_0^k\chi^m)=
\begin{cases}
t_0^k\chi^m,&m\in kF,\\
0,&m\in k\Delta-kF.
\end{cases}
\]
这个满射诱导环面簇的自然嵌入 $\mathbb P_F\hookrightarrow\mathbb P_\Delta$。若 $S_\Delta$ 由 $\Delta\cap M$ 生成，这个嵌入可以描述得特别简单：可把 $\mathbb P_\Delta$ 看成嵌入一个射影空间，其齐次坐标与 $\Delta\cap M$ 中的点一一对应。此时，只需把对应于 $(\Delta-F)\cap M$ 中各点的坐标全部置零，便得到 $\mathbb P_F$。

若从一个以原点为内点的整多胞形 $\Delta$ 出发，其极多胞形 $\Delta^\circ$ 未必是整的。\hyperref[sec:fano-reflexive]{第 3.5 节}将研究 Batyrev 的一类特殊多胞形，即极多胞形也是整的\emph{自反多胞形}。这类多胞形在许多镜像对称构造中起重要作用。\index{自反多胞形|hyperpage}

再回到任意 $n$ 维整多胞形 $\Delta$。我们注意到，环面簇 $\mathbb P=\mathbb P_\Delta$ 自然带有一个丰沛线丛：$\mathbb P=\operatorname{Proj}(S_\Delta)$ 带有丰沛的 $\mathcal O_{\mathbb P}(1)$。从环面簇的观点，这个线丛可描述如下。定义函数 $\phi_\Delta:N_{\mathbb R}\to\mathbb R$：
\[
\phi_\Delta(v)=\min_{m\in\Delta}\langle m,v\rangle .
\]
可以证明，$\phi_\Delta$ 是 $\mathbb P$ 上某个除子 $D_\Delta$ 的支撑函数。而且，$\phi_\Delta$ 严格上凸，$\mathcal O_{\mathbb P}(D_\Delta)\simeq\mathcal O_{\mathbb P}(1)$，与 $D_\Delta$ 对应的多胞形 $\Delta_{D_\Delta}$ 恰好是 $\Delta$，参见 \SourceCite{Batyrev4,Oda}。所以，由多胞形构造出的环面簇，是附带选定丰沛除子的射影簇。

若选定 $M$ 的一组基，就能更具体地观察 $\mathbb P_\Delta$ 及其射影嵌入。这等于为环面 $T_N$ 选取坐标 $t_1,\ldots,t_n$。当 $m\in M$ 写为 $m=(a_1,\ldots,a_n)$ 时，有 $\chi^m=\prod_{i=1}^{n}t_i^{a_i}$，因此可以写 $t^m$ 来代替 $\chi^m$。给定多胞形 $\Delta$，选择 $k$ 使 $kD_\Delta$ 十分丰沛。格点 $k\Delta\cap M=\{m_1,\ldots,m_s\}$ 给出单项式 $t^{m_i}$，从而定义映射 $T_N\to\mathbb P^{s-1}$，把 $t\in T_N$ 送到 $(t^{m_1},\ldots,t^{m_s})$。$\mathbb P_\Delta$ 就是此映射的像在 $\mathbb P^{s-1}$ 中的闭包。

另外，对任意 $k\geq0$，有 Laurent 多项式的向量空间
\begin{equation}\label{eq:laurent-space}
L(k\Delta)=\left\{f:f=\sum_{m\in k\Delta\cap M}a_mt^m,\ a_m\in\mathbb C\right\}.
\end{equation}
每个 $f\in L(k\Delta)$ 都定义了 $T_N$ 中由 $f=0$ 给出的仿射超曲面 $Z_f$；其在 $\mathbb P_\Delta$ 中的闭包 $\overline Z_f$ 是相应的除子。这样，把 $M$ 中的元素视为 Laurent 单项式的指数向量，便可用坐标研究 $\mathbb P_\Delta$。有时人们从 Laurent 多项式 $f$ 出发，将 $\Delta$ 定义为其指数向量的凸包，这就是 $f$ 的\emph{Newton 多胞形}。我们使用以下记号：
\begin{equation}\label{eq:lattice-points}
\begin{aligned}
l(k\Delta)&=|k\Delta\cap M|=\dim L(k\Delta),\\
l^*(k\Delta)&=|\{m\in k\Delta\cap M:m\text{ 不在 }k\Delta\text{ 的任何刻面上}\}|.
\end{aligned}
\end{equation}

\SourcePage{59}
数 $l(k\Delta)$ 和 $l^*(k\Delta)$ 具有许多有趣的性质，并在镜像对称的多种情境中出现。

最后对前面定义的多胞形环 $S_\Delta$ 作几点说明。首先，对 $f\in L(k\Delta)$ 而言，多项式 $t_0^kf$ 恰好是 $S_\Delta$ 中次数为 $k$ 的元素。这给出同构
\begin{equation}\label{eq:polytope-ring-sections}
(S_\Delta)_k\simeq L(k\Delta)\simeq
H^0(\mathbb P,\mathcal O_{\mathbb P}(kD_\Delta)),
\end{equation}
其中第二个同构由式~\eqref{eq:toric-sections} 得到。

还可将 $S_\Delta$ 解释为仿射坐标环。令 $\check\sigma_\Delta\subset\mathbb R\times M_{\mathbb R}$ 为 $\{1\}\times\Delta$ 上的锥。注意 $(k,m)\in\check\sigma_\Delta\cap(\mathbb Z\times M)$ 当且仅当 $m\in k\Delta$。因此，若 $\sigma_\Delta\subset\mathbb R\times N_{\mathbb R}$ 是 $\check\sigma_\Delta$ 的对偶，则 $S_\Delta$ 正是由 $\sigma_\Delta$ 给出的仿射环面簇的坐标环。由 \SourceCite{Danilov} 可知 $S_\Delta$ 是 Cohen--Macaulay 的。

\subsection{齐次坐标环}\label{subsec:homogeneous-ring}

紧接在式~\eqref{eq:laurent-space} 之前的讨论中，我们在环面簇的环面上引入了坐标 $t_1,\ldots,t_n$。这些坐标非常有用；但若能像射影空间上一样拥有整体坐标，就更方便了。为此，我们定义环面簇的齐次坐标环。

若 $X=X_\Sigma$ 由 $N_{\mathbb R}$ 中的扇 $\Sigma$ 给出，就为每个 $\rho\in\Sigma(1)$ 引入变量 $x_\rho$，并考虑多项式环
\[
S=\mathbb C[x_\rho:\rho\in\Sigma(1)].
\]
$S$ 中的单项式写作 $x^D=\prod_\rho x_\rho^{a_\rho}$，其中 $D=\sum_\rho a_\rho D_\rho$ 是 $X$ 上有效的环面不变除子；这里用到了对应 $\rho\leftrightarrow D_\rho$。规定 $x^D$ 的次数为
\[
\deg(x^D)=[D]\in A_{n-1}(X).
\]
于是 $S$ 按 Chow 群 $A_{n-1}(X)$ 分次。给定除子类 $\alpha\in A_{n-1}(X)$，以 $S_\alpha$ 表示 $S$ 中次数为 $\alpha$ 的分次部分。我们也常把变量写成 $x_1,\ldots,x_r$，其中 $x_i$ 对应锥生成元 $v_i$，$r=|\Sigma(1)|$；这时 $S=\mathbb C[x_1,\ldots,x_r]$。

\begin{definition}\label{def:toric-homogeneous-ring}
环 $S$ 连同上面定义的分次，称为 $X$ 的\emph{齐次坐标环}。
\end{definition}

\hyperref[sec:toric-examples]{第 3.7 节}将看到，$\mathbb P^n$ 的齐次坐标环就是带通常分次的 $\mathbb C[x_0,\ldots,x_n]$。对 $\mathbb P^n\times\mathbb P^m$，坐标环则是按双齐次多项式的相应次数分次的环。在与镜像对称有关的物理理论中，变量 $x_\rho\in S$ 对应某些可观测量，尽管单个可观测量本身并没有物理意义。所使用的是由它们组成的某些单项式，再模去 $S$ 的一个理想；这就是\emph{手征环}。

注意，分次环 $S$ 只依赖于 $\Sigma$ 的一维骨架。因此，可能存在许多具有同一一维骨架的扇，从而有许多具有同一坐标环的环面簇。\hyperref[sec:gkz-decomposition]{第~\ref*{sec:gkz-decomposition} 节}将进一步阐明这一点。

为在 $S$ 中体现整个扇 $\Sigma$，定义理想 $B(\Sigma)$。对 $\sigma\in\Sigma$，令
\[
\widehat x_\sigma=\prod_{\rho\not\subset\sigma}x_\rho
\]
为不属于 $\sigma$ 的变量之积，再定义理想
\[
B(\Sigma)=(\widehat x_\sigma:\sigma\in\Sigma)\subset S.
\]

\SourcePage{60}
对 $\mathbb P^n$，这个理想就是“无关”理想 $\langle x_0,\ldots,x_n\rangle$。一般而言，$B(\Sigma)$ 唯一地确定 $\Sigma$。

$S$ 的各分次部分还有一个很好的上同调解释。若除子 $D=\sum_\rho a_\rho D_\rho$ 不是 Cartier 除子，它就不能确定线丛，但仍给出一个自反层 $\mathcal O_X(D)$。这是 $X$ 上有理函数层的子层：它在开集 $U\subset X$ 上的截面，是在 $U$ 上极点阶数不超过 $D$ 所规定阶数的有理函数。关于 Weil 除子所确定的层，详见 \SourceCite[第 1 节附录]{Reid1}。如式~\eqref{eq:polytope-ring-sections} 所示，有同构
\[
H^0(X,\mathcal O_X(D))\simeq L(\Delta_D),
\]
这里 $\Delta_D$ 由式~\eqref{eq:divisor-polytope} 的第二行定义。由此容易看出，映射
$t^m\mapsto\prod_\rho x_\rho^{\langle m,v_\rho\rangle+a_\rho}$
诱导同构
\begin{equation}\label{eq:toric-graded-sections}
H^0(X,\mathcal O_X(D))\simeq S_\alpha,
\end{equation}
其中 $\alpha=[D]\in A_{n-1}(X)$。

在射影情形，可将它与式~\eqref{eq:polytope-ring-sections} 的同构结合。于是，若 $D$ 丰沛且 $\alpha=[D]$，就得到环同构
\[
S_{\Delta_D}\simeq\bigoplus_{k=0}^{\infty}S_{k\alpha};
\]
这里 $S_{\Delta_D}$ 是\hyperref[subsec:toric-from-polytopes]{第 3.2.1 小节}的多胞形环。因此，$S$ 包含了来自 $X$ 上所有可能丰沛除子 $D$ 的多胞形环 $S_{\Delta_D}$。

还可以用齐次坐标描述 $X$ 上的微分形式。\hyperref[chap:hodge-yukawa]{第 5 章}计算 $X$ 中丰沛超曲面的 Yukawa 耦合时，将采用这一描述。

\subsection{由齐次坐标构造环面簇}\label{subsec:toric-from-homogeneous}

下面说明怎样把射影空间的通常构造
$\mathbb P^n=(\mathbb C^{n+1}-\{0\})/\mathbb C^*$
推广到任意环面簇。如\hyperref[sec:cones-fans]{第 3.1 节}一样，假定 $\Sigma$ 的一维锥张成 $N_{\mathbb R}$；若 $\Sigma$ 完备，这一条件自动成立。定义群
\[
G=\operatorname{Hom}_{\mathbb Z}(A_{n-1}(X),\mathbb C^*).
\]
对正合列~\eqref{eq:toric-chow-sequence} 施加 $\operatorname{Hom}_{\mathbb Z}(-,\mathbb C^*)$，得到正合列
\begin{equation}\label{eq:toric-g-sequence}
1\longrightarrow G\longrightarrow(\mathbb C^*)^{\Sigma(1)}
\longrightarrow T_N\longrightarrow1.
\end{equation}
这里用到了
$\operatorname{Hom}_{\mathbb Z}(M,\mathbb C^*)
=N\otimes\mathbb C^*=T_N$。
所以，$G$ 同构于 $(\mathbb C^*)^{r-n}$ 与一个有限 Abel 群的乘积；当且仅当 $A_{n-1}(X)$ 有挠时，才出现这个有限群。于是，$G\to(\mathbb C^*)^{\Sigma(1)}$ 给出 $G$ 在 $\mathbb C^{\Sigma(1)}=\operatorname{Spec}(S)$ 上的作用。若 $a=(a_\rho)\in\mathbb C^{\Sigma(1)}$、$g\in G$，则
\[
g\cdot a=(g([D_\rho])a_\rho).
\]
按 $G$ 的定义，$g:A_{n-1}(X)\to\mathbb C^*$。另外，理想 $B(\Sigma)\subset S$ 在 $\mathbb C^{\Sigma(1)}$ 中定义簇 $Z(\Sigma)$。这些数据按下述方式确定 $X$。

\begin{theorem}\label{thm:toric-quotient}
若 $\Sigma$ 的一维锥张成 $N_{\mathbb R}$，则：
\begin{enumerate}[label=(\roman*)]
\item $X$ 是 $\mathbb C^{\Sigma(1)}-Z(\Sigma)$ 关于 $G$ 的范畴商；
\item $X$ 是 $\mathbb C^{\Sigma(1)}-Z(\Sigma)$ 关于 $G$ 的几何商，当且仅当 $X$ 单纯。
\end{enumerate}
\end{theorem}

\SourcePage{61}
在单纯情形，这个结果最早由 \SourceCite{Audin} 证明；不过，这一构造的想法曾多次独立出现。有关证明和参考文献清单，参见 \SourceCite{Cox}。

定理~\ref{thm:toric-quotient} 给出的 $X$ 的描述，也使\hyperref[sec:cones-fans]{第 3.1 节}提到的轨道闭包与 $\Sigma$ 中锥之间反转包含关系的对应变得简单：对每个 $\sigma\in\Sigma$，考虑线性子空间
\[
\{x\in\mathbb C^{\Sigma(1)}:x_\rho=0
\text{ 对所有 }\rho\in\sigma\}.
\]
它在 $X_\Sigma$ 中下降为一个子簇，容易验证这就是与 $\sigma$ 对应的轨道闭包。特别地，与边 $v_i$ 对应的环面除子 $D_i$，正是方程 $x_i=0$ 所定义的超曲面。

定理~\ref{thm:toric-quotient} 中的簇 $Z(\Sigma)$ 有简单结构。若子集 $\mathcal S\subset\Sigma(1)$ 不是 $\Sigma$ 中某个锥的全体一维锥，但 $\mathcal S$ 的每个真子集都是某个锥的全体一维锥，就称 $\mathcal S$ 为\emph{本原集合}。这时
\begin{equation}\label{eq:toric-primitive-z}
Z(\Sigma)=\bigcup_{\mathcal S}V(\mathcal S),
\end{equation}
其中并集遍历所有本原集合 $\mathcal S$，$V(\mathcal S)$ 是对 $\rho\in\mathcal S$ 令 $a_\rho=0$ 所定义的子空间。证明见 \SourceCite{BC}。

若 $X$ 单纯，定理~\ref{thm:toric-quotient} 表明我们的确有了齐次坐标：$X$ 中一点的坐标可取为
$a=(a_\rho)\in\mathbb C^{\Sigma(1)}-Z(\Sigma)$；另一个
$b\in\mathbb C^{\Sigma(1)}-Z(\Sigma)$ 表示 $X$ 中同一点，当且仅当存在 $g\in G$，使 $b=g\cdot a$。特别地，若 $f\in S$ 是齐次的，即 $f\in S_\alpha$ 对某个 $\alpha$ 成立，那么 $f=0$ 就在 $X$ 中定义一个超曲面。更一般地，$S$ 的任意齐次理想 $I$ 都定义 $X$ 中的一个子簇 $V(I)$。环面簇的\IndexLocation{indexloc:1370:61}理想与簇的对应关系见 \SourceCite{Cox}。

在非单纯情形，$\mathbb C^{\Sigma(1)}-Z(\Sigma)$ 与 $X$ 的关系不那么直接，但 $S$ 的齐次理想仍能确定 $X$ 的子簇。还有一种值得注意的现象，即使在单纯情形也会出现：齐次方程 $f=0$ 给出的超曲面 $Y\subset X$ 未必是 Cartier 除子。这意味着，虽然 $Y$ 可由整体方程 $f=0$ 定义，在 $X$ 上的局部却未必能由单个方程定义。\index{齐次坐标环|hyperpage}\index{本原集合|hyperpage}

\section{Kähler 锥与辛几何}\label{sec:kahler-symplectic}

本节描述\IndexLocation{indexloc:1375:61}单纯环面簇的 Kähler 锥，并解释如何利用辛几何构造环面簇。

\subsection{环面簇的 Kähler 锥}\label{subsec:toric-kahler-cone}

对于光滑射影环面簇 $X$，有 $A_{n-1}(X)\simeq H^2(X,\mathbb Z)$，从而 $H^2(X,\mathbb R)=H^{1,1}(X,\mathbb R)$。更一般地，若 $X$ 单纯且射影，则 $A_{n-1}(X)\otimes\mathbb R\simeq H^2(X,\mathbb R)$，所以仍有 $H^2(X,\mathbb R)=H^{1,1}(X,\mathbb R)$。回顾\hyperref[chap:introduction]{第 1 章}，$H^{1,1}(X,\mathbb R)$ 中的 Kähler 锥由 $X$ 上所有 Kähler 类组成；轨形 情形下 Kähler 的定义见\hyperref[app:singular-varieties]{附录 A}。在上述同构下，我们得到 $A_{n-1}(X)\otimes\mathbb R$ 中的一个锥，也称为 Kähler 锥。

令 $A^+_{n-1}(X)\otimes\mathbb R$ 为由除子类 $[D_i]$（$v_i\leftrightarrow D_i$，$i=1,\ldots,r$）生成的锥。Kähler 锥包含在其中，描述如下。设
$a=\sum_{i=1}^{r}a_i[D_i]\in A^+_{n-1}(X)\otimes\mathbb R$。
因为 $X$ 单纯，对每个 $\sigma\in\Sigma$ 都可找到 $m_\sigma\in M_{\mathbb R}$，使得在 $v_i\in\sigma$ 时 $\langle m_\sigma,v_i\rangle=-a_i$。若对所有 $\sigma\in\Sigma$ 和 $1\leq i\leq r$ 均有 $\langle m_\sigma,v_i\rangle\geq-a_i$，则称 $a$ 为\emph{凸的}。

\SourcePage{62}
定义
\begin{equation}\label{eq:cpl-cone}
\operatorname{cpl}(\Sigma)=
\{a\in A^+_{n-1}(X)\otimes\mathbb R:a\text{ 是凸的}\}.
\end{equation}
由于 $X$ 射影且单纯，$\operatorname{cpl}(\Sigma)$ 是一个 $(r-n)$ 维凸锥。而且，$a$ 属于 $\operatorname{cpl}(\Sigma)$ 的内部，当且仅当 $a$ 严格凸，即只要 $\sigma$ 为 $n$ 维且 $v_i\notin\sigma$，就有 $\langle m_\sigma,v_i\rangle>-a_i$，参见 \SourceCite{OdP}。于是可作如下描述。

\begin{proposition}\label{prop:toric-kahler-cone}
若 $X=X_\Sigma$ 射影且单纯，则 $X$ 的 Kähler 锥等于式~\eqref{eq:cpl-cone} 所定义的锥 $\operatorname{cpl}(\Sigma)$ 的内部。
\end{proposition}
\begin{proof}
由 \SourceCite{Baily2}，Kähler 锥包含在 $\operatorname{cpl}(\Sigma)$ 的内部；反向包含关系由 \SourceCite[第 4 节]{AGM1} 证明。
\end{proof}

由于 $\operatorname{cpl}(\Sigma)$ 是强凸多面锥，命题~\ref{prop:toric-kahler-cone} 表明，射影单纯环面簇的 Kähler 锥具有特别良好的结构。

\subsection{Mori 锥}\label{subsec:toric-mori-cone}

若 $X_\Sigma$ 单纯且完备，它的\emph{Mori 锥} $M(X_\Sigma)$ 是
$A_1(X_\Sigma)\otimes\mathbb R\simeq H_2(X_\Sigma,\mathbb R)$
中有效一维圈组成的锥。有时也记为
$NE(X_\Sigma)_{\mathbb R}$ 或 $\overline{NE}(X_\Sigma)_{\mathbb R}$。
在射影情形，我们关注 $M(X_\Sigma)$，是因为它与 Kähler 锥对偶。因此，由命题~\ref{prop:toric-kahler-cone}，$\operatorname{cpl}(\Sigma)$ 与 $M(X_\Sigma)$ 互为对偶。

正如所料，可以用扇 $\Sigma$ 的环面数据显式描述 $M(X_\Sigma)$。为此，取 $(n-1)$ 维锥 $\sigma\in\Sigma(n-1)$；相应的轨道闭包 $C_\sigma\subset X_\Sigma$ 是一条曲线。根据 Reid 的一个观察 \SourceCite[命题 1.6]{Reid3}，$M(X_\Sigma)$ 由这些 $C_\sigma$ 生成。

可以把这一描述写得更明确。令 $v_1,\ldots,v_{n-1}$ 为 $\sigma$ 的整生成元。注意 $\sigma$ 恰好包含在两个 $n$ 维锥中：第一个由 $\sigma$ 与另一个本原整生成元 $v_n$ 张成；第二个则由 $\sigma$ 与另一生成元 $v_{n+1}$ 张成。因此有线性关系
\begin{equation}\label{eq:mori-wall-relation}
\sum_{i=1}^{n+1}\lambda_i v_i=0.
\end{equation}
由于 $v_n$、$v_{n+1}$ 分居 $\sigma$ 两侧，可设 $\lambda_1,\ldots,\lambda_{n+1}$ 为互素整数，且 $\lambda_n,\lambda_{n+1}>0$。于是 $\lambda_i$ 在下列空间中给出一个关系 $\lambda_\sigma$：
\begin{equation}\label{eq:toric-relation-space}
\Lambda_{\mathbb Q}=
\left\{(\lambda_\rho)_{\rho\in\Sigma(1)}:
\lambda_\rho\in\mathbb Q,\ 
\sum_\rho\lambda_\rho v_\rho=0\right\}.
\end{equation}
对式~\eqref{eq:toric-chow-sequence} 取对偶并张量 $\mathbb Q$，得到自然同构
\[
A_1(X_\Sigma)\otimes\mathbb Q\simeq\Lambda_{\mathbb Q}.
\]
在这个同构下，曲线 $C_\sigma\in A_1(X_\Sigma)\otimes\mathbb Q$ 与关系 $\lambda_\sigma\in\Lambda_{\mathbb Q}$ 的联系如下。

\begin{lemma}\label{lem:mori-wall-curve}
设 $\sigma\in\Sigma(n-1)$。那么存在正常数 $c_\sigma$，使上述同构将 $C_\sigma$ 对应到 $c_\sigma\lambda_\sigma$；这里 $\lambda_\sigma$ 是由式~\eqref{eq:mori-wall-relation} 确定的关系。
\end{lemma}

\SourcePage{63}
\begin{proof}
首先，交积使 $A_1(X_\Sigma)\otimes\mathbb Q$ 与 $A_{n-1}(X_\Sigma)\otimes\mathbb Q$ 对偶；而借助 $\mathbb Q^{\Sigma(1)}$ 上的点积以及式~\eqref{eq:toric-chow-sequence}、\eqref{eq:toric-relation-space}，$\Lambda_{\mathbb Q}$ 也与 $A_{n-1}(X_\Sigma)\otimes\mathbb Q$ 对偶。由于同构
$A_1(X_\Sigma)\otimes\mathbb Q\simeq\Lambda_{\mathbb Q}$
与这两个对偶相容，只须找到 $c_\sigma>0$，使对所有 $\rho\in\Sigma(1)$ 有
\[
D_\rho\cdot C_\sigma=c_\sigma D_\rho\cdot\lambda_\sigma.
\]
若 $\rho$ 不属于 $v_i$ 中的任何一个，易见等式两边都是 $0$。若 $\rho$ 是某个 $v_i$，则如通常那样，将 $\rho$ 与 $\rho\cap N$ 的唯一生成元等同，此时 $D_{v_i}\cdot\lambda_\sigma=\lambda_i$。而 \SourceCite[命题 2.7]{Reid3} 表明
\begin{equation}\label{eq:reid-wall-intersection}
\lambda_j D_{v_i}\cdot C_\sigma
=\lambda_i D_{v_j}\cdot C_\sigma,
\qquad 1\leq i,j\leq n+1.
\end{equation}
已知 $\lambda_{n+1}>0$；又因 $\sigma$ 与 $v_{n+1}$ 确定 $\Sigma$ 的一个 $n$ 维锥，$D_{v_{n+1}}\cap C_\sigma$ 只含一个点，即该 $n$ 维锥所对应的不动点。因此
$D_{v_{n+1}}\cdot C_\sigma>0$，从而
$c_\sigma=(D_{v_{n+1}}\cdot C_\sigma)/\lambda_{n+1}>0$。
用式~\eqref{eq:reid-wall-intersection} 很容易验证 $c_\sigma$ 具有所需性质。
\end{proof}

注意，若 $X_\Sigma$ 在 $C_\sigma$ 的某个邻域光滑，则引理~\ref{lem:mori-wall-curve} 中的 $c_\sigma=1$，见 \SourceCite[第 99 页]{Fulton3}。

引理~\ref{lem:mori-wall-curve} 的一个有用推论，是下面的丰沛性判别准则。
\begin{corollary}\label{cor:toric-ample-wall}
除子 $D=\sum_{\rho\in\Sigma(1)}a_\rho D_\rho$ 丰沛，当且仅当它是 Cartier 除子，且对每个 $\sigma\in\Sigma(n-1)$，均满足
$(a_\rho)_{\rho\in\Sigma(1)}\cdot\lambda_\sigma>0$。
\end{corollary}
\begin{proof}
设 $D$ 是 Cartier 除子。由于 $C_\sigma$ 生成 Mori 锥，$D$ 丰沛当且仅当对所有 $\sigma\in\Sigma(n-1)$ 有 $D\cdot C_\sigma>0$。由引理~\ref{lem:mori-wall-curve} 立即得到结论。
\end{proof}

若 $X_\Sigma$ 光滑，Mori 锥 $M(X_\Sigma)$ 还有另一种描述，它涉及\hyperref[subsec:toric-from-homogeneous]{第 3.2.4 小节}定义的\IndexLocation{indexloc:1088:63}本原集合。详情见 \SourceCite{Batyrev1}，相应的丰沛性判别准则见 \SourceCite{Batyrev2}。\index{Mori锥@Mori 锥|hyperpage}

\subsection{辛几何}\label{subsec:toric-symplectic}

辛几何在镜像对称中起重要作用。这里回顾一些基本定义；\SourceCite{Audin} 是一本很好的辛几何参考书。

\emph{辛流形}是赋有闭的、非退化二形式 $\omega$ 的实流形。辛结构把函数转化为向量场：若 $f$ 是流形上的 $C^\infty$ 函数，则流形上存在唯一的向量场 $X_f$，使任意向量场 $X$ 都满足 $\omega(X,X_f)=X(f)$。我们称 $X_f$ 为 $f$ 的\emph{Hamilton 向量场}。\index{辛流形|hyperpage}

对我们而言，最重要的例子是赋有辛形式
$\omega=\sum_{i=1}^{r}dx_i\wedge dy_i$
的 $\mathbb C^r$，其中 $z_i=x_i+iy_i$。$\mathbb C^r$ 的一个特殊性质是，$U(1)^r=(S^1)^r$ 在其上的自然作用是辛的，即对所有 $g\in U(1)^r$ 都有 $g^*\omega=\omega$。

事实上，$U(1)^r$ 的作用还是 Hamilton 作用，含义如下。$U(1)^r$ 的 Lie 代数是 $\mathbb R^r$。元素 $\lambda=(\lambda_1,\ldots,\lambda_r)\in\mathbb R^r$ 在指数映射下对应
$\exp(i\lambda)=(\exp(i\lambda_1),\ldots,\exp(i\lambda_r))\in U(1)^r$。
给定 $\lambda$，在 $\mathbb C^r$ 上得到流
$v\mapsto\exp(it\lambda)\cdot v$，其向量场为
\[
X_\lambda=\sum_{i=1}^{r}\lambda_i
\left(-y_i\frac{\partial}{\partial x_i}
+x_i\frac{\partial}{\partial y_i}\right).
\]

\SourcePage{64}
所谓 Hamilton 作用，就是说对 Lie 代数中的每个 $\lambda$，向量场 $X_\lambda$ 都是 $\mathbb C^r$ 上某个函数的 Hamilton 向量场。

我们用\emph{矩映射}
\[
\mu:\mathbb C^r\longrightarrow(\mathbb R^r)^*
\]
来证明这一点。它定义为
$\mu(z_1,\ldots,z_r)=\tfrac12(|z_1|^2,\ldots,|z_r|^2)$；这里使用的是与 $\mathbb R^r$ 的标准基对偶的 $(\mathbb R^r)^*$ 的基。关键在于，$\lambda=(\lambda_1,\ldots,\lambda_r)\in\mathbb R^r$ 可视为映射 $\lambda:(\mathbb R^r)^*\to\mathbb R$，将它与矩映射复合就得到 $\lambda\circ\mu:\mathbb C^r\to\mathbb R$。简单计算可得
\[
X_\lambda=X_{\lambda\circ\mu},
\]
这便证明了上述作用是 Hamilton 作用。\index{矩映射|hyperpage}

一般而言，矩映射与 Hamilton 作用在辛几何中起重要作用。\hyperref[chap:gromov-witten]{第~\ref*{chap:gromov-witten} 章}给出 Gromov--Witten 不变量的辛几何定义时，我们还会讨论辛几何的其他方面。

\subsection{通过辛约化构造环面簇}\label{subsec:toric-reduction}

下面说明环面簇如何通过称为\emph{辛约化}的过程与辛几何联系起来。设 $X=X_\Sigma$ 由 $N_{\mathbb R}\simeq\mathbb R^n$ 中的扇 $\Sigma$ 确定。如\hyperref[subsec:toric-from-homogeneous]{第 3.2.3 小节}所述，有群 $G=\operatorname{Hom}_{\mathbb Z}(A_{n-1}(X),\mathbb C^*)$。$G$ 的极大紧子群是
\[
G_{\mathbb R}=\operatorname{Hom}_{\mathbb Z}(A_{n-1}(X),U(1)),
\]
它的 Lie 代数为 $\mathfrak g_{\mathbb R}=\operatorname{Hom}_{\mathbb Z}(A_{n-1}(X),\mathbb R)$。因此 $\mathfrak g_{\mathbb R}^*=A_{n-1}(X)\otimes\mathbb R$。

令 $r=|\Sigma(1)|$，则包含映射 $G\subset(\mathbb C^*)^r$ 诱导 $G_{\mathbb R}\subset U(1)^r$。由此得到 $G_{\mathbb R}$ 在 $\mathbb C^r$ 上的 Hamilton 作用，其矩映射显然是复合映射
\[
\mu_\Sigma:\mathbb C^r\xrightarrow{\mu}(\mathbb R^r)^*
\xrightarrow{p}\mathfrak g_{\mathbb R}^*
=A_{n-1}(X)\otimes\mathbb R\simeq\mathbb R^{r-n}.
\]
这里 $\mu$ 是 $U(1)^r$ 的矩映射，$p$ 来自正合列
\begin{equation}\label{eq:moment-exact-sequence}
0\longrightarrow M_{\mathbb R}\longrightarrow(\mathbb R^r)^*
\xrightarrow{p}A_{n-1}(X)\otimes\mathbb R\longrightarrow0,
\end{equation}
该正合列由式~\eqref{eq:toric-chow-sequence} 与 $\mathbb R$ 作张量积得到。映射 $\mu_\Sigma$ 很容易显式描述。首先注意，式~\eqref{eq:moment-exact-sequence} 中的 $(b_1,\ldots,b_r)\in\mathbb R^r$ 在 $M_{\mathbb R}$ 的像上为零，当且仅当在 $N_{\mathbb R}$ 中有 $\sum_{i=1}^r b_iv_i=0$。所有这样的 $(b_1,\ldots,b_r)$ 构成维数为 $r-n$ 的子空间 $\Lambda_{\mathbb R}\subset\mathbb R^r$；注意 $\Lambda_{\mathbb R}=\Lambda_{\mathbb Q}\otimes\mathbb R$，其中 $\Lambda_{\mathbb Q}$ 见式~\eqref{eq:toric-relation-space}。选取 $\Lambda_{\mathbb R}$ 的一组基 $(b_{1,j},\ldots,b_{r,j})$，$1\leq j\leq r-n$，就得到同构 $A_{n-1}(X)\otimes\mathbb R\simeq\mathbb R^{r-n}$，在此同构下 $p$ 变为
\[
p(x_1,\ldots,x_r)=
\left(\sum_{i=1}^r b_{i,1}x_i,\ldots,\sum_{i=1}^r b_{i,r-n}x_i\right).
\]
因此 $\mu_\Sigma$ 的公式为
\begin{equation}\label{eq:toric-moment-map}
\mu_\Sigma(z_1,\ldots,z_r)=\frac12
\left(\sum_{i=1}^r b_{i,1}|z_i|^2,\ldots,
\sum_{i=1}^r b_{i,r-n}|z_i|^2\right).
\end{equation}
我们也看到，$\mu_\Sigma$ 在 $G_{\mathbb R}$ 的每个轨道上取常值。

当 $X$ 射影且单纯时，前面已看到 $A_{n-1}(X)\otimes\mathbb R\simeq H^{1,1}(X,\mathbb R)$。由命题~\ref{prop:toric-kahler-cone}，$X$ 的 Kähler 锥可与 $\operatorname{cpl}(\Sigma)\subset A_{n-1}(X)\otimes\mathbb R$ 的内部等同。于是可如下构造 $X$。

\begin{theorem}\label{thm:toric-symplectic-reduction}
若环面簇 $X=X_\Sigma$ 射影且单纯，而 $a\in A_{n-1}(X)\otimes\mathbb R$ 是 Kähler 类，则 $\mu_\Sigma^{-1}(a)\subset\mathbb C^r-Z(\Sigma)$，且自然映射
\[
\mu_\Sigma^{-1}(a)/G_{\mathbb R}
\longrightarrow(\mathbb C^r-Z(\Sigma))/G=X
\]
\SourcePage{65}
是轨形微分同胚。此外，将 $\mathbb C^r$ 上的辛形式 $\omega$ 限制到 $\mu_\Sigma^{-1}(a)$ 后，它下降为 $\mu_\Sigma^{-1}(a)/G_{\mathbb R}$ 上的辛形式；其上同调类通过上述微分同胚对应于 $a\in H^2(X,\mathbb R)$。
\end{theorem}
\begin{proof}
在定理的陈述中，要注意 $\omega$ 限制到 $\mu_\Sigma^{-1}(a)$ 后并不是辛形式。取商恰好使它变为非退化的。这个过程称为辛约化。由于 $G_{\mathbb R}$ 可以有有限稳定子群，商 $\mu_\Sigma^{-1}(a)/G_{\mathbb R}$ 自然带有辛轨形结构，见\hyperref[app:singular-varieties]{附录 A}。

采用式~\eqref{eq:cpl-cone} 的记号，将 $a\in A_{n-1}(X)\otimes\mathbb R$ 写为 $a=\sum_{i=1}^r a_i[D_i]$。因为 $a$ 是 Kähler 类，命题~\ref{prop:toric-kahler-cone} 表明 $a$ 位于 $\operatorname{cpl}(\Sigma)$ 的内部，即 $a$ 严格凸。考虑由不等式 $\langle m,v_i\rangle\geq-a_i$ 定义的凸集 $\Delta\subset M_{\mathbb R}$。虽然 $\Delta$ 未必是整的，甚至未必是有理的，关键在于 $a$ 的严格凸性意味着多胞形 $\Delta$ 与 $\Sigma$ 在\emph{组合意义下对偶}。也就是说，当 $\sigma$ 遍历 $\Sigma$ 中的锥时，$\Delta$ 的各真面由
\[
\{m\in\Delta:\langle m,v_i\rangle=-a_i
\text{ 对所有 }v_i\in\sigma\}
\]
给出。当 $a_i$ 为整数时，这是\hyperref[subsec:toric-polytopes]{第 3.2.1 小节}讨论的丰沛性的一项性质。$a_i$ 为有理数的情形容易推出，再用连续性论证即可证明一般情形。

知道 $X$ 的扇在组合意义下与 $\Delta$ 对偶之后，就可用 \SourceCite[附录 1，定理 1.4]{Guillemin} 得出
\[
\mu_\Sigma^{-1}(a)/G_{\mathbb R}\simeq X.
\]
此外，定理的最后一个断言来自 \SourceCite[附录 2，式 (1.6)]{Guillemin}；Guillemin 使用的 $\lambda_i$ 在我们的记号下为 $-a_i$。Guillemin 的论证是在光滑情形下给出的，不过稍作修改也适用于 $X$ 单纯的情形。
\end{proof}

利用辛约化构造环面簇，在物理中有一些有趣的推论。特别地，某些称为\emph{规范线性 $\sigma$ 模型}的物理理论，以通过辛约化描述的环面簇为出发点。详情参见\hyperref[sec:gauged-linear-sigma-models]{附录 B.5}。\index{辛约化|hyperpage}\index{规范线性sigma模型@规范线性 $\sigma$ 模型|hyperpage}

\section{GKZ 分解}\label{sec:gkz-decomposition}

\hyperref[chap:moduli]{第~\ref*{chap:moduli} 章}研究 Kähler 模时，我们会看到需要扩大\hyperref[chap:introduction]{第 1 章}定义的 Kähler 模空间。大致的想法是：由一次翻越联系起来的两个 Calabi--Yau 流形，其 Kähler 锥是“相邻”的；这使我们能够将相应的 Kähler 模空间粘合起来。

在环面几何中，“相邻”的 Kähler 锥这一想法自然出现在 \emph{GKZ 分解}中。为了解它的由来，注意由式~\eqref{eq:toric-chow-sequence}，环面簇 $X=X_\Sigma$ 的 Chow 群 $A_{n-1}(X)$ 只依赖于 $\Sigma$ 的一维锥集合 $\Sigma(1)$。因此，具有相同一维锥的两个环面簇，其除子的 Chow 群典范同构。于是可以问：它们各自的 Kähler 锥在与 $\mathbb R$ 作张量积后的 Chow 群中处于什么位置？\index{GKZ分解@GKZ 分解|hyperpage}

\subsection{分解}\label{subsec:gkz-cones}

为精确表述这一想法，在 $N_{\mathbb R}$ 中固定由强凸有理一维锥构成的有限集合 $\Xi$，令 $\Sigma$ 是满足 $\Sigma(1)=\Xi$ 的扇。因为 $A_{n-1}(X)\otimes\mathbb R$ 只依赖于 $\Xi$，将此向量空间记为 $A(\Xi)$。类似地，有效除子类构成的锥 $A^+_{n-1}(X)\otimes\mathbb R$ 也只依赖于 $\Xi$，将它记为 $A^+(\Xi)$。因此 $A^+(\Xi)$ 是 $A(\Xi)\simeq\mathbb R^{r-n}$ 中的锥，其中 $r=|\Xi|$。

\SourcePage{66}
若 $\Sigma$ 是满足 $\Sigma(1)=\Xi$ 的射影单纯扇，则由命题~\ref{prop:toric-kahler-cone}，其 Kähler 锥的闭包是 $\operatorname{cpl}(\Sigma)$。用上述记号，就得到 $(r-n)$ 维锥 $\operatorname{cpl}(\Sigma)\subset A^+(\Xi)$。

现在让 $\Sigma$ 遍历所有这样的扇，便得到位于 $A^+(\Xi)$ 中的一族锥。直观上或许希望这些锥填满整个 $A^+(\Xi)$，但实际并非如此。为得到余下的部分，必须考虑满足 $\Sigma(1)\subset\Xi$ 的射影单纯扇。此时 $A_{n-1}(X_\Sigma)\otimes\mathbb R$ 可能不同于 $A(\Xi)$，但仍可如下定义锥 $\operatorname{cpl}(\Sigma)\subset A^+(\Xi)$。给定 $a=\sum_{i=1}^r a_i[D_i]\in A^+(\Xi)$，对 $\sigma\in\Sigma$，可找到 $m_\sigma\in M_{\mathbb R}$，使得只要 $v_i\in\sigma$ 且 $v_i$ 生成 $\Sigma(1)$ 的一个元素，就有 $\langle m_\sigma,v_i\rangle=-a_i$。若对所有 $\sigma\in\Sigma$ 及 $1\leq i\leq r$，都有 $\langle m_\sigma,v_i\rangle\geq-a_i$，则称 $a$ 为\emph{$\Sigma$ 凸的}。仿照式~\eqref{eq:cpl-cone}，定义
\[
\operatorname{cpl}(\Sigma)=
\{a\in A^+(\Xi):a\text{ 是 }\Sigma\text{ 凸的}\}.
\]
由 \SourceCite{OdP}，此锥仍有维数 $r-n$。值得注意的是，这些锥以非常规整的方式填满 $A^+(\Xi)$，形成 \SourceCite{OdP} 所称的 \emph{Gelfand--Kapranov--Zelevinsky 分解}，简称 \emph{GKZ 分解}。精确描述如下。

\begin{theorem}\label{thm:gkz-fan}
设 $\Xi$ 是 $N_{\mathbb R}$ 中由强凸有理一维锥构成的有限集合。当 $\Sigma$ 遍历所有满足 $\Sigma(1)\subset\Xi$ 的射影单纯扇时，各锥 $\operatorname{cpl}(\Sigma)$ 及其各面在 $A(\Xi)$ 中组成一个扇，其支集为 $A^+(\Xi)$。
\end{theorem}

\hyperref[sec:toric-examples]{第 3.7 节}将给出一个 GKZ 分解的例子，其中有两个锥 $\operatorname{cpl}(\Sigma)$ 和 $\operatorname{cpl}(\Sigma')$，满足 $\Sigma(1)=\Xi$ 而 $\Sigma'(1)\subsetneq\Xi$。

一般而言，可以问 GKZ 分解中的两个锥何时有一个共同的余维数为一的面。答案很漂亮：当 $N_{\mathbb R}$ 中对应的两个扇通过细分，即增加或减去一条一维锥，或通过一次“翻越”联系起来时，就会出现这种情形；详情见 \SourceCite{OdP}。\hyperref[subsec:enlarged-kahler-moduli]{第~\ref*{subsec:enlarged-kahler-moduli} 小节}将讨论翻越。还应提到，非单纯扇对应于 GKZ 分解中某些维数小于 $r-n$ 的锥。

\subsection{次级扇与 Gale 变换}\label{subsec:secondary-fan}

GKZ 分解可以扩充为一个完备扇，其支集是整个 $A(\Xi)$。具体做法如下。首先，把 $\Xi$ 中每条一维锥与它的本原整生成元相对应，将 $\Xi$ 视为 $N_{\mathbb R}$ 的一个子集。然后定义子集 $\Xi^+\subset N_{\mathbb R}\times\mathbb R$：
\begin{equation}\label{eq:extended-rays}
\Xi^+=(\Xi\cup\{0\})\times\{1\}.
\end{equation}
我们关心的是 $\Xi^+$ 的\emph{次级扇}。其定义类似于 GKZ 分解的定义：将 GKZ 分解中出现的射影单纯扇替换为凸包 $\operatorname{Conv}(\Xi^+)$ 的\emph{\IndexLocation{indexloc:1226:66}正则三角剖分} $\mathcal T$。$\mathcal T$ 中各单形的顶点必须属于 $\Xi^+$；所谓正则，是指 $\mathcal T$ 可由一个\IndexLocation{indexloc:0182:66}严格凸函数确定，正如射影单纯扇 $\Sigma$ 可由严格凸支撑函数 $a\in\operatorname{cpl}(\Sigma)$ 确定一样。与 $\mathcal T$ 对应的是锥 $C(\mathcal T)\subset A(\Xi)$，它由确定同一三角剖分 $\mathcal T$ 的严格凸支撑函数及其极限构成。各 $C(\mathcal T)$ 是次级扇的极大锥；次级扇是 $A(\Xi)$ 中的完备扇。\index{次级扇|hyperpage}\index{Gale变换@Gale 变换|hyperpage}

给定三角剖分 $\mathcal T$，锥 $C(\mathcal T)$ 还有如下描述。由式~\eqref{eq:toric-chow-sequence} 可见，$\Xi^+\subset N_{\mathbb R}\times\mathbb R$ 给出
\begin{equation}\label{eq:gale-exact-sequence}
0\longrightarrow M_{\mathbb R}\oplus\mathbb R
\longrightarrow\mathbb R^{\Xi^+}\longrightarrow A(\Xi)
\longrightarrow0.
\end{equation}

\SourcePage{67}
对每个 $v\in\Xi^+$，令 $e_v^*\in A(\Xi)$ 为 $\mathbb R^{\Xi^+}$ 中对应于 $v$ 的标准基向量在式~\eqref{eq:gale-exact-sequence} 的满射下的像。元素集合 $\{e_v^*:v\in\Xi^+\}$ 称为 $\Xi^+$ 的 \emph{Gale 变换}。读者可以验证，若 $\Sigma(1)=\Xi$ 且 $v=(\rho,1)\in\Xi^+-\{(0,1)\}$，则 $e_v^*=[D_\rho]$，即 $\rho\in\Sigma(1)$ 对应的环面除子在 $A(\Xi)=A_{n-1}(X_\Sigma)\otimes\mathbb R$ 中的类。

设 $\mathcal T$ 是上述 $\operatorname{Conv}(\Xi^+)$ 的正则三角剖分。每个单形 $\sigma\in\mathcal T$ 给出锥 $C(\sigma)\subset A(\Xi)$，其生成元是那些 $v$ 不是 $\sigma$ 顶点的 $e_v^*$。可以证明
\[
C(\mathcal T)=\bigcap_{\sigma\in\mathcal T}C(\sigma).
\]
若三角剖分 $\mathcal T$ 的每个极大单形都以 $(0,1)$ 为顶点，则因为 $(0,1)\in\operatorname{Conv}(\Xi^+)\subset N_{\mathbb R}\times\{1\}\simeq N_{\mathbb R}$ 是内点，从顶点 $(0,1)$ 向 $\mathcal T$ 中各单形的面作锥，便得到一个完备扇 $\Sigma$。所得扇满足 $\Sigma(1)\subset\Xi$，且 $C(\mathcal T)=\operatorname{cpl}(\Sigma)$。这个过程可以反过来进行，因此次级扇含有 GKZ 扇作为子扇。次级扇的更完整描述见 \SourceCite{GKZ2}，它与 GKZ 分解的关系见 \SourceCite{OdP}。

由式~\eqref{eq:gale-exact-sequence} 可见，对偶空间 $A(\Xi)^*$ 为
\[
A(\Xi)^*=
\left\{\lambda=(\lambda_v)\in\mathbb R^{\Xi^+}:
\sum_{v\in\Xi^+}\lambda_v\,v=0\right\}.
\]
这是 $\Xi^+$ 中各元素的线性关系构成的向量空间，类似于式~\eqref{eq:toric-relation-space} 定义的 $\Lambda_{\mathbb Q}$。注意，$A(\Xi)^*$ 中有一个显然的整数格，即系数全为整数的关系构成的格。此外，Gale 变换中的每个 $e_v^*\in A(\Xi)$ 都是 $A(\Xi)^*$ 上的线性泛函，容易算得 $e_v^*(\lambda)=\lambda_v$。\hyperref[chap:moduli]{第~\ref*{chap:moduli} 章}将按以下方式使用这一事实：固定 $A(\Xi)^*$ 的一组基，将每个基向量视为一行 $\lambda=(\lambda_v)$，这样便得到一个矩阵；若在 $A(\Xi)$ 中使用对偶基，则这个矩阵的各列就是各 $e_v^*$。这种描述使 Gale 变换易于处理。

一个特别好的情形是 $|\Xi^+|=n+3$，其中 $n$ 是格 $N$ 和 $M$ 的秩。此时 $A(\Xi)$ 的维数为 $2$，Gale 变换 $\{e_v^*\}$ 是平面上的一组向量。这时容易看到，次级扇就是以 Gale 变换中的各向量所张成的一维锥为一维锥的完备二维扇；\SourceCite{DHSS} 明确指出了这一点。

\subsection{与辛约化的关系}\label{subsec:gkz-reduction}

本节最后考察矩映射与 GKZ 分解的关系。首先注意，矩映射 $\mu_\Sigma$ 只依赖于 $\Sigma(1)$，因为正合列~\eqref{eq:moment-exact-sequence} 由 $\Sigma(1)$ 确定。因此可将矩映射写为 $\mu_\Xi:\mathbb C^\Xi\to A(\Xi)$。它的像是 $A^+(\Xi)$，所以有
\[
\mu_\Xi:\mathbb C^\Xi\longrightarrow A^+(\Xi).
\]
由于作用在 $\mathbb C^\Xi$ 上的群 $G$ 也只依赖于 $\Xi$，将它记为 $G(\Xi)$。

设 GKZ 分解中有一个 $(r-n)$ 维锥 $\operatorname{cpl}(\Sigma)\subset A^+(\Xi)$。取其内部的 $a$，形成商
\[
\mu_\Xi^{-1}(a)/G(\Xi)_{\mathbb R}.
\]
若 $\Sigma(1)=\Xi$ 且 $\Sigma$ 射影、单纯，由定理~\ref{thm:toric-symplectic-reduction}，这个商就是 $X_\Sigma$。当 $\Sigma(1)$ 是 $\Xi$ 的真子集时，上述定理不再适用，但以下推广告诉我们，得到的仍是 $X_\Sigma$。

\SourcePage{68}
\begin{theorem}\label{thm:gkz-symplectic-reduction}
若 $\Sigma$ 是满足 $\Sigma(1)\subset\Xi$ 的射影单纯扇，且 $a\in A^+(\Xi)$ 位于 $\operatorname{cpl}(\Sigma)$ 的内部，则有自然的轨形微分同胚
\[
\mu_\Xi^{-1}(a)/G(\Xi)_{\mathbb R}\simeq X_\Sigma.
\]
\end{theorem}
\begin{proof}
除了群 $G(\Xi)$ 与矩映射 $\mu_\Xi$，还有由 $\Sigma$ 得到的 $G(\Sigma)$ 和 $\mu_\Sigma$。为联系二者，令 $C=\Xi-\Sigma(1)$，于是 $\mathbb Z^\Xi=\mathbb Z^{\Sigma(1)}\oplus\mathbb Z^C$。投影 $\mathbb Z^\Xi\to\mathbb Z^{\Sigma(1)}$ 诱导正合列
\begin{equation}\label{eq:gkz-group-sequence}
1\longrightarrow G(\Sigma)\longrightarrow G(\Xi)
\longrightarrow(\mathbb C^*)^C\longrightarrow1,
\end{equation}
使得相对于分解 $\mathbb C^\Xi=\mathbb C^{\Sigma(1)}\oplus\mathbb C^C$，$G(\Sigma)$ 在第二个因子上的作用是平凡的。此外，可通过以下交换图联系矩映射：
\[
\begin{array}{rccccc}
\mu_\Xi:&\mathbb C^\Xi&\longrightarrow&\mathbb R^\Xi&\longrightarrow&A(\Xi)\\
&\big\downarrow&&\big\downarrow&&\big\downarrow\\
\mu_\Sigma:&\mathbb C^{\Sigma(1)}&\longrightarrow&\mathbb R^{\Sigma(1)}
&\longrightarrow&A_{n-1}(X_\Sigma)\otimes\mathbb R,
\end{array}
\]
其中竖直箭头是自然投影。

若 $a=\sum_{i=1}^r a_i[D_i]$ 位于 $\operatorname{cpl}(\Sigma)$ 的内部，则 \SourceCite[命题 3.1]{OdP} 的证明表明 $a$ 严格凸，即当 $\sigma\in\Sigma$ 为 $n$ 维、$v_i$ 不是 $\sigma$ 的生成元时，有 $\langle m_\sigma,v_i\rangle>-a_i$。令 $a_0$ 为 $a$ 在 $A_{n-1}(X_\Sigma)\otimes\mathbb R$ 中的像，则 $a_0$ 也严格凸。由定理~\ref{thm:toric-symplectic-reduction}，只须证明
\[
\mu_\Sigma^{-1}(a_0)/G(\Sigma)_{\mathbb R}
\simeq\mu_\Xi^{-1}(a)/G(\Xi)_{\mathbb R}.
\]

由于 $\Sigma$ 完备，每个 $v_i\in C$ 都位于某个 $n$ 维锥 $\sigma\in\Sigma$ 中；这里仍按惯例将一维锥与其生成元等同。于是 $v_i=\sum_j c_{ij}v_j$，其中 $c_{ij}\geq0$，求和遍历那些 $v_j$ 为 $\sigma$ 生成元的 $j$；本证明余下部分沿用这一约定。严格凸条件 $\langle m_\sigma,v_i\rangle>-a_i$ 给出不等式
\[
a_i>\sum_j c_{ij}a_j,\qquad v_i\in C.
\]
设 $\mathbb C^\Xi$ 的坐标为 $z_1,\ldots,z_r$，选择编号使 $v_1,\ldots,v_k\in\Sigma(1)$，$v_{k+1},\ldots,v_r\in C$。不难证明，$(z_1,\ldots,z_r)\in\mu_\Xi^{-1}(a)$，当且仅当 $(z_1,\ldots,z_k)\in\mu_\Sigma^{-1}(a_0)$，并且
\begin{equation}\label{eq:gkz-extra-coordinates}
\frac12|z_i|^2=\frac12\sum_j c_{ij}|z_j|^2
+a_i-\sum_j c_{ij}a_j,
\qquad k+1\leq i\leq r.
\end{equation}

定义映射 $\phi:\mu_\Sigma^{-1}(a_0)\to\mu_\Xi^{-1}(a)$：将 $(z_1,\ldots,z_k)\in\mu_\Sigma^{-1}(a_0)$ 送到 $(z_1,\ldots,z_r)\in\mu_\Xi^{-1}(a)$，其中 $z_{k+1},\ldots,z_r$ 取为满足式~\eqref{eq:gkz-extra-coordinates} 的正实数；根据上述不等式，该式右边为正。于是 $\phi$ 对 $G(\Sigma)_{\mathbb R}$ 与 $G(\Xi)_{\mathbb R}$ 的作用是等变的，而且由式~\eqref{eq:gkz-group-sequence}，$G(\Xi)_{\mathbb R}$ 的每个轨道都与 $\phi$ 的像相交。最后，设 $z,w\in\mu_\Sigma^{-1}(a_0)$、$g\in G(\Xi)_{\mathbb R}$ 满足 $g\phi(z)=\phi(w)$。由于 $g$ 不改变绝对值，式~\eqref{eq:gkz-extra-coordinates} 表明 $g$ 在 $z_{k+1},\ldots,z_r$ 上的作用平凡。这些数又都非零，因此由式~\eqref{eq:gkz-group-sequence} 得 $g\in G(\Sigma)_{\mathbb R}$。所以 $\phi$ 在轨道之间诱导双射，容易看出它是轨形微分同胚。
\end{proof}

定理~\ref{thm:gkz-symplectic-reduction} 表明，矩映射 $\mu_\Xi$ 不仅能构造一个环面簇，而且能构造所有满足 $\Sigma(1)\subset\Xi$ 的射影单纯环面簇 $X_\Sigma$。这与\emph{\IndexLocation{indexloc:0924:68}多重镜像}现象有关，第\hyperref[chap:mirror-constructions]{4 章}将对此展开讨论。\SourceCite{Witten5} 的\IndexLocation{indexloc:0354:68}规范线性 $\sigma$ 模型也与此相关。
\SourcePage{69}
在这些模型中，元素 $a\in A(\Xi)$ 是一个参数；$a$ 所处的位置不同，就得到不同的物理理论。若 $a$ 位于某个锥 $\operatorname{cpl}(\Sigma)$ 的内部，则理论涉及相应的环面簇 $X_\Sigma$；若它位于 $A^+(\Xi)$ 之外，则得到很不一样的理论，例如 Landau--Ginzburg 理论。物理文献将这种现象称为\emph{相}，例如参见 \SourceCite{MP1}。

\section{Fano 簇与自反多胞形}\label{sec:fano-reflexive}

射影空间的反典范类是丰沛的。更一般地，任何具有此性质的光滑完备簇都称为 \emph{Fano 簇}。为了研究镜像对称，还需考虑有奇点的 Fano 簇。为给出定义，回顾\hyperref[app:singular-varieties]{附录 A}：Cohen--Macaulay 簇 $V$ 有对偶化层 $\widehat\Omega_V^n$，其中 $n=\dim V$；而 $\widehat\Omega_V^n$ 是线丛，当且仅当 $V$ 是 Gorenstein 的。

Fano 性意味着对偶化层的对偶是丰沛的。特别地，这说明我们可能希望 $\widehat\Omega_V^n$ 是线丛，也就是希望这样的簇为 Gorenstein 簇。因此有如下定义。
\begin{definition}\label{def:fano}
若完备的 $n$ 维 Gorenstein 簇 $V$ 的线丛 $\widehat\Omega_V^n$ 的对偶是丰沛的，则称 $V$ 为 \emph{Fano 簇}。
\end{definition}

对环面簇，可以如下刻画这些概念。首先回顾，任意环面簇 $X$ 上的对偶化层都有简单描述
\[
\widehat\Omega_X^n=\mathcal O_X\left(-\sum_\rho D_\rho\right),
\]
见 \SourceCite{Danilov,Fulton3,Oda}。用典范除子 $K_X$ 表示，就是 $K_X=-\sum_\rho D_\rho$。因此，$X$ 是 Gorenstein 的，当且仅当 $\sum_\rho D_\rho$ 是 Cartier 除子。于是可如下刻画 Fano 环面簇。
\begin{lemma}\label{lem:toric-fano}
完备的 $n$ 维环面簇 $X$ 是 Fano 簇，当且仅当 $\sum_\rho D_\rho$ 是丰沛 Cartier 除子。
\end{lemma}

由于每个丰沛除子都来自一个多胞形，Fano 环面簇上的丰沛除子 $-K_X$ 确定了一个多胞形 $\Delta$，它具有一些很特殊的性质。这就引出了 Batyrev 的\emph{自反多胞形}概念 \SourceCite{Batyrev4}。精确定义如下。
\begin{definition}\label{def:reflexive-polytope}
若 $n$ 维整多胞形 $\Delta\subset M_{\mathbb R}\simeq\mathbb R^n$ 满足以下两个条件，则称 $\Delta$ 为\emph{自反}的：
\begin{enumerate}[label=(\roman*)]
\item $\Delta$ 的每个刻面 $\Gamma$ 都位于形如 $\{m\in M_{\mathbb R}:\langle m,v_\Gamma\rangle=-1\}$ 的支撑仿射超平面上，其中 $v_\Gamma\in N$；
\item $\operatorname{Int}(\Delta)\cap M=\{0\}$。
\end{enumerate}
\end{definition}

自反多胞形具有非常漂亮的组合对偶性。设 $\Delta$ 为整多胞形，$\Delta^\circ$ 是\hyperref[subsec:toric-from-polytopes]{第 3.2.2 小节}定义的极多胞形。除了 $(\Delta^\circ)^\circ=\Delta$ 之外，\SourceCite{Batyrev4} 还证明，$\Delta$ 与 $\Delta^\circ$ 的基本对偶性如下。
\begin{lemma}\label{lem:reflexive-duality}
$\Delta$ 自反，当且仅当 $\Delta^\circ$ 自反。
\end{lemma}

自反多胞形在这里之所以有趣，是因为以下结果刻画了何时 $\mathbb P_\Delta$ 是 Fano 簇。
\begin{proposition}\label{prop:reflexive-fano}
$\Delta$ 自反，当且仅当 $\mathbb P_\Delta$ 是 Fano 簇，且 $\Delta$ 是其反典范除子对应的多胞形。
\end{proposition}

\SourcePage{70}
\begin{proof}
设 $\Delta$ 自反，则定义~\ref{def:reflexive-polytope} 中的 $v_\Gamma$ 恰好是 $\mathbb P_\Delta$ 的扇的一维锥生成元。容易看出，支撑函数 $\phi_\Delta$ 在每个 $v_\Gamma$ 上都取值 $-1$。因此，$\mathbb P_\Delta$ 上相应的丰沛除子是 $\sum_\rho D_\rho=-K_{\mathbb P_\Delta}$，所以 $\mathbb P_\Delta$ 是 Fano 簇。逆命题也同样容易证明。
\end{proof}

Fano 环面簇最简单的例子是 $\mathbb P^n$。接着考虑加权射影空间，此时答案稍微有趣一些。
\begin{lemma}\label{lem:weighted-projective-fano}
设 $X=\mathbb P(q_0,\ldots,q_n)$ 是使用归一化权重的加权射影空间，令 $q=\sum_{i=0}^n q_i$。则 $X$ 是 Fano 簇，当且仅当对每个 $i$ 都有 $q_i\mid q$。
\end{lemma}
\begin{proof}
$X=\mathbb P(q_0,\ldots,q_n)$ 的扇的一维锥生成元为 $v_0,\ldots,v_n$，在 $n$ 维格 $N$ 中满足 $\sum_{i=0}^n q_iv_i=0$。扇的极大锥由 $\{v_0,\ldots,v_n\}$ 的各 $n$ 元子集生成。首先判定除子 $D=\sum_{i=0}^n D_i$ 何时是 Cartier 除子。

对每个 $i$，容易看出存在唯一的 $m_i\in M\otimes\mathbb Q$，使对所有 $j\ne i$ 都有 $\langle m_i,v_j\rangle=-1$。于是 $D$ 是 Cartier 除子，当且仅当每个 $m_i$ 都属于 $M$。不过，关系 $\sum_{i=0}^n q_iv_i=0$ 意味着
\[
\langle m_i,v_i\rangle=\frac{\sum_{j\ne i}q_j}{q_i}
=\frac q{q_i}-1.
\]
由于 $v_0,\ldots,v_n$ 生成 $N$，可见 $m_i\in M$ 当且仅当 $q_i\mid q$；对所有 $i$ 同时要求这一条件，就得到上述判别。一旦知道 $D$ 是 Cartier 除子，\hyperref[sec:cones-fans]{第 3.1 节}的丰沛性准则便化为不等式 $\langle m_i,v_i\rangle>-1$，这在此处显然成立。
\end{proof}

\hyperref[sec:toric-examples]{第 3.7 节}将给出更多 Fano 环面簇的例子。在任意给定维数中，模去幺模变换后，自反多胞形只有有限多个；这意味着 $n$ 维 Fano 环面簇的同构类也只有有限多个，见 \SourceCite{Batyrev4}。低维光滑 Fano 环面簇已得到分类，见 \SourceCite{Batyrev6,Oda}；同时，人们正在尝试分类所有四维自反多胞形，见 \SourceCite{KS3,Skarke}。

由于自反多胞形成对出现为 $\Delta$、$\Delta^\circ$，我们得到在某种意义下对偶的 Fano 环面簇 $\mathbb P_\Delta$、$\mathbb P_{\Delta^\circ}$。\hyperref[chap:mirror-constructions]{第 4 章}将利用这些环面簇构造 Calabi--Yau 超曲面的“对偶”族，它们在镜像对称中十分重要。\index{Fano簇@Fano 簇|hyperpage}

\section{环面簇的自同构}\label{sec:toric-automorphisms}

本节描述完备单纯环面簇 $X$ 的自同构群。这些结果源于 Demazure \SourceCite{Demazure}，他在 $X$ 光滑时计算了 $\operatorname{Aut}(X)$。\SourceCite{Cox} 证明，Demazure 对 $\operatorname{Aut}(X)$ 的描述在单纯情形同样成立。

若 $X$ 由扇 $\Sigma$ 给出，则定理~\ref{thm:toric-quotient} 给出商表示 $X=(\mathbb C^{\Sigma(1)}-Z(\Sigma))/G$，其中 $G=\operatorname{Hom}_{\mathbb Z}(A_{n-1}(X),\mathbb C^*)$。还需回顾齐次坐标环 $S=\mathbb C[x_\rho:\rho\in\Sigma(1)]$，其分次由 $A_{n-1}(X)$ 给出。

我们将通过寻找与 $G$ 的作用交换的 $\mathbb C^{\Sigma(1)}-Z(\Sigma)$ 的自同构，来构造 $X=(\mathbb C^{\Sigma(1)}-Z(\Sigma))/G$ 的自同构。一类显然的自同构由“大环面” $(\mathbb C^*)^{\Sigma(1)}$ 中的元素给出。正合列~\eqref{eq:toric-g-sequence}
\[
1\longrightarrow G\longrightarrow(\mathbb C^*)^{\Sigma(1)}
\longrightarrow T_N\longrightarrow1
\]
\SourcePage{71}
表明，这一类给出来自环面的 $X$ 的自同构，即 $T_N\subset\operatorname{Aut}(X)$。

下一类自同构来自\emph{根}。$X$ 的一个根是一个有序对 $(x_\rho,x^D)$，其中 $x_\rho$ 是我们的一个变量，$x^D\in S$ 是不同于 $x_\rho$、但与它同次的单项式，即在 $A_{n-1}(X)$ 中有 $\deg(x^D)=\deg(x_\rho)$。注意，$x_\rho$ 不可能整除 $x^D$，否则 $x^D/x_\rho$ 就是一个次数为 $0$ 的非平凡单项式；由于 $X$ 完备，这样的单项式不存在。

根 $(x_\rho,x^D)$ 如下确定 $\mathbb C^{\Sigma(1)}$ 的自同构。将 $\mathbb C^{\Sigma(1)}$ 的一点写为 $(x_\rho,\mathbf x)$，其中 $\mathbf x$ 是以 $\Sigma(1)-\{\rho\}$ 为指标的向量。于是得到单参数自同构族
\begin{equation}\label{eq:toric-root-automorphism}
y_\lambda(x_\rho,\mathbf x)=
(x_\rho+\lambda\mathbf x^D,\mathbf x),
\end{equation}
其中 $\lambda\in\mathbb C$，$\mathbf x^D$ 是单项式 $x^D$ 在 $\mathbf x$ 处的值；由于 $x_\rho$ 不整除 $x^D$，这个取值有意义。可以证明，$y_\lambda$ 是 $\mathbb C^{\Sigma(1)}-Z(\Sigma)$ 的自同构，并与 $G$ 的作用交换，因此下降为 $\operatorname{Aut}(X)$ 中的元素。

最后一类自同构来自扇 $\Sigma$ 的对称，即置换 $\Sigma$ 中各锥的 $N$ 的自同构。特别地，它们置换一维锥，因此给出 $\mathbb C^{\Sigma(1)}-Z(\Sigma)$ 的自同构；这些自同构与 $G$ 的作用交换，因而也下降到 $X$。

\begin{theorem}\label{thm:toric-automorphism-group}
若 $X$ 是单纯完备环面簇，则来自环面、根以及扇对称的上述三类自同构生成 $\operatorname{Aut}(X)$。此外，来自环面与根的前两类生成 $\operatorname{Aut}(X)$ 的单位元连通分支，而且
\[
\begin{aligned}
\dim\operatorname{Aut}(X)
&=n+\text{根的个数}\\
&=n+\sum_\rho\bigl(\dim S_{\deg(x_\rho)}-1\bigr).
\end{aligned}
\]
\end{theorem}

例如，$\mathbb P^n$ 的\IndexLocation{indexloc:1369:71}齐次坐标环是通常的分次环 $\mathbb C[x_0,\ldots,x_n]$。这意味着，根恰好是所有 $i\ne j$ 的有序对 $(x_i,x_j)$。这样的有序对有 $(n+1)^2-(n+1)$ 个，因此由定理~\ref{thm:toric-automorphism-group}，$\dim\operatorname{Aut}(\mathbb P^n)=(n+1)^2-1$。这当然就是 $\operatorname{PGL}(n+1,\mathbb C)$ 的维数。文献 \SourceCite{Cox,Demazure,Oda} 给出了 $\operatorname{Aut}(X)$ 作为代数群的结构的更多细节。

下面介绍 \SourceCite{AGM1} 给出的一种计算 $\operatorname{Aut}(X)$ 维数的几何方法。假设\IndexLocation{indexloc:1375:71}单纯环面簇 $X$ 是 Gorenstein 的，则反典范类 $-K_X=\sum_\rho D_\rho$ 是 Cartier 除子，给出多胞形
\[
\Delta=\Delta_{-K_X}=
\{m\in M_{\mathbb R}:\langle m,v_\rho\rangle\geq-1\}.
\]
注意，$0$ 在 $\Delta$ 的内部，因此 $\dim\Delta=n$，不过 $\Delta$ 未必是整的。

\begin{proposition}\label{prop:automorphism-facets}
若 $X$ 单纯、完备且 Gorenstein，而 $\Delta$ 如上定义，则
\[
\dim\operatorname{Aut}(X)=n+\sum_\Gamma l^*(\Gamma),
\]
其中求和遍历 $\Delta$ 的所有刻面 $\Gamma$。
\end{proposition}
\begin{proof}
回顾式~\eqref{eq:lattice-points}，$l^*(\Gamma)$ 是刻面 $\Gamma$ 的相对内部中的格点数。为证明命题，设 $(x_\rho,x^D)$ 是一个根，写 $x^D=\prod_{\rho'\ne\rho}x_{\rho'}^{a_{\rho'}}$。由正合列~\eqref{eq:toric-chow-sequence}，$\deg(x_\rho)=\deg(x^D)$ 当且仅当存在 $\widetilde m\in M$，使
\begin{equation}\label{eq:root-lattice-pairings}
\begin{aligned}
\langle\widetilde m,v_\rho\rangle&=-1,\\
\langle\widetilde m,v_{\rho'}\rangle&=a_{\rho'},\qquad\rho'\ne\rho.
\end{aligned}
\end{equation}

\SourcePage{72}
由于 $\Delta$ 由不等式 $\langle m,v_\rho\rangle\geq-1$ 及 $\rho'\ne\rho$ 时的 $\langle m,v_{\rho'}\rangle\geq-1$ 定义，式~\eqref{eq:root-lattice-pairings} 表明 $\widetilde m\in\Delta$，因为对 $\rho'\ne\rho$ 有 $a_{\rho'}\geq0>-1$。事实上，后面这些严格不等式与 $\langle\widetilde m,v_\rho\rangle=-1$ 意味着，$\widetilde m$ 位于 $\Delta$ 中由 $\langle m,v_\rho\rangle=-1$ 定义的刻面的相对内部。反过来，给定这个刻面相对内部的格点 $\widetilde m$，取
\[
x^D=\prod_{\rho'\ne\rho}x_{\rho'}^{\langle\widetilde m,v_{\rho'}\rangle},
\]
容易验证 $(x_\rho,x^D)$ 是一个根。
\end{proof}

举例说明这种方法：$\mathbb P^4$ 是 Fano 簇，因此 $\mathbb P^4=\mathbb P_\Delta$，其中 $\Delta=\Delta_{-K_{\mathbb P^4}}$ 自反。我们知道自同构群 $\operatorname{Aut}(\mathbb P^4)$，但 $X^\circ=\mathbb P_{\Delta^\circ}$ 的自同构如何？这里 $\Delta^\circ$ 是 $\Delta$ 的对偶多胞形。\hyperref[sec:quintic-revisited]{第 4.2 节}将说明，单形 $\Delta^\circ$ 的顶点为
\[
\begin{gathered}
(-1,-1,-1,-1),\\
(1,0,0,0),\\
(0,1,0,0),\\
(0,0,1,0),\\
(0,0,0,1).
\end{gathered}
\]
注意，这个多胞形的刻面没有内部格点。由于 $\Delta^\circ$ 是 $X^\circ$ 的反典范类对应的多胞形，命题~\ref{prop:automorphism-facets} 表明 $X^\circ$ 没有根。因此，$\operatorname{Aut}(X^\circ)$ 的单位元连通分支恰好是 $X^\circ$ 的环面。特别地，虽然原来的环面簇 $X=\mathbb P^4$ 有很大的自同构群，它的对偶 $X^\circ$ 的自同构群却相对较小。\index{自同构群|hyperpage}

\section{例题}\label{sec:toric-examples}

本节给出三个例题，说明目前已讨论的环面簇理论的不同方面。对不太熟悉环面几何的读者，仔细完成这些例题的细节是一项很好的练习。

\begin{example}\label{ex:toric-projective-space}
环面簇最基本的例子是 $\mathbb P^n$。此时 $N=\mathbb Z^n$，一维锥的生成元是标准基 $e_1,\ldots,e_n$ 以及 $e_0=-\sum_{i=1}^n e_i$；仍按惯例不区分一维锥与其生成元。由 $\{e_0,\ldots,e_n\}$ 的所有真子集所生成的锥组成一个光滑扇。正合列~\eqref{eq:toric-chow-sequence} 变为
\[
0\longrightarrow\mathbb Z^n\longrightarrow\mathbb Z^{n+1}
\longrightarrow\mathbb Z\longrightarrow0,
\]
其中 $(a_0,\ldots,a_n)\in\mathbb Z^{n+1}$ 映到 $a_0+\cdots+a_n$。因此，\IndexLocation{indexloc:1369:72}齐次坐标环就是赋有通常分次的 $\mathbb C[x_0,\ldots,x_n]$，定理~\ref{thm:toric-quotient} 给出的 $\mathbb P^n$ 的构造就是通常的 $\mathbb P^n=(\mathbb C^{n+1}-\{0\})/\mathbb C^*$。

矩映射为 $\mu(z_0,\ldots,z_n)=\tfrac12\sum_{i=0}^n|z_i|^2$，所以对 $r>0$，$\mu^{-1}(r)$ 是一个球面，而定理~\ref{thm:toric-symplectic-reduction} 的构造就是纤维为 $S^1$ 的 Hopf 纤维化 $S^{2n+1}\to\mathbb P^n$。

最后，对应于 $e_0$ 的除子 $D_0$ 是丰沛的，其对应多胞形为单形
\[
\Delta_{D_0}=\{(c_1,\ldots,c_n):c_i\geq0,\ c_1+\cdots+c_n\leq1\}.
\]
\end{example}

\begin{example}\label{ex:weighted-resolution}
下一个例题涉及\IndexLocation{indexloc:1432:72}加权射影空间及其环面消解。在 $N=\mathbb Z^4$ 中考虑向量
\[
v_0=(-1,-2,-2,-2),\quad v_1=e_1,\quad v_2=e_2,
\quad v_3=e_3,\quad v_4=e_4.
\]
\SourcePage{73}
这些向量满足关系 $v_0+v_1+2v_2+2v_3+2v_4=0$。使用由 $\{v_0,v_1,v_2,v_3,v_4\}$ 的真子集构成的扇，得到加权射影空间 $\mathbb P(1,1,2,2,2)$；由引理~\ref{lem:weighted-projective-fano}，它是 Fano 簇。正合列~\eqref{eq:toric-chow-sequence} 为
\[
0\longrightarrow\mathbb Z^4\longrightarrow\mathbb Z^5
\longrightarrow\mathbb Z\longrightarrow0,
\]
其中 $(a_0,a_1,a_2,a_3,a_4)\in\mathbb Z^5$ 映到 $a_0+a_1+2a_2+2a_3+2a_4$。因此，齐次坐标环是 $\mathbb C[x_0,x_1,x_2,x_3,x_4]$，其中 $x_0,x_1$ 的次数为 $1$，$x_2,x_3,x_4$ 的次数为 $2$。

为消解 $\mathbb P(1,1,2,2,2)$ 的奇点，取由 $v_0,v_1$ 生成的奇异锥，加入生成元
\[
v_5=\tfrac12(v_0+v_1)=(0,-1,-1,-1).
\]
这给出新扇 $\Sigma$：原扇中每个包含 $v_0,v_1$ 的四维锥，都被分成 $\Sigma$ 中的两个锥。$\Sigma$ 是光滑的，因此对应的环面簇 $X=X_\Sigma$ 是 $\mathbb P(1,1,2,2,2)$ 的奇点消解。

对 $X$，可以看到群 $G=(\mathbb C^*)^2$ 在 $\mathbb C^{\Sigma(1)}=\mathbb C^6$ 上的作用为
\[
\begin{aligned}
&(\lambda,\mu)(x_0,x_1,x_2,x_3,x_4,x_5)\\
&\qquad=(\lambda x_0,\lambda x_1,\mu x_2,\mu x_3,\mu x_4,
\lambda^{-2}\mu x_5).
\end{aligned}
\]
此外，仅有的两个\IndexLocation{indexloc:1087:73}本原集合是 $\{v_0,v_1\}$ 与 $\{v_2,v_3,v_4,v_5\}$。由式~\eqref{eq:toric-primitive-z}，
$Z(\Sigma)=\{x_0=x_1=0\}\cup\{x_2=x_3=x_4=x_5=0\}$。因此
\[
X=(\mathbb C^6-Z(\Sigma))/(\mathbb C^*)^2.
\]

还可以描述矩映射与 GKZ 分解。因为 $v_0,v_1,v_2,v_3,v_4,v_5$ 之间的关系由 $v_2+v_3+v_4+v_5=0$ 及 $v_0+v_1-2v_5=0$ 生成，式~\eqref{eq:toric-moment-map} 表明矩映射为
\[
\begin{aligned}
&\mu_\Sigma(z_0,z_1,z_2,z_3,z_4,z_5)\\
&\quad=\tfrac12\bigl(|z_2|^2+|z_3|^2+|z_4|^2+|z_5|^2,
|z_0|^2+|z_1|^2-2|z_5|^2\bigr).
\end{aligned}
\]
$A_3^+(X_\Sigma)\otimes\mathbb R\subset A_3(X_\Sigma)\otimes\mathbb R\simeq\mathbb R^2$ 的 GKZ 分解特别简单，因为生成元包含在 $\Sigma(1)$ 中的射影单纯扇只有两个：$X$ 的扇 $\Sigma$，以及加权射影空间 $\mathbb P(1,1,2,2,2)$ 原来的扇 $\Sigma'$。留给读者作为练习，验证 $A_3^+(X_\Sigma)\otimes\mathbb R$ 的 GKZ 分解如下：
\begin{figure}[H]
\refstepcounter{equation}\label{eq:weighted-gkz-picture}
\centering
\makebox[\textwidth]{\makebox[0pt][l]{(\theequation)}\hfill
\raisebox{-.5\height}{\includegraphics[width=0.56067\textwidth]{3-1-gkz-decomposition.png}}\hfill\mbox{}}
\end{figure}

转向次级扇，注意 $\Xi=\{v_0,\ldots,v_5\}$，因此由式~\eqref{eq:extended-rays}，$\Xi^+$ 由七个点 $(v_0,1),\ldots,(v_5,1),(0,1)$ 组成。它们之间的所有关系构成的空间有一组基，
\SourcePage{74}
由以下矩阵的各行给出：
\[
\begin{pmatrix}
0&0&1&1&1&1&-4\\
1&1&0&0&0&-2&0
\end{pmatrix}.
\]
由\hyperref[subsec:secondary-fan]{第 3.4.2 小节}，在 $A(\Xi)\simeq\mathbb R^2$ 的对偶基下，此矩阵的各列给出次级扇各边的生成元。事实上，从式~\eqref{eq:weighted-gkz-picture} 可见，上述矩阵的前六列给出 GKZ 分解，而最后一列通过沿 $r_1$ 轴负半轴插入一条边，将其扩充为完备扇。
\end{example}

\begin{example}\label{ex:cube-octahedron}
设 $M=\mathbb Z^3$，考虑以原点为中心、顶点为 $(\pm1,\pm1,\pm1)$ 的立方体 $\Delta\subset M_{\mathbb R}$。它给出环面簇 $X=\mathbb P_\Delta$。为描述 $X$ 的扇，注意极多胞形 $\Delta^\circ\subset N_{\mathbb R}$ 是以 $\pm e_1,\pm e_2,\pm e_3$ 为顶点的八面体。因此，法扇由八面体的各面构成，得到扇 $\Sigma$，其三维锥就是 $\mathbb R^3$ 的各卦限。容易推出 $X=\mathbb P^1\times\mathbb P^1\times\mathbb P^1$。此外，来自 $\Delta$ 的除子 $\mathcal O_X(1)$ 是反典范除子 $-K_X$。特别地，$X$ 是 Fano 簇。

$\Delta$ 与 $\Delta^\circ$ 的图形如下：
\begin{figure}[H]
\centering
\includegraphics[width=0.75174\textwidth]{3-2-cube-octahedron.png}
\end{figure}

容易验证，立方体 $\Delta\subset M_{\mathbb R}$ 与八面体 $\Delta^\circ\subset N_{\mathbb R}$ 是\hyperref[sec:fano-reflexive]{第 3.5 节}意义下的一对对偶\IndexLocation{indexloc:1222:74}自反多胞形。特别地，$\Delta^\circ$ 给出“对偶”环面簇 $X^\circ=\mathbb P_{\Delta^\circ}$，它由 $\Delta^\circ$ 的法扇确定；该法扇就是在 $M_{\mathbb R}$ 中向立方体 $\Delta$ 的各面作锥所得的扇。因此得到一对“对偶”环面簇 $X$ 与 $X^\circ$。有趣的是，$X$ 光滑，而 $X^\circ$ 却有相当严重的奇点。事实上，$\Sigma^\circ$ 的三维锥甚至不是单纯的：它们都是无限延伸的四棱锥。不过，$\Delta^\circ$ 自反，所以知道 $X^\circ$ 是 Fano 簇，而且 $\Delta^\circ\subset N_{\mathbb R}$ 是反典范除子 $-K_{X^\circ}$ 对应的多胞形。

还应注意，$\Delta$ 与 $\Delta^\circ$ 的格点数及格点几何分布也很不相同。对 $\Delta^\circ\subset N_{\mathbb R}$，属于 $N$ 的格点只有原点与各顶点；而 $\Delta\subset M_{\mathbb R}$ 中的格点要多得多，因为各棱的中点与各面的中心也属于 $M$。

\SourcePage{75}
还可以描述 $X^\circ$ 的 GKZ 分解。分解的极大锥对应于一维锥生成元位于 $\Sigma^\circ(1)=\{\text{立方体的顶点}\}$ 中的射影单纯扇。对原扇作单纯细分时，需要沿一条对角线将立方体的每个面分成两个三角形。每个面有两种分法，因此 GKZ 分解有 $2^6=64$ 个锥。不过，\hyperref[chap:mirror-constructions]{第 4 章}将说明，为了研究镜像对称，需要使用一维锥生成元位于 $\Delta\cap M$ 中的射影单纯扇来形成 GKZ 分解，而 $\Delta\cap M$ 远大于仅由立方体顶点组成的集合。因此，所得 GKZ 分解的锥数远多于 $2^6$。

这个例题表明，可以从一个光滑环面簇，例如 $\mathbb P^1\times\mathbb P^1\times\mathbb P^1$，出发，通过“对偶”过程得到有奇点的对象。\hyperref[chap:quintic]{第 2 章}中五次三维簇及其镜像出现类似情形，绝非偶然。\hyperref[chap:mirror-constructions]{第 4 章}将详细研究这一现象。
\end{example}

\chapter{镜像对称的构造}\label{chap:mirror-constructions}
\SourcePage{76}

本章介绍镜像流形的各种构造。其中大部分构造涉及环面簇，但并非全部如此。特别地，我们将从环面几何的角度解释五次超曲面的镜像对称。

\section{Batyrev 镜像构造}\label{sec:batyrev-mirror}

\hyperref[chap:introduction]{第 1 章}从超共形场论的角度讨论了镜像对称为何合理，却没有说明如何寻找给定 Calabi--Yau 三维簇的镜像。物理中最早的镜像流形显式例子由 Greene 与 Plesser \SourceCite{GPl} 给出，他们将镜像构造成加权射影空间中 Fermat 超曲面的商。\hyperref[chap:quintic]{第 2 章}给出的五次三维簇的镜像，就是这种构造的一个例子。Roan 观察到，环面几何提供了描述这类镜像的自然框架 \SourceCite{Roan1}。随后，Batyrev 找到了一种漂亮的环面几何描述，大大推广了 Greene--Plesser 构造。

Batyrev 的构造关键地使用了\hyperref[sec:fano-reflexive]{第 3.5 节}介绍的自反多胞形。那里已看到，自反多胞形 $\Delta$ 对应于 Fano 环面簇，而且 $\Delta$ 有一个同样自反的对偶 $\Delta^\circ$。本节将解释，$\Delta$ 与 $\Delta^\circ$ 之间的对偶性，如何导出某些与 $\mathbb P_\Delta$、$\mathbb P_{\Delta^\circ}$ 密切相关的环面簇中 Calabi--Yau 超曲面族之间的对偶性。\index{Batyrev镜像构造@Batyrev 镜像构造|hyperpage}

\subsection{环面簇中的 Calabi--Yau 超曲面}\label{subsec:toric-cy-hypersurfaces}

Batyrev 构造的第一步，是观察到环面簇中的 Calabi--Yau 超曲面自然来自自反多胞形。自反多胞形 $\Delta$ 给出环面簇 $\mathbb P_\Delta$。不过，对于镜像对称，$\mathbb P_\Delta$ 的奇点可能过于严重。我们需要能在其上进行 Hodge 理论的对象，因此将使用细化 $\mathbb P_\Delta$ 的扇的单纯扇。回顾，$\mathbb P_\Delta$ 的扇是 $\Delta$ 的法扇，由向 $\Delta^\circ\subset N_{\mathbb R}$ 的各真面作出的锥组成。此外，因为 $\Delta$ 自反，法扇的一维锥生成元就是 $\Delta^\circ$ 的各顶点，它们属于 $\Delta^\circ\cap N-\{0\}$。我们将考虑 $N_{\mathbb R}$ 中的以下扇。

\begin{definition}\label{def:projective-subdivision}
给定自反多胞形 $\Delta\subset M_{\mathbb R}$，若 $N_{\mathbb R}$ 中的扇 $\Sigma$ 具有以下性质，则称它为\emph{射影细分}：
\begin{itemize}
\item $\Sigma$ 细化 $\Delta$ 的法扇；
\item $\Sigma(1)\subset\Delta^\circ\cap N-\{0\}$；
\item $X_\Sigma$ 射影且单纯。
\end{itemize}
此外，若 $\Sigma(1)=\Delta^\circ\cap N-\{0\}$，则称 $\Sigma$ 为\emph{极大}的。
\end{definition}

可以证明，极大射影细分 $\Sigma$ 存在 \SourceCite{OdP}。这样的 $\Sigma$ 有多种选择。如\hyperref[sec:gkz-decomposition]{第 3.4 节}所述，每种选择给出
\SourcePage{77}
GKZ 分解中的一个不同的锥，也给出相关规范线性 $\sigma$ 模型中的一个不同的相，见\hyperref[app:physical-theories]{附录 B}。\hyperref[subsec:multiple-mirrors]{第 4.1.4 小节}将更详细地讨论这一点。

给定射影细分 $\Sigma$，得到环面簇 $X_\Sigma$ 以及双有理态射 $f:X_\Sigma\to\mathbb P_\Delta$，它们具有以下良好性质。
\begin{lemma}\label{lem:projective-subdivision-properties}
若 $\Sigma$ 是射影细分，则：
\begin{enumerate}[label=(\roman*)]
\item $X_\Sigma$ 是 Gorenstein 轨形；
\item $\Delta$ 是 $-K_{X_\Sigma}$ 对应的多胞形；
\item $-K_{X_\Sigma}$ 是半丰沛的，并且是大的；具体说，$-K_{X_\Sigma}$ 由整体截面生成，且 $(-K_{X_\Sigma})^n>0$；
\item 映射 $f:X_\Sigma\to\mathbb P_\Delta$ 是保典范的，即 $f^*(K_{\mathbb P_\Delta})=K_{X_\Sigma}$。
\end{enumerate}
此外，若 $\Sigma$ 是极大射影细分，则 $X_\Sigma$ 具有终端奇点。
\end{lemma}
\begin{proof}
因为 $\Delta^\circ$ 自反，它的每个刻面 $\Gamma^\circ$ 都由 $\langle m,v\rangle=-1$ 定义，其中某个 $m=m_{\Gamma^\circ}\in M$。于是当 $v_\rho\in\Gamma^\circ$ 时，有 $\langle m,v_\rho\rangle=-1$，根据\hyperref[sec:cones-fans]{第 3.1 节}的准则，容易推出 $-K_{X_\Sigma}=\sum_\rho D_\rho$ 是 Cartier 除子。因此 $X=X_\Sigma$ 是 Gorenstein 的；又因 $\Sigma$ 单纯，它也是轨形。此外，等式 $\langle m,v_\rho\rangle=-1$ 表明，与除子 $-K_X$ 对应的多胞形恰好是 $\Delta$。由对偶性，各 $m$ 就是 $\Delta$ 的顶点，\hyperref[subsec:toric-polytopes]{第 3.2.1 小节}表明 $-K_X$ 由整体截面生成。最后，由 \SourceCite[第 5.3 节]{Fulton3}，交数 $(-K_{X_\Sigma})^n=n!\operatorname{vol}(\Delta)$；由于 $\Delta$ 是 $n$ 维的，这个数为正。

为证明引理余下的部分，研究向 $\Sigma$ 加入新的一维锥生成元如何影响 $X$ 的典范类。回顾定义~\ref{def:fano} 后的讨论，对任意环面簇，$K_X$ 是 Weil 除子 $-\sum_{\rho\in\Sigma(1)}D_\rho$。取锥 $\sigma\in\Sigma$，在其相对内部加入新的本原整生成元 $v'\in N$，所得环面簇记为 $X'$。已知 $\sigma$ 的生成元 $v_\rho$ 位于 $\Delta^\circ$ 的某个刻面 $\Gamma^\circ$ 上，因此 $\langle m,v_\rho\rangle=-1$，其中 $m=m_{\Gamma^\circ}$ 如上。关键计算是：若 $g:X'\to X$ 为自然映射，则
\begin{equation}\label{eq:toric-canonical-discrepancy}
K_{X'}=g^*(K_X)-\bigl(\langle m,v'\rangle+1\bigr)D',
\end{equation}
其中 $D'$ 是 $X'$ 上对应于 $v'$ 的除子。这个等式应理解为 Weil 除子的线性等价关系。例如，可用 \SourceCite[第 4 节]{Reid4} 的方法证明式~\eqref{eq:toric-canonical-discrepancy}。只需考察其中 $D'$ 的系数：在仿射开集 $X_\sigma\subset X$ 上，$K_X$ 与 $\operatorname{div}(\chi^m)$ 一致，因此 $g^*K_X$ 中 $D'$ 的系数为 $\langle m,v'\rangle$。

若 $v'\in\Delta^\circ\cap N-\{0\}$，则 $v'$ 位于刻面 $\Gamma^\circ$ 上，所以 $\langle m,v'\rangle=-1$，由式~\eqref{eq:toric-canonical-discrepancy} 立即得到 $K_{X'}=g^*(K_X)$。因为 $\Sigma(1)\subset\Delta^\circ\cap N$，$X_\Sigma$ 由 $\mathbb P_\Delta$ 通过依次加入 $\Delta^\circ\cap N$ 中的新顶点得到，容易推出 $K_{X_\Sigma}=f^*(K_{\mathbb P_\Delta})$。

另一方面，若 $\Sigma$ 极大，则 $\Sigma(1)=\Delta^\circ\cap N-\{0\}$。因此，任何新加入的本原整生成元 $v'$ 必定位于 $\Delta^\circ$ 之外；按自反的定义，$0$ 是唯一的内部格点。于是式~\eqref{eq:toric-canonical-discrepancy} 中 $\langle m,v'\rangle+1<0$。由\hyperref[app:singular-varieties]{附录 A} 的讨论，可见 $X$ 具有终端奇点。
\end{proof}

下面说明如何利用 $X_\Sigma$ 和 $\mathbb P_\Delta$ 上的反典范线性系构造 Calabi--Yau 簇。\index{射影细分|hyperpage}\index{保典范态射|hyperpage}

\SourcePage{78}
\begin{proposition}\label{prop:toric-cy-construction}
若 $\Delta$ 是 $n$ 维自反多胞形，则一般成员 $\overline V\in|-K_{\mathbb P_\Delta}|$ 是 $n-1$ 维 Calabi--Yau 簇。此外，若 $\Sigma$ 是射影细分，$X=X_\Sigma$，则：
\begin{enumerate}[label=(\roman*)]
\item 一般成员 $V\in|-K_X|$ 是 Calabi--Yau 轨形；
\item 若 $\Sigma$ 极大，则一般成员 $V\in|-K_X|$ 是极小 Calabi--Yau 轨形。
\end{enumerate}
\end{proposition}
\begin{proof}
先考虑 $\overline V$。根据\hyperref[sec:math-preview]{第 1.4 节}给出的 Calabi--Yau 簇的定义，必须证明 $\overline V$ 具有典范奇点、平凡的典范层，而且对 $0<k<n-1$，上同调 $H^k(\overline V,\mathcal O_{\overline V})=0$。

由于 Fano 环面簇 $\mathbb P_\Delta$ 是 Gorenstein 的，它至多具有典范奇点，见 \SourceCite{Batyrev4}。于是 Bertini 定理的一个版本保证，一般成员 $\overline V\in|-K_{\mathbb P_\Delta}|$ 也至多具有典范奇点，见 \SourceCite[定理 1.13]{Reid1}。还应注意
\[
\widehat\Omega_{\overline V}^{n-1}
\simeq\widehat\Omega_{\mathbb P_\Delta}^{n}
(-K_{\mathbb P_\Delta})\otimes\mathcal O_{\overline V}
\simeq\mathcal O_{\overline V},
\]
其中第一个同构是伴随公式；因为 $\mathbb P_\Delta$ 是 Cohen--Macaulay 的，且 $-K_{\mathbb P_\Delta}$ 是 Cartier 除子，该公式成立。最后一步是证明 $0<k<n-1$ 时 $H^k(\overline V,\mathcal O_{\overline V})=0$。由于
$\mathcal O_{\mathbb P_\Delta}(-\overline V)\simeq\mathcal O_{\mathbb P_\Delta}(K_{\mathbb P_\Delta})=\widehat\Omega_{\mathbb P_\Delta}^n$，得到正合列
\[
0\longrightarrow\widehat\Omega_{\mathbb P_\Delta}^n
\longrightarrow\mathcal O_{\mathbb P_\Delta}
\longrightarrow\mathcal O_{\overline V}\longrightarrow0,
\]
进而得到长正合列
\[
\cdots\longrightarrow H^k(\mathbb P_\Delta,\mathcal O_{\mathbb P_\Delta})
\longrightarrow H^k(\overline V,\mathcal O_{\overline V})
\longrightarrow H^{k+1}(\mathbb P_\Delta,\widehat\Omega_{\mathbb P_\Delta}^n)
\longrightarrow\cdots.
\]
然而，对 $k>0$ 有 $H^k(\mathbb P_\Delta,\mathcal O_{\mathbb P_\Delta})=0$；由 Serre--Grothendieck 对偶性，对 $k<n-1$ 还有
\[
H^{k+1}(\mathbb P_\Delta,\widehat\Omega_{\mathbb P_\Delta}^n)
\simeq H^{n-k-1}(\mathbb P_\Delta,\mathcal O_{\mathbb P_\Delta})^*=0.
\]
这给出所需的 $H^k(\overline V,\mathcal O_{\overline V})$ 的消失，因此 $\overline V$ 是 Calabi--Yau 簇。

现在令 $\Sigma$ 是射影细分，置 $X=X_\Sigma$。由引理~\ref{lem:projective-subdivision-properties}，$-K_X$ 半丰沛，所以线性系 $|-K_X|$ 没有基点。此外，$\Sigma$ 单纯，所以 $X$ 是轨形；将 Bertini 定理应用于定义 $X$ 的局部商的各不动点集，便保证一般成员 $V\in|-K_X|$ 是 $X$ 的子轨形。上面的全部论证仍然成立，因此 $V$ 是 Calabi--Yau 轨形。

最后设 $\Sigma$ 极大。根据定义~\ref{def:calabi-yau-variety}，$V$ 为极小 Calabi--Yau 簇的条件是它具有 Gorenstein、$\mathbb Q$ 析因、终端奇点。因为 $V$ 已是 Gorenstein 轨形，它的奇点自动为 Gorenstein 且 $\mathbb Q$ 析因。因此只需证明 $V$ 具有终端奇点。不过，由引理~\ref{lem:projective-subdivision-properties}，环境空间 $X$ 具有终端奇点，像 \SourceCite[定理 1.13]{Reid1} 的证明那样应用 Bertini 定理即可。
\end{proof}

由于 $f:X_\Sigma\to\mathbb P_\Delta$ 保典范，上面构造的 Calabi--Yau 超曲面 $V\subset X_\Sigma$ 是 $\mathbb P_\Delta$ 中一般反典范超曲面 $\overline V$ 经 $f$ 的严格变换。还可以证明，诱导映射 $V\to\overline V$ 也保典范。

当 $\Sigma$ 是极大射影细分时，簇 $V\subset X_\Sigma$ 是 $\overline V\subset\mathbb P_\Delta$ 在 Mori 理论意义下的极小模型。Batyrev 称 $V$ 为 $\overline V$ 的\emph{极大射影保典范部分消解}，简称 \emph{MPCP 消解}。

当 $V$ 是三维簇时，情形尤其好。由\hyperref[app:singular-varieties]{附录 A}，三维 Gorenstein 轨形的终端奇点实际上是光滑点。因此，对极大射影细分，命题~\ref{prop:toric-cy-construction} 构造的极小 Calabi--Yau 三维簇 $V\subset X_\Sigma$ 都是光滑的。\index{MPCP消解@MPCP 消解|hyperpage}

\SourcePage{79}
\subsection{环面 Calabi--Yau 超曲面的镜像}\label{subsec:batyrev-pair}

现在可以描述 Batyrev 镜像构造。固定自反多胞形 $\Delta$，令 $\Sigma$ 是定义~\ref{def:projective-subdivision} 中其法扇的一个极大射影细分。这样得到由 $V\in|-K_{X_\Sigma}|$ 给出的 Calabi--Yau 超曲面族 $V\subset X_\Sigma$。如\hyperref[chap:introduction]{第 1 章}所述，这个族中的每个成员都带有一个复化 Kähler 锥，从而得到族 $\{(V,\omega)\}$，它由 Calabi--Yau 超曲面及其复化 Kähler 类的一切可能配对组成。

不过，$\Delta$ 自反意味着 $\Delta^\circ$ 自反，因此可对 $\Delta^\circ$ 重复上述构造。对 $\Delta^\circ$ 在 $M_{\mathbb R}$ 中的法扇的一个极大射影细分 $\Sigma^\circ$，$|-K_{X_{\Sigma^\circ}}|$ 给出 Calabi--Yau 超曲面族 $V^\circ\subset X_{\Sigma^\circ}$。加入各复化 Kähler 锥之后，得到族 $\{(V^\circ,\omega^\circ)\}$，称为原族 $\{(V,\omega)\}$ 的 \emph{Batyrev 镜像}。

物理学认为，这两个族是物理学定义~\ref{def:physics-mirror-pair} 意义下的镜像对。更精确地说，有以下猜想。
\begin{physicsconjecture}\label{conj:batyrev-physics}
族 $\{(V,\omega)\}$ 与 $\{(V^\circ,\omega^\circ)\}$ 诱导同构的超共形场论；它们的 $N=2$ 超共形表示除\hyperref[sec:physics]{第 1.1 节}讨论的符号改变外完全相同。
\end{physicsconjecture}

另一种说法是，族 $\{(V,\omega)\}$ 与 $\{(V^\circ,\omega^\circ)\}$ 分别表示 $N=2$ 超共形场论（SCFT）模空间\emph{同一个}连通分支中的开子集。严格说来，虽然 Batyrev 镜像构造适用于任意维数的自反多胞形，这个猜想只适用于 $V$ 与 $V^\circ$ 为三维簇的情形，此时 $\Delta$ 是四维的。

物理中尚未证明这个镜像对称猜想，不过 \SourceCite{MP2} 作过尝试。在某些加权射影空间中具有适当次数的 Fermat 超曲面这一特殊情形，Greene--Plesser 构造 \SourceCite{GPl} 从物理上证明了 $\{(V,\omega)\}$ 与 $\{(V^\circ,\omega^\circ)\}$ 构成镜像对。一般情形仍未解决。

由于 SCFT 缺少严格基础，物理学猜想~\ref{conj:batyrev-physics} 并非通常数学意义下的猜想。不过，可以把这种 SCFT 相等关系的数学推论精确地表述出来。特别地，\hyperref[chap:introduction]{第 1 章}已看到，镜像对称意味着下列断言应当成立：
\begin{itemize}
\item 上同调群 $H^{1,1}(V)$，相应地 $H^{1,1}(V^\circ)$，应同构于 $H^{n-2,1}(V^\circ)$，相应地 $H^{n-2,1}(V)$。回顾，$V$ 与 $V^\circ$ 的维数为 $n-1$，其中 $n=\dim\Delta$。
\item 根据镜像对称的图~\eqref{eq:mirror-foliations}，$V$，相应地 $V^\circ$，的 Kähler 模应通过镜像映射局部同构于 $V^\circ$，相应地 $V$，的复模。
\item $V$，相应地 $V^\circ$，的 A 模型关联函数应通过镜像映射，对应于经适当归一化的 $V^\circ$，相应地 $V$，的 B 模型关联函数。
\end{itemize}
数学家面临的挑战是证明这些推论。当 $V^\circ$ 是 $V$ 的 Batyrev 镜像时，已有实质性进展。\hyperref[subsec:batyrev-hodge]{第 4.1.3 小节}将解决第一项；\hyperref[chap:moduli]{第~\ref*{chap:moduli} 章}研究环面 Calabi--Yau 超曲面的复模与 Kähler 模时，将处理第二项。第三项则要等到\hyperref[chap:quantum-cohomology]{第~\ref*{chap:quantum-cohomology} 章}，那里将定义\IndexLocation{indexloc:0748:79}数学镜像对，并提出 Hodge 理论的\IndexLocation{indexloc:0795:79}环面镜像猜想。\hyperref[chap:mirror-theorem]{第~\ref*{chap:mirror-theorem} 章}将讨论若干特殊情形的证明。

\SourcePage{80}
\subsection{Batyrev 镜像的 Hodge 数}\label{subsec:batyrev-hodge}

要证明两个有限维向量空间同构，只须证明它们的维数相同。因此，上面的第一项要求比较环面 Calabi--Yau 超曲面 $V$ 与其 Batyrev 镜像 $V^\circ$ 的 Hodge 数。

仍假定 $\Delta$ 是 $n$ 维自反多胞形，$V$ 是 $X_\Sigma$ 中的一般反典范超曲面，其中 $\Sigma$ 是 $\Delta$ 的法扇的极大射影细分。对 $\Delta^\circ$ 应用这个过程，得到 Batyrev 镜像族 $V^\circ$。\SourceCite{Batyrev4,AGM1} 证明，$V$ 与 $V^\circ$ 的 Hodge 数有如下关系。
\begin{theorem}\label{thm:batyrev-hodge}
若 $V^\circ$ 是 $V$ 的 Batyrev 镜像，则
\[
h^{1,1}(V)=h^{n-2,1}(V^\circ),\qquad
h^{n-2,1}(V)=h^{1,1}(V^\circ).
\]
\end{theorem}
\begin{proof}
当 $n=2$ 或 $3$ 时，结论分别由椭圆曲线与 K3 曲面的 Hodge 数得到。以下设 $n\geq4$。
先定义 $H^{1,1}(V)$ 与 $H^{n-2,1}(V)$ 中的一些有趣的子空间。对每个 $v_i\in\Sigma(1)$，令 $D_i\subset X_\Sigma$ 为相应的 $T_N$ 不变除子。这些除子在超曲面 $V\subset X_\Sigma$ 上限制为除子 $D_i'$；其上同调类生成一个复数域上的子空间，记为
\begin{equation}\label{eq:toric-hodge-subspace}
H_{\mathrm{toric}}^{1,1}(V)\subset H^{1,1}(V).
\end{equation}
再考虑 $H^{n-2,1}(V)\simeq H^1(V,T_V)$。首先注意，虽然 $V$ 可能有奇点，它却是至多具有终端奇点的 Gorenstein 轨形；因此，根据命题~\ref{prop:terminal-gorenstein-orbifold-deformations}，$H^1(V,T_V)$ 分类 $V$ 的无穷小形变。由超曲面 $V\in|-K_{X_\Sigma}|$ 确定的无穷小形变给出子空间，记为
\begin{equation}\label{eq:polynomial-hodge-subspace}
H_{\mathrm{poly}}^{n-2,1}(V)\subset H^{n-2,1}(V).
\end{equation}
我们断言，这些子空间的维数可用 $\Delta$ 与 $\Delta^\circ$ 如下计算：
\begin{equation}\label{eq:toric-hodge-number}
h_{\mathrm{toric}}^{1,1}(V)
=l(\Delta^\circ)-n-1-\sum_{\Gamma^\circ}l^*(\Gamma^\circ),
\end{equation}
\begin{equation}\label{eq:polynomial-hodge-number}
h_{\mathrm{poly}}^{n-2,1}(V)
=l(\Delta)-n-1-\sum_\Gamma l^*(\Gamma),
\end{equation}
其中求和分别遍历 $\Delta^\circ$ 与 $\Delta$ 的所有刻面 $\Gamma^\circ$ 和 $\Gamma$；$l$、$l^*$ 如式~\eqref{eq:lattice-points} 定义。一旦证明这个断言，便立即得到等式
\[
h_{\mathrm{toric}}^{1,1}(V)=h_{\mathrm{poly}}^{n-2,1}(V^\circ),
\qquad
h_{\mathrm{toric}}^{1,1}(V^\circ)=h_{\mathrm{poly}}^{n-2,1}(V),
\]
因为由 $V$ 换到 $V^\circ$，正好交换了式~\eqref{eq:toric-hodge-number}、\eqref{eq:polynomial-hodge-number} 中 $\Delta$ 与 $\Delta^\circ$ 的角色。还应提到，可以构造自然同构 $H_{\mathrm{toric}}^{1,1}(V)\simeq H_{\mathrm{poly}}^{n-2,1}(V^\circ)$，称为\emph{单项式--\IndexLocation{indexloc:0910:80}除子镜像映射}，见 \SourceCite{AGM1}。

先证明式~\eqref{eq:toric-hodge-number}。设 $v_i\in\Sigma(1)$ 是 $\Delta^\circ$ 的某个刻面的内部格点。环面爆破的标准描述表明，$f(D_i)$ 是一个点，其中 $f:X_\Sigma\to\mathbb P_\Delta$ 为自然映射。立即可知，在 $X_\Sigma$ 中，$D_i$ 与一般的 $V\in|-K_{X_\Sigma}|$ 不相交。因此这些元素对 $H_{\mathrm{toric}}^{1,1}(V)$ 没有贡献。置
\[
\Sigma(1)'=\{v_i\in\Sigma(1):
v_i\text{ 不在 }\Delta^\circ\text{ 的任何刻面的相对内部}\},
\]
则 $H_{\mathrm{toric}}^{1,1}(V)$ 由 $v_i\in\Sigma(1)'$ 对应的 $D_i'$ 生成。要找出各 $D_i'$ 之间的关系，注意 $m\in M$ 给出 $T_N$ 的特征标 $\chi^m$，从而给出
\SourcePage{81}
$X_\Sigma$ 上的有理函数。式~\eqref{eq:toric-principal-divisor} 已给出 $\chi^m$ 的除子；将该公式限制到 $V$，得到线性等价关系
\begin{equation}\label{eq:toric-restricted-divisor-relation}
\sum_{v_i\in\Sigma(1)'}\langle m,v_i\rangle D_i'\sim0.
\end{equation}
\SourceCite{AGM1} 证明，这些就是各 $D_i'$ 之间的全部关系。因此
\[
H_{\mathrm{toric}}^{1,1}(V)
\simeq\bigl(\mathbb Z^{\Sigma(1)'}/M\bigr)\otimes\mathbb C,
\]
其中嵌入 $M\to\mathbb Z^{\Sigma(1)'}$ 按式~\eqref{eq:toric-restricted-divisor-relation} 的要求，由 $m\mapsto(\langle m,v_i\rangle)$ 给出。于是
\[
h_{\mathrm{toric}}^{1,1}(V)=|\Sigma(1)'|-n.
\]
用式~\eqref{eq:lattice-points} 的记号，因为 $\Sigma(1)=\Delta^\circ\cap N-\{0\}$，可见 $|\Sigma(1)|=l(\Delta^\circ)-1$；这里用到了射影细分 $\Sigma$ 的极大性。去掉各刻面的内部格点后，得到
$|\Sigma(1)'|=l(\Delta^\circ)-1-\sum_{\Gamma^\circ}l^*(\Gamma^\circ)$，式~\eqref{eq:toric-hodge-number} 随即成立。

现在考虑 $h_{\mathrm{poly}}^{n-2,1}(V)$。我们将描述来自族 $V\in|-K_{X_\Sigma}|$ 的模。这里先作非正式的讨论，\hyperref[chap:moduli]{第~\ref*{chap:moduli} 章}会作仔细处理。引理~\ref{lem:projective-subdivision-properties} 表明，$\Delta$ 是 $-K_{X_\Sigma}$ 对应的多胞形。因此，Laurent 多项式 $f\in L(\Delta)$ 给出仿射超曲面 $Z_f\subset T_N$，而且可设 $V=\overline Z_f\subset X_\Sigma$。由于将 $f$ 乘以非零常数不改变 $V$，变动 $f$ 得到一个维数为 $l(\Delta)-1$ 的多项式形变参数空间。不过，必须计入 $X_\Sigma$ 的自同构。因为 $\Delta$ 对应于 $-K_{X_\Sigma}$，命题~\ref{prop:automorphism-facets} 表明，自同构群的维数为
\[
\dim\operatorname{Aut}(X_\Sigma)=n+\sum_\Gamma l^*(\Gamma).
\]
现在容易推出，多项式模的维数由式~\eqref{eq:polynomial-hodge-number} 给出。

要完成定理的证明，还需理解 $V$ 的“非环面”除子类与“非多项式”形变。Batyrev \SourceCite{Batyrev4} 证明，在这两种情形中，“修正项”均由 $\Delta$ 与 $\Delta^\circ$ 的余维数为二的面确定。公式为
\begin{equation}\label{eq:batyrev-h11}
\begin{aligned}
h^{1,1}(V)={}&l(\Delta^\circ)-n-1
-\sum_{\Gamma^\circ}l^*(\Gamma^\circ)\\
&+\sum_{\Theta^\circ}l^*(\Theta^\circ)l^*(\widehat\Theta^\circ),
\end{aligned}
\end{equation}
\begin{equation}\label{eq:batyrev-hn21}
\begin{aligned}
h^{n-2,1}(V)={}&l(\Delta)-n-1-\sum_\Gamma l^*(\Gamma)\\
&+\sum_\Theta l^*(\Theta)l^*(\widehat\Theta),
\end{aligned}
\end{equation}
其中第一个等式里的 $\Gamma^\circ$ 指 $\Delta^\circ$ 的余维数为一的面，$\Theta^\circ$ 指其余维数为二的面，$\widehat\Theta^\circ$ 指 $\Delta$ 中与 $\Theta^\circ$ 对偶的面。第二个等式的含义类似。如前所述，将这些公式分别应用于 $V$ 和 $V^\circ$，就交换了 $\Delta$ 与 $\Delta^\circ$ 的角色，因此定理立即由式~\eqref{eq:batyrev-h11}、\eqref{eq:batyrev-hn21} 推出。

式~\eqref{eq:batyrev-h11} 的修正项来自与一般 $V$ 的交可约的除子 $D_i$。前面已看到，可以忽略那些 $v_i$ 位于刻面 $\Gamma^\circ$ 相对内部的 $D_i$。一般地，设 $v_i$ 位于某个面 $\Theta^\circ\subset\Delta^\circ$ 的相对内部。若 $\widehat\Theta^\circ\subset\Delta$ 表示对偶面，则 $\mathbb P_{\widehat\Theta^\circ}\subset\mathbb P_\Delta$ 是轨道闭包，并且在爆破映射 $f:X_\Sigma\to\mathbb P_\Delta$ 下，有 $f(D_i)=\mathbb P_{\widehat\Theta^\circ}$。若 $\overline V=f(V)$ 是 $\mathbb P_\Delta$ 中相应的反典范超曲面，则 $D_i'=V\cap D_i$ 与 $\overline V\cap\mathbb P_{\widehat\Theta^\circ}$ 有相同的不可约分支数。

回顾，面 $\Theta^\circ$ 与其对偶 $\widehat\Theta^\circ$ 满足 $\dim\Theta^\circ+\dim\widehat\Theta^\circ=n-1$。因此，若 $\Theta^\circ$ 的余维数至少为三，则 $\widehat\Theta^\circ$ 的维数至少为二，所以 $\mathbb P_{\widehat\Theta^\circ}$ 的
\SourcePage{82}
维数至少为二。Bertini 定理表明 $\overline V\cap\mathbb P_{\widehat\Theta^\circ}$ 不可约。因此，当 $v_i$ 位于 $\Delta^\circ$ 的余维数至少为三的面中时，$D_i'$ 不可约。

另一方面，若 $\Theta^\circ$ 的余维数为二，则 $\widehat\Theta^\circ$ 的维数为一，$\mathbb P_{\widehat\Theta^\circ}$ 是一条曲线。在这种情形，\SourceCite[第 5.3 节]{Fulton3} 表明
\[
\overline V\cdot\mathbb P_{\widehat\Theta^\circ}
=\operatorname{vol}(\widehat\Theta^\circ)
=l^*(\widehat\Theta^\circ)+1,
\]
这里 $\operatorname{vol}$ 表示归一化体积；第二个等式成立，是因为 $\widehat\Theta^\circ$ 为一维。这说明 $\overline V\cap\mathbb P_{\widehat\Theta^\circ}$ 由 $l^*(\widehat\Theta^\circ)+1$ 个点组成，从而 $D_i'$ 有 $l^*(\widehat\Theta^\circ)+1$ 个连通分支。可以证明，这些分支在 $H^{1,1}(V)$ 中给出线性无关的类。它们的和 $D_i'$ 已经属于 $H_{\mathrm{toric}}^{1,1}(V)$，因此，$\Theta^\circ$ 的每个内部格点增加 $l^*(\widehat\Theta^\circ)$ 个新的类。因为 $\Sigma$ 极大，使用了 $\Theta^\circ$ 的全部内部格点，这就解释了式~\eqref{eq:batyrev-h11} 中的修正项。

至于式~\eqref{eq:batyrev-hn21} 的修正项，我们将解释这个公式为什么合理；完整证明参见 \SourceCite{Batyrev4,Voisin3}。令 $Z_f=V\cap T_N\subset T_N$ 为 $V$ 所确定的仿射超曲面，则补集 $V-Z_f=\sum_i V\cap D_i$ 是轨形意义下的正规交叉除子。因此，Gysin 谱序列给出正合列
\[
\begin{aligned}
\longrightarrow\bigoplus_i H^{n-3}(V\cap D_i,\mathbb C)(-1)
&\longrightarrow H^{n-1}(V,\mathbb C)\\
&\longrightarrow\operatorname{Gr}_{n-1}^W H^{n-1}(Z_f,\mathbb C)
\longrightarrow0.
\end{aligned}
\]
这些群都带有纯 Hodge 结构；由此可以证明
\begin{equation}\label{eq:gysin-hodge-number}
\begin{aligned}
h^{n-2,1}(V)={}&h^{n-2,1}
\bigl(\operatorname{Gr}_{n-1}^W H^{n-1}(Z_f,\mathbb C)\bigr)\\
&+\sum_i h^{n-3,0}(V\cap D_i).
\end{aligned}
\end{equation}
根据 Danilov 与 Khovanskii 的工作 \SourceCite[5.9]{DK}，右边第一项可计算为
\[
\begin{aligned}
&h^{n-2,1}\bigl(\operatorname{Gr}_{n-1}^W H^{n-1}(Z_f,\mathbb C)\bigr)\\
&\qquad=l^*(2\Delta)-(n+1)l^*(\Delta)-\sum_\Gamma l^*(\Gamma).
\end{aligned}
\]
不过，自反性意味着 $l^*(\Delta)=1$，而 $l^*(2\Delta)=l(\Delta)$ 也来自自反性。因此上式简化为
\[
\begin{aligned}
&h^{n-2,1}\bigl(\operatorname{Gr}_{n-1}^W H^{n-1}(Z_f,\mathbb C)\bigr)\\
&\qquad=l(\Delta)-n-1-\sum_\Gamma l^*(\Gamma).
\end{aligned}
\]
余下只须证明，式~\eqref{eq:gysin-hodge-number} 右边第二项等于式~\eqref{eq:batyrev-hn21} 中的修正项。

理解它的一种方式如下。设 $v_i$ 位于面 $\widehat\Theta\subset\Delta^\circ$ 的相对内部，令 $\Theta\subset\Delta$ 表示对偶面。如上，这给出轨道闭包 $\mathbb P_\Theta\subset\mathbb P_\Delta$，且 $f(D_i)=\mathbb P_\Theta$。于是有映射 $V\cap D_i\to\overline V\cap\mathbb P_\Theta$，进一步诱导自然映射
\[
H^{n-3,0}(\overline V\cap\mathbb P_\Theta)
\longrightarrow H^{n-3,0}(V\cap D_i).
\]
大多数情形下，$h^{n-3,0}(\overline V\cap\mathbb P_\Theta)=0$；但当 $\Theta$ 单纯，且 $\overline V\cap\mathbb P_\Theta$ 的维数为 $n-3$ 时，\SourceCite[5.5]{DK} 给出公式
\begin{equation}\label{eq:face-holomorphic-forms}
h^{n-3,0}(\overline V\cap\mathbb P_\Theta)=l^*(\Theta).
\end{equation}
注意，$\dim(\overline V\cap\mathbb P_\Theta)=n-3$ 等价于 $\Theta$ 是余维数为二的面。因此，若 $\Theta\subset\Delta$ 是单纯的余维数为二的面，则 $\widehat\Theta$ 的 $l^*(\widehat\Theta)$ 个内部格点 $v_i$ 中的每一个，都贡献 $l^*(\Theta)$ 个全纯 $(n-3)$ 形式。这样，修正项 $\sum_\Theta l^*(\Theta)l^*(\widehat\Theta)$ 就自然出现了。遗憾的是，这个论证并不适用于一般情形。为实际计算式~\eqref{eq:gysin-hodge-number} 右边的和，需要将 $D_i$ 分解为轨道 $O$ 的不交并，并用
\SourcePage{83}
\SourceCite{DK} 的方法计算各 $V\cap O$ 的适当 $h^{p,q}$ 数。如上所述，这一论证的细节可在 \SourceCite{Batyrev4,Voisin3} 中找到。
\end{proof}

从历史上看，人们最初将大量 Calabi--Yau 簇族构造成加权射影空间中的拟光滑超曲面；最终找到了 $7555$ 个例子 \SourceCite{KS,KS2}。将这些例子的 Hodge 数绘成图，所得图形显示出醒目的\emph{部分对称性}，这是最早支持镜像对称的一些证据。\SourceCite{Batyrev4} 发表后，\SourceCite{CdK} 核查了表中全部 $7555$ 个簇：一般 Calabi--Yau 超曲面的 Newton 多胞形，经平移使唯一内部格点位于原点后，的确都是自反的。因此这些例子都是 Batyrev 镜像构造的特殊情形。\hyperref[subsec:nonreflexive-weights]{第 4.4.3 小节}将进一步介绍所用的方法。

定理~\ref{thm:batyrev-hodge} 给出的 Hodge 数相等关系，例如说明 $H^{1,1}(V)$ 同构于 $H^1(V^\circ,T_{V^\circ})\simeq H^{n-2,1}(V^\circ)$，但并未指定同构。根据镜像对称，联系 $V$ 的 Kähler 模与 $V^\circ$ 的复模的镜像映射，应至少定义在 $V$ 的 Kähler 模空间的一个开子集上；其导数应给出典范同构 $H^{1,1}(V)\simeq H^{n-2,1}(V^\circ)$。此外，在上述证明中还提到了单项式--除子镜像映射，即自然同构 $H_{\mathrm{toric}}^{1,1}(V)\simeq H_{\mathrm{poly}}^{n-2,1}(V^\circ)$。\hyperref[chap:moduli]{第~\ref*{chap:moduli} 章}将更详细地讨论此同构，并解释它与镜像映射的关系。\index{单项式--除子镜像映射|hyperpage}\index{Hodge数@Hodge 数|hyperpage}

还应提到，根据式~\eqref{eq:mirror-cohomology-forms}，镜像对称暗示更一般地，当 $V^\circ$ 是 $V$ 的 Batyrev 对偶时，应有 Hodge 数相等关系
\begin{equation}\label{eq:full-mirror-hodge}
h^{p,q}(V)=h^{n-1-p,q}(V^\circ).
\end{equation}
回顾，此时 $V$ 与 $V^\circ$ 的维数为 $d=n-1$。定理~\ref{thm:batyrev-hodge} 包含了 $(p,q)=(1,1)$ 与 $(n-2,1)$ 的情形，而 $q=0$ 的情形是平凡的。对三维簇，这些就是所需的全部情形；但在维数至少为四时，式~\eqref{eq:full-mirror-hodge} 是否成立仍是一个未解决的问题。

当 $V\subset X_\Sigma$ 是我们正在考虑的这类环面 Calabi--Yau 超曲面时，\SourceCite{BB2,BD} 定义了\emph{弦论 Hodge 数} $h_{\mathrm{st}}^{p,q}(V)$。当 $V$ 光滑或 $q=1$ 时，这些数与通常的 Hodge 数 $h^{p,q}(V)$ 相同。更重要的是，当 $V^\circ$ 是 $V$ 的 Batyrev 对偶时，对所有 $p+q=\dim V$，这些数满足
\[
h_{\mathrm{st}}^{p,q}(V)=h_{\mathrm{st}}^{n-1-p,q}(V^\circ).
\]
因此，在光滑情形且 $p+q=\dim V$ 时，式~\eqref{eq:full-mirror-hodge} 成立；不过，如前所述，若 $V$ 或 $V^\circ$ 是有奇点的极小 Calabi--Yau 簇，这个问题仍未解决。\index{弦论Hodge数@弦论 Hodge 数|hyperpage}

\subsection{多重镜像}\label{subsec:multiple-mirrors}

在 Batyrev 镜像构造中，$\Sigma$ 与 $\Sigma^\circ$ 的极大射影细分可能有不止一种选择。这似乎导致多个镜像族。例如，若 $\Sigma$ 与 $\Sigma'$ 都满足 $\Sigma(1)=\Sigma'(1)=\Delta^\circ\cap N-\{0\}$，便得到两个族 $V\subset X_\Sigma$ 与 $V'\subset X_{\Sigma'}$。$V^\circ\subset X_{\Sigma^\circ}$ 能否同时为这两个族的镜像？答案是肯定的，这是 \SourceCite{AGM2} 讨论的“多重镜像”现象的一部分。\hyperref[chap:moduli]{第~\ref*{chap:moduli} 章}将详细讨论多重镜像。不过，若限于式~\eqref{eq:toric-hodge-subspace} 中的环面 Kähler 模，大致的想法是：$\Sigma$ 与 $\Sigma'$ 对应于 GKZ 分解中不同的锥，而族 $V$ 与 $V'$ 的 Kähler 模给出\emph{同一个}整体 Kähler 模空间中互不相交的开集。\index{多重镜像|hyperpage}

\SourcePage{84}
也可以在另一侧作不同选择。此时，极大射影细分 $\Sigma^\circ$ 与 $\Sigma'^\circ$ 给出族 $V^\circ$ 与 $V'^\circ$。我们希望这些族的复模对应于族 $V$ 的 Kähler 模。因此 $V^\circ$ 与 $V'^\circ$ 应有\emph{相同的}复模。若限于式~\eqref{eq:polynomial-hodge-subspace} 中的多项式形变，这一点来自定理~\ref{thm:batyrev-hodge} 的证明。具体说，$V^\circ$ 的多项式模对应于 $L(\Delta^\circ)$ 的一个开子集模去 $X_{\Sigma^\circ}$ 的自同构群的作用。不过，$V'^\circ$ 也由同一空间 $L(\Delta^\circ)$ 中的多项式定义；利用 \SourceCite{Cox} 的结果，可以看到 $\operatorname{Aut}(X_{\Sigma^\circ})$ 与 $\operatorname{Aut}(X_{\Sigma'^\circ})$ 有同样的单位元连通分支。因此，至少在局部，这些多项式模空间看起来相同。\hyperref[chap:moduli]{第~\ref*{chap:moduli} 章}将仔细解释这些内容。

\section{重访五次三维簇}\label{sec:quintic-revisited}

现在回到\hyperref[chap:quintic]{第 2 章}讨论的五次三维簇 $V\subset\mathbb P^4$。由于 $V$ 是 Fano 环面簇 $\mathbb P^4$ 中的反典范超曲面，可用 Batyrev 构造描述它的镜像族。

$\mathbb P^4$ 上的反典范类 $D_0+\cdots+D_4$ 对应于\IndexLocation{indexloc:1177:84}\IndexLocation{indexloc:1223:84}自反多胞形 $\Delta\subset M_{\mathbb R}\simeq\mathbb R^4$，其顶点为
\begin{equation}\label{eq:quintic-polytope-vertices}
\begin{gathered}
(-1,-1,-1,-1),\\
(4,-1,-1,-1),\\
(-1,4,-1,-1),\\
(-1,-1,4,-1),\\
(-1,-1,-1,4).
\end{gathered}
\end{equation}
可以把 $\Delta$ 视为 $\mathbb P^4$ 中五次多项式的 Newton 多胞形。为看出这一点，考虑如式~\eqref{eq:laurent-space} 那样由 $\Delta$ 确定的 Laurent 多项式，例如
\[
f=\sum_{i=1}^s a_it^{m_i},\qquad
\Delta\cap M=\{m_1,\ldots,m_s\}.
\]
再注意，由式~\eqref{eq:toric-graded-sections} 给出的齐次化
\[
F(x_0,x_1,x_2,x_3,x_4)=
x_0x_1x_2x_3x_4\,f(x_1/x_0,x_2/x_0,x_3/x_0,x_4/x_0)
\]
就是 $x_0,x_1,x_2,x_3,x_4$ 的任意五次多项式。

为寻找 $V$ 的镜像族，使用对偶多胞形 $\Delta^\circ\subset N_{\mathbb R}$。容易看到，它的顶点为
\begin{equation}\label{eq:quintic-polar-vertices}
\begin{gathered}
(-1,-1,-1,-1),\\
(1,0,0,0),\\
(0,1,0,0),\\
(0,0,1,0),\\
(0,0,0,1).
\end{gathered}
\end{equation}
环面簇 $\mathbb P_{\Delta^\circ}$ 由 $\Delta^\circ$ 在 $M_{\mathbb R}$ 中的法扇确定，该扇由向 $\Delta$ 的各面作出的锥组成。因此，一维锥生成元是式~\eqref{eq:quintic-polytope-vertices} 中的向量，它们生成 $M$ 中指数为 $125$ 的子格 $M_1$。商 $M/M_1$ 为
\[
G=\left\{(a_0,a_1,a_2,a_3,a_4)\in(\mathbb Z_5)^5:
\sum_{i=0}^4a_i=0\pmod5\right\}/\mathbb Z_5,
\]
其中 $\mathbb Z_5\subset(\mathbb Z_5)^5$ 是对角子群。由于 $M_{\mathbb R}=(M_1)_{\mathbb R}$，可将 $\Delta^\circ$ 的法扇视为 $(M_1)_{\mathbb R}$ 中的扇。其生成元由 $M_1$ 的一组 $\mathbb Z$ 基以及这些基向量之和的相反数组成。这说明，使用格 $M_1$ 时，该法扇就是 $\mathbb P^4$ 的标准扇。若再换到更大的
\SourcePage{85}
格 $M$，环面几何的标准结果表明，$\mathbb P_{\Delta^\circ}$ 是 $\mathbb P^4$ 模去 $M/M_1\simeq G$ 的商，其中 $G$ 在 $\mathbb P^4$ 上的作用与\hyperref[chap:quintic]{第 2 章}所述完全相同。

$\mathbb P_{\Delta^\circ}$ 的齐次坐标环 $S$ 可描述如下。它是多项式环 $S=\mathbb C[x_0,x_1,x_2,x_3,x_4]$，由 Chow 群 $A_3(\mathbb P_{\Delta^\circ})$ 分次。由式~\eqref{eq:toric-chow-sequence} 得到正合列
\[
0\longrightarrow\mathbb Z^4\longrightarrow\mathbb Z^5
\longrightarrow\mathbb Z\oplus G\longrightarrow0,
\]
其中第一个映射的矩阵以式~\eqref{eq:quintic-polytope-vertices} 的各行为行，第二个映射将 $(a_0,a_1,a_2,a_3,a_4)\in\mathbb Z^5$ 送到
\[
\begin{aligned}
\bigl[a_0,a_1,a_2,a_3,a_4\bigr]
={}&\left(\sum_{i=0}^4a_i,
(-a_1-a_2-a_3-a_4,a_1,a_2,a_3,a_4)\right)\!\in\mathbb Z\oplus G.
\end{aligned}
\]
因此 $A_3(\mathbb P_{\Delta^\circ})\simeq\mathbb Z\oplus G$，而 $S$ 的分次规定单项式 $x_0^{a_0}\cdots x_4^{a_4}$ 的次数为 $[a_0,a_1,a_2,a_3,a_4]\in\mathbb Z\oplus G$。

于是，$\mathbb P_{\Delta^\circ}$ 的反典范类的次数为 $[1,1,1,1,1]=(5,0)\in\mathbb Z\oplus G$，$S$ 中仅有的这个次数的单项式是 $x_i^5$ 及 $x_0x_1x_2x_3x_4$。因此得到超曲面 $\overline V^\circ\subset\mathbb P_{\Delta^\circ}$，由方程
\begin{equation}\label{eq:quintic-general-mirror}
a_0x_0^5+a_1x_1^5+a_2x_2^5+a_3x_3^5+a_4x_4^5
+a_5x_0x_1x_2x_3x_4=0
\end{equation}
定义。具体说，$\overline V^\circ$ 就是 $\mathbb P^4$ 中由上述方程定义的超曲面模去 $G$ 的商。注意，式~\eqref{eq:quintic-general-mirror} 中的单项式对应于 $\Delta^\circ$ 中的整点，即原点与式~\eqref{eq:quintic-polar-vertices} 中的各点。此外，关系
\[
(x_0^5)(x_1^5)(x_2^5)(x_3^5)(x_4^5)
=(x_0x_1x_2x_3x_4)^5
\]
反映了式~\eqref{eq:quintic-polar-vertices} 中五个向量的和等于原点的五倍。

方程~\eqref{eq:quintic-general-mirror} 不同于\hyperref[chap:quintic]{第 2 章}讨论的单参数簇族。不过，回顾我们关心的是这些超曲面的模。$\mathbb P_{\Delta^\circ}$ 的\IndexLocation{indexloc:1355:85}自同构群是环面 $T_M$；利用这些自同构，可以约化到\hyperref[chap:quintic]{第 2 章}考虑的单参数族
\[
x_0^5+x_1^5+x_2^5+x_3^5+x_4^5+\psi x_0x_1x_2x_3x_4=0.
\]
\hyperref[chap:moduli]{第~\ref*{chap:moduli} 章}将给出完成此约化的系统方法。

最后，为得到五次超曲面 $V$ 的镜像，在 $M_{\mathbb R}$ 中选取\IndexLocation{indexloc:1096:85}极大射影细分 $\Sigma^\circ$。这样的扇细化 $\Delta^\circ$ 的法扇，并满足 $\Sigma^\circ(1)=\Delta\cap M-\{0\}$。于是镜像族由超曲面 $V^\circ\subset X_{\Sigma^\circ}$ 给出，它们是 $\overline V^\circ\subset\mathbb P_{\Delta^\circ}$ 的严格变换。

注意，$M\cap\Delta-\{0\}$ 除顶点外还有大量点，因此 $\Sigma^\circ$ 有许多选择。找出所有这样的 $\Sigma^\circ$ 这一计算问题，远远超出了当时计算机的能力。\SourceCite{Morrison2} 的附录描述了 $\Sigma^\circ$ 的一个具体选择。

\section{环面完全交}\label{sec:toric-complete-intersections}

推广\hyperref[sec:batyrev-mirror]{第 4.1 节}研究的环面超曲面，也可以问环面簇中的 Calabi--Yau 完全交是否具有镜像对称。本节通过介绍 Batyrev 与 Borisov \SourceCite{Borisov1,BB1} 的两种镜像对称构造来回答这个问题；这些构造涉及 nef 分拆与自反 Gorenstein 锥。

\SourcePage{86}
\subsection{nef 分拆}\label{subsec:nef-partitions}

设 $X=\mathbb P_\Delta$ 是自反多胞形 $\Delta$ 对应的 $n$ 维环面簇。令 $D_\rho$ 表示 $\rho\in\Sigma(1)$ 对应的环面除子，其中 $\Sigma$ 是 $\Delta$ 的法扇。给定将 $\Sigma(1)=I_1\cup\cdots\cup I_k$ 分成 $k$ 个互不相交的子集的分拆，得到除子 $E_j=\sum_{\rho\in I_j}D_\rho$，满足 $-K_X=E_1+\cdots+E_k$。
\begin{definition}\label{def:nef-partition}
若对每个 $j$，$E_j$ 都是由整体截面生成的 Cartier 除子，则称分解 $\Sigma(1)=I_1\cup\cdots\cup I_k$ 为 \emph{nef 分拆}。
\end{definition}

若 $\Delta_j$ 是 $E_j$ 对应的多胞形，则 nef 分拆意味着
\[
\Delta=\Delta_1+\cdots+\Delta_k.
\]
不严格地说，分解 $-K_X=E_1+\cdots+E_k$ 与 $\Delta=\Delta_1+\cdots+\Delta_k$ 也称为 nef 分拆。

$\mathcal O(E_1),\ldots,\mathcal O(E_k)$ 的一般截面定义一个维数为 $d=n-k$ 的完全交 $\overline V\subset X$。利用 $-K_X=E_1+\cdots+E_k$ 与伴随公式，看到 $\overline V\subset\mathbb P_\Delta$ 是 Calabi--Yau 簇。如\hyperref[sec:batyrev-mirror]{第 4.1 节}，$\Delta$ 的法扇的极大射影细分 $\Sigma'$ 给出映射 $X_{\Sigma'}\to\mathbb P_\Delta$，而 $\overline V\subset\mathbb P_\Delta$ 的严格变换 $V\subset X_{\Sigma'}$ 仍是完全交，并且成为极小 Calabi--Yau 簇，见 \SourceCite[命题 4.15]{BB1}。

为描述 $V$ 的镜像族，仍将 $\Sigma(1)$ 中的元素视为一维锥生成元。在自反情形，它们恰好是极多胞形 $\Delta^\circ\subset N_{\mathbb R}$ 的顶点。于是 nef 分拆 $\Sigma(1)=I_1\cup\cdots\cup I_k$ 给出多胞形
\[
\nabla_j=\operatorname{Conv}(\{0\}\cup I_j)\subset N_{\mathbb R}.
\]
立即得到
\[
\Delta^\circ=\operatorname{Conv}(\nabla_1\cup\cdots\cup\nabla_k).
\]
\SourceCite{Borisov1} 的基本结果是：Minkowski 和
\[
\nabla=\nabla_1+\cdots+\nabla_k
\]
仍是自反多胞形，而且 $\nabla=\nabla_1+\cdots+\nabla_k$ 也是 nef 分拆。按 \SourceCite{BB1} 的术语，这称为\emph{对偶 nef 分拆}。

使用对偶 nef 分拆，得到 Calabi--Yau 完全交 $\overline V^\circ\subset\mathbb P_\nabla$；选取 $\nabla$ 的法扇的极大射影细分 $\Sigma^\circ$，得到极小 Calabi--Yau 完全交族 $V^\circ\subset X_{\Sigma^\circ}$。猜想为
\[
V^\circ\text{ 是 }V\text{ 的镜像族。}
\]
这个猜想与 \SourceCite{LTe} 提出的射影空间完全交的镜像对称构造一致。

作为支持此猜想的证据，Batyrev 与 Borisov 证明
\begin{equation}\label{eq:complete-intersection-hodge}
h^{p,1}(V)=h^{d-p,1}(V^\circ),
\end{equation}
其中 $d=n-k$ 是 $V$ 与 $V^\circ$ 的维数。更一般地，使用\hyperref[subsec:batyrev-hodge]{第 4.1.3 小节}讨论的弦论 Hodge 数 $h_{\mathrm{st}}^{p,q}$，可以证明
\begin{equation}\label{eq:complete-intersection-string-hodge}
h_{\mathrm{st}}^{p,q}(V)=h_{\mathrm{st}}^{d-p,q}(V^\circ),
\end{equation}
见 \SourceCite{BB2}。由于 $h^{p,1}(V)=h_{\mathrm{st}}^{p,1}(V)$，可见式~\eqref{eq:complete-intersection-hodge} 是式~\eqref{eq:complete-intersection-string-hodge} 的特殊情形。\index{nef分拆@nef 分拆|hyperpage}\index{完全交|hyperpage}

\SourcePage{87}
\subsection{Gorenstein 锥}\label{subsec:gorenstein-cones}

研究 Gorenstein 锥，可推广 nef 分拆的组合对偶性。其定义如下。设有格 $\overline M$ 及其对偶 $\overline N$。
\begin{definition}\label{def:gorenstein-cone}
设 $\sigma\subset\overline M_{\mathbb R}$ 是有理多面锥。
\begin{enumerate}[label=(\roman*)]
\item 若 $\sigma$ 强凸，且存在 $n_\sigma\in\overline N$，使 $\sigma$ 的每个本原整生成元 $v$ 都满足 $\langle v,n_\sigma\rangle=1$，则称 $\sigma$ 为 \emph{Gorenstein 锥}。
\item 若 $\sigma$ 及其对偶 $\check\sigma$ 都是 Gorenstein 锥，则称 $\sigma$ 为\emph{自反 Gorenstein 锥}。
\end{enumerate}
\end{definition}
显然，$\sigma$ 是自反 Gorenstein 锥，当且仅当 $\check\sigma$ 是。此外，若 $\sigma$ 自反，则它必有满维数，这意味着 $n_\sigma\in\overline N$ 唯一确定。$\check\sigma$ 相应地确定 $m_\sigma\in\overline M$，整数
\[
r=\langle m_\sigma,n_\sigma\rangle
\]
称为 $\sigma$，或 $\check\sigma$，的\emph{指数}。

先解释自反 Gorenstein 锥与自反多胞形对偶性的关系。给定以原点为内点的整多胞形 $\Delta\subset M_{\mathbb R}$，将格 $M$ 扩充为 $\overline M=M\times\mathbb Z$，相应地有 $\overline N=N\times\mathbb Z$。令 $\sigma\subset\overline M_{\mathbb R}$ 为由元素 $\{(v,1):v\in\Delta\}$ 张成的锥。换言之，
\begin{equation}\label{eq:polytope-gorenstein-cone}
\sigma=\{(\lambda v,\lambda):v\in\Delta,\ \lambda\geq0\}.
\end{equation}
注意，$\sigma$ 是 Gorenstein 锥。简单论证表明，$\sigma$ 自反，当且仅当 $\Delta$ 自反。此外，$\Delta$ 自反时，指数 $r=1$，并且 $\check\sigma$ 由 $\Delta^\circ$ 构造的方式，与 $\sigma$ 由 $\Delta$ 构造的方式相同。

也可用自反 Gorenstein 锥描述 nef 分拆的对偶性。给定自反多胞形 $\Delta\subset M_{\mathbb R}$ 及 nef 分拆 $\Delta=\Delta_1+\cdots+\Delta_k$，定义格 $\overline M=M\times\mathbb Z^k$，其对偶格为 $\overline N=N\times\mathbb Z^k$。令 $\sigma$ 为锥
\begin{equation}\label{eq:nef-gorenstein-cone}
\sigma=\left\{\left(\sum_{i=1}^k\lambda_iv_i,
\lambda_1,\ldots,\lambda_k\right):
v_i\in\Delta_i,\ \lambda_i\geq0\right\}.
\end{equation}
注意，它推广了式~\eqref{eq:polytope-gorenstein-cone}。

\begin{proposition}\label{prop:nef-cone-duality}
$\sigma$ 是指数为 $k$ 的自反 Gorenstein 锥。它的对偶锥 $\check\sigma$，由对偶 nef 分拆 $\nabla=\nabla_1+\cdots+\nabla_k$ 应用式~\eqref{eq:nef-gorenstein-cone} 构造得到。
\end{proposition}
\begin{proof}
证明见 \SourceCite{BB1}。
\end{proof}

当 $\sigma$ 如式~\eqref{eq:nef-gorenstein-cone} 定义时，对偶锥 $\check\sigma$ 具有有趣的几何意义。这个锥是某个扇 $\overline\Sigma$ 的支集，其对应环面簇为 $X=\mathbb P_\Delta$ 上向量丛
\begin{equation}\label{eq:nef-vector-bundle}
\mathcal O(-E_1)\oplus\cdots\oplus\mathcal O(-E_k)
\end{equation}
的总空间。环面簇 $X$ 可视为经零截面嵌入 $X_{\overline\Sigma}$，因此 $V$ 可以构造成 $X_{\overline\Sigma}$ 的子簇。$\overline\Sigma$ 也可扩充为完备扇，其对应环面簇是 $X$ 上的射影丛
\[
\mathbb P\bigl(\mathcal O\oplus\mathcal O(-E_1)\oplus\cdots\oplus\mathcal O(-E_k)\bigr).
\]
给出完全交 $V$ 的 $\mathcal O(E_1),\ldots,\mathcal O(E_k)$ 的各截面，也确定了向量丛 $\mathcal O(E_1)\oplus\cdots\oplus\mathcal O(E_k)$ 的一个整体截面。相应的超曲面为
\begin{equation}\label{eq:cayley-hypersurface}
Y'\subset\mathbb P\bigl(\mathcal O(E_1)\oplus\cdots\oplus\mathcal O(E_k)\bigr)
\end{equation}
这个超曲面与 $V$ 密切相关。这就是 \SourceCite{BB3,GKZ2} 描述的 \emph{Cayley 技巧}。\index{Gorenstein锥@Gorenstein 锥|hyperpage}\index{Cayley技巧@Cayley 技巧|hyperpage}

\SourcePage{88}
式~\eqref{eq:cayley-hypersurface} 的超曲面 $Y'$ 不是 Calabi--Yau 簇，但它是 Schimmrigk 意义下的\emph{广义 Calabi--Yau 簇}，见 \SourceCite{BB3,Schimmrigk}。可用任意自反 Gorenstein 锥 $\sigma$ 如下构造广义 Calabi--Yau 簇。由 $\sigma$ 得到半群 $\sigma\cap\overline M$，并将 $m\in\sigma\cap\overline M$ 的次数定义为 $\langle m,n_\sigma\rangle$。这使 $\mathbb C[\sigma\cap\overline M]$ 成为分次环，给出环面簇
\[
\mathbb P_\sigma=\operatorname{Proj}\bigl(\mathbb C[\sigma\cap\overline M]\bigr).
\]
若 $\mathcal O_{\mathbb P_\sigma}(1)$ 为 $\mathbb P_\sigma$ 上的重言层，则它的一个整体截面定义子簇
\[
Y\subset\mathbb P_\sigma,
\]
Batyrev 与 Borisov 称其为\emph{与 $\sigma$ 对应的广义 Calabi--Yau 簇}。当然，对偶锥 $\check\sigma$ 也是自反 Gorenstein 锥，因此得到对偶族
\[
Y^\circ\subset\mathbb P_{\check\sigma}.
\]
猜想它在适当意义下是 $Y\subset\mathbb P_\sigma$ 的镜像，见 \SourceCite[猜想 2.17]{BB3}。

这个构造包含了目前考虑过的 Calabi--Yau 超曲面与完全交。更精确地说：
\begin{itemize}
\item 当 $\sigma$ 由式~\eqref{eq:polytope-gorenstein-cone} 给出时，$\mathbb P_\sigma$ 恰好是环面簇 $\mathbb P_\Delta$；而 $\mathbb C[\sigma\cap\overline M]$ 就是\hyperref[subsec:toric-from-polytopes]{第 3.2.2 小节}的环 $S_\Delta$。超曲面 $Y\subset\mathbb P_\sigma$ 恰好是\hyperref[sec:batyrev-mirror]{第 4.1 节}研究的 Calabi--Yau 超曲面 $V\subset\mathbb P_\Delta$。
\item 当 $\sigma$ 由式~\eqref{eq:nef-gorenstein-cone} 给出时，存在双有理保典范态射
\[
\mathbb P\bigl(\mathcal O(E_1)\oplus\cdots\oplus\mathcal O(E_k)\bigr)
\longrightarrow\mathbb P_\sigma,
\]
使 $Y\subset\mathbb P_\sigma$ 是 $Y'\subset\mathbb P(\mathcal O(E_1)\oplus\cdots\oplus\mathcal O(E_k))$ 的像。
\end{itemize}
此外，在这两种情形中，镜像对称都化为 Gorenstein 锥的对偶性。

{\setlength{\emergencystretch}{1em}
除了这些已知例子，Gorenstein 锥构造也有助于解释通常镜像对称表述失效的一些情形。\emph{刚性 Calabi--Yau 三维簇}的情形最为明显。它们没有非平凡形变，所以 $H^{2,1}(V)=0$。通常意义下的镜像 $V^\circ$ 应有 $H^{1,1}(V^\circ)=0$，但这对 Kähler 流形是不可能的。然而，在某些情形中，人们仍找到了“镜像”。例如，设 $E$ 为唯一具有阶为 $3$ 的自同构的椭圆曲线，令 $V=E^3/G$，其中 $G$ 由该自同构的对角作用生成。可以证明 $V$ 是一个有奇点的刚性 Calabi--Yau 簇；\SourceCite{CDP} 算出的镜像是 $\mathbb P^8$ 中七维 Fermat 超曲面的一个有限商。\SourceCite{BB3} 表明，这与 Gorenstein 锥的对偶性一致。\par}\index{刚性Calabi--Yau三维簇@刚性 Calabi--Yau 三维簇|hyperpage}

\section{Voisin--Borcea 构造}\label{sec:voisin-borcea}

本节描述 Voisin 与 Borcea 各自独立给出的另一种镜像族构造。其灵感来自 Nikulin 关于带对合的 K3 曲面的工作。他找到了 $75$ 个这样的曲面族，其不变量显示出醒目的部分对称性 \SourceCite{Nikulin}。利用这些曲面，Voisin 与 Borcea 构造了猜想具有镜像对称性的 Calabi--Yau 三维簇族。我们将勾勒一些证明，完整细节请参见原始论文 \SourceCite{Voisin1,Borcea2}。\index{Voisin--Borcea构造@Voisin--Borcea 构造|hyperpage}

\SourcePage{89}
\subsection{由 K3 曲面与椭圆曲线构造三维簇}\label{subsec:voisin-threefolds}

首先利用某些带对合的 K3 曲面与椭圆曲线构造 Calabi--Yau 三维簇。更精确地说，考虑以下资料：
\begin{itemize}
\item 一个 K3 曲面 $S$ 与对合 $i:S\to S$，它在 $H^{2,0}(S)$ 上作用为 $-1$。$i$ 的不动点集是光滑曲线 $C_1,\ldots,C_N$ 的不交并。我们明确假设 $N>0$。
\item 一条椭圆曲线 $E$ 与对合 $j:E\to E$，其商为 $\mathbb P^1$。将 $j$ 的四个不动点记为 $p_1,p_2,p_3,p_4$。
\end{itemize}
$S$ 的对合 $i$ 与 $E$ 的对合 $j$ 诱导 $E\times S$ 的对合 $k$，定义为 $k(e,s)=(j(e),i(s))$，其中 $(e,s)\in E\times S$。$k$ 的不动点集由 $4N$ 条曲线 $C_{r,s}=\{p_r\}\times C_s$ 组成。沿这 $4N$ 条曲线爆破 $E\times S$，得到三维簇 $\widetilde{E\times S}$，其光滑例外除子为 $D=\bigcup D_{r,s}$；对合 $k$ 自然提升为 $\widetilde{E\times S}$ 的对合 $\widetilde k$。置
\begin{equation}\label{eq:voisin-threefold}
V=\widetilde{E\times S}/\widetilde k.
\end{equation}
可将 $D=\bigcup D_{r,s}$ 视为 $V$ 上具有 $4N$ 个不可约分支的除子。于是可以证明以下结论。
\begin{lemma}\label{lem:voisin-cy}
$V$ 是光滑 Calabi--Yau 三维簇。
\end{lemma}

令 $g_s$ 表示曲线 $C_s\subset S$ 的亏格，置 $N'=\sum_{s=1}^N g_s$。$V$ 的 Hodge 分量计算如下。
\begin{theorem}\label{thm:voisin-hodge}
$h^{1,1}(V)=11+5N-N'$，且 $h^{2,1}(V)=11+5N'-N$。
\end{theorem}
\begin{note}
由这些公式可见，若同样构造一个 $V^\circ$，但交换 $N$ 与 $N'$ 的角色，则它便是镜像流形的一个候选者。这正是 Voisin 与 Borcea 所做的。
\end{note}
\begin{proof}
先给出 $H^{1,1}(V)$ 的显式分解。首先，除子 $D=\bigcup D_{r,s}$ 的各分支在 $H^{1,1}(V)$ 中给出 $4N$ 个类。其次，光滑曲面 $T=S/i$ 给出一些类。将每条曲线 $C_s\subset S$ 与其在 $T$ 中的像等同，通过计算 Euler 示性数来求 $h^{1,1}(T)$：
\begin{equation}\label{eq:voisin-euler-quotient}
\begin{aligned}
e(T)&=e\left(T-\bigcup_s C_s\right)+\sum_s e(C_s)\\
&=\tfrac12e\left(S-\bigcup_s C_s\right)+\sum_s e(C_s)\\
&=\tfrac12\left(24-\sum_s(2-2g_s)\right)+\sum_s(2-2g_s)\\
&=12+N-N'.
\end{aligned}
\end{equation}
以 $i^*$ 表示 $i$ 诱导的 $S$ 的上同调上的作用。考察 $S$ 的上同调中 $i^*$ 不变的部分，立即可得 $T$ 的上同调满足 $h^0(T)=h^4(T)=1$，以及 $h^1(T)=h^3(T)=h^{2,0}(T)=h^{0,2}(T)=0$。结合式~\eqref{eq:voisin-euler-quotient} 得到
\begin{equation}\label{eq:voisin-quotient-h11}
h^{1,1}(T)=10+N-N'.
\end{equation}
$H^{1,1}(T)$ 中的上同调类可经自然投影 $\pi_1:V\to T$ 拉回到 $V$；将所得 $H^{1,1}(V)$ 中的元素与 $H^{1,1}(T)$ 等同。

\SourcePage{90}
最后，$H^{1,1}(E)$ 中的元素经自然投影 $\widetilde{E\times S}\to E$，拉回为 $H^{1,1}(\widetilde{E\times S})$ 中的 $\widetilde k$ 不变上同调类，因此也可与 $V$ 上的上同调类等同。

刚才描述的三类上同调类在 $\mathbb C$ 上线性无关，并张成 $H^{1,1}(V)$。因此得到直和分解
\begin{equation}\label{eq:voisin-h11-decomposition}
H^{1,1}(V)=\bigoplus_{r,s}\mathbb C\cdot[D_{r,s}]
\oplus H^{1,1}(T)\oplus H^{1,1}(E).
\end{equation}
此分解意味着 $h^{1,1}(V)=4N+(10+N-N')+1=11+5N-N'$，证明了所需的 $h^{1,1}(V)$ 公式。

对定理的第二个断言，将 $H^{2,1}(V)$ 与 $H^{2,1}(\widetilde{E\times S})$ 的 $\widetilde k$ 不变部分等同，再次列出三类上同调类。首先，由构造，$D_{r,s}$ 是 $C_{r,s}$ 上的射影丛，因此给出 $\widetilde{E\times S}$ 中以 $C_{r,s}$ 为参数的一族曲线。有相应的 Abel--Jacobi 映射 $H^{1,0}(C_{r,s})\to H^{2,1}(\widetilde{E\times S})$，见 \SourceCite{CG}；将其像与 $H^{1,0}(C_{r,s})$ 等同。其次，令 $H^{1,1}(S)^-$ 表示 $i^*$ 在 $H^{1,1}(S)$ 上作用的 $-1$ 特征子空间。由于 $j$ 在 $H^1(E)$ 上诱导 $-1$，而 $i$ 在 $H^{2,0}(S)$ 上作用为 $-1$，由 Künneth 分解，还得到来自 $H^{1,1}(S)^-\otimes H^{1,0}(E)$ 与 $H^{2,0}(S)\otimes H^{0,1}(E)$ 的类。刚才描述的三类仍在 $\mathbb C$ 上线性无关，并张成 $H^{2,1}(V)$。注意 $H^{1,0}(E)$ 与 $H^{0,1}(E)$ 都是一维的，因此可以非典范地写为
\begin{equation}\label{eq:voisin-h21-decomposition}
H^{2,1}(V)\simeq\bigoplus_{r,s}H^{1,0}(C_{r,s})
\oplus H^{1,1}(S)^-\oplus H^{2,0}(S).
\end{equation}
所以 $h^{2,1}(V)=4N'+h^{1,1}(S)^-+1$。不过，若 $H^{1,1}(S)^+$ 表示 $i^*$ 在 $H^{1,1}(S)$ 上作用的 $+1$ 特征子空间，则式~\eqref{eq:voisin-quotient-h11} 意味着
\begin{equation}\label{eq:voisin-eigenspace-h11}
h^{1,1}(S)^+=10+N-N',\qquad
h^{1,1}(S)^-=10+N'-N.
\end{equation}
所需的 $h^{2,1}(V)$ 公式立即随之得到。
\end{proof}

后面还需要对合 $i$ 的如下不变量。交叉形式 $\langle\cdot,\cdot\rangle$ 在 $H^2(S,\mathbb Z)$ 上是幺模的，但在 $H^2(S,\mathbb Z)^+$ 或 $H^2(S,\mathbb Z)^-$ 上未必如此。不过，利用 Nikulin 关于 K3 格的对合的结果 \SourceCite{AN,Dolgachev2,Nikulin}，可以证明存在整数 $a$，使
\begin{equation}\label{eq:voisin-discriminant}
\begin{aligned}
2^a&=\langle\cdot,\cdot\rangle\text{ 在 }H^2(S,\mathbb Z)^+
\text{ 上的判别式绝对值}\\
&=\langle\cdot,\cdot\rangle\text{ 在 }H^2(S,\mathbb Z)^-
\text{ 上的判别式绝对值}.
\end{aligned}
\end{equation}
$a$ 有如下显式公式。
\begin{lemma}\label{lem:voisin-discriminant}
$a=12-N-N'$。
\end{lemma}
\begin{proof}
关键步骤是构造正合列
\[
0\longrightarrow H^2(T,\mathbb Z)
\longrightarrow H^2(S,\mathbb Z)^+
\longrightarrow(\mathbb Z_2)^{N-1}\longrightarrow0,
\]
其中嵌入由拉回诱导。这里省略构造。$H^2(S,\mathbb Z)^+$ 上的形式 $\langle\cdot,\cdot\rangle$，限制到 $H^2(T,\mathbb Z)$ 后为其幺模形式的两倍，因此该正合列意味着判别式绝对值为 $2^{b_2(T)-2(N-1)}$。现在由式~\eqref{eq:voisin-quotient-h11} 即得引理。
\end{proof}

\SourcePage{91}
\subsection{Voisin--Borcea 镜像}\label{subsec:voisin-mirror}

现在描述 Calabi--Yau 三维簇 $V=\widetilde{E\times S}/\widetilde k$ 的镜像族 $V^\circ$。如\hyperref[subsec:voisin-threefolds]{第 4.4.1 小节}已指出，定理~\ref{thm:voisin-hodge} 表明，若存在同类型的镜像族，$N$ 与 $N'$ 的角色就必定互换。构造族 $V^\circ$ 需要以下技术性引理。
\begin{lemma}\label{lem:voisin-hyperbolic-plane}
若 $N'>0$ 且 $(N,N')\ne(5,1)$，则 $H^2(S,\mathbb Z)^-$ 含有一个双曲平面。
\end{lemma}
\SourceCite{Voisin1} 中的陈述实际上稍强一些。证明使用前面提到的 Nikulin 关于 K3 格的对合的结果。

假设 $N'>0$ 且 $(N,N')\ne(5,1)$，选择双曲平面 $P\subset H^2(S,\mathbb Z)^-$，以 $r_P$ 表示关于 $P$ 的反射。也就是说，$P$ 是 $r_P$ 的 $+1$ 特征子空间，$P^\perp$ 是其 $-1$ 特征子空间。置 $i_1^*=r_P\cdot i^*$。仍利用 Nikulin 关于 K3 格的对合的结果，可以证明 $i_1^*$ 是一个对合，它来自某个 K3 曲面 $S_1$ 上的对合 $i_1$。此外，简单计算表明
\begin{equation}\label{eq:voisin-mirror-lattice}
H^2(S_1,\mathbb Z)^-=P\oplus H^2(S,\mathbb Z)^+.
\end{equation}

现在由 $(S_1,i_1)$ 与带对合 $j_1$ 的椭圆曲线 $E_1$，按构造 $V$ 的相同方法构造 Calabi--Yau 三维簇 $V^\circ$。如 \SourceCite{Voisin1} 所述，$(E_1,j_1)$ 有一个自然选择。设 $V^\circ$ 的不变量为 $(N_1,N_1')$，则有以下结果。
\begin{proposition}\label{prop:voisin-mirror-invariants}
$N_1=N'$，$N_1'=N$。
\end{proposition}
\begin{proof}
设 $H^2(S_1,\mathbb Z)^-$ 上交叉形式的判别式绝对值为 $2^{a_1}$。因为 $P\subset H^2(S,\mathbb Z)$ 幺模，式~\eqref{eq:voisin-discriminant}、\eqref{eq:voisin-mirror-lattice} 意味着 $a_1=a$，所以由引理~\ref{lem:voisin-discriminant}，
\[
12-N_1-N_1'=12-N-N'.
\]
式~\eqref{eq:voisin-mirror-lattice} 还意味着 $h^{1,1}(S_1)^-=h^{1,1}(S)^+$。结合式~\eqref{eq:voisin-eigenspace-h11} 得到
\[
10+N_1'-N_1=10+N-N'.
\]
立即得到所需等式 $N_1=N'$ 与 $N_1'=N$。
\end{proof}

猜想所得 Calabi--Yau 三维簇族 $V^\circ$ 是 $V$ 的镜像族。作为支持 $V$ 与 $V^\circ$ 的镜像关系的证据，由定理~\ref{thm:voisin-hodge} 与命题~\ref{prop:voisin-mirror-invariants} 得到 $h^{1,1}(V)=h^{2,1}(V^\circ)$ 及 $h^{2,1}(V)=h^{1,1}(V^\circ)$。此外，Voisin 与 Borcea 还给出了其他支持镜像关系的证据。

第一项证据来自 Voisin，涉及 $V$ 与 $V^\circ$ 的复模和 Kähler 模之间的关系。$V$ 的复模可描述如下。先考虑来自 K3 曲面 $S$ 的模，考察区域
\[
D=\{\omega\in\mathbb P(H^2(S,\mathbb C)^-):
\omega\cdot\omega=0,\ \omega\cdot\overline\omega>0\}.
\]
它几乎是具有满足 $i'^*=i^*$ 的对合 $i'$ 的 K3 曲面的周期域。更精确地说，一般点 $\omega\in D$ 对应于带对合 $i'$ 的 K3 曲面 $S'$，满足 $i'^*=i^*$ 及 $\mathbb C\cdot\omega=H^{2,0}(S')$；有序对 $(S,i)$ 的模空间是某个明确的第二纲子集 $U\subset D$ 的离散商。固定椭圆曲线 $(E,j)$，便得到 $(S,i)$ 的模空间到 $V$ 的复模空间中的嵌入。因此，可将 $D$ 基本上视为比 $V$ 的复模空间某个叶状结构的一片叶稍大的空间。\index{对合|hyperpage}\index{双曲平面|hyperpage}

\SourcePage{92}
转向 $V$ 的 Kähler 模，定义区域
\[
D'=\{\alpha\in H^2(S,\mathbb C)^+:(\operatorname{Re}\alpha)^2>0\}.
\]
注意，$D'$ 可与 $H^2(T,\mathbb C)$ 的一个子集等同；粗略地说，它比 $T$ 的 Kähler 模空间稍大。此外，经 $\pi_1:V\to T$ 拉回后，$T$ 的 Kähler 模中的元素可加到 $V$ 的复化 Kähler 类上，得到新的复化 Kähler 类。因此，$T$ 的 Kähler 模空间可与 $V$ 的 Kähler 模空间某个叶状结构的各叶等同，而区域 $D'$ 比 $V$ 的 Kähler 模中的这片叶稍大。这与刚才对复模所作的处理很相似。

K3 曲面 $S$ 给出区域 $D$ 与 $D'$，它们分别与 $V$ 的复模和 Kähler 模相关。镜像一侧有上面构造的曲面 $S_1$，以 $D_1$、$D_1'$ 表示 $S_1$ 的相应区域。若 $V$ 与 $V^\circ$ 的确是镜像，则 $D$ 与 $D_1'$ 之间以及 $D'$ 与 $D_1$ 之间存在关系，将与镜像对称相一致。Voisin 证明了如下结果。
\begin{theorem}\label{thm:voisin-domains}
$D\simeq D_1'$，且 $D'\simeq D_1$。
\end{theorem}

Voisin 给出的另一类证据是 Yukawa 耦合的渐近行为。在 $H^{1,1}(V)$ 上有杯积给出的耦合；如\hyperref[chap:introduction]{第 1 章}定义，$H^1(V,T_V)$ 上的 Yukawa 耦合为
\begin{equation}\label{eq:voisin-yukawa}
\langle\theta_1,\theta_2,\theta_3\rangle
=\int_V\Omega\wedge
(\nabla_{\theta_1}\nabla_{\theta_2}\nabla_{\theta_3}\Omega),
\end{equation}
其中 $\Omega$ 是 $V$ 上的全纯三形式，见式~\eqref{eq:b-yukawa}。同构 $H^1(V,T_V)\simeq H^{2,1}(V)$ 在 $H^{2,1}(V)$ 上给出一个耦合，也称为 Yukawa 耦合。

利用式~\eqref{eq:voisin-h11-decomposition}、\eqref{eq:voisin-h21-decomposition} 的分解，Voisin 仔细描述了 $H^{1,1}(V)$ 与 $H^{2,1}(V^\circ)$ 上的耦合，并证明 $H^{1,1}(V)$ 上的杯积可渐近地与 $H^{2,1}(V^\circ)$ 上的 Yukawa 耦合等同。

为解释最后这个断言，回顾\hyperref[chap:introduction]{第 1 章}：$H^{2,1}(V^\circ)$ 上的 B 模型关联函数恰好就是上述 Yukawa 耦合，而 $H^{1,1}(V)$ 上的 A 模型关联函数由以下公式给出：
\[
\begin{aligned}
\langle\omega_1,\omega_2,\omega_3\rangle
={}&\int_V\omega_1\wedge\omega_2\wedge\omega_3\\
&+\sum_{\substack{\beta\ne0\\\beta\text{ 有效}}}n_\beta
\left(\int_\beta\omega_1\right)
\left(\int_\beta\omega_2\right)
\left(\int_\beta\omega_3\right)
\frac{e^{2\pi i\int_\beta\omega}}{1-e^{2\pi i\int_\beta\omega}},
\end{aligned}
\]
见式~\eqref{eq:a-yukawa}。现在设 $V$ 的复化 Kähler 类趋于无穷，即对所有有效曲线 $\beta$，都有 $\int_\beta\operatorname{Im}(\omega)\to\infty$；物理文献称之为“大半径极限”。于是
\[
\langle\omega_1,\omega_2,\omega_3\rangle
\longrightarrow\int_V\omega_1\wedge\omega_2\wedge\omega_3,
\]
右边就是杯积。在这种情形，称 A 模型关联函数渐近地变为杯积。

由于镜像流形 $V^\circ$ 的 B 模型关联函数与原流形 $V$ 的 A 模型关联函数一致，预期 B 模型关联函数也应具有类似行为。这正是 Voisin 证明的：在适当假设下，$V^\circ$ 上的 Yukawa 耦合有一个极限，与 $V$ 上的杯积一致，见 \SourceCite[第 3 节]{Voisin1}。

\SourcePage{93}
\subsection{与 Batyrev 构造的关系}\label{subsec:nonreflexive-weights}

Borcea 为 $V$ 与其 Voisin--Borcea 镜像 $V^\circ$ 的镜像对称提供了另一类证据：他证明，此构造的某些情形是 Batyrev 镜像对偶性的特殊情形。例如，考虑加权射影空间 $\mathbb P(3,3,2,2,2)$ 中的反典范超曲面
\begin{equation}\label{eq:voisin-weighted-original}
u^4+v^4=y^6+z^6+w^6,
\end{equation}
其中变量 $u,v,y,z,w$ 的权重依次为 $3,3,2,2,2$。为将它与式~\eqref{eq:voisin-threefold} 的三维簇联系起来，注意 $x^2=u^4+v^4$ 在 $\mathbb P(2,1,1)$ 中定义椭圆曲线 $E$，而 $x^2=y^6+z^6+w^6$ 在 $\mathbb P(3,1,1,1)$ 中定义（奇异的）K3 曲面 $S$。两种情形中，$x\mapsto-x$ 都是一个对合。Borcea 证明，Calabi--Yau 三维簇 $V=\widetilde{E\times S}/\widetilde k$ 是上述 $\mathbb P(3,3,2,2,2)$ 中超曲面的消解。还应注意，在 $S$ 中，对合的不动点集由 $y^6+z^6+w^6=0$ 给出，这是亏格为 $10$ 的光滑曲线。因此，此时 $(N,N')=(1,10)$。

利用 Voisin--Borcea 构造，族~\eqref{eq:voisin-weighted-original} 有镜像族 $V_{\mathrm{VB}}^\circ$，其不变量为 $(N,N')=(10,1)$。在 \SourceCite{Borcea2} 中，Borcea 证明，这个族包含加权射影空间 $\mathbb P(25,25,20,16,14)$ 中的反典范超曲面
\begin{equation}\label{eq:voisin-weighted-mirror}
u^4+v^4=y^5+yz^5+zw^6.
\end{equation}
此时，$E$ 仍是 $\mathbb P(2,1,1)$ 中的 $x^2=u^4+v^4$，但 $S_1$ 是 $\mathbb P(25,10,8,7)$ 中由 $x^2=y^5+yz^5+zw^6$ 定义的奇异 K3 曲面。如前，式~\eqref{eq:voisin-weighted-mirror} 的消解为 $V_{\mathrm{VB}}^\circ=\widetilde{E\times S_1}/\widetilde k$。对合的不动点集在 $S_1$ 中为 $y^5+yz^5+zw^6=0$。消解 $S_1$ 后，不动点集成为一条椭圆曲线加上九条 $\mathbb P^1$，这解释了为何 $(N,N')=(10,1)$。

基本断言是：$V_{\mathrm{VB}}^\circ$ 与式~\eqref{eq:voisin-weighted-original} 的 Batyrev 镜像一致，后者记为 $V_{\mathrm{Bat}}^\circ$。Borcea 在 \SourceCite{Borcea2} 中用\emph{\IndexLocation{indexloc:0075:93}双棱锥构造}证明了这一点；我们将利用 \SourceCite{CdK} 的\emph{广义 Berglund--Hübsch \IndexLocation{indexloc:0371:93}转置法则}给出更初等的证明。式~\eqref{eq:voisin-weighted-original} 与式~\eqref{eq:voisin-weighted-mirror} 之间的对偶性并不显然，因为 $\mathbb P(3,3,2,2,2)$ 是 Fano 环面簇，而 $\mathbb P(25,25,20,16,14)$ 不是，见引理~\ref{lem:weighted-projective-fano}。那么如何应用 Batyrev 对偶性？

答案涉及物理文献中的一种方法：把 Batyrev 镜像表示成非 Fano 加权射影空间中的超曲面。这里不描述一般情形，而说明它对 $\mathbb P(q)$、$q=(3,3,2,2,2)$ 如何起作用；更多细节见 \SourceCite[第 3.1 节]{CdK}。从环面几何角度看 $\mathbb P(q)$，有
\[
M=\{a\in\mathbb Z^5:q\cdot a=0\},\qquad
N=\mathbb Z^5/\mathbb Z q.
\]
$\mathbb P(q)$ 的扇的一维锥生成元为 $w_i=[e_i]\in N$，其中 $[e_i]$ 表示 $\mathbb Z^5$ 中第 $i$ 个标准基向量 $e_i$ 的陪集。此外，反典范除子的多胞形为
\[
\begin{aligned}
\Delta&=\{a\in M_{\mathbb R}:\langle a,w_i\rangle\geq-1\}\\
&=\{a\in M_{\mathbb R}:a\cdot e_i\geq-1\}.
\end{aligned}
\]
由于 $\mathbb P(q)$ 是 Fano 簇，$\Delta$ 自反。

\SourcePage{94}
现在考虑 $\Delta\cap M$ 中的以下各点：
\begin{equation}\label{eq:voisin-exponent-rows}
\begin{aligned}
a_0&=(3,-1,-1,-1,-1),\\
a_1&=(-1,3,-1,-1,-1),\\
a_2&=(-1,-1,4,0,-1),\\
a_3&=(-1,-1,-1,4,0),\\
a_4&=(-1,-1,-1,-1,5).
\end{aligned}
\end{equation}
这些点张成以 $0$ 为内点的单形。因此，某个系数全为正的 $a_i$ 的线性组合等于零，可以验证
\begin{equation}\label{eq:voisin-dual-weights}
25a_0+25a_1+20a_2+16a_3+14a_4=0.
\end{equation}
而且，其他所有关系都是这个关系的倍数。因此，$\widehat q=(25,25,20,16,14)$ 是各 $a_i$ 之间唯一的系数为正整数且互素的关系。它给出加权射影空间 $\mathbb P(\widehat q)=\mathbb P(25,25,20,16,14)$，具有
\[
\widehat M=\{a\in\mathbb Z^5:\widehat q\cdot a=0\},\qquad
\widehat N=\mathbb Z^5/\mathbb Z\widehat q.
\]
$\mathbb P(\widehat q)$ 的扇的一维锥生成元为 $\widehat w_i=[e_i]\in\widehat N$。

由于 $\Delta$ 自反，$\mathbb P(q)=\mathbb P_\Delta$ 是 Fano 簇；极多胞形 $\Delta^\circ$ 给出 Fano 环面簇 $\mathbb P_{\Delta^\circ}$，记为 $\mathbb P(q)^\circ$。$\mathbb P(q)^\circ$ 的扇位于 $M_{\mathbb R}$ 中，各锥由 $\Delta$ 的各面确定。为将 $\mathbb P(q)^\circ$ 与非 Fano 加权射影空间 $\mathbb P(\widehat q)$ 联系起来，使用映射 $\phi:\widehat N\to M$，将 $\widehat w_i\in\widehat N$ 送到 $a_i\in M$。由式~\eqref{eq:voisin-dual-weights}，容易看出 $\phi$ 良定义，并诱导同构 $\widehat N_{\mathbb R}\simeq M_{\mathbb R}$。因为 $\phi$ 把一维锥生成元 $\widehat w_i$ 送到 $a_i\in\Delta\cap M$，后者也可作为一维锥生成元，故能找到与 $\mathbb P(\widehat q)$ 的扇相容的极大射影细分 $\Sigma^\circ$。这意味着，将 $\Sigma^\circ$ 经 $\phi$ 拉回到 $\widehat N_{\mathbb R}$ 后所得的扇，细化 $\mathbb P(\widehat q)$ 的扇。由此得到环面爆破 $\widehat X\to\mathbb P(\widehat q)$，而 $\phi$ 诱导态射
\[
\phi_*:\widehat X\longrightarrow X_{\Sigma^\circ}.
\]
根据标准理论，$\phi_*$ 是有限态射，其次数为 $\phi(\widehat N)\subset M$ 的指数。令 $d=q\cdot(1,1,1,1,1)=12$、$\widehat d=\widehat q\cdot(1,1,1,1,1)=100$，则指数由下式计算：
\[
[M:\phi(\widehat N)]
=\frac{\det(A+\mathbf1)}{d\widehat d}
=\frac{2400}{12\cdot100}=2,
\]
其中 $A$ 是以式~\eqref{eq:voisin-exponent-rows} 中各向量为行的 $5\times5$ 矩阵，$\mathbf1$ 是所有元素均为 $1$ 的 $5\times5$ 矩阵。\SourceCite[第 3.1 节]{CdK} 陈述了这个公式但未给证明，不过论证并不困难。因此 $\phi_*:\widehat X\to X_{\Sigma^\circ}$ 是一个 $2:1$ 覆盖。

下一步是将 $V_{\mathrm{VB}}^\circ\subset\widehat X$ 与反典范超曲面 $V_{\mathrm{Bat}}^\circ\subset X_{\Sigma^\circ}$ 的拉回等同。自反单形 $\Delta^\circ\subset N_{\mathbb R}$ 的顶点为各 $w_i$，它们给出 Laurent 多项式 $f=\sum_{i=0}^4 c_it^{w_i}\in L(\Delta^\circ)$。它虽不是 $L(\Delta^\circ)$ 中最一般的 Laurent 多项式，却有正确的 Newton 多胞形。因此可设 $V_{\mathrm{Bat}}^\circ$ 由 $f=0$ 定义。为看出这在 $\widehat X$ 中给出什么，需要研究与 $\phi$ 对偶的映射 $\phi^t:N\to\widehat M$。将 $\mathbb Z^5$ 的元素视为列向量，并使用 $\widehat N$ 与 $M$ 的自然基及坐标，则复合映射
\[
\mathbb Z^5\longrightarrow\widehat N\xrightarrow{\phi}M
\longrightarrow\mathbb Z^5
\]
\SourcePage{95}
就是乘以 $A^t$，其中 $A$ 仍是以式~\eqref{eq:voisin-exponent-rows} 中各向量为行的矩阵。取对偶后，复合映射
\[
\mathbb Z^5\longrightarrow N\xrightarrow{\phi^t}\widehat M
\longrightarrow\mathbb Z^5
\]
就是乘以 $A$。这意味着 $\phi^t(w_i)=\widehat a_i$，其中 $\widehat a_i$ 是 $A$ 的第 $i$ 列。因此，$X_{\Sigma^\circ}$ 上的 $f=0$ 拉回为 $\widehat f=0$，其中 $\widehat f=\sum_{i=0}^4c_it^{\widehat a_i}$。最后一步是关于变量 $u,v,y,z,w$ 作齐次化。留给读者验证，这只需将 $\widehat a_i$ 的每个分量加 $1$。因此，得到的指数向量为矩阵
\[
A+\mathbf1=
\begin{pmatrix}
4&0&0&0&0\\
0&4&0&0&0\\
0&0&5&1&0\\
0&0&0&5&1\\
0&0&0&0&6
\end{pmatrix}
\]
的各列。用 $u,v,y,z,w$ 可将 $\widehat f=0$ 写为
\[
\widehat f=c_0u^4+c_1v^4+c_2y^5+c_3yz^5+c_4zw^6=0.
\]
式~\eqref{eq:voisin-weighted-mirror} 是这种形式的方程，因此可设 $V_{\mathrm{VB}}^\circ$ 由 $\widehat f=0$ 定义。

上述论证证明了 $\phi_*$ 诱导 $2:1$ 映射 $\phi_*:V_{\mathrm{VB}}^\circ\to V_{\mathrm{Bat}}^\circ$。这看起来不完全符合期望，因为我们想让二者相同。幸而有一个简单解释。在处理 $V_{\mathrm{VB}}^\circ=\widetilde{E\times S_1}/\widetilde k$ 时，将 $E$ 模去一个二阶点生成的子群，得到次数为二的同源映射 $E\to E'$，它与对合交换。因此有 $2:1$ 覆盖
\[
V_{\mathrm{VB}}^\circ\longrightarrow V_{\mathrm{VB}}'^\circ
=\widetilde{E'\times S_1}/\widetilde k.
\]
可以证明，$\phi_*:V_{\mathrm{VB}}^\circ\to V_{\mathrm{Bat}}^\circ$ 恰好是这种映射，所以可将 $V_{\mathrm{Bat}}^\circ$ 与 $V_{\mathrm{VB}}'^\circ$ 等同。证明使用 Borcea 的\IndexLocation{indexloc:0075:95}双棱锥构造，见 \SourceCite{Borcea2} 中式 (29) 后的说明。现在有一个关键观察：因为 $E$ 与 $E'$ 互为形变，$V_{\mathrm{VB}}^\circ$ 与 $V_{\mathrm{VB}}'^\circ$ 也互为形变。因此，就\emph{镜像族}而言，式~\eqref{eq:voisin-weighted-original} 的 Batyrev 镜像就是式~\eqref{eq:voisin-weighted-mirror}。这表明，至少在 $(N,N')=(1,10)$ 的情形，Voisin--Borcea 构造与 Batyrev 镜像对偶性一致。

论文 \SourceCite{Borcea2} 还列出了一些可用 Batyrev 镜像对偶性解释 Voisin--Borcea 构造的 $(N,N')$。不过，并非每个 Calabi--Yau 簇都是环面超曲面，所以可能并非所有 $(N,N')$ 都能这样处理。尽管如此，两种差别很大的\IndexLocation{indexloc:0799:95}\IndexLocation{indexloc:0804:95}镜像构造，即 Batyrev 构造与 Voisin--Borcea 构造，在式~\eqref{eq:voisin-weighted-original} 这样的情形中一致，是镜像对称确实存在的有力迹象。另一方面，两种构造的差异也表明，我们离一个真正普遍适用的寻找镜像的方法还有很远。

最后谈谈\hyperref[subsec:batyrev-hodge]{第 4.1.3 小节}提到的 $7555$ 个加权射影空间的列表。与 $\mathbb P(25,25,20,16,14)$ 一样，表中大部分加权射影空间\emph{不是} Fano 的。列表只显示出部分对称性，人们因而担心那些“缺失的镜像”。这个困难的解决使用了 Batyrev 镜像对偶性，以及上述用非 Fano 加权射影空间解释 Batyrev 对偶的方法的一个变体。从这个角度看，“缺失的镜像”对应于那些不能用加权射影空间表示的 Batyrev 对偶。细节与例子参见 \SourceCite{CdK}。


\chapter{Hodge 理论与 Yukawa 耦合}\label{chap:hodge-yukawa}
\SourcePage{96}

本章回顾并发展 Hodge 理论中的一些概念，将其用于计算 Calabi--Yau 流形 $V$ 的 Yukawa 耦合。为此，必须先理解 $V$ 上的周期积分所满足的 Picard--Fuchs 微分方程。我们将说明求出这些方程的两种方法。

\hyperref[sec:hodge-theory]{第 5.1 节}回顾 Hodge 结构变分、Picard--Fuchs 方程以及边界行为；\hyperref[sec:maximal-unipotent]{第 5.2 节}定义\emph{最大幺幂边界点}，它们在确定镜像映射时起重要作用。\hyperref[sec:griffiths-dwork]{第 5.3 节}解释 Griffiths 与 Dwork 发展的方法，用射影空间上有理微分形式的留数表示 $V$ 上的上同调类。\hyperref[sec:hodge-examples]{第 5.4 节}给出两个例题。\hyperref[sec:hypergeometric]{第 5.5 节}研究计算 Picard--Fuchs 方程的另一种方法，涉及 Gelfand、Kapranov 与 Zelevinsky 的\emph{广义超几何方程}。然后在\hyperref[sec:yukawa-couplings]{第~\ref*{sec:yukawa-couplings} 节}将这些内容用于计算 Yukawa 耦合。

本章有两个出人意料的方面：\hyperref[subsec:nef-gkz]{第~\ref*{subsec:nef-gkz} 小节}的 GKZ 方程组以上同调为值的解 $\widetilde I$，将在\hyperref[chap:mirror-theorem]{第~\ref*{chap:mirror-theorem} 章}的镜像定理讨论中出现；\hyperref[subsec:normalized-yukawa-connection]{第~\ref*{subsec:normalized-yukawa-connection} 小节}归一化 Yukawa 耦合与 Gauss--Manin 联络之间的关系，将在\hyperref[chap:quantum-cohomology]{第~\ref*{chap:quantum-cohomology} 章}讨论的 A 模型 Hodge 结构变分中发挥突出作用。

\section{Hodge 理论}\label{sec:hodge-theory}

本节回顾镜像对称中使用的 Hodge 理论。我们遵循 \SourceCite{Morrison2} 的思路，同时推广到参数空间维数更高的族。关于 Hodge 结构变分的一般参考文献有 \SourceCite{Griffiths1,Griffiths2,Schmid,CaK}；Calabi--Yau 变分的具体内容可参见 \SourceCite{BG}。

\subsection{Hodge 结构及其变分}\label{subsec:hodge-variations}

若 $X$ 是 $n$ 维光滑复射影簇，则 $H^k(X,\mathbb C)$ 有 Hodge 分解
\[
H^k(X,\mathbb C)\simeq\bigoplus_{p+q=k}H^{p,q}(X),
\]
其中相对于 $H^k(X,\mathbb R)$ 所确定的实结构，有 $\overline{H^{p,q}(X)}=H^{q,p}(X)$。再加入来自 $H^k(X,\mathbb Z)$ 的整数格，就得到一个\emph{权为 $k$ 的 Hodge 结构}。

我们常用 \emph{Hodge 滤过} $F^\bullet(X)=\{F^p(X)\}_{p=0}^n$ 来表述 $H^k(X,\mathbb C)$ 上的 Hodge 结构，其中
\[
F^p(X)=\bigoplus_{a\geq p}H^{a,k-a}(X).
\]
\SourcePage{97}
注意，$F^p(X)\oplus\overline{F^{k-p+1}(X)}=H^k(X,\mathbb C)$，且
\[
H^{p,q}(X)=F^p(X)\cap\overline{F^q(X)}.
\]
还有 $H^{p,q}(X)\simeq F^p(X)/F^{p+1}(X)$，其中 $p+q=k$。因此，$H^k(X,\mathbb C)$ 上的 Hodge 结构由整数结构与 Hodge 滤过唯一确定。

最后，$X$ 上的 Kähler 类 $\omega$ 在 $H^k(X,\mathbb C)$ 中确定\emph{本原上同调群} $H_0^k(X,\mathbb C)$，也记为 $PH^k(X)$。配对
\[
Q(\alpha,\beta)=(-1)^{k(k-1)/2}
\int_X\omega^{n-k}\wedge\alpha\wedge\beta,
\qquad\alpha,\beta\in H_0^k(X,\mathbb C),
\]
使本原上同调成为\emph{极化 Hodge 结构}。这意味着
\begin{equation}\label{eq:hodge-polarization}
Q(i^{p-q}\alpha,\overline\alpha)>0,
\qquad0\ne\alpha\in H_0^{p,q}(X),\quad p+q=k.
\end{equation}
关于极化 Hodge 结构的更多内容，见 \SourceCite{CaK,Griffiths2}。

子空间 $H^{p,q}(X)$ 在族中并不全纯地变化，但 $F^p(X)$ 的相对版本具有很好的性质。设有相对维数为 $n$ 的光滑态射 $\pi:\mathcal X\to S$，其中 $S$ 光滑且拟射影。对任意 $t\in S$，置 $X_t=\pi^{-1}(t)$。各上同调群 $H^k(X_t,\mathbb C)$ 拼成 $S$ 上的局部自由层 $\mathcal F^0=R^k\pi_*\mathbb C\otimes\mathcal O_S$。各空间 $F^p(X_t)$ 也拼成局部自由子层 $\mathcal F^p$。

由构造，$\mathcal F^0$ 含有局部系统，即局部常值层 $R^k\pi_*\mathbb C$。它唯一确定 $\mathcal F^0$ 上的平坦联络 $\nabla$，称为 \emph{Gauss--Manin 联络}；其平坦截面，也称水平截面，恰好是局部系统 $R^k\pi_*\mathbb C$ 的截面。具体地，$\nabla:\mathcal F^0\to\mathcal F^0\otimes\Omega_S^1$ 定义为
\[
\nabla(s\otimes f)=s\otimes df,
\]
其中 $s$ 是 $R^k\pi_*\mathbb C$ 的截面，$f$ 是 $S$ 上的函数。Gauss--Manin 联络满足 \emph{Griffiths 横截性}
$\nabla(\mathcal F^p)\subset\mathcal F^{p-1}\otimes\Omega_S^1$。

若置 $\mathcal H=\mathcal F^0$，则 $\mathcal H$ 有局部常值子层 $\mathcal H_{\mathbb C}=R^k\pi_*\mathbb C$；它又有整数截面构成的子层 $\mathcal H_{\mathbb Z}$，即 $R^k\pi_*\mathbb Z\to R^k\pi_*\mathbb C$ 的像。称 $(\mathcal H,\nabla,\mathcal H_{\mathbb Z},\mathcal F^\bullet)$ 为 \emph{Hodge 结构变分}。\index{Hodge结构@Hodge 结构|hyperpage}\index{Hodge结构变分@Hodge 结构变分|hyperpage}\index{Gauss--Manin联络@Gauss--Manin 联络|hyperpage}\index{Griffiths横截性@Griffiths 横截性|hyperpage}

\subsection{Picard--Fuchs 方程}\label{subsec:picard-fuchs}

本节余下部分研究中维上同调群 $H^n(X_t,\mathbb C)$ 上的 Hodge 结构变分，其中 $n=\dim X_t$。固定 $p\in S$ 处 $\mathcal F^n$ 的局部截面 $\Omega(t)$，令 $\mathcal D$ 为 $S$ 上的线性微分算子层。在 $p$ 处的局部解析坐标 $z_1,\ldots,z_r$ 下，微分算子环为
\[
\mathcal D=\mathbb C\{z_1,\ldots,z_r\}
\left[\frac\partial{\partial z_1},\ldots,
\frac\partial{\partial z_r}\right],
\]
其中 $\mathbb C\{z_1,\ldots,z_r\}$ 是 $z_1,\ldots,z_r$ 的收敛幂级数环。Gauss--Manin 联络 $\nabla$ 确定 $\mathcal O_S$ 同态 $\phi:\mathcal D\to\mathcal F^0$，对 $S$ 上的向量场 $X_i$，其规则为
\[
\phi(X_1\cdots X_\ell)
=\nabla_{X_1}\cdots\nabla_{X_\ell}\Omega(t).
\]
联络同时赋予 $\mathcal F^0$ 一个 $\mathcal D$ 模结构。

左理想 $I=\ker\phi$ 由零化 $\Omega(t)$ 的微分算子组成，称为 \emph{Picard--Fuchs 理想}。此理想在理论中起重要作用。例如，后面会看到，知道 $I$ 就能计算 Yukawa 耦合。若上述映射 $\phi$ 满射，\SourceCite{BG} 证明这发生在 Calabi--Yau 三维簇上，
\SourcePage{98}
则 $\mathcal F^0$ 的 $\mathcal D$ 模结构完全由 Picard--Fuchs 理想 $I$ 确定。

对 $D\in I$，方程 $D\cdot y=0$ 称为 \emph{Picard--Fuchs 方程}；若 $\mathcal O_S$ 的局部截面 $y$ 对所有 $D\in I$ 都满足 $D\cdot y=0$，则称其为一个\emph{解}。由于对局部常值的同调 $n$ 圈 $\gamma_t$ 有
\begin{equation}\label{eq:period-derivative}
\frac\partial{\partial z_i}\int_{\gamma_t}\Omega(t)
=\int_{\gamma_t}\nabla_{\partial/\partial z_i}\Omega(t),
\end{equation}
所以 $\Omega$ 的周期 $y=\int_{\gamma_t}\Omega(t)$ 是 Picard--Fuchs 方程的解。我们常将周期写为 $y=\langle g(t),\Omega(t)\rangle$，其中 $g(t)\in H^n(X_t,\mathbb C)$ 是 $\gamma_t$ 的 Poincaré 对偶，因此 $g(t)$ 是 $\nabla$ 的平坦截面。注意，微分算子 $D$ 属于 Picard--Fuchs 理想，当且仅当对所有周期 $y=\langle g,\Omega\rangle$ 都有 $D\cdot y=0$。

虽然这里只在局部构造了 Picard--Fuchs 方程，但若限于代数微分形式的周期，通过解析延拓，它们可作为整个 $S$ 上的代数微分方程来理解。

\subsection{Hodge 结构的退化}\label{subsec:hodge-degenerations}

设光滑族 $\pi:\mathcal X\to S$ 可补成平坦族 $\overline\pi:\overline{\mathcal X}\to\overline S$，其中 $\overline S$ 是 $S$ 的光滑紧化，边界除子 $D=\bigcup_j D_j=\overline S-S$ 具有正规交叉。

$S$ 上的丛 $\mathcal F^0$ 在 $\overline S$ 上有一个\emph{典范延拓} $\overline{\mathcal F}^0$，见 \SourceCite{Deligne1}。Gauss--Manin 联络不一定延拓为 $\overline{\mathcal F}^0$ 上的通常联络，因为它可能产生奇点。这些奇点很温和，是\IndexLocation{indexloc:1225:98}正则奇点；也就是说，$\nabla$ 延拓为映射
$\overline\nabla:\overline{\mathcal F}^0\to\overline{\mathcal F}^0\otimes\Omega_{\overline S}^1(\log D)$。
层 $\Omega_{\overline S}^1(\log D)$ 是 $\overline S$ 上有理一形式层的子层。在 $\overline S$ 的局部坐标 $z_1,\ldots,z_r$ 中，若 $D$ 由 $z_1\cdots z_\ell=0$ 定义，则 $\Omega_{\overline S}^1(\log D)$ 由
\[
\frac{dz_1}{z_1},\ldots,\frac{dz_\ell}{z_\ell},
dz_{\ell+1},\ldots,dz_r
\]
生成。

还需考虑单值性。设 $\gamma_j=\gamma_j(u)$ 是围绕边界除子 $D_j$ 的一条小回路，以 $\gamma_j(0)=\gamma_j(1)=t\in S$ 为基点。上同调类 $\eta\in H^n(X_t,\mathbb C)$ 可以唯一提升为 $[0,1]$ 上的 $\nabla$ 平坦截面 $\eta(u)\in H^n(X_{\gamma_j(u)},\mathbb C)$，满足 $\eta(0)=\eta$。\emph{\IndexLocation{indexloc:0901:98}\IndexLocation{indexloc:0902:98}\IndexLocation{indexloc:0904:98}\IndexLocation{indexloc:0905:98}单值性变换}，也称 \emph{Picard--Lefschetz 变换}，
\[
\mathcal T_j:H^n(X_t,\mathbb C)\longrightarrow H^n(X_t,\mathbb C)
\]
定义为 $\mathcal T_j(\eta)=\eta(1)$。在共轭意义下，它只依赖于 $D_j$，而不依赖于 $\gamma_j$ 的选择。\IndexLocation{indexloc:0900:98}单值性定理 \SourceCite{Landman} 表明，对某个整数 $m>0$，有
\[
(\mathcal T_j^m-I)^{n+1}=0.
\]
因此 $\mathcal T_j$ 是拟幺幂的，幺幂指数至多为 $n+1$。特别地，$\mathcal T_j$ 的特征值都是单位根。\index{单值性|hyperpage}

\subsection{极限混合 Hodge 结构}\label{subsec:limiting-hodge}

假设单值性幺幂，且在点 $p\in\overline S$ 处，$D=\overline S-S$ 由 $z_1\cdots z_r=0$ 定义；$z_1,\ldots,z_r$ 仍为 $p$ 处的局部坐标。因此可设 $S=(\Delta^*)^r$、$\overline S=\Delta^r$、$p=0$。令 $N_j=\log\mathcal T_j$，其中 $\mathcal T_j$ 为逆时针绕 $(\Delta^*)^r$ 的第 $j$ 个因子一周所得的单值性。在这种情形，$\mathcal F^0$ 的典范延拓可描述如下。设 $s$ 是 $(\Delta^*)^r$ 上 $\mathcal H=\mathcal F^0$ 的平坦多值截面，则
\[
\exp\left(-\frac1{2\pi i}\sum_j\log(z_j)N_j\right)s
\]
\SourcePage{99}
是单值的，并延拓为典范延拓 $\overline{\mathcal F}^0$ 的截面。容易看出，这个性质刻画了 $\overline{\mathcal F}^0$。

此外，幂零轨道定理 \SourceCite{Schmid} 意味着丛 $\mathcal F^p$ 延拓为 $\overline{\mathcal F}^p$，它们是典范延拓的子丛。在 $0\in\Delta^r$，$\overline{\mathcal F}^\bullet$ 给出极限 Hodge 滤过 $\overline{\mathcal F}_0^\bullet=F_{\mathrm{lim}}^\bullet$。可以证明 $N_j(F_{\mathrm{lim}}^p)\subset F_{\mathrm{lim}}^{p-1}$，因此得到线性映射
\[
N_j:F_{\mathrm{lim}}^p/F_{\mathrm{lim}}^{p+1}
\longrightarrow F_{\mathrm{lim}}^{p-1}/F_{\mathrm{lim}}^p.
\]
若 $\delta_j=z_j\partial/\partial z_j$，则由 Griffiths 横截性，
$\overline\nabla_{\delta_j}(\overline{\mathcal F}^p)\subset\overline{\mathcal F}^{p-1}$。于是 $\overline\nabla_{\delta_j}$ 诱导线性映射
\[
\overline\nabla_{\delta_j}:F_{\mathrm{lim}}^p/F_{\mathrm{lim}}^{p+1}
\longrightarrow F_{\mathrm{lim}}^{p-1}/F_{\mathrm{lim}}^p.
\]
利用幂零轨道定理，可见这两个映射满足关系
\begin{equation}\label{eq:connection-monodromy-residue}
\overline\nabla_{\delta_j}=-\frac1{2\pi i}N_j.
\end{equation}

在 $0\in\Delta^r$ 上还有自然的整数结构。设 $s$ 是 $(\Delta^*)^r$ 上 $\mathcal H_{\mathbb Z}$ 的平坦多值截面。如上，
\[
\widetilde s=\exp\left(-\frac1{2\pi i}\sum_j\log(z_j)N_j\right)s
\]
是典范延拓 $\overline{\mathcal F}^0$ 的截面，$\widetilde s(0)$ 是 $0$ 上的整元素。若用 $\exp(\tfrac1{2\pi i}\sum_j\log(z_j)N_j)$ 将其传回 $H^n(X_t,\mathbb C)$，就得到 $H^n(X_t,\mathbb Z)$ 所确定的通常整数结构。

另一个重要部分是\emph{单值性权滤过}
\[
W_\bullet:\quad W_0\subset W_1\subset\cdots\subset W_{2n-1}
\subset W_{2n}=H^n(X_t,\mathbb C),
\]
它由 $N=\sum_j a_jN_j$、$a_j>0$ 在 $H^n(X_t,\mathbb C)$ 上的作用定义。例如，
\begin{equation}\label{eq:monodromy-weight-filtration}
\begin{aligned}
W_0&=\operatorname{im}(N^n),\\
W_1&=\operatorname{im}(N^{n-1})\cap\ker N,\\
W_2&=\operatorname{im}(N^{n-2})\cap\ker N
+\operatorname{im}(N^{n-1})\cap\ker(N^2),\\
&\ \vdots\\
W_{2n-1}&=\ker(N^n).
\end{aligned}
\end{equation}
其他 $W_k$ 的公式可在 \SourceCite{Griffiths2} 中找到。单值性权滤过的主要性质为：
\begin{itemize}
\item $N(W_k)\subset W_{k-2}$；
\item $N^k$ 诱导同构
$N^k:W_{n+k}/W_{n+k-1}\simeq W_{n-k}/W_{n-k-1}$；
\item $F_{\mathrm{lim}}^\bullet$ 在 $W_k/W_{k-1}$ 上诱导权为 $k$ 的 Hodge 结构。
\end{itemize}
最后一项是说，$(W_\bullet,F_{\mathrm{lim}}^\bullet)$ 是\emph{混合 Hodge 结构}。更完整的讨论见 \SourceCite{Griffiths2}。单值性权滤过在\hyperref[sec:maximal-unipotent]{第 5.2 节}中将起重要作用。\index{极限混合Hodge结构@极限混合 Hodge 结构|hyperpage}\index{单值性权滤过|hyperpage}

\subsection{单值性与 Picard--Fuchs 方程}\label{subsec:monodromy-pf}

为简单起见，设 $\dim S=1$，令 $z$ 为以 $p\in\overline S-S$ 为中心的圆盘 $\Delta\subset\overline S$ 上的坐标。选择 $\overline{\mathcal F}^0$ 在 $\Delta$ 上的一组基 $\omega_1,\ldots,\omega_r$，则联络 $\overline\nabla$ 完全由其\emph{\IndexLocation{indexloc:0174:99}联络矩阵} $(\Gamma_{ij})$ 确定，其中 $\overline\nabla_{d/dz}\omega_i=\sum_j\Gamma_{ij}\omega_j$。

\SourcePage{100}
因为 $\overline\nabla$ 有正则奇点，$\Gamma_{ij}$ 在 $z=0$ 处至多有一阶极点。由此得到留数矩阵
\begin{equation}\label{eq:connection-residue}
\operatorname{Res}(\nabla)=\operatorname{Res}_{z=0}(\Gamma_{ij}),
\end{equation}
根据 \SourceCite{Deligne1}，它具有以下性质：
\begin{itemize}
\item $\operatorname{Res}(\nabla)$ 的特征值 $\lambda$ 是有理数，满足 $0\leq\lambda<1$；
\item $\exp(-2\pi i\operatorname{Res}(\nabla))$ 与单值性 $\mathcal T$ 共轭；
\item $\mathcal T$ 幺幂，当且仅当 $\operatorname{Res}(\nabla)$ 幂零。
\end{itemize}
实际计算时，可能不知道 $\overline{\mathcal F}^0$ 的一组基。改为假定 $\omega_1,\ldots,\omega_r$ 是去心圆盘 $\Delta^*$ 上的一组基。若联络矩阵在 $z=0$ 处至多有一阶极点，留数矩阵 $\operatorname{Res}(\nabla)$ 仍给出有用的信息：
\begin{itemize}
\item 当 $\lambda$ 遍历 $\operatorname{Res}(\nabla)$ 的特征值时，单值性 $\mathcal T$ 的特征值为 $\exp(-2\pi i\lambda)$；
\item $\mathcal T$ 幺幂，当且仅当 $\operatorname{Res}(\nabla)$ 的特征值全为整数；
\item 若 $\operatorname{Res}(\nabla)$ 的任意两个不同特征值之差都不是整数，则 $\mathcal T$ 与 $\exp(-2\pi i\operatorname{Res}(\nabla))$ 共轭。
\end{itemize}

为看出此时 Picard--Fuchs 方程的形式，令 $\delta=z\,d/dz$。对 $\overline{\mathcal F}^n$ 的固定截面 $\Omega(z)$，设在 $\overline{\mathcal F}^0$ 中有关系
\[
\nabla_\delta^m\Omega+f_1(z)\nabla_\delta^{m-1}\Omega
+\cdots+f_m(z)\Omega=0,
\]
其中各 $f_i(z)$ 都在 $z=0$ 处解析。如前所述，这意味着 $\delta^m+f_1(z)\delta^{m-1}+\cdots+f_m(z)$ 属于 Picard--Fuchs 理想，而周期 $y=\int_{\gamma_z}\Omega(z)$ 是 Picard--Fuchs 方程
\begin{equation}\label{eq:regular-pf-equation}
\delta^my+f_1(z)\delta^{m-1}y+\cdots+f_m(z)y=0
\end{equation}
的解；由于单值性对 $\gamma_z$ 的影响，这些解可能是多值的。

现在像 \SourceCite{CL} 那样，将它与常微分方程联系起来。按照微分方程标准教材中的定义，式~\eqref{eq:regular-pf-equation} 在这里具有正则奇点。其解由\emph{指标方程}控制；这个方程通过将式~\eqref{eq:regular-pf-equation} 中的 $\delta$ 替换为辅助变量 $\lambda$，并在 $z=0$ 处取系数而得到：
\[
\lambda^m+f_1(0)\lambda^{m-1}+\cdots+f_m(0)=0.
\]
设 $\lambda$ 是指标方程的一个根，而且不存在根 $\lambda'$ 使 $\lambda'-\lambda$ 为正整数。则 Picard--Fuchs 方程~\eqref{eq:regular-pf-equation} 在 $z=0$ 附近有形如 $y_0=z^\lambda f_0(z)$ 的解，$f_0$ 在 $0$ 处解析且不为零。因此，将 $y_0$ 绕 $z=0$ 解析延拓一周，得到 $\exp(2\pi i\lambda)y_0$。此外，若 $\lambda$ 是重数 $\mu>1$ 的根，则经典的 Frobenius 方法除 $y_0=z^\lambda f_0(z)$ 外，还给出其他解 $y_1,\ldots,y_{\mu-1}$；当 $1\leq i\leq\mu-1$ 时，将 $y_i$ 绕 $z=0$ 解析延拓一周，得到 $\exp(2\pi i\lambda)(y_i+y_{i-1})$。这意味着当 $z\to0$ 时，有渐近行为
\[
y_i\sim\frac c{i!}z^\lambda
\left(\frac{\log z}{2\pi i}\right)^i,
\qquad0\leq i\leq\mu-1,
\]
其中 $c$ 是非零常数。数 $\mu$ 与指标方程的友矩阵的 Jordan 标准形有关。

上述 Picard--Fuchs 方程来自 $\mathcal F^n$ 的截面 $\Omega$。若改用 $\Omega'=g(z)\Omega$，其中 $g(z)$ 在 $S$ 上亚纯，则得到新的 Picard--Fuchs 方程，
\SourcePage{101}
其解为 $\widetilde y=g(z)y$，$y$ 是式~\eqref{eq:regular-pf-equation} 的解。在这种情形，可以证明新方程仍有正则奇点，指标方程为
\[
(\lambda-\ell)^m+f_1(0)(\lambda-\ell)^{m-1}
+\cdots+f_m(0)=0,
\]
其中 $g(z)$ 在 $z=0$ 处的 Laurent 展开为 $g(z)=a_\ell z^\ell+\cdots$，$a_\ell\ne0$。

下一节将看到，在适当类型的边界点处，Picard--Fuchs 方程与单值性之间的关系尤其紧密。\index{指标方程|hyperpage}\index{Picard--Fuchs方程@Picard--Fuchs 方程|hyperpage}

\section{最大幺幂单值性}\label{sec:maximal-unipotent}

本节将发现，五次三维簇的镜像具有某种特殊的单值性行为，它对镜像对称至关重要。

\subsection{一维模}\label{subsec:one-dimensional-mum}

首先在复模空间 $S$ 为一维的 Calabi--Yau 三维簇 $V$ 的情形中讨论最大幺幂单值性。\hyperref[chap:quintic]{第 2 章}描述的五次三维簇的镜像就是一个例子。$S$ 上未必有万有族，但可在 $S$ 的有限覆盖上局部构造族。只要从 $\mathcal F^3$ 的代数截面出发，所得 Picard--Fuchs 方程就具有代数微分方程的意义，因此可解析延拓到整个 $S$。仍令 $\overline S$ 是 $S$ 的紧化，设边界点 $p\in\overline S-S$ 具有\emph{最大幺幂单值性}。这意味着单值性 $\mathcal T:H^3(V,\mathbb C)\to H^3(V,\mathbb C)$ 幺幂，而单值性定理给出 $(\mathcal T-I)^4=0$，所以“最大幺幂”意味着 $(\mathcal T-I)^3\ne0$。因此 $N=\log\mathcal T$ 也满足 $N^3\ne0$。

首要目标是解释上一节的理论如何应用于此。我们将看到，最大幺幂单值性对周期与 Picard--Fuchs 方程有一些很强的推论。首先，因为 $V$ 的复模是一维的，所以 $h^{2,1}(V)=1$，进而 $h^3(V)=4$。由于 $\mathcal T$ 最大幺幂，可找到 $H^3(V,\mathbb C)$ 的一组基 $g_0,g_1,g_2,g_3$，使 $\mathcal T$ 的矩阵为
\begin{equation}\label{eq:mum-jordan-matrix}
\mathcal T=
\begin{pmatrix}
1&1&0&0\\
0&1&1&0\\
0&0&1&1\\
0&0&0&1
\end{pmatrix}.
\end{equation}
每个 $g_i$ 都对偶于同调类 $\gamma_i\in H_3(V,\mathbb C)$。给定 $\overline{\mathcal F}^3$ 在 $z=0$ 附近的一组固定局部基 $\Omega$，希望描述周期 $y_i=\int_{\gamma_i}\Omega$ 在 $0$ 附近的行为。

先遵循 \SourceCite{Morrison2} 的方法。任意 $g\in H^3(V,\mathbb C)$ 都延拓为 $\mathcal F^0$ 的多值平坦截面 $g(z)$，而根据幂零轨道定理，$\exp(-\tfrac{\log z}{2\pi i}N)g(z)$ 在 $z=0$ 处单值且解析。$g=g_0$ 时，$\mathcal T(g_0)=g_0$ 意味着 $N(g_0)=0$，所以 $\exp(-\tfrac{\log z}{2\pi i}N)g_0(z)=g_0(z)$。因此 $g_0(z)$ 在 $z=0$ 处单值。由于
\[
y_0=\int_{\gamma_0}\Omega=\langle g_0(z),\Omega(z)\rangle,
\]
$y_0$ 在 $z=0$ 处单值且解析。后面还将看到 $y_0$ 在 $z=0$ 处不为零。再考虑 $y_1,y_2,y_3$。注意，对 $i\geq1$，有 $\mathcal T(g_i)=g_i+g_{i-1}$，且 $y_i=\langle g_i(z),\Omega(z)\rangle$，所以将 $y_i$ 绕 $z=0$
\SourcePage{102}
解析延拓一周，得到 $y_i+y_{i-1}$。此外，由于 $\exp(-\tfrac{\log z}{2\pi i}N)g_i(z)$ 单值，容易得到当 $z\to0$ 时形如
\[
y_i\sim\frac c{i!}\left(\frac{\log z}{2\pi i}\right)^i,
\qquad0\leq i\leq3,
\]
的渐近公式，其中 $c$ 是非零常数。

这使人想到\hyperref[subsec:monodromy-pf]{第 5.1.5 小节}讨论的微分方程理论。因此，下面解释此时 Picard--Fuchs 方程的形式。关键结果如下。
\begin{proposition}\label{prop:mum-indicial-equation}
固定复模空间 $S$ 为一维的 Calabi--Yau 三维簇 $V$，并在 $S$ 的一个边界点处取局部坐标 $z$，其中 Gauss--Manin 联络在该点具有正则奇点。设 $\mathcal F^3$ 的局部截面 $\Omega$ 给出 Picard--Fuchs 方程
\[
\delta^4y+f_1(z)\delta^3y+f_2(z)\delta^2y
+f_3(z)\delta y+f_4(z)y=0,
\]
其中 $\delta=z\,d/dz$，各 $f_i$ 在 $z=0$ 处解析。则：
\begin{enumerate}[label=(\roman*)]
\item 单值性 $\mathcal T$ 幺幂，当且仅当指标方程的根全为整数；
\item 单值性 $\mathcal T$ 最大幺幂，当且仅当指标方程对某个整数 $\ell$ 具有形式 $(\lambda-\ell)^4=0$。
\end{enumerate}
\end{proposition}
\begin{proof}
上述 Picard--Fuchs 方程意味着，在 $\mathcal F^0$ 中有
\begin{equation}\label{eq:mum-connection-relation}
\nabla_\delta^4\Omega+f_1(z)\nabla_\delta^3\Omega
+f_2(z)\nabla_\delta^2\Omega+f_3(z)\nabla_\delta\Omega
+f_4(z)\Omega=0.
\end{equation}
根据 \SourceCite{BG}，还知道 $\Omega,\nabla_\delta\Omega,\nabla_\delta^2\Omega,\nabla_\delta^3\Omega$ 在一般点处必定线性无关。特别地，可设它们在 $z=0$ 的去心邻域 $\Delta^*$ 中构成 $\mathcal F^0$ 的一组基；因为 $\Omega$ 未必是 $\overline{\mathcal F}^3$ 的截面，$z=0$ 处未必仍是一组基。

利用基 $\Omega,\nabla_\delta\Omega,\nabla_\delta^2\Omega,\nabla_\delta^3\Omega$，可计算 $\nabla$ 的联络矩阵。对这些基向量应用 $\nabla_{d/dz}$，并使用式~\eqref{eq:mum-connection-relation}，得到
\[
\begin{aligned}
\nabla_{d/dz}\Omega&=\frac1z\nabla_\delta\Omega,\\
\nabla_{d/dz}(\nabla_\delta\Omega)&=\frac1z\nabla_\delta^2\Omega,\\
\nabla_{d/dz}(\nabla_\delta^2\Omega)&=\frac1z\nabla_\delta^3\Omega,\\
\nabla_{d/dz}(\nabla_\delta^3\Omega)&=-\frac{f_4(z)}z\Omega
-\frac{f_3(z)}z\nabla_\delta\Omega\\
&\quad-\frac{f_2(z)}z\nabla_\delta^2\Omega
-\frac{f_1(z)}z\nabla_\delta^3\Omega.
\end{aligned}
\]
联络矩阵在 $z=0$ 处至多有一阶极点，因此可用\hyperref[subsec:monodromy-pf]{第 5.1.5 小节}发展的理论。在此基下，$\nabla$ 的\IndexLocation{indexloc:1231:102}留数矩阵为
\begin{equation}\label{eq:mum-companion-residue}
\operatorname{Res}(\nabla)=
\begin{pmatrix}
0&1&0&0\\
0&0&1&0\\
0&0&0&1\\
-f_4(0)&-f_3(0)&-f_2(0)&-f_1(0)
\end{pmatrix},
\end{equation}
它是指标方程的友矩阵。

于是，$\mathcal T$ 幺幂，当且仅当 $\operatorname{Res}(\nabla)$ 的特征值全为整数，也当且仅当指标方程的根全为整数。这证明了命题的第一部分。

\SourcePage{103}
对第二部分，设指标方程为 $(\lambda-\ell)^4=0$，其中 $\ell$ 为整数。将 $\Omega$ 替换为 $z^{-\ell}\Omega$，根据\hyperref[sec:hodge-theory]{第 5.1 节}末尾的讨论，可设指标方程为 $\lambda^4=0$。于是式~\eqref{eq:mum-companion-residue} 变为
\[
\operatorname{Res}(\nabla)=
\begin{pmatrix}
0&1&0&0\\
0&0&1&0\\
0&0&0&1\\
0&0&0&0
\end{pmatrix}.
\]
没有两个不同的根之差为整数，因此 $\operatorname{Res}(\nabla)$ 的性质表明，$\mathcal T$ 与 $\exp(-2\pi i\operatorname{Res}(\nabla))$ 共轭，从而 $\mathcal T$ 最大幺幂。

反过来，设 $\mathcal T$ 最大幺幂。只须找出指标方程为 $\lambda^4=0$ 的 Picard--Fuchs 方程。取 $\overline{\mathcal F}^3$ 的局部基 $\Omega$。我们断言，$\Omega,\nabla_\delta\Omega,\nabla_\delta^2\Omega,\nabla_\delta^3\Omega$ 在 $z=0$ 附近构成 $\overline{\mathcal F}^0$ 的一组基。

利用\hyperref[subsec:limiting-hodge]{第 5.1.4 小节}描述的单值性权滤过证明此断言。由式~\eqref{eq:monodromy-weight-filtration}，$W_6=H^3(V,\mathbb C)$，$W_5=\ker(N^3)$，其中 $N=\log\mathcal T$。式~\eqref{eq:mum-jordan-matrix} 给出 $\dim(W_6/W_5)=1$。由于 $F_{\mathrm{lim}}^\bullet$ 在 $W_6/W_5$ 上诱导权为六的 Hodge 结构，且 $F_{\mathrm{lim}}^6=F_{\mathrm{lim}}^5=F_{\mathrm{lim}}^4=0$，必有 $W_6/W_5=H^{3,3}$。$W_6/W_5$ 与 $F_{\mathrm{lim}}^3$ 均为一维，因此 $F_{\mathrm{lim}}^3\cap W_5=\{0\}$。这给出直和分解
\begin{equation}\label{eq:mum-top-splitting}
F_{\mathrm{lim}}^3\oplus W_5=H^3(V,\mathbb C).
\end{equation}

$\overline{\mathcal F}^3$ 的局部基 $\Omega$ 给出非零元素 $\Omega(0)\in F_{\mathrm{lim}}^3$。由于 $W_5=\ker(N^3)$，式~\eqref{eq:mum-top-splitting} 意味着 $N^3(\Omega(0))\ne0$；再由式~\eqref{eq:mum-jordan-matrix}，$\Omega(0),N(\Omega(0)),N^2(\Omega(0)),N^3(\Omega(0))$ 是 $H^3(V,\mathbb C)$ 的一组基。迭代式~\eqref{eq:connection-monodromy-residue} 得
\[
\nabla_\delta^j(\Omega)(0)
=\left(-\frac1{2\pi i}\right)^jN^j(\Omega(0))
\pmod{F_{\mathrm{lim}}^{4-j}}.
\]
再利用 $N(F_{\mathrm{lim}}^p)\subset F_{\mathrm{lim}}^{p-1}$ 与 $\dim F_{\mathrm{lim}}^p=4-p$，容易推出 $\Omega,\nabla_\delta\Omega,\nabla_\delta^2\Omega,\nabla_\delta^3\Omega$ 在 $z=0$ 处线性无关。因此，它们在 $z=0$ 的某个邻域构成 $\overline{\mathcal F}^0$ 的一组基。

特别地，可以用这组基表示 $\nabla_\delta^4\Omega$，得到形如式~\eqref{eq:mum-connection-relation} 的 Picard--Fuchs 方程。还需证明 $f_1(0)=f_2(0)=f_3(0)=f_4(0)=0$。这已经很容易：如上，用这组基计算 $\nabla$ 的联络矩阵与留数矩阵，得到 $\operatorname{Res}(\nabla)$ 的公式~\eqref{eq:mum-companion-residue}。因为此矩阵使用 $\overline{\mathcal F}^0$ 的一组基构造，\hyperref[sec:hodge-theory]{第 5.1 节}表明，$\mathcal T$ 幺幂意味着 $\operatorname{Res}(\nabla)$ 幂零。因此指标方程为 $\lambda^4=0$，命题得证。
\end{proof}

关于命题~\ref{prop:mum-indicial-equation} 的证明，有两点说明。首先，可以证明 $W_{2k}=W_{2k+1}$，且对 $0\leq p\leq3$ 有
\begin{equation}\label{eq:mum-hodge-tate-splitting}
F_{\mathrm{lim}}^p\oplus W_{2p-1}
=F_{\mathrm{lim}}^p\oplus W_{2p-2}=H^3(V,\mathbb C),
\end{equation}
这推广了式~\eqref{eq:mum-top-splitting}。由此可见，此时混合 Hodge 结构 $(W_\bullet,F_{\mathrm{lim}}^\bullet)$ 是 \emph{Hodge--Tate 结构}。更一般地，\SourceCite{Deligne2} 讨论了 Hodge--Tate 结构与镜像对称的关系。

第二点涉及命题~\ref{prop:mum-indicial-equation} 前的讨论：那里断言解 $y_0=\langle g_0,\Omega\rangle$ 在 $z=0$ 处不为零。为理解这一点，先注意，由式~\eqref{eq:monodromy-weight-filtration} 与
\[
\langle N(a),b\rangle+\langle a,N(b)\rangle=0,
\qquad a,b\in H^3(V,\mathbb C),
\]
容易推出 $\langle W_0,W_5\rangle=0$。

\SourcePage{104}
由于 $g_0\in W_0$，而 $\Omega(0)$ 张成 $F_{\mathrm{lim}}^3$，由式~\eqref{eq:mum-top-splitting} 以及杯积的非退化性，必有 $\langle g_0,\Omega(0)\rangle\ne0$。所以 $y_0(0)=\langle g_0,\Omega(0)\rangle$ 不为零。

举例说明命题~\ref{prop:mum-indicial-equation} 的用法：考虑\hyperref[chap:quintic]{第 2 章}的五次镜像族。由式~\eqref{eq:quintic-pf}，Picard--Fuchs 方程为
\begin{equation}\label{eq:quintic-mum-pf}
\begin{aligned}
0={}&\left(z\frac d{dz}\right)^4y
+\frac{2\cdot5^5z}{1+5^5z}\left(z\frac d{dz}\right)^3y
+\frac{7\cdot5^4z}{1+5^5z}\left(z\frac d{dz}\right)^2y\\
&+\frac{2\cdot5^4z}{1+5^5z}\left(z\frac d{dz}\right)y
+\frac{24\cdot5z}{1+5^5z}y.
\end{aligned}
\end{equation}
\hyperref[sec:hodge-examples]{第 5.4 节}将推导这个方程。这里的变量 $z$，在\hyperref[chap:quintic]{第 2 章}记为 $x$，等于 $z=\psi^{-5}$，其中 $\psi$ 是五次镜像的定义方程~\eqref{eq:quintic-pencil} 中的参数。因为指标方程为 $\lambda^4=0$，命题~\ref{prop:mum-indicial-equation} 表明 $z=0$ 处有最大幺幂单值性。

另一方面，五次镜像族还有两个边界点 $z=-5^{-5},\infty$。容易看出，$z=\infty$ 处的单值性不是最大幺幂的。在式~\eqref{eq:quintic-mum-pf} 中作坐标变换 $w=1/z$，算得指标方程的根为 $\{1/5,2/5,3/5,4/5\}$；由命题~\ref{prop:mum-indicial-equation}，无穷远点处的单值性甚至不是幺幂的。在适合 $z=-5^{-5}$ 的坐标中，指标方程的根为整数，因此单值性幺幂。不过，根并非全相等，所以根据该命题，单值性不可能最大幺幂。事实上，可以证明此时 $N^2=0$。

\subsection{$r$ 维模空间}\label{subsec:higher-dimensional-mum}

下面推广最大幺幂单值性的定义。对维数为 $d$ 的 Calabi--Yau 流形 $V$，令 $S$ 为 $V$ 的完整复模空间；$V$ 有奇点的情形将在\hyperref[chap:moduli]{第~\ref*{chap:moduli} 章}考虑。则局部形变空间 $S$ 是维数 $r=h^{d-1,1}(V)$ 的光滑流形，见 \SourceCite{Bogomolov,Tian,Todorov}。如\hyperref[sec:hodge-theory]{第 5.1 节}，令 $\overline S$ 是 $S$ 的紧化，边界除子 $D=\overline S-S$ 具有正规交叉。我们重点考察 $r=\dim S$ 个边界除子 $D_1,\ldots,D_r$ 的交点 $p\in\overline S-S$；\hyperref[chap:moduli]{第~\ref*{chap:moduli} 章}会更详细研究 $S$ 的自然紧化。我们关心以下类型的边界点。

\begin{definition}\label{def:maximal-unipotent-point}
\SourceCite{Morrison3} 若点 $p=D_1\cap\cdots\cap D_r$ 满足以下条件，则称 $p$ 为\emph{最大幺幂边界点}：
\begin{enumerate}[label=(\roman*)]
\item 单值性变换 $\mathcal T_i$ 全部幺幂；
\item 置 $N_i=\log\mathcal T_i$，令 $N=\sum_i a_iN_i$ 为所有 $a_i>0$ 的线性组合。若 $W_\bullet$ 是由 $N$ 确定的单值性权滤过~\eqref{eq:monodromy-weight-filtration}，则 $\dim W_0=\dim W_1=1$，$\dim W_2=1+r$；
\item 取 $W_2$ 的一组基 $g_0,\ldots,g_r$，使 $g_0$ 张成 $W_0$，并对 $1\leq i,j\leq r$ 用 $N_i(g_j)=m_{ij}g_0$ 定义 $m_{ij}$。则矩阵 $(m_{ij})$ 可逆。
\end{enumerate}
\end{definition}

此定义可解释如下。每个 $g_0\in W_0$ 给出 Picard--Fuchs 方程的解 $y_0=\langle g_0,\Omega\rangle$，它在最大幺幂边界点 $p$ 处解析。因此 $\dim W_0=1$ 是说，在相差常数倍的意义下，Picard--Fuchs 方程在 $p$ 处有唯一的全纯解。还应注意，每个 $g_i\in W_2$ 给出解 $\langle g_i,\Omega\rangle$，它沿各除子 $D_j$ 至多按对数增长。因此 $\dim W_2=1+r$ 意味着有 $r$ 个线性无关的、沿 $D_j$ 按对数增长的解。最后，条件 (iii) 中矩阵 $(m_{ij})$ 的可逆性使我们可以换到
\SourcePage{105}
新基
\begin{equation}\label{eq:mum-logarithmic-basis}
g_k'=\sum_{j=1}^r g_jm^{jk},\qquad
(m^{jk})=(m_{jk})^{-1},
\end{equation}
满足 $N_j(g_k')=\delta_{jk}g_0$。于是 $y_k=\langle g_k',\Omega\rangle$ 是 Picard--Fuchs 方程的解；它在 $j\ne k$ 时 $D_j$ 的一般点处全纯，而沿 $D_k$ 按对数增长。更精确地说，在 $p$ 附近可以局部写为
\[
y_k=y_0\frac{\log z_k}{2\pi i}+\text{全纯函数}.
\]
实际处理最大幺幂边界点时，寻找的正是这种情况。\hyperref[chap:moduli]{第~\ref*{chap:moduli} 章}将看到，镜像映射由函数 $y_0,y_1,\ldots,y_r$ 构造。

理解条件 (iii) 的另一种方式来自镜像对称。设 $V$ 与 $V^\circ$ 互为镜像。如\hyperref[sec:mirror-map]{第~\ref*{sec:mirror-map} 节}将看到，镜像映射把 $V^\circ$ 的 Kähler 模空间与 $V$ 的复模空间中最大幺幂边界点的一个邻域等同。在镜像一侧，\hyperref[chap:quantum-cohomology]{第~\ref*{chap:quantum-cohomology} 章}将用\IndexLocation{indexloc:1111:105}量子上同调在 $V^\circ$ 的 Kähler 模空间上构造平坦联络，称为\emph{A 模型联络}。局部地，此联络位于去心多圆盘 $(\Delta^*)^r$ 上，其坐标轴对应于 $V^\circ$ 的 Kähler 锥的一组生成元 $T_i$，$1\leq i\leq r$。每条轴产生此平坦联络的一个单值性变换 $\widetilde{\mathcal T}_i$，其对数 $\widetilde N_i=\log\widetilde{\mathcal T}_i$ 是 $\bigoplus_pH^{p,p}(V^\circ)$ 上与 $-T_i$ 作杯积的自同态。猜想镜像映射将 A 模型联络与 Gauss--Manin 联络等同，使得：
\begin{itemize}
\item $N_i$ 对应于 $\widetilde N_i=\mathord{\cdot}\cup(-T_i)$，因此 $N=\sum_i a_iN_i$ 对应于与丰沛类 $D=\sum_i a_iT_i$ 的相反数作杯积；
\item $g_0$ 对应于一个点的上同调类；
\item $g_i$ 对应于曲线 $C_i$ 的 Poincaré 对偶，其中各 $C_i$ 是 $H_2(V^\circ,\mathbb Z)$ 的一组合适的基。
\end{itemize}
于是，在镜像映射下，矩阵 $(m_{ij})$ 对应于交数矩阵 $(-T_i\cdot C_j)$；根据 Poincaré 对偶性，它必定可逆。因此，从镜像对称的角度看，定义~\ref{def:maximal-unipotent-point} 的第三个条件完全自然。

在最大幺幂单值性点处，同构 $N^d:W_{2d}/W_{2d-1}\simeq W_0$ 意味着 $N^d\ne0$；由单值性定理 $N^{d+1}=0$，所以 $N$ 是具有最大幂零指数的幂零算子。特别地，对复模为一维的三维簇，容易推出定义~\ref{def:maximal-unipotent-point} 等价于 $N^3\ne0$。更一般地，对维数为 $d$、复模为一维的 Calabi--Yau 簇，定义~\ref{def:maximal-unipotent-point} 等价于 $N^d\ne0$，见 \SourceCite{Morrison2}。

最后提到，权滤过的整数结构如何与定义~\ref{def:maximal-unipotent-point} 相联系。设在最大幺幂边界点处，$g_0,g_1,\ldots,g_r$ 是 $W_2$ 的整数格的一组 $\mathbb Z$ 基，其中 $g_0$ 生成 $W_0$ 的整数格。容易证明，由 $N_j(g_k)=m_{jk}g_0$ 定义的数 $m_{jk}$ 是整数。在 \SourceCite{Morrison6} 中，Morrison 猜想矩阵 $(m_{jk})$ 在 $\mathbb Z$ 上可逆；等价地，式~\eqref{eq:mum-logarithmic-basis} 中的 $m^{jk}$ 都是整数。这个\emph{\IndexLocation{indexloc:0637:105}\IndexLocation{indexloc:0972:105}整性猜想}已对五次三维簇的镜像得到验证 \SourceCite{Morrison2}；\SourceCite{Morrison6} 还讨论了猜想的其他情形。\index{最大幺幂单值性|hyperpage}\index{最大幺幂边界点|hyperpage}

\SourcePage{106}
\section{Griffiths--Dwork 方法}\label{sec:griffiths-dwork}

现在转向 Griffiths 与 Dwork 提出的计算 Picard--Fuchs 理想的方法。这种方法最初用于射影超曲面，近来又推广到加权射影空间中的超曲面，乃至更一般的射影环面簇中的超曲面。

\subsection{射影超曲面}\label{subsec:projective-residues}

从次数为 $\ell$ 的光滑超曲面 $V\subset\mathbb P^n$ 开始。有留数映射
\[
\operatorname{Res}:H^n(\mathbb P^n-V)\longrightarrow H^{n-1}(V),
\]
稍后将描述它。$H^n(\mathbb P^n-V)$ 中的元素可用 $\mathbb P^n$ 上沿 $V$ 有极点的全纯 $n$ 形式表示。为此定义
\[
\Omega_0=\sum_j(-1)^j x_j\,dx_0\wedge\cdots
\wedge\widehat{dx_j}\wedge\cdots\wedge dx_n.
\]
注意，$\Omega_0$ 是 $\Omega_{\mathbb P^n}^n(n+1)$ 的截面，而这个层是平凡层。这意味着，在相差常数倍的意义下，$\Omega_0$ 是唯一的次数为 $n+1$ 的齐次全纯 $n$ 形式。

若 $V$ 由次数为 $\ell$ 的齐次方程 $f=0$ 定义，则根据 \SourceCite{Griffiths1}，$H^n(\mathbb P^n-V)$ 中的元素可用形式
\[
\frac{P\Omega_0}{f^k},\qquad\deg P=k\ell-(n+1),
\]
表示，其中 $P$ 也齐次。对 $V$ 中任意拓扑 $(n-1)$ 圈 $\gamma$，令 $T(\gamma)$ 为 $\gamma$ 上的管，即包含在 $\mathbb P^n-V$ 中、位于 $\gamma$ 上的圆周丛。留数定义为
\begin{equation}\label{eq:residue-tube}
\int_\gamma\operatorname{Res}\frac{P\Omega_0}{f^k}
=\int_{T(\gamma)}\frac{P\Omega_0}{f^k}.
\end{equation}
向 $P\Omega_0/f^k$ 加上恰当形式，不改变式~\eqref{eq:residue-tube} 的右边，因此 $\operatorname{Res}$ 在上同调类上良定义。由于 $\ell H\sim V$，其中 $H$ 为超平面类，故 $\operatorname{Res}(P\Omega_0/f^k)\cdot H=0$。定义本原上同调
\[
PH^{n-1}(V)=\{\eta\in H^{n-1}(V):\eta\cdot H=0\}.
\]
若 $V$ 的维数为奇数，本原上同调等于全部中维上同调。一般而言，$\operatorname{Res}$ 的像包含在 $PH^{n-1}(V)$ 中；根据 \SourceCite{Griffiths1}，映射 $\operatorname{Res}:H^n(\mathbb P^n-V)\to PH^{n-1}(V)$ 满射。下文将遵循通常做法，常在记号中省去“$\operatorname{Res}$”，将 $P\Omega_0/f^k$ 视为 $V$ 上的上同调类或微分形式。

根据 \SourceCite{Griffiths1}，$P\Omega_0/f^k$ 实际位于 $F^{n-k}PH^{n-1}(V)$ 中。若让定义多项式 $f$ 与分子 $P$ 依赖参数 $s$ 来得到族 $V$，则计算得
\[
\nabla_{d/ds}\left(\frac{P\Omega_0}{f^k}\right)
=\frac{(-kPf'+fP')\Omega_0}{f^{k+1}}
\in\mathcal F^{n-k-1},
\]
其中撇号表示关于 $s$ 求导。这验证了射影超曲面的 Griffiths 横截性。

\SourcePage{107}
\SourceCite{Griffiths1} 给出了一个有效过程，用于模去恰当形式进行计算。对每个 $i$，令 $G_i$ 为次数 $(k-1)\ell-n$ 的齐次多项式，置
\[
\eta=\frac1{f^{k-1}}\sum_{i<j}(-1)^{i+j}(x_iG_j-x_jG_i)
\,dx_0\wedge\cdots\wedge\widehat{dx_i}\wedge\cdots
\wedge\widehat{dx_j}\wedge\cdots\wedge dx_n.
\]
则
\begin{equation}\label{eq:griffiths-pole-reduction}
d\eta=
\frac{(k-1)\bigl(\sum_jG_j\frac{\partial f}{\partial x_j}\bigr)\Omega_0}{f^k}
-\frac{\bigl(\sum_j\frac{\partial G_j}{\partial x_j}\bigr)\Omega_0}{f^{k-1}}.
\end{equation}
关键在于：利用式~\eqref{eq:griffiths-pole-reduction}，当 $P$ 位于 Jacobian 理想
\[
J(f)=\left(\frac{\partial f}{\partial x_0},\ldots,
\frac{\partial f}{\partial x_n}\right)
\]
中时，任何形式 $P\Omega_0/f^k$ 都可模去恰当形式，约化为极点阶数较低的形式。反过来也成立：若 $P\Omega_0/f^k$ 可模去恰当形式，约化为 $F^{n-k+1}$ 中的元素，则 $P$ 位于 Jacobian 理想。换言之，映射 $P\mapsto P\Omega_0/f^k$ 在相应的分次商上定义同构
\begin{equation}\label{eq:jacobian-hodge-isomorphism}
(S/J(f))_{k\ell-(n+1)}\simeq PH^{n-k,k-1}(V),
\end{equation}
其中 $S=\mathbb C[x_0,\ldots,x_n]$，下标表示次数为 $k\ell-(n+1)$ 的分次部分。因此，选择 $(S/J(f))_{k\ell-(n+1)}$ 的一组基之后，任何 $\deg P=k\ell-(n+1)$ 的形式 $P\Omega_0/f^k$，都可以模去恰当形式与极点阶数较低的形式，约化为来自这组基的形式的线性组合。

下面计算 Picard--Fuchs 方程。为简单起见，在 $V$ 为 $\mathbb P^n$ 中 Calabi--Yau 超曲面的情形说明方法，这意味着 $\ell=n+1$。此时 $\omega=\operatorname{Res}(\Omega_0/f)$ 是 $V$ 上的全纯 $(n-1)$ 形式。\emph{Griffiths--Dwork 方法}按如下步骤进行：
\begin{itemize}
\item 设 $V$ 由依赖参数 $s\in S$ 的单个方程 $f$ 定义，其中 $S$ 是一维的。当 $s$ 变化时，$\omega=\operatorname{Res}(\Omega_0/f)$ 是 $f$ 所定义的超曲面族 $V$ 上的全纯 $(n-1)$ 形式。假定 $f$ 对 $s$ 的依赖为多项式。
\item 为 $V$ 的本原上同调选择一组基，用一组便于计算的形式 $\omega_i=P_i\Omega_0/f^{k_i}$ 表示，$1\leq i\leq r$。这里 $r=\dim PH^{n-1}(V)$；$V$ 的维数为偶数时 $r=h^{n-1}(V)-1$，为奇数时 $r=h^{n-1}(V)$。
\item 反复对 $\omega$ 关于 $s$ 求导，得到 $\mathcal F^0$ 的截面
\[
\omega,\nabla_{d/ds}\omega,\ldots,\nabla_{d/ds}^r\omega.
\]
这些截面都可模去恰当形式，用选定的基表示。方法如下：从最高阶极点向下处理，在模 $J(f)$ 的意义下，寻找 $\nabla_{d/ds}^i\omega$ 的分子多项式与选定基的分子之间的关系；最方便的是在环 $\mathbb C(s)[x_0,\ldots,x_n]$ 中使用 Gröbner 基方法。再用式~\eqref{eq:griffiths-pole-reduction}，将求导所得的形式模去恰当形式，表示为基元素的线性组合加上极点阶数较低的形式 $\eta$。线性组合的系数是参数 $s$ 的有理函数。然后对 $\eta$ 重复这个过程。
\SourcePage{108}
\item 由于有 $r+1$ 个截面，而 $\mathcal F^0$ 的秩为 $r$，$\omega,\nabla_{d/ds}\omega,\ldots,\nabla_{d/ds}^r\omega$ 之间必定存在一个系数属于 $\mathbb C(s)$ 的线性关系。这个关系给出 Picard--Fuchs 理想中的一个元素，所得常微分方程就是 Picard--Fuchs 方程。
\end{itemize}
下一节将把 Griffiths--Dwork 方法用于五次三维簇的镜像。此外，当超曲面的定义方程 $f$ 依赖多个参数时，容易调整这个方法，得到 Picard--Fuchs 理想中的元素。

遵循 \SourceCite{Dolgachev2}，不难把 Griffiths--Dwork 方法推广到加权射影空间中的超曲面。若加权射影空间为 $\mathbb P(q_0,\ldots,q_n)$，有
\[
\Omega_0=\sum_{j=0}^n(-1)^{j-1}q_jx_j\,dx_0\wedge\cdots
\wedge\widehat{dx_j}\wedge\cdots\wedge dx_n
\]
其余修改都很直接。下一节将给出说明这个方法的例子。
\index{Griffiths--Dwork方法@Griffiths--Dwork 方法|hyperpage}\index{本原上同调|hyperpage}\index{留数映射|hyperpage}

\subsection{环面超曲面}\label{subsec:toric-residues}

现在说明 Griffiths--Dwork 方法如何应用于环面簇中的丰沛超曲面。我们遵循 \SourceCite{BC} 的处理。设 $X=X_\Sigma$ 是一个 $n$ 维射影单纯环面簇，记 $\Sigma(1)=\{\rho_1,\ldots,\rho_r\}$，$S=\mathbb C[x_1,\ldots,x_r]$ 为~\hyperref[subsec:homogeneous-ring]{第 3.2.3 小节}定义的 $X$ 的齐次坐标环。取 $\beta\in A_{n-1}(X)$，并假定 $\beta$ 是 $X$ 上某个丰沛线丛的第一 Chern 类。于是 $S_\beta$ 是次数为 $\beta$ 的齐次多项式空间，这里的齐次性采用~\hyperref[subsec:homogeneous-ring]{第 3.2.3 小节}的含义：回顾 $S$ 按 $A_{n-1}(X)$ 分次。一般的 $f\in S_\beta$ 定义一个拟光滑超曲面 $V=\{f=0\}$，其除子类为 $\beta$。

Griffiths--Dwork 方法的思路可以自然地推广到这里。首先需要定义合适的形式 $\Omega_0$。可取 $\Omega_0$ 为 $\widehat\Omega_X^n$ 的次数为 $\beta_0=\sum_i[D_i]$ 的齐次截面，因为
\[
\widehat\Omega_X^n=\mathcal O_X\Bigl(-\sum_iD_i\Bigr).
\]
这里 $D_i$ 是对应于 $\rho_i$ 的环面除子。若记 $\beta_i=\deg x_i$，则 $\beta_0=\sum_i\beta_i$。

为给出 $\Omega_0$ 的显式公式，固定 $M$ 的一组基 $e_1,\ldots,e_n$。对指标集 $I=\{i_1,\ldots,i_n\}\subset\{1,\ldots,r\}$，定义
\[
\det(e_I)=\det\bigl(\langle e_j,\rho_{i_k}\rangle\bigr)_{1\leq j,k\leq n},
\]
其中也用 $\rho_i$ 表示相应射线的本原生成元。置
\[
dx_I=dx_{i_1}\wedge\cdots\wedge dx_{i_n},\qquad
\widehat x_I=\prod_{i\notin I}x_i.
\]
于是定义
\begin{equation}\label{eq:toric-volume-form}
\Omega_0=\sum_{|I|=n}\det(e_I)\widehat x_I\,dx_I.
\end{equation}
注意，$\Omega_0$ 的次数为 $\beta_0$。此外，如 \SourceCite{BC} 所解释的，这个表达式确实在 $X$ 上定义了一个带有相应扭曲的 $n$ 形式，而不只是 $\mathbb C^r$ 上的形式。于是，$X$ 上沿 $V$ 有极点的 $n$ 形式可写成
\[
\frac{P\Omega_0}{f^k},\qquad \deg P=k\beta-\beta_0.
\]

\SourcePage{109}
为降低极点阶数，对每个 $1\leq i\leq r$，需要一个次数为 $\beta_0-\beta_i$ 的 $(n-1)$ 形式 $\Omega_i$；这里不讨论其精确定义。关键是以下公式：
\begin{equation}\label{eq:toric-pole-reduction}
d\left(\frac{P\Omega_i}{f^k}\right)
=\frac{\bigl(f\frac{\partial P}{\partial x_i}
-kP\frac{\partial f}{\partial x_i}\bigr)\Omega_0}{f^{k+1}}.
\end{equation}
在这里，定义 Jacobian 理想
\[
J(f)=\left\langle\frac{\partial f}{\partial x_1},\ldots,
\frac{\partial f}{\partial x_r}\right\rangle,
\]
可以结合 Gröbner 基方法与 Jacobian 理想，降低极点阶数，从而将任意形式模去恰当形式后，用一组给定的基表示。

在环面情形中，我们采用一个稍有不同的本原上同调概念：
\[
PH^{n-1}(V)=H^{n-1}(V)\big/\operatorname{im}H^{n-1}(X).
\]
当 $V$ 为奇数维时，这就是通常的上同调。同样有留数映射
\[
\operatorname{Res}:H^n(X-V)\longrightarrow PH^{n-1}(V).
\]
在经典情形中它是同构，但在环面情形中有正合列
\[
0\longrightarrow H^{n-2}(X)\xrightarrow{\,\smile[V]\,}H^n(X)
\longrightarrow H^n(X-V)\xrightarrow{\operatorname{Res}}
PH^{n-1}(V)\longrightarrow0.
\]
由于 $H^n(X)=H^{n/2,n/2}(X)$，式~\eqref{eq:jacobian-hodge-isomorphism} 相应地变为同构
\[
(S/J(f))_{k\beta-\beta_0}\simeq PH^{n-k,k-1}(V),
\qquad k\ne n/2+1.
\]
而当 $k=n/2+1$ 时，有正合列
\[
\begin{split}
0\longrightarrow H^{n-2}(X)\longrightarrow H^n(X)
\longrightarrow (S/J(f))_{(n/2+1)\beta-\beta_0}\\
\longrightarrow PH^{n/2-1,n/2}(V)\longrightarrow0.
\end{split}
\]
（见 \SourceCite{BC}）。注意，只有在 $n$ 为偶数时才需要考虑这个正合列；此时超曲面 $V$ 为奇数维，例如 Calabi--Yau 三维簇。

也可以稍微改变 $J(f)$，从而避免使用这个正合列。用以下理想商代替 Jacobian 理想：
\[
J_1(f)=\left\langle x_1\frac{\partial f}{\partial x_1},\ldots,
x_r\frac{\partial f}{\partial x_r}\right\rangle:(x_1\cdots x_r).
\]
根据 \SourceCite{BC}，若 $f$ 非退化，即 $V$ 拟光滑且与 $X$ 的每个轨道拟横截地相交，则对所有 $k$ 都有
\[
(S/J_1(f))_{k\beta-\beta_0}\simeq PH^{n-k,k-1}(V).
\]
为解释理想 $J_1(f)$ 的出现，注意算子 $x_i\partial/\partial x_i$ 在环面作用下不变，而取理想商是使 $J(f)\subset J_1(f)$ 成立所必需的。此外，$J_1(f)$ 也自然出现在规范线性 $\sigma$ 模型的研究中，见 \SourceCite{MP1}。

因此，使用 $J(f)$（$n$ 为奇数时）或 $J_1(f)$（任意 $n$），就能在任何射影单纯环面簇中，对丰沛超曲面应用 Griffiths--Dwork 方法来计算 Picard--Fuchs 方程。

遗憾的是，这个方法要求 $V$ 丰沛且 $X$ 单纯。例如，在 Batyrev 镜像构造的情形中，回顾~\hyperref[sec:batyrev-mirror]{第 4.1 节}：从自反多胞形 $\Delta$ 出发，$\overline V$ 是 $\mathbb P_\Delta$ 中的丰沛反典范超曲面；对 $\Delta$ 的法扇作极大射影细分 $\Sigma$，则 $V\subset X_\Sigma$ 是 $\overline V$ 的严格变换，也是一个极小 Calabi--Yau 簇。
\SourcePage{110}
虽然 $\overline V$ 丰沛，$\mathbb P_\Delta$ 却未必单纯；而虽然 $X_\Sigma$ 单纯，$V$ 却未必丰沛。所以，这个方法可能对两者都不适用。\index{Jacobian理想@Jacobian 理想|hyperpage}

\section{例题}\label{sec:hodge-examples}

现在给出两个例子，说明如何使用上一节的方法计算 Picard--Fuchs 方程。

\begin{example}\label{ex:quintic-pf-method}
首先遵循 \SourceCite{Morrison1} 的思路，考虑五次镜像族。由式~\eqref{eq:quintic-pencil}，五次镜像由方程
\begin{equation}\label{eq:quintic-pf-family}
f=x_1^5+x_2^5+x_3^5+x_4^5+x_5^5+\psi x_1x_2x_3x_4x_5=0
\end{equation}
模去式~\eqref{eq:greene-plesser-group} 中给出的群 $G$ 的作用而得到。更精确地说，如\hyperref[sec:quintic-revisited]{第 4.2 节}所述，若 $W\subset\mathbb P^4$ 是式~\eqref{eq:quintic-pf-family} 定义的超曲面，则 $V^\circ$ 是
$W/G\subset\mathbb P^4/G=\mathbb P_{\Delta^\circ}$ 的消解，其中 $\Delta^\circ$ 的顶点由式~\eqref{eq:quintic-polar-vertices} 给出。再由\hyperref[chap:mirror-constructions]{第~\ref*{chap:mirror-constructions} 章}的式~\eqref{eq:batyrev-hn21}，有
\begin{equation}\label{eq:example-polynomial-hodge}
\begin{split}
h^{2,1}(V^\circ)
&=l(\Delta^\circ)-5-\sum_{\Gamma^\circ}l^*(\Gamma^\circ)
+\sum_{\Theta^\circ}l^*(\Theta^\circ)l^*(\widehat\Theta^\circ)\\
&=h_{\mathrm{poly}}^{2,1}(V^\circ)
+\sum_{\Theta^\circ}l^*(\Theta^\circ)l^*(\widehat\Theta^\circ).
\end{split}
\end{equation}
这里 $\Gamma^\circ$ 遍历 $\Delta^\circ$ 的余维 $1$ 的面，$\Theta^\circ$ 遍历其余维 $2$ 的面，$\widehat\Theta^\circ$ 是 $\Delta$ 中与 $\Theta^\circ$ 对偶的一维面。容易算出，$\Delta^\circ$ 的整点只有式~\eqref{eq:quintic-polar-vertices} 所列的五个顶点和原点 $(0,0,0,0)$，共六个，因而
\[
h^{2,1}(V^\circ)=h_{\mathrm{poly}}^{2,1}(V^\circ)=6-5-0+0=1.
\]
所以 $V^\circ$ 的复模空间为一维，来自式~\eqref{eq:quintic-pf-family} 中参数 $\psi$ 的变化。回顾\hyperref[chap:quintic]{第~\ref*{chap:quintic} 章}，$z=\psi^{-5}$ 是模空间的一个局部参数；\hyperref[chap:quintic]{第~\ref*{chap:quintic} 章}中记作 $x=\psi^{-5}$ 的变量，现在改记为 $z$。

因为 $V^\circ$ 是 $\mathbb P_{\Delta^\circ}$ 的一个环面爆破中的非丰沛除子，无法对它应用 Griffiths--Dwork 方法。幸而 $W/G$ 丰沛，且 $\mathbb P_{\Delta^\circ}$ 单纯，所以改取 $V^\circ=W/G$，因为消解不改变 Picard--Fuchs 方程。这样便可应用\hyperref[sec:griffiths-dwork]{第 5.3 节}的方法，只须用 $\mathbb P^4-W$ 上的 $G$ 不变形式 $P\Omega_0/f^k$ 表示 $V^\circ$ 上的上同调类。

从 $G$ 不变三形式 $\omega_1=\psi\Omega_0/f$ 开始，它已在\hyperref[chap:quintic]{第 2 章}的式~\eqref{eq:quintic-residue} 中出现。目标是找出
$\omega_1,\delta\omega_1,\ldots,\delta^4\omega_1$ 之间的关系，其中
$\delta=z\,d/dz=-\frac15\psi\,d/d\psi$。为此，需要一组上同调的基。考虑以下 $G$ 不变形式：
\[
\omega_j=(-1)^{j-1}\psi^j(x_1x_2x_3x_4x_5)^{j-1}
\frac{\Omega_0}{f^j},\qquad j=1,\ldots,4.
\]
选择因子 $\psi^j$，是为了使 $\omega_j$ 在~\hyperref[sec:quintic-mirror]{第 2.2 节}所述作用
$x_1\mapsto\mu^{-1}x_1$、$\psi\mapsto\mu\psi$（$\mu^5=1$）下不变。这个变换诱导同构 $V^\circ_\psi\simeq V^\circ_{\mu\psi}$；特别地，这说明 Picard--Fuchs 方程可以用 $z=\psi^{-5}$ 改写。
注意，$\operatorname{Res}(\omega_j)\in F^{4-j}H^3(V^\circ)$，一次简单的 Gröbner 基计算表明，这个类在 $H^{4-j,j-1}(V^\circ)$ 中的投影非零。由于 $h^{3,0}=h^{2,1}=h^{1,2}=h^{0,3}=1$，$\omega_j$（$j=1,2,3,4$）的留数组成 $H^3(V^\circ)$ 的一组基。

\SourcePage{111}
因此，对 $k=1,\ldots,4$，任何极点阶数为 $k$ 的 $G$ 不变形式，都可以模去恰当形式和极点阶数较低的形式，约化为 $\omega_k$ 的倍数；而具有任意阶极点的 $G$ 不变形式，都可以模去恰当形式，表示为这四个形式的线性组合。可以结合式~\eqref{eq:griffiths-pole-reduction} 中的极点约化公式与 Gröbner 基方法，完成这些计算。

现在把这个过程用于 $\omega_1,\delta\omega_1,\ldots,\delta^4\omega_1$。对任意 $j$，有
\begin{equation}\label{eq:quintic-basis-derivative}
\delta\omega_j=-\frac j5(\omega_j+\omega_{j+1}).
\end{equation}
由此立刻得到
\begin{equation}\label{eq:quintic-derivative-basis}
\begin{pmatrix}\omega_1\\\delta\omega_1\\\delta^2\omega_1\\\delta^3\omega_1\end{pmatrix}
=\begin{pmatrix}
1&0&0&0\\
-1/5&-1/5&0&0\\
1/25&3/25&2/25&0\\
-1/125&-7/125&-12/125&-6/125
\end{pmatrix}
\begin{pmatrix}\omega_1\\\omega_2\\\omega_3\\\omega_4\end{pmatrix}.
\end{equation}
这样便处理了 $\omega_1,\delta\omega_1,\delta^2\omega_1,\delta^3\omega_1$。至于 $\delta^4\omega_1$，反复进行 Gröbner 基计算并降低极点阶数，模去恰当形式后得到
\[
\delta^4\omega_1
=\frac{5z}{1+5^5z}\omega_1+\frac{75z}{1+5^5z}\omega_2
\quad+\frac{250z}{1+5^5z}\omega_3+\frac{300z}{1+5^5z}\omega_4.
\]
因此，
\[
\delta^4\omega_1
=\begin{pmatrix}
\dfrac{5z}{1+5^5z}&\dfrac{75z}{1+5^5z}&
\dfrac{250z}{1+5^5z}&\dfrac{300z}{1+5^5z}
\end{pmatrix}
\begin{pmatrix}\omega_1\\\omega_2\\\omega_3\\\omega_4\end{pmatrix}.
\]
将式~\eqref{eq:quintic-derivative-basis} 代入，便得到 Picard--Fuchs 方程~\eqref{eq:quintic-mum-pf}。
\end{example}

\begin{example}\label{ex:weighted-pf-method}
现在考虑 $\mathbb P(1,1,2,2,2)$ 的一个爆破中的 Calabi--Yau \IndexLocation{indexloc:1020:111}超曲面的镜像。稍后将看到，镜像的复模空间是二维的，这意味着 Picard--Fuchs 方程是偏微分方程。

对于 $\mathbb P(1,1,2,2,2)$ 的标准扇，一维锥的生成元是
\begin{equation}\label{eq:weighted-pf-rays}
\{(-1,-2,-2,-2),(1,0,0,0),(0,1,0,0),(0,0,1,0),(0,0,0,1)\}.
\end{equation}
这些点是一个\IndexLocation{indexloc:1222:111}自反多胞形 $\Delta^\circ\subset N_{\mathbb R}$ 的顶点。根据~\hyperref[sec:quintic-revisited]{第 4.2 节}的讨论，$\mathbb P_{\Delta^\circ}$ 是加权射影空间 $\mathbb P(1,1,2,2,2)$ 对某个有限群 $H$ 的商。
$\Delta^\circ$ 中只有七个整点：五个顶点、点
\[
(0,-1,-1,-1)=\tfrac12(-1,-2,-2,-2)+\tfrac12(1,0,0,0)
\]
以及原点。由式~\eqref{eq:example-polynomial-hodge}，可得
$h^{2,1}(V^\circ)=h_{\mathrm{poly}}^{2,1}(V^\circ)=2$。所以 $V^\circ$ 有二维复模；而且这些模是多项式模，也就是说，可以通过改变 $V^\circ$ 的定义方程来检测它们。

如前所述，可以假定 $V^\circ$ 是 $\mathbb P_{\Delta^\circ}$ 中的反典范超曲面。因为它丰沛，且 $\mathbb P_{\Delta^\circ}$ 单纯，所以 Griffiths--Dwork 方法适用。

\SourcePage{112}
由于 $\Delta^\circ$ 是式~\eqref{eq:weighted-pf-rays} 中各点的凸包，计算得极多胞形 $\Delta$ 是以下各点的凸包：
\begin{equation}\label{eq:weighted-pf-polar}
\begin{split}
\{&(-1,-1,-1,-1),(7,-1,-1,-1),(-1,3,-1,-1),\\
&(-1,-1,3,-1),(-1,-1,-1,3)\}.
\end{split}
\end{equation}
按给定顺序将 $\Delta$ 的顶点记为 $v_i$。这些 $v_i$ 生成 $\mathbb P_{\Delta^\circ}$ 的扇的一维锥。注意，$\mathbb P_\Delta=\mathbb P(1,1,2,2,2)$；若 $\Sigma$ 是 $\Delta$ 的法扇的唯一极大射影细分，则 $X_\Sigma$ 是 $\mathbb P(1,1,2,2,2)$ 的一次爆破，具体由向式~\eqref{eq:weighted-pf-rays} 确定的扇中插入射线 $(0,-1,-1,-1)$ 得到。由此可知，$V^\circ$ 是本例第一句话中所描述的族。

$\mathbb P_{\Delta^\circ}$ 的齐次坐标环就是 $\mathbb C[x_1,\ldots,x_5]$，按 $A_3(\mathbb P_{\Delta^\circ})$ 分次；这里这个群含有挠元。需要确定次数
$\beta_0=\sum_{i=1}^5\deg(x_i)\in A_3(\mathbb P_{\Delta^\circ})$ 的单项式。具体而言，它们是在
$\mathbb P_{\Delta^\circ}=\mathbb P(1,1,2,2,2)/H$ 中次数为 $8$ 的 $H$ 不变单项式。最简便的办法是使用同构~\eqref{eq:toric-graded-sections}；由此得出，所求单项式为
\[
\prod_i x_i^{\langle v,v_i\rangle+1},\qquad v\in\Delta^\circ\cap N.
\]
由于已经知道 $\Delta^\circ\cap N$ 的各点，所以得到单项式
\[
x_1^8,\quad x_2^8,\quad x_3^4,\quad x_4^4,\quad x_5^4,
\quad x_1^4x_2^4,\quad x_1x_2x_3x_4x_5.
\]
它们恰好是加权射影空间中权为 $8$ 的 $H$ 不变单项式。从加权射影空间的角度看，这正是 Greene--Plesser 轨形构造。一般地，可以证明：对允许 Fermat 超曲面的加权射影空间，Batyrev 镜像总与 Greene--Plesser 构造一致（见 \SourceCite{BK2}）。

因此，$\mathbb P_{\Delta^\circ}$ 中的一般反典范超曲面由这些 $H$ 不变单项式的线性组合给出。利用环面作用缩放变量 $x_i$，可以将反典范超曲面的方程写成
\begin{equation}\label{eq:weighted-pf-family}
f=z_2x_1^8+x_2^8+z_1x_3^4+x_4^4+x_5^4
+x_1x_2x_3x_4x_5+x_1^4x_2^4=0,
\end{equation}
其中 $z_1,z_2$ 是参数。之所以这样选择，是为了使 $z_1=z_2=0$ 为\IndexLocation{indexloc:0769:112}最大幺幂边界点。记 $\delta_i=z_i\,d/dz_i$。Hodge 数为
$h^{3,0}=h^{0,3}=1$，以及前面算出的 $h^{2,1}=h^{1,2}=2$。选择以下基：
\[
\begin{aligned}
\omega_1&=\frac{\Omega_0}{f},\\
\omega_2&=\delta_1\omega_1=-\frac{z_1x_3^4\Omega_0}{f^2},\\
\omega_3&=\delta_2\omega_1=-\frac{z_2x_1^8\Omega_0}{f^2},\\
\omega_4&=\delta_1^2\omega_1=\omega_2+\frac{2z_1^2x_3^8\Omega_0}{f^3},\\
\omega_5&=\delta_1\delta_2\omega_1=\frac{2z_1z_2x_1^8x_3^4\Omega_0}{f^3},\\
\omega_6&=\delta_1^2\delta_2\omega_1
=\omega_5-\frac{6z_1^2z_2x_1^8x_3^8\Omega_0}{f^4}.
\end{aligned}
\]
将各种 $\delta_i^j\omega_1$ 用这组基表示，就能得到 Picard--Fuchs 方程。

\SourcePage{113}
例如，计算得
\[
\delta_2^2\omega_1=\omega_3+\frac{2z_2^2x_1^{16}\Omega_0}{f^3}.
\]
相应地，尝试利用 Jacobian 理想，将 $x_1^{16}$ 约化为只含 $x_3^8$ 和 $x_1^8x_3^4$ 的表达式。使用 Gröbner 基方法，得到关系
\[
z_2(1-4z_2)x_1^{16}-z_1^2x_3^8+4z_1z_2x_1^8x_3^4\in J(f).
\]
事实上，容易验证
\[
z_2(1-4z_2)x_1^{16}-z_1^2x_3^8+4z_1z_2x_1^8x_3^4
=-\frac1{16}\left(\sum_{i=1}^3A_ix_i\frac{\partial f}{\partial x_i}\right),
\]
其中
\[
\begin{aligned}
A_1&=2(4z_2-1)x_1^8+x_1x_2x_3x_4x_5-4x_1^4x_2^4,\\
A_2&=2x_1^8,\\
A_3&=4z_1x_3^4-16z_2x_1^8-x_1x_2x_3x_4x_5.
\end{aligned}
\]
将它与式~\eqref{eq:toric-pole-reduction} 结合，便能降低 $\delta_2^2\omega_1$ 的极点阶数。反复进行这个过程，继续降低极点阶数，最终得到 Picard--Fuchs 方程
\begin{equation}\label{eq:weighted-pf-second}
(1-4z_2)\delta_2^2\omega_1
-z_2(\delta_1^2-4\delta_1\delta_2-\delta_1+2\delta_2)\omega_1=0.
\end{equation}

\SourceCite{CdFKM} 用其他方法详细研究了这个例子，所用方程~\eqref{eq:weighted-pf-second} 的形式略有不同（另见 \SourceCite{HKTY1}）。如前所述，可以证明 $z_1=z_2=0$ 是一个最大幺幂边界点。为了详细验证，还需要另一个 Picard--Fuchs 方程。这里暂且只指出，从式~\eqref{eq:weighted-pf-second} 的形式很快可得 $N_2^2=0$，所以 $N_2$ 单独不能定义最大幺幂边界点的单值性权滤过。但是，任意正线性组合 $a_1N_1+a_2N_2$ 都满足
$(a_1N_1+a_2N_2)^3\ne0$，并且确实定义了所需类型的单值性权滤过。

我们将在例题~\ref{ex:weighted-gkz} 中再次讨论这个例子，并在下一章继续讨论。注意，具有正规交叉边界的模空间光滑紧化仍然只有一个最大幺幂边界点。各对边界除子的其他交点都不是最大幺幂的。
\end{example}

\section{超几何方程}\label{sec:hypergeometric}

现在描述 Gelfand、Kapranov 与 Zelevinsky 的超几何方程组，并讨论镜像对称中自然出现的几个例子。

\subsection{$\mathcal A$ 方程组}\label{subsec:a-system}

设 $\mathcal A=\{v_1,\ldots,v_k\}\subset\mathbb Z^{n+1}$ 是位于同一个整仿射超平面中的 $k>n+1$ 个点。将 $v_i$ 的坐标记为 $v_{ij}$。对每个 $v_i\in\mathcal A$ 引入变量 $\lambda_i$，并固定向量
$\widehat\beta=(\widehat\beta_1,\ldots,\widehat\beta_{n+1})\in\mathbb C^{n+1}$。这些数据确定一个关于 $\lambda_i$ 的微分方程组，称为\emph{超几何方程的 $\mathcal A$ 方程组}。用 $\partial_i$ 表示微分算子 $\partial/\partial\lambda_i$，用 $\delta_i$ 表示对数导数 $\lambda_i\partial_i=\lambda_i\partial/\partial\lambda_i$。

$\mathcal A$ 方程组由两种微分算子构成，记为 $Z_j$ 与 $\Box_\ell$。算子 $Z_j$ 定义为
\begin{equation}\label{eq:gkz-euler-operators}
Z_j=\left(\sum_i v_{ij}\delta_i\right)-\widehat\beta_j,
\qquad j=1,\ldots,n+1.
\end{equation}

\SourcePage{114}
为描述 $\Box_\ell$，还需要一些记号。令 $\Lambda$ 为 $\mathcal A$ 中各元素的关系格：
\begin{equation}\label{eq:gkz-relation-lattice}
\Lambda=\left\{\ell=(\ell_i)\in\mathbb Z^k:\sum_i\ell_iv_i=0\right\}.
\end{equation}
根据关于 $\mathcal A$ 的假设，$\Lambda\ne\{0\}$。然后定义
\begin{equation}\label{eq:gkz-box-operators}
\Box_\ell=\prod_{\ell_i>0}\partial_i^{\ell_i}
-\prod_{\ell_i<0}\partial_i^{-\ell_i},\qquad\ell\in\Lambda.
\end{equation}

\begin{definition}\label{def:gkz-system}
\emph{超几何方程的 $\mathcal A$ 方程组}是微分方程组
$Z_j\Phi=0$（$1\leq j\leq n+1$）与 $\Box_\ell\Phi=0$（$\ell\in\Lambda$），其中 $Z_j$ 与 $\Box_\ell$ 分别由式~\eqref{eq:gkz-euler-operators} 与式~\eqref{eq:gkz-box-operators} 定义。
\end{definition}
根据 Gelfand、Kapranov 与 Zelevinsky 的工作 \SourceCite{GKZ1}，这个方程组通常也称为 \emph{GKZ 方程组}。\index{超几何方程|hyperpage}\index{GKZ方程组@GKZ 方程组|hyperpage}

\hypertarget{source:114-batch3}{}
\subsection{Calabi--Yau 环面超曲面}\label{subsec:toric-gkz}

现在设 $\Delta\subset M_{\mathbb R}$ 是一个 $n$ 维自反多胞形，$X_\Sigma$ 是对应于 $\Delta$ 的法扇的一个极大射影细分 $\Sigma$ 的环面簇。考虑 $X_\Sigma$ 中的反典范超曲面族；由命题~\ref{prop:toric-cy-construction}，这些超曲面是极小 Calabi--Yau 子簇。我们的任务是为这个族求出 Picard--Fuchs 方程。如前所述，这些方程与细分的选择无关，而且与 $\mathbb P_\Delta$ 中反典范超曲面族 $V$ 的方程相同；由命题~\ref{prop:toric-cy-construction}，后者也是 Calabi--Yau 簇族。因此，只需研究 $\mathbb P_\Delta$ 中的这个 Calabi--Yau 超曲面族。

将 $\Delta\cap M$ 中的点 $m_i$ 排序，使 $m_0=0$。由于 $\Delta$ 自反，反典范丛的截面可以与以下 Laurent 多项式等同：
\begin{equation}\label{eq:toric-gkz-laurent}
f=\sum_i\lambda_it^{m_i}=\lambda_0+\sum_{i>0}\lambda_it^{m_i}\in L(\Delta),
\end{equation}
其中 $\lambda_i$ 是参数。于是，Calabi--Yau 簇族 $V\subset\mathbb P_\Delta$ 由方程 $f=0$ 定义，其中 $f\in L(\Delta)$。

为了将它与超几何方程联系起来，令
\begin{equation}\label{eq:toric-gkz-configuration}
\mathcal A=(\Delta\cap M)\times\{1\}\subset M\times\mathbb Z\simeq\mathbb Z^{n+1},
\end{equation}
并置 $\widehat\beta=(0,\ldots,0,-1)$。注意，$\ell\in\Lambda$ 同时意味着 $\sum_i\ell_im_i=0$ 与 $\sum_i\ell_i=0$。后一个等式说明，式~\eqref{eq:gkz-box-operators} 中的 $\Box_\ell$ 是齐次微分算子。

我们断言，$\mathcal A$ 方程组中的方程是 Picard--Fuchs 方程；也就是说，$V$ 上某些微分形式的周期满足超几何方程的 $\mathcal A$ 方程组。

遵循 \SourceCite{Batyrev3}，使用 $n$ 形式
\begin{equation}\label{eq:toric-gkz-rational-form}
\omega=\frac1f\eta,\qquad
\eta=\frac{dt_1}{t_1}\wedge\cdots\wedge\frac{dt_n}{t_n},
\end{equation}
其中 $(t_1,\ldots,t_n)$ 是环面 $T\subset\mathbb P_\Delta$ 的坐标。对 $\ell=(\ell_i)\in\Lambda$ 以及一个指标子集 $I$，令 $\ell_I=\sum_{i\in I}\ell_i$，计算得
\begin{equation}\label{eq:gkz-form-derivatives}
\left(\prod_{i\in I}\partial_i^{\ell_i}\right)\omega
=\frac{(-1)^{\ell_I}(\ell_I)!\prod_{i\in I}t^{\ell_im_i}}{f^{\ell_I+1}}\eta.
\end{equation}
这里在求导的指标集 $I$ 上取 $\ell_i\geq0$。由于 $\ell\in\Lambda$ 意味着
$\sum_{\ell_i>0}\ell_im_i=\sum_{\ell_i<0}(-\ell_i)m_i$，故有
\[
\prod_{\ell_i>0}t^{\ell_im_i}=\prod_{\ell_i<0}t^{-\ell_im_i}.
\]
\SourcePage{115}
再由式~\eqref{eq:gkz-form-derivatives} 和
$\sum_{\ell_i>0}\ell_i=\sum_{\ell_i<0}(-\ell_i)$，容易推出 $\Box_\ell\omega=0$。若 $\Phi$ 是 $\omega$ 的任意周期，立即得到 $\Box_\ell\Phi=0$。

同样容易证明，对 $1\leq j\leq n+1$ 有 $Z_j\Phi=0$。考虑 $(n-1)$ 形式
\[
\eta_j=\frac{dt_1}{t_1}\wedge\cdots\wedge
\widehat{\frac{dt_j}{t_j}}\wedge\cdots\wedge\frac{dt_n}{t_n},
\qquad 1\leq j\leq n.
\]
计算得
\[
d\left(\frac1f\eta_j\right)
=\frac{(-1)^j}{f^2}\left(\sum_i\lambda_im_{ij}t^{m_i}\right)\eta,
\]
其中 $m_{ij}$ 是 $m_i$ 的坐标。因为对 $1\leq j\leq n$ 有 $\widehat\beta_j=0$，容易得出 $Z_j\omega$ 对 $j<n+1$ 是恰当形式，因而这些 $j$ 都满足 $Z_j\Phi=0$。至于 $j=n+1$，有 $\widehat\beta_{n+1}=-1$，而 $(m_i,1)\in\mathcal A$ 的最后一个坐标总为 $1$，所以
\[
Z_{n+1}\omega=\left(\sum_i\lambda_i\partial_i\right)\omega+\omega
=-\frac{\sum_i\lambda_it^{m_i}}{f^2}\eta+\omega=0,
\]
正如所需。

式~\eqref{eq:toric-gkz-rational-form} 中的 $\omega$ 定义在 $T-Z_f$ 上，其中 $T\subset\mathbb P_\Delta$ 是环面，$Z_f=V\cap T$ 是由 $f\in L(\Delta)$ 确定的仿射超曲面。于是 $\omega$ 的留数是 $Z_f$ 上的 $(n-1)$ 形式。\SourceCite{Batyrev3} 更详细地解释了这一仿射观点。不过，我们需要的是 $V$ 上周期的 Picard--Fuchs 方程，而不仅是仿射部分 $Z_f$ 上的方程。这可以借助 $\mathbb P_\Delta$ 的齐次坐标环 $\mathbb C[x_1,\ldots,x_r]$ 实现。$x_i$ 与 $t_j$ 之间的关系为
\[
t_j=\prod_{i=1}^r x_i^{\langle e_j,\rho_i\rangle},
\]
在这个代换下，有
\[
\Omega_0=x_1\cdots x_r\eta,
\]
其中 $\Omega_0$ 由式~\eqref{eq:toric-volume-form} 定义。这些事实的证明见 \SourceCite{BC}。$\Omega_0$ 的公式容易理解：形式 $\eta$ 在 $\mathbb P_\Delta-T$ 的各分支 $D_1,\ldots,D_r$ 上有一阶极点。另一方面，对应于 $0\in\Delta$ 的 $\mathcal O_{\mathbb P_\Delta}(\sum_{j=1}^rD_j)$ 的截面 $x_1\cdots x_r$，在每个环面除子 $D_j$ 上恰有一阶零点；由
$\widehat\Omega_{\mathbb P_\Delta}^n\simeq\mathcal O_{\mathbb P_\Delta}(-\sum_{j=1}^rD_j)$，便得到 $\Omega_0=x_1\cdots x_r\eta$。因此，在上述变量代换下，有
\[
\omega=\frac{\eta}{f(t_1,\ldots,t_n)}
=\frac{x_1\cdots x_r\eta}{x_1\cdots x_r f(t_1,\ldots,t_n)}
=\frac{\Omega_0}{f(x_1,\ldots,x_r)},
\]
其中 $f(x_1,\ldots,x_r)$ 是按式~\eqref{eq:toric-graded-sections} 定义的 $f(t_1,\ldots,t_n)$ 的齐次化。因此，$\omega$ 是~\hyperref[sec:griffiths-dwork]{第 5.3 节}考虑过的形式之一，其留数是 $V$ 上的全纯 $(n-1)$ 形式。我们得出结论：$\mathcal A$ 超几何方程是 Picard--Fuchs 方程。这里还有一个技术问题，即 $\mathbb P_\Delta$ 可能不是单纯的；不过，一旦拉回到一个消解上，就没有困难。

$\mathcal A$ 超几何方程定义在以 $\lambda_i$ 为坐标的空间 $L(\Delta)$ 上。然而，我们关心的是模空间上的微分方程，而模空间是 $L(\Delta)$ 的一个商；更准确地说，这个商给出模的多项式部分。\hyperref[chap:moduli]{第~\ref*{chap:moduli} 章}将详细讨论此问题，现在只考察模的两个较明显的方面。

\SourcePage{116}
首先，环面 $T$ 在 $L(\Delta)$ 上有自然作用
\[
\nu\cdot f=\nu\cdot\sum_i\lambda_it^{m_i}
=\sum_i\nu^{m_i}\lambda_it^{m_i},\qquad \nu\in T.
\]
这个作用来自 $T$ 在自身上的自然作用。因此，对每个 $\nu\in T$，$f\in L(\Delta)$ 与 $\nu\cdot f$ 在 $\mathbb P_\Delta$ 中定义同构的超曲面，从而给出相同的模点。

第二个作用是 $\mathbb C^*$ 的作用，对 $c\in\mathbb C^*$ 定义为 $f\mapsto c\cdot f$。由于缩放给出同一个超曲面，$f$ 与 $c\cdot f$ 也给出相同的模点。合并这两个作用，便得到 $T\times\mathbb C^*$ 在 $L(\Delta)$ 上的作用，其商映射到模空间。

因为 $\omega=\eta/f$ 由 $f\in L(\Delta)$ 确定，$T\times\mathbb C^*$ 也作用于 $\omega$ 的周期积分。可以证明，方程 $Z_j\Phi=0$（$1\leq j\leq n$）是周期在 $T$ 的作用下不变的无穷小形式；同样，$Z_{n+1}\omega=0$ 来自缩放下的不变性。细节见 \SourceCite{HLY1}。

若限于 $\lambda_i\ne0$（对所有 $i$）的子集
$(\mathbb C^*)^{\Delta\cap M}\subset L(\Delta)$，那么
$(\mathbb C^*)^{\Delta\cap M}/(T\times\mathbb C^*)$ 很容易描述，因为它是环面的商。此时式~\eqref{eq:gkz-relation-lattice} 中的格 $\Lambda$ 就发挥作用。选取 $\Lambda$ 的一组基 $\ell_j=(\ell_{ji})$，定义函数
\[
z_j=\prod_i\lambda_i^{\ell_{ji}}.
\]
容易验证，$z_j$ 在 $T\times\mathbb C^*$ 的作用下不变，且给出商空间的坐标。我们的目标之一，是找到 $\Lambda$ 的一组基，使 $z_j$ 成为所考察的最大幺幂边界点附近复模空间的局部坐标。这样，在最大幺幂边界点附近解析的周期，就会是 $z_j$ 的解析函数。\hyperref[chap:moduli]{第~\ref*{chap:moduli} 章}将再回到这个问题。

我们希望用这些变量 $z_j$ 写出 $\mathcal A$ 超几何方程。先将方程改写成更便于处理模的形式。问题在于，$\omega$ 在 $\mathbb C^*$ 作用 $\lambda_i\mapsto c\lambda_i$ 下并不保持不变。因此，选取形式
\begin{equation}\label{eq:normalized-toric-period-form}
\widetilde\omega=\lambda_0\omega=\frac{\lambda_0}{f}\eta,
\end{equation}
它关于 $\lambda_i$ 的齐次次数现在是 $0$。因为 $\omega$ 的周期满足 $\mathcal A$ 方程组，所以 $\widetilde\omega$ 的周期 $\widetilde\Phi$ 满足
\[
Z_j\lambda_0^{-1}\widetilde\Phi=0,\qquad
\Box_\ell\lambda_0^{-1}\widetilde\Phi=0,
\]
其中 $1\leq j\leq n+1$、$\ell\in\Lambda$。

方程 $Z_j\lambda_0^{-1}\widetilde\Phi=0$ 是 $\widetilde\Phi$ 在 $T\times\mathbb C^*$ 作用下不变的无穷小形式。因此，一旦按 $T\times\mathbb C^*$ 的作用取商，余下的任务就是将方程
$\Box_\ell\lambda_0^{-1}\widetilde\Phi=0$ 改写成 $z_j$ 坐标。为此，反复使用微分算子恒等式
\begin{equation}\label{eq:logarithmic-operator-commutation}
\delta_i\lambda_i^p=\lambda_i^p(\delta_i+p).
\end{equation}
以五次三维簇的镜像说明这个方法。这里，$\Lambda$ 的相关生成元为 $\ell=(1,1,1,1,1,-5)$；回顾 $\Lambda$ 由
$\mathcal A=(\Delta\cap M)\times\{1\}$ 的元素之间的关系组成。它给出坐标
\begin{equation}\label{eq:quintic-gkz-coordinate}
z=\lambda_1\cdots\lambda_5/\lambda_0^5.
\end{equation}
\SourcePage{117}
在这种情形下，对式~\eqref{eq:toric-gkz-laurent} 作齐次化，得到
\[
\lambda_1x_1^5+\lambda_2x_2^5+\lambda_3x_3^5+\lambda_4x_4^5+\lambda_5x_5^5
+\lambda_0x_1x_2x_3x_4x_5=0.
\]
置 $\lambda_1=\cdots=\lambda_5=1$、$\lambda_0=\psi$，就得到式~\eqref{eq:quintic-pencil}。此时上面的模坐标为 $z=\psi^{-5}$，与\hyperref[chap:quintic]{第~\ref*{chap:quintic} 章}选取的坐标一致；回顾，我们现在用 $z$ 而不是 $x$ 表示模坐标。必须把 $\mathcal A$ 超几何方程
\[
\Box_\ell\lambda_0^{-1}=(\partial_1\cdots\partial_5-\partial_0^5)\lambda_0^{-1}=0
\]
改写成关于 $z$ 的方程。在 $\Box_\ell\lambda_0^{-1}$ 的左侧乘以
$\lambda_0\cdots\lambda_5=z\lambda_0^6$，得到
\[
\delta_1\cdots\delta_5-z\lambda_0^6\partial_0^5\lambda_0^{-1}
=\delta_1\cdots\delta_5-z\lambda_0^5\delta_0(\lambda_0^{-1}\delta_0)^4\lambda_0^{-1}.
\]
从右到左反复应用式~\eqref{eq:logarithmic-operator-commutation}，得到方程
\[
\delta_1\cdots\delta_5-z(\delta_0-5)(\delta_0-4)\cdots(\delta_0-1)=0.
\]
现在用 $\delta$ 表示 $z\,d/dz$。限于仅依赖 $z$ 的函数时，由式~\eqref{eq:quintic-gkz-coordinate} 得恒等式
\[
\delta_1=\cdots=\delta_5=\delta,\qquad\delta_0=-5\delta.
\]
最终得到
\begin{equation}\label{eq:quintic-gkz-fifth}
\delta^5+z(5\delta+5)(5\delta+4)\cdots(5\delta+1)=0.
\end{equation}
这种方法的局限立刻显现出来：五阶算子~\eqref{eq:quintic-gkz-fifth} 并不生成 Picard--Fuchs 理想，因为还缺少四阶算子~\eqref{eq:quintic-mum-pf}。

一般而言，从 $\mathcal A$ 方程组得到的 $z_j$ 坐标中的微分方程组，其解包含周期，但通常还会有其他解。问题在于，找到一个更大的微分方程组，其共同解恰好由周期组成。为此，有几种策略。

求出全部 Picard--Fuchs 方程的一种方法，是从 $\mathcal A$ 方程组的超几何方程出发，再添加其他 Picard--Fuchs 方程。通常，所添加的是用 Griffiths--Dwork 方法找到的低阶方程。一些例子见 \SourceCite{BKK}。

另一种常用方法，是分解来自 $\mathcal A$ 方程组的某个 Picard--Fuchs 方程。如~\hyperref[subsec:picard-fuchs]{第 5.1.2 小节}所述，只要分解得到的低阶算子 $D$ 零化所有周期，它就是 Picard--Fuchs 算子。幸而在许多情形中，只需检验 $D$ 零化一个周期 $y$ 就够了。例如，假设 $D\cdot y=0$，且把 $y$ 沿整个模空间解析延拓，会得到 Picard--Fuchs 方程解的一组基 $y_i$。那么 $D\cdot y=0$ 蕴含对所有 $i$ 都有 $D\cdot y_i=0$。因此，$D$ 零化所有周期，因而是 Picard--Fuchs 算子。

以五次三维簇的镜像说明这一过程。从五阶方程~\eqref{eq:quintic-gkz-fifth} 出发，把 $z$ 交换到右边，得到
\[
\delta^5+(5\delta)(5\delta-1)\cdots(5\delta-4)z=0.
\]
于是从左侧提出因子 $\delta$，得到四阶方程
\begin{equation}\label{eq:quintic-gkz-factor}
\delta^4+5(5\delta-1)(5\delta-2)(5\delta-3)(5\delta-4)z=0,
\end{equation}
\SourcePage{118}
它等价于式~\eqref{eq:quintic-mum-pf}。\SourceCite{HKTY1} 描述了这种分解的系统方法。为了证明它是一个 Picard--Fuchs 方程，考虑式~\eqref{eq:y0-hypergeometric} 中的
\[
y_0(z)=\sum_{n=0}^\infty\frac{(5n)!}{(n!)^5}(-1)^nz^n.
\]
由后文的式~\eqref{eq:toric-euler-period}，将看出 $y_0$ 确实是一个周期；容易检验，它是式~\eqref{eq:quintic-gkz-factor} 的解。进一步，\SourceCite{CdGP} 证明，$y_0$ 的解析延拓给出周期的一组基。因此，式~\eqref{eq:quintic-gkz-factor} 是五次三维簇镜像的 Picard--Fuchs 方程。

要在一般情形中实施这个策略，就需要知道一个满足 $\mathcal A$ 方程组的显式周期。\hyperref[chap:moduli]{第~\ref*{chap:moduli} 章}的式~\eqref{eq:toric-euler-period} 将给出这样的一个周期，取自 \SourceCite[命题 14.6]{Batyrev3}。实际上，式~\eqref{eq:toric-euler-period} 中的周期 $y$ 还给出一种略有不同的生成 Picard--Fuchs 方程的方法：利用 $y$ 的系数之间的递推关系，推出微分方程。\SourceCite{BvS} 将这个技巧用于若干单参数族。其不足是，一般尚不知道式~\eqref{eq:toric-euler-period} 的解析延拓是否给出解的一组基，尽管对 Calabi--Yau 三维簇，人们预期确实如此。更多细节见 \SourceCite{BvS}。

遗憾的是，因式分解可能仍不能给出全部缺失的 Picard--Fuchs 算子，因为 $\mathcal A$ 方程组还有第二个不完整性的来源，即自同构。GKZ 方程 $Z_j\Phi=0$ 反映的是在 $\mathbb P_\Delta$ 的环面作用下的不变性；而要得到 $f=0$ 的模，必须取商于 $\operatorname{Aut}(\mathbb P_\Delta)$ 的作用。我们在~\hyperref[sec:toric-automorphisms]{第 3.6 节}看到，$\operatorname{Aut}(\mathbb P_\Delta)$ 还包含来自根的自同构。这些自同构使某些坐标 $\lambda_i$ 无效，即对应于 $\Delta$ 的刻面内部点的那些坐标。与其添加更多方程，可以直接略去这些 $\lambda_i$；但须记住，模空间上可能还保留离散的等同关系，一个例子见 \SourceCite{CFKM}。\hyperref[chap:moduli]{第~\ref*{chap:moduli} 章}定义环面超曲面 $V$ 的\emph{简化模空间} $\mathcal M_{\mathrm{simp}}$ 时，我们将采用这一方法。另一种密切相关的方法，是把这些额外自同构的无穷小形式加入 $\mathcal A$ 方程组，见 \SourceCite{HKTY1,HLY1}。它们补充了方程 $Z_j\Phi=0$；如刚才所述，后者是环面所诱导自同构的无穷小形式。

上述 GKZ 方程组的解~\eqref{eq:toric-euler-period}，只是众多精妙的\IndexLocation{indexloc:0615:118}超几何方程解公式之一。这些公式推广了经典的超几何幂级数。\SourceCite{GKZ1,HLY1} 讨论了这一点；在\hyperref[sec:mirror-map]{第~\ref*{sec:mirror-map} 节}讨论镜像映射时，我们还将给出一些显式例子与求解方法。

再回到已经考虑过的一个例子，说明上述讨论。
\SourceExampleNumber{5.5.2}{1}
\begin{example}\label{ex:weighted-gkz}
在~\hyperref[sec:hodge-examples]{第 5.4 节}，讨论了 $\mathbb P(1,1,2,2,2)$ 中次数为 $8$ 的\IndexLocation{indexloc:1020:118}\IndexLocation{indexloc:1027:118}超曲面的镜像 $V^\circ$。这里，将式~\eqref{eq:toric-gkz-laurent} 齐次化，得到形如
\[
\lambda_1x_1^8+\lambda_2x_2^8+\lambda_3x_3^4+\lambda_4x_4^4+\lambda_5x_5^4
+\lambda_6x_1^4x_2^4+\lambda_0x_1x_2x_3x_4x_5=0
\]
的方程。这些单项式已在例题~\ref{ex:weighted-pf-method} 中算出。利用由式~\eqref{eq:weighted-pf-rays} 确定的自反多胞形 $\Delta^\circ$，可以算出 $\Lambda$ 的一组基为
$\ell_1=(0,0,1,1,1,1,-4)$ 与 $\ell_2=(1,1,0,0,0,-2,0)$。于是得到模坐标
\[
z_1=\frac{\lambda_3\lambda_4\lambda_5\lambda_6}{\lambda_0^4},\qquad
z_2=\frac{\lambda_1\lambda_2}{\lambda_6^2}.
\]
\SourcePage{119}
注意，令 $\lambda_0=\lambda_2=\lambda_4=\lambda_5=\lambda_6=1$，就有 $z_1=\lambda_3$、$z_2=\lambda_1$，与式~\eqref{eq:weighted-pf-family} 一致。由
\[
\lambda_0\lambda_1\lambda_2\Box_{\ell_2}\lambda_0^{-1}
=\lambda_0\lambda_1\lambda_2(\partial_1\partial_2-\partial_6^2)\lambda_0^{-1}=0,
\]
按上面讨论的方法，得到
\begin{equation}\label{eq:weighted-gkz-second}
\delta_{z_2}^2-z_2(\delta_{z_1}-2\delta_{z_2})(\delta_{z_1}-2\delta_{z_2}-1)=0,
\end{equation}
其中 $\delta_{z_1}=z_1\partial/\partial z_1$、$\delta_{z_2}=z_2\partial/\partial z_2$，并使用了在 $z_1,z_2$ 的函数上成立的恒等式
\begin{equation}\label{eq:weighted-logarithmic-identities}
\delta_1=\delta_2=\delta_{z_2},\qquad\delta_6=\delta_{z_1}-2\delta_{z_2}.
\end{equation}
读者应注意，简单的代数运算就能把式~\eqref{eq:weighted-gkz-second} 变为先前用 Griffiths--Dwork 方法求出的 Picard--Fuchs 方程~\eqref{eq:weighted-pf-second}；其中出现的 $\delta_i$，就是这里的 $\delta_{z_i}$。

类似地，利用式~\eqref{eq:weighted-logarithmic-identities}，以及
$\delta_3=\delta_4=\delta_5=\delta_{z_1}$、$\delta_0=-4\delta_{z_1}$，可以算出 GKZ 方程
\[
\lambda_0\lambda_3\lambda_4\lambda_5\lambda_6\Box_{\ell_1}\lambda_0^{-1}
=\lambda_0\lambda_3\lambda_4\lambda_5\lambda_6
(\partial_3\partial_4\partial_5\partial_6-\partial_0^4)\lambda_0^{-1}=0
\]
给出方程
\[
\delta_{z_1}^3(\delta_{z_1}-2\delta_{z_2})
-(4\delta_{z_1})(4\delta_{z_1}-1)(4\delta_{z_1}-2)(4\delta_{z_1}-3)z_1=0.
\]
提出左侧因子，便得到三阶方程
\[
\delta_{z_1}^2(\delta_{z_1}-2\delta_{z_2})
-4(4\delta_{z_1}-1)(4\delta_{z_1}-2)(4\delta_{z_1}-3)z_1=0,
\]
它等价于
\begin{equation}\label{eq:weighted-gkz-third}
\delta_{z_1}^2(\delta_{z_1}-2\delta_{z_2})
-4z_1(4\delta_{z_1}+1)(4\delta_{z_1}+2)(4\delta_{z_1}+3)=0.
\end{equation}
解析周期~\eqref{eq:toric-period-local-series} 满足这个方程；而且，如 \SourceCite{CdFKM} 所示，这个周期的解析延拓给出解的一组基。由本例之前的讨论，式~\eqref{eq:weighted-gkz-third} 确实是一个 Picard--Fuchs 方程。

两个 Picard--Fuchs 算子~\eqref{eq:weighted-gkz-second} 与~\eqref{eq:weighted-gkz-third} 的阶数分别为 $2$ 与 $3$，所确定方程组的解空间维数至多为 $6$。因为 $h^3(V^\circ)=6$，我们已经找到生成全部 Picard--Fuchs 方程的算子。
\end{example}

\subsection{nef 完全交}\label{subsec:nef-gkz}\index{nef完全交@nef 完全交|hyperpage}

至今已讨论 Picard--Fuchs 方程与 GKZ 方程组，但还没有多谈它们的解。这要等到\hyperref[sec:mirror-map]{第~\ref*{sec:mirror-map} 节}：在那里，将利用 Picard--Fuchs 方程的解构造镜像映射。不过，我们觉得值得先在这里介绍一个由\IndexLocation{indexloc:1349:119}环面完全交产生的有趣的 GKZ 方程组的解。

设 $X=\mathbb P_\Delta$ 是对应于\IndexLocation{indexloc:1218:119}自反多胞形 $\Delta\subset M_{\mathbb R}$ 的 $n$ 维环面簇，法扇为 $\Sigma$。为方便叙述，假定 $X$ 光滑；这个假设实际上可以适当减弱。我们考虑~\hyperref[subsec:nef-partitions]{第 4.3.1 小节}所研究完全交的一种推广。
\begin{definition}\label{def:nef-complete-intersection}
$X$ 上的一个 \emph{nef 完全交}，是完全交 $Y=\bigcap_{i=1}^kY_i$，其中每个 $Y_i$ 都是由整体截面生成的线丛 $\mathcal L_i$ 的一个一般截面的零点集，并且 $-(K_X+\sum_{i=1}^k\mathcal L_i)$ 为 nef。
\end{definition}

\SourcePage{120}
注意，最后一个条件略强于 $-K_Y$ 为 nef。我们要求对所有曲线 $C\subset X$ 都有
\[
\left(K_X+\sum_{i=1}^k\mathcal L_i\right)\cdot C\leq0.
\]
根据伴随公式，$-K_Y$ 为 nef 只要求这个条件对 $C\subset Y$ 中的曲线成立。$Y$ 与 $X$ 的有效曲线锥可以不同，我们将在\hyperref[subsec:enlarged-kahler-moduli]{第~\ref*{subsec:enlarged-kahler-moduli} 小节}讨论这一点。

注意，~\hyperref[subsec:nef-partitions]{第 4.3.1 小节}定义的 nef 分拆
$-K_X=\sum_{i=1}^kE_i$，自动定义一个 nef 完全交。这时 $\mathcal L_i=\mathcal O(E_i)$，所以
$-(K_X+\sum_{i=1}^k\mathcal L_i)$ 是平凡的，显然为 nef。

现在置 $\overline N=N\oplus\mathbb Z^k$，并对 $1\leq i\leq k$ 选取 $a_{i\rho}$，使
\[
\mathcal L_i=\mathcal O\left(\sum_{\rho\in\Sigma(1)}a_{i\rho}D_\rho\right).
\]
考虑点集
\begin{equation}\label{eq:nef-gkz-configuration}
\mathcal A=\{(v_\rho,a_{1\rho},\ldots,a_{k\rho}):\rho\in\Sigma(1)\}
\cup\{(0,e_1),\ldots,(0,e_k)\}\subset\overline N,
\end{equation}
其中 $e_i$ 是 $\mathbb Z^k$ 的第 $i$ 个标准基向量。注意，这推广了式~\eqref{eq:toric-gkz-configuration}；后者对应于 $k=1$、$\mathcal L_1=\mathcal O(\sum_\rho D_\rho)$ 的情形。

如往常一样，令 $r=|\Sigma(1)|$。对每个 $\rho\in\Sigma(1)$，置
$w_\rho=(v_\rho,a_{1\rho},\ldots,a_{k\rho})$；对 $1\leq i\leq k$，置 $u_i=(0,e_i)$。于是式~\eqref{eq:nef-gkz-configuration} 可写成
$\mathcal A=\{w_\rho,\rho\in\Sigma(1);\ u_i,1\leq i\leq k\}$。

点集 $\mathcal A$ 与光滑环面簇 $X$ 有如下关系。
\begin{lemma}\label{lem:nef-relation-lattice}
设 $\Lambda$ 是 $\mathcal A$ 的元素之间的关系格。则：
\begin{enumerate}[label=(\roman*)]
\item 对每个 $\beta\in H_2(X,\mathbb Z)$，在 $\overline N$ 中有关系
\[
\sum_\rho(D_\rho\cdot\beta)w_\rho-
\sum_i(c_1(\mathcal L_i)\cdot\beta)u_i=0.
\]
因此 $\beta$ 给出格 $\Lambda$ 的一个元素。
\item 所诱导的映射 $H_2(X,\mathbb Z)\to\Lambda$ 是同构。
\end{enumerate}
\end{lemma}
\begin{proof}
先注意，由正合列~\eqref{eq:toric-chow-sequence}，对所有 $m\in M$，在 $H^2(X,\mathbb Z)$ 中有
$\sum_\rho\langle m,v_\rho\rangle[D_\rho]=0$。因此，当 $\beta\in H_2(X,\mathbb Z)$ 时，
$\sum_\rho\langle m,v_\rho\rangle(D_\rho\cdot\beta)=0$。这对所有 $m$ 成立，故必有
$\sum_\rho(D_\rho\cdot\beta)v_\rho=0$（在 $N$ 中）。再结合
$c_1(\mathcal L_i)=\sum_\rho a_{i\rho}D_\rho$，以及 $w_\rho,u_i$ 的定义，容易得到所需关系。

为证明第二部分，考虑交换图
\[
\begin{array}{ccccccc}
0&\longrightarrow&\Lambda&\longrightarrow&\mathbb Z^{\mathcal A}&\longrightarrow&\overline N\\
&&\downarrow&&\downarrow&&\downarrow\\
0&\longrightarrow&H_2(X,\mathbb Z)&\longrightarrow&\mathbb Z^{\Sigma(1)}&\longrightarrow&N.
\end{array}
\]
上行按 $\Lambda$ 的定义正合，下行则是对正合列~\eqref{eq:toric-chow-sequence} 取对偶所得，因而也正合。中间的竖直映射将 $w_\rho$ 送到 $v_\rho$，将 $u_i$ 送到 $0$；最右侧的竖直映射是自然投影。最左侧的竖直映射由交换图确定。若 $N'$ 是 $\mathbb Z^{\Sigma(1)}\to N$ 的像，则容易看出，$\mathbb Z^{\mathcal A}\to\overline N$ 的像为 $N'\oplus\mathbb Z^k$。由此作图追踪，可以证明 $\Lambda\to H_2(X,\mathbb Z)$ 是同构。还可以检验，它的逆映射恰好就是引理陈述中描述的 $H_2(X,\mathbb Z)\to\Lambda$。
\end{proof}

为简单起见，现在假定 Kähler 锥 $K(X)$ 为单纯锥。再假定对偶 Mori 锥的本原射线生成元，是 $H_2(X,\mathbb Z)$ 的一组整基。\hyperref[subsec:enlarged-kahler-moduli]{第~\ref*{subsec:enlarged-kahler-moduli} 小节}将再讨论与这个假设有关的问题。这样可以选取一组整同调基
\[
\beta_1,\ldots,\beta_{r-n}\in H_2(X,\mathbb Z),
\]
使它们生成
\SourcePage{121}
Mori 锥 $M(X)\subset H_2(X,\mathbb R)$ 的整类 $M(X)_{\mathbb Z}$；这个锥由 $X$ 上有效曲线的同调类生成。在本小节余下部分，固定 $\Sigma(1)$ 的元素次序；不过，继续将 $\Sigma(1)$ 的生成元记为 $v_\rho$ 往往仍很方便。固定这一排序后，按式~\eqref{eq:nef-gkz-configuration} 中的次序，把变量 $\lambda_1,\ldots,\lambda_{r+k}$ 对应于 $\mathcal A$ 的各个生成元。对应于 $w_\rho\in\mathcal A$ 的变量，也常记为 $\lambda_\rho$。

置 $\overline T=\overline N\otimes\mathbb C^*$，引入 $\overline T$ 不变变量
\begin{equation}\label{eq:nef-gkz-q-coordinates}
q_j=(-1)^{\sum_{i=1}^k c_1(\mathcal L_i)\cdot\beta_j}
\frac{\prod_\rho\lambda_\rho^{D_\rho\cdot\beta_j}}
{\prod_{i=1}^k\lambda_{r+i}^{c_1(\mathcal L_i)\cdot\beta_j}},
\qquad 1\leq j\leq r-n.
\end{equation}
令 $T_1,\ldots,T_{r-n}$ 是与 $\beta_1,\ldots,\beta_{r-n}$ 对偶的 $K(X)$ 的生成元，定义
\[
q^\beta=\prod_jq_j^{T_j\cdot\beta},
\qquad \beta\in M(X)_{\mathbb Z}\subset H_2(X,\mathbb Z).
\]
\hyperref[chap:quantum-cohomology]{第~\ref*{chap:quantum-cohomology} 章}讨论量子上同调时，将更频繁地使用这个记号。还置 $t_j=\log q_j$，于是 $q_j=e^{t_j}$。

受 \SourceCite{Givental4} 启发，考虑上同调值表达式
\[
\widetilde I=e^{\sum_jt_jT_j}\sum_{\beta\in M(X)_{\mathbb Z}}q^\beta
\frac{\left(\prod_{i=1}^k\prod_{m=1}^{c_1(\mathcal L_i)\cdot\beta}
(c_1(\mathcal L_i)+m)\right)
\left(\prod_\rho\prod_{m=-\infty}^{0}(D_\rho+m)\right)}
{\prod_\rho\prod_{m=-\infty}^{D_\rho\cdot\beta}(D_\rho+m)}.
\]
指数 $e^{\sum_jt_jT_j}$ 是 $H^*(X)[t_1,\ldots,t_{r-n}]$ 中的多项式，求和则可视为
$H^*(X)[[q_1,\ldots,q_{r-n}]]$ 中的元素。注意，$c_1(\mathcal L_i)\cdot\beta\geq0$，因为 $\mathcal L_i$ 由整体截面生成；分子与分母中的因子 $D_\rho+m$ 除有限个以外全部抵消。这种略显繁琐的定义方式，是为了处理 $D_\rho\cdot\beta<0$ 的可能性。由于 $-(K_X+\sum_{i=1}^k\mathcal L_i)$ 为 nef，猜想当 $q_j$ 足够小时，$\widetilde I$ 的级数收敛。

令 $\mathcal V=\bigoplus_{i=1}^k\mathcal L_i$，其 Euler 类为
$\operatorname{Euler}(\mathcal V)=\prod_{i=1}^kc_1(\mathcal L_i)$。\hyperref[chap:mirror-theorem]{第~\ref*{chap:mirror-theorem} 章}将看到，上同调值形式函数
\begin{equation}\label{eq:nef-givental-iv}
I_{\mathcal V}=\operatorname{Euler}(\mathcal V)\widetilde I
\end{equation}
在 Givental 版本的环面完全交镜像定理中起重要作用，见式~\eqref{eq:toric-complete-intersection-i-function} 与定理~\ref{thm:givental-toric-mirror}。

在 $X=\mathbb P^4$、$k=1$、$\mathcal V=\mathcal L_1=\mathcal O_{\mathbb P^4}(5H)$ 的特殊情形，完全交就是五次三维簇。此时
\begin{equation}\label{eq:quintic-tilde-i}
\widetilde I=e^{t_1H}\sum_{d=0}^\infty e^{dt_1}
\frac{\prod_{m=1}^{5d}(5H+m)}{\prod_{m=1}^{d}(H+m)^5},
\end{equation}
而 $I_{\mathcal V}=5H\widetilde I$ 变为
\begin{equation}\label{eq:quintic-iv-hypergeometric}
I_{\mathcal V}=e^{t_1H}5H\sum_{d=0}^\infty e^{dt_1}
\frac{\prod_{m=1}^{5d}(5H+m)}{\prod_{m=1}^{d}(H+m)^5}.
\end{equation}
最后这个公式已在\hyperref[chap:quintic]{第~\ref*{chap:quintic} 章}作为式~\eqref{eq:iv-simple} 出现。

函数 $\widetilde I$ 与\IndexLocation{indexloc:0615:121}超几何方程有如下关系。令
$Q=\prod_{i=1}^k\lambda_{r+i}^{-1}$。
\begin{proposition}\label{prop:nef-i-gkz}
形式函数 $Q\widetilde I$ 满足与式~\eqref{eq:nef-gkz-configuration} 对应的 $\mathcal A$ 方程组，其中
$\widehat\beta=-\sum_{i=1}^ku_i=(0,\ldots,0,-1,\ldots,-1)$。
\end{proposition}
\begin{proof}
\SourcePage{122}
这里 $\widehat\beta$ 指的是定义式~\eqref{eq:gkz-euler-operators} 中算子 $Z_j$ 所用的向量。如~\hyperref[subsec:toric-gkz]{第 5.5.2 小节}所讨论的，对所有 $j$，$Z_j(\widetilde I)=0$ 等价于形式函数 $\widetilde I$ 的 $\overline T$ 不变性。$q_j$ 按构造在 $\overline T$ 下不变，$\lambda_i$ 在 $\widetilde I$ 中也仅通过 $q_j$ 出现，所以 $\widetilde I$ 是 $\overline T$ 不变的。乘以 $Q$ 后，其权为 $-\sum_i u_i$，从而对所有 $j$ 都有 $Z_j(Q\widetilde I)=0$。

现在考虑 $\beta\in H_2(X,\mathbb Z)\simeq\Lambda$ 所对应的算子 $\Box_\beta$。由引理~\ref{lem:nef-relation-lattice}，$\beta$ 给出关系
$\sum_\rho(D_\rho\cdot\beta)w_\rho-\sum_i(c_1(\mathcal L_i)\cdot\beta)u_i=0$。因此，$\Box_\beta$ 可写为
\begin{equation}\label{eq:nef-gkz-box}
\begin{split}
\Box_\beta={}&
\prod_{D_\rho\cdot\beta>0}\partial_\rho^{D_\rho\cdot\beta}
\prod_{c_1(\mathcal L_i)\cdot\beta<0}\partial_{r+i}^{-c_1(\mathcal L_i)\cdot\beta}\\
&-\prod_{D_\rho\cdot\beta<0}\partial_\rho^{-D_\rho\cdot\beta}
\prod_{c_1(\mathcal L_i)\cdot\beta>0}\partial_{r+i}^{c_1(\mathcal L_i)\cdot\beta},
\end{split}
\end{equation}
其中 $\partial_\rho=\partial/\partial\lambda_\rho$、$\partial_{r+i}=\partial/\partial\lambda_{r+i}$。需要证明 $\Box_\beta(Q\widetilde I)=0$，等价地，证明
$(Q^{-1}\Box_\beta Q)\widetilde I=0$。

如往常一样，置 $\delta_\rho=\lambda_\rho\partial_\rho$、
$\delta_{r+i}=\lambda_{r+i}\partial_{r+i}$，并置
$\delta_{q_j}=q_j\partial/\partial q_j=\partial/\partial t_j$。将前述方法用于式~\eqref{eq:nef-gkz-q-coordinates}，对 $\overline T$ 不变函数得到恒等式
\begin{equation}\label{eq:nef-logarithmic-identities}
\begin{aligned}
\delta_\rho&=\sum_j(D_\rho\cdot\beta_j)\delta_{q_j},\\
\delta_{r+i}&=-\sum_j(c_1(\mathcal L_i)\cdot\beta_j)\delta_{q_j}.
\end{aligned}
\end{equation}
另外，计算得
\begin{equation}\label{eq:nef-i-derivative}
\delta_{q_j}\left(e^{\sum_jt_jT_j}q^\beta\right)
=(T_j+T_j\cdot\beta)\left(e^{\sum_jt_jT_j}q^\beta\right).
\end{equation}
结合式~\eqref{eq:nef-logarithmic-identities} 与~\eqref{eq:nef-i-derivative}，得到
\begin{equation}\label{eq:nef-i-lambda-derivatives}
\begin{aligned}
\delta_\rho\left(e^{\sum_jt_jT_j}q^\beta\right)
&=(D_\rho+D_\rho\cdot\beta)\left(e^{\sum_jt_jT_j}q^\beta\right),\\
\delta_{r+i}\left(e^{\sum_jt_jT_j}q^\beta\right)
&=-(c_1(\mathcal L_i)+c_1(\mathcal L_i)\cdot\beta)
\left(e^{\sum_jt_jT_j}q^\beta\right).
\end{aligned}
\end{equation}
这里使用了除子恒等式 $D=\sum_j(D\cdot\beta_j)T_j$，分别取 $D=D_\rho$ 与 $D=c_1(\mathcal L_i)$；该恒等式成立，是因为 $T_j$ 与 $\beta_j$ 互为对偶基。

$\Box_\beta$ 用 $\partial_\rho,\partial_{r+i}$ 定义。为了改用对数微分算子，考虑
\[
\Box'_\beta=
\left(\prod_{D_\rho\cdot\beta>0}\lambda_\rho^{D_\rho\cdot\beta}
\prod_{c_1(\mathcal L_i)\cdot\beta<0}\lambda_{r+i}^{-c_1(\mathcal L_i)\cdot\beta}\right)
Q^{-1}\Box_\beta Q.
\]
记 $s_\beta=\sum_i c_1(\mathcal L_i)\cdot\beta$，则
\begin{equation}\label{eq:nef-normalized-box}
\begin{aligned}
\Box'_\beta={}&
\prod_{D_\rho\cdot\beta>0}\delta_\rho(\delta_\rho-1)\cdots
(\delta_\rho-D_\rho\cdot\beta+1)\\
&\quad\times\prod_{c_1(\mathcal L_i)\cdot\beta<0}
(\delta_{r+i}-1)(\delta_{r+i}-2)\cdots
(\delta_{r+i}+c_1(\mathcal L_i)\cdot\beta)\\
&-(-1)^{s_\beta}q^\beta
\prod_{D_\rho\cdot\beta<0}\delta_\rho(\delta_\rho-1)\cdots
(\delta_\rho+D_\rho\cdot\beta+1)\\
&\quad\times\prod_{c_1(\mathcal L_i)\cdot\beta>0}
(\delta_{r+i}-1)(\delta_{r+i}-2)\cdots
(\delta_{r+i}-c_1(\mathcal L_i)\cdot\beta).
\end{aligned}
\end{equation}
为写出其作用后的表达式，置
$a_\rho=D_\rho\cdot\beta$、$b_i=c_1(\mathcal L_i)\cdot\beta$，以及
$\varepsilon_\beta=(-1)^{-\sum_{b_i<0}b_i}$。利用式~\eqref{eq:nef-normalized-box} 与~\eqref{eq:nef-i-lambda-derivatives}，计算得
\begin{equation}\label{eq:nef-gkz-cancellation}
\begin{split}
\Box'_\beta\widetilde I
={}&\varepsilon_\beta e^{\sum_jt_jT_j}\Biggl(\\
&\sum_{\beta'\in M(X)_{\mathbb Z}}q^{\beta'}
\frac{\left(\prod_i\prod_{m=1}^{c_1(\mathcal L_i)\cdot\beta'-\min(b_i,0)}
(c_1(\mathcal L_i)+m)\right)
\left(\prod_\rho\prod_{m=-\infty}^0(D_\rho+m)\right)}
{\prod_\rho\prod_{m=-\infty}^{D_\rho\cdot\beta'-\max(a_\rho,0)}(D_\rho+m)}\\
-&\sum_{\beta''\in M(X)_{\mathbb Z}}q^{\beta''+\beta}
\frac{\left(\prod_i\prod_{m=1}^{c_1(\mathcal L_i)\cdot\beta''+\max(b_i,0)}
(c_1(\mathcal L_i)+m)\right)
\left(\prod_\rho\prod_{m=-\infty}^0(D_\rho+m)\right)}
{\prod_\rho\prod_{m=-\infty}^{D_\rho\cdot\beta''+\min(a_\rho,0)}(D_\rho+m)}
\Biggr).
\end{split}
\end{equation}

\SourcePage{123}
在 $\Box'_\beta\widetilde I$ 的这个表达式中，若 $\beta'=\beta''+\beta$，第一个和中对应于 $\beta'\in M(X)_{\mathbb Z}$ 的项，就与第二个和中对应于 $\beta''\in M(X)_{\mathbb Z}$ 的项抵消。抵消后剩下两类项：第一类来自第一个和，其中 $\beta'\in M(X)_{\mathbb Z}$ 而 $\beta'-\beta\notin M(X)_{\mathbb Z}$；第二类来自第二个和，其中 $\beta''\in M(X)_{\mathbb Z}$ 而 $\beta''+\beta\notin M(X)_{\mathbb Z}$。

先考察第一类。设 $\beta'\in M(X)_{\mathbb Z}$，且 $\beta'-\beta\notin M(X)_{\mathbb Z}$。那么 $e^{\sum_jt_jT_j}q^{\beta'}$ 出现在式~\eqref{eq:nef-gkz-cancellation} 的第一个和中，而不出现在第二个和中。目标是证明它的系数为零。假设 $\beta'-\beta\notin M(X)_{\mathbb Z}$，意味着存在一个丰沛除子 $D$，使 $D\cdot(\beta'-\beta)<0$。不妨假定 $D$ 是 $T$ 不变的，并写成 $D=\sum_\rho a_\rho D_\rho$，其中 $a_\rho\geq0$。令
\[
S=\{\rho\in\Sigma(1):D_\rho\cdot(\beta'-\beta)<0\}.
\]
由于 $D\cdot(\beta'-\beta)<0$，有 $S\ne\varnothing$。

我们断言，在 $H^*(X)$ 中有 $\prod_{\rho\in S}D_\rho=0$。为此，只需证明 $\bigcap_{\rho\in S}D_\rho=\varnothing$。假定该交集非空；那么它至少包含一个环面作用的不动点。若 $\sigma$ 是 $\Sigma$ 中对应的 $n$ 维锥，则对所有 $\rho\in S$ 都有 $\rho\subset\sigma$。$D=\sum a_\rho D_\rho$ 是 Cartier 除子，所以存在 $m_\sigma\in M$，使对 $\sigma$ 的所有射线 $\rho$，都有 $\langle m_\sigma,v_\rho\rangle=-a_\rho$；特别地，这对 $S$ 中的所有元素成立。这里 $v_\rho\in N$ 是射线 $\rho$ 的生成元。又因为 $D$ 丰沛，~\hyperref[sec:cones-fans]{第 3.1 节}末的丰沛性判据给出：对所有 $\rho\in\Sigma(1)$，都有 $\langle m_\sigma,v_\rho\rangle\geq-a_\rho$。于是 $D$ 与除子 $D'=D+\operatorname{div}(\chi^{m_\sigma})$ 线性等价。将 $D$ 换成 $D'$ 后，可以假定
\begin{equation}\label{eq:nef-gkz-divisor-contradiction}
D=\sum_{\rho\notin S}a_\rho D_\rho,\qquad a_\rho\geq0.
\end{equation}
两边与 $\beta'-\beta$ 相交，左边为负，右边非负，产生矛盾。断言得证。

现在考察式~\eqref{eq:nef-gkz-cancellation} 第一个和中 $e^{\sum_jt_jT_j}q^{\beta'}$ 的系数，并固定 $\rho\in S$。由定义，$D_\rho\cdot(\beta'-\beta)<0$；于是，从式~\eqref{eq:nef-gkz-cancellation} 可见，$D_\rho$ 是该系数的一个因子。这对所有 $\rho\in S$ 成立，所以系数是 $\prod_{\rho\in S}D_\rho=0$ 的倍数，因而为零。

剩下第二类，即 $\beta''\in M(X)_{\mathbb Z}$ 且 $\beta''+\beta\notin M(X)_{\mathbb Z}$ 的项。论证完全类似，予以省略。
\end{proof}

命题~\ref{prop:nef-i-gkz} 的一个直接推论，是形式函数 $I_{\mathcal V}=\operatorname{Euler}(\mathcal V)\widetilde I$（由式~\eqref{eq:nef-givental-iv} 定义）也是 $\mathcal A$ 方程组的解。

命题~\ref{prop:nef-i-gkz} 为式~\eqref{eq:nef-gkz-configuration} 中的 $\mathcal A$ 提供了有趣的 GKZ 方程组的解。也就是说，把 $\widetilde I$ 在 $H^*(X)$ 的任意一组基下展开，系数都满足 GKZ 方程。那么，它与 Picard--Fuchs 方程有什么联系？对于与 nef 分拆对应的 Calabi--Yau 完全交 $V\subset X$，~\hyperref[subsec:nef-partitions]{第 4.3.1 小节}构造了其镜像族 $V^\circ$；可以证明，$\mathcal A$ 方程组中的微分方程包含在镜像族 $V^\circ$ 的 Picard--Fuchs 方程中。这推广了~\hyperref[subsec:toric-gkz]{第 5.5.2 小节}的超曲面情形。（读者可能会担心，超曲面情形使用了 $\widehat\beta=(0,\ldots,0,-1)$，而完全交的归一化形式函数 $\widetilde I$ 的 Euler 权为 $\widehat\beta=(0,\ldots,0)$。这一差别并不重要，来自式~\eqref{eq:normalized-toric-period-form} 中全纯形式 $\widetilde\omega$ 的选择。）

\SourcePage{124}
举例说明：取 $X=\mathbb P^4$、$\mathcal V=\mathcal L_1=\mathcal O_{\mathbb P^4}(5H)$。用 $H^*(\mathbb P^4)$ 的通常基展开式~\eqref{eq:quintic-tilde-i}，得到
\begin{equation}\label{eq:quintic-i-expansion}
\widetilde I=y_0+y_1H+y_2H^2+y_3H^3+y_4H^4.
\end{equation}
在变量识别 $z=-q_1$ 下，$y_0,\ldots,y_4$ 都是五阶 GKZ 方程~\eqref{eq:quintic-gkz-fifth} 的解。

不过，已经看到，GKZ 方程组只是全部 Picard--Fuchs 方程的一部分。因此，虽然 $\widetilde I$ 满足归一化后的 GKZ 方程组，却可能不满足全部 Picard--Fuchs 方程。例如，在~\hyperref[subsec:toric-gkz]{第 5.5.2 小节}，式~\eqref{eq:quintic-gkz-fifth} 分解得到四阶方程，即五次镜像的 Picard--Fuchs 方程。\hyperref[subsec:hypergeometric-mirror-map]{第~\ref*{subsec:hypergeometric-mirror-map} 小节}将证明，式~\eqref{eq:quintic-i-expansion} 中只有 $y_0,\ldots,y_3$ 是 Picard--Fuchs 方程的解；$y_4$ 是多余的。因此考虑 $I_{\mathcal V}=\operatorname{Euler}(\mathcal V)\widetilde I$，将式~\eqref{eq:quintic-i-expansion} 变为
\begin{equation}\label{eq:quintic-iv-period-expansion}
I_{\mathcal V}=5H(y_0+y_1H+y_2H^2+y_3H^3),
\end{equation}
因为 $H^5=0$。\hyperref[subsec:hypergeometric-mirror-map]{第~\ref*{subsec:hypergeometric-mirror-map} 小节}的理论将说明，$y_0,\ldots,y_3$ 是五次镜像 Picard--Fuchs 方程解的一组基。还将看到，将 Frobenius 方法用于 Picard--Fuchs 方程，会自然产生式~\eqref{eq:quintic-tilde-i} 中 $\widetilde I$ 的公式。

更一般地，当 $\mathcal L_i$ 确定的完全交是 Calabi--Yau 簇时，式~\eqref{eq:nef-givental-iv} 中的
$I_{\mathcal V}=\prod_i c_1(\mathcal L_i)\widetilde I$ 应满足镜像的 Picard--Fuchs 方程。\SourceCite{Givental4} 无证明地提出了这一断言。在显式例子中，检验 $I_{\mathcal V}$ 的系数满足 $V^\circ$ 的 Picard--Fuchs 方程，往往并不困难；例如，五次三维簇的镜像确实如此。在这些例子中，通常通过考察各种上同调类空间的维数，并使用与 Picard--Fuchs 方程组对应的 $\mathcal D$ 模的不可约性，给出论证。

\hyperref[chap:mirror-theorem]{第~\ref*{chap:mirror-theorem} 章}首次出现式~\eqref{eq:quintic-iv-hypergeometric} 与~\eqref{eq:quintic-iv-period-expansion} 时，将采用 \SourceCite{LLY} 中更一般的形式。具体地，对 $\mathcal O_{\mathbb P^4}(5)$，式~\eqref{eq:hypergeometric-euler-data-function} 化为
\[
HG[\mathcal I(\widehat P)](t_1)
=e^{-pt_1/\hbar}\sum_{d\geq0}
\frac{\prod_{m=0}^{5d}(5p-m\hbar)}
{\prod_{m=1}^{d}\prod_{k=0}^{4}(p-\lambda_k-m\hbar)}e^{dt_1}.
\]
这个公式取值于等变上同调，$p$ 是等变超平面类。令 $\hbar=-1$，再令 $\lambda_k\to0$ 以回到通常上同调，得到
\[
\lim_{\lambda\to0}HG[\mathcal I(\widehat P)](t_1)
=e^{t_1H}5H\sum_{d\geq0}
\frac{\prod_{m=1}^{5d}(5H+m)}{\prod_{m=1}^{d}(H+m)^5}e^{dt_1},
\]
这恰好是此时的 $I_{\mathcal V}$。\SourceCite{LLY} 的作者选用“$HG$”这一记号，是为了强调这些函数的超几何解释。

这里引入的某些对象，将在后续章节中起重要作用。例如，\hyperref[subsec:enlarged-kahler-moduli]{第~\ref*{subsec:enlarged-kahler-moduli} 小节}将说明式~\eqref{eq:nef-gkz-q-coordinates} 中的 $q_j$，如何给出 Calabi--Yau 超曲面 $V\subset X$ 的 Kähler 模空间上的自然坐标。如前所述，\hyperref[chap:mirror-theorem]{第~\ref*{chap:mirror-theorem} 章}讨论镜像定理时，函数 $I_{\mathcal V}$ 也将占据显著位置。特别是例题~\ref{ex:toric-fano-and-hirzebruch-quantum-ring} 将说明，命题~\ref{prop:nef-i-gkz} 蕴含环面完全交量子上同调环中的一些非平凡关系。
\index{nef完全交@nef 完全交|hyperpage}\index{GKZ方程组@GKZ 方程组|hyperpage}\index{上同调值超几何级数|hyperpage}

\SourcePage{125}
\section{Yukawa 耦合}\label{sec:yukawa-couplings}\index{Yukawa耦合@Yukawa 耦合|hyperpage}

本节集中讨论 Calabi--Yau 三维簇这一特殊情形。如\hyperref[chap:introduction]{第~\ref*{chap:introduction} 章}，\emph{Yukawa 耦合}定义为
\[
\langle\theta_1,\theta_2,\theta_3\rangle
=\int_V\Omega\wedge(\nabla_{\theta_1}\nabla_{\theta_2}\nabla_{\theta_3}\Omega),
\]
其中 $\theta_i\in H^1(V,T_V)$，见式~\eqref{eq:b-yukawa}。将 $H^1(V,T_V)$ 视为模空间的切空间，并在最大幺幂边界点 $p$ 附近选取模坐标 $z_1,\ldots,z_r$，则 $\partial/\partial z_i$ 给出切空间的一组基。不过，Gauss--Manin 联络 $\nabla$ 在 $p$ 处有正则奇点，因此使用 $\delta_i=z_i\partial/\partial z_i$ 更方便。我们因而关心
\[
\langle\delta_i,\delta_j,\delta_k\rangle
=\int_V\Omega\wedge(\nabla_{\delta_i}\nabla_{\delta_j}\nabla_{\delta_k}\Omega)
\]
的计算。

\subsection{一维情形}\label{subsec:yukawa-one-dimensional}

先考虑复模空间为一维的 Calabi--Yau 三维簇 $V$。假定在一个最大幺幂边界点附近，坐标为 $z$。对 $\delta=z\,d/dz$，可将 Picard--Fuchs 方程写成
\[
\nabla_\delta^4\Omega+f_1(z)\nabla_\delta^3\Omega
+f_2(z)\nabla_\delta^2\Omega+f_3(z)\nabla_\delta\Omega+f_4(z)\Omega=0.
\]
这个方程可以确定 Yukawa 耦合，差一个常数因子。\hyperref[chap:quintic]{第~\ref*{chap:quintic} 章}已对五次镜像做过这项计算，过程很容易推广。Griffiths 横截性给出
$\int_V\Omega\wedge\nabla_\delta^2\Omega=0$；求导两次，便得到
\[
\int_V\nabla_\delta\Omega\wedge\nabla_\delta^3\Omega
=-\frac12\int_V\Omega\wedge\nabla_\delta^4\Omega,
\]
正如式~\eqref{eq:omega-identity} 附近的计算。令 $Y=\int_V\Omega\wedge\nabla_\delta^3\Omega$ 表示 Yukawa 耦合，则
\[
\begin{aligned}
\delta Y
&=\int_V\Omega\wedge\nabla_\delta^4\Omega
+\int_V\nabla_\delta\Omega\wedge\nabla_\delta^3\Omega\\
&=\frac12\int_V\Omega\wedge\nabla_\delta^4\Omega\\
&=-\frac12f_1(z)\int_V\Omega\wedge\nabla_\delta^3\Omega.
\end{aligned}
\]
最后一个等式来自 Picard--Fuchs 方程，以及由 Griffiths 横截性得到的
$\int_V\Omega\wedge\nabla_\delta^i\Omega=0$（$i\leq2$）。因此
\[
\delta Y=-\frac12f_1(z)Y.
\]
由此可解出 $Y$，差一个常数因子，正如五次镜像的情形。

\subsection{一般情形}\label{subsec:yukawa-general}

一般情形基本上使用相同的方法，只是代数运算更繁杂。为简单起见，假定 $V$ 是复模空间为二维的 Calabi--Yau 三维簇。给定一个最大幺幂边界点，选取坐标 $z_1,z_2$，使 $z_i=0$ 定义边界除子。置 $\delta_i=z_i\partial/\partial z_i$，定义
\[
K^{ab}=\int_V\Omega\wedge\nabla_{\delta_1}^{a}\nabla_{\delta_2}^{b}\Omega,
\qquad a+b=3.
\]
\SourcePage{126}
于是 $K^{ab}$ 是下列 Yukawa 耦合之一：
\[
\langle\delta_1,\delta_1,\delta_1\rangle,\quad
\langle\delta_1,\delta_1,\delta_2\rangle,\quad
\langle\delta_1,\delta_2,\delta_2\rangle,\quad
\langle\delta_2,\delta_2,\delta_2\rangle.
\]
另外，Griffiths 横截性在这里成为
\begin{equation}\label{eq:yukawa-transversality}
\int_V\Omega\wedge\nabla_{\delta_1}^{a}\nabla_{\delta_2}^{b}\Omega=0,
\qquad a+b\leq2.
\end{equation}
反复求导，可得到将
\begin{equation}\label{eq:yukawa-fourth-integrals}
W^{ab}=\int_V\Omega\wedge\nabla_{\delta_1}^{a}\nabla_{\delta_2}^{b}\Omega,
\qquad a+b=4,
\end{equation}
与 Yukawa 耦合 $K^{ab}$ 的偏导数联系起来的恒等式。再结合 Picard--Fuchs 方程，便得到一个关于 Yukawa 耦合的可解的一阶微分方程组。更多细节见 \SourceCite{HKTY1}。

为了具体说明，再回到一个熟悉的例子。
\SourceExampleNumber{5.6.2}{1}
\begin{example}\label{ex:weighted-yukawa}
$\mathbb P(1,1,2,2,2)$ 中次数为 $8$ 的\IndexLocation{indexloc:1020:126}\IndexLocation{indexloc:1030:126}\IndexLocation{indexloc:1033:126}超曲面的镜像 $V^\circ$，有二维复模空间，坐标 $z_1,z_2$ 取自例题~\ref{ex:weighted-gkz}。在那里求出了两个 Picard--Fuchs 方程~\eqref{eq:weighted-gkz-second} 与~\eqref{eq:weighted-gkz-third}，经代数运算可写成
\begin{equation}\label{eq:weighted-yukawa-pf2}
0=(1-4z_2)\delta_2^2-z_2(\delta_1^2-4\delta_1\delta_2-\delta_1+2\delta_2),
\end{equation}
\begin{equation}\label{eq:weighted-yukawa-pf3}
0=(1-256z_1)\delta_1^3-2\delta_1^2\delta_2
-8z_1(48\delta_1^2+22\delta_1+3).
\end{equation}
这些方程立即给出 $K^{ab}$ 之间的一些关系。例如，由式~\eqref{eq:weighted-yukawa-pf3} 与 Griffiths 横截性~\eqref{eq:yukawa-transversality}，容易得出
\[
0=(1-256z_1)K^{30}-2K^{21},
\]
即
\[
K^{21}=\frac12(1-256z_1)K^{30}.
\]
类似地，在式~\eqref{eq:weighted-yukawa-pf2} 的左侧分别乘以 $\delta_1,\delta_2$，再用 $K^{30}$ 解出其余耦合，得到
\[
\begin{aligned}
K^{12}&=\frac{z_2(-1+512z_1)}{1-4z_2}K^{30},\\
K^{03}&=\frac{z_2(1-256z_1+4z_2-3072z_1z_2)}{2(1-4z_2)^2}K^{30}.
\end{aligned}
\]
因此，只需确定一个 Yukawa 耦合 $K^{30}$。为此，求出它满足的一阶偏微分方程。首先，在式~\eqref{eq:weighted-yukawa-pf3} 的左侧乘以 $\delta_1$。按上面的方法，并利用 $\delta_1z_1=z_1(\delta_1+1)$，容易得到
\begin{equation}\label{eq:weighted-yukawa-fourth-relation}
(1-256z_1)W^{40}-2W^{31}-640z_1K^{30}=0,
\end{equation}
其中 $W^{ab}$ 见式~\eqref{eq:yukawa-fourth-integrals}。另外，反复对式~\eqref{eq:yukawa-transversality} 求导，得到恒等式
\[
\begin{aligned}
W^{40}&=2\delta_1K^{30},\\
W^{31}&=\frac32\delta_1K^{21}+\frac12\delta_2K^{30},\\
W^{22}&=\delta_1K^{12}+\delta_2K^{21}.
\end{aligned}
\]
\SourcePage{127}
\SourceCite{HKTY1} 给出了系统地写出全部这类恒等式的方法。将这些恒等式、上面的 $K^{21}$ 公式与式~\eqref{eq:weighted-yukawa-fourth-relation} 结合，得到
\[
(1-256z_1)\delta_1K^{30}-2\delta_2K^{30}-512z_1K^{30}=0.
\]
还可以在式~\eqref{eq:weighted-yukawa-pf2} 的左侧乘以 $\delta_1^2$，推出第二个方程：
\[
(1-4z_2)W^{22}-z_2(W^{40}-4W^{31}-K^{30}+2K^{21})=0.
\]
按相同的方法，得到
\[
-256z_1z_2\delta_1K^{30}
+\frac12(1-256z_1+1024z_1z_2)\delta_2K^{30}=0.
\]
（这里发生了大量抵消！）现在已有两个关于 $K^{30}$ 的一阶线性偏微分方程。经过代数运算，并利用 $\delta_i=z_i\partial/\partial z_i$，可改写为
\[
\begin{aligned}
\frac1{K^{30}}\frac{\partial K^{30}}{\partial z_1}
&=\frac{512(1-256z_1+1024z_1z_2)}{D(z_1,z_2)},\\
\frac1{K^{30}}\frac{\partial K^{30}}{\partial z_2}
&=\frac{512^2z_1^2}{D(z_1,z_2)},
\end{aligned}
\]
其中
\[
D(z_1,z_2)=(1-256z_1)^2-512^2z_1^2z_2.
\]
因此，容易得出 Yukawa 耦合 $K^{30}$ 的公式
\[
K^{30}=\frac{c}{D(z_1,z_2)},\qquad c\ne0.
\]
由此立即得到 $K^{21},K^{12},K^{03}$ 的公式。这些答案与 \SourceCite{HKTY1} 附录中的结果一致。

所有 Yukawa 耦合 $K^{ab}$ 的分母都涉及 $D(z_1,z_2)$ 与 $1-4z_2$。\hyperref[subsec:two-parameter-moduli]{第~\ref*{subsec:two-parameter-moduli} 小节}将证明，在 $V^\circ$ 的复模空间中，方程 $D(z_1,z_2)=0$ 与 $1-4z_2=0$ 定义判别式轨迹的两个分支；这里判别式轨迹由所有使定义方程~\eqref{eq:weighted-pf-family} 变为奇异的点 $(z_1,z_2)$ 组成。因此，这些多项式出现在 Yukawa 耦合的分母中，并不令人意外。
\end{example}

\subsection{归一化 Yukawa 耦合与 Gauss--Manin 联络}\label{subsec:normalized-yukawa-connection}\index{Gauss--Manin联络@Gauss--Manin 联络!归一化 Yukawa 耦合|hyperpage}

接下来的任务，是讨论归一化 Yukawa 耦合。如将看到的，这需要仔细选取 $3$ 形式 $\Omega$ 与模坐标。

先注意，$\Omega$ 远不是唯一的：可以乘以模空间上的任意处处非零函数。如\hyperref[chap:quintic]{第~\ref*{chap:quintic} 章}，采用以下归一化。在\IndexLocation{indexloc:0766:127}最大幺幂边界点，有 $\dim W_0=1$。由\IndexLocation{indexloc:0907:127}单值性权滤过 $W_\bullet$ 的定义，$W_0$ 在单值性作用下不变。取 $g_0$ 为 $W_0\cap H^3(V,\mathbb Z)$ 的生成元，则 $g_0$ 除符号外唯一确定。于是，$y_0=\langle g_0,\Omega\rangle$ 是在边界点解析的唯一周期，差一个常数因子。如~\hyperref[sec:quintic-b-model]{第 2.4 节}，将 $\Omega$ 换成 $\Omega/y_0$ 来归一化 $3$ 形式。\hyperref[chap:quintic]{第~\ref*{chap:quintic} 章}的论证表明，这会把 Yukawa 耦合 $Y$ 换成 $Y/y_0^2$，与式~\eqref{eq:yukawa-normalized} 前的讨论完全对应。与 \SourceCite{Morrison2} 一样，我们指出 $Y/y_0^2$ 与 $g_0$ 的选择无关。

归一化的下一步，是选取典范的模变量。这里将发现 Yukawa 耦合与 Gauss--Manin 联络 $\nabla$ 之间一种很有趣的关系。
\SourcePage{128}
为简单起见，先讨论复模为一维的 Calabi--Yau 三维簇，并假定处于最大幺幂边界点附近。如上，有 $3$ 形式 $\Omega$ 与周期 $y_0=\langle g_0,\Omega\rangle$。以 $\Omega/y_0$ 代替 $\Omega$ 归一化，同样将 $Y$ 变为 $Y/y_0^2$；归一化后，有 $\langle g_0,\Omega\rangle=1$。以下命题表明，只要选取适当局部坐标，这就会使 Picard--Fuchs 方程的形式发生显著变化。
\begin{proposition}\label{prop:normalized-yukawa-pf}
设 $V$ 是复模为一维的 Calabi--Yau 三维簇，假定所处的最大幺幂边界点满足~\hyperref[subsec:higher-dimensional-mum]{第 5.2.2 小节}的整性猜想。则存在局部坐标 $q$，使归一化 $3$ 形式 $\Omega$ 的 Picard--Fuchs 方程为
\[
\nabla_\delta^2\left(\frac{\nabla_\delta^2\Omega}{Y}\right)=0,
\]
其中 $\delta=2\pi iq\,d/dq$，$Y$ 是定义为
\[
Y=-\int_V\Omega\wedge\nabla_\delta^3\Omega
\]
的 Yukawa 耦合。
\end{proposition}
\begin{proof}
证明基于 \SourceCite{Deligne2} 的思想。先在最大幺幂边界点选取局部坐标 $z$。在 $z=0$ 的邻域，$\mathcal H=\mathcal F^0$ 上有 Hodge 滤过 $\mathcal F^\bullet$，而权滤过 $W_\bullet$ 给出 $\mathcal H$ 的滤过 $\mathcal W_\bullet$。在 $z=0$ 上，由式~\eqref{eq:mum-hodge-tate-splitting}，$(W_\bullet,F_{\mathrm{lim}}^\bullet)$ 是 Hodge--Tate 型的。如 \SourceCite[第 6 节]{Deligne2} 所指出的，这意味着在 $0$ 的邻域有
\[
\mathcal F^p\oplus\mathcal W_{2p-1}
=\mathcal F^p\oplus\mathcal W_{2p-2}=\mathcal H.
\]
因此，$(W_\bullet,\mathcal F^\bullet)$ 是 Hodge--Tate 型混合 Hodge 结构变分；对每个 $p$，$\mathcal F^p$ 在 $\operatorname{Gr}_{2p}^W$ 上诱导一维的 $(p,p)$ 型纯 Hodge 结构。置
$\mathcal I^{p,p}=\mathcal W_{2p}\cap\mathcal F^p$，则有自然同构
\[
\mathcal I^{p,p}\simeq\operatorname{Gr}_{2p}^W.
\]
由于 $\mathcal W_\bullet$ 平坦，$\nabla$ 在 $\operatorname{Gr}_{2p}^W$ 上诱导联络 $\nabla_{2p}^W$。又有
$N(\mathcal W_{2p})\subset\mathcal W_{2p-2}=\mathcal W_{2p-1}$，其中 $N$ 是单值性的对数，故 $N$ 在 $\operatorname{Gr}_{2p}^W$ 上作用为零。于是 $\nabla_{2p}^W$ 在 $\operatorname{Gr}_{2p}^W$ 上的单值性平凡，自然延拓到 $z=0$。$\operatorname{Gr}_{2p}^W$ 为一维，且有自然的 $\mathbb Z$ 结构，因此在 $0$ 的邻域可以找到张成它的整值 $\nabla_{2p}^W$ 平坦截面，除 $\pm1$ 外唯一。借助上述同构，这给出 $\mathcal I^{p,p}$ 的一个截面，记为 $e_{3-p}$。

截面 $e_0,e_1,e_2,e_3$ 有一些很好的性质。首先，$N$ 是配对
$\langle\alpha,\beta\rangle=\int_V\alpha\wedge\beta$ 的无穷小自同构。结合式~\eqref{eq:monodromy-weight-filtration} 与 $\mathcal W_4=\mathcal W_5$，得到
$\langle\mathcal W_0,\mathcal W_4\rangle=\langle\mathcal W_2,\mathcal W_2\rangle=\{0\}$。因此，只有
$\langle e_0,e_3\rangle$ 与 $\langle e_1,e_2\rangle$ 可以非零，并可在 $\operatorname{Gr}_\bullet^W$ 中计算。配对在 $\mathbb Z$ 上完美，故它们必为 $\pm1$。可以选取 $e_0,e_1,e_2,e_3$，使
\begin{equation}\label{eq:normalized-integral-pairings}
\langle e_0,e_3\rangle=-1,\qquad\langle e_1,e_2\rangle=1.
\end{equation}
这个符号选择的原因很快就会清楚。

\SourcePage{129}
现在令 $\delta=2\pi iz\,d/dz$。因为 $e_{3-p}$ 映到 $\operatorname{Gr}_{2p}^W$ 的一个 $\nabla_{2p}^W$ 平坦截面，$\nabla(e_{3-p})$ 属于 $\mathcal W_{2p-1}=\mathcal W_{2p-2}$。另一方面，$e_{3-p}$ 是 $\mathcal F^p$ 的截面，所以 $\nabla_\delta(e_{3-p})$ 属于 $\mathcal F^{p-1}$。因此，它是 $\mathcal I^{p-1,p-1}$ 的截面，从而对某些函数 $Y_1,Y_2,Y_3$，有
\[
\nabla_\delta e_0=Y_1e_1,\quad
\nabla_\delta e_1=Y_2e_2,\quad
\nabla_\delta e_2=Y_3e_3,\quad
\nabla_\delta e_3=0.
\]
由于 $\langle e_0,e_2\rangle$ 恒为零，
\[
\begin{aligned}
0=\delta\langle e_0,e_2\rangle
&=\langle\nabla_\delta e_0,e_2\rangle+\langle e_0,\nabla_\delta e_2\rangle\\
&=\langle Y_1e_1,e_2\rangle+\langle e_0,Y_3e_3\rangle=Y_1-Y_3,
\end{aligned}
\]
最后一步用了式~\eqref{eq:normalized-integral-pairings}，故 $Y_1=Y_3$。

$\nabla$ 有正则奇点，而 $\delta=2\pi iz\,d/dz$，所以 $Y_3$ 在 $0$ 处全纯。下面证明 $Y_3(0)=\pm1$。由式~\eqref{eq:connection-monodromy-residue}，在 $0$ 上，把它视为
$F_{\mathrm{lim}}^1/F_{\mathrm{lim}}^2\to F_{\mathrm{lim}}^0/F_{\mathrm{lim}}^1$ 的映射时，有
\[
\overline\nabla_\delta=-N.
\]
于是 $\overline\nabla_\delta(e_2)=-N(e_2)$。又由
$\nabla_\delta(e_2)=Y_3e_3$，有
$\overline\nabla_\delta(e_2)=Y_3(0)e_3$，因而
$N(e_2)=-Y_3(0)e_3$。但在 $0$ 上，$e_3$ 是 $W_0$ 的一组整基，$e_2,e_3$ 是 $W_2$ 的一组整基。由~\hyperref[subsec:higher-dimensional-mum]{第 5.2.2 小节}末的整性猜想，$N(e_2)=\pm e_3$。所以 $Y_3(0)=\pm1$，正如所需。

必要时将 $e_1,e_2$ 换成 $-e_1,-e_2$，便可以假定 $N(e_2)=-e_3$；注意，式~\eqref{eq:normalized-integral-pairings} 不变。于是 $Y_3(0)=1$，函数
\[
q=\exp\left(\int Y_3(z)\frac{dz}{z}\right)
\]
给出 $0$ 处的局部坐标，满足
\[
Y_3(z)\,2\pi iq\frac d{dq}=2\pi iz\frac d{dz}.
\]
重新将 $\delta$ 定义为 $2\pi iq\,d/dq$，则 $\nabla_\delta(e_2)=e_3$。因为 $Y_1=Y_3$，也有 $\nabla_\delta(e_0)=e_1$。因此，在基 $e_0,e_1,e_2,e_3$ 下，以 $q$ 为局部坐标，$\nabla_\delta$ 的\IndexLocation{indexloc:0174:129}联络矩阵为
\begin{equation}\label{eq:normalized-yukawa-matrix}
\begin{pmatrix}
0&0&0&0\\1&0&0&0\\0&Y&0&0\\0&0&1&0
\end{pmatrix},
\end{equation}
其中 $Y=Y_2$。最后将把 $Y$ 与归一化 Yukawa 耦合等同。

下一步证明，$e_0$ 就是归一化 $3$ 形式 $\Omega$。$e_0$ 是 $\mathcal F^3$ 的处处非零截面，因此对某个全纯函数 $a(q)$，有 $\Omega=a(q)e_0$。上述矩阵给出
\[
\nabla_\delta^j\Omega=\delta^j(a)e_0\pmod{e_1,e_2,e_3},
\]
所以 $a(q)$ 满足与 $\Omega$ 相同的 Picard--Fuchs 方程。在最大幺幂边界点，除常数倍外只有一个全纯解；由于 $\Omega$ 已归一化，$\langle g_0,\Omega\rangle=1$ 就是一个解。这说明 $a(q)$ 是常数 $a$，故 $\Omega=ae_0$。回顾，$g_0$ 是 $W_0$ 的整基，所以可以假定 $e_3=g_0$。配对为反对称，因而
\[
1=\langle g_0,\Omega\rangle=\langle e_3,ae_0\rangle
=-a\langle e_0,e_3\rangle=a,
\]
\SourcePage{130}
其中最后一个等号来自式~\eqref{eq:normalized-integral-pairings}。这就证明了 $e_0=\Omega$。由此很容易得到所需的 Picard--Fuchs 方程：
\begin{equation}\label{eq:normalized-yukawa-pf}
\nabla_\delta^2\left(\frac{\nabla_\delta^2\Omega}{Y}\right)=0.
\end{equation}
证明的最后一步是说明 $Y$ 就是 Yukawa 耦合。这很容易，因为已经证明 $\Omega=e_0$；于是式~\eqref{eq:normalized-yukawa-matrix} 给出 $\nabla_\delta^3\Omega=\delta(Y)e_2+Ye_3$，再由式~\eqref{eq:normalized-integral-pairings} 得到
\[
\int_V\Omega\wedge\nabla_\delta^3\Omega
=\int_V e_0\smile\bigl(\delta(Y)e_2+Ye_3\bigr)
=\langle e_0,\delta(Y)e_2+Ye_3\rangle=-Y.
\]
命题得证。
\end{proof}

证明过程中如此仔细地处理了所有符号，最后 $Y$ 的公式却出现一个负号，这可能令人觉得奇怪。不过，这个负号完全自然，因为在等式
\begin{equation}\label{eq:normalized-yukawa-sign}
Y=-\int_V\Omega\wedge\nabla_\delta^3\Omega
\end{equation}
中，右端就是使 $H^3(V,\mathbb C)$ 上的 Hodge 结构极化的配对
$Q(\alpha,\beta)=(-1)^{3(3-1)/2}\int_V\alpha\smile\beta$
（见式~\eqref{eq:hodge-polarization}）。按照镜像对称，$Y$ 应对应镜像 $V^\circ$ 上的某个 A 模型关联函数；后者也可以用 \emph{A 模型 Hodge 结构变分}上的极化来定义，我们将在\hyperref[sec:a-model-vhs]{第~\ref*{sec:a-model-vhs} 节}讨论这一点。比较这两个极化 Hodge 结构变分时，就会明白为什么需要这个符号。用\hyperref[subsec:normalized-yukawa-mirror]{第 5.6.4 小节}的术语来说，式~\eqref{eq:normalized-yukawa-sign} 是一维模空间情形的\emph{归一化 Yukawa 耦合}。

\SourceCite{Deligne2} 中对命题~\ref{prop:normalized-yukawa-pf} 的证明比上述证明更为精巧。Deligne 从原来的局部坐标 $z$ 出发，固定 $1\leq i\leq3$，考虑 Hodge 结构的扩张
\begin{equation}\label{eq:hodge-extension-coordinates}
0\longrightarrow\operatorname{Gr}_{2i-2}^W
\longrightarrow\mathcal W_{2i}/\mathcal W_{2i-4}
\longrightarrow\operatorname{Gr}_{2i}^W\longrightarrow0.
\end{equation}
在每一点，所得扩张由 $\mathbb C^*$ 中的一个数分类。因此，当 $(\mathcal W_\bullet,\mathcal F^\bullet)$ 在 $z=0$ 的一个邻域内变化时，扩张就由一个函数 $q_i(z)$ 分类。前面记为 $Y_1,Y_2,Y_3$ 的函数可用 $q_i$ 表示为
\[
Y_{4-i}=\frac{z}{q_i}\frac{dq_i}{dz},\qquad1\leq i\leq3.
\]
（\SourceCite{Deligne2} 的第 7 节中出现一个负号，因为那里处理的是对偶情形。）所以命题~\ref{prop:normalized-yukawa-pf} 中的特殊坐标 $q$ 就是 $q_1=q_3$，它分别分类 $i=1,3$ 时的扩张~\eqref{eq:hodge-extension-coordinates}。

命题~\ref{prop:normalized-yukawa-pf} 可以很好地推广到 Calabi--Yau 三维簇具有 $r$ 维模空间的情形。假定\hyperref[subsec:higher-dimensional-mum]{第 5.2.2 小节}的整性猜想成立，那么在一个最大幺幂边界点附近，就可以如下描述 Gauss--Manin 联络。我们把该边界点视为 $0\in\Delta^r$。

与前面一样，$I^{p,p}=\mathcal W_{2p}\cap\mathcal F^p\simeq\operatorname{Gr}_{2p}^W$，不过现在 $I^{1,1}$ 和 $I^{2,2}$ 的秩为 $r$。命题证明中使用的基 $e_0,e_1,e_2,e_3$ 相应地推广为以下各空间的基：
\[
e_0\in I^{3,3},\quad e_j\in I^{2,2}\ (1\leq j\leq r),\quad
 e^k\in I^{1,1}\ (1\leq k\leq r),\quad e^0\in I^{0,0},
\]
使得
\begin{equation}\label{eq:multimoduli-integral-pairings}
\langle e_0,e^0\rangle=-1,\qquad
\langle e_j,e^k\rangle=\delta_{jk},\quad1\leq j,k\leq r,
\end{equation}
\SourcePage{131}
其余配对全部为零。注意，式~\eqref{eq:normalized-integral-pairings} 正是 $r=1$ 的特殊情形。此外，与命题~\ref{prop:normalized-yukawa-pf} 的证明中一样，可以假定这些基元素延拓到 $0$，并且在那里具有整性。

证明命题~\ref{prop:normalized-yukawa-pf} 所用的方法（或 \SourceCite{Deligne2} 中更强的方法）于是表明，可以找到 $0\in\Delta^r$ 的特殊局部坐标 $q_1,\ldots,q_r$，使得对 $\delta_j=2\pi i q_j\partial/\partial q_j$，有
\begin{equation}\label{eq:multimoduli-connection}
\begin{aligned}
\nabla_{\delta_i}e^0&=0,\\
\nabla_{\delta_i}e^k&=\delta_{ik}e^0,&&1\leq k\leq r,\\
\nabla_{\delta_i}e_j&=\sum_{k=1}^rY_{ijk}e^k,&&1\leq j\leq r,\\
\nabla_{\delta_i}e_0&=e_i,
\end{aligned}
\end{equation}
其中 $Y_{ijk}$ 是一些函数。这说明了矩阵~\eqref{eq:normalized-yukawa-matrix} 如何推广到 $r$ 维模空间的情形。

由命题~\ref{prop:normalized-yukawa-pf} 的论证，知道 $e_0$ 就是归一化 $3$ 形式 $\Omega$，于是式~\eqref{eq:multimoduli-connection} 表明
\begin{equation}\label{eq:multimoduli-yukawa}
Y_{ijk}=-\int_V\Omega\wedge\nabla_{\delta_k}\nabla_{\delta_i}\nabla_{\delta_j}\Omega.
\end{equation}
在\hyperref[subsec:normalized-yukawa-mirror]{第 5.6.4 小节}中，将把 $Y_{ijk}$ 定义为\emph{归一化 Yukawa 耦合}。上述讨论于是说明，Gauss--Manin 联络由归一化 Yukawa 耦合唯一确定。

下面证明，函数 $Y_{ijk}$ 都来自同一个“势函数”。引入变量 $t_j$，使 $q_j=e^{2\pi i t_j}$，于是 $\partial/\partial t_j=2\pi i q_j\partial/\partial q_j=\delta_j$。有如下引理。
\begin{lemma}\label{lem:yukawa-potential}\index{Yukawa 势函数|hyperpage}
存在全纯函数 $\Phi(t_1,\ldots,t_r)$，使得对所有 $i,j,k$，
\[
Y_{ijk}=\frac{1}{(2\pi i)^3}\frac{\partial^3\Phi}{\partial t_i\partial t_j\partial t_k}.
\]
\end{lemma}
\begin{proof}
关键是 $\nabla$ 为平坦联络，因此复合映射
\[
\mathcal H\xrightarrow{\nabla}\mathcal H\otimes\Omega_S^1
\xrightarrow{\nabla}\mathcal H\otimes\Omega_S^2
\]
为零（这里 $S=(\Delta^*)^r$）。利用式~\eqref{eq:multimoduli-connection}，计算得
\[
0=\nabla^2(e_0)=\sum_{i,j,k=1}^rY_{ijk}e^k\otimes dt_i\wedge dt_j.
\]
这表明，对所有 $i,j,k$，$Y_{ijk}=Y_{jik}$。此外，对 $1\leq j\leq r$，又可计算出
\[
0=\nabla^2(e_j)=\sum_{i,k=1}^rY_{ijk}e^0\otimes dt_k\wedge dt_i
+\sum_{i,k,\ell=1}^r\frac{\partial Y_{ijk}}{\partial t_\ell}e^k\otimes dt_\ell\wedge dt_i.
\]
这首先表明，对所有 $i,j,k$，$Y_{ijk}=Y_{kji}$；其次表明，对所有 $i,j,k,\ell$，
\[
\frac{\partial Y_{ijk}}{\partial t_\ell}
=\frac{\partial Y_{\ell jk}}{\partial t_i}.
\]
所以 $Y_{ijk}$ 关于 $i,j,k$ 完全对称。再利用上述偏导数条件，很容易推出 $\Phi$ 的存在性。
\end{proof}

\SourcePage{132}
注意，如果用式~\eqref{eq:multimoduli-connection} 定义 $\nabla$，那么势函数 $\Phi$ 的存在性等价于 $\nabla$ 的平坦性。

这样就证明了：在满足整性猜想的最大幺幂边界点 $0\in\Delta^r$ 附近，Calabi--Yau 三维簇族的 Gauss--Manin 联络由一个函数 $\Phi$ 确定。由于 $\Phi$ 是 $t_j=\frac1{2\pi i}\log q_j$ 的函数，它在 $(\Delta^*)^r$ 上是多值的；不过，它的三阶导数就是归一化 Yukawa 耦合，而后者是单值的，并全纯延拓到 $\Delta^r$。此外，利用 $\Phi$，还可以显式写出 $H^3(V,\mathbb C)$ 上的 Hodge 结构变分的许多有趣性质。例如，在式~\eqref{eq:a-model-flat-sections-threefold} 中将看到，可以用势函数 $\Phi$ 写出一组多值平坦截面的基。

真正令人惊讶的一点是，镜像一侧存在 $\Phi$ 的一个自然对应物，称为 \emph{Gromov--Witten 势}。将在\hyperref[sec:a-model-vhs]{第~\ref*{sec:a-model-vhs} 节}和\hyperref[sec:mirror-symmetry-revisited]{第~\ref*{sec:mirror-symmetry-revisited} 节}详细探究这种对应关系。

\subsection{归一化 Yukawa 耦合与镜像映射}\label{subsec:normalized-yukawa-mirror}
命题~\ref{prop:normalized-yukawa-pf} 表明，只要使用归一化 $3$ 形式及正确的局部坐标，Gauss--Manin 联络就由 Yukawa 耦合确定。现在出现了一个奇迹：\emph{那个神秘的坐标 $q$ 正是镜像映射！}镜像映射的一般定义要到\hyperref[sec:mirror-map]{第~\ref*{sec:mirror-map} 节}才给出，不过，沿用\hyperref[chap:quintic]{第 2 章}的做法，就能理解一维模空间的情形。回顾该章公式~\eqref{eq:mirror-map-periods}，它将镜像映射写成
\begin{equation}\label{eq:mirror-map-period-ratio}
\widehat q=\exp\left(2\pi i\frac{\int_{\gamma_1}\Omega}{\int_{\gamma_0}\Omega}\right),
\end{equation}
其中 $\int_{\gamma_0}\Omega$ 为全纯函数，而绕 $z=0$ 一周时，$\int_{\gamma_1}\Omega/\int_{\gamma_0}\Omega$ 增加 $1$。令 $g_0,g_1$ 为 $\gamma_0,\gamma_1$ 的 Poincaré 对偶，就得到 $g_0\in W_0$、$g_1\in W_2$ 以及 $N(g_1)=g_0$，其中 $N$ 为单值性的对数。所以 $g_0,g_1$ 满足最大幺幂单值性定义中的条件，而 $\widehat q$ 的公式变成
\begin{equation}\label{eq:mirror-map-normalized-period}
\widehat q=\exp\bigl(2\pi i\langle g_1,\Omega\rangle\bigr),
\end{equation}
因为 $\Omega$ 的归一化意味着 $\langle g_0,\Omega\rangle=1$。

可以如下证明，这就是命题~\ref{prop:normalized-yukawa-pf} 中的坐标 $q$。取命题证明中构造的基 $e_0,e_1,e_2,e_3$；该证明还给出了 $g_0=e_3$、$\Omega=e_0$ 和 $\nabla_\delta(e_2)=e_3$。由此很容易看出
\[
g_1=-e_2+\frac{\log q}{2\pi i}e_3
\]
为平坦截面，这里的 $q$ 是命题~\ref{prop:normalized-yukawa-pf} 中的坐标。逆时针绕原点一周，$\frac{\log q}{2\pi i}$ 增加 $1$，所以 $\mathcal T(g_1)=g_1+g_0$，从而 $N(g_1)=g_0$。按定义，$e_2,e_3$ 分别属于 $W_2,W_0$，于是 $g_1$ 是 $W_2$ 中满足 $N(g_1)=g_0$ 的平坦截面。这就是所需的 $g_1$，由此得到
\[
\langle g_1,\Omega\rangle=\frac1{2\pi i}\log q.
\]
因此，坐标 $q$ 就是式~\eqref{eq:mirror-map-normalized-period} 定义的镜像映射 $\widehat q$。

更一般地，在\hyperref[chap:moduli]{第~\ref*{chap:moduli} 章}给出镜像映射的一般定义以后，将会看到：对具有 $r$ 维模空间的 Calabi--Yau 三维簇，镜像映射正是前面推广命题~\ref{prop:normalized-yukawa-pf} 时，在式~\eqref{eq:multimoduli-integral-pairings} 和式~\eqref{eq:multimoduli-connection} 的讨论中出现的特殊坐标 $q_1,\ldots,q_r$。

\SourcePage{133}
暂且接受这一结论，将 $q_1,\ldots,q_r$ 称为\emph{镜像坐标}；利用这些坐标，终于可以给出归一化 Yukawa 耦合的正式定义。
\begin{definition}\label{def:normalized-yukawa}
设有一个 Calabi--Yau 三维簇族，它在某个\IndexLocation{indexloc:0766:133}最大幺幂边界点满足整性猜想。若 $q_1,\ldots,q_r$ 是上述镜像坐标，$\Omega$ 是归一化 $3$ 形式，则\emph{归一化 Yukawa 耦合}定义为
\[
Y_{ijk}=-\int_V\Omega\wedge\nabla_{\delta_k}\nabla_{\delta_i}\nabla_{\delta_j}\Omega,
\qquad \delta_j=2\pi i q_j\frac{\partial}{\partial q_j}.
\]
\end{definition}
\index{Yukawa耦合@Yukawa 耦合!归一化|hyperpage}\index{镜像坐标|hyperpage}
$Y_{ijk}$ 也常称为 \emph{B 模型关联函数}。有一个情形，可以显式计算归一化 Yukawa 耦合。

\SourceExampleNumber{5.6.4}{1}
\begin{example}\label{ex:quintic-normalized-yukawa}
令 $V^\circ$ 为五次三维簇的镜像。由 \SourceCite{Morrison2}，知道 $V^\circ$ 满足整性猜想，所以它的归一化 Yukawa 耦合 $Y$ 有定义。我们断言，$Y$ 由下式给出：
\begin{equation}\label{eq:quintic-normalized-yukawa}
Y=\frac{5}{(1+5^5z)y_0(z)^2}\left(\frac qz\frac{dz}{dq}\right)^3,
\end{equation}
其中 $z$ 是通常的模变量，且
\[
\begin{aligned}
y_0(z)&=\sum_{n=0}^\infty\frac{(5n)!}{(n!)^5}(-1)^nz^n,\\
y_1(z)&=y_0(z)\log(-z)+5\sum_{n=1}^\infty\frac{(5n)!}{(n!)^5}
\left[\sum_{j=n+1}^{5n}\frac1j\right](-1)^nz^n,\\
q&=\exp(y_1/y_0).
\end{aligned}
\]
注意，式~\eqref{eq:quintic-normalized-yukawa} 正是\hyperref[chap:quintic]{第 2 章}五次三维簇镜像对称陈述~\eqref{eq:quintic-mirror-identity} 中出现的 B 模型关联函数。

\hyperref[subsec:hypergeometric-mirror-map]{第~\ref*{subsec:hypergeometric-mirror-map} 小节}将证明 $y_0,y_1$ 是\IndexLocation{indexloc:1176:133}\IndexLocation{indexloc:1178:133}\IndexLocation{indexloc:1180:133}五次镜像的 Picard--Fuchs 方程的解。\SourceCite{CdGP} 用公式 $q=\exp(y_1/y_0)$ 定义镜像映射，而 \SourceCite{Morrison2} 附录 C 证明了它与式~\eqref{eq:mirror-map-period-ratio} 的定义一致。因此 $q$ 就是镜像坐标。

若 $\Omega$ 为式~\eqref{eq:quintic-residue} 定义的 $3$ 形式，那么将\hyperref[subsec:yukawa-one-dimensional]{第 5.6.1 小节}的方法应用于\hyperref[ex:quintic-pf-method]{例题~\ref*{ex:quintic-pf-method}} 中的 Picard--Fuchs 方程~\eqref{eq:quintic-mum-pf}，便得到
\[
\int_{V^\circ}\Omega\wedge\nabla_{z\frac d{dz}}\nabla_{z\frac d{dz}}\nabla_{z\frac d{dz}}\Omega
=\frac1{1+5^5z}\times\text{某个常数},
\]
与\hyperref[chap:quintic]{第~\ref*{chap:quintic} 章}的结果相同。归一化就是用 $\widetilde\Omega=\Omega/\langle g_0,\Omega\rangle$ 替代 $\Omega$。由于 $y_0$ 和 $\langle g_0,\Omega\rangle$ 都是 Picard--Fuchs 方程的全纯解，$y_0$ 是 $\langle g_0,\Omega\rangle$ 的一个非零常数倍，因此，如\hyperref[subsec:normalized-yukawa-connection]{第~\ref*{subsec:normalized-yukawa-connection} 小节}所解释的，
\[
\int_{V^\circ}\widetilde\Omega\wedge\nabla_{z\frac d{dz}}\nabla_{z\frac d{dz}}\nabla_{z\frac d{dz}}\widetilde\Omega
=\frac1{(1+5^5z)y_0(z)^2}\times\text{另一个常数}.
\]
最后，从 $z$ 换到镜像坐标 $q$ 时，有
\[
\delta=2\pi i q\frac d{dq}=2\pi i\frac qz\frac{dz}{dq}\,z\frac d{dz},
\]
\SourcePage{134}
所以由 Griffiths 横截性，归一化 Yukawa 耦合为
\begin{equation}\label{eq:quintic-yukawa-undetermined}
\begin{aligned}
Y&=-\int_{V^\circ}\widetilde\Omega\wedge\nabla_\delta\nabla_\delta\nabla_\delta\widetilde\Omega\\
&=\frac{c}{(1+5^5z)y_0(z)^2}\left(\frac qz\frac{dz}{dq}\right)^3,
\end{aligned}
\end{equation}
其中 $c$ 是某个非零常数。因此，只要证明 $c=5$，就得到式~\eqref{eq:quintic-normalized-yukawa}。在 $q=0$ 处取上式的值，得到 $Y(0)=c$，所以只需证明 $Y(0)=5$。

用\hyperref[chap:quintic]{第~\ref*{chap:quintic} 章}的记号，常数 $c$ 被写成 $(2\pi i)^3c_2$；在式~\eqref{eq:quintic-mirror-yukawa} 后面的讨论中，我们令 $c_2=5/(2\pi i)^3$，以得到镜像对称所预言的正确首项。不过，在本章中，可以通过计算归一化 Yukawa 耦合，\emph{预先预测这个常数}。更确切地说，将结合本章前面建立的 Hodge 理论与 \SourceCite{CdGP,Morrison2} 中的显式计算，求出 $Y(0)$。

需要使用命题~\ref{prop:normalized-yukawa-pf} 证明中的记号。回顾，$e_0=\widetilde\Omega$ 是典范延拓 $\widetilde{\mathcal H}$ 的一组基 $e_0,e_1,e_2,e_3$ 的一部分，并且这组基在 $q=0$ 处是整结构的一组基。

现在考察 $q=0$ 处的 $\overline\nabla_\delta^3(e_0)$。首先，由式~\eqref{eq:connection-monodromy-residue}，知道在 $q=0$ 处，极限 Hodge 滤过的分次部分上有
\[
\overline\nabla_\delta=-N,
\]
所以模去 $F_{\mathrm{lim}}^1$ 后，$\overline\nabla_\delta^3(e_0)(0)=(-N)^3(e_0(0))$。另一方面，由联络矩阵~\eqref{eq:normalized-yukawa-matrix}，还有 $\overline\nabla_\delta^3(e_0)=Ye_3+\delta(Y)e_2$。合起来便得到
\[
N^3(e_0(0))=-Y(0)e_3(0).
\]
令 $s_i$ 为 $\mathcal H$ 的平坦截面，使得 $\exp(-t_1N)s_i$ 在 $q=0$ 处的值为 $e_i(0)$。由于 $e_i$ 在 $q=0$ 处是一组整基，$s_i$ 就是 $\mathcal H_{\mathbb Z}$ 的一组整基。由前一等式得
\begin{equation}\label{eq:quintic-cubic-monodromy}
N^3(s_0)=-Y(0)s_3.
\end{equation}
最后，使用 \SourceCite{CdGP} 中\IndexLocation{indexloc:1174:134}\IndexLocation{indexloc:1180:134}五次镜像族单值性的显式计算。如 \SourceCite{Morrison2} 附录 C 所解释的，$\mathcal H$ 有整平坦截面 $\beta^2,\alpha_2$，它们属于一组 $\mathbb Z$ 基，满足 $\langle\alpha_2,\beta^2\rangle=1$ 及 $N^3(\beta^2)=-5\alpha_2$。（\SourceCite{Morrison2} 中写成正号，因为那里的边界点是 $\infty$，这里却是 $0$；把 $\infty$ 变到 $0$ 的反演映射反转了相应单值性群生成元的定向。）

比较式~\eqref{eq:quintic-cubic-monodromy} 与 $N^3(\beta^2)=-5\alpha_2$，知道 $s_3,\alpha_2$ 是 $W_0$ 的整生成元，因此必要时同时将 $\beta^2,\alpha_2$ 乘以 $-1$，即可假定 $\alpha_2=s_3$。由于 $N^3:W_6/W_5\simeq W_0$，且 $N^3(Y(0)\beta^2-5s_0)=0$，推知 $Y(0)\beta^2-5s_0\in W_5$。于是 $\langle W_5,W_0\rangle=0$ 表明
$\langle Y(0)\beta^2-5s_0,s_3\rangle=0$。另一方面，已经知道 $\langle\alpha_2,\beta^2\rangle=1$，而式~\eqref{eq:normalized-integral-pairings} 给出 $\langle s_0,s_3\rangle=-1$。立刻得到 $Y(0)=5$，这正是所需的结论。
\end{example}

关于刚才的例题，还想补充两点。首先，$Y(0)=5$ 与式~\eqref{eq:quintic-cubic-monodromy} 表明
\[
N^3(s_0)=-5s_3.
\]
从镜像对称的角度看，这一等式完全合乎情理。在\hyperref[sec:a-model-vhs]{第~\ref*{sec:a-model-vhs} 节}研究五次三维簇的 A 模型联络时，将看到该联络的单值性 $N$ 自然同构于
\SourcePage{135}
与 $-H$ 的杯积，其中 $H$ 是五次三维簇的超平面截面。所以上述单值性等式与 $H^3=5$ 一致。

其次，证明式~\eqref{eq:quintic-normalized-yukawa} 需要对五次镜像族有\emph{极为详尽的了解}。除本章讨论的计算外，还需要明确的单值性信息，才能验证整性猜想、计算镜像映射并确定 $Y(0)$。

这番讨论说明：归一化 Yukawa 耦合的定义虽然有很好的理论性质，实际计算却十分困难。更常见的做法是，只对 Yukawa 耦合进行部分归一化，即将 $3$ 形式 $\Omega$ 归一化，并换成镜像坐标，而暂不确定在 $0$ 处的值。在这种情形中，只能将 $Y_{ijk}$ 确定到一个常数；不过，由引理~\ref{lem:yukawa-potential}，需要处理的常数只有一个。在表述镜像定理的环面版本时，将采取这一做法。

本节从 Yukawa 耦合开始，随后讨论 Picard--Fuchs 方程及联络矩阵，最后引出了镜像映射和归一化 Yukawa 耦合的概念。这一发展颇为出人意料，而后面的故事还会更加精彩。在\hyperref[chap:quantum-cohomology]{第~\ref*{chap:quantum-cohomology} 章}，将研究由量子上同调构造的 \emph{A 模型联络}。特别是，例题~\ref{ex:a-model-yukawa-connection} 将说明：若 Calabi--Yau 三维簇 $V^\circ$ 满足 $h^2(V^\circ)=1$，那么中间上同调 $\bigoplus_pH^{p,p}(V^\circ)$ 有一组基，使 A 模型联络的矩阵\emph{恰好}就是式~\eqref{eq:normalized-yukawa-matrix}，只是 $Y$ 此时为\hyperref[sec:three-point]{第 1.2 节}讨论的 A 模型关联函数。\hyperref[chap:quantum-differential-equations]{第~\ref*{chap:quantum-differential-equations} 章}中的情形更加丰富；例题~\ref{ex:cy-quantum-equation-relation} 从\emph{量子微分方程}的角度讨论 $V^\circ$，并导出算子
\[
\left(\hbar\frac d{dt_1}\right)^2
\left(\frac{\left(\hbar\frac d{dt_1}\right)^2}{Y}\right),
\]
这里 $Y$ 仍是 $V^\circ$ 的 A 模型关联函数。它与 Picard--Fuchs 方程~\eqref{eq:normalized-yukawa-pf} 极为相似（特别是令 $q=e^{2\pi i t_1}$，从而 $\delta=d/dt_1$ 时）。最后，当 Calabi--Yau 三维簇 $V$ 为环面超曲面，$V^\circ$ 为其 Batyrev 镜像时，\hyperref[subsec:toric-mirror-conjecture]{第~\ref*{subsec:toric-mirror-conjecture} 小节}的 Hodge \IndexLocation{indexloc:0795:135}理论\IndexLocation{indexloc:0791:135}环面镜像猜想粗略地断言：在镜像映射给出的坐标变换下，$H^3(V,\mathbb C)$ 上的 Gauss--Manin 联络与 $\bigoplus_pH^{p,p}(V^\circ)$ 上的 A 模型联络一致。一个直接推论就是两侧 Yukawa 耦合的相等。

不过，要达到这一步，还需要做许多工作。首先，必须仔细定义镜像映射。这需要详细研究复结构模及 Kähler 模。复结构模空间有许多紧化，因此可能有许多最大幺幂边界点可供选择；Kähler 一侧的情形同样复杂。\hyperref[chap:moduli]{第~\ref*{chap:moduli} 章}将探究这些问题。

另一个尚缺的要素，是仔细定义上面提到的 A 模型关联函数或三点函数。为此，需要深入研究 Gromov--Witten 不变量。\hyperref[chap:gromov-witten]{第~\ref*{chap:gromov-witten} 章}将讨论这些不变量及其与枚举几何的关系。再将这些内容与\hyperref[chap:quantum-cohomology]{第~\ref*{chap:quantum-cohomology} 章}的量子上同调结合，就具备了陈述\IndexLocation{indexloc:0841:135}镜像对称数学版本所需的全部要素。\hyperref[chap:mirror-theorem]{第~\ref*{chap:mirror-theorem} 章}的主要主题，就是证明其中一些特殊情形。


\chapter{模空间}\label{chap:moduli}
\SourcePage{136}
本章的目标是描述复结构模空间与 Kähler 模空间，并解释如何在 Calabi--Yau 流形 $V$ 的 Kähler 模空间与其镜像 $V^\circ$ 的复结构模空间之间构造镜像映射。

本章的讨论将揭示，复结构模与 Kähler 模具有出人意料的丰富内容。我们将特别关注环面簇中的 Calabi--Yau 超曲面，因为在这一情形，可以显式描述这两类模空间的形态，尽管其中有些结论仍属猜想。实际上，只研究这些模空间的多项式部分和环面部分，但它们已经足够有趣。

复结构模与 Kähler 模之间的密切联系，通过镜像映射建立起来；镜像映射将在本章末定义。困难在于，不了解 Kähler 模就无法充分理解复结构模，反之亦然。此外，理解两者都需要了解镜像映射，而要充分理解镜像映射，又必须等到\hyperref[chap:quantum-cohomology]{第~\ref*{chap:quantum-cohomology} 章}和\hyperref[chap:mirror-theorem]{第~\ref*{chap:mirror-theorem} 章}研究量子上同调及量子微分方程以后。因此，我们的讨论虽然表面上依次从复结构模走向 Kähler 模，再走向镜像映射，实际上却会在这些主题之间往返穿插。

\section{复结构模}\label{sec:complex-moduli}\index{复结构模空间|hyperpage}
本节讨论 Calabi--Yau 流形的复结构模。\SourceCite{Viehweg} 表明，\emph{极化} Calabi--Yau 流形 $(V,\mathcal L)$ 的模可由光滑拟射影簇参数化。我们将关注该模空间的一个连通分支 $\mathcal M$。

目标是为 $\mathcal M$ 找到一个方便使用的紧化 $\overline{\mathcal M}$。我们采取实用的立场，因此会灵活理解“方便”的含义。最低限度的要求是：在与镜像对称有关的最大幺幂边界点附近，紧化边界除子具有正规交叉。由奇点消解，具有正规交叉边界的紧化总是存在，却远不唯一。不过，为什么这样的紧化应具有最大幺幂单值性，目前并不清楚。

先讨论最大幺幂单值性，随后主要关注\hyperref[sec:batyrev-mirror]{第 4.1 节}所述 Batyrev 构造得到的 Calabi--Yau 超曲面的模空间，再详细说明一些例子。

\subsection{最大幺幂边界点}\label{subsec:mum-boundary-points}
设在去心多圆盘 $(\Delta^*)^r\subset\Delta^r$ 上有一个 Calabi--Yau 流形族，且在原点的单值性为最大幺幂。这就是所谓的\emph{最大幺幂边界点}。令 $z_1,\ldots,z_r$ 为 $\Delta^r$ 的坐标，$g_0,\ldots,g_r$ 如最大幺幂单值性的定义~\ref{def:maximal-unipotent-point} 所述。
\SourcePage{137}
再用下式定义全纯函数 $q_1,\ldots,q_r$ \SourceCite{Morrison6}：
\begin{equation}\label{eq:mum-canonical-coordinates}
\frac1{2\pi i}\log q_k=\frac1{\langle g_0,\Omega\rangle}
\sum_{j=1}^r\langle g_j,\Omega\rangle m^{jk},
\end{equation}
其中 $(m^{jk})$ 为定义~\ref{def:maximal-unipotent-point} 中矩阵 $(m_{ij})$ 的逆矩阵。

计算 $q_k$ 在原点附近的行为。由 $N_i g_j=m_{ij}g_0$，注意到沿 $z_i=0$ 有
$\langle g_j,\Omega\rangle\sim m_{ij}\langle g_0,\Omega\rangle\log(z_i)/(2\pi i)$。
把所有 $z_i$ 对应的渐近式代入式~\eqref{eq:mum-canonical-coordinates}，得到在原点附近 $q_k\sim cz_k$，其中 $c\ne0$ 为某个常数。因此，$(q_1,\ldots,q_r)$ 在最大幺幂边界点附近定义局部坐标。此外，这些坐标几乎唯一。若把 $g_0,\ldots,g_r$ 换成 $g'_j=\sum_{k=0}^r c_{jk}g_k$，就得到新的局部坐标 $(q'_1,\ldots,q'_r)$；可以证明
\begin{equation}\label{eq:mum-coordinate-basis-change}
q'_k=e^{2\pi i(c^{k0}/c_{00})}q_k,\qquad
c^{k0}=\sum_{j=1}^r c_{j0}m'^{jk},
\end{equation}
其中 $(m'^{jk})$ 是由 $N_j(g'_k)=m'_{jk}g'_0$ 确定的矩阵的逆矩阵。\hyperref[sec:mirror-map]{第~\ref*{sec:mirror-map} 节}将讨论如何唯一确定 $q_k$。

遗憾的是，上述“唯一性”假定始终使用 $(\Delta^*)^r$ 的同一个正规交叉紧化。因此，需要确定改用另一紧化时，$q_k$ 如何变化。只需局部处理此事，所以考虑全纯映射 $\phi:\Delta^r\to\Delta^r$，在源和目标的去心部分采用环面坐标；假定它在稠密环面上为同构，并把坐标超平面映到坐标超平面组成的边界内。必要时缩小圆盘后，很容易看出：可以在源上选择坐标 $w=(w_1,\ldots,w_r)$，在目标上选择坐标 $z=(z_1,\ldots,z_r)$，使映射为
\begin{equation}\label{eq:mum-monomial-coordinate-change}
z_j=\phi_j(w)=\prod_i w_i^{a_{ij}},
\end{equation}
其中 $a_{ij}$ 为非负整数，矩阵 $(a_{ij})$ 在 $\mathbb Z$ 上可逆。

我们断言，源中的 $w=0$ 是最大幺幂边界点，当且仅当目标中的 $z=0$ 是最大幺幂边界点。为此，先比较绕除子 $z_i=0$ 与绕除子 $w_i=0$ 的单值性。分别用 $T_i^z,N_i^z$ 表示绕除子 $z_i=0$ 的单值性及其对数，类似地用 $T_i^w,N_i^w$ 表示绕除子 $w_i=0$ 的单值性及其对数。沿一条绕 $w_i=0$ 一周的路径，例如
\[
\theta\longmapsto(p_1,\ldots,p_{i-1},p_i e^{2\pi i\theta},p_{i+1},\ldots,p_r),
\]
代入式~\eqref{eq:mum-monomial-coordinate-change} 后，可以看到：与 $\phi$ 复合所得路径绕 $z_j=0$ 共 $a_{ij}$ 周。因此，对每个 $i$，得到
\[
T_i^w=\prod_j(T_j^z)^{a_{ij}}.
\]
现在假定 $T_j^z$ 为幺幂算子。由于它们彼此交换，$T_i^w$ 也为幺幂；反向结论则由 $(a_{ij})$ 的可逆性推出。于是可以取对数，得到
\begin{equation}\label{eq:mum-log-coordinate-change}
N_i^w=\sum_j a_{ij}N_j^z.
\end{equation}

\SourcePage{138}
为便于后面使用，注意该式可以写成
\begin{equation}\label{eq:mum-log-matrix-change}
\begin{pmatrix}N_1^w\\\vdots\\N_r^w\end{pmatrix}
=A\begin{pmatrix}N_1^z\\\vdots\\N_r^z\end{pmatrix},\qquad A=(a_{ij}).
\end{equation}
接着比较单值性权滤过。在源上，它由 $N=\sum_i a_iN_i^w$ 确定，其中 $a_i>0$。不过，由式~\eqref{eq:mum-log-coordinate-change}，有
\[
N=\sum_i a_iN_i^w=\sum_j\left(\sum_i a_i a_{ij}\right)N_j^z.
\]
已假定对所有 $i,j$ 都有 $a_{ij}\geq0$；由于 $(a_{ij})$ 可逆，$\sum_i a_i a_{ij}>0$。这说明在目标上也可使用 $N$，所以源和目标具有相同的单值性权滤过。由此很容易看出，在 $w=0$ 处满足定义~\ref{def:maximal-unipotent-point}，当且仅当在 $z=0$ 处满足它，断言得证。

在这种情况下，用 $q_i^z,q_i^w$ 表示不同局部紧化中由式~\eqref{eq:mum-canonical-coordinates} 确定的坐标。计算得
\begin{equation}\label{eq:mum-q-monomial-change}
q_j^z=\prod_i(q_i^w)^{a_{ij}}.
\end{equation}
注意它与式~\eqref{eq:mum-monomial-coordinate-change} 的相似性。

所以，一个最大幺幂边界点通过式~\eqref{eq:mum-monomial-coordinate-change} 可以产生\emph{许多}其他最大幺幂边界点，进而产生许多组局部坐标 $(q_1,\ldots,q_r)$；它们都通过式~\eqref{eq:mum-q-monomial-change} 联系起来。在\hyperref[subsec:large-radius-limits]{第 6.2.1 小节}将看到，Kähler 一侧的大半径极限点的局部坐标有\emph{完全相同}的不确定性。

回到整体情形，设 $\overline{\mathcal M}$ 是复结构模空间 $\mathcal M$ 的一个光滑正规交叉紧化。若 $\overline{\mathcal M}$ 有一个最大幺幂边界点，上述讨论就表明：通过爆破，可以得到具有多个最大幺幂边界点的其他紧化。$\overline{\mathcal M}$ 本身也可能是一个爆破所得空间。事实上，$\overline{\mathcal M}$ 可能有多个最大幺幂边界点，在爆缩到一个奇异紧化时，这些点会被识别为同一点。

整理这些可能性的一种办法，是在最大幺幂边界点之间引入等价关系。依照 \SourceCite{Morrison6}，作如下定义。
\begin{definition}\label{def:equivalent-mum-points}
设 $\overline{\mathcal M}$ 是模空间 $\mathcal M$ 的一个光滑正规交叉紧化。若存在包含 $p,q$ 的连通集 $\Xi\subset\overline{\mathcal M}-\mathcal M$，并且在 $\Xi$ 的一个邻域上存在由 $\overline{\mathcal M}$ 的余切丛中、沿边界有对数极点的形式组成的局部系统 $\mathcal L$，使得在 $p,q$ 各自附近，该局部系统由 $1$ 形式 $d\log q_k$ 张成，其中 $q_k$ 如式~\eqref{eq:mum-canonical-coordinates} 所述，则称两个最大幺幂边界点 $p,q\in\overline{\mathcal M}$ \emph{等价}。
\end{definition}
\index{最大幺幂边界点!等价关系|hyperpage}
很容易看出，爆破产生的多个最大幺幂边界点，在这一意义下彼此等价。

\IndexLocation{indexloc:0830:138}镜像对称的一个基本思想是：Calabi--Yau 流形 $V$ 的复结构模空间中，最大幺幂边界点的每个等价类，应当对应于其镜像 $V^\circ$ 的 Kähler 模空间中一个适当定义的
\SourcePage{139}
大半径极限点等价类（将在\hyperref[subsec:large-radius-limits]{第 6.2.1 小节}讨论）。

实际处理中，从一个选定的 $\overline{\mathcal M}$ 出发，通常会尽可能进行爆缩，直到得到一个无法继续爆缩的光滑正规交叉紧化。这样可以使等价类尽可能小。\hyperref[subsec:two-parameter-moduli]{第 6.1.4 小节}将给出一个例子。

还要指出，定义~\ref{def:equivalent-mum-points} 的等价关系，也可在 $\mathcal M$ 的更一般紧化上定义。\SourceCite{Morrison3} 已经做了这项工作。

最后，注意我们还没有讨论最大幺幂边界点的\emph{存在性}。这仍然是猜想，尽管迄今计算的每个例子都有这样的点。很快还会看到，环面超曲面的模空间具有自然出现的紧化，其中有一些特殊边界点；猜想这些点是最大幺幂的（\SourceCite{HLY2} 给出了一个很好的部分结果）。一般而言，镜像对称预言 $\mathcal M$ 存在十分特殊的紧化，它们不仅具有最大幺幂边界点，而且反映镜像的 Kähler 模结构。本章后面还会进一步讨论这一点。

\subsection{环面超曲面的复结构模}\label{subsec:toric-complex-moduli}\index{多项式模空间|hyperpage}
研究 Calabi--Yau 环面超曲面的复结构模。回顾\hyperref[sec:batyrev-mirror]{第 4.1 节}，对自反多胞形 $\Delta$，反典范超曲面 $\overline V\subset\mathbb P_\Delta$ 是 Calabi--Yau，尽管可能有奇点。若 $\Sigma$ 是 $\Delta$ 的法扇的一个极大射影细分，已经证明 $\overline V$ 的严格变换仍为反典范超曲面，从而得到一族极小 Calabi--Yau 超曲面 $V\subset X=X_\Sigma$。所以 $V$ 是一个至多具有终端奇点的 Gorenstein 轨形；当 $\dim V=3$ 时，它是光滑的。

不处理 $V$ 的整个复结构模空间 $\mathcal M$，而集中研究其中的“多项式”部分 $\mathcal M_{\mathrm{poly}}\subset\mathcal M$，即那些可实现为 $X$ 中超曲面的复结构。后面将看到，若 $\Delta$ 的任何余维二的面都没有内部格点，那么 $\mathcal M_{\mathrm{poly}}$ 就是整个复结构模空间。

为构造多项式模空间，从形如 $\sum_{m\in\Delta\cap M}\lambda_m t^m$ 的 Laurent 多项式构成的向量空间 $L(\Delta\cap M)$ 出发。与\hyperref[subsec:toric-polytopes]{第 3.2.1 小节}一样，$t^m$ 是对应于 $m$ 的特征标。（\hyperref[chap:toric]{第~\ref*{chap:toric} 章}中，将 $L(\Delta\cap M)$ 记作 $L(\Delta)$；这里希望更明确地体现其对格的依赖。）由式~\eqref{eq:toric-sections}，有自然同构
\[
H^0(X,\mathcal O_X(-K_X))\simeq L(\Delta\cap M).
\]
可以把 $L(\Delta\cap M)$ 中的一个多项式看作定义 $X$ 中一个反典范超曲面的仿射部分。此外，$-K_X=\sum_\rho D_\rho$ 表明
\[
\Delta\cap M=\{m\in M:\langle m,v_\rho\rangle\geq-1\text{，对所有 }\rho\}.
\]
由于典范微分的拉回仍是典范微分，\IndexLocation{indexloc:1355:139}自同构群 $\operatorname{Aut}(X)$ 自然作用于射影空间 $\mathbb P(L(\Delta\cap M))$，所以多项式模空间是商
\[
\mathcal M_{\mathrm{poly}}=\mathbb P(L(\Delta\cap M))/\operatorname{Aut}(X).
\]
要使这一表述精确，需将 $\mathbb P(L(\Delta\cap M))$ 换成对应于拟光滑超曲面的子集，并考虑商的存在性。后一问题很重要，因为 $\operatorname{Aut}(X)$ 未必为约化群。不过，如 \SourceCite{BC} 第 13 节所指出的，可以找到一个非空不变开集
\SourcePage{140}
$U\subset\mathbb P(L(\Delta\cap M))$，使几何商 $\mathcal M_{\mathrm{poly}}$ 存在。因此，至少在一般位置上可以构造多项式模空间。

若能更清楚地了解这一模空间，当然更好。因此，接着讨论 \SourceCite{AGM1} 引入的\emph{简化多项式模空间} $\mathcal M_{\mathrm{simp}}$。令 $(\Delta\cap M)_0$ 为 $\Delta\cap M$ 中不属于 $\Delta$ 任何刻面内部的点所组成的集合，于是 $L((\Delta\cap M)_0)$ 是形如
\[
\sum_{m\in(\Delta\cap M)_0}\lambda_m t^m
\]
的 Laurent 多项式构成的向量空间。由于环面 $T=T_N$ 通过 $v\cdot t^m=v^m t^m$（$v\in T$）作用于 $t^m$，便在 $L((\Delta\cap M)_0)$ 上诱导出作用
\begin{equation}\label{eq:toric-moduli-coefficient-action}
v\cdot\lambda_m=v^m\lambda_m,\qquad v\in T.
\end{equation}
因此，$T$ 作用于 $\mathbb P(L((\Delta\cap M)_0))$；可以验证，在包含映射 $T\subset\operatorname{Aut}(X)$ 下，这与 $\operatorname{Aut}(X)$ 在 $\mathbb P(L(\Delta\cap M))$ 上的作用相容。

现在将简化多项式模空间 $\mathcal M_{\mathrm{simp}}$ 定义为
\[
\mathcal M_{\mathrm{simp}}=\mathbb P(L((\Delta\cap M)_0))/T.
\]
与前面一样，实际上须限制到 $\mathbb P(L((\Delta\cap M)_0))$ 中对应于拟光滑超曲面的子集，并考虑商的存在性。幸运的是，后一问题不太棘手，因为 $T$ 是约化群，而式~\eqref{eq:toric-moduli-coefficient-action} 给出了作用的线性化。因此，不仅能够构造 $\mathcal M_{\mathrm{simp}}$，还可用 GIT 商将它紧化 \SourceCite{Mumford1}：
\[
\overline{\mathcal M}_{\mathrm{simp}}=\mathbb P(L((\Delta\cap M)_0))/\!/T.
\]
不过，对我们的目的而言，最有趣的紧化来自 \SourceCite{KSZ1} 的 \emph{Chow 商}。后面将看到，这种紧化具有由次级扇给出的自然环面结构。

基本断言是：显然定义的映射 $\phi:\mathcal M_{\mathrm{simp}}\to\mathcal M_{\mathrm{poly}}$ 至少在一般位置上是局部同构。这就是 \SourceCite{AGM1} 中的\emph{占优猜想}，现在加以证明。
\begin{proposition}\label{prop:simplified-moduli-etale}\index{简化多项式模空间|hyperpage}
映射 $\phi:\mathcal M_{\mathrm{simp}}\to\mathcal M_{\mathrm{poly}}$ 在一般位置上为平展映射。
\end{proposition}
\begin{proof}
用 $\operatorname{Aut}_0(X)$ 表示 $\operatorname{Aut}(X)$ 的单位元所在连通分支。如 \SourceCite{Cox} 所指出的，$\operatorname{Aut}_0(X)$ 在 $\operatorname{Aut}(X)$ 中具有有限指数，因此只需证明自然映射
\[
\mathbb P(L((\Delta\cap M)_0))/T\longrightarrow
\mathbb P(L(\Delta\cap M))/\operatorname{Aut}_0(X)
\]
为占优映射。

令 $I=(\Delta\cap M)-(\Delta\cap M)_0$，即 $\Delta$ 的刻面内部格点所组成的集合。由命题~\ref{prop:automorphism-facets} 的证明，每个 $m\in I$ 确定一个自同构的单参数族 $y_{m,u}:X\to X$，其中参数 $u\in\mathbb C$。为描述 $y_{m,u}$，使用\hyperref[subsec:homogeneous-ring]{第 3.2.3 小节}的齐次坐标环 $S=\mathbb C[x_\rho]$。由于 $m$ 位于某个刻面的内部，存在 $\rho$，使 $\langle m,v_\rho\rangle=-1$，而对所有 $\rho'\ne\rho$ 都有 $\langle m,v_{\rho'}\rangle>-1$。置
$\mathbf x^D=\prod_{\rho'\ne\rho}x_{\rho'}^{\langle m,v_{\rho'}\rangle}$，相应的自同构单参数族为
\[
y_{m,u}(x_\rho,\mathbf x)=(x_\rho+u\mathbf x^D,\mathbf x).
\]
\SourcePage{141}
这里与式~\eqref{eq:toric-root-automorphism} 一样，将 $\mathbb C^{\Sigma(1)}$ 的点写成 $(x_\rho,\mathbf x)$，其中 $\mathbf x$ 是以 $\Sigma(1)-\{\rho\}$ 为指标的向量。

考虑具有 $|I|$ 个参数的自同构族
\begin{equation}\label{eq:toric-root-product}
y_{\mathbf u}=\prod_{m_i\in I}y_{m_i,u_i}.
\end{equation}
为定义式~\eqref{eq:toric-root-product} 的右端，需要给 $I$ 选定一个次序，因为 $y_{m_i,u_i}$ 未必彼此交换。将包含映射 $T\subset\operatorname{Aut}_0(X)$ 与映射 $\mathbf u\mapsto y_{\mathbf u}$ 结合，得到映射
\[
T\times\mathbb C^I\longrightarrow\operatorname{Aut}_0(X),
\]
很容易看出它为占优映射。因此，只需考虑 $T\times\mathbb C^I$ 在 $\mathbb P(L(\Delta\cap M))$ 上的作用；只要证明在一般位置上，$\mathbb C^I$ 的轨道维数为 $|I|$，且与 $\mathbb P(L((\Delta\cap M)_0))$ 横截相交，命题就成立。

先计算 $y_{m,u}$ 在 Laurent 单项式 $t^\alpha\in L(\Delta\cap M)$ 上的作用。为此，需将 $t^\alpha$ 映到单项式 $\prod_\rho x_\rho^{\langle\alpha,v_\rho\rangle+1}$（见\hyperref[subsec:homogeneous-ring]{第 3.2.3 小节}），施以 $y_{m,u}$，再映回 $L(\Delta\cap M)$。若 $\langle m,v_\rho\rangle=-1$ 确定 $m$ 所在的刻面，则计算得
\[
y_{m,u}(t^\alpha)=
\begin{cases}
t^\alpha,&\langle\alpha,v_\rho\rangle=-1,\\
t^\alpha+u(\langle\alpha,v_\rho\rangle+1)t^{\alpha+m}
+\text{高阶项},&\langle\alpha,v_\rho\rangle>-1.
\end{cases}
\]
可以验证，$\langle\alpha,v_\rho\rangle>-1$ 蕴含 $\alpha+m\in\Delta\cap M$。\hyperref[ex:degree-four-moduli]{例题~\ref*{ex:degree-four-moduli}} 将给出一次具体计算，包括上面省略的高阶项。

由于 $\Delta$ 自反，原点 $0$ 属于 $(\Delta\cap M)_0$，所以 $1=t^0\in L((\Delta\cap M)_0)$。于是很容易得到
\[
y_{\mathbf u}(1)=1+\sum_{m_i\in I}u_i t^{m_i}+\text{高阶项}.
\]
这说明，$\mathbb C^I$ 作用于 $1$ 所得轨道的维数正确；而且它在 $1$ 处的切空间显然与 $L((\Delta\cap M)_0)$ 互补。由于 $y_{\mathbf u}(1)$ 的常数项恒为 $1$，它恒不为零；射影化后上述一阶结论仍然成立，命题随即得证。
\end{proof}

这一命题表明，在一般位置上，$\mathcal M_{\mathrm{simp}}$ 是 $\mathcal M_{\mathrm{poly}}$ 的有限覆盖。了解 $X$ 的自同构群后，就能显式计算作用于 $\mathcal M_{\mathrm{simp}}$ 的有限群。\SourceCite{CFKM} 已在一个例子中完成这一计算；后面的\hyperref[ex:degree-four-moduli]{例题~\ref*{ex:degree-four-moduli}} 将在稍简单的情形中说明这种技巧。

接着计算 $\mathcal M_{\mathrm{poly}}$ 与整个复结构模空间 $\mathcal M$ 的维数，并确定两者何时相等。
\begin{proposition}\label{prop:toric-moduli-dimensions}
设 $V\subset X=X_\Sigma$ 为极小 Calabi--Yau 环面超曲面，其中 $\Sigma$ 是 $n$ 维自反多胞形 $\Delta$ 的法扇的一个极大射影细分。则：
\begin{enumerate}[label=(\roman*)]
\item $V$ 的多项式模空间 $\mathcal M_{\mathrm{poly}}$ 的维数为
\[
\dim\mathcal M_{\mathrm{poly}}=l(\Delta)-n-1-\sum_\Gamma l^*(\Gamma),
\]
其中 $l(\Delta)=|\Delta\cap M|$，求和遍及 $\Delta$ 的所有刻面 $\Gamma$，$l^*(\Gamma)$ 为 $\Gamma$ 相对内部的格点数。
\SourcePage{142}
\item $V$ 的复结构模空间 $\mathcal M$ 的维数为
\[
\dim\mathcal M=l(\Delta)-n-1-\sum_\Gamma l^*(\Gamma)
+\sum_\Theta l^*(\Theta)l^*(\widehat\Theta^\circ),
\]
其中第二个求和遍及 $\Delta$ 的所有余维二的面，$\widehat\Theta^\circ$ 为 $\Delta^\circ$ 中与 $\Theta$ 对偶的一维面。
\item 若 $\Delta\cap M$ 中没有余维二的面的内部格点，则 $\mathcal M=\mathcal M_{\mathrm{poly}}$。
\end{enumerate}
\end{proposition}
\begin{proof}
$\mathcal M_{\mathrm{simp}}$ 的维数很容易计算：由于 $\Delta$ 的维数为 $n$，$T$ 在 $\mathbb P(L((\Delta\cap M)_0))$ 上的作用，在一般位置上具有有限稳定子。因此
\[
\dim\mathcal M_{\mathrm{simp}}=|(\Delta\cap M)_0|-n-1,
\]
这也就是 $\mathcal M_{\mathrm{poly}}$ 的维数。由 $l(\Delta)=|(\Delta\cap M)_0|+\sum_\Gamma l^*(\Gamma)$，得到命题的第一部分。

对于第二部分，注意 $V$ 为极小 Calabi--Yau 轨形，所以由命题~\ref{prop:minimal-cy-orbifold-deformations}，它的无穷小形变由 $H^1(V,\Theta_V)\simeq H^1(V,\widehat\Omega_V^{n-2})$ 分类（记住 $V$ 的维数为 $n-1$）。此外，由命题~\ref{prop:terminal-gorenstein-orbifold-codimension}，$V$ 的奇点集余维至少为 $4$，所以 Ran 的结果 \SourceCite{Ran2} 表明这些形变无障碍。于是 $\mathcal M$ 的维数为 $h^{n-2,1}(V)$。式~\eqref{eq:batyrev-hn21} 计算了这个 Hodge 数，从而证明了第二部分。最后一部分直接由前两部分推出，命题得证。
\end{proof}

这一命题与\hyperref[chap:mirror-constructions]{第 4 章}的定理~\ref{thm:batyrev-hodge} 有很好的联系。在一般点 $V$ 处，$\mathcal M$ 的切空间为 $H^1(V,\Theta_V)\simeq H^{n-2,1}(V)$；\hyperref[sec:batyrev-mirror]{第 4.1 节}中记为 $H_{\mathrm{poly}}^{n-2,1}(V)\subset H^{n-2,1}(V)$ 的空间，就是 $\mathcal M_{\mathrm{poly}}$ 在 $V$ 处的切空间。因此，命题~\ref{prop:toric-moduli-dimensions} 的第一部分给出了式~\eqref{eq:polynomial-hodge-number} 的严格证明。

现在对 $\mathcal M_{\mathrm{poly}}$ 有了更多了解，下一项任务就是将它紧化，并考察能否找到一些\IndexLocation{indexloc:0766:142}最大幺幂边界点。实际上，将使用 $\mathcal M_{\mathrm{simp}}$，因为它有一个特别好的紧化。如前所述，采用 \SourceCite{KSZ1} 的 Chow 商，并取其正规化：
\begin{equation}\label{eq:moduli-chow-quotient}\index{Chow 商!正规化|hyperpage}
\overline{\mathcal M}_{\mathrm{simp}}=\mathbb P(L((\Delta\cap M)_0))/\!/_{\mathrm{Ch}}T.
\end{equation}
它是一个环面簇，先描述其中的环面。注意，$L((\Delta\cap M)_0)$ 包含环面 $(\mathbb C^*)^{(\Delta\cap M)_0}$。位似给出自然映射 $\mathbb C^*\to(\mathbb C^*)^{(\Delta\cap M)_0}$（射影化时需要它）；而式~\eqref{eq:toric-moduli-coefficient-action} 表明，$T$ 的作用来自由特征标 $t^m$（$m\in(\Delta\cap M)_0$）给出的映射 $T\to(\mathbb C^*)^{(\Delta\cap M)_0}$。所以有环面
\begin{equation}\label{eq:moduli-quotient-torus}
T_0=(\mathbb C^*)^{(\Delta\cap M)_0}/\operatorname{im}(T\times\mathbb C^*),
\end{equation}
它包含在上述商中。利用特殊点 $0\in(\Delta\cap M)_0$，可以去掉 $\mathbb C^*$。具体地，令 $\Xi=(\Delta\cap M)_0-\{0\}$，则
\[
T_0=(\mathbb C^*)^\Xi/\operatorname{im}(T).
\]
$T_0$ 的单参数子群构成格
\begin{equation}\label{eq:moduli-quotient-lattice}
N_0=(\mathbb Z^\Xi/N)/\text{挠元},
\end{equation}
其中当 $\Xi=\{m_1,\ldots,m_s\}$ 时，包含映射 $N\to\mathbb Z^\Xi$ 为 $u\mapsto(\langle u,m_1\rangle,\ldots,\langle u,m_s\rangle)$。

\SourcePage{143}
格 $N_0$ 很重要，因为\hyperref[sec:gkz-decomposition]{第 3.4 节}定义的 $\Xi=(\Delta\cap M)_0-\{0\}$ 的 GKZ 分解，自然位于向量空间 $N_0\otimes\mathbb R$ 中。此外，由定理~\ref{thm:gkz-fan}，该扇的极大锥是所有 $M_{\mathbb R}$ 中射影、单纯且满足 $\Sigma^\circ(1)\subset\Xi$ 的扇 $\Sigma^\circ$ 对应的锥 $\operatorname{cpl}(\Sigma^\circ)$。

这与 $V$ 的 Batyrev 镜像密切相关。回顾这一构造：若 $\Delta^\circ\subset N_{\mathbb R}$ 为极多胞形，$\Sigma^\circ$ 为 $\Delta^\circ$ 的法扇的一个极大射影细分，那么 Batyrev 镜像就是反典范超曲面族 $V^\circ\subset X_{\Sigma^\circ}$。按定义~\ref{def:projective-subdivision}，$\Sigma^\circ$ 射影、单纯，细分 $\Delta^\circ$ 的法扇，并满足 $\Sigma^\circ(1)=\Delta\cap M-\{0\}$。不过，在\hyperref[subsec:enlarged-kahler-moduli]{第 6.2.3 小节}将看到，可以将最后一个条件换成
$\Sigma^\circ(1)=(\Delta\cap M)_0-\{0\}=\Xi$，而不改变 $V^\circ$。这样，$V^\circ\subset X_{\Sigma^\circ}$ 就给出 $\Xi$ 的 GKZ 分解中的一个锥 $\operatorname{cpl}(\Sigma^\circ)$。

因此，每一种 Batyrev 镜像 $V^\circ\subset X_{\Sigma^\circ}$ 的选择，都在 GKZ 分解中给出一个特殊极大锥。在对应于 GKZ 分解的环面簇中，这些锥给出环面边界上的特殊点。它们可能是奇异点；不过，很快将看到，猜想其消解能够给出所需的最大幺幂边界点。此外，作为\hyperref[subsec:enlarged-kahler-moduli]{第 6.2.3 小节}的预告，还应指出：$\operatorname{cpl}(\Sigma^\circ)$ 的内部就是 $X_{\Sigma^\circ}$ 的 Kähler 锥，并且自然同构于 $V^\circ$ 的 Kähler 锥的环面部分。这是 $V$ 的复结构模与其镜像 $V^\circ$ 的 Kähler 模之间同构关系的开端。

现在可以仔细描述 Chow 商~\eqref{eq:moduli-chow-quotient}。定理~\ref{thm:gkz-fan} 中 $\Xi$ 的 GKZ 分解并不完备，但如\hyperref[sec:gkz-decomposition]{第 3.4 节}所述，可以自然扩充为一个完备扇，称为 $\Xi$ 的\emph{次级扇}。为便于后面使用，注意次级扇由以下点集构造：
\begin{equation}\label{eq:moduli-secondary-configuration}\index{次级扇|hyperpage}
\Xi^+=(\Xi\cup\{0\})\times\{1\}=(\Delta\cap M)_0\times\{1\}\subset M\times\mathbb Z.
\end{equation}
\SourceCite{KSZ2} 表明，正规化的 Chow 商 $\overline{\mathcal M}_{\mathrm{simp}}$ 就是由 $\Xi$ 的次级扇确定的环面簇。\SourceCite{KSZ1} 探究了 Chow 商与各种 GIT 商之间的关系。还应指出，$\overline{\mathcal M}_{\mathrm{simp}}$ 有一个由 \SourceCite{GKZ2} 的\emph{\IndexLocation{indexloc:1242:143}次级多胞形}给出的自然射影嵌入。

环面紧化 $\overline{\mathcal M}_{\mathrm{simp}}$ 包含一些特殊的环面边界点，对应于镜像族的不同选择。不过，$\overline{\mathcal M}_{\mathrm{simp}}$ 不一定光滑，因此讨论最大幺幂边界点以前，需要先消解奇点。于是选取次级扇的一个细分，使它给出光滑环面簇。所得 $\mathcal M_{\mathrm{simp}}$ 的光滑紧化记为 $\widetilde{\mathcal M}_{\mathrm{simp}}$。

这一光滑紧化具有如下性质：每个 Batyrev 镜像 $V^\circ\subset X_{\Sigma^\circ}$ 都在 $\widetilde{\mathcal M}_{\mathrm{simp}}$ 中给出有限多个特殊边界点，对应于细分 $\operatorname{cpl}(\Sigma^\circ)$ 的锥。猜想这些正是所寻求的最大幺幂边界点。
\begin{conjecture}\label{conj:toric-mum-boundary}
设 $\widetilde{\mathcal M}_{\mathrm{simp}}$ 是 Chow 商 $\overline{\mathcal M}_{\mathrm{simp}}$ 的一个光滑环面消解，$V^\circ\subset X_{\Sigma^\circ}$ 是 $V$ 的一个 Batyrev 镜像。则 $\widetilde{\mathcal M}_{\mathrm{simp}}$ 中对应于细分 $\operatorname{cpl}(\Sigma^\circ)$ 的锥的点，全部都是最大幺幂边界点。此外，这些点在定义~\ref{def:equivalent-mum-points} 的等价关系下全部等价。
\end{conjecture}
\SourceCite{BK1} 曾给出这一猜想的一个不完整证明。若 $H^{1,1}(V)$ 在杯积下生成 $\bigoplus_pH^{p,p}(V)$，那么只要证明
\SourcePage{144}
式~\eqref{eq:toric-complete-intersection-i-function} 定义的形式函数 $I_V$ 的系数确实是镜像族 $V^\circ$ 的周期，正如 \SourceCite{Givental4} 所断言的（见\hyperref[subsec:nef-gkz]{第 5.5.3 小节}的讨论），就很容易推出这一猜想。特别是，这适用于 $V$ 为三维簇的情形。可以证明，猜想~\ref{conj:toric-mum-boundary} 对后面的\hyperref[ex:weighted-complex-moduli]{例题~\ref*{ex:weighted-complex-moduli}} 和\hyperref[ex:degree-four-moduli]{例题~\ref*{ex:degree-four-moduli}} 成立。例如，前一个例题的证明可由论文 \SourceCite{CdFKM} 中的计算推出。

猜想~\ref{conj:toric-mum-boundary} 的动机之一是：在镜像对称下，绕某个边界除子的单值性的对数，应与同细分 $\operatorname{cpl}(\Sigma^\circ)$ 的锥的相应生成元的负元作杯积识别起来；最大幺幂边界点的定义~\ref{def:maximal-unipotent-point} 后面的讨论已经提到这一点。\hyperref[chap:quantum-cohomology]{第~\ref*{chap:quantum-cohomology} 章}讨论 A 模型 Hodge 结构变分时，将从更深入的角度理解它。

\emph{本章余下部分都假定猜想~\ref{conj:toric-mum-boundary} 成立。}

猜想~\ref{conj:toric-mum-boundary} 描述的最大幺幂边界点来自次级扇中一些十分特殊的锥。GKZ 分解中还有其他锥：它们对应的 $\Sigma^\circ$ 不使用 $\Xi$ 的全部点，或者不细分 $\Delta^\circ$ 的法扇；将 GKZ 分解扩充为次级扇时，还会得到一些锥。猜想~\ref{conj:toric-mum-boundary} 没有说明 $\overline{\mathcal M}_{\mathrm{simp}}$ 中对应于这些锥的边界点的性质。不过，在\hyperref[subsec:enlarged-kahler-moduli]{第 6.2.3 小节}将看到，所有这些锥都有有趣的物理解释——每个锥都给出一种不同的物理理论。

还应提到，\emph{Gröbner 扇}给出了次级扇的一个典范细分 \SourceCite{MR,Sturmfels1}。它给出 $\mathcal M_{\mathrm{simp}}$ 的一个紧化，是 $\overline{\mathcal M}_{\mathrm{simp}}$ 的爆破，因而奇点可能更少。\SourceCite{HLY1,HLY2} 使用了这一紧化，并证明存在\emph{最大退化}边界点，即在局部存在唯一的全纯周期积分，常数倍除外。这是最大幺幂性的一个必要条件，因此应由猜想~\ref{conj:toric-mum-boundary} 推出。

目前的讨论着重于爆破 $\overline{\mathcal M}_{\mathrm{simp}}$，似乎 $\overline{\mathcal M}_{\mathrm{simp}}$ 是 $\mathcal M_{\mathrm{simp}}$ 的“极小”紧化。不过，在某些情况下，将次级扇中的若干锥合并成一个较大的锥是完全自然的。这会给出一个紧化 $\overline{\mathcal M}'_{\mathrm{simp}}$，使 $\overline{\mathcal M}_{\mathrm{simp}}$ 是 $\overline{\mathcal M}'_{\mathrm{simp}}$ 的爆破。\hyperref[ex:degree-four-moduli]{例题~\ref*{ex:degree-four-moduli}} 将展示这种现象；\hyperref[subsec:enlarged-kahler-moduli]{第 6.2.3 小节}将仔细研究应如何以及何时这样做。

还要注意，除了存在多个镜像，原来的族 $V\subset X_\Sigma$ 也有多个版本，因为极大射影细分 $\Sigma$ 可能有许多种选择。我们对 $\mathcal M_{\mathrm{simp}}$ 的描述没有提到 $\Sigma$，原因很明确：这些族具有相同的多项式模。变化的是 $V$ 的 Kähler 锥，因为每个 $\Sigma$ 都对应于\hyperref[subsec:enlarged-kahler-moduli]{第 6.2.3 小节}所讨论的 $V$ 的扩大 Kähler 模空间的一部分。\hyperref[chap:mirror-constructions]{第~\ref*{chap:mirror-constructions} 章}中，$\Sigma$ 和 $\Sigma^\circ$ 的多种选择似乎是个麻烦，因为它们使我们无法典范地选择镜像。不过，现在逐渐了解到：“多个镜像”的现象，对从整体上理解复结构模与 Kähler 模至关重要。

紧化 $\overline{\mathcal M}_{\mathrm{simp}}$ 的环面边界包含刚才讨论的最大幺幂边界点。不过，$\overline{\mathcal M}_{\mathrm{simp}}-\mathcal M_{\mathrm{simp}}$ 中也有不在环面边界上的点，因为环面 $T_0\subset\overline{\mathcal M}_{\mathrm{simp}}$ 中有许多点对应于奇异超曲面 $V\subset X_\Sigma$。实际上，由于 $V$ 和 $X_\Sigma$ 可能
\SourcePage{145}
有奇点，通常用\emph{非退化超曲面}来表述更方便。对 $V\subset X_\Sigma$，这意味着：对每个环面轨道 $\mathcal O\subset X_\Sigma$，交集 $V\cap\mathcal O$ 或者为空，或者光滑且余维为一。由于 $X_\Sigma$ 是轨形，很容易推出 $X_\Sigma$ 的每个非退化超曲面都是子轨形。

在 $L((\Delta\cap M)_0)$ 中，\emph{判别式轨迹} $\mathcal D$ 是所有使 $f=0$ 不是非退化超曲面的 Laurent 多项式 $f$ 构成的集合。由 \SourceCite[第 10 章]{GKZ2}，$\mathcal D$ 的定义方程由\emph{主 $\mathcal A$ 行列式}给出。置 $\mathcal A=(\Delta\cap M)_0=\Xi\cup\{0\}$，则它是 $f$ 的系数的多项式，定义为
\[
E_{\mathcal A}(f)=\operatorname{Res}_{\mathcal A}
\left(t_1\frac{\partial f}{\partial t_1},\ldots,t_n\frac{\partial f}{\partial t_n},f\right),
\]
其中 $\operatorname{Res}_{\mathcal A}$ 为 \SourceCite[第 8 章]{GKZ2} 定义的稀疏结式，$t_1,\ldots,t_n$ 为环面 $T$ 的坐标。这个多项式实际上可约，并具有因式分解
\[
E_{\mathcal A}(f)=\prod_{\Gamma\subset\Delta}
D_{\mathcal A\cap\Gamma}(f_\Gamma)^{m(\Gamma)}.
\]
右端可如下理解。$\mathcal A$ 判别式 $D_{\mathcal A}$ 定义了 $L((\Delta\cap M)_0)$ 中那些在环面 $T$ 上具有奇异点的多项式所组成集合的 Zariski 闭包（因此，若该集合的余维大于 $1$，则 $D_{\mathcal A}=1$）。对每个面 $\Gamma\subset\Delta$，再将 $f_\Gamma$ 定义为 $f$ 中对应于 $\Gamma\cap\mathcal A$ 的各项之和。最后，上式的指数 $m(\Gamma)$ 是正整数，在良好情形下等于 $1$。

由这一描述可知，判别式轨迹可能有许多连通分支。幸运的是，许多 $D_{\mathcal A\cap\Gamma}$ 常常等于 $1$。例如，后面的\hyperref[ex:weighted-complex-moduli]{例题~\ref*{ex:weighted-complex-moduli}} 中，判别式轨迹恰好有两个分支。还应指出，$\mathcal D$ 恰好也是下面将描述的 GKZ 超几何方程组的判别式轨迹。

从 $L((\Delta\cap M)_0)$ 下降到 Chow 商时，判别式轨迹的像仍记为 $\mathcal D$。判别式的\emph{主分支}是由 $D_{\mathcal A}=0$ 定义的超曲面。当 $V$ 为三维簇时，该分支的一般点参数化具有一个结点的超曲面 $V$（三维结点在物理文献中常称为\emph{锥形奇点}）。判别式的其他非空分支，则参数化带有可能更复杂的奇点的超曲面。

除刚才描述的判别式轨迹外，$\overline{\mathcal M}_{\mathrm{simp}}$ 的环面边界也参数化奇异超曲面，尽管参数值 $\lambda_i$ 可能没有良好定义。从 GIT 商的角度看，问题在于这些参数值可能对应于非稳定超曲面。因此，作为 $\overline{\mathcal M}_{\mathrm{simp}}$ 中的对象，超曲面本身可能没有良好定义。例如，五次镜像族~\eqref{eq:quintic-pencil} 在 $\psi\to\infty$ 时的极限，可以取为 $x_1x_2x_3x_4x_5=0$；然而，进行坐标变换 $x_5\mapsto x_5/\psi$ 后，极限却为
\[
x_1^5+x_2^5+x_3^5+x_4^5+x_1x_2x_3x_4x_5=0.
\]

\subsection{超几何方程与模空间}\label{subsec:hypergeometric-moduli}\index{GKZ方程组@GKZ 方程组!模空间|hyperpage}
\hyperref[sec:hypergeometric]{第 5.5 节}定义了 GKZ 超几何方程，并且以相当朴素的方式讨论了它们在 $V$ 的模空间上的行为。现在，对复结构模有了更多了解，可以更清楚地认识 GKZ 方程组。

\hyperref[sec:hypergeometric]{第 5.5 节}构造超几何方程时使用了 $\Delta\cap M$。这里，由于要在模空间 $\overline{\mathcal M}_{\mathrm{simp}}$ 上工作，将使用 $(\Delta\cap M)_0$
\SourcePage{146}
来代替它。按照式~\eqref{eq:toric-gkz-configuration}，这意味着使用从式~\eqref{eq:moduli-secondary-configuration} 中的点集
\[
(\Delta\cap M)_0\times\{1\}=\Xi^+
\]
构造的 GKZ 方程组。与\hyperref[chap:hodge-yukawa]{第~\ref*{chap:hodge-yukawa} 章}一样，将 $(\Delta\cap M)_0$ 的点记为 $m_i$，其中 $m_0=0$。于是 $\Xi^+$ 中各点的关系模为
\begin{equation}\label{eq:moduli-relation-lattice}
M_0=\left\{\ell=(\ell_i)\in\mathbb Z^{\Xi^+}:\sum_i\ell_i m_i=0,\ \sum_i\ell_i=0\right\}.
\end{equation}
由式~\eqref{eq:moduli-quotient-torus}，知道 $M_0$ 可识别为环面 $T_0$ 的特征标群，所以式~\eqref{eq:moduli-quotient-lattice} 中的 $N_0$ 是 $M_0$ 的对偶。GKZ 超几何方程组由算子 $\Box_\ell$（$\ell\in M_0$）及 $Z_j$（$j=1,\ldots,n+1$）组成。$\Box_\ell,Z_j$ 的精确定义见\hyperref[sec:hypergeometric]{第 5.5 节}。

\hyperref[chap:hodge-yukawa]{第~\ref*{chap:hodge-yukawa} 章}用 $(\mathbb C^*)^{\Delta\cap M}/\operatorname{im}(T\times\mathbb C^*)$ 粗略近似模空间。子空间 $(\mathbb C^*)^{(\Delta\cap M)_0}/\operatorname{im}(T\times\mathbb C^*)$ 恰好就是环面 $T_0\subset\overline{\mathcal M}_{\mathrm{simp}}$。$f\in L((\Delta\cap M)_0)$ 在由 $T_0$ 参数化的族上给出相对 $n$ 形式
\[
\widetilde\omega=\frac{\lambda_0}{f}\frac{dt_1}{t_1}\wedge\cdots\wedge\frac{dt_n}{t_n},
\]
已经证明它满足方程 $\Box_\ell\lambda_0^{-1}\widetilde\omega=0$ 及 $Z_j\lambda_0^{-1}\widetilde\omega=0$。后一个方程只是表达 $\widetilde\omega$ 在 $T$ 和 $\mathbb C^*$ 的作用下不变（这是因为 $\widetilde\omega$ 中额外加入了因子 $\lambda_0$），所以它下降为 $\overline{\mathcal M}_{\mathrm{simp}}$ 所参数化的族上的有理相对 $n$ 形式。因此，$\overline{\mathcal M}_{\mathrm{simp}}$ 上的 GKZ 方程组为
\[
\Box_\ell\lambda_0^{-1}\widetilde\omega=0,\qquad\ell\in M_0.
\]
为使用这些方程，需要以 $\overline{\mathcal M}_{\mathrm{simp}}$ 上的局部坐标表达它们。每个 $\ell=(\ell_i)\in M_0$ 给出
\begin{equation}\label{eq:moduli-monomial-coordinate}
z=\prod_i\lambda_i^{\ell_i},
\end{equation}
它在 $T\times\mathbb C^*$ 下不变，所以选取 $M_0$ 的一组基，就给出环面 $T_0$ 的坐标 $z_1,\ldots,z_r$。\hyperref[sec:hypergeometric]{第 5.5 节}的方法于是能将方程 $\Box_\ell\lambda_0^{-1}\widetilde\omega=0$ 写成 $z_i$ 的方程，不过该节没有指定应使用 $M_0$ 的哪组基。

现在，在猜想~\ref{conj:toric-mum-boundary} 的前提下，已经确切知道如何选择坐标。固定 $\overline{\mathcal M}_{\mathrm{simp}}$ 的一个光滑环面消解 $\widetilde{\mathcal M}_{\mathrm{simp}}$，令 $\sigma$ 为消解的扇中的一个极大锥，它细分来自 $V$ 的某个 Batyrev 镜像 $V^\circ\subset X_{\Sigma^\circ}$ 的锥 $\operatorname{cpl}(\Sigma^\circ)$。由于 $\widetilde{\mathcal M}_{\mathrm{simp}}$ 光滑，$\sigma$ 的生成元构成 $N_0$ 的一组基；$M_0$ 中的对偶基给出坐标 $z_1,\ldots,z_r$，它们是 $\widetilde{\mathcal M}_{\mathrm{simp}}$ 中对应于 $\sigma$ 的点的局部坐标。

后面的\hyperref[subsec:two-parameter-moduli]{第 6.1.4 小节}将给出一个很好的例子。还应指出，上述特殊坐标 $z_1,\ldots,z_r$ 是镜像映射的一部分；\hyperref[sec:mirror-map]{第~\ref*{sec:mirror-map} 节}将进一步详细讨论。

\subsection{例题}\label{subsec:two-parameter-moduli}
现在给出两个 Calabi--Yau \IndexLocation{indexloc:0099:146}环面超曲面复结构模的例子。
\SourceExampleNumber{6.1.4}{1}
\begin{example}\label{ex:weighted-complex-moduli}
首先考虑 $\mathbb P(1,1,2,2,2)$ 的环面爆破中反典范\IndexLocation{indexloc:1020:146}超曲面的镜像族，前面
\SourcePage{147}
\hyperref[sec:hodge-examples]{第 5.4 节}、\hyperref[sec:hypergeometric]{第 5.5 节}及\hyperref[sec:yukawa-couplings]{第 5.6 节}已研究过它。沿用\hyperref[ex:weighted-pf-method]{例题~\ref*{ex:weighted-pf-method}} 的记号，自反多胞形 $\Delta^\circ$ 的顶点为
\[
(-1,-2,-2,-2),\ (1,0,0,0),\ (0,1,0,0),\ (0,0,1,0),\ (0,0,0,1).
\]
要研究 $V^\circ\subset X_{\Sigma^\circ}$ 的复结构模，其中 $\Sigma^\circ$ 是 $\Delta^\circ$ 的法扇的一个极大射影细分。

先证明 $\overline{\mathcal M}=\overline{\mathcal M}_{\mathrm{poly}}=\overline{\mathcal M}_{\mathrm{simp}}$。\hyperref[ex:weighted-pf-method]{例题~\ref*{ex:weighted-pf-method}} 指出，$\Delta^\circ\cap N$ 由 $\Delta^\circ$ 的五个顶点，以及
$(0,-1,-1,-1)=\frac12(-1,-2,-2,-2)+\frac12(1,0,0,0)$ 和 $(0,0,0,0)$ 组成。这给出 $\Delta^\circ$ 的一条具有内部格点的边；与之对偶的 $\Delta$ 的二维面没有内部格点。$\Delta^\circ$ 的余维二的面也没有内部格点，所以由命题~\ref{prop:toric-moduli-dimensions}，$\overline{\mathcal M}=\overline{\mathcal M}_{\mathrm{poly}}$。

还注意到，$\Delta^\circ$ 的刻面没有内部格点。因此，由命题~\ref{prop:automorphism-facets}，$\operatorname{Aut}(X_{\Sigma^\circ})$ 没有根；这意味着 $X_{\Sigma^\circ}$ 的环面 $T^\circ$ 就是自同构群的单位元所在连通分支。不过，由 $\Delta^\circ$ 的法扇的自同构，得到
$\operatorname{Aut}(X_{\Sigma^\circ})/T^\circ\simeq\mathbb Z/2\mathbb Z\times S_3$
（交换 $\Delta^\circ$ 的前两个顶点给出 $\mathbb Z/2\mathbb Z$，置换其余三个顶点给出 $S_3$）。需证明这些自同构在 $\overline{\mathcal M}_{\mathrm{simp}}$ 上平凡作用。注意 $f\in L((\Delta^\circ\cap N)_0)$ 可写为
\[
f=\lambda_1t_1^{-1}t_2^{-2}t_3^{-2}t_4^{-2}+\lambda_2t_1
+\lambda_3t_2+\lambda_4t_3+\lambda_5t_4
+\lambda_6t_2^{-1}t_3^{-1}t_4^{-1}+\lambda_0.
\]
交换 $t_2,t_3$ 的自同构（$S_3$ 中的一个）在 $L((\Delta^\circ\cap N)_0)$ 上诱导一个交换 $\lambda_3,\lambda_4$ 的自同构。不过，可以验证，$(1,\lambda_4/\lambda_3,\lambda_3/\lambda_4,1)\in T^\circ$ 在 $f$ 上具有相同效果，所以这个自同构在 $(T^\circ\times\mathbb C^*)$ 轨道上平凡作用。其余自同构类似处理，因而 $\overline{\mathcal M}_{\mathrm{poly}}=\overline{\mathcal M}_{\mathrm{simp}}$。

下一步构造次级扇，它由 $\Xi^+=(\Delta^\circ\cap N)_0\times\{1\}$ 构造。因此
\[
\begin{split}
\Xi^+=\{&(-1,-2,-2,-2,1),\ (1,0,0,0,1),\ (0,1,0,0,1),\\
&(0,0,1,0,1),\ (0,0,0,1,1),\ (0,-1,-1,-1,1),\ (0,0,0,0,1)\}.
\end{split}
\]
由式~\eqref{eq:moduli-relation-lattice}，$M_0$ 为 $\Xi^+$ 各点的\IndexLocation{indexloc:0701:147}线性关系格。以下矩阵的各行给出一组基：
\begin{equation}\label{eq:weighted-moduli-relations}
\begin{pmatrix}0&0&1&1&1&1&-4\\1&1&0&0&0&-2&0\end{pmatrix}.
\end{equation}
也得到 $N_0$ 的对偶基。为描述次级扇，使用\hyperref[sec:gkz-decomposition]{第 3.4 节}：它表明，在秩为二的情形中，次级扇是 $N_0\otimes\mathbb R$ 中的完备扇，其一维锥由矩阵~\eqref{eq:weighted-moduli-relations} 的各列生成。因此，在 GKZ 分解~\eqref{eq:weighted-gkz-picture} 中加入由 $(-1,0)$ 生成的锥，并补成完备扇，就得到次级扇。

回顾式~\eqref{eq:weighted-gkz-picture}，知道 $V^\circ$ 有唯一的 Batyrev 镜像，即\hyperref[ex:weighted-resolution]{例题~\ref*{ex:weighted-resolution}} 描述的 $\mathbb P(1,1,2,2,2)$ 的环面爆破中的超曲面族。GKZ 分解中对应的锥为第一象限，因此所需局部坐标 $z_1,z_2$ 由矩阵~\eqref{eq:weighted-moduli-relations} 的行给出。利用式~\eqref{eq:moduli-monomial-coordinate}，各行给出
\[
z_1=\frac{\lambda_3\lambda_4\lambda_5\lambda_6}{\lambda_0^4},\qquad
z_2=\frac{\lambda_1\lambda_2}{\lambda_6^2}.
\]
\SourcePage{148}
如\hyperref[sec:hypergeometric]{第 5.5 节}所述，可利用 $T^\circ\times\mathbb C^*$，将一般的 $f\in L((\Delta\cap M)_0)$ 化为
\begin{equation}\label{eq:weighted-moduli-laurent}
f=z_2t_1^{-1}t_2^{-2}t_3^{-2}t_4^{-2}+t_1+z_1t_2+t_3+t_4
+t_2^{-1}t_3^{-1}t_4^{-1}+1.
\end{equation}
因此，式~\eqref{eq:weighted-pf-family} 首次使用的这种写法并非偶然：按猜想~\ref{conj:toric-mum-boundary}，它恰好给出最大幺幂边界点所需的形式。

在环面边界以外，本例中的判别式轨迹 $\mathcal D\subset\overline{\mathcal M}_{\mathrm{simp}}$ 有两个分支：由 $D_{\mathcal A}=0$（$\mathcal A=(\Delta\cap M)_0$）定义的主分支，以及与 $\Delta$ 的边 $\Gamma$ 对应的分支 $D_{\Gamma\cap\mathcal A}=0$，其中
\[
\Gamma\cap M=\{(-1,-2,-2,-2),(1,0,0,0),(0,-1,-1,-1)\}
=\{m_1,m_2,m_6\}.
\]
为确定 $D_{\mathcal A}$，将 $f$ 写为式~\eqref{eq:weighted-moduli-laurent}，并确定使方程 $f=t_i\partial f/\partial t_i=0$ 在 $T^\circ$ 中有解的 $z_1,z_2$ 的值。可以手算完成，结果是
\[
D_{\mathcal A}(f)=512^2z_1^2z_2-(1-256z_1)^2.
\]
再看边 $\Gamma$，有 $f_\Gamma=z_2t^{m_1}+t^{m_2}+t^{m_6}$。由于 $m_1+m_2=2m_6$，$t^{m_1+m_2}=(t^{m_6})^2$，计算得 $D_{\mathcal A\cap\Gamma}(f_\Gamma)=1-4z_2$。Picard--Fuchs 方程~\eqref{eq:weighted-pf-second} 在 $1-4z_2=0$ 时退化，因此它选出了判别式轨迹的一个分支。对于另一个分支 $D_{\mathcal A}=0$，\hyperref[ex:weighted-yukawa]{例题~\ref*{ex:weighted-yukawa}} 中计算的 Yukawa 耦合表明，GKZ 方程组也在该分支上退化。
\end{example}

\SourceExampleNumber{6.1.4}{2}
\begin{example}\label{ex:degree-four-moduli}
考虑自反多胞形 $\Delta$，它是 $\mathbb Z^4=M$ 中以下各点的凸包：
\begin{equation}\label{eq:second-moduli-polytope}
\begin{gathered}
(-1,-2,-3,-7),\ (1,0,0,0),\ (0,1,0,0),\\
(0,0,1,0),\ (0,0,0,1),\ (0,0,-1,-2).
\end{gathered}
\end{equation}
Batyrev 构造给出镜像族 $V\subset X_\Sigma$ 和 $V^\circ\subset X_{\Sigma^\circ}$。这里的目标是描述 $V$ 的复结构模。将在\hyperref[subsec:enlarged-kahler-moduli]{第 6.2.3 小节}考虑 $V^\circ$ 的 Kähler 模。

首先观察到，$\Delta$ 的格点是式~\eqref{eq:second-moduli-polytope} 的六个顶点，加上原点 $(0,0,0,0)$ 和点 $(0,0,0,-1)$；后一点位于由除第五个顶点以外的所有顶点张成的刻面的内部。由于余维二的面没有内部格点，命题~\ref{prop:toric-moduli-dimensions} 表明 $\overline{\mathcal M}=\overline{\mathcal M}_{\mathrm{poly}}$。刻面内部恰好只有一个格点，所以同一命题还表明模空间为二维。

刻面有内部格点，还意味着 $\operatorname{Aut}(X)$ 大于环面 $T$。不过，仍有 $\overline{\mathcal M}_{\mathrm{poly}}=\overline{\mathcal M}_{\mathrm{simp}}$，现在来证明。先描述由 $(0,0,0,-1)$ 产生的额外自同构在空间 $L(\Delta\cap M)$ 上的作用。虽然可以直接计算，但用齐次坐标来做颇有启发。回顾，$X_\Sigma$ 的扇 $\Sigma$ 满足 $\Sigma(1)=\Delta^\circ\cap N-\{0\}$。可以验证 $\Delta^\circ$ 是以下各点的凸包：
\[
\begin{gathered}
(13,-1,-1,-1),\ (1,-1,3,-1),\ (-1,6,-1,-1),\ (-1,-1,3,-1),\\
(-1,-1,-1,-1),\ (-1,-1,-1,1),\ (-1,0,3,-1).
\end{gathered}
\]
$\Delta^\circ$ 还有许多其他格点。为简化问题，回顾超曲面 $V$ 都是 $\mathbb P_\Delta$ 中反典范超曲面的拉回。因此，可以使用与法扇本身对应的齐次坐标
\SourcePage{149}
进行计算。将按上述次序使用齐次坐标 $x_1,\ldots,x_7$。

由\hyperref[sec:toric-automorphisms]{第 3.6 节}，知道 $\operatorname{Aut}(X_\Sigma)$ 唯一的根来自 $(0,0,0,-1)$，因为它是 $\Delta$ 的刻面内部唯一的格点。包含它的刻面与 $\Delta^\circ$ 的顶点 $(-1,-1,-1,1)$ 对偶，因此很容易看出，相应的根为 $(x_6,x_1x_2x_3x_4x_5x_7)$。自同构的相应单参数族为
\[
y_u(x_6)=x_6+u x_1x_2x_3x_4x_5x_7,\qquad y_u(x_i)=x_i\quad(i\ne6).
\]
回顾命题~\ref{prop:automorphism-facets} 之前的讨论：Laurent 单项式 $t^m\in L(\Delta\cap M)$ 在齐次坐标环中给出单项式
\[
t^m=\prod_i x_i^{\langle m,v_i\rangle+1}.
\]
这里 $v_i$ 为上面列出的 $\Delta^\circ$ 的顶点。将 $\Delta\cap M$ 的点记为：式~\eqref{eq:second-moduli-polytope} 中依次列出的点为 $m_1,\ldots,m_6$，$m_7=(0,0,0,-1)$，$m_0=(0,0,0,0)$。于是
\[
\begin{aligned}
t^{m_0}&=x_1x_2x_3x_4x_5x_6x_7,\\
t^{m_1}&=x_4^2x_5^{14},&t^{m_2}&=x_1^{14}x_2^2,\\
t^{m_3}&=x_3^7x_7,&t^{m_4}&=x_2^4x_4^4x_7^4,\\
t^{m_5}&=x_6^2,&t^{m_6}&=x_1^4x_3^4x_5^4,\\
t^{m_7}&=x_1^2x_2^2x_3^2x_4^2x_5^2x_7^2.
\end{aligned}
\]
自同构的单参数族 $y_u$ 在这些单项式上的作用为
\[
\begin{aligned}
y_u(t^{m_0})&=t^{m_0}+u t^{m_7},\\
y_u(t^{m_5})&=t^{m_5}+2u t^{m_0}+u^2t^{m_7},\\
y_u(t^{m_i})&=t^{m_i},\qquad i\ne0,5.
\end{aligned}
\]
现在证明，自然映射 $\overline{\mathcal M}_{\mathrm{simp}}\to\overline{\mathcal M}_{\mathrm{poly}}$ 在一般位置上是一一的。第一步是说明保持 $\Delta$ 的格自同构在 $\overline{\mathcal M}_{\mathrm{simp}}$ 上平凡作用；这可用\hyperref[ex:weighted-complex-moduli]{例题~\ref*{ex:weighted-complex-moduli}} 的方法完成。所以只需考虑 $\operatorname{Aut}(X_\Sigma)$ 的单位元所在连通分支，它由 $T$ 和自同构 $y_u$ 生成。取 $\overline{\mathcal M}_{\mathrm{simp}}$ 中一点，以 $f=\sum_{i=0}^6\lambda_i t^{m_i}$ 表示，并施以 $y_u$。$y_u(f)$ 中 $t^{m_7}$ 的系数为 $u\lambda_0+u^2\lambda_5$。令该系数等于零，就能找到 $\overline{\mathcal M}_{\mathrm{simp}}$ 中与 $f$ 在 $\overline{\mathcal M}_{\mathrm{poly}}$ 中有相同像的所有点。有两个解：$u=0$ 和 $u=-\lambda_0/\lambda_5$。非平凡解 $u=-\lambda_0/\lambda_5$ 给出
\[
y_u(f)=-\lambda_0t^{m_0}+\sum_{i=1}^6\lambda_i t^{m_i}.
\]
不过，它与 $f$ 表示 $\overline{\mathcal M}_{\mathrm{simp}}$ 的同一点，因为可通过 $(-1,-1,-1,-1,-1)\in T\times\mathbb C^*$ 的作用由 $f$ 得到它。因此，$\overline{\mathcal M}_{\mathrm{simp}}=\overline{\mathcal M}_{\mathrm{poly}}$。

所以可用次级扇紧化 $V$ 的复结构模空间；这里次级扇由 $\Xi^+=(\Delta\cap M)_0\times\{1\}$ 构造。为线性关系格 $M_0$ 选取一组基，由以下矩阵的各行给出：
\[
\begin{pmatrix}-2&-2&-4&1&0&7&0\\1&1&2&0&1&-3&-2\end{pmatrix}.
\]
与上一例一样，可从矩阵的各列读出次级扇的一维锥。因此，$\overline{\mathcal M}_{\mathrm{simp}}$ 是与
\SourcePage{150}
下图所示的扇对应的环面簇：
\begin{figure}[H]
\centering
\refstepcounter{equation}\label{eq:second-moduli-secondary-fan}
\makebox[\textwidth]{\makebox[0pt][l]{(\theequation)}\hfill
\raisebox{-.5\height}{\includegraphics[width=0.53975\textwidth]{6-1-secondary-fan.png}}\hfill\mbox{}}
\end{figure}
在\hyperref[subsec:enlarged-kahler-moduli]{第 6.2.3 小节}计算镜像 $V^\circ$ 的 Kähler 模时，将看到 GKZ 分解由锥 $\sigma_1,\sigma_2,\sigma_3$ 组成（记住，GKZ 分解是次级扇的一个子扇）。不过，只有 $\sigma_1,\sigma_2$ 具有 $\operatorname{cpl}(\Sigma^\circ)$ 的形式，其中 $\Sigma^\circ$ 是 $\Delta^\circ$ 的法扇的适当射影细分。因此，只有 $\sigma_1,\sigma_2$ 给出 $\overline{\mathcal M}_{\mathrm{simp}}$ 的最大幺幂边界点。

这两个锥中，只有 $\sigma_2$ 光滑，所以要将理论应用于 $\sigma_1$，必须对它细分。这会在一个消解 $\widetilde{\mathcal M}_{\mathrm{simp}}$ 中给出多个最大幺幂边界点。不过，还有另一种处理办法：将 $\sigma_1,\sigma_2$ \emph{合并}成一个锥 $\sigma=\sigma_1\cup\sigma_2$。将 $\sigma$ 与式~\eqref{eq:second-moduli-secondary-fan} 中其余锥 $\sigma_3,\sigma_4,\sigma_5$ 合用，得到一个爆缩
$\overline{\mathcal M}_{\mathrm{simp}}\to\overline{\mathcal M}'_{\mathrm{simp}}$，而 $\overline{\mathcal M}'_{\mathrm{simp}}$ 中对应于 $\sigma$ 的点 $p$ 光滑。可以证明，$p$ 是最大幺幂边界点；并且在爆破 $\widetilde{\mathcal M}_{\mathrm{simp}}\to\overline{\mathcal M}'_{\mathrm{simp}}$ 下，$\widetilde{\mathcal M}_{\mathrm{simp}}$ 的所有最大幺幂边界点都映到 $p$，在定义~\ref{def:equivalent-mum-points} 的意义下彼此等价。因此，从\IndexLocation{indexloc:0772:150}最大幺幂单值性的角度看，$\overline{\mathcal M}'_{\mathrm{simp}}$ 是最适宜采用的紧化。

\hyperref[subsec:enlarged-kahler-moduli]{第 6.2.3 小节}将从镜像流形 $V^\circ$ 的 Kähler 模空间的角度重新讨论本例。特别是，将了解到为何应将 $\sigma_1,\sigma_2$ 合并为一个锥 $\sigma$ 的系统性原因。
\end{example}

\section{Kähler 模}\label{sec:kahler-moduli}\index{Kähler 模空间|hyperpage}
现在转向 Kähler 模。维数至少为 $2$ 的 Calabi--Yau 流形 $V$ 具有 Kähler 锥 $K(V)$，并且如定义~\ref{def:kahler-moduli} 所述，具有复化 Kähler 空间
\[
K_{\mathbb C}(V)=\{\omega\in H^2(V,\mathbb C):\operatorname{Im}(\omega)\in K(V)\}/\operatorname{im}H^2(V,\mathbb Z).
\]
\emph{复化 Kähler 模空间}则是商
\[
\mathcal{KM}(V)=K_{\mathbb C}(V)/\operatorname{Aut}(V).
\]
不会造成混淆时，将 $\mathcal{KM}(V)$ 简记为 $\mathcal{KM}$。本节的目标是理解 $\mathcal{KM}$ 的结构，特别是找到一个良好的紧化，其中包含与\hyperref[sec:complex-moduli]{第 6.1 节}研究的最大幺幂边界点相对应的点。与该节一样，在构造中保持一定灵活性会更方便。
\SourcePage{151}
\hyperref[sec:mirror-map]{第~\ref*{sec:mirror-map} 节}将讨论镜像映射。猜想它把 $\mathcal{KM}(V)$ 的\emph{大半径极限点}（将在\hyperref[subsec:large-radius-limits]{第 6.2.1 小节}定义）的一个邻域，映到 $V^\circ$ 的复结构模空间中一个最大幺幂边界点的邻域。这将把读者在本节与\hyperref[sec:complex-moduli]{第 6.1 节}之间发现的相似之处联系起来。镜像映射可用来将这些邻域与 SCFT 模空间的子集识别起来。一开始就必须强调，镜像映射只在局部定义，所以目前试图把本节构造的整个紧化，与 SCFT 模空间的一部分识别起来，会有风险。另一方面，将在例子中看到，本节构造的紧化确实为理解 SCFT 模空间的整体结构提供了有用的几何框架，不过还须由镜像映射加以“修正”。如何给出 SCFT 模空间的精确数学描述，仍是一个有待解开的精彩谜题。

\subsection{大半径极限点}\label{subsec:large-radius-limits}
研究 Kähler 模的第一步，是为复化 Kähler 空间 $K_{\mathbb C}(V)$ 找到一些良好的部分紧化。为此，固定一个满维单纯锥 $\sigma$，其内部位于 Kähler 锥 $K(V)$ 中，并令
\[
\mathcal D_\sigma=(H^2(V,\mathbb R)+i\operatorname{Int}(\sigma))/\operatorname{im}H^2(V,\mathbb Z)
\subset K_{\mathbb C}(V).
\]
注意，$\mathcal D_\sigma$ 是 $K_{\mathbb C}(V)$ 的开子集。进一步假定 $\sigma$ 由 $H^2(V,\mathbb Z)/\text{挠元}$ 的一组基 $T_1,\ldots,T_r$ 生成。$T_i$ 位于 Kähler 锥的闭包中。于是
$\operatorname{Int}(\sigma)=\{t_1T_1+\cdots+t_rT_r:t_1,\ldots,t_r>0\}$，而映射
\begin{equation}\label{eq:large-radius-exponential-coordinates}
t_1T_1+\cdots+t_rT_r\longmapsto(q_1,\ldots,q_r)
=(e^{2\pi i t_1},\ldots,e^{2\pi i t_r})
\end{equation}
诱导双全纯同构
\[
\mathcal D_\sigma\simeq(\Delta^*)^r.
\]
利用该同构及包含映射 $(\Delta^*)^r\subset\Delta^r$，得到部分紧化
\[
\mathcal D_\sigma\subset\mathcal D_\sigma^-\simeq\Delta^r.
\]
特别地，原点 $0\in\Delta^r$ 成为一个特殊边界点，记为 $0\in\mathcal D_\sigma^-$。

对 $\omega\in H^2(V,\mathbb R)+i\operatorname{Int}(\sigma)$，当 $\omega=\sum_jt_jT_j$ 的所有 $\operatorname{Im}(t_j)$ 都趋于 $+\infty$ 时，在映射~\eqref{eq:large-radius-exponential-coordinates} 下，$\omega$ 的像趋于 $0$。因此，称 $0\in\mathcal D_\sigma^-$ 为\emph{大半径极限点}。这是 Kähler 一侧与最大幺幂边界点对应的概念。注意，每个大半径极限点都具有上述典范坐标 $q_1,\ldots,q_r$。
\index{大半径极限点|hyperpage}

对于\hyperref[subsec:nef-gkz]{第 5.5.3 小节}定义的 nef 完全交，命题~\ref{prop:nef-i-gkz} 中已出现坐标 $q_1,\ldots,q_r$ 的一种版本。那里令 $q_j=e^{t_j}$，并对 $\beta\in H_2(V,\mathbb Z)$ 置
\[
q^\beta=\prod_jq_j^{\int_\beta T_j}
=\exp\left(\int_\beta(t_1T_1+\cdots+t_rT_r)\right)
\]
（见式~\eqref{eq:nef-gkz-q-coordinates} 后面的讨论）。若将 $t_j$ 换成 $2\pi i t_j$，就恰好得到上面的变量 $t_j,q_j$。所以，除多出的因子 $2\pi i$ 外，命题~\ref{prop:nef-i-gkz} 的形式函数 $\widetilde I$ 内在地定义于复化 Kähler 空间 $K_{\mathbb C}(V)$ 。由于 $\widetilde I$ 与镜像 $V^\circ$ 的 Picard--Fuchs 方程有关，这已经提示了\IndexLocation{indexloc:0830:151}镜像对称给出的复结构模空间与 Kähler 模空间之间的关系。

\SourcePage{152}
定义了大半径极限点以后，接着注意到，$\sigma$ 的不同选择给出许多这样的点。将 $\mathcal D_\sigma^-$ 上相应的坐标记为 $q_1^\sigma,\ldots,q_r^\sigma$，就可以研究改变 $\sigma$ 时这些坐标如何变化。例如，设有两个这样的锥，且 $\tau\subset\sigma$。这给出包含关系 $\mathcal D_\tau\subset\mathcal D_\sigma$。为联系相应的部分紧化，将 $\tau$ 的生成元用 $\sigma$ 的生成元表示为
\[
T_i^\tau=\sum_{j=1}^r a_{ij}T_j^\sigma,\qquad a_{ij}\geq0.
\]
则包含映射 $\mathcal D_\tau\subset\mathcal D_\sigma$ 延拓为映射 $\mathcal D_\tau^-\to\mathcal D_\sigma^-$，其表达式为
\[
q_j^\sigma=\prod_{i=1}^r(q_i^\tau)^{a_{ij}}.
\]
这令人想起式~\eqref{eq:mum-log-coordinate-change} 和式~\eqref{eq:mum-q-monomial-change}，而且相似之处比预想的更深。首先，\hyperref[subsec:mum-boundary-points]{第 6.1.1 小节}也有锥，即由单值性算子的对数生成的锥；式~\eqref{eq:mum-log-coordinate-change} 来自一个锥包含于另一个锥。此外，\hyperref[chap:quantum-cohomology]{第~\ref*{chap:quantum-cohomology} 章}研究量子上同调时，将看到这里锥 $\sigma$ 的生成元 $T_i$ 自然对应于 A 模型联络的单值性的对数。因此，两种情形极其相似，这正是镜像对称的预言。

大半径极限点与最大幺幂边界点之间，一个可能的差别在于：前者的局部坐标唯一，后者却因式~\eqref{eq:mum-coordinate-basis-change} 而不唯一。如\hyperref[subsec:mirror-map-definition]{第~\ref*{subsec:mirror-map-definition} 小节}将说明的，假定式~\eqref{eq:mum-canonical-coordinates} 中 $g_0,\ldots,g_r$ 具有整性，就能消除这一差别。这正好对应于上面假定 $T_1,\ldots,T_r$ 具有整性。

接着讨论特殊情形的 Kähler 模：$\operatorname{Aut}(V)=\{1\}$，且 $K(V)$ 是 $H^2(V,\mathbb R)$ 中一个有理多面锥的内部。此时 Kähler 模空间为 $\mathcal{KM}=K_{\mathbb C}(V)$，可以如下构造它的环面部分紧化。闭包 $\overline{K(V)}$ 关于格 $N=H^2(V,\mathbb Z)/\text{挠元}$ 为有理多面锥，所以得到\IndexLocation{indexloc:1354:152}仿射环面簇 $X_{\overline{K(V)}}$，简记为 $X_K$。同构
\begin{equation}\label{eq:kahler-torus-isomorphism}
H^2(V,\mathbb C)/N\simeq N\otimes\mathbb C^*=T_N=X_K\text{ 的环面}
\end{equation}
由 $K_{\mathbb C}(V)$ 的定义给出开集的自然包含映射
\begin{equation}\label{eq:kahler-toric-inclusion}
\mathcal{KM}=K_{\mathbb C}(V)\subset T_N\subset X_K.
\end{equation}
利用环面簇的边界 $B=X_K-T_N$，将部分紧化 $\overline{\mathcal{KM}}$ 定义为 $\mathcal{KM}$，加上 $B$ 中所有属于 $\mathcal{KM}$ 在 $X_K$ 中的像的闭包的点。

锥 $\overline{K(V)}$ 确定一点 $0\in X_K$，它是环面作用的不动点，并且可能有很严重的奇点。可以验证 $0\in\overline{\mathcal{KM}}$，因为在式~\eqref{eq:kahler-torus-isomorphism} 下，沿 Kähler 锥内部的射线使虚部趋于无穷时，像趋于 $0$。所以 $\overline{\mathcal{KM}}$ 很可能有奇点。

用标准环面方法消解该奇点。令 $\Sigma_+$ 为支集 $\overline{K(V)}$ 上的一个扇，使环面簇 $X_{\Sigma_+}$ 光滑。将 $X_K$ 换成 $X_{\Sigma_+}$ 后，式~\eqref{eq:kahler-toric-inclusion} 仍然适用；可以完全按前面的方式定义 $\overline{\mathcal{KM}}$，使它的边界由 $X_{\Sigma_+}-T_N$ 中属于 $\mathcal{KM}$ 的像的闭包的点组成。

\SourcePage{153}
$\Sigma_+$ 的极大锥给出 $X_{\Sigma_+}$ 的特殊点。我们断言，这些点全都是 $\overline{\mathcal{KM}}$ 的大半径极限点。原因在于，极大锥都是满足 $\operatorname{Int}(\sigma)\subset K(V)$ 的单纯锥 $\sigma$；这正是定义大半径极限点时所用的锥。将 $\mathcal{KM}$ 视为 $X_{\Sigma_+}$ 的子集。仿射环面簇 $X_\sigma$ 就是 $\mathbb C^r$；在其中取坐标多圆盘 $U_\sigma=\{|q_j|<1\text{，对所有 }j\}$，由式~\eqref{eq:kahler-toric-inclusion} 得到
\[
\mathcal D_\sigma=U_\sigma\cap\mathcal{KM}.
\]
很容易推出：$U_\sigma\cap\overline{\mathcal{KM}}$ \emph{恰好}是前面讨论的部分紧化 $\mathcal D_\sigma^-$。

因此，得到一个良好的部分紧化，其中有有限多个大半径极限点。此外，若 $\omega\in H^2(V,\mathbb R)+iK(V)$ 的虚部在某个极大锥 $\sigma\in\Sigma_+$ 的全部生成元方向上都趋于 $+\infty$，那么 $\omega$ 在 $\mathcal{KM}$ 中的像就趋于 $\overline{\mathcal{KM}}$ 中相应的大半径极限点。

这幅图景很好，但仍十分特殊。对于一般的 Calabi--Yau 流形，自同构群未必平凡，Kähler 锥也未必为多面锥。$K(V)$ 的结构可能相当复杂；\hyperref[subsec:threefold-kahler-moduli]{第 6.2.2 小节}详细研究三维簇时将看到这一点。一般而言，所能期望的最好情形是\emph{锥猜想} \SourceCite{Morrison3}。为陈述它，令 $K_+$ 为 $\overline{K(V)}\cap H^2(V,\mathbb Q)$ 的凸包。由于 $K(V)$ 为开集，$K(V)\subset K_+$，所以 $K_+$ 由 $K(V)$ 及其边界中由有理点生成的部分组成。
\begin{conjecture}\label{conj:kahler-cone}
设 $V$ 为维数至少为 $2$ 的 Calabi--Yau 流形，Kähler 锥为 $K(V)$。则存在有理多面锥 $\Pi\subset K_+$，使得
\[
\operatorname{Aut}(V)\cdot\Pi=K_+.
\]
\end{conjecture}
\index{锥猜想|hyperpage}
\SourceCite{GM} 已在一个非平凡例子中验证了锥猜想。此外，关于 $\operatorname{Aut}(V)$ 在 Kähler 锥上的作用，一些有趣的有限性结论是猜想~\ref{conj:kahler-cone} 的推论。例如，它表明模去 $\operatorname{Aut}(V)$ 后，$\overline{K(V)}$ 只有有限多条边。\SourceCite{Borcea1} 在一个特定 Calabi--Yau 三维簇 $V$ 上验证了这一点；该三维簇具有非多面的 Kähler 锥及无限的 $\operatorname{Aut}(V)$。锥猜想还表明，某些类的 Calabi--Yau 三维簇仅具有有限多个代数纤维空间结构；这一结论已知成立 \SourceCite{Oguiso,OgP}。\SourceCite{Szendroi} 证明了锥猜想的另一个推论。\SourceCite{Kawamata2} 提出了锥猜想的相对版本，并在其中证明了低维 Calabi--Yau 纤维空间的情形。

对我们而言，锥猜想的关键性质是：对于\emph{任何}满足它的 Calabi--Yau 流形，都能构造 $\mathcal{KM}$ 的一个部分紧化。因此，假定 $V$ 满足猜想~\ref{conj:kahler-cone}。Looijenga 关于管状域（如 $H^2(V,\mathbb R)+iK(V)$）的商的部分紧化的工作 \SourceCite{Looijenga} 表明，可以找到一个 $\operatorname{Aut}(V)$ 不变的扇 $\Sigma_+$（可能无限），其支集为 $K_+$。还可以假定自同构群在 $\Sigma_+$ 上自由作用，且 $\Sigma_+$ 的极大锥具有上述 $\sigma$ 的形式。

根据这些数据，Looijenga 从扇 $\Sigma_+$ 确定的各层构造空间 $\widehat{\mathcal D}(\Sigma_+)$。困难之处是赋予 $\widehat{\mathcal D}(\Sigma_+)$ 正确的拓扑，使商 $\widehat{\mathcal D}(\Sigma_+)/\operatorname{Aut}(V)$ 成为正规解析空间。这个商称为\emph{半环面部分紧化}，就是所需的部分紧化 $\overline{\mathcal{KM}}$。\SourceCite{Morrison3} 描述了这一构造的具体方式。

\SourcePage{154}
当 $\operatorname{Aut}(V)$ 正常不连续地作用，且 $K(V)$ 是齐次自伴锥 \SourceCite{AMRT} 时，完全可以在环面几何的范围内完成这一构造。为便于说明，在这一特殊情形中考察它与前面构造的关系。其思想是：无限扇 $\Sigma_+$ 确定一个局部有限型概形 $X_{\Sigma_+}$，将式~\eqref{eq:kahler-toric-inclusion} 中的 $X_K$ 换为 $X_{\Sigma_+}$，得到开集的包含映射
\[
K_{\mathbb C}(V)\subset X_{\Sigma_+}.
\]
由于扇在 $\operatorname{Aut}(V)$ 下不变，这一包含映射关于 $\operatorname{Aut}(V)$ 等变，所以得到包含映射
\[
\mathcal{KM}\subset X_{\Sigma_+}/\operatorname{Aut}(V).
\]
\SourceCite{AMRT} 给出了 $X_{\Sigma_+}/\operatorname{Aut}(V)$ 的一幅很好的图示，\SourceCite{Morrison3} 也转载了它。利用 $X_{\Sigma_+}/\operatorname{Aut}(V)$ 的环面边界，完全按前面的方式得到部分紧化 $\overline{\mathcal{KM}}$。此外，因为在这一情形中，有效自同构群在扇上自由作用，每个极大锥都给出一个大半径极限点，其邻域同构于 $\mathcal D_\sigma^-$。（不过，由于 $\operatorname{Aut}(V)$ 的作用，$\Sigma_+$ 中无限多个不同的锥可能给出同一个大半径极限点。）如 \SourceCite{Morrison3} 所示，这幅图景在 Looijenga 考虑的更一般情形中仍然成立。

所有这些讨论的结果是：得到一个良好的部分紧化 $\overline{\mathcal{KM}}$，其中包含有限多个大半径极限点。不过，这样的紧化有许多种，因为给定一个紧化后，总能进一步细分扇，构造其他紧化。因此，与最大幺幂边界点的情形一样，依照 \SourceCite{Morrison6}，在大半径极限点之间引入如下等价关系。
\begin{definition}\label{def:equivalent-large-radius-points}\index{大半径极限点!等价关系|hyperpage}
设 $\overline{\mathcal{KM}}$ 为 Kähler 模空间 $\mathcal{KM}$ 的一个光滑正规交叉部分紧化。若存在包含 $p,q$ 的连通集 $\Xi\subset\overline{\mathcal{KM}}-\mathcal{KM}$，并且在 $\Xi$ 的一个邻域上存在由 $\overline{\mathcal{KM}}$ 的余切丛中、沿边界有对数极点的形式组成的局部系统 $\mathcal L$，使得在 $p,q$ 各自附近，该局部系统由 $1$ 形式 $d\log q_k$ 张成，其中 $q_k$ 如式~\eqref{eq:large-radius-exponential-coordinates} 所述，则称大半径极限点 $p,q\in\overline{\mathcal{KM}}$ \emph{等价}。
\end{definition}
很容易看出，进一步细分扇 $\Sigma_+$ 所产生的多个大半径极限点，在这一意义下彼此等价。

如\hyperref[sec:complex-moduli]{第 6.1 节}所解释的，镜像对称告诉我们：$V$ 的 Kähler 模应通过镜像映射，同构于镜像 $V^\circ$ 的复结构模的一部分，且使大半径极限点的每个等价类映到最大幺幂边界点的一个等价类。事实上，镜像映射应保持两个等价关系定义中提到的局部系统。$\overline{\mathcal{KM}}$ 的构造有很大的非唯一性，镜像复结构模的紧化也不唯一。因此，镜像猜想的一部分内容，就是能找到相容的紧化。

实际处理中，在 $X_{\Sigma_+}$ 光滑的假设下，选取锥尽可能大的 $\Sigma_+$。这样可以缩小等价类，与\hyperref[sec:complex-moduli]{第 6.1 节}处理复结构模的做法相似。不过，若允许 $\mathcal{KM}$ 的紧化有奇点，则存在一个典范的极小紧化；它由包含 $K_+$ 及其各个面的扇 $\Sigma_{\mathrm{SBB}}$ 确定。由于 $K_+$ 可能有无限多个面，环面方法不适用，但 Looijenga 的构造
\SourcePage{155}
仍然有效，所得极小紧化记为 $\overline{\mathcal{KM}}_{\mathrm{SBB}}$。这一构造推广了有界对称域的商的 Satake--Baily--Borel 紧化构造，这也解释了该记号。当 $\operatorname{Aut}(V)=1$，且 $\overline{K(V)}$ 为有理多面锥时，$\overline{\mathcal{KM}}_{\mathrm{SBB}}$ 就是前面由 $\overline{K(V)}$ 的仿射环面簇 $X_K$ 得到的紧化。

Kähler 模存在 Satake--Baily--Borel 型极小紧化，通过镜像对称意味着：镜像的复结构模也应有类似的极小紧化。\SourceCite{Morrison3} 提出了这一方向的精确猜想。

最后讨论一个读者可能已经想到的问题：为什么对复结构模作紧化，却只对 Kähler 模作\emph{部分}紧化？原因是，$\mathcal{KM}$ 实际上只是整个 Kähler 模空间的一小部分；要找到缺失的部分，就需要“越出 Kähler 锥” \SourceCite{Morrison5}。一般 Calabi--Yau 流形的相应理论尚未建立，但在三维簇及环面超曲面两个特殊情形中，已经进行了广泛研究。因此，接着讨论这两类十分有趣的簇的 Kähler 模。

\subsection{Calabi--Yau 三维簇的 Kähler 模}\label{subsec:threefold-kahler-moduli}\index{Kähler 锥!三维簇|hyperpage}
令 $V$ 为光滑 Calabi--Yau 三维簇。这里有两个目标：首先，描述关于 Kähler 锥 $K(V)$ 及其闭包 $\overline{K(V)}$ 的结构的已知结果；其次，解释如何扩大 Kähler 模空间。

描述 Kähler 锥时，一个重要要素是\emph{三次锥} $W^*\subset H^2(V,\mathbb R)$，它定义为满足 $D^3=0$ 的所有类 $D$ 所组成的集合。基本思想是：在 $W^*$ 以外，$\overline{K(V)}$ 很容易理解。将使用 Wilson 的工作 \SourceCite{Wilson2,Wilson1}，它将 $\overline{K(V)}$ 的结构与本原收缩联系起来。
\begin{definition}\label{def:primitive-contraction}
\emph{收缩}是双有理态射 $V\to\overline V$，其中 $\overline V$ 为射影正规三维簇，且 $\rho(\overline V)<\rho(V)$（通常以 $\rho$ 表示 Picard 数）。若收缩不能进一步分解为正规簇之间的非平凡双有理态射，则称它为\emph{本原收缩}。
\end{definition}
\index{本原收缩|hyperpage}
若 $\phi$ 为本原收缩，则可以证明：除倍数外，$V$ 上被 $\phi$ 收缩的 $1$ 圈的数值类 $C$ 唯一。这定义超平面 $\mathcal H=\{D\in H^2(V,\mathbb R):D\cdot C=0\}$；很容易看出，该超平面支撑 $\overline{K(V)}$ 的一个余维一的面。称这个面\emph{对应于}该本原收缩。这些面决定 $\overline{K(V)}$ 在三次锥 $W^*$ 以外的结构。更确切地说，有如下结果。
\begin{proposition}[\SourceCite{Wilson2}]\label{prop:kahler-cone-contractions}\index{三次锥|hyperpage}
设 $V$ 为光滑 Calabi--Yau 三维簇。
\begin{enumerate}[label=(\roman*)]
\item 在 $W^*$ 以外，锥 $\overline{K(V)}$ 局部为有理多面锥。
\item 在 $W^*$ 以外，$\overline{K(V)}$ 的余维一的面对应于本原双有理收缩。
\end{enumerate}
\end{proposition}
这仍允许 Kähler 锥具有相当复杂的结构。例如，$\overline{K(V)}$ 可能在 $W^*$ 以外有无限多个面，并向 $W^*$ 聚集。还应指出，$W^*$ 可能包含 $\overline{K(V)}$ 的余维一的面。这样的面由将 $V$ 映到一个较低维簇的线性系 $|D|$ 确定，而该线性系的一般成员 $D$ 的类位于这个面上。
\SourcePage{156}
接着描述如何扩大 $V$ 的 Kähler 模空间。思想是将 Kähler 锥 $K(V)$ 跨过 $\overline{K(V)}$ 的某些面延伸出去。将使用对应于本原收缩的面；回顾，这样的收缩 $\phi$ 分为三种类型：
\begin{itemize}
\item I 型：存在曲线 $C\subset V$，使 $\phi$ 在 $V-C$ 上为同构，并将 $C$ 映到一点。
\item II 型：存在曲面 $S\subset V$，使 $\phi$ 在 $V-S$ 上为同构，并将 $S$ 映到一点。
\item III 型：存在曲面 $S\subset V$，使 $\phi$ 在 $V-S$ 上为同构，并将 $S$ 映到一条曲线。
\end{itemize}
现在假定 $\overline{K(V)}$ 的一个面对应于 I 型本原收缩 $\phi:V\to\overline V$。有一条不必不可约的曲线 $C\subset V$ 被 $\phi$ 收缩。三维簇几何的一个基本定理表明，可以对 $C$ 作\emph{翻越} \SourceCite{Reid2,Kollar}。特别地，这意味着存在不同于 $V$ 的 Calabi--Yau 三维簇 $V'$、曲线 $C'\subset V'$ 及映射 $\phi':V'\to\overline V$，它收缩 $C'$，并在 $V'-C'$ 上为同构。所得双有理映射 $V\dashrightarrow V'$ 在余维一处为同构，所以有自然识别 $H^2(V)\simeq H^2(V')$。在这一识别下，$\overline{K(V)}$ 与 $\overline{K(V')}$ 沿各自的一个公共面相交；它们分别由超平面 $\{D:D\cdot C=0\}$ 及 $\{D':D'\cdot C'=0\}$ 支撑。
\index{翻越|hyperpage}

Kawamata 的结果 \SourceCite{Kawamata1} 表明：对于所有可从 $V$ 经一系列翻越得到的 Calabi--Yau 三维簇 $V'$（包括 $V$ 自身），锥 $\overline{K(V')}$ 构成其并集的凸包的一个多面分解的各个腔；该凸包称为\emph{可动锥} $\operatorname{Mov}(V)$。因为 $V$ 的所有双有理模型都由一系列翻越构造，$V$ 的双有理自同构群 $\operatorname{Bir}(V)$ 自然作用于 $\operatorname{Mov}(V)$。锥猜想的双有理版本是\emph{\IndexLocation{indexloc:0158:156}双有理锥猜想} \SourceCite{Morrison5}；它断言存在有理多面锥 $\Pi$，使
$\operatorname{Bir}(V)\cdot\Pi=\operatorname{Mov}(V)_+$，其中 $\operatorname{Mov}(V)_+$ 是 $\operatorname{Mov}(V)\cap H^2(V,\mathbb Q)$ 的凸包。用与前面类似的构造，可以适当细分 $\operatorname{Mov}(V)$，同时构造所有 $V'$ 的 Kähler 模空间的一个部分紧化。\SourceCite{Morrison5} 给出了更详细的说明；\hyperref[subsec:enlarged-kahler-moduli]{第 6.2.3 小节}将结合环面超曲面重新讨论这一思想。
\index{可动锥|hyperpage}

还可以进一步扩大 Kähler 模空间。设 $\phi$ 是 III 型收缩，其例外曲面为一个以亏格 $g\geq2$ 的曲线 $C$ 为底的直纹曲面 $S$，且 $\phi$ 沿直纹将 $S$ 收缩到 $C$。若对 $V$ 的复结构作一般形变，$S\to C$ 的一般纤维就不再保持全纯，但有 $2g-2$ 个纤维可以全纯形变。事实上，一个纤维的一阶形变障碍，就是 Kodaira--Spencer 类 $\rho\in H^1(\Theta_V)$ 在自然映射
$r:H^1(\Theta_V)\to H^0(K_C)$ 下的像；若 $\rho$ 一般，则 $r(\rho)$ 在 $2g-2$ 个不同的点为零。这些点上的纤维可以一阶形变 \SourceCite{Wilson2}。还可计算出没有更高阶障碍。于是可同时翻越这 $2g-2$ 条形变后的曲线。翻越所得 Calabi--Yau 三维簇属于同一个复结构模空间，但复结构不同。Kähler 结构也随之改变。所得 $H^2$ 上的同构可以沿形变追溯回原来的复结构，从而成为 $H^2(V)$ 的自同构。置 $E=[S]$，该自同构为
\[
D\longmapsto D+(D\cdot f)E,\qquad f=\text{纤维的类}.
\]

\SourcePage{157}
由于 $E\cdot f=-2$，这是关于超平面 $\{D:D\cdot f=0\}$ 的反射。用上述所有可能的 III 型收缩给出的反射，对 $\operatorname{Mov}(V)$ 作反射，就得到\emph{反射可动锥} \SourceCite{Morrison5}。若同时扩大双有理自同构群，使它包含上述反射，就得到一个更大的 Kähler 模空间，使用的正是反射可动锥及刚才描述的扩大自同构群。此外，锥猜想的适当版本还能使我们对这一扩大模空间作部分紧化。

反射可动锥给出的 Kähler 模空间，是跨越 $\overline{K(V)}$ 中对应于 I 型和 III 型收缩的墙构造的，不过只使用了将直纹曲面收缩到亏格 $g\geq2$ 的曲线的 III 型收缩。目前，跨过对应于 II 型及其余 III 型收缩的墙后，应如何扩大 Kähler 模空间，还没有方案。不过，在\hyperref[subsec:enlarged-kahler-moduli]{第 6.2.3 小节}将看到，对于 Calabi--Yau 环面超曲面，可以进一步扩大 Kähler 模空间。

最后讨论的主题是：改变 $V$ 的复结构时，Kähler 锥如何变化。\hyperref[sec:physics]{第 1.1 节}的简单图景似乎表明，可以保持 Kähler 结构不变，同时改变复结构。对于三维簇，这在一般情形下成立。更确切地说，Wilson \SourceCite{Wilson2} 证明：Calabi--Yau 三维簇 $V$ 的复结构处于一般位置时，Kähler 锥在复结构的小形变下保持不变。这种一般位置的刚性很强。例如，\SourceCite{Wilson3} 证明：若两个 Calabi--Yau 三维簇 $V_1,V_2$ 在各自的复结构模空间中处于一般位置，且互为辛形变，那么它们的 Kähler 锥相同；\SourceCite{Wilson4} 则表明，在三次锥 $W^*$ 以外，与 $\overline{K(V_1)}$ 的某个余维一的面对应的本原收缩的类型（I、II 或 III 型），与 $\overline{K(V_2)}$ 的相应面对应的本原收缩的类型相同。

然而，实际情形更为复杂，\hyperref[sec:physics]{第 1.1 节}图下的脚注已暗示了这一点。特别是，有些 $V$ 的形变会使 Kähler 锥变化。设 $\phi$ 为 III 型收缩，将直纹曲面 $S$ 映到亏格为 $g$ 的曲线 $C$。当 $g\geq2$ 时，前面可以用反射扩展可动锥。不过，底曲线满足 $g=1$ 时，上面的论证表明：在复结构的一般形变中，没有任何纤维随之形变。这意味着，对含有曲面 $S$（其类为 $E$）的特殊复结构，有效曲线锥严格包含一般复结构的有效曲线锥，因为一般复结构中不存在以类 $f$ 表示的有效曲线。对偶地，这一特殊复结构的 Kähler 锥变小。因此，Kähler 模空间可能依赖于复结构。

\subsection{环面超曲面的 Kähler 模}\label{subsec:enlarged-kahler-moduli}
设 $\Delta$ 为自反多胞形，$\Sigma$ 为 $\Delta$ 的法扇的一个极大射影细分。若 $V\subset X_\Sigma$ 为一般 Calabi--Yau 超曲面，回顾\hyperref[chap:mirror-constructions]{第~\ref*{chap:mirror-constructions} 章}对限制映射 $H^2(X_\Sigma)\to H^2(V)$ 的描述。它的核由除子类 $[D_\rho]$ 生成；这些除子对应于由 $\Delta^\circ$ 某个刻面的内部本原格点张成的一维锥 $\rho\in\Sigma(1)$。限制映射的像按定义为 $H_{\mathrm{toric}}^2(V)$。目标是为 $V$ 的“环面
\SourcePage{158}
部分”构造 Kähler 模空间。从某种意义上说，这一构造与\hyperref[subsec:toric-complex-moduli]{第 6.1.2 小节}中 $\overline{\mathcal M}_{\mathrm{simp}}$ 的构造互为“镜像”。

与\hyperref[subsec:toric-complex-moduli]{第 6.1.2 小节}一样，稍微改变问题，排除 $\Delta^\circ$ 的刻面内部格点。相应地，稍微修改极大射影细分的定义，以避开这些点。
\begin{definition}\label{def:simplified-projective-subdivision}\index{简化射影细分|hyperpage}
给定自反多胞形 $\Delta\subset M_{\mathbb R}$，若 $N_{\mathbb R}$ 中的扇 $\Sigma$ 具有以下性质，则称为\emph{简化射影细分}：
\begin{itemize}
\item $\Sigma$ 细分 $\Delta$ 的法扇。
\item $\Sigma(1)=(\Delta^\circ\cap N)_0-\{0\}$。
\item $X_\Sigma$ 射影且单纯。
\end{itemize}
\end{definition}
现在可以仿照\hyperref[sec:batyrev-mirror]{第 4.1 节}的 Batyrev 构造，将极大射影细分 $\Sigma$ 换成一个简化射影细分，仍记为 $\Sigma$。下面证明这样得到的 Calabi--Yau 簇，与 Batyrev 构造给出的为同一类。首先，给定极大射影细分 $\Sigma$，去掉由 $\Delta^\circ$ 的刻面内部格点生成的所有一维锥，就得到 $X_\Sigma$ 的一个爆缩的扇。所得扇未必单纯，所以可能需要在不加入新一维锥的情况下进一步细分，以得到简化射影细分 $\Sigma'$。得到交换图
\begin{equation}\label{eq:simplified-subdivision-diagram}
\begin{array}{ccc}
X_\Sigma&\overset{\phi}{\dashrightarrow}&X_{\Sigma'}\\
\scriptstyle\pi\searrow&&\swarrow\scriptstyle\pi'\\
&\mathbb P_\Delta&
\end{array}
\end{equation}
按构造，除 $\mathbb P_\Delta$ 中与 $\Delta^\circ$ 的那些含内部格点的刻面对应的点的上方以外，$\phi$ 为同构。一般 Calabi--Yau 超曲面 $\overline V\in|-K_{\mathbb P_\Delta}|$ 与这一有限集不交，因此它在 $\pi,\pi'$ 下的严格变换——即 $X_\Sigma,X_{\Sigma'}$ 中的一般 Calabi--Yau 超曲面——同构，正如所断言的。反过来更简单：任何简化射影细分都可细分为极大射影细分，仍得到同构的 Calabi--Yau 超曲面。

现在固定 $\Delta$ 的法扇的一个简化射影细分 $\Sigma$，考虑 $V\subset X_\Sigma$。由于排除了 $\Delta^\circ$ 的刻面内部格点，自然映射 $H^2(X_\Sigma)\to H^2(V)$ 为单射。其像已知为 $H_{\mathrm{toric}}^2(V)$，所以
\[
H^2(X_\Sigma)\simeq H_{\mathrm{toric}}^2(V).
\]
\hyperref[subsec:toric-kahler-cone]{第 3.3.3 小节}表明，$X_\Sigma$ 的 Kähler 锥的闭包为 $\overline{K(X_\Sigma)}=\operatorname{cpl}(\Sigma)$。因此，可以将 $\operatorname{cpl}(\Sigma)$ 视为位于 $H_{\mathrm{toric}}^2(V)$ 中。将 $V$ 的\emph{Kähler 锥的环面部分}定义为
\[
K(V)_{\mathrm{toric}}=K(V)\cap H_{\mathrm{toric}}^2(V),
\]
则有显然的包含关系 $\operatorname{cpl}(\Sigma)\subset\overline{K(V)}_{\mathrm{toric}}$，因为 $X_\Sigma$ 的丰沛类限制到 $V$ 后仍为丰沛类。

暂且假定总有 $\operatorname{cpl}(\Sigma)=\overline{K(V)}_{\mathrm{toric}}$。此时，很容易构造一个扩大的环面 Kähler 模空间，并用\hyperref[sec:gkz-decomposition]{第 3.4 节}的次级扇将它紧化。置 $\Xi=(\Delta^\circ\cap N)_0-\{0\}$，次级扇自然位于向量空间
$A(\Xi)\simeq H^2(X_\Sigma,\mathbb R)\simeq H_{\mathrm{toric}}^2(V,\mathbb R)$ 中（沿用\hyperref[sec:gkz-decomposition]{第~\ref*{sec:gkz-decomposition} 节}的记号）。次级扇及整上同调给出的格确定一个完备环面簇，暂且将它
\SourcePage{159}
记为 $\overline{\mathcal{KM}}_{\mathrm{toric}}$。对于每个简化射影细分 $\Sigma$，$\operatorname{cpl}(\Sigma)$ 是次级扇的一个极大锥；若 $\operatorname{cpl}(\Sigma)=\overline{K(V)}_{\mathrm{toric}}$，则 $\overline{\mathcal{KM}}_{\mathrm{toric}}$ 中相应的点就是\hyperref[subsec:large-radius-limits]{第 6.2.1 小节}意义下的大半径极限点（可能需要细分 $\operatorname{cpl}(\Sigma)$，以得到光滑空间）。

这幅图景的问题在于：包含关系 $\operatorname{cpl}(\Sigma)\subset\overline{K(V)}_{\mathrm{toric}}$ 中，\emph{不一定取等号}。环面 Kähler 锥的闭包 $\overline{K(V)}_{\mathrm{toric}}$ 可能严格大于 $\operatorname{cpl}(\Sigma)$。原因是，可能存在与 $\Sigma$ 不同的简化射影细分 $\Sigma'$，给出同构的 Calabi--Yau 超曲面。因此，$\overline{K(V)}_{\mathrm{toric}}$ 可能同时包含 $\operatorname{cpl}(\Sigma)$ 和 $\operatorname{cpl}(\Sigma')$。本小节后面将看到一个例子。

事实上，目前还不知道环面 Kähler 锥 $\overline{K(V)}_{\mathrm{toric}}$ 的精确结构。猜想它为有理多面锥；在所有已计算出 $\overline{K(V)}_{\mathrm{toric}}$ 的情形中，确实如此。下面的猜想~\ref{conj:toric-kahler-cone} 将仔细描述我们认为一般情形下 $\overline{K(V)}_{\mathrm{toric}}$ 应有的形态。猜想的实质是：$\overline{K(V)}_{\mathrm{toric}}$ 是某些锥 $\operatorname{cpl}(\Sigma')$ 的并集。若猜想正确，就意味着由次级扇构造的空间 $\overline{\mathcal{KM}}_{\mathrm{toric}}$，是扩大环面 Kähler 模空间的“真正”紧化的一个爆破。

本小节余下的计划是：先考察 $\overline{K(V)}_{\mathrm{toric}}$ 的结构，再构造一个扩大 Kähler 模空间，并将它紧化。

开始研究 $\overline{K(V)}_{\mathrm{toric}}$ 时，需要考察：$\Sigma$ 遍历所有简化射影细分时，各锥 $\operatorname{cpl}(\Sigma)$ 如何拼接。例如，如何从一个细分转到“邻近”的另一个细分？为此使用\emph{线性回路}的语言，现在给出定义。
\begin{definition}\label{def:linear-circuit}
给定 $\Xi=(\Delta^\circ\cap N)_0-\{0\}$，\emph{线性回路}是线性相关的子集 $\Xi'\subset\Xi$，使 $\Xi'$ 的任何真子集都线性无关。
\end{definition}
\index{线性回路|hyperpage}
更多细节可参阅 \SourceCite{GKZ2} 及其参考文献。要注意，\SourceCite{GKZ2} 的情形稍有不同：那里研究\emph{仿射回路}，使用仿射相关与三角剖分，而不是线性相关与扇。

下面用 $v_i$ 表示 $\Xi$ 的元素。给定线性回路 $\Xi'\subset\Xi$，固定其元素 $v_i\in\Xi'$ 之间的非平凡线性关系 $\sum_i b_i v_i=0$，并分解为 $\Xi'=\Xi'_+\cup\Xi'_-$；其中 $\Xi'_+$ 为在线性关系中具有正系数 $b_i$ 的元素 $v_i\in\Xi'$ 构成的集合，$\Xi'_-$ 类似定义。

对任何子集 $S\subset\Xi$，以 $C(S)$ 表示 $S$ 生成的凸锥。给定满足 $|\Xi'_+|\geq1$ 的回路 $\Xi'\subset\Xi$，将 $\Sigma_+(\Xi')$ 定义为极大维锥由 $C(\Xi'-\{v_i\})$（$v_i\in\Xi'_+$）组成的扇，其支集为 $C(\Xi')$。类似地，若 $|\Xi'_-|\geq1$，就得到扇 $\Sigma_-(\Xi')$。举例说明，设 $\Xi'=\{v_1,v_2,v_3,v_4\}$，线性关系为 $v_1+v_2=v_3+v_4$。此时，$C(\Xi')$ 是有四个生成元的锥。将它细分为一对单纯锥有两种办法：$\Sigma_+(\Xi')$ 给出一种，$\Sigma_-(\Xi')$ 给出另一种。

接着描述线性回路与扇 $\Sigma$ 如何相互作用。

\SourcePage{160}
\begin{definition}\label{def:circuit-supported-fan}\index{线性回路!扇的支撑|hyperpage}
设 $\Sigma$ 为 $N_{\mathbb R}$ 中的射影单纯扇，$\Sigma(1)\subset\Xi$；令 $\Xi'\subset\Xi$ 为线性回路。若满足下列条件，则称 $\Sigma$ \emph{支撑于} $\Xi'$（或称 $\Xi'$ \emph{由} $\Sigma$ \emph{支撑}）：
\begin{enumerate}[label=(\roman*)]
\item $\Sigma_+(\Xi')$ 是 $\Sigma$ 的子扇。
\item 设 $\sigma$ 为 $\Sigma_+(\Xi')$ 的一个极大维锥。若存在子集 $\Pi\subset\Xi-\Xi'$，使 $\sigma\cup\Pi$ 生成 $\Sigma$ 的一个极大维锥，则对 $\Sigma_+(\Xi')$ 的其他所有极大维锥 $\sigma'$，$\sigma'\cup\Pi$ 也生成 $\Sigma$ 的一个锥。
\end{enumerate}
\end{definition}
因此，支撑于 $\Xi'$ 意味着：$\Sigma$ 中与 $C(\Xi')$ 相接的锥，由 $\Sigma_+(\Xi')$ 以一种特别强的方式确定。本小节后面重新讨论\hyperref[ex:degree-four-moduli]{例题~\ref*{ex:degree-four-moduli}} 时，将给出一个支撑于回路的射影细分的例子。

定义~\ref{def:circuit-supported-fan} 使我们能够用线性回路修改扇 $\Sigma$。更确切地说，假定 $\Sigma$ 支撑于 $\Xi'$，且 $|\Xi'_-|\geq1$。将 $\Sigma$ 中由 $\sigma\cup\Pi$ 张成的单纯锥——其中 $\sigma$ 是 $\Sigma_+(\Xi')$ 的锥，$\Pi\subset\Xi-\Xi'$——换成由 $\sigma'\cup\Pi$ 张成的单纯锥，其中 $\sigma'$ 是 $\Sigma_-(\Xi')$ 的锥，就得到新的单纯扇 $\Sigma'=\operatorname{flip}_{\Xi'}(\Sigma)$。例如，设 $\Sigma$ 支撑于线性回路 $\Xi'=\{v_1,v_2,v_3,v_4\}$，且 $v_1+v_2=v_3+v_4$（即上述例子）。此时，$\Sigma$ 使用将 $C(\Xi')$ 分为两个单纯锥的一种方式，翻转所得的 $\Sigma'$ 则使用另一种方式。
\index{翻转|hyperpage}

从满足 $\Sigma(1)\subset\Xi$ 的单纯射影扇 $\Sigma$ 出发，翻转 $\Sigma'$ 显然仍单纯，并满足 $\Sigma'(1)\subset\Xi$。在所考虑的相邻射影腔之间的翻转中，利用 \SourceCite{OdP} 的技巧还能证明 $\Sigma'$ 也射影。因此，$\operatorname{cpl}(\Sigma')$ 是 GKZ 分解中的一个锥。不过，即使 $\Sigma$ 细分 $\Delta$ 的法扇（如简化射影细分的定义所要求的），$\Sigma'$ 也未必如此。但若 $\Xi'$ 位于 $\Delta^\circ$ 的一个真面中，则 $\Sigma'$ 也细分 $\Delta$ 的法扇。

接着考察翻转 $\Sigma'=\operatorname{flip}_{\Xi'}(\Sigma)$ 的几何意义。设 $\Sigma$ 为单纯射影扇，$\Sigma(1)\subset\Xi$，且 $|\Xi'_+|,|\Xi'_-|\geq2$。由定义立即可知，$X_{\Sigma'}$ 是从 $X_\Sigma$ 爆缩子簇 $\bigcap_{v_i\in\Xi'_-}D_i$，再爆破所得奇点，以 $X_{\Sigma'}$ 的子簇 $\bigcap_{v_i\in\Xi'_+}D_i$ 代替它而得到的。将 $X_\Sigma$ 到 $X_{\Sigma'}$ 的这一双有理变换称为\emph{广义翻越}，记为 $\operatorname{flop}_{\Xi'}$。例如，余维二的经典翻越可看作与回路 $\{v_1,v_2,v_3,v_4\}$、关系 $v_1+v_3=v_2+v_4$ 对应的广义翻越。

现在令 $\Xi'$ 为由 $\Sigma$ 支撑的线性回路，满足 $|\Xi'_+|,|\Xi'_-|\geq2$。再假定 $\operatorname{flop}_{\Xi'}$ 的基点集与 $V$ 不交，即对一般 Calabi--Yau 超曲面 $V\subset X_\Sigma$，有
\[
\bigcap_{v_i\in\Xi'_-}D_i\cap V=\varnothing.
\]
立即可知，$\operatorname{flop}_{\Xi'}$ 诱导 Calabi--Yau 超曲面之间的同构 $V\simeq V'$。将这种翻越称为\emph{\IndexLocation{indexloc:0335:160}平凡翻越}，相应的翻转 $\Sigma'=\operatorname{flip}_{\Xi'}(\Sigma)$ 称为 $\Sigma$ 的\emph{平凡翻转}。

当 $\Sigma'$ 为 $\Sigma$ 的平凡翻转时，得到映射
\[
H^2(X_{\Sigma'})\longrightarrow H^2(V')\simeq H^2(V),
\]
它把 Kähler 类映到 Kähler 类。因此，$\operatorname{cpl}(\Sigma')\subset\overline{K(V)}_{\mathrm{toric}}$。由此引出如下定义。给定 Calabi--Yau 环面
\SourcePage{161}
超曲面 $V\subset X_\Sigma$，考虑锥
\begin{equation}\label{eq:trivial-flip-kahler-union}\index{平凡翻转|hyperpage}
\widetilde{K(V)}_{\mathrm{toric}}=\bigcup_{\Sigma'}\operatorname{cpl}(\Sigma'),
\end{equation}
其中并集遍及所有可由 $\Sigma$ 经一系列平凡翻转得到的扇 $\Sigma'$。由上述讨论立即可知
\[
\widetilde{K(V)}_{\mathrm{toric}}\subset\overline{K(V)}_{\mathrm{toric}}.
\]
\begin{conjecture}\label{conj:toric-kahler-cone}\index{环面 Kähler 锥|hyperpage}
$\widetilde{K(V)}_{\mathrm{toric}}=\overline{K(V)}_{\mathrm{toric}}$。
\end{conjecture}
据我们所知，在所有已知例子中，$\widetilde{K(V)}_{\mathrm{toric}}$ 都与 $\overline{K(V)}_{\mathrm{toric}}$ 相等。有趣的是，在这些例子中，$\widetilde{K(V)}_{\mathrm{toric}}$ 总是单纯锥，尽管 $\overline{K(X_\Sigma)}$ 未必如此。下面重新讨论\hyperref[ex:degree-four-moduli]{例题~\ref*{ex:degree-four-moduli}} 时，将说明猜想~\ref{conj:toric-kahler-cone}，以及一些可能有助于最终证明它的思想。然而，\SourceCite{Szendroi2} 表明，猜想~\ref{conj:toric-kahler-cone} 一般情形下不成立。

\emph{本小节余下部分都假定猜想~\ref{conj:toric-kahler-cone} 成立。}

来探究这对 Kähler 模的环面部分意味着什么。将尚未扩大的 Kähler 模空间的环面部分定义为商
\begin{equation}\label{eq:toric-kahler-moduli}\index{环面 Kähler 模空间|hyperpage}
\mathcal{KM}_{\mathrm{toric}}
=\frac{\{\omega\in H_{\mathrm{toric}}^2(V,\mathbb C):
\operatorname{Im}(\omega)\in K(V)_{\mathrm{toric}}\}}
{\operatorname{im}H_{\mathrm{toric}}^2(V,\mathbb Z)}.
\end{equation}
严格地说，还应对 $V$ 的保持环面部分的自同构取商。不过，这些自同构保持有理多面锥 $\overline{K(V)}_{\mathrm{toric}}$，每个自同构都置换其本原射线生成元。因此，它们通过一个有限自同构群作用，真正的环面 Kähler 模空间就是式~\eqref{eq:toric-kahler-moduli} 的有限商。这类似于\hyperref[subsec:toric-complex-moduli]{第 6.1.2 小节}中 $\mathcal M_{\mathrm{poly}}$ 是 $\mathcal M_{\mathrm{simp}}$ 的有限商的情形。正如使用 $\mathcal M_{\mathrm{simp}}$ 更方便，这里考虑按上述方式定义的 $\mathcal{KM}_{\mathrm{toric}}$ 也更简单。

有了这一环面 Kähler 模空间的定义，可以看到，式~\eqref{eq:trivial-flip-kahler-union} 中的锥 $\operatorname{cpl}(\Sigma')$ 给出一个细分 $\overline{K(V)}_{\mathrm{toric}}$ 的扇。与\hyperref[subsec:large-radius-limits]{第 6.2.1 小节}一样，所得环面簇给出 $\mathcal{KM}_{\mathrm{toric}}$ 的一个部分紧化；若进一步细分以得到光滑部分紧化，就会得到许多大半径极限点。其中包括来自 $\operatorname{cpl}(\Sigma)$ 的点，但也可能有更多点，因为 $\overline{K(V)}_{\mathrm{toric}}$ 中可能出现其他锥 $\operatorname{cpl}(\Sigma')$。所有这些点都应对应于镜像一侧的最大幺幂边界点。

若要寻找 $\mathcal{KM}_{\mathrm{toric}}$ 的\emph{极小}部分紧化，显然应使用 $\overline{K(V)}_{\mathrm{toric}}$ 给出的仿射环面簇。在镜像一侧，本章后面将看到：需要将 $\overline{\mathcal M}_{\mathrm{simp}}$ 略作爆缩，才能得到简化复结构模空间 $\mathcal M_{\mathrm{simp}}$ 的相应极小紧化；\hyperref[ex:degree-four-moduli]{例题~\ref*{ex:degree-four-moduli}} 已暗示了这一点。

下面讨论线性回路的一些其他性质，它们在计算例子时很有用。首先观察到，线性回路给出对偶空间 $A(\Xi)^*$ 中有趣的元素。由于 $A(\Xi)=\mathbb R^\Xi/M_{\mathbb R}$，$A(\Xi)^*$ 是线性关系的向量空间
\[
A(\Xi)^*=\left\{\ell=(\ell_i)\in\mathbb R^\Xi:\sum_i\ell_i v_i=0\right\}.
\]
若 $\Xi'$ 为线性回路，有关系 $\sum_i b_i v_i=0$，由此得到分解 $\Xi'=\Xi'_+\cup\Xi'_-$。该关系还给出 $\ell_{\Xi'}=(b_i)\in A(\Xi)^*$；
\SourcePage{162}
固定分解 $\Xi'=\Xi'_+\cup\Xi'_-$ 后，它确定到一个正标量倍。将元素 $D\in A(\Xi)$ 视为除子类 $D=[\sum_i a_iD_i]$，则
\[
\ell_{\Xi'}(D)=\sum_i a_i b_i.
\]
现在可用 $\Sigma$ 所支撑的线性回路刻画 $\operatorname{cpl}(\Sigma)$。若 $\Sigma$ 支撑于 $\Xi'$，且 $\Sigma'=\operatorname{flip}_{\Xi'}(\Sigma)$，很容易证明：对 $D\in\operatorname{cpl}(\Sigma)$，有 $\ell_{\Xi'}(D)\geq0$；对 $D'\in\operatorname{cpl}(\Sigma')$，有 $\ell_{\Xi'}(D')\leq0$。因此，超平面 $\ell_{\Xi'}=0$ 将两个锥的内部分开。此外，可以证明
\[
\operatorname{cpl}(\Sigma)=\{D\in A(\Xi):\ell_{\Xi'}(D)\geq0
\text{，对所有由 }\Sigma\text{ 支撑的 }\Xi'\}.
\]
事实上，$\operatorname{cpl}(\Sigma)$ 的刻面与由 $\Sigma$ 支撑的线性回路 $\Xi'$ 一一对应。

还可用线性回路缩小 GKZ 分解。这很有用，因为我们真正感兴趣的，并不是那些 $\Sigma(1)$ 为 $\Xi=(\Delta^\circ\cap N)_0-\{0\}$ 的真子集所对应的锥 $\operatorname{cpl}(\Sigma)$。思想如下。令 $\Sigma$ 为简化射影细分，设 $\Xi'\subset\Xi$ 为任意满足 $|\Xi'_-|=1$ 的线性回路，不必由 $\Sigma$ 支撑。我们断言，$\operatorname{cpl}(\Sigma)$ 上有 $\ell_{\Xi'}\geq0$。为此，重标 $\Xi'$，使线性关系为 $b_0v_0=\sum_{i>0}b_i v_i$，所有系数都为正。由于 $v_0\in\Xi=\Sigma(1)$，某个极大锥 $\sigma\in\Sigma$ 含有 $v_0$。取 $D=[\sum_i a_iD_i]\in\operatorname{cpl}(\Sigma)$。因为 $\sigma$ 单纯，可以找到 $m_\sigma\in M_{\mathbb R}$，使对所有 $v_i\in\sigma$ 都有 $\langle m_\sigma,v_i\rangle=-a_i$；还注意到，对任意 $v_i\in\Xi$，都有 $\langle m_\sigma,v_i\rangle\geq-a_i$，因为 $D\in\operatorname{cpl}(\Sigma)$（见\hyperref[subsec:toric-kahler-cone]{第 3.3.1 小节}）。于是
\[
-b_0a_0=\langle m_\sigma,b_0v_0\rangle
=\left\langle m_\sigma,\sum_{i>0}b_i v_i\right\rangle
\geq-\sum_{i>0}a_i b_i,
\]
立即得到 $\ell_{\Xi'}(D)\geq0$。为使用这一结论，回顾锥
\[
A^+(\Xi)=\left\{\left[\sum_i a_iD_i\right]\in A(\Xi):a_i\geq0\text{，对所有 }i\right\}
\]
是 GKZ 分解的支集。将 \emph{Calabi--Yau 锥}\footnote{B. Sturmfels 向我们指出了这个锥的重要性。}定义为
\[
C_{\mathrm{CY}}=\{D\in A^+(\Xi):\ell_{\Xi'}(D)\geq0
\text{，对所有满足 }|\Xi'_-|=1\text{ 的 }\Xi'\}.
\]
关键在于，若 $\Sigma$ 为简化射影细分，由前一段知 $\operatorname{cpl}(\Sigma)\subset C_{\mathrm{CY}}$。所以，只需研究 GKZ 分解在锥 $C_{\mathrm{CY}}$ 上的限制。

下面用一个前面见过的例子，说明迄今建立的概念。
\begingroup
\SourceExampleNumber{6.1.4}{2}
\renewcommand{\theHexam}{6.1.4.revisited.\arabic{exam}}
\begin{example}[重访]\label{ex:revisited-toric-kahler}
\hyperref[ex:degree-four-moduli]{例题~\ref*{ex:degree-four-moduli}} 研究了由自反多胞形 $\Delta$ 得到的 Calabi--Yau 三维簇 $V\subset X_\Sigma$ 的复结构模；$\Delta$ 的顶点 $v_1,\ldots,v_6$ 为
\begin{equation}\label{eq:revisited-polytope}
\begin{gathered}
(-1,-2,-3,-7),\ (1,0,0,0),\ (0,1,0,0),\\
(0,0,1,0),\ (0,0,0,1),\ (0,0,-1,-2).
\end{gathered}
\end{equation}
现在要研究镜像 $V^\circ\subset X_{\Sigma^\circ}$ 的 Kähler 模。因为处理的是 $V^\circ$，需要交换 $N,M$、$\Delta,\Delta^\circ$ 等。

首先观察到，$K(V^\circ)=K(V^\circ)_{\mathrm{toric}}$。这由\hyperref[sec:batyrev-mirror]{第 4.1 节}的公式，特别是式~\eqref{eq:toric-hodge-subspace} 和式~\eqref{eq:batyrev-h11} 推出，因为它们表明 $h^2(V^\circ)=h_{\mathrm{toric}}^2(V^\circ)=2$。

\SourcePage{163}
$\Delta^\circ$ 的法扇的简化射影细分 $\Sigma^\circ$ 位于 $M_{\mathbb R}$，满足 $\Sigma^\circ(1)=(\Delta\cap M)_0-\{0\}$。这恰好是\hyperref[ex:degree-four-moduli]{例题~\ref*{ex:degree-four-moduli}} 的集合 $\Xi$；注意 $\Xi=\{v_1,v_2,v_3,v_4,v_5,v_6\}$。锥 $\operatorname{cpl}(\Sigma^\circ)$ 是 $\Xi$ 的 GKZ 分解的一部分。前面用 $\Xi^+=(\Xi\cup\{0\})\times\{1\}=(\Delta\cap M)_0\times\{1\}$ 构造了次级扇~\eqref{eq:second-moduli-secondary-fan}。改用 $\Xi$ 重复该构造，就给出 GKZ 分解。回顾 $A(\Xi)^*$ 是 $\Xi$ 各元素之间的线性关系向量空间。在本例中，以下矩阵的各行给出 $A(\Xi)^*$ 的一组基：
\[
\begin{pmatrix}-2&-2&-4&1&0&7\\1&1&2&0&1&-3\end{pmatrix}.
\]
使用 $A(\Xi)$ 的对偶基，该矩阵的各列给出向量 $e_{v_1}^*,\ldots,e_{v_6}^*\in A(\Xi)$。\hyperref[sec:gkz-decomposition]{第 3.4 节}于是表明，GKZ 分解如下图所示：
\begin{figure}[H]
\centering
\refstepcounter{equation}\label{eq:revisited-gkz-fan}
\makebox[\textwidth]{\makebox[0pt][l]{(\theequation)}\hfill
\raisebox{-.5\height}{\includegraphics[width=0.43375\textwidth]{6-2-gkz-fan.png}}\hfill\mbox{}}
\end{figure}
每条边都标出张成它的点 $e_{v_i}^*$。GKZ 分解有三个极大锥，记为 $\sigma_1,\sigma_2,\sigma_3$。$\Xi$ 的线性回路为
\begin{equation}\label{eq:revisited-circuits}\index{线性回路!例题|hyperpage}
\begin{gathered}
\{v_1,v_2,v_3,v_5,v_6\},\quad\{v_1,v_2,v_3,v_4,v_6\},\\
\{v_4,v_5,v_6\},\quad\{v_1,v_2,v_3,v_4,v_5\},
\end{gathered}
\end{equation}
相应的线性关系为
\begin{equation}\label{eq:revisited-circuit-relations}
\begin{aligned}
v_1+v_2+2v_3+v_5-3v_6&=0,\\
-2v_1-2v_2-4v_3+v_4+7v_6&=0,\\
v_4+2v_5+v_6&=0,\\
v_1+v_2+2v_3+3v_4+7v_5&=0.
\end{aligned}
\end{equation}
观察式~\eqref{eq:revisited-gkz-fan} 最容易找出这些线性回路。每个回路都可以如下得到：沿图中一条边确定的直线，将平面分为两个半平面。$\Xi'$ 由不在该直线上的顶点组成；$\Xi'_+$ 的点位于直线一侧，$\Xi'_-$ 的点位于另一侧。

很容易看出，唯一满足 $|\Xi'_-|=1$ 的线性回路 $\Xi'$（使用式~\eqref{eq:revisited-circuit-relations} 的关系或其相反关系），是按所写符号取的第一个关系。因此，Calabi--Yau 锥 $C_{\mathrm{CY}}$ 是 $A^+(\Xi)=\sigma_1\cup\sigma_2\cup\sigma_3$ 中由单个不等式 $a_1+a_2+2a_3+a_5-3a_6\geq0$ 确定的锥；与通常一样，$A^+(\Xi)$ 的一般点表示为 $[\sum_i a_iD_i]$。得到 $C_{\mathrm{CY}}=\sigma_1\cup\sigma_2$。

接着构造对应于 $\sigma_1,\sigma_2$ 的扇。先从 $\sigma_2$ 开始。一种做法是在 $\sigma_2$ 中取一个除子类，用丰沛性的凸性条件确定扇（\SourceCite{KS} 已将这一过程自动化）。不过，更容易的办法是认识到 $V^\circ$ 来自\IndexLocation{indexloc:1432:163}加权射影空间 $\mathbb P(1,1,2,3,7)$（最初正是这样构造 $V^\circ$ 的）。式~\eqref{eq:revisited-polytope} 的六个顶点中，前五个确定一个单纯形；取该单纯形各个面上的锥，给出
\SourcePage{164}
$M_{\mathbb R}$ 中的一个扇。对应环面簇为 $\mathbb P(1,1,2,3,7)$，它不是 Fano 簇。不过，加入由 $(0,0,-1,-2)$ 生成的一维锥来爆破 $\mathbb P(1,1,2,3,7)$ 后，就得到单纯扇 $\Sigma^\circ$，其极大维锥为
\[
\begin{gathered}
C(v_2,v_3,v_4,v_5),\ C(v_1,v_3,v_4,v_5),\ C(v_1,v_2,v_4,v_5),\ C(v_1,v_2,v_3,v_4),\\
C(v_2,v_3,v_5,v_6),\ C(v_1,v_3,v_5,v_6),\ C(v_1,v_2,v_5,v_6),\ C(v_1,v_2,v_3,v_6).
\end{gathered}
\]
这个扇显然细分 $\Delta^\circ$ 的法扇。还可验证，$\Sigma^\circ$ 支撑于线性回路

\[
\Xi'_1=\{v_1,v_2,v_3,v_5,v_6\},\qquad\Xi'_2=\{v_1,v_2,v_3,v_4,v_6\}.
\]
（符号由式~\eqref{eq:revisited-circuit-relations} 确定）。不等式 $\ell_{\Xi'_i}(D)\geq0$（$i=1,2$）定义 $\sigma_2$，因此 $\operatorname{cpl}(\Sigma^\circ)=\sigma_2$。

此外，$\Xi'_2$ 满足 $|(\Xi'_2)_+|=2$、$|(\Xi'_2)_-|=3$，所以得到扇 $\Sigma^{\circ\prime}=\operatorname{flip}_{\Xi'_2}(\Sigma^\circ)$ 及双有理映射 $\operatorname{flop}_{\Xi'_2}$。利用式~\eqref{eq:revisited-gkz-fan}，很容易验证 $\operatorname{cpl}(\Sigma^{\circ\prime})=\sigma_1$。还注意到，$\Xi'_2$ 包含在 $\Delta$ 的一个刻面中，该刻面由 $M_{\mathbb R}$ 中方程
\[
\langle m,(-1,-1,-1,1)\rangle=-1
\]
确定的超平面支撑，对应于顶点 $(-1,-1,-1,1)\in\Delta^\circ$（见\hyperref[ex:degree-four-moduli]{例题~\ref*{ex:degree-four-moduli}}）。由前面的讨论，$\Sigma^\circ,\Sigma^{\circ\prime}$ 都是简化射影细分。

现在证明，$\Sigma^{\circ\prime}$ 是 $\Sigma^\circ$ 的平凡翻转。因为二者都是简化射影细分，所以处于式~\eqref{eq:simplified-subdivision-diagram} 的情形。此外，可验证
$\operatorname{flop}_{\Xi'_2}:X_{\Sigma^\circ}\dashrightarrow X_{\Sigma^{\circ\prime}}$
除 $\mathbb P_{\Delta^\circ}$ 中与顶点 $(-1,-1,-1,1)\in\Delta^\circ$ 对应的点的上方外，为同构。一般 $\overline{V^\circ}\subset\mathbb P_{\Delta^\circ}$ 避开这一点，因此一般 $V^\circ\subset X_{\Sigma^\circ}$ 避开 $\operatorname{flop}_{\Xi'_2}$ 的基点集。所以，$X_{\Sigma^\circ},X_{\Sigma^{\circ\prime}}$ 中的一般 Calabi--Yau 超曲面同构，$\Sigma^{\circ\prime}$ 是 $\Sigma^\circ$ 的平凡翻转。

这样就证明了 $\widetilde{K(V^\circ)}_{\mathrm{toric}}=\sigma_1\cup\sigma_2$。下一步是证明它等于 $\overline{K(V^\circ)}_{\mathrm{toric}}$，即本例中猜想~\ref{conj:toric-kahler-cone} 成立。取 $D\in\overline{K(V^\circ)}_{\mathrm{toric}}$。先将 $V^\circ$ 视为 $X_{\Sigma^\circ}$ 的子簇，用其齐次坐标在 $V^\circ$ 上施加方程 $x_1=x_6=0$，定义曲线 $C\subset V^\circ$。再将 $V^\circ$ 视为 $X_{\Sigma^{\circ\prime}}$ 的子簇，用后者的齐次坐标及方程 $x_1=x_2=0$ 定义曲线 $C'\subset V^\circ$。因为 $D\in\overline{K(V^\circ)}_{\mathrm{toric}}$，有 $C\cdot D\geq0$ 和 $C'\cdot D\geq0$。可以证明，这两个不等式恰好定义 $\sigma_1\cup\sigma_2$。事实上，第一个不等式等价于 $\ell_{\Xi'_1}(D)\geq0$，其中 $\Xi'_1$ 为上述由 $\Sigma^\circ$ 支撑的线性回路；第二个不等式等价于 $\ell_{\Xi'_3}(D)\geq0$，其中 $\Xi'_3$ 是式~\eqref{eq:revisited-circuits} 的第三个回路，符号取式~\eqref{eq:revisited-circuit-relations} 所写的。因此，$\overline{K(V^\circ)}_{\mathrm{toric}}\subset\sigma_1\cup\sigma_2$，随即得到等号。

一般情形下，证明猜想~\ref{conj:toric-kahler-cone} 的一种思路，是构造类似上述 $C,C'$ 的曲线，使它们给出的线性泛函，与定义 $\widetilde{K(V)}_{\mathrm{toric}}$ 所用各细分的适当回路给出的泛函 $\ell_{\Xi'}$ 一致。

现在已得到 Kähler 锥的闭包
$\overline{K(V^\circ)}=\overline{K(V^\circ)}_{\mathrm{toric}}=\sigma_1\cup\sigma_2$，接着讨论 $V^\circ$ 的 Kähler 模空间。由\hyperref[subsec:large-radius-limits]{第 6.2.1 小节}，$\mathcal{KM}$ 的一个\IndexLocation{indexloc:0893:164}部分紧化由 $\overline{K(V^\circ)}$ 的仿射环面簇给出，很容易看出它光滑。此外，\IndexLocation{indexloc:0698:164}大半径极限点附近的自然坐标由 $\overline{K(V^\circ)}$ 的生成元确定。所需生成元 $D,D'$ 为式~\eqref{eq:revisited-gkz-fan} 的向量 $D=(-2,1)=e_{v_1}^*$、$D'=(1,0)=e_{v_4}^*$，所以 $D=D_1$ 且
\SourcePage{165}
$D'=D_4$，这里 $D_i$ 表示与 $v_i$ 对应的环面除子在 $V^\circ$ 上的限制。可以计算交数
\[
D^3=0,\qquad D^2\cdot D'=1,\qquad D\cdot D'^2=3,\qquad D'^3=9,
\]
由此得到 $V^\circ$ 上同调中的关系 $D^3=0$、$D'^2-3D\cdot D'=0$。下面将用到这些关系。

下一步是与镜像 $V$ 的复结构模比较。\hyperref[ex:degree-four-moduli]{例题~\ref*{ex:degree-four-moduli}} 用次级扇~\eqref{eq:second-moduli-secondary-fan} 紧化了 $\mathcal M_{\mathrm{simp}}$。现在知道，需要将锥 $\sigma_1,\sigma_2$ 合并为一个锥，这对应于爆缩原来定义的 $\overline{\mathcal M}_{\mathrm{simp}}$。这样，Kähler 锥 $\overline{K(V^\circ)}$ 就对应于紧化后的复结构模空间中的一个边界点；可以证明它为最大幺幂。事实上，使用锥 $\overline{K(V^\circ)}$ 的生成元提供的自然坐标，就得到两个单值性变换，其对数为 $N_1,N_2$。计算得 $N_1^3=0$、$N_2^2-3N_1N_2=0$。\footnote{计算使用了 \SourceCite{CdFKM,CFKM} 的方法；在这一用途下，它们效率更高。}结合相应周期空间的维数与非零最高阶项，足以证明极小爆缩中的相应边界点为最大幺幂。注意，它们与 $D,D'$ 满足的关系 $D^3=0$、$D'^2-3D\cdot D'=0$ 相似。这并非偶然，因为镜像映射应当将单值性的对数，与同 $\overline{K(V^\circ)}$ 中相应锥的生成元作杯积对应起来。于是得到本例中镜像对称如何起作用的一个显式例子。
\end{example}
\endgroup

接着转向 Kähler 模空间的整体描述。目标是使它与\hyperref[subsec:toric-complex-moduli]{第 6.1.2 小节}中镜像流形的简化多项式复结构模空间的整体描述，直接按定义相同。

更正式地，设 $V\subset X_\Sigma$ 为 Calabi--Yau 环面超曲面，$\Sigma$ 为自反多胞形 $\Delta$ 的法扇的一个简化射影细分，并且如通常一样，$\Xi=(\Delta^\circ\cap N)_0-\{0\}$。将\emph{紧化环面 Kähler 模空间} $\overline{\mathcal{KM}}_{\mathrm{toric}}$ 定义为 $A(\Xi)$ 中以下扇对应的环面簇：从 $\Xi^+=(\Delta^\circ\cap N)_0\times\{1\}$ 的次级扇出发，凡一组锥 $\operatorname{cpl}(\Sigma')$ 可通过一系列\IndexLocation{indexloc:0331:165}平凡翻转联系起来，就用它们的并集代替这一组锥（如猜想~\ref{conj:toric-kahler-cone} 所述）。

因此，紧化环面模空间与镜像族的紧化复结构模空间 $\overline{\mathcal M}_{\mathrm{simp}}$ 经极小爆缩后的空间，本质上相同。注意，按定义，紧化 Kähler 模空间包含与 Kähler 锥对应的仿射环面簇。\hyperref[subsec:large-radius-limits]{第 6.2.1 小节}已经说明，这个环面簇自然包含 Kähler 模空间。如前所述，忽略了自同构群的作用；在始终假定猜想~\ref{conj:toric-kahler-cone} 成立的前提下，它将给出 $\overline{\mathcal{KM}}_{\mathrm{toric}}$ 的一个有限商。

$\overline{\mathcal{KM}}_{\mathrm{toric}}$ 的定义有一个很好的特点：它自然地扩大 Kähler 模空间，同时将它紧化。\hyperref[subsec:threefold-kahler-moduli]{第 6.2.2 小节}中的一般 Calabi--Yau 三维簇，必须先扩大到可动锥，再进一步扩大到\IndexLocation{indexloc:0914:165}反射可动锥；即便如此，也只得到\IndexLocation{indexloc:0893:165}部分紧化。因此，虽然描述环面 Kähler 模空间花了一些功夫，Calabi--Yau 环面超曲面仍然比一般 Calabi--Yau 三维簇容易处理得多，更不用说任意维数的 Calabi--Yau 簇了。

\SourcePage{166}
紧化环面 Kähler 模空间与\hyperref[subsec:toric-reduction]{第~\ref*{subsec:toric-reduction} 小节}中讨论的矩映射，为规范线性 $\sigma$ 模型（GLSM）的相结构提供了数学描述；\hyperref[sec:gauged-linear-sigma-models]{附录 B.5} 将概述这一模型。GLSM 依赖于一些参数，可将它们与 $H^2(X_\Sigma,\mathbb R)$ 的元素识别起来。该空间自然识别为矩映射的目标空间。与讨论辛约化时一样，在上同调类 $D\in A^+(\Xi)$ 上方，对矩映射的一个纤维取商，得到辛商。若 $D$ 位于 $A^+(\Xi)$ 的 GKZ 分解中 $\operatorname{cpl}(\Sigma)$ 的内部，则定理~\ref{thm:gkz-symplectic-reduction} 已证明该辛商就是 $X_\Sigma$。对一般 $D\in A(\Xi)$，辛商仍有定义，但未必为射影环面簇。由此，可能的辛商的同构类与次级扇中的极大维锥一一对应。

GLSM 的额外要素是从辛商中得到一个子簇 $V$。在这里，对应于次级扇不同部分的不同辛商，称为 GLSM 的\emph{相}。若 $D$ 位于某个简化射影细分对应的 $\operatorname{cpl}(\Sigma)$ 的内部，则子簇 $V$ 为 Calabi--Yau 超曲面。一般而言，$V$ 甚至未必为超曲面。\hyperref[sec:gauged-linear-sigma-models]{附录 B.5} 的例子有两个相：相应辛商为 $\mathbb P^4$ 和 $\mathbb C^5/\mathbb Z_5$，对应的子簇 $V$ 为 Fermat 五次三维簇和原点。后一情形是一个 \emph{Landau--Ginzburg 轨形}的例子，\hyperref[sec:landau-ginzburg-models]{附录 B.4} 也将讨论它。

环面 Kähler 模空间的定义带来的一个结果是：若 $V\subset X_\Sigma$ 和 $V^\circ\subset X_{\Sigma^\circ}$ 为 Batyrev 镜像，则 $\overline{\mathcal{KM}}_{\mathrm{toric}}(V)$ 与 $\overline{\mathcal M}_{\mathrm{simp}}(V^\circ)$ 的相应极小爆缩，是同一个空间。这可能使读者觉得镜像映射就是恒等映射；不过，如例题~\ref{ex:weighted-extremal-transition} 将展示的，镜像映射并没有这么简单。真正的镜像映射将在\hyperref[sec:mirror-map]{第~\ref*{sec:mirror-map} 节}详细讨论。

\hypertarget{source:166-batch4}{}
\subsection{Kähler 模空间的边界}\label{subsec:kahler-moduli-boundary}
本小节的讨论仍然在很大程度上具有猜测性质。我们先考察 Calabi--Yau 三维簇，再专门讨论作为环面超曲面出现的情形。出发点是以下观察：常常可以将一个已知的 Calabi--Yau 三维簇爆缩为奇异簇，再将它光滑化，得到一个新的 Calabi--Yau 三维簇。更准确地说，\SourceCite{Morrison9} 给出了以下定义。

\begin{definition}\label{def:extremal-transition}\index{极端跃迁|hyperpage}
设 $V$ 为 Calabi--Yau 三维簇，$\phi:V\to\overline V$ 为双有理收缩。若 $\overline V$ 可以光滑化为 Calabi--Yau 三维簇 $\widetilde V$，则从 $V$ 到 $\widetilde V$ 的这一过程称为\emph{极端跃迁}。
\end{definition}

若 $\phi$ 为本原双有理收缩，为了强调这一点，有时称它为\emph{本原极端跃迁}。回忆\hyperref[subsec:threefold-kahler-moduli]{第 6.2.2 小节}，$\phi$ 分为 I、II、III 三型，因此，例如可以说一个极端跃迁属于 I 型。

在物理学文献中，I 型极端跃迁称为\emph{锥形奇点跃迁}，因为物理学家常将三维结点称为锥形奇点。通常所说的“Reid 的幻想”中出现的跃迁，正是这种锥形奇点跃迁。在 \SourceCite{Reid5} 中，Reid 猜想：若将所有可能的 $V$ 与 $\widetilde V$ 的复模空间，沿所有可能的 $\overline V$ 的模空间粘合起来，所得空间或许是连通的。

下面是取自 \SourceCite{GMS} 的一个漂亮的锥形奇点跃迁例子。

\SourcePage{167}
\SourceExampleNumber{6.2.4}{1}
\begin{example}\label{ex:quintic-conifold-transition}
设 $\overline V\subset\mathbb P^4$ 为一个一般的奇异五次三维簇，包含 $\mathbb P^4$ 中由 $x_3=x_4=0$ 定义的平面。这样的 $\overline V$ 的方程为
$x_3g(x_0,\ldots,x_4)+x_4h(x_0,\ldots,x_4)=0$，其中 $g,h$ 为一般四次多项式。超曲面 $\overline V$ 在 $x_3=x_4=g(x_0,\ldots,x_4)=h(x_0,\ldots,x_4)=0$ 的 16 个点处有结点。考虑 $\mathbb P^4\times\mathbb P^1$ 中由
\[
y_0x_4-y_1x_3=0,\qquad
y_0g(x_0,\ldots,x_4)+y_1h(x_0,\ldots,x_4)=0
\]
定义的光滑 Calabi--Yau 完全交 $V$；这里 $y_0,y_1$ 为 $\mathbb P^1$ 的齐次坐标，$x_0,\ldots,x_4$ 仍为 $\mathbb P^4$ 的齐次坐标。容易看出，$V$ 的 Kähler 锥同构于 $\mathbb P^4\times\mathbb P^1$ 的 Kähler 锥，即由 $\mathbb P^4$ 与 $\mathbb P^1$ 的超平面类的拉回生成。来自 $\mathbb P^4$ 的超平面类 $H$ 张成 Kähler 锥的一个面，且 $H^3=5\ne0$，所以投影 $\mathbb P^4\times\mathbb P^1\to\mathbb P^4$ 限制为本原收缩 $\phi:V\to\overline V$。由于 $\phi$ 收缩的是 16 个结点上方的曲线，它属于 I 型。奇异三维簇 $\overline V$ 可光滑化为光滑五次三维簇 $\widetilde V\subset\mathbb P^4$，从而完成这一锥形奇点跃迁。
\end{example}

我们组织从 $V$ 到 $\widetilde V$ 的极端跃迁的讨论时，所要提示的是：同时粘合 $V$ 与 $\widetilde V$ 的复模空间和 Kähler 模空间，比只粘合复模空间更自然。换言之，我们是在完整的 SCFT 模空间中工作。下面将看到，识别 $V$ 与 $V^\circ$ 的 SCFT 模空间，会带来出人意料的丰富结构。

在 SCFT 模空间中，一个极端跃迁自然可看作两步过程：先固定复结构，使 Kähler 结构发生形变；再固定 Kähler 结构，使复结构发生形变。具体地，根据\hyperref[subsec:threefold-kahler-moduli]{第 6.2.2 小节}的讨论，$\overline V$ 的丰沛类空间经 $\phi$ 拉回，成为 $\overline{K(V)}$ 的一个面 $F$。我们先在 $\overline{\mathcal{KM}}$ 中选取一条趋近于其边界中与 $F$ 关联的分支的路径，同时保持复结构不变。然后，将这一边界分支与 $V$ 的复模空间的乘积，粘合到 $\widetilde V$ 的紧化 Kähler 模空间与 $\overline V$ 的复模空间的乘积上。注意，$\overline V$ 的复模空间是 $\widetilde V$ 的复模空间边界的一个子集。接下来在 $\widetilde V$ 的复模空间中自由移动，将 $\overline V$ 光滑化，即可完成跃迁。

\SourceCite{Morrison9} 猜想，这一跃迁的“镜像”是 $V^\circ$ 的 SCFT 模空间中沿反方向进行的极端跃迁。镜像对称应将 $V$ 的 SCFT 模空间中的路径，与 $V^\circ$ 的 SCFT 模空间中的路径识别起来。若 $\widetilde V^\circ$ 为 $\widetilde V$ 的镜像，那么描述这一跃迁的路径应当如下：先在 $V^\circ$ 的复模空间中移动，固定 Kähler 结构；其极限奇异复结构，则来自 $\widetilde V^\circ$ 的一个本原收缩，可以表示为 $\widetilde V^\circ$ 的 Kähler 模空间中的路径。猜想从 $V^\circ$ 到 $\widetilde V^\circ$ 的跃迁，是一个极端跃迁的逆过程。

正如 \SourceCite{BKK,Morrison9} 所观察到的，Batyrev 镜像构造提供了这一机制的一个漂亮例子。

\SourceExampleNumber{6.2.4}{2}
\begin{example}\label{ex:toric-extremal-transition}
设 $\Delta$ 与 $\widetilde\Delta$ 为两个不同的自反多胞形，且 $\Delta\subset\widetilde\Delta$。为简便起见，假定二者都为四维。我们先用这些数据构造一个从 $V$ 到 $\widetilde V$ 的极端跃迁。

\SourcePage{168}
由于 $\widetilde\Delta^\circ\subset\Delta^\circ$，$\widetilde\Delta$ 的法扇的任一细分 $\widetilde\Sigma$，都可进一步细化为 $\Delta$ 的法扇的一个细分 $\Sigma$。这一细分给出态射 $\pi:X_\Sigma\to X_{\widetilde\Sigma}$。Calabi--Yau 超曲面 $V\subset X_\Sigma$ 由方程
\[
\sum_{m_i\in\Delta\cap M}\lambda_i t^{m_i}=0
\]
定义。类似地，由方程
\begin{equation}\label{eq:extremal-transition-polynomial}
\sum_{m_i\in\widetilde\Delta\cap M}\lambda_i t^{m_i}=0
\end{equation}
得到 Calabi--Yau 超曲面 $\widetilde V\subset X_{\widetilde\Sigma}$。令 $\phi=\pi|_V$，即得到双有理收缩 $\phi:V\to\overline V$；其中 $\overline V$ 由式~\eqref{eq:extremal-transition-polynomial} 定义，只是当 $m_i\in(\widetilde\Delta-\Delta)\cap M$ 时，所有系数 $\lambda_i$ 都为零。让式~\eqref{eq:extremal-transition-polynomial} 中全部 $\lambda_i$ 任意取值，就将 $\overline V$ 光滑化为 $\widetilde V$，完成极端跃迁。从 $V$ 到 $\widetilde V$ 的这一过程称为\emph{环面极端跃迁}。

关键的观察在于：由于 $\widetilde\Delta^\circ\subset\Delta^\circ$，同样的构造给出镜像之间的环面极端跃迁，不过方向是从 $\widetilde V^\circ$ 到 $V^\circ$。这恰好符合 Morrison 猜想的预言。
\end{example}

还应提到，\hyperref[ex:toric-extremal-transition]{例题~\ref*{ex:toric-extremal-transition}}的方法，已经应用于\hyperref[sec:batyrev-mirror]{第 4.1 节}讨论的\IndexLocation{indexloc:0109:168}加权射影空间超曲面所对应的 7555 个\IndexLocation{indexloc:1218:168}自反多胞形。可以证明 \SourceCite{CGGK,ACJM}，用环面极端跃迁连接模空间的各个分支后，所得空间的这 7555 个不可约分支全部位于同一个连通分支中。

一般地，设有一个从 $V$ 到 $\widetilde V$ 的极端跃迁。对于 SCFT 模空间中与 $V$、$\widetilde V$ 关联的两个模空间分支交会之处的结构，目前还没有良好的数学定义。现有的解释只有物理学中的解释：在 I 型极端跃迁的情形下，利用 II 型弦理论和无质量黑洞 \SourceCite{GMS}。II 型弦理论的模空间大于 SCFT 模空间；极端跃迁将这些更大的模空间粘合起来，物理学给出的描述如下。被 $V\to\overline V$ 收缩的曲线对应于黑洞。当我们沿极端跃迁趋近 $V$ 的 Kähler 模空间边界时，这些黑洞变为无质量。另一方面，II 型弦理论中的上同调类也对应于无质量粒子，而 $V$ 与 $\widetilde V$ 的 Hodge 数之差，正由这些无质量黑洞说明。基于 \SourceCite{KMP} 所解释的物理，对某些 III 型极端收缩也能给出类似说明。第 \hyperref[subsec:conifold-transitions-conclusion]{\ref*{subsec:conifold-transitions-conclusion}} 小节将简要回到这些思想。（这里所说的 II 型弦理论，在英文中使用小写的“type”；它与极端收缩的 I、II、III 型，即英文使用大写“Type”的分类，没有直接关系。）

接下来讨论取自 \SourceCite{CdFKM,KMP} 的一个极端跃迁例子。它揭示了有关模空间细致而丰富的结构。令人意外的是，这个例子我们已经很熟悉了。

\SourceExampleNumber{6.2.4}{3}
\begin{example}\label{ex:weighted-extremal-transition}
继续讨论\hyperref[ex:weighted-pf-method]{例题~\ref*{ex:weighted-pf-method}}和\hyperref[ex:weighted-complex-moduli]{例题~\ref*{ex:weighted-complex-moduli}}开始研究的镜像族，它来自 $\mathbb P(1,1,2,2,2)$ 的环面爆破中的反典范超曲面。
\SourcePage{169}
有自反多胞形 $\Delta$，其极多胞形 $\Delta^\circ$ 满足
\[
\begin{aligned}
\Xi=(\Delta^\circ\cap N)_0-\{0\}=\{&(-1,-2,-2,-2),(1,0,0,0),(0,1,0,0),\\
&(0,0,1,0),(0,0,0,1),(0,-1,-1,-1)\}.
\end{aligned}
\]
将这些点记为 $m_1,\ldots,m_6$。\hyperref[ex:weighted-resolution]{例题~\ref*{ex:weighted-resolution}}中的 GKZ 分解表明，满足 $\Sigma(1)=\Xi$ 的简化射影细分 $\Sigma$ 唯一。它给出 Calabi--Yau 三维簇 $V\subset X_\Sigma$。我们将构造一个从 $V$ 出发的极端跃迁，并研究这一跃迁在 $V$ 的 Kähler 模空间中的表现。

利用\hyperref[sec:batyrev-mirror]{第 4.1 节}的公式，可以算出
\[
h^{1,1}(V)=h^{1,1}_{\mathrm{toric}}(V)=2,\qquad
h^{2,1}(V)=86,\quad h^{2,1}_{\mathrm{poly}}(V)=83.
\]
因此 Batyrev 镜像 $V^\circ$ 满足
\[
h^{2,1}(V^\circ)=h^{2,1}_{\mathrm{poly}}(V^\circ)=2,\qquad
h^{1,1}(V^\circ)=86,\quad h^{1,1}_{\mathrm{toric}}(V^\circ)=83.
\]
令 $D_1,\ldots,D_6$ 为与 $\Sigma(1)=\Xi$ 关联的环面除子。由于 $H^2(X_\Sigma)\simeq H^2(V)$，既可将它们视为 $X_\Sigma$ 上的类，也可视为 $V$ 上的类。$V$ 的 Kähler 锥由 $D_1,D_3$ 的类张成。\hyperref[ex:weighted-complex-moduli]{例题~\ref*{ex:weighted-complex-moduli}}中矩阵的前六列给出线性等价关系
$D_2\sim D_1$、$D_4\sim D_5\sim D_3$、$D_6\sim D_3-2D_1$，它们对后续计算很有用。特别地，由此得到 $V\in|\sum_iD_i|=|4D_3|$。

接着寻找 $V$ 的双有理收缩。利用\hyperref[subsec:threefold-kahler-moduli]{第 6.2.2 小节}的理论，研究 Kähler 锥的各个面。使用 $X_\Sigma$ 中与 $\Sigma(1)=\Xi$ 对应的齐次坐标 $x_1,\ldots,x_6$ 很方便。可验证，$|D_1|$ 是由截面 $x_1,x_2$ 定义的线性系，它们给出 K3 纤维化 $V\to\mathbb P^1$，且 $D_1^3=0$。因此 $|D_1|$ 不给出双有理收缩。另一方面，$|D_3|$ 由截面 $x_1^2x_6,x_1x_2x_6,x_2^2x_6,x_3,x_4,x_5$ 定义。由于 $D_3^3=8$，得到本原收缩 $\phi=\phi_{|D_3|}$，其像为 $\mathbb P^5$ 中次数为 8 的三维簇 $\overline V$。事实上，$\overline V$ 是一个二次超曲面与一个四次超曲面的完全交。为看清这一点，在 $\mathbb P^5$ 上引入坐标 $y_0,\ldots,y_5$。按上述顺序用 $|D_3|$ 的截面定义 $\phi$，可见 $\overline V$ 位于秩为 3 的奇异二次超曲面 $y_0y_2=y_1^2$ 上；而定义 $V\subset X_\Sigma$ 的 $|4D_3|$ 的截面，由 $y_i$ 的一个四次多项式诱导。可以验证，$\phi$ 为 III 型本原收缩，将例外曲面 $x_6=0$ 收缩到 $\overline V$ 中由 $y_0=y_1=y_2=0$ 截出的轨迹；它实际上是一条平面四次曲线。例外除子 $x_6=0$ 是这条亏格为 3 的曲线上的直纹曲面。将 $\overline V$ 光滑化为 $\mathbb P^5$ 中光滑的 $(2,4)$ 完全交 $\widetilde V$，便完成了本原极端跃迁。$\widetilde V$ 的有关 Hodge 数为 $h^{1,1}(\widetilde V)=1$、$h^{2,1}(\widetilde V)=89$。

在 Kähler 模的语言中，这一跃迁描述如下：由于 $\overline V$ 来自 $V$ 的 Kähler 锥的一个面，显然应将 $\overline{\mathcal{KM}}(\widetilde V)$ 沿这一面对应的子集，接到 $\overline{\mathcal{KM}}(V)$ 上。不过，这还没有交代全部情形，因为 Kähler 锥的一个面的实余维数为一，而 $\overline{\mathcal{KM}}(\widetilde V)\subset\overline{\mathcal{KM}}(V)$ 的复余维数为一。因此，还需要进一步确定粘接 $\overline{\mathcal{KM}}(\widetilde V)$ 的精确轨迹。下面将看到，镜像会告诉我们该怎样做。

\SourcePage{170}
至此，SCFT 模空间的图景已近乎完整。迄今我们只看了 $V$ 的复模空间中多项式的部分，但 $V$ 的 Hodge 数表明，存在非多项式形变。例外除子 $x_6=0$ 在 $V$ 的一般复结构形变下不会继续存在；不过，由于它是亏格为 3 的曲线上的直纹曲面，有 $2g-2=4$ 条纤维会随复结构的一般形变而形变。在这一邻近复结构上，形变后的 $|D_3|$ 仍会诱导本原收缩，但这次属于 I 型。此外，这 4 条孤立曲线可以同时翻越，得到不同的复结构和 Kähler 结构。可以证明，这一翻越诱导 SCFT 模空间上的一个对合。

于是得到 $V$ 的 SCFT 模空间的如下描述。对于 86 维复模空间中的任一复结构，都有一个二维 Kähler 模空间，其中含有与本原收缩 $\phi_{|D_3|}$ 对应的一条特殊曲线。另外，考虑刚才构造的 SCFT 模空间上的翻越对合：83 维的多项式子空间逐点不动，而它在 Kähler 模空间上的作用非平凡，是\hyperref[subsec:threefold-kahler-moduli]{第 6.2.2 小节}所述的反射。$\widetilde V$ 的 89 维复模空间中，包含定义它的二次式秩为 3 的 $(2,4)$ 完全交所构成的 83 维子空间，以及二次式秩为 4 的 $(2,4)$ 完全交所构成的 86 维子空间。它们分别与 $V$ 的多项式模空间及复模空间识别；$\widetilde V$ 的 Kähler 模空间则与 $V$ 的 Kähler 模空间中的一条曲线识别，下面会描述这条曲线。

接着考察从 $\widetilde V^\circ$ 到 $V^\circ$ 的“对偶”极端跃迁。\hyperref[ex:weighted-pf-method]{例题~\ref*{ex:weighted-pf-method}}和\hyperref[ex:weighted-complex-moduli]{例题~\ref*{ex:weighted-complex-moduli}}已描述 $V^\circ$；$\widetilde V^\circ$ 则可用\hyperref[sec:toric-complete-intersections]{第 4.3 节}讨论的 Batyrev--Borisov 构造得到。这样，根据 \SourceCite{CdFKM} 的计算，可以显式写出从 $\widetilde V^\circ$ 到 $V^\circ$ 的极端跃迁。它诱导模空间上的映射 $\overline{\mathcal M}_{\mathrm{simp}}(\widetilde V^\circ)\to\overline{\mathcal M}_{\mathrm{simp}}(V^\circ)$。利用 $\overline{\mathcal M}_{\mathrm{simp}}$ 上的坐标 $z_1,z_2$，可证明 \SourceCite{CdFKM,KMP}，$\overline{\mathcal M}_{\mathrm{simp}}(\widetilde V^\circ)$ 沿\hyperref[ex:weighted-complex-moduli]{例题~\ref*{ex:weighted-complex-moduli}}中由 $D_{A\cap\Gamma}(f_\Gamma)=1-4z_2=0$ 定义的判别式轨迹分支粘接。\hyperref[subsec:toric-mirror-map]{第~\ref*{subsec:toric-mirror-map} 小节}将简要说明这一论证。判别式轨迹的出现并不意外，因为从 $\widetilde V^\circ$ 到 $V^\circ$ 的过程，首先收缩到一个奇异簇。

由于从 $\widetilde V^\circ$ 到 $V^\circ$ 的跃迁，经镜像对称与从 $V$ 到 $\widetilde V$ 的跃迁相容，现在终于能完成对 $V$ 的 SCFT 模空间的描述。即用镜像映射将 $\overline{\mathcal M}_{\mathrm{simp}}(V^\circ)$ 与 $\overline{\mathcal{KM}}(V)$ 识别后，根据上一段，$\widetilde V$ 的 Kähler 模空间应与 $V$ 的 Kähler 模空间中由 $1-4z_2=0$ 定义的曲线识别起来。

$V^\circ$ 的 SCFT 模空间具有与 $V$ 相同的复杂几何。此外，也能看出翻越对合的作用。由定理~\ref{thm:batyrev-hodge} 的证明，$V^\circ$ 构造为环面簇 $X_{\Sigma^\circ}$ 中的超曲面；该环面簇含有 3 个环面除子 $D_i$，每个限制到 $V^\circ$ 后都有两个分支 $D_i',D_i''$。除子类 $D_i'+D_i''$ 属于 $H^{1,1}_{\mathrm{toric}}(V^\circ)$，但单独的 $D_i',D_i''$ 不属于这个子空间。交换 $D_i'$ 与 $D_i''$（$i=1,2,3$）的对合，对应于翻越对合在 $V$ 的复模空间上的作用。它在 $H^{1,1}_{\mathrm{toric}}(V^\circ)$ 上的作用平凡。
\end{example}

\SourcePage{171}
这个例子对环面情形的镜像映射有若干有趣的启示。通过次级扇的构造，Batyrev 镜像 $V,V^\circ$ 使 $\overline{\mathcal{KM}}_{\mathrm{toric}}(V)$ 与 $\overline{\mathcal M}_{\mathrm{simp}}(V^\circ)$ 成为同一个模空间。这似乎暗示镜像映射是恒等映射。然而，仔细考察\hyperref[ex:weighted-extremal-transition]{例题~\ref*{ex:weighted-extremal-transition}}就知道并非如此。基本思想是：要使镜像对称与极端跃迁相容，镜像映射就必须不同于恒等映射。

具体说明如下。回到\hyperref[ex:weighted-extremal-transition]{例题~\ref*{ex:weighted-extremal-transition}}，$V$ 的 Kähler 锥闭包为 $\{t_1D_3+t_2D_1:t_1,t_2\geq0\}$。注意，这里的 $t_1,t_2$ 是\hyperref[chap:toric]{第~\ref*{chap:toric} 章}式~\eqref{eq:weighted-gkz-picture} 中记为 $r_1,r_2$ 的坐标。因此，Kähler 锥的面由 $t_j=0$ 定义。转到 Kähler 模空间时，让 $t_j$ 位于上半平面，则 $q_j=e^{2\pi it_j}$ 是 $\mathcal{KM}(V)$ 的相应局部坐标。因此，$\operatorname{Im}(t_j)=0$ 对应于 $V$ 的 Kähler 模空间中的 $|q_j|=1$。注意，这个集合在 $\overline{\mathcal{KM}}(V)$ 中的实余维数为一。

镜像一侧有 $V^\circ$ 的简化复模，它由同一个空间给出。\hyperref[ex:weighted-complex-moduli]{例题~\ref*{ex:weighted-complex-moduli}}中的自然坐标 $z_1,z_2$ 来自矩阵~\eqref{eq:weighted-moduli-relations}。该矩阵的行给出 $A^*(\Xi)$ 的一组基，以及 $A(\Xi)$ 的对偶基；GKZ 分解就在后一个空间中。这正是上面描述 $V$ 的 Kähler 锥时用到的基 $D_3,D_1$。因此，将两边都与 Chow 商识别之后，$\overline{\mathcal M}_{\mathrm{simp}}(V^\circ)$ 的局部坐标 $z_1,z_2$ 就是 $\overline{\mathcal{KM}}(V)$ 的坐标 $q_1,q_2$；\hyperref[subsec:toric-mirror-map]{第~\ref*{subsec:toric-mirror-map} 小节}将仔细证明这一点。

为说明镜像映射为何不是 $q_j=z_j$，使用\hyperref[ex:weighted-extremal-transition]{例题~\ref*{ex:weighted-extremal-transition}}中从 $V$ 到 $\widetilde V$ 的极端跃迁。如上所见，由 $D_3$ 生成的 Kähler 锥面给出\IndexLocation{indexloc:1089:171}本原收缩 $\phi=\phi_{|D_3|}:V\to\overline V$。在 Kähler 模中，这个面由 $\operatorname{Im}(t_2)=0$ 定义，也就是 $|q_2|=1$。若镜像映射为恒等映射，则从 $\widetilde V^\circ$ 到 $V^\circ$ 的“对偶”极端跃迁，应在 $|z_2|=1$ 这一轨迹的某处粘合。然而，\hyperref[ex:weighted-extremal-transition]{例题~\ref*{ex:weighted-extremal-transition}}说明，这条轨迹由 $1-4z_2=0$，即 $z_2=1/4$ 定义。两者不一致，表明镜像映射不能是恒等映射。

物理学家会用 Kähler 模空间“受到瞬子修正”来解释这一差异。我们以前已见过这种现象，最突出的是：瞬子，即从 Riemann 曲面到 $V$ 的全纯映射，使 A 模型关联函数不同于交数。\hyperref[sec:mirror-map]{第~\ref*{sec:mirror-map} 节}将最终解释镜像映射。

还有一点很有趣：尽管 $|z_2|=1$ 与 $1-4z_2=0$ 不相容，它们至少涉及同一个变量。这并非偶然，而是 Morrison 关于极端跃迁中复模与 Kähler 模识别位置的一般猜想 \SourceCite{Morrison9} 的一部分。

第 \hyperref[subsec:conifold-transitions-conclusion]{\ref*{subsec:conifold-transitions-conclusion}}、\hyperref[subsec:other-physics-topics]{\ref*{subsec:other-physics-topics}} 小节将讨论极端跃迁的更多方面。第 \hyperref[subsec:flag-cy-complete-intersections]{\ref*{subsec:flag-cy-complete-intersections}} 小节还会说明，极端跃迁可用于构造 Grassmann 簇和部分旗流形中的 Calabi--Yau 完全交的镜像。

\section{镜像映射}\label{sec:mirror-map}
本节描述 Calabi--Yau 流形 $V$ 的复模空间与其镜像流形 $V^\circ$ 的 Kähler 模空间之间的镜像映射。预期这些模空间彼此同构；\SourcePage{172}
镜像映射是一个显式的局部同构，将 $V^\circ$ 的 Kähler 模空间在大半径极限点处的一组自然坐标，与 $V$ 的复模空间在相应最大幺幂边界点附近的一组自然坐标联系起来。

先给出镜像映射的一般定义，再将它应用于环面超曲面，最后讨论几个显式例子。

\subsection{镜像映射的定义}\label{subsec:mirror-map-definition}\index{镜像映射!定义|hyperpage}
\SourceCite{CdGP} 为五次三维簇写出了镜像映射；其他 Kähler 模空间为一维的例子，基本上也用同样的方法处理 \SourceCite{Font,KT1,KT2,Morrison1}。更高维模空间的镜像映射见 \SourceCite{CdFKM,HKTY1,CFKM,HKTY2}，而 \SourceCite{BCOV2} 给出了镜像映射的物理论证。\hyperref[chap:quantum-cohomology]{第~\ref*{chap:quantum-cohomology} 章}将基于量子上同调，从数学上说明镜像映射的动机；\hyperref[chap:mirror-theorem]{第~\ref*{chap:mirror-theorem} 章}给出另一种计算它的方法，例题~\ref{ex:quintic-hypergeometric-mirror-map} 末尾将透露一点线索。镜像映射的纯数学定义归于 Morrison；下面遵循 \SourceCite{Morrison6} 的论述。

为定义镜像映射，先在复模空间 $\overline{\mathcal M}(V)$ 中取一个最大幺幂边界点。按定义~\ref{def:maximal-unipotent-point}，取 $W_2\subset H^d(V)$ 的一组基 $g_0,\ldots,g_r$；这里 $W_2$ 来自单值性权滤过，$d=\dim V$。如式~\eqref{eq:mum-canonical-coordinates}，定义 $q_k$ 为
\begin{equation}\label{eq:mirror-map-period-definition}
\frac{1}{2\pi i}\log q_k
=\frac{1}{\langle g_0,\Omega\rangle}\sum_{j=1}^r\langle g_j,\Omega\rangle m^{jk},
\end{equation}
其中 $(m^{jk})$ 为定义~\ref{def:maximal-unipotent-point} 中矩阵 $(m_{ij})$ 的逆矩阵。由\hyperref[subsec:mum-boundary-points]{第 6.1.1 小节}，当 $p$ 位于所取最大幺幂边界点的邻域中时，$q_i=q_i(p)$ 是全纯函数。

根据定义~\ref{def:maximal-unipotent-point} 之后的讨论，还可如下改写 $q_k$ 的公式。首先，$y_0=\langle g_0,\Omega\rangle$ 是在最大幺幂边界点处全纯的 Picard--Fuchs 方程解，并且在相差常数倍的意义下唯一。再令
\[
y_k=\sum_{j=1}^r\langle g_j,\Omega\rangle m^{jk}.
\]
根据\hyperref[chap:hodge-yukawa]{第~\ref*{chap:hodge-yukawa} 章}，$y_k$ 为 Picard--Fuchs 方程的一个解，满足
\[
y_k=y_0\frac{\log z_k}{2\pi i}+\widetilde y_k,
\]
其中 $\widetilde y_k$ 在边界点处全纯。于是式~\eqref{eq:mirror-map-period-definition} 可写为
$\frac{1}{2\pi i}\log q_k=y_k/y_0$，从而
\begin{equation}\label{eq:mirror-map-holomorphic-correction}
q_k=z_k\exp(2\pi i\,\widetilde y_k/y_0).
\end{equation}
注意，固定 $y_0$ 后，$y_k$ 在相差 $y_0$ 的常数倍的意义下唯一。因此，$q_k$ 在相差常数倍的意义下唯一。

\SourcePage{173}
现在令 $V^\circ$ 为 $V$ 的镜像，假定 $V^\circ$ 的 Kähler 锥满足锥猜想。再取\hyperref[subsec:large-radius-limits]{第 6.2.1 小节}中锥 $K_+$ 的一个细分 $\Sigma_+$，使它给出未扩大的 Kähler 模空间的光滑部分紧化 $\overline{\mathcal{KM}}(V^\circ)$。每个极大维锥 $\sigma\in\Sigma_+$ 都给出 $\overline{\mathcal{KM}}(V^\circ)$ 的一个大半径极限点。由式~\eqref{eq:large-radius-exponential-coordinates}，$\sigma$ 的生成元 $T_i$ 确定典范坐标 $q_i$；它们除次序外是唯一的。

由于 $\overline{\mathcal M}(V)$ 与 $\overline{\mathcal{KM}}(V^\circ)$ 有相同的维数 $r$，这是二者互为镜像的条件之一，可以在 $\overline{\mathcal M}(V)$ 的最大幺幂边界点附近取 $p$，定义
\begin{equation}\label{eq:mirror-map-local}
M(p)=(q_1(p),\ldots,q_r(p))\in\overline{\mathcal{KM}}(V^\circ).
\end{equation}
这里 $q_i(p)$ 是上面定义的全纯函数，而目标空间使用镜像一侧所选大半径极限点处的典范局部坐标。因此，定义镜像映射 $M$ 需要两类数据：
\begin{itemize}
\item 式~\eqref{eq:mirror-map-period-definition} 中的函数 $q_1,\ldots,q_r$；
\item 给出大半径极限点的锥 $\sigma$ 的生成元 $T_1,\ldots,T_r\in H^2(V^\circ)$。
\end{itemize}
第一项完全由 $V$ 的最大幺幂边界点确定，第二项由 $V^\circ$ 的大半径极限点确定。

选定函数 $q_1,\ldots,q_r$ 及锥生成元 $T_1,\ldots,T_r$ 后，式~\eqref{eq:mirror-map-period-definition}、\eqref{eq:mirror-map-local} 就定义了映射 $M$。对于多数选择，$M$ 并不是我们所要的映射；然而，镜像对称的核心正在于：若满足下列条件，
\begin{itemize}
\item $W_2$ 中的类 $g_0,g_1,\ldots,g_r$ 选取得正确，等价地，Picard--Fuchs 方程的解 $y_0,y_1,\ldots,y_r$ 选取得正确；
\item 在两个模空间的相容紧化中，最大幺幂边界点及大半径极限点选取得正确，
\end{itemize}
那么式~\eqref{eq:mirror-map-local} 给出的就是镜像映射。

为理解怎样才算“选取得正确”，需要进一步解释这两点。第一点回到式~\eqref{eq:mum-coordinate-basis-change}：将 $g_0,\ldots,g_r$ 换为 $g_j'=\sum_{k=0}^rc_{jk}g_k$ 后，所得新局部坐标 $(q_1',\ldots,q_r')$ 满足
\begin{equation}\label{eq:mirror-map-constant-ambiguity}
q_k'=e^{2\pi i(c^{k0}/c_{00})}q_k,\qquad
c^{k0}=\sum_{j=1}^rc_{j0}m'^{jk},
\end{equation}
其中 $(m'^{jk})$ 为由 $N_j(g_k')=m'_{jk}g_0'$ 确定的矩阵的逆矩阵。

处理这种非唯一性，有两种已知的方法。第一种使用整数结构 \SourceCite{Morrison6}。单值性权滤过中的空间 $W_0\subset W_2$ 定义在 $\mathbb Q$ 上。用 $W_i\cap H^d(V,\mathbb Z)$ 表示 $W_i$ 与 $H^d(V,\mathbb Z)$ 在 $H^d(V,\mathbb Q)$ 中的像的交，则可以取
$g_0\in W_0\cap H^d(V,\mathbb Z)$、$g_1,\ldots,g_r\in W_2\cap H^d(V,\mathbb Z)$，使 $g_0$ 在 $\mathbb Z$ 上张成前一个格，而 $g_0,g_1,\ldots,g_r$ 在 $\mathbb Z$ 上张成后一个格。称这样的 $g_0,g_1,\ldots,g_r$ 为 $W_0\subset W_2$ 的一组 $\mathbb Z$ 基。

给定这样一组 $\mathbb Z$ 基，容易证明，由 $N_j(g_k)=m_{jk}g_0$ 定义的数 $m_{jk}$ 是整数。\hyperref[subsec:higher-dimensional-mum]{第 5.2.2 小节}的整性猜想断言，矩阵 $(m_{jk})$ 在 $\mathbb Z$ 上可逆；等价地，式~\eqref{eq:mirror-map-period-definition} 中的 $m^{jk}$ 为整数。五次三维簇的镜像已验证这一猜想 \SourceCite{Morrison2}；\SourceCite{Morrison6} 讨论了其他情形。

\SourcePage{174}
下面假定整性猜想成立。设两组 $\mathbb Z$ 基 $g_0,g_1,\ldots,g_r$ 与 $g_0',g_1',\ldots,g_r'$ 满足 $g_j'=\sum_{k=0}^rc_{jk}g_k$，则 $c_{00}=\pm1$ 且 $c_{jk}\in\mathbb Z$。由于还假定 $m'^{jk}\in\mathbb Z$，式~\eqref{eq:mirror-map-constant-ambiguity} 给出 $q_k'=q_k$。因此，只要使用 $\mathbb Z$ 基，式~\eqref{eq:mirror-map-period-definition} 的函数除次序外就是唯一的。

\SourceExampleNumber{6.3.1}{1}
\begin{example}\label{ex:mirror-map-integral-connection}
设 $V$ 为光滑 Calabi--Yau 三维簇，在一个最大幺幂边界点处满足整性猜想；将该点看作 $0\in\Delta^r$。命题~\ref{prop:normalized-yukawa-pf} 之后的讨论说明，可以找到变量 $q_1,\ldots,q_r$，使 Gauss--Manin 联络 $\nabla$ 为
\begin{equation}\label{eq:mirror-map-flat-connection}
\begin{aligned}
\nabla_{\delta_i}e^0&=0,\\
\nabla_{\delta_i}e^k&=\delta_{ik}e^0,&&1\leq k\leq r,\\
\nabla_{\delta_i}e_j&=\sum_{k=1}^rY_{ijk}e^k,&&1\leq j\leq r,\\
\nabla_{\delta_i}e_0&=e_i,
\end{aligned}
\end{equation}
其中 $\delta_j=2\pi iq_j\partial/\partial q_j$，见式~\eqref{eq:multimoduli-connection}。下面说明，此时 $q_1,\ldots,q_r$ 确定镜像映射。

命题~\ref{prop:normalized-yukawa-pf} 的证明中已说明，$e_0$ 是归一化的三形式 $\widetilde\Omega$，平坦截面 $e^0=g_0$ 是 $W_0$ 的整数生成元，且 $\langle g_0,\widetilde\Omega\rangle=1$。此外，$e^1,\ldots,e^r$ 为 $W_2$ 的截面。写 $q_j=\exp(2\pi iu_j)$，容易看出
\[
g_j=-e^j+u_je^0,\qquad1\leq j\leq r
\]
是平坦截面。逆时针绕 $(\Delta^*)^r$ 的第 $i$ 个因子一周，将 $u_j$ 变为 $u_j+\delta_{ij}$，所以 $\mathcal T_i(g_j)=g_j+\delta_{ij}g_0$，进而 $N_i(g_j)=\delta_{ij}g_0$。

由于 $e^0,e^1,\ldots,e^r$ 为 $W_2$ 的截面，$g_0,g_1,\ldots,g_r$ 也是。我们还断言这些截面都是整的。对 $g_0$ 已知如此；对其他截面，注意
$\exp(-\sum_i u_iN_i)g_j=-e^j$。与命题~\ref{prop:normalized-yukawa-pf} 的证明一样，$e^j$ 在 $0\in\Delta^r$ 上为整；由\hyperref[subsec:limiting-hodge]{第 5.1.4 小节}的讨论，这意味着平坦截面 $g_j$ 也为整。因此，$g_0,g_1,\ldots,g_r$ 构成 $W_2$ 的一组平坦整数截面基，且 $g_0$ 张成 $W_0$。

现在可以计算镜像映射。由于 $N_i(g_j)=\delta_{ij}g_0$，矩阵 $(m_{ij})=(\delta_{ij})$ 为单位矩阵，因此式~\eqref{eq:mirror-map-period-definition} 的右端化为 $\langle g_k,\widetilde\Omega\rangle$；回忆 $\langle g_0,\widetilde\Omega\rangle=1$。而
\[
\langle g_j,\widetilde\Omega\rangle
=\langle-e^j+u_je^0,e_0\rangle=u_j,
\]
因为式~\eqref{eq:multimoduli-integral-pairings} 给出 $\langle e^j,e_0\rangle=0$（$1\leq j\leq r$）。由式~\eqref{eq:mirror-map-period-definition}，$q_k$ 确实是镜像映射的第 $k$ 个坐标。
\end{example}

环面情形还有第二种方法，能在不依赖整性猜想的情况下，选出式~\eqref{eq:mirror-map-period-definition} 中正确的 $g_0,g_1,\ldots,g_r$。基本思想是：环面几何自然给出 $z_i$ 的选择；对式~\eqref{eq:mirror-map-holomorphic-correction} 中 $\widetilde y_k/y_0$ 的常数项作适当归一化，就唯一选出 $q_k$。\hyperref[subsec:toric-mirror-map]{第~\ref*{subsec:toric-mirror-map} 小节}将说明这一方法。

\SourcePage{175}
算出 $q_1,\ldots,q_r$ 后，“选取得正确”的另一面，是寻找 $\overline{\mathcal{KM}}(V^\circ)$ 中对应的大半径极限点。这可能并不简单：如\hyperref[subsec:mum-boundary-points]{第 6.1.1 小节}和\hyperref[subsec:large-radius-limits]{第 6.2.1 小节}所见，紧化可能不唯一，还可能有许多不同的点的等价类可供选择。一种办法是寻找适当的极小紧化，尽可能爆缩边界；这或许能减少将最大幺幂边界点与大半径极限点配对的选择。不过，多数时候的问题是，必须相当完整地了解模空间的形态。事实上，$V^\circ$ 为 $V$ 的镜像所包含的要求之一，就是\emph{要充分理解 $\overline{\mathcal M}(V)$ 和 $\overline{\mathcal{KM}}(V^\circ)$，从而知道怎样将最大幺幂边界点与大半径极限点相匹配}。最好的例子是环面情形：两个模空间本质上相同，所以能够准确知道边界点如何匹配。\hyperref[subsec:toric-mirror-map]{第~\ref*{subsec:toric-mirror-map} 小节}将作更详细的讨论。

此外，正确选定大半径极限点后，得到一个由 $T_1,\ldots,T_r$ 生成的锥 $\sigma$，但这些生成元应如何排序？\hyperref[subsec:mirror-map-derivative]{第~\ref*{subsec:mirror-map-derivative} 小节}讨论镜像映射的导数时，将看到这基本上由\hyperref[chap:introduction]{第~\ref*{chap:introduction} 章}式~\eqref{eq:mirror-cohomology-forms} 中的同构 $H^{d-1,1}(V)\simeq H^{1,1}(V^\circ)$ 决定。

所有选择都正确完成后，关键思想是：从 $\overline{\mathcal M}(V)$ 到 $\overline{\mathcal{KM}}(V^\circ)$ 的镜像映射，应与 A 模型及 B 模型关联函数相容。进一步还猜想，这种相容性可推广到 Hodge 结构变分。$\overline{\mathcal M}(V)$ 上有一个来自 $V$ 的复结构变化的自然 Hodge 结构变分；它的渐近行为决定某个边界点是否最大幺幂。\hyperref[chap:quantum-cohomology]{第~\ref*{chap:quantum-cohomology} 章}将定义 $\overline{\mathcal{KM}}(V^\circ)$ 上的 A 模型 Hodge 结构变分，并证明它在大半径极限点处最大幺幂。镜像猜想断言，镜像映射与这些 Hodge 结构变分相容；这将蕴含相应关联函数的相容性。\hyperref[chap:mirror-theorem]{第~\ref*{chap:mirror-theorem} 章}将说明，在环面背景下这一猜想能证明到什么程度。

\subsection{镜像映射的导数}\label{subsec:mirror-map-derivative}
设 $V,V^\circ$ 为维数 $d\geq3$ 的 Calabi--Yau 流形，假定镜像映射已按\hyperref[subsec:mirror-map-definition]{第 6.3.1 小节}定义。接下来要联系以下三个主题：
\begin{itemize}
\item 镜像映射 $M$ 的导数；
\item \hyperref[chap:introduction]{第~\ref*{chap:introduction} 章}式~\eqref{eq:mirror-cohomology-forms} 中的同构 $H^{d-1,1}(V)\simeq H^{1,1}(V^\circ)$；
\item 已知 $q_1,\ldots,q_r$ 后，怎样“正确选取”$T_1,\ldots,T_r$。
\end{itemize}
为研究镜像映射，用复化 Kähler 空间
\[
K_{\mathbb C}(V^\circ)
=(H^2(V^\circ,\mathbb R)+iK(V^\circ))/\operatorname{im}H^2(V^\circ,\mathbb Z)
\]
代替 Kähler 模空间 $\mathcal{KM}(V^\circ)$ 很方便。若假定\hyperref[subsec:large-radius-limits]{第 6.2.1 小节}的锥猜想成立，前者是后者的离散覆盖。为避免处理这一商，只将镜像映射看作 $M:S\to K_{\mathbb C}(V^\circ)$；其中 $S$ 是 $V$ 的复模空间在一个最大幺幂边界点附近的内侧邻域。

\SourcePage{176}
$V^\circ$ 的一个大半径极限点，由 Kähler 锥闭包中整类 $T_1,\ldots,T_r\in H^2(V^\circ,\mathbb C)=H^{1,1}(V^\circ)$ 生成的锥 $\sigma$ 确定。它在所选大半径极限点附近给出坐标
\begin{equation}\label{eq:mirror-map-kahler-coordinates}
(e^{2\pi it_1},\ldots,e^{2\pi it_r})
\longleftrightarrow\left[\sum_jt_jT_j\right]\in K_{\mathbb C}(V^\circ).
\end{equation}
由式~\eqref{eq:mirror-map-local}，$p\in S$ 时 $M(p)=(q_1(p),\ldots,q_r(p))$；式~\eqref{eq:mirror-map-holomorphic-correction} 随即说明，$M$ 可延拓到边界点，并在该处为局部双全纯映射。

为描述 $M$ 的导数，需要理解 $S$ 和 $K_{\mathbb C}(V^\circ)$ 的切丛。先看后者。每个 $T\in H^2(V^\circ,\mathbb C)$ 都在 $H^2(V^\circ,\mathbb R)+iK(V^\circ)$ 上给出一个平移不变的向量场 $\partial/\partial T$，因而可将它看作 $K_{\mathbb C}(V^\circ)$ 上的向量场。由此得到平凡化
\[
T_{K_{\mathbb C}(V^\circ)}\simeq H^2(V^\circ,\mathbb C)\times K_{\mathbb C}(V^\circ).
\]
坐标~\eqref{eq:mirror-map-kahler-coordinates} 给出向量场 $\partial/\partial t_j$，它自然与 $\partial/\partial T_j$ 识别；在上述同构下，$\partial/\partial t_j$ 对应于 $T_j$。利用这一平凡化，局部双全纯映射 $M$ 诱导
\begin{equation}\label{eq:mirror-map-differential}
D_M:T_S\longrightarrow H^2(V^\circ,\mathbb C)\times K_{\mathbb C}(V^\circ).
\end{equation}
由于 $q_1,\ldots,q_r$ 可用作 $S$ 的局部坐标，向量场 $\delta_j=2\pi iq_j\partial/\partial q_j$ 在内部构成 $T_S$ 的一组基；$M$ 的上述公式说明，$D_M$ 将 $\delta_j$ 映到 $\partial/\partial t_j=\partial/\partial T_j$。以式~\eqref{eq:mirror-map-differential} 的语言说，就是映到 $T_j$。

再描述 $S$ 的切丛 $T_S$。$V$ 的复模空间为 $\mathcal M(V)$；在最大幺幂边界点 $p_0$ 附近的内侧邻域 $S$ 上，假定有万有族 $\pi:\mathcal V\to S$；严格地说，应先取 $S$ 的有限覆盖。由于 $V$ 为 Calabi--Yau，Kodaira--Spencer 映射为同构 \SourceCite{Bogomolov,Tian,Todorov}：
\[
\kappa:T_S\simeq R^1\pi_*\Theta_{\mathcal V/S},
\]
其中 $\Theta_{\mathcal V/S}$ 为相对切丛。

还有 $\mathcal H=R^d\pi_*\mathbb C\otimes_{\mathbb C}\mathcal O_S$ 的 Hodge 滤过 $\mathcal F^p$。我们最关心
$\mathcal F^d=\pi_*\Omega^d_{\mathcal V/S}$ 与
$\mathcal F^{d-1}/\mathcal F^d=R^1\pi_*\Omega^{d-1}_{\mathcal V/S}$。内乘给出映射
\begin{equation}\label{eq:mirror-map-contraction}
R^1\pi_*\Theta_{\mathcal V/S}
\longrightarrow\operatorname{Hom}(\mathcal F^d,\mathcal F^{d-1}/\mathcal F^d).
\end{equation}
众所周知，它与 $\kappa$ 的复合
\begin{equation}\label{eq:mirror-map-kodaira-spencer}
T_S\longrightarrow\operatorname{Hom}(\mathcal F^d,\mathcal F^{d-1}/\mathcal F^d)
\end{equation}
将向量场 $X$ 与全纯 $d$ 形式 $\Omega$ 映为 $[\nabla_X(\Omega)]$；方括号表示在商 $\mathcal F^{d-1}/\mathcal F^d$ 中的类。

由于 $V$ 为 Calabi--Yau，可以在 $S$ 上固定 $\mathcal F^d$ 的一个处处非零截面 $\Omega$。照常令 $\widetilde\Omega=\Omega/\langle g_0,\Omega\rangle$ 作归一化；$g_0$ 来自\IndexLocation{indexloc:0772:176}最大幺幂单值性的定义。$\widetilde\Omega$ 诱导 $\mathcal F^d\simeq\mathcal O_S$，因而式~\eqref{eq:mirror-map-kodaira-spencer} 给出同构
\begin{equation}\label{eq:mirror-map-infinitesimal-period}
T_S\simeq R^1\pi_*\Theta_{\mathcal V/S}
\simeq R^1\pi_*\Omega^{d-1}_{\mathcal V/S},
\end{equation}
由 $X\mapsto[\nabla_X(\widetilde\Omega)]$ 定义。它的含义如下：设 $p\in S$、$V=V_p$，将 $X$ 看作把 $V_p$ 变为 $V_\varepsilon$ 的\IndexLocation{indexloc:0209:176}无穷小形变，则在此形变下，$\widetilde\Omega_p$ 变为
$\widetilde\Omega_\varepsilon=\widetilde\Omega_p+\varepsilon\nabla_X(\widetilde\Omega)_p+O(\varepsilon^2)$。

\SourcePage{177}
现在可以联系前面提到的三点。首先，对 $p\in S$、$V=V_p$，结合式~\eqref{eq:mirror-map-differential} 与~\eqref{eq:mirror-map-infinitesimal-period}，得到同构
\begin{equation}\label{eq:mirror-map-hodge-differential}
H^1(V,\Omega_V^{d-1})\simeq T_{S,p}
\simeq H^2(V^\circ,\mathbb C)=H^{1,1}(V^\circ).
\end{equation}
因此，镜像映射的导数就是\hyperref[chap:introduction]{第~\ref*{chap:introduction} 章}同构~\eqref{eq:mirror-cohomology-forms} 的 $d$ 维版本。此外，前面已看到 $D_M$ 将 $\delta_j=2\pi iq_j\partial/\partial q_j$ 映为 $T_j$。在式~\eqref{eq:mirror-map-hodge-differential} 中，这表示在 $V=V_p$ 处，将 $[\nabla_{\delta_j}\widetilde\Omega]$ 映为 $T_j$，即由 $2\pi iq_j\partial/\partial q_j\mapsto T_j$ 定义相应同构。因此，知道式~\eqref{eq:mirror-map-hodge-differential}，就等于知道怎样“正确选取”$T_1,\ldots,T_r$。

最后给出一个\hyperref[chap:quantum-cohomology]{第~\ref*{chap:quantum-cohomology} 章}会用到的引理。它将镜像映射的导数与丛 $\mathcal H=\mathcal F^0=R^d\pi_*\mathbb C\otimes_{\mathbb C}\mathcal O_S$ 联系起来。

\begin{lemma}\label{lem:mirror-map-hodge-splitting}
设 $\widetilde\Omega$ 为最大幺幂边界点 $p_0$ 附近的归一化 $d$ 形式。则：
\begin{enumerate}[label=(\roman*)]
\item 在 $p_0$ 附近，映射
$\psi:\mathcal O_S\oplus T_S\to\mathcal F^{d-1}\subset\mathcal H$，
$\psi(f,X)=f\widetilde\Omega+\nabla_X(\widetilde\Omega)$，为与式~\eqref{eq:mirror-map-infinitesimal-period} 相容的同构。边界处使用对数切丛及 Hodge 丛的典范延拓。
\item 经 $M$ 识别底空间后，有丛同构
\[
\phi:\mathcal F^{d-1}\simeq
\bigl(H^0(V^\circ,\mathbb C)\oplus H^2(V^\circ,\mathbb C)\bigr)\times M(S),
\]
它将 $\widetilde\Omega$ 映为 $1\in H^0(V^\circ,\mathbb C)$；将 $T_S$ 经 $\psi$ 视为 $\mathcal F^{d-1}$ 的子丛后，$\phi$ 在其上的限制就是镜像映射的导数。
\end{enumerate}
\end{lemma}
\begin{proof}
最大幺幂单值性的定义给出 $W_0=W_1$ 且维数为 1。杯积与单值性权滤过相容，因此 $\operatorname{Gr}^W_{2d}$ 的维数为 1，且 $W_{2d-1}=W_{2d-2}$。用于证明式~\eqref{eq:mum-top-splitting} 的论证给出
\[
F^d_{\mathrm{lim}}\oplus W_{2d-1}
=F^d_{\mathrm{lim}}\oplus W_{2d-2}=H^d(V,\mathbb C),
\]
其中 $F^d_{\mathrm{lim}}$ 来自极限 Hodge 滤过。由此，$\operatorname{Gr}^W_{2d-2}$ 上的 Hodge 结构纯为 $(d-1,d-1)$ 型，进而
$F^{d-1}_{\mathrm{lim}}/F^d_{\mathrm{lim}}\simeq\operatorname{Gr}^W_{2d-2}$。

结合最大幺幂单值性的定义与杯积在 $\operatorname{Gr}^W_\bullet$ 上给出的对偶性，可知：若 $g\in W_{2d}$ 生成 $\operatorname{Gr}^W_{2d}$，则 $N_j(g)$ 构成 $\operatorname{Gr}^W_{2d-2}$ 的一组基。再结合式~\eqref{eq:connection-monodromy-residue}，便知转到典范延拓后，引理在 $p_0$ 处成立。于是 $\mathcal O_S\oplus T_S\simeq\mathcal F^{d-1}$ 立即成立；其余结论也随之得到。
\end{proof}

式~\eqref{eq:mirror-map-flat-connection} 提供了这一引理的具体说明：$e_0=\widetilde\Omega$、$e_i=\nabla_{\delta_i}(\widetilde\Omega)$。若用 $1$ 作 $\mathcal O_S$ 的基，用 $\delta_i$ 作 $T_S$ 的基，则引理第一部分的同构将 $1$ 映为 $\widetilde\Omega$，将 $\delta_i$ 映为 $e_i$。特别地，式~\eqref{eq:mirror-map-flat-connection} 显式包含了 $T_S$ 的一个副本。第二部分的同构将 $e_0$ 映为 $1\in H^0(V^\circ,\mathbb C)$，将 $e_i$ 映为 $T_i$（$1\leq i\leq r$）。\hyperref[chap:quantum-cohomology]{第~\ref*{chap:quantum-cohomology} 章}将把 $\phi$ 扩大为 $\mathcal H=\mathcal F^0$ 与由 $\bigoplus_pH^{2p}(V^\circ,\mathbb C)$ 给出的平凡丛之间的同构。我们将看到，这属于镜像对称的 Hodge 理论表述。

\SourcePage{178}
\subsection{环面超曲面的镜像映射}\label{subsec:toric-mirror-map}
对于环面超曲面，可以更加精确地说明镜像映射应是什么，也能给出处理式~\eqref{eq:mirror-map-constant-ambiguity} 中非唯一性的另一种方案。

照常固定\IndexLocation{indexloc:1218:178}自反多胞形 $\Delta$，令 $V\subset X_\Sigma$ 为 Calabi--Yau 超曲面，$\Sigma$ 为 $\Delta$ 的法扇的简化射影细分。再令 $V^\circ\subset X_{\Sigma^\circ}$ 为 $V$ 的一个 Batyrev 镜像，置
\[
\Xi=(\Delta\cap M)_0-\{0\},\qquad
\Xi^+=(\Xi\cup\{0\})\times\{1\}=(\Delta\cap M)_0\times\{1\}.
\]
注意，$\Sigma^\circ(1)=\Xi$，因为 $\Sigma^\circ$ 是 $\Delta^\circ$ 的法扇的简化射影细分。还要回忆，$\Xi$ 的 GKZ 分解与 $\Xi^+$ 的次级扇，都位于向量空间 $A(\Xi)$ 中。

为了确定镜像映射，还需要更多数据。锥 $\operatorname{cpl}(\Sigma^\circ)$ 是 $\Xi$ 的 GKZ 分解中的一个锥。将它细分为光滑锥 $\sigma$；于是每个 $\sigma$ 都是单纯锥，其生成元构成格的一组整基。

我们的基本主张是：每个锥 $\sigma$ 都给出了确定镜像映射所需的全部信息。换言之，$\sigma$ 确定相互对应的最大幺幂边界点与\IndexLocation{indexloc:0696:178}大半径极限点，并且\emph{不必假定\IndexLocation{indexloc:0636:178}整性猜想，也不必知道整数上同调的任何信息}，就能选定式~\eqref{eq:mirror-map-constant-ambiguity} 中的常数。此外，从 $\sigma$ 出发，还会自动“正确选取”$T_1,\ldots,T_r$。

先看\hyperref[subsec:toric-complex-moduli]{第 6.1.2 小节}用 $\Xi^+$ 的次级扇构造的简化多项式模空间，见式~\eqref{eq:moduli-secondary-configuration}。由于希望紧化光滑，采用次级扇的一个细化，其中包含上述细分 $\operatorname{cpl}(\Sigma^\circ)$ 的各锥 $\sigma$。将这一紧化记为 $\overline{\mathcal M}_{\mathrm{simp}}$。假定猜想~\ref{conj:toric-mum-boundary} 成立，本章余下部分都作此假定，则每个 $\sigma$ 都给出它的一个最大幺幂边界点。

固定一个 $\sigma$ 后，可如下得到该最大幺幂边界点处的自然坐标。$\sigma$ 的生成元构成 $A(\Xi)$ 中的格的一组基；\hyperref[subsec:toric-complex-moduli]{第 6.1.2 小节}将该格记为 $N_0$。对偶基 $\ell_1,\ldots,\ell_r$ 生成 $A(\Xi)^*$ 中的格，即式~\eqref{eq:moduli-relation-lattice} 中的 $M_0$。它是 $\Xi^+$ 的关系格，所以每个基元素可写为 $\ell_k=(\ell_{kj})$。由式~\eqref{eq:moduli-monomial-coordinate}，
\begin{equation}\label{eq:toric-mirror-moduli-coordinates}
z_k=\prod_j\lambda_j^{\ell_{kj}}
\end{equation}
为最大幺幂边界点处的局部坐标。回忆，$\lambda_i$ 为 Laurent 多项式
\begin{equation}\label{eq:toric-mirror-laurent-polynomial}
\sum_{m_i\in(\Delta\cap M)_0}\lambda_i t^{m_i}
=\lambda_0+\sum_{m_i\in\Xi}\lambda_i t^{m_i}
\end{equation}
的系数；仍取 $m_0=0$。这样书写局部坐标，强调我们所处理的是模；更准确地说是简化模，因为\IndexLocation{indexloc:0217:178}占优猜想意味着 $\overline{\mathcal M}_{\mathrm{simp}}$ 为真正多项式模空间的有限覆盖。

镜像映射将是式~\eqref{eq:toric-mirror-moduli-coordinates} 中坐标 $z_1,\ldots,z_r$ 的函数。下面将看到，每个 $\ell_k=(\ell_{kj})$ 都贡献一个符号
\begin{equation}\label{eq:toric-mirror-sign}
(-1)^{\ell_{k0}},\qquad k=1,\ldots,r,
\end{equation}
它对确定镜像映射有重要作用。

再看\hyperref[subsec:enlarged-kahler-moduli]{第 6.2.3 小节}的环面 Kähler 模空间 $\overline{\mathcal{KM}}_{\mathrm{toric}}(V^\circ)$。假定猜想~\ref{conj:toric-kahler-cone} 成立，则在相差有限群作用的意义下，
\SourcePage{179}
环面 Kähler 模的一个紧化由 $\Xi^+$ 的次级扇给出，其中需将对应于平凡翻转的锥合并。为了得到与 $V$ 的复模相容的紧化，使用上述次级扇细化给出的光滑爆破，记为 $\overline{\mathcal{KM}}_{\mathrm{toric}}(V^\circ)$。特别地，细分 $\operatorname{cpl}(\Sigma^\circ)$ 的锥 $\sigma$ 给出该空间的大半径极限点。

对每个 $\sigma$，可如下得到相应大半径极限点处的局部坐标 $q_1,\ldots,q_r$。由定理~\ref{thm:batyrev-hodge} 的证明，首先有自然同构
$A(\Xi)\simeq H^2_{\mathrm{toric}}(V^\circ,\mathbb R)$。若 $\sigma$ 的生成元为 $T_1,\ldots,T_r$，则映射
\[
(\Delta^*)^r\longrightarrow
\frac{H^2_{\mathrm{toric}}(V^\circ,\mathbb R)+iK_{\mathrm{toric}}(V^\circ)}
{\operatorname{im}H^2_{\mathrm{toric}}(V^\circ,\mathbb Z)},
\qquad(e^{2\pi it_1},\ldots,e^{2\pi it_r})\longmapsto\left[\sum_{j=1}^rt_jT_j\right]
\]
延拓为开嵌入 $\Delta^r\subset\overline{\mathcal{KM}}_{\mathrm{toric}}(V^\circ)$，原点映到所选大半径极限点。该点处的局部坐标于是为 $q_j=e^{2\pi it_j}$；定义镜像映射时将使用它们。

紧化 $\overline{\mathcal M}_{\mathrm{simp}}(V)$ 与 $\overline{\mathcal{KM}}_{\mathrm{toric}}(V^\circ)$ 相同；固定 $\sigma$ 给出相容的\IndexLocation{indexloc:0766:179}最大幺幂边界点与大半径极限点。前者有局部坐标 $z_1,\ldots,z_r$，后者有 $q_1,\ldots,q_r$。这些坐标其实相同，虽然这一点未必显然。原因是 $z_j$ 对应于模空间中环面 $T_0$ 的特征标群的一组基，沿用\hyperref[subsec:toric-complex-moduli]{第 6.1.2 小节}的记号；而 $q_j$ 由 $T_0$ 的单参数子群的一组基确定。两组基互为对偶，因而给出相同坐标。为使记号清楚，仍在 $\overline{\mathcal M}_{\mathrm{simp}}(V)$ 上使用 $z_j$，在 $\overline{\mathcal{KM}}_{\mathrm{toric}}(V^\circ)$ 上使用 $q_j$。

现在可以确定镜像映射。按式~\eqref{eq:mirror-map-period-definition} 定义 $q_k=q_k(z_1,\ldots,z_r)$，并如下消除其非唯一性。由式~\eqref{eq:mirror-map-holomorphic-correction}，有 $q_k(z_1,\ldots,z_r)\sim cz_k$，其中 $c\ne0$ 为常数。利用式~\eqref{eq:mirror-map-constant-ambiguity}，可使
\begin{equation}\label{eq:toric-mirror-normalization}
q_k(z_1,\ldots,z_r)\sim(-1)^{\ell_{k0}}z_k,
\end{equation}
其中符号来自式~\eqref{eq:toric-mirror-sign}。这显然唯一确定函数 $q_1,\ldots,q_r$，也就定义了唯一的镜像映射~\eqref{eq:mirror-map-local}。用式~\eqref{eq:mirror-map-holomorphic-correction} 前的记号，就是寻找 Picard--Fuchs 方程的解 $y_k$，使
\begin{equation}\label{eq:toric-mirror-normalized-period}
y_k=y_0\frac{\log\bigl((-1)^{\ell_{k0}}z_k\bigr)}{2\pi i}+\widetilde y_k,
\end{equation}
其中 $\widetilde y_k$ 全纯，且在 $z_1=\cdots=z_r=0$ 处为零。于是
$q_k=(-1)^{\ell_{k0}}z_k\exp(2\pi i\,\widetilde y_k/y_0)$。

这一公式虽然看起来有些特别，但符号 $(-1)^{\ell_{k0}}$ 已出现在例题~\ref{ex:quintic-toric-mirror-map} 中，下面加以说明。

\SourceExampleNumber{6.3.3}{1}
\begin{example}\label{ex:quintic-toric-mirror-map}
\hyperref[chap:quintic]{第~\ref*{chap:quintic} 章}讨论\IndexLocation{indexloc:1195:179}五次三维簇 $V$ 时，镜像映射从镜像 $V^\circ$ 的复模，映到 $V$ 的 Kähler 模。
\SourcePage{180}
用齐次坐标，$V^\circ$ 的方程为
\[
\lambda_1x_1^5+\lambda_2x_2^5+\lambda_3x_3^5+\lambda_4x_4^5+\lambda_5x_5^5
+\lambda_0x_1x_2x_3x_4x_5=0,
\]
而 $\Xi^+$ 的元素之间的关系为 $\ell=(1,1,1,1,1,-5)$。如式~\eqref{eq:quintic-gkz-coordinate}，$V^\circ$ 的模坐标为 $z=\lambda_1\cdots\lambda_5/\lambda_0^5$。该式之后的讨论表明，$z$ 正是\hyperref[chap:quintic]{第~\ref*{chap:quintic} 章}称为 $x$ 的坐标。$V$ 的 Kähler 模还有坐标 $q$。

例题~\ref{ex:quintic-hypergeometric-mirror-map} 及其后续讨论将证明，$V^\circ$ 的 Picard--Fuchs 方程有全纯解
\[
y_0(z)=\sum_{n=0}^\infty\frac{(5n)!}{(n!)^5}(-z)^n,
\]
以及解
\[
y_0\frac{\log(-z)}{2\pi i}
+\frac5{2\pi i}\sum_{n=1}^\infty\frac{(5n)!}{(n!)^5}
\left[\sum_{j=n+1}^{5n}\frac1j\right](-z)^n,
\]
具有式~\eqref{eq:toric-mirror-normalization} 所要求的形式。因此，五次三维簇的环面镜像映射为
\[
q=-z\exp\left(\frac5{y_0(z)}\sum_{n=1}^\infty\frac{(5n)!}{(n!)^5}
\left[\sum_{j=n+1}^{5n}\frac1j\right](-z)^n\right).
\]
这就是式~\eqref{eq:q-series-c1} 中的镜像映射。\hyperref[chap:quintic]{第~\ref*{chap:quintic} 章}通过计算 $V^\circ$ 的 Yukawa 耦合，并利用一般 $V$ 含有 2875 条直线，确定右端前面的负号。现在知道了式~\eqref{eq:toric-mirror-normalization}，就直接由 $(-1)^{\ell_{k0}}=\allowbreak(-1)^{-5}=-1$ 得到这个负号。

还可与\hyperref[subsec:mirror-map-definition]{第 6.3.1 小节}定义的镜像映射比较。由于五次镜像满足整性猜想 \SourceCite{Morrison2}，能取满足 $N(g_1)=g_0$ 的整类 $g_0,g_1$，使式~\eqref{eq:mirror-map-period-definition} 中的镜像映射为
\[
\frac1{2\pi i}\log q=\frac{\langle g_1,\Omega\rangle}{\langle g_0,\Omega\rangle}.
\]
它与上述环面定义一致，并非显然。不过，如\hyperref[ex:quintic-normalized-yukawa]{例题~\ref*{ex:quintic-normalized-yukawa}}所述，这由 \SourceCite{CdGP,Morrison2} 的显式计算得到。令 $\gamma_0,\gamma_1$ 分别为 $g_0,g_1$ 的 Poincaré 对偶，上式还意味着
\[
q=\exp\left(2\pi i\frac{\int_{\gamma_1}\Omega}{\int_{\gamma_0}\Omega}\right).
\]
这正是\hyperref[sec:quintic-mirror-map]{第 2.3 节}开头的定义。因此，对五次三维簇，各种不同的镜像映射定义都一致。
\end{example}

理解式~\eqref{eq:toric-mirror-normalization} 的一种方式，是改变 $\lambda_0$ 的符号：在式~\eqref{eq:toric-mirror-laurent-polynomial} 中用 $-\lambda_0$ 代替 $\lambda_0$，则由式~\eqref{eq:toric-mirror-moduli-coordinates}，每个 $z_k$ 乘上符号 $(-1)^{\ell_{k0}}$。若一开始这样安排，就用 $q_k\sim z_k$ 归一化镜像映射，而不必使用式~\eqref{eq:toric-mirror-normalization}。\SourceCite{BvS} 采用这种方法。另一方面，\SourceCite{HKTY1} 使用与我们相同的 $z_k$，随后引入新变量 $x_k=(-1)^{\ell_{k0}}z_k$，用 $x_k$ 定义镜像映射。

\SourcePage{181}
另一种理解是：式~\eqref{eq:toric-mirror-normalization} 等价于指定镜像映射的导数。利用\hyperref[subsec:mirror-map-derivative]{第 6.3.2 小节}的思想，最大幺幂边界点处的切空间可自然识别为定理~\ref{thm:batyrev-hodge} 的证明中记为 $H^1_{\mathrm{poly}}(V,T_V)\simeq H^{n-2,1}_{\mathrm{poly}}(V)$ 的空间；类似地，大半径极限点处的切空间为 $H^2_{\mathrm{toric}}(V^\circ)$。对于坐标 $z_k,q_k$，分别有自然基 $\partial/\partial z_k$ 和 $\partial/\partial q_k$。在这些坐标中，镜像映射的导数为同构
\[
H^{n-2,1}_{\mathrm{poly}}(V)\simeq H^2_{\mathrm{toric}}(V^\circ),\qquad
\frac\partial{\partial z_k}\longmapsto
\sum_j\frac{\partial q_j}{\partial z_k}\frac\partial{\partial q_j}.
\]
在边界点 $z_1=\cdots=z_r=0$，式~\eqref{eq:toric-mirror-normalization} 表明，导数将 $\partial/\partial z_k$ 映为 $(-1)^{\ell_{k0}}\partial/\partial q_k$。理清这些记号，就知道这一映射把单项式确定的 $V$ 的一阶形变，映为来自 $V^\circ$ 的环面除子的上同调类，至多相差符号 $(-1)^{\ell_{k0}}$。因此，它，或者更准确地说，不带这些符号的版本，称为\emph{单项式–除子镜像映射} \SourceCite{AGM1}。

\SourceCite[第 4 节]{AGM1} 的基本猜想之一，是单项式–除子镜像映射在相差一个阶为 2 的自同构的意义下，为镜像映射的导数。所谓阶为 2 的自同构，表示需要指定某些符号；这正是式~\eqref{eq:toric-mirror-normalization} 所做的事情。在我们的处理下，鉴于环面情形的镜像映射定义，这一猜想由定义成立。不过，要真正称得上“镜像映射”，这里描述的映射必须与\hyperref[subsec:mirror-map-definition]{第 6.3.1 小节}末尾提到的关联函数和 Hodge 结构变分相容。这在完全一般的情形中尚未得到证明，因此仍有工作要做。

最后还可以问，为何式~\eqref{eq:toric-mirror-normalization} 是正确的归一化方式。物理学的解释见 \SourceCite{MP1}。从数学看，符号 $(-1)^{\ell_{k0}}$ 最容易用超几何函数解释；下一小节将讨论它们。

最后，我们计算一个熟悉例子的部分镜像映射。

\SourceExampleNumber{6.3.3}{2}
\begin{example}\label{ex:weighted-boundary-mirror-map}
$\mathbb P(1,1,2,2,2)$ 的环面爆破中反典范超曲面的镜像 $V^\circ$，已经讨论多次，最近一次是\hyperref[subsec:kahler-moduli-boundary]{第 6.2.4 小节}。这里不计算整个镜像映射，而只研究与\hyperref[ex:weighted-extremal-transition]{例题~\ref*{ex:weighted-extremal-transition}}中的极端跃迁相关的部分。回忆，实际上有两个跃迁：从 $V$ 到 $\widetilde V$ 的原跃迁，以及从 $\widetilde V^\circ$ 到 $V^\circ$ 的“对偶”跃迁。前者涉及 $V$ 的 Kähler 锥的一个面，后者使用 $V^\circ$ 的判别式轨迹的一部分。问题在于：它们在镜像映射下怎样匹配？

使用\hyperref[ex:weighted-complex-moduli]{例题~\ref*{ex:weighted-complex-moduli}}给出的最大幺幂边界点附近的局部坐标 $z_1,z_2$，在轨迹 $z_1=0$ 上研究镜像映射。这里的计算首先见于 \SourceCite{AGM3}，\SourceCite{Morrison9} 作了进一步阐述。目标是假定 $z_1=0$，将 $q_2$ 算为 $z_2$ 的函数。由于 $z_2=\lambda_1\lambda_2/\lambda_6^2$，有 GKZ 算子 $\Box=\partial_1\partial_2-\partial_6^2$，其中 $\partial_i$ 表示对 $\lambda_i$ 的偏导。导出式~\eqref{eq:weighted-logarithmic-identities} 的计算同样适用，给出周期的 Picard--Fuchs 方程
\[
\delta^2y-2z_2\delta(2\delta+1)y=0.
\]
\SourcePage{182}
这里 $y=y(z_2)$、$\delta=z_2\,d/dz_2$。一般解为
\begin{equation}\label{eq:weighted-boundary-period}
y(z_2)=c_1+c_2\log\left(\frac{\sqrt{1-4z_2}+2z_2-1}{2z_2}\right).
\end{equation}
为了计算镜像映射，需要 $z_2=0$ 附近的解析解与对数解。解析解只有常数，因而取 $y_0(z_2)=1$。由于 $\sqrt{1-4z_2}$ 的 Taylor 展开为 $1-2z_2-2z_2^2+\cdots$，式~\eqref{eq:weighted-boundary-period} 中乘 $c_2$ 的函数，在 $z_2=0$ 处按对数增长。为了使 $q_2(z_2)=\exp(2\pi i\,y_1(z_2)/y_0(z_2))$ 具有式~\eqref{eq:toric-mirror-normalization} 预言的渐近行为，取 $c_1=1/2$、$c_2=1/(2\pi i)$。这给出
\[
y_1(z_2)=\frac1{2\pi i}\log\left(\frac{1-2z_2-\sqrt{1-4z_2}}{2z_2}\right),
\]
于是 $z_1=0$ 上的镜像映射为
\begin{equation}\label{eq:weighted-boundary-mirror-coordinate}
q_2(z_2)=\frac{1-2z_2-\sqrt{1-4z_2}}{2z_2}
=z_2+2z_2^2+5z_2^3+\cdots.
\end{equation}
在 $\overline{\mathcal M}_{\mathrm{simp}}(V^\circ)$ 中，\hyperref[ex:weighted-extremal-transition]{例题~\ref*{ex:weighted-extremal-transition}}说明，对偶极端跃迁发生于\IndexLocation{indexloc:0215:182}判别式轨迹的一部分 $1-4z_2=0$。这条轨迹与 $z_1=0$ 交于 $z_2=1/4$。代入式~\eqref{eq:weighted-boundary-mirror-coordinate}，得到 $q_2=1$。而 $q_2=\exp(2\pi it_2)$，所以 $t_2\in\mathbb Z$，特别地，$t_2$ 的虚部为零。坐标为 $q_1,q_2$ 的点，对应于复化 Kähler 类
\[
\begin{aligned}
t_1T_1+t_2T_2
&=\operatorname{Re}(t_1)T_1+\operatorname{Re}(t_2)T_2\\
&\quad+i\bigl(\operatorname{Im}(t_1)T_1+\operatorname{Im}(t_2)T_2\bigr)=B+i\omega.
\end{aligned}
\]
$\operatorname{Im}(t_2)=0$ 意味着 $\omega$ 位于原极端跃迁所发生的 Kähler 锥面。因此，恰好得到镜像对称预言的相容性；也看清了为何出现 $z_2=1/4$ 而非 $z_2=1$。若镜像映射只是 $q_2=z_2$，才会得到后一个值。

下一小节将回到这一例子，用超几何函数计算完整镜像映射。
\end{example}

\subsection{用超几何函数求镜像映射}\label{subsec:hypergeometric-mirror-map}
\hyperref[sec:hypergeometric]{第~\ref*{sec:hypergeometric} 节}和\hyperref[subsec:hypergeometric-moduli]{第~\ref*{subsec:hypergeometric-moduli} 小节}用 GKZ 超几何方程组，研究 Calabi--Yau \IndexLocation{indexloc:0099:182}环面超曲面 $V\subset X_\Sigma$ 的简化模空间上的 Picard--Fuchs 方程。这里继续这一讨论，用超几何方法构造镜像映射。沿用上一小节的记号；特别地，假定所处位置是\IndexLocation{indexloc:0766:182}最大幺幂边界点，其局部坐标 $z_1,\ldots,z_r$ 由式~\eqref{eq:toric-mirror-moduli-coordinates} 给出。

$\overline{\mathcal M}_{\mathrm{simp}}$ 上的 GKZ 方程组由\hyperref[subsec:hypergeometric-moduli]{第~\ref*{subsec:hypergeometric-moduli} 小节}讨论的微分算子 $\Box_\ell\circ\lambda_0^{-1}$ 组成，其中 $\ell$ 属于式~\eqref{eq:moduli-relation-lattice} 中 $\Xi^+$ 的关系格 $M_0$。对 $f=\lambda_0+\sum_{m_i\in\Xi}\lambda_i t^{m_i}$，还知道 $n$ 形式
\[
\widetilde\omega=\frac{\lambda_0}{f}\frac{dt_1}{t_1}\wedge\cdots\wedge\frac{dt_n}{t_n}
\]
是一个解，即对所有 $\ell$，有 $\Box_\ell(\lambda_0^{-1}\widetilde\omega)=0$。

形式 $\widetilde\omega$ 位于 $X_\Sigma-V$ 上，其在 $V$ 上的留数为全纯 $(n-1)$ 形式 $\Omega$。此外，在 $T\subset X_\Sigma$ 上还有由 $|t_1|=\cdots=|t_n|=1$ 定义的圈 $\gamma$，
\SourcePage{183}
它给出上同调类 $g\in H^{n-1}(V)$。于是
\[
\langle g,\Omega\rangle=\frac1{(2\pi i)^n}\int_\gamma
\frac{\lambda_0}{f}\frac{dt_1}{t_1}\wedge\cdots\wedge\frac{dt_n}{t_n}.
\]
这里加入常数 $1/(2\pi i)^n$，以便计算积分。由\hyperref[chap:hodge-yukawa]{第~\ref*{chap:hodge-yukawa} 章}，知道这一函数是 Picard--Fuchs 方程的解。上述积分是 \emph{Euler 积分}的一个例子；这类 GKZ 方程组的解在 \SourceCite{GKZ1} 中有详细研究。\index{Euler积分@Euler 积分|hyperpage}

为了看清这个解的形式，沿用 \SourceCite{Batyrev3} 的方法，注意到
\[
\frac{\lambda_0}{f}
=\frac1{1-\sum_{m_i\in\Xi}\lambda_i(-\lambda_0)^{-1}t^{m_i}}
=\sum_{k=0}^\infty
\left(\sum_{m_i\in\Xi}\lambda_i(-\lambda_0)^{-1}t^{m_i}\right)^k.
\]
设 $\Xi$ 有 $s$ 个元素。把这一恒等式代入上述积分，逐项积分，便得
\begin{equation}\label{eq:toric-euler-period}
\langle g,\Omega\rangle=
\sum_{\substack{\ell=(u_1,\ldots,u_s,u_0)\in M_0\\u_1,\ldots,u_s\geq0}}
\frac{(-u_0)!}{u_1!\cdots u_s!}(-\lambda_0)^{u_0}
\lambda_1^{u_1}\cdots\lambda_s^{u_s}.
\end{equation}
这是因为 $\ell\in M_0$ 蕴含 $\sum_{i=1}^su_i m_i=0$ 及 $-u_0=u_1+\cdots+u_s$。式~\eqref{eq:toric-euler-period} 中 $\lambda_0$ 前的负号至关重要。

下面用一个具体例子说明级数~\eqref{eq:toric-euler-period}。

\SourceExampleNumber{6.3.4}{1}
\begin{example}\leavevmode\phantomsection\label{ex:quintic-hypergeometric-mirror-map}%
对于五次三维簇的镜像 $V^\circ$，$\Xi^+$ 的元素间的关系只有 $\ell=(1,1,1,1,1,-5)$ 的整数倍。因此，对 $(n,n,n,n,n,-5n)$，$n\geq0$，求和，得到
\begin{equation}\label{eq:quintic-euler-period}
\begin{aligned}
\langle g,\Omega\rangle
&=\sum_{n=0}^\infty\frac{(-(-5n))!}{(n!)^5}
(-1)^{-5n}\lambda_0^{-5n}\lambda_1^n\lambda_2^n\lambda_3^n\lambda_4^n\lambda_5^n\\
&=\sum_{n=0}^\infty\frac{(5n)!}{(n!)^5}(-z)^n,
\end{aligned}
\end{equation}
其中照常取 $z=\lambda_1\lambda_2\lambda_3\lambda_4\lambda_5/\lambda_0^5$。这恰好是\hyperref[chap:quintic]{第~\ref*{chap:quintic} 章}给出的 $y_0(z)$ 的公式~\eqref{eq:y0-hypergeometric}。
\end{example}

一般地，可证明解~\eqref{eq:toric-euler-period} 在最大幺幂边界点处全纯。

\begin{proposition}\label{prop:toric-period-holomorphic}
将级数~\eqref{eq:toric-euler-period} 写成最大幺幂边界点处坐标 $z_1,\ldots,z_r$ 的函数，便得到该点某个邻域内处处非零的全纯函数。
\end{proposition}
\begin{proof}
第一步是用 $z_1,\ldots,z_r$ 写出式~\eqref{eq:toric-euler-period}，并证明只出现非负指数。$z_k$ 来自 $\Xi^+$ 的关系格的一组基 $\ell_k=(\ell_{kj})$，所以式~\eqref{eq:toric-euler-period} 中任一 $\ell$ 都形如 $\ell=n_1\ell_1+\cdots+n_r\ell_r$，其中 $n_1,\ldots,n_r$ 为整数。因此，该级数可写为
\begin{equation}\label{eq:toric-period-local-series}
\sum_{n_1,\ldots,n_r}
\frac{\bigl(-\sum_{k=1}^rn_k\ell_{k0}\bigr)!}
{\prod_{j=1}^s\bigl(\sum_{k=1}^rn_k\ell_{kj}\bigr)!}
\prod_{k=1}^r\bigl((-1)^{\ell_{k0}}z_k\bigr)^{n_k},
\end{equation}
其中求和遍及满足 $\sum_{k=1}^rn_k\ell_{kj}\geq0$，$j=1,\ldots,s$，的所有 $n_1,\ldots,n_r$。需要证明，这些不等式蕴含 $n_1,\ldots,n_r\geq0$；也就是，每个满足 $u_1,\ldots,u_s\geq0$ 的
\SourcePage{184}
$\ell=(u_1,\ldots,u_s,u_0)\in M_0$，都是 $\ell_k=(\ell_{kj})$ 的非负整数线性组合。

只要理清 $M_0$ 及其对偶 $N_0$ 的元素的各种解释，这一点其实很容易看出。在 $\mathbb R$ 上，有
$N_0\otimes\mathbb R=A(\Xi)=H^2_{\mathrm{toric}}(V^\circ,\mathbb R)$，而 $M_0\otimes\mathbb R=A(\Xi)^*$ 为其对偶。$X_{\Sigma^\circ}$ 上每个环面除子 $D_j$，在 $V^\circ$ 上给出仍记为 $D_j$ 的除子，因而有 $[D_j]\in H^2_{\mathrm{toric}}(V^\circ,\mathbb R)$。由于 $\Sigma^\circ(1)=\Xi$，$[D_1],\ldots,[D_s]$ 给出 $M_0$ 上的线性泛函。进一步，作为 $M_0$ 上的泛函，$[D_j]$ 取出 $\ell=(u_1,\ldots,u_s,u_0)$ 的第 $j$ 个坐标。于是 $u_1,\ldots,u_s\geq0$ 就是说，$\ell$ 位于各 $[D_j]$ 生成的锥的对偶锥中。各 $[D_j]$ 生成的锥正是 $A(\Xi)$ 的有效锥，也是 GKZ 分解的支集。指定所选最大幺幂边界点的锥 $\sigma$ 位于 $\operatorname{cpl}(\Sigma^\circ)$ 中，故位于有效锥中；因此，其对偶锥包含有效锥的对偶锥。于是 $\ell$ 位于 $\sigma$ 的对偶锥中，而这个锥由 $\ell_k=(\ell_{kj})$ 生成。

最后证明级数在原点的某个邻域收敛。展开 $(1+\cdots+1)^{u_1+\cdots+u_s}$，得到初等估计
\[
\frac{(u_1+\cdots+u_s)!}{u_1!\cdots u_s!}\leq s^{u_1+\cdots+u_s}.
\]
若像上面一样取 $u_j=\sum_{k=1}^rn_k\ell_{kj}$，则
$\sum_{j=1}^su_j\leq C\sum_{k=1}^rn_k$，其中适当的常数 $C$ 与各 $n_k$ 无关，因为 $n_1,\ldots,n_r\geq0$。级数的常数项显然为 $1$，故命题容易得到。
\end{proof}

这一命题表明，同一个级数~\eqref{eq:toric-euler-period} 在\emph{每个}最大幺幂边界点处都全纯。证明还提示了 GKZ 方程组与 GKZ 分解之间一个有趣的联系。更多细节见 \SourceCite{GKZ1}。

在最大幺幂边界点，Picard--Fuchs 方程只有一个全纯解，相差一个常数因子；这就是单值性权滤过满足 $\dim W_0=1$ 的含义。因此，命题~\ref{prop:toric-period-holomorphic} 描述了镜像映射公式~\eqref{eq:mirror-map-period-definition} 中的分母 $\langle g_0,\Omega\rangle$。为简便起见，以 $y_0$ 表示式~\eqref{eq:toric-period-local-series} 给出的级数。

为了求镜像映射的第 $k$ 个分量 $q_k$，需要找到满足式~\eqref{eq:toric-mirror-normalized-period} 的 Picard--Fuchs 方程解。不过，处理超几何函数时，去掉该式中的因子 $1/(2\pi i)$ 比较方便。因此，寻找满足
\begin{equation}\label{eq:hypergeometric-log-period}
y_k=y_0\log\bigl((-1)^{\ell_{k0}}z_k\bigr)+\widetilde y_k,
\qquad1\leq k\leq r,
\end{equation}
的解 $y_k$，其中 $\widetilde y_k$ 在原点全纯，且 $\widetilde y_k(0)=0$。在这一归一化下，$q_k$ 的公式成为
\begin{equation}\label{eq:hypergeometric-mirror-coordinate}
q_k=(-1)^{\ell_{k0}}z_k\exp(\widetilde y_k/y_0)
\sim(-1)^{\ell_{k0}}z_k,\qquad1\leq k\leq r,
\end{equation}
因为 $y_0$ 在原点非零，而 $\widetilde y_k$ 在原点为零。由式~\eqref{eq:toric-mirror-normalization}，这恰好符合要求。

求满足式~\eqref{eq:hypergeometric-log-period} 的 $y_k$，已有若干方法。先介绍 \SourceCite{HLY1} 中 Frobenius 方法的一个变体，它很适合\IndexLocation{indexloc:0615:184}超几何方程。其思想是，将式~\eqref{eq:toric-period-local-series} 中每个 $n_k$ 换成 $n_k+\rho_k$，其中 $\rho_k$ 为实参数。这样，诸如 $\bigl(\sum_{k=1}^rn_k\ell_{kj}\bigr)!$ 的阶乘，
\SourcePage{185}
换成经典的 Euler $\Gamma$ 函数
\[
\Gamma\left(\sum_{k=1}^r(n_k+\rho_k)\ell_{kj}+1\right).
\]
指数也随之改变，即 $\bigl((-1)^{\ell_{k0}}z_k\bigr)^{n_k}$ 变成
$\bigl((-1)^{\ell_{k0}}z_k\bigr)^{n_k+\rho_k}$。将这样由式~\eqref{eq:toric-period-local-series} 得到的级数记为 $y_0(z;\rho)$，其中 $z$ 代表 $z_1,\ldots,z_r$，$\rho$ 代表 $\rho_1,\ldots,\rho_r$。

若 $\partial_{\rho_k}=\partial/\partial\rho_k$，关键结果是
\[
y_k=\left.\partial_{\rho_k}y_0(z;\rho)\right|_{\rho=0},
\qquad1\leq k\leq r,
\]
具有所需的性质~\eqref{eq:hypergeometric-log-period}。所用的恒等式为
\[
\begin{aligned}
\partial_{\rho_k}\bigl((-1)^{\ell_{k0}}z_k\bigr)^{n_k+\rho_k}
&=\log\bigl((-1)^{\ell_{k0}}z_k\bigr)
\bigl((-1)^{\ell_{k0}}z_k\bigr)^{n_k+\rho_k},\\
[\Box_\ell,\partial_{\rho_k}]&=0.
\end{aligned}
\]
这里略去证明细节，改用几个例子说明实际计算。\index{Frobenius方法@Frobenius 方法|hyperpage}

\begingroup
\SourceExampleNumber{6.3.4}{1}
\renewcommand{\theHexam}{6.3.4.continued.\arabic{exam}}
\begin{example}[续]\label{ex:quintic-frobenius-continued}
回到前面开始讨论的五次镜像。式~\eqref{eq:quintic-euler-period} 给出的 $y_0(z)$ 意味着
\begin{equation}\label{eq:quintic-gamma-frobenius}
y_0(z;\rho)=\sum_{n=0}^\infty
\frac{\Gamma(5n+5\rho+1)}{\Gamma(n+\rho+1)^5}(-z)^{n+\rho}.
\end{equation}
于是 $y_0(z)=y_0(z;0)$，而上述理论说明，GKZ 方程的对数解为
\[
\begin{aligned}
y_1(z)&=\left.\partial_\rho y_0(z;\rho)\right|_{\rho=0}\\
&=y_0(z)\log(-z)+
\sum_{n=0}^\infty
\left.\partial_\rho\left(\frac{\Gamma(5n+5\rho+1)}
{\Gamma(n+\rho+1)^5}\right)\right|_{\rho=0}(-z)^n.
\end{aligned}
\]
为了化简，需要了解 $\Gamma$ 函数的导数。幸运的是，$\Gamma'(s)$ 在正整数处有以下公式：
\[
\Gamma'(n+1)=\Gamma(n+1)\left(-\gamma+\sum_{j=1}^n\frac1j\right),
\]
其中 $\gamma$ 为 Euler 常数。用归纳法及经典事实 $\Gamma'(1)=-\gamma$，可证明该式，见 \SourceCite[第 236 页]{WW}。由此容易得到
\[
y_1(z)=y_0(z)\log(-z)+5\sum_{n=1}^\infty
\frac{(5n)!}{(n!)^5}\left[\sum_{j=n+1}^{5n}\frac1j\right](-z)^n,
\]
式~\eqref{eq:hypergeometric-mirror-coordinate} 于是给出
\[
q=-z\exp\left(\frac5{y_0(z)}\sum_{n=1}^\infty
\frac{(5n)!}{(n!)^5}\left[\sum_{j=n+1}^{5n}\frac1j\right](-z)^n\right).
\]
这恰好是\hyperref[chap:quintic]{第~\ref*{chap:quintic} 章}式~\eqref{eq:y0-hypergeometric} 之后给出的五次三维簇镜像映射公式。

现在改用标准的 Frobenius 方法，它以递推关系为基础，类似于 \SourceCite{BvS} 的做法。为简化记号，令 $x=-z$。
\SourcePage{186}
采用这一变量，五次镜像族的 Picard--Fuchs 方程~\eqref{eq:quintic-pf} 可写为
\[
\delta^4y=5x(5\delta+1)(5\delta+2)(5\delta+3)(5\delta+4)y,
\]
其中 $\delta=x\,d/dx$，见式~\eqref{eq:quintic-gkz-fifth} 及随后的讨论。标准 Frobenius 方法 \SourceCite{CL} 使用函数
\[
\widetilde y(x;\rho)=\sum_{n=0}^\infty a_n(\rho)x^{n+\rho}.
\]
设 $L$ 为 Picard--Fuchs 方程给出的算子，需要选择 $a_n(\rho)$，使 $\widetilde y(x;\rho)$ 满足
\begin{equation}\label{eq:quintic-frobenius-inhomogeneous}
L(\widetilde y(x;\rho))=\rho^4x^\rho,
\end{equation}
因为 Picard--Fuchs 方程的\IndexLocation{indexloc:0621:186}指标方程为 $\rho^4=0$，对应\IndexLocation{indexloc:0772:186}最大幺幂单值性。对 $n>0$，得到递推关系
\[
(\rho+n)^4a_n(\rho)
=5(5\rho+5n-1)(5\rho+5n-2)(5\rho+5n-3)(5\rho+5n-4)a_{n-1}(\rho).
\]
选择 $a_0=1$ 的解，解出递推关系，得到
\[
a_n(\rho)=\frac{\prod_{m=1}^{5n}(5\rho+m)}
{\prod_{m=1}^n(\rho+m)^5}.
\]
于是
\begin{equation}\label{eq:quintic-product-frobenius}
\widetilde y(x;\rho)=\sum_{n=0}^\infty
\frac{\prod_{m=1}^{5n}(5\rho+m)}{\prod_{m=1}^n(\rho+m)^5}x^{n+\rho}.
\end{equation}
由式~\eqref{eq:quintic-frobenius-inhomogeneous} 容易得出，函数
\begin{equation}\label{eq:quintic-frobenius-basis}
y_i(x)=\left.\frac1{i!}\frac{\partial^i\widetilde y(x;\rho)}
{\partial\rho^i}\right|_{\rho=0},\qquad0\leq i\leq3,
\end{equation}
满足 Picard--Fuchs 方程，且给出 $x=0$ 附近解的一组基。注意，可把 $y_i(x)$ 看成 Taylor 级数
\begin{equation}\label{eq:quintic-frobenius-expansion}
\widetilde y(x;\rho)=y_0(x)+y_1(x)\rho+y_2(x)\rho^2+y_3(x)\rho^3+\cdots
\end{equation}
的系数。

读者可能会担心，这与前面使用 $\Gamma$ 函数的 Frobenius 方法不同。不过容易验证，$\widetilde y(x;\rho)$ 与式~\eqref{eq:quintic-gamma-frobenius} 定义的 $y_0(z;\rho)$ 满足
\[
y_0(z;\rho)=\widetilde y(-z;\rho)\frac{\Gamma(5\rho+1)}{\Gamma(\rho+1)^5}.
\]
经典公式，或 Mathematica，给出幂级数展开
\[
\frac{\Gamma(5\rho+1)}{\Gamma(\rho+1)^5}
=1+\frac{5\pi^2}{3}\rho^2+\cdots,
\]
从而
\[
\left.\frac{\partial\widetilde y(-z;\rho)}{\partial\rho}\right|_{\rho=0}
=\left.\frac{\partial y_0(z;\rho)}{\partial\rho}\right|_{\rho=0}.
\]
因此，两种 Frobenius 方法对对数解 $y_1$ 给出相同的公式。

更令人惊讶的是，从更形式化的角度看待这一计算时会发生什么。先将 $x$ 换成 $e^{t_1}$，于是 $\delta=x\,d/dx=d/dt_1$；再将 $\rho$ 换成
\SourcePage{187}
$\mathbb P^4$ 上同调环中的超平面类 $H$。同时将求和指标由 $n$ 改为 $d$，式~\eqref{eq:quintic-product-frobenius} 就成为函数
\[
\widetilde I=e^{t_1H}\sum_{d=0}^\infty
\frac{\prod_{m=1}^{5d}(5H+m)}{\prod_{m=1}^d(H+m)^5}e^{dt_1},
\]
这正是\hyperref[chap:hodge-yukawa]{第~\ref*{chap:hodge-yukawa} 章}式~\eqref{eq:quintic-tilde-i} 考察的函数。因此，从 Frobenius 方法看，乍看奇怪的上同调值函数 $\widetilde I$ 十分自然；感谢 \SourceCite{LLY} 的作者向我们解释这一点。按 $H^*(\mathbb P^4)$ 的通常基展开 $\widetilde I$，得到
\[
\widetilde I=y_0(x)+y_1(x)H+y_2(x)H^2+y_3(x)H^3+y_4(x)H^4.
\]
这就是在式~\eqref{eq:quintic-frobenius-expansion} 中用 $H$ 代替 $\rho$，因为 $H^5=0$。于是 $y_0,y_1,y_2,y_3$ 是式~\eqref{eq:quintic-frobenius-basis} 中的解，而 $y_4$ 不是。这证实了\hyperref[subsec:nef-gkz]{第~\ref*{subsec:nef-gkz} 小节}关于式~\eqref{eq:quintic-i-expansion} 所作的断言。进一步，若像式~\eqref{eq:quintic-iv-hypergeometric} 一样定义 $I_{\mathcal V}=5H\widetilde I$，则
\[
\begin{aligned}
I_{\mathcal V}
&=e^{t_1H}5H\sum_{d=0}^\infty
\frac{\prod_{m=1}^{5d}(5H+m)}{\prod_{m=1}^d(H+m)^5}e^{dt_1}\\
&=5H\bigl(y_0(x)+y_1(x)H+y_2(x)H^2+y_3(x)H^3\bigr),
\end{aligned}
\]
其中 $y_0,y_1,y_2,y_3$ 给出\IndexLocation{indexloc:1176:187}五次镜像 Picard--Fuchs 方程解的一组基。这些公式已经作为式~\eqref{eq:iv-simple} 和~\eqref{eq:iv-period-expansion} 出现在\hyperref[chap:quintic]{第~\ref*{chap:quintic} 章}中。用这种形式写出 Picard--Fuchs 方程的解，首先见于 \SourceCite{Givental1}。如\hyperref[subsec:nef-gkz]{第~\ref*{subsec:nef-gkz} 小节}末尾所述，$\widetilde I$ 在 \SourceCite{LLY} 中以 $\lim_{\lambda\to0}HG[\mathcal I(P)](t_1)$ 出现。上述求解超几何方程的讨论，也说明了为何采用“$HG$”记号。函数
$I_{\mathcal V}=\lim_{\lambda\to0}HG[\mathcal I(\widehat P)](t_1)$ 将在\hyperref[chap:mirror-theorem]{第~\ref*{chap:mirror-theorem} 章}镜像定理的处理中起重要作用。
\end{example}
\endgroup

\SourceExampleNumber{6.3.4}{2}
\begin{example}\label{ex:weighted-hypergeometric-mirror-map}
最后一个例子，回到 $\mathbb P(1,1,2,2,2)$ 的环面爆破中反典范\IndexLocation{indexloc:1020:187}\IndexLocation{indexloc:1029:187}超曲面的镜像 $V^\circ$。例题~\ref{ex:weighted-boundary-mirror-map} 在假定 $z_1=0$ 时计算了镜像映射。用\IndexLocation{indexloc:0616:187}超几何函数，现在能确定所有 $z_1,z_2$ 上的映射。例题~\ref{ex:octic-toric-mirror-theorem} 将再次讨论这一镜像映射，并进一步利用\hyperref[subsec:nef-gkz]{第~\ref*{subsec:nef-gkz} 小节}中本例的函数 $\widetilde I$。

解 $y_0(z_1,z_2)$ 很容易写出，因为格 $M_0$ 的生成元为矩阵~\eqref{eq:weighted-moduli-relations} 的两行
\[
(0,0,1,1,1,1,-4),\qquad(1,1,0,0,0,-2,0).
\]
于是
\[
(u_1,u_2,u_3,u_4,u_5,u_6,u_0)
=n_1(0,0,1,1,1,1,-4)+n_2(1,1,0,0,0,-2,0)
\]
给出式~\eqref{eq:toric-euler-period} 中的 $u_1=u_2=n_2$、
$u_3=u_4=u_5=n_1$、$u_6=n_1-2n_2$ 和 $u_0=-4n_1$，故式~\eqref{eq:toric-period-local-series} 给出级数
\[
y_0(z_1,z_2)=
\sum_{\substack{n_1,n_2\geq0\\n_1\geq2n_2}}
\frac{(4n_1)!}{(n_1!)^3(n_2!)^2(n_1-2n_2)!}z_1^{n_1}z_2^{n_2}.
\]
条件 $n_1\geq2n_2$ 表明 $y_0(0,z_2)=1$，与例题~\ref{ex:weighted-boundary-mirror-map} 的结果一致。

为了描述具有对数增长的解，按上述方式定义 $y_0(z_1,z_2;\rho_1,\rho_2)$，并计算
\[
y_k(z_1,z_2)=\left.\partial_{\rho_k}y_0(z_1,z_2;\rho_1,\rho_2)\right|_{\rho_1=\rho_2=0}.
\]
可以采用上一个例子的方法，
\SourcePage{188}
但必须适当处理条件 $n_1\geq2n_2$。这里微妙的一点是，$y_0(z_1,z_2)$ 实际上为级数
\[
y_0(z_1,z_2)=\sum_{n_1,n_2\geq0}
\frac{\Gamma(4n_1+1)}
{\Gamma(n_1+1)^3\Gamma(n_2+1)^2\Gamma(n_1-2n_2+1)}
z_1^{n_1}z_2^{n_2}.
\]
这与先前的级数相同，因为 Euler $\Gamma$ 函数在 $0,-1,-2,\ldots$ 处有单极点，其倒数在这些点为零。但由此构造 $y_0(z_1,z_2;\rho_1,\rho_2)$ 时，$n_1\geq2n_2$ 与 $n_1<2n_2$ 两类项都出现，对 $\rho_1,\rho_2$ 的导数中也都会出现。

用 $y_2=\left.\partial_{\rho_2}y_0\right|_{\rho_1=\rho_2=0}$ 说明这一过程。令 $\rho=\rho_2$，只需处理
\begin{equation}\label{eq:weighted-gamma-frobenius}
\sum_{n_1,n_2\geq0}
\frac{\Gamma(4n_1+1)}
{\Gamma(n_1+1)^3\Gamma(n_2+\rho+1)^2\Gamma(n_1-2n_2-2\rho+1)}
z_1^{n_1}z_2^{n_2+\rho}.
\end{equation}
对 $\rho$ 求导后置 $\rho=0$，得到 $y_0\log z_2+\widetilde y_2$，其中 $\widetilde y_2$ 来自式~\eqref{eq:weighted-gamma-frobenius} 的系数的导数。它有两类项，分别满足 $n_1\geq2n_2$ 和 $n_1<2n_2$。前者可用上一例题的方法处理，后者则需多做一点计算。使用 \SourceCite[第 239 页]{WW} 中的恒等式
\[
\Gamma(s)\Gamma(1-s)=\frac\pi{\sin(\pi s)},
\]
当 $n_1-2n_2<0$ 时，得到
\[
\left.\partial_\rho\bigl(\Gamma(n_1-2n_2-2\rho+1)^{-1}\bigr)\right|_{\rho=0}
=2(-1)^{n_1}(2n_2-n_1-1)!.
\]
因此，得到所需公式：
\[
\begin{aligned}
y_2(z_1,z_2)
&=y_0(z_1,z_2)\log z_2\\
&\quad+2\sum_{\substack{n_1,n_2\geq0\\n_1<2n_2}}
\frac{(4n_1)!(2n_2-n_1-1)!}{(n_1!)^3(n_2!)^2}
(-1)^{n_1}z_1^{n_1}z_2^{n_2}\\
&\quad+2\sum_{\substack{n_1,n_2\geq0\\n_1\geq2n_2}}
\frac{(4n_1)!}{(n_1!)^3(n_2!)^2(n_1-2n_2)!}
\left[\sum_{j=1}^{n_1-2n_2}\frac1j-\sum_{j=1}^{n_2}\frac1j\right]
z_1^{n_1}z_2^{n_2}.
\end{aligned}
\]
写成 $y_2=y_0\log z_2+\widetilde y_2$ 后，由式~\eqref{eq:hypergeometric-mirror-coordinate}，镜像映射有 $q_2=z_2\exp(\widetilde y_2/y_0)$。这给出了 $q_2$ 的显式公式，虽然相当复杂；也能为 $q_1$ 求出类似公式。至此，镜像映射已完全确定。

作为检验，像例题~\ref{ex:weighted-boundary-mirror-map} 一样置 $z_1=0$。所有 $n_1\geq2n_2$ 的求和项都消失；又因为 $y_0(0,z_2)=1$，得到
\begin{equation}\label{eq:weighted-boundary-log-series}
y_2(0,z_2)=\log z_2+
2\sum_{n_2=1}^\infty\frac{(2n_2-1)!}{(n_2!)^2}z_2^{n_2}.
\end{equation}
另一方面，例题~\ref{ex:weighted-boundary-mirror-map} 给出
\begin{equation}\label{eq:weighted-boundary-log-closed}
y_2(0,z_2)=\log\left(\frac{1-2z_2-\sqrt{1-4z_2}}{2z_2}\right).
\end{equation}
这是式~\eqref{eq:weighted-boundary-mirror-coordinate} 前一公式的表达式；因子 $1/(2\pi i)$ 消失，是因为本小节采用了不同的归一化方式。由式~\eqref{eq:weighted-boundary-period}，这两个公式相同，因为 $y_2(0,z_2)$ 的级数满足例题~\ref{ex:weighted-boundary-mirror-map} 的 Picard--Fuchs 方程
$\delta^2y-2z_2\delta(2\delta+1)y=0$，并具有相同的归一化。因此，虽然这里采用了与例题~\ref{ex:weighted-boundary-mirror-map} 很不同的方法求镜像映射，所得结果完全一致。
\end{example}

\SourcePage{189}
最后补充一点：上述两个例题介绍的镜像映射计算方法，已经用于大量环面超曲面，见 \SourceCite{HKTY1}。适当修改后，这一方法也已用于许多完全交，见 \SourceCite{BvS,HKTY2}。

\chapter{Gromov--Witten 不变量}\label{chap:gromov-witten}
\SourcePage{190}
如\hyperref[sec:three-point]{第~\ref*{sec:three-point} 节}所见，Calabi--Yau 三维簇 $V$ 上的 A 模型关联函数用不变量 $n_\beta$ 表示。朴素地说，$n_\beta$ 是 $V$ 上代表同调类 $\beta$ 的有理曲线的数目，尤其参见式~\eqref{eq:a-yukawa}。本章有两个目标：严格定义 $n_\beta$，并严格定义相关的 \emph{Gromov--Witten 不变量}。事实上，$n_\beta$ 将通过 Gromov--Witten 不变量来定义。\index{Gromov--Witten不变量@Gromov--Witten 不变量|hyperpage}

在以下一般情形中开展讨论。给定射影代数簇 $X$，固定同调类 $\beta\in H_2(X,\mathbb Z)$ 及 $X$ 上的圈 $Z_1,\ldots,Z_n$。基本问题是研究下列曲线集合的结构：
\begin{equation}\label{eq:gw-naive-curves}
C\subset X,\qquad C\text{ 的亏格为 }g,\quad[C]=\beta,\quad
C\cap Z_i\ne\varnothing\text{ 对所有 }i\text{ 成立},
\end{equation}
其中假定各 $Z_i$ 处于一般位置。我们将区分 \emph{Gromov--Witten 不变量}与 \emph{Gromov--Witten 类}：粗略地说，前者是式~\eqref{eq:gw-naive-curves} 中曲线只有有限条时的曲线数目，后者则是数目无限时得到的上同调类。\index{Gromov--Witten类@Gromov--Witten 类|hyperpage}

Gromov--Witten 不变量源于物理学中的\emph{与引力耦合的拓扑 $\sigma$ 模型}。特别地，亏格 $0$ 的 Gromov--Witten 不变量，有时称为\emph{树级}不变量，来自拓扑 $\sigma$ 模型，而后者是一种拓扑量子场论。尤其，Gromov--Witten 不变量的构造本来就要求它在 $X$ 的复结构形变下不变。\hyperref[sec:topological-qft]{附录 B.6} 简要介绍拓扑量子场论。

下面将看到，式~\eqref{eq:gw-naive-curves} 对 Gromov--Witten 不变量或类的描述实在过于朴素。要给出严格定义，需要相当精细的工具。代数几何与辛几何各提供了若干方法。第一个具有相当一般性的定义是在半正辛流形的框架中给出的：亏格 $0$ 的情形见 \SourceCite{RT1}，此前一些特殊情形已见于 \SourceCite{Ruan2}；更高亏格的情形见 \SourceCite{RT2}。此后，又出现了对一般辛流形的若干推广，如 \SourceCite{LTi3,Siebert1,Ruan3,FO}。在代数几何方面，早期首先为 Grassmann 簇构造了这些不变量 \SourceCite{BDW}，然后推广到齐性空间 \SourceCite{LTi1,FP}。一般的代数几何定义见 \SourceCite{LTi2} 以及 \SourceCite{BF,Behrend}。

\hyperref[sec:gw-algebraic-definition]{第~\ref*{sec:gw-algebraic-definition} 节}和\hyperref[sec:gw-symplectic-definition]{第~\ref*{sec:gw-symplectic-definition} 节}将介绍定义 Gromov--Witten 类的主要代数几何方法与辛几何方法。定义种类繁多，可能令人望而生畏；幸运的是，它们都满足\hyperref[subsec:gw-axioms]{第~\ref*{subsec:gw-axioms} 小节}陈述的 Gromov--Witten 类的基本公理。通常相信，各定义在共同的适用范围内彼此一致。\hyperref[subsec:gw-equivalence]{第~\ref*{subsec:gw-equivalence} 小节}将说明，某些情形中已经证明了这一点，但一般情形的等价性尚不清楚。因此，将所有这些对象都称作
\SourcePage{191}
“Gromov--Witten 不变量”，在某种意义上是对记号的宽松使用。幸运的是，只要所做的工作仅涉及 Gromov--Witten 公理，就没有问题；而在讨论各定义猜想中的等价性可能带来的推论时，我们会明确指出所使用的定义。

本章末尾将给出几个计算 Gromov--Witten 不变量的例子。接着，\hyperref[chap:quantum-cohomology]{第~\ref*{chap:quantum-cohomology} 章}将用这些不变量构造\emph{量子上同调环}。

文献中已有一些很好的 Gromov--Witten 理论入门资料。特别推荐 Voisin 的书 \SourceCite{Voisin3}、Morrison 的 Park City 讲义 \SourceCite{Morrison7}，以及 Fulton 和 Pandharipande 的 Santa Cruz 讲义 \SourceCite{FP}。

\section{代数几何中的定义}\label{sec:gw-algebraic-definition}
\index{Gromov--Witten类@Gromov--Witten 类!代数几何中的定义|hyperpage}
在 Gromov--Witten 理论的代数几何方法中，Kontsevich \SourceCite{Kontsevich2} 的关键观察是：式~\eqref{eq:gw-naive-curves} 中的曲线 $C\subset X$ 应当换成带 $n$ 个标记的曲线 $(C,p_1,\ldots,p_n)$ 与全纯映射 $f:C\to X$。这一观察与该概念在非线性 $\sigma$ 模型（\hyperref[sec:nonlinear-sigma-models]{附录 B.2}）以及与引力耦合的二维拓扑量子场论中的起源一致。于是，式~\eqref{eq:gw-naive-curves} 提示应研究映射
\begin{equation}\label{eq:gw-pointed-maps}
f:C\longrightarrow X,\qquad f_*[C]=\beta,\quad
f(p_i)\in Z_i,\quad i=1,\ldots,n.
\end{equation}
这些映射之间的同构概念是显然的，因而需要考虑它们的\emph{模}。为了得到紧的模空间，将允许某些亏格 $g$ 的可约曲线 $C$。带 $n$ 个标记的亏格 $g$ 稳定曲线的 Deligne--Mumford 粗模空间 $\overline M_{g,n}$，将在理论中发挥重要作用。\SourceCite{FP} 给出了关于 $\overline M_{g,n}$ 的很好介绍。

Gromov--Witten 类的粗略思想是：式~\eqref{eq:gw-pointed-maps} 中的映射给出 $\overline M_{g,n}$ 的一个子集，进而给出 $H^*(\overline M_{g,n},\mathbb Q)$ 中的上同调类。采用 $\mathbb Q$ 系数，是因为当 $n+2g\geq3$ 时，$\overline M_{g,n}$ 存在，并且是维数 $3g-3+n$ 的轨形。为了使这一表述更正式，令 $\alpha_i\in H^*(X)$ 为圈 $Z_i$ 的对偶上同调类。则 Gromov--Witten 类
\[
I_{g,n,\beta}(\alpha_1,\ldots,\alpha_n)\in H^*(\overline M_{g,n},\mathbb Q)
\]
直观上应是由式~\eqref{eq:gw-pointed-maps} 中出现的带标记曲线 $(C,p_1,\ldots,p_n)$ 的集合所代表的上同调类。因此，Gromov--Witten 类是一族映射
\begin{equation}\label{eq:gw-class-map}
I_{g,n,\beta}:H^*(X,\mathbb Q)^{\otimes n}
\longrightarrow H^*(\overline M_{g,n},\mathbb Q).
\end{equation}
\hyperref[sec:gw-properties]{第~\ref*{sec:gw-properties} 节}将研究这些映射的性质。

除了上同调类，还能定义数值不变量。具体如下：若 $I_{g,n,\beta}(\alpha_1,\ldots,\alpha_n)$ 如上所述，则得到有理数
\begin{equation}\label{eq:gw-numerical-invariant}
\langle I_{g,n,\beta}\rangle(\alpha_1,\ldots,\alpha_n)
=\int_{\overline M_{g,n}}I_{g,n,\beta}(\alpha_1,\ldots,\alpha_n).
\end{equation}
这就是一个 \emph{Gromov--Witten 不变量}。注意，除非 Gromov--Witten 类
$I_{g,n,\beta}(\alpha_1,\ldots,\alpha_n)$ 在 $H^*(\overline M_{g,n},\mathbb Q)$ 中含有最高次数的分量，式~\eqref{eq:gw-numerical-invariant} 为零。直观地，当式~\eqref{eq:gw-pointed-maps} 只含有限条曲线时，$I_{g,n,\beta}(\alpha_1,\ldots,\alpha_n)$ 应具有最高次数，此时曲线数目就等于
$\langle I_{g,n,\beta}\rangle(\alpha_1,\ldots,\alpha_n)$。不过，\hyperref[sec:computing-gw-i]{第~\ref*{sec:computing-gw-i} 节}将看到，Gromov--Witten 不变量可能是分数，甚至为负，因此它们的枚举意义并非总是直接明了。

\SourcePage{192}
到目前为止，所有讨论仍然极其朴素。要定义 Gromov--Witten 类，必须更仔细地研究映射~\eqref{eq:gw-pointed-maps} 的模。特别地，需要了解\emph{稳定映射}与\emph{虚拟基本类}。下面分别讨论这两个主题。

\subsection{稳定映射及其模}\label{subsec:stable-maps}
\index{稳定曲线|hyperpage}\index{带标记曲线|hyperpage}\index{稳定映射!模空间|hyperpage}
寻找 Gromov--Witten 类的代数定义时，\SourceCite{Kontsevich2} 的另一个关键思想，是构造映射~\eqref{eq:gw-pointed-maps} 的模空间的一个紧化。这通过\emph{稳定映射}实现，定义如下。\index{稳定映射|hyperpage}

\begin{definition}\label{def:stable-map}
带 $n$ 个标记的稳定映射，由连通的带标记曲线 $(C,p_1,\ldots,p_n)$ 与态射 $f:C\to X$ 组成，满足以下性质：
\begin{enumerate}[label=(\roman*)]
\item $C$ 的奇点只能是普通二重点。
\item $p_1,\ldots,p_n$ 是 $C$ 上互不相同、有次序的光滑点。
\item 若 $C_i$ 为 $C$ 的一个分支，$C_i\simeq\mathbb P^1$，且 $f$ 在 $C_i$ 上为常值，则 $C_i$ 至少含三个特殊点，即结点或标记点。
\item 若 $C$ 的算术亏格为 $1$ 且 $n=0$，则 $f$ 不是常值映射。
\end{enumerate}
\end{definition}
在定义的前两个条件成立的前提下，第三与第四个条件等价于：数据
$(f,C,p_1,\ldots,p_n)$ 只有有限多个自同构。无限多个自同构的存在，正是 $n+2g<3$ 时 $\overline M_{g,n}$ 不存在的原因。

如下定义稳定映射族。

\begin{definition}\label{def:stable-map-family}
设 $X$ 为射影代数簇，$S$ 为 $\mathbb C$ 上的概形。$S$ 上带 $n$ 个标记的稳定映射，是平坦的固有态射 $\mathcal C\to S$，连同 $n$ 个截面 $s_1,\ldots,s_n$ 及映射 $f:\mathcal C\to X$；对 $S$ 的每个几何点 $s$，$f$ 在 $\mathcal C$ 位于 $s$ 上的几何纤维上的限制 $f_s:\mathcal C_s\to X$，连同各截面 $s_i$ 的像，定义一个稳定映射。此外：
\begin{enumerate}[label=(\roman*)]
\item 若对 $S$ 的每个几何点 $s$，曲线 $\mathcal C_s$ 的算术亏格都为 $g$，则称 $f:\mathcal C\to X$ 的亏格为 $g$。
\item 给定同调类 $\beta\in H_2(X,\mathbb Z)$，若对 $S$ 的每个几何点 $s$，都有 $(f_s)_*[\mathcal C_s]=\beta$，则称 $f:\mathcal C\to X$ 的类为 $\beta$。
\end{enumerate}
\end{definition}
于是，给定 $\beta\in H_2(X,\mathbb Z)$，可定义反变函子
\[
\overline{\mathcal M}_{g,n}(X,\beta):
(\mathbb C\text{-概形})\longrightarrow(\text{集合}),
\]
将概形 $S$ 映为 $S$ 上所有亏格 $g$、类为 $\beta$、带 $n$ 个标记的稳定映射的同构类集合。

为了解决稳定映射的模问题，需要在某种意义下表示这一函子。最直接的方法，是考虑稳定映射 $f:C\to X$ 的\emph{粗模空间}。粗模空间的思想见 \SourceCite{Mumford1}。有以下基本存在性结果。\index{粗模空间|hyperpage}

\begin{theorem}\label{thm:stable-map-coarse-moduli}
若 $X$ 为射影簇，$\beta$ 为 $H_2(X,\mathbb Z)$ 中的同调类，则函子
$\overline{\mathcal M}_{g,n}(X,\beta)$ 有粗模空间 $\overline M_{g,n}(X,\beta)$，它是 $\mathbb C$ 上的射影概形。
\end{theorem}
这一结果首先由 Alexeev \SourceCite{Alexeev} 证明。\SourceCite{FP} 给出了 $\overline M_{g,n}(X,\beta)$ 的相当显式的构造。

模问题的一种更精细的方法涉及\emph{代数叠}。叠的语言见 \SourceCite{DM} 或 \SourceCite{BEFFGK}。
\SourcePage{193}
对本书的目的，只需指出，叠是从概形范畴到集合范畴、满足某些性质的反变函子；随着讨论展开，将明确说明其中一些性质。概形本身也是叠。对概形 $S$，可赋予由
\begin{equation}\label{eq:representable-scheme-functor}
h_S(X)=\operatorname{Hom}(X,S)
=\{\text{所有概形态射 }f:X\to S\}
\end{equation}
定义的可表示函子 $h_S$，其中 $X$ 为任意概形。\index{代数叠|hyperpage}

用这种语言，\SourceCite{Kontsevich2} 陈述了以下基本存在性结果。

\begin{theorem}\label{thm:stable-map-stack}
若 $X$ 为射影簇，$\beta$ 为 $H_2(X,\mathbb Z)$ 中的同调类，则
$\overline{\mathcal M}_{g,n}(X,\beta)$ 为 $\mathbb C$ 上的固有代数叠。此外，叠 $\overline{\mathcal M}_{0,n}(\mathbb P^r,\beta)$ 为光滑叠。
\end{theorem}
证明见 \SourceCite{BM}。由于这一结果，以下用 $\overline{\mathcal M}_{g,n}(X,\beta)$ 表示上述模叠。还要回忆，每个代数叠都有底层代数空间。对叠 $\overline{\mathcal M}_{g,n}(X,\beta)$，底层空间为粗模空间 $\overline M_{g,n}(X,\beta)$。有自然的叠态射
$\overline{\mathcal M}_{g,n}(X,\beta)\to\overline M_{g,n}(X,\beta)$，将 $S$ 上带 $n$ 个标记的稳定映射，映为其分类映射 $S\to\overline M_{g,n}(X,\beta)$。粗略地说，这一映射的存在表明，叠比粗模空间本身保留了更多信息。

叠是概形的推广，因此会经常用概形的语言讨论叠。我们希望读者能在几何上理解
$\overline{\mathcal M}_{g,n}(X,\beta)$，而不必深入涉及叠的技术问题；代价是引入一些不够精确的表述，希望不至于干扰阅读。想用叠的语言理解我们的确切意思的读者，需要稍多做一些工作。例如，有时把一个亏格 $g$、类为 $\beta$、带 $n$ 个标记的稳定映射
$f:(C,p_1,\ldots,p_n)\to X$ 称为 $\overline{\mathcal M}_{g,n}(X,\beta)$ 的一个\emph{点}。严格地说，函子并没有通常意义下的点，所以这一说法需要解释。按语境，它可能指关联于 $\overline{\mathcal M}_{g,n}(X,\beta)$ 的代数空间上的相应点，也可能指群胚
$\overline{\mathcal M}_{g,n}(X,\beta)(\operatorname{Spec}\mathbb C)$ 中的相应对象。

略说特殊情形 $g=0$、$X=\mathbb P^r$。此时可写 $\beta=d\ell$，其中 $\ell$ 为直线的类。由定理~\ref{thm:stable-map-stack}，
$\overline{\mathcal M}_{0,n}(\mathbb P^r,d\ell)$ 光滑，故其底层代数空间
$\overline M_{0,n}(\mathbb P^r,d\ell)$ 为轨形。可以证明，它是维数
$dr+r+d+n-3$ 的不可约正规射影簇。经常用 $\overline M_{g,n}(\mathbb P^r,d)$ 代替
$\overline M_{g,n}(\mathbb P^r,d\ell)$。

\subsection{从模空间到 Gromov--Witten 类}\label{subsec:moduli-to-gw}
\index{Gysin映射@Gysin 映射|hyperpage}\index{预期维数|hyperpage}
有了模空间 $\overline M_{g,n}(X,\beta)$，便得到以下自然映射：
\begin{equation}\label{eq:stable-map-evaluation-forgetting}
\begin{aligned}
\pi_1:\overline M_{g,n}(X,\beta)&\longrightarrow X^n,\\
\pi_2:\overline M_{g,n}(X,\beta)&\longrightarrow\overline M_{g,n}.
\end{aligned}
\end{equation}
这些映射容易定义。给定稳定映射 $f:(C,p_1,\ldots,p_n)\to X$，因为 $p_1,\ldots,p_n$ 是有次序的点列，得到良定义的 $n$ 元组
\[
(f(p_1),\ldots,f(p_n))\in X^n.
\]
这在集合层面给出了 $\pi_1$，不难验证它是态射。对于 $\pi_2$，先注意：若 $f:C\to X$ 为稳定映射，$C$ 并不必是 Deligne--Mumford 意义下的稳定曲线。
\SourcePage{194}
不过，若 $n+2g\geq3$，依次收缩 $C$ 的不稳定分支，就得到稳定曲线 $\widetilde C$。于是 $\pi_2$ 在集合层面将 $f:(C,p_1,\ldots,p_n)\to X$ 映为 $\widetilde C$ 的同构类。用 \SourceCite{Knudsen} 的技术可证明，它也是态射。

现在说明特殊情形下 Gromov--Witten 类的形式。假定 $X$ 光滑，且 $\overline M_{g,n}(X,\beta)$ 为维数
\begin{equation}\label{eq:stable-map-expected-dimension}
(1-g)(\dim X-3)-\int_\beta\omega_X+n
\end{equation}
的轨形，其中 $\omega_X$ 为 $X$ 的典范类。\hyperref[subsec:normal-cone]{第~\ref*{subsec:normal-cone} 小节}将说明，这是模空间的“预期维数”。在这些假设下，映射~\eqref{eq:stable-map-evaluation-forgetting} 给出自然映射
\[
\begin{aligned}
\pi_1^*:H^*(X,\mathbb Q)^{\otimes n}
&\longrightarrow H^*(\overline M_{g,n}(X,\beta),\mathbb Q),\\
\pi_{2*}:H_*(\overline M_{g,n}(X,\beta),\mathbb Q)
&\longrightarrow H_*(\overline M_{g,n},\mathbb Q),
\end{aligned}
\]
其中仍假定 $n+2g\geq3$。由于涉及的空间都是轨形，Poincaré 对偶（见\hyperref[app:singular-varieties]{附录 A}）与 $\pi_{2*}$ 诱导 Gysin 映射
\[
\pi_{2!}:H^*(\overline M_{g,n}(X,\beta),\mathbb Q)
\longrightarrow H^{2m+*}(\overline M_{g,n},\mathbb Q),
\]
其中 $m=(g-1)\dim X+\int_\beta\omega_X$。当 $n+2g\geq3$ 时，用公式
\begin{equation}\label{eq:gw-gysin-definition}
I_{g,n,\beta}(\alpha_1,\ldots,\alpha_n)
=\pi_{2!}\bigl(\pi_1^*(\alpha_1\otimes\cdots\otimes\alpha_n)\bigr)
\end{equation}
定义 Gromov--Witten 类。注意，这精确地定义了式~\eqref{eq:gw-pointed-maps} 中曲线集合的上同调类；回忆 $\alpha_i$ 是 $Z_i$ 的 Poincaré 对偶。容易验证，
$I_{g,n,\beta}(\alpha_1,\ldots,\alpha_n)$ 的次数为
\[
2m+\sum_{i=1}^n\deg\alpha_i
=2(g-1)\dim X+2\int_\beta\omega_X+\sum_{i=1}^n\deg\alpha_i.
\]
由式~\eqref{eq:gw-numerical-invariant}，还有 Gromov--Witten 不变量
\[
\langle I_{g,n,\beta}\rangle(\alpha_1,\ldots,\alpha_n)
=\int_{\overline M_{g,n}}I_{g,n,\beta}(\alpha_1,\ldots,\alpha_n),
\]
这里 $I_{g,n,\beta}(\alpha_1,\ldots,\alpha_n)$ 的维数须为 $2(3g-3+n)$，即 $\overline M_{g,n}$ 的实维数。利用式~\eqref{eq:gw-gysin-definition}，容易化简为
\begin{equation}\label{eq:gw-moduli-integral}
\langle I_{g,n,\beta}\rangle(\alpha_1,\ldots,\alpha_n)
=\int_{\overline M_{g,n}(X,\beta)}
\pi_1^*(\alpha_1\otimes\cdots\otimes\alpha_n).
\end{equation}
直观地，这恰好是式~\eqref{eq:gw-pointed-maps} 给出的数目。一个重要观察是：当 $n+2g<3$ 时，式~\eqref{eq:gw-moduli-integral} 仍有意义。因此，虽然 Gromov--Witten 类要求 $n+2g\geq3$，Gromov--Witten 不变量却对 $n,g\geq0$ 都有定义。\hyperref[chap:mirror-theorem]{第~\ref*{chap:mirror-theorem} 章}中，当语境已经明确 $n$ 时，经常采用简写
\[
\langle\alpha_1,\ldots,\alpha_n\rangle_{g,\beta}
=\langle I_{g,n,\beta}\rangle(\alpha_1,\ldots,\alpha_n).
\]

到目前为止，所有构造都假定 $\overline M_{g,n}(X,\beta)$ 为具有预期维数~\eqref{eq:stable-map-expected-dimension} 的光滑轨形。一个满足这些假设的良好例子是 $g=0$、$X=\mathbb P^r$。如上所述，
$\overline M_{0,n}(\mathbb P^r,d)$ 是维数 $rd+r+d+n-3$ 的轨形，而这正是本例的预期维数。完整细节见 \SourceCite{BM} 或 \SourceCite{FP}。

遗憾的是，定义~\eqref{eq:gw-gysin-definition} 和~\eqref{eq:gw-moduli-integral} 一般并不适用。例如，当 $g\geq1$ 时，甚至当 $g=0$ 而 $X\ne\mathbb P^r$ 的大多数流形情形中，空间 $\overline M_{g,n}(X,\beta)$
\SourcePage{195}
可能含有维数超过预期维数的分支。为了看清如何修正，需要强调\emph{基本类}的作用。上述情形中，有基本类
\[
\xi=[\overline M_{g,n}(X,\beta)]
\in H_*(\overline M_{g,n}(X,\beta),\mathbb Q),
\]
它在 Poincaré 对偶下对应于 $1\in H^*(\overline M_{g,n}(X,\beta),\mathbb Q)$。可用 $\xi$ 重写式~\eqref{eq:gw-gysin-definition}。首先，若
$e=\dim\overline M_{g,n}(X,\beta)$，则 Poincaré 对偶
$PD:H^*(\overline M_{g,n}(X,\beta),\mathbb Q)\to H_{2e-*}(\overline M_{g,n}(X,\beta),\mathbb Q)$ 为与 $\xi$ 的帽积。式~\eqref{eq:stable-map-evaluation-forgetting} 中的 $\pi_1,\pi_2$ 给出映射
\[
\pi:\overline M_{g,n}(X,\beta)\longrightarrow X^n\times\overline M_{g,n},
\]
满足 $\pi_i=p_i\circ\pi$，其中 $p_1,p_2$ 为自然投影。因此
\[
\begin{aligned}
\pi_{2!}\bigl(\pi_1^*(\alpha_1\otimes\cdots\otimes\alpha_n)\bigr)
&=\pi_{2!}PD^{-1}\bigl(\pi_1^*(\alpha_1\otimes\cdots\otimes\alpha_n)\cap\xi\bigr)\\
&=PD^{-1}\pi_{2*}\bigl(\pi_1^*(\alpha_1\otimes\cdots\otimes\alpha_n)\cap\xi\bigr)\\
&=PD^{-1}p_{2*}\pi_*\bigl(\pi^*p_1^*(\alpha_1\otimes\cdots\otimes\alpha_n)\cap\xi\bigr).
\end{aligned}
\]
最后两个等式中的 $PD$ 指 $\overline M_{g,n}$ 的 Poincaré 对偶。由投影公式及式~\eqref{eq:gw-gysin-definition}，Gromov--Witten 类为
\begin{equation}\label{eq:gw-fundamental-class-definition}
I_{g,n,\beta}(\alpha_1,\ldots,\alpha_n)
=PD^{-1}p_{2*}\bigl(p_1^*(\alpha_1\otimes\cdots\otimes\alpha_n)\cap\pi_*(\xi)\bigr),
\end{equation}
其中 $n+2g\geq3$。类似地，对 $n,g\geq0$，式~\eqref{eq:gw-moduli-integral} 中的 Gromov--Witten 不变量可用 $\xi$ 写为
\[
\langle I_{g,n,\beta}\rangle(\alpha_1,\ldots,\alpha_n)
=\int_\xi\pi_1^*(\alpha_1\otimes\cdots\otimes\alpha_n).
\]

后文中，用各分量表示式~\eqref{eq:stable-map-evaluation-forgetting} 中的 $\pi_1$ 会很方便。对 $1\leq i\leq n$，定义\emph{赋值映射}
\[
e_i:\overline{\mathcal M}_{g,n}(X,\beta)\longrightarrow X,
\]
将稳定映射 $f:(C,p_1,\ldots,p_n)\to X$ 映为 $f(p_i)$；等价地，它是 $\pi_1$ 与 $X^n$ 到第 $i$ 个因子的投影的复合。用此记号，Gromov--Witten 不变量为
\begin{equation}\label{eq:gw-evaluation-integral}
\langle I_{g,n,\beta}\rangle(\alpha_1,\ldots,\alpha_n)
=\int_\xi e_1^*(\alpha_1)\cup\cdots\cup e_n^*(\alpha_n).
\end{equation}
关键在于，只要有适当的“基本类”
$\xi\in H_*(\overline{\mathcal M}_{g,n}(X,\beta),\mathbb Q)$，公式~\eqref{eq:gw-fundamental-class-definition} 和~\eqref{eq:gw-evaluation-integral} 就有意义。具有所需性质的类 $\xi$ 通常称为\emph{虚拟基本类}。不过，很快将看到，其定义相当微妙。\index{虚拟基本类|hyperpage}

\subsection{法锥}\label{subsec:normal-cone}
\index{法锥!与Euler类的关系@与 Euler 类的关系|hyperpage}
先从\emph{法锥}说起，它是虚拟基本类的一个典型范例。设 $Y$ 为维数 $n$ 的光滑簇，$E$ 为其上的秩 $r$ 向量丛。给定截面 $s\in H^0(E)$，令 $Z=Z(s)\subset Y$ 为 $s$ 的零点概形。随着 $s$ 变化，$Z$ 的表现可能很糟糕，甚至维数也会改变。这与空间 $\overline{\mathcal M}_{g,n}(X,\beta)$ 所面临的问题类似。不过，本情形的处理方法是熟知的，见 \SourceCite{Fulton1}；更深入的处理见 \SourceCite{Fulton2}。
\SourcePage{196}
若 $\mathcal I$ 为 $Z$ 在 $Y$ 中的理想层，则 $Y$ 中 $Z$ 的\emph{法锥}定义为
\begin{equation}\label{eq:normal-cone-definition}
C_ZY=\operatorname{Spec}_Z\left(\bigoplus_{k=0}^\infty
\mathcal I^k/\mathcal I^{k+1}\right).
\end{equation}
法锥 $C_ZY$ 是 $Z$ 上纯维数为 $n$ 的仿射锥。与 $s$ 相乘给出的满射
$\mathcal O(E^*)\to\mathcal I$，诱导满射
\[
\bigoplus_k\operatorname{Sym}^k
\bigl(\mathcal O(E^*)/\mathcal I\mathcal O(E^*)\bigr)
\longrightarrow\bigoplus_k\mathcal I^k/\mathcal I^{k+1}.
\]
取 $\operatorname{Spec}$，得到 $C_ZY$ 到 $Z$ 上的丛 $E|_Z$ 的嵌入。法锥定义圈类
$[C_ZY]\in A_n(E|_Z)$。于是，类 $s^*[C_ZY]$ 位于 $A_{n-r}(Z)$。注意，若 $s$ 将 $Z$ 定义为局部完全交，$Z$ 的维数恰好为 $n-r$。因此，$n-r$ 为 $Z$ 的“预期维数”。关于 $s^*[C_ZY]$，还有以下熟知结果，可用 \SourceCite{Fulton2} 的结果与技术证明。\index{法锥|hyperpage}

\begin{lemma}\label{lem:normal-cone-euler-class}
若 $i:Z\hookrightarrow Y$ 为包含映射，则
$i_*(s^*[C_ZY])\in A_{n-r}(Y)$ 为 $E$ 的 Euler 类 $c_r(E)\cap[Y]$。
\end{lemma}
这一引理表明，$s^*[C_ZY]$ 是 Euler 类 $c_r(E)\cap[Y]$ 的精细化。随着截面 $s$ 变化，类 $s^*[C_ZY]$ 的行为也很好。还要指出，有时可把 $s^*[C_ZY]$ 算为关联的\emph{超额法丛}的 Euler 类。有时将用上同调而非 Chow 群开展工作；在这种情形中，秩 $r$ 丛 $E$ 的 Euler 类指最高 Chern 类 $c_r(E)\in H^{2r}(Y,\mathbb Z)$。这与拓扑学中的通行用法稍有不同，后者将 Euler 类写为同调类 $c_r(E)\cap[Y]$。

下面的例子说明法锥与稳定映射的关系。

\SourceExampleNumber{7.1.3}{1}
\begin{example}\label{ex:quintic-lines-normal-cone}
设 $V\subset\mathbb P^4$ 为光滑五次三维簇。由于一般 $V$ 恰好含 2875 条直线，Gromov--Witten 不变量的任何合理定义，都应在 $\ell$ 为 $V$ 中直线的类时满足
$\langle I_{0,0,\ell}\rangle=2875$。注意，$\langle I_{0,0,\ell}\rangle$ 是具有 $n=0$ 个变元的函数，也就是一个数。后面最终将证明这一点；目前先看法锥为何与本例相关。

首先注意，$\overline M_{0,0}(\mathbb P^4,\ell)$ 是 $\mathbb P^4$ 中直线的 Grassmann 簇 $G(2,5)$。因此，包含映射 $V\subset\mathbb P^4$ 诱导自然映射
$\overline M_{0,0}(V,\ell)\to\overline M_{0,0}(\mathbb P^4,\ell)=G(2,5)$，将稳定映射映为其在 $\mathbb P^4$ 中的像直线。还要注意 $\dim G(2,5)=6$。例题~\ref{ex:quintic-virtual-class} 将如下描述这一映射：若 $U$ 为 $G(2,5)$ 上的秩 $2$ 重言子丛，则直线 $\ell$ 处的纤维 $U_\ell$ 是 $\mathbb C^5$ 中射影化为 $\ell$ 的二维子空间。$V$ 的一个方程给出秩 $6$ 丛 $E=\operatorname{Sym}^5U^*$ 的截面 $s$，而 $\overline M_{0,0}(V,\ell)$ 正是 $s$ 的零点概形。

上述构造产生类 $s^*[C_ZY]\in A_0(\overline M_{0,0}(V,\ell))$，其中
$C_ZY$ 是 $\overline M_{0,0}(V,\ell)=Z(s)\subset Y=G(2,5)$ 的法锥。由下文的例题~\ref{ex:quintic-virtual-class}，这个零圈就是 $\overline M_{0,0}(V,\ell)$ 的虚拟基本类。这是合理的，因为 $V$ 的典范丛平凡，故由式~\eqref{eq:stable-map-expected-dimension}，
$\overline M_{0,0}(V,\ell)$ 的预期维数为 $0$，其虚拟基本类就应是零圈。

若 $V$ 为一般五次三维簇，可证明截面 $s$ 的零点轨迹为既约完全交，因而
$\overline M_{0,0}(V,\ell)$ 的实际维数确为 $0$。于是，一般情形中，
$\langle I_{0,0,\ell}\rangle$ 恰好是 $V$ 上直线的数目。
\SourcePage{197}
例题~\ref{ex:quintic-lines-localization} 将算出这一数目为 2875，与\hyperref[chap:quintic]{第~\ref*{chap:quintic} 章}一致。

另一方面，设 $V\subset\mathbb P^4$ 为由
\[
x_0^5+x_1^5+x_2^5+x_3^5+x_4^5=0
\]
定义的 Fermat 五次三维簇。这是含有无限多条直线的光滑 Calabi--Yau 三维簇。为看清这一点，令 $\zeta$ 为任意五次单位根。对任意 $i\ne j$，超平面
$x_i+\zeta x_j=0$ 与 $V$ 相交于一个以平面 Fermat 五次曲线为底的锥，给出 $V$ 中的一个单参数直线族。进一步可证明，按此方式得到的 50 个族包含 $V$ 上的全部直线，每条直线的法丛都是
$N=\mathcal O_{\mathbb P^1}(-3)\oplus\mathcal O_{\mathbb P^1}(1)$。于是，
$\overline M_{0,0}(V,\ell)$ 是一维的，有 50 个分支，每个分支都是平面五次曲线。此外，法丛满足 $h^0(N)=2$，由此可证明 $\overline M_{0,0}(V,\ell)$ 处处非既约。事实上，每个分支在一般点的重数为 $2$；而对应于同时属于不止一个锥的 375 条直线的 375 个特殊点，各点重数为 $5$。例如，由
$x_0+x_1=x_2+x_3=x_4=0$ 定义的直线便是这种情形。这些结果都在 \SourceCite{AKa} 中用形变理论证明。

虽然 $\overline M_{0,0}(V,\ell)$ 很复杂，其虚拟基本类仍是法锥构造给出的零圈
$s^*[C_ZY]\in A_0(\overline M_{0,0}(V,\ell))$，其中 $s$ 由 Fermat 五次方程诱导。
\SourceCite[例题 4.2]{ClK} 描述了法锥 $C=C_ZY$ 的结构，并给出更多细节。
在 $\overline M_{0,0}(V,\ell)$ 的 50 个分支各自上方，法锥 $C$ 都有一个纤维维数 $5$、重数 $2$ 的分支；在对应于同时属于不止一个锥的 375 条直线的各特殊点上方，则有一个纤维维数 $6$、重数 $5$ 的分支。因此，$C$ 的纯维数为 $6$，正如预期。再看虚拟基本类 $s^*[C]$，可证明：$C$ 中位于 $\overline M_{0,0}(V,\ell)$ 的一个分支上方的每个分支，对虚拟基本类贡献 $20$，这里需要计算平面五次曲线上\IndexLocation{indexloc:0280:197}超额法丛的 Euler 类；位于特殊点上方的每个分支则贡献 $5$。由式~\eqref{eq:gw-evaluation-integral}，Gromov--Witten 不变量
$\langle I_{0,0,\ell}\rangle$ 为这一零圈的次数，即
$20\times50+5\times375=2875$，正如预期。

可以验证，上述计算与 \SourceCite{AKa} 完全对应。事实上，我们的论述表明，这一对应来自动态相交理论，见 \SourceCite[第 11 章]{Fulton2}。虚拟基本类在某种意义上延伸了动态相交理论的适用范围。
\end{example}

\subsection{虚拟基本类}\label{subsec:virtual-fundamental-class}
\index{切—障碍复形|hyperpage}\index{Kuranishi映射@Kuranishi 映射|hyperpage}\index{障碍理论|hyperpage}\index{余切复形|hyperpage}\index{虚法锥|hyperpage}\index{预稳定曲线|hyperpage}\index{内蕴法层|hyperpage}
为了对任意 $g,n,X,\beta$ 定义 $\overline{\mathcal M}_{g,n}(X,\beta)$ 的虚拟基本类，需要比法锥构造更一般的工具。因此，不再把 $\overline{\mathcal M}_{g,n}(X,\beta)$ 实现为向量丛截面的零点概形，而是赋予它更多结构。在 \SourceCite{LTi2} 的方法中，这种结构称为\emph{完美切—障碍复形}；在 \SourceCite{BF} 的方法中，则称为\emph{完美障碍理论}。\index{完美障碍理论|hyperpage}

深入这些定义的细节前，先说明障碍为何与虚拟基本类有关。根本原因是：计算式~\eqref{eq:stable-map-expected-dimension}
给出的粗模空间 $\overline M_{g,n}(X,\beta)$ 的预期维数时，障碍自然出现。说明如下。

“预期维数”的思想常出现在代数几何中；也许在紧复流形
\SourcePage{198}
$M$ 的形变情形中最容易理解。无穷小形变由 $H^1(M,\Theta_M)$ 给出，障碍位于
$H^2(M,\Theta_M)$。Kuranishi 理论 \SourceCite{Kuranishi} 将 $M$ 的模空间局部描述为全纯障碍映射
\begin{equation}\label{eq:kuranishi-obstruction-map}
U\longrightarrow H^2(M,\Theta_M)
\end{equation}
的零点轨迹，其中 $U\subset H^1(M,\Theta_M)$ 为开集。因此，$M$ 的模空间的预期维数为
\[
\dim H^1(M,\Theta_M)-\dim H^2(M,\Theta_M).
\]

对 $\overline M_{g,n}(X,\beta)$，推理类似。令 $f:(C,p_1,\ldots,p_n)\to X$ 为稳定映射。叠
$\overline{\mathcal M}_{g,n}(X,\beta)$ 在 $f$ 处的切空间是超 $\operatorname{Ext}$ 群
\[
\operatorname{Ext}_C^1
\left(f^*\Omega_X^1\longrightarrow\Omega_C^1\left(\sum_{i=1}^np_i\right),
\mathcal O_C\right),
\]
而障碍位于
\[
\operatorname{Ext}_C^2
\left(f^*\Omega_X^1\longrightarrow\Omega_C^1\left(\sum_{i=1}^np_i\right),
\mathcal O_C\right).
\]
这里沿用 \SourceCite{LTi2}。\SourceCite{Kontsevich2} 则用相应的切形变复形，将切空间与障碍空间写成
\[
\mathbb H^1(C,\Theta'_C\longrightarrow f^*\Theta_X),
\qquad
\mathbb H^2(C,\Theta'_C\longrightarrow f^*\Theta_X),
\]
其中 $\Theta_X$ 为 $X$ 的切层，$\Theta'_C$ 为在 $p_1,\ldots,p_n$ 处为零的向量场所组成的层；在结点处应采用相应的导出形变复形。

将看到，式~\eqref{eq:kuranishi-obstruction-map} 在这一框架中有对应形式。因此，
$\overline M_{g,n}(X,\beta)$ 的预期维数为
\[
\begin{aligned}
&\dim\operatorname{Ext}_C^1
\left(f^*\Omega_X^1\to\Omega_C^1\left(\sum_{i=1}^np_i\right),\mathcal O_C\right)\\
&\qquad-\dim\operatorname{Ext}_C^2
\left(f^*\Omega_X^1\to\Omega_C^1\left(\sum_{i=1}^np_i\right),\mathcal O_C\right).
\end{aligned}
\]
考察 $\operatorname{Ext}$ 的长正合列，其开头为
\begin{equation}\label{eq:stable-map-ext-sequence}
\begin{aligned}
0&\longrightarrow
\operatorname{Ext}_C^0
\left(f^*\Omega_X^1\to\Omega_C^1\left(\sum_{i=1}^np_i\right),\mathcal O_C\right)\\
&\longrightarrow\operatorname{Ext}_C^0
\left(\Omega_C^1\left(\sum_{i=1}^np_i\right),\mathcal O_C\right)
\longrightarrow\operatorname{Ext}_C^0(f^*\Omega_X^1,\mathcal O_C)\\
&\longrightarrow\operatorname{Ext}_C^1
\left(f^*\Omega_X^1\to\Omega_C^1\left(\sum_{i=1}^np_i\right),\mathcal O_C\right)
\longrightarrow\cdots.
\end{aligned}
\end{equation}
注意，$\operatorname{Ext}_C^0(\Omega_C^1(\sum_{i=1}^np_i),\mathcal O_C)$ 为
$(C,p_1,\ldots,p_n)$ 的无穷小自同构空间；由于 $C$ 本身不必稳定，这些自同构可能存在。
此外，$\operatorname{Ext}_C^0(f^*\Omega_X^1,\mathcal O_C)$ 为固定 $C$ 时映射 $f$ 的一阶形变空间。
由 $f$ 的稳定性，可知连接这两项的映射是单射。因此，
$\operatorname{Ext}_C^0(f^*\Omega_X^1\to\Omega_C^1(\sum_{i=1}^np_i),\mathcal O_C)=0$。
于是，预期维数可改写为
\[
\chi(f^*\Theta_X)+
\dim\operatorname{Ext}_C^1\left(\Omega_C^1\left(\sum_{i=1}^np_i\right),\mathcal O_C\right)
-\dim\operatorname{Ext}_C^0\left(\Omega_C^1\left(\sum_{i=1}^np_i\right),\mathcal O_C\right),
\]
其中使用同构 $\operatorname{Ext}_C^i(f^*\Omega_X^1,\mathcal O_C)\simeq H^i(f^*\Theta_X)$。
Hirzebruch--\allowbreak Riemann--\allowbreak Roch 定理给出
\[
\chi(f^*\Theta_X)=-\int_\beta\omega_X+(1-g)\dim X.
\]
另外两个 $\operatorname{Ext}$ 项中，
\[
\operatorname{Ext}_C^1\left(\Omega_C^1\left(\sum_{i=1}^np_i\right),\mathcal O_C\right)
\]
为 $(C,p_1,\ldots,p_n)$ 的一阶形变空间，而
\[
\operatorname{Ext}_C^0\left(\Omega_C^1\left(\sum_{i=1}^np_i\right),\mathcal O_C\right)
\]
为其无穷小自同构空间。
回忆 $(C,p_1,\ldots,p_n)$ 本身不必稳定。因此，两者维数之差为带标记曲线模叠的维数；在稳定范围内，这就是 $\overline M_{g,n}$ 的轨形切空间的维数 $3g-3+n$。故预期维数为
\[
\begin{aligned}
-\int_\beta\omega_X+(1-g)\dim X+(3g-3+n)
&=(1-g)(\dim X-3)-\int_\beta\omega_X+n,
\end{aligned}
\]
与式~\eqref{eq:stable-map-expected-dimension} 一致。

\SourcePage{199}
现在可以精确定义虚拟基本类。如前所述，\SourceCite{LTi2} 与 \SourceCite{BF} 给出了两种方法。
预期二者在共同适用范围内定义相同的虚拟基本类。证明应当直接，但尚未见于文献。因此，有时会针对关于虚拟基本类的一个论断给出两个证明，各对应一种方法。此外，使用特定方法时，会明确说明采用哪一种。为了强调虚拟基本类的思想，所用记号不区分这些方法。

先介绍 Li 与 Tian 的方法，更多细节见 \SourceCite{LTi2}。令
\[
\mathcal F:(\mathbb C\text{-概形})\longrightarrow(\text{集合})
\]
为反变模函子，将概形 $S$ 映为 $S$ 上某种固定类型的几何对象的所有同构类的集合。
稳定映射的同构类函子 $\overline{\mathcal M}_{g,n}(X,\beta)$ 是其一例。

设 $S$ 为 $\mathbb C$ 上的仿射概形，$\mathcal N$ 为 $\mathcal O_S$ 模。用 $S_{\mathcal N}$ 表示 $S$ 由 $\mathcal N$ 给出的平凡扩张，即
$S_{\mathcal N}=\operatorname{Spec}(\Gamma(\mathcal O_S)\oplus\Gamma(\mathcal N))$，其中
$\Gamma(\mathcal O_S)\oplus\Gamma(\mathcal N)$ 是 $\Gamma(\mathcal O_S)$ 由 $\Gamma(\mathcal N)$ 给出的平凡环扩张。特别地，得到无穷小扩张 $S\hookrightarrow S_{\mathcal N}$。
$\mathcal F$ 的\emph{切函子} $T\mathcal F$ 是以下函子族。对如上所述的元素
$\alpha\in\mathcal F(S)$，有反变函子
\[
T\mathcal F(\alpha):(\mathcal O_S\text{-模})\longrightarrow(\text{集合}),
\]
将 $\mathcal N$ 映为 $\mathcal F(S_{\mathcal N})$ 中在自然“限制”映射
$\mathcal F(S_{\mathcal N})\to\mathcal F(S)$ 下限制为 $\alpha$ 的所有元素的集合。
更精确的描述见 \SourceCite{LTi2}。因此，切函子记录 $\alpha$ 到 $S$ 的各平凡无穷小扩张的所有“扩张”的数据。组成 $T\mathcal F$ 的各函子满足自然的基变换性质。

$\mathcal F$ 的\emph{切—障碍复形}由函子复形
\[
T^1\mathcal F\longrightarrow T^2\mathcal F
\]
组成，其中箭头为零自然变换，$T^1\mathcal F$ 为上述切函子，$T^2\mathcal F$ 给出
$\mathcal F$ 的适当障碍理论。这里不进一步展开障碍理论，只指出：对数据 $(\alpha,S,\mathcal N)$，有障碍类
\[
\operatorname{ob}(\alpha,S\hookrightarrow S',\mathcal N)
\in\Gamma\bigl(T^2\mathcal F(\alpha)\otimes_{\mathcal O_S}\mathcal N\bigr),
\]
其消失，是存在将 $\alpha$ 延拓的元素 $\widetilde\alpha\in\mathcal F(S')$ 的充要条件。不过，障碍理论还包含更多信息；更多细节见 \SourceCite{LTi2}。

若对上述每个 $(\alpha,S)$，至少局部于 $S$，存在局部自由 $\mathcal O_S$ 模层的两项复形
$\mathcal E^1\to\mathcal E^2$，使对任意 $\mathcal N$，
$T^i\mathcal F(\alpha)(\mathcal N)$ 为诱导复形
$\mathcal E^\bullet\otimes_{\mathcal O_S}\mathcal N$ 的第 $i$ 个层上同调，则称切—障碍复形为\emph{完美的}。

假定 $\mathcal F$ 可由概形 $Z$ 表示。按式~\eqref{eq:representable-scheme-functor} 的记号，即
$\mathcal F=h_Z$，其中 $1_Z\in h_Z(Z)=\operatorname{Hom}(Z,Z)$ 对应于泛对象
$\alpha\in\mathcal F(Z)$。为了易于说明，再假定 $\mathcal F$ 有完美切—障碍复形，且定义中的
$\mathcal E^i$ 都由整个 $Z$ 上的局部自由层诱导，而不仅仅局部存在。
\SourcePage{200}
仍用 $\mathcal E^i$ 表示这些层。下面的例题~\ref{ex:zero-scheme-obstruction-theory} 给出这种完美障碍理论的简单例子。
对 $i=1,2$，令 $E^i$ 为满足 $\mathcal O(E^i)=\mathcal E^i$ 的向量丛。
Li 与 Tian 用障碍理论构造一个远非唯一的 \emph{Kuranishi 映射}
\[
F:\widehat E^1\longrightarrow E^2,
\]
其中 $\widehat E^1$ 为 $E^1$ 沿零截面的形式完备化。这一 Kuranishi 映射 $F$ 是通常映射~\eqref{eq:kuranishi-obstruction-map} 的形式相对版本。
$F$ 的特征性质是：在每个 $z\in Z$ 上的纤维中，$F$ 的零点轨迹与 $Z$ 在 $z$ 处的完备化形式同构。
更精确地说，$F$ 本身的零点轨迹形式同构于 $Z\times Z$ 沿对角线的完备化 $\widehat Z$。
特别地，Kuranishi 映射给出 $\widehat Z$ 到 $\widehat E^1$ 的嵌入。

现在构造\emph{虚法锥}便很简单。考虑 $\widehat Z$ 在 $\widehat E^1$ 中的法锥，它是
$\widehat Z$ 上的锥。\emph{虚法锥} $C^{\mathcal E^\bullet}$ 是该锥在
$Z\subset\widehat Z$ 上的限制。注意，通过 $Z\subset\widehat Z$，$Z$ 作为 $E^1$ 的零截面嵌入。
虚法锥的维数等于 $E^1$ 的秩。再用 Kuranishi 映射，得到
$C^{\mathcal E^\bullet}$ 到 $E^2$ 的嵌入。\emph{虚拟基本类}定义为
$s^*[C^{\mathcal E^\bullet}]$，其中 $s$ 为 $E^2$ 的零截面。立即可见，它是 $Z$ 上维数
$\operatorname{rk}(E^1)-\operatorname{rk}(E^2)$ 的圈类。
这一差值可识别为 $Z$ 的预期维数；更准确地说，这也成为预期维数的一种定义。
由完美切—障碍复形的定义，显然这一差值完全由切—障碍复形决定，而与 $\mathcal E^i$ 的选择无关。
不那么显然、但由 \SourceCite{LTi2} 证明的是：虚拟基本类本身也与 $\mathcal E^i$ 和 Kuranishi 映射的选择无关。

\SourceExampleNumber{7.1.4}{1}
\begin{example}\label{ex:zero-scheme-obstruction-theory}
为了将这些内容具体化，再回到光滑概形 $Y$、向量丛 $E$ 与截面 $s\in H^0(E)$ 的情形。
令 $\mathcal F$ 为由 $Z=Z(s)$ 表示的函子。换言之，
\[
\mathcal F(S)=\{\text{所有满足 }f^*(s)=0\text{ 的态射 }f:S\to Y\}.
\]
给定 $\mathcal F(S)$ 的元素 $f:S\to Y$，我们断言：由
\[
f^*\Theta_Y\longrightarrow f^*\mathcal O(E)
\]
的核与余核定义的复形 $T^\bullet\mathcal F(f)$，满足完美切—障碍复形所需的条件。
上述复形中的映射，是截面 $s:Y\to E$ 的微分与自然投影 $T_E\to E$ 的复合；后者将 $E$ 的切丛沿零截面投影到零截面的法丛，而这个法丛同构于 $E$ 本身。
它显然由全局定义在 $Z$ 上的复形诱导，即
\begin{equation}\label{eq:zero-scheme-tangent-obstruction-complex}
\Theta_Y|_Z\longrightarrow\mathcal O(E)|_Z,
\end{equation}
这个复形充当上面的 $\mathcal E^\bullet$。按照 \SourceCite{LTi2} 的构造，虚法锥
$C^{\mathcal E^\bullet}$ 与法锥 $C_ZY$ 相同，其到 $E$ 的嵌入也正是刚才描述的嵌入。
于是，虚拟基本类就是 $s^*[C_ZY]$。这表明，虚拟基本类推广了\hyperref[subsec:normal-cone]{第~\ref*{subsec:normal-cone} 小节}的法锥构造。
\end{example}

到目前为止，只在函子 $\mathcal F$ 可由概形表示、且复形 $\mathcal E^\bullet$ 全局存在的特殊情形中构造了虚拟基本类。
\SourcePage{201}
\SourceCite{LTi2} 表明，这些假设可以放宽。特别地，对于
$\overline{\mathcal M}_{g,n}(X,\beta)$ 这种具有粗模空间的不可表示函子，也能应用这一构造定义虚拟基本类。

为了看清对 $\overline{\mathcal M}_{g,n}(X,\beta)$ 如何实现，考虑仿射概形 $S$ 上亏格 $g$、
带 $n$ 个标记的稳定映射族 $f:\mathcal C\to X$，见定义~\ref{def:stable-map-family}。
以 $D\subset\mathcal C$ 表示各截面 $s_i:S\to\mathcal C$ 的像的并。
Li 与 Tian 考察的切—障碍复形，将 $f$ 对应于函子复形
\begin{equation}\label{eq:stable-map-relative-obstruction-complex}
T^i\overline{\mathcal M}_{g,n}(X,\beta)(f)(\mathcal N)
=\operatorname{Ext}_{\mathcal C/S}^i
\left(f^*\Omega_X^1\longrightarrow\Omega_{\mathcal C/S}(D),
\mathcal O_{\mathcal C}\otimes_{\mathcal O_S}\mathcal N\right),
\end{equation}
其中 $i=1,2$。Li 与 Tian 证明，这是完美复形。因此，将
$\overline{\mathcal M}_{g,n}(X,\beta)$ 的\emph{虚拟基本类}定义为上述完美切—障碍复形的虚拟基本类。
这一类记作 $[\overline{\mathcal M}_{g,n}(X,\beta)]^{\mathrm{virt}}$，其维数为
$\operatorname{rk}(E^1)-\operatorname{rk}(E^2)$。由式~\eqref{eq:stable-map-ext-sequence} 紧前面的表达式，可见
$[\overline{\mathcal M}_{g,n}(X,\beta)]^{\mathrm{virt}}$ 具有预期维数~\eqref{eq:stable-map-expected-dimension}。

下面简要介绍 Behrend 与 Fantechi \SourceCite{BF} 的相关虚拟基本类构造。不会深入描述这一构造，因为这需要大幅转入叠的理论。叠的语言仍参见 \SourceCite{DM} 或 \SourceCite{BEFFGK}，代数叠的相交理论见 \SourceCite{Vistoli}。

Deligne--Mumford 叠是对角线非分歧的代数叠。\SourceCite{BF} 的构造始于以下观察：
任何 Deligne--Mumford 叠 $\mathcal M$ 都有一个\emph{内蕴法锥} $C_{\mathcal M}$。
它是 $\mathcal M$ 上纯维数为 $0$ 的叠。可如下局部描述内蕴法锥。
由于 $\mathcal M$ 有由概形给出的平展开覆盖，可以在一个这样的概形 $U$ 上工作。
选取嵌入 $U\hookrightarrow W$，其中 $W$ 光滑。令 $\mathcal I$ 为 $U$ 在 $W$ 中的理想层，
照常用式~\eqref{eq:normal-cone-definition} 的右端定义法锥 $C_UW$。微分映射
$\mathcal I\to\Omega_W^1$ 将 $f$ 映为 $df$，诱导映射
\[
\bigoplus_k\mathcal I^k/\mathcal I^{k+1}
\longrightarrow\bigoplus_k\operatorname{Sym}^k
(\Omega_W^1/\mathcal I\Omega_W^1).
\]
取 $\operatorname{Spec}$，得到映射
\[
T_W|_U\longrightarrow C_UW,
\]
因为 $T_W=\operatorname{Spec}(\bigoplus_{k=0}^\infty\operatorname{Sym}^k\Omega_W^1)$。
它给出一个作用，从而可取叠商；内蕴法锥 $C_U$ 定义为 $C_UW$ 关于 $T_W|_U$ 的叠商。
$C_U$ 与 $W$ 的选择无关，所以选取平展覆盖 $\mathcal M=\bigcup_iU_i$，其中各 $U_i$ 为概形时，
局部内蕴法锥 $C_{U_i}$ 能粘合成 $C_{\mathcal M}$。\index{内蕴法锥|hyperpage}

在上述构造中，把法锥 $C_UW$ 换成法层，也得到\emph{内蕴法层} $\mathcal N_{\mathcal M}$ 的构造。
事实上，$C_{\mathcal M}$ 嵌入 $\mathcal N_{\mathcal M}$。这一嵌入局部由满射
\[
\bigoplus_k\operatorname{Sym}^k(\mathcal I/\mathcal I^2)
\longrightarrow\bigoplus_k\mathcal I^k/\mathcal I^{k+1}
\]
诱导。

现在，$\mathcal M$ 有余切复形 $L^\bullet_{\mathcal M}$ \SourceCite{Illusie}。
粗略地说，$h^0(L^\bullet_{\mathcal M})$ 控制形变，$h^{-1}(L^\bullet_{\mathcal M})$ 控制障碍。
在 \SourceCite{BF} 中，$\mathcal M$ 的\emph{障碍理论}粗略指 $\mathcal M$ 上的层复形
$\mathcal E^\bullet$ 及态射 $\mathcal E^\bullet\to L^\bullet_{\mathcal M}$，
它在 $h^0$ 上为同构，在 $h^{-1}$ 上为满射。更精确的处理应在导出范畴中进行。
注意，这与 \SourceCite{LTi2} 中的切—障碍复形类似：
\SourcePage{202}
切空间 $T^1\mathcal F$ 固定，但障碍 $T^2\mathcal F$ 有一定的灵活性。

若 $\mathcal E^0\to\mathcal E^1$ 为 $\mathcal M$ 上 Abel 层的两项复形，可取
$\mathcal E^0$ 在 $\mathcal E^1$ 上的作用的叠商。上面构造内蕴法锥时，已经见过这种商的例子。
现在，若 $\mathcal E^\bullet$ 为任意长度的复形，按照 \SourceCite{BF}，
可将 $h^1/h^0(\mathcal E^\bullet)$ 定义为：$\mathcal E^1\to\mathcal E^2$ 的核，
关于 $\mathcal E^{-1}\to\mathcal E^0$ 的余核的叠商。

\begin{proposition}\label{prop:intrinsic-normal-sheaf}
\[
h^1/h^0\bigl((L^\bullet_{\mathcal M})^\vee\bigr)
\simeq\mathcal N_{\mathcal M}.
\]
\end{proposition}
\begin{proof}
这由 \SourceCite[第 66--69 页]{BF} 得到。符号 $(L^\bullet_{\mathcal M})^\vee$ 表示导出范畴中的“对偶”。
\end{proof}

现在可以构造虚拟基本类。若 $h^1/h^0((\mathcal E^\bullet)^\vee)$ 在 $\mathcal M$ 上光滑，
则障碍理论 $\mathcal E^\bullet$ 是\emph{完美的}。由命题~\ref{prop:intrinsic-normal-sheaf} 及完美障碍理论的定义，
$\mathcal N_{\mathcal M}$ 嵌入 $h^1/h^0((\mathcal E^\bullet)^\vee)$，因而 $C_{\mathcal M}$ 也嵌入其中。
令 $C$ 为 $(E^{-1})^*$ 与 $C_{\mathcal M}$ 在 $h^1/h^0((\mathcal E^\bullet)^\vee)$ 上的纤维积，其中
$\mathcal O(E^{-1})=\mathcal E^{-1}$。这是向量丛 $(E^{-1})^*$ 中的普通锥。
\emph{虚拟基本类}定义为 $C$ 与 $(E^{-1})^*$ 的零截面的交。

下面用 \SourceCite{BF} 的方法继续前面的例子。

\begingroup
\SourceExampleNumber{7.1.4}{1}
\renewcommand{\theHexam}{7.1.4.revisited.\arabic{exam}}
\begin{example}[再访]\label{ex:zero-scheme-bf-obstruction}
回到向量丛 $E$ 的截面 $s$ 的零点轨迹 $Z\subset Y$。令 $\mathcal E^\bullet$ 为
式~\eqref{eq:zero-scheme-tangent-obstruction-complex} 的对偶复形
$\mathcal O(E^*)|_Z\to\Omega_Y^1|_Z$，两项的次数分别为 $-1$ 和 $0$。
由于 $Y$ 光滑，有 $L_Y^\bullet\simeq\Omega_Y^1$。
截面 $s$ 给出的方程诱导 $\mathcal O(E^*)|_Z\to\mathcal I/\mathcal I^2$，
连同次数 $0$ 上 $\Omega_Y^1|_Z$ 的自然映射，得到 $\mathcal E^\bullet\to L_Z^\bullet$。
这使 $\mathcal E^\bullet$ 成为完美障碍理论。
于是，$h^1/h^0((\mathcal E^\bullet)^\vee)$ 是 $E|_Z$ 关于 $\Theta_Y|_Z$ 的叠商，
可与式~\eqref{eq:zero-scheme-tangent-obstruction-complex} 比较。
接着把虚拟基本类算为丛 $E|_Z$ 中锥 $C$ 与零截面的交。
由内蕴法锥的定义，$C=C_ZY$，到 $E|_Z$ 的嵌入也与本小节开头所给的嵌入相同。
因此，虚拟基本类确实等于 $s^*[C_ZY]$。
\end{example}
\endgroup

现在将这一构造用于 $\overline{\mathcal M}_{g,n}(X,\beta)$。
沿用 \SourceCite{Behrend}，定义虚拟基本类的一个相对版本，它更容易应用。
令 $\mathfrak M_{g,n}$ 为满足定义~\ref{def:stable-map} 中条件 (i)、(ii) 的带 $n$ 个标记的亏格 $g$ 曲线的叠，
这些曲线称为\emph{预稳定曲线}。于是有叠映射
\begin{equation}\label{eq:stable-map-prestable-forgetting}
F:\overline{\mathcal M}_{g,n}(X,\beta)\longrightarrow\mathfrak M_{g,n},
\end{equation}
它忘掉映射，但不作稳定化。上述构造的相对版本表明，$F$ 有内蕴法锥 $C_F$ 与法层 $\mathcal N_F$。
再考虑叠映射
\[
\pi_{n+1}:\overline{\mathcal M}_{g,n+1}(X,\beta)
\longrightarrow\overline{\mathcal M}_{g,n}(X,\beta),
\]
它忘掉最后一个点 $p_{n+1}$，并收缩因此出现的任何不稳定分支。
如\hyperref[subsec:gravitational-correlator-definition]{第~\ref*{subsec:gravitational-correlator-definition} 小节}所说明、并由 \SourceCite{BM} 证明，泛\IndexLocation{indexloc:1282:202}稳定映射在叠的意义下由 $\pi_{n+1}$ 与\IndexLocation{indexloc:0276:202}赋值映射
\[
e_{n+1}:\overline{\mathcal M}_{g,n+1}(X,\beta)\longrightarrow X
\]
组成。

\SourcePage{203}
所需完美障碍理论为
$\mathcal E^\bullet=(R\pi_{n+1*}e_{n+1}^*\Theta_X)^\vee$。
于是，$C_F$ 嵌入
\[
h^1/h^0((\mathcal E^\bullet)^\vee)=h^1/h^0(R\pi_{n+1*}e_{n+1}^*\Theta_X).
\]
同样，可将它拉回为普通向量丛中的普通锥，与零截面相交，得到\emph{相对虚拟基本类}
\[
[\overline{\mathcal M}_{g,n}(X,\beta)/\mathfrak M_{g,n},
R\pi_{n+1*}e_{n+1}^*\Theta_X]
\]
位于双变 Chow 群
$A^*(\overline{\mathcal M}_{g,n}(X,\beta)\to\mathfrak M_{g,n})$ 中。
双变 Chow 群见 \SourceCite{Fulton2} 与 \SourceCite{Vistoli}。
最后，将 $\overline{\mathcal M}_{g,n}(X,\beta)$ 的\emph{虚拟基本类}定义为
\[
[\overline{\mathcal M}_{g,n}(X,\beta)]^{\mathrm{virt}}
=[\mathfrak M_{g,n}]\cap
[\overline{\mathcal M}_{g,n}(X,\beta)/\mathfrak M_{g,n},
R\pi_{n+1*}e_{n+1}^*\Theta_X],
\]
它位于 $A_*(\overline{\mathcal M}_{g,n}(X,\beta))$ 中。

注意，上述各种虚拟基本类定义都位于 Chow 群或双变 Chow 群中。
本章余下部分经常将虚拟基本类与其关联的同调类识别，这一过程会舍弃部分信息。
这样做是为了与\hyperref[sec:gw-symplectic-definition]{第~\ref*{sec:gw-symplectic-definition} 节}讨论的辛几何版本的虚拟基本类比较，
也为了在\hyperref[chap:localization]{第~\ref*{chap:localization} 章}及后续章节中避免使用等变 Chow 群。

\subsection{计算虚拟基本类}\label{subsec:computing-virtual-class}
\index{虚拟基本类!计算|hyperpage}\index{Euler类@Euler 类!虚拟基本类|hyperpage}
这里不穷尽虚拟基本类的计算与运算，而是介绍四种理解得较为清楚的情形。
首先是一种特别容易求出 $[\overline M_{0,n}(X,\beta)]^{\mathrm{virt}}$ 的情况。

\begin{definition}\label{def:convex-variety}
\SourceCite{KoM1} 若光滑射影簇 $X$ 对所有亏格 $0$ 的稳定映射 $f:C\to X$，都满足
$H^1(C,f^*T_X)=0$，则称 $X$ 为\emph{凸的}。
\end{definition}
\index{凸簇|hyperpage}
当 $X$ 凸时，可证明 $\overline M_{0,n}(X,\beta)$ 是具有预期维数的轨形 \SourceCite{FP}。
特别地，有基本类 $[\overline M_{0,n}(X,\beta)]\in H_*(\overline M_{0,n}(X,\beta),\mathbb Q)$。
不过，凸性还意味着到 $X$ 的亏格 $0$ 稳定映射的模函子无障碍。
使用上述任一种方法，都化为平凡丛中的平凡相交，故本情形中
$[\overline M_{0,n}(X,\beta)]^{\mathrm{virt}}=[\overline M_{0,n}(X,\beta)]$。
射影空间以及更一般的齐性空间，都是凸的。

\hyperref[subsec:moduli-to-gw]{第~\ref*{subsec:moduli-to-gw} 小节}讨论了 $\overline M_{g,n}(X,\beta)$ 是具有预期维数的轨形的情形。
对这样的 $X$，仍有
\[
[\overline M_{g,n}(X,\beta)]^{\mathrm{virt}}=[\overline M_{g,n}(X,\beta)].
\]
原因是：用 \SourceCite{LTi2} 的语言，切—\IndexLocation{indexloc:1328:203}障碍复形由关联模叠
$\overline{\mathcal M}_{g,n}(X,\beta)$ 的切层诱导；用 \SourceCite{BF} 的语言，
完美障碍理论由叠的余切层诱导。与上一段类似，得到所需等式。

接着计算五次三维簇的一个虚拟基本类。

\SourceExampleNumber{7.1.5}{1}
\begin{example}\label{ex:quintic-virtual-class}
设 $V\subset\mathbb P^4$ 为光滑五次三维簇，$d>0$。本例目标是构造虚拟基本类
$[\overline M_{0,0}(V,d\ell)]^{\mathrm{virt}}$。
先证明自然映射 $\overline M_{0,0}(V,d\ell)\to\overline M_{0,0}(\mathbb P^4,d)$ 为嵌入。
为此，令 $s$ 为 $\mathcal O_{\mathbb P^4}(5)$ 上在 $V$ 上消失的截面，考虑函子
\begin{equation}\label{eq:quintic-stable-map-zero-functor}
F(S)=\{f:\mathcal C\to\mathbb P^4:
f\text{ 为 }S\text{ 上亏格 }0\text{、次数 }d\text{ 的稳定映射，且 }f^*(s)=0\}/\sim,
\end{equation}
其中 $\sim$ 表示 $S$ 上稳定映射的同构。这是关联于
$\overline M_{0,0}(V,d\ell)$ 的函子，因此得到嵌入
$\overline M_{0,0}(V,d\ell)\hookrightarrow\overline M_{0,0}(\mathbb P^4,d)$。

\SourcePage{204}
可如下描述这一嵌入。令 $\mathcal V_d$ 为叠
$\overline{\mathcal M}_{0,0}(\mathbb P^4,d)$ 上的秩 $5d+1$ 向量丛，
其在稳定映射 $f:C\to\mathbb P^4$ 处的纤维为 $H^0(C,f^*\mathcal O_{\mathbb P^4}(5))$。
定义 $V$ 的截面 $s$ 确定 $\overline{\mathcal M}_{0,0}(\mathbb P^4,d)$ 上
$\mathcal V_d$ 的截面 $\widetilde s$；由式~\eqref{eq:quintic-stable-map-zero-functor}，
$\overline{\mathcal M}_{0,0}(V,d\ell)$ 是 $\widetilde s$ 的零点轨迹。
这里使用叠，是因为 $\mathcal V_d$ 未必是轨形 $\overline M_{0,0}(\mathbb P^4,d)$ 上的向量丛。
这样表示 $\overline{\mathcal M}_{0,0}(V,d\ell)$，说明我们几乎处于\hyperref[subsec:normal-cone]{第~\ref*{subsec:normal-cone} 小节}的情形：
有向量丛 $E=\mathcal V_d$，底空间 $Y=\overline{\mathcal M}_{0,0}(\mathbb P^4,d)$ 光滑，
这里是叠意义下的光滑，而非作为簇光滑，且 $Z(\widetilde s)=\overline{\mathcal M}_{0,0}(V,d\ell)$
是截面 $\widetilde s$ 的零点轨迹。于是，法锥构造给出 $\widetilde s^*[C_ZY]$。
它是零圈，因为 $\mathcal V_d$ 的秩为 $5d+1$，由式~\eqref{eq:stable-map-expected-dimension}，
$\overline M_{0,0}(\mathbb P^4,d)$ 的维数也是 $5d+1$。

我们断言，$\widetilde s^*[C_ZY]$ 就是虚拟基本类
$[\overline M_{0,0}(V,d\ell)]^{\mathrm{virt}}$。采用 Li 与 Tian 的方法证明。
由例题~\ref{ex:zero-scheme-obstruction-theory}，$\widetilde s^*[C_ZY]$ 为下列映射的核与余核
给出的完美切—障碍复形的虚拟基本类，参见式~\eqref{eq:zero-scheme-tangent-obstruction-complex}：
\begin{equation}\label{eq:quintic-zero-obstruction-complex}
\Theta_{\overline{\mathcal M}_{0,0}(\mathbb P^4,d)}
\big|_{\overline{\mathcal M}_{0,0}(V,d\ell)}
\longrightarrow
\mathcal V_d\big|_{\overline{\mathcal M}_{0,0}(V,d\ell)}.
\end{equation}
因此，只需证明式~\eqref{eq:quintic-zero-obstruction-complex} 的核与余核，与
式~\eqref{eq:stable-map-relative-obstruction-complex} 为
$\overline{\mathcal M}_{0,0}(V,d\ell)$ 定义的切—障碍复形一致。

为此，令 $f:\mathcal C\to V$ 为概形 $S$ 上的稳定映射，$\pi:\mathcal C\to S$ 为结构映射。
考虑复形的短正合列
\[
\begin{array}{ccccccccc}
0&\longrightarrow&f^*\mathcal O_{\mathbb P^4}(-5)&\longrightarrow&
f^*\Omega_{\mathbb P^4}^1&\longrightarrow&f^*\Omega_V^1&\longrightarrow&0\\
&&\downarrow&&\downarrow&&\downarrow&&\\
0&\longrightarrow&0&\longrightarrow&
\Omega_{\mathcal C/S}^1&=&\Omega_{\mathcal C/S}^1&\longrightarrow&0.
\end{array}
\]
再考察相对 $\operatorname{Ext}$ 的长正合列：
\begin{equation}\label{eq:quintic-relative-ext-sequence}
\begin{aligned}
0&\longrightarrow
\operatorname{Ext}_{\mathcal C/S}^1(f^*\Omega_V^1\to\Omega_{\mathcal C/S}^1,\mathcal O_{\mathcal C})\\
&\longrightarrow
\operatorname{Ext}_{\mathcal C/S}^1(f^*\Omega_{\mathbb P^4}^1\to\Omega_{\mathcal C/S}^1,\mathcal O_{\mathcal C})\\
&\longrightarrow
\operatorname{Ext}_{\mathcal C/S}^0(f^*\mathcal O_{\mathbb P^4}(-5),\mathcal O_{\mathcal C})\\
&\longrightarrow
\operatorname{Ext}_{\mathcal C/S}^2(f^*\Omega_V^1\to\Omega_{\mathcal C/S}^1,\mathcal O_{\mathcal C})
\longrightarrow0.
\end{aligned}
\end{equation}
末尾的 $\to0$ 来自 $\mathbb P^4$ 的凸性。

考察式~\eqref{eq:quintic-relative-ext-sequence} 的非零项。第一项是
$\overline{\mathcal M}_{0,0}(V,d\ell)$ 的切空间拉回到 $S$，
第二项是 $\overline{\mathcal M}_{0,0}(\mathbb P^4,d)$ 的切空间拉回到 $S$，
第三项是 $\mathcal V_d$ 拉回到 $S$，第四项为障碍空间 $T^2$。
由式~\eqref{eq:quintic-relative-ext-sequence} 的正合性，可将
式~\eqref{eq:quintic-zero-obstruction-complex} 的核识别为
$\operatorname{Ext}_{\mathcal C/S}^1(f^*\Omega_V^1\to\Omega_{\mathcal C/S}^1,\mathcal O_{\mathcal C})$，
将余核识别为
$\operatorname{Ext}_{\mathcal C/S}^2(f^*\Omega_V^1\to\Omega_{\mathcal C/S}^1,\mathcal O_{\mathcal C})$。
这正是式~\eqref{eq:stable-map-relative-obstruction-complex} 的切—障碍复形，断言得证。

细心的读者会发现，这一断言稍微宽松地使用了记号。丛 $\mathcal V_d$ 及其 Euler 类
仅定义在叠 $\overline{\mathcal M}_{0,0}(\mathbb P^4,d)$ 上，而非空间 $\overline M_{0,0}(\mathbb P^4,d)$ 上，
但我们却在粗模空间 $\overline M_{0,0}(V,d\ell)$ 上定义虚拟基本类。
这是因为一旦在叠上得到 Euler 类，就能将其推到粗模空间，而不必要求
$\overline M_{0,0}(V,d)$ 上存在相关的丛。全书中将继续采用这种宽松记法。

等式 $[\overline M_{0,0}(V,d\ell)]^{\mathrm{virt}}=\widetilde s^*[C_ZY]$ 有一个很好的推论。
令 $i:\overline{\mathcal M}_{0,0}(V,d\ell)\hookrightarrow
\overline{\mathcal M}_{0,0}(\mathbb P^4,d)$ 为上述嵌入。引理~\ref{lem:normal-cone-euler-class} 蕴含
\begin{equation}\label{eq:quintic-virtual-class-pushforward}
i_*([\overline{\mathcal M}_{0,0}(V,d\ell)]^{\mathrm{virt}})
=i_*(\widetilde s^*[C_ZY])
=c_{5d+1}(\mathcal V_d)\cap[\overline{\mathcal M}_{0,0}(\mathbb P^4,d)].
\end{equation}
\SourcePage{205}
这一等式成立于 $\overline{\mathcal M}_{0,0}(\mathbb P^4,d)$ 上；
推到 $\overline M_{0,0}(\mathbb P^4,d)$ 后，得到代表丛 $\mathcal V_d$ 的 Euler 类的零圈。
由此将得到一个很好的 Gromov--Witten 不变量公式，例题~\ref{ex:quintic-gw-euler-integral} 将研究它。

按例题~\ref{ex:quintic-lines-normal-cone} 的记号，$G(2,5)$ 上的丛
$\operatorname{Sym}^5U^*$ 就是这里考虑的 $\overline M_{0,0}(\mathbb P^4,1)$ 上的丛
$\mathcal V_1$。因此，前一例题的嵌入
$\overline M_{0,0}(V,\ell)\hookrightarrow G(2,5)$，由上述
$\overline{\mathcal M}_{0,0}(V,d\ell)$ 的描述得到。
\end{example}

计算虚拟基本类的第三个方面涉及“超额维数”。在相交理论中，当一个交的维数大于预期维数时，
就出现超额维数。在某些良好情形下，可用所谓超额法丛计算相交类。
清楚的讨论见 \SourceCite[第 6.3 节]{Fulton2}。

“超额维数”的思想也适用于当前情形。将
$\overline M_{g,n}(X,\beta)$ 的\emph{超额维数}定义为实际维数
$\dim\overline M_{g,n}(X,\beta)$ 与预期维数~\eqref{eq:stable-map-expected-dimension} 之差。
若超额维数为正，等式
$[\overline M_{g,n}(X,\beta)]^{\mathrm{virt}}=[\overline M_{g,n}(X,\beta)]$ 显然不可能成立，
但有时仍能用 Euler 类计算虚拟基本类。
若 $F:\overline{\mathcal M}_{g,n}\to\mathfrak M_{g,n}$ 为式~\eqref{eq:stable-map-prestable-forgetting}
中的叠映射，则有以下结果。\index{超额维数|hyperpage}

\begin{proposition}\label{prop:excess-obstruction-euler-class}
若定义 $[\overline M_{g,n}(X,\beta)]^{\mathrm{virt}}$ 时使用的层
$R^1\pi_{n+1*}e_{n+1}^*\Theta_X$ 局部自由，其秩 $e$ 等于超额维数，则
\[
[\overline M_{g,n}(X,\beta)]^{\mathrm{virt}}
=c_e(R^1\pi_{n+1*}e_{n+1}^*\Theta_X)
\cap\bigl([\overline{\mathcal M}_{g,n}(X,\beta)/\mathfrak M_{g,n}]
\cap[\mathfrak M_{g,n}]\bigr),
\]
其中 $[\overline{\mathcal M}_{g,n}(X,\beta)/\mathfrak M_{g,n}]
\in A^*(\overline{\mathcal M}_{g,n}(X,\beta)\to\mathfrak M_{g,n})$
是 \SourceCite{Fulton2} 意义下的基本类。
\end{proposition}
\begin{proof}
见 \SourceCite{Behrend,Getzler1}。
\end{proof}

讨论亏格 $1$ 的 Gromov--Witten 不变量时，将在例题~\ref{ex:plane-genus-one} 使用这一命题。
注意，即使 $\overline M_{g,n}(X,\beta)$ 光滑，命题~\ref{prop:excess-obstruction-euler-class} 的假设也未必成立。
原因在于：$\overline M_{g,n}(X,\beta)$ 中映射的障碍并不位于 $H^1(C,f^*\Theta_X)$ 本身，
而是位于它的商
$\operatorname{Ext}_C^2(f^*\Omega_X^1\to\Omega_C^1(\sum_i p_i),\mathcal O_C)$；
这一空间处理起来更微妙。\hyperref[subsec:nodal-quintic-double-covers]{第~\ref*{subsec:nodal-quintic-double-covers} 小节}将给出说明。

第四个、也是最后一个观察，涉及 $\overline M_{g,n}(X,\beta)$ 与
$\overline M_{g,n-1}(X,\beta)$ 的虚拟基本类之间的相容性。考虑交换图
\begin{figure}[H]
\centering
\includegraphics[width=0.568\textwidth]{7-1-forgetting-diagram.png}
\end{figure}
其中 $\phi_i$ 忘掉第一个标记点并作稳定化，$\pi_i$ 忘掉最后一个标记点，
$e_i:\overline{\mathcal M}_{g,i}(X,\beta)\to X$ 在最后一个点处赋值。
由基变换，
\[
R^1\pi_{n+1*}e_{n+1}^*\Theta_X
\simeq\phi_n^*R^1\pi_{n*}e_n^*\Theta_X.
\]
\SourcePage{206}
由此得到
\begin{equation}\label{eq:virtual-class-forgetting-compatibility}
[\overline M_{g,n}(X,\beta)]^{\mathrm{virt}}
=\phi_n^*[\overline M_{g,n-1}(X,\beta)]^{\mathrm{virt}}.
\end{equation}
换言之，第 $n$ 个标记点的形变理论无障碍，见 \SourceCite[公理 IV]{Behrend}。
相容性~\eqref{eq:virtual-class-forgetting-compatibility} 将在\hyperref[chap:quantum-differential-equations]{第~\ref*{chap:quantum-differential-equations} 章}中有用。
应指出，忘却映射 $\phi_n:\overline M_{g,n}(X,\beta)\to\overline M_{g,n-1}(X,\beta)$ 并非总存在。
问题是，当 $\beta=0$ 时，忘掉一个标记点并作稳定化，可能使整条曲线消失。
例如，$\phi_3:\overline M_{0,3}(X,0)\to\overline M_{0,2}(X,0)$ 不存在，
因为 $\overline M_{0,3}(X,0)=\overline M_{0,3}\times X$，而 $\overline M_{0,2}(X,0)=\varnothing$。
幸运的是，不难证明：若 $n+2g\geq4$，或者 $\beta\ne0$ 且 $n\geq1$，则 $\phi_n$ 存在。
对 $\pi_n$ 也一样。

\subsection{定义 Gromov--Witten 类与不变量}\label{subsec:gw-algebraic-classes}
\index{Gromov--Witten不变量@Gromov--Witten 不变量!代数几何中的定义|hyperpage}
有了虚拟基本类 $\xi=[\overline M_{g,n}(X,\beta)]^{\mathrm{virt}}$，终于能严格定义
Gromov--Witten 类与不变量。回忆，由式~\eqref{eq:stable-map-evaluation-forgetting}，
有映射 $\pi:\overline M_{g,n}(X,\beta)\to X^n\times\overline M_{g,n}$；
以 $p_i$，$i=1,2$，表示到 $X^n\times\overline M_{g,n}$ 第 $i$ 个因子的投影。

\begin{definition}\label{def:gw-algebraic-invariants}
设 $\beta\in H_2(X,\mathbb Z)$ 为同调类，$\alpha_1,\ldots,\alpha_n\in H^*(X,\mathbb Q)$ 为上同调类。
\begin{enumerate}[label=(\roman*)]
\item 若 $n+2g\geq3$，则 Gromov--Witten 类 $I_{g,n,\beta}(\alpha_1,\ldots,\alpha_n)$ 为
$H^*(\overline M_{g,n},\mathbb Q)$ 中由下式定义的上同调类：
\[
I_{g,n,\beta}(\alpha_1,\ldots,\alpha_n)
=PD^{-1}p_{2*}\bigl(p_1^*(\alpha_1\otimes\cdots\otimes\alpha_n)\cap\pi_*(\xi)\bigr),
\]
其中 $\xi=[\overline M_{g,n}(X,\beta)]^{\mathrm{virt}}$ 为
$\overline M_{g,n}(X,\beta)$ 的虚拟基本类，$PD$ 为 Poincaré 对偶。
\item 对 $n,g\geq0$，Gromov--Witten 不变量
$\langle I_{g,n,\beta}\rangle(\alpha_1,\ldots,\alpha_n)$ 为由
\[
\langle I_{g,n,\beta}\rangle(\alpha_1,\ldots,\alpha_n)
=\int_\xi e_1^*(\alpha_1)\cup\cdots\cup e_n^*(\alpha_n)
\]
定义的有理数，其中 $e_i$ 如式~\eqref{eq:gw-evaluation-integral} 所述。
\end{enumerate}
\end{definition}
注意，这些定义与式~\eqref{eq:gw-fundamental-class-definition} 和~\eqref{eq:gw-evaluation-integral} 相同，
只是 $\xi$ 改为虚拟基本类。此外，若 $n+2g\geq3$，可证明
\[
\langle I_{g,n,\beta}\rangle(\alpha_1,\ldots,\alpha_n)
=\int_{\overline M_{g,n}}I_{g,n,\beta}(\alpha_1,\ldots,\alpha_n).
\]
因此，$n+2g\geq3$ 时，Gromov--Witten 不变量由相应的 Gromov--Witten 类确定。

\hyperref[sec:gw-properties]{第~\ref*{sec:gw-properties} 节}将研究 Gromov--Witten 类与不变量的性质，\hyperref[sec:computing-gw-i]{第~\ref*{sec:computing-gw-i} 节}将计算一些例子。
还应说明：虽然这里用 $\mathbb Q$ 上的同调与上同调处理 Gromov--Witten 类，
许多文献使用的却是 Chow 群。\SourceCite{Behrend} 使用\hyperref[subsec:virtual-fundamental-class]{第~\ref*{subsec:virtual-fundamental-class} 小节}先前讨论的相对虚拟基本类版本，
给出了类似的 Gromov--Witten 不变量定义。

本节以 \SourceCite{Kontsevich2} 给出的五次三维簇某些 Gromov--Witten 不变量的公式结束。
这一结果将在\hyperref[chap:mirror-theorem]{第~\ref*{chap:mirror-theorem} 章}发挥重要作用。

\SourcePage{207}
\SourceExampleNumber{7.1.6}{1}
\begin{example}\label{ex:quintic-gw-euler-integral}
设 $V\subset\mathbb P^4$ 为光滑五次三维簇，仍以 $\ell\in H_2(V,\mathbb Z)$ 表示直线的类。
例题~\ref{ex:quintic-virtual-class} 已说明，虚拟基本类
$[\overline M_{0,0}(V,d\ell)]^{\mathrm{virt}}$ 为零圈。
因此，由定义~\ref{def:gw-algebraic-invariants}，
\[
\langle I_{0,0,d\ell}\rangle
=\int_{[\overline M_{0,0}(V,d\ell)]^{\mathrm{virt}}}1
=\text{零圈 }[\overline M_{0,0}(V,d\ell)]^{\mathrm{virt}}\text{ 的次数}.
\]
可通过嵌入 $i:\overline M_{0,0}(V,d\ell)\hookrightarrow\overline M_{0,0}(\mathbb P^4,d)$，
在 $\overline M_{0,0}(\mathbb P^4,d)$ 上计算这一次数。
由式~\eqref{eq:quintic-virtual-class-pushforward}，
\[
i_*[\overline M_{0,0}(V,d\ell)]^{\mathrm{virt}}
=c_{5d+1}(\mathcal V_d)\cap[\overline M_{0,0}(\mathbb P^4,d)],
\]
其中 $\mathcal V_d$ 为 $\overline{\mathcal M}_{0,0}(\mathbb P^4,d)$ 上的向量丛，
在亏格 $0$ 的稳定映射 $f:C\to\mathbb P^4$ 处的纤维为 $H^0(C,f^*\mathcal O_{\mathbb P^4}(5))$。
由于 $c_{5d+1}(\mathcal V_d)$ 为最高次数的上同调类，
零圈 $c_{5d+1}(\mathcal V_d)\cap[\overline M_{0,0}(\mathbb P^4,d)]$ 的次数，
由在 $\overline M_{0,0}(\mathbb P^4,d)$ 的基本类上积分给出。因此
\[
\langle I_{0,0,d\ell}\rangle
=\int_{\overline M_{0,0}(\mathbb P^4,d)}c_{5d+1}(\mathcal V_d).
\]
这一公式首次见于 \SourceCite[第 2.2 节]{Kontsevich2}。
这里稍宽松地把 $\mathcal V_d$ 写成定义在 $\overline M_{0,0}(\mathbb P^4,d)$ 上的丛；
如前所述，其 Euler 类在 $\overline M_{0,0}(\mathbb P^4,d)$ 上给出良定义的零圈。
\hyperref[subsec:kontsevich-localization]{第~\ref*{subsec:kontsevich-localization} 小节}将说明如何用等变上同调与局部化计算这一积分。
\end{example}

\section{辛几何中的定义}\label{sec:gw-symplectic-definition}
\index{Gromov--Witten不变量@Gromov--Witten 不变量!辛几何中的定义|hyperpage}
本节简要概述 Gromov--Witten 不变量的辛几何表述所用的思想与技术。
这些定义以辛几何中 $J$ 全纯映射的理论为基础，带来了对量子上同调的最早严格处理。
部分基本思想见 \SourceCite{Witten3}；对半正辛流形的亏格 $0$ 辛 Gromov--Witten 不变量的首次显式处理，
则归功于 \SourceCite{Ruan2} 与 \SourceCite{RT1}。\SourceCite{MS} 可作为这一方法的入门书。
半正情形的更高亏格不变量，先在 \SourceCite{RT1} 中为固定复结构定义，
随后在 \SourceCite{RT2} 中推广到变化的复结构。
现在已有若干一般辛流形的 Gromov--Witten 不变量定义，见 \SourceCite{LTi3,Siebert1,Ruan3,FO}。
本节概述省略的许多细节，均可查阅这些文献。

\subsection{辛流形与 $J$ 全纯映射}\label{subsec:j-holomorphic-maps}
\index{几乎复结构|hyperpage}\index{几乎Kähler流形@几乎 Kähler 流形|hyperpage}
设 $M$ 为实维数 $2n$ 的紧辛流形。即 $M$ 上有闭 $2$ 形式 $\omega$，使 $\omega^n$ 在每一点非零。
每个 Kähler 流形都有由 Kähler 形式给出的自然辛结构；辛几何中的一个有力思想是，
任何辛流形 $(M,\omega)$ 都是\emph{几乎 Kähler} 的。
特别地，Kähler 流形的复结构与 Kähler 度量，可以换成几乎复结构及与 $\omega$ 关联的 Riemann 度量。

辛流形上的几乎复结构可满足若干条件。下面的定义正式给出其中两个条件。

\SourcePage{208}
\begin{definition}\label{def:tamed-compatible-complex-structure}
设 $(M,\omega)$ 为辛流形，$J$ 为 $M$ 上的几乎复结构。若
\[
\omega(X,JX)>0,\qquad
p\in M,\quad0\ne X\in T_pM,
\]
则称 $J$ \emph{受 $\omega$ 驯服}。
此外，若 $J$ 受 $\omega$ 驯服，且对 $M$ 上所有向量场 $X,Y$ 都有
\[
\omega(JX,JY)=\omega(X,Y),
\]
则称 $J$ \emph{与 $\omega$ 相容}。
\end{definition}
\index{驯服的几乎复结构|hyperpage}\index{相容的几乎复结构|hyperpage}
相容性等价于 $g(X,Y)=\omega(X,JY)$ 是 $M$ 上的 Riemann 度量。
满足上述任一条件的 $J$ 总存在，而且彼此之间都可通过形变连接。
例如，这意味着：在辛流形上，可用任意受 $\omega$ 驯服的 $J$，
赋予切丛 $TM$ 复结构，从而定义它的 Chern 类 $c_i(TM)$。
若 $J$ 受 $\omega$ 驯服，则
\[
\mu_J(X,Y)=\frac12\bigl(\omega(X,JY)+\omega(Y,JX)\bigr)
\]
定义一个度量；当 $J$ 与 $\omega$ 相容时，它就是上述 $g$。

使用\IndexLocation{indexloc:0048:208}受驯服的几乎复结构，其好处是更容易得到“一般”情形。
另一方面，若施加更强的相容性条件，往往能作出更精确的陈述。

这里的主要处理沿用 \SourceCite{Ruan2} 开创的方法，也讨论 \SourceCite{RT1,RT2} 中对它的推广；
这些工作为\emph{半正}辛流形发展了 Gromov--Witten 不变量理论。
半正是指：对每个可由映射 $f:S^2\to M$ 代表的同调类 $\beta\in H_2(M,\mathbb Z)$，
不会同时出现
\[
\int_\beta\omega>0,\qquad
3-\tfrac12\dim_{\mathbb R}M\leq\int_\beta c_1(TM)<0.
\]
例如，每个 Calabi--Yau 流形在其自然辛结构下都是半正的，$\mathbb P^r$ 也是如此。
还要注意，每个复维数不超过 $3$ 的光滑射影簇都是半正的。
提醒读者，文献中“半正”有时采用稍有不同的含义。\index{半正辛流形|hyperpage}

接下来讨论辛流形 $M$ 上的曲线。在代数几何中，用方程定义曲线及其他子簇。
在辛几何中这一做法不适用，因此改令 $C$ 为 Riemann 曲面，考虑
$C^\infty$ 映射 $f:C\to M$，它在以下意义下“全纯”。

\begin{definition}\label{def:j-holomorphic-map}
设 $(M,\omega)$ 为辛流形，$J$ 为受 $\omega$ 驯服的几乎复结构。
$J$ 全纯映射是 Riemann 曲面 $C$ 到 $M$ 的 $C^\infty$ 映射 $f:C\to M$，
其微分 $df:TC\to TM$ 相对于 $C$ 的复结构与 $M$ 的几乎复结构，是复向量丛映射。
\end{definition}
\index{J全纯映射@$J$ 全纯映射|hyperpage}
注意，映射 $f:C\to M$ 有一个可作用于它的 $\overline\partial$ 算子，定义为
\begin{equation}\label{eq:symplectic-cauchy-riemann-operator}
\overline\partial_J=d+J\circ d\circ j_C,
\end{equation}
其中 $j_C$ 为 $C$ 的可积几乎复结构。
$J$ 全纯性等价于 Cauchy--Riemann 方程 $\overline\partial_Jf=0$；
对 $\overline\partial_J$ 的定义施加 $J$ 并使用 $J^2=-1$，即可验证。
特别地，若 $J$ 可积，本情形中的 $J$ 全纯映射就是参数化的复曲线。

\SourcePage{209}
辛 Gromov--Witten 不变量的基本思想，类似于式~\eqref{eq:gw-pointed-maps} 的做法：
希望计数曲线
\begin{equation}\label{eq:symplectic-gw-pointed-maps}
f:C\to M,\qquad f_*[C]=\beta,\quad f(p_i)\in Z_i,\quad i=1,\ldots,n,
\end{equation}
其中，像先前一样，$p_1,\ldots,p_n$ 为 $C$ 的固定点，
$\beta\in H_2(M,\mathbb Z)$ 为同调类，$Z_1,\ldots,Z_n$ 为 $M$ 上处于一般位置的圈。
在 \SourceCite{Ruan2,RT2} 的方法中，通过要求 $f:C\to M$ 对固定的一般 $J$ 为 $J$ 全纯映射，
使这一表述精确化；实际上，还需要扰动 Cauchy--Riemann 方程。
一般性的作用，是确保有关模空间具有预期维数。所得不变量依赖 $\omega$，不依赖所选 $J$。
这与代数情形一致，因为\hyperref[sec:gw-properties]{第~\ref*{sec:gw-properties} 节}将看到，\hyperref[sec:gw-algebraic-definition]{第~\ref*{sec:gw-algebraic-definition} 节}的 Gromov--Witten 不变量在复结构形变下不变。
从镜像对称看，这是合理的：这些不变量出现在 A 模型关联函数中，而后者仅依赖 Kähler 模。

辛几何方法的一个优点是：选用一般几乎复结构时，代数情形中的许多困难，
例如模空间奇异或维数不符，便消失了。
事实上，与代数情形相比，式~\eqref{eq:symplectic-gw-pointed-maps} 已更接近 Gromov--Witten 不变量的真正定义。
不过，在较新的辛 Gromov--Witten 不变量方法中，可去掉半正假设，同时构造虚拟基本类的对应物。
这时，也不再需要实际模空间达到预期维数。

为了描述这些方法，必须了解有关模空间。下面转入这一主题。

\subsection{模空间}\label{subsec:symplectic-map-moduli}
\index{J全纯映射@$J$ 全纯映射!模空间|hyperpage}\index{扰动的J全纯映射@扰动的 $J$ 全纯映射|hyperpage}
先沿用 \SourceCite{Ruan2} 的处理，介绍性论述见 \SourceCite{MS}。
这里的思想，是研究类为 $\beta\in H_2(M,\mathbb Z)$ 的 $J$ 全纯映射 $f:C\to M$。
固定 $C$ 的复结构与 $M$ 的几乎复结构 $J$。出于技术原因，需要假定 $f$ 是\emph{简单的}，
即非常值且不能分解为 $C\to C'\to M$，其中 $C\to C'$ 的次数大于 $1$。
若 $f$ 不是简单的，称它为\emph{\IndexLocation{indexloc:0653:209}多重覆盖}。令
\[
\mathcal M(C,J,\beta)=
\{f:C\to M:f\text{ 为简单的 }J\text{ 全纯映射，且 }f_*[C]=\beta\}.
\]
第一个结果是：当 $J$ 为受 $\omega$ 驯服的一般几乎复结构时，
$\mathcal M(C,J,\beta)$ 自然为实流形，其维数为
\begin{equation}\label{eq:fixed-curve-map-index}
(1-g)\dim_{\mathbb R}M+2\int_\beta c_1(TM).
\end{equation}
为了证明，将 $\mathcal M(C,J,\beta)$ 表示为两个无穷维流形之间某个映射的纤维。
对一般 $J$，这一映射具有 Fredholm 线性化，隐函数定理便表明纤维
$\mathcal M(C,J,\beta)$ 为有限维流形。
纤维维数由 Fredholm 指数给出，而 Hirzebruch--\allowbreak Riemann--\allowbreak Roch 定理给出式~\eqref{eq:fixed-curve-map-index}。
顺便指出，类似分析也用于构造 Donaldson 理论与 Seiberg--Witten 理论中的模空间。

流形 $\mathcal M(C,J,\beta)$ 对 $J$ 的依赖也很好。
更精确地说，随着 $J$ 变化，只要 $J$ 足够一般，各 $\mathcal M(C,J,\beta)$ 就彼此协边。
不过，由于 $\mathcal M(C,J,\beta)$ 通常非紧，这一结果没有想象中那么强。
\hyperref[subsec:instanton-numbers]{第~\ref*{subsec:instanton-numbers} 小节}猜想~\ref{conj:symplectic-instantons} 之后的讨论，将着重指出非紧性带来的一个困难。\index{协边|hyperpage}

\SourcePage{210}
接着介绍 Ruan 与 Tian \SourceCite{RT1} 对这些空间采用的另一种方法。
这里扰动 $J$ 全纯映射的概念；部分原因是为了处理多重覆盖。
而且，正如下文将指出的，扰动能得到紧化模空间。

回忆，$J$ 全纯意味着 Cauchy--Riemann 方程 $\overline\partial_Jf=0$。
在积空间 $C\times M$ 上，令
\[
\nu:p_1^*TC\longrightarrow p_2^*TM
\]
为共轭线性映射，其中 $p_i$，$i=1,2$，为到 $C\times M$ 第 $i$ 个因子的投影。
考虑\emph{非齐次 Cauchy--Riemann 方程}
\[
(\overline\partial_Jf)(p)=\nu(p,f(p)),\qquad p\in C.
\]
若 $f:C\to M$ 满足此方程，称它为\emph{扰动的 $J$ 全纯映射}。
假定 $J,\nu$ 都一般，考虑集合
\[
\mathcal M(C,J,\nu,\beta)
=\{f:C\to M:\overline\partial_Jf=\nu,\ f_*[C]=\beta\}.
\]
注意，不再要求 $f$ 简单。这是采用扰动的 $J$ 全纯映射的优点之一。
与齐次情形 $\nu=0$ 的关系，见 \SourceCite[第 2.10 节]{RT1}。
可证明，$\mathcal M(C,J,\nu,\beta)$ 自然为流形，维数为式~\eqref{eq:fixed-curve-map-index}。

到目前为止，曲线 $C$ 的复结构固定。\SourceCite{RT1} 中这一情形产生的不变量，
可用来定义包含所有亏格的拓扑量子场论（\hyperref[sec:topological-qft]{附录 B.6}）。
若允许复结构变化，用物理语言说就是拓扑理论\emph{与引力耦合}，
便得到先前所研究的模空间 $\overline M_{g,n}(X,\beta)$ 的辛几何对应物。
由于需要空间 $\overline M_{g,n}$，必须假定 $2g+n\geq3$。
不过，$\overline M_{g,n}$ 的奇点要求谨慎处理。
在 \SourceCite{RT2} 中，Ruan 与 Tian 定义空间 $\mathcal M^\mu_{g,n}(J,\nu,\beta)_I$，
粗略地说，它由满足 $f_*[C]=\beta$ 的所有扰动的 $J$ 全纯映射
\[
f:(C,p_1,\ldots,p_n)\longrightarrow M
\]
组成。这里 $C$ 光滑，亏格为 $g$；上标 $\mu$ 提醒我们，实际上是在
$\overline M_{g,n}$ 的有限覆盖 $\overline M^\mu_{g,n}$ 上工作。
通常通过固定一个级结构来实现这一覆盖，它是获得泛曲线所必需的。
下标 $I$ 则表示限制到 $\overline M^\mu_{g,n}$ 中由自同构群平凡的带 $n$ 个标记曲线组成的子集。

与先前一样，$\mathcal M^\mu_{g,n}(J,\nu,\beta)_I$ 自然为流形，维数为
\begin{equation}\label{eq:varying-curve-map-index}
2(1-g)\left(\tfrac12\dim_{\mathbb R}M-3\right)
+2\int_\beta c_1(TM)+2n.
\end{equation}
当 $M$ 为射影簇 $X$ 时，有 $\int_\beta c_1(TX)=-\int_\beta\omega_X$。
因此，这一维数恰为 $\overline M_{g,n}(X,\beta)$ 的预期实维数，
可与式~\eqref{eq:stable-map-expected-dimension} 比较。

这样，从扰动的 $J$ 全纯映射构造了两类空间：
$\mathcal M(C,J,\nu,\beta)$ 与 $\mathcal M^\mu_{g,n}(J,\nu,\beta)_I$。
各空间都有称为 \emph{Gromov--Uhlenbeck 紧化}的自然紧化。
由于 $(M,\omega)$ 半正，可证明这些紧化通过增加有限多个实余维至少 $2$ 的层得到。
下面将看到，这两个空间给出两种不同但密切相关的 Gromov--Witten 不变量。\index{Gromov--Uhlenbeck紧化@Gromov--Uhlenbeck 紧化|hyperpage}

\SourcePage{211}
现在介绍 \SourceCite{LTi3} 提出的第三种模空间方法。
这一方法使用与 $\omega$ 相容的几乎复结构 $J$，并引入\emph{稳定 $C^\ell$ 映射}。
它们是稳定映射 $f:C\to M$ 的 $C^\ell$ 版本。
曲线 $C$ 与稳定映射定义中的完全相同，标记点仍为光滑点。
区别在于，$f$ 不再是代数态射，而满足以下条件：
对 $C$ 的每个分支 $C_i$，$f$ 与正规化映射 $\widetilde C_i\to C_i$ 的复合是 $C^\ell$ 的。
此外，不仅对 $f$ 为常值的分支施加稳定性条件，而且对满足 $f_*[C_i]=0$ 的分支都施加这一条件。
$C$ 的双全纯同构给出 $C^\ell$ 映射之间自然的等价概念。

在这一情形中，还有相应的 $J$ 全纯稳定映射概念：
要求各映射 $\widetilde C_i\to C_i\to M$ 为 $J$ 全纯映射。

若固定亏格 $g$、标记点数目 $n$ 与同调类 $\beta=f_*[C]$，所得稳定 $C^\ell$ 映射的等价类集合记为
$\overline{\mathcal F}^{\ell}_\beta(M,g,n)$。
这一空间有自然拓扑，是要研究的模空间。\SourceCite{Ruan3,Siebert1} 的方法中也使用类似空间。

\subsection{辛 Gromov--Witten 不变量}\label{subsec:symplectic-gw-invariants}
\index{伪流形|hyperpage}\index{虚邻域|hyperpage}\index{有向Euler类@有向 Euler 类|hyperpage}
空间 $\mathcal M(C,J,\nu,\beta)$ 与 $\mathcal M^\mu_{g,n}(J,\nu,\beta)_I$ 有良好的紧化，
意味着它们几乎具有基本类。问题在于，边界 $\partial$ 是否有邻域 $U$，使 $\partial$ 为 $U$ 的形变收缩，
见 \SourceCite[注 4.3]{RT1}。若有正确的基本类，便可用类似式~\eqref{eq:gw-fundamental-class-definition}
的公式定义 Gromov--Witten 不变量。

由于没有一个良好的基本类，有两种处理方法：
\begin{itemize}
\item 用横截性论证，将定义化为有向交中的点数。这是 \SourceCite{RT1,RT2} 的方法。
\item 如\hyperref[sec:gw-algebraic-definition]{第~\ref*{sec:gw-algebraic-definition} 节}的虚拟基本类讨论一样，在更大的空间中工作。
在 \SourceCite{Ruan3} 中，Ruan 使用\emph{虚邻域}，通过将模空间嵌入一个 V 流形得到它。
\SourceCite{Siebert1} 有相关方法；\SourceCite{LTi3} 中也以类似方式使用上面定义的
空间 $\overline{\mathcal F}^{\ell}_\beta(M,g,n)$。
\end{itemize}
第二种方法更有力，因为它适用于所有辛流形，而非仅限半正辛流形。
不过，这里的讨论主要采用第一种方法，同时简要提及 \SourceCite{LTi3} 的一些进一步思想，
作为更一般技术的代表。

先从 $\mathcal M(C,J,\nu,\beta)$ 给出的不变量开始。目标是定义\emph{混合 Gromov--Witten 不变量}
\begin{equation}\label{eq:mixed-gw-invariant}
\Phi_{g,\beta}(\alpha_1,\ldots,\alpha_k\mid\beta_1,\ldots,\beta_l),
\end{equation}
其中 $\alpha_i,\beta_j\in H^*(M,\mathbb Z)$ 为满足
\begin{equation}\label{eq:mixed-gw-degree-condition}
\sum_{i=1}^k\deg\alpha_i+\sum_{j=1}^l\deg\beta_j
=(1-g)\dim_{\mathbb R}M+2\int_\beta c_1(TM)+2l
\end{equation}
的上同调类。在 \SourceCite{RT1} 中，Ruan 与 Tian 使用同调类，因此这里的 $\alpha_i,\beta_j$
与他们所用的类互为对偶。粗略思想如下：固定 $C$ 上的点 $p_1,\ldots,p_k$。
混合不变量~\eqref{eq:mixed-gw-invariant} 应是 $\mathcal M(C,J,\nu,\beta)$ 中满足
$f(p_i)$ 位于 $\alpha_i$ 的对偶圈、$f(C)$ 遇到 $\beta_j$ 的对偶圈的映射 $f:C\to M$ 的数目。
\index{混合Gromov--Witten不变量@混合 Gromov--Witten 不变量|hyperpage}

\SourcePage{212}
为了精确定义，用 $M$ 的\emph{伪流形}表示 $\alpha_i,\beta_j$ 的对偶圈。
这里，伪流形指连续映射 $F:Y\to M$，其中 $Y$ 分层，每个较低层在相邻较高层中至少有余维 $2$，
而 $F$ 在每一层上光滑。若 $(F_i,Y_i)$ 与 $(G_j,Z_j)$ 分别表示 $\alpha_i,\beta_j$ 的对偶圈，则
\[
F=\prod_{i=1}^kF_i\times\prod_{j=1}^lG_j:
\prod_{i=1}^kY_i\times\prod_{j=1}^lZ_j\longrightarrow M^{k+l}
\]
为伪流形，它表示与上同调类
$\alpha_1\otimes\cdots\otimes\alpha_k\otimes\beta_1\otimes\cdots\otimes\beta_l
\in H^*(M^{k+l},\mathbb Z)$ 对偶的圈。

另一方面，有固定点 $p_1,\ldots,p_k\in C$。它们给出\IndexLocation{indexloc:0276:212}赋值映射
\[
e:\mathcal M(C,J,\nu,\beta)\times C^l\longrightarrow M^{k+l},
\]
将 $(f,q_1,\ldots,q_l)$ 映为
$(f(p_1),\ldots,f(p_k),f(q_1),\ldots,f(q_l))\in M^{k+l}$。

于是，有两个到 $M^{k+l}$ 的映射。对一般 $J,\nu$，可证明两者表示横截相交的圈；
特别地，交点出现在光滑部分。次数条件~\eqref{eq:mixed-gw-degree-condition} 还表明，
$\operatorname{Im}(F)$ 与 $\operatorname{Im}(e)$ 表示互补维数的圈。
故按原像计数的相交集只有有限多个点。
由于所有对象都有自然定向，每个相交点 $P$ 都有由定向确定的符号 $\epsilon(P)=\pm1$。

\begin{definition}\label{def:mixed-gw-invariant}
设 $(M,\omega)$ 为半正辛流形，$\beta\in H_2(M,\mathbb Z)$ 为同调类。
给定满足式~\eqref{eq:mixed-gw-degree-condition} 的
$\alpha_i,\beta_j\in H^*(M,\mathbb Z)$，混合 Gromov--Witten 不变量定义为
\[
\Phi_{g,\beta}(\alpha_1,\ldots,\alpha_k\mid\beta_1,\ldots,\beta_l)
=\sum_{P\in F\times_{M^{k+l}}e}\epsilon(P).
\]
\end{definition}

Ruan 与 Tian \SourceCite{RT1} 证明，
$\Phi_{g,\beta}(\alpha_1,\ldots,\alpha_k\mid\beta_1,\ldots,\beta_l)$
与所选 $J,\nu$、点 $p_1,\ldots,p_k\in C$、$C$ 的复结构，
以及表示 $\alpha_i,\beta_j$ 的伪流形无关。他们还讨论了\IndexLocation{indexloc:0546:212}混合不变量的许多性质与应用。

可将这一构造与\hyperref[sec:gw-algebraic-definition]{第~\ref*{sec:gw-algebraic-definition} 节}的代数不变量联系起来。
设 $(M,\omega)$ 是带 Kähler 形式的光滑射影簇。令
\[
\pi:\overline M_{g,k+l}(M,\beta)\longrightarrow\overline M_{g,k}
\]
将稳定映射忘掉映射本身及最后 $l$ 个标记点，并收缩因此出现的不稳定分支。
这是式~\eqref{eq:stable-map-evaluation-forgetting} 中 $\pi_2$ 的推广；$\pi_2$ 对应 $l=0$。
令 $[pt]$ 表示 $\overline M_{g,k}$ 上一点的上同调类，
它代表带固定标记点 $(p_1,\ldots,p_k)$ 的固定曲线 $C$。
则不变量 $\Phi_{g,\beta}(\alpha_1,\ldots,\alpha_k\mid\beta_1,\ldots,\beta_l)$ 在思想上就是
\[
\int_{[\overline M_{g,k+l}(M,\beta)]^{\mathrm{virt}}}
\left(\prod_{i=1}^ke_i^*(\alpha_i)\right)\cup
\left(\prod_{j=1}^le_{k+j}^*(\beta_j)\right)\cup\pi^*[pt].
\]

不过，我们更关心\hyperref[sec:gw-algebraic-definition]{第~\ref*{sec:gw-algebraic-definition} 节}所讨论类型的 Gromov--Witten 不变量。
因此，转而研究由 $\mathcal M^\mu_{g,n}(J,\nu,\beta)_I$ 构造的不变量。
它们是\emph{辛 Gromov--Witten 不变量}
\begin{equation}\label{eq:symplectic-gw-invariant}
\Psi_{g,n,\beta}(\gamma;\alpha_1,\ldots,\alpha_n),
\end{equation}
\SourcePage{213}
为满足
\[
\deg\gamma+\sum_{i=1}^n\deg\alpha_i
=2(1-g)\left(\tfrac12\dim_{\mathbb R}M-3\right)
+2\int_\beta c_1(TM)+2n
\]
的上同调类 $\gamma\in H^*(\overline M_{g,n},\mathbb Q)$、
$\alpha_i\in H^*(M,\mathbb Z)$ 定义。
朴素思想是：若 $Z$ 为 $\overline M_{g,n}$ 中与 $\gamma$ 对偶的圈，
$Z_i$ 为 $M$ 中与 $\alpha_i$ 对偶的圈，则
$\Psi_{g,n,\beta}(\gamma;\alpha_1,\ldots,\alpha_n)$ 应是满足
$(C,p_1,\ldots,p_n)\in Z\subset\overline M_{g,n}$、
$f_*[C]=\beta$、$f(p_i)\in Z_i$ 的扰动 $J$ 全纯映射
$f:(C,p_1,\ldots,p_n)\to M$ 的数目。

式~\eqref{eq:symplectic-gw-invariant} 的精确定义类似混合情形，但细节更复杂：
并非 $H_*(\overline M_{g,n},\mathbb Z)$ 中所有类都由所需的伪流形表示，
而使用 $\mathcal M^\mu_{g,n}(J,\nu,\beta)_I$ 意味着是在有限覆盖
$\overline M^\mu_{g,n}\to\overline M_{g,n}$ 上工作。
这样先定义依赖有限覆盖选择的不变量
$\Psi^\mu_{g,n,\beta}(\gamma;\alpha_1,\ldots,\alpha_n)$。
再除以覆盖次数，得到与 $\mu$ 无关的
$\Psi_{g,n,\beta}(\gamma;\alpha_1,\ldots,\alpha_n)$。
这里不展开细节，完整定义见 \SourceCite{RT2}。

注意，固定 $\alpha_1,\ldots,\alpha_n$ 后，对所有 $\gamma$ 知道
$\Psi_{g,n,\beta}(\gamma;\alpha_1,\ldots,\alpha_n)$，
就唯一确定类
$\Psi_{g,n,\beta}(\alpha_1,\ldots,\alpha_n)\in H^*(\overline M_{g,n},\mathbb Q)$，满足
\[
\int_{\overline M_{g,n}}\gamma\cup
\Psi_{g,n,\beta}(\alpha_1,\ldots,\alpha_n)
=\Psi_{g,n,\beta}(\gamma;\alpha_1,\ldots,\alpha_n).
\]
称 $\Psi_{g,n,\beta}(\alpha_1,\ldots,\alpha_n)$ 为\emph{辛 Gromov--Witten 类}。
与代数情形相同，辛类可看成映射
\begin{equation}\label{eq:symplectic-gw-class-map}
\Psi_{g,n,\beta}:H^*(M)^{\otimes n}\longrightarrow H^*(\overline M_{g,n},\mathbb Q).
\end{equation}
后面将看到，代数 Gromov--Witten 类与辛 Gromov--Witten 类具有相似性质。

混合 Gromov--Witten 不变量与辛 Gromov--Witten 不变量之间有一些很好的关系。
例如，混合不变量固定了复结构，因而下式并不令人意外：
\[
\Phi_{g,\beta}(\alpha_1,\ldots,\alpha_n\mid)
=\Psi_{g,n,\beta}([pt];\alpha_1,\ldots,\alpha_n).
\]
为了陈述一般关系，令 $K_{k,l}$ 为 $\overline M_{g,k+l}$ 中以下集合的闭包所给出的圈：
\[
\{(C,p_1,\ldots,p_{k+l})\in\overline M_{g,k+l}:
(C,p_1,\ldots,p_k)\text{ 为 }\overline M_{g,k}\text{ 中的固定点}\}.
\]
则可证明
\[
\Phi_{g,\beta}(\alpha_1,\ldots,\alpha_k\mid\beta_1,\ldots,\beta_l)
=\Psi_{g,k+l,\beta}([K_{k,l}];
\alpha_1,\ldots,\alpha_k,\beta_1,\ldots,\beta_l),
\]
这意味着所有混合不变量都是辛不变量的特殊情形。
一个尤其相关的情况是 $g=0$。由于 $\overline M_{0,3}$ 只有一点，
故
\begin{equation}\label{eq:genus-zero-mixed-symplectic}
\Phi_{0,\beta}(\alpha_1,\alpha_2,\alpha_3\mid\alpha_4,\ldots,\alpha_n)
=\Psi_{0,n,\beta}(1;\alpha_1,\ldots,\alpha_n).
\end{equation}
对此公式的很好说明见 \SourceCite{Voisin3}。

本节最后简要提及 \SourceCite{LTi3} 方法中的一些思想。
令 $\overline{\mathcal F}^{\ell}_\beta(M,g,n)$ 为\hyperref[subsec:symplectic-map-moduli]{第~\ref*{subsec:symplectic-map-moduli} 小节}末尾描述的模空间。
\SourcePage{214}
在其上如下定义广义丛 $E$。给定 $C$ 光滑的稳定 $C^\ell$ 映射 $f:C\to M$，
令 $\Lambda_f^{0,1}$ 为
\[
\operatorname{Hom}(TC,f^*TM)
\]
中满足 $\nu\circ j_C=-J\circ\nu$ 的连续截面 $\nu$ 的空间。
这些截面应看作 $C$ 上取值于 $f^*(TM)$ 的 $(0,1)$ 形式。
注意，Cauchy--Riemann 算子 $\overline\partial_Jf$ 属于 $\Lambda_f^{0,1}$，
这可用式~\eqref{eq:symplectic-cauchy-riemann-operator}、$J^2=-1$ 与 $j_C^2=-1$ 立即算出。
当 $C$ 不光滑时，$\Lambda_f^{0,1}$ 的定义需稍作修改，细节见 \SourceCite{LTi3}；
修改与定义 $\overline{\mathcal F}^{\ell}_\beta(M,g,n)$ 的等价关系相容。
结果是，各 $\Lambda_f^{0,1}$ 成为 $\overline{\mathcal F}^{\ell}_\beta(M,g,n)$ 上广义丛 $E$ 的纤维。

由构造，$\overline\partial_J$ 诱导 $E$ 的全局截面。
于是，我们得到无穷维空间上一个无穷维丛的截面。
思想是寻找一种对象，使它发挥有限维情形中丛的 Euler 类的作用。
\SourceCite{LTi3} 证明，当 $\ell\geq2$ 时，这一截面赋予 $E$ \emph{广义 Fredholm 轨形丛}结构，
这一概念由 \SourceCite{LTi3} 引入。主要结果是：这些数据确定一个\emph{有向 Euler 类}，
它是丛的底空间上的同调类。
这个类支于 $E$ 的截面 $\overline\partial_J$ 的一个精心选取的扰动的零点轨迹。
所研究情形中，该同调类的维数恰好是预期维数~\eqref{eq:varying-curve-map-index}。
这一有向 Euler 类充当虚拟基本类，容易按熟悉的方法定义辛 Gromov--Witten 不变量。

\SourceCite{LTi3} 的 Gromov--Witten 不变量推广了 \SourceCite{RT2} 的不变量，
后者只在 $M$ 半正时定义。注意，不必假定 $2g+n\geq3$。
\SourceCite{Siebert1} 也有相关构造。
人们相信各种辛 Gromov--Witten 不变量彼此一致，但截至本书写作时，尚无已写出的证明。
下一节将讨论：各种辛不变量与\hyperref[sec:gw-algebraic-definition]{第~\ref*{sec:gw-algebraic-definition} 节}的代数 Gromov--Witten 不变量，在多大程度上相同。

\section{Gromov--Witten 类的性质}\label{sec:gw-properties}
接下来的任务，是研究 Gromov--Witten 类，以及前两节各种定义之间的关系。
先描述所有定义共有的性质。

\subsection{Gromov--Witten 类的公理}\label{subsec:gw-axioms}
\index{Gromov--Witten类@Gromov--Witten 类!公理|hyperpage}\index{分裂公理|hyperpage}\index{约化公理|hyperpage}\index{除子公理|hyperpage}\index{基本类公理|hyperpage}\index{点映射公理|hyperpage}\index{形变公理|hyperpage}\index{复合律|hyperpage}
在 \SourceCite{KoM1} 中，Kontsevich 与 Manin 为 Gromov--Witten 类提出了一组公理。
现在已知，\hyperref[sec:gw-algebraic-definition]{第~\ref*{sec:gw-algebraic-definition} 节}的代数类 $I_{g,n,\beta}$ 与\hyperref[sec:gw-symplectic-definition]{第~\ref*{sec:gw-symplectic-definition} 节}的辛类 $\Psi_{g,n,\beta}$ 都满足这些公理。
代数情形的证明见 \SourceCite{BM,LTi2,BF,Behrend}；
辛情形的证明见 \SourceCite{RT1,RT2,LTi3}。
还将加入一条形变公理；它在代数情形中由 \SourceCite{LTi2} 证明，
在辛情形中则由辛不变性得到。

为简便起见，为 $I_{g,n,\beta}$ 陈述公理。
回忆\hyperref[sec:gw-algebraic-definition]{第~\ref*{sec:gw-algebraic-definition} 节}：若 $X$ 为光滑射影簇，则对 $g\geq0$、$n\geq0$、$n+2g\geq3$，
Gromov--Witten 类为映射
\[
I_{g,n,\beta}:H^*(X,\mathbb Q)^{\otimes n}
\longrightarrow H^*(\overline M_{g,n},\mathbb Q).
\]
\SourcePage{215}
此外，对 $g,n\geq0$，Gromov--Witten 不变量为映射
\[
\langle I_{g,n,\beta}\rangle:H^*(X,\mathbb Q)^{\otimes n}\longrightarrow\mathbb Q.
\]
当 $n+2g\geq3$ 时，它们满足关系
\[
\langle I_{g,n,\beta}\rangle(\alpha_1,\ldots,\alpha_n)
=\int_{\overline M_{g,n}}I_{g,n,\beta}(\alpha_1,\ldots,\alpha_n).
\]

理解下面的公理时，应回忆式~\eqref{eq:gw-pointed-maps} 对 Gromov--Witten 类的朴素解释：
若 $Z_i$ 为 $X$ 中与 $\alpha_i$ 对偶的圈，
则 $I_{g,n,\beta}(\alpha_1,\ldots,\alpha_n)$ 应是由那些
$(C,p_1,\ldots,p_n)\in\overline M_{g,n}$ 给出的上同调类，其中可以找到 $f$ 满足
\begin{equation}\label{eq:gw-axiom-naive-stable-map}
f:(C,p_1,\ldots,p_n)\longrightarrow X\text{ 为稳定映射},\quad
f(p_i)\in Z_i,\quad f_*[C]=\beta.
\end{equation}
带着这一想法，进入各条公理。

\begin{axiom*}{线性}
$I_{g,n,\beta}$ 对每个变元都是线性的。
\end{axiom*}
朴素地说，这是因为圈的和就是它们的并，计入各自的重数。

\begin{axiom*}{有效性}
对光滑射影簇 $X$，若 $\beta$ 不是有效类，则 $I_{g,n,\beta}=0$。
\end{axiom*}
这是合理的，因为只要 $f:C\to X$ 全纯，$f_*[C]$ 就有效。
对辛流形 $(M,\omega)$，这一公理换成以下观察：只要 $\int_\beta\omega<0$，不变量就消失。

\begin{axiom*}{次数}
若 $\alpha_1,\ldots,\alpha_n$ 为齐次上同调类，则
$I_{g,n,\beta}(\alpha_1,\ldots,\alpha_n)\in H^*(\overline M_{g,n},\mathbb Q)$ 的次数为
\[
2(g-1)\dim X+2\int_\beta\omega_X+\sum_{i=1}^n\deg\alpha_i.
\]
\end{axiom*}
当模空间 $\overline M_{g,n}(X,\beta)$ 具有预期维数~\eqref{eq:stable-map-expected-dimension}，
并有良好的基本类时，这由式~\eqref{eq:gw-gysin-definition} 得到。
一般地，由于\IndexLocation{indexloc:1415:215}虚拟基本类具有预期次数，次数公理由定义~\ref{def:gw-algebraic-invariants} 得到。

次数公理意味着，$I_{g,n,\beta}(\alpha_1,\ldots,\alpha_n)$ 为最高次数的类，当且仅当
\begin{equation}\label{eq:gw-top-degree-condition}
\sum_{i=1}^n\deg\alpha_i
=2(1-g)\dim X-2\int_\beta\omega_X+2(3g-3+n).
\end{equation}
因此，考虑 Gromov--Witten 不变量 $\langle I_{g,n,\beta}\rangle(\alpha_1,\ldots,\alpha_n)$ 时，
总可假定式~\eqref{eq:gw-top-degree-condition} 成立，否则不变量为零。
由定义~\ref{def:gw-algebraic-invariants}，这对所有 $n,g\geq0$ 都有效。
文献中，有时用“代数次数” $\tfrac12\deg\alpha_i$ 代替 $\deg\alpha_i$；
如此，余维 $i$ 的子簇 $Z$ 的对偶类 $[Z]$ 具有次数 $i$。
处理偶次数上同调时，这一约定最为有用。

\begin{axiom*}{等变性}
对称群 $S_n$ 自然作用于 $H^*(X,\mathbb Q)^{\otimes n}$ 与 $H^*(\overline M_{g,n},\mathbb Q)$。
在前者上采用分次置换作用，在后者上通过置换 $(C,p_1,\ldots,p_n)$ 的标记点 $p_i$ 作用。
映射
\[
I_{g,n,\beta}:H^*(X,\mathbb Q)^{\otimes n}
\longrightarrow H^*(\overline M_{g,n},\mathbb Q)
\]
是 $S_n$ 等变的。
\end{axiom*}
这一性质背后的直观理由，应由式~\eqref{eq:gw-axiom-naive-stable-map} 看出。
对 Gromov--Witten 不变量，等变性意味着
\[
\begin{aligned}
&\langle I_{g,n,\beta}\rangle(\alpha_1,\ldots,\alpha_i,\alpha_{i+1},\ldots,\alpha_n)\\
&\quad=(-1)^{\deg\alpha_i\deg\alpha_{i+1}}
\langle I_{g,n,\beta}\rangle(\alpha_1,\ldots,\alpha_{i+1},\alpha_i,\ldots,\alpha_n).
\end{aligned}
\]

\SourcePage{216}
\begin{axiom*}{基本类}
若 $n+2g\geq4$，忘掉最后一点得到自然映射
$\pi_n:\overline M_{g,n}\to\overline M_{g,n-1}$。
若 $[X]\in H^0(X,\mathbb Q)$ 为 $X$ 的基本类的对偶，即单位类，则
\[
I_{g,n,\beta}(\alpha_1,\ldots,\alpha_{n-1},[X])
=\pi_n^*I_{g,n-1,\beta}(\alpha_1,\ldots,\alpha_{n-1}).
\]
\end{axiom*}
用式~\eqref{eq:gw-axiom-naive-stable-map} 理解，这是合理的：
$f(p_n)\in X$ 对 $p_n$ 不施加条件。
对 Gromov--Witten 不变量，基本类公理蕴含
\begin{equation}\label{eq:gw-fundamental-class-vanishing}
\langle I_{g,n,\beta}\rangle(\alpha_1,\ldots,\alpha_{n-1},[X])=0.
\end{equation}
因为这一不变量只有在
$I_{g,n,\beta}(\alpha_1,\ldots,\alpha_{n-1},[X])$ 为最高次数类时才可能非零。
但 $\overline M_{g,n-1}$ 的维数更小，故相应次数的
$I_{g,n-1,\beta}(\alpha_1,\ldots,\alpha_{n-1})$ 为零。
注意，只要稳定映射的忘却映射 $\pi_n$ 有定义，式~\eqref{eq:gw-fundamental-class-vanishing} 就成立。
\hyperref[subsec:computing-virtual-class]{第~\ref*{subsec:computing-virtual-class} 小节}末尾的讨论表明，$n+2g\geq4$，或者 $\beta\ne0$ 且 $n\geq1$，都保证这一点。

\begin{axiom*}{除子}
若 $n+2g\geq4$，取与上一公理相同的
$\pi_n:\overline M_{g,n}\to\overline M_{g,n-1}$。若 $\alpha_n\in H^2(X,\mathbb Q)$，则
\[
\pi_{n!}I_{g,n,\beta}(\alpha_1,\ldots,\alpha_{n-1},\alpha_n)
=\left(\int_\beta\alpha_n\right)
I_{g,n-1,\beta}(\alpha_1,\ldots,\alpha_{n-1}).
\]
\end{axiom*}
用式~\eqref{eq:gw-axiom-naive-stable-map} 看其理由：
令稳定映射 $f:(C,p_1,\ldots,p_n)\to X$ 满足 $f_*[C]=\beta$，
且 $f(p_i)\in Z_i$，$i=1,\ldots,n-1$。
则 $f(p_n)$ 必须位于 $f(C)\cap Z_n$，其相交数为 $\beta\cdot Z_n$，
即 $p_n$ 按重数计有 $\int_\beta\alpha_n$ 种选择。由此得到上述公式。

对 Gromov--Witten 不变量，若 $\alpha_n\in H^2(X,\mathbb Q)$ 且
式~\eqref{eq:gw-top-degree-condition} 成立，则
\[
\langle I_{g,n,\beta}\rangle(\alpha_1,\ldots,\alpha_{n-1},\alpha_n)
=\left(\int_\beta\alpha_n\right)
\langle I_{g,n-1,\beta}\rangle(\alpha_1,\ldots,\alpha_{n-1}).
\]
如基本类公理，只要稳定映射的 $\pi_n$ 有定义，这一等式就成立；
由\hyperref[subsec:computing-virtual-class]{第~\ref*{subsec:computing-virtual-class} 小节}，$n+2g\geq4$，或者 $\beta\ne0$ 且 $n\geq1$，都足够。

\begin{axiom*}{点映射}
当 $\beta=0$ 时，若 $g=0$ 且各 $\alpha_i$ 为齐次上同调类，则
\[
I_{0,n,0}(\alpha_1,\ldots,\alpha_n)=
\begin{cases}
\displaystyle\left(\int_X\alpha_1\cup\cdots\cup\alpha_n\right)
[\overline M_{0,n}],&\displaystyle\sum_{i=1}^n\deg\alpha_i=2\dim X,\\
0,&\text{其他情形},
\end{cases}
\]
其中 $[\overline M_{0,n}]$ 表示其基本类的对偶，即单位类。
\end{axiom*}
为了看清这一点，注意满足 $f_*[C]=0$ 的映射必为常值。
因此，式~\eqref{eq:gw-axiom-naive-stable-map} 化为 $f(C)\in Z_1\cap\cdots\cap Z_n$。
当 $\sum_{i=1}^n\deg\alpha_i=2\dim X$ 时，由次数公理，该类的次数为 $0$，
从而得到上述公式；本情形中的其他类都为零。

这对 Gromov--Witten 不变量 $\langle I_{0,n,0}\rangle$ 有一个很好的推论。
将条件 $\sum_{i=1}^n\deg\alpha_i=2\dim X$ 与式~\eqref{eq:gw-top-degree-condition} 比较，
可见 $n=3$ 很特殊。更精确地，
\begin{equation}\label{eq:gw-degree-zero-three-point}
\langle I_{0,n,0}\rangle(\alpha_1,\ldots,\alpha_n)
=\begin{cases}
\displaystyle\int_X\alpha_1\cup\alpha_2\cup\alpha_3,&n=3,\\
0,&\text{其他情形}.
\end{cases}
\end{equation}
遗憾的是，这种推理仅适用于 $g=0$。
问题在于，$\overline M_{g,n}(X,0)=\overline M_{g,n}\times X$，
只有 $g=0$ 时才具有预期维数。
对所有 $g$ 都有 $I_{g,n,0}$ 的公式，但它们可能相当复杂。
后面的式~\eqref{eq:genus-one-constant} 将给出 $g=1$ 的情形；一般情形的困难见 \SourceCite[第 2.2.5 节]{KoM1}。

\SourcePage{217}
\begin{axiom*}{分裂}
设 $g=g_1+g_2$、$n=n_1+n_2$，且 $n_i+2g_i\geq2$。
对亏格分别为 $g_1,g_2$ 的稳定曲线
$(C_1,p_1,\ldots,p_{n_1+1})$ 与 $(C_2,q_1,\ldots,q_{n_2+1})$，
将 $p_{n_1+1}$ 与 $q_1$ 识别，把 $C_1\cup C_2$ 粘合为曲线 $C$，
得到亏格 $g$ 的稳定曲线 $(C,p_1,\ldots,p_{n_1},q_2,\ldots,q_{n_2+1})$。
因此，有映射
\begin{equation}\label{eq:gw-splitting-gluing-map}
\varphi:\overline M_{g_1,n_1+1}\times\overline M_{g_2,n_2+1}
\longrightarrow\overline M_{g,n}.
\end{equation}
分裂公理断言：
\[
\begin{aligned}
\varphi^*I_{g,n,\beta}(\alpha_1,\ldots,\alpha_n)
={}&\sum_{\beta_1+\beta_2=\beta}\sum_{i,j}g^{ij}\\
&\quad\cdot I_{g_1,n_1+1,\beta_1}(\alpha_1,\ldots,\alpha_{n_1},T_i)\\
&\quad\otimes I_{g_2,n_2+1,\beta_2}(T_j,\alpha_{n_1+1},\ldots,\alpha_n).
\end{aligned}
\]
其中 $T_i$ 为 $H^*(X,\mathbb Q)$ 的齐次基，$(g^{ij})$ 是矩阵 $(g_{ij})$ 的逆，
$g_{ij}=\int_XT_i\cup T_j$。等价地，对角线在 $H^*(X\times X,\mathbb Q)$ 中的上同调类为
$\sum_{i,j}g^{ij}T_i\otimes T_j$。由有效性公理，上述求和只有有限项。
\end{axiom*}
这一公理陈述起来稍复杂。为看清公式的来源，首先注意：
式~\eqref{eq:gw-axiom-naive-stable-map} 中映射在 $\varphi$ 下的逆像，由映射
\[
f:(C_1\cup C_2,p_1,\ldots,p_{n_1},p,q,p_{n_1+1},\ldots,p_n)\longrightarrow X
\]
组成，满足 $f_*[C_1]+f_*[C_2]=\beta$、$f(p_i)\in Z_i$ 及 $f(p)=f(q)$。
第一个条件对应于上述公式中 $\beta_1+\beta_2=\beta$ 的分解；
最后一个条件表示 $(f(p),f(q))$ 应位于 $X\times X$ 的对角线上。
由于对角线的类为 $\sum_{i,j}g^{ij}T_i\otimes T_j$，得到上述公式。

\begin{axiom*}{约化}
将最后两个标记点粘合，得到自然映射
$\psi:\overline M_{g-1,n+2}\to\overline M_{g,n}$。
用上一公理的对角线类及 $(g^{ij})$，有
\[
\psi^*I_{g,n,\beta}(\alpha_1,\ldots,\alpha_n)
=\sum_{i,j}g^{ij}I_{g-1,n+2,\beta}(\alpha_1,\ldots,\alpha_n,T_i,T_j).
\]
\end{axiom*}
这是合理的，因为式~\eqref{eq:gw-axiom-naive-stable-map} 中映射在 $\psi$ 下的逆像，
由亏格 $g-1$ 的映射 $f:(C,p_1,\ldots,p_{n+2})\to X$ 组成，满足
$f_*[C]=\beta$、$f(p_i)\in Z_i$，$i=1,\ldots,n$，及 $f(p_{n+1})=f(p_{n+2})$。
最后一个条件要求 $(f(p_{n+1}),f(p_{n+2}))$ 位于 $X\times X$ 的对角线上。
如上，对角线类为 $\sum_{i,j}g^{ij}T_i\otimes T_j$，公式随之成立。

\begin{axiom*}{形变}
设 $f:\mathcal X\to T$ 为光滑固有映射，底空间 $T$ 连通；对 $t\in T$，置 $X_t=f^{-1}(t)$。
于是，对每个 $t\in T$ 与 $\beta_t\in H_2(X_t,\mathbb Z)$，有映射
\[
I^{X_t}_{g,n,\beta_t}:H^*(X_t,\mathbb Q)^{\otimes n}
\longrightarrow H^*(\overline M_{g,n},\mathbb Q).
\]
若 $\beta_t$ 是同调局部系统 $H_2(X_t,\mathbb Z)$ 的局部常值截面，
$\alpha_1,\ldots,\alpha_n$ 是 $H^*(X_t,\mathbb Q)$ 的局部常值截面，则
$I^{X_t}_{g,n,\beta_t}(\alpha_1,\ldots,\alpha_n)$ 不随 $t$ 改变。
\end{axiom*}
其直观理由是：Gromov--Witten 类应在复结构形变下不变。
采用 $I_{g,n,\beta}$ 的辛几何定义时，这一点很清楚。

这些公理中，分裂公理与\IndexLocation{indexloc:0506:217}约化公理也许最重要。
\hyperref[chap:quantum-cohomology]{第~\ref*{chap:quantum-cohomology} 章}将看到，它们与量子上同调的结合律证明密切相关。
在 Ruan 与 Tian 的论文中，这两条公理合称\emph{复合律}。

除上述公理外，\SourceCite{KoM1} 还包含一条\emph{\IndexLocation{indexloc:0504:217}动机公理}：
Gromov--Witten 类应由
\SourcePage{218}
$X^n\times\overline M_{g,n}$ 的 Chow 环中适当次数的类诱导。
由于代数 Gromov--Witten 类用虚拟基本类
$[\overline{\mathcal M}_{g,n}(X,\beta)]^{\mathrm{virt}}$ 定义，这条公理随之成立。

应强调，上述各公理的直观解释不是证明。
不过，在辛情形中，某些“论证”确实与实际证明相当接近，
但复合律除外，见 \SourceCite{RT1,RT2}。
另一方面，代数情形的证明 \SourceCite{LTi2,Behrend} 更为精细。
如 \SourceCite{BM} 所解释，上述公理归结为虚拟基本类
$[\overline{\mathcal M}_{g,n}(X,\beta)]^{\mathrm{virt}}$ 之间五种关键的相容性。

\subsection{各种定义的等价性}\label{subsec:gw-equivalence}
\index{Gromov--Witten不变量@Gromov--Witten 不变量!各种定义的等价性|hyperpage}
\hyperref[sec:gw-algebraic-definition]{第~\ref*{sec:gw-algebraic-definition} 节}和\hyperref[sec:gw-symplectic-definition]{第~\ref*{sec:gw-symplectic-definition} 节}介绍了 Gromov--Witten 类与不变量的各种定义。
代数情形中，有两个定义，因为 \SourceCite{LTi2} 与 \SourceCite{BF,Behrend}
给出的虚拟基本类构造略有不同。如\hyperref[subsec:virtual-fundamental-class]{第~\ref*{subsec:virtual-fundamental-class} 小节}所述，二者应等价，但证明尚未写出。

辛情形中，\SourceCite{LTi3} 的不变量推广了 \SourceCite{RT2} 的不变量，
此外还有 \SourceCite{Siebert1} 给出的辛不变量。
同样，预期这些定义给出相同不变量，但尚未证明。

一个很好的进展是，代数不变量与辛不变量的一致性证明已有实质推进。
特别地，\SourceCite{LTi4} 证明了 \SourceCite{LTi2} 定义的代数不变量与
\SourceCite{LTi3} 定义的辛不变量一致；
\SourceCite{Siebert2} 证明了 \SourceCite{Behrend} 的代数不变量与
\SourceCite{Siebert1} 的辛不变量一致。
因此，似乎很可能所有定义在共同适用范围内都一致。

下面是一种已知这种一致性成立的情形。
用 $\langle I_{g,n,\beta}\rangle^{\mathrm{LRT}}$ 表示：
由 \SourceCite{LTi2} 的代数构造，或由 \SourceCite{LTi3} 的辛构造定义的 Gromov--Witten 不变量，
后者推广了 \SourceCite{RT2}。
又用 $\langle I_{g,n,\beta}\rangle^{\mathrm{BFS}}$ 表示：
由 \SourceCite{Behrend} 的代数构造，或由 \SourceCite{Siebert1} 的辛构造定义的 Gromov--Witten 不变量；
前者使用 \SourceCite{BF} 中的相对虚拟基本类。于是有以下结果。

\begin{theorem}\label{thm:homogeneous-gw-equivalence}
设 $X$ 为光滑射影簇，某个线性代数群在其上可迁地作用。
取 $\alpha_1,\ldots,\alpha_n\in H^*(X,\mathbb Q)$，则
\[
\langle I_{0,n,\beta}\rangle^{\mathrm{LRT}}(\alpha_1,\ldots,\alpha_n)
=\langle I_{0,n,\beta}\rangle^{\mathrm{BFS}}(\alpha_1,\ldots,\alpha_n).
\]
此外，当次数条件~\eqref{eq:gw-top-degree-condition} 成立、各类以一般位置的圈表示时，
这些不变量按重数计数满足式~\eqref{eq:gw-pointed-maps} 的带标记映射 $f:\mathbb P^1\to X$。
\end{theorem}
\begin{proof}
为检验等式，可以采用 $\langle I_{0,n,\beta}\rangle^{\mathrm{LRT}}$
与 $\langle I_{0,n,\beta}\rangle^{\mathrm{BFS}}$ 各自的代数定义。
\SourceCite{FP} 的引理 13 表明，计数定义给出代数 Gromov--Witten 不变量。
特别地，虚拟基本类等于基本类；采用虚拟基本类的任一种代数定义都如此。
群作用保证能找到与 $\alpha_i$ 对偶、处于一般位置的圈 $Z_i$。

在辛几何方面，Li 与 Tian 在 \SourceCite{LTi1} 中利用可迁群作用证明：
采用 $X$ 给定的复结构 $J_0$ 并取 $\nu=0$，得到相同数目。
对 \SourceCite{Siebert1} 的辛定义，也有类似结论。
\end{proof}

定理~\ref{thm:homogeneous-gw-equivalence} 断言这些不变量相等，其证明还表明：
它们确实计数相应类型的有理曲线。下一节将看到，一般 Gromov--Witten 不变量的\IndexLocation{indexloc:0534:218}枚举意义，
通常更难理解。

\SourcePage{219}
最后还应提到，除了迄今讨论的代数与辛 Gromov--Witten 不变量外，
还有一些相关不变量也值得研究。在第 \hyperref[chap:localization]{\ref*{chap:localization}}、\hyperref[chap:quantum-differential-equations]{\ref*{chap:quantum-differential-equations}} 章中，我们将遇到其中两种：
\emph{等变 Gromov--Witten 不变量}与\emph{引力关联数}。
它们在 Givental 关于镜像对称的工作中，以及相关思想和后续发展中，都起着关键作用。
\index{Gromov--Witten不变量@Gromov--Witten 不变量!等变不变量|hyperpage}
\index{引力关联数|hyperpage}

\hypertarget{source:219-batch5}{}
\section{计算 Gromov--Witten 不变量，I}\label{sec:computing-gw-i}
既然已经了解 Gromov--Witten 类与不变量的定义和性质，现在就来计算一些例子。
本节将讨论各种不同空间 $X$ 的 Gromov--Witten 不变量
$\langle I_{g,n,\beta}\rangle(\alpha_1,\ldots,\alpha_n)$。
大多数例子的亏格为 $g=0$ 或 $1$。鉴于 Gromov--Witten 不变量的直观解释~\eqref{eq:gw-pointed-maps}，
我们还将探讨它们的枚举意义。我们会发现，这些数与枚举几何之间的关系有时颇为复杂。

\hyperref[chap:quantum-cohomology]{第~\ref*{chap:quantum-cohomology} 章}将看到，量子上同调是由不变量 $\langle I_{0,n,\beta}\rangle$ 构造出来的。
该章还将介绍 \emph{Gromov--Witten 势}，它有一些非常有趣的性质，
尤其能帮助解释本节的若干例子。因此，本节实际上是讨论 Gromov--Witten 不变量的计算及其枚举意义的两节中的第一节。

\subsection{树级 Gromov--Witten 类}\label{subsec:tree-level-gw}
\index{树级 Gromov--Witten 不变量|hyperpage}\index{重构定理|hyperpage}
\index{Gromov--Witten类@Gromov--Witten 类!树级|hyperpage}
\index{重构定理|hyperpage}
当 $g=0$ 时，各类 $I_{0,n,\beta}$ 构成一个\emph{树级 Gromov--Witten 类系统}。
“树级”一词是说，亏格为零的稳定曲线是由若干 $\mathbb P^1$ 构成的树
（对应于量子场论中的树形 Feynman 图，这一术语就源自那里）。
由于始终要求 $n+2g\geq3$，树级类在 $n\geq3$ 时定义。
模空间 $\overline M_{0,n}$ 是维数为 $n-3$ 的轨形，其上同调已得到充分了解。
因此，对于 $I_{0,n,\beta}$，我们能说出许多事情。特别地，有下面这个 Kontsevich 与 Manin 的\emph{重构定理}。

\begin{theorem}\label{thm:gw-reconstruction}
设 $X$ 为光滑射影簇，且 $H^*(X,\mathbb Q)$ 由 $H^2(X,\mathbb Q)$ 生成。
再假设我们已经知道所有满足 $-\int_\beta\omega_X\leq\dim X+1$、$\deg\alpha_3=2$ 的
$\beta\in H_2(X,\mathbb Z)$ 所对应的 Gromov--Witten 不变量
$\langle I_{0,3,\beta}\rangle(\alpha_1,\alpha_2,\alpha_3)$。
那么，对所有 $\beta\in H_2(X,\mathbb Z)$，都能确定全部树级 Gromov--Witten 类
$I_{0,n,\beta}(\alpha_1,\ldots,\alpha_n)$。
\end{theorem}
\begin{proof}
先证明所有树级类都能由 Gromov--Witten 不变量重构。对 $n$ 作归纳；
由于 $\overline M_{0,3}$ 是一点，可假设 $n>3$。

为描述 $\overline M_{0,n}$ 的上同调，注意式~\eqref{eq:gw-splitting-gluing-map} 中映射 $\varphi$ 的像是
$\overline M_{0,n}$ 上的一个除子。对点 $p_1,\ldots,p_n$ 作置换，还能得到映射 $\varphi$ 的其他版本，
从而得到 $\overline M_{0,n}$ 上的其他除子。由 \SourceCite{Keel}，上同调环
$H^*(\overline M_{0,n},\mathbb Q)$ 由这些除子的类生成。
利用这一点，\SourceCite{KoM1} 证明：$H^*(\overline M_{0,n},\mathbb Q)$ 中不处于最高次数的上同调类，
由它与这些除子的交唯一确定。

现取一个类 $I_{0,n,\beta}(\alpha_1,\ldots,\alpha_n)$。若它处于最高次数，
它就是 $\langle I_{0,n,\beta}\rangle(\alpha_1,\ldots,\alpha_n)$ 乘以一点的类。
若不处于最高次数，则由它与上述除子的交确定。
而与 $\varphi$ 的像所给的除子相交时，得到 $\varphi^*I_{0,n,\beta}(\alpha_1,\ldots,\alpha_n)$。
由\IndexLocation{indexloc:0507:219}分裂公理，它由 $n$ 较小的 Gromov--Witten 类构成。
\SourcePage{220}
此外，等变性公理表明，与其他除子的交也具有同样的性质。
所以 $I_{0,n,\beta}(\alpha_1,\ldots,\alpha_n)$ 由 $n$ 较小的 Gromov--Witten 类确定。
由归纳假设，这些类由 Gromov--Witten 不变量确定，第一步证毕。

下一步是研究不变量 $\langle I_{0,n,\beta}\rangle$。
其思路是：上述除子之间的某些线性关系，给出不同 $n$、$\beta$ 所对应的
$\langle I_{0,n,\beta}\rangle$ 之间的二次关系。这些关系使我们能将任意 Gromov--Witten 不变量
用定理陈述中列出的不变量表示出来。尤其，定理中的不等式很容易解释。
若各 $\alpha_i$ 为齐次类，则 $\langle I_{0,3,\beta}\rangle(\alpha_1,\alpha_2,\alpha_3)=0$，除非
\[
\sum_i\deg\alpha_i=2\dim[\overline M_{0,3}(X,\beta)]^{\mathrm{virt}}
=2\left(-\int_\beta\omega_X+\dim X\right).
\]
结合显然成立的不等式 $\sum_i\deg\alpha_i\leq2\dim X+2$
（仍假设 $\deg\alpha_3=2$），就得到所需不等式。
论证的完整细节见 \SourceCite{KoM1}。下面说明当 $X=\mathbb P^2$ 时，这一方法如何运作。
\end{proof}

\subsection{射影平面的亏格零不变量}\label{subsec:plane-genus-zero}
\index{有理曲线!射影平面的计数|hyperpage}
\index{Gromov--Witten不变量@Gromov--Witten 不变量!射影平面|hyperpage}
我们将计算平面 $\mathbb P^2$ 的树级 Gromov--Witten 不变量。
这里 $\beta=d[\ell]$，其中 $\ell\subset\mathbb P^2$ 是一条直线，$d\geq0$。
为简化记号，用 $\langle I_{0,n,d}\rangle$ 代替 $\langle I_{0,n,\beta}\rangle$。
先计算 $\langle I_{0,3,1}\rangle([pt],[pt],[\ell])$。这很容易！
给定映射 $f:(\mathbb P^1,p_1,p_2,p_3)\to\mathbb P^2$，
使 $f(p_1)$ 与 $f(p_2)$ 是预先指定的点，且 $f(p_3)\in\ell$，便得到下图：
\begin{figure}[H]\centering
\includegraphics[width=0.46\textwidth]{7-4-plane-line.png}
\end{figure}
由于次数为 $1$，$f$ 的像就是由这两个点确定的直线。
又因为 $f(p_1)$、$f(p_2)$ 已固定，而 $f(p_3)$ 由交点 $\ell\cap f(\mathbb P^1)$ 确定，
我们得到一个\emph{唯一的映射}。随即容易看出，即使允许任意源曲线 $C$，
它也是唯一满足这些标记点条件的亏格零稳定映射 $f:(C,p_1,p_2,p_3)\to\mathbb P^2$。
此外，直接可见，这个稳定映射在 Gromov--Witten 不变量的定义~\eqref{eq:gw-moduli-integral} 中出现的重数为 $1$。
因此
\[
\langle I_{0,3,1}\rangle([pt],[pt],[\ell])=1.
\]
而且，依照定理~\ref{thm:gw-reconstruction}，这是唯一需要另行计算的不变量；
因为很容易检验，$\langle I_{0,3,1}\rangle([pt],[pt],[\ell])$ 是满足该定理条件的唯一非平凡不变量。

那么，其他不变量是什么？由\IndexLocation{indexloc:0505:220}点映射公理，可假设 $d\geq1$。
为计算 $\langle I_{0,n,d}\rangle(\alpha_1,\ldots,\alpha_n)$，可假设每个 $\alpha_i$ 都是
$[\mathbb P^2]$、$[\ell]$ 或 $[pt]$。又由式~\eqref{eq:gw-top-degree-condition}，有
\begin{equation}\label{eq:plane-degree-condition}
\sum_{i=1}^n\deg\alpha_i=6d-2+2n.
\end{equation}
当 $n=0,1$ 时，此方程无解；而当 $n=2,3$ 时，可以检验：不计置换，唯一的不变量是
\[
\langle I_{0,2,1}\rangle([pt],[pt])
=\langle I_{0,3,1}\rangle([pt],[pt],[\ell])=1.
\]
\SourcePage{221}
第一个等式由除子公理得到；更简单地说，$\langle I_{0,2,1}\rangle([pt],[pt])=1$，
因为在 $\mathbb P^2$ 中，经过两个点的直线唯一。

若 $n\geq4$，则由基本类公理，只要某个 $\alpha_i=[\mathbb P^2]$，就有
$\langle I_{0,n,d}\rangle(\alpha_1,\ldots,\alpha_n)=0$。
类似地，若 $n\geq4$ 且 $\alpha_n=[\ell]$，则由除子公理，
\[
\langle I_{0,n,d}\rangle(\alpha_1,\ldots,\alpha_{n-1},[\ell])
=d\langle I_{0,n-1,d}\rangle(\alpha_1,\ldots,\alpha_{n-1}),
\]
从而可以归纳计算。尚待计算的只有不变量 $\langle I_{0,n,d}\rangle([pt],\ldots,[pt])$。
由于 $\deg[pt]=4$，式~\eqref{eq:plane-degree-condition} 给出 $n=3d-1$。
因此，我们要计算的是
\begin{equation}\label{eq:plane-nd}
N_d=\langle I_{0,3d-1,d}\rangle(\underbrace{[pt],\ldots,[pt]}_{3d-1\text{ 次}}),\qquad d\geq1.
\end{equation}
由于 $\mathbb P^2$ 上有可迁的线性群作用，定理~\ref{thm:homogeneous-gw-equivalence} 表明，
$N_d$ 是 $\mathbb P^2$ 中经过 $3d-1$ 个一般位置点的 $d$ 次有理曲线的数目。

这些数 $N_d$ 也可定义为 $d$ 次有理平面曲线的 Severi 簇的次数。
然而，传统的代数几何方法此前只能给出最初几个数。
这里的新进展是：将 $N_d$ 解释为 Gromov--Witten 不变量，可以得到递推公式
\begin{equation}\label{eq:plane-recursion}
N_d=\sum_{\substack{d=d_1+d_2\\d_1,d_2>0}}N_{d_1}N_{d_2}
\left(d_1^2d_2^2\binom{3d-4}{3d_1-2}-d_1^3d_2\binom{3d-4}{3d_1-1}\right).
\end{equation}
由于 $N_1=\langle I_{0,2,1}\rangle([pt],[pt])=1$，此关系给出
\[
N_2=1,\quad N_3=12,\quad N_4=620,\quad N_5=87304,\quad\ldots.
\]
这说明稳定映射往往比嵌入曲线更容易处理，甚至在回答关于嵌入曲线的经典问题时也是如此！

式~\eqref{eq:plane-recursion} 的证明采用定理~\ref{thm:gw-reconstruction} 证明中概述的策略。
从下面的 Gromov--Witten 类出发：
\begin{equation}\label{eq:plane-xi}
\xi=I_{0,3d,d}(\underbrace{[pt],\ldots,[pt]}_{3d-2\text{ 次}},[\ell],[\ell])
\in H^*(\overline M_{0,3d},\mathbb Q).
\end{equation}
由次数公理，该类的次数为 $2\dim\overline M_{0,3d}-2$，所以它与一个除子的交是良定义的有理数。

幸运的是，$\overline M_{0,3d}$ 有许多有趣的除子。
用集合 $\mathcal I=\{1,\ldots,3d-4,p,q,r,s\}$ 为 $\overline M_{0,3d}$ 中稳定曲线的 $3d$ 个标记点编号。
用式~\eqref{eq:plane-xi} 的记号，这意味着 $\alpha_1=\cdots=\alpha_{3d-4}=\alpha_p=\alpha_q=[pt]$，
而 $\alpha_r=\alpha_s=[\ell]$。现在将 $\mathcal I$ 分成不相交的集合 $A$ 与 $B$。
与分裂公理中一样，这给出自然映射
\[
\varphi_{A,B}:\overline M_{0,|A|+1}\times\overline M_{0,|B|+1}\longrightarrow\overline M_{0,3d}.
\]
$\varphi_{A,B}$ 的像是 $\overline M_{0,3d}$ 上的一个除子。此外，还有下面这个基本的除子线性等价关系：
\begin{equation}\label{eq:plane-boundary-relation}
\sum_{\substack{r,s\in A\\p,q\in B}}\operatorname{Im}(\varphi_{A,B})
\sim\sum_{\substack{p,r\in A\\q,s\in B}}\operatorname{Im}(\varphi_{A,B}).
\end{equation}
\SourcePage{222}
这个关系来自自然映射 $\phi:\overline M_{0,3d}\to\overline M_{0,4}$，
它遗忘最初的 $3d-4$ 个点，并收缩不稳定分支。在此映射下，式~\eqref{eq:plane-boundary-relation}
的左边是 $\phi$ 在 $\overline M_{0,4}$ 中下面这条稳定曲线上的纤维：
\begin{figure}[H]\centering
\includegraphics[width=0.31250\textwidth]{7-4-four-point-boundary.png}
\end{figure}
类似地，式~\eqref{eq:plane-boundary-relation} 的右边是 $\phi$ 在另一条稳定曲线上的纤维，
其中 $p,r$ 位于同一分支，$q,s$ 位于另一分支。交比给出
$M_{0,4}\simeq\mathbb P^1\setminus\{0,1,\infty\}$，上图中的稳定曲线是
$\overline M_{0,4}\simeq\mathbb P^1$ 的三个边界点之一。
由于这些边界点显然线性等价，$\phi$ 在这些点上的纤维也线性等价，式~\eqref{eq:plane-boundary-relation} 因而成立。

证明的下一步，是将式~\eqref{eq:plane-boundary-relation} 的两边分别与式~\eqref{eq:plane-xi}
定义的 Gromov--Witten 类 $\xi$ 相交。先看左边。$\xi$ 与 $\operatorname{Im}(\varphi_{A,B})$
的交对应于拉回 $\varphi_{A,B}^*\xi$，可用分裂公理计算。
在\hyperref[sec:gw-properties]{第~\ref*{sec:gw-properties} 节}中，我们只对指标集 $\mathcal I$
的一个特定分划陈述了分裂公理；但由等变性公理，它适用于所有分划。
因此，由于 $\mathbb P^2\times\mathbb P^2$ 中的对角线类是
\[
\Delta=[pt]\otimes[\mathbb P^2]+[\ell]\otimes[\ell]+[\mathbb P^2]\otimes[pt],
\]
分裂公理与等变性公理表明，交数 $\xi\cdot\operatorname{Im}(\varphi_{A,B})$ 由下式给出：
\begin{equation}\label{eq:plane-splitting-intersection}
\begin{aligned}
\sum_{d=d_1+d_2}\Big(&\langle I_{0,|A|+1,d_1}\rangle(\alpha_i,\ i\in A,[pt])
\cdot\langle I_{0,|B|+1,d_2}\rangle(\alpha_i,\ i\in B,[\mathbb P^2])\\
&+\langle I_{0,|A|+1,d_1}\rangle(\alpha_i,\ i\in A,[\ell])
\cdot\langle I_{0,|B|+1,d_2}\rangle(\alpha_i,\ i\in B,[\ell])\\
&+\langle I_{0,|A|+1,d_1}\rangle(\alpha_i,\ i\in A,[\mathbb P^2])
\cdot\langle I_{0,|B|+1,d_2}\rangle(\alpha_i,\ i\in B,[pt])\Big),
\end{aligned}
\end{equation}
因为 $\varphi_{A,B}^*\xi$ 是最高次数的类。

虽然公式看起来复杂，但由 Gromov--Witten 不变量的次数条件，多数项为零。
例如，考虑不变量 $\langle I_{0,|A|+1,d_1}\rangle(\alpha_i,\ i\in A,\gamma)$，
其中 $\gamma$ 是 $[pt]$、$[\ell]$、$[\mathbb P^2]$ 之一。
由对 $A$ 的条件，$\alpha_r=\alpha_s=[\ell]$，其余 $\alpha_i=[pt]$。
所以式~\eqref{eq:gw-top-degree-condition} 表明，此不变量为零，除非
\[
|A|=3d_1+2-\tfrac12\deg\gamma.
\]
由于 $|A|$ 固定，这个方程确定了 $d_1$ 与 $\deg\gamma$。
因此，式~\eqref{eq:plane-splitting-intersection} 至多只有一项非零。

现假设 $|A|,|B|\geq3$。由基本类公理，
$\langle I_{0,|A|+1,d_1}\rangle(\alpha_i,\ i\in A,[\mathbb P^2])
=\langle I_{0,|B|+1,d_2}\rangle(\alpha_i,\ i\in B,[\mathbb P^2])=0$。
所以非零项必须满足 $\gamma=[\ell]$，进而 $|A|=3d_1+1$、$|B|=3d_2-1$。
于是式~\eqref{eq:plane-splitting-intersection} 化为
\[
\langle I_{0,3d_1+2,d_1}\rangle(\underbrace{[pt],\ldots,[pt]}_{3d_1-1\text{ 次}},[\ell],[\ell],[\ell])
\cdot\langle I_{0,3d_2,d_2}\rangle(\underbrace{[pt],\ldots,[pt]}_{3d_2-1\text{ 次}},[\ell]).
\]
\SourcePage{223}
由除子公理与 $N_d$ 的定义，上式化为 $d_1^3d_2N_{d_1}N_{d_2}$。因此
\begin{equation}\label{eq:plane-boundary-positive}
\xi\cdot\operatorname{Im}(\varphi_{A,B})=d_1^3d_2N_{d_1}N_{d_2},\qquad |A|=3d_1+1.
\end{equation}
还要注意，此时 $d_1>0$、$d_2>1$。

另一方面，若 $|B|=2$，请读者证明，式~\eqref{eq:plane-splitting-intersection} 的非零项为
\[
\langle I_{0,3d-1,d-1}\rangle(\underbrace{[pt],\ldots,[pt]}_{3d-4\text{ 次}},[\ell],[\ell],[\ell])
\cdot\langle I_{0,3,1}\rangle([pt],[pt],[\ell]).
\]
利用除子公理，它化为 $(d-1)^3N_{d-1}N_1$，
即 $d_1=d-1$、$d_2=1$ 时的式~\eqref{eq:plane-boundary-positive}。
最后，若 $|A|=2$，则非零项为
\[
\langle I_{0,3,0}\rangle([\ell],[\ell],[\mathbb P^2])
\cdot\langle I_{0,3d-1,d}\rangle(\underbrace{[pt],\ldots,[pt]}_{3d-2\text{ 次}},[pt]).
\]
由\IndexLocation{indexloc:0505:223}点映射公理，乘积中的第一个因子为 $1$，第二个因子为 $N_d$。于是
\begin{equation}\label{eq:plane-boundary-zero}
\xi\cdot\operatorname{Im}(\varphi_{A,B})=N_d,\qquad |A|=2,\ d_1=0,\ d_2=d.
\end{equation}
现在容易看出，类 $\xi$ 与式~\eqref{eq:plane-boundary-relation} 左边的交为
\[
N_d+\sum_{\substack{d=d_1+d_2\\d_1,d_2>0}}N_{d_1}N_{d_2}d_1^3d_2\binom{3d-4}{3d_1-1}.
\]
为说明原因，考虑指标集 $\mathcal I=\{1,\ldots,3d-4,p,q,r,s\}$ 的分划 $A,B$，
其中 $r,s\in A$、$p,q\in B$。当 $A,B$ 变化时，式~\eqref{eq:plane-boundary-zero}
恰好出现一次，即 $A=\{r,s\}$ 时。
而当 $d=d_1+d_2$、$d_1,d_2>0$ 时，式~\eqref{eq:plane-boundary-positive}
恰好出现 $\binom{3d-4}{3d_1-1}$ 次；因为 $|A|=3d_1+1$ 意味着，
除 $r,s$ 之外，$A$ 还有 $3d_1-1$ 个元素，选自 $\mathcal I$ 最初的 $3d-4$ 个元素。

为完成式~\eqref{eq:plane-recursion} 的证明，可用类似论证说明，
$\xi$ 与式~\eqref{eq:plane-boundary-relation} 右边的交为
\[
\sum_{\substack{d=d_1+d_2\\d_1,d_2>0}}N_{d_1}N_{d_2}d_1^2d_2^2\binom{3d-4}{3d_1-2}.
\]
细节留给读者，这是运用这些公理的一道很好的练习。
由线性等价关系~\eqref{eq:plane-boundary-relation}，$\xi$ 与两边的交相同，
递推关系~\eqref{eq:plane-recursion} 立即成立。

在上述由 Kontsevich 给出的论证中，请注意：将式~\eqref{eq:plane-boundary-relation}
与 $\xi$ 相交，一边产生了 $N_d$，另一边却没有。
这对式~\eqref{eq:plane-recursion} 的证明至关重要，而这一现象有更广泛的适用性。
事实上，\IndexLocation{indexloc:1208:223}重构定理的证明也以类似方式进行：将式~\eqref{eq:plane-boundary-relation}
与各种 Gromov--Witten 类相交，说明如何用定理~\ref{thm:gw-reconstruction}
中列出的特殊不变量表示一般的树级不变量。

还可用其他方法证明 $N_d$ 的递推关系。例如，\hyperref[chap:quantum-cohomology]{第~\ref*{chap:quantum-cohomology} 章}将看到，量子上同调这一套工具
如何使上述论证自动化并加以简化。尤其，最初的严格证明 \SourceCite{RT1}
使用了量子上同调的辛几何表述。
\SourcePage{224}
也能给出更几何化的论证，省去 Gromov--Witten 类的形式体系，而遵循上面所做的基本思路，见 \SourceCite{FP}。
另外，还可以用其他代数几何方法给出完全自足的论证，见 \SourceCite{Ran1,CH1,CH2,Vakil3}。
关于数 $N_d$ 的推广，见 \SourceCite{DI,EK1}。

最后还应提到，这些方法对其他具有可迁线性群作用的簇，也能给出很好的结果。
例如，\SourceCite{FP} 计算了 $\mathbb P^3$ 与光滑二次三维簇的树级 Gromov--Witten 不变量，
$\mathbb P^n$ 的情形则在 \SourceCite{Vakil3,GPa} 中讨论。其他例子见 \SourceCite{DI}。

\subsection{Calabi--Yau 三维簇的亏格零不变量}\label{subsec:cy-genus-zero}
\index{Clemens猜想@Clemens 猜想|hyperpage}
我们将讨论 Calabi--Yau 三维簇的树级 Gromov--Witten 不变量。
这里 $V$ 表示一个光滑 Calabi--Yau 三维簇。对非零 $\beta\in H_2(V,\mathbb Z)$，置 $N_\beta=\langle I_{0,0,\beta}\rangle$。
\begin{proposition}\label{prop:cy-gw-determined}
$V$ 的亏格零 Gromov--Witten 不变量完全由各 $N_\beta$ 及 $V$ 上的交数确定。
\end{proposition}
\begin{proof}
只需考虑各 $\alpha_i$ 都是齐次类的不变量 $\langle I_{0,n,\beta}\rangle(\alpha_1,\ldots,\alpha_n)$。
由\IndexLocation{indexloc:0498:224}次数公理，可假设 $\sum_i\deg\alpha_i=2n$。
由于每个 $\alpha_i$ 或为基本类，或满足 $\deg\alpha_i\geq2$，需考虑两种情形：
\begin{itemize}
\item 某个 $\alpha_i$ 是基本类；
\item 所有 $i$ 都满足 $\deg\alpha_i=2$。
\end{itemize}
在第一种情形，由基本类公理、\IndexLocation{indexloc:0501:224}等变性公理及点映射公理，Gromov--Witten 不变量是一个交数，或为 $0$。
在第二种情形，当 $\beta\ne0$ 时，除子公理给出
\[
\langle I_{0,n,\beta}\rangle(\alpha_1,\ldots,\alpha_n)=N_\beta\int_\beta\alpha_1\cdots\int_\beta\alpha_n.
\]
而当 $\beta=0$ 时，式~\eqref{eq:gw-degree-zero-three-point} 表明，$n\ne3$ 时 $\langle I_{0,n,0}\rangle=0$，且
\[
\langle I_{0,3,0}\rangle(\alpha_1,\alpha_2,\alpha_3)=\int_V\alpha_1\cup\alpha_2\cup\alpha_3.
\]
命题证毕。
\end{proof}
\SourceExampleNumber{7.4.3}{1}
\begin{example}\label{ex:quintic-gw-types}
设 $V\subset\mathbb P^4$ 是一个一般的五次三维簇，取直线 $\ell\subset V$
（这样的直线共有 $2875$ 条）。用 $H\in H^2(V,\mathbb Z)$ 表示超平面类。
那么，不计置换，唯一可能非零的 Gromov--Witten 不变量是
\[
\langle I_{0,3,0}\rangle(\alpha_1,\alpha_2,1),\qquad\langle I_{0,n,d}\rangle(H,\ldots,H).
\]
其中，为简化记号，用 $\langle I_{0,n,d}\rangle$ 代替 $\langle I_{0,n,d[\ell]}\rangle$，而 $n,d\geq0$ 任意。
由基本类公理，$\langle I_{0,3,0}\rangle(\alpha_1,\alpha_2,1)=\int_V\alpha_1\cup\alpha_2$。
置 $N_d=\langle I_{0,0,d}\rangle$，则当 $d>0$ 时，除子公理给出
$\langle I_{0,n,d}\rangle(H,\ldots,H)=d^nN_d$。
最后，由点映射公理，$\langle I_{0,3,0}\rangle(H,H,H)=\int_VH^3=5$。
\end{example}
对于五次三维簇 $V$，上述例子将求 Gromov--Witten 不变量的问题归结为确定
$N_d=\langle I_{0,0,d}\rangle$。由例题~\ref{ex:quintic-gw-euler-integral}，
\[
N_d=\langle I_{0,0,d}\rangle=\int_{\overline M_{0,0}(\mathbb P^4,d)}c_{5d+1}(\mathcal V_d).
\]
\SourcePage{225}
当然，我们希望用 $V$ 上的有理曲线来计算 $N_d$。
由于 Clemens 猜想以及多重覆盖的存在，这仍是一个未解决的问题。
我们将分别讨论这两点，并说明它们与计算 $N_d$ 的问题有何关系。

如\hyperref[chap:quintic]{第~\ref*{chap:quintic} 章}所讨论，镜像对称的一个引人注目的方面，是它对五次三维簇 $V$ 上的有理曲线作出的预测。
困难在于，我们一般并不知道 $V$ 上给定次数的有理曲线是否只有有限多条。
这就引出了如下的 Clemens 猜想。
\index{Clemens猜想@Clemens 猜想|hyperpage}
\begin{conjecture}\label{conj:clemens}
设 $V\subset\mathbb P^4$ 是一个一般的五次三维簇。则对每个次数 $d\geq1$：
\begin{enumerate}[label=(\roman*)]
\item $V$ 中 $d$ 次不可约有理曲线 $C\subset V$ 只有有限多条。
\item 遍历所有次数时，这些曲线彼此不相交。
\item 若 $f:\mathbb P^1\to C$ 是不可约有理曲线 $C\subset V$ 的正规化，
则法丛 $N_f$ 同构于 $\mathcal O_{\mathbb P^1}(-1)\oplus\mathcal O_{\mathbb P^1}(-1)$。
\end{enumerate}
\end{conjecture}
注意，这里并未断言有理曲线都是光滑的。
那样的断言是错误的，因为一般 $V$ 中存在具有六个结点的有理平面五次曲线 $C\subset V$，见 \SourceCite{Vainsencher}。
目前已知，$V$ 上次数不超过 $9$ 的有理曲线只有有限多条。
而且，正如猜想预测的，它们都彼此不相交，且除上述平面五次曲线外，都非奇异。
此外，光滑曲线的法丛都是 $\mathcal O_{\mathbb P^1}(-1)\oplus\mathcal O_{\mathbb P^1}(-1)$。
这些结果均在 \SourceCite{Katz2,JK1} 中证明。
\hyperref[subsec:nodal-quintic-double-covers]{第~\ref*{subsec:nodal-quintic-double-covers} 小节}将说明，$V$ 中的奇异五次曲线也具有预测的法丛，因此猜想~\ref{conj:clemens} 对 $d\leq9$ 成立。

但即使 Clemens 猜想对所有 $d$ 都成立，仍有\emph{多重覆盖}的问题。
设 $C\subset V$ 是一条 $d/k$ 次不可约有理曲线，其正规化映射为 $f:\mathbb P^1\to V$。
将 $f$ 与任意 $k$ 次映射 $g:\mathbb P^1\to\mathbb P^1$ 复合，便得到稳定映射
$f\circ g:\mathbb P^1\to V$，满足 $(f\circ g)_*[\mathbb P^1]=k(d/k)=d$。
由于 $g$ 在一个族中变化，$k>1$ 时就得到正维的稳定映射族。
这说明 $\overline M_{0,0}(V,d)$ 为正维，尽管由于 $V$ 是 Calabi--Yau 三维簇，预期维数~\eqref{eq:stable-map-expected-dimension} 为 $0$。

我们需要确定上述多重覆盖族对虚拟基本类 $[\overline M_{0,0}(V,d)]^{\mathrm{virt}}$，从而对 $N_d$ 的贡献。
\hyperref[subsec:instanton-numbers]{第~\ref*{subsec:instanton-numbers} 小节}将看到：
若 $C\subset V$ 是 $d/k$ 次光滑有理曲线，法丛为
$N\simeq\mathcal O_{\mathbb P^1}(-1)\oplus\mathcal O_{\mathbb P^1}(-1)$，
则其 $k$ 次多重覆盖对 $N_d$ 的贡献为 $k^{-3}$。

暂且作一个朴素假设：对所有 $d$，$V$ 上的 $d$ 次有理曲线只有有限多条，
且全都光滑、彼此不相交、法丛为 $\mathcal O_{\mathbb P^1}(-1)\oplus\mathcal O_{\mathbb P^1}(-1)$。
若 $n_d$ 是这样的 $d$ 次曲线的总数，则上述讨论提示
\begin{equation}\label{eq:quintic-multiple-cover}
N_d=\sum_{k\mid d}n_{d/k}k^{-3},
\end{equation}
因为每条 $d/k$ 次有理曲线对 $N_d$ 都贡献 $k^{-3}$。
利用\IndexLocation{indexloc:0499:225}除子公理，式~\eqref{eq:quintic-multiple-cover} 给出
\begin{equation}\label{eq:quintic-three-divisors}
\langle I_{0,3,d}\rangle(H,H,H)=d^3\langle I_{0,0,d}\rangle=d^3N_d
=\sum_{k\mid d}n_{d/k}\left(\frac d k\right)^3=\sum_{k\mid d}n_kk^3,
\end{equation}
\SourcePage{226}
其中 $d\geq1$。这使我们能与\hyperref[chap:quintic]{第~\ref*{chap:quintic} 章}的公式建立如下联系。
回忆式~\eqref{eq:quintic-a}，A 模型关联函数为
\begin{equation}\label{eq:quintic-a-instanton}
\langle H,H,H\rangle=5+\sum_{d=1}^{\infty}n_dd^3\frac{q^d}{1-q^d}.
\end{equation}
由于此式右边可以写成
\[
5+\sum_{d=1}^{\infty}\left(\sum_{k\mid d}n_kk^3\right)q^d,
\]
上面关于 $\langle I_{0,3,d}\rangle(H,H,H)$ 的公式表明
\begin{equation}\label{eq:quintic-a-gw}
\langle H,H,H\rangle=\sum_{d=0}^{\infty}\langle I_{0,3,d}\rangle(H,H,H)q^d,
\end{equation}
这里用到了例题~\ref{ex:quintic-gw-types} 的 $\langle I_{0,3,0}\rangle(H,H,H)=5$。

这是一个重要进展：虽然上述多数等式依赖各种假设，但给出 $\langle H,H,H\rangle$
的式~\eqref{eq:quintic-a-gw} 有严格的定义。这意味着，我们终于扎实地理解了五次三维簇的 A 模型关联函数。
此外，\hyperref[chap:mirror-theorem]{第~\ref*{chap:mirror-theorem} 章}证明式~\eqref{eq:quintic-mirror-identity} 中 $\langle H,H,H\rangle$ 的精彩公式时，将使用这个严格定义的 $\langle H,H,H\rangle$。

从上述讨论可以得出两个结论。第一，Gromov--Witten 不变量 $N_d=\langle I_{0,0,d}\rangle$
正是表述五次三维簇的镜像对称所需要的量。
第二，我们仍有\emph{大量}工作要做，因为还需要在不依赖上述所有假设的情况下，将 $N_d$ 与 $V$ 上的有理曲线联系起来。
因此，应研究式~\eqref{eq:quintic-multiple-cover} 真正意味着什么。我们很快会看到，答案比预想的复杂。

\subsection{Calabi--Yau 三维簇的瞬子数}\label{subsec:instanton-numbers}
\index{瞬子数|hyperpage}\index{多重覆盖公式|hyperpage}
\index{瞬子数|hyperpage}\index{多重覆盖|hyperpage}
为更好地理解 $N_d$ 与五次三维簇上有理曲线的关系，我们采取如下策略：
先定义\emph{瞬子数} $n_d$，使式~\eqref{eq:quintic-multiple-cover} 自动成立；
再处理更精细的问题，将 $n_d$ 与 $V$ 上 $d$ 次有理曲线的实际数目联系起来。

考虑如下的一般情形。设 $V$ 是光滑 Calabi--Yau 三维簇，$\beta\in H_2(V,\mathbb Z)$ 是非零有效类。
那么，容易看出存在唯一的有理数 $n_\beta$，称为\emph{瞬子数}，使
\begin{equation}\label{eq:instanton-definition}
N_\beta=\langle I_{0,0,\beta}\rangle=\sum_{\beta=k\gamma}n_\gamma k^{-3}.
\end{equation}
求和遍历满足 $\beta=k\gamma$ 的所有正整数 $k$ 与非零有效类 $\gamma\in H_2(V,\mathbb Z)$。
$n_\beta$ 的存在唯一性如下证明：若 $\beta$ 在 $H_2(V,\mathbb Z)$ 中不可整除，直接置
$n_\beta=\langle I_{0,0,\beta}\rangle$；否则，利用式~\eqref{eq:instanton-definition} 归纳定义 $n_\beta$。
这证明了存在性，唯一性也随即得到。

有了瞬子数 $n_\beta$ 的这一定义，首先要说明式~\eqref{eq:instanton-definition} 中因子 $k^{-3}$ 的由来。
纯粹从形式上看，可以说它给出了漂亮的公式。例如，设 $\beta\in H_2(V,\mathbb Z)$
是非零有效类，$\alpha_1,\alpha_2,\alpha_3\in H^2(V,\mathbb C)$。
\SourcePage{227}
由式~\eqref{eq:instanton-definition} 与除子公理，
\begin{equation}\label{eq:instanton-three-divisors}
\begin{aligned}
\langle I_{0,3,\beta}\rangle(\alpha_1,\alpha_2,\alpha_3)
&=N_\beta\int_\beta\alpha_1\int_\beta\alpha_2\int_\beta\alpha_3\\
&=\left(\sum_{\beta=k\gamma}n_\gamma k^{-3}\right)
\int_\beta\alpha_1\int_\beta\alpha_2\int_\beta\alpha_3\\
&=\sum_{\beta=k\gamma}n_\gamma\int_\gamma\alpha_1\int_\gamma\alpha_2\int_\gamma\alpha_3.
\end{aligned}
\end{equation}
最后一个等式来自 $\beta=k\gamma$ 时的 $k^{-1}\int_\beta u=\int_\gamma u$。
注意，这推广了式~\eqref{eq:quintic-three-divisors}。

\hyperref[chap:quantum-cohomology]{第~\ref*{chap:quantum-cohomology} 章}将严格表述 A 模型关联函数这一数学概念。对于 Calabi--Yau 三维簇 $V$，关联函数为
\[
\langle\alpha_1,\alpha_2,\alpha_3\rangle
=\sum_\beta\langle I_{0,3,\beta}\rangle(\alpha_1,\alpha_2,\alpha_3)q^\beta,
\]
其中 $\alpha_i\in H^2(V,\mathbb C)$，而 $q^\beta=e^{2\pi i\int_\beta\omega}$，
$\omega$ 是复化 Kähler 类。利用点映射公理与式~\eqref{eq:instanton-three-divisors}，得到
\begin{equation}\label{eq:cy-a-correlation-instanton}
\langle\alpha_1,\alpha_2,\alpha_3\rangle
=\int_V\alpha_1\cup\alpha_2\cup\alpha_3
+\sum_{\substack{\beta\ne0\\\beta\text{ 有效}}}n_\beta\frac{q^\beta}{1-q^\beta}
\int_\beta\alpha_1\int_\beta\alpha_2\int_\beta\alpha_3.
\end{equation}
推导方式与式~\eqref{eq:quintic-a-gw} 类似。此公式在镜像对称中起重要作用，
并与量子上同调密切相关。\hyperref[subsec:a-model-correlation]{第~\ref*{subsec:a-model-correlation} 小节}将更详细地讨论它。

式~\eqref{eq:instanton-definition} 中的因子 $k^{-3}$ 还有更深刻的理由。
\hyperref[subsec:cy-genus-zero]{第~\ref*{subsec:cy-genus-zero} 小节}已略作提及，现在作更充分的说明。
设 $C\subset V$ 为光滑有理曲线，法丛为
$N\simeq\mathcal O_{\mathbb P^1}(-1)\oplus\mathcal O_{\mathbb P^1}(-1)$，
并置 $\gamma=[C]\in H_2(V,\mathbb Z)$。
我们要计算 $C$ 的 $k$ 次多重覆盖对 Gromov--Witten 不变量
$N_\beta=\langle I_{0,0,\beta}\rangle$ 的贡献，其中 $\beta=k\gamma$。

固定同构 $f:\mathbb P^1\simeq C$。则 $k$ 次覆盖 $g:\mathbb P^1\to\mathbb P^1$
给出稳定映射 $f\circ g\in\overline M_{0,0}(V,\beta)$。
由于 $g$ 自然属于 $\overline M_{0,0}(\mathbb P^1,k)$，得到嵌入
$\overline M_{0,0}(\mathbb P^1,k)\hookrightarrow\overline M_{0,0}(V,\beta)$。
记这个嵌入的像为 $M_{k,C}(V)$。我们断言，$M_{k,C}(V)$ 实际上是 $\overline M_{0,0}(V,\beta)$ 的一个连通分支。

为说明这一点，考虑在连通参数空间 $S$ 上、同调类为 $\beta$ 的任意亏格零稳定映射族
$f:\mathcal C\to V$，使 $0\in S$ 上的稳定映射属于 $M_{k,C}(V)$。
记结构态射为 $\pi:\mathcal C\to S$。关于 $C$ 的法丛 $N$ 的假设保证存在解析收缩映射
$\rho:V\to V'$，将 $C$ 收缩为一点 $p$，且在 $C$ 与 $p$ 之外是同构。
考虑复合映射 $\widetilde f=\rho\circ f$，注意 $\widetilde f(\pi^{-1}(0))$ 是点 $p$。
由 \SourceCite[引理 1.5]{CKM}，容易得到，对所有 $s\in S$，$\widetilde f(\pi^{-1}(s))$ 都是一点。
由于 $f_*[\pi^{-1}(s)]=\beta$，该点也必须是 $p$。断言得证。

虚拟基本类 $[\overline M_{0,0}(V,\beta)]^{\mathrm{virt}}$ 是零维的；
按定义，$\langle I_{0,0,\beta}\rangle$ 是它的次数。
曲线 $C$ 对 $\langle I_{0,0,\beta}\rangle$ 的贡献如下。
\begin{theorem}\label{thm:multiple-cover}
在上述关于 $C$ 的假设下，设 $\beta=k\gamma=k[C]$，
并令 $M_{k,C}(V)$ 为 $\overline M_{0,0}(V,\beta)$ 中包含 $C$ 的所有 $k$ 次多重覆盖的分支。
\SourcePage{228}
则虚拟基本类 $[\overline M_{0,0}(V,\beta)]^{\mathrm{virt}}$ 中支撑在 $M_{k,C}(V)$ 上的部分，
对 $N_\beta=\langle I_{0,0,\beta}\rangle$ 的贡献为 $k^{-3}$。
\end{theorem}
\begin{proof}
从历史上看，这一结果最初表述为以下断言：$C$ 的 $k$ 次多重覆盖对
$\langle I_{0,3,\beta}\rangle(\alpha_1,\alpha_2,\alpha_3)$ 的贡献是
\begin{equation}\label{eq:multiple-cover-three-point}
\int_\gamma\alpha_1\int_\gamma\alpha_2\int_\gamma\alpha_3,
\end{equation}
其中 $\alpha_i\in H^2(V,\mathbb C)$。
这一断言与定理~\ref{thm:multiple-cover} 的等价性可由式~\eqref{eq:instanton-three-divisors} 的推导得到。
\SourceCite{CdGP} 写出三点函数~\eqref{eq:quintic-a-instanton} 的方式，已经隐含地提出了这一猜想；
Aspinwall 与 Morrison 在 \SourceCite{AM1} 中给出了证明概要。
现在必须将他们的论证视为不完整的，因为它没有与后来任何一种 Gromov--Witten 不变量的定义作直接比较。

下面概述 Voisin 的证明 \SourceCite{Voisin2}。
与\hyperref[sec:gw-symplectic-definition]{第~\ref*{sec:gw-symplectic-definition} 节}一样，
将 $V$ 的复结构形变为一般的几乎复结构 $J$，再加入一般的非齐次项 $\nu$。
考虑扰动的 $J$ 全纯映射 $g:\mathbb P^1\to V$，满足 $g_*[\mathbb P^1]=\beta$。
用该节的记号，它们构成空间 $\mathcal M(\mathbb P^1,J,\nu,\beta)$。
若 $J$ 足够接近原复结构、$\nu$ 足够小，而 $U$ 是 $V$ 中 $C$ 的一个足够小的邻域，
则由像位于 $U$ 中的映射构成的子集 $\mathcal M^U(\mathbb P^1,J,\nu,\beta)$
是 $\mathcal M(\mathbb P^1,J,\nu,\beta)$ 的一个分支（未必连通）。
固定三个点 $p_1,p_2,p_3\in\mathbb P^1$，得到赋值映射
\[
\pi_1:\mathcal M^U(\mathbb P^1,J,\nu,\beta)\longrightarrow V^3,
\qquad g\longmapsto(g(p_1),g(p_2),g(p_3)).
\]
此映射的像为六维，具有自然定向，并能紧化为边界维数至多为四的空间。
因此，$\pi_1$ 的像给出 $H_6(V^3,\mathbb Z)$ 中一个良定义的同调类。
\SourceCite{Voisin2} 的主要结果是：由于 $C$ 的嵌入是刚性的，该类恰好是 $C\times C\times C$ 的同调类。

假定这一结果，取一般位置的 $\Gamma_1,\Gamma_2,\Gamma_3$，分别为
$\alpha_1,\alpha_2,\alpha_3$ 的 Poincaré 对偶。
由于 $C\times C\times C$ 与 $\Gamma_1\times\Gamma_2\times\Gamma_3$ 相交的点数
由式~\eqref{eq:multiple-cover-three-point} 给出（记得 $[C]=\gamma$），
$\Gamma_1\times\Gamma_2\times\Gamma_3$ 与 $\pi_1$ 的像相交也得到同样的数。
于是，辛 Gromov--Witten 不变量的计数定义表明，
对 $\langle I_{0,3,\beta}\rangle(\alpha_1,\alpha_2,\alpha_3)$ 的贡献正是式~\eqref{eq:multiple-cover-three-point}。

定理~\ref{thm:multiple-cover} 的一些较新证明直接处理 $N_\beta$，说明 $C$ 的 $k$ 次覆盖贡献 $k^{-3}$。
\hyperref[subsec:multiple-cover-localization]{第~\ref*{subsec:multiple-cover-localization} 小节}将介绍 Kontsevich 的证明方法；我们将详细处理 $d=2$ 的情形，
一般情形则参见 \SourceCite{Manin2}。
不过，该小节还将给出 Pandharipande 的一个简化论证。
此定理的其他证明见 \SourceCite{LLY,Gathmann2}；\hyperref[sec:quintic-mirror-theorem]{第~\ref*{sec:quintic-mirror-theorem} 节}将讨论 \SourceCite{LLY} 的证明。
\end{proof}
接下来讨论两个例子，说明瞬子数 $n_\beta$ 的枚举意义有时相当微妙。
即使对于五次三维簇，事情也并非如人们所希望的那样简单。

\SourceExampleNumber{7.4.4}{1}
\begin{example}\label{ex:quintic-instanton-enumerative}
设 $V\subset\mathbb P^4$ 是一个一般的五次三维簇。
将瞬子数 $n_\beta$ 记为 $n_d$，其中 $\beta=d[\ell]$，$\ell\subset V$ 为直线。
这些数正是式~\eqref{eq:quintic-multiple-cover}、\eqref{eq:quintic-a-instanton} 中出现的数，
并由镜像对称通过\hyperref[chap:quintic]{第~\ref*{chap:quintic} 章}的公式计算。
\SourcePage{229}
当 $d\leq9$ 时，$n_d$ 与 $d$ 次有理曲线的关系尤其好：
\begin{equation}\label{eq:quintic-instanton-small-degree}
d\leq9\quad\Longrightarrow\quad n_d=V\text{ 中 }d\text{ 次有理曲线的数目}.
\end{equation}
由于 Clemens 猜想对 $d\leq9$ 成立，此式可由定理~\ref{thm:multiple-cover} 得到。

下面更详细地解释原因。由于猜想~\ref{conj:clemens} 在 $d\leq9$ 时成立，
式~\eqref{eq:quintic-instanton-small-degree} 的右边有限。
用 $n'_d$ 表示 $V$ 中 $d$ 次有理曲线的数目。
如\hyperref[subsec:cy-genus-zero]{第~\ref*{subsec:cy-genus-zero} 小节}所述，
$d\leq9$ 时，$V$ 中所有 $d$ 次有理曲线都光滑（$d=5$ 时有例外），
且法丛为 $\mathcal O_{\mathbb P^1}(-1)\oplus\mathcal O_{\mathbb P^1}(-1)$，见 \SourceCite{JK1}。

为证明式~\eqref{eq:quintic-instanton-small-degree}，先假设 $d\ne5$，
并取 $\overline M_{0,0}(V,d)$ 中的稳定映射 $f:C'\to V$。
由于次数不超过 $9$ 的所有曲线彼此不相交，$C=f(C')$ 必为 $d/k$ 次光滑曲线，
其中 $k$ 是 $f:C'\to C$ 的次数。
用定理~\ref{thm:multiple-cover} 的记号，$f$ 位于 $\overline M_{0,0}(V,d)$ 的分支 $M_{k,C}(V)$ 中。
这说明 $\overline M_{0,0}(V,d)$ 是这些分支的不交并，于是定理~\ref{thm:multiple-cover} 给出
\[
N_d=\sum_{k\mid d}n'_{d/k}k^{-3}.
\]
由 $n_d$ 的唯一性，$n'_d=n_d$，所以式~\eqref{eq:quintic-instanton-small-degree} 对 $d\ne5$ 得证。

再设 $d=5$。由 \SourceCite{Vainsencher}，$V$ 中有 $17\,601\,000$ 条具有六个结点的五次有理曲线。
现取 $\overline M_{0,0}(V,5)$ 中的稳定映射 $f:C'\to V$。
此时，$f(C')$ 或为直线，或为五次光滑曲线，或为具有六个结点的五次有理曲线。
因此，$\overline M_{0,0}(V,5)$ 分解为三种相应的分支；定理~\ref{thm:multiple-cover} 适用于前两种。
余下只需考虑 $f:C'\to C$，其中 $C$ 是 $V$ 中具有六个结点的有理曲线。
容易看出，$f$ 是正规化映射 $f:\mathbb P^1\to C$。
此外，\hyperref[subsec:nodal-quintic-double-covers]{第~\ref*{subsec:nodal-quintic-double-covers} 小节}将证明 $f$ 的法丛为
$N_f=\mathcal O_{\mathbb P^1}(-1)\oplus\mathcal O_{\mathbb P^1}(-1)$。
这表明，$\overline M_{0,0}(V,5)$ 中对应于 $C$ 的分支是一个光滑点，
且它对虚拟基本类的贡献为 $1$，与定理~\ref{thm:multiple-cover} 相符。
于是，上一段的论证给出 $n'_5=n_5$，完成式~\eqref{eq:quintic-instanton-small-degree} 的证明。

基于式~\eqref{eq:quintic-instanton-small-degree}，自然会希望：
只要证明更高次数 $d$ 的 Clemens 猜想，就能自动推出相应瞬子数是\IndexLocation{indexloc:1193:229}\IndexLocation{indexloc:1199:229}\IndexLocation{indexloc:1200:229}五次三维簇上 $d$ 次有理曲线的数目。
长期以来，许多人相信这一点；然而，R. Pandharipande 最近的一个观察表明，这种希望过于朴素。
事实上，对于 $d=10$——猜想~\ref{conj:clemens} 中第一个尚未解决的情形——有下面关于 $n_{10}$ 的猜想公式：
\begin{equation}\label{eq:quintic-degree-ten}
\begin{gathered}
\text{若 Clemens 猜想对所有 }d\leq10\text{ 成立，则}\\
n_{10}=6\times17\,601\,000+V\text{ 中 }10\text{ 次有理曲线的数目}.
\end{gathered}
\end{equation}
\hyperref[subsec:nodal-quintic-double-covers]{第~\ref*{subsec:nodal-quintic-double-covers} 小节}将仔细证明此式。
关键在于，$V$ 中那 $17\,601\,000$ 条具有六个结点的五次曲线的二重覆盖，造成了复杂情形。
更确切地说，该小节将证明，每条这样的曲线 $C$ 对 $N_{10}$ 的贡献是 $6\tfrac18$，而非 $\tfrac18$。
我们会看到，这是因为二重覆盖 $f:C'\to C$ 有两种：
一种经由正规化映射分解；另一种的 $C'$ 可约，且 $f$ 在六个结点中某一个的邻域内是局部同构。
\SourcePage{230}
假设猜想~\ref{conj:clemens} 对 $d\leq10$ 成立，且每条上述 $C$ 对 $N_{10}$ 都贡献 $6\tfrac18$，
则式~\eqref{eq:quintic-degree-ten} 可如下推导。首先，由定义，
\[
N_{10}=n_1 10^{-3}+n_2 5^{-3}+n_5 2^{-3}+n_{10}.
\]
再利用 $d\leq9$ 时的 $n_d=n'_d$，以及证明式~\eqref{eq:quintic-instanton-small-degree} 时的方法，得到
\[
N_{10}=n_1 10^{-3}+n_2 5^{-3}+(n_5-17\,601\,000)2^{-3}
+17\,601\,000\cdot6\tfrac18+n'_{10},
\]
其中 $n'_{10}$ 是 $10$ 次有理曲线的数目。
将这两个等式相减，容易得到式~\eqref{eq:quintic-degree-ten}。
\end{example}

\hyperref[chap:mirror-theorem]{第~\ref*{chap:mirror-theorem} 章}证明五次三维簇 $V$ 的镜像定理后，我们就能按需要计算任意多个瞬子数 $n_d$。
但从前一个例子可见，这离知道 $V$ 上 $d$ 次有理曲线的数目仍很远。
首先，对于 $d\geq10$，证明 Clemens 猜想是一个非平凡的问题
（关于 $10\leq d\leq24$ 的一些初步进展，见 \SourceCite{JK2}）；
然后还要面对这些曲线并非全都光滑的困难。
如上所见，这类曲线的多重覆盖，并非总按定理~\ref{thm:multiple-cover} 的方式作出贡献。
事实上，只要 $d\geq10$ 是 $5$ 的倍数，就可能出现这一问题。
所以，即使\hyperref[chap:quintic]{第~\ref*{chap:quintic} 章}的公式已经全部得到证明，五次三维簇仍提出了许多未解决的问题。

下面给出一个瞬子数更微妙的例子。这里将看到，除多重覆盖之外，虚拟基本类还会以其他方式发挥作用。
\SourceExampleNumber{7.4.4}{2}
\begin{example}\label{ex:local-plane-instantons}
设 $\widehat V\subset\mathbb P(1,1,1,6,9)$ 是由下式定义的 $18$ 次超曲面：
\[
x_0^{18}+x_1^{18}+x_2^{18}+x_3^3+x_4^2=0.
\]
可以证明，$\widehat V$ 在 $(0,0,0,-1,1)$ 有唯一奇点，解消此奇点得到 Calabi--Yau 三维簇 $V$。
例外集是 $V$ 中的一个除子 $E\simeq\mathbb P^2$。
一般情形也是如此；而 $\mathbb P^2\subset V$ 表明，$V$ 有无穷多条各个次数的有理曲线。

为计算 Gromov--Witten 不变量，需要描述 $V$ 上的一些类。
依照 \SourceCite{CFKM}，令 $L,H$ 分别为由 $1$ 次与 $3$ 次方程定义的 $V$ 上的除子类。
由于爆破，有 $3L+E\in|H|$；可以证明，$L,H$ 张成 $H^2(V,\mathbb Z)$ 并生成 Kähler 锥。
此外，$\ell=L\cap E$ 是 $E\simeq\mathbb P^2$ 中的直线，
$h=L_1\cap L_2$ 是椭圆曲线，其中 $L_i\in|L|$ 为一般成员。
于是 $\ell,h$ 生成 $H^4(V,\mathbb Z)\simeq H_2(V,\mathbb Z)$。

在这一情形中，用 $\langle I_{0,3,r,s}\rangle$ 表示 $\langle I_{0,3,[rh+s\ell]}\rangle$。
则式~\eqref{eq:instanton-three-divisors} 变成
\[
\langle I_{0,3,r,s}\rangle(\alpha_1,\alpha_2,\alpha_3)
=\sum_{k\mid\gcd(r,s)}n\!\left(\frac r k,\frac s k\right)
\int_{\frac r kh+\frac s k\ell}\alpha_1
\int_{\frac r kh+\frac s k\ell}\alpha_2
\int_{\frac r kh+\frac s k\ell}\alpha_3,
\]
其中 $\alpha_1,\alpha_2,\alpha_3\in H^2(V,\mathbb C)$，
$n(r,s)$ 是类为 $C\sim rh+s\ell$ 的有理曲线 $C\subset V$ 的瞬子数。
为看看会得到怎样的数，考虑
\[
\langle I_{0,3,0,2}\rangle(L,L,L)=n(0,1)+n(0,2)2^3.
\]
由 \SourceCite{CFKM}，镜像对称预测 $n(0,1)=3$、$n(0,2)=-6$，
从而 $\langle I_{0,3,0,2}\rangle(L,L,L)=-45$。
这些数也可由定理~\ref{thm:givental-toric-mirror} 得到。因此，我们得到一个负的 Gromov--Witten 不变量。
\SourcePage{231}
这些数可作如下解释。首先，$C\sim\ell$ 表明 $C$ 是 $\mathbb P^2$ 中的直线。
这类 $C$ 的模空间是对偶的 $\mathbb P^2$。
为计算它对 $n(0,1)$ 的贡献，回忆\hyperref[subsec:normal-cone]{第~\ref*{subsec:normal-cone} 小节}
关于虚拟基本类的讨论。在命题~\ref{prop:excess-obstruction-euler-class} 的情形下，
构造归结为求障碍丛 Euler 类（即最高 Chern 类）的次数。
又因为这里处理的是嵌入曲线，直线 $C\subset\mathbb P^2$ 的障碍空间为 $H^1(N)$，
其中 $N$ 是 $C$ 在 $V$ 中的法丛。但 $H^1(N)$ 通过配对
\[
H^1(N)\otimes H^0(N)\longrightarrow H^1(\wedge^2N)
\simeq H^1(T_C^*)=\mathbb C
\]
与 $H^0(N)$ 对偶，所以全局的障碍丛是 $T^*_{\mathrm{moduli}}=T^*_{\mathbb P^2}$。
因此瞬子数为
\[
n(0,1)=\int_{\mathbb P^2}c_2(T^*_{\mathbb P^2})
=(-1)^{\dim\mathbb P^2}\chi(\mathbb P^2)=3.
\]
类似地，$C\sim2\ell$ 表明 $C$ 是 $\mathbb P^2$ 中的二次曲线。
二次曲线的模空间是 $\mathbb P^5$。由于它光滑，与上面相同的计算给出
\[
n(0,2)=\int_{\mathbb P^5}c_5(T^*_{\mathbb P^5})
=(-1)^{\dim\mathbb P^5}\chi(\mathbb P^5)=-6.
\]
这就是 \SourceCite{CFKM} 中的论证，但它并不完整：满足 $f_*[C]=2\ell$ 的稳定映射
$f:C'\to V$ 的模空间，与这里所用的二次曲线 Hilbert 概形不同。

$n(0,1)$ 也可在辛几何框架中计算，见 \SourceCite{Ruan1}。
对 $n(0,2)$ 的朴素计算得到正确数值，提示非孤立曲线可能存在一种多重覆盖公式的推广，
使我们能用浸入曲线，而不是稳定映射，更广泛地从几何上描述 $n_\beta$。
下面的猜想~\ref{conj:symplectic-instantons} 涉及其中一些微妙之处。
关于这里研究的 Calabi--Yau 三维簇 $V$ 的进一步细节与计算，见 \SourceCite{CFKM}。
\end{example}

我们将以一种定义 $n_\beta$ 的猜想性辛几何方法结束瞬子数的讨论。
在 $V$ 上固定一个辛结构（例如 Kähler 形式定义的结构）。
回忆\hyperref[subsec:symplectic-map-moduli]{第~\ref*{subsec:symplectic-map-moduli} 小节}，
$\mathcal M(\mathbb P^1,J,\beta)$ 是简单的 $J$ 全纯映射 $f:\mathbb P^1\to V$ 的空间，
其中 $J$ 是一般的几乎复结构，受 $V$ 的辛形式驯服。
由 \SourceCite[第 2 章]{MS}，每个 $f\in\mathcal M(\mathbb P^1,J,\beta)$
在一般点为浸入，意思是 $\mathbb P^1$ 有稠密开集 $U$，在其上 $df$ 单射，且对所有 $u\in U$，
$f^{-1}(f(u))=\{u\}$。
$\operatorname{PGL}(2,\mathbb C)$ 在 $\mathcal M(\mathbb P^1,J,\beta)$ 上自然作用，
我们将此作用的一个轨道定义为 $V$ 上类为 $\beta$ 的一条 \emph{$J$ 全纯曲线}。

由式~\eqref{eq:fixed-curve-map-index}，$\mathcal M(\mathbb P^1,J,\beta)$ 的实维数为 $6$。
这正是 $\operatorname{PGL}(2,\mathbb C)$ 的维数；由于作用自由，各轨道就是
$\mathcal M(\mathbb P^1,J,\beta)$ 的连通分支。
此外，每个分支 $\mathcal C$ 有自然定向，因此可定义符号 $\epsilon(\mathcal C)=+1$ 或 $-1$，
分别对应于映射 $\phi:\operatorname{PGL}(2,\mathbb C)\simeq\mathcal C$ 保持或不保持定向；
这里 $\phi$ 由群作用于固定的 $f_0\in\mathcal C$ 给出。于是提出如下猜想。
\begin{conjecture}\label{conj:symplectic-instantons}
设 $(V,\omega)$ 是带 Kähler 形式的 Calabi--Yau 三维簇。
固定一般的、受 $\omega$ 驯服的几乎复结构 $J$。则：
\SourcePage{232}
\begin{enumerate}[label=(\roman*)]
\item 对任意同调类 $\beta\in H_2(V,\mathbb Z)$，类为 $\beta$、亏格为零的浸入 $J$ 全纯曲线只有有限多条。
它们都是嵌入的，且任意两条都不相交（即使 $\beta$ 不同）。定义整数
\[
n'_\beta=\sum_{\mathcal C\in\pi_0\mathcal M(\mathbb P^1,J,\beta)}\epsilon(\mathcal C),
\]
其中求和遍历 $\mathcal M(\mathbb P^1,J,\beta)$ 的所有连通分支 $\mathcal C$。
\item 亏格零的所有 $J$ 全纯映射 $C\to V$，都是嵌入的多重覆盖。
\item 对每个 $\beta$，$n_\beta$ 可表示为
\[
n_\beta=n'_\beta+\sum_{\substack{\beta=k\gamma\\k>1}}
\sum_{\mathcal C\in\pi_0\mathcal M(\mathbb P^1,J,\gamma)}c(k,\mathcal C),
\]
其中 $c(k,\mathcal C)$ 是由 $\mathcal C$ 的 $k$ 次多重覆盖确定的整数。
\end{enumerate}
\end{conjecture}
此猜想说，$n_\beta$ 应为类 $\beta$ 中 $J$ 全纯曲线的带符号条数，
再加上当 $\beta$ 可被 $k$ 整除时，$k$ 次多重覆盖的某些贡献。
特别地，猜想~\ref{conj:symplectic-instantons} 推出每个 $n_\beta$ 都是整数。
$\langle I_{0,0,\beta}\rangle$ 的辛定义表明 $n_\beta$ 是辛不变量；
但如下面所解释，$n'_\beta$ 可能不是。

复杂性主要仍来自多重覆盖。
如\hyperref[sec:gw-symplectic-definition]{第~\ref*{sec:gw-symplectic-definition} 节}所述，
空间 $\mathcal M(\mathbb P^1,J,\beta)$ 不紧，而 Gromov--Uhlenbeck 紧化需要加入多重覆盖。
因此，甚至不能预先断定猜想~\ref{conj:symplectic-instantons} 中等式的右边是有限和。

读者也许会问，为何不更简单地猜想总有 $n_\beta=n'_\beta$？
事实上，我们不知道任何违反这个等式的例子。另一方面，\SourceCite{Taubes} 中有类似情形：
它分析了辛四维流形中法丛平凡的椭圆曲线的多重覆盖。
在那里，二重覆盖有一个相应的离散不变量，其值会改变该覆盖对 Gromov--Witten 不变量的贡献。
此外，所有 $J$ 构成的空间有一些称为\emph{墙}的实余维一轨迹；
当 $J$ 穿过一堵墙时，这种贡献的变化为整数。
穿墙时出现额外的类为 $2\beta$ 的 $J$ 全纯曲线，使 Gromov--Witten 不变量保持不变。
而二重覆盖的多重覆盖所产生的贡献变化，又由类 $2\beta$ 中这些新曲线的多重覆盖的贡献抵消。
我们感谢 \SourceCite{IP} 的作者解释 Taubes 的工作及他们自己的论文；
该论文研究 Taubes 定义的不变量与 Gromov--Witten 不变量的联系。
这番讨论促使我们提出猜想~\ref{conj:symplectic-instantons}。

可以设想，Calabi--Yau 三维簇上的有理曲线也出现类似情形：
所有 $J$ 的空间中可能存在墙；若 $V$ 的给定复结构位于墙上，
那么“一般”可能有两种含义，取决于“一般”的 $J$ 位于墙的哪一侧。
正如上一段所说的 $g=1$ 情形，穿墙可能使某些嵌入的 $J$ 全纯曲线变成多重覆盖。
我们希望为每个 $k>1$ 和嵌入的 $J$ 全纯曲线 $\mathcal C$ 定义离散不变量 $c(k,\mathcal C)$，
使穿墙时 $n'_\beta$ 的损失由 $c(k,\mathcal C)$ 的增加抵消，
从而猜想~\ref{conj:symplectic-instantons} 中的总和保持不变。
\SourcePage{233}
在一种情形下，猜想~\ref{conj:symplectic-instantons} 的一部分可以证明。
假设 $\beta$ 是\emph{本原}的，即不是 $H_2(V,\mathbb Z)$ 中另一个类的某个大于 $1$ 的整数倍。
\begin{proposition}\label{prop:primitive-instantons}
若 $\beta\in H_2(V,\mathbb Z)$ 非零且本原，则 $n_\beta=n'_\beta$，
其中 $n'_\beta$ 如猜想~\ref{conj:symplectic-instantons} 所定义。
\end{proposition}
\begin{proof}
这由 \SourceCite{Ruan2} 得到。采用 Li--Ruan--Tian 的 Gromov--Witten 不变量定义，
则用\hyperref[subsec:gw-equivalence]{第~\ref*{subsec:gw-equivalence} 小节}的记号，
$n_\beta=\langle I_{0,0,\beta}\rangle^{\mathrm{LRT}}$。
由式~\eqref{eq:instanton-three-divisors}，可以从 $\langle I_{0,3,\beta}\rangle^{\mathrm{LRT}}$ 恢复 $n_\beta$。
由于 $V$ 半正，可用不变量 $\Psi_{0,3,\beta}$ 代替 $\langle I_{0,3,\beta}\rangle^{\mathrm{LRT}}$。
取 $\alpha_1,\alpha_2,\alpha_3\in H^2(V,\mathbb Z)$。
用\hyperref[subsec:symplectic-gw-invariants]{第~\ref*{subsec:symplectic-gw-invariants} 小节}
定义的辛 Gromov--Witten 不变量 $\Psi,\Phi$，由式~\eqref{eq:genus-zero-mixed-symplectic}，
\[
\Psi_{0,3,\beta}([pt];\alpha_1,\alpha_2,\alpha_3)
=\Phi_{0,\beta}(\alpha_1,\alpha_2,\alpha_3\mid\,).
\]
可以用 \SourceCite{Ruan2} 的方法计算 $\Phi_{0,\beta}(\alpha_1,\alpha_2,\alpha_3\mid\,)$，
因为 $\beta$ 本原意味着所有满足 $f_*[\mathbb P^1]=\beta$ 的 $J$ \IndexLocation{indexloc:0655:233}全纯映射 $f:\mathbb P^1\to V$ 都是简单的。
所以 $\mathcal M(\mathbb P^1,J,\beta)$ 给出了类 $\beta$ 中的所有 $J$ 全纯 $\mathbb P^1$。
考虑\IndexLocation{indexloc:0276:233}赋值映射
\[
\mathcal M(\mathbb P^1,J,\beta)\times_{\operatorname{PGL}(2,\mathbb C)}(\mathbb P^1)^3
\longrightarrow V^3,\quad(f,p_1,p_2,p_3)\longmapsto(f(p_1),f(p_2),f(p_3)).
\]
此映射的像的闭包为六维圈 $\Gamma$。
由 \SourceCite[第 7.2 节]{MS}，$\Phi_{0,\beta}(\alpha_1,\alpha_2,\alpha_3\mid\,)$
为交数 $\Gamma\cdot(\Gamma_1\times\Gamma_2\times\Gamma_3)$，
其中与定理~\ref{thm:multiple-cover} 的证明一样，$\Gamma_i$ 是与 $\alpha_i$ 对偶的四维圈。

已知 $\mathcal M(\mathbb P^1,J,\beta)$ 为各连通分支 $\mathcal C$ 的并。
固定 $f_0\in\mathcal C$，让 $\operatorname{PGL}(2,\mathbb C)$ 作用于 $f_0$，得到保持定向的同构
$\mathcal C\simeq\epsilon(\mathcal C)\operatorname{PGL}(2,\mathbb C)$，
其中 $\epsilon(\mathcal C)$ 决定如何给群定向。
据此，上述映射变成
\[
f_0\times f_0\times f_0:\epsilon(\mathcal C)(\mathbb P^1)^3
\simeq\mathcal C\times_{\operatorname{PGL}(2,\mathbb C)}(\mathbb P^1)^3\longrightarrow V^3.
\]
由于 $f_{0*}[\mathbb P^1]=\beta$，其像与 $\Gamma_1\times\Gamma_2\times\Gamma_3$ 的交数为
$\epsilon(\mathcal C)\int_\beta\alpha_1\int_\beta\alpha_2\int_\beta\alpha_3$。
对所有分支求和，就得到所需结果。
\end{proof}
类似地，对于五次三维簇，猜想~\ref{conj:symplectic-instantons} 对所有满足 $d\leq9$、$d\ne5$ 的次数 $d$ 成立。
原因是：取可积的 $J$，就已能达到足够一般的情形，
然后用例题~\ref{ex:quintic-instanton-enumerative} 的论证得到 $n_d=n'_d$，
这里的 $n'_d$ 如猜想~\ref{conj:symplectic-instantons} 所定义。

\subsection{亏格一不变量}\label{subsec:genus-one-invariants}
\index{Gromov--Witten不变量@Gromov--Witten 不变量!亏格一|hyperpage}\index{椭圆曲线!计数|hyperpage}\index{固定j不变量@固定 $j$ 不变量|hyperpage}\index{退化瞬子|hyperpage}
\index{Gromov--Witten不变量@Gromov--Witten 不变量!亏格一|hyperpage}
当 $g=1$ 时，得到\emph{椭圆 Gromov--Witten 不变量}
$\langle I_{1,n,\beta}\rangle(\alpha_1,\ldots,\alpha_n)$。
这些不变量最初在 \SourceCite{BCOV1} 中讨论，此后已在各种背景下得到研究。
下面的例子说明，确定 Gromov--Witten 不变量的枚举意义，往往是一项微妙的工作。
\SourceExampleNumber{7.4.5}{1}
\begin{example}\label{ex:plane-genus-one}
先研究射影平面 $\mathbb P^2$ 的椭圆 Gromov--Witten 不变量。
若 $\ell\subset\mathbb P^2$ 为直线，与往常一样，用 $\langle I_{1,n,d}\rangle$ 代替 $\langle I_{1,n,d[\ell]}\rangle$。
\SourcePage{234}
首先考虑 $\langle I_{1,1,0}\rangle$，它出人意料地并非平凡。
这里处理恒定映射（次数为零），但\hyperref[sec:gw-properties]{第~\ref*{sec:gw-properties} 节}的点映射公理不适用。
不过，如 \SourceCite{KoM1} 所解释，对每个亏格都有相应版本。
对 $g=1$ 和任意 $X$，当 $\deg\alpha=2$ 时，得到
\begin{equation}\label{eq:genus-one-constant}
\langle I_{1,1,0}\rangle(\alpha)=-\frac1{24}\int_Xc_{\dim X-1}(T_X)\cup\alpha.
\end{equation}
问题在于，$n$ 点、亏格一、次数零的映射 $f:C\to X$ 的空间是
$\overline M_{1,n}\times X$，它并不具有预期维数。
所以必须像\hyperref[subsec:computing-virtual-class]{第~\ref*{subsec:computing-virtual-class} 小节}一样计算虚拟基本类。
可用命题~\ref{prop:excess-obstruction-euler-class} 完成这一计算，例题~\ref{ex:genus-one-gravitational-euler} 将详细说明，
式~\eqref{eq:genus-one-constant} 随即得到。
将它用于 $X=\mathbb P^2$，得到
\[
\langle I_{1,1,0}\rangle([\ell])=-\frac1{24}\int_{\mathbb P^2}c_1(T_{\mathbb P^2})\cup[\ell]=-\frac18.
\]
所以，Gromov--Witten 不变量既可能为分数，也可能为负！
例题~\ref{ex:genus-one-gravitational-euler} 在将 Gromov--Witten 不变量推广为\emph{引力关联数}之后，将讨论一个有趣的相关计算。

幸运的是，其他椭圆不变量具有很好的枚举意义。对 $d\geq1$，定义
\[
N_d^1=\langle I_{1,3d,d}\rangle(\underbrace{[pt],\ldots,[pt]}_{3d\text{ 次}}).
\]
可以证明，$N_d^1$ 是 $\mathbb P^2$ 中经过 $3d$ 个一般点的 $d$ 次椭圆曲线的数目。
于是 $N_1^1=N_2^1=0$、$N_3^1=1$，因为平面中经过九个一般点的光滑椭圆曲线唯一。
当 $d>3$ 时，$N_d^1$ 所计数的曲线都奇异。
容易证明：一旦知道 $\langle I_{1,1,0}\rangle([\ell])$ 与所有 $d\geq1$ 的 $N_d^1$，
就确定了全部椭圆 Gromov--Witten 不变量。

关于 $N_d^1$ 的主要结果是各种递推关系；一个例子见 \SourceCite{Getzler1}。
\SourceCite{EHX} 用 Virasoro 约束的方法预测了一个漂亮的递推式，\hyperref[subsec:virasoro-conjecture]{第~\ref*{subsec:virasoro-conjecture} 小节}将简要讨论这一方法。
此公式由 \SourceCite{Pandharipande2} 证明：若 $N_d$ 是式~\eqref{eq:plane-nd} 给出的有理曲线数目，则
\begin{equation}\label{eq:plane-genus-one-recursion}
N_d^1=\frac1{12}\binom d3N_d+
\sum_{\substack{d=d_1+d_2\\d_1,d_2>0}}
\frac{3d_1^2d_2-2d_1d_2}{9}N_{d_1}N_{d_2}^1\binom{3d-1}{3d_1-1}.
\end{equation}
\SourceCite{Pandharipande2} 的证明，使用 Getzler 在
$H^4(\overline M_{1,4},\mathbb Q)$ 中发现的余维二分层之间一个关系的几何构造，
然后应用\hyperref[chap:quantum-cohomology]{第~\ref*{chap:quantum-cohomology} 章}将讨论的 WDVV 方程。
由此可按需要计算任意多个 $N_d^1$。一个初等证明见 \SourceCite{Ran3}。

在迄今讨论的椭圆不变量中，我们允许 $j$ 不变量变化。
但也可研究 $\mathbb P^2$ 中 $j$ 不变量\emph{固定}的椭圆曲线。
更准确地，令 $E_{d,j}$ 为 $\mathbb P^2$ 中经过 $3d-1$ 个一般点、具有给定 $j$ 不变量的
$d$ 次不可约椭圆曲线数目。
\SourcePage{235}
\SourceCite{Ionel,Pandharipande1} 证明
\begin{equation}\label{eq:plane-fixed-j}
E_{d,j}=\begin{cases}
\binom{d-1}{2}N_d,&j\ne0,1728,\\
\frac13\binom{d-1}{2}N_d,&j=0,\\
\frac12\binom{d-1}{2}N_d,&j=1728,
\end{cases}
\end{equation}
其中 $N_d$ 如上。

试用 Gromov--Witten 类进行这一计算。为简便，假设 $j\ne0,1728$。
由于希望经过 $\mathbb P^2$ 中的 $3d-1$ 个点，从 Gromov--Witten 类
\[
\xi=I_{1,3d-1,d}(\underbrace{[pt],\ldots,[pt]}_{3d-1\text{ 次}})
\]
出发。由式~\eqref{eq:gw-top-degree-condition}，这个类的次数为 $6d-4$；
因 $\overline M_{1,3d-1}$ 维数为 $3d-1$，它比最高次数低二。
固定 $j$，得到由所有具有此 $j$ 不变量的带标记曲线构成的除子
$Z_j\subset\overline M_{1,3d-1}$。
朴素地，我们可能预期 $E_{d,j}$ 是 $\xi$ 与 $Z_j$ 的交，即
\begin{equation}\label{eq:plane-fixed-j-naive}
E_{d,j}=\xi\cdot Z_j.
\end{equation}
如下计算 $\xi\cdot Z_j$。
自然映射 $\overline M_{1,3d-1}\to\overline M_{1,1}\simeq\mathbb P^1$
遗忘最后的 $3d-2$ 个点，并收缩不稳定分支。由于 $j\ne0,1728$，$Z_j$ 就是 $j$ 上的纤维。
而 $\infty$ 上的纤维由所有映到结点曲线的 $\overline M_{1,3d-1}$ 中的元素组成。
容易看出，这个纤维是映射
$\psi:\overline M_{0,3d+1}\to\overline M_{1,3d-1}$ 的像，$\psi$ 将最后两个点粘合。
所以 $Z_j\sim\operatorname{Im}(\psi)$。
这意味着
\[
\begin{aligned}
\xi\cdot Z_j=\xi\cdot\operatorname{Im}(\psi)=\frac12\Big(&
\langle I_{0,3d+1,d}\rangle([pt],\ldots,[pt],[pt],[\mathbb P^2])\\
&+\langle I_{0,3d+1,d}\rangle([pt],\ldots,[pt],[\ell],[\ell])\\
&+\langle I_{0,3d+1,d}\rangle([pt],\ldots,[pt],[\mathbb P^2],[pt])\Big).
\end{aligned}
\]
这是第一次使用约化公理。由此，基本类公理与除子公理容易给出
\[
\xi\cdot Z_j=\xi\cdot\operatorname{Im}(\psi)=\frac12d^2N_d.
\]
与式~\eqref{eq:plane-fixed-j} 比较，可见式~\eqref{eq:plane-fixed-j-naive} 是错误的。
原因在于，$E_{d,j}$ 只计数 $\mathbb P^2$ 中的不可约曲线，
而 $\xi\cdot Z_j$ 还包括来自 $\overline M_{1,3d-1}$ 边界的可约曲线。
事实上，\SourceCite{Ionel} 的论证先进行上述计算，再在辛几何背景下仔细计入边界，来求 $E_{d,j}$。
还应提到，用\hyperref[sec:gw-symplectic-definition]{第~\ref*{sec:gw-symplectic-definition} 节}
定义的混合不变量与辛不变量，有
\[
\xi\cdot Z_j=\frac12\Phi_{1,d}([pt]\mid\underbrace{[pt],\ldots,[pt]}_{3d-2\text{ 次}})
=\Psi_{1,3d-1,d}(Z_j;\underbrace{[pt],\ldots,[pt]}_{3d-1\text{ 次}}),
\]
尤其参见\hyperref[subsec:symplectic-gw-invariants]{第~\ref*{subsec:symplectic-gw-invariants} 小节}中混合不变量与辛不变量的关系。
这些不变量及相关椭圆不变量的\IndexLocation{indexloc:0534:235}枚举意义，见 \SourceCite[注记 2.19]{RT2}。
\end{example}
\SourcePage{236}
\SourceExampleNumber{7.4.5}{2}
\begin{example}\label{ex:space-genus-one}
在 $\mathbb P^3$ 中取直线 $\ell$，并定义
\[
\begin{aligned}
N_{a,b}&=\langle I_{0,a+b,(a+2b)/4}\rangle
(\underbrace{[\ell],\ldots,[\ell]}_{a\text{ 次}},\underbrace{[pt],\ldots,[pt]}_{b\text{ 次}}),\\
N_{a,b}^1&=\langle I_{1,a+b,(a+2b)/4}\rangle
(\underbrace{[\ell],\ldots,[\ell]}_{a\text{ 次}},\underbrace{[pt],\ldots,[pt]}_{b\text{ 次}}),
\end{aligned}
\]
其中假设 $(a+2b)/4$ 为整数。
由定理~\ref{thm:homogeneous-gw-equivalence}，$N_{a,b}$ 是 $\mathbb P^3$ 中
经过 $a$ 条一般直线及 $b$ 个一般点的 $d=(a+2b)/4$ 次有理曲线数目。
$N_{a,b}$ 的递推关系见 \SourceCite{FP}。

椭圆不变量 $N_{a,b}^1$ 也可用 \SourceCite{Getzler1} 附录中的递推关系计算。
但与亏格零情形不同，它们不是枚举不变量。
例如，\SourceCite{Getzler1} 附录算出的\emph{所有} $N_{a,b}^1$ 都是负数，许多还是分数。
不过，若取线性组合
\[
E_{a,b}=N_{a,b}^1+\frac{2d-1}{12}N_{a,b},
\]
其中仍有 $d=(a+2b)/4$，则 \SourceCite{GeP} 表明，$E_{a,b}$ 是
$\mathbb P^3$ 中经过 $a$ 条一般直线及 $b$ 个一般点的 $d$ 次椭圆曲线数目。
计算 $N_{a,b}^1$ 的其他方法，以及对 $\mathbb P^n$ 的推广，见 \SourceCite{Vakil3,GPa}。

还应注意，也能研究 $\mathbb P^3$ 中 $j$ 不变量固定的椭圆曲线。
与 $\mathbb P^2$ 一样，可得到这类曲线的数目与 $N_{a,b}$ 之间的关系式。
精确公式见 \SourceCite{Ionel}。
\end{example}
\SourceExampleNumber{7.4.5}{3}
\begin{example}\label{ex:quintic-genus-one}
最后，回到一般的五次三维簇 $V\subset\mathbb P^4$。
从式~\eqref{eq:cy-a-correlation-instanton} 的椭圆版本开始。
令 $n_d$ 为 $V$ 上 $d$ 次有理曲线的瞬子数，令 $e_d$ 为 $V$ 上 $d$ 次“椭圆曲线”的类似瞬子数
（下面解释引号的含义）。假设这些数有限，并将 $n_d,e_d$ 理解为计数适当的刚性嵌入曲线。
那么有公式
\begin{equation}\label{eq:quintic-genus-one-instanton}
\langle I_{1,1,d}\rangle(H)=\sum_{k\mid d}\sigma\!\left(\frac d k\right)ke_k
+\frac1{12}\sum_{k\mid d}kn_k,
\end{equation}
其中 $\sigma(d/k)=\sum_{l\mid(d/k)}l$ 是 $d/k$ 的所有正因子之和。

此公式通过物理推理得到，见 \SourceCite{BCOV1}，但迄今尚无数学证明。
不过，可作如下直观解释。先看式~\eqref{eq:quintic-genus-one-instanton} 中第一个和。
设 $f:(C,p)\to V$ 为稳定映射，其像 $E=f(C)$ 为椭圆曲线，
$C\to E$ 的次数为 $d/k$，$E\subset V$ 的次数为 $k$。
那么 $E$ 是这 $e_k$ 条椭圆曲线之一；由于 $f(p)\in E\cap H$，$f(p)$ 有 $k$ 种选择。
给定 $f(p)$，$p\in C$ 有 $d/k$ 种选择，但因 $C\to E$ 是覆盖，它们彼此同构。
因此，只需计数 $E$ 的 $d/k$ 次覆盖空间。这相当于计数 $\mathbb Z^2$ 中指数为 $d/k$ 的子格，
众所周知，其数目为 $\sigma(d/k)$。

式~\eqref{eq:quintic-genus-one-instanton} 中的第二个和更微妙。例如，考虑 $k=d$。
那 $n_d$ 条 $d$ 次“有理曲线”对 $\langle I_{1,1,d}\rangle(H)$ 如何作出贡献？
取 $d$ 次有理曲线 $C_0\subset V$。如往常一样，$f(p)\in C_0\cap H$ 有 $d$ 种可能。
给定 $C_0$ 与 $f(p)$，我们想构造次数为 $1$ 的亏格一稳定映射 $f:(C,p)\to C_0$。
这似乎不可能，
\SourcePage{237}
但稳定映射的定义允许 $C$ 可约。
物理文献把这种情形称为\emph{退化瞬子}。
这里 $C=E\cup C_0$，$E$ 是在一点与 $C_0$ 相接的椭圆曲线，
$f:C\to C_0$ 在 $C_0$ 上为恒等映射，在 $E$ 上为恒定映射。
一旦选取映到 $f(p)$ 的点 $p\in C$，就得到我们所关心的这类稳定映射的一个族。

若第二个和中的 $k\mid d$ 是 $d$ 的真因子，情形更复杂，主要原因是多重覆盖。
与亏格零情形一样，需要处理有理曲线的多重覆盖；
此外，还有由椭圆曲线到 $\mathbb P^1$ 的熟知二次映射产生的多重覆盖。
物理计算 \SourceCite{BCOV1} 提示，答案仍为 $1/12$；
当底层有理曲线的法丛为 $\mathcal O_{\mathbb P^1}(-1)\oplus\mathcal O_{\mathbb P^1}(-1)$ 时，
这一点已得到数学证明，见 \SourceCite{GPa}。
所以，若 Clemens 猜想~\ref{conj:clemens} 的假设对所有整除 $d$ 的次数成立，
并且椭圆曲线也满足类似假设，彼此不相交且与有理曲线不相交，
则式~\eqref{eq:quintic-genus-one-instanton} 对该次数 $d$ 成立。

还有一个复杂之处：即使 $e_d$ 很可能有限，它也可能计数 $V$ 上光滑椭圆曲线之外的对象。
这就是我们为“椭圆曲线”加引号的原因。
最好的例子是五次三维簇上四次“椭圆曲线”的数目 $e_4$。
\SourceCite{BCOV1} 将此数算为
\[
e_4=3721431625.
\]
然而，由 \SourceCite{ES2}，一般五次三维簇上的光滑四次椭圆曲线数目是
\[
3718024750.
\]
它与 $e_4$ 不同。如 \SourceCite{ES2} 所解释，差别来自 $e_4$ 允许\emph{奇异}曲线。
具体地，$V$ 含有若干双结点的四次平面曲线，每条都产生一个对 $e_4$ 有贡献的稳定映射。
这些奇异曲线如下计数。若 $\ell$ 是 $V$ 上的 $2875$ 条直线之一，
则包含 $\ell$ 的任意平面 $P\subset\mathbb P^4$，与 $V$ 的交为 $\ell\cup C$，其中 $C$ 是平面四次曲线。
一般情况下，$C$ 光滑；但在有限多个特殊平面 $P$ 上，它出现两个结点。
\SourceCite{Vainsencher} 证明，这样的特殊平面 $P$ 共有 $1185$ 个。
由于
\[
3721431625=3718024750+1185\times2875,
\]
这些数是相符的。这又说明，在理解 Gromov--Witten 不变量的\IndexLocation{indexloc:0534:237}枚举意义时，必须十分仔细。
\end{example}

\begingroup\linespread{1.22}\selectfont
\subsection{其他例题}\label{subsec:other-gw-examples}
除已经列出的例子外，关于 Gromov--Witten 不变量与枚举几何，还有\emph{大量}其他工作。
例子的清单已经很长，下面列出一些供进一步阅读的文献。

\looseness=-1 镜像对称的枚举方面可追溯到 1991 年：\SourceCite{CdGP} 对五次三维簇上的有理曲线作出了惊人的预测。
随后，大批论文对 Calabi--Yau 流形作出枚举预测，
包括物理论文 \SourceCite{BCOV1,BCOV2,CdFKM,CFKM,HKTY1,HKTY2,BKK,AM2,Font,KT1,KT2}，
以及数学论文 \SourceCite{Katz2,Morrison1,ES2,ES1,LTe,BvS,Voisin1,Meurer,HSS,Givental2,LLY,Givental4,BCFKvS2}。
关于镜像对称枚举方面的一篇很好的综述，见 \SourceCite{Morrison6}。
\SourcePage{238}
虽然 Calabi--Yau 流形的枚举几何仍有许多未解决问题，但不少较新的工作已经转向其他类型的代数簇。
这部分源于\hyperref[sec:gw-algebraic-definition]{第~\ref*{sec:gw-algebraic-definition} 节}
所述任意光滑射影簇的 Gromov--Witten 不变量定义，它促成了上述许多例子与论文。
此外，还有对有理曲面 \SourceCite{CM,GoP,Vakil1}、Del Pezzo 曲面 \SourceCite{Vakil2}、
$\mathbb P^n$ 的爆破 \SourceCite{Gathmann1}、凸簇的爆破 \SourceCite{Gathmann2}、
完全交 \SourceCite{Beauville1}，以及 $\mathbb P^2$ 中带结点的亏格二曲线 \SourceCite{KQR} 的应用。
与此同时，人们也受到启发，探索不用 Gromov--Witten 理论能做些什么。
由此产生了 Crauder 与 Miranda 的论文 \SourceCite{CM}，Ran 的论文 \SourceCite{Ran1,Ran3}，
以及 Caporaso 与 Harris 的论文 \SourceCite{CH1,CH2,CH3}（综述见 \SourceCite{Caporaso}）。
五次三维簇及其他相关 Calabi--Yau 三维簇上的 Clemens 猜想也得到研究，见 \SourceCite{JK1,JK2,EJS}。
综观这些论文，近年来对枚举几何的兴趣迅速增长。
这种兴趣可能起源于镜像对称，但如今在代数几何中已有自身的发展动力。

\enlargethispage{\baselineskip}
\looseness=-1 这一长章至此结束。我们已经看到，Gromov--Witten 不变量与辛几何建立了非平凡的联系，
并在当代代数几何中发挥重要作用。本章的一个启示是，曲线计数比预想的更加深刻、丰富。
然而，故事尚未讲完：我们还需要了解量子上同调。
\hyperref[chap:quantum-cohomology]{第~\ref*{chap:quantum-cohomology} 章}将看到，它会进一步解释这里的例子，
并补全理解镜像对称的数学表述所需的一整套思想。
\par\endgroup

\SourcePage{240}
\chapter{量子上同调}\label{chap:quantum-cohomology}
\index{量子上同调|hyperpage}
本章将看到，量子上同调实际上有两种，即所谓\emph{大量子上同调环}与\emph{小量子上同调环}。
我们还将研究 \emph{Gromov--Witten 势}，它可视为树级 Gromov--Witten 不变量的指数型生成函数。
Gromov--Witten 势满足的 WDVV 方程，自然导出大量子上同调环的结合律，
并对\hyperref[chap:gromov-witten]{第~\ref*{chap:gromov-witten} 章}研究的例子产生一些有趣的推论。
另一个重要主题是 Dubrovin 形式体系，它说明量子上同调如何在上同调环上给出 Hodge 结构变分。

量子上同调是理解镜像对称的数学版本所需的最后一项要素。
如本章将看到，它给出三点函数的清晰表述。
此外，它导出上一段提到的 Hodge 结构变分，这是 Kähler 模上的变分。
镜像对称的一个关键思想是：经由镜像映射，这个 Hodge 结构变分对应于
镜像族的复模空间上的几何 Hodge 结构变分。
本章末尾定义数学上的镜像对并讨论镜像猜想时，将详细解释这一点。

从历史上看，量子上同调的思想最初在物理文献中以 Calabi--Yau 三维簇的\emph{手征环}的形式出现。
手征环被理解为通常上同调环的形变，见 \SourceCite{LVW}，
在大半径极限点与通常上同调环一致。
量子乘积也可在\IndexLocation{indexloc:1336:240}拓扑量子场论中作更一般的理解，见 \SourceCite{Witten1,Witten2}；\hyperref[sec:topological-qft]{附录 B.6} 将讨论这一点。
量子上同调最初的数学构造适用于半正辛流形，见 \SourceCite{RT1,MS}。

本章采用的量子上同调方法几乎完全是代数几何的。
介绍性的论述见 \SourceCite{Aluffi,FP,MS,Morrison7,Voisin3}。

\section{小量子上同调}\label{sec:small-quantum-cohomology}
\index{量子上同调!小量子上同调|hyperpage}
光滑射影簇 $X$ 的小量子上同调环，用\hyperref[chap:gromov-witten]{第~\ref*{chap:gromov-witten} 章}的 Gromov--Witten 不变量 $\langle I_{0,3,\beta}\rangle$ 定义。

\subsection{定义与性质}\label{subsec:small-quantum-definition}
取 $H^*(X,\mathbb Q)$ 的齐次基 $T_0=1,T_1,\ldots,T_m$，令
$g_{ij}=\int_XT_i\cup T_j$，则 $(g^{ij})=(g_{ij})^{-1}$ 为其逆矩阵。
回忆\hyperref[chap:moduli]{第~\ref*{chap:moduli} 章}，若 $\omega\in H^2(X,\mathbb C)$ 的虚部为 Kähler 类，则 $\omega$ 是一个复化 Kähler 类。
\SourcePage{241}
\begin{definition}\label{def:small-quantum-product}
设 $\omega$ 是光滑射影簇 $X$ 上的复化 Kähler 类。
对 $a,b\in H^*(X,\mathbb C)$，定义
\[
a*b=\sum_{i,j}\sum_{\beta\in H_2(X,\mathbb Z)}
\langle I_{0,3,\beta}\rangle(a,b,T_i)g^{ij}q^\beta T_j,
\qquad q^\beta=e^{2\pi i\int_\beta\omega}.
\]
称 $a*b$ 为 $a,b$ 的\emph{小量子乘积}。
\end{definition}
\index{量子乘积!小量子乘积|hyperpage}
按此写法，这一定义确实不够精确，因为对 $\beta$ 的求和可能无限，因而需要收敛性结果。
\hyperref[subsec:formal-small-quantum]{第~\ref*{subsec:formal-small-quantum} 小节}将说明几种赋予
定义~\ref{def:small-quantum-product} 严格意义的方法：以适当的形式表达式代替 $q^\beta$。
事实上，这才是通常的做法。不过，我们仍愿先给出上面这个不完整的定义，
因为它明确显示小量子乘积依赖 Kähler 模，而这是量子上同调的核心思想之一。
请读者暂且“悬置疑虑”，直到该小节；届时会看到，本节与下一节的所有断言，
经适当解释后，都是完全严格的。

采用物理学中“升指标”的记号，定义
\[
T^i=\sum_jg^{ij}T_j,
\]
于是 $T^0,\ldots,T^m$ 构成 $T_0,\ldots,T_m$ 的对偶基，满足
$\int_XT^i\cup T_j=\delta^i_j$。这样，小量子乘积可更简洁地写为
\begin{equation}\label{eq:small-quantum-product}
a*b=\sum_i\sum_{\beta\in H_2(X,\mathbb Z)}\langle I_{0,3,\beta}\rangle(a,b,T_i)q^\beta T^i.
\end{equation}
本章与下一章都使用此记号。

\hyperref[chap:gromov-witten]{第~\ref*{chap:gromov-witten} 章}为有理上同调类定义了 Gromov--Witten 不变量。
在上面的 $*$ 公式中，我们通过线性性将定义扩张到复系数，
使 $\langle I_{0,3,\beta}\rangle(a,b,T_i)$ 对 $a,b\in H^*(X,\mathbb C)$ 有意义。
还要注意，对 $\beta\in H_2(X,\mathbb Z)$ 与齐次 $a,b$，
次数条件~\eqref{eq:gw-top-degree-condition} 表明定义~\ref{def:small-quantum-product} 中的
Gromov--Witten 不变量 $\langle I_{0,3,\beta}\rangle(a,b,T_i)$ 为零，除非
\begin{equation}\label{eq:quantum-three-point-degree}
\deg a+\deg b+\deg T_i=2\dim X-2\int_\beta\omega_X,
\end{equation}
其中 $\omega_X$ 是 $X$ 的典范类。因此，上述定义只对这些 $\beta$ 求和。
也可检验，二元运算 $*$ 不依赖所选的基 $T_i$。

对于 Calabi--Yau 流形，有如下猜想。
\begin{conjecture}\label{conj:small-quantum-convergence}
对于 Calabi--Yau 流形，只要 $\omega$ 的虚部足够大，定义~\ref{def:small-quantum-product} 中的和便收敛。
\end{conjecture}
暂且假设此猜想成立。如前所述，\hyperref[subsec:formal-small-quantum]{第~\ref*{subsec:formal-small-quantum} 小节}
将通过重新解释 $q^\beta$，说明如何避免收敛问题。
幸运的是，对一些射影流形，已知上述求和性质良好，
因而能在 Kähler 模的意义下严格定义小量子乘积。
\begin{proposition}\label{prop:quantum-finite-fano}
设 $X$ 是光滑射影簇，满足以下条件之一：
\begin{enumerate}[label=(\roman*)]
\item $X$ 是 Fano 簇；
\item $X$ 上有半单 Lie 群的可迁作用。
\end{enumerate}
则对任意 $a,b\in H^*(X,\mathbb C)$，定义~\ref{def:small-quantum-product} 中的和为有限和。
\end{proposition}
\SourcePage{242}
\begin{proof}
需证明：$X$ 上满足式~\eqref{eq:quantum-three-point-degree} 且可由亏格零稳定映射表示的类 $\beta$ 只有有限多个。
先设 $X$ 为 Fano 簇。忽略有限的挠部分后，这些 $\beta$ 是 $H_2(X,\mathbb R)$ 中的格点。
Fano 条件意味着 $-\omega_X$ 位于 Kähler 锥的内部，
所以线性泛函 $\beta\mapsto-\int_\beta\omega_X$ 在有效 $\beta$ 生成的锥的闭包上严格为正（原点除外）。
于是，对任意常数 $C$，满足 $-\int_\beta\omega_X<C$ 的 $\beta$ 只有有限多个。
由式~\eqref{eq:quantum-three-point-degree}，得到 Fano 情形的结论。

再设 $X$ 为齐次空间 $G/P$。此时 $X$ 的上同调已知；
由亏格零稳定映射表示的类 $\beta$，是有限多个类 $\beta_1,\ldots,\beta_r$ 的非负整系数线性组合，
每个 $\beta_i$ 都是非恒定映射 $\mu_i:\mathbb P^1\to X$ 的像类。
由式~\eqref{eq:quantum-three-point-degree}，只需证明
$\int_{\beta_i}\omega_X<0$，等价地，$\int_{\beta_i}c_1(T_X)>0$。

可迁群作用表明，$T_X$ 由全局截面生成，因此 $\mu_i^*(T_X)$ 也由全局截面生成。
它是 $\mathbb P^1$ 上的向量丛，所以
$\mu_i^*(T_X)\simeq\bigoplus_j\mathcal O_{\mathbb P^1}(a_j)$，其中所有 $a_j\geq0$。
若能证明某个 $a_j>0$，便有 $\int_{\beta_i}c_1(T_X)>0$。
否则，$\mu_i^*(T_X)$ 平凡，而由于 $T_{\mathbb P^1}\simeq\mathcal O_{\mathbb P^1}(2)$，
自然映射 $T_{\mathbb P^1}\to\mu_i^*(T_X)$ 必为零。
这与 $\mu_i$ 非恒定矛盾，命题得证。
\end{proof}
命题的第二部分说明，对任意射影空间、Grassmann 簇、二次超曲面、旗流形，
以及这些簇的任意乘积，定义~\ref{def:small-quantum-product} 中的和都是有限的。

接着证明 $H^*(X,\mathbb C)$ 在小量子乘积下成为环。
\begin{theorem}\label{thm:small-quantum-ring}
$H^*(X,\mathbb C)$ 在小量子乘积下是有单位元的超交换环。
\end{theorem}
\begin{proof}
所谓超交换，是指
\[
a*b=(-1)^{\deg a\deg b}b*a,
\]
与杯积类似。这是因为\hyperref[sec:gw-properties]{第~\ref*{sec:gw-properties} 节}的等变性公理给出
$\langle I_{0,3,\beta}\rangle(a,b,T_i)=(-1)^{\deg a\deg b}\langle I_{0,3,\beta}\rangle(b,a,T_i)$。
至于单位元，取 $T_0=1=[X]\in H^0(X,\mathbb C)$ 为 $X$ 的基本类。
它是杯积的单位元，下面证明也是 $*$ 的单位元。
先注意，当 $\beta\ne0$ 时，$\langle I_{0,3,\beta}\rangle(a,T_0,T_i)=0$，
这是等变性公理与\IndexLocation{indexloc:0502:242}基本类公理的直接推论。因此
\[
\begin{aligned}
a*T_0&=\sum_i\langle I_{0,3,0}\rangle(a,T_0,T_i)T^i\\
&=\sum_i\left(\int_Xa\cup T_0\cup T_i\right)T^i=a\cup T_0=a.
\end{aligned}
\]
第二个等式由\IndexLocation{indexloc:0505:242}点映射公理得到，第三个等式由 $T^i$ 的定义得到。
所以 $T_0$ 是 $*$ 的单位元。

最后一步是证明 $*$ 满足结合律。
\hyperref[sec:gw-properties]{第~\ref*{sec:gw-properties} 节}的\IndexLocation{indexloc:0507:242}分裂公理在证明中起关键作用。
给定 $H^*(X,\mathbb C)$ 中的齐次类 $a,b,c$，先考虑 $(a*b)*c$。
用定义~\ref{def:small-quantum-product} 展开，并利用
\SourcePage{243}
$T^i=\sum_jg^{ij}T_j$，得到
\[
\sum_\beta\sum_{t,u}\left(\sum_{\beta=\beta_1+\beta_2}\sum_{r,s}
g^{rs}\langle I_{0,3,\beta_1}\rangle(a,b,T_r)
\langle I_{0,3,\beta_2}\rangle(T_s,c,T_t)\right)g^{tu}q^\beta T_u.
\]
现考虑 Gromov--Witten 类 $I_{0,4,\beta}(a,b,c,T_t)$。
分裂公理描述它沿式~\eqref{eq:gw-splitting-gluing-map} 定义的映射
\[
\varphi:\overline M_{0,3}\times\overline M_{0,3}\longrightarrow\overline M_{0,4}\simeq\mathbb P^1
\]
的拉回。注意，除非 $I_{0,4,\beta}(a,b,c,T_t)\in H^0(\overline M_{0,4})$，此拉回为零；
下面假设它确实位于此空间。由于 $\varphi$ 是 $\mathbb P^1$ 中一点的嵌入，可将
\begin{equation}\label{eq:small-associativity-left}
\varphi^*I_{0,4,\beta}(a,b,c,T_t)
\end{equation}
视为一个数。分裂公理于是说，式~\eqref{eq:small-associativity-left}
恰好是上面显示式的大圆括号内的量。

转而考虑 $a*(b*c)$，得到
\[
\sum_\beta\sum_{t,u}\left(\sum_{\beta=\beta_1+\beta_2}\sum_{r,s}
g^{rs}\langle I_{0,3,\beta_1}\rangle(b,c,T_r)
\langle I_{0,3,\beta_2}\rangle(a,T_s,T_t)\right)g^{tu}q^\beta T_u.
\]
由等变性公理，
$\langle I_{0,3,\beta_2}\rangle(a,T_s,T_t)=
(-1)^{\deg a\deg T_s}\langle I_{0,3,\beta_2}\rangle(T_s,a,T_t)$；
而由式~\eqref{eq:quantum-three-point-degree}，有
$\deg a+\deg T_s+\deg T_t\equiv0\pmod2$。因此
\[
\langle I_{0,3,\beta_2}\rangle(a,T_s,T_t)
=(-1)^{\deg a(\deg a+\deg T_t)}\langle I_{0,3,\beta_2}\rangle(T_s,a,T_t).
\]
所以 $a*(b*c)$ 表达式中括号内的量为
\begin{equation}\label{eq:small-associativity-right}
(-1)^{\deg a(\deg a+\deg T_t)}\widetilde\varphi^*I_{0,4,\beta}(b,c,a,T_t),
\end{equation}
其中 $\widetilde\varphi$ 与 $\varphi$ 类似。
注意 $\widetilde\varphi$ 是 $\mathbb P^1$ 中另一个点的嵌入，故与 $\varphi$ 线性等价。

比较 $(a*b)*c$ 与 $a*(b*c)$ 的表达式，可见，只要证明
式~\eqref{eq:small-associativity-left} 与式~\eqref{eq:small-associativity-right} 相等，定理立即成立。
所以只需证明
\[
I_{0,4,\beta}(a,b,c,T_t)
=(-1)^{\deg a(\deg a+\deg T_t)}I_{0,4,\beta}(b,c,a,T_t).
\]
这是等变性公理与由式~\eqref{eq:gw-top-degree-condition} 得到的
$\deg a+\deg b+\deg c+\deg T_t\equiv0\pmod2$ 的直接推论。
\end{proof}
注意，上述证明与\hyperref[subsec:plane-genus-zero]{第~\ref*{subsec:plane-genus-zero} 小节}
证明 $\mathbb P^2$ 的 Gromov--Witten 不变量递推关系时所用的方法，有许多共同之处。
这里出现了奇数次上同调类，因此还需处理符号。
关于结合律的直观解释，见 \SourceCite{Morrison7}。

除赋予环结构外，小量子乘积还有其他性质。
先描述 $a*b$ 的次数。下一小节的例子将说明，即使 $a,b$ 齐次，$a*b$ 也未必齐次。
不过，$a*b$ 中出现的次数并非任意。
\SourcePage{244}
\begin{proposition}\label{prop:small-quantum-degree}
若 $a,b\in H^*(X,\mathbb C)$ 齐次，则其小量子乘积为齐次类之和 $a*b=\sum_lc_l$，且
\[
\deg c_l\equiv\deg a+\deg b\pmod2.
\]
此外，若 $\omega_X\in H^2(X,\mathbb Z)$ 的类可被整数 $r$ 整除，则
$\deg c_l\equiv\deg a+\deg b\pmod{2r}$。
\end{proposition}
\begin{proof}
由于只有 $\deg T_i+\deg T_j=2\dim X$ 时 $g^{ij}$ 才可能非零，
\IndexLocation{indexloc:1133:244}次数条件~\eqref{eq:quantum-three-point-degree} 表明，在定义~\ref{def:small-quantum-product} 中，
$T_j$ 的系数非零只可能在
\[
\deg a+\deg b=\deg T_j-2\int_\beta\omega_X
\]
时发生。命题立即成立。
\end{proof}
还可描述小量子乘积与 $X$ 上交配对的关系。
记此配对为 $g(a,b)=\int_Xa\cup b$。
\begin{proposition}\label{prop:quantum-frobenius-pairing}
对所有 $a,b,c\in H^*(X,\mathbb C)$：
\begin{enumerate}[label=(\roman*)]
\item $g(a*b,c)=g(a,b*c)=\sum_{\beta\in H_2(X,\mathbb Z)}\langle I_{0,3,\beta}\rangle(a,b,c)q^\beta$。
\item $\int_Xa*b=g(a,b)$。
\end{enumerate}
\end{proposition}
\begin{proof}
容易证明 $g(T_i*T_j,T_k)=\sum_\beta\langle I_{0,3,\beta}\rangle(T_i,T_j,T_k)q^\beta$。
转而考虑 $g(T_i,T_j*T_k)$，杯积的超交换性给出
\[
g(T_i,T_j*T_k)=(-1)^{\deg T_i(\deg T_j+\deg T_k)}g(T_j*T_k,T_i),
\]
而由上述等式，$g(T_j*T_k,T_i)=\sum_\beta\langle I_{0,3,\beta}\rangle(T_j,T_k,T_i)q^\beta$。
第一部分随即由\IndexLocation{indexloc:0501:244}等变性公理得到。为证明第二部分，注意
\[
\int_Xa*b=\int_XT_0\cup(a*b)=g(T_0,a*b)=g(T_0*a,b)=g(a,b),
\]
因为 $T_0$ 是杯积与小量子乘积共同的单位元。
\end{proof}
此命题表明，$H^*(X,\mathbb C)$ 在 $*$ 与 $g$ 下是一个 \emph{Frobenius 代数}，见定义~\ref{def:commutative-frobenius-algebra}。
\index{Frobenius代数@Frobenius 代数|hyperpage}
命题~\ref{prop:quantum-frobenius-pairing} 也与\IndexLocation{indexloc:1129:244}\IndexLocation{indexloc:1144:244}小量子上同调的另一种略有不同的方法有关。
该方法从\emph{三点函数} $\langle a,b,c\rangle$ 出发，定义为
\begin{equation}\label{eq:a-three-point-function}
\index{三点函数!量子上同调|hyperpage}
\langle a,b,c\rangle=\sum_{\beta\in H_2(X,\mathbb Z)}\langle I_{0,3,\beta}\rangle(a,b,c)q^\beta.
\end{equation}
也称 $\langle a,b,c\rangle$ 为 \emph{A \IndexLocation{indexloc:0185:244}模型\IndexLocation{indexloc:0195:244}关联函数}。
容易看出，$a,b$ 的小量子乘积是满足
$g(a*b,c)=\langle a,b,c\rangle$（对所有 $c\in H^*(X,\mathbb C)$）的唯一元素 $a*b\in H^*(X,\mathbb C)$。
文献中常采用这种方式定义量子上同调，见 \SourceCite{MS,Morrison7,Voisin3}。
还要注意，式~\eqref{eq:small-quantum-product} 与式~\eqref{eq:a-three-point-function} 给出
\begin{equation}\label{eq:small-product-from-correlations}
a*b=\sum_i\langle a,b,T_i\rangle T^i,
\end{equation}
\SourcePage{245}
这再次证明，量子上同调由三点函数确定。
最后还应提到，命题~\ref{prop:quantum-frobenius-pairing} 给出
\begin{equation}\label{eq:three-point-quantum-integral}
\langle a,b,c\rangle=\int_Xa*b*c,
\end{equation}
所以三点函数能完全用小量子乘积表示。

对于 Calabi--Yau 流形，小量子乘积尤其好。
\begin{corollary}\label{cor:cy-quantum-hodge}
若 $V$ 为 Calabi--Yau 流形，则小量子乘积满足 $\deg(a*b)=\deg a+\deg b$。
此外，若 $a\in H^{p,q}(V)$、$b\in H^{r,s}(V)$，则 $a*b\in H^{p+r,q+s}(V)$。
\end{corollary}
\begin{proof}
第一个断言立即由命题~\ref{prop:small-quantum-degree} 得到，因为 $\omega_V$ 平凡。
设 $V$ 的维数为 $d$。对于第二个断言，由命题~\ref{prop:quantum-frobenius-pairing} 的第一部分，只需证明
\[
\langle I_{0,3,\beta}\rangle(a,b,c)=0,
\qquad c\in H^{u,v}(V),\quad(u,v)\ne(d-p-r,d-q-s).
\]
但由于 $\overline M_{0,3}$ 是一点，定义~\ref{def:gw-algebraic-invariants} 给出
\[
\langle I_{0,3,\beta}\rangle(a,b,c)=\int_{\pi_{1*}(\xi)}a\otimes b\otimes c,
\]
其中 $\pi_1:\overline M_{0,3}(V,\beta)\to V^3$ 是\IndexLocation{indexloc:0276:245}赋值映射，$\xi$ 为\IndexLocation{indexloc:1415:245}虚拟基本类。
由于 $\xi$ 是预期维数的代数类，由式~\eqref{eq:stable-map-expected-dimension}，
$\pi_{1*}(\xi)\in H_{2d}(V^3)$ 是代数圈的同调类。
上述积分只有在 $a\otimes b\otimes c$ 的型为 $(d,d)$ 时才可能非零，结论随之成立。
\end{proof}
推论~\ref{cor:cy-quantum-hodge} 表明：当 $V$ 为 Calabi--Yau 流形时，
$\bigoplus_{p=0}^dH^{p,p}(V)$ 是赋予小量子乘积后的 $H^*(V,\mathbb C)$ 的子环。
本章稍后研究 A 侧 Hodge \IndexLocation{indexloc:0019:245}结构变分时，将用到这一点。

\subsection{例题}\label{subsec:small-quantum-examples}
\label{subsec:a-model-correlation}
现在计算\IndexLocation{indexloc:1129:245}\IndexLocation{indexloc:1133:245}\IndexLocation{indexloc:1137:245}\IndexLocation{indexloc:1138:245}小量子上同调环的几个例子。
\SourceExampleNumber{8.1.2}{1}
\begin{example}\label{ex:projective-quantum-ring}
对 $\mathbb P^r$，使用上同调基 $H^i$，$0\leq i\leq r$，
其中 $H$ 为超平面类，指数表示杯积。我们断言
\begin{equation}\label{eq:projective-quantum-product}
H^i*H^j=\begin{cases}
H^{i+j},&i+j\leq r,\\
qH^{i+j-r-1},&i+j\geq r+1,
\end{cases}
\end{equation}
其中 $q=e^{2\pi i\int_\ell\omega}$，$\ell$ 是 $\mathbb P^r$ 中的直线。
为证明此式，取 $\beta=d[\ell]$，$d\geq0$。
由式~\eqref{eq:quantum-three-point-degree}，不变量
$\langle I_{0,3,d}\rangle(H^i,H^j,H^k)$（这里用 $d$ 代替 $\beta$）为零，除非
\begin{equation}\label{eq:projective-quantum-degree}
i+j+k=r+d(r+1).
\end{equation}
由于 $i,j,k\leq r$，$d$ 只能为 $0$ 或 $1$；进一步，若 $i+j\leq r$，则 $d=0$。
因此容易得到，当 $i+j\leq r$ 时，$H^i*H^j=H^{i+j}$。
特别地，对于 $0\leq i\leq r$，$H^i$ 对杯积与小量子乘积具有相同含义。

为完成证明，只需证明 $H*H^r=qH^0$。
此时由式~\eqref{eq:projective-quantum-degree}，$d=1$、$k=r$。
于是容易得到
\[
H*H^r=\langle I_{0,3,1}\rangle(H,H^r,H^r)qH^0.
\]
\SourcePage{246}
与\hyperref[subsec:plane-genus-zero]{第~\ref*{subsec:plane-genus-zero} 小节}一样，
$\langle I_{0,3,1}\rangle(H,H^r,H^r)=\langle I_{0,3,1}\rangle(H,[pt],[pt])=1$，
因为 $\mathbb P^r$ 中两个点确定唯一一条直线，它在第三个点与 $H$ 相交。
这证明了式~\eqref{eq:projective-quantum-product}。
还要注意，由于 $\omega_{\mathbb P^r}=\mathcal O_{\mathbb P^r}(-r-1)$ 可被 $r+1$ 整除，
此式是命题~\ref{prop:small-quantum-degree} 的一个很好的例子。

由于 $H^0=[\mathbb P^r]$ 是单位元，且在小量子乘积下 $H^{r+1}=qH^0$，
$\mathbb P^r$ 的小量子环同构于
$\mathbb C[H]/\langle H^{r+1}-q\rangle$。
另一方面，通常的上同调环同构于 $\mathbb C[H]/\langle H^{r+1}\rangle$。
当 $\omega$ 的虚部增大时，$q\to0$，可见小量子乘积在极限中成为杯积。
\end{example}

\SourceExampleNumber{8.1.2}{2}
\begin{example}\label{ex:batyrev-quantum-ring}
\SourceCite{Batyrev2} 为光滑射影环面簇 $X$ 定义了一个“量子上同调环”。
结果表明，若 $X$ 为 Fano 簇，这个环与小量子上同调环一致；否则可能不同。
为避免混淆，记 Batyrev 的环为 $H_\omega^*(X)$。
\SourceCite{QR} 证明，对于射影空间上的某些射影丛，$H_\omega^*(X)$ 与小量子上同调环一致。
例题~\ref{ex:toric-fano-and-hirzebruch-quantum-ring} 将解释为何它们在 Fano 情形一致（推广 \SourceCite{QR} 的结果），
并给出非 Fano 情形的一个例子，表明一般情况下 $H_\omega^*(X)$ 确实与小量子上同调环不同。
该例题还将说明，无论如何，$H_\omega^*(X)$ 与小量子上同调环密切相关。
现在讨论 $H_\omega^*(X)$ 的定义和性质。
我们会看到，它的性质与预期的量子上同调性质相似（至少在 Fano 情形下应当如此）。

置 $X=X_\Sigma$，其中 $\Sigma$ 是 $N_{\mathbb R}$ 中给出 $X$ 的扇。
设 $v_1,\ldots,v_s\in N$ 是 $\Sigma$ 的一维锥的本原整生成元；
每个 $v_i$ 对应环面不变除子 $D_i\subset X$。
先回忆 $X$ 的复系数上同调环的通常描述。
给每个 $v_i,D_i$ 配一个变量 $x_i$，得到多项式环 $\mathbb C[x_1,\ldots,x_s]$。
考虑其中两个理想：
\[
\begin{aligned}
P(\Sigma)&=\left\langle\sum_{i=1}^s\langle m,v_i\rangle x_i:m\in M\right\rangle,\\
SR(\Sigma)&=\langle x_{i_1}\cdots x_{i_k}:\{v_{i_1},\ldots,v_{i_k}\}\not\subset\sigma
\text{，对所有 }\sigma\in\Sigma\rangle.
\end{aligned}
\]
$SR(\Sigma)$ 是 $\Sigma$ 的 \emph{Stanley--Reisner 理想}。
众所周知，见 \SourceCite{Fulton3,Oda}，将 $[D_i]\in H^2(X,\mathbb C)$ 映到 $x_i$，诱导同构
\begin{equation}\label{eq:toric-cohomology-sr}
\index{Stanley--Reisner理想@Stanley--Reisner 理想|hyperpage}
H^*(X,\mathbb C)\simeq\mathbb C[x_1,\ldots,x_s]/(P(\Sigma)+SR(\Sigma)).
\end{equation}
还需要 Batyrev 对此作的一个修改，见 \SourceCite{Batyrev2}。
\hyperref[subsec:homogeneous-ring]{第~\ref*{subsec:homogeneous-ring} 小节}定义：
若 $\mathcal P=\{v_{i_1},\ldots,v_{i_k}\}$ 不包含在任何 $\sigma\in\Sigma$ 中，
但 $\mathcal P$ 的每个真子集都包含在某个锥中，则称 $\mathcal P$ 为\emph{本原集合}。
于是容易看出
\begin{equation}\label{eq:sr-primitive-collections}
SR(\Sigma)=\langle x_{i_1}\cdots x_{i_k}:\{v_{i_1},\ldots,v_{i_k}\}\text{ 为本原集合}\rangle.
\end{equation}
现在考虑量子上同调。对于 $\mathbb P^r$，我们从杯积的 $\mathbb C[h]/\langle h^{r+1}\rangle$
过渡到量子乘积的 $\mathbb C[h]/\langle h^{r+1}-q\rangle$。
为将此推广到环面簇 $X$，先引入一些记号。
设 $\mathcal P=\{v_{i_1},\ldots,v_{i_k}\}$ 为本原集合。
则和 $v_{\mathcal P}=v_{i_1}+\cdots+v_{i_k}$ 位于某个锥 $\sigma\in\Sigma$ 中，
\SourcePage{247}
所以可写成 $v_{\mathcal P}=c_1v_{j_1}+\cdots+c_lv_{j_l}$，
其中 $v_{j_1},\ldots,v_{j_l}$ 为 $\sigma$ 的生成元，$c_1,\ldots,c_l\geq0$。
这给出 $v_1,\ldots,v_s$ 之间的非平凡关系
\begin{equation}\label{eq:quantum-primitive-relation}
v_{i_1}+\cdots+v_{i_k}-c_1v_{j_1}-\cdots-c_lv_{j_l}=0.
\end{equation}
这些关系按如下方式对应于同调类。
由于 $X$ 光滑，$\operatorname{Pic}(X)=A_{n-1}(X)=H^2(X,\mathbb Z)$，
且光滑情形中的 $A_{n-1}(X)$ 无挠，因此\hyperref[chap:toric]{第~\ref*{chap:toric} 章}正合列~\eqref{eq:toric-chow-sequence} 的对偶为
\[
0\longrightarrow H_2(X,\mathbb Z)\longrightarrow\mathbb Z^s\longrightarrow N\longrightarrow0.
\]
这里 $(\lambda_1,\ldots,\lambda_s)\in\mathbb Z^s$ 映到 $\sum_{i=1}^s\lambda_iv_i\in N$。
于是有自然同构
\begin{equation}\label{eq:toric-relation-homology}
\left\{(\lambda_1,\ldots,\lambda_s)\in\mathbb Z^s:\sum_{i=1}^s\lambda_iv_i=0\right\}
\simeq H_2(X,\mathbb Z).
\end{equation}
特别地，关系~\eqref{eq:quantum-primitive-relation} 对应于一个同调类，
记为 $\beta(\mathcal P)\in H_2(X,\mathbb Z)$，参见引理~\ref{lem:mori-wall-curve}。
\SourceCite[定理 2.15]{Batyrev1} 中，Batyrev 证明 $\beta(\mathcal P)$ 是有效的一维圈类。

现给定复化 Kähler 类 $\omega\in H^{1,1}(X)=H^2(X,\mathbb C)$，定义\emph{量子 Stanley--Reisner 理想}
\begin{equation}\label{eq:quantum-sr-primitive}
SR_\omega(\Sigma)=\langle x_{i_1}\cdots x_{i_k}
-q^{\beta(\mathcal P)}x_{j_1}^{c_1}\cdots x_{j_l}^{c_l}:\mathcal P\text{ 为本原集合}\rangle,
\end{equation}
其中，如通常一样，$q^{\beta(\mathcal P)}=e^{2\pi i\int_{\beta(\mathcal P)}\omega}$。
Batyrev 随后定义环
\begin{equation}\label{eq:batyrev-quantum-ring}
H_\omega^*(X)=\mathbb C[x_1,\ldots,x_s]/(P(\Sigma)+SR_\omega(\Sigma)),
\end{equation}
见 \SourceCite[定理 5.3]{Batyrev2}。
称 $H_\omega^*(X)$ 为 \emph{Batyrev 量子环}。
\index{Batyrev量子环@Batyrev 量子环|hyperpage}\index{Stanley--Reisner理想@Stanley--Reisner 理想!量子理想|hyperpage}
为了解环~\eqref{eq:batyrev-quantum-ring} 与通常上同调环~\eqref{eq:toric-cohomology-sr} 的关系，
让 $\omega$ 的虚部趋于无穷大。
由于 $\beta(\mathcal P)$ 是有效圈，$\int_{\beta(\mathcal P)}\omega$ 的虚部也增大，
从而 $q^{\beta(\mathcal P)}\to0$。
因此，$\omega$ 的虚部趋于无穷大时，量子 Stanley--Reisner 理想~\eqref{eq:quantum-sr-primitive}
成为通常的 Stanley--Reisner 理想~\eqref{eq:sr-primitive-collections}。
这说明 $X$ 的通常上同调环是 $H_\omega^*(X)$ 的一个极限。

检验 $X=\mathbb P^r$ 的情形。
锥的生成元为 $e_0,\ldots,e_r$，其中 $e_1,\ldots,e_r$ 生成格，$e_0=-\sum_{i=1}^re_i$。
所以 $H^*(\mathbb P^r,\mathbb C)$ 的描述中的理想 $P(\Sigma)$ 为
$\langle x_1-x_0,\ldots,x_r-x_0\rangle$。
此外，$\mathcal P=\{e_0,\ldots,e_r\}$ 是唯一的本原集合。
由于 $e_0+\cdots+e_r=0$，式~\eqref{eq:quantum-primitive-relation} 中的 $c_j$ 都为 $0$。
因此，量子 Stanley--Reisner 理想~\eqref{eq:quantum-sr-primitive} 为
$SR_\omega(\Sigma)=\langle x_0\cdots x_r-q^{\beta(\mathcal P)}\rangle$。
结合 $P(\Sigma)$ 与 $SR_\omega(\Sigma)$，环~\eqref{eq:batyrev-quantum-ring} 化为
$\mathbb C[x_0]/\langle x_0^{r+1}-q^{\beta(\mathcal P)}\rangle$。
还需确定 $\beta(\mathcal P)\in H_2(\mathbb P^r,\mathbb Z)$。
因为它有效，而关系 $e_0+\cdots+e_r=0$ 生成 $e_i$ 之间全部关系构成的 $\mathbb Z$ 模，
必有 $\beta(\mathcal P)=[\ell]$，另见引理~\ref{lem:mori-wall-curve} 及其证明之后的一句话。
因此，重新得到前面对 $\mathbb P^r$ \IndexLocation{indexloc:1129:247}小量子上同调的描述。

环 $H_\omega^*(X)$ 还有其他描述。若
\begin{equation}\label{eq:toric-quantum-general-relation}
\sum_{i=1}^sa_iv_i=\sum_{i=1}^sb_iv_i,\qquad a_i,b_i\in\mathbb Z_{\geq0},
\end{equation}
\SourcePage{248}
则由式~\eqref{eq:toric-relation-homology}，
$a-b=(a_1-b_1,\ldots,a_s-b_s)$ 给出元素 $\beta(a,b)\in H_2(X,\mathbb Z)$。
\SourceCite{Batyrev2} 中，Batyrev 证明量子 Stanley--Reisner 理想也可写为
\begin{equation}\label{eq:quantum-sr-all-relations}
SR_\omega(\Sigma)=\langle x_1^{a_1}\cdots x_s^{a_s}
-q^{\beta(a,b)}x_1^{b_1}\cdots x_s^{b_s}:a_i,b_i\text{ 满足式~\eqref{eq:toric-quantum-general-relation}}\rangle.
\end{equation}
这一描述说明，$H_\omega^*(X)$ 只依赖 $\Sigma$ 的一维锥与复化 Kähler 类 $\omega$。
所以，具有相同一维骨架的不同扇，量子上同调密切相关。
例如，设 $X_1,X_2$ 是光滑环面簇，它们的扇有相同一维骨架，
并假设它们的 Kähler 锥在\hyperref[sec:gkz-decomposition]{第~\ref*{sec:gkz-decomposition} 节}
描述的 GKZ 分解中有共同的墙。
那么 $X_1,X_2$ 通过小双有理变换联系；若该变换保典范，则是一条翻越。
因此，它们的通常上同调环可能不同构。
不过，利用 Batyrev 量子环，可以按以下方式从一个过渡到另一个：
\[
\begin{array}{ccccc}
\text{通常}&&\text{量子}\quad\text{量子}&&\text{通常}\\[-2pt]
H^*(X_1,\mathbb C)&\longleftarrow&H_\omega^*(X_1)\longleftrightarrow H_\omega^*(X_2)
&\longrightarrow&H^*(X_2,\mathbb C).
\end{array}
\]
左边箭头表示：当 $\omega$ 的虚部在 $X_1$ 的 Kähler 锥中趋于无穷大时，
$X_1$ 的 Batyrev 量子环趋于通常上同调环。
右边箭头类似，只是 $\omega$ 的虚部在 $X_2$ 的 Kähler 锥中趋于无穷大。
中间箭头则对应于 $\omega$ 穿过这堵墙。
只要保留相同的一维锥，式~\eqref{eq:quantum-sr-all-relations} 就允许改变扇，这使上述过渡成为可能。

可如下理解 $SR_\omega(\Sigma)$ 的两种描述：
式~\eqref{eq:quantum-sr-all-relations} 只用一维骨架，说明如何将 $H_\omega^*(X)$ 扩展到整个 GKZ 分解；
式~\eqref{eq:quantum-sr-primitive} 则通过本原集合使用整个扇，
说明如何在任意一个 Kähler 锥中取极限。进一步讨论与例子见 \SourceCite{Batyrev2}。

最后，在令 $\omega$ 的虚部趋于无穷大时，我们一直以很朴素的方式取理想 $SR_\omega(\Sigma)$ 的极限。
正确做法当然是使用单项式序与 Gröbner 基。
事实上，$\omega$ 的虚部给出一个单项式序，见 \SourceCite{Batyrev2} 定理 5.3 的证明。
此外，按所给单项式序将所有可能的 $\operatorname{Im}\omega$ 分组，
就得到 GKZ 分解的一个细分，称为 \emph{Gröbner 扇}。
\SourceCite{Sturmfels1} 描述了这个扇，\SourceCite{HLY1,HLY2} 将它用于镜像对称，研究超几何方程的解。
\index{Gröbner扇@Gröbner 扇|hyperpage}
\end{example}

\SourceExampleNumber{8.1.2}{3}
\begin{example}\label{ex:cy-small-quantum}
最后，考虑光滑 Calabi--Yau 三维簇 $V$ 的小量子上同调。
先证明，大多数情况下，小量子乘积与杯积一致。
\begin{lemma}\label{lem:cy-quantum-cup}
若 $V$ 是光滑 Calabi--Yau 三维簇，则只要 $\deg a\ne2$ 或 $\deg b\ne2$，就有 $a*b=a\cup b$。
\end{lemma}
\begin{proof}
先注意，当 $\beta\ne0$，且某个 $\deg\alpha_i\ne2$ 时，
$\langle I_{0,3,\beta}\rangle(\alpha_1,\alpha_2,\alpha_3)=0$。
\hyperref[subsec:cy-genus-zero]{第~\ref*{subsec:cy-genus-zero} 小节}已对五次三维簇证明了这一点，
该证明也适用于任意 Calabi--Yau 三维簇。引理随即得到。
\end{proof}
还需计算 $H^2(V,\mathbb C)$ 上的量子乘积。先设 $V$ 为五次三维簇。
如往常一样，用 $\langle I_{0,3,d}\rangle$ 代替 $\langle I_{0,3,\beta}\rangle$。
一旦知道 $H*H$，就唯一确定小量子乘积，其中
\SourcePage{249}
$H\in H^2(V,\mathbb Z)$ 是超平面类。
由推论~\ref{cor:cy-quantum-hodge}，$H*H\in H^4(V,\mathbb C)$，
然后定义~\ref{def:small-quantum-product} 与 $\int_VH^3=5$ 给出
\[
H*H=\frac15\sum_{d=0}^{\infty}\langle I_{0,3,d}\rangle(H,H,H)q^dH^2.
\]
还有由式~\eqref{eq:a-three-point-function}、\eqref{eq:three-point-quantum-integral} 定义的三点函数，
这里变成
\begin{equation}\label{eq:quintic-quantum-correlation}
\langle H,H,H\rangle=\int_VH*H*H
=\sum_{d=0}^{\infty}\langle I_{0,3,d}\rangle(H,H,H)q^d.
\end{equation}
\hyperref[subsec:cy-genus-zero]{第~\ref*{subsec:cy-genus-zero} 小节}中已看到，
此式右边就是\hyperref[chap:quintic]{第~\ref*{chap:quintic} 章}研究的 A \IndexLocation{indexloc:0187:249}\IndexLocation{indexloc:0188:249}\IndexLocation{indexloc:1183:249}模型\IndexLocation{indexloc:0195:249}关联函数 $\langle H,H,H\rangle$。
所以，五次三维簇的三点函数由小量子乘积给出。
如式~\eqref{eq:small-product-from-correlations} 所示，反过来也成立；结合上述两个公式，得到
\[
H*H=\langle H,H,H\rangle\frac15H^2.
\]
因此，量子乘积用三点函数表示出来。

对于任意 Calabi--Yau 三维簇 $V$，若限制到 $a,b,c\in H^2(V,\mathbb C)$，情形类似。
由式~\eqref{eq:a-three-point-function}、\eqref{eq:three-point-quantum-integral}，三点函数为
\begin{equation}\label{eq:cy-quantum-correlation-gw}
\begin{aligned}
\langle a,b,c\rangle&=\int_Va*b*c=\sum_\beta\langle I_{0,3,\beta}\rangle(a,b,c)q^\beta\\
&=\langle I_{0,3,0}\rangle(a,b,c)+\sum_{\beta\ne0}\langle I_{0,0,\beta}\rangle q^\beta
\int_\beta a\int_\beta b\int_\beta c.
\end{aligned}
\end{equation}
\hyperref[subsec:instanton-numbers]{第~\ref*{subsec:instanton-numbers} 小节}已介绍瞬子数 $n_\beta$，满足
\begin{equation}\label{eq:cy-instanton-mobius}
\langle I_{0,0,\beta}\rangle=\sum_{\beta=k\gamma}n_\gamma k^{-3}.
\end{equation}
如式~\eqref{eq:cy-a-correlation-instanton} 一样，由此得到
\begin{equation}\label{eq:cy-quantum-correlation-instanton}
\langle a,b,c\rangle=\int_Va*b*c=\int_Va\cup b\cup c
+\sum_{\substack{\beta\ne0\\\beta\text{ 有效}}}n_\beta\frac{q^\beta}{1-q^\beta}\int_\beta a\int_\beta b\int_\beta c.
\end{equation}
光滑 Calabi--Yau 三维簇三点函数的这一经典公式，最初对五次三维簇给出，见 \SourceCite{CdGP}，
并经常出现在文献中，见 \SourceCite{Witten5,AGM3,Morrison5,Voisin3}。
通过式~\eqref{eq:cy-instanton-mobius} 定义瞬子数 $n_\beta$ 时，这一公式具有明确意义。
但由\hyperref[subsec:instanton-numbers]{第~\ref*{subsec:instanton-numbers} 小节}可知，
$n_\beta$ 未必计数 $V$ 同调类为 $\beta$ 的有理曲线。
不过，猜想~\ref{conj:symplectic-instantons} 断言，$n_\beta$ 应与类 $\beta$ 中按定向带符号计数的 $J$ 全纯曲线数目密切相关。

能用量子上同调表达 Calabi--Yau 三维簇的三点函数，令人十分满意。
此外，如\hyperref[subsec:small-quantum-definition]{第~\ref*{subsec:small-quantum-definition} 小节}所见，
反过来也成立：可以用三点函数构造量子上同调。
\end{example}

\SourceExampleNumber{8.1.2}{4}
\begin{example}\label{ex:flop-quantum-invariance}
\index{翻越!量子上同调|hyperpage}
本例研究翻越对 Calabi--Yau 三维簇 $V$ 的\IndexLocation{indexloc:1129:249}\IndexLocation{indexloc:1134:249}\IndexLocation{indexloc:1139:249}\IndexLocation{indexloc:1142:249}\IndexLocation{indexloc:1202:249}小量子上同调的影响。
用\hyperref[subsec:threefold-kahler-moduli]{第~\ref*{subsec:threefold-kahler-moduli} 小节}的术语，设 $\phi:V\to\overline V$ 为本原 I 型收缩，
给出翻越 $V\dashrightarrow V'$，两边被收缩的曲线分别为 $C\subset V$ 与 $C'\subset V'$。
\SourcePage{250}
此时有自然识别 $H^2(V)\simeq H^2(V')$，在该识别下，Kähler 锥 $K(V)$ 与 $K(V')$ 沿共同的面相接。

为比较 $V,V'$ 的量子上同调，设可约曲线 $C$ 由同调类为 $\beta$、
彼此不相交且刚性嵌入的若干 $\mathbb P^1$ 构成。
进一步假设，$C$ 与所有其他类不为 $\beta$ 的曲线不相交。

在这一情形下，已知 $V,V'$ 有相同的 Hodge 数，但交形式可能不同，
所以通常上同调环可能不同构。不过，其小量子上同调环密切相关，见 \SourceCite{Witten5,AGM3,Morrison5}。
利用式~\eqref{eq:cy-quantum-correlation-instanton}，将 $V$ 上的三点函数写成
\[
\begin{aligned}
\langle a,b,c\rangle={}&\int_Va\cup b\cup c
+n_\beta\frac{q^\beta}{1-q^\beta}\int_\beta a\int_\beta b\int_\beta c\\
&+\sum_{\substack{\gamma\ne\beta,0\\\gamma\text{ 有效}}}n_\gamma\frac{q^\gamma}{1-q^\gamma}
\int_\gamma a\int_\gamma b\int_\gamma c.
\end{aligned}
\]
过渡到 $V'$ 时，右边第一项可能变化，第二项也可能变化，第三项则不变。
这是因为模空间 $\overline M_{0,3}(V,\gamma)$ 与 $V'$ 上相应的模空间同构，且这些同构保持虚拟基本类。
上面列出的文献中一个简单计算表明：只要通过 $H^2(V)\simeq H^2(V')$ 变换 $a,b,c$，
并通过该映射的 Poincaré 对偶变换 $\beta$，那么从 $V$ 过渡到 $V'$ 时，前两项之和
\[
\int_Va\cup b\cup c+n_\beta\frac{q^\beta}{1-q^\beta}\int_\beta a\int_\beta b\int_\beta c
\]
经相应解析延拓后\emph{保持不变}。

从形式角度看，在上述识别下，$V,V'$ 的三点函数给出相同的表达式；
容易看出，采用适当形式系数的小量子上同调环也相应一致。
不过，$V$ 的小量子乘积只有在 $\omega$ 的虚部足够深入 $K(V)$ 时才猜想收敛，$K(V')$ 的情形类似。
在 $H^2(V)\simeq H^2(V')$ 下，两个 Kähler 锥的内部不相交，
所以这个识别不再直接给出两者在数值系数下的量子上同调环同构，因为此时是在互不相交的 $\omega$ 集合上取值。
得到的是 $V$ 的三点函数到 $V'$ 的三点函数的解析延拓。
这与例题~\ref{ex:batyrev-quantum-ring} 研究翻越对环面簇量子上同调的影响时十分相似。
与那里一样，量子上同调使我们能将 $V$ 的通常上同调环形变为 $V'$ 的上同调环。

这些讨论都针对对应于 I 型收缩的 Kähler 锥的面。
\SourceCite{Wilson4} 研究了 III 型收缩的影响。
\end{example}

小量子乘积一直是活跃的研究领域。除刚才讨论的例子外，还能描述其他有趣\IndexLocation{indexloc:1140:250}簇的小量子上同调，包括：
\begin{itemize}
\item \textbf{Grassmann 簇。}Grassmann 簇的小量子环最初由 \SourceCite{Vafa2} 描述，
随后 \SourceCite{Witten6,ST} 给出了数学处理。
其他文献见 \SourceCite{MS,FP,Bertram,BCF}。
有关主题包括 Giambelli 与 Pieri 公式的量子版本、与 Verlinde 代数的联系，以及 Vafa--Intriligator 公式。
\SourcePage{251}
\item \textbf{旗流形。}旗流形的小量子环与 Toda 格有引人入胜的关系，最初在 \SourceCite{GK} 中描述。
其他文献包括 \SourceCite{AS,Kim1,Givental3,CiocanF,Kostant}。
\item \textbf{Fano 流形。}关于 Fano 流形的小量子环，已经有大量工作。
这类流形包括 Grassmann 簇、旗流形，以及某些射影超曲面与完全交。
相关论文包括 \SourceCite{ST,Beauville1,TX,CJ}。
\item \textbf{有理曲面。}\SourceCite{CM} 描述了一般有理曲面的小量子环。
对三次曲面，量子乘积的结合律基本上等价于存在 $27$ 条直线，且每条与另外 $10$ 条相交。
\end{itemize}

\subsection{系数}\label{subsec:formal-small-quantum}
\index{量子上同调!系数环|hyperpage}
\index{Novikov环@Novikov 环|hyperpage}
在迄今的讨论中，定义小量子乘积时使用复系数上同调。
这是因为，若 $\omega\in H^2(X,\mathbb C)$ 是复化 Kähler 类，则
$q^\beta=e^{2\pi i\int_\beta\omega}\in\mathbb C$。
但尚不知定义~\ref{def:small-quantum-product} 中的求和是否收敛。
避免这一问题的一种方法，是将 $q^\beta$ 视为形式变量，而非数值。
有几种方式可使这一点精确化。

对于 Fano 流形，由命题~\ref{prop:quantum-finite-fano}，定义~\ref{def:small-quantum-product} 中的和有限。
此时，可取系数环为半群环
\[
R=\mathbb Q[q^\beta:\beta\in H_2(X,\mathbb Z)].
\]
也就是说，每个 $\beta\in H_2(X,\mathbb Z)$ 给出形式符号 $q^\beta$，满足
$q^{\beta_1}q^{\beta_2}=q^{\beta_1+\beta_2}$；$R$ 的元素是各 $q^\beta$ 的有限有理系数和。
例如，对 $\mathbb P^r$，得到关于 $q$ 的 Laurent 多项式环 $R=\mathbb Q[q,q^{-1}]$。
对任意 Fano 流形，小量子乘积定义在 $H^*(X,R)$ 上。

若 $X$ 为任意射影流形，这种办法可能失效，因为定义~\ref{def:small-quantum-product} 中的和可能无限。
但由\IndexLocation{indexloc:0500:251}有效性公理，可假设 $\beta$ 有效。
更准确地说，可假设 $\beta$ 位于整 Mori 锥 $M(X)_{\mathbb Z}$ 中。
考虑
\[
\mathcal R=\mathbb Q[[q^\beta:\beta\in M(X)_{\mathbb Z}]],
\]
其中双括号表示允许无限和。两个这样的元素的乘积为
\[
\left(\sum_{\beta\in M(X)_{\mathbb Z}}a_\beta q^\beta\right)
\left(\sum_{\beta\in M(X)_{\mathbb Z}}b_\beta q^\beta\right)
=\sum_{\beta\in M(X)_{\mathbb Z}}
\left(\sum_{\beta=\beta_1+\beta_2}a_{\beta_1}b_{\beta_2}\right)q^\beta.
\]
这个和有意义，因为对任意 $\beta\in M(X)_{\mathbb Z}$，括号内的和有限：
Mori 锥的自由部分由强凸锥中的格点构成。
因此 $\mathcal R$ 是环，可视为整 Mori 锥的形式半群环。
以五次三维簇为例，$\mathcal R=\mathbb Q[[q]]$ 是关于 $q$ 的形式幂级数环。
一般地，容易看出小量子乘积定义在 $H^*(X,\mathcal R)$ 上。

最后，若 $(M,\omega)$ 为辛流形，可用定义~\ref{def:small-quantum-product} 定义小量子乘积，
但必须使用以下系数环的上同调：
\SourcePage{252}
\emph{Novikov 环} $\Lambda(\omega,\mathbb Q)$。
它由所有形式和 $\sum_{\beta\in H_2(M,\mathbb Z)}a_\beta q^\beta$，$a_\beta\in\mathbb Q$ 构成，
满足对任意 $C\in\mathbb R$，集合
\[
\left\{\beta:a_\beta\ne0\text{ 且 }\int_\beta\omega<C\right\}
\]
有限。直接可证，$\Lambda(\omega,\mathbb Q)$ 为环，
且小量子乘积定义在 $H^*(M,\Lambda(\omega,\mathbb Q))$ 上。

也可使用其他系数。在代数情形，可用 $\mathcal R$ 在适当自同构群下的不变子环代替 $\mathcal R$，见 \SourceCite{Morrison7}。
在辛情形，往往将 Novikov 环限制到球面同调类，见 \SourceCite{MS}，另见 \SourceCite{Getzler1}。

\section{大量子上同调}\label{sec:big-quantum-cohomology}
\index{量子上同调!大量子上同调|hyperpage}
\hyperref[sec:small-quantum-cohomology]{第~\ref*{sec:small-quantum-cohomology} 节}
讨论的量子上同调，并未用尽 Gromov--Witten 不变量 $\langle I_{0,n,\beta}\rangle$ 提供的枚举信息。
原因在于，小量子乘积由 $\langle I_{0,3,\beta}\rangle$ 定义，而更有趣的不变量常在 $n\geq3$ 时出现。
例如，\hyperref[chap:gromov-witten]{第~\ref*{chap:gromov-witten} 章}已看到，对 $\mathbb P^2$，数
\[
N_d=\langle I_{0,3d-1,d}\rangle(\underbrace{[pt],\ldots,[pt]}_{3d-1\text{ 次}})
\]
有一些显著性质。但 $\mathbb P^2$ 的小量子乘积只用到
$N_1=\langle I_{0,3,1}\rangle([pt],[pt],[\ell])=1$。
本节探讨一种称为\emph{大量子上同调}的量子上同调版本，它考虑全部 Gromov--Witten 不变量。

\subsection{作为超流形的 $H^*(X,\mathbb C)$}\label{subsec:cohomology-supermanifold}
\index{超流形|hyperpage}\index{超交换|hyperpage}
定义大量子乘积之前，需讨论 $H^*(X,\mathbb C)$ 的超流形结构。
固定有理上同调 $H^*(X,\mathbb Q)$ 的齐次基 $T_0=1,T_1,\ldots,T_m$。
对每个 $T_i$，引入与它次数相同的变量 $t_i$；这些变量是超流形 $H^*(X,\mathbb C)$ 的局部坐标。
$t_i$ 超交换，即 $t_it_j=(-1)^{\deg t_i\deg t_j}t_jt_i$；若 $t_i$ 的次数为奇数，则 $t_i^2=0$。
所以环 $\mathbb C[t_0,\ldots,t_m]$ 是一个关于偶变量的多项式环，与一个关于奇变量的外代数的张量积。
对不含 $t_i$ 的任意单项式 $t^\alpha$，定义偏导数算子 $\partial/\partial t_i$ 为
\[
\frac\partial{\partial t_i}(t_i^k\cdot t^\alpha)=kt_i^{k-1}\cdot t^\alpha.
\]
于是可以计算任意 $F\in\mathbb C[t_0,\ldots,t_m]$ 的 $\partial F/\partial t_i$，并且
\[
\frac{\partial^2F}{\partial t_i\partial t_j}
=(-1)^{\deg t_i\deg t_j}\frac{\partial^2F}{\partial t_j\partial t_i}.
\]
这些都容易扩展到形式幂级数环 $\mathbb C[[t_0,\ldots,t_m]]$。
\SourcePage{253}
关于超流形的更多内容，请参阅 \SourceCite{Manin3}。
如 \SourceCite[第 2 讲]{Alvarez} 所讨论，研究费米子系统时，超交换变量自然出现在量子物理中。

\subsection{Gromov--Witten 势}\label{subsec:gw-potential}
\index{Gromov--Witten势@Gromov--Witten 势|hyperpage}
定义大量子乘积的关键工具是 \emph{Gromov--Witten 势}。
如上，固定有理上同调 $H^*(X,\mathbb Q)$ 的基 $T_0=1,T_1,\ldots,T_m$，
置 $\gamma=\sum_{i=0}^mt_iT_i$，其中 $t_i$ 为超交换变量，$\deg t_i=\deg T_i$。
\begin{definition}\label{def:gw-potential}
设 $\omega$ 是光滑射影簇 $X$ 上的复化 Kähler 类。
则 \emph{Gromov--Witten 势}为形式和
\[
\Phi(\gamma)=\sum_{n=0}^{\infty}\sum_{\beta\in H_2(X,\mathbb Z)}
\frac1{n!}\langle I_{0,n,\beta}\rangle(\gamma^n)q^\beta,
\qquad q^\beta=e^{2\pi i\int_\beta\omega}.
\]
\end{definition}
在此和中，$\beta=0$ 时隐含 $n\geq3$，因为 $n\leq2$ 时 $\overline M_{0,n}(X,0)$ 不存在。
文献中 Gromov--Witten 势的定义略有差异：一些作者对\emph{所有} $\beta$ 都要求 $n\geq3$，从而截断级数。
当希望将 Gromov--Witten 势与 Gromov--Witten 类联系起来时，这一做法尤其合适。

为使定义有意义，需解释 $\langle I_{0,n,\beta}\rangle(\gamma^n)$ 的含义。
由于超交换性，这一点稍微微妙。思路是展开 $\gamma^n=(\sum_{i=0}^mt_iT_i)^n$，
将 $T_i$ 与 $t_i$ 都视为超交换变量。
采用多重指标记号：指数向量 $\alpha=(a_0,\ldots,a_m)$ 给出总次数
$|\alpha|=a_0+\cdots+a_m$ 的单项式 $T^\alpha,t^\alpha$。
还置 $\alpha!=a_0!\cdots a_m!$。
一个有用的观察是，$\gamma$ 的每项 $t_iT_i$ 都为偶，因此可将 $t_iT_i$ 作为整体交换。
于是，Gromov--Witten 不变量的符号规则给出
\begin{equation}\label{eq:gw-potential-super-expansion}
\frac1{n!}\langle I_{0,n,\beta}\rangle(\gamma^n)
=\sum_{|\alpha|=n}\epsilon(\alpha)\langle I_{0,n,\beta}\rangle(T^\alpha)\frac{t^\alpha}{\alpha!}.
\end{equation}
因子 $\epsilon(\alpha)=\pm1$ 是由等式
\begin{equation}\label{eq:gw-potential-super-sign}
(t_0T_0)^{a_0}\cdots(t_mT_m)^{a_m}
=\epsilon(\alpha)T_0^{a_0}\cdots T_m^{a_m}t_0^{a_0}\cdots t_m^{a_m}
\end{equation}
确定的符号。

将式~\eqref{eq:gw-potential-super-expansion} 代入 $\Phi$ 的定义，
希望将 $\Phi$ 视为\hyperref[subsec:cohomology-supermanifold]{第~\ref*{subsec:cohomology-supermanifold} 小节}
讨论的环 $\mathbb C[[t_0,\ldots,t_m]]$ 中的形式幂级数。
不过，小量子上同调的收敛问题在这里仍然存在。
为求严格，可能需要用\hyperref[subsec:formal-small-quantum]{第~\ref*{subsec:formal-small-quantum} 小节}
的环 $R,\mathcal R$ 或 $\Lambda(\omega,\mathbb Q)$ 代替 $\mathbb C$。
所以，在下面的讨论中，用 $\mathcal C$ 按需要表示
\[
\mathbb C[[t_0,\ldots,t_m]],\quad R[[t_0,\ldots,t_m]],\quad
\mathcal R[[t_0,\ldots,t_m]],\quad\Lambda(\omega,\mathbb Q)[[t_0,\ldots,t_m]]
\]
中的一个环。

由于 Gromov--Witten 势 $\Phi$ 包含全部可能的亏格零不变量，
可将它视为树级 Gromov--Witten 不变量的生成函数。
考虑到阶乘与不交换变量，更准确地说，$\Phi$ 是
\SourcePage{254}
$\langle I_{0,n,\beta}\rangle$ 的\emph{超交换指数型生成函数}。
从物理角度看，$\Phi$ 自然解释为亏格零自由能。

Gromov--Witten 势有时称为亏格零 Gromov--Witten 势。
\emph{亏格 $g$ 的 Gromov--Witten 势}可用亏格 $g$ 的不变量，以类似方式定义。

椭圆曲线 $E$ 给出一个简单的 Gromov--Witten 势例子。
由于不存在非恒定映射 $\mathbb P^1\to E$，只需考虑由点映射公理确定的
Gromov--Witten 不变量 $\langle I_{0,3,0}\rangle$。
取 $T_0=1\in H^0(E,\mathbb C)$、$T_3=[pt]\in H^2(E,\mathbb C)$，
再取 $T_1,T_2\in H^1(E,\mathbb C)$，满足 $\int_ET_1\cup T_2=1$。
于是计算得到
\begin{equation}\label{eq:elliptic-gw-potential}
\Phi=\tfrac12t_0^2t_3-t_0t_1t_2.
\end{equation}
负号来自式~\eqref{eq:gw-potential-super-expansion} 中的 $\epsilon(\alpha)$。
\hyperref[sec:computing-gw-ii]{第~\ref*{sec:computing-gw-ii} 节}将讨论一些不那么平凡的 Gromov--Witten 势例子。

利用 Gromov--Witten 势，如下定义大量子乘积。
\begin{definition}\label{def:big-quantum-product}
设 $\Phi$ 是光滑射影簇 $X$ 的 Gromov--Witten 势。定义
\[
T_i*T_j=\sum_k\frac{\partial^3\Phi}{\partial t_i\partial t_j\partial t_k}T^k,
\]
其中 $g^{kl}$ 如定义~\ref{def:small-quantum-product}，$T^k=\sum_lg^{kl}T_l$。
线性扩张后，得到上同调 $H^*(X,\mathcal C)$ 上的\emph{大量子乘积}。
\end{definition}
\index{量子乘积!大量子乘积|hyperpage}
注意，$n<3$ 的 Gromov--Witten 不变量是 $\Phi$ 中关于 $t_i$ 次数至多为二的项的系数，
所以在 Gromov--Witten 势的定义中保留或删去它们，并不影响大量子乘积的定义。

在椭圆曲线例子~\eqref{eq:elliptic-gw-potential} 中，$\Phi=\tfrac12t_0^2t_3-t_0t_1t_2$。
使用前面的上同调基 $T_0,T_1,T_2,T_3$，可计算量子乘积
\[
T_1*T_2=\frac{\partial^3\Phi}{\partial t_1\partial t_2\partial t_0}T_3
=\frac{\partial^3}{\partial t_1\partial t_2\partial t_0}
\left(\tfrac12t_0^2t_3-t_0t_1t_2\right)T_3=T_3,
\]
因为 $t_1,t_2$ 为奇变量。容易看出，这个例子中的大量子乘积与通常的杯积一致。
下面的引理有助于处理 Gromov--Witten 势的\IndexLocation{indexloc:0567:254}三阶偏导数。
\begin{lemma}\label{lem:gw-potential-third-derivatives}
对所有 $i,j,k$，
\[
\frac{\partial^3\Phi}{\partial t_i\partial t_j\partial t_k}
=\sum_{n=0}^{\infty}\sum_{\beta\in H_2(X,\mathbb Z)}
\frac1{n!}\langle I_{0,n+3,\beta}\rangle(T_i,T_j,T_k,\gamma^n)q^\beta.
\]
\end{lemma}
\begin{proof}
设指数向量 $\alpha=(a_0,\ldots,a_m)$ 满足 $|\alpha|=n$，
且 $\eta=\alpha-e_i-e_j-e_k$ 的各分量非负，$e_i,e_j,e_k$ 为标准基向量。
若能证明
\[
\frac{\partial^3}{\partial t_i\partial t_j\partial t_k}
\left(\epsilon(\alpha)\langle I_{0,n,\beta}\rangle(T^\alpha)\frac{t^\alpha}{\alpha!}\right)
=\epsilon(\eta)\langle I_{0,n,\beta}\rangle(T_i,T_j,T_k,T^\eta)\frac{t^\eta}{\eta!},
\]
\SourcePage{255}
则引理由式~\eqref{eq:gw-potential-super-expansion} 得到。
先设 $t_i,t_j,t_k$ 都为奇变量。可假设 $a_i=a_j=a_k=1$。
利用 $\langle I_{0,n,\beta}\rangle$ 与 $t_i$ 的符号规则，可写成
\[
\epsilon(\alpha)\langle I_{0,n,\beta}\rangle(T^\alpha)\frac{t^\alpha}{\alpha!}
=\epsilon(\eta+e_i+e_j+e_k)\langle I_{0,n,\beta}\rangle(T_k,T_j,T_i,T^\eta)
\,t_kt_jt_i\frac{t^\eta}{\eta!}.
\]
求三次偏导数，得到
\[
\frac{\partial^3}{\partial t_i\partial t_j\partial t_k}
\left(\epsilon(\alpha)\langle I_{0,n,\beta}\rangle(T^\alpha)\frac{t^\alpha}{\alpha!}\right)
=\epsilon(\eta+e_i+e_j+e_k)\langle I_{0,n,\beta}\rangle(T_k,T_j,T_i,T^\eta)\frac{t^\eta}{\eta!}.
\]
因此只需证明
\[
\epsilon(\eta+e_i+e_j+e_k)\langle I_{0,n,\beta}\rangle(T_k,T_j,T_i,T^\eta)
=\epsilon(\eta)\langle I_{0,n,\beta}\rangle(T_i,T_j,T_k,T^\eta).
\]
由式~\eqref{eq:gw-potential-super-sign}，$\epsilon(\eta)=(-1)^{s(s+1)/2}$，
其中 $s$ 是 $t^\eta$ 中奇变量的个数（可假设奇变量只出现一次幂）。
又由于 $t_i,t_j,t_k$ 为奇变量，式~\eqref{eq:gw-top-degree-condition} 给出 $s+3\equiv0\pmod2$。
所需等式随即成立。其他情形类似，留给读者。
\end{proof}

\subsection{WDVV 方程与结合律}\label{subsec:wdvv-associativity}
\index{WDVV方程@WDVV 方程|hyperpage}
下一个目标，是证明 $H^*(X,\mathcal C)$ 在大量子乘积下成为环。
与\hyperref[sec:small-quantum-cohomology]{第~\ref*{sec:small-quantum-cohomology} 节}一样，难点是结合律。
但这里出现了一个显著现象：大量子上同调的结合律等价于 $\Phi$ 满足的某个偏微分方程，
即所谓 WDVV（Witten--Dijkgraaf--Verlinde--Verlinde）方程，见 \SourceCite{Witten2,DVV1,DVV2}。
\begin{theorem}\label{thm:wdvv}
Gromov--Witten 势 $\Phi$ 对所有 $i,j,k,l$ 满足方程
\[
\begin{aligned}
\sum_{a,b}\frac{\partial^3\Phi}{\partial t_i\partial t_j\partial t_a}g^{ab}
\frac{\partial^3\Phi}{\partial t_b\partial t_k\partial t_l}
={}&(-1)^{\deg t_i(\deg t_j+\deg t_k)}\\
&\cdot\sum_{a,b}\frac{\partial^3\Phi}{\partial t_j\partial t_k\partial t_a}g^{ab}
\frac{\partial^3\Phi}{\partial t_b\partial t_i\partial t_l}.
\end{aligned}
\]
\end{theorem}
\begin{proof}
这里只概述主要思想，详细证明见 \SourceCite{KoM1}。
为简便，再假设 $H^*(X,\mathbb C)$ 只有偶数次上同调，这样就不必处理符号。

证明采用与 $\mathbb P^2$ 的 Gromov--Witten 不变量递推式~\eqref{eq:plane-recursion}
相同的策略。将式~\eqref{eq:plane-xi} 的类 $\xi$ 换成 Gromov--Witten 类
\[
I_{0,n+4,\beta}(\gamma^n,T_i,T_j,T_k,T_l)\in H^*(\overline M_{0,n+4},\mathcal C),
\]
其中 $\gamma=\sum_{i=0}^mt_iT_i$。
将在两个线性等价除子上对此类求值，下面描述它们。
用集合 $\mathcal I=\{1,\ldots,n,i,j,k,l\}$ 标记 $I_{0,n+4,\beta}$ 的 $n+4$ 个自变量
（这是对记号的轻微滥用）。
与式~\eqref{eq:plane-boundary-relation} 前的讨论一样，
将 $\mathcal I$ 分成不相交的 $A,B$，就得到一个除子
$\operatorname{Im}(\varphi_{A,B})\subset\overline M_{0,n+4}$。
\SourcePage{256}
定义
\[
D(i,j\mid k,l)=\sum_{\substack{i,j\in A\\k,l\in B}}\operatorname{Im}(\varphi_{A,B}),
\quad
D(j,k\mid i,l)=\sum_{\substack{j,k\in A\\i,l\in B}}\operatorname{Im}(\varphi_{A,B}).
\]
类似式~\eqref{eq:plane-boundary-relation} 的做法，可得到线性等价关系
\begin{equation}\label{eq:wdvv-boundary-equivalence}
D(i,j\mid k,l)\sim D(j,k\mid i,l).
\end{equation}
现考虑 WDVV 方程左边的和。利用 $T^k=\sum_lg^{kl}T_l$ 与引理~\ref{lem:gw-potential-third-derivatives}，可写为
\[
\sum\frac1{n_1!n_2!}\langle I_{0,n_1+3,\beta_1}\rangle(T_i,T_j,T_a,\gamma^{n_1})
g^{ab}\langle I_{0,n_2+3,\beta_2}\rangle(T_b,T_k,T_l,\gamma^{n_2})q^{\beta_1+\beta_2},
\]
其中求和遍历 $n_1,n_2,\beta_1,\beta_2,a,b$。
将此和分成较小的和，使每个小和中的 $n_1+n_2=n$ 与 $\beta_1+\beta_2=\beta$ 固定。
由交换性，对应于固定 $n,\beta$ 的小和可写成
\[
\sum\frac1{n_1!n_2!}\langle I_{0,n_1+3,\beta_1}\rangle(\gamma^{n_1},T_i,T_j,T_a)
g^{ab}\langle I_{0,n_2+3,\beta_2}\rangle(T_b,\gamma^{n_2},T_k,T_l)q^\beta,
\]
其中现在对满足 $n_1+n_2=n$、$\beta_1+\beta_2=\beta$ 的 $n_1,n_2,\beta_1,\beta_2,a,b$ 求和。
利用式~\eqref{eq:gw-potential-super-expansion} 展开并应用分裂公理，可证明此小和化为
\[
\int_{D(i,j\mid k,l)}\frac1{n!}I_{0,n+4,\beta}(\gamma^n,T_i,T_j,T_k,T_l)q^\beta,
\]
其中 $D(i,j\mid k,l)$ 是上面定义的除子。因此，WDVV 方程的左边为
\[
\sum_{n=0}^{\infty}\sum_{\beta\in H_2(X,\mathbb Z)}
\int_{D(i,j\mid k,l)}\frac1{n!}I_{0,n+4,\beta}(\gamma^n,T_i,T_j,T_k,T_l)q^\beta.
\]
对方程右边进行同样的处理，得到
\[
\sum_{n=0}^{\infty}\sum_{\beta\in H_2(X,\mathbb Z)}
\int_{D(j,k\mid i,l)}\frac1{n!}I_{0,n+4,\beta}(\gamma^n,T_i,T_j,T_k,T_l)q^\beta.
\]
由线性等价关系~\eqref{eq:wdvv-boundary-equivalence}，两者相等，定理立即成立。
\end{proof}
现在容易证明，大量子上同调构成环。
\begin{theorem}\label{thm:big-quantum-ring}
$H^*(X,\mathcal C)$ 在大量子乘积下是单位元为 $T_0=1$ 的交换环。
\end{theorem}
\begin{proof}
由\hyperref[sec:gw-properties]{第~\ref*{sec:gw-properties} 节}的\IndexLocation{indexloc:0501:256}等变性公理，大量子乘积超交换。
下面证明，$T_0=1$ 不仅是杯积的单位元，也是
\SourcePage{257}
大量子乘积的单位元。这很容易：在引理~\ref{lem:gw-potential-third-derivatives} 中取 $T_0=[X]$，得到
\begin{equation}\label{eq:gw-potential-unit-derivative}
\begin{aligned}
\frac{\partial^3\Phi}{\partial t_0\partial t_j\partial t_k}
&=\sum_{n=0}^{\infty}\sum_{\beta\in H_2(X,\mathbb Z)}
\frac1{n!}\langle I_{0,n+3,\beta}\rangle(T_0,T_j,T_k,\gamma^n)q^\beta\\
&=\langle I_{0,3,0}\rangle(T_0,T_j,T_k)=\int_XT_j\cup T_k=g_{jk}.
\end{aligned}
\end{equation}
第二、第三个等式分别由基本类公理与点映射公理得到。
此式给出 $T_0*T_j=T_j$，故 $T_0$ 为单位元。

最后证明结合律。利用 $T^k=\sum_lg^{kl}T_l$ 与大量子乘积的定义，得到
\[
(T_i*T_j)*T_k=\sum_{a,b}\sum_{l,s}
\frac{\partial^3\Phi}{\partial t_i\partial t_j\partial t_a}g^{ab}
\frac{\partial^3\Phi}{\partial t_b\partial t_k\partial t_l}g^{ls}T_s.
\]
以同样方式展开 $(T_j*T_k)*T_i$，并应用定理~\ref{thm:wdvv} 的 WDVV 方程，立即得到
\[
(T_i*T_j)*T_k=(-1)^{\deg T_i(\deg T_j+\deg T_k)}(T_j*T_k)*T_i.
\]
由于 $*$ 超交换，结合律随之成立。
\end{proof}
此定理说明，WDVV 方程推出大量子乘积的结合律。
反过来也成立，所以结合律等价于 WDVV 方程。
这是\hyperref[sec:dubrovin-formalism]{第~\ref*{sec:dubrovin-formalism} 节}将讨论的 Dubrovin 形式体系的一部分。

还应提到，大量子上同调与交乘积 $g(a,b)=\int_Xa\cup b$ 配合良好。
类似命题~\ref{prop:quantum-frobenius-pairing}，可以证明大量子乘积满足
\begin{equation}\label{eq:big-quantum-frobenius}
g(a*b,c)=g(a,b*c),\qquad a,b,c\in H^*(X,\mathcal C).
\end{equation}
所以 $H^*(X,\mathcal C)$ 在大量子乘积下是 Frobenius 代数，并有公式
\[
\int_XT_i*T_j*T_k=\frac{\partial^3\Phi}{\partial t_i\partial t_j\partial t_k}.
\]
大量子乘积有 Künneth 公式，见 \SourceCite{KoM2,Kaufmann}。
此外，\hyperref[subsec:gravitational-correlator-definition]{第~\ref*{subsec:gravitational-correlator-definition} 小节}将讨论一种称为\emph{引力量子乘积}的大量子乘积推广。

\section{计算 Gromov--Witten 不变量，II}\label{sec:computing-gw-ii}
现在重访\hyperref[sec:computing-gw-i]{第~\ref*{sec:computing-gw-i} 节}的若干例子，
看看它们与大量子上同调及 Gromov--Witten 势有何关系。
\subsection{Gromov--Witten 势的结构}\label{subsec:gw-potential-structure}
\index{Gromov--Witten势@Gromov--Witten 势!齐次性|hyperpage}\index{Gromov--Witten势@Gromov--Witten 势!经典部分和量子部分|hyperpage}
如下将看到，Gromov--Witten 势的显式公式可能相当复杂。
不过，$\Phi$ 具有一种结构，能使公式更容易理解。
\SourcePage{258}
首先说明：给各变量赋予适当的次数后，$\Phi$ 是齐次的。沿用上一节的记号，我们有上同调基 $T_i$ 与超交换变量 $t_i$。由式~\eqref{eq:gw-potential-super-expansion}，$\Phi$ 的一个典型项形如
\[
\langle I_{0,n,\beta}\rangle(T_0^{a_0},\ldots,T_m^{a_m})
\frac{t_0^{a_0}}{a_0!}\cdots\frac{t_m^{a_m}}{a_m!}q^\beta,
\]
其中 $a_0+\cdots+a_m=n$。另一方面，次数条件~\eqref{eq:gw-top-degree-condition} 给出
\[
a_0\deg T_0+\cdots+a_m\deg T_m
=2\dim X-2\int_\beta\omega_X+2(3\cdot0-3+n).
\]
利用 $a_0+\cdots+a_m=n$，可将它改写为
\[
a_0(\deg T_0-2)+\cdots+a_m(\deg T_m-2)+2\int_\beta\omega_X=2\dim X-6.
\]
因此，只要给 $t_0,\ldots,t_m$ 和 $q^\beta$ 赋予如下次数，Gromov--Witten 势 $\Phi$ 就是次数为 $2\dim X-6$ 的齐次式：
\begin{equation}\label{eq:gw-potential-grading}
\deg t_i=\deg T_i-2,\qquad \deg q^\beta=2\int_\beta\omega_X.
\end{equation}
也就是说，$t_i$ 的次数要从先前的次数中减去 $2$，因而奇偶性不变。至于 $q^\beta$，要这样赋予次数，就必须把它们看作形式符号，例如\hyperref[subsec:formal-small-quantum]{第~\ref*{subsec:formal-small-quantum} 小节}中的环 $R$、$\mathcal R$ 或 $\Lambda(\omega,\mathbb Q)$ 中的符号，而不能仅看作数值。

Gromov--Witten 势的齐次性在实际计算中十分有用。它说明，势中的许多项因次数不符而为零。应用于\IndexLocation{indexloc:1116:258}量子上同调时，齐次性还说明，$T_i*T_j$ 中 $T_\ell q^\beta$ 的系数具有次数 $\deg T_\ell-(\deg T_i+\deg T_j+2\int_\beta\omega_X)$。正如我们在命题~\ref{prop:small-quantum-degree} 的证明中对小量子乘积所注意到的那样。（不过，如果只关心量子乘积，通常直接使用次数条件比使用齐次性更方便。）

作为一个简单的例子，考虑式~\eqref{eq:elliptic-gw-potential} 中椭圆曲线的 Gromov--Witten 势 $\Phi=\tfrac12t_0^2t_3-t_0t_1t_2$。只要 $t_0,t_1,t_2,t_3$ 的次数分别为 $-2,-1,-1,0$，它就是次数为 $-4$ 的齐次式。

$\Phi$ 还有一个有用的性质：由\IndexLocation{indexloc:0505:258}点映射公理，对应于 $\beta=0$ 的各项化为 $\tfrac16\langle I_{0,3,0}\rangle(\gamma^3)=\tfrac16\int_X\gamma^3$。因此，可将 Gromov--Witten 势分解为“经典”项与“量子”项（这里采用 \SourceCite{FP} 的术语）：
\[
\begin{aligned}
\Phi&=\Phi_{\mathrm{classical}}+\Phi_{\mathrm{quantum}}\\
&=\frac16\int_X\gamma^3+
\sum_n\sum_{\beta\in H_2(X,\mathbb Z)\setminus\{0\}}
\frac1{n!}\langle I_{0,n,\beta}\rangle(\gamma^n)q^\beta.
\end{aligned}
\]
注意，由\IndexLocation{indexloc:0502:258}基本类公理，变量 $t_0$ 只出现在 $\Phi$ 的\IndexLocation{indexloc:0553:258}经典部分 $\Phi_{\mathrm{classical}}$ 中。

$\Phi$ 的另一个好性质是：由\IndexLocation{indexloc:0499:258}除子公理，对应于除子的变量在 $\Phi_{\mathrm{quantum}}$ 中以指数形式出现。为说明这一点，将基元素 $T_i$ 重新排列，使得 $T_1,\ldots,T_r$ 是 $H^2(X,\mathbb Q)$ 的一组基，并写成
\SourcePage{259}
$\gamma=t_0T_0+\delta+\epsilon$，其中 $\delta=\sum_{i=1}^rt_iT_i$，$\epsilon=\sum_{i=r+1}^mt_iT_i$。随后，利用二项式定理和除子公理，得到
\begin{equation}\label{eq:gw-potential-divisor-exponential}
\begin{aligned}
\Phi&=\frac16\int_X\gamma^3+\sum_n\sum_{\beta\ne0}
\frac1{n!}\langle I_{0,n,\beta}\rangle\bigl((\delta+\epsilon)^n\bigr)q^\beta\\
&=\frac16\int_X\gamma^3+\sum_{a,b}\sum_{\beta\ne0}
\frac1{a!b!}\langle I_{0,b,\beta}\rangle(\epsilon^b)
\left(\int_\beta\delta\right)^a q^\beta\\
&=\frac16\int_X\gamma^3+\sum_b\sum_{\beta\ne0}
\frac1{b!}\langle I_{0,b,\beta}\rangle(\epsilon^b)e^{\int_\beta\delta}q^\beta.
\end{aligned}
\end{equation}
\SourceCite{KoM1} 的命题 4.4 讨论了此公式；它对量子上同调还有一些推论，将在\hyperref[subsec:quantum-product-dependence]{第~\ref*{subsec:quantum-product-dependence} 小节}讨论。

\subsection{重访 $\mathbb P^2$}\label{subsec:plane-gw-potential}
\index{Gromov--Witten势@Gromov--Witten 势!射影平面|hyperpage}
对 $\mathbb P^2$，首先计算 Gromov--Witten 势。上同调基为 $T_0=[\mathbb P^2]=1$、$T_1=[\ell]$ 和 $T_2=[pt]$，因此，由式~\eqref{eq:gw-potential-grading}，$t_0,t_1,t_2$ 的次数为 $-2,0,2$，而 $q=q^\ell$ 的次数为 $-6$。照例将 Gromov--Witten 不变量写为 $\langle I_{0,n,d}\rangle$。用此记号，$N_d=\langle I_{0,3d-1,d}\rangle(T_2^{3d-1})$ 是通过 $\mathbb P^2$ 中 $3d-1$ 个一般位置点的 $d$ 次有理曲线数目。

我们知道，$\Phi$ 是次数为 $-2$ 的齐次式，其经典部分为 $\tfrac12(t_0^2t_2+t_0t_1^2)$。量子部分的一个典型项为
\[
\langle I_{0,a+b,d}\rangle(T_1^a,T_2^b)\frac{t_1^a}{a!}\frac{t_2^b}{b!}q^d,
\]
由齐次性，$2b+(-6)d=-2$。这给出 $b=3d-1$，所以该项可写成
\[
\langle I_{0,3d-1,d}\rangle(T_2^{3d-1})
\frac{(dt_1)^a}{a!}\frac{t_2^{3d-1}}{(3d-1)!}q^d
=N_d\frac{(dt_1)^a}{a!}\frac{t_2^{3d-1}}{(3d-1)!}q^d,
\]
其中因子 $d^a$ 来自除子公理。（这重新得到了\hyperref[subsec:plane-genus-zero]{第~\ref*{subsec:plane-genus-zero} 小节}中的一个计算，不过这里的推导更快。）因此，$\mathbb P^2$ 的 Gromov--Witten 势为
\[
\Phi=\frac12(t_0^2t_2+t_0t_1^2)+\sum_{d=1}^{\infty}
N_de^{dt_1}\frac{t_2^{3d-1}}{(3d-1)!}q^d.
\]
现在可以做些有趣的计算。对指标 $(i,j,k,l)=(1,1,2,2)$，定理~\ref{thm:wdvv} 的 WDVV 方程成为
\begin{equation}\label{eq:plane-wdvv}
\Phi_{222}+\Phi_{111}\Phi_{122}=\Phi_{112}^2,
\end{equation}
其中下标表示偏导数。因为式~\eqref{eq:gw-potential-unit-derivative} 给出 $\Phi_{0jk}=g_{jk}$，所以上式容易算出。为利用式~\eqref{eq:plane-wdvv}，注意，$N_d$ 在 $\Phi_{222}$ 中出现在如下项里：
\[
N_d\frac1{(3d-4)!}t_2^{3d-4}q^d.
\]
这提示我们考察 $\Phi_{112}^2$ 和 $\Phi_{111}\Phi_{122}$ 中含 $t_2^{3d-4}q^d$ 的项。在 $\Phi_{112}$ 中，不含 $t_1$ 的项为
\[
N_{d_1}d_1^2\frac1{(3d_1-2)!}t_2^{3d_1-2}q^{d_1},
\]
\SourcePage{260}
所以 $\Phi_{112}^2$ 中含 $t_2^{3d-4}q^d$ 的项为
\[
\sum_{d_1+d_2=d}N_{d_1}N_{d_2}d_1^2d_2^2
\frac1{(3d_1-2)!}\frac1{(3d_2-2)!}t_2^{3d-4}q^d.
\]
类似地，$\Phi_{111}\Phi_{122}$ 中含 $t_2^{3d-4}q^d$ 的项为
\[
\sum_{d_1+d_2=d}N_{d_1}N_{d_2}d_1^3d_2
\frac1{(3d_1-1)!}\frac1{(3d_2-3)!}t_2^{3d-4}q^d.
\]
这一计算留给读者验证。由此，式~\eqref{eq:plane-wdvv} 立即给出递推关系
\[
N_d=\sum_{d_1+d_2=d}N_{d_1}N_{d_2}
\left(d_1^2d_2^2\binom{3d-4}{3d_1-2}
-d_1^3d_2\binom{3d-4}{3d_1-1}\right).
\]

$\mathbb P^2$ 的大量子上同调也容易确定。例如，由于 $\Phi_{011}=1$，有
\[
T_1*T_1=T_2+\Phi_{111}T_1+\Phi_{112}T_0.
\]
另一方面，$\mathbb P^2$ 的小量子乘积满足 $T_1*T_1=T_2$，可见两种量子上同调有明显差别。还要注意，大量子乘积的结合律与\hyperref[sec:dubrovin-formalism]{第~\ref*{sec:dubrovin-formalism} 节}讨论的 WDVV 方程实质上等价。实际使用中，结合律往往更方便。例如，展开 $(T_1*T_1)*T_2=T_1*(T_1*T_2)$，并比较 $T_0$ 的系数，就直接得到式~\eqref{eq:plane-wdvv}。因此，$N_d$ 的递推关系是大量子乘积结合律的直接推论。正如 \SourceCite{FP} 所说：“量子形式体系使人不再非得想出巧妙的办法不可。只需写出结合律方程，便可从中读出计数信息。”\SourceCite{DI} 讨论了这样做的一种系统方法；对 $\mathbb P^3$ 和三维二次超曲面的应用见 \SourceCite{FP}。\SourceCite{Kresch} 还更深入地研究了结合律方程及其与 Kontsevich 第一重构定理（定理~\ref{thm:gw-reconstruction}）的关系。

对 $\mathbb P^2$，还可以将 $N_d$ 推广为 $N_d(a,b,c)$：它计数通过 $a$ 个一般位置点、与 $b$ 条一般位置直线相切，并在另外 $c$ 条一般位置直线上各自指定的一般位置点处与这些直线相切的 $d$ 次有理曲线；这里 $a+b+2c=3d-1$。\SourceCite{EK1} 研究了这些不变量，并导出了 $N_d(a,b,c)$ 的递推关系，推广了 \SourceCite{DI} 给出的 $N_d(a,b,0)$ 的递推关系。此外，可以用 $N_d(a,b,c)$ 定义 $\mathbb P^2$ 大量子上同调的一种推广，称为\emph{接触上同调环} \SourceCite{EK2}。\index{接触上同调环|hyperpage}

\subsection{重访五次三维簇}\label{subsec:quintic-gw-potential}
\index{Gromov--Witten势@Gromov--Witten 势!五次三维簇|hyperpage}
首先研究五次三维簇 $V$ 的 Gromov--Witten 势。$H^*(V,\mathbb C)$ 的一组上同调基为 $T_0=[V]$、$T_1=[H]$、$T_2=[\ell]$、$T_3=[pt]$ 与 $\gamma_1,\ldots,\gamma_{204}$，其中 $\gamma_i$ 是 $H^3(V,\mathbb C)$ 的一组基。对应的超交换变量记为 $t_0,t_1,t_2,t_3$ 和 $u_1,\ldots,u_{204}$。再令 $q_{ij}=\int_V\gamma_i\cup\gamma_j$。于是，$\Phi$ 的经典部分为
\[
\Phi_{\mathrm{classical}}=\frac12t_0^2t_3+t_0t_1t_2+
\frac56t_1^3-\sum_{i<j}q_{ij}t_0u_iu_j,
\]
其中负号来自式~\eqref{eq:gw-potential-super-expansion} 中的 $\epsilon(\alpha)$。
\SourcePage{261}
再看量子部分。我们知道，$\Phi$ 和 $q^\beta$ 的次数均为 $0$，而 $t_0$ 不出现在 $\Phi_{\mathrm{quantum}}$ 中。由式~\eqref{eq:gw-potential-grading}，$t_2,t_3$ 和 $u_i$ 的次数都是正的，所以 $\Phi_{\mathrm{quantum}}$ 只依赖于 $t_1$。又因为 $T_1$ 是除子（超平面截面 $H$）的类，式~\eqref{eq:gw-potential-divisor-exponential} 给出
\begin{equation}\label{eq:quintic-quantum-potential}
\Phi_{\mathrm{quantum}}=\sum_{d\ge1}\langle I_{0,0,d}\rangle e^{dt_1}q^d,
\end{equation}
其中照例用 $\langle I_{0,0,d}\rangle$ 代替 $\langle I_{0,0,\beta}\rangle$。因此，五次三维簇的 Gromov--Witten 势为
\[
\Phi=\frac12t_0^2t_3+t_0t_1t_2+\frac56t_1^3
-\sum_{i<j}q_{ij}t_0u_iu_j+\sum_{d\ge1}\langle I_{0,0,d}\rangle(e^{t_1}q)^d.
\]
由此，描述 $V$ 的大量子乘积相当简单。除 $T_1*T_1$ 外，所有基元素对的乘积都与杯积一致；而这个例外的乘积为
\[
T_1*T_1=\Phi_{111}T_2.
\]
此外，有
\[
\int_V T_1*T_1*T_1=\Phi_{111},
\]
并且，由于 $T_1$ 是 $H$ 的类，$\Phi_{111}$ 由如下公式给出：
\[
\Phi_{111}=5+\sum_{d\ge1}\langle I_{0,0,d}\rangle d^3(e^{t_1}q)^d
=5+\sum_{d\ge1}\langle I_{0,3,d}\rangle(H,H,H)(e^{t_1}q)^d.
\]
在这些公式中，将 $e^{t_1}q$ 换成 $q$，就恰好恢复五次三维簇的小量子上同调。

更一般地，对任意光滑 Calabi--Yau 三维簇，Gromov--Witten 势都可写成特别简洁的形式。首先注意，对于齐次插入类，$\beta\ne0$ 蕴含
\begin{equation}\label{eq:cy-nonzero-degree-insertions}
\langle I_{0,n,\beta}\rangle(\alpha_1,\ldots,\alpha_n)=0,
\quad\text{除非所有 $\alpha_i$ 的次数都为 $2$。}
\end{equation}
若选定的 $H^2(V,\mathbb C)$ 的基为 $T_1,\ldots,T_r$，式~\eqref{eq:cy-nonzero-degree-insertions} 说明 $\Phi_{\mathrm{quantum}}$ 是 $t_1,\ldots,t_r$ 的函数；由式~\eqref{eq:gw-potential-divisor-exponential}，这些变量以指数形式出现。因此，量子部分可写成
\[
\Phi_{\mathrm{quantum}}=\sum_{\beta\ne0}\langle I_{0,0,\beta}\rangle
 e^{\int_\beta\delta}q^\beta,
\]
其中 $\delta=\sum_{i=1}^rt_iT_i$。若去掉原来的因子 $q^\beta$，并令
$q^\beta=e^{\int_\beta\delta}=\prod_{i=1}^re^{t_i\int_\beta T_i}$，Gromov--Witten 势就由如下简洁公式给出：
\begin{equation}\label{eq:cy-total-gw-potential}
\Phi=\Phi_{\mathrm{classical}}+\sum_{\beta\ne0}\langle I_{0,0,\beta}\rangle q^\beta.
\end{equation}
利用式~\eqref{eq:cy-nonzero-degree-insertions} 和~\eqref{eq:cy-total-gw-potential}，得到
\begin{equation}\label{eq:cy-potential-three-derivatives}
\begin{aligned}
\frac{\partial^3\Phi}{\partial t_i\partial t_j\partial t_k}
&=\int_V T_i\cup T_j\cup T_k+
\sum_{\beta\ne0}\langle I_{0,0,\beta}\rangle
\int_\beta T_i\int_\beta T_j\int_\beta T_k\,q^\beta\\
&=\langle I_{0,3,0}\rangle(T_i,T_j,T_k)+
\sum_{\beta\ne0}\langle I_{0,3,\beta}\rangle(T_i,T_j,T_k)q^\beta\\
&=\langle T_i,T_j,T_k\rangle.
\end{aligned}
\end{equation}
\SourcePage{262}
其中第二个等式使用点映射公理和除子公理，第三个等式使用式~\eqref{eq:cy-quantum-correlation-gw}。这说明，将 Gromov--Witten 势改写为式~\eqref{eq:cy-total-gw-potential} 后，定义~\ref{def:big-quantum-product} 就成为小量子乘积的定义。因此，对 Calabi--Yau 三维簇，大量子上同调与小量子上同调实质上没有区别。还可以验证，此时 WDVV 方程归结为杯积的结合律。这与 $\mathbb P^2$ 的情形很不相同：在那里，大量子乘积与 WDVV 方程包含的计数信息远多于小量子上同调。

在\hyperref[subsec:quantum-product-dependence]{第~\ref*{subsec:quantum-product-dependence} 小节}将看到，对一般的簇，只要在 $T_i$ 的次数不为 $2$ 时令参数 $t_i=0$，大量子乘积与小量子乘积也实质上没有区别。

\section{Dubrovin 形式体系}\label{sec:dubrovin-formalism}
\index{Dubrovin 形式体系|hyperpage}\index{挠率|hyperpage}\index{平坦联络|hyperpage}
Gromov--Witten 势的代数结果是大量子乘积。本节将探讨 $\Phi$ 的主要几何结果，即 \emph{Dubrovin 联络}。\hyperref[sec:a-model-vhs]{第~\ref*{sec:a-model-vhs} 节}将利用 Dubrovin 形式体系，在 Calabi--Yau 流形的 Kähler 模空间上构造一种自然出现的 Hodge 结构变分。\index{Dubrovin 联络|hyperpage}

我们将 Dubrovin 联络的定义从 Gromov--Witten 势 $\Phi$ 推广到任意势函数 $F$；为此，依照 \SourceCite{Dubrovin1} 发展更一般的理论。

照例，取 $H^*(X,\mathbb C)$ 的基 $T_0=1,\ldots,T_m$，以及相应的超交换变量 $t_0,\ldots,t_m$。交配对的矩阵记为 $(g_{kl})$，其逆矩阵记为 $(g^{kl})$。令 $F\colon H^*(X,\mathbb C)\to\mathbb C$ 是关于 $t_i$ 的偶形式幂级数。在较好的情形，$F$ 在超流形 $H^*(X,\mathbb C)$ 的某个开集上收敛。给定 $F$ 后，定义张量 $A_{ijk}$ 为
\begin{equation}\label{eq:dubrovin-cubic-tensor}
A_{ijk}=\frac{\partial^3F}{\partial t_i\partial t_j\partial t_k}.
\end{equation}
若用下标表示偏导数，则须注意 $A_{ijk}=F_{kji}$。按物理学中通常的“升指标”约定，令
\begin{equation}\label{eq:dubrovin-raised-tensor}
A_{ij}^k=\sum_lA_{ijl}g^{lk}.
\end{equation}
现在可定义两个基本对象：二元运算
\[
T_i*T_j=\sum_kA_{ij}^kT_k,
\]
以及 \emph{Dubrovin 联络}
\[
\nabla^\lambda_{\frac{\partial}{\partial t_i}}
\left(\frac{\partial}{\partial t_j}\right)
=\lambda\sum_kA_{ij}^k\frac{\partial}{\partial t_k},
\]
其中 $\lambda$ 是非零复数。

二元运算 $*$ 与联络 $\nabla^\lambda$ 关系密切：一者的性质可以转化为另一者的性质。下面简要说明这种关系的三个重要方面。
\begin{itemize}
\item \textbf{挠率与交换律。}式~\eqref{eq:dubrovin-cubic-tensor} 和~\eqref{eq:dubrovin-raised-tensor} 容易推出恒等式
\[
A_{ij}^k=(-1)^{\deg t_i\deg t_j}A_{ji}^k.
\]
\SourcePage{263}
利用此式，不难证明联络 $\nabla^\lambda$ 的挠率为零（在超流形意义下），且二元运算 $*$ 超交换。
\item \textbf{曲率与结合律。}计算联络 $\nabla^\lambda$ 的曲率（在超流形意义下）可知，$(\nabla^\lambda)^2$ 是两项之和。由式~\eqref{eq:dubrovin-cubic-tensor} 和~\eqref{eq:dubrovin-raised-tensor}，有
\[
\frac{\partial}{\partial t_l}A_{ij}^k
=(-1)^{\deg t_i\deg t_l}\frac{\partial}{\partial t_i}A_{lj}^k,
\]
所以第一项为零；第二项为零，当且仅当
\[
\sum_aA_{ij}^aA_{ak}^l
=(-1)^{\deg t_i(\deg t_j+\deg t_k)}\sum_aA_{jk}^aA_{ai}^l.
\]
我们认出，这就是 WDVV 方程。于是容易得到
\[
\begin{aligned}
\text{$\nabla^\lambda$ 是平坦联络}
&\Longleftrightarrow\text{$F$ 满足 WDVV 方程}\\
&\Longleftrightarrow\text{$*$ 满足结合律。}
\end{aligned}
\]
\item \textbf{单位元。}我们希望 $T_0$ 是 $*$ 的单位元。容易得到等价关系
\[
\text{对所有 $i,j$，$A_{0ij}=g_{ij}$}
\Longleftrightarrow\text{对所有 $i,j$，$A_{0i}^j=\delta_{ij}$，}
\]
随后直接可证
\[
\begin{aligned}
\text{$T_0$ 是 $*$ 的单位元}
&\Longleftrightarrow\text{对所有 $i,j$，$A_{0ij}=g_{ij}$}\\
&\Longleftrightarrow\text{对所有 $i$，$\nabla^\lambda_{\frac{\partial}{\partial t_0}}
\left(\frac{\partial}{\partial t_i}\right)=\lambda\frac{\partial}{\partial t_i}$。}
\end{aligned}
\]
\end{itemize}
这些断言的证明见 \SourceCite{Dubrovin1,KoM1,Manin3}。

从现在起，始终假定 $F$ 满足 WDVV 方程，且对所有 $i,j$ 都有 $A_{0ij}=g_{ij}$。这意味着 Dubrovin 联络 $\nabla^\lambda$ 平坦且无挠，而 $\nabla^\lambda_{\frac{\partial}{\partial t_0}}$ 在 $\frac{\partial}{\partial t_i}$ 上的作用是乘以 $\lambda$。就二元运算 $*$ 而言，$H^*(X,\mathbb C)$ 有如下性质。
\begin{proposition}\label{prop:dubrovin-frobenius}
若 $F$ 满足 WDVV 方程，且对所有 $i,j$ 都有 $A_{0ij}=g_{ij}$，则 $H^*(X,\mathbb C)$ 在 $*$ 下是以 $T_0$ 为单位元的 Frobenius 代数。
\end{proposition}
\begin{proof}
假设说明，$*$ 给出以 $T_0$ 为单位元的超交换代数。因此只须证明它满足 Frobenius 性质。式~\eqref{eq:dubrovin-cubic-tensor} 还给出
\[
A_{ijk}=(-1)^{\deg t_i(\deg t_j+\deg t_k)}A_{jki}.
\]
我们又知道 $F$ 是关于 $t_i$ 的偶函数。这些事实蕴含
\[
g(T_i*T_j,T_k)=g(T_i,T_j*T_k).
\]
所以 $H^*(X,\mathbb C)$ 在 $*$ 下是 Frobenius 代数。
\end{proof}
此前研究的势函数 $\Phi$ 与 $\Phi_{\mathrm{small}}$ 都满足命题~\ref{prop:dubrovin-frobenius} 的条件。这就是大量子乘积与小量子乘积使 $H^*(X,\mathbb C)$ 成为 Frobenius 代数的原因。

Dubrovin 联络的构造在超流形 $M=H^*(X,\mathbb C)$ 的切丛上定义了 $\nabla^\lambda$。若将 $\frac{\partial}{\partial t_i}$ 与 $T_i$ 等同，就得到平凡丛
\SourcePage{264}
$H^*(X,\mathbb C)\times M$ 上的联络 $\nabla^\lambda$。容易验证，这个新联络由
\[
\nabla^\lambda_{\frac{\partial}{\partial t_i}}(T_j)=\lambda T_i*T_j
\]
给出。因此，可以将 Dubrovin 联络看作由量子乘法给出的联络。下一节将用到这种观点。

Dubrovin 形式体系还有若干方面，下面继续讨论。首先注意，到目前为止的理论并未使用\hyperref[subsec:gw-potential-structure]{第~\ref*{subsec:gw-potential-structure} 小节}所述 $\Phi$ 的齐次性。用 Dubrovin 联络的语言来表述，齐次性意味着存在一个偶向量场 $E$，称为 \emph{Euler 向量场}，它与 $\nabla^\lambda$ 有良好的相容关系。\index{Euler 向量场|hyperpage}

其次，$\nabla^\lambda$ 是由一个参数 $\lambda$ 参数化的联络族。若将 $\lambda$ 看作 $\mathbb P^1$ 的一个仿射坐标，\SourceCite{KoM1,Manin3} 说明，$\nabla^\lambda$ 可延拓为 $H^*(X,\mathbb C)\times\mathbb P^1$ 切丛上的联络 $\widehat\nabla$，在 $\lambda=0,\infty$ 处有奇点。当 $X=\mathbb P^2$ 时，$\widehat\nabla$ 的 Laplace 变换与第六 Painlevé 方程有一些令人意外的联系，见 \SourceCite{Manin3}。

超流形 $H^*(X,\mathbb C)$，连同度量 $g$ 及满足 WDVV 方程的偶势函数 $F$，给出的结构是 \emph{Frobenius 流形} 的一个例子 \SourceCite{Dubrovin1,Dubrovin2}。\index{Frobenius 流形|hyperpage}这有一些有趣的推论。例如，考虑 $F$ 的幂级数展开。忽略关于 $t_i$ 的次数不超过 $2$ 的项，可将 $F$ 写成
\[
F=\sum_{n=3}^{\infty}\sum_{|\alpha|=n}c_\alpha t^\alpha.
\]
随后，通过等式
\[
c_\alpha=\frac{\epsilon(\alpha)}{\alpha!}\langle I_n\rangle(T^\alpha),\qquad n\ge3,
\]
定义 $\langle I_n\rangle\colon H^*(X,\mathbb C)^{\otimes n}\to\mathbb C$，其中采用式~\eqref{eq:gw-potential-super-expansion} 的记号。于是可将 $F$ 写成
\[
F=\sum_{n=3}^{\infty}\frac1{n!}\langle I_n\rangle(\gamma^n),
\]
这是一个熟悉的形式。\SourceCite{Manin3} 中，Manin 证明了一个引人注目的结果：对 $n\ge3$，存在映射
\[
I_n\colon H^*(X,\mathbb C)^{\otimes n}\longrightarrow
H^*(\overline M_{0,n},\mathbb C),
\]
它们满足\IndexLocation{indexloc:0507:264}分裂公理和\IndexLocation{indexloc:0501:264}等变性公理，并按通常的方式确定 $\langle I_n\rangle$，即
\[
\langle I_n\rangle(\alpha_1,\ldots,\alpha_n)
=\int_{\overline M_{0,n}}I_n(\alpha_1,\ldots,\alpha_n).
\]
这些 $I_n$（$n\ge3$）组成 Manin 所称的\emph{上同调场论}，而 $\langle I_n\rangle$ 是它的\emph{\IndexLocation{indexloc:0127:264}\IndexLocation{indexloc:0183:264}关联函数}。\index{上同调场论|hyperpage}结论是：$H^*(X,\mathbb C)$ 上任意满足 WDVV 方程的偶函数 $F$，都会给出一种与 Gromov--Witten 类惊人相似的结构。还要指出，\hyperref[subsec:wdvv-associativity]{第~\ref*{subsec:wdvv-associativity} 小节}提到的量子上同调 Künneth 公式，更一般地也适用于上同调场论 \SourceCite{KoM2,Kaufmann}。还可以定义更高亏格的上同调场论 \SourceCite{Manin3,KK}。
\SourcePage{265}
最后，讨论 Dubrovin 形式体系时，务必记住，函数 $F$ 可以是形式幂级数。因此，联络 $\nabla^\lambda$ 也可能是形式联络；为此，它们有时称为\emph{形式 Dubrovin 联络}。这种情形下，必须使用\hyperref[subsec:formal-small-quantum]{第~\ref*{subsec:formal-small-quantum} 小节}讨论的一种系数环，才能定义上同调中相应的环结构。

\section{A 型 Hodge 结构变分}\label{sec:a-model-vhs}
\index{Hodge结构变分@Hodge 结构变分!A 型|hyperpage}
本节将说明，小量子上同调如何在\hyperref[chap:moduli]{第~\ref*{chap:moduli} 章}研究的 Kähler 模空间上给出自然的 Hodge 结构变分。这将使我们能够在\hyperref[sec:mirror-symmetry-revisited]{第~\ref*{sec:mirror-symmetry-revisited} 节}为镜像对建立精确的数学定义。
\subsection{大量子上同调与小量子上同调的关系}\label{subsec:quantum-product-dependence}
为研究\hyperref[subsec:a-model-connection]{第~\ref*{subsec:a-model-connection} 小节}定义的 A 模型联络，须先理解大量子上同调与小量子上同调之间的关系。

沿用通常的设定，取 $H^*(X,\mathbb Q)$ 的上同调基 $T_0=1,T_1,\ldots,T_m$，其中 $T_1,\ldots,T_r$ 张成 $H^2(X,\mathbb Q)$。定义~\ref{def:big-quantum-product} 中的大量子乘积是关于 $t_0,\ldots,t_m$ 和 $q^\beta=\exp(2\pi i\int_\beta\omega)$ 的形式幂级数。首先注意，变量 $t_1,\ldots,t_r$ 在此级数中以指数形式出现。更准确地说，定义~\ref{def:big-quantum-product}、引理~\ref{lem:gw-potential-third-derivatives} 和式~\eqref{eq:gw-potential-divisor-exponential} 的推导方法，容易给出如下引理。
\begin{lemma}\label{lem:big-quantum-divisor-exponential}
令 $\delta=\sum_{i=1}^rt_iT_i$ 与 $\epsilon=t_0T_0+\sum_{i=r+1}^mt_iT_i$，则大量子乘积为
\[
T_i*T_j=\sum_k\sum_{n=0}^{\infty}\sum_\beta\frac1{n!}
\langle I_{0,n+3,\beta}\rangle(T_i,T_j,T_k,\epsilon^n)
 e^{\int_\beta\delta}q^\beta T^k.
\]
\end{lemma}
利用此引理，可用几种方式比较大量子上同调与小量子上同调。最简单的是在 $T_i*T_j$ 的公式中令变量 $t_0,\ldots,t_m=0$。用引理~\ref{lem:big-quantum-divisor-exponential} 的记号，就是令 $\delta=\epsilon=0$，得到
\[
T_i*T_j\big|_{\delta=\epsilon=0}
=\sum_k\sum_\beta\langle I_{0,3,\beta}\rangle(T_i,T_j,T_k)q^\beta T^k.
\]
由于
\begin{equation}\label{eq:small-quantum-product-comparison}
T_i*_{\mathrm{small}}T_j=\sum_k\sum_\beta
\langle I_{0,3,\beta}\rangle(T_i,T_j,T_k)q^\beta T^k,
\end{equation}
可知
\[
T_i*T_j\big|_{\delta=\epsilon=0}=T_i*_{\mathrm{small}}T_j.
\]
因此，小量子乘积可看作大量子乘积的限制。

不过，为使用\hyperref[sec:dubrovin-formalism]{第~\ref*{sec:dubrovin-formalism} 节}的 Dubrovin 形式体系，我们采取稍有不同的方法，即“限制到 $H^2(X)$”。其含义如下：大量子乘积用 $t_0,\ldots,t_m$ 与 $q^\beta$ 在 $H^*(X)$ 上形式地定义，所以在 $T_i*T_j$ 的公式中“限制到 $H^2(X)$”是指令 $t_0=t_{r+1}=\cdots=
\SourcePage{266}
t_m=0$。这等价于在引理~\ref{lem:big-quantum-divisor-exponential} 中令 $\epsilon=0$，得到
\begin{equation}\label{eq:restricted-big-quantum-product}
T_i*T_j\big|_{\epsilon=0}=\sum_k\sum_\beta
\langle I_{0,3,\beta}\rangle(T_i,T_j,T_k)e^{\int_\beta\delta}q^\beta T^k.
\end{equation}
将它称为\emph{限制的大量子乘积}。它与小量子乘积~\eqref{eq:small-quantum-product-comparison} 相似，但尚不完全相同。注意到
\[
e^{\int_\beta\delta}q^\beta=e^{\int_\beta(\delta+2\pi i\omega)},
\]
这种相似性就更为明显。因此，式~\eqref{eq:restricted-big-quantum-product} 和~\eqref{eq:small-quantum-product-comparison} 中的指数项都涉及 $H^2(X)$ 中某个元素在 $\beta$ 上的积分，分别为 $\delta+2\pi i\omega$ 和 $2\pi i\omega$。

由此，可以如下使两个公式一致：令 $\omega=0$，并设
\[
\mathbf q^\beta=\exp\left(\int_\beta\delta\right)
=\exp\left(\int_\beta\sum_{i=1}^rt_iT_i\right),
\]
则限制的大量子乘积成为
\begin{equation}\label{eq:restricted-big-quantum-q}
T_i*T_j\big|_{\epsilon=0}=\sum_k\sum_\beta
\langle I_{0,3,\beta}\rangle(T_i,T_j,T_k)\mathbf q^\beta T^k.
\end{equation}
它与式~\eqref{eq:small-quantum-product-comparison} 很好地对应。从形式角度看，将式~\eqref{eq:restricted-big-quantum-q} 中的 $\mathbf q^\beta$ 看作适当系数环中依赖于 $\beta,t_1,\ldots,t_r$ 的形式符号，再使用将 $\mathbf q^\beta$ 映为 $q^\beta$ 的系数环同态，就得到与式~\eqref{eq:small-quantum-product-comparison} 相同的形式级数。

另一方面，将式~\eqref{eq:restricted-big-quantum-q} 和~\eqref{eq:small-quantum-product-comparison} 看作 $H^2(X)$ 上关于变量 $t_1,\ldots,t_r$ 的形式函数时，也可以联系起来。只要在式~\eqref{eq:small-quantum-product-comparison} 中令
\begin{equation}\label{eq:quantum-coordinate-rescaling}
\omega=\frac1{2\pi i}\sum_{i=1}^rt_iT_i,
\end{equation}
这些函数就相同。因此，式~\eqref{eq:small-quantum-product-comparison} 与~\eqref{eq:restricted-big-quantum-q} 的收敛假设也相同。采用这些约定后，无论在形式意义还是解析意义下，大量子乘积限制到 $H^2(X)$ 就是小量子乘积。

\subsection{A 模型联络}\label{subsec:a-model-connection}
\index{复化 Kähler 空间|hyperpage}
本章余下部分中，令 $V$ 是维数 $d\ge3$ 的 Calabi--Yau 流形。回顾\hyperref[chap:introduction]{第~\ref*{chap:introduction} 章}，$V$ 的复化 Kähler 空间是商空间
\[
K_{\mathbb C}(V)=\{\omega\in H^2(V,\mathbb C)\colon
\operatorname{Im}(\omega)\text{ 为 Kähler 类}\}/\operatorname{im}H^2(V,\mathbb Z),
\]
其中 $\operatorname{im}H^2(V,\mathbb Z)$ 是自然映射 $H^2(V,\mathbb Z)\to H^2(V,\mathbb C)$ 的像。下面说明，小量子乘积 $*_{\mathrm{small}}$ 如何在 $K_{\mathbb C}(V)$ 上给出一个有趣的联络。

照例，使用 $H^*(V,\mathbb Q)$ 的上同调基 $T_0,\ldots,T_m$。此外，假定其中位于 $H^2(V,\mathbb Q)$ 的基元素 $T_1,\ldots,T_r$ 满足：
\begin{itemize}
\item $T_1,\ldots,T_r$ 是 $\operatorname{im}H^2(V,\mathbb Z)$ 的一组基；
\item $T_1,\ldots,T_r$ 位于 Kähler 锥的闭包中。
\end{itemize}
由于 $H^{2,0}(V)=0$，容易看出可以找到这样的 $T_1,\ldots,T_r$。

现在引入对应于 $T_1,\ldots,T_r$ 的变量 $u_1,\ldots,u_r$，并令
\begin{equation}\label{eq:a-model-kahler-coordinates}
\omega=\sum_{i=1}^ru_iT_i\in H^2(V,\mathbb R)+iK(V).
\end{equation}
\SourcePage{267}
于是，对任意 $\beta\in H_2(V,\mathbb Z)$，函数 $q^\beta=e^{2\pi i\int_\beta\omega}$ 在 $K_{\mathbb C}(V)$ 上良好定义，因为 $T_i\in\operatorname{im}H^2(V,\mathbb Z)$。\hyperref[chap:moduli]{第~\ref*{chap:moduli} 章}曾用 $t_i$ 而非 $u_i$；很快就会明白这里使用 $u_i$ 的原因。

现在令 $*_{\mathrm{small}}$ 为 $H^*(V,\mathbb C)$ 上的小量子乘积。这需要一些收敛假设，留待稍后讨论。考虑 $K_{\mathbb C}(V)$ 上的平凡丛 $\mathcal H=H^*(V,\mathbb C)\times K_{\mathbb C}(V)$。按如下规则定义 $\mathcal H$ 上的 \emph{A 模型联络} $\nabla$：\index{A 模型联络|hyperpage}
\begin{equation}\label{eq:a-model-connection}
\nabla_{\frac{\partial}{\partial u_i}}(T_j)=T_i*_{\mathrm{small}}T_j,
\qquad 1\le i\le r,\quad0\le j\le m.
\end{equation}
因此，$\nabla$ 由小量子乘积构造，不过其中只用 $H^2(V,\mathbb C)$ 的元素进行乘法。用小量子乘积的定义展开 $\nabla$，得到
\begin{equation}\label{eq:a-model-connection-expanded}
\nabla_{\frac{\partial}{\partial u_i}}(T_j)
=\sum_k\sum_{\beta\in H_2(V,\mathbb Z)}
\langle I_{0,3,\beta}\rangle(T_i,T_j,T_k)q^\beta T^k.
\end{equation}
注意，此联络定义在 $K_{\mathbb C}(V)$ 上，且 $q^\beta$ 是这个空间上的非常值函数。因此，$\nabla$ 记录了当 $\omega$ 改变时，小量子乘积的变化。

严格地说，$\nabla$ 只是一个形式联络，因为式~\eqref{eq:a-model-connection-expanded} 中级数的收敛性仍是猜想。不过，为便于阐述，我们将把 $\nabla$ 看作真正的联络。\hyperref[subsec:a-model-asymptotics]{第~\ref*{subsec:a-model-asymptotics} 小节}将为 $\nabla$ 提出精确的收敛假设。

联络 $\nabla$ 有如下良好性质。
\begin{proposition}\label{prop:a-model-flat}
A 模型联络 $\nabla$ 平坦。
\end{proposition}
\begin{proof}
只要谨慎使用\hyperref[sec:dubrovin-formalism]{第~\ref*{sec:dubrovin-formalism} 节}的理论，这一结论就由 Dubrovin 联络 $\nabla^\lambda$ 的平坦性得到。

我们对取 $q^\beta=1$ 的 Gromov--Witten 势应用 Dubrovin 形式体系。这在 $H^*(V,\mathbb C)$ 上给出平坦联络 $\nabla^\lambda$，定义为
\begin{equation}\label{eq:dubrovin-big-quantum-connection}
\nabla^\lambda_{\frac{\partial}{\partial t_i}}(T_j)=\lambda T_i*T_j,
\end{equation}
其中 $*$ 为大量子乘积，$t_0,\ldots,t_m$ 为 $H^*(V,\mathbb C)$ 的通常变量。像\hyperref[subsec:quantum-product-dependence]{第~\ref*{subsec:quantum-product-dependence} 小节}那样，将 $\nabla^\lambda$ 限制到 $H^2(V,\mathbb C)$，则断言
\begin{equation}\label{eq:a-model-dubrovin-identification}
\nabla=\nabla^\lambda,\qquad\lambda=\frac1{2\pi i}.
\end{equation}
一旦证明此式，命题便由 $\nabla^\lambda$ 的平坦性立即得到。

为证明式~\eqref{eq:a-model-dubrovin-identification}，注意，由式~\eqref{eq:restricted-big-quantum-q}，限制的 Dubrovin 联络可写成
\[
\nabla^\lambda_{\frac{\partial}{\partial t_i}}(T_j)
=\lambda\sum_k\sum_\beta\langle I_{0,3,\beta}\rangle(T_i,T_j,T_k)\mathbf q^\beta T^k,
\]
其中 $t_1,\ldots,t_r$ 是 $H^2(V,\mathbb C)$ 的变量。由\hyperref[subsec:quantum-product-dependence]{第~\ref*{subsec:quantum-product-dependence} 小节}，在代换 $\mathbf q^\beta\mapsto q^\beta$ 下，这成为
\begin{equation}\label{eq:restricted-dubrovin-connection}
\nabla^\lambda_{\frac{\partial}{\partial t_i}}(T_j)=\lambda T_i*_{\mathrm{small}}T_j,
\end{equation}
前提是，像式~\eqref{eq:quantum-coordinate-rescaling} 那样，以
\[
\omega=\frac1{2\pi i}\sum_{i=1}^rt_iT_i
\]
定义 $q^\beta$。
\SourcePage{268}
A 模型联络用变量 $u_1,\ldots,u_r$ 定义；将上式与式~\eqref{eq:a-model-kahler-coordinates} 比较，得到 $u_j=(1/2\pi i)t_j$。用复化 Kähler 空间上的坐标 $q_1,\ldots,q_r$ 表示，就是
\[
q_j=\exp(2\pi iu_j)=\exp(t_j),
\]
因而
\begin{equation}\label{eq:a-model-coordinate-derivatives}
\frac{\partial}{\partial t_j}=q_j\frac{\partial}{\partial q_j},
\qquad
\frac{\partial}{\partial u_j}=2\pi iq_j\frac{\partial}{\partial q_j}
=2\pi i\frac{\partial}{\partial t_j}.
\end{equation}
于是，式~\eqref{eq:a-model-connection} 和~\eqref{eq:restricted-dubrovin-connection} 立即给出式~\eqref{eq:a-model-dubrovin-identification}，命题得证。
\end{proof}
第 \hyperref[chap:quantum-differential-equations]{\ref*{chap:quantum-differential-equations}}、\hyperref[chap:mirror-theorem]{\ref*{chap:mirror-theorem}} 章中，将使用 $t_1,\ldots,t_r$ 而非定义 A 模型联络时所用的 $u_1,\ldots,u_r$，因为这样更方便。不过，本章余下部分仍将使用 $u_1,\ldots,u_r$，因为它们与\hyperref[chap:moduli]{第~\ref*{chap:moduli} 章}的讨论联系更直接。

此外，\hyperref[chap:quantum-differential-equations]{第~\ref*{chap:quantum-differential-equations} 章}将遇到 $\nabla^\lambda$ 的一种扭曲形式，即 \emph{Givental 联络} $\nabla^g$，定义为\index{Givental联络@Givental 联络|hyperpage}
\[
\nabla^g_{\frac{\partial}{\partial t_i}}\left(\sum_{j=0}^ma_jT_j\right)
=\hbar\sum_{j=0}^m\frac{\partial a_j}{\partial t_i}T_j
-\sum_{j=0}^ma_jT_i*_{\mathrm{small}}T_j,
\]
其中 $\hbar$ 为参数。目前可将 $\hbar$ 看作非零复数，不过，\hyperref[subsec:equivariant-giventals-connection]{第~\ref*{subsec:equivariant-giventals-connection} 小节}将说明，$\hbar$ 在 $\mathbb C^*$ 等变上同调中有内在含义。用 Dubrovin 联络 $\nabla^\lambda$ 表示，则有
\[
\nabla^g=\hbar\nabla^\lambda,\qquad\lambda=-\hbar^{-1}.
\]
严格地说，由于因子 $\hbar$，$\nabla^g$ 不是联络，但讨论其平坦截面仍有意义；这些截面就是 $\nabla^{-\hbar^{-1}}$ 的平坦截面。

\hypertarget{source:268-batch6}{}
\subsection{A 模型联络的渐近性质}\label{subsec:a-model-asymptotics}\index{单值性!A 模型|hyperpage}
定义 A 模型联络 $\nabla$ 时，我们使用了位于 $V$ 的 Kähler 锥闭包内的整基 $T_1,\ldots,T_r$。为研究 $\nabla$ 的渐近性质，令 $\sigma$ 为这些类生成的单纯锥。容易看出 $\operatorname{Int}(\sigma)\subset K(V)$，因此
\[
D_\sigma=\bigl(H^2(V,\mathbb R)+i\operatorname{Int}(\sigma)\bigr)/\operatorname{im}H^2(V,\mathbb Z)\subset K_{\mathbb C}(V)
\]
是 $K_{\mathbb C}(V)$ 的开子集。

如\hyperref[subsec:large-radius-limits]{第~\ref*{subsec:large-radius-limits} 小节}所述，映射
\[
u_1T_1+\cdots+u_rT_r\longmapsto(q_1,\ldots,q_r)=(e^{2\pi iu_1},\ldots,e^{2\pi iu_r})
\]
诱导双全纯同构 $D_\sigma\simeq(\Delta^*)^r$。利用包含关系 $(\Delta^*)^r\subset\Delta^r$，按\hyperref[chap:moduli]{第~\ref*{chap:moduli} 章}的术语，可以把 $0\in\Delta^r$ 看作大半径极限点。此外，假设\hyperref[subsec:large-radius-limits]{第~\ref*{subsec:large-radius-limits} 小节}的\hyperref[conj:kahler-cone]{锥猜想}成立，那么这一大半径极限点自然位于 $V$ 的 Kähler 模空间某个光滑紧化的边界上。

对于小量子乘积，我们希望作如下收敛假设：只要 $\omega$ 在上述映射下的像位于 $0\in\Delta^r$ 的某个固定邻域 $U$ 中，式~\eqref{eq:a-model-connection-expanded} 的级数就收敛。因此，A 模型联络定义在 $U\cap(\Delta^*)^r$ 上。为使叙述简便，以下假设 $\nabla$ 定义在整个 $(\Delta^*)^r$ 上。
\SourcePage{269}
现在应用\hyperref[chap:hodge-yukawa]{第~\ref*{chap:hodge-yukawa} 章}的方法，研究 A 模型联络的单值性。在 $D_\sigma$ 的部分紧化 $\Delta^r$ 中，补集 $\Delta^r-D_\sigma=\Delta^r-(\Delta^*)^r$ 是正规交叉除子 $\bigcup_{j=1}^rD_j$。令 $\mathcal T_j$ 表示绕 $D_j$ 的单值性，它由围绕 $(\Delta^*)^r$ 第 $j$ 个因子的原点一周的回路确定。这些单值性变换可按如下方式计算。
\begin{proposition}\label{prop:a-model-regular-monodromy}\index{正则奇点|hyperpage}
Calabi--Yau 流形 $V$ 的 A 模型联络 $\nabla$ 沿 $\Delta^r-D_\sigma$ 具有正则奇点。此外，单值性变换 $\mathcal T_j$ 幺幂；在相差共轭的意义下，它的对数 $N_j=\log\mathcal T_j$ 由与 $-T_j$ 作杯积给出。
\end{proposition}
\begin{proof}
由式~\eqref{eq:a-model-coordinate-derivatives}，$\Delta^r$ 的坐标 $q_j=e^{2\pi iu_j}$ 满足
\[
\frac{\partial}{\partial u_j}=2\pi iq_j\frac{\partial}{\partial q_j}.
\]
于是，式~\eqref{eq:a-model-connection-expanded} 和 $T^l=\sum_sg^{ls}T_s$ 给出
\begin{equation}\label{eq:a-model-log-poles}
\nabla_{\frac{\partial}{\partial q_j}}(T_k)
=\frac1{2\pi iq_j}\sum_{l,s}\sum_{\beta\in H_2(V,\mathbb Z)}
\langle I_{0,3,\beta}\rangle(T_j,T_k,T_l)q^\beta g^{ls}T_s.
\end{equation}
因为 $\beta$ 有效，而且 $T_1,\ldots,T_r$ 位于 Kähler 锥的闭包内，
\[
q^\beta=q_1^{\int_\beta T_1}\cdots q_r^{\int_\beta T_r}
\]
的指数均非负。实际上，在固定 $j$ 的联络矩阵中，只要 $\beta\ne0$，就可以假设 $q_j$ 的指数为正。为说明这一点，设 $\int_\beta T_j=0$。因为 $\beta\ne0$，除子公理给出
\[
\langle I_{0,3,\beta}\rangle(T_j,T_k,T_l)
=\left(\int_\beta T_j\right)\langle I_{0,2,\beta}\rangle(T_k,T_l)=0.
\]
因此，可以把式~\eqref{eq:a-model-log-poles} 的右端看作一个级数，其中每个非零有效 $\beta$ 对应的非零项，在所固定的 $q_j$ 中具有正指数。

由式~\eqref{eq:a-model-log-poles}，$\nabla$ 相对于基 $T_0,\ldots,T_m$ 的联络矩阵，沿 $\Delta^r-D_\sigma$ 至多有对数极点。由 \SourceCite[定理 4.1]{Deligne1}，这些点为正则奇点。

现在固定 $1\le j\le r$。为计算单值性 $\mathcal T_j$，在 $\Delta^r$ 中取由 $q_i=\text{常数}$（$i\ne j$）给出的“切片”。这样就可把 $\nabla$ 看作以 $q_j$ 为坐标的 $\Delta^*$ 上的联络；其联络矩阵由固定 $j$ 后的式~\eqref{eq:a-model-log-poles} 给出。为计算该矩阵在 $q_j=0$ 处的留数，注意，若 $\beta\ne0$ 有效且相应项非零，则
\begin{equation}\label{eq:a-model-positive-vanishing}
q^\beta\longrightarrow0\qquad\text{当 }q_j\longrightarrow0,
\end{equation}
因为 $q^\beta$ 在 $q_j$ 中具有正指数。

于是，由式~\eqref{eq:a-model-positive-vanishing} 和~\eqref{eq:a-model-log-poles}，$\nabla_{\partial/\partial q_j}$ 在 $q_j=0$ 处的留数矩阵为
\[
\frac1{2\pi i}\left(\sum_l\langle I_{0,3,0}\rangle(T_j,T_k,T_l)g^{ls}\right)_{k,s}
=\frac1{2\pi i}\left(\sum_lg^{ls}\int_VT_j\cup T_k\cup T_l\right)_{k,s}.
\]
将此矩阵记为 $\operatorname{Res}_{q_j=0}(\nabla)$。由 $g^{ls}$ 的定义，容易看出，这就是与 $T_j$ 作杯积的线性映射所对应的矩阵的 $1/(2\pi i)$ 倍。

为完成证明，使用\hyperref[subsec:monodromy-pf]{第~\ref*{subsec:monodromy-pf} 小节}中式~\eqref{eq:connection-residue} 后第二组分项列出的事实。关键在于，与 $T_j$ 作杯积是幂零算子，故
\SourcePage{270}
$\operatorname{Res}_{q_j=0}(\nabla)$ 也幂零。其唯一特征值因而为 $0$，这意味着 $\mathcal T_j$ 的矩阵与 $\exp(-2\pi i\operatorname{Res}_{q_j=0}(\nabla))$ 共轭。因此，$\mathcal T_j$ 幺幂，它的对数 $N_j$ 与以 $-T_j$ 作杯积的算子共轭。
\end{proof}
尽管 $N_j$ 与以 $-T_j$ 作杯积的算子共轭，解释这一说法时仍须谨慎。理解 $\mathcal H$ 的典范延拓以后，情形就会更清楚。

因为 $\nabla$ 具有\IndexLocation{indexloc:1225:270}正则奇点，$(\Delta^*)^r$ 上的丛 $\mathcal H$ 在 $\Delta^r$ 上有典范延拓 $\overline{\mathcal H}$。理解 $\overline{\mathcal H}$ 的关键思想如下：若 $s$ 是 $\mathcal H$ 的一个 $\nabla$ 平坦多值截面，则 $\widetilde s=\exp(-\sum_ju_jN_j)s$ 是 $\mathcal H$ 的单值截面，并自然延拓为 $\overline{\mathcal H}$ 的截面 \SourceCite{CaK}。因而可用这些 $\widetilde s$ 在 $0\in\Delta^r$ 的某个邻域内平凡化 $\overline{\mathcal H}$。（我们的符号与 \SourceCite{CaK} 不同，因为这里使用逆时针的单值性生成元。）

一种常见的表述方式是使用联络 $\nabla^c$，其定义为
\begin{equation}\label{eq:a-model-canonical-connection}
\nabla^c=\nabla+\sum_jN_j\,du_j.
\end{equation}
简单计算表明，$\nabla^c$ 平坦截面恰为 $\exp(-\sum_ju_jN_j)s$，其中 $s$ 为 $\nabla$ 平坦截面。因此，$\nabla^c$ 延拓为 $\overline{\mathcal H}$ 上的联络；其平坦截面在边界点 $0\in\Delta^r$ 的某个邻域内平凡化典范延拓。

上同调基 $T_k$ 可以看作 $(\Delta^*)^r$ 上 $\mathcal H$ 的截面，并将此丛平凡化。接下来要把 $T_k$ 与由 $\nabla^c$ 平坦截面给出的平凡化联系起来。
\begin{proposition}\label{prop:a-model-canonical-flat-basis}\index{典范延拓|hyperpage}
对每个 $T_k$，存在唯一的 $\nabla^c$ 平坦截面 $\widetilde s_k$，满足 $\widetilde s_k=T_k+\text{次数更高的项}$ 且 $\widetilde s_k(0)=T_k$。此外，$N_j$ 作用于 $\widetilde s_k$ 的矩阵，等于以 $-T_j$ 与 $T_k$ 作杯积所得的矩阵。
\end{proposition}
\begin{proof}
先证明 $\dim V=3$ 的情形。为简便起见，仅考虑偶数次上同调，其基可写为
\[
T_0=1,T_1,\ldots,T_r,T^1,\ldots,T^r,T^0,
\]
其中 $T^j\in H^4(V,\mathbb Q)$ 与 $T_j\in H^2(V,\mathbb Q)$ 对偶（$1\le j\le r$），而 $T^0\in H^6(V,\mathbb Q)$ 与 $T_0$ 对偶。这些类都是整类。

先用 Gromov--Witten 势描述 $\nabla$ 的一组平坦截面基。\hyperref[subsec:quintic-gw-potential]{第~\ref*{subsec:quintic-gw-potential} 小节}末尾给出了 $\Phi$ 的一个简洁公式，在这里可写为
\begin{equation}\label{eq:a-model-potential}
\begin{aligned}
\Phi&=\frac{(2\pi i)^3}{6}\int_V\left(\sum_{j=1}^ru_jT_j\right)^3
+\sum_{\beta\ne0}\langle I_{0,0,\beta}\rangle q^\beta\\
&=\frac{(2\pi i)^3}{6}\int_V\left(\sum_{j=1}^ru_jT_j\right)^3+\Phi_{\mathrm{hol}},
\end{aligned}
\end{equation}
其中 $q^\beta=q_1^{\int_\beta T_1}\cdots q_r^{\int_\beta T_r}$，$q_j=\exp(2\pi iu_j)$。与命题~\ref{prop:a-model-regular-monodromy} 的证明一样，求和可限制在非零有效 $\beta$ 上，它们满足所有 $\int_\beta T_j\ge0$；因为 $\operatorname{Int}(\sigma)$ 包含 Kähler 类，这些指数不可能全为零。因此，$\Phi_{\mathrm{hol}}$ 关于 $q_j$ 全纯，并在 $0\in\Delta^r$ 处为零。
\SourcePage{271}
用 $\Phi$ 表示时，容易验证 A 模型联络由下列公式给出：
\begin{equation}\label{eq:a-model-threefold-connection}
\begin{aligned}
\nabla_{\frac{\partial}{\partial u_i}}(T^0)&=0,\\
\nabla_{\frac{\partial}{\partial u_i}}(T^k)&=\delta_{ik}T^0,&&1\le k\le r,\\
\nabla_{\frac{\partial}{\partial u_i}}(T_k)&=\frac1{(2\pi i)^3}\sum_l\frac{\partial^3\Phi}{\partial u_i\partial u_k\partial u_l}T^l,&&1\le k\le r,\\
\nabla_{\frac{\partial}{\partial u_i}}(T_0)&=T_i,
\end{aligned}
\end{equation}
其中用到了式~\eqref{eq:a-model-coordinate-derivatives}。直接计算可知，$\mathcal H$ 的一组 $\nabla$ 平坦基为
\begin{equation}\label{eq:a-model-flat-sections-threefold}
\begin{aligned}
s^0&=T^0,\\
s^k&=T^k-u_kT^0,&&1\le k\le r,\\
s_k&=T_k-\frac1{(2\pi i)^3}\sum_l\frac{\partial^2\Phi}{\partial u_k\partial u_l}T^l
+\frac1{(2\pi i)^3}\frac{\partial\Phi}{\partial u_k}T^0,&&1\le k\le r,\\
s_0&=T_0-\sum_lu_lT_l
-\frac1{(2\pi i)^3}\sum_l\left(\frac{\partial\Phi}{\partial u_l}-\sum_ju_j\frac{\partial^2\Phi}{\partial u_j\partial u_l}\right)T^l\\
&\quad+\frac1{(2\pi i)^3}\left(2\Phi-\sum_ju_j\frac{\partial\Phi}{\partial u_j}\right)T^0.
\end{aligned}
\end{equation}
这些截面在 $(\Delta^*)^r$ 上多值，因为 $u_j$ 出现在 $\Phi$ 的“拓扑”部分 $\frac{(2\pi i)^3}{6}\int_V(\sum_{j=1}^ru_jT_j)^3$ 中。

下一步计算单值性。基本观察是，沿 $(\Delta^*)^r$ 的第 $j$ 个切片绕一周，会把 $u_j$ 变成 $u_j+1$。将此代入上述平坦截面，并把 $\Phi$ 的拓扑部分的偏导数按 $u_j$ 的幂展开，即可计算 $\mathcal T_j$ 的作用。取对数得到
\begin{equation}\label{eq:a-model-monodromy-flat-basis}
\begin{aligned}
N_j(s^0)&=0,\\
N_j(s^k)&=-\delta_{jk}s^0,&&1\le k\le r,\\
N_j(s_k)&=-\sum_l\left.\frac{\partial^3(\Phi/(2\pi i)^3)}{\partial u_j\partial u_k\partial u_l}\right|_{q=0}s^l\\
&=-\sum_l\left(\int_VT_j\cup T_k\cup T_l\right)s^l,&&1\le k\le r,\\
N_j(s_0)&=-s_j.
\end{aligned}
\end{equation}
这些公式中的系数，显然恰为以 $-T_j$ 与给定上同调基作杯积所得矩阵的各项。

知道 $N_j$ 对 $\nabla$ 平坦截面的作用以后，便可计算施加 $\exp(-\sum_ju_jN_j)$ 的结果。这样得到下列 $\nabla^c$ 平坦截面：
\begin{equation}\label{eq:a-model-canonical-sections-threefold}
\begin{aligned}
\widetilde s^0&=T^0,\\
\widetilde s^k&=T^k,&&1\le k\le r,\\
\widetilde s_k&=T_k-\frac1{(2\pi i)^3}\sum_l\frac{\partial^2\Phi_{\mathrm{hol}}}{\partial u_k\partial u_l}T^l
+\frac1{(2\pi i)^3}\frac{\partial\Phi_{\mathrm{hol}}}{\partial u_k}T^0,&&1\le k\le r,\\
\widetilde s_0&=T_0-\frac1{(2\pi i)^3}\sum_l\frac{\partial\Phi_{\mathrm{hol}}}{\partial u_l}T^l
+\frac2{(2\pi i)^3}\Phi_{\mathrm{hol}}T^0.
\end{aligned}
\end{equation}
\SourcePage{272}
这些 $\nabla^c$ 平坦截面显然具有所需形式。在 $0\in\Delta^r$ 处也具有所需的值，因为 $\Phi_{\mathrm{hol}}$ 及其关于 $\partial/\partial u_k=2\pi iq_k\partial/\partial q_k$ 的各阶导数，都在 $0$ 处消失。最后，算子 $\exp(-\sum_ju_jN_j)$ 与 $N_j$ 交换，所以 $N_j$ 相对于式~\eqref{eq:a-model-flat-sections-threefold} 和~\eqref{eq:a-model-canonical-sections-threefold} 的矩阵相同。由式~\eqref{eq:a-model-monodromy-flat-basis}，$N_j$ 对式~\eqref{eq:a-model-canonical-sections-threefold} 的作用具有所需矩阵。

这就证明了 Calabi--Yau 三维簇偶数次上同调这一特殊情形。当 $V$ 的维数为任意 $d\ge3$ 时，需要使用\hyperref[chap:quantum-differential-equations]{第十章}\hyperref[subsec:a-model-gravitational-flat-sections]{第~\ref*{subsec:a-model-gravitational-flat-sections} 小节}对 $\nabla$ 平坦截面的描述。完整证明将在推论~\ref{cor:a-model-canonical-section-triangular} 中给出，故暂时推迟。
\end{proof}
此命题表明，在 $\Delta^r$ 上，可自然地将典范延拓 $\overline{\mathcal H}$ 等同于平凡丛 $H^*(V,\mathbb C)\times\Delta^r$。特别地，$0\in\Delta^r$ 上的纤维是 $H^*(V,\mathbb C)$。利用这一点，现在可以说明单值性对数 $N_j$ 的作用。由命题~\ref{prop:a-model-canonical-flat-basis}，$N_j$ 对 $\nabla^c$ 平坦基 $\widetilde s_k$ 的作用，与以 $-T_j$ 与 $T_k$ 作杯积相同。因为 $\widetilde s_k(0)=T_k$，在 $0$ 上的纤维中就有
\begin{equation}\label{eq:a-model-monodromy-boundary-cup}
N_j(\alpha)=-T_j\cup\alpha,\qquad\alpha\in H^*(V,\mathbb C).
\end{equation}
不过，离开 $0\in\Delta^r$ 后，此式可能不再成立。例如，当 $\dim V=3$ 时，命题~\ref{prop:a-model-canonical-flat-basis} 证明中的公式意味着，若把 $T_k$（$1\le k\le r$）看作 $\overline{\mathcal H}$ 的截面，则
\[
N_j(T_k)=-T_j\cup T_k-\frac1{(2\pi i)^3}\frac{\partial^2\Phi_{\mathrm{hol}}}{\partial u_j\partial u_k}T^0.
\]
在 $\Delta^r$ 的大多数点上，此式与 $-T_j\cup T_k$ 不同。

也可用命题~\ref{prop:a-model-canonical-flat-basis} 给 $\mathcal H$ 赋予整数结构。先取 $H^*(V,\mathbb C)$ 上的自然整数结构，并将其看作 $0\in\Delta^r$ 上的纤维中的整数结构。若 $\widetilde s=\exp(-\sum_ju_jN_j)s$ 在 $0$ 处的值是 $H^*(V,\mathbb C)$ 中的整类，就称 $\nabla$ 平坦截面 $s$ 为整截面。这些整平坦截面在 $(\Delta^*)^r$ 上形成局部常值层 $\mathcal H_{\mathbb Z}$。与 $\mathbb R$ 作张量积得到实向量空间的局部常值层 $\mathcal H_{\mathbb R}$。特别地，$\mathcal H$ 上有复共轭作用。

一个有趣的结果是，若把 $T_k$ 看作 $\mathcal H$ 的截面，那么在大多数纤维中，它既不是整的，也不是实的。例如，当 $\dim V=3$ 时，由所选上同调基可验证，式~\eqref{eq:a-model-flat-sections-threefold} 定义的 $\nabla$ 平坦截面都是整的。但这些公式也给出，例如
\[
T^k=s^k+u_ks^0,\qquad1\le k\le r.
\]
虽然 $s^k$ 和 $s^0$ 都整，$T^k$ 通常连实的都不是，更不用说整的。

还要注意，$\mathcal H_{\mathbb Z}$ 由 $\mathcal H$ 在 $0\in\Delta^r$ 上的行为决定。因此，$\mathcal H$ 的整数结构可能依赖大半径极限点的选择。我们猜想不会出现这种情形，即所有这些点都给出同一个 $\mathcal H_{\mathbb Z}$。只要作出保证 $\mathcal H_{\mathbb Z}$ 存在所需的适当收敛假设，证明应当不难。

既然已知单值性，也就可以确定 A 模型联络的单值性权滤过。
\begin{proposition}\label{prop:a-model-weight-filtration}\index{单值性权滤过|hyperpage}
对于维数 $d\ge3$ 的 Calabi--Yau 流形 $V$，$\nabla$ 在 $0\in\Delta^r$ 处的单值性权滤过 $W_\bullet$ 为
\[
W_i=\bigoplus_{j\ge2d-i}H^j(V,\mathbb C).
\]
\end{proposition}
\SourcePage{273}
\begin{proof}
滤过 $W_\bullet$ 由 $N=\sum_{j=1}^ra_jN_j$ 决定，其中 $N_j=\log\mathcal T_j$ 且 $a_j>0$。由式~\eqref{eq:a-model-monodromy-boundary-cup}，$N_j$ 在 $0$ 上的纤维中是与 $-T_j$ 作杯积的算子。因此，除去一个负号，$N$ 就是在此纤维中与 Kähler 类 $\omega=\sum_{j=1}^ra_jT_j$ 作杯积。由强 Lefschetz 定理，
\[
\cup\omega^k:H^{d-k}(V,\mathbb C)\simeq H^{d+k}(V,\mathbb C)
\]
是同构。因此，\SourceCite{Griffiths2} 中计算权滤过的过程，容易给出上述 $W_i$ 的描述。实际上，刻画单值性权滤过的一个关键性质，是同构 $N^k:W_{d+k}/W_{d+k-1}\simeq W_{d-k}/W_{d-k-1}$，这里正是强 Lefschetz 定理。
\end{proof}
$0$ 处的权滤过 $W_\bullet$ 可借助 $\nabla^c$ 传播为丛 $\overline{\mathcal H}$ 的滤过 $\mathcal W_\bullet$。这样得到如下结果。
\begin{corollary}\label{cor:a-model-weight-bundles}
$\nabla$ 在 $\Delta^r$ 上的单值性权滤过 $\mathcal W_\bullet$ 由平凡丛
\[
\mathcal W_i=\bigoplus_{j\ge2d-i}H^j(V,\mathbb C)\times\Delta^r
\]
给出。
\end{corollary}
\begin{proof}
在 $0$ 处，$W_i$ 由满足 $\deg T_k\ge2d-i$ 的 $T_k$ 决定。因此，在 $\Delta^r$ 上，$\mathcal W_i$ 由这些 $T_k$ 所对应的 $\widetilde s_k$ 决定（记号同命题~\ref{prop:a-model-canonical-flat-basis}）。因为 $\widetilde s_k=T_k+\text{次数更高的项}$，推论随即得到。
\end{proof}
此外，利用命题~\ref{prop:a-model-weight-filtration}，还可以按如下方式刻画 $\nabla$ 的单值性。
\begin{corollary}\label{cor:a-model-maximal-unipotent}\index{最大幺幂单值性!A 模型|hyperpage}
A 模型联络 $\nabla$ 在 $0\in\Delta^r$ 处\IndexLocation{indexloc:0010:273}具有最大幺幂单值性。
\end{corollary}
\begin{proof}
已知单值性幺幂，只须检验定义~\ref{def:maximal-unipotent-point} 的条件 (ii) 和 (iii)。因为 $H^1(V,\mathbb C)=0$，上一命题给出 $\dim W_0=\dim W_1=\dim H^{2d}(V,\mathbb C)=1$，以及 $\dim W_2=\dim(H^{2d}(V,\mathbb C)\oplus H^{2d-2}(V,\mathbb C))=1+r$。此外，若 $g_0\in H^{2d}(V,\mathbb C)$ 是一个点的 Poincaré 对偶，而 $g_1,\ldots,g_r\in H^{2d-2}(V,\mathbb C)$ 在杯积下与 $T_1,\ldots,T_r\in H^2(V,\mathbb C)$ 对偶，则在 $0\in\Delta^r$ 上的纤维中，
\[
N_i(g_j)=-T_i\cup g_j=-\left(\int_VT_i\cup g_j\right)g_0=-\delta_{ij}g_0,
\]
其中第一个等号来自式~\eqref{eq:a-model-monodromy-boundary-cup}。矩阵 $(\delta_{ij})$ 显然可逆，故单值性最大幺幂。
\end{proof}
这个推论有助于说明，为什么\hyperref[chap:hodge-yukawa]{第~\ref*{chap:hodge-yukawa} 章}如此强调最大幺幂单值性。在\IndexLocation{indexloc:0830:273}镜像对称下，A 模型联络应对应于描述镜像复结构变分的 Gauss--Manin 联络（精确表述见\hyperref[sec:mirror-symmetry-revisited]{第~\ref*{sec:mirror-symmetry-revisited} 节}）。A 模型联络具有良好的边界点，其单值性最大幺幂，故镜像的 Gauss--Manin 联络也应具有相同性质。

A 模型联络还满足一个与\hyperref[subsec:mirror-map-definition]{第~\ref*{subsec:mirror-map-definition} 小节}所讨论的整性猜想相类似的命题。为说明这一点，注意，在推论~\ref{cor:a-model-maximal-unipotent} 的证明中，若选取 $T_1,\ldots,T_r$ 为 $H^2(V,\mathbb Z)$ 的一组基，则 $g_0,g_1,\ldots,g_r$ 是 $H^{2d}(V,\mathbb Z)\oplus H^{2d-2}(V,\mathbb Z)$ 的一组基。由 $N(g_j)=-\delta_{ij}g_0$，整性猜想随即成立。不难预料，A 模型联络的这一行为也是镜像一侧整性猜想的动机之一。
\SourcePage{274}
关于 A 模型联络的最后一点说明，涉及 $H^*(V)$ 的基 $T_0,\ldots,T_m$。这些 $T_j$ 可以看作典范延拓 $\overline{\mathcal H}$ 的截面，但它们对于 $\nabla$ 和 $\nabla^c$ 都不平坦，而且除 $0\in\Delta^r$ 处以外也不整。不过，若考察它们在权滤过 $\mathcal W_\bullet$ 的分次部分上的行为，就会看到另一种情形。
\begin{corollary}\label{cor:a-model-graded-integral-basis}
若 $[T_j]$ 表示在 $\operatorname{Gr}^{\mathcal W}$ 上诱导的截面，则这些 $[T_j]$ 构成 $\operatorname{Gr}^{\mathcal W}$ 的典范延拓的一组整平坦截面基。
\end{corollary}
\begin{proof}
推论~\ref{cor:a-model-weight-bundles} 和 $\nabla$ 的定义容易推出，$[T_j]$ 对于 $\operatorname{Gr}^{\mathcal W}$ 上的诱导联络平坦。还可验证，单值性在 $\operatorname{Gr}^{\mathcal W}$ 上平凡，故 $\nabla$ 和 $\nabla^c$ 诱导同一个联络。结合 $T_j$ 在 $0\in\Delta^r$ 上的整数性，就知 $[T_j]$ 整。它们显然构成一组基，证明完成。
\end{proof}
此推论有趣之处如下。\hyperref[subsec:normalized-yukawa-connection]{第~\ref*{subsec:normalized-yukawa-connection} 小节}研究了 Calabi--Yau 三维簇族在最大幺幂边界点处的 Hodge 结构变分。回顾命题~\ref{prop:normalized-yukawa-pf} 的证明，会看到我们先取 $\operatorname{Gr}^{\mathcal W}$ 的一组整平坦截面基，再提升为 $\mathcal H$ 的一组基。由推论~\ref{cor:a-model-graded-integral-basis}，在 A 模型情形，这正是 $T_j$ 所起的作用。

最后两个推论表明，A 模型联络与 Calabi--Yau 三维簇族的 Gauss--Manin 联络十分相似。定义 A 模型 Hodge 滤过以后，这种相似性还将进一步加深；这就是接下来要讨论的主题。

\subsection{A 模型 Hodge 滤过}\label{subsec:a-model-hodge-filtration}
接下来要把 A 模型联络组成一个 Hodge 结构变分。为此，需要一个起 Hodge 滤过作用的对象。照例假设 $V$ 是维数 $\dim V=d\ge3$ 的 Calabi--Yau 流形，并令
\begin{equation}\label{eq:a-model-hodge-filtration}\index{Hodge 滤过!A 模型|hyperpage}
F^p=\bigoplus_{a\le d-p}H^{a,b}(V).
\end{equation}
这给出递减滤过 $H^*(V,\mathbb C)=F^0\supset F^1\supset\cdots$。在丛 $\mathcal H=H^*(V,\mathbb C)\times K_{\mathbb C}(V)$ 中，子丛
\[
\mathcal F^p=F^p\times K_{\mathbb C}(V)\subset\mathcal H
\]
确定 $\mathcal F^0=\mathcal H$ 的一个滤过。

我们的基本断言是，滤过 $\mathcal F^\bullet$ 与 A 模型联络 $\nabla$ 组成权为 $d$ 的 Hodge 结构变分，至少在复化 Kähler 空间 $K_{\mathbb C}(V)$ 的某个大半径极限点附近如此。为作出更精确的表述，回顾整类 $T_1,\ldots,T_r$ 生成的锥的内部 $\operatorname{Int}(\sigma)\subset K(V)$。这给出 $(\Delta^*)^r\subset K_{\mathbb C}(V)$，所讨论的大半径极限点就是 $0\in\Delta^r$。

\hyperref[subsec:a-model-asymptotics]{第~\ref*{subsec:a-model-asymptotics} 小节}在 $\mathcal H$ 上定义了整数结构 $\mathcal H_{\mathbb Z}$，从而特别给出了 $\mathcal H$ 上的复共轭作用。于是，证明 $(\mathcal H,\nabla,\mathcal H_{\mathbb Z},\mathcal F^\bullet)$ 是 Hodge 结构变分，归结为证明以下两个性质：
\begin{itemize}
\item 对每个 $p$，$\nabla(\mathcal F^p)\subset\mathcal F^{p-1}\otimes\Omega^1_{K_{\mathbb C}(V)}$。
\item 对每个 $p$，$\mathcal F^p\oplus\overline{\mathcal F^{d-p+1}}=\mathcal H$。
\end{itemize}
\SourcePage{275}
第一个性质容易证明，因为由\hyperref[cor:cy-quantum-hodge]{推论~\ref*{cor:cy-quantum-hodge}}，小量子乘积与 Hodge 分解相容。论证如下。给定 $\alpha\in H^{a,b}(V)$，其中 $a\le d-p$，可把 $\alpha$ 看作 $\mathcal F^p$ 的截面。因为 $T_i\in H^{1,1}(V)$，有
\[
\nabla_{\frac{\partial}{\partial u_i}}(\alpha)=T_i*_{\mathrm{small}}\alpha\in H^{a+1,b+1}(V).
\]
所以 $\nabla_{\partial/\partial u_i}(\alpha)$ 是 $\mathcal F^{p-1}$ 的截面。这就证明了 $\nabla$ 的 Griffiths 横截性。

不过，证明第二个性质 $\mathcal F^p\oplus\overline{\mathcal F^{d-p+1}}=\mathcal H$ 需要更多工作。为此，引入如下配对：若 $\alpha\in H^k(V,\mathbb C)$、$\beta\in H^l(V,\mathbb C)$，则
\begin{equation}\label{eq:a-model-polarizing-form}
S(\alpha,\beta)=\begin{cases}
(-1)^{k(k+1)/2}\displaystyle\int_V\alpha\cup\beta,&k+l=2d,\\
0,&k+l\ne2d.
\end{cases}
\end{equation}
若把 $\alpha$、$\beta$ 看作 $\mathcal H$ 的截面，上式仍有意义，因此可把 $S$ 看作定义在 $\mathcal H$ 上的配对。

配对 $S$ 具有下列有用性质。
\begin{lemma}\label{lem:a-model-pairing-properties}
设 $S$ 为上述 $\mathcal H$ 上的双线性型，则：
\begin{enumerate}[label=(\roman*)]
\item $S$ 的奇偶性为 $(-1)^d$，即 $S(\alpha,\beta)=(-1)^dS(\beta,\alpha)$。
\item $S$ 关于 $\nabla$ 平坦，即对 $\mathcal H$ 的任意截面 $\alpha,\beta$，以及所有 $1\le i\le r$，有
\[
\frac{\partial}{\partial u_i}S(\alpha,\beta)
=S\left(\nabla_{\frac{\partial}{\partial u_i}}\alpha,\beta\right)
+S\left(\alpha,\nabla_{\frac{\partial}{\partial u_i}}\beta\right).
\]
\item 单值性对数 $N_j$ 是 $S$ 的无穷小自同构，即对任意 $\alpha,\beta\in H^*(V,\mathbb C)$，有
\[
S(N_j(\alpha),\beta)+S(\alpha,N_j(\beta))=0.
\]
\item $S$ 在 $\mathcal H_{\mathbb Z}$ 上取整数值。
\item 对所有 $p$，$S(\mathcal F^p,\mathcal F^{d-p+1})=0$。
\end{enumerate}
\end{lemma}
\begin{proof}
对于第一部分，设 $\alpha\in H^k(V,\mathbb C)$、$\beta\in H^l(V,\mathbb C)$ 且 $k+l=2d$。则 $\alpha\cup\beta=(-1)^{kl}\beta\cup\alpha$，由 $S$ 的定义容易得到所需等式 $S(\alpha,\beta)=(-1)^dS(\beta,\alpha)$。

为证明第二部分，可设 $\alpha=T_j$、$\beta=T_l$。用 $\nabla_{\partial/\partial u_i}$ 与 $T_i$ 的量子乘积之间的定义关系，问题归结为证明
\[
S(T_i*_{\mathrm{small}}T_j,T_l)+S(T_j,T_i*_{\mathrm{small}}T_l)=0.
\]
可设 $\deg T_j=k$、$\deg T_l=2d-k-2$，否则上式平凡成立。于是 $\deg(T_i*_{\mathrm{small}}T_j)=k+2$，故
\[
\begin{aligned}
S(T_i*_{\mathrm{small}}T_j,T_l)
&=(-1)^{(k+2)(k+3)/2}\int_V(T_i*_{\mathrm{small}}T_j)\cup T_l\\
&=-(-1)^{k(k+1)/2}\int_V(T_i*_{\mathrm{small}}T_j)*_{\mathrm{small}}T_l,
\end{aligned}
\]
其中第二个等号使用了命题~\ref{prop:quantum-frobenius-pairing} 给出的 $\int_V\alpha*_{\mathrm{small}}\beta=\int_V\alpha\cup\beta$。类似地，
\[
S(T_j,T_i*_{\mathrm{small}}T_l)=(-1)^{k(k+1)/2}\int_VT_j*_{\mathrm{small}}(T_i*_{\mathrm{small}}T_l).
\]
\SourcePage{276}
因为 $T_i$ 的次数为偶数，$T_i*_{\mathrm{small}}T_j=T_j*_{\mathrm{small}}T_i$。结合小量子乘积的结合律，便得到所需恒等式。

对于第三部分，注意 (ii) 意味着对任意 $a,b\in H^*(V,\mathbb C)$，有 $S(\mathcal T_j(a),\mathcal T_j(b))=S(a,b)$，因为 $\mathcal T_j(a)$ 和 $\mathcal T_j(b)$ 分别由 $a$、$b$ 经平行移动得到。幺幂的保双线性型变换的对数，是无穷小的保双线性型变换，所需结论由此得到。

为证明 (iv)，令 $s_1,s_2$ 为整 $\nabla$ 平坦截面。由 (ii)，$S(s_1,s_2)$ 为常数。令 $\widetilde s_i=\exp(-\sum_ju_jN_j)s_i$；由 (iii)，$S(\widetilde s_1,\widetilde s_2)=S(s_1,s_2)$ 也为常数，所以 $S(s_1,s_2)=S(\widetilde s_1(0),\widetilde s_2(0))$。由 $\mathcal H_{\mathbb Z}$ 的定义，$\widetilde s_i(0)$ 是整类，故 $S(s_1,s_2)\in\mathbb Z$。

最后，设 $\alpha\in H^{a,b}(V)$、$a\le d-p$，以及 $\beta\in H^{e,f}(V)$、$e\le p-1=d-(d-p+1)$。则 $\alpha$、$\beta$ 分别给出 $\mathcal F^p$、$\mathcal F^{d-p+1}$ 的截面。$\alpha\cup\beta$ 的 Hodge 型为 $(a+e,b+f)$，而 $a+e\le d-1$，所以它不是 $(d,d)$ 型，$S(\alpha,\beta)=0$ 随即得到。
\end{proof}
证明 $\nabla$ 和 $\mathcal F^\bullet$ 给出 Hodge 结构变分的关键步骤之一，是证明 $S$ 极化某个混合 Hodge 结构。所讨论的结构来自式~\eqref{eq:a-model-hodge-filtration} 定义的滤过 $F^\bullet$，以及命题~\ref{prop:a-model-weight-filtration} 给出的权滤过。容易验证
\begin{equation}\label{eq:a-model-mixed-hodge-pieces}
F^p\cap\overline{F^{k-p}}\cap W_k=H^{d-p,d-k+p}(V).
\end{equation}
因为 $\operatorname{Gr}^W_k=H^{2d-k}(V,\mathbb C)$，可给 $\operatorname{Gr}^W_k$ 赋予权为 $k$ 的纯 Hodge 结构，使 $(\operatorname{Gr}^W_k)^{p,k-p}=H^{d-p,d-k+p}(V)$。因此，$(W_\bullet,F^\bullet)$ 是混合 Hodge 结构。

还要回顾，在 $0\in\Delta^r$ 处，单值性对数 $N_j$ 是与 $-T_j$ 作杯积的算子。因此，$N=\sum_ja_jN_j$（$a_j>0$）是与 $-\omega$ 作杯积的算子，其中 $\omega=\sum_ja_jT_j$ 为 Kähler 类。命题~\ref{prop:a-model-weight-filtration} 的证明已说明，$W_\bullet$ 是 $N$ 的单值性权滤过。于是有 \SourceCite{CKS2} 中的下列观察。
\begin{proposition}\label{prop:a-model-polarized-mixed-hodge}\index{混合 Hodge 结构!极化|hyperpage}
对任意 $N=\sum_ja_jN_j$，若所有 $a_j>0$，则混合 Hodge 结构 $(W_\bullet,F^\bullet)$ 由 $-N$ 和 $S$ 极化。
\end{proposition}
\begin{proof}
先回顾 \SourceCite{CaK} 中 $(W_\bullet,F^\bullet)$ 被 $-N$ 和 $S$ 极化的含义。定义首先要求 $W_\bullet$ 是 $N$ 的单值性权滤过（已在上文说明），$N(F^p)\subset F^{p-1}$（由定义显然成立），且 $N$ 是 $S$ 的无穷小自同构（由引理~\ref{lem:a-model-pairing-properties} 的证明得到）。此外，若 $k\ge d$，“本原部分” $P_k\subset\operatorname{Gr}^W_k$ 定义为
\[
P_k=\ker\bigl(N^{k-d+1}:\operatorname{Gr}^W_k\longrightarrow\operatorname{Gr}^W_{2d-k-2}\bigr).
\]
极化混合 Hodge 结构定义的最后一项要求，$P_k$ 上的纯 Hodge 结构由 $S(\,\cdot\,,(-N)^{k-d}(\,\cdot\,))$ 极化（这里的负号源于我们的 $N$ 是 \SourceCite{CaK} 所用 $N$ 的相反数）。因此，只须证明，对任意非零 $\alpha\in P_k^{p,k-p}$，有
\begin{equation}\label{eq:a-model-primitive-positivity}
S\bigl(i^{p-(k-p)}\alpha,(-N)^{k-d}(\overline\alpha)\bigr)>0.
\end{equation}
关键在于，以上安排使此式成为本原上同调的通常极化性质。为说明这一点，回顾 $\operatorname{Gr}^W_k=H^{2d-k}(V,\mathbb C)$，且 $(\operatorname{Gr}^W_k)^{p,k-p}=H^{d-p,d-k+p}(V)$。
\SourcePage{277}
与命题~\ref{prop:a-model-weight-filtration} 的证明一样，因 $N$ 是与 $-\omega$ 作杯积的算子，容易看出 $P_k\subset\operatorname{Gr}^W_k$ 恰为关于 $\omega$ 的通常本原上同调 $H^{2d-k}_0(V,\mathbb C)$。因此，若 $\alpha\in P_k^{p,k-p}=H^{d-p,d-k+p}_0(V)$ 非零，通常的 Hodge--Riemann 双线性关系~\eqref{eq:hodge-polarization} 给出
\[
(-1)^{(2d-k)(2d-k-1)/2}\int_V\omega^{d-(2d-k)}\wedge i^{(d-p)-(d-k+p)}\alpha\wedge\overline\alpha>0.
\]
利用 $S$ 的定义，经过简单代数运算，此式化为式~\eqref{eq:a-model-primitive-positivity}。
\end{proof}
现在终于可以证明，我们得到了 Hodge 结构变分。
\begin{theorem}\label{thm:a-model-hodge-variation}\index{极化!A 模型|hyperpage}
设 $V$ 为维数 $d\ge3$ 的 Calabi--Yau 流形。则存在 $0\in\Delta^r$ 的邻域 $U$，使 $(\mathcal H,\nabla,\mathcal H_{\mathbb Z},\mathcal F^\bullet)$ 在 $U\cap(\Delta^*)^r$ 上是由双线性型 $S$ 极化的权 $d$ Hodge 结构变分。
\end{theorem}
\begin{proof}
证明思路与定义 $\mathcal H_{\mathbb Z}$ 时相似：当时用 $0\in\Delta^r$ 上的整数结构确定 $\mathcal H$ 的整数结构。这里用命题~\ref{prop:a-model-polarized-mixed-hodge} 在 $0\in\Delta^r$ 上给出的极化混合 Hodge 结构，证明 $(\mathcal H,\nabla,\mathcal H_{\mathbb Z},\mathcal F^\bullet)$ 在 $0$ 附近是 Hodge 结构变分。

证明将使用 \SourceCite{CaK,CKS1} 的强有力工具，并沿用其中的记号。我们感谢 \SourceCite{CaK} 的作者解释了他们的工作与此情形的联系。令 $\check D$ 为分类 $H^*(V,\mathbb C)$ 上滤过 $\widetilde F^\bullet$ 的旗流形，这些滤过对所有 $p$ 都满足 $\dim\widetilde F^p=\sum_{a\le d-p}h^{a,b}(V)$ 和 $S(\widetilde F^p,\widetilde F^{d-p+1})=0$。在 $\check D$ 内有开集 $D$，分类那些在 $H^*(V,\mathbb C)$ 上诱导权 $d$ 极化 Hodge 结构的滤过。

式~\eqref{eq:a-model-hodge-filtration} 定义的 $F^\bullet$ 位于 $\check D$ 中，因为引理~\ref{lem:a-model-pairing-properties} 给出 $S(F^p,F^{d-p+1})=0$。再考虑滤过族 $\exp(-\sum_ju_jN_j)F^\bullet$；因为 $N_j$ 是 $S$ 的无穷小自同构，这个族也位于 $\check D$ 中。按 \SourceCite{CaK} 的定义，只要满足
\begin{itemize}
\item $N_j(F^p)\subset F^{p-1}$；
\item 当 $\operatorname{Im}(u_j)\gg0$ 时，$\exp(-\sum_ju_jN_j)F^\bullet\in D$，
\end{itemize}
这个族就是幂零轨道。

第一条已在滤过 $F^\bullet$ 的定义~\eqref{eq:a-model-hodge-filtration} 后的讨论中证明。对于第二条，令 $N=\sum_ja_jN_j$，其中 $a_j>0$。则 $N$ 给出命题~\ref{prop:a-model-weight-filtration} 的权滤过 $W_\bullet$，而由命题~\ref{prop:a-model-polarized-mixed-hodge}，$(W_\bullet,F^\bullet)$ 被 $-N$ 和 $S$ 极化。由 \SourceCite[命题 (4.66)]{CKS1}，$\exp(-\sum_ju_jN_j)F^\bullet$ 是幂零轨道。（应当说明，\SourceCite{CaK,CKS1} 将幂零轨道写成 $\exp(\sum_ju_jN_j)F^\bullet$，因为他们的 $N_j$ 是我们的 $-N_j$。）

更具体地说，刚定义的幂零轨道给出 $\mathcal H$ 的滤过 $\widetilde{\mathcal F}^\bullet$。\SourceCite[第 74 页]{CaK} 阐述的关键思想如下：把 $H^*(V,\mathbb C)$ 的滤过 $F^\bullet$ 看作位于\IndexLocation{indexloc:0005:277}典范延拓 $\overline{\mathcal H}$ 在 $0\in\Delta^r$ 上的纤维中（这里沿用\hyperref[subsec:a-model-asymptotics]{第~\ref*{subsec:a-model-asymptotics} 小节}的记号）。用 $\overline{\mathcal H}$ 的 $\nabla^c$ 平坦截面传播 $F^p$，得到子丛 $\widetilde{\mathcal F}^p$，其中 $\nabla^c$ 由式~\eqref{eq:a-model-canonical-connection} 定义。换言之，$\nabla^c$ 平坦截面 $\widetilde s$ 属于 $\widetilde{\mathcal F}^p$，当且仅当 $\widetilde s(0)\in F^p$。

因此，在 $(\Delta^*)^r$ 中充分靠近 $0$ 的地方，$(\mathcal H,\nabla,\mathcal H_{\mathbb Z},\widetilde{\mathcal F}^\bullet)$ 是来自幂零轨道的 Hodge 结构变分。另一方面，滤过 $\mathcal F^\bullet$ 满足 Griffiths 横截性，并在 $0\in\Delta^r$ 处与 $\widetilde{\mathcal F}^\bullet$ 一致。正如 Deligne 在 \SourceCite[第 2.3 节]{Deligne2} 中指出的，这些条件与 \SourceCite[定理 (2.8)]{CaK} 一起，意味着 $(\mathcal H,\nabla,\mathcal H_{\mathbb Z},\mathcal F^\bullet)$ 在 $(\Delta^*)^r$ 中 $0$ 的某个邻域内是 Hodge 结构变分。定理得证。
\end{proof}
\SourcePage{278}
按 \SourceCite{Morrison7} 的术语，称 $(\mathcal H,\nabla,\mathcal H_{\mathbb Z},\mathcal F^\bullet)$ 为 \emph{A 模型 Hodge 结构变分}。由我们的构造方式，直接可证，A 模型变分的极限 Hodge 结构恰为命题~\ref{prop:a-model-polarized-mixed-hodge} 中的混合 Hodge 结构 $(W_\bullet,F^\bullet)$。\index{A 模型 Hodge 结构变分|hyperpage}

为研究镜像对称，需要考虑 A 模型 Hodge 结构变分的一个有趣子变分。令 $\mathcal H^{\mathrm{middle}}$ 为 $\bigoplus_{p=0}^dH^{p,p}(V)$ 给出的丛；这是 $H^*(V,\mathbb C)$ 的 Hodge 分解的中间部分。由\hyperref[cor:cy-quantum-hodge]{推论~\ref*{cor:cy-quantum-hodge}}，小量子乘积与 Hodge 分解相容，故 $\nabla$ 保持 $\mathcal H^{\mathrm{middle}}$。将 $\nabla$ 在其上的限制记为 $\nabla^{\mathrm{middle}}$。还要注意，$\bigoplus_{p=0}^dH^{p,p}(V)$ 具有自然的实结构；通过\hyperref[subsec:a-model-asymptotics]{第~\ref*{subsec:a-model-asymptotics} 小节}的过程，得到实结构 $\mathcal H^{\mathrm{middle}}_{\mathbb R}$。最后，略微滥用记号，仍用 $\mathcal F^\bullet$ 表示 A 模型变分在 $\mathcal H^{\mathrm{middle}}$ 上诱导的滤过。

由这些定义以及定理~\ref{thm:a-model-hodge-variation}，容易得到实 Hodge 结构变分
\begin{equation}\label{eq:middle-a-model-variation}
(\mathcal H^{\mathrm{middle}},\nabla^{\mathrm{middle}},\mathcal H^{\mathrm{middle}}_{\mathbb R},\mathcal F^\bullet).
\end{equation}
这是\emph{中间 A 模型 Hodge 结构变分}，不过通常也将式~\eqref{eq:middle-a-model-variation} 简称为 A 模型变分。上述所有结果都适用于它。特别地，它由交配对
\begin{equation}\label{eq:middle-a-model-polarization}
S(\alpha,\beta)=(-1)^p\int_V\alpha\cup\beta,
\qquad\alpha\in H^{p,p}(V),\quad\beta\in H^{d-p,d-p}(V)
\end{equation}
极化。这是将式~\eqref{eq:a-model-polarizing-form} 限制到 $\bigoplus_{p=0}^dH^{p,p}(V)$ 的结果。还应说明，当 $d=3$ 时，$\mathcal H^{\mathrm{middle}}$ 具有自然的整数结构，因为这时 $H^{1,1}(V)=H^2(V,\mathbb C)$、$H^{2,2}(V)=H^4(V,\mathbb C)$。因此，三维簇情形可将式~\eqref{eq:middle-a-model-variation} 替换为 $(\mathcal H^{\mathrm{middle}},\nabla^{\mathrm{middle}},\mathcal H^{\mathrm{middle}}_{\mathbb Z},\mathcal F^\bullet)$。

还可描述式~\eqref{eq:middle-a-model-variation} 在 $0\in\Delta^r$ 处的行为。由推论~\ref{cor:a-model-maximal-unipotent}，这里的单值性最大幺幂；$0\in\Delta^r$ 上的极限混合 Hodge 结构为 $(W_\bullet,F^\bullet)$，其中
\[
W_k=\bigoplus_{2j\ge2d-k}H^{j,j}(V),\qquad
F^p=\bigoplus_{j\le d-p}H^{j,j}(V).
\]
这是分别限制命题~\ref{prop:a-model-weight-filtration} 和式~\eqref{eq:a-model-hodge-filtration} 对 $W_\bullet$、$F^\bullet$ 的描述所得的结果。此外，这些公式意味着，对所有 $i$，有 $W_{2i}=W_{2i+1}$；对所有 $p$，有 $F^p\oplus W_{2p-2}=\bigoplus_{j=0}^dH^{j,j}(V)$。它们恰为混合 Hodge 结构是 Hodge--Tate 结构的条件~\eqref{eq:mum-hodge-tate-splitting}。因此得到如下命题。
\begin{proposition}\label{prop:a-model-hodge-tate}\index{Hodge--Tate结构@Hodge--Tate 结构|hyperpage}
变分 $(\mathcal H^{\mathrm{middle}},\nabla^{\mathrm{middle}},\mathcal H^{\mathrm{middle}}_{\mathbb R},\mathcal F^\bullet)$ 的极限混合 Hodge 结构为 Hodge--Tate 结构。
\end{proposition}
本节最后讨论 Calabi--Yau 三维簇的三点函数及其与 A 模型联络的关系。
\SourceExampleNumber{8.5.4}{1}
\begin{example}\label{ex:a-model-yukawa-connection}
设 $V$ 是 Calabi--Yau 三维簇。\hyperref[subsec:small-quantum-definition]{第~\ref*{subsec:small-quantum-definition} 小节}对 $a,b,c\in H^2(V,\mathbb C)$ 定义了三点函数（也称 A \IndexLocation{indexloc:0188:278}模型\IndexLocation{indexloc:0195:278}关联函数）$\langle a,b,c\rangle$，并证明它可表示为
\[
\langle a,b,c\rangle=\int_Va*_{\mathrm{small}}b*_{\mathrm{small}}c
\]
\SourcePage{279}
（见式~\eqref{eq:a-three-point-function} 和~\eqref{eq:three-point-quantum-integral}）。我们将从两个方面讨论 $\langle a,b,c\rangle$ 与 A 模型联络 $\nabla^{\mathrm{middle}}$ 的关系，以下将后者简记为 $\nabla$。

先说明如何用 A 模型联络表示 $\langle a,b,c\rangle$。关键观察是，由 $\nabla$ 的定义，
\[
\nabla_{\frac{\partial}{\partial u_i}}(T_j*_{\mathrm{small}}T_k)
=T_i*_{\mathrm{small}}T_j*_{\mathrm{small}}T_k
+H^*(V,\mathbb C)\text{ 中次数较低的项}.
\]
因为 $1=T_0$，容易得到
\[
\int_V\nabla_{\frac{\partial}{\partial u_i}}\nabla_{\frac{\partial}{\partial u_j}}\nabla_{\frac{\partial}{\partial u_k}}1
=\int_VT_i*_{\mathrm{small}}T_j*_{\mathrm{small}}T_k=\langle T_i,T_j,T_k\rangle.
\]
此式可写成
\begin{equation}\label{eq:a-model-yukawa-derivatives}
\langle T_i,T_j,T_k\rangle
=\int_V1\cup\nabla_{\frac{\partial}{\partial u_i}}\nabla_{\frac{\partial}{\partial u_j}}\nabla_{\frac{\partial}{\partial u_k}}1,
\end{equation}
它与\hyperref[sec:yukawa-couplings]{第~\ref*{sec:yukawa-couplings} 节}定义的 B 模型关联函数，也就是 Yukawa 耦合的公式
\[
\int_V\Omega\wedge\nabla_{\theta_1}\nabla_{\theta_2}\nabla_{\theta_3}\Omega
\]
十分相似。在两种情形中，关联函数都以完全相同的方式，由联络（A 模型联络或 Gauss--Manin 联络）、$\mathcal F^3$ 的一个截面（$1$ 或 $\Omega$），以及交配对构造。很快会看到，这种相似性是镜像对称的重要组成部分。

第二个方面是，$\nabla$ 不仅如上所述确定 $\langle a,b,c\rangle$，反过来也成立。这来自命题~\ref{prop:a-model-canonical-flat-basis} 的证明。用式~\eqref{eq:a-model-potential} 给出的 Gromov--Witten 势 $\Phi$，将式~\eqref{eq:cy-potential-three-derivatives} 替换为
\begin{equation}\label{eq:a-model-yukawa-potential}
\frac1{(2\pi i)^3}\frac{\partial^3\Phi}{\partial u_i\partial u_j\partial u_k}=\langle T_i,T_j,T_k\rangle,
\end{equation}
因而式~\eqref{eq:a-model-threefold-connection} 可写为
\begin{equation}\label{eq:a-model-yukawa-full-connection}
\begin{aligned}
\nabla_{\frac{\partial}{\partial u_i}}(T^0)&=0,\\
\nabla_{\frac{\partial}{\partial u_i}}(T^k)&=\delta_{ik}T^0,&&1\le k\le r,\\
\nabla_{\frac{\partial}{\partial u_i}}(T_k)&=\sum_l\langle T_i,T_k,T_l\rangle T^l,&&1\le k\le r,\\
\nabla_{\frac{\partial}{\partial u_i}}(T_0)&=T_i.
\end{aligned}
\end{equation}
所以 $\nabla$ 完全由 A 模型关联函数 $\langle a,b,c\rangle$ 决定。

一个特殊情形是 $h^2(V)=1$。令 $H=T_1$ 生成 $H^2(V,\mathbb C)$，$C=T^1\in H^4(V,\mathbb C)$ 为对偶生成元，并令 $u=u_1$。推广例题~\ref{ex:cy-small-quantum} 的计算，得到
\[
H*_{\mathrm{small}}H=\langle H,H,H\rangle C.
\]
记三点函数为 $Y(q)=\langle H,H,H\rangle$，则上式成为
\begin{equation}\label{eq:one-parameter-quantum-yukawa}
H*_{\mathrm{small}}H=Y(q)C.
\end{equation}
\SourcePage{280}
由式~\eqref{eq:a-model-yukawa-full-connection}，A 模型联络的矩阵为
\[
\begin{pmatrix}0&0&0&0\\1&0&0&0\\0&Y(q)&0&0\\0&0&1&0\end{pmatrix}.
\]
这个矩阵~\eqref{eq:normalized-yukawa-matrix} 曾出现在命题~\ref{prop:normalized-yukawa-pf} 的证明中，当时研究的是复模空间一维的 Calabi--Yau 三维簇。特别地，该命题说明，在最大幺幂边界点处，只要 $Y$ 是归一化 Yukawa 耦合，并以镜像映射作为局部坐标，Gauss--Manin 联络就具有上述形式。与前面一样，这种相似性并非偶然，而是镜像猜想的一部分。

此例也略微预示了\hyperref[chap:quantum-differential-equations]{第十章}的内容。上述矩阵意味着，$1=T_0$ 满足方程
\[
\nabla_{\frac{\partial}{\partial u}}^2
\left(\frac1{Y(q)}\nabla_{\frac{\partial}{\partial u}}^2\right)1=0.
\]
按\hyperref[sec:quantum-cohomology-relations]{第~\ref*{sec:quantum-cohomology-relations} 节}的语言，这是一条\emph{量子微分方程}。在例题~\ref{ex:cy-quantum-equation-relation} 中将看到，它意味着 $V$ 的小量子上同调环中的关系
\[
H*_{\mathrm{small}}H*_{\mathrm{small}}\left(\frac{H*_{\mathrm{small}}H}{Y(q)}\right)=0.
\]
这个关系是平凡的，因为左端属于 $H^8(V)$，而该群显然为零。另一方面，从上述量子微分方程可提取 $Y(q)$（相差一个常数），所以量子微分方程包含了量子上同调的十分非平凡的信息。\hyperref[chap:quantum-differential-equations]{第十章}将证明，一般的量子微分方程同样给出量子上同调中的关系；还将通过例题说明，其中许多关系非平凡。
\end{example}

\section{镜像猜想}\label{sec:mirror-symmetry-revisited}
\hyperref[chap:quantum-cohomology]{第八章}最后要给出镜像对称的一个数学版本。为简便起见，主要考虑局部版本。其基本思想是，在复模空间与 Kähler 模空间中适当选取的边界点附近，\hyperref[chap:moduli]{第~\ref*{chap:moduli} 章}讨论的镜像映射应当把 $V$ 的 Gauss--Manin 联络变成镜像 $V^\circ$ 的 A 模型联络 $\nabla^{\mathrm{middle}}$。先在三维情形给出其精确含义。

\subsection{Calabi--Yau 三维簇的镜像对称}\label{subsec:threefold-mirror-symmetry}
设 $V$、$V^\circ$ 是光滑 Calabi--Yau 三维簇。与\hyperref[sec:mirror-map]{第~\ref*{sec:mirror-map} 节}一样，还假设 $V$ 的复模空间 $\mathcal M(V)$ 和 $V^\circ$ 的 Kähler 模空间 $\mathcal{KM}(V^\circ)$ 具有适当相容的紧化。也就是说，具有：
\begin{itemize}
\item $V$ 的复模空间的某个光滑紧化 $\overline{\mathcal M}(V)$ 上的最大幺幂边界点 $p_0$。
\item 镜像 $V^\circ$ 的 Kähler 模空间的相容光滑紧化 $\overline{\mathcal{KM}}(V^\circ)$ 上的\IndexLocation{indexloc:0696:280}大半径极限点 $q_0$。
\item 镜像映射~\eqref{eq:mirror-map-local}，将 $p_0\in\overline{\mathcal M}(V)$ 的某个邻域映到 $q_0\in\overline{\mathcal{KM}}(V^\circ)$ 的某个邻域。假设 $V$ 满足\hyperref[subsec:higher-dimensional-mum]{第~\ref*{subsec:higher-dimensional-mum} 小节}的整性猜想；由\hyperref[subsec:mirror-map-definition]{第~\ref*{subsec:mirror-map-definition} 小节}，这就唯一确定镜像映射。
\end{itemize}
\SourcePage{281}
如\hyperref[subsec:mirror-map-definition]{第~\ref*{subsec:mirror-map-definition} 小节}所述，将 $p_0$ 与 $q_0$ 配对以及确定镜像映射时，需要作出若干选择。以下假设这些选择已完成。

这里还有两个极化 Hodge 结构变分。在 $V$ 的复模空间上，由 $H^3(V,\mathbb C)$ 给出的丛 $\mathcal H^V$ 具有 Gauss--Manin 联络及 Hodge 滤过。为避免混淆，将 Gauss--Manin 联络记为 $\nabla^{\mathrm{GM}}$。在镜像一侧，$V^\circ$ 的 Kähler 模空间上，由 $H^{\mathrm{middle}}(V^\circ)=\bigoplus_{p=0}^3H^{p,p}(V^\circ)=\bigoplus_{p=0}^3H^{2p}(V^\circ,\mathbb C)$ 给出的丛 $\mathcal H^{V^\circ}$，具有 A 模型联络 $\nabla^{\mathrm{middle}}$ 及 Hodge 滤过。此外，$\mathcal H^V$ 和 $\mathcal H^{V^\circ}$ 分别有整数结构 $\mathcal H^V_{\mathbb Z}$、$\mathcal H^{V^\circ}_{\mathbb Z}$。最后，还有如下极化：
\begin{itemize}
\item 因为每个 $\alpha\in H^3(V,\mathbb C)$ 都本原，由式~\eqref{eq:hodge-polarization}，双线性型
\[
Q(\alpha,\beta)=(-1)^{3(3-1)/2}\int_V\alpha\cup\beta=-\int_V\alpha\cup\beta
\]
极化 $(\mathcal H^V,\nabla^{\mathrm{GM}},\mathcal H^V_{\mathbb Z},\mathcal F^\bullet)$。
\item 式~\eqref{eq:middle-a-model-polarization} 在 $\bigoplus_{p=0}^3H^{p,p}(V^\circ)$ 上定义了双线性型
\[
S(\alpha,\beta)=(-1)^p\int_{V^\circ}\alpha\cup\beta,\qquad
\alpha\in H^{p,p}(V^\circ),\quad\beta\in H^{3-p,3-p}(V^\circ).
\]
由定理~\ref{thm:a-model-hodge-variation}，$S$ 极化 $(\mathcal H^{V^\circ},\nabla^{\mathrm{middle}},\mathcal H^{V^\circ}_{\mathbb Z},\mathcal F^\bullet)$。
\end{itemize}
除此以外，还有几个特定截面：
\begin{itemize}
\item 对于 $V^\circ$，$\mathcal H^{V^\circ}$ 具有 $W_0$ 的整生成元 $[pt]$，其中 $[pt]\in H^{3,3}(V^\circ,\mathbb Z)$ 为一个点的类。还有特定截面 $1=T_0\in H^0(V^\circ,\mathbb C)$，由式~\eqref{eq:a-model-hodge-filtration}，它生成 $\mathcal F^3$。注意 $S(1,[pt])=1$。
\item 对于 $V$，丛 $\mathcal H^V$ 具有最大幺幂边界点定义中的截面 $g_0$。它是 $W_0$ 的整生成元，且在相差 $\pm1$ 的意义下唯一。此外，还有 $\mathcal F^3$ 的截面 $\Omega$。为得到典范的对象，令 $y_0=\langle g_0,\Omega\rangle$。一旦固定 $g_0$，$\Omega/y_0$ 就是满足 $Q(\Omega/y_0,g_0)=1$ 的唯一 $\mathcal F^3$ 截面。
\end{itemize}
于是终于可以如下精确定义 Calabi--Yau 三维簇的镜像对称。
\begin{definition}\label{def:mathematical-mirror-pair}
给定如上的 Calabi--Yau 三维簇 $V$、$V^\circ$，若镜像映射可提升为 $p_0\in\overline{\mathcal M}(V)$ 和 $q_0\in\overline{\mathcal{KM}}(V^\circ)$ 的邻域上，丛 $\mathcal H^V$ 与 $\mathcal H^{V^\circ}$ 的同构，就称 $(V,V^\circ)$ 为\emph{数学镜像对}。此外，此同构应保持来自 $\nabla^{\mathrm{GM}}$ 和 $\nabla^{\mathrm{middle}}$ 的极化 Hodge 结构变分，并将 $\mathcal H^V$ 的截面 $\Omega/y_0$、$g_0$ 分别映到 $\mathcal H^{V^\circ}$ 的截面 $1$、$[pt]$。
\end{definition}
\index{数学镜像对|hyperpage}\index{镜像对!数学镜像对|hyperpage}
借用 \SourceCite{Morrison7} 的术语，数学镜像对 $(V,V^\circ)$ 也可称为\emph{基于 Hodge 理论的镜像对}。定义~\ref{def:mathematical-mirror-pair} 只是局部表述；利用\hyperref[chap:moduli]{第~\ref*{chap:moduli} 章}讨论的复模空间和 Kähler 模空间，也可以给出更具整体性的定义。

接下来讨论定义~\ref{def:mathematical-mirror-pair}，并给出一些等价表述，随后再讨论镜像猜想的含义。

首先，定义~\ref{def:mathematical-mirror-pair} 假设 $V$ 满足整性猜想。\hyperref[subsec:a-model-asymptotics]{第~\ref*{subsec:a-model-asymptotics} 小节}末尾已指出，$V^\circ$ 上的 A 模型 Hodge 结构变分总满足整性猜想。因此，对 $V$ 的这一假设
\SourcePage{282}
是合理的，因为定义~\ref{def:mathematical-mirror-pair} 中的丛同构保持整数结构。

其次，与定义~\ref{def:mathematical-mirror-pair} 相容的丛同构 $\mathcal H^V\simeq\mathcal H^{V^\circ}$ 至多有一个。为说明这一点，注意，由 \SourceCite{BG}，反复对 $\widetilde\Omega=\Omega/y_0$ 施加 $\nabla^{\mathrm{GM}}_{\delta_i}$ 就能生成 $\mathcal H^V$。镜像一侧的 Lefschetz 分解表明，$H^{\mathrm{middle}}(V^\circ)$ 在杯积下由 $H^{1,1}(V^\circ)$ 生成；至少在 $q_0$ 附近，对小量子乘积也如此。这容易推出，反复对 $T_0=1$ 施加 $\nabla^{\mathrm{middle}}_{\delta_i}$ 就能生成 $\mathcal H^{V^\circ}$。所需的唯一性随即得到。

下面用更具体的方式讨论定义~\ref{def:mathematical-mirror-pair}。可以使用关联函数，也可以使用势函数，这将分别导出下面的定理~\ref{thm:mirror-correlation-equivalence} 和推论~\ref{cor:mirror-potential-equivalence}。

先\IndexLocation{indexloc:0757:282}用关联函数表述数学镜像对的含义。$V$ 上有 B 模型关联函数，也称归一化 Yukawa 耦合。回顾命题~\ref{prop:normalized-yukawa-pf} 后的讨论，在 $p_0$ 处存在局部坐标 $q_1,\ldots,q_r$，使 $\nabla^{\mathrm{GM}}$ 具有特别简单的形式~\eqref{eq:multimoduli-connection}；例题~\ref{ex:mirror-map-integral-connection} 表明，$q_j$ 确定镜像映射。若 $q_j=\exp(2\pi iu_j)$，则由式~\eqref{eq:multimoduli-yukawa}，$V$ 的归一化 Yukawa 耦合为
\begin{equation}\label{eq:b-model-normalized-correlations}\index{Yukawa耦合@Yukawa 耦合!归一化|hyperpage}
\begin{aligned}
Y_{ijk}&=-\int_V\widetilde\Omega\wedge\nabla_{\frac{\partial}{\partial u_i}}\nabla_{\frac{\partial}{\partial u_j}}\nabla_{\frac{\partial}{\partial u_k}}\widetilde\Omega\\
&=Q\left(\widetilde\Omega,\nabla_{\frac{\partial}{\partial u_i}}\nabla_{\frac{\partial}{\partial u_j}}\nabla_{\frac{\partial}{\partial u_k}}\widetilde\Omega\right),
\end{aligned}
\end{equation}
其中 $\widetilde\Omega=\Omega/y_0$ 为归一化三形式，$\nabla=\nabla^{\mathrm{GM}}$ 为 Gauss--Manin 联络，最后一个等号来自极化 $Q$ 的定义。这些就是 $V$ 的 B 模型关联函数。

在镜像一侧，存在 $T_1,\ldots,T_r\in H^2(V^\circ,\mathbb Z)$，使
\[
(q_1,\ldots,q_r)=(e^{2\pi iu_1},\ldots,e^{2\pi iu_r})
\longmapsto\left[\sum_ju_jT_j\right]\in K_{\mathbb C}(V^\circ)
\]
给出大半径极限点 $q_0$ 处的局部坐标。利用这些坐标，式~\eqref{eq:three-point-quantum-integral} 给出 A 模型关联函数
\begin{equation}\label{eq:a-model-mirror-correlations}
\langle T_i,T_j,T_k\rangle=\int_{V^\circ}T_i*_{\mathrm{small}}T_j*_{\mathrm{small}}T_k.
\end{equation}
用 A 模型联络 $\nabla=\nabla^{\mathrm{middle}}$ 表示时，式~\eqref{eq:a-model-yukawa-derivatives} 给出
\begin{equation}\label{eq:a-model-correlations-polarization}
\begin{aligned}
\langle T_i,T_j,T_k\rangle
&=\int_{V^\circ}1\cup\nabla_{\frac{\partial}{\partial u_i}}\nabla_{\frac{\partial}{\partial u_j}}\nabla_{\frac{\partial}{\partial u_k}}1\\
&=S\left(1,\nabla_{\frac{\partial}{\partial u_i}}\nabla_{\frac{\partial}{\partial u_j}}\nabla_{\frac{\partial}{\partial u_k}}1\right),
\end{aligned}
\end{equation}
其中最后一个等号使用了极化 $S$ 的定义。
\begin{theorem}\label{thm:mirror-correlation-equivalence}
若 $V$ 满足整性猜想，且镜像映射如上定义，则 $(V,V^\circ)$ 是定义~\ref{def:mathematical-mirror-pair} 意义下的数学镜像对，当且仅当 $V$ 的 B 模型关联函数 $Y_{ijk}$ 和 $V^\circ$ 的 A 模型关联函数 $\langle T_i,T_j,T_k\rangle$ 满足
\[
Y_{ijk}=\langle T_i,T_j,T_k\rangle\qquad\text{对所有 }i,j,k.
\]
\end{theorem}
\SourcePage{283}
\begin{proof}
一个方向平凡：若有同构将 $\widetilde\Omega$ 映到 $1$、$\nabla^{\mathrm{GM}}$ 映到 $\nabla^{\mathrm{middle}}$、$Q$ 映到 $S$，则关联函数相等由式~\eqref{eq:b-model-normalized-correlations} 和~\eqref{eq:a-model-correlations-polarization} 立即得到。

反过来，设关联函数按所需方式相同。用 $q_j$ 同时作为 $V$ 的复模空间与 $V^\circ$ 的 Kähler 模空间的局部坐标。由命题~\ref{prop:normalized-yukawa-pf} 后的讨论，丛 $\mathcal H^V$ 具有基 $e_0,e_j$（$1\le j\le r$）、$e^j$（$1\le j\le r$）、$e^0$，其中 $e_0=\widetilde\Omega$ 是特定截面。类似地，由命题~\ref{prop:a-model-canonical-flat-basis} 的证明，$\mathcal H^{V^\circ}$ 具有基 $T_0,T_j$（$1\le j\le r$）、$T^j$（$1\le j\le r$）、$T^0$，其中 $T_0=1$ 为特定截面。令 $e_i\mapsto T_i$、$e^i\mapsto T^i$，显然定义了一个丛同构。

为完成证明，须检验联络相容，且极化与整数结构都被保持。关联函数相同的假设，结合式~\eqref{eq:a-model-threefold-connection} 和~\eqref{eq:multimoduli-connection}，立即给出联络相容。极化也同样容易处理：使用双线性型 $Q$、$S$ 的定义，以及式~\eqref{eq:multimoduli-integral-pairings} 中 $e_i,e^i$ 的描述和命题~\ref{prop:a-model-canonical-flat-basis} 中 $T_i,T^i$ 的描述即可。

还需考虑整数结构。这一点稍微微妙，因为上述基不是整的。不过，这些基可延拓到 $\Delta^r$ 上的典范延拓，并在 $0\in\Delta^r$ 上成为整的。\hyperref[subsec:a-model-asymptotics]{第~\ref*{subsec:a-model-asymptotics} 小节}定义 $\mathcal H^{V^\circ}$ 的整数结构时规定：$\nabla^{\mathrm{middle}}$ 平坦截面 $s$ 为整截面，当且仅当 $\exp(-\sum_ju_jN_j)s$ 在 $0$ 上整。对于 $\mathcal H^V$ 也如此，于是容易得到，同构保持整数结构。

\end{proof}
这个定理很好地将数学镜像对的定义与\hyperref[chap:introduction]{第~\ref*{chap:introduction} 章}引入的“经典”镜像对称概念联系起来。这种方法的一个不足是，当 $V$ 的复模空间维数 $r$ 较大时，可能需要考虑很多关联函数，计算会变得相当复杂。使用势函数是一种简化办法。

回顾这一过程。在 $V$ 上，引理~\ref{lem:yukawa-potential} 证明了势函数 $\Phi^{\mathrm{GM}}$ 的存在。如\hyperref[subsec:normalized-yukawa-mirror]{第~\ref*{subsec:normalized-yukawa-mirror} 小节}所述，它是变量 $u_j$ 的函数，其中镜像映射为 $q_j=e^{2\pi iu_j}$。在这些变量中，引理~\ref{lem:yukawa-potential} 说明，归一化 Yukawa 耦合~\eqref{eq:b-model-normalized-correlations} 为
\begin{equation}\label{eq:b-model-potential-correlations}
Y_{ijk}=\frac1{(2\pi i)^3}\frac{\partial^3\Phi^{\mathrm{GM}}}{\partial u_i\partial u_j\partial u_k}
\qquad\text{对所有 }i,j,k.
\end{equation}
转到 $V^\circ$，由式~\eqref{eq:a-model-potential}，Gromov--Witten 势 $\Phi^{\mathrm{GW}}$ 等于
\[
\Phi^{\mathrm{GW}}=\frac{(2\pi i)^3}{6}\int_{V^\circ}\left(\sum_{j=1}^ru_jT_j\right)^3
+\sum_{\beta\ne0}\langle I_{0,0,\beta}\rangle q^\beta,
\]
其中 $q^\beta=q_1^{\int_\beta T_1}\cdots q_r^{\int_\beta T_r}$、$q_j=\exp(2\pi iu_j)$。类似地，式~\eqref{eq:a-model-yukawa-potential} 给出
\[
\langle T_i,T_j,T_k\rangle=\frac1{(2\pi i)^3}\frac{\partial^3\Phi^{\mathrm{GW}}}{\partial u_i\partial u_j\partial u_k}
\qquad\text{对所有 }i,j,k.
\]
因此，定理~\ref{thm:mirror-correlation-equivalence} 立即给出数学镜像对的势函数表述。
\SourcePage{284}
\begin{corollary}\label{cor:mirror-potential-equivalence}\index{镜像对称!势函数|hyperpage}
若 $V$ 满足整性猜想，且镜像映射如上定义，则 $(V,V^\circ)$ 是数学镜像对，当且仅当
\[
\Phi^{\mathrm{GM}}=\Phi^{\mathrm{GW}},
\]
其中允许两边相差关于 $u_j$ 的至多二次项。
\end{corollary}
\hyperref[subsec:quintic-mirror-symmetry]{第~\ref*{subsec:quintic-mirror-symmetry} 小节}将说明定理~\ref{thm:mirror-correlation-equivalence} 和推论~\ref{cor:mirror-potential-equivalence} 如何适用于五次三维簇。

在上述对定义~\ref{def:mathematical-mirror-pair} 及其后果的讨论中，应始终记住：要使 A 模型联络 $\nabla^{\mathrm{middle}}$ 成为真正的联络，需要收敛假设。有两种方式可以避开这一困难。第一种方式是注意，定理~\ref{thm:mirror-correlation-equivalence} 中的等式
\[
Y_{ijk}=\langle T_i,T_j,T_k\rangle
\]
作为关于形式幂级数的陈述完全有意义。因此，若能证明此等式，因为已知 $Y_{ijk}$ 收敛，$\langle T_i,T_j,T_k\rangle$ 的收敛性随即得到。于是\hyperref[sec:a-model-vhs]{第~\ref*{sec:a-model-vhs} 节}的全部理论都适用。类似地，在推论~\ref{cor:mirror-potential-equivalence} 的情形，只须证明形式幂级数等式 $\Phi^{\mathrm{GM}}=\Phi^{\mathrm{GW}}$，其余结论便都成立。

另一种可能更有趣的方式，是按 \SourceCite[第 7.2 节]{Morrison7} 的主张，建立“形式退化的 Hodge 结构变分”理论。这要求将\hyperref[sec:a-model-vhs]{第~\ref*{sec:a-model-vhs} 节}的一切发展为形式版本，再将定义~\ref{def:mathematical-mirror-pair} 替换为相应的形式陈述。

既然已有数学镜像对的严格定义，下一步似乎应当给出一个\emph{镜像猜想}，断言每个 Calabi--Yau 三维簇 $V$ 都有一个镜像 $V^\circ$，使 $(V,V^\circ)$ 满足定义~\ref{def:mathematical-mirror-pair}。遗憾的是，这行不通。刚性 Calabi--Yau 三维簇 $V$（只有平凡形变，故 $H^{2,1}(V)=0$）不可能具有该定义意义下的镜像 $V^\circ$，因为这会导致 $H^{1,1}(V^\circ)\simeq H^{2,1}(V)=0$，这是不可能的。因此，最多只能希望一个镜像猜想断言某些 Calabi--Yau 三维簇具有镜像。从这个角度看，\hyperref[chap:mirror-constructions]{第~\ref*{chap:mirror-constructions} 章}的每一种镜像构造，都导出一个镜像猜想。\hyperref[subsec:toric-mirror-conjecture]{第~\ref*{subsec:toric-mirror-conjecture} 小节}将详细探讨 Batyrev 镜像构造中的情形。

以推论~\ref{cor:mirror-potential-equivalence} 提示的一个有趣问题结束数学镜像对的讨论。给定任意函数 $\Phi(u_1,\ldots,u_r)$，使
\[
Y_{ijk}=\frac1{(2\pi i)^3}\frac{\partial^3\Phi}{\partial u_i\partial u_j\partial u_k}
\]
关于 $q_j=\exp(2\pi iu_j)$ 全纯，就可以用式~\eqref{eq:multimoduli-connection} 在某个向量丛上定义联络，并得到显然的 Hodge 滤过。利用\hyperref[subsec:a-model-hodge-filtration]{第~\ref*{subsec:a-model-hodge-filtration} 小节}的方法，可证明这给出 Hodge 结构变分。可能的 $\Phi$ 很多，故以此方式得到的变分也很多。但其中哪些来自几何，也就是来自 Calabi--Yau 三维簇族？换言之，Calabi--Yau 变分的势函数有什么特殊之处？镜像对称给出了部分答案：若 $\Phi$ 来自 Calabi--Yau 三维簇族，就应当是镜像的 Gromov--Witten 势
\SourcePage{285}
$\Phi^{\mathrm{GW}}$。这意味着势函数 $\Phi$ 具有非常特殊的形式。因此，镜像对称也帮助我们理解，哪些抽象的 Hodge 结构变分来自几何这一一般问题。

\subsection{五次三维簇的镜像对称}\label{subsec:quintic-mirror-symmetry}
接下来比较\hyperref[chap:quintic]{第~\ref*{chap:quintic} 章}给出的五次三维簇镜像对称的初步版本，与定义~\ref{def:mathematical-mirror-pair} 中更精细的表述。

\hyperref[chap:quintic]{第~\ref*{chap:quintic} 章}将五次三维簇的镜像对称表述为
\begin{equation}\label{eq:quintic-mirror-correlations-revisited}
5+\sum_{d=1}^\infty n_dd^3\frac{q^d}{1-q^d}
=\frac5{(1+5^5z)y_0(z)^2}\left(\frac qz\frac{dz}{dq}\right)^3,
\end{equation}
其中 $n_d$ 是五次三维簇的瞬子数，$z$ 是五次镜像的模坐标（\hyperref[chap:quintic]{第二章}记为 $x$），使五次镜像族在 $z=0$ 处的单值性最大幺幂。此外，
\[
\begin{aligned}
y_0(z)&=\sum_{n=0}^\infty\frac{(5n)!}{(n!)^5}(-1)^nz^n,\\
y_1(z)&=y_0(z)\log(-z)+5\sum_{n=1}^\infty\frac{(5n)!}{(n!)^5}
\left[\sum_{j=n+1}^{5n}\frac1j\right](-1)^nz^n,\\
q&=\exp(y_1/y_0).
\end{aligned}
\]
于是有如下良好事实。
\begin{theorem}\label{thm:quintic-mirror-correlation-equivalence}
取 $V^\circ$ 为五次三维簇、$V$ 为五次镜像时，定义~\ref{def:mathematical-mirror-pair} 与式~\eqref{eq:quintic-mirror-correlations-revisited} 等价。
\end{theorem}
\begin{proof}
根据现有知识，这容易看出。由式~\eqref{eq:cy-quantum-correlation-instanton}，式~\eqref{eq:quintic-mirror-correlations-revisited} 左端是五次三维簇的 A 模型关联函数 $\langle H,H,H\rangle$；例题~\ref{ex:quintic-normalized-yukawa} 则证明，右端是五次镜像的归一化 Yukawa 耦合。因此，该等式意味着定义 A 模型联络的级数收敛；如前所述，由 \SourceCite{Morrison2}，五次镜像族满足整性猜想。定理遂由定理~\ref{thm:mirror-correlation-equivalence} 得到。
\end{proof}
换言之，一旦证明\hyperref[chap:quintic]{第二章}中的公式，定义~\ref{def:mathematical-mirror-pair} 给出的五次三维簇镜像对称的完整版本，就成为直接推论。不过，不应因刚才的证明简短而低估其中的工作。对式~\eqref{eq:quintic-mirror-correlations-revisited} 两端所作的论断，背后都有大量工作。对于 A 模型，定义 $n_d$ 需要\hyperref[chap:gromov-witten]{第七章}的 Gromov--Witten 不变量定义，而理解其计数意义也并非平凡，见\hyperref[subsec:instanton-numbers]{第~\ref*{subsec:instanton-numbers} 小节}。随后还须完成\hyperref[sec:a-model-vhs]{第~\ref*{sec:a-model-vhs} 节}的 A 模型 Hodge 理论。对于 B 模型，\hyperref[chap:hodge-yukawa]{第五章}研究了镜像族的 Hodge 理论；例题~\ref{ex:quintic-normalized-yukawa} 表明，为证明右端是归一化 Yukawa 耦合，需要非常具体的单值性信息。

下面按推论~\ref{cor:mirror-potential-equivalence} 的方式，说明五次三维簇的镜像对称与势函数的关系。这很重要，因为\hyperref[chap:mirror-theorem]{第十一章}正是以此方式证明五次三维簇的镜像对称。

先描述 A 模型和 B 模型的势函数。它们是 $q$ 的多值函数，不过不再用 $u_1=(1/2\pi i)\log q$ 表示，而改用 $t_1=\log q$。如\hyperref[subsec:quantum-product-dependence]{第~\ref*{subsec:quantum-product-dependence} 小节}所述，变量
\SourcePage{286}
$t_1=2\pi iu_1$ 更适合\hyperref[chap:mirror-theorem]{第十一章}的讨论。对于 A 模型，这意味着将式~\eqref{eq:a-model-potential} 替换为 Gromov--Witten 势
\[
\Phi^{\mathrm{GW}}=\frac56t_1^3+\sum_{d=1}^\infty\langle I_{0,0,d}\rangle q^d.
\]
对于 B 模型，令 $Y$ 为式~\eqref{eq:quintic-mirror-correlations-revisited} 右端的归一化 Yukawa 耦合。若把 $Y$ 看作 $t_1=\log q$ 的函数，则势函数为三阶导数等于 $Y$ 的函数 $\Phi^{\mathrm{GM}}(t_1)$。为描述 $\Phi^{\mathrm{GM}}$，令 $z=-x$。将 $y_0,y_1$ 看作 $x$ 的函数时，它们由式~\eqref{eq:quintic-frobenius-basis} 给出；该式还给出 $y_2,y_3$，使 $y_0,\ldots,y_3$ 构成五次镜像 Picard--Fuchs 方程（以 $x$ 为变量）的一组解基。势函数由如下简洁公式给出。
\begin{proposition}\label{prop:quintic-mirror-potential}
五次镜像族的势函数为
\[
\Phi^{\mathrm{GM}}=\frac52\left(\frac{y_1}{y_0}\frac{y_2}{y_0}-\frac{y_3}{y_0}\right),
\]
其中借助镜像映射 $q=\exp(y_1/y_0)$，将其看作 $t_1=\log q$ 的函数。
\end{proposition}
\begin{proof}
须证明
\[
\frac{d^3}{dt_1^3}\frac52\left(\frac{y_1}{y_0}\frac{y_2}{y_0}-\frac{y_3}{y_0}\right)=Y.
\]
\hyperref[chap:quintic]{第二章}已证明此式，为完整起见，再回顾论证。$y_i$ 满足例题~\ref{ex:quintic-pf-method} 推导的 Picard--Fuchs 方程。将 $y_i$ 替换为 $y_i/y_0$，将 $x$ 替换为 $q=e^{t_1}$，命题~\ref{prop:normalized-yukawa-pf} 表明该方程成为
\[
\frac{d^2}{dt_1^2}\left(\frac{d^2y/dt_1^2}{Y}\right)=0.
\]
还要注意，由镜像映射的定义，$t_1=y_1/y_0$。利用此方程以及 $Y$、$y_i$ 的幂级数展开，式~\eqref{eq:period-derivative-identities} 中已证明
\[
5\frac{d^2}{dt_1^2}\frac{y_2}{y_0}=Y,\qquad
5\frac{d^2}{dt_1^2}\frac{y_3}{y_0}=Yt_1.
\]
于是，式~\eqref{eq:lly-yukawa} 后的论证给出所需等式。
\end{proof}
结合推论~\ref{cor:mirror-potential-equivalence}，得到如下定理。以下在镜像映射的源端用 $t_1=\log x$，在目标端用 $\Psi(t_1)=\log q=y_1/y_0$。
\begin{theorem}\label{thm:quintic-mirror-potential-equivalence}
取 $V^\circ$ 为五次三维簇、$V$ 为五次镜像时，定义~\ref{def:mathematical-mirror-pair} 等价于方程
\[
\Phi^{\mathrm{GW}}(\Psi(t_1))=\frac52\left(\frac{y_1}{y_0}\frac{y_2}{y_0}-\frac{y_3}{y_0}\right),
\]
其中镜像映射为 $\Psi(t_1)=y_1/y_0$，势函数取上述归一化。
\end{theorem}
采用这种形式表述定理~\ref{thm:quintic-mirror-potential-equivalence}，是因为上述方程就是 \SourceCite{LLY} 中的“\IndexLocation{indexloc:0864:286}\IndexLocation{indexloc:0867:286}镜像定理”，也是\hyperref[chap:mirror-theorem]{第十一章}定理~\ref{thm:lly-quintic-mirror} 证明五次三维簇\IndexLocation{indexloc:1197:286}镜像对称所用的形式。

\SourcePage{287}
\subsection{更高维情形}\label{subsec:higher-dimensional-mirror}\index{镜像对称!高维情形|hyperpage}
下面考虑 $V$、$V^\circ$ 是维数 $d>3$ 的光滑 Calabi--Yau 流形的情形。与前面一样，假设有对应的点 $p_0\in\overline{\mathcal M}(V)$、$q_0\in\overline{\mathcal{KM}}(V^\circ)$。困难出现在考察两个极化 Hodge 结构变分时。将会看到，情形不如三维簇那么理想。

先看 $V$。虽然 $H^d(V,\mathbb C)$ 上有良好的 Hodge 结构变分，但极化~\eqref{eq:hodge-polarization} 只适用于本原上同调，而本原上同调不一定是整个 $H^d(V,\mathbb C)$。为处理这一点，固定 Kähler 形式 $\omega$，并考虑 Lefschetz 分解。每个 $\alpha\in H^d(V,\mathbb C)$ 都可写成 $\sum_j\omega^j\cup\alpha_{d-2j}$，其中 $\alpha_{d-2j}\in H^{d-2j}_0(V,\mathbb C)$ 本原。于是有双线性型
\[
Q(\alpha,\beta)=\sum_j(-1)^{d(d-1)/2}(-1)^j
\int_V\omega^{2j}\cup\alpha_{d-2j}\cup\beta_{d-2j}.
\]
一般而言，此型只在 $\mathbb R$ 上定义（因为 $\omega$ 未必整），但由式~\eqref{eq:hodge-polarization}，它仍是极化。因此，在 $V$ 上选定 Kähler 结构以后，就得到极化实 Hodge 结构变分 $(\mathcal H^V,\nabla^{\mathrm{GM}},\mathcal H^V_{\mathbb R},\mathcal F^\bullet)$。

镜像一侧要考虑 $\bigoplus_{p=0}^dH^{p,p}(V^\circ)$ 上的 A 模型 Hodge 结构变分。虽然已有式~\eqref{eq:middle-a-model-polarization} 的良好极化型 $S$，但关键是：为定义 $\nabla^{\mathrm{middle}}$，须知道 $H^{p,p}(V^\circ)$ 如何嵌入 $H^{2p}(V^\circ,\mathbb C)$，因为量子上同调定义在后者上。这意味着必须固定 $V^\circ$ 的复结构。此外，$\bigoplus_{p=0}^dH^{p,p}(V^\circ)$ 一般只在 $\mathbb R$ 上定义。因此，选定 $V^\circ$ 的复结构以后，就得到极化实 Hodge 结构变分 $(\mathcal H^{V^\circ},\nabla^{\mathrm{middle}},\mathcal H^{V^\circ}_{\mathbb R},\mathcal F^\bullet)$。

还有 $\mathcal H^V$ 的特定截面 $\widetilde\Omega,g_0$，以及 $\mathcal H^{V^\circ}$ 的特定截面 $1=T_0,[pt]$。据此如下定义高维镜像对。
\begin{definition}\label{def:higher-mathematical-mirror-pair}
给定维数 $d>3$ 的 Calabi--Yau 流形 $V$、$V^\circ$，若能在 $V$ 上找到 Kähler 结构、在 $V^\circ$ 上找到复结构，使镜像映射提升为 $p_0\in\overline{\mathcal M}(V)$ 与 $q_0\in\overline{\mathcal{KM}}(V^\circ)$ 的邻域上丛 $\mathcal H^V$ 与 $\mathcal H^{V^\circ}$ 的同构，就称 $(V,V^\circ)$ 为\emph{数学镜像对}。此外，此同构保持来自 $\nabla^{\mathrm{GM}}$ 与 $\nabla^{\mathrm{middle}}$ 的极化实 Hodge 结构变分，并将 $\mathcal H^V$ 的截面 $\widetilde\Omega,g_0$ 分别映到 $\mathcal H^{V^\circ}$ 的截面 $1,[pt]$。
\end{definition}
与三维情形一样，数学镜像对有时也称为基于 Hodge 理论的镜像对。此定义有趣之处在于所需的额外数据：$V$ 上的 Kähler 结构与 $V^\circ$ 上的复结构。将这些数据加到 $\mathcal M(V)$ 与 $\mathcal{KM}(V^\circ)$ 上，意味着使用的是完整 SCFT 模空间上的镜像映射，如\hyperref[chap:introduction]{第一章}的式~\eqref{eq:mirror-foliations} 所示。还要注意，与 $p_0,q_0$ 的十分具体的选择相比，此定义对额外数据如何选择说得比较含糊。因此，或许最好将定义~\ref{def:higher-mathematical-mirror-pair} 看作初步版本。

另一个不确定之处是，A 模型变分在\IndexLocation{indexloc:0696:287}大半径极限点 $q_0$ 总为 Hodge--Tate，而\IndexLocation{indexloc:0772:287}最大幺幂单值性的定义，并不意味着 $H^d(V,\mathbb C)$ 的 Hodge 结构变分在对应点 $p_0$ 处为 Hodge--Tate。因此，定义~\ref{def:higher-mathematical-mirror-pair} 可能还需要加强。
\SourcePage{288}
尽管有这些限制，此定义与 $d$ 点函数仍有良好关系。$d$ 点函数定义如下。设镜像映射由 $q_1,\ldots,q_r$ 给出，照例同时用这些 $q_j$ 作为复模空间与 Kähler 模空间的局部坐标。令 $q_j=\exp(2\pi iu_j)$；为简化记号，令 $\delta_j=\partial/\partial u_j$。在这一约定下，$V$ 的 B 模型关联函数为
\[
Y_{i_1\cdots i_d}=\frac{(-1)^{d(d-1)/2}}{(2\pi i)^d}\int_V\widetilde\Omega\wedge\nabla_{\delta_{i_1}}\cdots\nabla_{\delta_{i_d}}\widetilde\Omega,
\qquad1\le i_1,\ldots,i_d\le r.
\]
类似地，$V^\circ$ 的 A 模型关联函数为
\[
\langle T_{i_1},\ldots,T_{i_d}\rangle=\int_{V^\circ}T_{i_1}*_{\mathrm{small}}\cdots*_{\mathrm{small}}T_{i_d}.
\]
这些关联函数具有如下关系。
\begin{proposition}\label{prop:higher-mirror-correlations}
若 $(V,V^\circ)$ 是定义~\ref{def:higher-mathematical-mirror-pair} 意义下的数学镜像对，则
\[
Y_{i_1\cdots i_d}=\langle T_{i_1},\ldots,T_{i_d}\rangle
\qquad\text{对所有 }1\le i_1,\ldots,i_d\le r.
\]
\end{proposition}
\begin{proof}
与定理~\ref{thm:mirror-correlation-equivalence} 的证明一样，定义~\ref{def:higher-mathematical-mirror-pair} 给出
\begin{equation}\label{eq:higher-mirror-pairing-equality}
Q(\widetilde\Omega,\nabla_{\delta_{i_1}}\cdots\nabla_{\delta_{i_d}}\widetilde\Omega)
=S(1,\nabla_{\delta_{i_1}}\cdots\nabla_{\delta_{i_d}}1),
\end{equation}
其中左端的 $\nabla$ 表示 Gauss--Manin 联络，右端表示 A 模型联络。

为理解左端，先注意 $\widetilde\Omega$ 是 $(d,0)$ 形式，故本原；类似地，$\nabla_{\delta_{i_1}}\cdots\nabla_{\delta_{i_d}}\widetilde\Omega$ 的 $(0,d)$ 分量也本原。由 $Q$ 的定义，不难得到式~\eqref{eq:higher-mirror-pairing-equality} 左端为 B 模型关联函数 $Y_{i_1\cdots i_d}$。

再看右端，因为所有 $T_{i_j}$ 的次数均为 $2$，
\[
\nabla_{\delta_{i_1}}\cdots\nabla_{\delta_{i_d}}1
=T_{i_1}*_{\mathrm{small}}\cdots*_{\mathrm{small}}T_{i_d}
+\text{次数小于 }2d\text{ 的项}.
\]
又因 $1\in H^0(V^\circ,\mathbb C)$，由 $S$ 的定义，右端为
\[
\begin{aligned}
\int_{V^\circ}1\cup\nabla_{\delta_{i_1}}\cdots\nabla_{\delta_{i_d}}1
&=\int_{V^\circ}\nabla_{\delta_{i_1}}\cdots\nabla_{\delta_{i_d}}1\\
&=\int_{V^\circ}T_{i_1}*_{\mathrm{small}}\cdots*_{\mathrm{small}}T_{i_d},
\end{aligned}
\]
恰为 A 模型关联函数 $\langle T_{i_1},\ldots,T_{i_d}\rangle$。
\end{proof}
虽然 $d>3$ 时关联函数有良好对应，但保持这些关联函数，不再等价于保持 Hodge 结构变分。这与三维情形很不同。通过一个例题说明其精确含义。
\SourceExampleNumber{8.6.3}{1}
\begin{example}\label{ex:fourfold-mirror-correlations}
设 $V$、$V^\circ$ 为 Calabi--Yau 四维簇。在 $V^\circ$ 上，只考虑 $\bigoplus_{p=0}^4H^{p,p}(V^\circ)$ 中的类。选取上同调基 $T_i\in H^{1,1}(V^\circ)$（$1\le i\le r$）、$T_k\in H^{2,2}(V^\circ)$（$r+1\le k\le s$）。令 $T^l\in H^{2,2}(V^\circ)$ 为对偶基，使 $\langle T_k,T^l\rangle=\delta_{kl}$。\IndexLocation{indexloc:0498:288}次数公理、\IndexLocation{indexloc:0502:288}基本类公理及\IndexLocation{indexloc:0499:288}除子公理表明，对于 $\beta\ne0$，只需考虑的非零 Gromov--Witten 不变量为 $\langle I_{0,3,\beta}\rangle(T_i,T_j,T_k)$，其中 $1\le i,j\le r$、$r+1\le k\le s$。定义
\[
Y^A_{ijk}=\sum_\beta\langle I_{0,3,\beta}\rangle(T_i,T_j,T_k)q^\beta,
\]
\SourcePage{289}
其中上标 $A$ 表示“A 模型”。于是，全部计数信息都包含在
\[
\nabla_{\delta_i}(T_j)=T_i*_{\mathrm{small}}T_j=\sum_{k=r+1}^sY^A_{ijk}T^k
\]
中，$\nabla$ 为 A 模型联络。用上述四点函数 $\langle T_i,T_j,T_a,T_b\rangle$ 表示，注意到
\begin{equation}\label{eq:fourfold-four-point-factorization}
\begin{aligned}
\langle T_i,T_j,T_a,T_b\rangle
&=\int_{V^\circ}T_i*_{\mathrm{small}}T_j*_{\mathrm{small}}T_a*_{\mathrm{small}}T_b\\
&=\int_{V^\circ}(T_i*_{\mathrm{small}}T_j)\cup(T_a*_{\mathrm{small}}T_b)\\
&=\sum_{k,l=r+1}^sY^A_{ijk}Y^A_{abl}\langle T^k,T^l\rangle,
\end{aligned}
\end{equation}
其中中间的等号来自命题~\ref{prop:quantum-frobenius-pairing}。

现在转向四维簇 $V$ 上的 B 模型。假设 $V$ 在所选的最大幺幂边界点处满足整性猜想，并且极限混合 Hodge 结构为 Hodge--Tate。与前面一样，可找到坐标 $q_1,\ldots,q_r$，使取 $e_0=\widetilde\Omega$ 后，$e_i=\nabla_{\delta_i}(e_0)$（$1\le i\le r$）给出 $\mathcal W_6\cap\mathcal F^3$ 的一组基。再令 $e^k$ 为 $\mathcal W_4\cap\mathcal F^2$ 的某组基，就可写出
\[
\nabla_{\delta_i}(e_j)=\sum_{k=r+1}^sY^B_{ijk}e^k
\]
，其中 $Y^B_{ijk}$ 为一些函数，$B$ 当然表示“B 模型”。

设 $(V,V^\circ)$ 为数学镜像对。由\hyperref[subsec:mirror-map-derivative]{第~\ref*{subsec:mirror-map-derivative} 小节}，镜像映射的导数为 $e_i\mapsto T_i$（$1\le i\le r$）。应用命题~\ref{prop:higher-mirror-correlations} 得到四点函数相等，从而可由 B 模型计算 A 模型关联函数 $\langle T_i,T_j,T_a,T_b\rangle$。但如式~\eqref{eq:fourfold-four-point-factorization} 所示，这不能再直接给出 $V^\circ$ 的 Gromov--Witten 不变量；这些不变量编码在 $Y^A_{ijk}$ 中。

如 \SourceCite{Morrison7} 指出的，真正希望得到的是等式 $Y^A_{ijk}=Y^B_{ijk}$，它能告诉我们如何用 B 模型找到 Gromov--Witten 不变量。但这需要完整同构 $\mathcal H^V\simeq\mathcal H^{V^\circ}$，特别要用到 $e^k\mapsto T^k$（$r+1\le k\le s$）。后者等价于知道同构 $H^{2,2}(V)\simeq H^{2,2}(V^\circ)$。因此，需要\hyperref[chap:introduction]{第一章}镜像对称同构~\eqref{eq:mirror-cohomology-forms} 的四维版本。

不过，若能自然选取同构 $H^{2,2}(V)\simeq H^{2,2}(V^\circ)$，应当就能给出定义~\ref{def:higher-mathematical-mirror-pair} 的精确版本，从而推出 $Y^B_{ijk}=Y^A_{ijk}$。这样就可根据 $V$ 的 Hodge 理论预测 $V^\circ$ 的 Gromov--Witten 不变量。\SourceCite{Morrison7} 讨论了这种四维镜像对称方法，并给出一些例题。此外，若同构 $H^{2,2}(V)\simeq H^{2,2}(V^\circ)$ 保持极限混合 Hodge 结构的实结构，那么定义~\ref{def:higher-mathematical-mirror-pair} 等价于对所有 $i,j,k$ 都有 $Y^B_{ijk}=Y^A_{ijk}$。因此，四维簇情形也有定理~\ref{thm:mirror-correlation-equivalence} 的类似命题，只是稍微复杂一些。
\end{example}
\SourcePage{290}
有了 $d>3$ 时镜像对的定义，就可以像 \SourceCite{Morrison7} 那样，尝试提出一般的镜像猜想。不过，由于刚性 Calabi--Yau 流形的存在，而且缺少镜像的一般描述，这样的猜想必然比较含糊。因此，这里不提出一般猜想，下面只讨论环面情形。

\subsection{环面镜像猜想}\label{subsec:toric-mirror-conjecture}
刚才指出，提出一般镜像猜想的困难之一，是给定 $V$ 后不知道如何找到 $V^\circ$。环面情形更为理想，因为可以使用\hyperref[chap:mirror-constructions]{第四章}的镜像构造。为简便起见，只考虑环面超曲面，此时 $V^\circ$ 为 $V$ 的 Batyrev 镜像。将会看到，得到的猜想在某些方面更具体，在另一些方面却不那么具体。

假设处于\hyperref[sec:batyrev-mirror]{第~\ref*{sec:batyrev-mirror} 节}的情形。固定维数至少为 $4$ 的自反多胞形 $\Delta$，令 $\Sigma$、$\Sigma^\circ$ 分别为 $\Delta$、$\Delta^\circ$ 的法扇的极大射影细分。得到极小 Calabi--Yau 环面超曲面 $V\subset X_\Sigma$ 及其 Batyrev 对偶 $V^\circ\subset X_{\Sigma^\circ}$。

为提出猜想，\hyperref[subsec:threefold-mirror-symmetry]{第~\ref*{subsec:threefold-mirror-symmetry}} 和\hyperref[subsec:higher-dimensional-mirror]{第~\ref*{subsec:higher-dimensional-mirror} 小节}提示，应从 $V$、$V^\circ$ 为三维簇的情形开始。此时得到如下\emph{环面三维簇镜像猜想}。
\begin{conjecture}\label{conj:toric-threefold-mirror}
若 $V\subset X_\Sigma$ 及其 Batyrev 对偶 $V^\circ\subset X_{\Sigma^\circ}$ 如上，且 $\dim V=\dim V^\circ=3$，则 $(V,V^\circ)$ 为定义~\ref{def:mathematical-mirror-pair} 意义下的数学镜像对。
\end{conjecture}
这是一个良好的猜想，但实际中，它与通常所谓环面镜像猜想不完全相同。原因在于，环面情形使用不同的模空间及镜像映射定义。

回顾这一过程，现在设 $\dim V=d\ge3$ 任意。\hyperref[chap:moduli]{第六章}在\hyperref[subsec:toric-complex-moduli]{第~\ref*{subsec:toric-complex-moduli} 小节}构造了简化模空间 $\overline{\mathcal M}_{\mathrm{simp}}(V)$，并在\hyperref[subsec:enlarged-kahler-moduli]{第~\ref*{subsec:enlarged-kahler-moduli} 小节}构造了环面 Kähler 模空间 $\overline{\mathcal{KM}}_{\mathrm{toric}}(V^\circ)$。为得到边界点，取由 $H^2_{\mathrm{toric}}(V^\circ,\mathbb Z)$ 的基 $T_1,\ldots,T_{\widetilde r}$ 生成的锥 $\sigma$，这些基类位于 Kähler 锥的闭包内。\hyperref[chap:moduli]{第六章}说明，$\sigma$ 给出最大幺幂边界点 $p_0\in\overline{\mathcal M}_{\mathrm{simp}}(V)$（假设猜想~\ref{conj:toric-mum-boundary} 成立），以及大半径极限点 $q_0\in\overline{\mathcal{KM}}_{\mathrm{toric}}(V^\circ)$。锥 $\sigma$ 还分别确定 $p_0$、$q_0$ 附近的显式局部坐标 $z_1,\ldots,z_{\widetilde r}$、$q_1,\ldots,q_{\widetilde r}$。此外，如\hyperref[subsec:toric-mirror-map]{第~\ref*{subsec:toric-mirror-map} 小节}所述，可在 $p_0,q_0$ 附近完全显式地确定环面版本的镜像映射。

为提出猜想，再回到 $V$、$V^\circ$ 为三维簇的情形。由命题~\ref{prop:toric-cy-construction} 证明后的说明，它们光滑。为进一步简化，先假设
\begin{equation}\label{eq:toric-divisor-surjectivity}
H^{1,1}(X_{\Sigma^\circ})\longrightarrow H^{1,1}(V^\circ)
\end{equation}
满射。按\hyperref[subsec:batyrev-hodge]{第~\ref*{subsec:batyrev-hodge} 小节}的记号，这意味着 $H^{1,1}_{\mathrm{toric}}(V^\circ)=H^{1,1}(V^\circ)$；由定理~\ref{thm:batyrev-hodge} 的证明，它又等价于 $H^{2,1}_{\mathrm{poly}}(V)=H^{2,1}(V)$。因此，上一段中的模空间正是\hyperref[subsec:threefold-mirror-symmetry]{第~\ref*{subsec:threefold-mirror-symmetry} 小节}所考虑的空间 $\overline{\mathcal{KM}}(V^\circ)$ 与 $\overline{\mathcal M}(V)$。

仍有联络 $\nabla^{\mathrm{GM}}$、$\nabla^{\mathrm{middle}}$ 及相应的 Hodge 结构变分。此外，如\hyperref[subsec:threefold-mirror-symmetry]{第~\ref*{subsec:threefold-mirror-symmetry} 小节}所述，镜像映射的导数、截面 $\widetilde\Omega$ 和 $1$、以及极化，给出唯一同构 $\mathcal H^V\simeq\mathcal H^{V^\circ}$。但即使这个同构保持联络 $\nabla^{\mathrm{GM}}$
\SourcePage{291}
和 $\nabla^{\mathrm{middle}}$，也不能断言它保持整数结构 $\mathcal H^V_{\mathbb Z}$ 和 $\mathcal H^{V^\circ}_{\mathbb Z}$。问题如下：定义~\ref{def:mathematical-mirror-pair} 中的镜像映射要求整性猜想，因而几乎由定义就保持整数结构。而\hyperref[subsec:toric-mirror-map]{第~\ref*{subsec:toric-mirror-map} 小节}给出的环面镜像映射可能与它不同，特别可能不保持整数结构。换句话说，尚不知道\hyperref[subsec:batyrev-hodge]{第~\ref*{subsec:batyrev-hodge}} 和\hyperref[subsec:toric-mirror-map]{第~\ref*{subsec:toric-mirror-map} 小节}的单项式—除子镜像映射是否保持整数结构。实际上，据我们所知，它甚至可能不保持实结构。

当然，猜想认为两种镜像映射相同，且 $V$ 满足整性猜想。但为了提出不依赖这些假设的猜想，只考虑仅涉及复数的结构。因此，对于 $V$，有 $(\mathcal H^V,\nabla^{\mathrm{GM}},\mathcal F^\bullet)$ 给出的\emph{复 Hodge 结构变分}。它由型 $Q$“极化”，这意味着第一条 Hodge--Riemann \IndexLocation{indexloc:0610:291}双线性关系的加强形式，即对所有 $p$，$(\mathcal F^p)^\perp=\mathcal F^{d-p+1}$。类似地，对于 $V^\circ$，有 $(\mathcal H^{V^\circ},\nabla^{\mathrm{middle}},\mathcal F^\bullet)$，它如\hyperref[subsec:threefold-mirror-symmetry]{第~\ref*{subsec:threefold-mirror-symmetry} 小节}所述被 $S$ 极化。

还要注意，由命题~\ref{prop:a-model-weight-filtration}，$\mathcal H^{V^\circ}$ 上的\IndexLocation{indexloc:0907:291}单值性权滤过在 $0\in\Delta^r$ 处满足 $\operatorname{Gr}^{\mathcal W}_4=H^2(V^\circ,\mathbb C)=H^{1,1}(V^\circ)$。类似地，在最大幺幂边界点处，$\mathcal H^V$ 上有 $\operatorname{Gr}^{\mathcal W}_4=H^{2,1}(V)$。这是混合 Hodge 结构为 Hodge--Tate 的结果；如式~\eqref{eq:mum-hodge-tate-splitting} 附近所述，在三维情形，它来自最大幺幂单值性。

在这种情形，得到\emph{基于 Hodge 理论的\IndexLocation{indexloc:0791:291}\IndexLocation{indexloc:0795:291}\IndexLocation{indexloc:0797:291}环面镜像猜想}的下列特殊情形。
\begin{conjecture}\label{conj:toric-hodge-mirror-surjective}
设 $V$、$V^\circ$ 如上为互为 Batyrev 镜像的三维 Calabi--Yau 环面超曲面，并假设式~\eqref{eq:toric-divisor-surjectivity} 满射。则镜像映射提升为 $p_0\in\overline{\mathcal M}(V)$ 和 $q_0\in\overline{\mathcal{KM}}(V^\circ)$ 的邻域上，丛 $\mathcal H^V$ 与 $\mathcal H^{V^\circ}$ 的同构。此外，此同构保持来自 $\nabla^{\mathrm{GM}}$ 与 $\nabla^{\mathrm{middle}}$ 的极化复 Hodge 结构变分，并将 $\mathcal H^V$ 的截面 $\widetilde\Omega,g_0$ 分别映到 $\mathcal H^{V^\circ}$ 的截面 $1,[pt]$。最后，在 $p_0,q_0$ 处的 $\operatorname{Gr}^{\mathcal W}_4$ 上，诱导同构就是单项式—除子镜像映射。
\end{conjecture}
\index{镜像猜想!环面情形|hyperpage}
与定理~\ref{thm:mirror-correlation-equivalence} 类似，此时又可用 A 模型、B 模型关联函数表述镜像对称。与前面一样，$V^\circ$ 上有通常的 A 模型关联函数 $\langle T_i,T_j,T_k\rangle$。在 $V$ 上，定义\emph{环面 B 模型关联函数}为
\[
Y_{ijk}=\int_V\widetilde\Omega\wedge\nabla_{\delta_i}\nabla_{\delta_j}\nabla_{\delta_k}\widetilde\Omega,
\]
其中 $\delta_j=2\pi iq_j\partial/\partial q_j$，$\widetilde\Omega=\Omega/y_0$ 为归一化三形式。它与\hyperref[subsec:normalized-yukawa-mirror]{第~\ref*{subsec:normalized-yukawa-mirror} 小节}定义的归一化 Yukawa 耦合有两点不同：
\begin{itemize}
\item 这里的 $y_0$ 是 $\Omega$ 的 Picard--Fuchs 方程在最大幺幂边界点处的任意非零全纯解。因此，只在相差常数的意义下知道 $y_0$，从而也只在此意义下知道 $\widetilde\Omega$。
\item 使用的是环面版本的镜像映射，它可能与定义归一化 Yukawa 耦合时所用的镜像映射不同。
\end{itemize}
因此，$Y_{ijk}$ 只在相差乘法常数的意义下确定；与定义~\ref{def:normalized-yukawa} 相比，上式没有负号，原因就在这里。不过，如\hyperref[subsec:normalized-yukawa-mirror]{第~\ref*{subsec:normalized-yukawa-mirror} 小节}指出的，除非极其详细地研究单值性，我们的计算方法也只能
\SourcePage{292}
在相差常数的意义下确定 $Y_{ijk}$。幸而，此时仍有势函数（引理~\ref{lem:yukawa-potential} 的证明仍适用），故本质上只须处理一个常数。

由此容易给出定理~\ref{thm:mirror-correlation-equivalence} 的一个版本，断言猜想~\ref{conj:toric-hodge-mirror-surjective} 等价于
\begin{equation}\label{eq:toric-mirror-correlations}
Y_{ijk}=\langle T_i,T_j,T_k\rangle\qquad\text{对所有 }i,j,k,
\end{equation}
其中允许将所有 $Y_{ijk}$ 同乘一个与 $i,j,k$ 无关的常数 $c$，使等式成立。这个证明直接，细节略去。

以下例题说明，所涉及模空间为二维时，猜想~\ref{conj:toric-hodge-mirror-surjective} 具有怎样的形式。
\SourceExampleNumber{8.6.4}{1}
\begin{example}\label{ex:octic-toric-mirror-correlations}
一个多次考虑过的例子，是 $\mathbb P(1,1,2,2,2)$ 中八次超曲面的环面解消 $V$ 及其 Batyrev 镜像 $V^\circ$。已知 $V$ 的环面 Kähler 模空间为二维（这里也等于全部 Kähler 模空间的维数），故式~\eqref{eq:toric-divisor-surjectivity} 满射。先写出 $V$ 的一个 A 模型关联函数。

在例题~\ref{ex:weighted-extremal-transition} 中看到，$V$ 的 Kähler 锥由记为 $D_3,D_1$ 的类生成。由 \SourceCite{CdFKM}，这些类可描述如下：$D_3=H$ 是 $\mathbb P(1,1,2,2,2)$ 上由二次方程定义的丰沛类的严格变换；$D_1=L$ 是由一次方程定义的除子的严格变换。若 $E\subset V$ 为映射 $V\to\overline V$ 的例外轨迹，则 $H\sim2L+E$。如该例指出的，$H^3=8$、$L^3=0$。

由 $H,L$ 生成的锥 $\sigma$（这里是整个 Kähler 锥）确定 $V$ 的大半径极限点和 $V^\circ$ 的最大幺幂边界点。$V$ 的 Kähler 模空间的局部坐标为 $q_j=\exp(2\pi iu_j)$，其中 $\sigma=\{u_1H+u_2L:u_1,u_2\ge0\}$。若 $h,\ell$ 为 $H_2(V,\mathbb Z)$ 中的对偶基，则 $\beta\in H_2(V,\mathbb Z)$ 可写为 $\beta=ah+b\ell$。因为 $H^3=8$、$\int_\beta H=a$，式~\eqref{eq:cy-quantum-correlation-instanton} 给出
\begin{equation}\label{eq:octic-a-model-instants}
\langle H,H,H\rangle=8+\sum_{(a,b)\ne(0,0)}n(a,b)\frac{a^3q_1^aq_2^b}{1-q_1^aq_2^b},
\end{equation}
其中 $n(a,b)$ 为\hyperref[subsec:instanton-numbers]{第~\ref*{subsec:instanton-numbers} 小节}定义的瞬子数 $n_{ah+b\ell}$。粗略地说，$n(a,b)$ 是 $V$ 中同调类为 $ah+b\ell$ 的有理曲线数；但例题~\ref{ex:quintic-instanton-enumerative} 表明，瞬子数与有理曲线的关系比初看时更微妙。

利用猜想~\ref{conj:toric-hodge-mirror-surjective} 求出 $\langle H,H,H\rangle$ 的猜想公式，就能同时预测所有瞬子数 $n(a,b)$。先注意，按这里的记号，使用环面镜像猜想时交换了 $V$ 与 $V^\circ$ 的角色。希望这不会造成太多混淆。以下用 $\nabla$ 表示 $V^\circ$ 的 Gauss--Manin 联络。由式~\eqref{eq:toric-mirror-correlations}，环面镜像猜想预测
\[
\langle H,H,H\rangle=Y_{111}=\int_{V^\circ}\widetilde\Omega\wedge\nabla_{\delta_1}\nabla_{\delta_1}\nabla_{\delta_1}\widetilde\Omega,
\]
其中 $\delta_1=2\pi iq_1\partial/\partial q_1$。

为计算右端积分，先用 $V^\circ$ 的模坐标 $z_1,z_2$。在例题~\ref{ex:weighted-yukawa} 中，已利用这些变量计算 Yukawa
\SourcePage{293}
耦合
\[
K^{ij}=\int_{V^\circ}\Omega\wedge\nabla_{z_1\frac{\partial}{\partial z_1}}^i\nabla_{z_2\frac{\partial}{\partial z_2}}^j\Omega,
\qquad i+j=3.
\]
现在须将 $\Omega$ 除以 $y_0$ 归一化，再转到镜像坐标 $q_1,q_2$。因为
\[
\delta_1=2\pi iq_1\frac{\partial}{\partial q_1}
=2\pi i\left(\frac{q_1}{z_1}\frac{\partial z_1}{\partial q_1}z_1\frac{\partial}{\partial z_1}
+\frac{q_1}{z_2}\frac{\partial z_2}{\partial q_1}z_2\frac{\partial}{\partial z_2}\right),
\]
代入 $\langle H,H,H\rangle$ 的公式，得到
\[
\begin{aligned}
\langle H,H,H\rangle=\frac{(2\pi i)^3}{y_0^2}\Biggl[&
\left(\frac{q_1}{z_1}\frac{\partial z_1}{\partial q_1}\right)^3K^{30}
+3\left(\frac{q_1}{z_1}\frac{\partial z_1}{\partial q_1}\right)^2
\left(\frac{q_1}{z_2}\frac{\partial z_2}{\partial q_1}\right)K^{21}\\
&+3\left(\frac{q_1}{z_1}\frac{\partial z_1}{\partial q_1}\right)
\left(\frac{q_1}{z_2}\frac{\partial z_2}{\partial q_1}\right)^2K^{12}
+\left(\frac{q_1}{z_2}\frac{\partial z_2}{\partial q_1}\right)^3K^{03}\Biggr].
\end{aligned}
\]
例题~\ref{ex:weighted-yukawa} 已显式计算 $K^{ij}$，只差一个共同常数 $c\ne0$；例题~\ref{ex:weighted-hypergeometric-mirror-map} 又将 $y_0,q_2$ 显式计算为 $z_1,z_2$ 的函数。$q_1$ 也可这样计算，因此上述公式能将 $\langle H,H,H\rangle$ 表示为 $q_1,q_2$ 的函数。

先确定常数 $c$。由镜像映射的环面归一化~\eqref{eq:toric-mirror-normalization}，$q_j=z_j+\cdots$；由例题~\ref{ex:weighted-yukawa}，
\[
K^{30}=\frac c{(1-256z_1)^2-512^2z_1^2z_2},
\]
故 $K^{30}=c+\cdots$。又因 $y_0=1+\cdots$，上述公式给出
\[
\langle H,H,H\rangle=(2\pi i)^3c+\cdots.
\]
而 $\langle H,H,H\rangle=8+\cdots$，故必须有
\[
c=\frac8{(2\pi i)^3},
\]
这与\hyperref[chap:quintic]{第二章}出现的情形相似。取此 $c$，得到展开
\[
\begin{aligned}
\langle H,H,H\rangle={}&8+640\frac{q_1}{1-q_1}
+10032\frac{2^3q_1^2}{1-q_1^2}+640\frac{q_1q_2}{1-q_1q_2}\\
&+288384\frac{3^3q_1^3}{1-q_1^3}
+72224\frac{2^3q_1^2q_2}{1-q_1^2q_2}
+10979984\frac{4^3q_1^4}{1-q_1^4}+\cdots.
\end{aligned}
\]
只要有耐心及足够计算能力，便可计算任意多项。与式~\eqref{eq:octic-a-model-instants} 比较，得到瞬子数 $n(a,b)$。可证明，$a\ge1$ 时有 $n(a,b)=n(a,a-b)$，例如 $n(1,1)=n(1,0)=640$；$b>a$ 时，除 $n(0,1)=4$ 外，都有 $n(a,b)=0$。因此，$\langle H,H,H\rangle$ 的级数确定所有 $(a,b)\ne(0,1)$ 对应的 $n(a,b)$。此例及其他 Calabi--Yau 环面超曲面例子的进一步计算，见 \SourceCite{CdFKM,HKTY1}。

\hyperref[chap:mirror-theorem]{第十一章}的例题~\ref{ex:octic-toric-mirror-theorem} 将说明，上述公式来自定理~\ref{thm:givental-toric-mirror}。此外，利用 \SourceCite{CdFKM} 中详细的单值性计算，可以实际证明，此例\IndexLocation{indexloc:1031:293}满足\IndexLocation{indexloc:0636:293}整性猜想，且取 $c=8/(2\pi i)^3$ 后，$Y_{111}$ 就是\hyperref[subsec:normalized-yukawa-mirror]{第~\ref*{subsec:normalized-yukawa-mirror} 小节}严格意义下的归一化 Yukawa 耦合。
\end{example}
\SourcePage{294}
下面考虑 $V$ 及其 Batyrev 镜像 $V^\circ$ 仍为三维簇，但不再假设映射~\eqref{eq:toric-divisor-surjectivity} 满射的情形。这时须处理\hyperref[subsec:toric-complex-moduli]{第~\ref*{subsec:toric-complex-moduli} 小节}的简化模空间 $\overline{\mathcal M}_{\mathrm{simp}}(V)$ 和\hyperref[subsec:enlarged-kahler-moduli]{第~\ref*{subsec:enlarged-kahler-moduli} 小节}的环面 Kähler 模空间 $\overline{\mathcal{KM}}_{\mathrm{toric}}(V^\circ)$。

一种提出猜想的方式如下。完整模空间上有丛 $\mathcal H^V$、$\mathcal H^{V^\circ}$；限制后得到
\[
\begin{aligned}
\mathcal H^V_{\mathrm{simp}}&=\mathcal H^V\text{ 在 }\overline{\mathcal M}_{\mathrm{simp}}(V)\text{ 上的限制},\\
\mathcal H^{V^\circ}_{\mathrm{toric}}&=\mathcal H^{V^\circ}\text{ 在 }\overline{\mathcal{KM}}_{\mathrm{toric}}(V^\circ)\text{ 上的限制}.
\end{aligned}
\]
这些限制丛继承自然的极化 Hodge 结构变分。于是，基于 Hodge 理论的环面镜像猜想会断言，存在保持极化复 Hodge 结构变分的丛同构 $\mathcal H^V_{\mathrm{simp}}\simeq\mathcal H^{V^\circ}_{\mathrm{toric}}$，如猜想~\ref{conj:toric-hodge-mirror-surjective} 所述。不过，不将它正式列为猜想，因为这样的丛同构需要知道同构 $H^{2,1}(V)\simeq H^{1,1}(V^\circ)$，而环面情形只有
\[
H^{2,1}_{\mathrm{poly}}(V)\simeq H^{1,1}_{\mathrm{toric}}(V^\circ),
\]
它来自\hyperref[subsec:batyrev-hodge]{第~\ref*{subsec:batyrev-hodge}} 和\hyperref[subsec:toric-mirror-map]{第~\ref*{subsec:toric-mirror-map} 小节}的单项式—除子镜像映射。回顾，$H^{2,1}_{\mathrm{poly}}(V)$ 是 $H^{2,1}(V)\simeq H^1(V,\Theta_V)$ 的子空间，对应于改变 $V\subset X_\Sigma$ 的定义方程得到的形变。因此，问题在于还不理解如何定义 $\mathcal H^V_{\mathrm{simp}}\simeq\mathcal H^{V^\circ}_{\mathrm{toric}}$。

不过，若将这两个丛替换为略小的丛，就可给出相当精确的环面镜像猜想版本。先从 $V^\circ$ 开始。令 $i:V^\circ\hookrightarrow X_{\Sigma^\circ}$ 为包含映射，定义
\[
H^*_{\mathrm{toric}}(V^\circ)=\operatorname{im}\bigl(i^*:H^*(X_{\Sigma^\circ})\longrightarrow H^*(V^\circ)\bigr).
\]
因为 $p\ne q$ 时 $H^{p,q}(X_{\Sigma^\circ})=0$，$H^*_{\mathrm{toric}}(V^\circ)$ 由 $(p,p)$ 型上同调类组成。因此，限制到该空间得到子丛
\[
\begin{aligned}
H^*_{\mathrm{toric}}(V^\circ)\times\overline{\mathcal{KM}}_{\mathrm{toric}}(V^\circ)
&\subset H^{\mathrm{middle}}(V^\circ)\times\overline{\mathcal{KM}}_{\mathrm{toric}}(V^\circ)\\
&=\mathcal H^{V^\circ}_{\mathrm{toric}}.
\end{aligned}
\]
将此子丛记为 $\mathcal T\mathcal H^{V^\circ}_{\mathrm{toric}}$，其中 $\mathcal T$ 表示所用的环面子环。

当 $X_{\Sigma^\circ}$ 凸（例如 $X_{\Sigma^\circ}=\mathbb P^4$）且 $V^\circ$ 丰沛时，\SourceCite[命题 4]{Pandharipande3} 的论证表明，$H^*_{\mathrm{toric}}(V^\circ)$ 是 $H^*(V^\circ)$ 关于小量子乘积的子环。一般情形是否如此尚不知道，因此进一步假设
\begin{equation}\label{eq:toric-quantum-subring-assumption}
H^*_{\mathrm{toric}}(V^\circ)\text{ 是 }H^*(V^\circ)\text{ 关于 }*_{\mathrm{small}}\text{ 的子环}.
\end{equation}
在此假设下，A 模型联络 $\nabla^{\mathrm{middle}}$ 在 $\mathcal T\mathcal H^{V^\circ}_{\mathrm{toric}}$ 上诱导联络。断言实际上得到极化 Hodge 结构变分。为证明这一点，假设 $X_{\Sigma^\circ}$ 光滑，则有如下引理。
\begin{lemma}\label{lem:toric-cohomology-pairing-nondegenerate}
若 $V^\circ$ 是光滑环面簇 $X_{\Sigma^\circ}$ 中的 Calabi--Yau 环面超曲面，则杯积配对在 $H^*_{\mathrm{toric}}(V^\circ)$ 上的限制非退化。
\end{lemma}
\SourcePage{295}
\begin{proof}
只须证明
\[
\int_{V^\circ}i^*(\alpha)\cup i^*(\beta)=0\quad\text{对所有 }\beta\in H^*(X_{\Sigma^\circ})
\quad\Longrightarrow\quad i^*(\alpha)=0.
\]
由投影公式，
\[
\int_{V^\circ}i^*(\alpha)\cup i^*(\beta)=\int_{X_{\Sigma^\circ}}\alpha\cup\beta\cup[V^\circ].
\]
而 \SourceCite[(3.36)]{HLY1} 意味着，对于 $\alpha\in H^*(X_{\Sigma^\circ})$，有
\[
\alpha\cup[V^\circ]=0\text{ 于 }H^*(X_{\Sigma^\circ})
\quad\Longleftrightarrow\quad
 i^*(\alpha)=0\text{ 于 }H^*_{\mathrm{toric}}(V^\circ).
\]
这两个事实容易推出所需的非退化性。
\end{proof}
假设式~\eqref{eq:toric-quantum-subring-assumption} 成立，便可将 $\mathcal H^{V^\circ}_{\mathrm{toric}}$ 替换为 $\mathcal T\mathcal H^{V^\circ}_{\mathrm{toric}}$，重复\hyperref[sec:a-model-vhs]{第~\ref*{sec:a-model-vhs} 节}的构造。极化来自刚证明的引理；如\hyperref[subsec:a-model-asymptotics]{第~\ref*{subsec:a-model-asymptotics} 小节}所述，还得到 $\mathcal T\mathcal H^{V^\circ}_{\mathrm{toric}}$ 的整数结构，因为 $H^*_{\mathrm{toric}}(V^\circ)$ 自然在整数上定义。

镜像一侧需要找到对应的 Hodge 结构变分。回顾定理~\ref{thm:batyrev-hodge} 的证明，其中有正合列
\[
\longrightarrow\bigoplus_iH^1(V\cap D_i,\mathbb C)\longrightarrow H^3(V,\mathbb C)
\longrightarrow\operatorname{Gr}^W_3H^3(Z_f,\mathbb C)\longrightarrow0.
\]
这里 $D_i$ 是 $X_\Sigma$ 的环面除子，$Z_f=T\cap V$ 是环面 $T\subset X_\Sigma$ 中的仿射超曲面。在该证明中，最终得到
\[
h^{2,1}\bigl(\operatorname{Gr}^W_3H^3(Z_f,\mathbb C)\bigr)=h^{2,1}_{\mathrm{poly}}(V).
\]
这提示存在自然同构
\begin{equation}\label{eq:polynomial-affine-hodge-assumption}
H^{2,1}_{\mathrm{poly}}(V)\simeq H^{2,1}\bigl(\operatorname{Gr}^W_3H^3(Z_f,\mathbb C)\bigr),
\end{equation}
由 $H^{2,1}_{\mathrm{poly}}(V)\subset H^{2,1}(V)\to H^{2,1}(\operatorname{Gr}^W_3H^3(Z_f,\mathbb C))$ 诱导，最后一个映射来自上述正合列。因为其证明尚未写出，这里直接假设式~\eqref{eq:polynomial-affine-hodge-assumption} 是同构。

在此假设下，利用 $H^{2,1}_{\mathrm{poly}}(V)$ 构造 Hodge 结构变分的自然办法，就是使用模空间 $\overline{\mathcal M}_{\mathrm{simp}}(V)$ 上 $\operatorname{Gr}^W_3H^3(Z_f,\mathbb C)$ 的自然极化 Hodge 结构变分。将所得丛记为 $\mathcal T\mathcal H^V_{\mathrm{simp}}$；这时 $\mathcal T$ 表示取交所用的环面。现在可以表述三维簇的基于 Hodge 理论的\IndexLocation{indexloc:0791:295}\IndexLocation{indexloc:0795:295}\IndexLocation{indexloc:0797:295}环面镜像猜想。
\begin{conjecture}\label{conj:toric-hodge-mirror-general}\index{镜像猜想!Hodge 理论|hyperpage}
设 $V$、$V^\circ$ 如上为互为 Batyrev 镜像的三维 Calabi--Yau 环面超曲面，并假设式~\eqref{eq:toric-quantum-subring-assumption}、\eqref{eq:polynomial-affine-hodge-assumption} 成立，且 $X_{\Sigma^\circ}$ 光滑。则镜像映射提升为 $p_0\in\overline{\mathcal M}_{\mathrm{simp}}(V)$ 和 $q_0\in\overline{\mathcal{KM}}_{\mathrm{toric}}(V^\circ)$ 的邻域上，丛 $\mathcal T\mathcal H^V_{\mathrm{simp}}$ 与 $\mathcal T\mathcal H^{V^\circ}_{\mathrm{toric}}$ 的同构。此外，此同构保持来自 $\nabla^{\mathrm{GM}}$ 与 $\nabla^{\mathrm{middle}}$ 的极化复 Hodge 结构变分，并将 $\widetilde\Omega,g_0$ 分别映到 $1,[pt]$。最后，在 $p_0,q_0$ 处的 $\operatorname{Gr}^{\mathcal W}_4$ 上，诱导同构就是单项式—\IndexLocation{indexloc:0910:295}除子镜像映射。
\end{conjecture}
因为此猜想的表述包含许多假设，应将其看作“尚在推进中的工作”。但希望猜想~\ref{conj:toric-hodge-mirror-general} 已相当接近真实情形。一个令人鼓舞的事实是，式~\eqref{eq:toric-divisor-surjectivity} 满射时，此猜想化为猜想~\ref{conj:toric-hodge-mirror-surjective}。
\SourcePage{296}
最后考虑 $V$、$V^\circ$ 的维数大于 $3$ 的情形。尝试提出环面镜像猜想时，会立即遇到一些问题。主要问题是，$\dim V>3$ 时，极小 Calabi--Yau 簇 $V$ 可能有奇点，而我们只对光滑射影簇定义了 Gromov--Witten 不变量。这大概只是技术困难；最终希望 Gromov--Witten 理论能适用于更广泛的簇。此外，若丛同构 $\mathcal H^V_{\mathrm{simp}}\simeq\mathcal H^{V^\circ}_{\mathrm{toric}}$ 保持 Hodge 滤过，则会有
\begin{equation}\label{eq:toric-mirror-hodge-equality}
h^{p,q}(V)=h^{q,q}(V^\circ),\qquad p+q=\dim V.
\end{equation}
\hyperref[chap:mirror-constructions]{第四章}已证明 $q=1$ 的情形；对于其他 $q$，已知的最好结果是联系 $V$、$V^\circ$ 的弦理论 Hodge 数的公式~\eqref{eq:complete-intersection-string-hodge}
\[
h^{p,q}_{\mathrm{st}}(V)=h^{q,q}_{\mathrm{st}}(V^\circ),\qquad p+q=\dim V,
\]
见 \SourceCite{BD}。$V$、$V^\circ$ 光滑时，此式化为式~\eqref{eq:toric-mirror-hodge-equality}；但 $V$ 有奇点时，$h^{p,q}_{\mathrm{st}}(V)$ 可能不同于 $h^{p,q}(V)$。

当然，可假设 $V$、$V^\circ$ 光滑来避开这一困难，但这会限制猜想的适用范围。即便如此，仍需知道如何定义 $\mathcal H^V_{\mathrm{simp}}\simeq\mathcal H^{V^\circ}_{\mathrm{toric}}$。仅有 Hodge 数等式~\eqref{eq:toric-mirror-hodge-equality} 还不够，必须有实际同构 $H^{p,q}(V)\simeq H^{q,q}(V^\circ)$，其中 $p+q=\dim V$。这与\hyperref[subsec:higher-dimensional-mirror]{第~\ref*{subsec:higher-dimensional-mirror} 小节}遇到的问题相同。

因此，不在一般情形表述基于 Hodge 理论的环面镜像猜想。猜想~\ref{conj:toric-hodge-mirror-general} 的某个版本可能在高维成立，但要给出精确猜想，还需更多工作。按 \SourceCite[第 8.2 讲]{Morrison7} 中基于 Hodge 理论的镜像对称猜想的表述，最终版本应涉及某些 Hodge 结构子变分，例如 $\mathcal T\mathcal H^{V^\circ}_{\mathrm{toric}}$。不过，找到这个猜想的精确版本，仍是代数几何中的开放问题。

\subsection{结语}\label{subsec:mirror-conjecture-conclusion}
本节严格定义了数学镜像对（定义~\ref{def:mathematical-mirror-pair}、\ref{def:higher-mathematical-mirror-pair}），并考虑了基于 Hodge 理论的环面镜像猜想的不同版本（猜想~\ref{conj:toric-hodge-mirror-surjective}、\ref{conj:toric-hodge-mirror-general}）。这些表明，对于三维簇，已经较好地理解镜像对称的含义，虽然猜想~\ref{conj:toric-hodge-mirror-general} 仍需进一步研究。但进入高维以后，猜想中还有许多部分未予明确。即便在环面情形，距离基于 Hodge 理论的镜像猜想的确定版本，也还有很长的路。

不过，近期文献中已出现若干“镜像定理”，包括五次三维簇镜像对称的完整证明，以及 Calabi--Yau 环面超曲面和更一般的 Calabi--Yau 环面完全交的结果。后者使用\hyperref[sec:toric-complete-intersections]{第~\ref*{sec:toric-complete-intersections} 节}的 Batyrev--Borisov 镜像构造。\footnote{环面簇中 Calabi--Yau 完全交的讨论和例子，见 \SourceCite{BvS,LTe,HKTY2,Givental2,Givental4}，以及\hyperref[subsec:other-gw-examples]{第~\ref*{subsec:other-gw-examples} 小节}列出的参考文献。}此外，Givental 提出了适用于某些非 Calabi--Yau 环面完全交的“镜像定理”；\SourceCite{LLY} 的“镜像原理”则适用于许多超出 Calabi--Yau 三维簇及其镜像的情形。
\SourcePage{297}
\hyperref[chap:mirror-theorem]{第十一章}将看到，尽管这些“镜像定理”没有直接提及这里讨论的基于 Hodge 理论的镜像对称版本，它们仍有共同的基本思想：把计算 Gromov--Witten 不变量的困难问题转化为更简单的问题。维数大于 $3$ 时，尚不清楚这些结果与 Hodge 理论的关系。因此，正如在维数大于 $3$ 时还没有确定的“基于 Hodge 理论的镜像猜想”，也还没有该情形下确定的“镜像定理”。希望代数几何学家继续探索镜像对称的数学理论，证明更一般的“镜像定理”，提出更精确的“镜像猜想”，最终汇合为对镜像对称现象的广泛理解。

本章使我们更清楚地认识到，镜像对称的数学版本应具有怎样的形式。$\dim V>3$ 时没有精确的“镜像猜想”表述，固然令人遗憾，但这不应妨碍考察目前已经证明的“镜像定理”。带着这一目标，接下来\hyperref[chap:localization]{第九章}、\hyperref[chap:quantum-differential-equations]{第十章}将发展一些必要工具，然后在\hyperref[chap:mirror-theorem]{第十一章}研究若干“镜像定理”。

\chapter{局部化}\label{chap:localization}\index{局部化|hyperpage}
\SourcePage{298}
本章定义等变上同调，并说明强有力的局部化方法。这里的思想与技巧，不仅将完成\hyperref[chap:gromov-witten]{第七章}开始的部分工作，还将提供\hyperref[chap:mirror-theorem]{第十一章}证明镜像定理所需的一些工具。

先在\hyperref[sec:localization-theorem]{第~\ref*{sec:localization-theorem} 节}介绍等变上同调的定义和基本性质，并讨论局部化定理 \SourceCite{AB} 及其各种推论。基本思想是，当光滑簇具有群作用时，簇上的上同调类“局部化”为固定点轨迹上的类，从而容易开展计算。将会看到，固定点集的法丛的等变 Euler 类起着重要作用。作为应用，先描述 $\mathbb P^r$ 的局部化，再用它确定一般五次三维簇上的直线数。就计数几何而言，局部化首次由文献 \SourceCite{ES2} 利用 Hilbert 概形上的环面作用引入。

本章主要工作从\hyperref[sec:stable-map-localization]{第~\ref*{sec:stable-map-localization} 节}开始，将局部化应用于模空间 $\overline M_{0,n}(\mathbb P^r,d)$。射影空间 $\mathbb P^r$ 有代数环面 $T=(\mathbb C^*)^{r+1}$ 的自然作用，从而诱导 $T$ 在该模空间上的作用。按照 \SourceCite{Kontsevich2}，将学习如何用某些带标记的图 $\Gamma$ 描述固定点集的分支，以及如何计算相应固定点分支的等变法丛。然后，在\hyperref[subsec:multiple-cover-localization]{第~\ref*{subsec:multiple-cover-localization}} 和\hyperref[subsec:nodal-quintic-double-covers]{第~\ref*{subsec:nodal-quintic-double-covers} 小节}继续\hyperref[subsec:instanton-numbers]{第~\ref*{subsec:instanton-numbers} 小节}开始的 Calabi--Yau 三维簇 Gromov--Witten 不变量研究。先计算三维簇中刚性嵌入的光滑有理曲线的 $d$ 重覆盖所作的贡献，再专门考察五次三维簇，研究五次结点有理曲线的二重覆盖这一更微妙的问题。这样将严格证明\hyperref[chap:gromov-witten]{第七章}关于瞬子数 $n_{10}$ 的论断。

本章最后一节简要介绍等变 Gromov--Witten 不变量与虚拟基本类的局部化。

\section{局部化定理}\label{sec:localization-theorem}
先定义\emph{等变上同调}。这里给出拓扑构造；代数构造参见 \SourceCite{EG}。设 $G$ 为连通 Lie 群，其分类空间由主 $G$ 丛 $EG\to BG$ 给出，总空间 $EG$ 可缩。此 $G$ 丛在同伦等价意义下唯一确定。我们将使用代数环面 $G=(\mathbb C^*)^n$；此时 $BG=(\mathbb{CP}^\infty)^n$，可取
\[
EG=\pi_1^*S^\times\times_{BG}\cdots\times_{BG}\pi_n^*S^\times,
\]
其中 $\pi_i:BG\to\mathbb{CP}^\infty$ 为第 $i$ 个投影，$S$ 是 $\mathbb{CP}^\infty$ 上截面层为 $\mathcal O_{\mathbb{CP}^\infty}(-1)$ 的重言丛，$S^\times$ 表示去掉零截面。
\SourcePage{299}
\subsection{等变上同调}\label{subsec:equivariant-cohomology}
若拓扑空间 $X$ 具有 $G$ 作用，令 $X_G=X\times_GEG$；它是 $BG$ 上纤维为 $X$ 的丛。为以后使用，令 $i_X:X\hookrightarrow X_G$ 表示一条纤维的包含映射。
\begin{definition}\label{def:equivariant-cohomology}
$X$ 的\emph{等变上同调}定义为
\[
H_G^*(X)=H^*(X_G),
\]
其中 $H^*(X_G)$ 为 $X_G$ 的通常上同调。
\end{definition}
\index{等变上同调|hyperpage}
等变上同调具有通常上同调的许多性质，例如存在平坦等变拉回与固有等变推前。这些运算满足通常上同调中相应运算的常见性质 \SourceCite{AB}。注意 $H_G^*(\mathrm{point})=H^*(BG)$。沿 $X\to\mathrm{point}$ 拉回可见，一般的 $H_G^*(X)$ 是 $H^*(BG)$ 模，因此可将 $H^*(BG)$ 看作等变上同调的系数环。此外，$i_X:X\to X_G$ 诱导“遗忘”映射 $i_X^*:H_G^*(X)\to H^*(X)$。

主要考虑 $G=(\mathbb C^*)^n$ 的情形。令 $M(G)$ 为环面 $G$ 的特征标群。每个 $\rho\in M(G)$ 给出一维向量空间 $\mathbb C_\rho$，$G$ 通过 $\rho$ 作用于其上。若 $L_\rho=(\mathbb C_\rho)_G$ 为 $BG$ 上的相应线丛，则 $\rho\mapsto c_1(L_\rho)$ 定义同构 $\psi:M(G)\simeq H^2(BG,\mathbb Z)$，进而诱导复系数环同构 $\operatorname{Sym}(M(G))\otimes\mathbb C\simeq H^*(BG,\mathbb C)$。称 $\psi(\rho)$ 为 $\rho$ 的\emph{权}，以下未注明系数的上同调均取复系数。

特别地，若 $\rho_i$ 是由 $\rho_i(t_1,\ldots,t_n)=t_i$ 定义的 $G=(\mathbb C^*)^n$ 的特征标，令 $\lambda_i$ 表示其权。于是得到同构
\begin{equation}\label{eq:equivariant-coefficients}
H_G^*(\mathrm{point})=H^*(BG)\simeq\mathbb C[\lambda_1,\ldots,\lambda_n].
\end{equation}
将线丛 $L_{\rho_i}$ 记为 $\mathcal O(\lambda_i)$，则 $\lambda_i=c_1(\mathcal O(\lambda_i))$。在式~\eqref{eq:equivariant-coefficients} 的等同下，$i_{\mathrm{point}}^*$ 可看作将全部 $\lambda_i$ 映到 $0$ 的“非等变极限”。更一般地，对任意 $X$，映射 $i_X^*$ 都可看作非等变极限，特别会将全部 $\lambda_i$ 映到 $0$。

一般而言，对于 $G$ 空间 $X$，\emph{等变向量丛}是 $X$ 上的向量丛 $E$，使 $G$ 在 $X$ 上的作用提升为 $E$ 上的作用，并在纤维上为线性作用。此时 $E_G$ 是 $X_G$ 上的向量丛；\emph{等变 Chern 类} $c_k^G(E)\in H_G^*(X)$ 定义为通常 Chern 类 $c_k(E_G)$。若 $E$ 的秩为 $r$，最高 Chern 类 $c_r^G(E)$ 称为 $E$ 的\emph{等变 Euler 类}，记为
\[
\operatorname{Euler}_G(E)\in H_G^*(X).
\]
\hyperref[chap:gromov-witten]{第七章}中的 Euler 类是同调类，即最高 Chern 类的 Poincaré 对偶；从这里开始，将通常或等变 Euler 类看作上同调类更为方便。
\index{等变向量丛|hyperpage}\index{Chern 类!等变|hyperpage}\index{Euler类@Euler 类!等变|hyperpage}
\SourceExampleNumber{9.1.1}{1}
\begin{example}\label{ex:diagonal-equivariant-chern}
$G=(\mathbb C^*)^n$ 在 $\mathbb C^n$ 上的对角作用，给出 $X=\mathrm{point}$ 上的等变向量丛 $E$，满足 $E_G\simeq\mathcal O(\lambda_1)\oplus\cdots\oplus\mathcal O(\lambda_n)$。因此，$\lambda_1,\ldots,\lambda_n$ 是该\IndexLocation{indexloc:1434:299}表示的权。因为 $\lambda_i=c_1(\mathcal O(\lambda_i))$，有
\begin{equation}\label{eq:equivariant-chern-symmetric}
c_k^G(E)=\sigma_k(\lambda_1,\ldots,\lambda_n)\in\mathbb C[\lambda_1,\ldots,\lambda_n],
\end{equation}
其中 $\sigma_k$ 为第 $k$ 个初等对称函数。
\end{example}
\SourcePage{300}
\subsection{Atiyah 与 Bott 的定理}\label{subsec:atiyah-bott-localization}
接下来需要局部化定理。设环面 $T=(\mathbb C^*)^n$ 作用于紧复流形 $X$。由 \SourceCite{Iversen}，固定点轨迹 $X^T$ 是光滑连通分支 $Z_j$ 的并。令 $i_j:Z_j\hookrightarrow X$ 为包含映射，$N_j$ 为 $Z_j$ 在 $X$ 中的法丛。$N_j$ 是等变向量丛，故具有等变 Euler 类
\[
\operatorname{Euler}_T(N_j)\in H_T^*(Z_j).
\]
等变包含映射 $i_j$ 诱导 $i_j^*:H_T^*(X)\to H_T^*(Z_j)$。此外，$Z_j$ 是 $X$ 的子流形，故还有 Gysin 映射
\[
i_{j!}:H_T^*(Z_j)\longrightarrow H_T^*(X).
\]
这是因为诱导映射 $(Z_j)_T\to X_T$ 可以由子流形的嵌入作有限维近似。Gysin 映射具有如下性质：对任意 $\alpha\in H_T^*(Z_j)$，
\begin{equation}\label{eq:equivariant-self-intersection}
i_j^*\circ i_{j!}(\alpha)=\alpha\cup\operatorname{Euler}_T(N_j).
\end{equation}
这些断言的证明见 \SourceCite[第 6 章附录 A]{Audin}。最后，$H_T^*(Z_j)$ 是 $H^*(BT)\simeq\mathbb C[\lambda_1,\ldots,\lambda_n]$ 模。若
\[
\mathcal R_T\simeq\mathbb C(\lambda_1,\ldots,\lambda_n)
\]
是 $H^*(BT)$ 的分式域，就得到局部化 $H_T^*(Z_j)\otimes_{H^*(BT)}\mathcal R_T$。\SourceCite{AB} 的一个重要观察是，$\operatorname{Euler}_T(N_j)$ 在此局部化中可逆。

在上述设置下，有 \SourceCite{AB} 证明的如下\emph{局部化定理}。
\begin{proposition}\label{prop:atiyah-bott-localization}
存在同构
\[
H_T^*(X)\otimes\mathcal R_T\xrightarrow{\sim}\bigoplus_jH_T^*(Z_j)\otimes\mathcal R_T,
\]
由映射 $\alpha\mapsto(i_j^*(\alpha)/\operatorname{Euler}_T(N_j))_j$ 诱导。其逆映射由 $(\alpha_j)_j\mapsto\sum_ji_{j!}(\alpha_j)$ 诱导。特别地，对任意 $\alpha\in H_T^*(X)\otimes\mathcal R_T$，有
\[
\alpha=\sum_ji_{j!}\left(\frac{i_j^*(\alpha)}{\operatorname{Euler}_T(N_j)}\right).
\]
\end{proposition}
\index{局部化定理|hyperpage}\index{Atiyah--Bott 定理|hyperpage}
看看此定理对 $\mathbb P^r$ 给出什么。
\SourceExampleNumber{9.1.2}{1}
\begin{example}\label{ex:projective-equivariant-localization}
考虑 $T=(\mathbb C^*)^{r+1}$ 在 $X=\mathbb P^r$ 上的作用
\[
(t_0,\ldots,t_r)\cdot(x_0,\ldots,x_r)=(t_0^{-1}x_0,\ldots,t_r^{-1}x_r).
\]
选择逆元，是为了使 $(t_0,\ldots,t_r)$ 在齐次形式 $x_i\in H^0(\mathcal O_{\mathbb P^r}(1))$ 上以乘 $t_i$ 的方式作用。

先定义等变超平面类。$\mathbb P_T^r$ 是 $BT$ 上向量丛 $\mathcal O(-\lambda_0)\oplus\cdots\oplus\mathcal O(-\lambda_r)$ 的射影化。因此，$\mathbb P_T^r=\mathbb P(E_T^*)$，其中 $E^*$ 为例题~\ref{ex:diagonal-equivariant-chern} 中丛 $E$ 的对偶，$\mathbb P$ 表示射影化。这给出重言线丛 $\mathcal O_{\mathbb P_T^r}(1)$，令
\[
p=c_1(\mathcal O_{\mathbb P_T^r}(1))\in H_T^*(\mathbb P^r).
\]
可以验证，$p$ 是等变 Chern 类 $c_1^T(\mathcal O_{\mathbb P^r}(1))$，故称为\emph{等变超平面类}。
\SourcePage{301}
接下来计算 $\mathbb P^r$ 的等变上同调。射影丛上同调的标准公式，例如 \SourceCite[第 606 页]{GH}，给出
\[
(-p)^{r+1}-c_1(E_T^*)(-p)^r+c_2(E_T^*)(-p)^{r-1}+\cdots+(-1)^{r+1}c_{r+1}(E_T^*)=0.
\]
\SourceCite{GH} 中记为 $T$ 的丛，在这里是 $\mathcal O_{\mathbb P_T^r}(-1)$，故 $c_1(T)=-p$。由式~\eqref{eq:equivariant-chern-symmetric}，$c_k(E_T^*)=(-1)^kc_k^T(E)=(-1)^k\sigma_k(\lambda_0,\ldots,\lambda_r)$，因此在 $H_T^*(\mathbb P^r)$ 中有
\begin{equation}\label{eq:projective-equivariant-relation}
\prod_{j=0}^r(p-\lambda_j)=0.
\end{equation}
于是，$\mathbb P^r$ 的等变上同调为
\[
H_T^*(\mathbb P^r)=\mathbb C[p,\lambda_0,\ldots,\lambda_r]\Big/\left(\prod_{j=0}^r(p-\lambda_j)\right).
\]
形式地令 $\lambda_i=0$，再把 $p$ 换成通常超平面类 $H=i_{\mathbb P^r}^*p$，就恢复 $\mathbb P^r$ 通常上同调的描述。

为应用局部化定理，先理解 $T$ 作用的固定点。共有 $r+1$ 个固定点 $q_j$，照常排序，使 $q_j$ 的第 $j$ 个坐标非零，其余为 $0$。注意 $(q_j)_T=\mathbb P(\mathcal O(-\lambda_j))$，所以 $\mathcal O_{\mathbb P_T^r}(-1)$ 在 $(q_j)_T$ 上限制为 $\mathcal O(-\lambda_j)$；相应地，$\mathcal O_{\mathbb P_T^r}(1)$ 限制为 $\mathcal O(\lambda_j)$。因此
\begin{equation}\label{eq:projective-hyperplane-fixed-restriction}\index{等变超平面类|hyperpage}
i_j^*(p)=\lambda_j.
\end{equation}
换句话说，$T$ 在 $\mathbb Cq_j$ 上的作用具有权 $-\lambda_j$，因为按定义 $(t_0,\ldots,t_r)\cdot q_j=t_j^{-1}q_j$。更一般地，$H_T^*(\mathbb P^r)$ 的任何元素都可写为关于 $p$、系数在 $\mathbb C[\lambda_0,\ldots,\lambda_r]$ 中的多项式 $F(p)$，式~\eqref{eq:projective-hyperplane-fixed-restriction} 给出
\begin{equation}\label{eq:projective-polynomial-fixed-restriction}
i_j^*(F(p))=F(\lambda_j).
\end{equation}
因为 $q_j$ 是一个点，法丛 $N_j$ 就是切空间
\[
N_j=T_{q_j}\mathbb P^r=\operatorname{Hom}(\mathbb Cq_j,\mathbb C^{r+1}/\mathbb Cq_j).
\]
所以 $T$ 在 $N_j$ 上的表示具有权 $\lambda_j-\lambda_k$（$k\ne j$），故
\begin{equation}\label{eq:projective-fixed-euler}
\operatorname{Euler}_T(N_j)=\prod_{k\ne j}(\lambda_j-\lambda_k).
\end{equation}
现在计算出命题~\ref{prop:atiyah-bott-localization} 的同构为
\begin{equation}\label{eq:projective-localization-map}
p\longmapsto\left(\frac{\lambda_j}{\prod_{i\ne j}(\lambda_j-\lambda_i)}\right)_{j=0}^r.
\end{equation}
为描述逆映射，令
\begin{equation}\label{eq:projective-fixed-dual-class}
\phi_j=\prod_{k\ne j}(p-\lambda_k).
\end{equation}
于是式~\eqref{eq:projective-polynomial-fixed-restriction} 给出
\begin{equation}\label{eq:projective-dual-class-restriction}
i_k^*(\phi_j)=\begin{cases}\operatorname{Euler}_T(N_j),&k=j,\\0,&k\ne j.\end{cases}
\end{equation}
\SourcePage{302}
因此，命题~\ref{prop:atiyah-bott-localization} 的最后一条断言给出
\[
\phi_j=\sum_{k=0}^ri_{k!}\left(\frac{i_k^*(\phi_j)}{\prod_{l\ne k}(\lambda_k-\lambda_l)}\right)
=i_{j!}\left(\frac{i_j^*(\phi_j)}{\prod_{l\ne j}(\lambda_j-\lambda_l)}\right)=i_{j!}(1).
\]
利用 $H_T^*(q_j)=\mathbb C[\lambda_0,\ldots,\lambda_r]$ 及投影公式 $i_{j!}(i_j^*(\alpha)\cup1)=\alpha\cup i_{j!}(1)$，容易得到式~\eqref{eq:projective-localization-map} 的逆映射为
\[
(\alpha_j)_{j=0}^r\longmapsto\sum_{j=0}^r\alpha_j\cup\phi_j=\sum_{j=0}^r\alpha_j\phi_j.
\]
所以 $\phi_0,\ldots,\phi_r$ 构成 $H_T^*(\mathbb P^r)\otimes\mathcal R_T$ 作为 $\mathcal R_T$ 向量空间的一组基。

最后，取该空间中的任意元素，如上写为关于 $p$、系数在 $\mathcal R_T=\mathbb C(\lambda_0,\ldots,\lambda_r)$ 中的多项式 $F(p)$。再用式~\eqref{eq:projective-polynomial-fixed-restriction}，局部化定理的最后一条断言给出恒等式
\[
F(p)=\sum_{j=0}^rF(\lambda_j)\frac{\prod_{k\ne j}(p-\lambda_k)}{\prod_{k\ne j}(\lambda_j-\lambda_k)}.
\]
这表明，命题~\ref{prop:atiyah-bott-localization} 可以看作 Lagrange 插值公式的广泛推广。
\end{example}
通过向一个点推前，命题~\ref{prop:atiyah-bott-localization} 给出积分公式。对于具有 $T$ 作用的簇 $X$，平凡映射 $X\to\mathrm{point}$ 诱导等变投影 $\pi_X:X_T\to BT$。当 $X$ 完备时，推前 $\pi_{X!}$ 由沿纤维积分给出，在这里称为\emph{等变积分}。因此，常将 $\pi_{X!}$ 写为
\[
\int_{X_T}:H_T^*(X)\longrightarrow H^*(BT).
\]
利用等变积分，得到命题~\ref{prop:atiyah-bott-localization} 的如下推论。
\begin{corollary}\label{cor:localization-integration}
对于任意 $\alpha\in H_T^*(X)\otimes\mathcal R_T$，有
\[
\int_{X_T}\alpha=\sum_j\int_{(Z_j)_T}\frac{i_j^*(\alpha)}{\operatorname{Euler}_T(N_j)}.
\]
\end{corollary}
\index{等变积分|hyperpage}
命题~\ref{prop:atiyah-bott-localization} 和推论~\ref{cor:localization-integration} 给出\emph{局部化}这一概念：将一个类等变地限制到各固定点分支，并除以该分支法丛的等变 Euler 类。这样，一个类的等变积分的计算，就归结为其限制在全部固定点分支上的等变积分之和。

最终要将理论用于轨形 $\overline M_{0,n}(\mathbb P^r,d)$，故实际上需要推论~\ref{cor:localization-integration} 对光滑代数叠的稍加变化的版本。如\hyperref[chap:gromov-witten]{第七章}所述，代数叠的交理论参见 \SourceCite{Vistoli}，代数叠的介绍参见 \SourceCite{BEFFGK}。

设轨形 $X$ 为光滑代数叠的底层簇。即 $X$ 具有\hyperref[app:singular-varieties]{附录 A} 所讨论的局部商图 $U/H$，其中 $U$ 光滑，$H$ 为 $GL(n,\mathbb C)$ 中作用于 $U$ 的小群。在相应光滑代数叠上有 $T$ 作用时，可以在 $T$ 与 $X$ 上局部地将其实现为 $U$ 上的作用。这样，通过局部在 $U$ 中工作，仍能定义固定点轨迹 $Z_j$、包含映射 $i_j:Z_j\hookrightarrow X$ 和等变法丛。因此得到如下代数叠版本。
\SourcePage{303}
\begin{corollary}\label{cor:orbifold-localization-integration}
设 $X$ 为具有 $T$ 作用的光滑代数叠的底层轨形。若 $\alpha\in H_T^*(X)\otimes\mathcal R_T$，则
\[
\int_{X_T}\alpha=\sum_j\int_{(Z_j)_T}\frac{i_j^*(\alpha)}{a_j\operatorname{Euler}_T(N_j)},
\]
其中 $a_j$ 是 $Z_j$ 一般点所对应对象的自同构群的阶。
\end{corollary}

\subsection{Chern 类的多项式}\label{subsec:chern-polynomial-localization}
许多应用中，可更具体地表述推论~\ref{cor:localization-integration}。设具有 $T$ 作用的光滑簇 $X$ 上有等变向量丛 $E$，希望计算 $\int_XP(c_k(E))$，其中 $P$ 是 $E$ 的 Chern 类的多项式。将 $c_k(E)$ 替换为\hyperref[subsec:equivariant-cohomology]{第~\ref*{subsec:equivariant-cohomology} 小节}定义的等变 Chern 类 $c_k^T(E)\in H_T^*(X)$，并沿包含映射 $i_{\mathrm{point}}$ 拉回，就将问题置于等变背景中。为说明这给出 $\int_XP(c_k(E))$，考虑交换图
\begin{equation}\label{eq:equivariant-integral-forgetting-diagram}
\begin{array}{ccc}
X&\longrightarrow&\mathrm{point}\\
\big\downarrow\scriptstyle i_X&&\big\downarrow\scriptstyle i_{\mathrm{point}}\\
X_T&\longrightarrow&BT.
\end{array}
\end{equation}
它意味着 $i_{\mathrm{point}}^*\circ\int_{X_T}=\int_X\circ i_X^*$，故
\[
i_{\mathrm{point}}^*\int_{X_T}P(c_k^T(E))
=\int_Xi_X^*P(c_k^T(E))=\int_XP(c_k(E)),
\]
正如所断言的。关键是 $\int_{X_T}P(c_k^T(E))$ 为可用局部化计算的等变积分。更精确地，将此式与推论~\ref{cor:localization-integration} 结合，就能根据 $E$ 在 $X$ 的 $T$ 固定点集上的行为，计算 $\int_XP(c_k(E))$。

为更明确地说明，考察固定点集 $X^T$ 有限的情形。每个点 $Z_j\in X^T$ 上，限制 $E|_{Z_j}$ 分解为 $T$ 的特征标 $\chi_j^1,\ldots,\chi_j^s$。每个 $\chi_j^k$ 都是已定义的基本特征标 $\rho_1,\ldots,\rho_n$ 的线性组合。这里使用特征标的加法记号，即 $(\rho_1+\rho_2)(t)=\rho_1(t)\rho_2(t)$。将此线性组合写为 $\chi_j^k=\ell_j^k(\rho_1,\ldots,\rho_n)$，简记为 $\chi_j^k=\ell_j^k(\rho)$。

因为 $\rho_i$ 的权为 $\lambda_i$，$\chi_j^k=\ell_j^k(\rho)$ 的权为 $\ell_j^k(\lambda)$，即在线性组合 $\ell_j^k(\rho)$ 中将 $\rho_i$ 替换为 $\lambda_i$ 的结果。因此，$T$ 在 $E|_{Z_j}$ 上的作用具有权 $\ell_j^k(\lambda)$，从而
\[
i_j^*(c_k^T(E))=c_k^T(E|_{Z_j})=\sigma_k(\ell_j^1(\lambda),\ldots,\ell_j^s(\lambda)).
\]
此外，因 $Z_j$ 是点，法丛 $N_j$ 就是切空间 $T_{Z_j}X$。若 $T$ 在该切空间上的特征标为 $t_j^1(\rho),\ldots,t_j^d(\rho)$，其中 $d=\dim X$，则
\begin{equation}\label{eq:isolated-fixed-tangent-euler}
\operatorname{Euler}_T(N_j)=\prod_kt_j^k(\lambda).
\end{equation}
现在对推论~\ref{cor:localization-integration} 的右端施加 $i_{\mathrm{point}}^*$，沿用上述记号，得到如下陈述。
\SourcePage{304}
\begin{proposition}\label{prop:chern-polynomial-localization}
设 $T$ 作用于完备光滑簇 $X$，其固定点 $Z_j$ 孤立，且 $E$ 为 $X$ 上的等变向量丛。若 $P(c_k(E))$ 为 $E$ 的 Chern 类中上同调次数为 $2\dim X$ 的多项式，则
\[
\int_XP(c_k(E))=\sum_j\frac{P\bigl(\sigma_k(\ell_j^1(\lambda),\ldots,\ell_j^s(\lambda))\bigr)}{\prod_kt_j^k(\lambda)},
\]
其中 $\sigma_k$ 为第 $k$ 个初等对称函数。
\end{proposition}
这就是 \SourceCite{ES2} 用于计数应用的表述。若对一般分支 $Z_j$ 重复论证，就得到 \SourceCite{Kontsevich2} 中的表述。

也可将命题~\ref{prop:chern-polynomial-localization} 推广到光滑代数叠：在上述论证中，将推论~\ref{cor:localization-integration} 替换为推论~\ref{cor:orbifold-localization-integration}。结果是，若用轨形描述代数叠，则上述求和中的每一项，都须除以相应 $Z_j$ 对应对象的自同构群的阶。下面\hyperref[subsec:multiple-cover-localization]{第~\ref*{subsec:multiple-cover-localization} 小节}将举例说明。
\SourceExampleNumber{9.1.3}{1}
\begin{example}\label{ex:quintic-lines-localization}
用命题~\ref{prop:chern-polynomial-localization} 计算一般\IndexLocation{indexloc:1182:304}\IndexLocation{indexloc:1188:304}五次三维簇 $V\subset\mathbb P^4$ 上的直线数。例题~\ref{ex:quintic-lines-normal-cone} 曾将 $V$ 上的直线概形实现为 $\mathbb P^4$ 中直线的 Grassmann 簇 $G(2,5)$ 上向量丛 $\operatorname{Sym}^5U^*$ 的一个截面的零点轨迹。$U$ 为 Grassmann 簇上的秩 $2$ 重言子丛，直线 $\ell$ 上的纤维 $U_\ell$ 是射影化为 $\ell$ 的 $\mathbb C^5$ 的二维子空间。$V\subset\mathbb P^4$ 的方程诱导秩 $6$ 丛 $\operatorname{Sym}^5U^*$ 的截面 $s$；例题~\ref{ex:quintic-gw-euler-integral} 表明，直线数是 Euler 类 $c_6(\operatorname{Sym}^5U^*)$ 的次数，即 Gromov--Witten 不变量
\[
\langle I_{0,0,1}\rangle=\int_{G(2,5)}c_6(\operatorname{Sym}^5U^*).
\]
在命题~\ref{prop:chern-polynomial-localization} 中取 $E=\operatorname{Sym}^5U^*$，$P(c_k(E))$ 为其 Euler 类；使用 $(\mathbb C^*)^5$ 在 $\mathbb C^5$ 上的自然作用
\[
(t_1,\ldots,t_5)\cdot(x_1,\ldots,x_5)=(t_1^{-1}x_1,\ldots,t_5^{-1}x_5).
\]
它诱导 $T=(\mathbb C^*)^5$ 在 $G(2,5)$ 上的作用，共有 $10$ 个孤立固定点 $L_I$，对应 $\mathbb P^4$ 的 $10$ 条坐标直线。每个 $L_I$ 由 $I\subset\{1,\ldots,5\}$、$|I|=2$ 标记，并由所有 $j\notin I$ 对应的方程 $x_j=0$ 定义。$L_I$ 的法丛就是 $G(2,5)$ 在 $L_I$ 处的切空间，需要计算它的等变 Euler 类。与式~\eqref{eq:isolated-fixed-tangent-euler} 一样，将切空间的 $T$ 作用分解为特征标，把每个 $\rho_i$ 替换为 $\lambda_i$，然后相乘。

Grassmann 簇的切丛为 $\operatorname{Hom}(U,\mathbb C^5/U)$。关键是，在 $L_I\in G(2,5)$ 处，$U|_{L_I}$ 的特征标为 $-\rho_i$（$i\in I$）。因此，$T$ 在该切空间上的特征标为
\[
\{\rho_i-\rho_j:i\in I,\ j\notin I\},
\]
其等变 Euler 类为
\[
\prod_{i\in I,\,j\notin I}(\lambda_i-\lambda_j).
\]
\SourcePage{305}
另一个所需对象是 $i_I^*(\operatorname{Sym}^5U^*)$ 的等变 Euler 类，其中 $i_I:\{L_I\}\hookrightarrow G(2,5)$ 为包含映射。用 $U^*|_{L_I}$ 的特征标计算；如上，这些特征标为 $\rho_i$（$i\in I$）。因此，写 $I=\{i_1,i_2\}$，得到
\[
\operatorname{Euler}_T\bigl(i_I^*(\operatorname{Sym}^5U^*)\bigr)
=\prod_{a=0}^5\bigl(a\lambda_{i_1}+(5-a)\lambda_{i_2}\bigr).
\]
命题~\ref{prop:chern-polynomial-localization} 现在给出
\[
\int_{G(2,5)}c_6(\operatorname{Sym}^5U^*)
=\sum_{|I|=2}\frac{\prod_{a=0}^5(a\lambda_{i_1}+(5-a)\lambda_{i_2})}{\prod_{i\in I}\prod_{j\notin I}(\lambda_i-\lambda_j)},
\]
仍取 $I=\{i_1,i_2\}$。正如应有的那样，全部 $\lambda_i$ 消去；计算得到答案 $2875$，即一般五次三维簇的直线数。
\end{example}
刚完成的计算涉及将 $\lambda_i$ 的 $10$ 个有理式之和化简；每个有理式的分子、分母都是 $6$ 个关于 $\lambda_i$ 的线性二项式的乘积。手算可能繁琐，计算机却能很快完成。此方法还可迅速进行射影空间中直线、二次曲线和扭曲三次曲线的 Hilbert 概形上的计算，因为这些都是光滑空间，其环面作用也很清楚。例如，完全交 Calabi--Yau 三维簇上的扭曲三次曲线的计算见 \SourceCite{ES2}。这些计算需要两个要素：首先，为所考察的射影曲线找到光滑模空间及 $T$ 作用，使固定点轨迹和权都可计算；其次，将待求数目表示为某个等变向量丛的 Chern 类多项式的次数，而且该丛在固定点处的权可计算。

\section{\texorpdfstring{$\overline M_{0,n}(\mathbb P^r,d)$}{稳定映射模空间}中的局部化}\label{sec:stable-map-localization}
为了求五次三维簇上 $d$ 次有理曲线的数目，Kontsevich 的见解 \SourceCite{Kontsevich2} 是：使用\hyperref[subsec:stable-maps]{第~\ref*{subsec:stable-maps} 小节}描述的稳定映射模空间，就可应用\hyperref[subsec:chern-polynomial-localization]{第~\ref*{subsec:chern-polynomial-localization} 小节}的方法。空间 $\overline M_{0,n}(\mathbb P^r,d)$ 不光滑，但 $\overline{\mathcal M}_{0,n}(\mathbb P^r,d)$ 是光滑代数叠，故推论~\ref{cor:orbifold-localization-integration} 适用。

\subsection{Kontsevich 的方法}\label{subsec:kontsevich-localization}\index{稳定映射!固定点图|hyperpage}
$T=(\mathbb C^*)^{r+1}$ 在 $\mathbb P^r$ 上的自然作用，诱导 $\overline M_{0,n}(\mathbb P^r,d)$ 上的 $T$ 作用。固定点轨迹及其等变法丛已在 \SourceCite{Kontsevich2} 中计算；$g>0$ 的情形，包括虚法丛的局部化公式，见 \SourceCite{GPa}。以下简述主要思想和结果。

$\overline M_{0,n}(\mathbb P^r,d)$ 的 $T$ 固定点由稳定映射 $(f,C,p_1,\ldots,p_n)$ 构成，$C$ 的每个分支 $C_i$ 或被 $f$ 映到 $\mathbb P^r$ 的某个 $T$ 固定点，或多重覆盖一条坐标直线。此外，每个标记点 $p_i$、$C$ 的每个结点、以及 $f$ 的每个分歧点，都映到 $T$ 固定点。因此，若 $C_i$ 通过 $f$ 以次数 $d_i$ 覆盖一条坐标直线，则可选取二者的齐次坐标，使覆盖由 $(x_0,x_1)\mapsto(x_0^{d_i},x_1^{d_i})$ 给出。
\SourcePage{306}
立即可见，该模空间上 $T$ 作用的固定点分支 $Z_j$ 可用组合数据描述。令 $q_0,\ldots,q_r$ 为 $\mathbb P^r$ 的通常 $T$ 固定点。对稳定映射 $(f,C,p_1,\ldots,p_n)$ 关联一棵树 $\Gamma$，其顶点 $v$ 与 $f^{-1}(\{q_0,\ldots,q_r\})$ 的连通分支 $C_v$ 一一对应。每个 $C_v$ 或为 $C$ 中一个点，或为其不可约分支的连通并。$\Gamma$ 的边 $e$ 对应 $C$ 中被 $f$ 映满某条坐标直线 $\ell_e\subset\mathbb P^r$ 的不可约分支 $C_e$。

树 $\Gamma$ 还带有各种标记。对每个顶点 $v$，关联由 $f(C_v)=q_{i_v}$ 定义的数 $i_v$，以及由所有满足 $p_i\in C_v$ 的 $i$ 组成的集合 $S_v$。对每条边 $e$，关联映射 $f|_{C_e}:C_e\to\ell_e$ 的次数 $d_e$。于是，固定点集 $\overline M_{0,n}(\mathbb P^r,d)^T$ 的连通分支，与带有标记 $i_v,S_v,d_e$ 并满足以下条件的连通树一一对应：
\begin{itemize}
\item 若边 $e$ 含顶点 $v,v'$，则 $i_v\ne i_{v'}$，且 $\ell_e$ 为连接 $q_{i_v}$ 与 $q_{i_{v'}}$ 的坐标直线。
\item $\sum_ed_e=d$。
\item $\{1,\ldots,n\}=\coprod_vS_v$。
\end{itemize}
称数据 $(\Gamma,i_v,S_v,d_e)$ 为稳定映射的\emph{图}，通常将全部数据简记为 $\Gamma$。回顾，顶点 $v$ 的\emph{价数} $\operatorname{val}(v)$ 是与 $v$ 相接的边数。\index{稳定映射!固定点图|hyperpage}\index{价数|hyperpage}

具有固定图 $\Gamma$ 的稳定映射，自然定义子叠
\[
\overline{\mathcal M}_\Gamma\subset\overline{\mathcal M}_{0,n}(\mathbb P^r,d).
\]
为研究它，取 $(f,C,p_1,\ldots,p_n)\in\overline{\mathcal M}_\Gamma$。对于 $C_v$ 是曲线的顶点 $v$，$C_v$ 有 $n(v)=|S_v|+\operatorname{val}(v)$ 个特殊点：其上有 $|S_v|$ 个标记点，另有 $\operatorname{val}(v)$ 个结点，使 $C_v$ 与包含 $v$ 的边 $e$ 对应的不可约分支 $C_e$ 相交。实际上，$C_v$ 连同这些 $n(v)$ 个点组成稳定曲线，给出 $\overline M_{0,n(v)}$ 的元素。

更仔细的描述如下。给定一组 $n(v)$ 点亏格零稳定曲线 $\{C_v\}\in\prod_{\dim C_v=1}\overline M_{0,n(v)}$，每当边 $e$ 含顶点 $v,v'$，就用一条新曲线 $C_e\simeq\mathbb P^1$ 的 $0,\infty$ 将 $C_v$ 与 $C_{v'}$ 连接，从而构造曲线 $C$。这可以系统地完成。再定义 $f:C\to\mathbb P^r$，将每个曲线 $C_v$ 收缩到 $q_{i_v}$，并令 $f|_{C_e}$ 为 $(x_0,x_1)\mapsto(x_0^{d_e},x_1^{d_e})$。得到态射
\[
\psi_\Gamma:\prod_{\dim C_v=1}\overline M_{0,n(v)}\longrightarrow\overline{\mathcal M}_\Gamma.
\]
$n(v)$ 对所有顶点都有定义，而且 $\dim C_v=1$，当且仅当 $C_v$ 含有被 $f$ 收缩的分支，也当且仅当 $n(v)\ge3$。

将上述乘积定义为 $\overline M_\Gamma$；若没有收缩分支，就令 $\overline M_\Gamma$ 为一个点。上述 $\psi_\Gamma:\overline M_\Gamma\to\overline{\mathcal M}_\Gamma$ 是有限态射，并非同构。用代数叠语言说，有有限自同构群 $A_\Gamma$ 作用于 $\overline M_\Gamma$，其代数叠意义下的商为 $\overline{\mathcal M}_\Gamma$。群 $A_\Gamma$ 处于正合列
\begin{equation}\label{eq:fixed-graph-automorphisms}\index{稳定映射!自同构群|hyperpage}
1\longrightarrow\prod_e\mathbb Z/d_e\mathbb Z\longrightarrow A_\Gamma
\longrightarrow\operatorname{Aut}(\Gamma)\longrightarrow1,
\end{equation}
其中 $\operatorname{Aut}(\Gamma)$ 为保持标记的图自同构群。此列可解释如下：子群 $\prod_e\mathbb Z/d_e\mathbb Z$ 来自 $d_e$ 重覆盖的覆盖
\SourcePage{307}
变换。商映射如下产生：每个 $g\in A_\Gamma$ 将一组 $n(v)$ 点亏格零稳定曲线 $\{C_v\}\in\overline M_\Gamma$ 映到另一组 $\{C'_v\}$，使 $f=\psi_\Gamma(\{C_v\})$ 与 $f'=\psi_\Gamma(\{C'_v\})$ 是同构的稳定映射。该同构将 $f^{-1}(q_i)$ 与 $(f')^{-1}(q_i)$ 的连通分支等同。特别地，每个 $C_v$ 与某个 $C'_w$ 等同，且 $i_v=i_w$。容易看出，顶点集的置换 $v\mapsto w$ 与 $C_v$ 的选择无关，并诱导保持标记的 $\Gamma$ 自同构 $\alpha(g)$。映射 $g\mapsto\alpha(g)$ 定义式~\eqref{eq:fixed-graph-automorphisms} 中的 $A_\Gamma\to\operatorname{Aut}(\Gamma)$，该列的正合性随即可见。

因此，将推论~\ref{cor:orbifold-localization-integration} 应用于 $\overline M_{0,n}(\mathbb P^r,d)$，此时与固定点分支 $\overline M_\Gamma$ 对应的项，相应项的分母因子 $a_\Gamma$ 就是 $A_\Gamma$ 的阶。
\SourceExampleNumber{9.2.1}{1}
\begin{example}\label{ex:fixed-graphs-p1-degree2}
计算标准 $T$ 作用在 $\overline M_{0,0}(\mathbb P^1,2)$ 上的固定点分支。可能的图 $\Gamma$ 为：
\begin{figure}[H]\centering
\includegraphics[width=.785\textwidth]{fixed-graphs-p1-degree2.png}
\end{figure}
图中 $d_e$ 放在边 $e$ 上方，$i_v$ 放在顶点 $v$ 下方。没有标记点，故全部 $S_v$ 为空。

得到固定点集的三个分支，相应 $T$ 固定稳定映射分别如下：
\begin{enumerate}[label=(\roman*)]
\item 曲线 $C=C_1\cup C_2$ 有两个分支。$f$ 在每个分支上的限制 $f|_{C_i}:C_i\xrightarrow{\sim}\mathbb P^1$ 为同构，且 $f(C_1\cap C_2)=q_1$。
\item 曲线 $C=C_1\cup C_2$ 有两个分支。$f$ 在每个分支上的限制 $f|_{C_i}:C_i\xrightarrow{\sim}\mathbb P^1$ 为同构，且 $f(C_1\cap C_2)=q_0$。
\item 曲线 $C$ 不可约，可在 $C\simeq\mathbb P^1$ 上取坐标 $(z_0,z_1)$，使 $f(z_0,z_1)=(z_0^2,z_1^2)$。
\end{enumerate}
每种情形都有 $A_\Gamma\simeq\mathbb Z_2$，故 $a_\Gamma=2$。在 (i)、(ii) 中，非平凡自同构来自 $\Gamma$ 的显然自同构；在 (iii) 中，它来自二重覆盖。
\end{example}
\SourceExampleNumber{9.2.1}{2}
\begin{example}\label{ex:fixed-graphs-p2-marked-degree2}
计算标准 $T$ 作用在 $\overline M_{0,1}(\mathbb P^2,2)$ 上的固定点分支。$\mathbb P^2$ 的 $T$ 作用有三个固定点
\[
q_0=(1,0,0),\qquad q_1=(0,1,0),\qquad q_2=(0,0,1).
\]
相应的可能图为：
\begin{figure}[H]\centering
\includegraphics[width=.792\textwidth]{fixed-graphs-p2-marked-degree2.png}
\end{figure}
只要 $S_v$ 非空，就在顶点 $v$ 上方加入标记 $S_v$。

第二个图对应的稳定映射，其源曲线为 $C=C_1\cup C_2\cup C_3$，$C_2$ 分别与 $C_1,C_3$ 相交于一个点，没有其他交点。$f$ 将 $C_1$ 同构地映到连接 $q_i,q_j$ 的直线，将 $C_2$ 映到点 $q_j$，并将 $C_3$
\SourcePage{308}
同构地映到连接 $q_j,q_k$ 的直线。标记点 $p_1$ 是 $C_2$ 的光滑点。$f$ 稳定，因为 $C_2$ 含三个特殊点 $p_1,C_1\cap C_2,C_2\cap C_3$。从图的中间顶点 $v$ 可立即识别收缩分支的存在，因为 $n(v)=1+2=3$。
\end{example}
为在 $\overline M_{0,n}(\mathbb P^r,d)$ 中开展计算，使用推论~\ref{cor:orbifold-localization-integration}，故须描述固定点分支的法丛。将 $\overline{\mathcal M}_\Gamma$ 的法丛记为 $N_\Gamma$。其等变 Euler 类计算的参考文献包括 \SourceCite{Kontsevich2,Givental4,GPa}；其中 \SourceCite{GPa} 还推广到 $\overline M_{g,n}(\mathbb P^r,d)$，需要等变虚拟基本类的\IndexLocation{indexloc:1417:308}局部化公式。

回顾\hyperref[subsec:virtual-fundamental-class]{第~\ref*{subsec:virtual-fundamental-class} 小节}，代数叠 $\overline{\mathcal M}_{0,n}(\mathbb P^r,d)$ 在稳定映射 $(f,C,p_1,\ldots,p_n)$ 处的切空间是超 Ext 群
\[
\operatorname{Ext}^1_C(f^*\Omega^1_{\mathbb P^r}\longrightarrow\Omega_C^1(D),\mathcal O_C),
\]
其中 $D=\sum_ip_i$。用 Ext 的长正合列理解此群，得到
\begin{equation}\label{eq:stable-map-tangent-exact-localization}
\begin{aligned}
0\longrightarrow\operatorname{Hom}(\Omega_C^1(D),\mathcal O_C)
&\longrightarrow H^0(C,f^*T_{\mathbb P^r})\\
&\longrightarrow\operatorname{Ext}^1(f^*\Omega^1_{\mathbb P^r}\to\Omega_C^1(D),\mathcal O_C)\\
&\longrightarrow\operatorname{Ext}^1(\Omega_C^1(D),\mathcal O_C)\longrightarrow0.
\end{aligned}
\end{equation}
同时，$\operatorname{Ext}^2(f^*\Omega^1_{\mathbb P^r}\to\Omega_C^1(D),\mathcal O_C)=0$，因为 $\mathbb P^r$ 凸，故 $H^1(C,f^*T_{\mathbb P^r})=0$。对更一般的 $g,X$，须用计算切空间—\IndexLocation{indexloc:1328:308}障碍复形的类似正合列替换式~\eqref{eq:stable-map-tangent-exact-localization}，如\hyperref[subsec:virtual-fundamental-class]{第~\ref*{subsec:virtual-fundamental-class} 小节}所述。

在 $\overline{\mathcal M}_\Gamma\subset\overline M_{0,n}(\mathbb P^r,d)$ 的 $T$ 固定点处，可由此正合列求切空间的 $T$ 作用的权。再考察各项的 $T$ 作用，并限制到“移动部分”，就得到法丛 $N_\Gamma$ 的权。将这些权相乘，照常得到其等变 Euler 类。

给出 $\operatorname{Euler}_T(N_\Gamma)$ 的公式前，需要一些记号。按 \SourceCite{Kontsevich2}，\emph{旗} $F$ 是一对 $(v,e)$，其中 $e$ 为含顶点 $v$ 的边。令 $i(F)=i_v$，$j(F)=i_{v'}$，其中 $v'$ 为 $e$ 的另一个顶点。对旗 $F=(v,e)$，令
\[
\omega_F=\frac{\lambda_{i(F)}-\lambda_{j(F)}}{d_e}.
\]
它是 $H^2(BT)$ 中对应于如下权的元素：$T$ 在边 $e$ 对应的分支 $C_e$ 上、位于 $q_{i_v}$ 上方的点 $p_F$ 处切空间上的作用的权。当 $n(v)\ge3$ 时，令 $e_F$ 为 $\overline M_\Gamma$ 上某线丛的第一 Chern 类，其纤维为顶点 $v$ 对应曲线分支在 $p_F$ 处的余切空间。

若 $\operatorname{val}(v)=1$，将含 $v$ 的唯一旗记为 $F(v)$。若 $\operatorname{val}(v)=2$，将两面旗记为 $F_1(v),F_2(v)$。\index{旗!稳定映射图|hyperpage}
\begin{theorem}\label{thm:fixed-graph-normal-euler}\index{等变 Euler 类!法丛|hyperpage}
法丛 $N_\Gamma$ 的等变 Euler 类，是旗、顶点、边的贡献之积。更精确地，
\[
\operatorname{Euler}_T(N_\Gamma)=e_\Gamma^{\mathrm F}e_\Gamma^{\mathrm v}e_\Gamma^{\mathrm e},
\]
\SourcePage{309}
其中三者定义为
\[
\begin{aligned}
e_\Gamma^{\mathrm F}&=\frac{\prod_{F=(v,e),\,n(v)\ge3}(\omega_F-e_F)}
{\prod_{F=(v,e)}\prod_{j\ne i(F)}(\lambda_{i(F)}-\lambda_j)},\\
e_\Gamma^{\mathrm v}&=\prod_v\prod_{j\ne i_v}(\lambda_{i_v}-\lambda_j)
\frac{\prod_{\operatorname{val}(v)=2,\,S_v=\varnothing}(\omega_{F_1(v)}+\omega_{F_2(v)})}
{\prod_{\operatorname{val}(v)=1,\,S_v=\varnothing}\omega_{F(v)}},\\
e_\Gamma^{\mathrm e}&=\prod_{e=(v,v')}\left[
\frac{(-1)^{d_e}(d_e!)^2(\lambda_i-\lambda_j)^{2d_e}}{d_e^{2d_e}}
\prod_{\substack{a+b=d_e,\ a,b\ge0\\k\ne i,j}}
\left(\frac{a\lambda_i+b\lambda_j}{d_e}-\lambda_k\right)\right],
\end{aligned}
\]
边项中 $i=i_v$、$j=i_{v'}$。
\end{theorem}
\begin{proof}
完整证明见 \SourceCite{Kontsevich2,GPa}。这里只简述涉及下面引理~\ref{lem:node-smoothing-tangent-tensor} 的一小部分证明，此引理以后还会有用。

先注意，$\operatorname{Hom}(\Omega_C^1(D),\mathcal O_C)$ 是 $(C,p_1,\ldots,p_n)$ 的无穷小自同构空间，而 $H^0(C,f^*T_{\mathbb P^r})$ 是映射 $f$ 的一阶形变空间。因此，
\[
H^0(C,f^*T_{\mathbb P^r})/\operatorname{Hom}(\Omega_C^1(D),\mathcal O_C)
\]
是整个稳定映射模空间在 $f$ 处切空间的子空间 $S_{\Gamma,f}$，对应于固定带标记源曲线、仅改变映射的形变。结合式~\eqref{eq:stable-map-tangent-exact-localization}，得到正合列
\[
0\longrightarrow S_{\Gamma,f}\longrightarrow T_f\overline M_{0,n}(\mathbb P^r,d)
\xrightarrow{\kappa}\operatorname{Ext}^1(\Omega_C^1(D),\mathcal O_C)\longrightarrow0.
\]
此外，由 Ext 的局部到整体谱序列，还有映射
\[
\mu:\operatorname{Ext}^1(\Omega_C^1(D),\mathcal O_C)
\longrightarrow H^0(C,\mathcal E xt^1(\Omega_C^1(D),\mathcal O_C)).
\]
$C$ 是曲线，谱序列在 $E_2$ 退化，故 $\mu$ 满射。$\mathcal E xt^1(\Omega_C^1(D),\mathcal O_C)$ 是支集位于 $C$ 的结点的挠层，因为 $\Omega_C^1(D)$ 在光滑点处局部自由。立即得到
\[
\mathcal E xt^1(\Omega_C^1(D),\mathcal O_C)\simeq\mathcal E xt^1(\Omega_C^1,\mathcal O_C),
\]
因而可将 $\mu$ 看作满射
\[
\mu:\operatorname{Ext}^1(\Omega_C^1(D),\mathcal O_C)
\longrightarrow H^0(C,\mathcal E xt^1(\Omega_C^1,\mathcal O_C)).
\]
由 \SourceCite{DM}，右端是忘掉 $f$ 和全部标记点后，曲线 $C$ 的一阶光滑化空间。尽管此处 $C$ 未必稳定，\SourceCite{DM} 的论证仍然适用。

这说明，$\kappa^{-1}(\mu^{-1}(\{0\}))\subset T_f\overline M_{0,n}(\mathbb P^r,d)$ 是保持全部结点的一阶形变空间。由于来自 $T_f\overline M_\Gamma$ 的形变显然保持这些结点，因此 $T_f\overline M_\Gamma\subset\kappa^{-1}(\mu^{-1}(\{0\}))$。于是 $\mu\circ\kappa$ 诱导满射
\[
N_{\Gamma,f}\longrightarrow H^0(C,\mathcal E xt^1(\Omega_C^1,\mathcal O_C)),
\]
它与右端的自然 $T$ 作用相容。

现给出 \SourceCite{Kontsevich2} 的一个重要引理，显式计算此空间。令 $q_i$ 为 $C$ 的结点，$C_i^1,C_i^2$ 为含该点的两个分支。
\SourcePage{310}
\begin{lemma}\label{lem:node-smoothing-tangent-tensor}\index{结点!光滑化|hyperpage}
存在与 $T$ 作用相容的自然同构
\[
H^0(C,\mathcal E xt^1(\Omega_C^1,\mathcal O_C))
\simeq\bigoplus_i\bigl(T_{q_i}C_i^1\otimes T_{q_i}C_i^2\bigr).
\]
\end{lemma}
\begin{proof}
计算纯属局部，故令 $C$ 为局部曲面 $S$ 内由 $xy=0$ 定义的局部解析曲线，$(x,y)$ 为 $S$ 的局部坐标。将 $C$ 的分支记为 $C_1,C_2$，令 $q=(0,0)$。用标准正合列计算 $\mathcal E xt^1$：
\begin{equation}\label{eq:node-conormal-sequence}
0\longrightarrow I_C/I_C^2\longrightarrow\Omega_S^1|_C\longrightarrow\Omega_C^1\longrightarrow0,
\end{equation}
其中第一个映射将 $[f]$ 映到 $df|_C$，第二个映射是到 $C$ 的拉回。施加 $\mathcal E xt^*(\,\cdot\,,\mathcal O_C)$ 的长正合列，得到
\begin{equation}\label{eq:node-ext-cokernel}
\mathcal E xt^1(\Omega_C^1,\mathcal O_C)
\simeq\operatorname{coker}\left(T_S|_C\xrightarrow h(I_C/I_C^2)^*\right),
\end{equation}
其中 $(I_C/I_C^2)^*=\mathcal H om(I_C/I_C^2,\mathcal O_C)$。计算可知，$h(\partial/\partial x)$ 将 $I_C/I_C^2$ 的生成元 $xy$ 映到 $y$，而 $h(\partial/\partial y)$ 将 $xy$ 映到 $x$。

内在地，这给出同构
\begin{equation}\label{eq:node-local-tensor-isomorphism}
T_qC_1\otimes T_qC_2\simeq\mathcal E xt^1(\Omega_C^1,\mathcal O_C)_q,
\end{equation}
借助式~\eqref{eq:node-ext-cokernel} 的等同，定义为
\begin{equation}\label{eq:node-local-tensor-map}
v_1\otimes v_2\longmapsto[f\longmapsto v_1v_2(f)(q)].
\end{equation}
为理解此式，将 $v_1,v_2$ 看作 $S$ 上向量场在 $C_1,C_2$ 上的限制，且每个 $v_i$ 与 $C_i$ 相切。将所得函数 $v_1v_2f$ 在 $q$ 处取值。给定向量场 $v_1,v_2$，映射 $f\mapsto v_1v_2(f)(q)$ 给出 $(I_C/I_C^2)^*|_q$ 的元素，并通过式~\eqref{eq:node-ext-cokernel} 确定余核中的类。容易看出，这个类仅依赖点 $q$ 处的 $v_1\otimes v_2$。用局部坐标 $x,y$ 表示，式~\eqref{eq:node-local-tensor-map} 将生成元 $(\partial/\partial x)\otimes(\partial/\partial y)$ 映到把 $xy$ 送到 $1$ 的 $(I_C/I_C^2)^*$ 元素的类。因此，它确实定义了同构~\eqref{eq:node-local-tensor-isomorphism}。

该同构的自然性说明，它与两侧的自然 $T$ 作用相容。
\end{proof}
在定理~\ref{thm:fixed-graph-normal-euler} 的证明中，引理~\ref{lem:node-smoothing-tangent-tensor} 在 $\operatorname{Euler}_T(N_\Gamma)$ 的公式中用于两处：
\begin{itemize}
\item $e_\Gamma^{\mathrm F}$ 中的 $\omega_F-e_F$ 项，来自收缩分支与以正次数映到 $\mathbb P^r$ 的分支相遇的结点。
\item $e_\Gamma^{\mathrm v}$ 中的 $\omega_{F_1(v)}+\omega_{F_2(v)}$ 项，来自两个都以正次数映到 $\mathbb P^r$ 的分支相遇的结点。
\end{itemize}
为完成证明，还须理解 $e_\Gamma^{\mathrm F},e_\Gamma^{\mathrm v},e_\Gamma^{\mathrm e}$ 的其他项。例如，可证明：
\begin{itemize}
\item $e_\Gamma^{\mathrm v}$ 的分母来自 $T$ 在 $\operatorname{Hom}(\Omega_C^1,\mathcal O_C)$ 中无穷小自同构上的作用的权。
\item $e_\Gamma^{\mathrm F}$ 的分母，以及 $e_\Gamma^{\mathrm v}$ 的第一个因子，分别来自 $\mathbb P^r$ 在 $q_{i(F)}$、$q_{i_v}$ 处切空间的贡献。
\end{itemize}
这些由式~\eqref{eq:stable-map-tangent-exact-localization} 不难得到。不过，还没有讨论 $e_\Gamma^{\mathrm e}$ 中的项，所以要证明定理，仍有许多工作。如前所述，完整细节见 \SourceCite{Kontsevich2,GPa}。
\end{proof}

\SourcePage{311}
定理~\ref{thm:fixed-graph-normal-euler} 给出的 $\operatorname{Euler}_T(N_\Gamma)$ 公式中，三个因子 $e_\Gamma^{\mathrm F},e_\Gamma^{\mathrm v},e_\Gamma^{\mathrm e}$ 之间显然有许多约去的项。因此，在具体例题中，这个公式未必是计算效率最高的形式。定理的优点，是给出了容易自动化的 $\operatorname{Euler}_T(N_\Gamma)$ 计算算法。下面的第~\ref{subsec:multiple-cover-localization} 小节将详细计算一个非平凡例题，并验证这些公式的正确性。

应当指出，引理~\ref{lem:node-smoothing-tangent-tensor} 中的等同
\[
H^0(C,\mathcal E xt^1(\Omega_C^1,\mathcal O_C))
\simeq\bigoplus_i(T_{q_i}C_i^1\otimes T_{q_i}C_i^2)
\]
已隐含地出现在 \SourceCite{Mumford2,HM} 中。下面第~\ref{subsec:multiple-cover-localization} 小节还将使用这一引理。\hyperref[chap:mirror-theorem]{第~\ref*{chap:mirror-theorem} 章}所述镜像定理的证明中，它也起着核心作用。

在更一般的空间 $\overline M_{0,n}(X,\beta)$ 上计算时，引理右侧对应的 $T$ 表示仍自然出现，尽管引理按原有表述不能直接适用。这源于切空间与虚拟基本类之间的联系；二者都由 $\operatorname{Ext}^*(f^*\Omega_X^1\to\Omega_C^1(D),\mathcal O_C)$ 的长正合列产生。

\SourceExampleNumber{9.2.1}{3}
\begin{example}\label{ex:quintic-quartics-localization}\index{五次三维簇!有理四次曲线|hyperpage}
这里建立的理论最早的应用之一，是计算五次三维簇的 Gromov--Witten 不变量 $N_4=\langle I_{0,0,4}\rangle$ \SourceCite{Kontsevich2}。回忆例题~\ref{ex:quintic-gw-euler-integral} 中的公式，
\[
N_4=\int_{\overline M_{0,0}(\mathbb P^4,4)}\operatorname{Euler}(\mathcal V_4),
\]
其中 $\mathcal V_4$ 是 $\overline M_{0,0}(\mathbb P^4,4)$ 上的向量丛，在 $(f,C)$ 处的纤维为 $H^0(C,f^*\mathcal O_{\mathbb P^4}(5))$。

对于 $T=(\mathbb C^*)^5$ 在 $\mathbb P^4$ 上的作用，可以列出所有可能的图 $\overline M_\Gamma$，并计算相应的等变 Euler 类 $\operatorname{Euler}_T(N_\Gamma)$。还可以将上同调类显式表示成 Čech 上闭链，计算 $\operatorname{Euler}(\mathcal V_4)$ 的限制的权。第~\ref{subsec:multiple-cover-localization} 小节将解释这一做法。再施加命题~\ref{prop:chern-polynomial-localization}，所得计算规模很大；不过，正如 \SourceCite{Kontsevich2} 所说，“在 Sun 上计算 5 分钟”便得到
\[
N_4=\frac{15517926796875}{64}.
\]
用瞬子数表示，有 $N_4=n_1\,4^{-3}+n_2\,2^{-3}+n_4$。代入已知值 $n_1=2875$ 和 $n_2=609250$，Kontsevich 得到
\[
n_4=242467530000.
\]
这是首次严格计算出 $n_4$，证实了镜像对称的预言 \SourceCite{CdGP}。正如例题~\ref{ex:quintic-instanton-enumerative} 所说明的，这意味着一般的五次三维簇恰有 $242467530000$ 条有理四次曲线。
\end{example}

\subsection{多重覆盖与 Gromov--Witten 不变量}\label{subsec:multiple-cover-localization}
设 $V$ 是 Calabi--Yau 三维簇，$C\subset V$ 是光滑的无穷小刚性有理曲线。掌握局部化以后，我们便可以介绍 Kontsevich 的方法，计算 $C$ 的 $d$ 次覆盖对 Gromov--Witten 不变量 $\langle I_{0,0,d[C]}\rangle$ 的贡献。根据定理~\ref{thm:multiple-cover}，答案是 $1/d^3$。\index{多重覆盖|hyperpage}

固定同构 $\mathbb P^1\simeq C$。刚性假设意味着 $C\subset V$ 的法丛为 $N\simeq\mathcal O_{\mathbb P^1}(-1)\oplus\mathcal O_{\mathbb P^1}(-1)$。
\SourcePage{312}
虚拟基本类 $[\overline M_{0,0}(V,d[C])]^{\mathrm{virt}}$ 是一个零维圈；按定义，$\langle I_{0,0,d[C]}\rangle$ 是它的次数。我们要计算其中支集位于 $C$ 的 $d$ 次多重覆盖簇 $M_{d,C}(V)$ 上的那一部分的贡献。在定理~\ref{thm:multiple-cover} 前的讨论中，已证明 $M_{d,C}(V)$ 是 $\overline M_{0,0}(V,d[C])$ 的一个连通分支。

同构 $\mathbb P^1\simeq C$ 诱导同构 $M_{d,C}(V)\simeq\overline M_{0,0}(\mathbb P^1,d)$，计算就在后一个空间上进行。容易看出，$\dim\overline M_{0,0}(\mathbb P^1,d)=2d-2$。证明的第一步，是将 $[\overline M_{0,0}(V,d[C])]^{\mathrm{virt}}$ 在该空间上的限制，表示为适当向量丛的 Euler 类。

具体如下。考虑稳定映射 $f\in\overline M_{0,0}(\mathbb P^1,d)$。拉回定义 $C$ 在 $V$ 中法丛的正合列，得到
\[
0\longrightarrow f^*T_{\mathbb P^1}\longrightarrow f^*T_V
\longrightarrow f^*(\mathcal O_{\mathbb P^1}(-1)\oplus\mathcal O_{\mathbb P^1}(-1))\longrightarrow0.
\]
命题~\ref{prop:excess-obstruction-euler-class} 的假设成立，因为
\[
h^1(C,f^*T_V)=h^1(C,f^*(\mathcal O_{\mathbb P^1}(-1)\oplus\mathcal O_{\mathbb P^1}(-1)))=2d-2,
\]
最后一个等式来自 Riemann--Roch 定理。因此，该命题说明，虚拟基本类的限制就是秩为 $2d-2$ 的障碍丛 $R^1\pi_{1*}e_1^*(\mathcal O_{\mathbb P^1}(-1)\oplus\mathcal O_{\mathbb P^1}(-1))$ 的 Euler 类。这里的万有稳定映射由曲线族 $\pi_1:\overline M_{0,1}(\mathbb P^1,d)\to\overline M_{0,0}(\mathbb P^1,d)$ 及万有映射 $e_1:\overline M_{0,1}(\mathbb P^1,d)\to\mathbb P^1$ 给出。

于是，\hyperref[chap:gromov-witten]{第七章}的定理~\ref{thm:multiple-cover} 可以改述如下。
\begin{theorem}\label{thm:multiple-cover-localization}\index{多重覆盖公式|hyperpage}
设 $C\subset V$ 是刚性嵌入 Calabi--Yau 三维簇的光滑有理曲线。则 $C$ 的 $d$ 次多重覆盖对 Gromov--Witten 不变量 $\langle I_{0,0,d[C]}\rangle$ 的贡献为
\[
\int_{\overline M_{0,0}(\mathbb P^1,d)}
\operatorname{Euler}\bigl(R^1\pi_{1*}e_1^*(\mathcal O_{\mathbb P^1}(-1)\oplus\mathcal O_{\mathbb P^1}(-1))\bigr)=\frac1{d^3}.
\]
\end{theorem}
\begin{proof}
证明的思路，是对 $T=(\mathbb C^*)^2$ 在 $\mathbb P^1$ 上的通常作用，使用代数叠的局部化公式，即推论~\ref{cor:orbifold-localization-integration}，计算积分。

先考虑 $d=2$。例题~\ref{ex:fixed-graphs-p1-degree2} 已确定 $T$ 在 $\overline M_{0,0}(\mathbb P^1,2)$ 上作用的三个固定点分支。将该例中的图依次记为 $\Gamma_1,\Gamma_2,\Gamma_3$。对于每个图 $\Gamma_i$，计算 $T$ 在法丛 $N_{\Gamma_i}$ 上，以及在障碍丛 $R^1\pi_{1*}e_1^*(\mathcal O_{\mathbb P^1}(-1)\oplus\mathcal O_{\mathbb P^1}(-1))$ 限制到 $Z_{\Gamma_i}$ 上的作用的权。

从 $\Gamma_1$ 开始。它是下面这个带标号的图：
\begin{figure}[H]
\centering\includegraphics[width=.171\textwidth]{multiple-cover-graph-chain.png}
\end{figure}
相应的固定点分支 $Z_{\Gamma_1}$ 只有一点，对应稳定映射 $f_1:C_1\cup C_2\to\mathbb P^1$，其中 $f_1|_{C_i}$ 是同构，且 $f_1(p)=q_1$，$\{p\}=C_1\cap C_2$。

先计算等变 Euler 类。图中有 4 个旗、3 个顶点和 2 条边。利用定理~\ref{thm:fixed-graph-normal-euler} 的公式，得到
\begin{align*}
e_{\Gamma_1}^{\mathrm F}&=\frac1{(\lambda_0-\lambda_1)^2(\lambda_1-\lambda_0)^2},\\
e_{\Gamma_1}^{\mathrm v}&=((\lambda_0-\lambda_1)^2(\lambda_1-\lambda_0))
\frac{2\lambda_1-2\lambda_0}{(\lambda_0-\lambda_1)^2},\\
e_{\Gamma_1}^{\mathrm e}&=(\lambda_0-\lambda_1)^4.
\end{align*}
\SourcePage{313}
立即得出
\begin{equation}\label{eq:degree-two-chain-normal-euler}
\operatorname{Euler}_T(N_{\Gamma_1})=2(\lambda_1-\lambda_0)^2.
\end{equation}
为进一步理解这些公式的来源，下面用更初等的方法计算切空间的权，验证此式。

为描述 $\overline M_{0,0}(\mathbb P^1,2)$ 在 $f_1$ 处的二维切空间，考虑子叠 $\overline M(\mathbb P^1,2)_1\subset\overline M_{0,0}(\mathbb P^1,2)$，由零标记稳定映射 $f:C_1\cup C_2\to\mathbb P^1$ 构成，其中 $f|_{C_i}$ 的次数为 $1$。注意 $f_1\in\overline M(\mathbb P^1,2)_1$。存在态射 $\overline M(\mathbb P^1,2)_1\to\mathbb P^1$，将 $f$ 映到 $f(C_1\cap C_2)$。由于存在交换 $C_1,C_2$ 的自同构，该态射作为代数叠的映射，其次数为 $2$。固定点 $Z_{\Gamma_1}=\{f_1\}$ 由此映到 $q_1$。该自同构对应图 $\Gamma_1$ 的一个自同构，在施加局部化公式时，会在分母中额外贡献因子 $2$。

子叠在 $f_1$ 处的切空间，给出 $T_{f_1}\overline M_{0,0}(\mathbb P^1,2)$ 的一个切方向，所做的只是移动点 $q_1$。因此，$T$ 在这个方向上的作用的权，等于它在 $\mathbb P^1$ 于 $q_1$ 处切空间上的权。用通常方法容易算得此权为 $\lambda_1-\lambda_0$。

为求剩下的权，对定理~\ref{thm:fixed-graph-normal-euler} 的证明中用到的映射作适当的图表追踪，可知商空间
\[
T_{f_1}\overline M_{0,0}(\mathbb P^1,2)/T_{f_1}\overline M(\mathbb P^1,2)_1
\]
自然同构于
\[
H^0(C_1\cup C_2,\mathcal E xt^1(\Omega_{C_1\cup C_2}^1,\mathcal O_{C_1\cup C_2})).
\]
引理~\ref{lem:node-smoothing-tangent-tensor} 又说明，这个商同构于 $T_pC_1\otimes T_pC_2$。

所有这些同构都保持自然的 $T$ 作用。而每个 $T_pC_i$ 作为 $T$ 表示都同构于 $T_{q_1}\mathbb P^1$，所以剩下的权为 $2(\lambda_1-\lambda_0)$。于是切空间的两个权是 $\lambda_1-\lambda_0$ 和 $2(\lambda_1-\lambda_0)$，等变 Euler 类为 $2(\lambda_1-\lambda_0)^2$，与式~\eqref{eq:degree-two-chain-normal-euler} 一致。

现在考虑障碍丛 $R^1\pi_{1*}e_1^*(\mathcal O_{\mathbb P^1}(-1)\oplus\mathcal O_{\mathbb P^1}(-1))$。计算其 $T$ 权，只须求 $R^1\pi_{1*}e_1^*\mathcal O_{\mathbb P^1}(-1)$ 的权，再重复一次答案。$f_1^*\mathcal O_{\mathbb P^1}(-1)$ 在 $C_i$ 上的限制为 $\mathcal O_{C_i}(-1)$，$i=1,2$。正规化映射诱导短正合列
\[
0\longrightarrow f_1^*\mathcal O_{\mathbb P^1}(-1)
\longrightarrow\mathcal O_{C_1}(-1)\oplus\mathcal O_{C_2}(-1)
\longrightarrow f_1^*\mathcal O_{q_1}(-1)\longrightarrow0.
\]
取上同调，得到一维空间
\[
H^1(C_1\cup C_2,f_1^*\mathcal O_{\mathbb P^1}(-1))
\simeq H^0(\mathbb P^1,\mathcal O_{q_1}(-1)).
\]
$T$ 在 $\mathcal O_{\mathbb P^1}(-1)$ 于 $q_1$ 处的作用为 $-\rho_1$，所以权为 $-\lambda_1$。因此，障碍丛在 $f_1$ 处的两个权为 $-\lambda_1,-\lambda_1$。

与 $\Gamma_2$ 对应的固定点 $Z_{\Gamma_2}$ 处的计算完全类似。只需交换 $\lambda_0,\lambda_1$，得到
\begin{align*}
\text{切空间的权：}&\quad\lambda_0-\lambda_1,\;2(\lambda_0-\lambda_1),\\
\text{障碍丛的权：}&\quad-\lambda_0,\;-\lambda_0.
\end{align*}
\SourcePage{314}
最后考虑 $\Gamma_3$。其中有 2 个旗、2 个顶点和 1 条边。由定理~\ref{thm:fixed-graph-normal-euler}，
\begin{align*}
e_{\Gamma_3}^{\mathrm F}&=\frac1{(\lambda_0-\lambda_1)(\lambda_1-\lambda_0)},\\
e_{\Gamma_3}^{\mathrm v}&=\frac{(\lambda_1-\lambda_0)(\lambda_0-\lambda_1)}
{((\lambda_0-\lambda_1)/2)((\lambda_1-\lambda_0)/2)},\\
e_{\Gamma_3}^{\mathrm e}&=\frac{2^2(\lambda_0-\lambda_1)^4}{2^4}.
\end{align*}
所以 $\operatorname{Euler}_T(N_{\Gamma_3})=(\lambda_0-\lambda_1)(\lambda_1-\lambda_0)$。也可像 $\Gamma_1$ 那样，通过显式计算得到这一结果。

还需要固定点 $\Gamma_3$ 处障碍丛的权。这里的固定点由二次映射 $f_3:\mathbb P^1\to\mathbb P^1$，$f_3(z_0,z_1)=(z_0^2,z_1^2)$ 表示。沿用 $\Gamma_1$ 的方法，只须计算 $T$ 在 $H^1(\mathbb P^1,f_3^*\mathcal O_{\mathbb P^1}(-1))$ 上作用的权，再将结果重复一次。

考虑由 Čech 上闭链
\[
\frac1{z_0z_1}\in Z^1(U_0\cap U_1,f_3^*\mathcal O_{\mathbb P^1}(-1))
\]
表示的非零上同调类，其中 $U_i=\{z\in\mathbb P^1:z_i\ne0\}$。变量 $x_i\in H^0(\mathbb P^1,\mathcal O_{\mathbb P^1}(1))$ 的权为 $\lambda_i$；由 $f_3$ 有 $x_i=z_i^2$，所以 $z_i$ 的权为 $\lambda_i/2$。于是 $1/(z_0z_1)$ 的权为 $-(\lambda_0+\lambda_1)/2$。因此，障碍丛的两个权均为 $-(\lambda_0+\lambda_1)/2$。

施加局部化公式所需的所有量现已齐备。积分
\[
\int_{\overline M_{0,0}(\mathbb P^1,2)}\operatorname{Euler}
\bigl(R^1\pi_{1*}e_1^*(\mathcal O_{\mathbb P^1}(-1)\oplus\mathcal O_{\mathbb P^1}(-1))\bigr)
\]
等于
\begin{align*}
&\frac{\lambda_1\lambda_1}{2(\lambda_1-\lambda_0)(2(\lambda_1-\lambda_0))}
+\frac{\lambda_0\lambda_0}{2(\lambda_0-\lambda_1)(2(\lambda_0-\lambda_1))}\\
&\qquad+\frac{(- (\lambda_0+\lambda_1)/2)(-(\lambda_0+\lambda_1)/2)}
{2(\lambda_1-\lambda_0)(\lambda_0-\lambda_1)}=\frac18.
\end{align*}
注意，每项分母中额外的因子 $2$，来自每个固定点处的阶为 $2$ 的自同构群，见例题~\ref{ex:fixed-graphs-p1-degree2}。这就完成了 $d=2$ 情形的证明。

对于一般的 $d$，\SourceCite{Manin2} 给出了一种证明，将上述方法与巧妙的组合求和结合起来。这里采用 R. Pandharipande 的简化方法证明任意 $d$ 的情形。其思路，是改变 $\mathcal O_{\mathbb P^1}(-1)\oplus\mathcal O_{\mathbb P^1}(-1)$ 上的环面作用，以简化计算。先将环面限制到子环面 $T'=\mathbb C^*\times\{1\}\subset T=(\mathbb C^*)^2$，记 $\lambda_0$ 在 $T'$ 上的限制为 $h$。在 $\overline M_{0,0}(\mathbb P^1,d)$ 上使用限制所得的作用，但在两个负线丛的直和上使用另一个 $T'$ 作用。

具体地，赋予 $\mathcal O_{\mathbb P^1}(-1)\oplus\mathcal O_{\mathbb P^1}(-1)$ 如下的作用：
\[
t\cdot(l(x_0,x_1),m(x_0,x_1))=(l(tx_0,x_1),m(x_0,t^{-1}x_1)).
\]
这里 $l,m$ 关于 $(x_0,x_1)$ 都是 $-1$ 次齐次的。注意，这个作用与 $T'$ 在 $\mathbb P^1$ 上的作用相容。
\SourcePage{315}
现在利用该作用和推论~\ref{cor:orbifold-localization-integration}，计算障碍丛的 Euler 类在 $\overline M_{0,0}(\mathbb P^1,d)$ 上的积分。容易看出，$T'$ 与 $T$ 作用的固定点集相同。设图 $\Gamma$ 对应一个固定点分支 $\overline M_\Gamma$。我们断言：若 $\Gamma$ 不止一条边，则它在局部化公式中的贡献为零。

证明并不困难。设 $f:C\to\mathbb P^1$ 是该分支中的一般稳定映射，将 $C$ 分解为分支 $C_i$，结点集记为 $\{r_j\}$。若图不止一条边，则至少有一个结点。正规化正合列为
\begin{align*}
0\longrightarrow f^*(\mathcal O_{\mathbb P^1}(-1)\oplus\mathcal O_{\mathbb P^1}(-1))
&\longrightarrow\bigoplus_i(f|_{C_i})^*(\mathcal O_{\mathbb P^1}(-1)\oplus\mathcal O_{\mathbb P^1}(-1))\\
&\longrightarrow\bigoplus_j f^*(\mathcal O_{\mathbb P^1}(-1)\oplus\mathcal O_{\mathbb P^1}(-1))_{r_j}
\longrightarrow0.
\end{align*}
取上同调，结点处纤维截面的空间，模去各收缩分支上的常值截面的像后，嵌入 $H^1(C,f^*(\mathcal O_{\mathbb P^1}(-1)\oplus\mathcal O_{\mathbb P^1}(-1)))$。

由于 $f\in\overline M_\Gamma$，每个 $f(r_j)$ 都是 $q_0$ 或 $q_1$。若 $f(r_j)=q_1$，则 $1/x_1$ 是负线丛在 $q_1$ 处的局部基，所以 $T'$ 在第一个纤维因子上的作用平凡。类似地，若 $f(r_j)=q_0$，以 $1/x_0$ 为局部基可知，$T'$ 在第二个纤维因子上的作用平凡。具有至少两条边的树必有价数至少为 $2$ 的顶点，因而上述商中至少保留一个零权。这一零权乘入障碍丛的等变 Euler 类的限制，故固定点分支 $\overline M_\Gamma$ 的贡献为零，断言得证。

因此，只需考虑对应一个 $d$ 次映射的单边图：
\begin{figure}[H]
\centering\includegraphics[width=.103\textwidth]{multiple-cover-graph-edge.png}
\end{figure}
对于该图，$\overline M_\Gamma$ 只有一点，自同构群为 $A_\Gamma=\mathbb Z/d\mathbb Z$；相应映射仍为 $f(z_0,z_1)=(z_0^d,z_1^d)$。

使用前面 $d=2$ 时的开覆盖 $U_0,U_1$，计算障碍丛的限制的权。一般的 $d$ 情形中，$H^1(\mathbb P^1,f^*(\mathcal O_{\mathbb P^1}(-1)\oplus\mathcal O_{\mathbb P^1}(-1)))$ 的一组基可由下列 Čech 上闭链给出：
\begin{equation}\label{eq:multiple-cover-cech-basis}
\left(\frac1{z_0^iz_1^{d-i}},0\right),\quad
\left(0,\frac1{z_0^iz_1^{d-i}}\right),\qquad1\le i\le d-1.
\end{equation}
$x_0,x_1$ 的权为 $h,0$；由 $f$ 有 $x_i=z_i^d$，故 $z_0,z_1$ 的权为 $h/d,0$。于是式~\eqref{eq:multiple-cover-cech-basis} 中两类上闭链的 $T'$ 权分别为 $-ih/d$ 和 $(d-i)h/d$，$1\le i\le d-1$。因此，障碍丛的等变 Euler 类在局部化公式中贡献
\[
(-1)^{d-1}((d-1)!)^2h^{2d-2}/d^{2d-2}.
\]
\SourcePage{316}
还须计算法丛 $N_\Gamma$ 的等变 Euler 类。在定理~\ref{thm:fixed-graph-normal-euler} 中取 $\lambda_0=h,\lambda_1=0$，得到
\begin{align*}
e_\Gamma^{\mathrm F}&=\frac1{(h-0)(0-h)}=-h^{-2},\\
e_\Gamma^{\mathrm v}&=\frac{(h-0)(0-h)}{((h-0)/d)((0-h)/d)}=d^2,\\
e_\Gamma^{\mathrm e}&=\frac{(-1)^d(d!)^2(h-0)^{2d}}{d^{2d}}
=\frac{(-1)^d(d!)^2h^{2d}}{d^{2d}}.
\end{align*}
由此可知
\[
\operatorname{Euler}_{T'}(N_\Gamma)=(-1)^{d-1}(d!)^2h^{2d-2}/d^{2d-2}.
\]
最后，因为处理的是代数叠，还要计入 $\overline M_\Gamma$ 的自同构因子 $a_\Gamma$。正合列~\eqref{eq:fixed-graph-automorphisms} 说明，这里 $a_\Gamma=d$，故分母中额外出现因子 $d$。综合以上结果，推论~\ref{cor:orbifold-localization-integration} 给出所求积分
\[
\frac{(-1)^{d-1}((d-1)!)^2h^{2d-2}/d^{2d-2}}
{d\cdot(-1)^{d-1}(d!)^2h^{2d-2}/d^{2d-2}}=\frac1{d^3}.
\]
这就证明了定理。
\end{proof}

\SourceCite[推论 3.5]{LLY} 给出了定理~\ref{thm:multiple-cover-localization} 的完全不同的证明。其思想与\hyperref[sec:quintic-mirror-theorem]{第~\ref*{sec:quintic-mirror-theorem} 节}所述镜像定理的证明相似，例题~\ref{ex:lly-multiple-cover-proof} 将介绍这一证明。

\subsection{六结点五次曲线的二重覆盖与 \texorpdfstring{$n_{10}$}{n10}}\label{subsec:nodal-quintic-double-covers}
\hyperref[subsec:instanton-numbers]{第~\ref*{subsec:instanton-numbers} 小节}研究五次三维簇 $V\subset\mathbb P^4$ 上的有理曲线时，已说明：当 $d\le9$ 时，瞬子数 $n_d$ 就是 $V$ 上次数为 $d$ 的有理曲线数。论证中的关键，是定理~\ref{thm:multiple-cover-localization}，它告诉我们刚性嵌入的光滑有理曲线的多重覆盖，如何对 $N_d=\langle I_{0,0,d}\rangle$ 作出贡献。不过，R. Pandharipande 向我们指出，该定理并不适用于 $V$ 上次数为 $5$ 的结点有理曲线的二重覆盖。这与 \SourceCite{Katz3,BL} 所发现的结点曲线多重覆盖的异常现象相似。下面将证明，每条这样的曲线对 $N_{10}=\langle I_{0,0,10}\rangle$ 的贡献为 $6+\frac18$，而非简单套用该定理所得的 $2^{-3}=\frac18$。\index{结点曲线!二重覆盖|hyperpage}

先解释为什么 $V$ 含有次数为 $5$ 的结点有理曲线。若 $P\subset\mathbb P^4$ 是平面，则 $C=V\cap P$ 是算术亏格为 $6$ 的平面五次曲线。我们希望在这些 $C$ 中找出有理曲线。若 $C$ 有 $6$ 个结点，它便是有理曲线。注意，$C$ 在 $p$ 处有普通结点，当且仅当 $P$ 在 $p$ 处与 $V$ 相切，且相切的二次项非退化。要求 $P$ 与 $V$ 相切，是关于 $P$ 的一个条件，故在 $6$ 个点处相切应给出 $6$ 个条件。曲线 $C$ 由射影平面 $P\subset\mathbb P^4$ 的 Grassmann 簇 $G(3,5)$ 参数化，而该簇的维数为 $6$，所以预期一般五次三维簇包含有限且非零数目的六结点有理平面五次曲线。这种朴素的维数计数可以严格化；已知其带重数的曲线数为 $17\,601\,000$ \SourceCite{Vainsencher}；这一计数的完全严格证明见 S.~Kleiman 与 R.~Piene, \emph{Node polynomials for families: methods and applications}, Math.\ Nachr.\ \textbf{271} (2004), 69--90。下面将证明，一般情形下所有重数都为 $1$，从而一般五次三维簇恰有 $17\,601\,000$ 条六结点平面五次曲线。

设 $C_n$ 是一般五次三维簇 $V$ 上这样的一条结点有理平面五次曲线。$C_n$ 的一个\emph{$d$ 次多重覆盖}，是指亏格零稳定映射 $f:C\to C_n$，满足 $f_*[C]=d[C_n]$，在 $V$ 中的次数为 $5d$。
\SourcePage{317}
也可将它看作通过 $C_n$ 分解的映射 $f:C\to V$。我们要找出这些多重覆盖对 Gromov--Witten 不变量 $N_{5d}$ 的贡献。对定理~\ref{thm:multiple-cover-localization} 的公式作朴素解释，会预期贡献为 $1/d^3$；我们将说明，这在 $d=2$ 时已不成立。论证分三步：
\begin{itemize}
\item 先计算 $C_n$ 的正规化映射的法丛。
\item 再证明 $C_n$ 的二重覆盖构成 $\overline M_{0,0}(V,10)$ 的若干连通分支的并。
\item 最后证明虚拟基本类 $[\overline M_{0,0}(V,10)]^{\mathrm{virt}}$ 限制到这些分支后的总次数为 $6+\frac18$。
\end{itemize}
先处理第一步，它将说明 $C_n$ 刚性嵌入 $V$。\hyperref[subsec:instanton-numbers]{第~\ref*{subsec:instanton-numbers} 小节}已使用过这个事实。

\begin{lemma}\label{lem:nodal-quintic-normal-bundle}\index{法丛!结点曲线的正规化|hyperpage}
若 $V$ 是一般五次三维簇，$g:\mathbb P^1\to C_n$ 是 $C_n\subset V$ 的正规化，则 $g$ 在 $V$ 中的法丛为 $N_g\simeq\mathcal O_{\mathbb P^1}(-1)\oplus\mathcal O_{\mathbb P^1}(-1)$。
\end{lemma}
\begin{proof}
$N_g$ 是 $\mathbb P^1$ 上的秩二向量丛，所以存在 $a,b$，使 $N_g\simeq\mathcal O_{\mathbb P^1}(a)\oplus\mathcal O_{\mathbb P^1}(b)$。由于 $V$ 是 Calabi--Yau 簇，伴随公式给出 $a+b=-2$。因此，$a=b=-1$ 当且仅当 $h^0(\mathbb P^1,N_g)=0$。

下面将 \SourceCite[第 152--153 页]{Katz2} 的论证适用于此情形。令 $\mathbb P=\mathbb P(H^0(\mathbb P^4,\mathcal O_{\mathbb P^4}(5)))$ 为五次超曲面的射影参数空间，$M$ 为 $\mathbb P^4$ 中亏格零不可约有理平面五次曲线的模空间。$M$ 纤维化于 $G(3,5)$ 上，纤维同构于 $\overline M_{0,0}(\mathbb P^2,5)$ 的开子集，故 $M$ 不可约且维数为 $20$；六结点曲线构成其中的开集。令
\[
I=\{(D,V)\in M\times\mathbb P:D\subset V\}
\]
为关联对应。由 Riemann--Roch 定理，$h^0(D,\mathcal O_D(5))=20$。每个 $D\in M$ 都是 $\mathbb P^4$ 中的完全交，故为射影正规。于是可算出投影 $\pi_1:I\to M$ 的纤维是 $105$ 维射影空间。所以 $I$ 不可约，维数为 $125$，与 $\mathbb P$ 的维数相同。

现在令
\[
I_1=\{(D,V)\in I:h^0(\mathbb P^1,N_g)\ge1\},
\]
其中 $g:\mathbb P^1\to D\subset V$ 是正规化映射，在上述六结点开集及光滑 $V$ 上考察此法丛。由前面对 $N_g$ 的描述，$(D,V)\in I\setminus I_1$ 当且仅当 $N_g\simeq\mathcal O_{\mathbb P^1}(-1)\oplus\mathcal O_{\mathbb P^1}(-1)$。稍后将证明这个补集非空。假设如此，由 $I$ 的不可约性，便有 $\dim I_1\le124$，所以 $\pi_2(I_1)$ 的闭包是 $\mathbb P$ 的真子簇。对补集中的 $V$，引理的断言对 $V$ 上所有这类平面有理五次曲线都成立。

因此，只需给出一个例子。设 $(z_0,z_1)$ 为 $\mathbb P^1$ 的齐次坐标，考虑下列映射的像曲线 $D$：
\begin{equation}\label{eq:nodal-quintic-parametrization}
\begin{aligned}
(z_0,z_1)\longmapsto\bigl(&z_0^5+z_1^5,\;
 z_0z_1(z_0^3+z_0^2z_1+z_0z_1^2),\\
 &z_0z_1(z_0^3+3z_0^2z_1+2z_0z_1^2+z_1^3),\;0,\;0\bigr).
\end{aligned}
\end{equation}
曲线 $D$ 位于五次超曲面 $F=0$ 中，其中
\begin{equation}\label{eq:nodal-quintic-equation}
\begin{aligned}
F(x_0,\ldots,x_4)={}&x_0^3x_1^2+3x_0^2x_1^3+20x_0x_1^4+61x_1^5
 -x_0^3x_1x_2+x_0^2x_1^2x_2\\
&-24x_0x_1^3x_2-25x_1^4x_2-3x_0^2x_1x_2^2
 +2x_0x_1^2x_2^2-40x_1^3x_2^2\\
&+x_0^2x_2^3+4x_0x_1x_2^3+35x_1^2x_2^3
 -x_0x_2^4-10x_1x_2^4+x_2^5,
\end{aligned}
\end{equation}
\SourcePage{318}
$x_0,\ldots,x_4$ 是 $\mathbb P^4$ 的坐标。包含 $D$ 的一般五次超曲面的方程形如
\begin{equation}\label{eq:quintic-containing-nodal-curve}
F+x_3G+x_4H=0,
\end{equation}
其中 $G,H$ 都是四次齐次多项式。

由法丛正合列
\begin{equation}\label{eq:nodal-normalization-tangent-sequence}
0\longrightarrow T_{\mathbb P^1}\longrightarrow g^*T_V\longrightarrow N_g\longrightarrow0,
\end{equation}
可知，$h^0(\mathbb P^1,N_g)=0$ 当且仅当 $h^0(\mathbb P^1,g^*T_V)=3$。由 Riemann--Roch 定理，这又等价于 $H^1(\mathbb P^1,g^*T_V)=0$。我们略加推广 \SourceCite{Katz1} 中的计算方法来求这个上同调群。设 $f:C\to\mathbb P^4$，$C$ 为可能可约的亏格零有理曲线，令 $\mathcal L=f^*\mathcal O_{\mathbb P^4}(1)$。若 $\widetilde F=0$ 是 $V$ 的方程，记 $\widetilde F_j=f^*(\partial\widetilde F/\partial x_j)\in H^0(C,4\mathcal L)$。考虑正合列
\[
0\longrightarrow\mathcal O_{\mathbb P^4}\longrightarrow\mathcal O_{\mathbb P^4}(1)^5\longrightarrow T_{\mathbb P^4}\longrightarrow0
\]
及
\[
0\longrightarrow T_V\longrightarrow T_{\mathbb P^4}|_V\longrightarrow\mathcal O_V(5)\longrightarrow0.
\]
由 $f$ 拉回并结合这两个正合列，得到映射
\[
\Phi_V:H^0(C,\mathcal L)^5\longrightarrow H^0(C,5\mathcal L),\qquad
(s_0,\ldots,s_4)\longmapsto\sum_{j=0}^4s_j\widetilde F_j.
\]
简单的图表追踪给出 $H^1(C,f^*T_V)\simeq\operatorname{coker}\Phi_V$。这里用了 $h^1(C,\mathcal O_C)=h^1(C,f^*T_{\mathbb P^4})=0$；后一个等式来自 $\mathbb P^4$ 的凸性。

这样就能用线性代数和计算机直接求法丛。例如，取 $G=x_0^4,H=x_1^4$，$\Phi_V$ 的矩阵有整数元素；它是 $30\times26$ 矩阵，计算所得秩为 $26$。相应的五次超曲面有奇点，但形如式~\eqref{eq:quintic-containing-nodal-curve} 的一般五次超曲面光滑。由秩的半连续性，也能找到满足条件的光滑 $V$。
\end{proof}

该引理说明，一般 $V$ 上的每条六结点有理平面五次曲线的重数均为 $1$，所以 $V$ 恰有 $17\,601\,000$ 条这样的曲线。

下面处理二重覆盖研究的第二步，证明这些覆盖给出 $\overline M_{0,0}(V,10)$ 的 $7$ 个连通分支。令 $M_{2,C_n}(V)\subset\overline M_{0,0}(V,10)$ 为分类所有零标记亏格零稳定二重覆盖 $f:C\to C_n$ 的子空间。
\begin{proposition}\label{prop:nodal-double-cover-components}\index{稳定映射!二重覆盖|hyperpage}
$M_{2,C_n}(V)$ 由模空间 $\overline M_{0,0}(V,10)$ 的 $7$ 个连通分支构成，且：
\begin{enumerate}[label=(\roman*)]
\item 一个分支同构于 $\overline M_{0,0}(\mathbb P^1,2)$，由通过正规化映射分解的二重覆盖组成。
\item 其余 $6$ 个分支是 $\overline M_{0,0}(V,10)$ 的光滑点，对应二重覆盖 $f:C_1\cup C_2\to C_n$，其中 $f|_{C_i}$ 是正规化映射，且 $f$ 在 $C_n$ 的一个结点上方为局部同构。
\end{enumerate}
\end{proposition}
\SourcePage{319}
\begin{proof}
先考虑通过正规化映射 $g:\mathbb P^1\to C_n$ 分解的二重覆盖。令 $\mathcal M_1$ 为这类稳定二重覆盖的代数叠，并考虑相应的模空间 $M_1\subset\overline M_{0,0}(V,10)$。它们恰是可写成 $f=g\circ h$ 的映射，其中 $h:C\to\mathbb P^1$ 属于 $\overline M_{0,0}(\mathbb P^1,2)$。我们断言，$M_1$ 是一个分支。为此，在 $f\in M_1$ 处对 $\overline M_{0,0}(V,10)$ 作无穷小分析。由于 $\overline M_{0,0}(\mathbb P^1,2)$ 在轨形意义下光滑、完备且二维，只需证明环境模空间在 $f$ 处的切空间也为二维。照例，这个切空间是 $\operatorname{Ext}^1(f^*\Omega_V^1\to\Omega_C^1,\mathcal O_C)$。

先设 $(f,C)\in M_1$ 且 $C$ 不可约。那么 $C\simeq\mathbb P^1$，相应的 $\operatorname{Ext}^*$ 长正合列化为
\[
0\longrightarrow H^0(\mathbb P^1,T_{\mathbb P^1})\longrightarrow H^0(\mathbb P^1,f^*T_V)
\longrightarrow\operatorname{Ext}^1(f^*\Omega_V^1\to\Omega_{\mathbb P^1}^1,\mathcal O_{\mathbb P^1})\longrightarrow0,
\]
并给出同构 $\operatorname{Ext}^2(f^*\Omega_V^1\to\Omega_{\mathbb P^1}^1,\mathcal O_{\mathbb P^1})\simeq H^1(\mathbb P^1,f^*T_V)$。由引理~\ref{lem:nodal-quintic-normal-bundle} 和式~\eqref{eq:nodal-normalization-tangent-sequence}，
\[
g^*T_V\simeq\mathcal O_{\mathbb P^1}(2)\oplus\mathcal O_{\mathbb P^1}(-1)\oplus\mathcal O_{\mathbb P^1}(-1),
\]
从而
\[
f^*T_V\simeq\mathcal O_{\mathbb P^1}(4)\oplus\mathcal O_{\mathbb P^1}(-2)\oplus\mathcal O_{\mathbb P^1}(-2).
\]
立即可知上述 $\operatorname{Ext}^1$ 的维数为 $2$。

也可能有 $(f,C)\in M_1$，其中 $C$ 有两个分支，每个都以次数 $1$ 映满 $C_n$。此时同样可算出 $\operatorname{Ext}^1(f^*\Omega_V^1\to\Omega_C^1,\mathcal O_C)$ 为二维，细节略去。命题的第 (i) 项得证。

现在设 $f:C\to C_n$ 是不能通过正规化 $g$ 分解的二次覆盖。写 $C=\bigcup_iC_i$，记 $f$ 在 $C_i$ 上的次数为 $d_i$，则 $\sum_i d_i=2$。最多只有两个分支以正次数映到 $C_n$。若有一棵分支组成的树被 $f$ 收缩，即树上所有 $d_i=0$，它最多与另两个分支相遇，与稳定性矛盾。因此，$C$ 只能有一或两个分支。若只有一个，则 $C$ 正规，$f$ 必通过 $g$ 分解。所以 $C=C_1\cup C_2$ 是两个光滑有理分支的并，且 $d_1=d_2=1$。每个 $f|_{C_i}$ 因而可与正规化映射 $g$ 等同，只需描述如何将二者粘合。

若 $f$ 将 $C$ 的结点映到 $C_n$ 的光滑点，显然 $f$ 通过 $g$ 分解。若结点映到 $C_n$ 的结点，则 $f$ 通过 $g$ 分解，当且仅当 $C$ 在结点附近的两个分支映到 $C_n$ 在像结点处的同一分支。因此，不通过 $g$ 分解时，两个分支必须分别映到像结点的两个不同分支。换言之，$f$ 在该结点上方为局部同构。

于是，恰有 $6$ 个不通过正规化分解的二重覆盖，每个结点对应一个。记结点为 $p_1,\ldots,p_6$，记 $f_i:C=C_1\cup C_2\to C_n\subset V$ 为如下覆盖：在每个 $C_j$ 上为正规化映射，将 $C$ 的结点映到 $p_i$，并使 $C_1,C_2$ 分别映到 $p_i$ 处的两个不同分支。

以上说明，$M_{2,C_n}(V)$ 由这六个映射及通过正规化分解的二重覆盖分支 $\overline M_{0,0}(\mathbb P^1,2)$ 组成。为完成证明，只需说明 $f_i$ 是环境模空间的光滑点，等价地，其切空间 $\operatorname{Ext}^1(f_i^*\Omega_V^1\to\Omega_C^1,\mathcal O_C)$ 的维数为零。
\SourcePage{320}
由通常的正合列，只须证明映射
\begin{equation}\label{eq:nodal-cover-smoothing-surjectivity}
\operatorname{Ext}^1(\Omega_C^1,\mathcal O_C)\longrightarrow H^1(C,f_i^*T_V)
\end{equation}
满射。令 $K$ 为 $f_i^*\Omega_V^1\to\Omega_C^1$ 的核。用局部坐标容易看出，$K$ 是秩为 $2$ 的局部自由层；式~\eqref{eq:nodal-cover-smoothing-surjectivity} 的满射性等价于 $H^1(C,K^*)=0$。下面计算这个上同调群。

采用间接论证。考虑关联对应
\begin{equation}\label{eq:nodal-quintic-point-incidence}
\mathcal I=\{(p,D,V)\in\mathbb P^4\times M\times\mathbb P:
 p\text{ 是 }D\subset V\text{ 的一个结点}\},
\end{equation}
其中 $M$ 与引理~\ref{lem:nodal-quintic-normal-bundle} 的证明中一样，是亏格零不可约有理平面五次曲线的模空间。

在固定点处有结点的不可约有理曲线组成的集合不可约。这是因为，例如在 $(1,0,0,0,0)$ 处有结点的曲线，可以用如下映射参数化：
\[
\begin{aligned}
(z_0,z_1)\longmapsto\bigl(&\alpha_0(z_0,z_1),\;z_0z_1\alpha_1(z_0,z_1),\;z_0z_1\alpha_2(z_0,z_1),\\
 &z_0z_1\alpha_3(z_0,z_1),\;z_0z_1\alpha_4(z_0,z_1)\bigr),
\end{aligned}
\]
这里 $\alpha_0$ 为五次齐次式，$\alpha_i$ 为三次齐次式，$1\le i\le4$。在这参数化中，任意将 $(1,0)$、$(0,1)$ 选为映到结点的两点。式~\eqref{eq:nodal-quintic-parametrization} 是其特例。所有 $\alpha_i$ 的参数空间不可约；限制到平面曲线时，再选取包含该点的平面及其中的三个坐标。平面的参数空间亦不可约，故所需结论成立。

利用这一点，可证明式~\eqref{eq:nodal-quintic-point-incidence} 定义的 $\mathcal I$ 不可约：先投影到 $\mathbb P^4\times M$，再投影到 $\mathbb P^4$，应用上一段的结论。

此后，引理~\ref{lem:nodal-quintic-normal-bundle} 的证明所用的论证，将问题归结为找出一个 $(p,D,V)\in\mathcal I$，满足 $H^1(C,K^*)=0$。取 $D$ 为式~\eqref{eq:nodal-quintic-equation} 给出的曲线，$p=(1,0,0,0,0)$。由参数化~\eqref{eq:nodal-quintic-parametrization}，该结点是 $(1,0)$ 和 $(0,1)$ 的共同像。用 Macaulay \SourceCite{BS}，可将 $C$ 显式构造为 $\mathbb P^2$ 中的两条直线，将 $f$ 构造为显式映射 $\mathbb P^2\to\mathbb P^4$ 在 $C$ 上的限制，再追踪上述构造，得到层 $K^*$。计算给出所需的 $H^1(C,K^*)=0$。虽然 Macaulay 在非零特征下计算，但由上同调的上半连续性，仍可推出一般情形下的消失。
\end{proof}

最后，计算虚拟基本类 $[\overline M_{0,0}(V,10)]^{\mathrm{virt}}$ 限制到命题~\ref{prop:nodal-double-cover-components} 的各分支后的次数。
\begin{theorem}\label{thm:nodal-quintic-double-cover-contribution}\index{五次三维簇!六结点五次曲线|hyperpage}
若 $C_n$ 是一般五次三维簇 $V\subset\mathbb P^4$ 上的六结点有理平面五次曲线，则 $C_n$ 的二重覆盖对 Gromov--Witten 不变量 $N_{10}=\langle I_{0,0,10}\rangle$ 的贡献为 $6+\frac18$。
\end{theorem}
\begin{proof}
先考虑通过正规化分解的二重覆盖所构成的分支 $\overline M_{0,0}(\mathbb P^1,2)$。直接沿用定理~\ref{thm:multiple-cover-localization} 的证明即可；切空间与障碍空间的权均不变，故该分支对 $N_{10}$ 的贡献为 $1/8$。

再看由 $f_1,\ldots,f_6$ 表示的六个分支。关键在于，它们都是具有预期维数的光滑点。由定义~\ref{def:convex-variety} 后的讨论，虚拟基本类在各点的限制为通常的点的基本类，即系数为 $1$ 的点。因此，六个分支各贡献 $1$，定理得证。
\end{proof}

\SourcePage{321}
\section{等变 Gromov--Witten 不变量}\label{sec:equivariant-gw}\index{Gromov--Witten不变量@Gromov--Witten 不变量!等变不变量|hyperpage}
现在简要讨论等变 Gromov--Witten 不变量。为简便起见，先考虑 $g=0$ 且 $X$ 为凸簇的情形。

\subsection{凸簇的亏格零不变量}\label{subsec:convex-equivariant-gw}
设 $X$ 光滑且凸。如\hyperref[subsec:computing-virtual-class]{第~\ref*{subsec:computing-virtual-class} 小节}所述，这意味着 $[\overline M_{0,n}(X,\beta)]^{\mathrm{virt}}=[\overline M_{0,n}(X,\beta)]$。于是，可以如下定义亏格零等变 Gromov--Witten 不变量。

设代数群 $G$ 在 $X$ 上作用，并保持曲线类 $\beta$。则 $h\in G$ 按公式 $h\cdot f=\tau_h\circ f$ 自然地作用于 $\overline M_{0,n}(X,\beta)$，其中 $\tau_h:X\to X$ 是与 $h$ 对应的自同构。在\hyperref[sec:stable-map-localization]{第~\ref*{sec:stable-map-localization} 节}中，已经对 $X=\mathbb P^r$、$G=(\mathbb C^*)^{r+1}$ 的情形详细分析了这个作用。

由于 $[\overline M_{0,n}(X,\beta)]^{\mathrm{virt}}=[\overline M_{0,n}(X,\beta)]$，\hyperref[subsec:gw-algebraic-classes]{第~\ref*{subsec:gw-algebraic-classes} 小节}给出的 Gromov--Witten 不变量定义简化为
\[
\langle I_{0,n,\beta}\rangle(\alpha_1,\ldots,\alpha_n)
=\int_{\overline M_{0,n}(X,\beta)}e_1^*(\alpha_1)\cup\cdots\cup e_n^*(\alpha_n),
\]
其中 $\alpha_i\in H^*(X)$，而
\[
e_i:\overline M_{0,n}(X,\beta)\longrightarrow X
\]
是\IndexLocation{indexloc:0276:321}赋值映射，将稳定映射 $f:(C,p_1,\ldots,p_n)\to X$ 映为 $f(p_i)$（见\hyperref[subsec:moduli-to-gw]{第~\ref*{subsec:moduli-to-gw} 小节}）。

要给出它的等变版本，先注意到，对于 $G$ 在 $\overline M_{0,n}(X,\beta)$ 和 $X$ 上的作用，$e_i$ 是等变的。因此有等变拉回映射
\[
e_i^*:H_G^*(X)\longrightarrow H_G^*(\overline M_{0,n}(X,\beta)).
\]
此外，$X$ 为凸簇，所以 $\overline M_{0,n}(X,\beta)$ 是轨形，这意味着存在等变积分
\[
\int_{\overline M_{0,n}(X,\beta)_G}:
H_G^*(\overline M_{0,n}(X,\beta))\longrightarrow H^*(BG).
\]
定义所需的各个要素现在都已具备。
\begin{definition}\label{def:equivariant-gw}
设 $X$ 光滑且凸，$\alpha_1,\ldots,\alpha_n\in H_G^*(X)$。则表达式
\[
\langle I_{0,n,\beta}\rangle_G(\alpha_1,\ldots,\alpha_n)
=\int_{\overline M_{0,n}(X,\beta)_G}e_1^*(\alpha_1)\cup\cdots\cup e_n^*(\alpha_n)
\in H^*(BG)
\]
称为\emph{亏格零等变 Gromov--Witten 不变量}。
\end{definition}

这个定义与通常的 Gromov--Witten 不变量有如下联系。由\hyperref[subsec:equivariant-cohomology]{第~\ref*{subsec:equivariant-cohomology} 小节}，有“\IndexLocation{indexloc:0538:321}非等变极限”映射 $i_{\mathrm{point}}^*:H^*(BG)\to H^*(\mathrm{point})=\mathbb C$，以及 $i_X^*:H_G^*(X)\to H^*(X)$。于是，恒等式
\[
i_{\mathrm{point}}^*\langle I_{0,n,\beta}\rangle_G(\alpha_1,\ldots,\alpha_n)
=\langle I_{0,n,\beta}\rangle(i_X^*\alpha_1,\ldots,i_X^*\alpha_n)
\]
可由\hyperref[subsec:chern-polynomial-localization]{第~\ref*{subsec:chern-polynomial-localization} 小节}的交换图~\eqref{eq:equivariant-integral-forgetting-diagram} 证明。因此，等变 Gromov--Witten 不变量是通常 Gromov--Witten 不变量的精细化，包含来自群作用的更多信息。

接下来，设 $E$ 是光滑凸簇 $X$ 上的\IndexLocation{indexloc:0255:321}等变向量丛。如果光滑子簇 $Y\subset X$ 由 $E$ 的一个一般
\SourcePage{322}
截面的零点集定义，那么即使 $Y$ 本身没有群作用，也能为 $Y$ 的某些 Gromov--Witten 不变量定义等变版本。这里不讨论一般情形，只说明一个对我们尤其重要的例子。
\SourceExampleNumber{9.3.1}{1}
\begin{example}\label{ex:equivariant-quintic-gw}
设 $V$ 为五次三维簇。由例题~\ref{ex:quintic-gw-euler-integral}，有 Gromov--Witten 不变量
\[
N_d=\langle I_{0,0,d}\rangle
=\int_{\overline M_{0,0}(\mathbb P^4,d)}\operatorname{Euler}(\mathcal V_d),
\]
其中 $\mathcal V_d$ 是 $\overline M_{0,0}(\mathbb P^4,d)$ 上的向量丛，其在 $(f,C)\in\overline M_{0,0}(\mathbb P^4,d)$ 处的纤维为 $H^0(C,f^*\mathcal O_{\mathbb P^4}(5))$。由于 $T=(\mathbb C^*)^5$ 自然地作用于 $\overline M_{0,0}(\mathbb P^4,d)$ 和 $\mathcal V_d$，这个不变量的等变版本为
\[
\langle I_{0,0,d}\rangle_T
=\int_{\overline M_{0,0}(\mathbb P^4,d)_T}\operatorname{Euler}_T(\mathcal V_d)
\in H^*(BT).
\]
对于例题~\ref{ex:quintic-quartics-localization} 中 $N_4$ 的计算，可以这样理解：先应用推论~\ref{cor:orbifold-localization-integration} 计算等变 Gromov--Witten 不变量 $\langle I_{0,0,4}\rangle_T$，再取非等变极限，得到所需的 $N_4$。
\end{example}

以下不会显式使用等变 Gromov--Witten 不变量，尽管\hyperref[chap:mirror-theorem]{第十一章}的一些公式隐含着它们。不过，\hyperref[chap:mirror-theorem]{第十一章}确实会采用上述策略：先处理某些等变类，再在非等变极限中得到有意义的结果。

\subsection{虚拟基本类的局部化}\label{subsec:virtual-localization}\index{虚拟基本类!局部化|hyperpage}
要为任意光滑射影簇 $X$ 和任意亏格 $g\ge0$ 定义等变 Gromov--Witten 不变量，需要虚拟基本类的等变版本。当 $X$ 具有 $\mathbb C^*$ 作用时，\SourceCite{GPa} 已利用 $X$ 的等变 Chow 群 $A_*^{\mathbb C^*}(X)$ 建立了这一理论。

这个理论有两个关键点。第一个，正如所预期的，是定义等变虚拟基本类
\[
[\overline M_{g,n}(X,\beta)]_{\mathbb C^*}^{\mathrm{virt}}
\in A_*^{\mathbb C^*}(\overline M_{g,n}(X,\beta)).
\]
由此容易得到等变 Gromov--Witten 类。第二个关键要素是局部化公式。这也不意外，因为已经看到，局部化在计算 Gromov--Witten 不变量的贡献中发挥着重要作用。

基本思想是，\SourceCite{BF,Behrend} 的构造将等变虚拟基本类
\[
[Y]_{\mathbb C^*}^{\mathrm{virt}}\in A_*^{\mathbb C^*}(Y)
\]
赋予任何带有 $\mathbb C^*$ 作用和 $\mathbb C^*$ 等变完美障碍理论的代数概形 $Y$；相应的构造也适用于 Deligne--Mumford 叠。将它应用于 $Y=\overline M_{g,n}(X,\beta)$ 时，由 $X$ 上的作用诱导 $\mathbb C^*$ 作用，而定义通常虚拟基本类所用的完美障碍理论也自然继承 $\mathbb C^*$ 作用，从而给出
\[
[\overline M_{g,n}(X,\beta)]_{\mathbb C^*}^{\mathrm{virt}}
\in A_*^{\mathbb C^*}(\overline M_{g,n}(X,\beta)).
\]

此外，\SourceCite{GPa} 说明，等变虚拟基本类有很好的局部化性质。在上述 $Y$ 带有 $\mathbb C^*$ 作用及 $\mathbb C^*$ 等变完美障碍理论的情形，固定点轨迹的分支 $Z_j$ 继承作用权为零的完美障碍理论，因而得到 $A_*^{\mathbb C^*}(Z_j)$
\SourcePage{323}
中的等变虚拟基本类。$Z_j$ 的\emph{虚法丛} $N_j^{\mathrm{virt}}$ 来自障碍理论所确定的虚切空间的“移动部分”。于是，等变虚拟基本类的局部化公式 \SourceCite{GPa} 为
\[
[Y]_{\mathbb C^*}^{\mathrm{virt}}
=\sum_j i_{j!}\left(\frac{[Z_j]_{\mathbb C^*}^{\mathrm{virt}}}
{\operatorname{Euler}_{\mathbb C^*}(N_j^{\mathrm{virt}})}\right),
\]
其中，如往常一样，$\operatorname{Euler}_{\mathbb C^*}(N_j^{\mathrm{virt}})$ 是 $N_j^{\mathrm{virt}}$ 的等变 Euler 类。

将这一公式推广到 $G=(\mathbb C^*)^n$ 是直接的。这个局部化公式在定理~\ref{thm:givental-toric-mirror} 中起着关键作用。\SourceCite{GPa} 所述的局部化方法对 Gromov--Witten 不变量还有一些很好的应用，包括例题~\ref{ex:quintic-genus-one} 中提到的亏格一 Gromov--Witten 不变量的多重覆盖贡献。

\chapter{量子微分方程}\label{chap:quantum-differential-equations}\index{量子微分方程|hyperpage}
\SourcePage{324}
本章将介绍 Gromov--Witten 不变量的一种推广，称为\emph{引力关联数}，并用它们定义\emph{量子微分方程}。在\hyperref[sec:gravitational-correlators]{第~\ref*{sec:gravitational-correlators} 节}中，定义引力关联数并研究它们的性质。\hyperref[sec:giventals-connection]{第~\ref*{sec:giventals-connection} 节}随后利用引力关联数，显式描述上同调丛相对于\hyperref[chap:quantum-cohomology]{第八章}所构造的 A 模型联络的平坦截面。借助这些平坦截面，还将在\hyperref[sec:quantum-cohomology-relations]{第~\ref*{sec:quantum-cohomology-relations} 节}中定义一个取值于 $H^*(X)$、定义在 $H^*(X)$ 上的形式函数 $J$。这个函数可看作由某些较为重要的引力关联数组成的生成函数。特别地，Calabi--Yau 三维簇的 A 模型关联函数可以完全从 $J$ 中提取出来。然后，将量子微分算子定义为零化 $J$ 的微分算子。将证明，这些算子给出量子上同调中的关系，这也说明了 $J$ 本身的重要性。此函数在 Givental 对镜像定理的研究方法中同样起着核心作用；这一方法将在\hyperref[sec:givental-mirror-approach]{第~\ref*{sec:givental-mirror-approach} 节}讨论。

\section{引力关联数}\label{sec:gravitational-correlators}\index{引力关联数|hyperpage}
本节讨论引力关联数及其性质。它们是光滑代数簇的不变量，并包含\hyperref[chap:gromov-witten]{第七章}所研究的 Gromov--Witten 不变量。在\hyperref[subsec:gravitational-correlator-definition]{第~\ref*{subsec:gravitational-correlator-definition} 小节}中定义引力关联数，在\hyperref[subsec:gravitational-correlator-properties]{第~\ref*{subsec:gravitational-correlator-properties} 小节}中研究它们的性质；最后给出一系列例子。

\subsection{关联数的定义}\label{subsec:gravitational-correlator-definition}
定义引力关联数之前，需要回顾稳定映射的一些事实。给定光滑代数簇 $X$ 和 $\beta\in H_2(X,\mathbb Z)$，由\hyperref[subsec:stable-maps]{第~\ref*{subsec:stable-maps} 小节}有模叠 $\overline{\mathcal M}_{g,n}(X,\beta)$，及其对应的粗模空间 $\overline M_{g,n}(X,\beta)$。还有赋值映射
\[
e_i:\overline{\mathcal M}_{g,n}(X,\beta)\longrightarrow X,
\qquad i=1,\ldots,n,
\]
由 $e_i(f,p_1,\ldots,p_n)=f(p_i)$ 给出，以及映射
\[
\pi_{n+1}:\overline{\mathcal M}_{g,n+1}(X,\beta)
\longrightarrow\overline{\mathcal M}_{g,n}(X,\beta),
\]
它通过忘却最后一个标记点并收缩不稳定分支来定义。

由于 $\overline M_{g,n}(X,\beta)$ 不是 $\overline{\mathcal M}_{g,n}(X,\beta)$ 的细模空间，它不能作为一个万有曲线的参数空间。不过，我们断言，$\pi_{n+1}$ 是万有曲线在代数叠意义下的版本，而 $e_{n+1}:\overline{\mathcal M}_{g,n+1}(X,\beta)\to X$ 是万有稳定映射。为说明这一点，仍用 $e_{n+1}$、$\pi_{n+1}$ 表示粗模空间 $\overline M_{g,n+1}(X,\beta)$
\SourcePage{325}
和 $\overline M_{g,n}(X,\beta)$ 上对应的映射。对任何 $(f,C,p_1,\ldots,p_n)\in\overline M_{g,n}(X,\beta)$，有自然同构
\begin{equation}\label{eq:universal-stable-curve-coarse-fiber}
C/\operatorname{Aut}(f)\simeq\pi_{n+1}^{-1}(f),
\end{equation}
其中 $\operatorname{Aut}(f)$ 是稳定映射 $f$ 的有限自同构群。这个同构如下定义：若 $p\in C$ 是不同于 $p_1,\ldots,p_n$ 的光滑点，就将它映为 $(f:C\to X,p_1,\ldots,p_n,p)$。若 $p=p_i$，则将它映为稳定映射 $(f:C\cup\mathbb P^1\to X,p_1,\ldots,p_i',\ldots,p_n,p_{n+1}')$，其中 $\mathbb P^1$ 在 $p_i$ 处附着于 $C$，$p_i'$ 和 $p_{n+1}'$ 是 $\mathbb P^1$ 上两个互异且不同于附着点的点。最后，$p$ 为结点的情形留给读者处理。由于可将 $f$ 与 $e_{n+1}$ 在 $\pi_{n+1}^{-1}(f)$ 上的限制认同，我们的断言得到证明。在对应的代数叠表述中，自同构已被叠结构计入，所以在代数叠意义下，$\pi_{n+1}$ 就是万有曲线。

式~\eqref{eq:universal-stable-curve-coarse-fiber} 对 $\pi_{n+1}$ 纤维的描述表明，它具有典范截面 $s_i:\overline M_{g,n}(X,\beta)\to\overline M_{g,n+1}(X,\beta)$，定义为
\[
s_i(f,C,p_1,\ldots,p_n)
=(f:C\cup\mathbb P^1\to X,p_1,\ldots,p_i',\ldots,p_n,p_{n+1}').
\]
按式~\eqref{eq:universal-stable-curve-coarse-fiber} 的描述，此映射就是将 $(f,C,p_1,\ldots,p_n)$ 送到 $p_i\in C$ 的映射。

现在令 $\mathcal L_i$ 表示“第 $i$ 个标记点处的余切线”，即在稳定映射 $(f:C\to X,p_1,\ldots,p_n)$ 处以余切空间 $T_{p_i}^*C$ 为纤维的线丛。严格地说，令 $\omega_{n+1}$ 为 $\pi_{n+1}$ 的相对对偶化层，则 $\mathcal L_i=s_i^*\omega_{n+1}$，这里使用代数叠上的万有曲线及其截面。注意，$\mathcal L_i$ 是叠 $\overline{\mathcal M}_{g,n}(X,\beta)$ 上的线丛，却不一定是粗模空间 $\overline M_{g,n}(X,\beta)$ 上的局部自由层。除了线丛 $\mathcal L_i$ 和赋值映射 $e_i$，还需要\IndexLocation{indexloc:1406:325}虚拟基本类 $[\overline M_{g,n}(X,\beta)]^{\mathrm{virt}}$。由\hyperref[subsec:virtual-fundamental-class]{第~\ref*{subsec:virtual-fundamental-class} 小节}，其同调次数为 $2(1-g)(d-3)-2\int_\beta\omega_X+2n$，其中 $d$ 是 $X$ 的复维数，$\omega_X$ 为典范类。现在可以作出定义。
\begin{definition}\label{def:gravitational-correlator}
给定类 $\gamma_1,\ldots,\gamma_n\in H^*(X)$，以及非负整数 $d_i$（$i=1,\ldots,n$），不变量
\[
\langle\tau_{d_1}\gamma_1,\ldots,\tau_{d_n}\gamma_n\rangle_{g,\beta}
=\int_{[\overline M_{g,n}(X,\beta)]^{\mathrm{virt}}}
\prod_{i=1}^n\bigl(c_1(\mathcal L_i)^{d_i}\cup e_i^*(\gamma_i)\bigr)
\]
称为\emph{引力关联数}。
\end{definition}

注意，只有在 $\overline M_{g,n}(X,\beta)$ 有定义时，引力关联数才有定义。这会限制三元组 $(g,n,\beta)$ 的取值。

定义~\ref{def:gravitational-correlator} 中的积分将 $\mathcal L_i$ 当作模空间 $\overline M_{g,n}(X,\beta)$ 上的线丛。严格来说，这个积分应在叠 $\overline{\mathcal M}_{g,n}(X,\beta)$ 上求值。不过，如果认为被积式已推前到 $\overline M_{g,n}(X,\beta)$ 上，那么在粗模空间上求积分就完全有意义。这里以及本书余下部分，都采用这一约定。它在思想上类似于例题~\ref{ex:quintic-virtual-class} 中对丛 $\mathcal V_d$ 及其 Euler 类所采用的约定。

引力关联数这一名称来自物理学中的类似术语。如\hyperref[sec:nonlinear-sigma-models]{附录 B.2} 所述，非线性 sigma 模型将算子 $\mathcal O_\gamma$（作为 BRST 上同调类）赋予每个上同调类 $\gamma$。算子 $\mathcal O_\gamma$ 又产生其他物理算子 $\mathcal O_{i,\gamma}$，其中整数 $i\ge0$；这些算子称为\emph{引力后裔}。引力关联数与引力后裔
\SourcePage{326}
之间的关系由方程
\begin{equation}\label{eq:gravitational-physical-normalization}
\langle\tau_{d_1}\gamma_1,\ldots,\tau_{d_n}\gamma_n\rangle_{g,\beta}
=\frac{\langle\mathcal O_{d_1,\gamma_1},\ldots,\mathcal O_{d_n,\gamma_n}\rangle_{g,\beta}}
{\prod_{j=1}^n d_j!}
\end{equation}
给出。式~\eqref{eq:gravitational-physical-normalization} 的左端是定义~\ref{def:gravitational-correlator} 中的数学引力关联数。右端的表达式 $\langle\mathcal O_{d_1,\gamma_1},\ldots,\mathcal O_{d_n,\gamma_n}\rangle_{g,\beta}$ 是物理关联函数 $\langle\mathcal O_{d_1,\gamma_1},\ldots,\mathcal O_{d_n,\gamma_n}\rangle$ 的一部分：这一部分使用亏格为 $g$、满足 $f_*[\Sigma]=\beta$ 的映射 $f:\Sigma\to X$，以及引力后裔场来计算。\hyperref[sec:general-field-theories]{附录 B.1} 概述了物理关联函数的思想。式~\eqref{eq:gravitational-physical-normalization} 右端分母中的因子纯属约定。

一个重要观察是，Gromov--Witten 不变量恰好是所有 $d_i=0$ 的引力关联数。若将 $\tau_0\gamma$ 简写为 $\gamma$，则按\hyperref[chap:gromov-witten]{第七章}的记号，
\[
\langle\gamma_1,\ldots,\gamma_n\rangle_{g,\beta}
=\langle I_{g,n,\beta}\rangle(\gamma_1,\ldots,\gamma_n).
\]
本章对 Gromov--Witten 不变量使用较简便的记号 $\langle\gamma_1,\ldots,\gamma_n\rangle_{g,\beta}$。还经常将 $\tau_d1$ 简写为 $\tau_d$，其中 $1\in H^0(X)$ 是 $X$ 的基本类所对应的上同调单位。

引力关联数还有一些有用的组合，可产生形式函数。先回顾\hyperref[chap:quantum-cohomology]{第八章}的记号：取 $H^*(X,\mathbb Q)$ 的基 $T_0=1,T_1,\ldots,T_m$。又如\hyperref[subsec:gw-potential]{第~\ref*{subsec:gw-potential} 小节}中定义 Gromov--Witten 势时一样，令 $\gamma=\sum_i t_iT_i$。若 $\omega$ 是 $X$ 上的复化 Kähler 类，则定义\emph{亏格 $g$ 耦合}
\begin{equation}\label{eq:genus-g-gravitational-coupling}
\begin{aligned}
\langle\!\langle\tau_{d_1}\gamma_1,\ldots,\tau_{d_n}\gamma_n\rangle\!\rangle_g
=\sum_{k=0}^{\infty}\sum_\beta\frac1{k!}
\langle\tau_{d_1}\gamma_1,\ldots,\tau_{d_n}\gamma_n,
\underbrace{\gamma,\ldots,\gamma}_{k\text{ 次}}\rangle_{g,\beta}\,q^\beta,
\end{aligned}
\end{equation}
其中，如往常一样，$q^\beta=e^{2\pi i\int_\beta\omega}$。在不会混淆时，有时将右端求和中的引力关联数简写为
\[
\langle\tau_{d_1}\gamma_1,\ldots,\tau_{d_n}\gamma_n,(\gamma)^k\rangle_{g,\beta}.
\]

类似于\hyperref[subsec:formal-small-quantum]{第~\ref*{subsec:formal-small-quantum} 小节}和\hyperref[subsec:gw-potential]{第~\ref*{subsec:gw-potential} 小节}中的处理，可以把式~\eqref{eq:genus-g-gravitational-coupling} 的系数解释为属于某个环 $\widehat C$；视具体情形，可取
\[
\begin{gathered}
\mathbb C[[t_0,\ldots,t_m]],\quad R[[t_0,\ldots,t_m]],\\
\mathcal R[[t_0,\ldots,t_m]],\quad
\Lambda(\omega,\mathbb Q)[[t_0,\ldots,t_m]].
\end{gathered}
\]
因此，引力关联数及亏格 $g$ 耦合都是定义良好的形式对象。

用\emph{亏格 $g$ 引力 Gromov--Witten 势}重新表述这些内容，会更清楚。对 $d\ge0$ 和 $0\le j\le m$，引入超交换变量 $t_d^j$，使 $\deg t_d^j=\deg T_j$，其中 $t_0^j$ 就是通常的变量 $t_j$。随后修改 Gromov--Witten 势 $\Phi$ 的定义（定义~\ref{def:gw-potential}），得到亏格 $g$ 引力 Gromov--Witten 势 $\Phi_g^{\mathrm{grav}}$。具体地，将表达式 $\gamma=\sum_i t_iT_i$ 替换为形式表达式
\[
\gamma=\sum_{d=0}^{\infty}\sum_{j=0}^m t_d^j\tau_dT_j,
\]
然后用亏格 $g$ 映射代替亏格零映射，再如\hyperref[subsec:gw-potential]{第~\ref*{subsec:gw-potential} 小节}中那样展开。更准确地说，有如下定义。
\SourcePage{327}
\begin{definition}\label{def:gravitational-gw-potential}\index{Gromov--Witten势@Gromov--Witten 势!引力势|hyperpage}
设 $\omega$ 是光滑射影簇 $X$ 上的复化 Kähler 类。则\emph{亏格 $g$ 引力 Gromov--Witten 势}为形式和
\[
\Phi_g^{\mathrm{grav}}(\gamma)
=\sum_{n=0}^{\infty}\sum_{\beta\in H_2(X,\mathbb Z)}
\frac1{n!}\langle\gamma^n\rangle_{g,\beta}\,q^\beta,
\]
其中 $q^\beta=e^{2\pi i\int_\beta\omega}$。
\end{definition}
在上述求和中，只对使 $\overline M_{g,n}(X,\beta)$ 存在的 $(n,\beta)$ 求和。注意，
\begin{equation}\label{eq:gravitational-potential-derivatives}
\langle\!\langle\tau_{d_1}T_{i_1},\ldots,\tau_{d_n}T_{i_n}\rangle\!\rangle_g
=\left.\frac{\partial^n\Phi_g^{\mathrm{grav}}}
{\partial t_{d_1}^{i_1}\cdots\partial t_{d_n}^{i_n}}
\right|_{t_d^j=0\text{，所有 }d>0},
\end{equation}
这里，对于奇变量的偏导数，采用\hyperref[subsec:cohomology-supermanifold]{第~\ref*{subsec:cohomology-supermanifold} 小节}中的符号约定。特殊地，当 $g=0$ 且所有 $d_i=0$ 时，可写成
\[
\langle\!\langle T_{j_1},\ldots,T_{j_n}\rangle\!\rangle_0
=\frac{\partial^n\Phi}{\partial t_{j_1}\cdots\partial t_{j_n}},
\]
其中 $\Phi$ 是通常的 Gromov--Witten 势。

立即可知，大量子乘积（定义~\ref{def:big-quantum-product}）可改写为
\begin{equation}\label{eq:big-quantum-product-couplings}
T_i*T_j=\sum_k\langle\!\langle T_i,T_j,T_k\rangle\!\rangle_0\,T^k,
\end{equation}
这里利用 $g^{kl}$ 升指标，令 $T^k=\sum_l g^{kl}T_l$。如\hyperref[chap:quantum-cohomology]{第八章}所述，这个乘积通过 Dubrovin 形式体系由 Gromov--Witten 势 $\Phi$ 得到。若将 $\Phi$ 替换为 $\Phi_0^{\mathrm{grav}}$，便在 $H^*(X)$ 上得到一个新的结合乘积，定义为
\[
T_i*_gT_j=\sum_k\frac{\partial^3\Phi_0^{\mathrm{grav}}}
{\partial t_i\partial t_j\partial t_k}\,T^k.
\]
这就是\emph{引力量子乘积}。其中的系数可视具体情形属于以下某个环：
\[
\mathbb C[[t_d^j]],\qquad R[[t_d^j]],\qquad
\mathcal R[[t_d^j]],\qquad\Lambda(\omega,\mathbb Q)[[t_d^j]].
\]

\subsection{关联数的性质}\label{subsec:gravitational-correlator-properties}
刚刚定义的引力关联数具有与\hyperref[sec:gw-properties]{第~\ref*{sec:gw-properties} 节}中 Gromov--Witten 不变量的性质相似的性质。这里只陈述将要用到的那些。它们来自虚拟基本类的性质 \SourceCite{BM,Behrend}。
\begin{axiom*}{线性}
$\langle\tau_{d_1}\gamma_1,\ldots,\tau_{d_n}\gamma_n\rangle_{g,\beta}$ 对每个变元都是线性的。
\end{axiom*}

\begin{axiom*}{有效性}
对光滑射影簇 $X$，若 $\beta$ 不是有效类，则 $\langle\tau_{d_1}\gamma_1,\ldots,\tau_{d_n}\gamma_n\rangle_{g,\beta}=0$。
\end{axiom*}

\begin{axiom*}{次数}
对齐次类 $\gamma_i$，引力关联数
\[
\langle\tau_{d_1}\gamma_1,\ldots,\tau_{d_n}\gamma_n\rangle_{g,\beta}
\]
只有在
\[
\sum_{i=1}^n(2d_i+\deg\gamma_i)
=2(1-g)\dim X-2\int_\beta\omega_X+2(3g-3+n)
\]
时才可能非零。
\end{axiom*}
这容易由定义~\ref{def:gravitational-correlator} 和虚拟基本类的维数公式得到。

\SourcePage{328}
\begin{axiom*}{基本类}
如\hyperref[chap:gromov-witten]{第七章}中一样，$\pi_n:\overline M_{g,n}(X,\beta)\to\overline M_{g,n-1}(X,\beta)$ 是忘却最后一点并收缩不稳定分支的映射。若 $n+2g\ge4$，或 $\beta\ne0$ 且 $n\ge1$，则此映射有定义。令 $\mathcal L_i'$ 为 $\overline M_{g,n-1}(X,\beta)$ 上第 $i$ 个标记点处的余切线丛，则有恒等式
\begin{equation}\label{eq:cotangent-line-forgetting-boundary}
c_1(\mathcal L_i)=\pi_n^*c_1(\mathcal L_i')
+\widetilde D_{(i,n\mid1,\ldots,i-1,i+1,\ldots,n-1)},
\end{equation}
其中 $\widetilde D_{(i,n\mid1,\ldots,i-1,i+1,\ldots,n-1)}$ 表示如下稳定映射轨迹的闭包：源曲线有两个分支 $C_1,C_2$，$C_1$ 上有标记点 $\{i,n\}$，$C_2$ 上有标记点 $\{1,\ldots,i-1,i+1,\ldots,n-1\}$，而且 $f$ 在 $C_1$ 上为常值。又如式~\eqref{eq:virtual-class-forgetting-compatibility} 所述，$\overline M_{g,n}(X,\beta)$ 的虚拟基本类是 $\overline M_{g,n-1}(X,\beta)$ 的虚拟基本类的拉回。将这一事实、式~\eqref{eq:cotangent-line-forgetting-boundary} 和 \SourceCite{Witten3} 中的计算结合，得到
\[
\begin{aligned}
&\langle\tau_{d_1}\gamma_1,\ldots,\tau_{d_{n-1}}\gamma_{n-1},1\rangle_{g,\beta}\\
&\quad=\sum_{i=1}^{n-1}\langle\tau_{d_1}\gamma_1,\ldots,
\tau_{d_{i-1}}\gamma_{i-1},\tau_{d_i-1}\gamma_i,
\tau_{d_{i+1}}\gamma_{i+1},\ldots,\tau_{d_{n-1}}\gamma_{n-1}\rangle_{g,\beta},
\end{aligned}
\]
其中任何含有 $\tau_{-1}$ 的项均取为零。
\end{axiom*}

\begin{axiom*}{除子}
设 $D$ 为除子。式~\eqref{eq:cotangent-line-forgetting-boundary} 及虚拟基本类的相容性给出
\[
\begin{aligned}
&\langle\tau_{d_1}\gamma_1,\ldots,\tau_{d_{n-1}}\gamma_{n-1},D\rangle_{g,\beta}\\
&=\left(\int_\beta D\right)
\langle\tau_{d_1}\gamma_1,\ldots,\tau_{d_{n-1}}\gamma_{n-1}\rangle_{g,\beta}\\
&\quad+\sum_{j=1}^{n-1}\langle\tau_{d_1}\gamma_1,\ldots,
\tau_{d_{j-1}}\gamma_{j-1},\tau_{d_j-1}(D\cup\gamma_j),
\tau_{d_{j+1}}\gamma_{j+1},\ldots,\tau_{d_{n-1}}\gamma_{n-1}\rangle_{g,\beta}.
\end{aligned}
\]
如同基本类公理，这条公理在 $n+2g\ge4$，或 $\beta\ne0$ 且 $n\ge1$ 时成立。
\end{axiom*}

\begin{axiom*}{分裂}
为便于陈述这条公理，稍微改变观点，将 Gromov--Witten 类的定义推广，使它也包含引力项。沿用定义~\ref{def:gw-algebraic-invariants} 的记号，若 $2g+n\ge3$，$\gamma_1,\ldots,\gamma_n\in H^*(X,\mathbb Q)$，则称由
\begin{equation}\label{eq:gravitational-class-definition}
\begin{aligned}
I_{g,n,\beta}(\tau_{d_1}\gamma_1,\ldots,\tau_{d_n}\gamma_n)
=PD^{-1}p_{2*}\pi_*\left((p_1\circ\pi)^*(\gamma_1\otimes\cdots\otimes\gamma_n)
\cap\left(\left(\prod_{i=1}^n c_1(\mathcal L_i)^{d_i}\right)\cup\xi\right)\right)
\end{aligned}
\end{equation}
给出的 $H^*(\overline M_{g,n},\mathbb Q)$ 中的元素为\emph{引力类}。其中，$\xi$ 是虚拟基本类，
\[
PD:H^*(\overline M_{g,n},\mathbb Q)
\simeq H_{6g-6+2n-*}(\overline M_{g,n},\mathbb Q)
\]
是 Poincaré 对偶同构。这里仍采用先前的约定：先将括号内的表达式看作叠 $\overline{\mathcal M}_{g,n}(X,\beta)$ 上的类，再推前到粗模空间 $\overline M_{g,n}$。在这一记号下，注意
\begin{equation}\label{eq:gravitational-class-integral}
\langle\tau_{d_1}\gamma_1,\ldots,\tau_{d_n}\gamma_n\rangle_{g,\beta}
=\int_{\overline M_{g,n}}I_{g,n,\beta}(\tau_{d_1}\gamma_1,\ldots,\tau_{d_n}\gamma_n).
\end{equation}
如\hyperref[sec:gw-properties]{第~\ref*{sec:gw-properties} 节}中一样，考虑将两条带标记曲线的粘合点接在一起的映射
\[
\varphi:\overline M_{g_1,n_1+1}\times\overline M_{g_2,n_2+1}
\longrightarrow\overline M_{g,n},
\qquad g=g_1+g_2,\quad n=n_1+n_2.
\]
\SourcePage{329}
在稳定映射的源曲线上施行这种粘合，便形成新的稳定映射。分裂公理断言
\[
\begin{aligned}
&\varphi^*I_{g,n,\beta}(\tau_{d_1}\gamma_1,\ldots,\tau_{d_n}\gamma_n)\\
&=\sum_{i,\,\beta=\beta_1+\beta_2}
I_{g_1,n_1+1,\beta_1}(\tau_{d_1}\gamma_1,\ldots,\tau_{d_{n_1}}\gamma_{n_1},T_i)\otimes\\
&\hspace{5em}I_{g_2,n_2+1,\beta_2}(T^i,\tau_{d_{n_1+1}}\gamma_{n_1+1},\ldots,\tau_{d_n}\gamma_n).
\end{aligned}
\]
\end{axiom*}

\begin{axiom*}{伸缩}\index{伸缩公理|hyperpage}
最后一条公理在 Gromov--Witten 理论中没有对应物，它是恒等式
\[
\langle\tau_1,\tau_{d_1}\gamma_1,\ldots,\tau_{d_n}\gamma_n\rangle_{g,\beta}
=(2g-2+n)\langle\tau_{d_1}\gamma_1,\ldots,\tau_{d_n}\gamma_n\rangle_{g,\beta}.
\]
为解释这一点，将万有曲线 $C_{g,n}(X,\beta)\to\overline M_{g,n}(X,\beta)$ 与基本类公理中使用的映射 $\pi_{n+1}:\overline M_{g,n+1}(X,\beta)\to\overline M_{g,n}(X,\beta)$ 认同。按这一认同，新增标记点对应的余切线丛在 $\pi_{n+1}$ 的纤维上，即在稳定映射 $(f,C,p_1,\ldots,p_n)$ 的源曲线 $C$ 上，限制为允许在各 $p_i$ 处有至多一阶极点的对偶化微分层 $\omega_C(p_1+\cdots+p_n)$。于是，由虚拟基本类相对于 $\pi_{n+1}$ 的相容性，以及该层在每个纤维上的次数为 $2g-2+n$，得到伸缩公理。
\end{axiom*}

虽然引力关联数看上去比 Gromov--Witten 不变量复杂，但 \SourceCite{Dubrovin2} 已注意到，在亏格零情形中，可以用 Gromov--Witten 不变量 $\langle\gamma_1,\ldots,\gamma_n\rangle_{0,\beta}$ 计算所有亏格零引力关联数。Dubrovin 的证明使用除子公理和下面的引理~\ref{lem:genus-zero-topological-recursion}。这里简要介绍 R. Pandharipande 的另一种方法。

其思想是先增加三个辅助标记点 $p_{n+1}$、$p_{n+2}$ 和 $p_{n+3}$。接着，将引力关联数定义中的被积式拉回到模空间 $\overline M_{0,n+3}(X,\beta)$，再与下列类作杯积：
\[
\phi=e_{n+1}^*(H)\cup e_{n+2}^*(H)
\]
其中 $H$ 是 $X$ 上的丰沛类。忘却前 $n-1$ 个标记点和映射，便得到 $s:\overline M_{0,n+3}(X,\beta)\to\overline M_{0,4}$。然后将 $\phi$ 分别与式~\eqref{eq:wdvv-boundary-equivalence} 中线性等价关系
\[
D(n,n+1\mid n+2,n+3)\sim D(n,n+3\mid n+1,n+2)
\]
两端在 $s$ 下的拉回作杯积。使用除子公理、基本类公理和分裂公理，就得到用亏格零 Gromov--Witten 不变量递归计算亏格零引力关联数的方法。下面的例题~\ref{ex:projective-line-gravitational} 将概述这一方法。将引力关联数约化的另一种方法见 \SourceCite{KoM3}。

\subsection{关联数的例题}\label{subsec:gravitational-correlator-examples}
现在给出几个例子，说明引力关联数的计算。为方便起见，当 $X=\mathbb P^r$、$\beta=d[\ell]$，其中 $\ell\subset\mathbb P^r$ 是直线时，有时用 $\overline M_{g,n}(\mathbb P^r,d)$ 代替 $\overline M_{g,n}(\mathbb P^r,\beta)$，用 $\langle\tau_{d_1}\gamma_1,\ldots,\tau_{d_n}\gamma_n\rangle_{g,d}$ 代替 $\langle\tau_{d_1}\gamma_1,\ldots,\tau_{d_n}\gamma_n\rangle_{g,\beta}$。
\SourceExampleNumber{10.1.3}{1}
\begin{example}\label{ex:projective-line-gravitational}
将计算 $X=\mathbb P^1$ 的一些引力关联数。如\hyperref[subsec:moduli-to-gw]{第~\ref*{subsec:moduli-to-gw} 小节}所述，$\overline M_{0,n}(\mathbb P^1,d)$ 的维数为\IndexLocation{indexloc:1279:329}预期维数 $2d+n-2$，因此计算中不必引入虚拟基本类。

先计算次数为 $1$ 的不变量。注意，$\overline M_{0,1}(\mathbb P^1,1)\simeq\mathbb P^1$，同构将 $f:(C,p)\to\mathbb P^1$ 映为 $f(p)$。在此同构下，$\mathcal L_1$ 对应于 $\mathbb P^1$ 的余切丛，因此得到一点引力
\SourcePage{330}
关联数 $\langle\tau_1\rangle_{0,1}=-2$。令 $H$ 为 $\mathbb P^1$ 的超平面类，也就是一个点的类。则有 Gromov--Witten 不变量
\[
\langle H\rangle_{0,1}=\langle H,H,H\rangle_{0,1}
=\langle I_{0,3,1}\rangle(H,H,H)
=\langle I_{0,3,1}\rangle(pt,pt,pt)=1.
\]
第一个等号来自除子公理，而最后一个等号是显然的，因为存在唯一的 $\mathbb P^1$ 自同构，将 $0,1,\infty$ 映为任意三个互异的点 $p_1,p_2,p_3$。

可以用两种方法计算引力关联数 $\langle H,\tau_1\rangle_{0,1}$。第一种，由伸缩公理，
\[
\langle H,\tau_1\rangle_{0,1}=-\langle H\rangle_{0,1}=-1.
\]
另一种方法是利用除子公理，得到
\[
\langle H,\tau_1\rangle_{0,1}
=1\cdot\langle\tau_1\rangle_{0,1}+\langle H\rangle_{0,1}=-2+1=-1.
\]

接着，利用\hyperref[subsec:gravitational-correlator-properties]{第~\ref*{subsec:gravitational-correlator-properties} 小节}末尾所述的约化过程，计算任意 $d$ 的不变量 $\langle\tau_{2d},1\rangle_{0,d}$ 和 $\langle\tau_{2d-1}H,1\rangle_{0,d}$。为求后者，先注意基本类公理给出
\[
\langle\tau_1H,1\rangle_{0,1}=\langle H\rangle_{0,1}=1.
\]
对一般 $d$，增加三个标记点，并选 $H$ 作为丰沛类。考虑类
\[
c_1(\mathcal L_1)^{2d-1}\cup e_1^*(H)\cup e_3^*(H)\cup e_4^*(H),
\]
将它分别与线性等价关系 $D(2,3\mid4,5)\sim D(2,5\mid3,4)$ 两端在映射 $s:\overline M_{0,5}(\mathbb P^1,d)\to\overline M_{0,4}$ 下的拉回作杯积。然后在 $\overline M_{0,5}(\mathbb P^1,d)$ 上积分并使用分裂公理。可以看到，许多项为零，非零项给出
\begin{equation}\label{eq:projective-line-descendant-splitting}
\begin{aligned}
&\langle\tau_{2d-1}H,1,H,H\rangle_{0,d}\langle1,H,1\rangle_{0,0}\\
&\quad=\langle\tau_{2d-1}H,1,1,1\rangle_{0,d-1}\langle H,H,H\rangle_{0,1}.
\end{aligned}
\end{equation}
由 Gromov--Witten 理论的点映射公理，$\langle1,H,1\rangle_{0,0}=1$；又如上所述，$\langle H,H,H\rangle_{0,1}=1$。于是，除子公理和基本类公理将式~\eqref{eq:projective-line-descendant-splitting} 化为
\[
d^2\langle\tau_{2d-1}H,1\rangle_{0,d}
=\langle\tau_{2d-3}H,1\rangle_{0,d-1}.
\]
结合此前对 $\langle\tau_1H,1\rangle_{0,1}$ 的计算，此递推式给出
\begin{equation}\label{eq:projective-line-descendant-factorial}
\langle\tau_{2d-1}H,1\rangle_{0,d}=\frac1{(d!)^2}.
\end{equation}

转而考虑 $\langle\tau_{2d},1\rangle_{0,d}$，基本类公理给出
\[
\langle\tau_2,1\rangle_{0,1}=\langle\tau_1\rangle_{0,1}=-2.
\]
对一般 $d$，考虑类 $c_1(\mathcal L_1)^{2d}\cup e_3^*(H)\cup e_4^*(H)$，如上将它与线性等价关系作杯积，再积分并使用分裂公理。得到
\[
\begin{aligned}
&\langle\tau_{2d},1,H,H\rangle_{0,d}\langle1,H,1\rangle_{0,0}\\
&\quad=\langle\tau_{2d},1,1,1\rangle_{0,d-1}\langle H,H,H\rangle_{0,1},
\end{aligned}
\]
由此可得
\[
d^2\langle\tau_{2d},1\rangle_{0,d}
+2d\langle\tau_{2d-1}H,1\rangle_{0,d}
=\langle\tau_{2d-2},1\rangle_{0,d-1}.
\]
\SourcePage{331}
这里使用了除子公理和基本类公理。结合 $\langle\tau_2,1\rangle_{0,1}$ 的公式和式~\eqref{eq:projective-line-descendant-factorial}，可解此递推式，得到
\begin{equation}\label{eq:projective-line-descendant-harmonic}
\langle\tau_{2d},1\rangle_{0,d}
=\frac{-2}{(d!)^2}\left(1+\frac12+\cdots+\frac1d\right).
\end{equation}
式~\eqref{eq:projective-line-descendant-factorial} 和~\eqref{eq:projective-line-descendant-harmonic} 的这一计算归功于 Pandharipande。这些公式隐含在 \SourceCite{Givental2} 中，例题~\ref{ex:projective-line-i-j-equality} 将说明这一点。
\end{example}

\SourceExampleNumber{10.1.3}{2}
\begin{example}\label{ex:cy-threefold-gravitational}
设 $V$ 为 Calabi--Yau 三维簇。则虚拟基本类 $[\overline M_{0,n}(V,\beta)]^{\mathrm{virt}}$ 的维数为 $n$。将计算形如 $\langle\tau_d\gamma,1\rangle_{0,\beta}$ 的全部引力关联数。次数公理表明，$\langle\tau_d\gamma,1\rangle_{0,\beta}$ 只有在 $2d+\deg\gamma=4$ 时才可能非零。因此，只需计算 $\langle\tau_2,1\rangle_{0,\beta}$、对除子 $D$ 的 $\langle\tau_1D,1\rangle_{0,\beta}$，以及对曲线类 $C$ 的 $\langle C,1\rangle_{0,\beta}=0$。还可假设 $\beta\ne0$。

回顾\hyperref[subsec:instanton-numbers]{第~\ref*{subsec:instanton-numbers} 小节}，有 Gromov--Witten 不变量
\[
\langle I_{0,0,\beta}\rangle=\sum_{k\mid\beta}n_{\beta/k}k^{-3},
\]
其中 $n_\beta$ 是该小节定义的瞬子数。如\hyperref[chap:gromov-witten]{第七章}中一样，令 $N_\beta=\langle I_{0,0,\beta}\rangle$。

现在计算所需的引力关联数。先考虑 $\langle\tau_2,1\rangle_{0,\beta}$。由基本类公理和伸缩公理，
\[
\langle\tau_2,1\rangle_{0,\beta}
=\langle\tau_1\rangle_{0,\beta}=-2N_\beta.
\]
类似地，若 $D$ 为除子，则
\[
\langle\tau_1D,1\rangle_{0,\beta}
=\langle D\rangle_{0,\beta}
=\left(\int_\beta D\right)N_\beta.
\]
最后，若 $C$ 为曲线类，基本类公理给出 $\langle C,1\rangle_{0,\beta}=0$。将在\hyperref[subsec:cy-threefold-j-function]{第~\ref*{subsec:cy-threefold-j-function} 小节}中使用这些计算。
\end{example}

下面的例子中，虚拟基本类起着关键作用。
\SourceExampleNumber{10.1.3}{3}
\begin{example}\label{ex:genus-one-gravitational-euler}
设 $X$ 为任意复维数 $d$ 的光滑射影簇。将计算 $\langle\tau_1\rangle_{1,0}$。满足 $f_*[C]=0$ 的亏格一稳定映射 $f:(C,p)\to X$ 是常值映射。因此，$(C,p)$ 是带一个标记点的稳定曲线，容易得到同构
\[
\overline M_{1,1}(X,0)\simeq X\times\overline M_{1,1},
\]
将 $f:(C,p)\to X$ 与 $(f(C),(C,p))$ 认同。预期维数为 $1$，而 $\overline M_{1,1}(X,0)$ 的实际维数是 $d+1$，所以超额维数为 $d$。由于 $f^*\Theta_X\simeq\mathcal O_C^d$，有 $h^1(C,f^*\Theta_X)=d$。这说明命题~\ref{prop:excess-obstruction-euler-class} 的假设成立。按该命题的记号，需要计算 $R^1\pi_{2*}e_2^*\Theta_X$ 的第 $d$ 个 Chern 类。由上述同构，容易得到
\[
R^1\pi_{2*}e_2^*\Theta_X
\simeq p_2^*R^1\pi_{2*}\mathcal O_{\overline{\mathcal M}_{1,2}}\otimes p_1^*\Theta_X,
\]
其中 $p_i$ 是 $X\times\overline M_{1,1}$ 的两个投影，而右端的 $\pi_2$ 是通常的忘却第二点的映射 $\overline M_{1,2}\to\overline M_{1,1}$。为简便起见，令 $\mathcal L=R^1\pi_{2*}\mathcal O_{\overline{\mathcal M}_{1,2}}$，这是 $\overline{\mathcal M}_{1,1}$ 上的线丛。于是，命题~\ref{prop:excess-obstruction-euler-class}
\SourcePage{332}
和 Chern 类的标准事实表明，虚拟基本类 $\xi=[\overline M_{1,1}(X,0)]^{\mathrm{virt}}$ 为
\begin{equation}\label{eq:genus-one-constant-virtual-class}
\begin{aligned}
\xi&=c_d(p_2^*\mathcal L\otimes p_1^*\Theta_X)\cap[\overline M_{1,1}(X,0)]\\
&=\bigl(p_1^*c_d(X)+p_2^*c_1(\mathcal L)\cup p_1^*c_{d-1}(X)\bigr)
\cap[\overline M_{1,1}(X,0)],
\end{aligned}
\end{equation}
其中 $c_k(X)=c_k(\Theta_X)$ 是切丛的 Chern 类。

在此情形也能描述线丛 $\mathcal L_1$。若 $\widetilde{\mathcal L}_1$ 为 $\overline{\mathcal M}_{1,1}$ 上相应的丛，则 $\mathcal L_1=p_2^*(\widetilde{\mathcal L}_1)$，因为映射 $\overline M_{1,1}(X,0)\to\overline M_{1,1}$ 不收缩任何曲线。还已知
\[
\int_{\overline M_{1,1}}c_1(\widetilde{\mathcal L}_1)=\frac1{24}
\qquad\text{\SourceCite{Mumford2}}.
\]
现在可以计算
\[
\begin{aligned}
\langle\tau_1\rangle_{1,0}
&=\int_{[\overline M_{1,1}(X,0)]^{\mathrm{virt}}}c_1(\mathcal L_1)\\
&=\int_{X\times\overline M_{1,1}}
p_2^*c_1(\widetilde{\mathcal L}_1)\cup
\bigl(p_1^*c_d(X)+p_2^*c_1(\mathcal L)\cup p_1^*c_{d-1}(X)\bigr)\\
&=\left(\int_{\overline M_{1,1}}c_1(\widetilde{\mathcal L}_1)\right)
\left(\int_Xc_d(X)\right)
=\frac1{24}\int_Xc_d(X).
\end{aligned}
\]
第二个等号使用式~\eqref{eq:genus-one-constant-virtual-class}，第三个等号则来自次数原因所给出的 $c_1(\widetilde{\mathcal L}_1)\cup c_1(\mathcal L)=0$。

文献中通常以另一种方式陈述式~\eqref{eq:genus-one-constant-virtual-class}。由 Serre 对偶，$\mathcal L=R^1\pi_{2*}\mathcal O_{\overline{\mathcal M}_{1,2}}$ 对偶于相对对偶化层 $\omega_2$ 沿 $\pi_2:\overline{\mathcal M}_{1,2}\to\overline{\mathcal M}_{1,1}$ 的推前 $\pi_{2*}\omega_2$。习惯上将该层记为 $\mathcal H$。于是，如果把虚拟基本类看作上同调类，式~\eqref{eq:genus-one-constant-virtual-class} 成为
\[
[\overline M_{1,1}(X,0)]^{\mathrm{virt}}
=p_1^*c_d(X)-p_2^*c_1(\mathcal H)\cup p_1^*c_{d-1}(X).
\]
\SourceCite{BCOV1} 解释了这个公式，\SourceCite{Katz3} 在稳定映射的背景中重新诠释了它；另见 \SourceCite{Getzler1}。

现在重新计算 $\langle\tau_1\rangle_{1,0}$，将得到一个令人意外而又很好的结果。先注意 $\mathcal L_1=p_2^*(\widetilde{\mathcal L}_1)$，故式~\eqref{eq:gravitational-class-definition} 所定义的引力类 $I_{1,1,0}(\tau_1)$ 化为
\[
\begin{aligned}
I_{1,1,0}(\tau_1)
&=PD^{-1}p_{2*}\bigl((c_1(\mathcal L_1)\cup p_1^*(1))\cap\xi\bigr)\\
&=c_1(\widetilde{\mathcal L}_1)\cup PD^{-1}p_{2*}(p_1^*(1)\cap\xi)\\
&=c_1(\widetilde{\mathcal L}_1)\cup I_{1,1,0}(1).
\end{aligned}
\]
第二个等号来自投影公式，而 $I_{1,1,0}(1)$ 是 $H^0(\overline M_{1,1})$ 中的 Gromov--Witten 类。于是，式~\eqref{eq:gravitational-class-integral} 给出
\[
\begin{aligned}
\langle\tau_1\rangle_{1,0}
&=\int_{\overline M_{1,1}}c_1(\widetilde{\mathcal L}_1)\cup I_{1,1,0}(1)\\
&=\deg(I_{1,1,0}(1))\int_{\overline M_{1,1}}c_1(\widetilde{\mathcal L}_1)
=\frac1{24}\deg(I_{1,1,0}(1)),
\end{aligned}
\]
其中 $\deg(I_{1,1,0}(1))$ 使用同构 $\deg:H^0(\overline M_{1,1})\simeq\mathbb C$ 定义，该同构将 $1$ 映为 $1$。

为求 $\deg(I_{1,1,0}(1))$，令 $\psi:\overline M_{0,3}\to\overline M_{1,1}$ 为将前两个标记点粘合以产生亏格一结点曲线的映射。由于 $\overline M_{0,3}$ 是一个点，有
\[
\deg(I_{1,1,0}(1))=\int_{\overline M_{0,3}}\psi^*I_{1,1,0}(1).
\]
\SourcePage{333}
结合上述两个方程，以及\hyperref[chap:gromov-witten]{第七章}的约化公理和点映射公理，得到
\[
\begin{aligned}
\langle\tau_1\rangle_{1,0}
&=\frac1{24}\int_{\overline M_{0,3}}\psi^*I_{1,1,0}(1)\\
&=\frac1{24}\int_{\overline M_{0,3}}\sum_{i,j}g^{ij}I_{0,3}(1,T_i,T_j)
=\frac1{24}\sum_{i,j}g^{ij}g_{ij}.
\end{aligned}
\]
若 $T_i,T_j$ 都是奇类，则 $g^{ij}=-g^{ji}$；若至少有一个是偶类，则 $g^{ij}=g^{ji}$。因此得到 $\langle\tau_1\rangle_{1,0}=\frac1{24}\sum_i(-1)^ib_i(X)$，其中 $b_i(X)$ 是 $X$ 的第 $i$ 个 Betti 数。比较 $\langle\tau_1\rangle_{1,0}$ 的两种计算，就得到熟悉的公式
\[
\int_Xc_d(X)=\sum_{i=0}^{2d}(-1)^ib_i(X).
\]
\end{example}

\subsection{Virasoro 猜想}\label{subsec:virasoro-conjecture}\index{Virasoro猜想@Virasoro 猜想|hyperpage}
本节最后简要介绍一个关于\IndexLocation{indexloc:0459:333}\IndexLocation{indexloc:0473:333}引力关联数的引人关注的猜想。将\emph{引力势}定义为无穷和
\[
\Phi_{\mathrm{grav}}=\sum_{g=0}^{\infty}\kappa^{2g-2}\Phi_g^{\mathrm{grav}},
\]
其中 $\kappa$ 是参数。再定义\emph{配分函数}
\[
Z=e^{\Phi_{\mathrm{grav}}}.
\]
\emph{Virasoro 猜想}断言，存在一个由变量 $t_d^j$ 的形式微分算子组成的 Virasoro 代数（式~\eqref{eq:virasoro-central-extension}），其中这些算子零化 $Z$ \SourceCite{EHX}。

更准确地说，该 Virasoro 代数包含变量 $t_d^j$ 的显式微分算子 $L_n$，$n\ge-1$，满足 $[L_n,L_m]=(n-m)L_{n+m}$。其中，方程 $L_{-1}Z=0$ 来自基本类公理 \SourceCite{Witten3}；由\IndexLocation{indexloc:0461:333}次数公理、\IndexLocation{indexloc:0463:333}除子公理与\IndexLocation{indexloc:0462:333}伸缩公理的组合，可以推出方程 $L_0Z=0$ \SourceCite{Hori}。猜想认为，对所有 $n\ge-1$，均有 $L_nZ=0$。

有时可以将这个微分算子代数扩张为完整 Virasoro 代数 $L_n$、$n\in\mathbb Z$ 的中心扩张。\hyperref[sec:conformal-field-theories]{附录 B.3} 描述了这种中心扩张，尤其见式~\eqref{eq:virasoro-central-extension}。这个表示的中心荷是 $X$ 的 Euler 示性数 $\chi(X)$。猜想并不声称 $n<-1$ 时也有 $L_nZ=0$。事实上，将交换关系 $[L_n,L_{-n}]=2nL_0+(n^3-n)\chi(X)/12$ 与 $Z\ne0$ 结合，在 $\chi(X)\ne0$ 时可知，对任意 $n>1$，$L_nZ=L_{-n}Z=0$ 都不可能同时成立。

当 $X$ 是一个点时，此猜想成为 Witten 猜想 \SourceCite{Witten3}，Kontsevich 已证明它 \SourceCite{Kontsevich1}。Kontsevich 证明中的一个关键思想是，将配分函数表示为\emph{\IndexLocation{indexloc:0760:333}矩阵积分}，即 Hermitian 矩阵空间上的积分；这导出 Virasoro 代数的一种内在作用。在物理学中，有理由认为，更一般的配分函数也可类似地用\emph{\IndexLocation{indexloc:0761:333}矩阵模型}描述，从而导出 Virasoro 代数的自然作用。

此方法已用于预测射影空间的有理曲线和椭圆曲线 Gromov--Witten 不变量。对于 $\mathbb P^2$，它给出有理曲线不变量的式~\eqref{eq:plane-recursion}，以及椭圆曲线不变量的式~\eqref{eq:plane-genus-one-recursion}。
\SourcePage{334}
\SourceCite{EHX} 提出的 Virasoro 代数只有在 $H^*(X)$ 中每个类的 Hodge 型均为 $(p,p)$ 时才能成立（见 \SourceCite{Borisov2}）。在一般情形中，S. Katz 提出了一对经过修改的 Virasoro 代数；\SourceCite{EJX} 表明，它们在算例中给出了正确的预测。陈述这个猜想时，需要将次数公理替换为使用 Hodge 型的类似公理。

\SourceCite{EX} 还表明，Virasoro 猜想加上一个额外假设，可以导出亏格 $g$ 引力后裔之间的关系，称为\emph{拓扑递推关系}（TRR）。这是因为这些关系将含 $\tau_k$ 的引力后裔约化为含 $\tau_{k-1}$ 的后裔，从而得到递推。\SourceCite{EX} 中描述的 TRR 足以确定 $\mathbb P^2$ 的 $g=2$ Gromov--Witten 不变量。E. Getzler 向我们指出，\SourceCite{Getzler1,Getzler2} 的方法可用于给出同一条 TRR，因而证明 \SourceCite{EX} 所得 TRR 的正确性。下面的引理~\ref{lem:genus-zero-topological-recursion} 将给出 $g=0$ 的一条 TRR。

为研究这个猜想，\SourceCite{LiuT} 将 $L_nZ$ 写为 $\bigl(\sum_{g\ge0}\kappa^{2g-2}\Psi_{g,n}\bigr)Z$，其中 $\Psi_{g,n}$ 是变量 $t_d^j$ 和 $q^\beta$ 的形式表达式，而不依赖 $\kappa$。可以证明，$\Psi_{g,n}$ 只依赖亏格 $\le g$ 的引力关联数。\SourceCite{LiuT} 的主要结果是：若 $H^*(X)$ 只有 $(p,p)$ 型，则对所有 $n\ge-1$，有 $\Psi_{0,n}=0$。

\section{Givental 联络}\label{sec:giventals-connection}\index{Givental联络@Givental 联络|hyperpage}
在\hyperref[sec:dubrovin-formalism]{第~\ref*{sec:dubrovin-formalism} 节}中，已经说明，超流形 $H^*(X)$ 上满足命题~\ref{prop:dubrovin-frobenius} 条件的势函数，都自然导出 $H^*(X,\mathbb C)$ 的切丛上的一族形式平坦联络 $\nabla^\lambda$。当势函数为 Gromov--Witten 势时，得到的联络尤其值得研究，它们与\hyperref[sec:a-model-vhs]{第~\ref*{sec:a-model-vhs} 节}所研究的 A 模型联络密切相关。

这里实际研究的是 Dubrovin 联络的一个扭曲版本，称为\emph{Givental 联络} $\nabla^g$。在\hyperref[subsec:giventals-flat-sections]{第~\ref*{subsec:giventals-flat-sections} 小节}中，将说明如何用引力关联数写出一组 $\nabla^g$ 平坦截面的基，并讨论它与\hyperref[chap:quantum-cohomology]{第八章}的 A 模型联络的关系。\hyperref[subsec:equivariant-giventals-connection]{第~\ref*{subsec:equivariant-giventals-connection} 小节}还将解释 Givental 联络中的参数 $\hbar$ 与等变上同调的关系。

本节构造的平坦截面将用于\hyperref[sec:quantum-cohomology-relations]{第~\ref*{sec:quantum-cohomology-relations} 节}中定义 Givental 函数 $J_X$；它又将在该节导出量子微分方程及\IndexLocation{indexloc:1124:334}量子上同调中的关系。

\subsection{平坦截面}\label{subsec:giventals-flat-sections}
给定光滑射影簇 $X$，令 $\Phi$ 表示 Gromov--Witten 势，$*$ 表示大量子乘积。这里不直接使用\hyperref[sec:dubrovin-formalism]{第~\ref*{sec:dubrovin-formalism} 节}所定义的 $\Phi$ 的 Dubrovin 联络，而是遵循 \SourceCite{Givental2}，对它略作修改。如该节一样，将 $\partial/\partial t_i$ 与 $T_i$ 认同，并将联络看作定义在 $H^*(X,\mathbb C)$ 上的平凡上同调丛上。于是，定义\emph{Givental 联络} $\nabla^g$ 为
\begin{equation}\label{eq:giventals-connection}
\nabla^g_{\frac{\partial}{\partial t_i}}\left(\sum_j a_jT_j\right)
=\hbar\sum_j\frac{\partial a_j}{\partial t_i}T_j
-\sum_j a_jT_i*T_j,
\end{equation}
其中 $\hbar$ 为参数。在\hyperref[subsec:equivariant-giventals-connection]{第~\ref*{subsec:equivariant-giventals-connection} 小节}中，将把 $\hbar$ 与 $H^2(B\mathbb C^*)$ 的一个生成元认同，这将在\hyperref[subsec:quintic-mirror-theorem-statement]{第~\ref*{subsec:quintic-mirror-theorem-statement} 小节}中有用。

虽然由于 $\hbar$ 的存在，$\nabla^g$ 并不满足通常联络的 Leibniz 法则，但仍使用联络的标准术语，尤其使用平坦截面的概念。如\hyperref[subsec:a-model-connection]{第~\ref*{subsec:a-model-connection} 小节}所述，Givental 联络与 Dubrovin 联络 $\nabla^\lambda$ 的关系为
\SourcePage{335}
$\nabla^g=\hbar\nabla^{-\hbar^{-1}}$。因此，式~\eqref{eq:giventals-connection} 的平坦截面正是参数 $\lambda=-\hbar^{-1}$ 时 Dubrovin 联络的平坦截面。

对每个指标 $a=0,\ldots,m$，定义形式截面
\begin{equation}\label{eq:giventals-flat-section-series}
s_a=T_a+\sum_{n=0}^{\infty}\hbar^{-(n+1)}\sum_j
\langle\!\langle\tau_nT_a,T_j\rangle\!\rangle_0\,T^j.
\end{equation}
注意，$\langle\!\langle\tau_nT_a,T_j\rangle\!\rangle_0$ 是变量 $t_0,\ldots,t_m$ 的形式幂级数。若以形式符号 $c$ 表示 $c_1(\mathcal L_1)$，则可将式~\eqref{eq:giventals-flat-section-series} 形式地写为
\begin{equation}\label{eq:giventals-flat-section-symbolic}
s_a=T_a+\sum_j\left\langle\!\left\langle\frac{T_a}{\hbar-c},T_j\right\rangle\!\right\rangle_0T^j,
\end{equation}
因为将分母按 $c/\hbar$ 的幂形式展开，再以 $\tau_n$ 替换 $c^n$，就得到式~\eqref{eq:giventals-flat-section-series}。截面 $s_a$ 在这里具有很好的性质。
\begin{proposition}\label{prop:giventals-flat-sections}
截面 $s_a$ 构成 $\nabla^g$ 平坦截面的一组基。由式~\eqref{eq:giventals-connection}，这等价于方程
\[
\hbar\frac{\partial s_a}{\partial t_i}=T_i*s_a,
\qquad a,i=0,\ldots,m,
\]
其中 $*$ 表示大量子乘积。
\end{proposition}

证明之前，需要 \SourceCite{Witten3} 的如下结果。按\hyperref[subsec:virasoro-conjecture]{第~\ref*{subsec:virasoro-conjecture} 小节}的术语，这个引理是一条拓扑递推关系。
\begin{lemma}\label{lem:genus-zero-topological-recursion}\index{拓扑递推关系|hyperpage}
对 $d_1,d_2,d_3\ge0$、$0\le j_1,j_2,j_3\le m$，有
\[
\begin{aligned}
&\langle\!\langle\tau_{d_1+1}T_{j_1},\tau_{d_2}T_{j_2},\tau_{d_3}T_{j_3}\rangle\!\rangle_0\\
&\quad=\sum_a\langle\!\langle\tau_{d_1}T_{j_1},T_a\rangle\!\rangle_0
\langle\!\langle T^a,\tau_{d_2}T_{j_2},\tau_{d_3}T_{j_3}\rangle\!\rangle_0,
\end{aligned}
\]
其中，如往常一样，$T^a$ 是 $T_a$ 相对于杯积的对偶基。
\end{lemma}
\begin{proof}
对 $n\ge3$，考虑从 $\overline M_{0,n}(X,\beta)$ 到一点的映射，记为 $\pi:\overline M_{0,n}(X,\beta)\to\overline M_{0,3}$，以便利用稳定映射空间的函子性质。这个映射将 $(f,C,p_1,\ldots,p_n)$ 送到从 $C$ 删除 $p_4,\ldots,p_n$ 并收缩全部不稳定分支所得的曲线。令 $\mathcal L_1'$ 为 $\overline M_{0,3}$ 上按 $\mathcal L_1$ 的方式定义的（平凡）丛。计算得到
\begin{equation}\label{eq:cotangent-line-three-point-boundary}
c_1(\pi^*\mathcal L_1')
=c_1(\mathcal L_1)-\sum_{K\sqcup L=\{4,\ldots,n\}}
D_{(\{1\}\cup K\mid\{2,3\}\cup L)},
\end{equation}
其中 $D_{(I\mid J)}$ 表示源曲线有两个分支、标记点 $I$ 位于 $f$ 为常值的一个分支、标记点 $J$ 位于另一个分支的稳定映射轨迹的闭包，并对各可能的曲线类分解取和，如图所示：
\begin{figure}[H]
\centering
\includegraphics[width=0.285\textwidth]{gravitational-recursion-boundary.png}
\end{figure}
\SourcePage{336}
现在将式~\eqref{eq:cotangent-line-three-point-boundary} 乘以 $\prod_a\bigl(c_1(\mathcal L_a)^{d_a}\cup e_a^*(T_{j_a})\bigr)$，对虚拟基本类积分，并使用分裂公理，得到
\begin{equation}\label{eq:genus-zero-topological-recursion-coefficients}
\begin{aligned}
&\langle\tau_{d_1+1}T_{j_1},\tau_{d_2}T_{j_2},\tau_{d_3}T_{j_3},
\tau_{d_4}T_{j_4},\ldots,\tau_{d_n}T_{j_n}\rangle_{0,\beta}\\
&=\sum_{\substack{K\sqcup L=\{4,\ldots,n\}\\\beta=\beta_1+\beta_2}}
\sum_a\pm\langle\tau_{d_1}T_{j_1},\tau_{d_{k_1}}T_{j_{k_1}},\ldots,T_a\rangle_{0,\beta_1}\\
&\hspace{5em}\times\langle T^a,\tau_{d_2}T_{j_2},\tau_{d_3}T_{j_3},
\tau_{d_{l_1}}T_{j_{l_1}},\ldots\rangle_{0,\beta_2},
\end{aligned}
\end{equation}
其中 $K=\{k_1,\ldots\}$、$L=\{l_1,\ldots\}$。式~\eqref{eq:genus-zero-topological-recursion-coefficients} 的符号来自方程左右两端奇上同调类次序的可能改变。

按式~\eqref{eq:genus-g-gravitational-coupling} 展开后，容易看出，式~\eqref{eq:genus-zero-topological-recursion-coefficients} 正是所需等式
\[
\begin{aligned}
&\langle\!\langle\tau_{d_1+1}T_{j_1},\tau_{d_2}T_{j_2},\tau_{d_3}T_{j_3}\rangle\!\rangle_0\\
&=\sum_a\langle\!\langle\tau_{d_1}T_{j_1},T_a\rangle\!\rangle_0
\langle\!\langle T^a,\tau_{d_2}T_{j_2},\tau_{d_3}T_{j_3}\rangle\!\rangle_0
\end{aligned}
\]
中的一个系数。注意，导致式~\eqref{eq:genus-zero-topological-recursion-coefficients} 中符号的奇上同调类重新排列，恰好与比较引理右端变量次序和左端变量次序时的重新排列相同。
\end{proof}

现在可证明命题~\ref{prop:giventals-flat-sections}，即证明 $s_a$ 为平坦截面。
\begin{proof}[命题~\ref{prop:giventals-flat-sections} 的证明]
先观察到
\[
\frac{\partial}{\partial t_i}\langle\!\langle\tau_{d_1}\gamma_1,\ldots,\tau_{d_n}\gamma_n\rangle\!\rangle_0
=\langle\!\langle T_i,\tau_{d_1}\gamma_1,\ldots,\tau_{d_n}\gamma_n\rangle\!\rangle_0.
\]
它由式~\eqref{eq:genus-g-gravitational-coupling} 得到，论证类似于引理~\ref{lem:gw-potential-third-derivatives} 的证明。利用此式，从式~\eqref{eq:giventals-flat-section-series} 计算出
\[
\hbar\frac{\partial s_a}{\partial t_i}
=\sum_{n=0}^{\infty}\sum_j\hbar^{-n}
\langle\!\langle T_i,\tau_nT_a,T_j\rangle\!\rangle_0T^j,
\]
而式~\eqref{eq:giventals-flat-section-series} 和~\eqref{eq:big-quantum-product-couplings} 给出
\begin{equation}\label{eq:giventals-section-product}
\begin{aligned}
T_i*s_a
&=\sum_j\langle\!\langle T_i,T_a,T_j\rangle\!\rangle_0T^j\\
&\quad+\sum_{n=0}^{\infty}\sum_{j,c}\hbar^{-(n+1)}
\langle\!\langle\tau_nT_a,T_j\rangle\!\rangle_0
\langle\!\langle T_i,T^j,T_c\rangle\!\rangle_0T^c.
\end{aligned}
\end{equation}
引理~\ref{lem:genus-zero-topological-recursion} 表明，式~\eqref{eq:giventals-section-product} 右端的第二项为
\[
\sum_{n=0}^{\infty}\sum_c\hbar^{-(n+1)}
\langle\!\langle T_i,\tau_{n+1}T_a,T_c\rangle\!\rangle_0T^c.
\]
将它代回式~\eqref{eq:giventals-section-product}，就得到所需结论。
\end{proof}

命题~\ref{prop:giventals-flat-sections} 的一个推论是
\[
\hbar\frac{\partial s_a}{\partial t_0}=s_a,
\]
\SourcePage{337}
因为 $T_0$ 是大量子乘积的单位元。因此，$s_a$ 通过乘法因子 $e^{t_0/\hbar}$ 依赖 $t_0$。鉴于这一点，经常忽略 $t_0$ 的依赖，因为只要插入因子 $e^{t_0/\hbar}$ 就能将它恢复。

以下有时会限制丛及联络 $\nabla^g$ 的底流形。特别地，在\hyperref[sec:quantum-cohomology-relations]{第~\ref*{sec:quantum-cohomology-relations} 节}中，将限制到子空间
\[
M=H^0(X,\mathbb C)\oplus H^2(X,\mathbb C)\subset H^*(X,\mathbb C).
\]
若 $T_1,\ldots,T_r$ 构成 $H^2(X,\mathbb Q)$ 的基，则 $M$ 的变量为 $t_0,t_1,\ldots,t_r$。为看清此时 $\nabla^g$ 的形式，先注意：由\hyperref[subsec:quantum-product-dependence]{第~\ref*{subsec:quantum-product-dependence} 小节}，将大量子乘积限制到 $H^2(X,\mathbb C)$，恰好得到 $*_{\mathrm{small}}$，只需在式~\eqref{eq:genus-g-gravitational-coupling} 中令 $q^\beta=1$，并使用式~\eqref{eq:restricted-big-quantum-product} 附近说明的约定。又因为 $T_0=1$ 同时是大量子上同调与小量子上同调的单位元，故 $\nabla^g$ 在 $M$ 上的限制为
\begin{equation}\label{eq:restricted-giventals-connection}
\nabla^g_{\frac{\partial}{\partial t_i}}\left(\sum_{j=0}^m a_jT_j\right)
=\hbar\sum_{j=0}^m\frac{\partial a_j}{\partial t_i}T_j
-\sum_{j=0}^m a_jT_i*_{\mathrm{small}}T_j,
\end{equation}
其中 $i=0,1,\ldots,r$。特别地，截面 $s_a$ 限制到 $M$ 后满足
\[
\hbar\frac{\partial s_a}{\partial t_i}=T_i*_{\mathrm{small}}s_a,
\qquad a=0,\ldots,m,\quad i=0,\ldots,r.
\]
限制到 $M$ 后，$s_a$ 的公式有很好的简化。
\begin{proposition}\label{prop:restricted-giventals-flat-sections}
令 $\delta=\sum_{i=1}^r t_iT_i$。则 $s_a$ 在 $M$ 上的限制可写为
\[
\begin{aligned}
s_a=e^{t_0/\hbar}\biggl(&e^{\delta/\hbar}\cup T_a\\
&+\sum_{\beta\ne0}\sum_{j=0}^m e^{\int_\beta\delta}
\left\langle\frac{e^{\delta/\hbar}\cup T_a}{\hbar-c},T_j\right\rangle_{0,\beta}T^j\biggr).
\end{aligned}
\]
\end{proposition}
\begin{proof}
前面已注意到，$t_0$ 通过乘法因子 $e^{t_0/\hbar}$ 出现，因此只需证明 $t_0=0$ 的情形。利用 $\delta=\sum_{i=1}^rt_iT_i$，引力耦合~\eqref{eq:genus-g-gravitational-coupling} 成为
\[
\langle\!\langle\tau_nT_a,T_j\rangle\!\rangle_0
=\sum_{k=0}^{\infty}\sum_\beta\frac1{k!}
\langle\tau_nT_a,T_j,(\delta)^k\rangle_{0,\beta},
\]
其中按上述说明，令 $q^\beta=1$。这里使用式~\eqref{eq:genus-g-gravitational-coupling} 后面引入的记号。

由于 $\delta$ 为除子，可以用除子公理简化这个表达式。当 $\beta\ne0$ 时，容易对 $k$ 归纳证明恒等式
\begin{equation}\label{eq:descendant-divisor-binomial}
\langle\tau_nT_a,T_j,(\delta)^k\rangle_{0,\beta}
=\sum_{\mu+\nu=k}\frac{k!}{\mu!\nu!}
\left(\int_\beta\delta\right)^\mu
\langle\tau_{n-\nu}(T_a\cup\delta^\nu),T_j\rangle_{0,\beta},
\end{equation}
其中，当 $n-\nu<0$ 时，将 $\tau_{n-\nu}(T_a\cup\delta^\nu)$ 取为零。对于 $\beta=0$，需要另一个恒等式，因为 $\overline M_{0,2}(X,0)$ 不存在。此时，对 $k\ge1$，利用除子公理得到
\begin{equation}\label{eq:degree-zero-descendant-divisors}
\langle\tau_nT_a,T_j,(\delta)^k\rangle_{0,0}
=\langle\tau_{n-k+1}(T_a\cup\delta^{k-1}),T_j,\delta\rangle_{0,0}.
\end{equation}
若 $k>n+1$，此式显然为零；若 $k=n+1$，则由 Gromov--Witten 不变量的点映射公理，它化为积分 $\int_XT_a\cup T_j\cup\delta^k$。最后，若 $k<n+1$，有同构 $\overline M_{0,3}(X,0)\simeq X$，将 $(f,C,p_1,p_2,p_3)$ 映为 $f(C)\in X$。此外，$\overline M_{0,4}(X,0)\to\overline M_{0,3}(X,0)$ 就是投影 $X\times\overline M_{0,4}\to X$。
\SourcePage{338}
因此容易看出，$\mathcal L_1$ 在 $\overline M_{0,3}(X,0)$ 上平凡，其 Chern 类为零。于是，当 $k<n+1$ 时，式~\eqref{eq:degree-zero-descendant-divisors} 为零，综上可得
\begin{equation}\label{eq:degree-zero-descendant-kronecker}
\langle\tau_nT_a,T_j,(\delta)^k\rangle_{0,0}
=\left(\int_XT_a\cup T_j\cup\delta^k\right)\delta_{k,n+1}.
\end{equation}
利用式~\eqref{eq:descendant-divisor-binomial} 和~\eqref{eq:degree-zero-descendant-kronecker}，$\langle\!\langle\tau_nT_a,T_j\rangle\!\rangle_0$ 的公式简化为
\[
\begin{aligned}
&\frac1{(n+1)!}\int_XT_a\cup T_j\cup\delta^{n+1}\\
&\quad+\sum_{\beta\ne0}\sum_{k=0}^{\infty}\sum_{\mu+\nu=k}
\frac1{\mu!\nu!}\left(\int_\beta\delta\right)^\mu
\langle\tau_{n-\nu}(T_a\cup\delta^\nu),T_j\rangle_{0,\beta}\\
&=\frac1{(n+1)!}\int_XT_a\cup T_j\cup\delta^{n+1}
+\sum_{\beta\ne0}\sum_{\nu=0}^{\infty}
e^{\int_\beta\delta}\frac1{\nu!}
\langle\tau_{n-\nu}(T_a\cup\delta^\nu),T_j\rangle_{0,\beta}.
\end{aligned}
\]
将此表达式代入 $s_a$ 的定义，得到
\begin{equation}\label{eq:restricted-flat-section-expanded}
\begin{aligned}
s_a={}&T_a+\sum_j\left(\sum_{n=0}^{\infty}\hbar^{-(n+1)}
\frac1{(n+1)!}\int_XT_a\cup T_j\cup\delta^{n+1}\right)T^j\\
&+\sum_{\beta\ne0}\sum_j e^{\int_\beta\delta}
\left(\sum_{\nu=0}^{\infty}\sum_{n\ge\nu}\hbar^{-(n+1)}\frac1{\nu!}
\langle\tau_{n-\nu}(T_a\cup\delta^\nu),T_j\rangle_{0,\beta}\right)T^j.
\end{aligned}
\end{equation}
这个式子看起来复杂，但由恒等式
\begin{equation}\label{eq:poincare-dual-basis-reconstruction}
\sum_j\left(\int_X\phi\cup T_j\right)T^j=\phi,
\qquad\phi\in H^*(X,\mathbb C),
\end{equation}
式~\eqref{eq:restricted-flat-section-expanded} 第一行连同 $T_a$ 简化为 $T_a\cup e^{\delta/\hbar}$。第二行括号内的表达式则可令 $k=n-\nu$ 来简化；由于 $\hbar^{-(n+1)}=\hbar^{-(k+1)}\hbar^{-\nu}$，它成为
\[
\sum_{k=0}^{\infty}\hbar^{-(k+1)}
\langle\tau_k(T_a\cup e^{\delta/\hbar}),T_j\rangle_{0,\beta}.
\]
按式~\eqref{eq:giventals-flat-section-symbolic} 引入的\IndexLocation{indexloc:0476:338}符号记法，它等于
\[
\left\langle\frac{T_a\cup e^{\delta/\hbar}}{\hbar-c},T_j\right\rangle_{0,\beta},
\]
由此立即得到命题。
\end{proof}

\subsection{A 模型联络的平坦截面}\label{subsec:a-model-gravitational-flat-sections}
设 $V$ 是复维数 $d$ 的 Calabi--Yau 流形。在\hyperref[subsec:a-model-connection]{第~\ref*{subsec:a-model-connection} 小节}中，按方程
\[
\nabla_{\frac{\partial}{\partial u_j}}(T_k)=T_j*_{\mathrm{small}}T_k,
\qquad1\le j\le r,\quad0\le k\le m,
\]
定义了 A 模型联络，其中 $u_j=\frac1{2\pi i}t_j$。若将 Givental 联络 $\nabla^g$ 限制到 $H^2(V,\mathbb C)$，则式~\eqref{eq:restricted-giventals-connection} 表明，在 $\hbar=-2\pi i$ 时，$-2\pi i\nabla$ 恰好就是 $\nabla^g$。本小节余下部分固定使用这一 $\hbar$ 值。

于是，\hyperref[subsec:giventals-flat-sections]{第~\ref*{subsec:giventals-flat-sections} 小节}中定义的截面 $s_a$ 限制后成为 A 模型联络 $\nabla$ 的平坦截面。为使用这些截面，引入如下通用记号：给定上同调类 $T\in H^*(V,\mathbb C)$，令
\[
s(T)=e^{\delta/\hbar}\cup T
+\sum_{\beta\ne0}\sum_{j=0}^m e^{\int_\beta\delta}
\left\langle\frac{e^{\delta/\hbar}\cup T}{\hbar-c},T_j\right\rangle_{0,\beta}T^j.
\]
\SourcePage{339}
由命题~\ref{prop:restricted-giventals-flat-sections}，$s(T_a)$ 就是截面 $s_a$；容易推出，对任意 $T\in H^*(V,\mathbb C)$，$s(T)$ 都是 $\nabla$ 的形式平坦截面。沿用\hyperref[subsec:quantum-product-dependence]{第~\ref*{subsec:quantum-product-dependence} 小节}的处理，令 $q^\beta=e^{\int_\beta\delta}$，于是可写成
\begin{equation}\label{eq:a-model-flat-section-gravitational}
s(T)=e^{\delta/\hbar}\cup T
+\sum_{\beta\ne0}\sum_{j=0}^m q^\beta
\left\langle\frac{e^{\delta/\hbar}\cup T}{\hbar-c},T_j\right\rangle_{0,\beta}T^j.
\end{equation}
利用这一公式，容易确定 $\nabla$ 的单值性。假设处于\hyperref[subsec:a-model-asymptotics]{第~\ref*{subsec:a-model-asymptotics} 小节}的情形：$T_1,\ldots,T_r$ 是位于 $V$ 的 Kähler 锥闭包中的整数类，变量 $q_j=e^{t_j}=e^{2\pi iu_j}$ 使 $\nabla$ 成为 $(\Delta^*)^r$ 中原点附近的联络，并采用该小节同样的收敛假设。

令 $\mathcal T_j$ 是绕 $(\Delta^*)^r$ 中第 $j$ 个因子一周所给出的\IndexLocation{indexloc:0901:339}\IndexLocation{indexloc:0902:339}单值性变换。
\begin{theorem}\label{thm:a-model-gravitational-monodromy}
单值性 $\mathcal T_j$ 及其对数 $N_j$ 如下作用于 A 模型联络的平坦截面：
\[
\mathcal T_j(s(T))=s(e^{-T_j}\cup T),
\qquad N_j(s(T))=-s(T_j\cup T).
\]
\end{theorem}
\begin{proof}
由于 $q_j=e^{t_j}$，$\mathcal T_j$ 将 $t_j$ 变为 $t_j+2\pi i$。上面已知 $\nabla$ 是 $\hbar=-2\pi i$ 时 Givental 联络的一个倍数，因此可将 $\mathcal T_j$ 的作用写为 $t_j\mapsto t_j-\hbar$。

现在考察此作用对式~\eqref{eq:a-model-flat-section-gravitational} 的影响。因为 $\delta=\sum_{i=1}^rt_iT_i$，$\mathcal T_j$ 将 $\delta$ 变为 $\delta+2\pi iT_j=\delta-\hbar T_j$。它不影响 $q^\beta=e^{\int_\beta\delta}$，因为 $T_j$ 和 $\beta$ 都是整数类。不过，$\delta$ 也出现在 $e^{\delta/\hbar}$ 中，此处 $\mathcal T_j$ 的作用为
\[
e^{\delta/\hbar}\longmapsto
e^{(\delta-\hbar T_j)/\hbar}
=e^{\delta/\hbar-T_j}=e^{\delta/\hbar}\cup e^{-T_j}.
\]
最后一个等号使用了 $T_j$ 的次数为偶数、因此与 $\delta$ 交换这一事实。将它代入式~\eqref{eq:a-model-flat-section-gravitational}，立即得到 $\mathcal T_j(s(T))$ 的公式；取对数即得 $N_j(s(T))$ 的公式。
\end{proof}

定理~\ref{thm:a-model-gravitational-monodromy} 表明，单值性对数 $N_j$ 本质上是与 $-T_j$ 作杯积。更准确地说，$N_j(s(T_a))=-s(T_j\cup T_a)$ 意味着：与 $-T_j$ 作杯积相对于上同调基 $\{T_a\}$ 的矩阵，恰好是 $N_j$ 相对于平坦截面基 $\{s(T_a)\}$ 的矩阵。这是\hyperref[chap:quantum-cohomology]{第八章}中经常使用的事实。

还可用定理~\ref{thm:a-model-gravitational-monodromy} 研究 $\nabla$ 的典范延拓。回顾\hyperref[subsec:limiting-hodge]{第~\ref*{subsec:limiting-hodge} 小节}：由于 $s(T)$ 是 $\nabla$ 的平坦截面，截面
\[
\widetilde s(T)=\exp\left(-\frac1{2\pi i}\sum_j\log(q_j)N_j\right)s(T)
\]
自然延伸为典范延拓上的截面。于是，由定理~\ref{thm:a-model-gravitational-monodromy} 得到以下公式。
\begin{corollary}\label{cor:a-model-canonical-gravitational-section}
对 $T\in H^*(V,\mathbb C)$，截面 $\widetilde s(T)$ 为
\[
\widetilde s(T)=T+\sum_{\beta\ne0}\sum_{j=0}^m q^\beta
\left\langle\frac{T}{\hbar-c},T_j\right\rangle_{0,\beta}T^j.
\]
\end{corollary}
\SourcePage{340}
\begin{proof}
证明很简单。先注意，$\log(q_j)=t_j$、$\hbar=-2\pi i$ 给出
\[
\begin{aligned}
\widetilde s(T)
&=\exp\left(\hbar^{-1}\sum_jt_jN_j\right)s(T)\\
&=s\left(\exp\left(\hbar^{-1}\sum_jt_j(-T_j)\right)\cup T\right)
=s(e^{-\delta/\hbar}\cup T),
\end{aligned}
\]
第二个等号来自 $N_j(s(T))=s((-T_j)\cup T)$。又有 $\delta=\sum_jt_jT_j$。根据式~\eqref{eq:a-model-flat-section-gravitational}，
\[
\begin{aligned}
s(e^{-\delta/\hbar}\cup T)
&=e^{\delta/\hbar}\cup(e^{-\delta/\hbar}\cup T)\\
&\quad+\sum_{\beta\ne0}\sum_{j=0}^m q^\beta
\left\langle\frac{e^{\delta/\hbar}\cup(e^{-\delta/\hbar}\cup T)}{\hbar-c},T_j\right\rangle_{0,\beta}T^j\\
&=T+\sum_{\beta\ne0}\sum_{j=0}^m q^\beta
\left\langle\frac{T}{\hbar-c},T_j\right\rangle_{0,\beta}T^j.
\end{aligned}
\]
结合这些公式，立即得到推论。
\end{proof}

注意，$\widetilde s(T)$ 的公式只包含 $q^\beta=\prod_iq_i^{\int_\beta T_i}$。如往常一样，可假定 $\beta$ 有效；由于 $T_i$ 位于 Kähler 锥闭包中，有 $\int_\beta T_i\ge0$。因此，$\widetilde s(T)$ 全纯延伸到 $0$。

利用上面的推论，现在终于可以证明\hyperref[chap:quantum-cohomology]{第八章}的命题~\ref{prop:a-model-canonical-flat-basis}。将该命题重新陈述如下。
\begin{corollary}\label{cor:a-model-canonical-section-triangular}
对每个齐次类 $T\in H^*(V,\mathbb C)$，截面 $\widetilde s(T)$ 满足
\[
\widetilde s(T)=T+\text{次数更高的项},\qquad\widetilde s(T)(0)=T.
\]
此外，$N_j$ 作用于 $\widetilde s(T_a)$ 的矩阵，等于与 $-T_j$ 作杯积作用于 $T_a$ 的矩阵。
\end{corollary}
\begin{proof}
最后一个断言立即来自定理~\ref{thm:a-model-gravitational-monodromy} 证明后的说明。对于其他断言，先将推论~\ref{cor:a-model-canonical-gravitational-section} 的符号记法展开，得到
\[
\widetilde s(T)=T+\sum_{\beta\ne0}\sum_{j=0}^m\sum_{n=0}^{\infty}
q^\beta\hbar^{-(n+1)}\langle\tau_nT,T_j\rangle_{0,\beta}T^j.
\]
由于 $V$ 为 Calabi--Yau，\IndexLocation{indexloc:0461:340}次数公理意味着，除满足
\[
2n+\deg T+\deg T_j=2d-2,\qquad d=\dim V
\]
的 $j$ 外，所有项均为零。又因 $\deg T^j=2d-\deg T_j$，有 $\deg T^j=\deg T+2n+2$，故上述求和中任何非零项的次数均大于 $\deg T$。

最后，要证明所有 $q_i=0$ 时 $\widetilde s(T)=T$。前面已指出，可以将求和限制到使
\[
q^\beta=\prod_iq_i^{\int_\beta T_i}
\]
各指数均非负的 $\beta\ne0$。因此，只需证明，每个这样的有效 $\beta$ 至少有一个指数为正。否则，就有 $\int_\beta T_i=0$ 对所有 $i$ 成立；因为 $T_i$ 构成 $H^2(V,\mathbb C)$ 的基，这意味着 $\beta$ 与每个 Kähler 类的配对均为零，与非零有效曲线类的正性矛盾。
\end{proof}

\SourcePage{341}
首次在式~\eqref{eq:giventals-flat-section-symbolic} 中引入符号表达式 $T/(\hbar-c)$ 时，它可能显得有些人为。不过，它对定理~\ref{thm:a-model-gravitational-monodromy} 和推论~\ref{cor:a-model-canonical-gravitational-section} 的证明所带来的简化，已充分说明其价值。在\hyperref[sec:quantum-cohomology-relations]{第~\ref*{sec:quantum-cohomology-relations} 节}中，还会经常遇到这个记号。

\hyperref[chap:quantum-cohomology]{第八章}已研究过 Calabi--Yau 三维簇的 A 模型联络的平坦截面。下面比较那里的公式与这里得到的公式。
\SourceExampleNumber{10.2.2}{1}
\begin{example}\label{ex:threefold-flat-section-comparison}
设 $V$ 为 Calabi--Yau 三维簇。在\hyperref[subsec:a-model-asymptotics]{第~\ref*{subsec:a-model-asymptotics} 小节}中，利用式~\eqref{eq:a-model-potential} 的 Gromov--Witten 势
\[
\begin{aligned}
\Phi&=\frac{(2\pi i)^3}{6}\int_V\left(\sum_{j=1}^ru_jT_j\right)^3
+\sum_{\beta\ne0}N_\beta q^\beta\\
&=\frac{(2\pi i)^3}{6}\int_V\left(\sum_{j=1}^ru_jT_j\right)^3+\Phi_{\mathrm{hol}}
\end{aligned}
\]
构造了 $\nabla$ 的平坦截面（如往常一样，$N_\beta=\langle I_{0,0,\beta}\rangle$）。这里限制到偶上同调。在命题~\ref{prop:a-model-canonical-flat-basis} 的证明中，先用式~\eqref{eq:a-model-flat-sections-threefold} 写出平坦截面，再在式~\eqref{eq:a-model-canonical-sections-threefold} 中，对它们施加 $\exp(-\sum_j u_jN_j)=\exp(-\frac1{2\pi i}\sum_j\log(q_j)N_j)$。特别地，得到公式
\begin{equation}\label{eq:threefold-canonical-unit-potential}
\widetilde s_0=T_0-\frac1{(2\pi i)^3}\sum_j
\frac{\partial\Phi_{\mathrm{hol}}}{\partial u_j}T^j
+\frac2{(2\pi i)^3}\Phi_{\mathrm{hol}}T^0.
\end{equation}
其中 $T_0=1$，$T_j$（$1\le j\le r$）构成 $H^2(V,\mathbb C)$ 的基，$T^j\in H^4(V,\mathbb C)$（$1\le j\le r$）及 $T^0\in H^6(V,\mathbb C)$ 为对偶类。

按\hyperref[subsec:a-model-asymptotics]{第~\ref*{subsec:a-model-asymptotics} 小节}的术语，$\widetilde s_0$ 是 $\nabla^c$ 平坦截面，且已证明其在 $0$ 的值为 $T_0$。推论~\ref{cor:a-model-canonical-gravitational-section} 表明，截面 $\widetilde s(T_0)$ 也对 $\nabla^c$ 平坦，而推论~\ref{cor:a-model-canonical-section-triangular} 表明，它在 $0$ 的值也是 $T_0$。因此 $\widetilde s_0=\widetilde s(T_0)$。

为看清它的含义，用上述上同调基写出 $\widetilde s(T_0)$。由于 $T_0=1$，推论~\ref{cor:a-model-canonical-gravitational-section} 的公式成为
\begin{equation}\label{eq:threefold-canonical-unit-descendants}
\begin{aligned}
\widetilde s(T_0)=T_0+\sum_{\beta\ne0}q^\beta\biggl(&
\left\langle\frac1{\hbar-c},T^0\right\rangle_{0,\beta}T_0
+\sum_j\left\langle\frac1{\hbar-c},T^j\right\rangle_{0,\beta}T_j\\
&+\sum_j\left\langle\frac1{\hbar-c},T_j\right\rangle_{0,\beta}T^j
+\left\langle\frac1{\hbar-c},T_0\right\rangle_{0,\beta}T^0\biggr).
\end{aligned}
\end{equation}
比较式~\eqref{eq:threefold-canonical-unit-potential} 和~\eqref{eq:threefold-canonical-unit-descendants} 的系数。对于 $T^0$ 的系数，先注意
\[
\left\langle\frac1{\hbar-c},T_0\right\rangle_{0,\beta}
=\sum_{n=0}^{\infty}\hbar^{-(n+1)}\langle\tau_n,1\rangle_{0,\beta}
=\hbar^{-3}\langle\tau_2,1\rangle_{0,\beta},
\]
因为由次数公理，其余项均为零。比较 $T^0$ 的系数，并使用上述 $\Phi$ 的公式，得到恒等式
\[
\frac2{(2\pi i)^3}N_\beta
=\hbar^{-3}\langle\tau_2,1\rangle_{0,\beta}.
\]
由于 $\hbar=-2\pi i$，有 $\langle\tau_2,1\rangle_{0,\beta}=-2N_\beta$，恰好与例题~\ref{ex:cy-threefold-gravitational} 一致。对于式~\eqref{eq:threefold-canonical-unit-descendants} 求和中的其他项，由次数公理和\IndexLocation{indexloc:0464:341}基本类公理可知，$T_0,T_j$ 的系数均为零。最后，式~\eqref{eq:threefold-canonical-unit-potential} 和~\eqref{eq:threefold-canonical-unit-descendants} 中 $T^j$ 系数相等，等价于
\[
\langle\tau_1,T_j\rangle_{0,\beta}
=-\left(\int_\beta T_j\right)N_\beta,
\]
它来自\IndexLocation{indexloc:0463:341}除子公理与\IndexLocation{indexloc:0462:341}伸缩公理。
\end{example}

\SourcePage{342}
\subsection{参数 \texorpdfstring{$\hbar$}{hbar}}\label{subsec:equivariant-giventals-connection}
本节最后简要解释，参数 $\hbar$ 如何自然地来自等变上同调。以下讨论取自 \SourceCite{Givental2}。

考虑 $\mathbb P^1$ 上的 $\mathbb C^*$ 作用 $t\cdot(x_0,x_1)=(t^{-1}x_0,x_1)$，令 $\hbar\in H^2(B\mathbb C^*)$ 为相应的生成元，故 $H^*(B\mathbb C^*)=\mathbb C[\hbar]$。可通过代换 $\lambda_0\mapsto\hbar$、$\lambda_1\mapsto0$，将它与例题~\ref{ex:projective-equivariant-localization} 中的 $(\mathbb C^*)^2$ 作用比较。于是
\[
H_{\mathbb C^*}^*(\mathbb P^1)\simeq\mathbb C[p,\hbar]/(p(p-\hbar)).
\]
令 $\mathcal R_{\mathbb C^*}=\mathbb C(\hbar)$ 为 $H^*(B\mathbb C^*)$ 的分式域，这是\hyperref[subsec:atiyah-bott-localization]{第~\ref*{subsec:atiyah-bott-localization} 小节}的记号。则 $\{p/\hbar,(\hbar-p)/\hbar\}$ 构成 $H_{\mathbb C^*}^*(\mathbb P^1)\otimes_{H^*(B\mathbb C^*)}\mathcal R_{\mathbb C^*}$ 作为 $\mathcal R_{\mathbb C^*}$ 向量空间的一组基。这组基尤其适合\hyperref[subsec:atiyah-bott-localization]{第~\ref*{subsec:atiyah-bott-localization} 小节}的局部化定理。（注意，也可将 $\mathcal R_{\mathbb C^*}$ 替换为较小的环 $\mathbb C[\hbar,\hbar^{-1}]$。）

\hyperref[chap:mirror-theorem]{第十一章}将使用如下重要思想。给定稳定映射 $f:(\mathbb P^1,p_1,\ldots,p_n)\to X$，令 $\beta=f_*[\mathbb P^1]\in H_2(X,\mathbb Z)$，则 $f$ 的图像自然给出 $\overline M_{0,n}(X\times\mathbb P^1,(\beta,1))$ 的一个元素。此空间相对于 $\overline M_{0,n}(X,\beta)$ 的一个优点是：$\mathbb C^*$ 通过在 $X$ 上的平凡作用和上述 $\mathbb P^1$ 作用，作用于 $X\times\mathbb P^1$。

由于这个 $\mathbb C^*$ 作用，可使用\hyperref[sec:equivariant-gw]{第~\ref*{sec:equivariant-gw} 节}中定义的等变 Gromov--Witten 不变量。为简化叙述，假定 $X$ 光滑且凸。于是 $X\times\mathbb P^1$ 也凸，定义~\ref{def:equivariant-gw} 适用。给定 $\gamma\in H_{\mathbb C^*}^*(X\times\mathbb P^1)$，得到\IndexLocation{indexloc:0536:342}等变不变量
\[
\begin{aligned}
\langle I_{0,n,(\beta,1)}\rangle_{\mathbb C^*}(\gamma,\ldots,\gamma)
&=\int_{\overline M_{0,n}(X\times\mathbb P^1,(\beta,1))_{\mathbb C^*}}
e_1^*(\gamma)\cup\cdots\cup e_n^*(\gamma)\\
&\in H^*(B\mathbb C^*).
\end{aligned}
\]
将它看作分式域 $\mathcal R_{\mathbb C^*}=\mathbb C(\hbar)$ 中的元素较为方便。利用这些不变量，定义等变势函数
\[
\mathcal F^1(\gamma)=\sum_{n=0}^{\infty}\sum_{\beta\in H_2(X)}\frac1{n!}
\langle I_{0,n,(\beta,1)}\rangle_{\mathbb C^*}(\gamma,\ldots,\gamma)q^\beta.
\]
为将它写成形式幂级数，引入一些变量。注意
\[
\begin{aligned}
H_{\mathbb C^*}^*(X\times\mathbb P^1)\otimes_{H^*(B\mathbb C^*)}\mathcal R_{\mathbb C^*}
=H^*(X)\otimes_{\mathbb C}H_{\mathbb C^*}^*(\mathbb P^1)
\otimes_{H^*(B\mathbb C^*)}\mathcal R_{\mathbb C^*}.
\end{aligned}
\]
因此，若 $T_i$ 为 $H^*(X)$ 的一组上同调基，则 $\{T_i\otimes p/\hbar,T_i\otimes(\hbar-p)/\hbar\}$ 是局部化上同调在 $\mathcal R_{\mathbb C^*}$ 上的一组基。然后将
\[
\gamma=\sum_jt_jT_j\otimes p/\hbar
+\sum_i\tau_iT_i\otimes(\hbar-p)/\hbar
\]
代入 $\mathcal F^1$ 的公式。可将 $t_j$、$\tau_i$ 看作取值于 $\mathcal R_{\mathbb C^*}$ 的变量。

用变量 $t_j,\tau_i$ 表示 $\mathcal F^1$ 后，会出现一些很好的性质。先考虑二阶导数
\[
\Phi_{ij}=\frac{\partial^2\mathcal F^1}{\partial\tau_i\partial t_j}.
\]
将 $\Phi=(\Phi_{ij})$ 看作取矩阵值的形式函数。Givental 表明，$\Phi$ 的转置给出 $\nabla^g$ 的\IndexLocation{indexloc:0440:342}平坦截面。
\begin{proposition}\label{prop:graph-space-flat-matrix}
若 $\Phi^*$ 表示 $\Phi$ 的转置，则对所有 $j$，
\[
\hbar\frac{\partial}{\partial t_j}\Phi^*=T_j*\Phi^*,
\]
其中 $s_j$ 是式~(10.15) 所定义的平坦截面。
\end{proposition}
\SourcePage{343}
\begin{proof}
证明类似于命题~\ref{prop:giventals-flat-sections} 的证明，但因为这里处理的是等变 Gromov--Witten 不变量而非引力关联数，稍为简单。详细证明见 \SourceCite{Givental2}。
\end{proof}

$\Phi_{ij}$ 还有一个很好的方面，即可自然地用引力关联数解释。为此，令
\[
\psi_{ij}=g_{ij}+\left\langle\!\left\langle\frac{T_j}{\hbar-c},T_i\right\rangle\!\right\rangle_0.
\]
由式~\eqref{eq:giventals-flat-section-symbolic}，这就是将 $s_j$ 在对偶基中展开后 $T^i$ 的系数，其中，如往常一样，$\{T^i\}$ 为 $\{T_i\}$ 的对偶基。可以将 $\psi_{ij}$ 看作 $t_i$ 和 $\hbar$ 的函数，即 $\psi_{ij}=\psi_{ij}(t_0,\ldots,t_m,\hbar)$。它与上述函数 $\Phi_{ij}$ 的关系如下。
\begin{lemma}\label{lem:graph-space-descendant-factorization}
\[
\Phi_{ij}=-\sum_{k,\ell}\hbar^{-2}
\psi_{ik}(\tau_0,\ldots,\tau_m,-\hbar)g^{k\ell}
\psi_{j\ell}(t_0,\ldots,t_m,\hbar).
\]
\end{lemma}
\begin{proof}
\SourceCite{Givental2} 证明了这个引理，这里只简述证明。关键是，如\hyperref[chap:localization]{第九章}中一样，用局部化计算左端。$X\times\mathbb P^1$ 上的 $\mathbb C^*$ 作用诱导稳定映射空间上的作用。可以看到，固定点轨迹由稳定映射 $f:C\to X\times\mathbb P^1$ 组成，其中 $C=C'\cup C_0\cup C_\infty$，满足 $f(C_0)\subset X\times\{0\}$、$f(C_\infty)\subset X\times\{\infty\}$，而对某点 $x\in X$，有 $f(C')\subset\{x\}\times\mathbb P^1$。分支 $C_0$、$C_\infty$ 可为空。分支 $C'$ 分别在一个点与 $C_0$、$C_\infty$ 相交，将引理~\ref{lem:node-smoothing-tangent-tensor} 应用于这两个结点，便得到 $\psi_{ij}$ 定义中的 $\hbar-c$ 和 $-\hbar-c$ 因子。
\end{proof}

最后应说明，这些思想也可用于给出命题~\ref{prop:giventals-flat-sections} 的另一证明：再次使用局部化，从引理~\ref{lem:graph-space-descendant-factorization} 和命题~\ref{prop:graph-space-flat-matrix} 推出该结果。

\section{量子上同调中的关系}\label{sec:quantum-cohomology-relations}
本节用式~\eqref{eq:giventals-flat-section-series} 中的平坦截面 $s_a$ 研究量子上同调。将 $\nabla^g$ 看作底空间 $M=H^0(X,\mathbb C)\oplus H^2(X,\mathbb C)$ 上平凡丛 $M\times H^*(X,\mathbb C)$ 的形式联络。因此，变量为 $t_0,t_1,\ldots,t_r$，$\nabla^g$ 由式~\eqref{eq:restricted-giventals-connection} 中的小量子乘积给出。

\subsection{Givental 函数 \texorpdfstring{$J$}{J}}\label{subsec:giventals-j-function}\index{Givental $J$ 函数|hyperpage}\index{J函数@$J$ 函数|hyperpage}
遵循 Givental，定义函数
\begin{equation}\label{eq:giventals-j-pairing-definition}
J=\sum_j\langle s_j,1\rangle T^j,
\end{equation}
其中 $\langle\alpha,\beta\rangle=\int_X\alpha\cup\beta$ 是上同调上通常的交配对。利用式~\eqref{eq:giventals-flat-section-series}，也可将 $J$ 的定义写成
\begin{equation}\label{eq:giventals-j-descendant-series}
J=J(t_0,\ldots,t_m,\hbar^{-1})
=1+\sum_{n=0}^{\infty}\sum_{a=0}^m\hbar^{-(n+1)}
\langle\!\langle\tau_nT_a,1\rangle\!\rangle_0T^a.
\end{equation}
\hyperref[subsec:quantum-product-dependence]{第~\ref*{subsec:quantum-product-dependence} 小节}已说明，变量 $t_1,\ldots,t_r$ 在\IndexLocation{indexloc:1122:343}大量子上同调中以指数 $q_j=e^{t_j}$ 出现。因此，将小量子乘积看作大量子乘积的限制时，给出 $*_{\mathrm{small}}$ 的级数是
\SourcePage{344}
$q_j$ 的形式幂级数。为简便起见，以下将 $*_{\mathrm{small}}$ 简记为 $*$。

设 $P(\hbar\partial/\partial t,e^t,\hbar)$ 是以下量的形式幂级数：
\[
\hbar\frac{\partial}{\partial t_0},\ldots,
\hbar\frac{\partial}{\partial t_r},e^{t_0},\ldots,e^{t_r},\hbar.
\]
为明确起见，可以且将假定，指数项都放在微分算子左侧，尽管下面的定理~\ref{thm:quantum-differential-relations} 并不需要这一条件。

若 $H^*(X,\widehat C)$ 表示系数为 $\widehat C=\mathbb C[[q_1,\ldots,q_r]]$ 的上同调，则以 $P(T,q,0)$ 表示对 $P(\hbar\partial/\partial t,e^t,\hbar)$ 作代换
\[
\hbar\frac{\partial}{\partial t_j}\longmapsto T_j,
\qquad e^{t_j}\longmapsto q_j,\qquad\hbar\longmapsto0
\]
并将微分算子的复合替换为量子乘积后，所得的 $H^*(X,\widehat C)$ 中的形式幂级数。若 $P$ 还含 $e^{t_0}$，则在系数环中相应增添标量参数 $q_0$。关键结果如下。
\begin{theorem}\label{thm:quantum-differential-relations}
设 $P(\hbar\partial/\partial t,e^t,\hbar)J=0$，其中 $P$ 是按上述记号表示的形式幂级数。则关系 $P(T,q,0)=0$ 在小量子上同调中成立。
\end{theorem}
\begin{proof}
以下证明由 B. Kim 提出。引入\emph{对偶 Givental 联络} $\widetilde\nabla^g$，定义为
\[
\widetilde\nabla^g_{\frac{\partial}{\partial t_i}}\left(\sum_j a_jT_j\right)
=\hbar\sum_j\frac{\partial a_j}{\partial t_i}T_j
+\sum_j a_jT_i*T_j.
\]
简单计算可验证，对变量 $t_0,\ldots,t_m$ 的取值于 $H^*(X,\mathbb C)$ 的函数 $G,H$，有恒等式
\begin{equation}\label{eq:dual-giventals-pairing-derivative}
\hbar\frac{\partial}{\partial t_i}\langle G,H\rangle
=\left\langle\nabla^g_{\frac{\partial}{\partial t_i}}G,H\right\rangle
+\left\langle G,\widetilde\nabla^g_{\frac{\partial}{\partial t_i}}H\right\rangle.
\end{equation}
又可算得
\begin{equation}\label{eq:dual-giventals-classical-symbol}
\widetilde\nabla^g_{\frac{\partial}{\partial t_{i_1}}}\cdots
\widetilde\nabla^g_{\frac{\partial}{\partial t_{i_k}}}1
=T_{i_1}*\cdots*T_{i_k}+\hbar A(T,q,\hbar),
\end{equation}
其中 $A$ 是变量 $T_i,q_i,\hbar$ 的形式幂级数。

表达式 $PJ$ 定义为 $\sum_jP\langle s_j,1\rangle T^j$。由于 $T^0,\ldots,T^m$ 线性无关，$PJ=0$ 意味着，对所有 $j$，均有 $P(\hbar\partial/\partial t,e^t,\hbar)\langle s_j,1\rangle=0$。命题~\ref{prop:giventals-flat-sections} 又给出 $\nabla^g_{\partial/\partial t_i}s_j=0$。重复应用式~\eqref{eq:dual-giventals-pairing-derivative} 和~\eqref{eq:dual-giventals-classical-symbol}，可得
\[
\begin{aligned}
0=P\langle s_j,1\rangle
&=\langle s_j,P(T,q,0)\rangle
+\hbar\langle s_j,B(T,q,\hbar)\rangle,\\[-2pt]
&\hspace{10em}j=0,\ldots,m,
\end{aligned}
\]
其中 $B(T,q,\hbar)$ 是形式幂级数。由于 $s_j$ 构成一组基，且交配对非退化，先得到 $P(T,q,0)+\hbar B(T,q,\hbar)=0$。这是 $\hbar$ 的非负幂级数恒等式，因此可令 $\hbar=0$，得到 $P(T,q,0)=0$，并有
\[
\langle s_j,P(T,q,0)\rangle=0,\qquad j=0,\ldots,m.
\]
所需关系便在 $H^*(X,\widehat C)$ 的小量子乘积下成立。
\end{proof}

这个定理导出如下定义。
\begin{definition}\label{def:quantum-differential-operator}\index{量子微分方程|hyperpage}
满足定理~\ref{thm:quantum-differential-relations} 假设的微分算子 $P(\hbar\partial/\partial t,e^t,\hbar)$ 称为\emph{\IndexLocation{indexloc:1151:344}\IndexLocation{indexloc:1154:344}量子微分算子}。方程 $P(\hbar\partial/\partial t,e^t,\hbar)Y=0$（由 $Y=J$ 满足）称为\emph{量子微分方程}。
\end{definition}

\SourcePage{345}
也可以如下理解定理~\ref{thm:quantum-differential-relations}。令 $P_{\widetilde\nabla^g}$ 表示对 $P(\hbar\partial/\partial t,e^t,\hbar)$ 作代换
\[
\hbar\frac{\partial}{\partial t_i}\longmapsto
\widetilde\nabla^g_{\frac{\partial}{\partial t_i}}
\]
所得的表达式。由于 $s_j$ 是 $\nabla^g$ 的平坦截面，式~\eqref{eq:dual-giventals-pairing-derivative} 给出
\begin{equation}\label{eq:j-annihilator-dual-connection}
P\langle s_j,1\rangle=\langle s_j,P_{\widetilde\nabla^g}1\rangle.
\end{equation}
因此，对于 Givental 函数 $J$，有等价关系
\begin{equation}\label{eq:j-annihilator-equivalences}
PJ=0\ \Longleftrightarrow\ P\langle s_j,1\rangle=0\text{ 对所有 }j
\ \Longleftrightarrow\ P_{\widetilde\nabla^g}1=0.
\end{equation}
第一个等价关系来自定理~\ref{thm:quantum-differential-relations} 的证明；第二个来自式~\eqref{eq:j-annihilator-dual-connection}，因为 $s_j$ 成基。这解释了式~\eqref{eq:giventals-j-pairing-definition} 中“$1$”的作用。

注意，若 $P(T,q)=0$ 是小量子上同调中的关系，则 $P(\hbar\partial/\partial t,e^t)$ 不一定是量子微分算子（见例题~\ref{ex:cy-quantum-equation-relation}）。不过，若此关系的次数至多为二，且指数项放在微分算子的左边，就容易看出，$P(\hbar\partial/\partial t,e^t)$ 是量子微分算子。此外，下面的例题~\ref{ex:second-order-quantum-relations} 将说明，某些特殊情形下，可直接由量子上同调中的关系导出量子微分方程。最后还应说明，\hyperref[chap:mirror-theorem]{第十一章}的例题~\ref{ex:toric-fano-and-hirzebruch-quantum-ring} 将利用定理~\ref{thm:quantum-differential-relations} 计算某些环面簇的量子上同调。

接着要简化 $J$ 的公式，以便计算量子微分算子的一些例子。以下引理记录了两个经常使用的公式。
\begin{lemma}\label{lem:giventals-j-two-formulas}
$X$ 的 Givental $J$ 函数由以下两个公式给出：
\[
\begin{aligned}
J&=e^{(t_0+\delta)/\hbar}\left(1+\sum_{\beta\ne0}\sum_{a=0}^m q^\beta
\left\langle\frac{T_a}{\hbar-c},1\right\rangle_{0,\beta}T^a\right)\\
&=e^{(t_0+\delta)/\hbar}\left(1+\sum_{\beta\ne0}q^\beta PD^{-1}e_{1*}
\left(\frac1{\hbar-c}\cap[\overline M_{0,2}(X,\beta)]^{\mathrm{virt}}\right)\right),
\end{aligned}
\]
其中 $\delta=\sum_{i=1}^rt_iT_i$，$q^\beta=e^{\int_\beta\delta}$。第二个等式中，$PD$ 是 Poincaré 对偶，$e_1:\overline M_{0,2}(X,\beta)\to X$ 是第一个标记点处的赋值映射。
\end{lemma}
\begin{proof}
回顾 $J$ 是用平坦截面 $s_a$ 构造的，而我们限制在 $M=H^0(X,\mathbb C)\oplus H^2(X,\mathbb C)$ 上，因此命题~\ref{prop:restricted-giventals-flat-sections} 给出
\[
s_a=e^{t_0/\hbar}\left(e^{\delta/\hbar}\cup T_a
+\sum_{\beta\ne0}\sum_{j=0}^m q^\beta
\left\langle\frac{e^{\delta/\hbar}\cup T_a}{\hbar-c},T_j\right\rangle_{0,\beta}T^j\right).
\]
\SourcePage{346}
故 Givental 函数 $J$ 为
\begin{equation}\label{eq:j-function-flat-section-expansion}
\begin{aligned}
J&=\sum_{a=0}^m\langle s_a,1\rangle T^a\\
&=e^{t_0/\hbar}\biggl(\sum_{a=0}^m\langle e^{\delta/\hbar}\cup T_a,1\rangle T^a\\
&\hspace{4em}+\sum_{a=0}^m\sum_{\beta\ne0}\sum_{j=0}^m q^\beta
\left\langle\frac{e^{\delta/\hbar}\cup T_a}{\hbar-c},T_j\right\rangle_{0,\beta}
\langle T^j,1\rangle T^a\biggr).
\end{aligned}
\end{equation}
第二行的求和容易简化：
\[
\sum_{a=0}^m\langle e^{\delta/\hbar}\cup T_a,1\rangle T^a
=\sum_{a=0}^m\langle e^{\delta/\hbar},T_a\rangle T^a=e^{\delta/\hbar},
\]
最后一个等号来自式~\eqref{eq:poincare-dual-basis-reconstruction}。第三行更简单，因为
\[
\langle T^j,1\rangle=
\begin{cases}1,&T_j=T_0=1,\\0,&\text{其他情形}.
\end{cases}
\]
故第三行只有 $j=0$ 的项有贡献。因此
\begin{equation}\label{eq:j-function-exponential-insertion}
J=e^{t_0/\hbar}\left(e^{\delta/\hbar}
+\sum_{\beta\ne0}\sum_{a=0}^m q^\beta
\left\langle\frac{e^{\delta/\hbar}\cup T_a}{\hbar-c},1\right\rangle_{0,\beta}T^a\right).
\end{equation}
还可进一步简化。虽然上同调基 $\{T_a\}$ 及其对偶基 $\{T^a\}$ 出现在式~\eqref{eq:j-function-exponential-insertion} 中，$J$ 并不依赖所用的基，故可用任意其他基替换 $\{T_a\}$。甚至可取非齐次基；由于使用形式系数，可取基 $\{e^{-\delta/\hbar}\cup T_a\}$。它的对偶基为 $\{e^{\delta/\hbar}\cup T^a\}$，因此
\[
\begin{aligned}
J=e^{t_0/\hbar}\biggl(&e^{\delta/\hbar}\\
&+\sum_{\beta\ne0}\sum_{a=0}^m q^\beta
\left\langle\frac{e^{\delta/\hbar}\cup(e^{-\delta/\hbar}\cup T_a)}{\hbar-c},1\right\rangle_{0,\beta}
e^{\delta/\hbar}\cup T^a\biggr),
\end{aligned}
\]
简化后便是所需的第一个公式。

为证明第二个公式，注意
\begin{equation}\label{eq:descendant-j-evaluation-integral}
\left\langle\frac{T_a}{\hbar-c},1\right\rangle_{0,\beta}
=\int_{[\overline M_{0,2}(X,\beta)]^{\mathrm{virt}}}\frac{e_1^*(T_a)}{\hbar-c}.
\end{equation}
由 $e_1$ 的投影公式，此式等于
\[
\int_XT_a\cup PD^{-1}\left(e_{1*}
\left(\frac1{\hbar-c}\cap[\overline M_{0,2}(X,\beta)]^{\mathrm{virt}}\right)\right).
\]
将 $\phi=PD^{-1}\bigl(e_{1*}((1/(\hbar-c))\cap[\overline M_{0,2}(X,\beta)]^{\mathrm{virt}})\bigr)$ 代入式~\eqref{eq:poincare-dual-basis-reconstruction}，再用式~\eqref{eq:descendant-j-evaluation-integral}，得到
\[
\begin{aligned}
\sum_{a=0}^m\left\langle\frac{T_a}{\hbar-c},1\right\rangle_{0,\beta}T^a
=PD^{-1}\left(e_{1*}\left(\frac1{\hbar-c}\cap[\overline M_{0,2}(X,\beta)]^{\mathrm{virt}}\right)\right).
\end{aligned}
\]
这表明，引理中的第二个公式立即来自第一个。
\end{proof}

\SourcePage{347}
引理~\ref{lem:giventals-j-two-formulas} 的一个重要推论，结合基本类公理所给出的 $\langle T_a,1\rangle_{0,\beta}=0$，是
\begin{equation}\label{eq:j-function-asymptotic-normalization}
J=e^{(t_0+\delta)/\hbar}\bigl(1+o(\hbar^{-1})\bigr),
\end{equation}
如 \SourceCite{Givental4} 所指出的。这个事实对\hyperref[sec:givental-mirror-approach]{第~\ref*{sec:givental-mirror-approach} 节}中定理~\ref{thm:givental-projective-mirror} 的陈述至关重要。

下面给出两个例题，既给出 $J$ 的公式，也说明量子微分方程如何导出量子上同调中的关系。第~\ref{subsec:cy-threefold-j-function} 小节还将给出更多例题。
\SourceExampleNumber{10.3.1}{1}
\begin{example}\label{ex:projective-line-j-function}
计算 $X=\mathbb P^1$ 的 $J$。若 $H$ 是超平面截面类（这里就是点类），则 $\delta=t_1H$；对于 $\beta=d\in H_2(\mathbb P^1,\mathbb Z)=\mathbb Z$，有 $\int_\beta\delta=dt_1$。由引理~\ref{lem:giventals-j-two-formulas}，
\[
J=e^{(t_0+t_1H)/\hbar}\left(1+\sum_{d=1}^{\infty}q^d
\left(\left\langle\frac H{\hbar-c},1\right\rangle_{0,d}1
+\left\langle\frac1{\hbar-c},1\right\rangle_{0,d}H\right)\right).
\]
利用式~\eqref{eq:projective-line-descendant-factorial} 和~\eqref{eq:projective-line-descendant-harmonic}，它成为
\[
\begin{aligned}
J=e^{(t_0+t_1H)/\hbar}\biggl(1+\sum_{d=1}^{\infty}q^d\biggl(&
\frac{\hbar^{-2d}}{(d!)^2}1\\
&-\frac{2\hbar^{-(2d+1)}}{(d!)^2}
\left(1+\frac12+\cdots+\frac1d\right)H\biggr)\biggr)\\
=e^{(t_0+t_1H)/\hbar}\biggl(1+\sum_{d=1}^{\infty}
&\left(\frac q{\hbar^2}\right)^d\frac1{(d!)^2}
\left(1-2\frac H\hbar\left(1+\frac12+\cdots+\frac1d\right)\right)\biggr).
\end{aligned}
\]
直接验证可知，$(\hbar\,d/dt_1)^2-e^{t_1}$ 零化 $J$，因而是量子微分算子。于是，定理~\ref{thm:quantum-differential-relations} 在代换 $\hbar\,d/dt_1\mapsto H$、$e^{t_1}\mapsto q$ 后，给出小量子上同调关系 $H^2=q$，与\hyperref[subsec:a-model-correlation]{第~\ref*{subsec:a-model-correlation} 小节}的例题~\ref{ex:projective-quantum-ring} 一致。

还可如下得到一个很有趣的 $J$ 公式。由于 $H^2=0$，上式可写成
\[
J=e^{(t_0+t_1H)/\hbar}\sum_{d=0}^{\infty}q^d
\frac1{\left(d!\hbar^{d-1}\left(\sum_{j=1}^d1/j\right)H+d!\hbar^d\right)^2},
\]
再利用一次 $H^2=0$，得到
\[
J=e^{(t_0+t_1H)/\hbar}\sum_{d=0}^{\infty}q^d
\frac1{\bigl((H+\hbar)(H+2\hbar)\cdots(H+d\hbar)\bigr)^2}.
\]
\hyperref[subsec:givental-nef-projective-mirror]{第~\ref*{subsec:givental-nef-projective-mirror} 小节}将说明，这是 Givental 版本的镜像定理的一个特例。
\end{example}

\SourceExampleNumber{10.3.1}{2}
\begin{example}\label{ex:second-order-quantum-relations}
考虑 $n$ 维射影空间 $X=\mathbb P^n$，令 $H$ 为 $\mathbb P^n$ 的超平面类。取上同调基 $1,H,\ldots,H^n$，其中变量 $t_0,t_1$ 对应于类 $1,H$。令 $q=e^{t_1}$，则由例题~\ref{ex:projective-quantum-ring}，小量子乘积满足 $H^{n+1}=e^{t_1}$。

证明这个关系导出一个量子微分方程。由于当 $i\le n$ 时，量子幂 $H^i$ 不依赖任何 $t_0,t_1$，推导式~\eqref{eq:dual-giventals-classical-symbol} 的分析在此给出更精确的恒等式
\[
\left(\widetilde\nabla^g_{\frac d{dt_1}}\right)^{n+1}1=H^{n+1}.
\]
\SourcePage{348}
又因小量子上同调中 $H^{n+1}=e^{t_1}$，此式成为
\[
\left(\widetilde\nabla^g_{\frac d{dt_1}}\right)^{n+1}1=e^{t_1}.
\]
由式~\eqref{eq:j-annihilator-equivalences}，得到
\begin{equation}\label{eq:projective-space-quantum-operator}
\left(\hbar\frac d{dt_1}\right)^{n+1}-e^{t_1}
\end{equation}
是量子微分算子。

也可反过来推导。命题~\ref{prop:projective-space-i-equals-j} 断言，$\mathbb P^n$ 的 Givental 函数 $J$ 由下式给出：
\begin{equation}\label{eq:projective-space-j-hypergeometric}\index{Givental $J$ 函数!射影空间|hyperpage}
J=e^{(t_0+t_1H)/\hbar}\sum_{d=0}^{\infty}e^{dt_1}
\frac1{\bigl((H+\hbar)(H+2\hbar)\cdots(H+d\hbar)\bigr)^{n+1}}.
\end{equation}
目前先注意这个公式的两点。第一，当 $n=1$ 时，它与例题~\ref{ex:projective-line-j-function} 的计算一致（令 $q=e^{t_1}$ 即可）；第二，对于一般的 $n$，可证明式~\eqref{eq:projective-space-j-hypergeometric} 右端被式~\eqref{eq:projective-space-quantum-operator} 的微分算子消去。因此，\hyperref[chap:mirror-theorem]{第十一章}证明式~\eqref{eq:projective-space-j-hypergeometric} 后，即可推出式~\eqref{eq:projective-space-quantum-operator} 是量子微分算子；再由定理~\ref{thm:quantum-differential-relations}，得到 $\mathbb P^n$ 小量子上同调中的关系 $H^{n+1}=e^{t_1}=q$。
\end{example}

\subsection{Calabi--Yau 三维簇}\label{subsec:cy-threefold-j-function}
本节余下部分均假设 $V$ 是 Calabi--Yau 三维簇。先解释如何用\hyperref[subsec:instanton-numbers]{第~\ref*{subsec:instanton-numbers} 小节}所讨论的 Gromov--Witten 不变量
\[
N_\beta=\langle I_{0,0,\beta}\rangle
\]
表示 Givental $J$ 函数。如例题~\ref{ex:threefold-flat-section-comparison} 一样，限制到偶上同调，将 $H^{\mathrm{even}}(V)$ 的上同调基排列成 $T_0=1$，$T_1,\ldots,T_r$ 生成 $H^2(V)$，而 $T^a\in H^4(V)$（$1\le a\le r$）与 $T^0\in H^6(V)$ 为对偶类。还令 $\delta=\sum_{i=1}^rt_iT_i$。

现在给出 $J$ 的公式。引理~\ref{lem:giventals-j-two-formulas} 给出
\[
J=e^{(t_0+\delta)/\hbar}\left(1+\sum_{\beta\ne0}q^\beta
\left(\hbar^{-2}\sum_{a=1}^r\langle\tau_1T_a,1\rangle_{0,\beta}T^a
+\hbar^{-3}\langle\tau_2,1\rangle_{0,\beta}T^0\right)\right).
\]
由次数公理与基本类公理容易得到此式。在例题~\ref{ex:cy-threefold-gravitational} 中，已证明
\[
\langle\tau_1T_a,1\rangle_{0,\beta}=\left(\int_\beta T_a\right)N_\beta,
\qquad\langle\tau_2,1\rangle_{0,\beta}=-2N_\beta,
\]
故上述公式化简为
\begin{equation}\label{eq:cy-threefold-j-invariants}
J=e^{(t_0+\delta)/\hbar}\left(1+\sum_{\beta\ne0}q^\beta
\left(\hbar^{-2}N_\beta\sum_{a=1}^r\left(\int_\beta T_a\right)T^a
-2\hbar^{-3}N_\beta T^0\right)\right).
\end{equation}
不过，式~\eqref{eq:poincare-dual-basis-reconstruction} 容易推出，$\beta\in H_2(V,\mathbb C)$ 的 Poincaré 对偶为 $\sum_{a=1}^r(\int_\beta T_a)T^a\in H^4(V,\mathbb C)$。略微滥用记号，仍用 $\beta$ 表示这个 $H^4(V,\mathbb C)$ 中的类，就可将 Calabi--Yau 三维簇的 Givental 函数公式改写为
\begin{equation}\label{eq:cy-threefold-j-poincare-dual}
J=e^{(t_0+\delta)/\hbar}\left(1+\hbar^{-2}\sum_{\beta\ne0}N_\beta q^\beta\beta
-2\hbar^{-3}\sum_{\beta\ne0}N_\beta q^\beta\,pt\right).
\end{equation}
\SourcePage{349}
若有独立的方法计算 $J$，则与上述公式比较系数，便可计算 $N_\beta$。Givental 在 \SourceCite{Givental4} 中断言，环面完全交具有这样的计算方法。\hyperref[subsec:givental-nef-projective-mirror]{第~\ref*{subsec:givental-nef-projective-mirror} 小节}的例题~\ref{ex:givental-quintic-mirror-potential} 将以五次三维簇说明这种方法。

下面解释 Givental $J$ 函数与 Gromov--Witten 势及小量子乘积的关系。结果如下。
\begin{proposition}\label{prop:cy-threefold-j-potential}\index{Givental $J$ 函数!Calabi--Yau 三维簇|hyperpage}
若 $\Phi$ 是 Calabi--Yau 三维簇 $V$ 的 Gromov--Witten 势，则 $V$ 的 Givental $J$ 函数为
\[
\begin{aligned}
J=e^{t_0/\hbar}\biggl(1+\hbar^{-1}\sum_{a=1}^rt_aT_a
&+\hbar^{-2}\sum_{a=1}^r\frac{\partial\Phi}{\partial t_a}T^a\\
&+\hbar^{-3}\left(\sum_{a=1}^rt_a\frac{\partial\Phi}{\partial t_a}-2\Phi\right)T^0\biggr).
\end{aligned}
\]
此外，对任意 $1\le i,j\le r$，有
\[
\frac{\partial^2J}{\partial t_i\partial t_j}
=\hbar^{-2}e^{(t_0+\delta)/\hbar}T_i*T_j.
\]
\end{proposition}
\begin{proof}
回顾式~\eqref{eq:cy-total-gw-potential}，Gromov--Witten 势为
\[
\Phi=\frac16\int_V\delta^3+\sum_{\beta\ne0}N_\beta q^\beta
=\frac16\int_V\delta^3+\Phi_{\mathrm{hol}},\qquad\delta=\sum_{a=1}^rt_aT_a.
\]
利用此式及式~\eqref{eq:cy-threefold-j-invariants}，容易得到
\[
J=e^{(t_0+\delta)/\hbar}\left(1+\hbar^{-2}\sum_{a=1}^r
\frac{\partial\Phi_{\mathrm{hol}}}{\partial t_a}T^a-2\hbar^{-3}\Phi_{\mathrm{hol}}T^0\right).
\]
将乘积展开，得到
\[
\begin{aligned}
J=e^{t_0/\hbar}\biggl(&1+\hbar^{-1}\delta
+\frac{\hbar^{-2}}2\delta^2+\frac{\hbar^{-3}}6\delta^3\\
&+\hbar^{-2}\sum_{a=1}^r\frac{\partial\Phi_{\mathrm{hol}}}{\partial t_a}T^a
+\hbar^{-3}\delta\sum_{a=1}^r\frac{\partial\Phi_{\mathrm{hol}}}{\partial t_a}T^a
-2\hbar^{-3}\Phi_{\mathrm{hol}}T^0\biggr).
\end{aligned}
\]
由此出发，利用以下恒等式即可得到所需的 $J$ 公式，细节留给读者：
\begin{equation}\label{eq:cy-threefold-potential-euler-identities}
\begin{aligned}
\delta\sum_{a=1}^r\frac{\partial\Phi_{\mathrm{hol}}}{\partial t_a}T^a
&=\left(\sum_{a=1}^rt_a\frac{\partial\Phi_{\mathrm{hol}}}{\partial t_a}\right)T^0,\\
\frac16\delta^3&=\frac16\left(\int_V\delta^3\right)T^0
=\frac13\sum_{a=1}^rt_a\frac\partial{\partial t_a}
\left(\frac16\int_V\delta^3\right)T^0,\\
\frac12\delta^2&=\frac12\sum_{a=1}^r\left(\int_V\delta^2\cup T_a\right)T^a
=\sum_{a=1}^r\frac\partial{\partial t_a}
\left(\frac16\int_V\delta^3\right)T^a.
\end{aligned}
\end{equation}
注意，第二行使用 Euler 公式，第三行使用式~\eqref{eq:poincare-dual-basis-reconstruction}。

为证明 $\partial^2J/\partial t_i\partial t_j$ 的公式，将命题陈述中的 $J$ 公式对 $t_i,t_j$ 求导，容易得到
\[
\begin{aligned}
\frac{\partial^2J}{\partial t_i\partial t_j}=e^{t_0/\hbar}\biggl(&
\hbar^{-2}\sum_{a=1}^r\frac{\partial^3\Phi}{\partial t_i\partial t_j\partial t_a}T^a\\
&+\hbar^{-3}\left(\sum_{a=1}^rt_a\frac{\partial^3\Phi}{\partial t_i\partial t_j\partial t_a}\right)T^0\biggr).
\end{aligned}
\]
\SourcePage{350}
利用类似于式~\eqref{eq:cy-threefold-potential-euler-identities} 第一行的恒等式，得到
\[
\begin{aligned}
\frac{\partial^2J}{\partial t_i\partial t_j}
&=\hbar^{-2}e^{(t_0+\delta)/\hbar}\sum_a
\frac{\partial^3\Phi}{\partial t_i\partial t_j\partial t_a}T^a\\
&=\hbar^{-2}e^{(t_0+\delta)/\hbar}T_i*T_j.
\end{aligned}
\]
第二个等号来自命题~\ref{prop:giventals-flat-sections} 证明后的讨论，那里解释了：小量子乘积就是大量子乘积在 $M=H^0(V)\oplus H^2(V)$ 上的限制。
\end{proof}

应用此命题，计算一个有趣的量子微分算子。
\SourceExampleNumber{10.3.2}{1}
\begin{example}\label{ex:cy-quantum-equation-relation}
设 Calabi--Yau 三维簇 $V$ 满足 $h^2(V)=1$。令 $H$ 为 $H^2(V)$ 的生成元，$C\in H^4(V)$ 为对偶生成元，用 $pt\in H^6(V,\mathbb Z)$ 表示点类，故 $H\cup C=pt$。如往常一样，变量 $t_0,t_1$ 对应于 $1,H$，并令 $q=e^{t_1}$。

\hyperref[subsec:a-model-hodge-filtration]{第~\ref*{subsec:a-model-hodge-filtration} 小节}研究了三点函数 $Y(q)=\langle H,H,H\rangle$，并在式~\eqref{eq:one-parameter-quantum-yukawa} 中指出关系 $H*H=Y(q)C$。利用命题~\ref{prop:cy-threefold-j-potential}，得到
\[
\left(\hbar\frac d{dt_1}\right)^2J
=e^{(t_0+t_1H)/\hbar}H*H=e^{(t_0+t_1H)/\hbar}Y(q)C.
\]
容易由此推出
\[
\begin{aligned}
\left(\hbar\frac d{dt_1}\right)^2
\left(\frac{(\hbar\,d/dt_1)^2J}{Y(q)}\right)
&=\left(\hbar\frac d{dt_1}\right)^2e^{(t_0+t_1H)/\hbar}C\\
&=e^{(t_0+t_1H)/\hbar}H\cup H\cup C=0,
\end{aligned}
\]
因为 $H^8(V)=0$。因此
\begin{equation}\label{eq:cy-threefold-quantum-operator}
\left(\hbar\frac d{dt_1}\right)^2\circ\frac1{Y(q)}\circ
\left(\hbar\frac d{dt_1}\right)^2
\end{equation}
是量子微分算子。由定理~\ref{thm:quantum-differential-relations}，它在 $V$ 的\IndexLocation{indexloc:1139:350}小量子上同调环中给出关系 $H*H*(H*H/Y)=0$。

由于 $H^8(V)=0$，这个关系本身相当平凡。不过，微分方程~\eqref{eq:cy-threefold-quantum-operator} 仍很有趣，主要是因为它与\hyperref[sec:yukawa-couplings]{第~\ref*{sec:yukawa-couplings} 节}为 Calabi--Yau 三维簇的 Yukawa 耦合所推导的 Picard--Fuchs 方程 \eqref{eq:normalized-yukawa-pf} 极其相似。这并非巧合，而是镜像对称的一个关键部分。下面的定理~\ref{thm:picard-fuchs-quantum-equivalence} 将进一步研究 Picard--Fuchs 方程与量子微分方程的关系。

还可用这个例题说明\hyperref[subsec:giventals-j-function]{第~\ref*{subsec:giventals-j-function} 小节}的评论：定理~\ref{thm:quantum-differential-relations} 的逆命题可能不成立，即量子上同调中的关系并非都给出量子微分方程。对于这里的 Calabi--Yau 三维簇 $V$，量子上同调中确有关系 $H*H*H*H=0$，
\SourcePage{351}
但 $(\hbar\,d/dt_1)^4J\ne0$（当 $Y$ 非常数时）。这是因为，用式~\eqref{eq:dual-giventals-pairing-derivative} 作简单计算，得到
\[
\left(\hbar\frac d{dt_1}\right)^4\langle s_j,1\rangle
=\hbar^2\frac{d^2Y}{dt_1^2}\langle s_j,C\rangle
+2\hbar\frac{dY}{dt_1}\langle s_j,pt\rangle.
\]
由于 $C$ 和 $pt$ 是线性无关的上同调类，且 $Y$ 非常数，容易推出 $(\hbar\,d/dt_1)^4\langle s_j,1\rangle$ 不可能对所有 $j$ 都为零。
\end{example}

本节最后解释，定理~\ref{thm:quantum-differential-relations} 与例题~\ref{ex:threefold-flat-section-comparison} 中讨论的 A 模型联络 $\nabla$ 有何关系。回顾，自然变量为 $u_j$，它与 $t_j$ 的关系是 $t_j=2\pi iu_j$。考虑算子
\[
P_\nabla=P\left(\nabla_{\frac\partial{\partial u_1}},\ldots,
\nabla_{\frac\partial{\partial u_r}},q_1,\ldots,q_r\right),
\qquad q_j=e^{2\pi iu_j}.
\]
假定 $P_\nabla$ 是联络项的多项式，其系数为 $q_j$ 的形式幂级数，并假定所有 $q_j$ 都放在联络项的左边。

\begingroup\linespread{1.18}\selectfont
算子 $P_\nabla$ 可作用于定义 $\nabla$ 的丛 $\mathcal H$ 的任意截面。特别地，有由 $1=T_0$ 给出的截面。若 $P_\nabla1=0$，则称 $P_\nabla$ 为 \emph{Picard--Fuchs 算子}。

这个术语源于镜像对称，因为在\hyperref[sec:mirror-symmetry-revisited]{第~\ref*{sec:mirror-symmetry-revisited} 节}讨论的镜像猜想中，截面 $1$ 对应于镜像上的归一化三形式 $\widetilde\Omega$，而通常的 Picard--Fuchs 算子恰好是消去 $\widetilde\Omega$ 的算子。

给定 $P_\nabla$，将它写成
\begin{equation}\label{eq:a-model-picard-fuchs-operator}
P_\nabla=\sum_\alpha A_\alpha(q)\nabla^\alpha,
\end{equation}
其中 $\alpha=(a_1,\ldots,a_r)$ 为多重指标，
\[
\nabla^\alpha=\nabla_{\frac\partial{\partial u_1}}^{a_1}\cdots
\nabla_{\frac\partial{\partial u_r}}^{a_r}.
\]
如往常一样，令 $|\alpha|=\sum_i a_i$，并以 $m=\max_\alpha\{|\alpha|:A_\alpha(q)\ne0\}$ 表示 $P_\nabla$ 的阶数。定义 $P_\nabla$ 的\emph{$\hbar$ 齐次化}为微分算子
\begin{equation}\label{eq:picard-fuchs-homogenization}
P(\hbar\partial/\partial t,e^t,\hbar)
=\sum_\alpha\left(\frac\hbar{2\pi i}\right)^{m-|\alpha|}
A_\alpha(e^t)\left(\hbar\frac\partial{\partial t}\right)^\alpha,
\end{equation}
其中 $(\hbar\partial/\partial t)^\alpha$ 的含义显而易见。（下面会看到为何需要 $2\pi i$。）还令
\[
P_m(T,q)=\sum_{|\alpha|=m}A_\alpha(q)T^\alpha,
\]
其中 $T^\alpha$ 是用 $T_1,\ldots,T_r$ 在 $H^*(V)$ 的小量子乘积下组成的乘积。可将 $P_m(T,q)$ 看作算子 $P_\nabla$ 的主部（特征形式）的量子版本。

这些对象的关系如下。
\begin{theorem}\label{thm:picard-fuchs-quantum-equivalence}\index{量子微分方程!Picard--Fuchs 算子|hyperpage}
设 $V$ 为 Calabi--Yau 三维簇，$P_\nabla$ 为式~\eqref{eq:a-model-picard-fuchs-operator} 中的 $m$ 阶算子。若 $P(\hbar\partial/\partial t,e^t,\hbar)$ 由式~\eqref{eq:picard-fuchs-homogenization} 定义，则：
\SourcePage{352}
\begin{enumerate}[label=(\roman*)]
\item $P_\nabla$ 是 Picard--Fuchs 算子，当且仅当 $P(\hbar\partial/\partial t,e^t,\hbar)$ 是量子微分算子。换言之，
\[
P_\nabla1=0\quad\Longleftrightarrow\quad
P(\hbar\partial/\partial t,e^t,\hbar)J=0.
\]
\item 若 $P_\nabla$ 是 Picard--Fuchs 算子，则在小量子上同调中 $P_m(T,q)=0$。
\end{enumerate}
\end{theorem}
\begin{proof}
先指出，此时量子微分方程具有一个特殊性质。令 $J$ 如式~\eqref{eq:cy-threefold-j-poincare-dual}，将 $\hbar$ 视为变量；以 $J_0$ 表示把其中的 $\hbar$ 替换为非零复数 $\hbar_0$ 后所得的表达式。断言
\begin{equation}\label{eq:quantum-operator-hbar-specialization}
P(\hbar\partial/\partial t,e^t,\hbar)J=0
\quad\Longleftrightarrow\quad
P(\hbar_0\partial/\partial t,e^t,\hbar_0)J_0=0.
\end{equation}
\looseness=-1 为说明这一点，先注意 $P(\hbar\partial/\partial t,e^t,\hbar)$ 关于 $\hbar$ 齐次，因此它等于 $\hbar^m\widetilde P$，其中算子 $\widetilde P$ 不含 $\hbar$。此外，$\hbar$ 在 $J$ 的系数中也是齐次出现的。将式~\eqref{eq:cy-threefold-j-invariants} 展开，并暂时除去不受这些 $t_1,\ldots,t_r$ 微分算子影响的公因子 $e^{t_0/\hbar}$，可写成
\[
J=1+\hbar^{-1}\sum_{j=1}^rA_jT_j
+\hbar^{-2}\sum_{j=1}^rB_jT^j+\hbar^{-3}CT^0,
\]
其中 $A_j,B_j,C$ 为 $q_1,\ldots,q_r$ 的函数。由 $P=\hbar^m\widetilde P$，$PJ=0$ 等价于
\[
\widetilde P1=\widetilde PA_j=\widetilde PB_j=\widetilde PC=0
\qquad\text{对所有 }j.
\]
由于 $\hbar_0\ne0$，这也等价于 $P(\hbar_0\partial/\partial t,e^t,\hbar_0)J_0=0$，从而得到式~\eqref{eq:quantum-operator-hbar-specialization}。

\looseness=-1 现在证明定理。由式~\eqref{eq:quantum-operator-hbar-specialization}，可将 $\hbar$ 特化为任意非零常数。故在证明余下部分中，不失一般性，取 $\hbar=2\pi i$。马上就会看到这是正确的选择。

由式~\eqref{eq:j-annihilator-equivalences}，已知
\begin{equation}\label{eq:quantum-operator-dual-annihilation}
P(\hbar\partial/\partial t,e^t,\hbar)J=0
\quad\Longleftrightarrow\quad P_{\widetilde\nabla^g}1=0,
\end{equation}
其中 $\widetilde\nabla^g$ 为对偶 Givental 联络，$P_{\widetilde\nabla^g}$ 由代换 $\hbar\partial/\partial t_j\mapsto\widetilde\nabla^g_{\partial/\partial t_j}$ 定义。

关键观察是：由于 $\hbar=2\pi i$，对偶 Givental 联络 $\widetilde\nabla^g$ 恰好是 $2\pi i\nabla$。因此
\[
\widetilde\nabla^g_{\frac\partial{\partial t_j}}
=2\pi i\nabla_{\frac\partial{\partial t_j}}
=\nabla_{\frac\partial{\partial u_j}},
\]
第二个等号使用了 $t_j=2\pi iu_j$。由于 $\hbar=2\pi i$，比较式~\eqref{eq:a-model-picard-fuchs-operator} 和~\eqref{eq:picard-fuchs-homogenization}，可见 $P_{\widetilde\nabla^g}$ 恰好为 $P_\nabla$。于是式~\eqref{eq:quantum-operator-dual-annihilation} 给出
\[
P(\hbar\partial/\partial t,e^t,\hbar)J=0\quad\Longleftrightarrow\quad P_\nabla1=0.
\]
这证明了第一部分，第二部分立即由定理~\ref{thm:quantum-differential-relations} 得到。
\end{proof}

\enlargethispage{2\baselineskip}
\looseness=-1 这个定理表明，Calabi--Yau 三维簇的 A 模型联络与小量子上同调有着深刻联系。不过，镜像对称所预言的、它与镜像的 B 模型之间的联系更深。下一章将证明镜像定理的几个特殊情形，最终将一直以来所研究的这些美妙数学联系起来。
\par\endgroup

% 第七批：原 PDF 物理页354起。
\chapter{镜像定理}\label{chap:mirror-theorem}\index{镜像定理|hyperpage}
\SourcePage{354}
本章旨在证明镜像定理的几个版本，特别着重于五次三维簇。还将涉及与证明技巧有关的一些有趣主题。如第一、二章所指出的，证明需要几个此前尚未讨论的新概念。

\hyperref[chap:quantum-cohomology]{第八章}说明，“镜像定理”有多种表述方式。\hyperref[sec:mirror-symmetry-revisited]{第~\ref*{sec:mirror-symmetry-revisited} 节}的版本用 Hodge 结构变分陈述。这里采取另一种观点，介绍两种密切相关的证明方法，两者都以稳定映射空间上的等变交理论为基础。历史上较早的方法由 Givental 提出 \SourceCite{Givental2}，其特点是使用\hyperref[chap:quantum-differential-equations]{第十章}定义的引力关联数。这些关联数导出\hyperref[subsec:giventals-j-function]{第~\ref*{subsec:giventals-j-function} 小节}定义的 Givental $J$ 函数；我们将看到，$J$ 在 Givental 版本的镜像定理中起重要作用。Givental 的方法可推广到环面完全交 \SourceCite{Givental4}、齐性空间 \SourceCite{Kim2}，并可部分推广到高亏格 \SourceCite{Givental5}。Lian、Liu 和 Yau 的方法 \SourceCite{LLY} 则着重于线性与非线性 $\sigma$ 模型之间的相互作用。在数学上，非线性 $\sigma$ 模型用稳定映射模空间的交理论理解，而线性 $\sigma$ 模型用射影空间 $N_d$ 上的交理论理解；$N_d$ 参数化 $\mathbb P^1$ 上次数 $d$ 的齐次式元组。主要结果用 $N_d$ 的等变交理论表述。这种方法也适用于某些非紧簇。我们先介绍 \SourceCite{LLY} 的方法，因为展开它所需的技术细节较少。

\hyperref[sec:quintic-mirror-theorem]{第~\ref*{sec:quintic-mirror-theorem} 节}概述 \SourceCite{LLY} 对五次三维簇镜像定理的证明，定理在这里陈述为定理~\ref{thm:lly-quintic-mirror}。由\hyperref[subsec:quintic-mirror-symmetry]{第~\ref*{subsec:quintic-mirror-symmetry} 小节}，它蕴含\hyperref[chap:quantum-cohomology]{第八章}用 Hodge 理论表述的镜像对称。在 \SourceCite{LLY} 的方法中，等变版本的 Gromov--Witten 不变量被编码为射影空间 $N_d$ 上的一列等变上同调类 $\widehat Q_d$。我们将证明，它们与等变上同调类 $\widehat P_d$ 有关；后者联系于给出镜像族周期的超几何函数。证明依靠直接比较 $\widehat P_d$ 与 $\widehat Q_d$。为此使用Euler 数据（定义~\ref{def:omega-euler-data}）及其关联关系（定义~\ref{def:linked-euler-data}）这一概念，它本身以乘性上同调类的粘合引理（引理~\ref{lem:euler-data-gluing}）为基础；还要使用镜像变换（定义~\ref{def:euler-data-mirror-transformation}）。也将解释此方法如何适用于\hyperref[subsec:multiple-cover-localization]{第~\ref*{subsec:multiple-cover-localization} 小节}中的多重覆盖计算。

\hyperref[sec:givental-mirror-approach]{第~\ref*{sec:givental-mirror-approach} 节}转向 Givental 的镜像定理方法。这里着重讨论完全交 $X\subset\mathbb P^n$，将其表示为 $\mathcal V=\bigoplus_{i=1}^{\ell}\mathcal O_{\mathbb P^n}(a_i)$（$a_i>0$）的一个截面的零点集。在此情形下，定义两个取形式上同调值的函数 $J_{\mathcal V}$ 和 $I_{\mathcal V}$。函数 $J_{\mathcal V}$ 与该完全交的 Givental $J$ 函数密切相关，而 $I_{\mathcal V}$
\SourcePage{355}
就是\hyperref[subsec:nef-gkz]{第~\ref*{subsec:nef-gkz} 小节}研究的函数。在\hyperref[subsec:givental-nef-projective-mirror]{第~\ref*{subsec:givental-nef-projective-mirror} 小节}，陈述定理~\ref{thm:givental-projective-mirror}，说明当 $\sum_{i=1}^{\ell}a_i\le n+1$ 时，如何用 $I_{\mathcal V}$ 计算 $J_{\mathcal V}$。对于 $\mathcal V=\mathcal O_{\mathbb P^4}(5)$，此镜像定理蕴含五次三维簇的镜像对称。更一般地，在\hyperref[subsec:quantum-hyperplane-section-principle]{第~\ref*{subsec:quantum-hyperplane-section-principle} 小节}考虑完全交 $X\subset Y$，它由 $Y$ 上一个向量丛 $\mathcal V$ 的截面消失所定义，其中 $\mathcal V$ 是线丛的直和。这里将给出 $J_{\mathcal V}$ 的一般定义，并讨论猜想上可以计算 $J_{\mathcal V}$ 的\emph{量子超平面截面原理}。然后证明定理~\ref{thm:givental-projective-mirror}：先在\hyperref[subsec:givental-projective-space-proof]{第~\ref*{subsec:givental-projective-space-proof} 小节}证明 $\mathcal V=0$ 的情形，再在\hyperref[subsec:givental-general-bundle-proof]{第~\ref*{subsec:givental-general-bundle-proof} 小节}证明一般 $\mathcal V$ 的情形。证明使用等变上同调中的局部化，并将它应用于 $J_{\mathcal V}$、$I_{\mathcal V}$ 的等变版本。当 $\mathcal V=0$ 时，\hyperref[chap:localization]{第九章}的局部化公式（推论~\ref{cor:orbifold-localization-integration}）导出 $J_{\mathcal V}$ 局部化的递推关系，将其显式求和便得到 $J_{\mathcal V}$。对于一般 $\mathcal V$，所得关系尚不是递推关系，因此需要更多工作。本章最后讨论 Givental 版本的环面完全交镜像定理。

\section{五次三维簇的镜像定理}\label{sec:quintic-mirror-theorem}
本节遵循 \SourceCite{LLY}，概述五次三维簇镜像定理的证明。
\subsection{定理的陈述}\label{subsec:quintic-mirror-theorem-statement}
设 $V\subset\mathbb P^4$ 为一般的五次三维簇。先回顾 Gromov--Witten 不变量 $\langle I_{0,0,d}\rangle$。固定正整数 $d$，考虑映射
\[
\begin{tikzcd}[column sep=large,row sep=large]
\overline M_{0,1}(\mathbb P^4,d)\arrow[r,"e_1"]\arrow[d,"\pi_1"']&\mathbb P^4\\
\overline M_{0,0}(\mathbb P^4,d)&
\end{tikzcd}
\]
其中 $\pi_1$ 忘掉标记点，$e_1$ 为通常的\IndexLocation{indexloc:0276:355}赋值映射。令 $\mathcal V=\mathcal O_{\mathbb P^4}(5)$，并令 $\mathcal V_d=\pi_{1*}e_1^*(\mathcal V)$，它是 $\overline M_{0,0}(\mathbb P^4,d)$ 上的向量丛。容易验证，它与例题~\ref{ex:quintic-virtual-class} 中定义的丛一致。于是上述 Gromov--Witten 不变量为
\[
\langle I_{0,0,d}\rangle=\int_{\overline M_{0,0}(\mathbb P^4,d)}c_{5d+1}(\mathcal V_d).
\]
\hyperref[chap:gromov-witten]{第七章}用 $N_d$ 表示它，但这里改用 $K_d$（与 \SourceCite{LLY} 一致），因为下面的 $N_d$ 将有另一含义。此外，$c_{5d+1}(\mathcal V_d)$ 为 $\mathcal V_d$ 的 Euler 类，记为 $\operatorname{Euler}(\mathcal V_d)$。用这一记号，上式写成
\[
K_d=\int_{\overline M_{0,0}(\mathbb P^4,d)}\operatorname{Euler}(\mathcal V_d).
\]
注意，这里将积分写在模空间 $\overline M_{0,0}(\mathbb P^4,d)$ 上，而不是代数叠 $\overline{\mathcal M}_{0,0}(\mathbb P^4,d)$ 上。这是因为，虽然 $\mathcal V_d$ 仅是代数叠上的向量丛，但它的 Euler 类在粗模空间上给出定义良好的有理零维闭链，因而上述积分完全有意义。以下常滥用记号，将 $\mathcal V_d$ 当作 $\overline M_{0,0}(\mathbb P^4,d)$ 上的向量丛。

\SourcePage{356}
由\hyperref[chap:gromov-witten]{第七章}还知道，$K_d=\sum_{k\mid d}n_{d/k}k^{-3}$，其中 $n_d$ 是\hyperref[subsec:instanton-numbers]{第~\ref*{subsec:instanton-numbers} 小节}讨论的瞬子数。此外，由\hyperref[subsec:quintic-gw-potential]{第~\ref*{subsec:quintic-gw-potential} 小节}，五次三维簇的 Gromov--Witten 势为
\[
\Phi(T)=\frac56T^3+\sum_{d>0}K_de^{dT},
\]
所以 $\Phi'''$ 为 A 模型关联函数。\hyperref[sec:quintic-mirror-theorem]{第~\ref*{sec:quintic-mirror-theorem} 节}的目标是证明：按照镜像对称的规定，$\Phi$ 由五次镜像族的 Hodge 理论决定，从而证明\hyperref[chap:quintic]{第二章}中给出的 $n_d$ 公式。

为陈述 \SourceCite{LLY} 所证明的结果，考虑取 $H^*(\mathbb P^4,\mathbb C)$ 中形式值的表达式
\[
e^{-Ht/\hbar}\sum_{d\ge0}
\frac{\prod_{m=0}^{5d}(5H-m\hbar)}{\prod_{m=1}^d(H-m\hbar)^5}e^{dt},
\]
其中 $H$ 为超平面类。按 $H$ 的幂展开，得到
\[
5H\left(y_0-y_1\frac H\hbar+y_2\frac{H^2}{\hbar^2}-y_3\frac{H^3}{\hbar^3}\right),
\]
其中 $y_i=y_i(t)$ 是一些函数。在\hyperref[chap:moduli]{第六章}中已经看到，$y_i$ 是五次镜像周期所满足的超几何微分方程
\[
\left[\left(\frac d{dt}\right)^4-5e^t
\left(5\frac d{dt}+1\right)\left(5\frac d{dt}+2\right)
\left(5\frac d{dt}+3\right)\left(5\frac d{dt}+4\right)\right]y=0
\]
的一组解基，用 Frobenius 方法即可验证。例题~\ref{ex:quintic-hypergeometric-mirror-map} 的后续讨论解释了这一点。注意，$y_i$ 与\hyperref[subsec:nef-gkz]{第~\ref*{subsec:nef-gkz} 小节}的函数 $I$ 的系数一致；此函数对应于集合 $\mathcal A=\Delta\cap M$，其中 $\Delta$ 是集合 \eqref{eq:quintic-polar-vertices} 的凸包，该集合由 $\mathbb P^4$ 标准扇的本原整数生成元组成。由命题~\ref{prop:nef-i-gkz}，知道 $y_i$ 满足与 $\mathcal A$ 对应的 GKZ 方程组。

令 $\Psi(t)=y_1/y_0$。\hyperref[subsec:givental-ij]{第~\ref*{subsec:givental-ij} 小节}已指出，镜像映射由 $q=\exp(\Psi(t))$ 定义。于是，\SourceCite{LLY} 的\emph{镜像定理}陈述如下。
\begin{theorem}\label{thm:lly-quintic-mirror}
若 $\Phi$ 为五次三维簇的 Gromov--Witten 势，则
\[
\Phi(\Psi(t))=\frac52\left(\frac{y_1}{y_0}\frac{y_2}{y_0}-\frac{y_3}{y_0}\right).
\]
\end{theorem}
这个五次簇\IndexLocation{indexloc:0788:356}镜像猜想的表述最早见于 \SourceCite{Kontsevich3}。例题~\ref{ex:givental-quintic-mirror-potential} 将说明，定理~\ref{thm:lly-quintic-mirror} 可由 \SourceCite{Givental2} 的结果推出，不过第一个完整而详细的证明见于 \SourceCite{LLY}。\hyperref[subsec:euler-data]{第~\ref*{subsec:euler-data} 小节}解释了，此版本的镜像定理如何导出“经典”镜像对称 \eqref{eq:quintic-mirror-identity}。还应回顾命题~\ref{prop:quintic-mirror-potential}：上述定理等式的右端，是五次镜像的 Hodge 结构变分的“势函数”。如\hyperref[subsec:quintic-mirror-symmetry]{第~\ref*{subsec:quintic-mirror-symmetry} 小节}所指出，定理~\ref{thm:lly-quintic-mirror} 意味着，\IndexLocation{indexloc:0759:356}五次镜像与五次三维簇构成定义~\ref{def:mathematical-mirror-pair} 意义下的数学镜像对。

\SourcePage{357}
本节余下部分用于证明定理~\ref{thm:lly-quintic-mirror}；\hyperref[sec:givental-mirror-approach]{第~\ref*{sec:givental-mirror-approach} 节}将概述此定理的第二种证明。
\subsection{稳定模空间的线性模型}\label{subsec:stable-moduli-linear-models}\index{稳定映射!线性模型|hyperpage}
下面引入稳定映射空间的一个“线性化”版本，它在后续讨论中起关键作用。在物理上，这类似于从非线性 $\sigma$ 模型转到线性 $\sigma$ 模型。先对 $\mathbb P^n$ 中次数 $\ell$ 的超曲面建立理论，再特化到 $\mathbb P^4$ 中的五次三维簇。

先紧化次数 $d$ 的映射 $\mathbb P^1\to\mathbb P^n$ 所组成的空间。通常的模空间 $\overline M_{0,0}(\mathbb P^n,d)$ 并不实现这种紧化，因为它的一般元素只能在相差源空间重新参数化的意义下确定这样的映射。因此改用空间
\[
M_d=\overline M_{0,0}(\mathbb P^1\times\mathbb P^n,(1,d)).
\]
它的一般元素是一个 $\mathbb P^1\to\mathbb P^1\times\mathbb P^n$ 的映射，像的次数为 $(1,d)$。关键在于，重新参数化源空间不会改变其像，而这个像是一个次数 $d$ 的映射 $\mathbb P^1\to\mathbb P^n$ 的图像。因此，$M_d$ 确实紧化了这些映射组成的空间。

第二种紧化如下。一个次数 $d$ 的映射 $\mathbb P^1\to\mathbb P^n$ 可由 $\mathbb P^1$ 上的 $n+1$ 个次数 $d$ 的齐次多项式 $(\alpha_0,\ldots,\alpha_n)$ 表示。这导出射影空间
\[
N_d=\mathbb P\bigl(H^0(\mathbb P^1,\mathcal O_{\mathbb P^1}(d))^{n+1}\bigr),
\]
其中 $\mathbb P$ 表示一维子空间组成的射影空间。\hyperref[sec:gauged-linear-sigma-models]{附录 B.5} 在 $n=4$ 的情形下解释了 $N_d$ 在线性 $\sigma$ 模型中的作用。

$M_d$ 和 $N_d$ 均有自然的环面作用。以下大量使用第 \hyperref[sec:localization-theorem]{\ref*{sec:localization-theorem}}、\hyperref[sec:stable-map-localization]{\ref*{sec:stable-map-localization}} 节，读者继续之前可以先复习。命题~\ref{prop:atiyah-bott-localization} 及其推论、例题~\ref{ex:projective-equivariant-localization} 和\hyperref[subsec:kontsevich-localization]{第~\ref*{subsec:kontsevich-localization} 小节}尤其重要。

令 $G=\mathbb C^*\times T$ 为环面，其中 $T=(\mathbb C^*)^{n+1}$。先描述要使用的 $\mathbb P^1\times\mathbb P^n$ 上的 $G$ 作用。令 $t\in\mathbb C^*$、$(t_0,\ldots,t_n)\in T$，则
\[
\begin{aligned}
&(t,t_0,\ldots,t_n)\cdot((w_0,w_1),(x_0,\ldots,x_n))\\
&\qquad=((t^{-1}w_0,w_1),(t_0^{-1}x_0,\ldots,t_n^{-1}x_n)).
\end{aligned}
\]
如例题~\ref{ex:projective-equivariant-localization} 一样，选择逆元是为了让经由 $G$ 作用拉回所诱导的齐次式上的作用具有正指数。

接着描述 $G$ 在 $M_d=\overline M_{0,0}(\mathbb P^1\times\mathbb P^n,(1,d))$ 上的作用，它由 $\mathbb P^1\times\mathbb P^n$ 上的 $G$ 作用诱导。这与\hyperref[sec:stable-map-localization]{第~\ref*{sec:stable-map-localization} 节}研究的 $(\mathbb C^*)^{r+1}$ 在 $\overline M_{0,n}(\mathbb P^r,d)$ 上的作用密切相关。回顾，我们已仔细描述了固定点轨迹的各分支。这里关心的是 $G$ 在 $M_d$ 上作用的固定点轨迹的分支；可得到类似描述，只需作以下修改：
\begin{itemize}
\item 一个 $G$ 不动稳定映射的收缩分支必须映到 $\mathbb P^1\times\mathbb P^n$ 的某个 $G$ 不动点。
\item 非收缩分支必须覆盖 $\mathbb P^1\times\mathbb P^n$ 中的某条 $G$ 不变曲线。这些曲线形如 $\mathbb P^1\times\{q_i\}$ 和 $\{p\}\times\ell$，其中 $q_i\in\mathbb P^n$ 为 $T$ 不动点，$p$ 为 $0$ 或 $\infty$，$\ell\subset\mathbb P^n$ 为坐标直线。
\item 若 $\sum'$ 表示对覆盖 $\mathbb P^1\times\{q_i\}$ 的分支求和，$\sum''$ 表示对覆盖 $\{p\}\times\ell$ 的分支求和，则 $\sum'd_i=1$、$\sum''d_i=d$，其中 $d_i$ 为各覆盖的次数。
\end{itemize}
\SourcePage{358}
这些 $G$ 固定点轨迹的分支，与引理~\ref{lem:graph-space-descendant-factorization} 的证明中出现的分支类似。

最后描述 $G$ 如何作用于 $N_d$。此处作用为
\[
\begin{aligned}
&(t,t_0,\ldots,t_n)\cdot(\alpha_0(w_0,w_1),\ldots,\alpha_n(w_0,w_1))\\
&\qquad=(t_0^{-1}\alpha_0(t^{-1}w_0,w_1),\ldots,
 t_n^{-1}\alpha_n(t^{-1}w_0,w_1)).
\end{aligned}
\]
这给出射影空间 $N_d$ 上的一个线性化 $G$ 作用。$N_d$ 的 $G$ 不动点为
\[
p_{i,r}=(0,\ldots,0,w_0^rw_1^{d-r},0,\ldots,0),
\]
其中非零单项式位于第 $i$ 个位置。以 $i_{p_{i,r}}$ 表示 $p_{i,r}$ 到 $N_d$ 的包含映射。

如\hyperref[subsec:equivariant-cohomology]{第~\ref*{subsec:equivariant-cohomology} 小节}一样，令 $\lambda_0,\ldots,\lambda_n$ 生成 $H^*(BT)$，并类似地引入 $H^*(B\mathbb C^*)$ 的生成元 $\hbar$。\SourceCite{LLY} 将这个生成元记为 $\alpha$，但这里采用 \SourceCite{Givental2} 的记号，以使全章记号一致（不过 \SourceCite{Givental2} 中 $\hbar$ 的符号不同，对应于 $\mathbb C^*$ 作用的不同约定）。

可用 $\lambda_i$、$\hbar$ 描述等变上同调 $H_G^*(N_d)$，答案类似于例题~\ref{ex:projective-equivariant-localization}。将 $p_{i,r}$ 视为 $H^0(\mathbb P^1,\mathcal O_{\mathbb P^1}(d))^{n+1}$ 中的向量，就使 $G$ 在此空间上的作用对角化。由 $G$ 在 $N_d$ 上的作用，作用于 $p_{i,r}$ 的权重显然为 $-(\lambda_i+r\hbar)$。因此，只要将权重 $\{\lambda_i\}$ 换成 $\{\lambda_i+r\hbar\}$，就得到例题~\ref{ex:projective-equivariant-localization} 中的同样公式。特别地，$(N_d)_G$ 是 $BG$ 上丛 $\bigoplus_{i,r}\mathcal O(-\lambda_i-r\hbar)$ 的射影化。若 $\kappa$ 为 $N_d$ 的等变超平面类，则与式~\eqref{eq:projective-equivariant-relation} 对应的关系为
\begin{equation}\label{eq:linear-model-equivariant-cohomology}
\prod_{i,r}(\kappa-\lambda_i-r\hbar)=0
\end{equation}
（在 $H_G^*(N_d)$ 中）。此外，作为式~\eqref{eq:projective-fixed-dual-class} 的特殊情形，令
\begin{equation}\label{eq:linear-model-fixed-point-class}
\phi_{p_{i,r}}=\prod_{(j,s)\ne(i,r)}\bigl(\kappa-(\lambda_j+s\hbar)\bigr),
\end{equation}
因为 $\lambda_j+s\hbar$ 是 $N_d$ 线性化 $G$ 作用的相应权重参数。回顾，将非零权重取逆后，$\phi_{p_{i,r}}$ 构成 $\mathcal R_GH_G^*(N_d)$ 作为 $\mathcal R_G$ 向量空间的一组基。

注意，每个类 $Q\in H_G^*(N_d)$ 都可由 $\kappa,\hbar,\lambda_i$ 的多项式表示。上述讨论结合式~\eqref{eq:projective-hyperplane-fixed-restriction} 表明，类 $i_{p_{i,r}}^*(Q)$ 就是在 $\kappa$ 中代入 $\lambda_i+r\hbar$ 后所得的 $H^*(BG)=\mathbb C[\hbar,\lambda_i]$ 中的元素。因此，也可合理地以 $Q(\lambda_i+r\hbar)$ 表示 $i_{p_{i,r}}^*(Q)$。由关系~\eqref{eq:linear-model-equivariant-cohomology}，此记号定义良好。当 $Q$ 属于上同调环的某个适当局部化（交换代数意义下）时，也使用类似记号。

作为局部化定理应用于这一情形的简单例子，证明
\begin{equation}\label{eq:linear-model-restriction-integral}
Q(\lambda_i+r\hbar)=i_{p_{i,r}}^*(Q)
=\int_{(N_d)_G}\phi_{p_{i,r}}\cup Q,
\end{equation}
其中 $\int_{(N_d)_G}$ 为推论~\ref{cor:localization-integration} 之前定义的等变积分。为证明此式，先注意，若 $N_{i,r}$ 为 $p_{i,r}$ 的法丛，$\operatorname{Euler}_G(N_{i,r})$ 为其等变 Euler 类，则由式~\eqref{eq:projective-fixed-euler}，
\SourcePage{359}
\begin{equation}\label{eq:linear-model-normal-euler}
\operatorname{Euler}_G(N_{i,r})
=\prod_{(j,s)\ne(i,r)}\bigl(\lambda_i+r\hbar-(\lambda_j+s\hbar)\bigr).
\end{equation}
此外，将式~\eqref{eq:projective-dual-class-restriction} 适配到 $N_d$，容易得到
\begin{equation}\label{eq:linear-model-fixed-class-restriction}
i_{p_{j,s}}^*(\phi_{p_{i,r}})=
\begin{cases}\operatorname{Euler}_G(N_{i,r}),&(j,s)=(i,r),\\0,&\text{其他情形}.
\end{cases}
\end{equation}
现在应用推论~\ref{cor:localization-integration}，得到
\[
\int_{(N_d)_G}\phi_{p_{i,r}}\cup Q
=\frac{i_{p_{i,r}}^*(\phi_{p_{i,r}}\cup Q)}{\operatorname{Euler}_G(N_{i,r})}
=i_{p_{i,r}}^*(Q).
\]

应指出，\hyperref[chap:localization]{第九章}中的局部化定理，使用的是\IndexLocation{indexloc:0247:359}等变上同调在交换代数意义下的局部化。如\hyperref[subsec:atiyah-bott-localization]{第~\ref*{subsec:atiyah-bott-localization} 小节}一样，以 $\mathcal R_T$ 表示 $H^*(BT)\simeq\mathbb C[\lambda_i]$ 的分式域，以 $\mathcal R_TH_G^*(N_d)$ 表示张量积 $H_G^*(N_d)\otimes_{H^*(BT)}\mathcal R_T$。这是将 $\lambda_i$ 的非零多项式取逆后得到的局部化。类似地，若 $\mathcal R_G$ 为 $H^*(BG)\simeq\mathbb C[\lambda_i,\hbar]$ 的分式域，则得到局部化 $\mathcal R_GH_G^*(N_d)$，其中 $\lambda_i$ 和 $\hbar$ 的非零多项式均取逆。

现在断言，存在自然的 $G$ 等变态射 $\varphi:M_d\to N_d$，描述如下。考虑双次数 $(1,d)$ 的映射 $f:\mathbb P^1\to\mathbb P^1\times\mathbb P^n$，它表示 $M_d$ 的一个元素。必要时重新参数化源空间 $\mathbb P^1$，可用以下形式的映射表示 $f\in M_d$：
\begin{equation}\label{eq:graph-space-standard-parametrization}
(w_0,w_1)\longmapsto\bigl((w_1,w_0),
(\alpha_0(w_0,w_1),\ldots,\alpha_n(w_0,w_1))\bigr),
\end{equation}
其中 $\alpha_i$ 是 $\mathbb P^1$ 坐标 $(w_0,w_1)$ 中次数 $d$ 的齐次式。对于这样的 $f$，定义 $\varphi(f)=(\alpha_0,\ldots,\alpha_n)$。概括地说，一个参数化的次数 $d$ 的映射，可由 $\mathbb P^1$ 上 $n+1$ 个无公因子的次数 $d$ 的齐次式模去共同非零标量倍来表示。

之所以看似奇怪地将 $(w_0,w_1)$ 换为 $(w_1,w_0)$，是因为要求 $\varphi$ 为 $\mathbb C^*$ 等变映射。对于 $t\in\mathbb C^*$，$t\cdot f$ 为
\[
(w_0,w_1)\longmapsto\bigl((t^{-1}w_1,w_0),(\alpha_0,\ldots,\alpha_n)\bigr)
=\bigl((w_1,tw_0),(\alpha_0,\ldots,\alpha_n)\bigr).
\]
为计算 $\varphi(t\cdot f)$，须将 $t\cdot f$ 写为标准形式~\eqref{eq:graph-space-standard-parametrization}：用 $t^{-1}w_0$ 替换 $w_0$ 即可。结果为
\[
\varphi(t\cdot f)
=(\alpha_0(t^{-1}w_0,w_1),\ldots,\alpha_n(t^{-1}w_0,w_1)),
\]
它等于 $t\cdot\varphi(f)$，如所断言。

对于任意稳定映射 $f$，$\varphi(f)=(\alpha_0,\ldots,\alpha_n)$ 的描述来自对收缩分支的研究。令 $p_1:\mathbb P^1\times\mathbb P^n\to\mathbb P^1$、$p_2:\mathbb P^1\times\mathbb P^n\to\mathbb P^n$ 为两个投影，并假设 $p_1\circ f$ 将一个分支收缩到点 $p\in\mathbb P^1$，而 $p_2\circ f$ 在此分支上的次数为 $d'$。此时可证明，$\alpha_i$ 有公因子 $L(w_0,w_1)^{d'}$，其中 $L$ 是在按式~\eqref{eq:graph-space-standard-parametrization} 的坐标交换对应于 $p$ 的点处消失的线性函数。下面定理~\ref{thm:geometric-euler-data-gluing} 的证明将使用此事实。$\varphi$ 的更明确公式见 \SourceCite{Givental2} 或 \SourceCite{LLY}。这些文献也证明了 $\varphi$ 是 $M_d$ 上的态射（\SourceCite{LLY} 所给的证明出自 J. Li）。\SourceCite{Givental2} 实际构造的是映射 $L_d\to N_d$，其中 $L_d$ 与 $M_d$ 密切相关；将在引理~\ref{lem:complete-intersection-double-construction} 的证明中定义 $L_d$。

\SourcePage{360}
\subsection{Euler 数据}\label{subsec:lly-euler-data}\index{Euler 数据|hyperpage}
固定整数 $\ell\ge1$，令 $\mathcal V=\mathcal O_{\mathbb P^n}(\ell)$。如第~\ref{subsec:quintic-mirror-theorem-statement} 小节一样，有图
\begin{equation}\label{eq:euler-data-universal-curve}
\begin{tikzcd}[column sep=large,row sep=large]
\overline M_{0,1}(\mathbb P^n,d)\arrow[r,"e_1"]\arrow[d,"\pi_1"']&\mathbb P^n\\
\overline M_{0,0}(\mathbb P^n,d)&
\end{tikzcd}
\end{equation}
并定义 $\mathcal V_d=\pi_{1*}e_1^*(\mathcal V)$。这是代数叠 $\overline{\mathcal M}_{0,0}(\mathbb P^n,d)$ 上的丛，不过与第~\ref{subsec:quintic-mirror-theorem-statement} 小节一样，滥用记号，将它视为 $\overline M_{0,0}(\mathbb P^n,d)$ 上的丛。

现在定义几个重要的等变上同调类。考虑自然投影 $\pi:M_d\to\overline M_{0,0}(\mathbb P^n,d)$，定义
\[
\chi_d^{\mathcal V}=\operatorname{Euler}_G(\pi^*(\mathcal V_d))\in H_G^*(M_d).
\]
利用上述映射 $\varphi:M_d\to N_d$，定义
\begin{equation}\label{eq:geometric-euler-data}
\widehat Q_d=\varphi_!(\chi_d^{\mathcal V})\in H_G^*(N_d),
\end{equation}
即 $\chi_d^{\mathcal V}$ 的等变推前。可将 $\widehat Q_d$ 看作 $\chi_d^{\mathcal V}$ 的“线性化”。接着定义
\begin{equation}\label{eq:hypergeometric-euler-data}
\widehat P_d=\prod_{m=0}^{\ell d}(\ell\kappa-m\hbar)\in H_G^*(N_d).
\end{equation}
$\widehat P_d$、$\widehat Q_d$ 的定义均对应于丛 $\mathcal V=\mathcal O_{\mathbb P^n}(\ell)$。在 $\mathcal V=\mathcal O_{\mathbb P^4}(5)$ 的特殊情形下，这两个等变上同调类将在定理~\ref{thm:lly-quintic-mirror} 的证明中起核心作用。

\SourceCite{LLY} 对比这里广得多的一类丛 $\mathcal V$ 定义了 $\widehat P_d$ 和 $\widehat Q_d$。若将 $\mathcal O_{\mathbb P^n}(\ell)$ 看作与超曲面相关联的丛，则 \SourceCite{LLY} 所考虑的丛，不但包含与 $\mathbb P^n$ 中完全交相关联的丛，也包含\emph{凸凹丛}。下面\hyperref[subsec:hypergeometric-mirror-transformations]{第~\ref*{subsec:hypergeometric-mirror-transformations} 小节}末的评论中将定义凸凹丛。

按照目前的构造，丛 $\mathcal V=\mathcal O_{\mathbb P^n}(\ell)$ 导出无穷列丛 $\{\mathcal V_d\}_{d=1}^{\infty}$。也定义 $d=0$ 的情形：以 $\mathcal V_0$ 表示 $\mathbb P^n$ 上的丛 $\mathcal V$。须注意，虽然这些丛位于不同空间上，它们之间仍有非平凡关系。为说明这一点，取 $0<r<d$，用拉回图定义 $\overline M(r,d-r)$：
\begin{equation}\label{eq:euler-data-gluing-fiber-product}
\begin{tikzcd}[column sep=large,row sep=large]
\overline M(r,d-r)\arrow[r,"p_1"]\arrow[d,"p_2"']&\overline M_{0,1}(\mathbb P^n,r)\arrow[d,"e_1"]\\
\overline M_{0,1}(\mathbb P^n,d-r)\arrow[r,"e_1"']&\mathbb P^n
\end{tikzcd}
\end{equation}
其中 $e_1$ 为通常的赋值映射，$p_i$ 为投影。可将 $\overline M(r,d-r)$ 看作由单标记稳定映射对 $((f_1,C_1,p_1),\allowbreak(f_2,C_2,p_2))$ 组成，它们满足 $f_1(p_1)=f_2(p_2)$。于是有图
\begin{equation}\label{eq:euler-data-gluing-map}
\begin{tikzcd}[column sep=large,row sep=large]
\overline M(r,d-r)\arrow[r,"\psi"]\arrow[d,"e_1"']&\overline M_{0,0}(\mathbb P^n,d)\\
\mathbb P^n&
\end{tikzcd}
\end{equation}
\SourcePage{361}
其中 $e_1$ 为赋值映射（由 $\overline M(r,d-r)$ 的定义可知它定义良好）。水平映射将 $((f_1,C_1,p_1),\allowbreak(f_2,C_2,p_2))$ 映到无标记稳定映射 $(f,C)$；这里 $C$ 由 $C_1\cup C_2$ 识别 $p_1$ 和 $p_2$ 后得到，且 $f|_{C_i}=f_i$（同样，由 $\overline M(r,d-r)$ 的定义可知它定义良好）。

当 $r=0$ 或 $r=d$ 时，须修改 $\overline M(r,d-r)$ 的定义。由于 $\overline M_{0,1}(\mathbb P^n,0)$ 为空，在图~\eqref{eq:euler-data-gluing-fiber-product} 中以恒等映射 $1_{\mathbb P^n}:\mathbb P^n\to\mathbb P^n$ 替换 $e_1:\overline M_{0,1}(\mathbb P^n,0)\to\mathbb P^n$。对于 $r=0$，得到 $\overline M(0,d)=\overline M_{0,1}(\mathbb P^n,d)$；在图~\eqref{eq:euler-data-gluing-fiber-product} 中，$p_1$ 识别为赋值映射 $e_1$，$p_2$ 为恒等映射。此外，图~\eqref{eq:euler-data-gluing-map} 中的 $\psi$ 识别为忘掉标记点的映射 $\pi_1:\overline M_{0,1}(\mathbb P^n,d)\to\overline M_{0,0}(\mathbb P^n,d)$。$r=d$ 的情形类似。

在这些准备下，得到以下\emph{粘合引理}。
\begin{lemma}\label{lem:euler-data-gluing}\index{Euler 数据!粘合引理|hyperpage}
若 $b$ 是向量丛的任意乘性等变特征类（例如等变 Euler 类），则在 $\overline M(r,d-r)$ 上有
\[
e_1^*b(\mathcal V)\,\psi^*b(\mathcal V_d)
=p_1^*\pi_1^*b(\mathcal V_r)\,p_2^*\pi_1^*b(\mathcal V_{d-r}).
\]
当 $r=0$ 或 $r=d$ 时，为使该式有意义，所有出现的 $\pi_1^*b(\mathcal V_0)$ 都须解释为 $b(\mathcal V_0)=b(\mathcal V)$。
\end{lemma}
\begin{proof}
当 $r=0,d$ 时，引理显然成立。设 $0<r<d$，并设 $((f_1,C_1,p_1),\allowbreak(f_2,C_2,p_2))$ 如上映到 $(f,C)$。有正合列
\[
0\longrightarrow f^*\mathcal V\longrightarrow
f_1^*\mathcal V\oplus f_2^*\mathcal V\longrightarrow\mathcal V|_p\longrightarrow0,
\qquad p=f_1(p_1)=f_2(p_2).
\]
于是得到正合列
\[
0\longrightarrow H^0(C,f^*\mathcal V)\longrightarrow
H^0(C_1,f_1^*\mathcal V)\oplus H^0(C_2,f_2^*\mathcal V)
\longrightarrow\mathcal V|_p\longrightarrow0.
\]
由此即可得到引理。
\end{proof}

粘合引理给出等变 Euler 类 $\operatorname{Euler}_G(\mathcal V_d)\in H_G^*(\overline M_{0,0}(\mathbb P^n,d))$ 之间的恒等式。由于 $\widehat Q_d\in H_G^*(N_d)$ 用 $\operatorname{Euler}_G(\mathcal V_d)$ 定义，自然会问这些类是否满足类似恒等式。$H_G^*(N_d)$ 中的类容易研究：固定点集很简单，仅由孤立点 $p_{i,r}$ 组成。因此，由\IndexLocation{indexloc:0739:361}局部化定理，类 $\widehat Q_d$ 被限制 $i_{p_{i,r}}^*(\widehat Q_d)$ 唯一确定。这提示我们寻找这些限制所满足的、类似于引理~\ref{lem:euler-data-gluing} 的恒等式。

下面的定义~\ref{def:omega-euler-data} 将使这一点明确。不过，在陈述定义之前，需要映射 $\overline{\phantom\alpha}:N_d\to N_d$，定义为
\[
(\alpha_0(w_0,w_1),\ldots,\alpha_n(w_0,w_1))
\longmapsto(\alpha_0(w_1,w_0),\ldots,\alpha_n(w_1,w_0)).
\]
此映射显然 $T$ 等变，但不 $\mathbb C^*$ 等变。\SourceCite{LLY} 证明，配合 $\mathbb C^*$ 参数反转，它在等变上同调上给出的变换 $\overline{\phantom\alpha}:H_G^*(N_d)\to H_G^*(N_d)$ 满足
\begin{equation}\label{eq:euler-data-involution}
\overline\kappa=\kappa-d\hbar,\qquad\overline\hbar=-\hbar,
\qquad\overline{\lambda_i}=\lambda_i.
\end{equation}
还回顾例题~\ref{ex:projective-equivariant-localization}：$G$（或 $T$）在 $\mathbb P^n$ 上作用的不动点记为 $q_i$。

\SourcePage{362}
\begin{definition}\label{def:omega-euler-data}\index{Euler 数据!定义|hyperpage}
设 $\Omega\in\mathcal R_TH_T^*(\mathbb P^n)$ 为可逆类。一列 $Q=\{Q_d\}$，其中 $Q_d\in\mathcal R_TH_G^*(N_d)$（$d>0$），若对所有 $r=0,\ldots,d$ 及 $i=0,\ldots,n$ 均有
\[
i_{q_i}^*(\Omega)i_{p_{i,r}}^*(Q_d)
=\overline{i_{p_{i,0}}^*(Q_r)}\,i_{p_{i,0}}^*(Q_{d-r}),
\]
则称它为一组 \emph{$\Omega$--Euler 数据}，其中 $Q_0=\Omega$。所有 $\Omega$--Euler 数据组成的集合记为 $\mathcal A^\Omega$。
\end{definition}
前面已指出，$i_{p_{i,r}}^*(Q_d)=Q_d(\lambda_i+r\hbar)$；由式~\eqref{eq:projective-polynomial-fixed-restriction}，有 $i_{q_i}^*(\Omega)=\Omega(\lambda_i)$。因此，上式可写成
\[
\Omega(\lambda_i)Q_d(\lambda_i+r\hbar)
=\overline{Q_r(\lambda_i)}\,Q_{d-r}(\lambda_i).
\]
注意，由于 $\Omega\in\mathcal R_TH_T^*(\mathbb P^n)$，有 $\overline\Omega=\Omega$，故 $r=0$ 时此式显然成立。下面将证明，由式~\eqref{eq:geometric-euler-data} 和~\eqref{eq:hypergeometric-euler-data} 定义的序列 $\widehat Q=\{\widehat Q_d\}$、$\widehat P=\{\widehat P_d\}$ 均为 $\Omega=\operatorname{Euler}_T(\mathcal V)=\operatorname{Euler}_T(\mathcal O_{\mathbb P^n}(\ell))=\ell p$ 的 Euler 数据，其中 $p$ 为等变超平面类。采用“Euler 数据”这一名称，部分是因为 $\widehat Q$ 用等变 Euler 类定义。

先给出 $\widehat P$ 是一组 $\ell p$--Euler 数据的简单证明。
\SourceExampleNumber{11.1.3}{1}
\begin{example}\label{ex:hypergeometric-euler-data-gluing}
先注意，$\Omega=\ell p$ 在 $\mathcal R_TH_T^*(\mathbb P^n)$ 中可逆。由局部化（命题~\ref{prop:atiyah-bott-localization}）即可看出这一点，因为限制 $i_{q_i}^*(\Omega)=\Omega(\lambda_i)=\ell\lambda_i$ 均可逆。还回顾式~\eqref{eq:hypergeometric-euler-data}，有 $\widehat P_d=\prod_{m=0}^{\ell d}(\ell\kappa-m\hbar)$。于是，$\widehat P=\{\widehat P_d\}$ 为 $\ell p$--Euler 数据，等价于恒等式
\[
\ell\lambda_i\prod_{m=0}^{\ell d}\bigl(\ell(\lambda_i+r\hbar)-m\hbar\bigr)
=\prod_{m=0}^{\ell r}(\ell\lambda_i+m\hbar)
\prod_{m=0}^{\ell(d-r)}(\ell\lambda_i-m\hbar),
\]
而此式立即可验证。
\end{example}
证明 $\widehat Q$ 是 $\ell p$--Euler 数据需要更多工作。
\begin{theorem}\label{thm:geometric-euler-data-gluing}
序列 $\widehat Q=\{\widehat Q_d\}=\{\varphi_!\chi_d^{\mathcal V}\}$ 是一组 $\ell p$--Euler 数据。
\end{theorem}
\begin{proof}
须证明，对所有 $r,i$，有
\begin{equation}\label{eq:geometric-euler-data-gluing-target}
\ell\lambda_i\widehat Q_d(\lambda_i+r\hbar)
=\overline{\widehat Q_r(\lambda_i)}\,\widehat Q_{d-r}(\lambda_i).
\end{equation}
先计算 $i_{p_{i,r}}^*(\widehat Q_d)$，它为
\begin{equation}\label{eq:geometric-euler-data-localization-integral}
\widehat Q_d(\lambda_i+r\hbar)
=\int_{(N_d)_G}\phi_{p_{i,r}}\widehat Q_d
=\int_{(M_d)_G}\varphi^*(\phi_{p_{i,r}})\chi_d^{\mathcal V}.
\end{equation}
第一个等号来自式~\eqref{eq:linear-model-restriction-integral}，第二个等号来自 $\widehat Q_d$ 的定义和投影公式。策略是用推论~\ref{cor:localization-integration} 的局部化公式计算式~\eqref{eq:geometric-euler-data-localization-integral} 右端的积分。

先确定 $M_d$ 中哪些 $G$ 不动分支对此积分有贡献。令 $i_\Gamma:\overline M_\Gamma\subset M_d$ 表示与图 $\Gamma$ 对应的分支。将推论~\ref{cor:localization-integration} 应用于此积分时，$\overline M_\Gamma$ 的贡献项含因子
\[
i_\Gamma^*(\varphi^*(\phi_{p_{i,r}}))
=(\varphi|_{\overline M_\Gamma})^*\bigl(i_{\varphi(\overline M_\Gamma)}^*(\phi_{p_{i,r}})\bigr).
\]
\SourcePage{363}
由于 $\varphi(\overline M_\Gamma)$ 必为某个不动点 $p_{j,s}$，由式~\eqref{eq:linear-model-fixed-point-class} 中 $\phi_{p_{i,r}}$ 的定义，当 $(j,s)\ne(i,r)$ 时，相应项为零。

剩下的是确定在 $\varphi$ 下映到 $p_{i,r}$ 的 $G$ 固定点集分支。为此使用式~\eqref{eq:euler-data-gluing-map} 中 $\psi$ 的一个变体。给定 $\overline M(r,d-r)$ 的一点，它由满足 $f_1(p_1)=f_2(p_2)=q\in\mathbb P^n$ 的 $f_j:(C_j,p_j)\to\mathbb P^n$ 表示。定义 $C=C_0\cup C_1\cup C_2$，其中 $C_0=\mathbb P^1$，将 $0\in\mathbb P^1$ 与 $p_1\in C_1$ 粘合，将 $\infty\in\mathbb P^1$ 与 $p_2\in C_2$ 粘合。然后定义 $f:C\to\mathbb P^1\times\mathbb P^n$ 为
\[
f|_{C_0}(z)=(z,q),\qquad f|_{C_1}(z)=(0,f_1(z)),
\qquad f|_{C_2}(z)=(\infty,f_2(z)).
\]
于是 $f\in M_d$，故这一过程定义映射
\begin{equation}\label{eq:euler-data-graph-gluing-map}
\widetilde\psi:\overline M(r,d-r)\longrightarrow M_d.
\end{equation}
令 $\{F_k\}$ 为 $\overline M_{0,1}(\mathbb P^n,k)$ 的 $T$ 不动分支中、标记点映到 $q_i$ 的分支集合。于是 $F_r\times F_{d-r}\subset\overline M(r,d-r)$；将其与 $\widetilde\psi$ 下的像识别，便写作 $F_r\times F_{d-r}\subset M_d$。

断言：遍历所有可能的 $F_r$ 和 $F_{d-r}$ 所得到的 $F_r\times F_{d-r}$，恰好就是在 $\varphi$ 下映到 $p_{i,r}$ 的 $G$ 固定点集分支。为证明此断言，取 $f\in F_r\times F_{d-r}$。按定义，$p_1\circ f$ 收缩 $C_1$、$C_2$，而 $p_2\circ f$ 在 $C_1$、$C_2$ 上的次数分别为 $r$、$d-r$。于是，第~\ref{subsec:stable-moduli-linear-models} 小节对 $\varphi(f)$ 的描述给出 $\varphi(f)=p_{i,r}$。利用 \SourceCite{LLY} 中更详细的 $\varphi$ 描述，还可证明 $F_r\times F_{d-r}$ 是 $G$ 固定点集的一个分支，且 $\varphi^{-1}(p_{i,r})$ 中的所有 $G$ 不动分支都由此得到。

因此，用局部化计算式~\eqref{eq:geometric-euler-data-localization-integral} 右端积分时，只需要形如 $F_r\times F_{d-r}$ 的分支。先处理 $r\ne0,d$ 的情形。以下记号较为有用。令 $\mathcal L_{k,1}$ 为 $\overline M_{0,1}(\mathbb P^n,k)$ 上的线丛，其在 $f:(C,p)\to\mathbb P^n$ 处的纤维是余切空间 $T_p^*(C)$，如\hyperref[subsec:gravitational-correlator-definition]{第~\ref*{subsec:gravitational-correlator-definition} 小节}开头所定义。对任意分支 $F_k$，以 $N(F_k)$ 表示 $F_k$ 在 $\overline M_{0,1}(\mathbb P^n,k)$ 中的法丛。最后，令 $\pi_1:\overline M_{0,1}(\mathbb P^n,k)\to\overline M_{0,0}(\mathbb P^n,k)$ 为图~\eqref{eq:euler-data-universal-curve} 中定义丛 $\mathcal V_k$ 时使用的映射。将推论~\ref{cor:orbifold-localization-integration} 应用于式~\eqref{eq:geometric-euler-data-localization-integral} 右端积分，断言
\begin{equation}\label{eq:geometric-euler-data-product-localization}
\begin{aligned}
\widehat Q_d(\lambda_i+r\hbar)
&=\int_{(M_d)_G}\varphi^*(\phi_{p_{i,r}})\chi_d^{\mathcal V}\\
&=-\hbar^{-2}(\ell\lambda_i)^{-1}i_{p_{i,r}}^*(\phi_{p_{i,r}})
\prod_{j\ne i}(\lambda_i-\lambda_j)\\
&\quad\times\sum_{F_r}\int_{(F_r)_G}
\frac{\pi_1^*\operatorname{Euler}_G(\mathcal V_r)}
{a_r\operatorname{Euler}_G(N(F_r))(\hbar-c_1^G(\mathcal L_{r,1}))}\\
&\quad\times\sum_{F_{d-r}}\int_{(F_{d-r})_G}
\frac{\pi_1^*\operatorname{Euler}_G(\mathcal V_{d-r})}
{a_{d-r}\operatorname{Euler}_G(N(F_{d-r}))(-\hbar-c_1^G(\mathcal L_{d-r,1}))}.
\end{aligned}
\end{equation}
其中 $a_k$ 为与 $F_k$ 相关联的有限自同构群的阶。（在 \SourceCite{LLY} 中，$a_k$ 已包含在轨形积分的定义中，因此未在公式里明确出现。）下面概述这个计算，省略许多细节。

\SourcePage{364}
固定一般稳定映射 $f\in F_r\times F_{d-r}$ 的自同构，只能来自 $f|_{C_1}\in F_r$ 和 $f|_{C_2}\in F_{d-r}$ 的自同构，所以局部化公式所需的自同构群的阶为 $a_ra_{d-r}$。这解释了式~\eqref{eq:geometric-euler-data-product-localization} 分母中的因子 $a_r$、$a_{d-r}$。

还需要 $F_r\times F_{d-r}$ 在 $M_d$ 中的法丛 $N(F_r\times F_{d-r})$ 的等变 Euler 类。这个计算是论证中技术性最强的部分，因为有若干项贡献于 $\operatorname{Euler}_G(N(F_r\times F_{d-r}))$。有些地方只作简略说明，更完整的证明见 \SourceCite{LLY}。为研究此法丛，令 $M^i(k)\subset\overline M_{0,1}(\mathbb P^n,k)$ 为将标记点映到 $q_i\in\mathbb P^n$ 的稳定映射集。注意，此嵌入的法丛是 $T_{q_i}\mathbb P^n$ 经由 $e_1$ 的拉回。考虑以下图：
\[
\begin{tikzcd}[column sep=1.5em,row sep=1.2em]
F_r\times F_{d-r}\arrow[r]&M^i(r)\times M^i(d-r)\arrow[r]\arrow[d]&
\overline M_{0,1}(\mathbb P^n,r)\times\overline M_{0,1}(\mathbb P^n,d-r)\\
&\overline M(r,d-r)\arrow[d,"\widetilde\psi"]&\\
&M_d&
\end{tikzcd}
\]
其中 $\widetilde\psi$ 来自式~\eqref{eq:euler-data-graph-gluing-map}。式~\eqref{eq:euler-data-gluing-fiber-product} 的定义表明，$e_1^{-1}(q_i)=M^i(r)\times M^i(d-r)\subset\overline M(r,d-r)$，其中 $e_1$ 是图~\eqref{eq:euler-data-gluing-map} 的赋值映射。因此，此嵌入的法丛也是 $T_{q_i}\mathbb P^n$ 的拉回。

计算 $F_r\times F_{d-r}\subset M_d$ 的法丛时，上图意味着，在 $F_r\times F_{d-r}$ 的等变 $K$ 群中，有
\[
N(F_r\times F_{d-r})=N(F_r)+N(F_{d-r})-2T_{q_i}\mathbb P^n
+T_{q_i}\mathbb P^n+N(\widetilde\psi).
\]
为清楚起见，此式省去了若干拉回。取等变 Euler 类，得到（仍省去拉回）
\[
\operatorname{Euler}_G(N(F_r\times F_{d-r}))
=\frac{\operatorname{Euler}_G(N(F_r))\operatorname{Euler}_G(N(F_{d-r}))}
{\prod_{j\ne i}(\lambda_i-\lambda_j)}\operatorname{Euler}_G(N(\widetilde\psi)),
\]
因为由式~\eqref{eq:projective-fixed-euler}，$\operatorname{Euler}_G(T_{q_i}\mathbb P^n)=\prod_{j\ne i}(\lambda_i-\lambda_j)$。这出现在局部化定理的分母中，因而又解释了式~\eqref{eq:geometric-euler-data-product-localization} 中的三个因子。

还须考虑映射 $\widetilde\psi$ 的法丛。先考虑引理~\ref{lem:node-smoothing-tangent-tensor} 研究过的\IndexLocation{indexloc:0940:364}结点形变。对于 $p_1=C_0\cap C_1$，引理要求考察 $T_{p_1}C_0\otimes T_{p_1}C_1$。由于 $C_0\simeq\mathbb P^1$ 为等变同构，而 $\mathcal L_{r,1}$ 由 $C_1$ 的余切丛定义，权重为 $\hbar$ 和 $-c_1^G(\mathcal L_{r,1})$。这解释了式~\eqref{eq:geometric-euler-data-product-localization} 分母中的 $\hbar-c_1^G(\mathcal L_{r,1})$；类似地，形变 $p_2=C_0\cap C_2$ 导出分母中的 $-\hbar-c_1^G(\mathcal L_{d-r,1})$。\footnote{这与引理~\ref{lem:graph-space-descendant-factorization} 有关，并解释了式~\eqref{eq:giventals-flat-section-symbolic} 等公式中的 $\hbar-c$。}最后研究 $f|_{C_0}$ 的形变。知道 $C_0=\mathbb P^1$，而 $\overline M(r,d-r)$ 和 $\widetilde\psi$ 的定义表明，此类映射分别将 $0$、$\infty$ 映到 $(0,q)$、$(\infty,q)$，其中 $q\in\mathbb P^n$。允许 $(0,q)$、$(\infty,q)$ 的第一坐标变化，就可形变这个映射。这些法向形变可识别为 $\mathbb P^1$ 的无穷小自同构模去保持 $0$、$\infty$ 的那些自同构。用 $\mathbb P^1$ 的标准坐标 $(z_0,z_1)$，它们由向量场 $z_0\partial/\partial z_1$、$z_1\partial/\partial z_0$ 给出，产生因子 $-\hbar^2$。由于此因子出现在局部化定理的分母中，便得到式~\eqref{eq:geometric-euler-data-product-localization} 中的因子 $-\hbar^{-2}$。

应用局部化公式还需要最后一个要素：确定 $\varphi^*(\phi_{p_{i,r}})\chi_d^{\mathcal V}=\varphi^*(\phi_{p_{i,r}})\pi^*(\operatorname{Euler}_G(\mathcal V_d))$ 在 $F_r\times F_{d-r}$ 上的限制。由于 $\varphi(F_r\times F_{d-r})=p_{i,r}$，$\varphi^*(\phi_{p_{i,r}})$ 的限制为 $i_{p_{i,r}}^*(\phi_{p_{i,r}})$，这就是式~\eqref{eq:geometric-euler-data-product-localization} 的一个因子。至于
\SourcePage{365}
$\pi^*(\operatorname{Euler}_G(\mathcal V_d))$ 的限制，粘合引理（引理~\ref{lem:euler-data-gluing}）意味着，在 $F_r\times F_{d-r}$ 上有恒等式
\[
\operatorname{Euler}_G(\mathcal V)\operatorname{Euler}_G(\mathcal V_d)
=\operatorname{Euler}_G(\mathcal V_r)\operatorname{Euler}_G(\mathcal V_{d-r}),
\]
为方便起见，省去了拉回。由于 $\operatorname{Euler}_G(\mathcal V)=\ell p$，可知 $\chi_d^{\mathcal V}$ 在 $F_r\times F_{d-r}$ 上的限制为
\[
(\ell p)^{-1}\pi_1^*\operatorname{Euler}_G(\mathcal V_r)
\pi_1^*\operatorname{Euler}_G(\mathcal V_{d-r}).
\]
这些因子出现在局部化公式的分子中，从而完成对式~\eqref{eq:geometric-euler-data-product-localization} 的解释。

对 $r=0$ 作类似分析，得到公式
\begin{equation}\label{eq:geometric-euler-data-endpoint-localization}
\widehat Q_d(\lambda_i)=-\hbar^{-1}i_{p_{i,0}}^*(\phi_{p_{i,0}})
\sum_{F_d}\int_{(F_d)_G}
\frac{\pi_1^*\operatorname{Euler}_G(\mathcal V_d)}
{a_d\operatorname{Euler}_G(N(F_d))(-\hbar-c_1^G(\mathcal L_{d,1}))}.
\end{equation}
现在可证明定理。在上式中以 $r$ 替换 $d$，并应用式~\eqref{eq:euler-data-involution}，得到
\[
\overline{\widehat Q_r(\lambda_i)}
=\hbar^{-1}\overline{i_{p_{i,0}}^*(\phi_{p_{i,0}})}
\sum_{F_r}\int_{(F_r)_G}
\frac{\pi_1^*\operatorname{Euler}_G(\mathcal V_r)}
{a_r\operatorname{Euler}_G(N(F_r))(\hbar-c_1^G(\mathcal L_{r,1}))},
\]
再在式~\eqref{eq:geometric-euler-data-endpoint-localization} 中以 $d-r$ 替换 $d$，得到
\[
\widehat Q_{d-r}(\lambda_i)
=-\hbar^{-1}i_{p_{i,0}}^*(\phi_{p_{i,0}})
\sum_{F_{d-r}}\int_{(F_{d-r})_G}
\frac{\pi_1^*\operatorname{Euler}_G(\mathcal V_{d-r})}
{a_{d-r}\operatorname{Euler}_G(N(F_{d-r}))(-\hbar-c_1^G(\mathcal L_{d-r,1}))}.
\]
将两式相乘，与式~\eqref{eq:geometric-euler-data-product-localization} 乘以 $\ell\lambda_i$ 后比较，所需等式~\eqref{eq:geometric-euler-data-gluing-target} 便归结为恒等式
\[
i_{p_{i,r}}^*(\phi_{p_{i,r}})\prod_{j\ne i}(\lambda_i-\lambda_j)
=\overline{i_{p_{i,0}}^*(\phi_{p_{i,0}})}\,i_{p_{i,0}}^*(\phi_{p_{i,0}}).
\]
回顾式~\eqref{eq:linear-model-fixed-class-restriction}，$i_{p_{i,r}}^*(\phi_{p_{i,r}})$ 正是式~\eqref{eq:linear-model-normal-euler} 所计算的 $\operatorname{Euler}_G(N_{i,r})$。因此，上式归结为
\[
\begin{aligned}
&\prod_{(j,s)\ne(i,r)}\bigl(\lambda_i-\lambda_j+(r-s)\hbar\bigr)
\prod_{j\ne i}(\lambda_i-\lambda_j)\\
&\qquad=\prod_{(j,s)\ne(i,0)}\bigl(\lambda_i-\lambda_j-s\hbar\bigr)
\prod_{(j,s)\ne(i,0)}\bigl(\lambda_i-\lambda_j+s\hbar\bigr).
\end{aligned}
\]
其中左端的 $s$ 从 $0$ 到 $d$；右端第一个乘积的 $s$ 从 $0$ 到 $d-r$，第二个乘积的 $s$ 从 $0$ 到 $r$；所有乘积中 $j$ 均从 $0$ 到 $n$。此恒等式容易验证，从而完成 $r\ne0,d$ 时的证明。

由于 $r=0$ 时 Euler 恒等式显然成立（见定义~\ref{def:omega-euler-data} 后的讨论），只须再考虑 $r=d$。此时可为 $\widehat Q_d(\lambda_i+d\hbar)$ 导出类似于式~\eqref{eq:geometric-euler-data-endpoint-localization} 中 $\widehat Q_d(\lambda_i)$ 的公式，比较这些公式便得到所需恒等式。细节见 \SourceCite{LLY}。
\end{proof}

\subsection{相关联的 Euler 数据}\label{subsec:linked-euler-data}
现在转向一个重要概念。令 $\mathcal S$ 为序列 $P=\{P_d\}_{d=1}^{\infty}$ 组成的集合，其中对所有 $d$，有 $P_d\in\mathcal R_GH_G^*(N_d)$；注意 $i_{p_{i,0}}^*$ 将 $\mathcal R_GH_G^*(N_d)$ 映到 $\mathcal R_G$。
\begin{definition}\label{def:linked-euler-data}\index{Euler 数据!关联关系|hyperpage}
若对所有 $i\ne j$、$d>0$，$i_{p_{i,0}}^*(P_d-Q_d)\in\mathcal R_G$ 在 $\hbar=(\lambda_i-\lambda_j)/d$ 处有定义且为零，则称序列 $P,Q\in\mathcal S$ \emph{相关联}。
\end{definition}
\SourcePage{366}
为使此定义精确，需要解释 $\hbar=(\lambda_i-\lambda_j)/d$ 的含义。定义子群 $G'\subset G$ 为
\[
G'=\{(t,t_0,\ldots,t_n):t^d=t_i/t_j\}.
\]
包含 $G'\hookrightarrow G$ 诱导映射 $BG'\to BG$，从而在上同调上诱导拉回 $\operatorname{res}_{ij}:H^*(BG)\to H^*(BG')$。由\hyperref[subsec:equivariant-cohomology]{第~\ref*{subsec:equivariant-cohomology} 小节}的讨论，容易看出 $\operatorname{res}_{ij}$ 来自特征标的限制。这意味着 $\operatorname{res}_{ij}(\hbar)=\operatorname{res}_{ij}((\lambda_i-\lambda_j)/d)$；有时简写为 $\hbar=(\lambda_i-\lambda_j)/d$，并理解为这些符号位于 $G'$ 等变上同调中。类似地，若空间 $X$ 上有 $G$ 作用，则有映射 $\operatorname{res}_{ij}:H_G^*(X)\to H_{G'}^*(X)$。于是得到复合
\begin{equation}\label{eq:linked-data-restriction-first}
H_G^*(N_d)\xrightarrow{i_{p_{i,0}}^*}H_G^*(p_{i,0})
\xrightarrow{\operatorname{res}_{ij}}H_{G'}^*(p_{i,0}).
\end{equation}
定义~\ref{def:linked-euler-data} 的条件意味着，在所有相关分母有定义的局部化上，对所有 $i\ne j$、$d$，有 $\operatorname{res}_{ij}\circ i_{p_{i,0}}^*(P_d-Q_d)=0$。为后续使用，注意上述复合等于
\begin{equation}\label{eq:linked-data-restriction-second}
H_G^*(N_d)\xrightarrow{\operatorname{res}_{ij}}H_{G'}^*(N_d)
\xrightarrow{i_{p_{i,0}}^*}H_{G'}^*(p_{i,0}),
\end{equation}
所以，相关联也意味着，在这样的局部化上，对所有 $i\ne j$、$d$，有 $i_{p_{i,0}}^*\circ\operatorname{res}_{ij}(P_d-Q_d)=0$。

子群 $G'$ 值得研究，因为它固定了直线 $\overline{q_iq_j}\subset\mathbb P^n$ 的多重覆盖所组成的空间。考虑 $N_d$ 中参数化给出的直线
\[
L_{i,j}=\{(0,\ldots,0,aw_0^d,0,\ldots,0,bw_1^d,0,\ldots,0)\},
\]
其中 $aw_0^d$ 在第 $j$ 个位置，$bw_1^d$ 在第 $i$ 个位置。注意，$p_{i,0}\in L_{i,j}$ 对应于 $a=0$。$L_{i,j}$ 的一般点，是直线 $\overline{q_iq_j}\subset\mathbb P^n$ 的一个次数 $d$ 的多重覆盖在 $\varphi$ 下的像。子群 $G'$ 逐点固定 $L_{i,j}$，所以 $H_{G'}^*(L_{i,j})=H^*(L_{i,j})\otimes H^*(BG')$。因此，若 $u\in L_{i,j}$ 为任意点，则\IndexLocation{indexloc:0249:366}限制映射 $H_{G'}^*(L_{i,j})\to H_{G'}^*(u)$ 不依赖于 $u$。于是图
\[
\begin{tikzcd}[column sep=large,row sep=large]
H_G^*(N_d)\arrow[r,"\operatorname{res}_{ij}"]&H_{G'}^*(N_d)\arrow[r]\arrow[dr,"i_u^*"']&H_{G'}^*(L_{i,j})\arrow[d]\\
&&H_{G'}^*(u)
\end{tikzcd}
\]
中的映射也不依赖于 $u$。由于式~\eqref{eq:linked-data-restriction-second} 是 $u=p_{i,0}$ 时的特殊情形，可知在其后讨论的“限制”映射 $i_{p_{i,0}}^*\circ\operatorname{res}_{ij}$ 中，可用任意 $u\in L_{i,j}$ 替换 $p_{i,0}$。因此，$P$ 与 $Q$ 相关联，当且仅当它们在 $\mathbb P^n$ 的坐标直线的多重覆盖上相等。

现在陈述一个关键结果。
\begin{theorem}\label{thm:linked-geometric-hypergeometric-data}
在 $\hbar=(\lambda_i-\lambda_j)/d$、$i\ne j$ 时，有
\[
i_{p_{i,0}}^*(\widehat Q_d)
=\prod_{m=0}^{\ell d}\left(\ell\lambda_i-m\frac{\lambda_i-\lambda_j}{d}\right).
\]
特别地，Euler 数据 $\widehat P$ 与 $\widehat Q$ 相关联。
\end{theorem}
\SourcePage{367}
\begin{proof}
先回顾 $\widehat P_d=\prod_{m=0}^{\ell d}(\ell\kappa-m\hbar)$。由式~\eqref{eq:linked-data-restriction-first}，可用复合 $\operatorname{res}_{ij}\circ i_{p_{i,0}}^*$ 理解 $\hbar=(\lambda_i-\lambda_j)/d$。于是 $\operatorname{res}_{ij}\circ i_{p_{i,0}}^*(\hbar)=(\lambda_i-\lambda_j)/d$，而第~\ref{subsec:stable-moduli-linear-models} 小节给出 $i_{p_{i,0}}^*(\kappa)=\lambda_i$。因此
\[
\operatorname{res}_{ij}\circ i_{p_{i,0}}^*(\widehat P_d)
=\prod_{m=0}^{\ell d}\left(\ell\lambda_i-m\frac{\lambda_i-\lambda_j}{d}\right).
\]
所以，计算出 $\operatorname{res}_{ij}\circ i_{p_{i,0}}^*(\widehat Q_d)$ 后，即可得到定理的最后一个断言。

上面已证明，$\widehat Q_d$ 的限制也可由任意 $u\in L_{i,j}$ 对应的复合 $i_u^*\circ\operatorname{res}_{ij}(\widehat Q_d)$ 表示。特别地，可选
\[
u=(0,\ldots,0,w_0^d,0,\ldots,0,w_1^d,0,\ldots,0)\in L_{i,j}\subset N_d,
\]
其中 $w_0^d$、$w_1^d$ 分别位于第 $j$、第 $i$ 个位置。

证明的策略是：由于 $\widehat Q_d=\varphi_!(\chi_d^{\mathcal V})$，可在 $M_d$ 上进行计算。对于上述 $u$，考虑映射 $f:\mathbb P^1\to\mathbb P^1\times\mathbb P^n$，定义为
\[
f(w_0,w_1)=\bigl((w_1,w_0),
(0,\ldots,0,w_0^d,0,\ldots,0,w_1^d,0,\ldots,0)\bigr).
\]
这是 $M_d$ 的一个元素，显然 $\varphi(f)=u$。关键事实是，映射 $\varphi:M_d\to N_d$ 在这些点的邻域中为局部同构。考虑交换图
\[
\begin{tikzcd}[column sep=large,row sep=large]
H_G^*(M_d)\arrow[r,"\operatorname{res}_{ij}"]\arrow[d,"\varphi_!"']&H_{G'}^*(M_d)\arrow[r,"i_f^*"]\arrow[d,"\varphi_!"']&H_{G'}^*(f)\arrow[d,"\varphi_!"']\\
H_G^*(N_d)\arrow[r,"\operatorname{res}_{ij}"']&H_{G'}^*(N_d)\arrow[r,"i_u^*"']&H_{G'}^*(u)
\end{tikzcd}
\]
右端的竖直映射为同构，所以只须计算 $i_f^*\circ\operatorname{res}_{ij}(\chi_d^{\mathcal V})$。

回顾 $\chi_d^{\mathcal V}$ 为 $\pi^*\mathcal V_d$ 的等变 Euler 类。因此，在 $f$ 上须计算 $G'$ 对 $H^0(\mathbb P^1,f^*\mathcal O_{\mathbb P^n}(\ell))$ 作用的权重。一组截面基为 $w_0^mw_1^{\ell d-m}$（$0\le m\le\ell d$）。由于 $f$ 的次数为 $d$，$w_0$ 的权重为 $\lambda_j/d$，$w_1$ 的权重为 $\lambda_i/d$。这类似于\hyperref[chap:localization]{第九章}式~\eqref{eq:multiple-cover-cech-basis} 之后的讨论。于是截面 $w_0^mw_1^{\ell d-m}$ 的权重为
\[
m\frac{\lambda_j}d+(\ell d-m)\frac{\lambda_i}d
=\ell\lambda_i-m\frac{\lambda_i-\lambda_j}d.
\]
将这些权重相乘，就得到 $\hbar=(\lambda_i-\lambda_j)/d$ 时 $i_{p_{i,0}}^*(\widehat Q_d)$ 的所需公式。
\end{proof}

还需要一个唯一性结果。
\begin{theorem}\label{thm:linked-euler-data-uniqueness}
设 $P,Q$ 为相关联的 $\Omega$--Euler 数据，且对所有 $0\le i\le n$、$d\ge1$，有
\[
\deg_\hbar i_{p_{i,0}}^*(P_d-Q_d)\le(n+1)d-2.
\]
则 $P=Q$。
\end{theorem}
\begin{proof}
只简述证明。完整论证需要\IndexLocation{indexloc:1207:367}互易引理 \SourceCite[引理 2.4]{LLY}，该引理研究作为 Euler 数据所带来的结果。证明从以下观察开始：$\{\kappa^s\}=\{1,\kappa,\kappa^2,\ldots\}$ 生成
\SourcePage{368}
$H_G^*(N_d)$ 这个 $H^*(BG)$ 模。由于交配对非退化，只须证明，对所有 $s\ge0$，有
\begin{equation}\label{eq:euler-data-uniqueness-integral}
\int_{(N_d)_G}\kappa^s(P_d-Q_d)=0.
\end{equation}
以 $L_s$ 表示此式左端的积分。Euler 数据的定义意味着 $P_d$、$Q_d$ 关于 $\hbar$ 为多项式，故 $L_s$ 也是 $\hbar$ 的多项式。这一点稍后很重要。

对 $d$ 归纳，假定对 $0\le j\le d-1$，有 $P_j=Q_j$（回顾 $P_0=Q_0=\Omega$）。用局部化（推论~\ref{cor:localization-integration}）计算 $L_s$，得到对不动点 $p_{i,r}$ 的求和。由于 $P$ 为 Euler 数据，当 $0<r<d$ 时，可用 $P_1,\ldots,P_{d-1}$ 表示 $i_{p_{i,r}}^*(P_d)$；对 $Q$ 也一样。由归纳假设，仅 $r=0,d$ 的项可能贡献于式~\eqref{eq:euler-data-uniqueness-integral} 的积分。利用 \SourceCite[引理 2.4]{LLY}，可证明 $L_s$ 的局部化公式归结为
\begin{equation}\label{eq:euler-data-uniqueness-localized-sum}
L_s=\sum_{i=0}^n\left(\frac{\lambda_i^sA_i(\hbar)}{\hbar^d}
+\frac{(\lambda_i+d\hbar)^sA_i(-\hbar)}{(-\hbar)^d}\right),
\end{equation}
其中
\begin{equation}\label{eq:euler-data-uniqueness-a-polynomial}
A_i(\hbar)=\frac{(-1)^d}{d!\prod_{j\ne i}(\lambda_i-\lambda_j)}
\frac{i_{p_{i,0}}^*(P_d-Q_d)}
{\prod_{j\ne i}\prod_{s=1}^d(\lambda_i-(\lambda_j+s\hbar))}.
\end{equation}
分母部分由以下公式解释：
\[
\operatorname{Euler}_G(N_{p_{i,0}})=i_{p_{i,0}}^*(\phi_{p_{i,0}})
=\prod_{(j,s)\ne(i,0)}(\lambda_i-(\lambda_j+s\hbar)).
\]
式~\eqref{eq:euler-data-uniqueness-localized-sum} 和~\eqref{eq:euler-data-uniqueness-a-polynomial} 的分子、分母均含 $\hbar$。但 $P,Q$ 相关联，所以对所有 $j\ne i$，$i_{p_{i,0}}^*(P_d-Q_d)$ 在 $\hbar=(\lambda_i-\lambda_j)/d$ 处为零。分母中的因子 $\lambda_i-(\lambda_j+d\hbar)$ 也在此处为零，故与分子中的一个因子约去。此外，利用 \SourceCite[引理 2.4]{LLY} 和归纳假设，还可证明式~\eqref{eq:euler-data-uniqueness-a-polynomial} 分母中其他含 $\hbar$ 的因子也都约去。因此，$A_i$ 是 $\hbar$ 的多项式。于是式~\eqref{eq:euler-data-uniqueness-a-polynomial} 及定理中的次数界给出
\[
\deg_\hbar A_i\le(n+1)d-2-nd=d-2.
\]
由于对所有 $s$，$L_s$ 均为 $\hbar$ 的多项式，式~\eqref{eq:euler-data-uniqueness-localized-sum} 意味着 $\hbar^d$ 必须整除
\begin{equation}\label{eq:euler-data-uniqueness-divisibility}
\sum_{i=0}^n\bigl(\lambda_i^sA_i(\hbar)
+(-1)^d(\lambda_i+d\hbar)^sA_i(-\hbar)\bigr).
\end{equation}
现在证明所有 $A_i=0$。否则，令 $m$ 为各 $A_i$ 中非零项的最小次数，并写成
\[
A_i=a_i\hbar^m+b_i\hbar^{m+1}+\cdots,
\qquad\text{至少有一个 }a_i\ne0.
\]
$A_i$ 的次数界意味着 $m+1<d$，故式~\eqref{eq:euler-data-uniqueness-divisibility} 中 $\hbar^m$、$\hbar^{m+1}$ 的系数必须为零。若 $d+m$ 为偶数，$\hbar^m$ 的系数消失给出 $\sum_{i=0}^n2\lambda_i^sa_i=0$，
\SourcePage{369}
由此得到所有 $a_i=0$，矛盾。若 $d+m$ 为奇数，则类似地使用 $\hbar^{m+1}$ 的系数。
\end{proof}

\subsection{超几何函数与镜像变换}\label{subsec:hypergeometric-mirror-transformations}
先引入一些取等变上同调值的形式函数。令 $\mathcal S_0$ 表示序列 $B=\{B_d\}_{d=1}^{\infty}$ 所组成的集合，其中对所有 $d$，有 $B_d\in\mathcal R_GH_G^*(\mathbb P^n)$。

现在为每个 $B\in\mathcal S_0$ 关联一个形式函数，它是超几何函数的推广。
\begin{definition}\label{def:euler-data-hypergeometric-function}\index{Euler 数据!超几何函数|hyperpage}
对 $B=\{B_d\}\in\mathcal S_0$ 和可逆的 $\Omega\in\mathcal R_TH_G^*(\mathbb P^n)$，定义
\[
HG[B](t)=e^{-pt/\hbar}\left(\Omega+\sum_{d=1}^{\infty}
\frac{B_de^{dt}}{\prod_{k=0}^n\prod_{m=1}^d(p-\lambda_k-m\hbar)}\right),
\]
其中 $p\in H_G^*(\mathbb P^n)$ 为 $\mathbb P^n$ 的等变超平面类。
\end{definition}
若非等变极限 $\lambda_i\to0$ 存在，则得到形式表达式
\[
\lim_{\lambda_i\to0}HG[B](t)
=e^{-Ht/\hbar}\left(\overline\Omega+\sum_{d=1}^{\infty}
\frac{\overline B_de^{dt}}{\prod_{k=0}^n\prod_{m=1}^d(H-m\hbar)}\right),
\]
其中 $H\in H^2(\mathbb P^n)$ 为通常的超平面类，$\overline B_d=\lim_{\lambda_i\to0}B_d$，$\overline\Omega=\lim_{\lambda_i\to0}\Omega$。这是取 $H_{\mathbb C^*}^*(\mathbb P^n)\otimes\mathbb C(\hbar)$ 中形式值的函数。由于 $\mathbb C^*$ 在 $\mathbb P^n$ 上平凡作用，得到取 $H^*(\mathbb P^n)\otimes\mathbb C(\hbar)$ 中形式值的函数，此时 $\hbar$ 仅为参数。

前面几章已见过这类形式函数，\hyperref[chap:quantum-differential-equations]{第十章}中 $X=\mathbb P^n$ 的 Givental $J$ 函数就是一个明显例子。此外，下面\hyperref[subsec:lly-quintic-mirror-proof]{第~\ref*{subsec:lly-quintic-mirror-proof} 小节}将说明，形式函数 $HG[B](t)$ 在五次三维簇镜像定理的证明中起关键作用。

接下来，要把 $HG[B](t)$ 与一直研究的 Euler 数据联系起来。给定可逆的 $\Omega\in\mathcal R_TH_G^*(\mathbb P^n)$，如定义~\ref{def:omega-euler-data}，以 $\mathcal A^\Omega$ 表示所有 $\Omega$--Euler 数据的集合。还有 $\mathcal A^\Omega\subset\mathcal S$，其中 $\mathcal S$ 为序列 $Q=\{Q_d\}_{d=1}^{\infty}$ 所组成的集合，对所有 $d\ge1$，有 $Q_d\in\mathcal R_GH_G^*(N_d)$。

$\mathcal S$ 与 $\mathcal S_0$ 有何关系并不明显：后者的整个序列都位于 $\mathcal R_GH_G^*(\mathbb P^n)$ 中，而前者各项则位于不同空间的等变上同调中。幸而，有自然方法从 $\mathcal S$ 转到 $\mathcal S_0$。对于每个 $d$，有 $G$ 等变映射
\begin{equation}\label{eq:euler-data-restriction-embedding}
I_d:\mathbb P^n\longrightarrow N_d,\qquad
(a_0,\ldots,a_n)\longmapsto(a_0w_1^d,\ldots,a_nw_1^d).
\end{equation}
注意，$I_d$ 将不动点 $q_i\in\mathbb P^n$ 映到不动点 $p_{i,0}\in N_d$。定义映射 $\mathcal I:\mathcal S\to\mathcal S_0$ 为 $\mathcal I(P)_d=I_d^*(P_d)$。

通过 $\mathcal I$，任意 $\Omega$--Euler 数据 $Q$ 给出 $\mathcal I(Q)\in\mathcal S_0$，进而确定形式超几何函数 $HG[\mathcal I(Q)](t)$。在一个熟悉的例子中具体计算。
\SourceExampleNumber{11.1.5}{1}
\begin{example}\label{ex:hypergeometric-euler-data-function}
对于 $\Omega=\ell p$，有式~\eqref{eq:hypergeometric-euler-data} 定义的 $\ell p$--Euler 数据 $\widehat P$。为计算 $\mathcal I(\widehat P)$，先注意 $I_d^*(\kappa)=p$，其中 $\kappa\in H_G^*(N_d)$、$p\in H_G^*(\mathbb P^n)$ 为等变超平面类。由于 $I_d:\mathbb P^n\to N_d$ 是射影空间的等变嵌入，这一点立即成立。

由 $\widehat P$ 的定义~\eqref{eq:hypergeometric-euler-data}，得到
\[
I_d^*(\widehat P_d)=I_d^*\left(\prod_{m=0}^{\ell d}(\ell\kappa-m\hbar)\right)
=\prod_{m=0}^{\ell d}(\ell p-m\hbar).
\]
\SourcePage{370}
由此得到形式函数
\begin{equation}\label{eq:hypergeometric-euler-data-function}
HG[\mathcal I(\widehat P)](t)
=e^{-pt/\hbar}\sum_{d=0}^{\infty}
\frac{\prod_{m=0}^{\ell d}(\ell p-m\hbar)}
{\prod_{m=1}^d\prod_{k=0}^n(p-\lambda_k-m\hbar)}e^{dt},
\end{equation}
除 $\hbar$ 的符号外，它与 \SourceCite{Givental2} 中为描述 $\mathbb P^n$ 的次数 $\ell$ 超曲面的量子上同调而研究的函数 $I_{\mathcal V}$ 一致。\hyperref[subsec:lly-quintic-mirror-proof]{第~\ref*{subsec:lly-quintic-mirror-proof} 小节}将说明，当 $\mathcal V=\mathcal O_{\mathbb P^4}(5)$ 时，非等变极限 $\lim_{\lambda_i\to0}HG[\mathcal I(\widehat P)](t)$ 给出\hyperref[subsec:hypergeometric-mirror-map]{第~\ref*{subsec:hypergeometric-mirror-map} 小节}研究的五次镜像 Picard--Fuchs 方程的解 $y_0,y_1,y_2,y_3$。
\end{example}

除 $\widehat P$ 外，另一组重要的 $\ell p$--Euler 数据是定理~\ref{thm:geometric-euler-data-gluing} 中的 $\widehat Q$。此时得到形式函数
\[
HG[\mathcal I(\widehat Q)](t)
=e^{-pt/\hbar}\left(\Omega+\sum_{d=1}^{\infty}
\frac{I_d^*(\widehat Q_d)e^{dt}}
{\prod_{k=0}^n\prod_{m=1}^d(p-\lambda_k-m\hbar)}\right).
\]
可如下计算 $I_d^*(\widehat Q_d)$。
\begin{lemma}\label{lem:geometric-euler-data-hypergeometric-bound}\index{Euler 数据!次数界|hyperpage}
令 $\phi^d=\prod_{k=0}^n\prod_{m=1}^d(p-\lambda_k-m\hbar)\in H_G^*(\mathbb P^n)$，表示 $HG[\mathcal I(\widehat Q)](t)$ 中 $e^{dt}$ 系数的分母。则：
\begin{enumerate}[label=(\roman*)]
\item 若 $\pi_1,e_1$ 如图~\eqref{eq:euler-data-universal-curve}，$\mathcal L_{d,1}$ 如定理~\ref{thm:geometric-euler-data-gluing} 的证明，则
\[
I_d^*(\widehat Q_d)=\phi^d e_{1!}\left(
\frac{\pi_1^*\operatorname{Euler}_G(\mathcal V_d)}
{\hbar(\hbar+c_1^G(\mathcal L_{d,1}))}\right).
\]
\item $\deg_\hbar I_d^*(\widehat Q_d)\le(n+1)d-2$。
\end{enumerate}
\end{lemma}
\begin{proof}
为证明第 (i) 部分，只须证明，对每个 $i=0,\ldots,n$，两端在 $q_i$ 处的局部化相等。对于左端，由 $I_d(q_i)=p_{i,0}$，有
\[
i_{q_i}^*I_d^*(\widehat Q_d)=i_{p_{i,0}}^*(\widehat Q_d)=\widehat Q_d(\lambda_i).
\]
结合式~\eqref{eq:geometric-euler-data-endpoint-localization}，得到
\[
i_{q_i}^*I_d^*(\widehat Q_d)=i_{p_{i,0}}^*(\phi_{p_{i,0}})
\sum_{F_d}\int_{(F_d)_G}
\frac{\pi_1^*\operatorname{Euler}_G(\mathcal V_d)}
{a_d\operatorname{Euler}_G(N(F_d))\hbar(\hbar+c_1^G(\mathcal L_{d,1}))}.
\]
另一方面，若如式~\eqref{eq:projective-fixed-dual-class} 令 $\phi_{q_i}=\prod_{j\ne i}(p-\lambda_j)$，则式~\eqref{eq:linear-model-normal-euler} 和 $\phi^d$ 的定义立即给出
\[
i_{p_{i,0}}^*(\phi_{p_{i,0}})=i_{q_i}^*(\phi_{q_i})i_{q_i}^*(\phi^d).
\]
因此，上式成为
\begin{equation}\label{eq:geometric-euler-data-restriction-localized}
\begin{aligned}
i_{q_i}^*I_d^*(\widehat Q_d)
&=i_{q_i}^*(\phi_{q_i})i_{q_i}^*(\phi^d)\\
&\quad\times\sum_{F_d}\int_{(F_d)_G}
\frac{\pi_1^*\operatorname{Euler}_G(\mathcal V_d)}
{a_d\operatorname{Euler}_G(N(F_d))\hbar(\hbar+c_1^G(\mathcal L_{d,1}))}.
\end{aligned}
\end{equation}
转向右端，令
\[
\alpha^d=\frac{\pi_1^*\operatorname{Euler}_G(\mathcal V_d)}
{\hbar(\hbar+c_1^G(\mathcal L_{d,1}))},
\]
\SourcePage{371}
于是须计算 $\phi^d e_{1!}(\alpha^d)$ 在 $q_i$ 处的局部化。将式~\eqref{eq:geometric-euler-data-localization-integral} 适配到 $\mathbb P^n$，得到
\[
i_{q_i}^*(\phi^d e_{1!}(\alpha^d))
=\int_{(\mathbb P^n)_G}\phi_{q_i}\phi^d e_{1!}(\alpha^d)
=\int_{\overline M_{0,1}(\mathbb P^n,d)_G}e_1^*(\phi_{q_i}\phi^d)\alpha^d.
\]
用局部化计算右端积分时，因子 $\phi_{q_i}$ 意味着只须考虑在 $e_1$ 下映到 $q_i$ 的固定点集分支。这些恰好是定理~\ref{thm:geometric-euler-data-gluing} 证明中使用的分支 $\{F_d\}$；将局部化公式应用于此积分，立即得到式~\eqref{eq:geometric-euler-data-restriction-localized} 的右端。这就证明了第 (i) 部分。

对于第 (ii) 部分，由于 $\widehat Q$ 为 Euler 数据，$\widehat Q_d$、进而 $I_d^*(\widehat Q_d)$ 均为 $\hbar$ 的多项式。不过，由 $\phi^d$ 的定义，有 $\deg_\hbar\phi^d=(n+1)d$。于是，$I_d^*(\widehat Q_d)$ 的所需次数界由第 (i) 部分得到。
\end{proof}
注意，引理第二部分所证明的 $\deg_\hbar I_d^*(\widehat Q_d)$ 的界，恰好就是\IndexLocation{indexloc:0270:371}唯一性定理（定理~\ref{thm:linked-euler-data-uniqueness}）所需的次数界。

最终目标是在五次三维簇的情形下，联系 $HG[\mathcal I(\widehat P)](t)$ 与 $HG[\mathcal I(\widehat Q)](t)$。但这两个函数并不相等，Euler 数据 $\widehat P$、$\widehat Q$ 也不相等。因此，需要一种将一组 Euler 数据变成另一组的方法。这导出以下定义。
\begin{definition}\label{def:euler-data-mirror-transformation}\index{Euler 数据!镜像变换|hyperpage}
可逆映射 $\mu:\mathcal A^\Omega\to\mathcal A^\Omega$ 若对所有 $P\in\mathcal A^\Omega$，$\mu(P)$ 与 $P$ 相关联，则称它为\emph{镜像变换}。
\end{definition}
马上将看到，第~\ref{subsec:quintic-mirror-theorem-statement} 小节讨论的镜像映射，可描述为此定义意义下的镜像变换。

目标是构造一类特殊镜像变换 $\mu$，使 $HG[\mathcal I(P)](t)$ 与 $HG[\mathcal I(\mu(P))](t)$ 容易联系。先考虑超几何函数 $HG[B](t)$ 的一些初等变换。
\begin{lemma}\label{lem:hypergeometric-formal-transformation}
给定任意 $B\in\mathcal S_0$，给定任意 $f,g\in e^t\mathcal R_T[[e^t]]$。则存在唯一 $\widetilde B\in\mathcal S_0$，使
\[
e^{f/\hbar}HG[B](t+g)=HG[\widetilde B](t).
\]
\end{lemma}
\begin{proof}
这是直接的形式幂级数计算，见 \SourceCite{LLY}。
\end{proof}
利用此引理，可描述我们关心的镜像变换类。
\begin{proposition}\label{prop:hypergeometric-mirror-transformation-lift}\index{Euler 数据!Lagrange 提升|hyperpage}
令 $f=a\hbar+b$，其中 $a,b,g\in e^t\mathcal R_T[[e^t]]$。则存在镜像变换 $\mu:\mathcal A^\Omega\to\mathcal A^\Omega$，使对所有 Euler 数据 $P\in\mathcal A^\Omega$，有
\[
e^{f/\hbar}HG[\mathcal I(P)](t+g)=HG[\mathcal I(\mu(P))](t).
\]
\end{proposition}
\begin{proof}
给定 $f,g$，引理~\ref{lem:hypergeometric-formal-transformation} 中的映射 $B\mapsto\widetilde B$ 定义变换 $\mu_0:\mathcal S_0\to\mathcal S_0$。用此记号，前一引理意味着，对 $B\in\mathcal S_0$，有
\[
e^{f/\hbar}HG[B](t+g)=HG[\mu_0(B)](t).
\]
为证明命题，只须证明 $\mu_0$ 可提升为满足 $\mathcal I(\mu(P))=\mu_0(\mathcal I(P))$ 的镜像变换 $\mu$。

\SourcePage{372}
回顾 $\mathcal I:\mathcal S\to\mathcal S_0$ 由 $\mathcal I(P)_d=I_d^*(P_d)$ 定义，$I_d$ 如式~\eqref{eq:euler-data-restriction-embedding}。为定义 $\mu$，须构造反方向的映射，做法如下。命题~\ref{prop:atiyah-bott-localization} 对 $N_d$ 所给的同构可写成
\[
\mathcal R_GH_G^*(N_d)\xrightarrow{\ \bigoplus i_{p_{i,r}}^*\ }
(\mathcal R_G)^{(N_d)^G},
\]
其中 $(N_d)^G$ 如往常一样表示 $G$ 固定点集。因此，为指定元素 $\omega\in\mathcal R_GH_G^*(N_d)$，只须指定所有 $i,r$ 对应的限制 $i_{p_{i,r}}^*(\omega)$。据此，用以下条件唯一定义 $\mathcal L:\mathcal S_0\to\mathcal S$：
\begin{equation}\label{eq:euler-data-lagrange-lift}
i_{p_{i,r}}^*(\mathcal L(B)_d)
=(i_{q_i}^*(\Omega))^{-1}\overline{i_{q_i}^*(B_r)}\,i_{q_i}^*(B_{d-r}),
\qquad0\le i\le n,\quad0\le r\le d,
\end{equation}
对任意 $B\in\mathcal S_0$ 均如此。这些等式中，$B_0=\Omega$。

映射 $\mathcal L$ 有一些良好性质。首先，$\mathcal I\circ\mathcal L=\operatorname{id}_{\mathcal S_0}$。为证明这一点，注意 $I_d(q_i)=p_{i,0}$；式~\eqref{eq:euler-data-lagrange-lift} 在 $r=0$ 时给出
\[
\begin{aligned}
i_{q_i}^*(\mathcal I(\mathcal L(B))_d)
&=i_{q_i}^*\bigl(I_d^*(\mathcal L(B)_d)\bigr)\\
&=i_{p_{i,0}}^*(\mathcal L(B)_d)\\
&=(i_{q_i}^*(\Omega))^{-1}\overline{i_{q_i}^*(B_0)}\,i_{q_i}^*(B_d)
=i_{q_i}^*(B_d),
\end{aligned}
\]
最后一个等号来自 $B_0=\Omega\in\mathcal R_TH_T^*(\mathbb P^n)$。由 $\mathcal R_TH_T^*(\mathbb P^n)$ 的局部化（并保留形式参数 $\hbar$），得到 $\mathcal I(\mathcal L(B))=B$。$\mathcal L$ 的第二个良好性质是：对任意 $B\in\mathcal S_0$，$\mathcal L(B)$ 满足定义~\ref{def:omega-euler-data} 中的 Euler 恒等式。因此，若所有 $d$ 都有 $\mathcal L(B)_d\in\mathcal R_TH_G^*(N_d)$，则 $\mathcal L(B)$ 是 Euler 数据。最后，由定义也立即可知，若 $Q\in\mathcal A^\Omega\subset\mathcal S$，则 $Q=\mathcal L\circ\mathcal I(Q)$。

既然有了 $\mathcal L:\mathcal S_0\to\mathcal S$ 和 $\mathcal I:\mathcal S\to\mathcal S_0$，映射 $\mu_0:\mathcal S_0\to\mathcal S_0$ 可延伸为 $\mu=\mathcal L\circ\mu_0\circ\mathcal I:\mathcal S\to\mathcal S$。用 \SourceCite{LLY} 的术语，$\mu$ 是 $\mu_0$ 的 \emph{Lagrange 提升}。如此定义 $\mu$ 后，$\mathcal I\circ\mu=\mu_0\circ\mathcal I$ 立即来自 $\mathcal I\circ\mathcal L=\operatorname{id}_{\mathcal S_0}$。因此，还须证明 $\mu$ 是镜像变换。给定 $P\in\mathcal A^\Omega$，$P$ 与 $\mu(P)$ 相关联这一条件，即 $i_{p_{i,0}}^*(P_d-\mu(P)_d)$ 在 $\hbar=(\lambda_i-\lambda_j)/d$ 处为零，立即来自证明引理~\ref{lem:hypergeometric-formal-transformation} 时所作的显式幂级数计算。最后，须证明 $\mu(P)$ 是 $\Omega$--Euler 数据。前一段表明，$\mu(P)$ 自动满足 Euler 条件，因此最后一步是证明 $\mu(P)_d\in\mathcal R_TH_G^*(N_d)$。证明见 \SourceCite[引理 2.15]{LLY}。
\end{proof}

现在几乎已有证明\IndexLocation{indexloc:0864:372}镜像定理所需的全部准备。不过，在开始证明之前，应解释\hyperref[subsec:lly-euler-data]{第~\ref*{subsec:lly-euler-data}}、\hyperref[subsec:linked-euler-data]{\ref*{subsec:linked-euler-data}}、\hyperref[subsec:hypergeometric-mirror-transformations]{\ref*{subsec:hypergeometric-mirror-transformations} 小节}与 \SourceCite{LLY} 的关系。一直使用的丛 $\mathcal O_{\mathbb P^n}(\ell)$ 是\emph{凸丛}\index{凸丛|hyperpage}的特殊情形；凸丛是指一个 $T$ 等变丛 $\mathcal V$，其 $\operatorname{Euler}_T(\mathcal V)$ 在 $\mathcal R_TH_T^*(\mathbb P^n)$ 中可逆，且对每个无标记亏格零稳定映射 $f:C\to\mathbb P^n$，有 $H^1(C,f^*\mathcal V)=0$。凸丛的其他例子有 $\mathcal V=\bigoplus_i\mathcal O_{\mathbb P^n}(\ell_i)$，其中 $\ell_i>0$。类似地，\SourceCite{LLY} 将满足相应条件、但要求对每个无标记亏格零稳定映射 $f$ 有 $H^0(C,f^*\mathcal V)=0$ 的丛 $\mathcal V$ 称为\emph{凹丛}\index{凹丛|hyperpage}。最后，凸丛与凹丛的直和称为\emph{凸凹丛}\index{凸凹丛|hyperpage}。

给定 $\mathbb P^n$ 上的凸凹丛 $\mathcal V$，如下定义 $\overline M_{0,0}(\mathbb P^n,d)$ 上的丛 $\mathcal V_d$。若 $\mathcal V$ 凸，则如图~\eqref{eq:euler-data-universal-curve} 后的讨论一样，取 $\mathcal V_d=\pi_{1*}e_1^*(\mathcal V)$。若 $\mathcal V$ 凹，则取 $\mathcal V_d=R^1\pi_{1*}e_1^*(\mathcal V)$。若 $\mathcal V$ 凸凹，则 $\mathcal V_d$ 为将这两种构造分别应用于凸部分、凹部分后所得丛的直和。

最后，设 $\mathcal V$ 为形如 $\bigoplus_i\mathcal O_{\mathbb P^n}(\ell_i)$（$\ell_i\ne0$）的凸凹丛。用 $\mathcal V_d$，仍如式~\eqref{eq:geometric-euler-data} 一样定义 $\widehat Q_d$。此外，还可
\SourcePage{373}
推广式~\eqref{eq:hypergeometric-euler-data}，对 $d>0$ 定义
\begin{equation}\label{eq:concavex-hypergeometric-euler-data}
\widehat P_d=
\prod_{\ell_i>0}\prod_{m=0}^{\ell_i d}(\ell_i\kappa-m\hbar)
\times\prod_{\ell_i<0}\prod_{m=1}^{-\ell_i d-1}(\ell_i\kappa+m\hbar).
\end{equation}
关键在于，对于 $\mathcal V=\mathcal O_{\mathbb P^n}(\ell)$ 已证明的 $\widehat Q$、$\widehat P$ 的结果，在 $\mathcal V=\bigoplus_i\mathcal O_{\mathbb P^n}(\ell_i)$ 时仍成立；此时取 $\Omega=\operatorname{Euler}_T(\mathcal V^+)/\operatorname{Euler}_T(\mathcal V^-)$，其中 $\mathcal V^+$、$\mathcal V^-$ 为凸、凹部分，。细节见 \SourceCite{LLY}。
% 第八批接续：原 PDF 物理页373中部。
\subsection{临界丛}\label{subsec:critical-bundles}\index{临界丛|hyperpage}
\hypertarget{source:373-batch8}{}
在证明五次三维簇 $V\subset\mathbb P^4$ 的镜像定理之前，需要解释 $V$ 的 Gromov--Witten 不变量 $K_d=\langle I_{0,0,d}\rangle$ 与此前讨论的联系。定理~\ref{thm:critical-quintic-hypergeometric-potential} 将表明，对丛 $\mathcal V=\mathcal O_{\mathbb P^4}(5)$，$K_d$ 自然出现在超几何函数 $HG[\mathcal I(\widehat Q)](t)$ 中。还将讨论凸凹丛的一类子类，即\emph{临界丛}；它们满足该定理的一个修改版本。

回顾\hyperref[subsec:quintic-mirror-theorem-statement]{第~\ref*{subsec:quintic-mirror-theorem-statement} 小节}，$V$ 的 Gromov--Witten 不变量 $K_d$ 由
\[
K_d=\int_{\overline M_{0,0}(\mathbb P^4,d)}\operatorname{Euler}(\mathcal V_d)
\]
给出。用 $K_d$ 可将五次三维簇的 Gromov--Witten 势 $\Phi$ 写成
\[
\Phi=\frac{5t^3}{6}+\sum_{d=1}^{\infty}K_de^{dt}.
\]
于是有如下结果。
\begin{theorem}\label{thm:critical-quintic-hypergeometric-potential}
令 $\mathcal V=\mathcal O_{\mathbb P^4}(5)$，并如式~\eqref{eq:geometric-euler-data} 定义 $\widehat Q_d$。则在非等变极限 $\lambda_i\to0$ 下，有
\[
\lim_{\lambda_i\to0}HG[\mathcal I(\widehat Q)](t)
=5H\left(1-t\frac H\hbar+\frac{\Phi'}5\frac{H^2}{\hbar^2}
-\frac{t\Phi'-2\Phi}5\frac{H^3}{\hbar^3}\right),
\]
其中 $H$ 为 $H^*(\mathbb P^4)$ 中的超平面类。
\end{theorem}
\begin{proof}
定义~\ref{def:euler-data-hypergeometric-function} 的非等变极限给出
\[
\lim_{\lambda_i\to0}HG[\mathcal I(\widehat Q)](t)
=e^{-Ht/\hbar}\left(5H+\sum_{d=1}^{\infty}
\frac{\lim_{\lambda_i\to0}I_d^*(\widehat Q_d)}{\prod_{m=1}^d(H-m\hbar)^5}e^{dt}\right).
\]
现在取此表达式中 $e^{dt}$ 的系数，并用 $H^*(\mathbb P^4)$ 的基 $1,H,\ldots,H^4$ 表示：
\[
e^{-Ht/\hbar}\frac{\lim_{\lambda_i\to0}I_d^*(\widehat Q_d)}{\prod_{m=1}^d(H-m\hbar)^5}
=A_0+A_1H+A_2H^2+A_3H^3+A_4H^4.
\]
将此等式乘以 $H^k$ 并在 $\mathbb P^4$ 上积分，即可求出 $A_i$。这给出
\[
A_{4-k}=\int_{\mathbb P^4}H^ke^{-Ht/\hbar}
\frac{\lim_{\lambda_i\to0}I_d^*(\widehat Q_d)}{\prod_{m=1}^d(H-m\hbar)^5}.
\]
\SourcePage{374}
以下公式如何推出本定理，留给读者验证：
\begin{equation}\label{eq:critical-quintic-integral-coefficients}
\int_{\mathbb P^4}H^ke^{-Ht/\hbar}
\frac{\lim_{\lambda_i\to0}I_d^*(\widehat Q_d)}{\prod_{m=1}^d(H-m\hbar)^5}
=\begin{cases}
\hbar^{-3}(2-dt)K_d,&k=0,\\
\hbar^{-2}dK_d,&k=1,\\
0,&k>1.
\end{cases}
\end{equation}
式~\eqref{eq:critical-quintic-integral-coefficients} 的 $k=0$ 情形在 \SourceCite{LLY} 中明确给出，而 $k\ge1$ 的情形则隐含于其中。

为证明式~\eqref{eq:critical-quintic-integral-coefficients}，使用引理~\ref{lem:geometric-euler-data-hypergeometric-bound} 给出的 $I_d^*(\widehat Q_d)$ 公式。令 $\lambda_i\to0$，取其非等变极限，得到
\[
\lim_{\lambda_i\to0}I_d^*(\widehat Q_d)
=\prod_{m=1}^d(H-m\hbar)^5\,
 e_{1!}\left(\frac{\pi_1^*\operatorname{Euler}(\mathcal V_d)}
 {\hbar(\hbar+c_1(\mathcal L_{d,1}))}\right).
\]
将其代入式~\eqref{eq:critical-quintic-integral-coefficients} 左端的积分，积分简化为
\begin{equation}\label{eq:critical-quintic-pushforward-integral}
\begin{aligned}
&\int_{\mathbb P^4}H^ke^{-Ht/\hbar}
 e_{1!}\left(\frac{\pi_1^*\operatorname{Euler}(\mathcal V_d)}
 {\hbar(\hbar+c_1(\mathcal L_{d,1}))}\right)\\
&\quad=\int_{\overline M_{0,1}(\mathbb P^4,d)}
 e_1^*H^k e^{-e_1^*Ht/\hbar}
 \frac{\pi_1^*\operatorname{Euler}(\mathcal V_d)}{\hbar(\hbar+c_1(\mathcal L_{d,1}))}\\
&\quad=\int_{\overline M_{0,0}(\mathbb P^4,d)}\operatorname{Euler}(\mathcal V_d)\,
 \pi_{1!}\left(e_1^*H^k\frac{e^{-e_1^*Ht/\hbar}}
 {\hbar(\hbar+c_1(\mathcal L_{d,1}))}\right).
\end{aligned}
\end{equation}
因为 $\operatorname{Euler}(\mathcal V_d)$ 的次数为最高次，这个积分恰为 $K_d$ 乘以被积式中第二个因子的零次分量。又因为 $\pi_1$ 的相对维数为 $1$，该分量就是
\[
\begin{cases}
\pi_{1!}\bigl(\hbar^{-3}(-e_1^*Ht-c_1(\mathcal L_{d,1}))\bigr),&k=0,\\
\pi_{1!}(\hbar^{-2}e_1^*H),&k=1,\\
0,&k>1
\end{cases}
\]
的零次值。在 $\pi_1$ 的一条纤维上积分后，它化为
\[
\begin{cases}
\hbar^{-3}(-dt+2),&k=0,\\
\hbar^{-2}d,&k=1,\\
0,&k>1.
\end{cases}
\]
这就得到式~\eqref{eq:critical-quintic-integral-coefficients}，从而证明了定理。
\end{proof}

事实上，这个定理在更一般的情形下也成立。考察上述证明可知，关键步骤之一是式~\eqref{eq:critical-quintic-pushforward-integral} 后的论证，其中用到了 $\operatorname{Euler}(\mathcal V_d)$ 在 $\overline M_{0,0}(\mathbb P^4,d)$ 上为最高次类这一事实。这导出如下定义。
\begin{definition}\label{def:critical-bundle}
若 $\mathbb P^n$ 上的凸凹丛 $\mathcal V$ 所诱导的丛 $\mathcal V_d$ 在 $\overline M_{0,0}(\mathbb P^n,d)$ 上，对所有 $d>0$ 都满足
\[
\operatorname{rank}\mathcal V_d=(n+1)d+n-3=\dim\overline M_{0,0}(\mathbb P^n,d),
\]
则称 $\mathcal V$ 为\emph{临界丛}。
\end{definition}
\SourcePage{375}
当 $\mathcal V$ 为临界丛时，$\operatorname{Euler}(\mathcal V_d)$ 是 $\overline M_{0,0}(\mathbb P^n,d)$ 上的最高次类，因而可以定义
\[
K_d=\int_{\overline M_{0,0}(\mathbb P^n,d)}\operatorname{Euler}(\mathcal V_d),
\]
推广五次三维簇情形的构造。在此情形下，由于 $\mathcal V$ 凸凹，引理~\ref{lem:geometric-euler-data-hypergeometric-bound} 仍适用；由于 $\mathcal V$ 临界，定理~\ref{thm:critical-quintic-hypergeometric-potential} 的论证也适用。因此，已证明
\begin{equation}\label{eq:critical-bundle-integral-coefficients}
\int_{\mathbb P^n}H^ke^{-Ht/\hbar}
\frac{\lim_{\lambda_i\to0}I_d^*(\widehat Q_d)}{\prod_{m=1}^d(H-m\hbar)^{n+1}}
=\begin{cases}
\hbar^{-3}(2-dt)K_d,&k=0,\\
\hbar^{-2}dK_d,&k=1,\\
0,&k>1.
\end{cases}
\end{equation}
细节见 \SourceCite{LLY}。下面将看到，这有一些有趣的应用。

\hyperref[subsec:lly-euler-data]{第~\ref*{subsec:lly-euler-data} 小节}考虑了形如 $\mathcal V=\mathcal O_{\mathbb P^n}(\ell)$ 的丛。读者容易验证，此类型中唯一的临界丛是 $\mathcal V=\mathcal O_{\mathbb P^4}(5)$。不过，对于 $\mathcal V=\bigoplus_i\mathcal O_{\mathbb P^n}(\ell_i)$，还有其他值得研究的临界丛。以下例子与我们的讨论尤其相关。
\SourceExampleNumber{11.1.6}{1}
\begin{example}\label{ex:critical-multiple-cover-bundle}
考虑 $\mathbb P^1$ 上的凹丛
\[
\mathcal V=\mathcal O_{\mathbb P^1}(-1)\oplus\mathcal O_{\mathbb P^1}(-1).
\]
这里，$\mathcal V_d$ 定义为
\[
\mathcal V_d=R^1\pi_{1*}e_1^*\bigl(\mathcal O_{\mathbb P^1}(-1)\oplus\mathcal O_{\mathbb P^1}(-1)\bigr).
\]
对代表 $\overline M_{0,0}(\mathbb P^1,d)$ 中一点的映射 $f:C\to\mathbb P^1$，容易验证
\[
\begin{aligned}
\operatorname{rank}\mathcal V_d
&=\dim H^1\bigl(C,f^*\mathcal O_{\mathbb P^1}(-1)\oplus f^*\mathcal O_{\mathbb P^1}(-1)\bigr)\\
&=2d-2=\dim\overline M_{0,0}(\mathbb P^1,d).
\end{aligned}
\]
因此 $\mathcal V$ 为临界丛。于是式~\eqref{eq:critical-bundle-integral-coefficients} 表明，$HG[\mathcal I(\widehat Q)](t)$ 的非等变极限涉及
\[
K_d=\int_{\overline M_{0,0}(\mathbb P^1,d)}
\operatorname{Euler}\!\left(R^1\pi_{1*}e_1^*
\bigl(\mathcal O_{\mathbb P^1}(-1)\oplus\mathcal O_{\mathbb P^1}(-1)\bigr)\right).
\]
这恰是定理~\ref{thm:multiple-cover-localization} 计算的积分；该定理计算了 Calabi--Yau 三维簇中刚性嵌入的光滑有理曲线的\IndexLocation{indexloc:0916:375}多重覆盖贡献。即将应用 \SourceCite{LLY} 的方法，给出该定理的第二个证明。
\end{example}
还应指出，\SourceCite{LLY} 给出了 $\mathbb P^n$ 上所有由线丛直和构成的临界丛的完整清单。
\subsection{定理~\ref{thm:lly-quintic-mirror} 的证明}\label{subsec:lly-quintic-mirror-proof}
现在可以证明\IndexLocation{indexloc:0867:375}五次三维簇的镜像定理了。在本节余下部分，均假设 $\mathcal V=\mathcal O_{\mathbb P^4}(5)$，并使用\hyperref[subsec:lly-euler-data]{第~\ref*{subsec:lly-euler-data} 小节}定义的 $\widehat P$、$\widehat Q$。证明的关键是构造两个\IndexLocation{indexloc:0268:375}镜像变换 $\mu$、$\nu$，使 $\mu(\widehat P)=\nu(\widehat Q)$。这将导出超几何函数 $HG[\mathcal I(\widehat P)](t)$ 与 $HG[\mathcal I(\widehat Q)](t)$ 之间的关系；取非等变极限 $\lambda_i\to0$，便得到所需结果。
\SourcePage{376}
为构造所需的镜像变换，首先更仔细地考察 $HG[\mathcal I(\widehat P)](t)$。取式~\eqref{eq:hypergeometric-euler-data-function} 在 $\lambda_i\to0$ 时的非等变极限，得到
\begin{equation}\label{eq:quintic-hypergeometric-period-expansion}
\begin{aligned}
\lim_{\lambda_i\to0}HG[\mathcal I(\widehat P)](t)
&=e^{-Ht/\hbar}\sum_{d=0}^{\infty}
\frac{\prod_{m=0}^{5d}(5H-m\hbar)}{\prod_{m=1}^d(H-m\hbar)^5}e^{dt}\\
&=5H\left(y_0-y_1\frac H\hbar+y_2\frac{H^2}{\hbar^2}-y_3\frac{H^3}{\hbar^3}\right).
\end{aligned}
\end{equation}
这些就是\hyperref[subsec:quintic-mirror-theorem-statement]{第~\ref*{subsec:quintic-mirror-theorem-statement} 小节}讨论的函数 $y_0,y_1,y_2,y_3$；它们出现在镜像映射 $\Psi(t)=y_1/y_0$ 中，也出现在定理~\ref{thm:lly-quintic-mirror} 的陈述中。

现在两次使用命题~\ref{prop:hypergeometric-mirror-transformation-lift}，构造镜像变换 $\mu$、$\nu$。对于 $\mu$，取
\begin{equation}\label{eq:quintic-first-mirror-transformation-parameters}
\begin{gathered}
f_1=-(\log y_0)\hbar+\frac{h_2}{y_0}\sum_{k=0}^4\lambda_k,\\
h_2=\sum_{d=1}^{\infty}\frac{(5d)!}{(d!)^5}
\left(\sum_{m=1}^d\frac1m\right)e^{dt},
\end{gathered}
\end{equation}
并取 $g_1=0$，随后命题~\ref{prop:hypergeometric-mirror-transformation-lift} 给出镜像变换 $\mu$，使
\begin{equation}\label{eq:quintic-first-mirror-transformation}
e^{f_1/\hbar}HG[\mathcal I(\widehat P)](t)
=HG[\mathcal I(\mu(\widehat P))](t).
\end{equation}
对于第二个镜像变换，取 $f_2=0$、$g_2=\Psi(t)-t\in e^t\mathcal R[[e^t]]$，其中 $\Psi(t)=y_1/y_0$ 是镜像映射。再次应用上述命题，得到镜像变换 $\nu$，使
\begin{equation}\label{eq:quintic-second-mirror-transformation}
HG[\mathcal I(\nu(\widehat Q))](t)
=HG[\mathcal I(\widehat Q)](\Psi(t)).
\end{equation}
镜像变换所得数据 $\mu(\widehat P)$ 与 $\nu(\widehat Q)$ 有如下关系。
\begin{theorem}\label{thm:quintic-transformed-euler-data-equality}
$\mu(\widehat P)=\nu(\widehat Q)$。
\end{theorem}
\begin{proof}
遵循 \SourceCite{LLY} 的证明。由定理~\ref{thm:linked-geometric-hypergeometric-data}，$\widehat P$、$\widehat Q$ 是相关联的 Euler 数据。此外，因为 $\mu$、$\nu$ 是镜像变换，$\widehat P,\mu(\widehat P)$ 及 $\widehat Q,\nu(\widehat Q)$ 也分别相关联。由传递性，$\mu(\widehat P)$、$\nu(\widehat Q)$ 是相关联的 Euler 数据。因此，只要证明次数界
\[
\deg_{\hbar}i_{p_{i,0}}^*\bigl(\mu(\widehat P)_d-\nu(\widehat Q)_d\bigr)\le5d-2,
\]
本定理就由定理~\ref{thm:linked-euler-data-uniqueness} 得到。由于 $i_{p_{i,0}}=I_d\circ i_{q_i}$，只须证明
\[
\frac{I_d^*(\mu(\widehat P)_d-\nu(\widehat Q)_d)}
{\prod_{k=0}^4\prod_{m=1}^d(p-\lambda_k-m\hbar)}
=0\pmod{\hbar^{-2}}.
\]
不过，由定义~\ref{def:euler-data-hypergeometric-function}，这等价于
\[
HG[\mathcal I(\nu(\widehat Q))](t)
=HG[\mathcal I(\mu(\widehat P))](t)\pmod{\hbar^{-2}}.
\]
再用式~\eqref{eq:quintic-first-mirror-transformation}、\eqref{eq:quintic-second-mirror-transformation}，问题归结为证明
\begin{equation}\label{eq:quintic-transformed-hypergeometric-congruence}
e^{f_1/\hbar}HG[\mathcal I(\widehat P)](t)
=HG[\mathcal I(\widehat Q)](\Psi(t))\pmod{\hbar^{-2}}.
\end{equation}
为此，首先按 $\hbar^{-1}$ 的幂展开 $HG[\mathcal I(\widehat P)](t)$，得到
\[
HG[\mathcal I(\widehat P)](t)
=5p\left(y_0-\hbar^{-1}\left(py_1+h_2\sum_{k=0}^4\lambda_k\right)+O(\hbar^{-2})\right).
\]
于是，由式~\eqref{eq:quintic-first-mirror-transformation-parameters} 及一些代数运算可得
\[
e^{f_1/\hbar}HG[\mathcal I(\widehat P)](t)
=5p-\hbar^{-1}5p^2\frac{y_1}{y_0}+O(\hbar^{-2}).
\]
\SourcePage{377}
再看式~\eqref{eq:quintic-transformed-hypergeometric-congruence} 的右端。首先，由引理~\ref{lem:geometric-euler-data-hypergeometric-bound} 的第 (ii) 部分，有 $\deg_{\hbar}I_d^*(\widehat Q_d)\le5d-2$。于是定义~\ref{def:euler-data-hypergeometric-function} 给出
\[
\begin{aligned}
HG[\mathcal I(\widehat Q)](t)
&=e^{-pt/\hbar}\bigl(5p+O(\hbar^{-2})\bigr)\\
&=5p-\hbar^{-1}5p^2t+O(\hbar^{-2}).
\end{aligned}
\]
因为 $\Psi=y_1/y_0$，有
\[
HG[\mathcal I(\widehat Q)](\Psi(t))
=5p-\hbar^{-1}5p^2\frac{y_1}{y_0}+O(\hbar^{-2}).
\]
比较上述 $e^{f_1/\hbar}HG[\mathcal I(\widehat P)](t)$ 与 $HG[\mathcal I(\widehat Q)](\Psi(t))$ 的表达式，就立即得到所需的式~\eqref{eq:quintic-transformed-hypergeometric-congruence}。定理得证。
\end{proof}
注意，以上证明给出的 $HG[\mathcal I(\widehat P)](t)$ 与 $HG[\mathcal I(\widehat Q)](t)$ 的展开，说明了选择镜像变换 $\mu$、$\nu$ 的依据：基本想法是确保同余式~\eqref{eq:quintic-transformed-hypergeometric-congruence} 成立。
\begin{corollary}\label{cor:quintic-hypergeometric-mirror-equality}
在非等变极限 $\lambda_i\to0$ 下，有
\[
\lim_{\lambda_i\to0}HG[\mathcal I(\widehat Q)](\Psi(t))
=\frac1{y_0}\lim_{\lambda_i\to0}HG[\mathcal I(\widehat P)](t).
\]
\end{corollary}
\begin{proof}
将定理~\ref{thm:quintic-transformed-euler-data-equality} 与式~\eqref{eq:quintic-first-mirror-transformation}、\eqref{eq:quintic-second-mirror-transformation} 合并，并取非等变极限，得到
\[
\lim_{\lambda_i\to0}HG[\mathcal I(\widehat Q)](\Psi(t))
=\lim_{\lambda_i\to0}e^{f_1/\hbar}HG[\mathcal I(\widehat P)](t).
\]
不过，由式~\eqref{eq:quintic-first-mirror-transformation-parameters}，知道
$f_1=-(\log y_0)\hbar+(h_2/y_0)\sum_{k=0}^4\lambda_k$。这给出 $\lim_{\lambda_i\to0}e^{f_1/\hbar}=1/y_0$，从而得到推论。
\end{proof}

\begin{proof}[定理~\ref{thm:lly-quintic-mirror} 的证明]
现在容易证明定理~\ref{thm:lly-quintic-mirror}。由定理~\ref{thm:critical-quintic-hypergeometric-potential}，有
\[
\begin{aligned}
&\lim_{\lambda_i\to0}HG[\mathcal I(\widehat Q)](\Psi(t))\\
&\quad=5H\left(1-\Psi(t)\frac H\hbar+\frac{\Phi'(\Psi(t))}{5}\frac{H^2}{\hbar^2}
-\frac{\Psi(t)\Phi'(\Psi(t))-2\Phi(\Psi(t))}{5}\frac{H^3}{\hbar^3}\right),
\end{aligned}
\]
又由式~\eqref{eq:quintic-hypergeometric-period-expansion}，有
\[
\lim_{\lambda_i\to0}HG[\mathcal I(\widehat P)](t)
=5H\left(y_0-y_1\frac H\hbar+y_2\frac{H^2}{\hbar^2}-y_3\frac{H^3}{\hbar^3}\right).
\]
把这些等式与推论~\ref{cor:quintic-hypergeometric-mirror-equality} 合并，得到
\[
\begin{aligned}
&5H\left(1-\Psi(t)\frac H\hbar+\frac{\Phi'(\Psi(t))}{5}\frac{H^2}{\hbar^2}
-\frac{\Psi(t)\Phi'(\Psi(t))-2\Phi(\Psi(t))}{5}\frac{H^3}{\hbar^3}\right)\\
&\qquad=5H\left(1-\frac{y_1}{y_0}\frac H\hbar
+\frac{y_2}{y_0}\frac{H^2}{\hbar^2}-\frac{y_3}{y_0}\frac{H^3}{\hbar^3}\right).
\end{aligned}
\]
比较系数，简单计算即得
\[
\Phi(\Psi(t))=\frac52\left(\frac{y_1}{y_0}\frac{y_2}{y_0}-\frac{y_3}{y_0}\right).
\]
这完成了五次三维簇\IndexLocation{indexloc:0864:377}\IndexLocation{indexloc:0867:377}镜像定理的证明。
\end{proof}
\SourcePage{378}
刚才的证明十分漂亮；若再认识到，本节先前的工作还能轻而易举地给出其他重要结果，就更能体会它的价值。\SourceCite{LLY} 给出了若干这样的例子，这里只介绍一个与我们的目标有关的例子。
\SourceExampleNumber{11.1.7}{1}
\begin{example}\label{ex:lly-multiple-cover-proof}
如例题~\ref{ex:critical-multiple-cover-bundle}，考虑 $\mathcal V=\mathcal O_{\mathbb P^1}(-1)\oplus\mathcal O_{\mathbb P^1}(-1)$。由式~\eqref{eq:critical-bundle-integral-coefficients}，知道
\begin{equation}\label{eq:lly-multiple-cover-integral}
\int_{\mathbb P^1}e^{-Ht/\hbar}
\frac{\lim_{\lambda_i\to0}I_d^*(\widehat Q_d)}{\prod_{m=1}^d(H-m\hbar)^2}
=\hbar^{-3}(2-dt)K_d,
\end{equation}
其中
\[
K_d=\int_{\overline M_{0,0}(\mathbb P^1,d)}
\operatorname{Euler}\!\left(R^1\pi_{1*}e_1^*
\bigl(\mathcal O_{\mathbb P^1}(-1)\oplus\mathcal O_{\mathbb P^1}(-1)\bigr)\right).
\]
可以如下计算 $K_d$。已知 $\widehat P$、$\widehat Q$ 是相关联的 Euler 数据，而引理~\ref{lem:geometric-euler-data-hypergeometric-bound} 给出 $\deg_{\hbar}I_d^*(\widehat Q_d)\le2d-2$。不过，在此情形下，式~\eqref{eq:concavex-hypergeometric-euler-data} 表明
\[
\widehat P_d=\prod_{m=1}^{d-1}(-\kappa+m\hbar)^2
=\prod_{m=1}^{d-1}(\kappa-m\hbar)^2.
\]
因此 $I_d^*(\widehat P_d)=\prod_{m=1}^{d-1}(p-m\hbar)^2$，其关于 $\hbar$ 的次数也不超过 $2d-2$。所以，此情形满足定理~\ref{thm:linked-euler-data-uniqueness} 的次数条件，从而 $\widehat P=\widehat Q$。于是
\[
\begin{aligned}
&\int_{\mathbb P^1}e^{-Ht/\hbar}
\frac{\lim_{\lambda_i\to0}I_d^*(\widehat Q_d)}{\prod_{m=1}^d(H-m\hbar)^2}\\
&\quad=\int_{\mathbb P^1}e^{-Ht/\hbar}
\frac{\lim_{\lambda_i\to0}I_d^*(\widehat P_d)}{\prod_{m=1}^d(H-m\hbar)^2}\\
&\quad=\int_{\mathbb P^1}e^{-Ht/\hbar}
\frac{\prod_{m=1}^{d-1}(H-m\hbar)^2}{\prod_{m=1}^d(H-m\hbar)^2}\\
&\quad=\int_{\mathbb P^1}e^{-Ht/\hbar}\frac1{(H-d\hbar)^2}
=\hbar^{-3}(2-dt)d^{-3},
\end{aligned}
\]
最后一个等号利用了 $H^*(\mathbb P^1)$ 中的 $H^2=0$。与式~\eqref{eq:lly-multiple-cover-integral} 比较，即得 $K_d=d^{-3}$；这完成了定理~\ref{thm:multiple-cover-localization} 的第二个证明。
\end{example}

这个例子表明，有必要拓宽我们对“镜像定理”的认识：导出五次三维簇镜像对称的同样思想与技巧，也会同时给出其他有趣结果。论文 \SourceCite{LLY} 是朝理解镜像对称的更广阔背景迈出的重要一步。

下一节将看到，与 \SourceCite{LLY} 相比，Givental 的镜像定理方法具有相似而又不同的结构。他考虑两个取上同调值的形式函数 $I_{\mathcal V}$（式~\eqref{eq:givental-complete-intersection-i-function}）和 $J_{\mathcal V}$（式~\eqref{eq:givental-general-j-function}），然后作变量变换，使一个唯一性定理适用。参见定理~\ref{thm:givental-projective-mirror} 及 \SourceCite{Givental2}。例题~\ref{ex:givental-quintic-mirror-potential} 将以五次三维簇说明此方法；\hyperref[subsec:givental-general-bundle-proof]{第~\ref*{subsec:givental-general-bundle-proof} 小节}末附近还会进一步比较这两种方法。
\section{Givental 的方法}\label{sec:givental-mirror-approach}
\SourcePage{379}
本节解释 Givental 对射影空间中完全交的镜像定理的表述（定理~\ref{thm:givental-projective-mirror}）。其方法涉及构造取形式上同调值的函数 $I_{\mathcal V}$、$J_{\mathcal V}$，其中 $J_{\mathcal V}$ 是\hyperref[sec:quantum-cohomology-relations]{第~\ref*{sec:quantum-cohomology-relations} 节}讨论的函数 $J$ 的一个变体。镜像定理则断言，作变量变换后，$I_{\mathcal V}=J_{\mathcal V}$。

\SourceCite{BDPP,Pandharipande3} 更详细地解释了 \SourceCite{Givental2} 的方法，也可作为本节的有用参考文献。五次簇的情形见 \SourceCite{Givental6}。这里将综合 \SourceCite{Givental2,Givental4,Kim2,Pandharipande3} 中的思想，概述 Givental 关于 $\mathbb P^n$ 中 nef 完全交的定理的一个略有不同的证明；读者可参考这些文献。\SourceCite{Givental2} 引入的一些所需思想，已在\hyperref[chap:localization]{第九章}和\hyperref[chap:quantum-differential-equations]{第十章}解释过。还将讨论环面完全交，并给出一些说明性的例子。
\subsection{射影空间中 nef 完全交的镜像定理}\label{subsec:givental-nef-projective-mirror}
考虑 $\mathbb P^n$ 中由 $\ell$ 个次数分别为 $a_1,\ldots,a_\ell$ 的超曲面构成的光滑完全交 $X$。令 $\mathcal V=\bigoplus_{i=1}^{\ell}\mathcal O_{\mathbb P^n}(a_i)$，于是 $X$ 是 $\mathcal V$ 的一个整体截面的零点集。还假设 $X$ 的反典范类为 nef；在此情形下，这等价于 $\sum_{i=1}^{\ell}a_i\le n+1$。

首先要严谨陈述定理~\ref{thm:givental-projective-mirror}，它联系了下文定义的两个取上同调值的形式函数 $I_{\mathcal V}$、$J_{\mathcal V}$。其中 $I_{\mathcal V}$ 较为初等。遵循 \SourceCite{Givental2,Kim2}，定义
\begin{equation}\label{eq:givental-complete-intersection-i-function}\index{I函数@$I$ 函数|hyperpage}
\begin{aligned}
I_{\mathcal V}&=I_{\mathcal V}(t_0,t_1,\hbar^{-1})\\
&=e^{(t_0+t_1H)/\hbar}\operatorname{Euler}(\mathcal V)
\sum_{d=0}^{\infty}e^{dt_1}
\frac{\prod_{i=1}^{\ell}\prod_{m=1}^{a_id}(a_iH+m\hbar)}
{\prod_{m=1}^d(H+m\hbar)^{n+1}},
\end{aligned}
\end{equation}
其中 $H\in H^2(\mathbb P^n)$ 是超平面类。记号 $I_{\mathcal V}(t_0,t_1,\hbar^{-1})$ 表示，有时把 $I_{\mathcal V}$ 看作按 $\hbar^{-1}$ 的幂展开的形式级数。若 $X=\mathbb P^n$，即 $\mathcal V=0$，则把 $I_{\mathcal V}$ 记为 $I_{\mathbb P^n}$。

另一方面，形式函数 $J_{\mathcal V}$ 的定义较为复杂。\hyperref[chap:quantum-differential-equations]{第十章}研究了 Givental 函数 $J=J_X$。若 $\dim X\ge3$，则 $H^2(X)$ 由 $i^*H$ 生成，于是引理~\ref{lem:giventals-j-two-formulas} 给出 $J_X$ 的两个公式：
\begin{equation}\label{eq:complete-intersection-ordinary-j-function}
\begin{aligned}
J_X&=e^{(t_0+t_1i^*H)/\hbar}
\left(1+\sum_{d=1}^{\infty}\sum_{a=0}^m e^{dt_1}
\left\langle\frac{T_a}{\hbar-c},1\right\rangle_{0,d}T^a\right)\\
&=e^{(t_0+t_1i^*H)/\hbar}
\left(1+\sum_{d=1}^{\infty}e^{dt_1}PD^{-1}e_{1*}
\left(\frac1{\hbar-c}\cap[\overline M_{0,2}(X,d)]^{\mathrm{virt}}\right)\right).
\end{aligned}
\end{equation}
这里，$T_a$ 是 $H^*(X)$ 的上同调基，$T^a$ 是其对偶基，$PD$ 是 Poincar\'e 对偶，$e_1:\overline M_{0,2}(X,d)\to X$ 是赋值映射，而
\begin{equation}\label{eq:complete-intersection-symbolic-descendants}
\left\langle\frac{T_a}{\hbar-c},1\right\rangle_{0,d}
=\sum_{k=0}^{\infty}\hbar^{-(k+1)}\langle\tau_kT_a,1\rangle_{0,d}
\end{equation}
是\hyperref[chap:quantum-differential-equations]{第十章}采用的引力关联数的符号记法。此前把 Givental 函数简记为 $J$，但现在须区分 $X$ 与 $\mathbb P^n$，故改记为 $J_X$。函数 $J_X$ 在量子微分方程理论中起核心作用（见定理~\ref{thm:quantum-differential-relations}）。
\SourcePage{380}
注意，不能直接比较 $I_{\mathcal V}$ 与 $J_X$，因为前者取值于 $H^*(\mathbb P^n)$，后者取值于 $H^*(X)$。为此，将引入一个取值于 $H^*(\mathbb P^n)$ 的 $J$ 函数变体 $J_{\mathcal V}$。然后，Givental 的定理将说明，当 $X$ 的反典范类为 nef 时，$I_{\mathcal V}$ 与 $J_{\mathcal V}$ 如何联系。

解释此定理的安排如下。先讨论 $\mathcal V=0$ 的情形；此时 Givental 的定理就是简单的等式 $I_{\mathbb P^n}=J_{\mathbb P^n}$（见下文命题~\ref{prop:projective-space-i-equals-j}）。然后定义 $\mathbb P^n$ 上一般丛所对应的 $J_{\mathcal V}$，并陈述 Givental 的镜像定理（定理~\ref{thm:givental-projective-mirror}）。证明将推迟至\hyperref[subsec:givental-projective-space-proof]{第~\ref*{subsec:givental-projective-space-proof}}、\hyperref[subsec:givental-general-bundle-proof]{\ref*{subsec:givental-general-bundle-proof} 小节}。还将解释 $J_{\mathcal V}$ 与 $J_X$ 的联系，并说明定理~\ref{thm:givental-projective-mirror} 蕴含五次三维簇的镜像对称。

此时可以联系 \SourceCite{LLY} 的方法。Givental 的函数 $I_{\mathcal V}$ 在 $t_0=0$ 时，除去 $\hbar$ 的符号变化，正是以 \SourceCite{LLY} 记号表示的\IndexLocation{indexloc:0945:380}非等变极限 $\lim_{\lambda_i\to0}HG[\mathcal I(\widehat P)](t_1)$。在将本节计算与 \SourceCite{LLY} 的相应计算比较时，不再明确重复指出这一符号差别。此外，一般情形下，定理~\ref{thm:givental-projective-mirror} 的内容是：变量变换后，$I_{\mathcal V}$ 与 $J_{\mathcal V}$ 相同。在 \SourceCite{LLY} 的方法中，对应陈述则是：作同样变量变换后，$HG[\mathcal I(\widehat P)](t_1)$ 与 $HG[\mathcal I(\widehat Q)](t_1)$ 相同。

现在可以陈述 Givental 版本镜像定理在 $X=\mathbb P^n$，即射影空间本身这一特殊情形下的结果。用式~\eqref{eq:givental-complete-intersection-i-function} 中 $\ell=0$ 定义 $I_{\mathbb P^n}$，并如式~\eqref{eq:complete-intersection-ordinary-j-function} 定义 $J_{\mathbb P^n}$。
\begin{proposition}\label{prop:projective-space-i-equals-j}
对于射影空间 $\mathbb P^n$，有 $I_{\mathbb P^n}=J_{\mathbb P^n}$。
\end{proposition}
将在\hyperref[subsec:givental-projective-space-proof]{第~\ref*{subsec:givental-projective-space-proof} 小节}证明此命题。不过，\hyperref[chap:quantum-differential-equations]{第十章}已验证了 $n=1$ 的情形，下面先回顾这一结果。
\SourceExampleNumber{11.2.1}{1}
\begin{example}\label{ex:projective-line-i-j-equality}
当 $X=\mathbb P^1$ 时，由\hyperref[subsec:gravitational-correlator-properties]{第~\ref*{subsec:gravitational-correlator-properties} 小节}的次数公理，式~\eqref{eq:complete-intersection-ordinary-j-function} 的第一行简化为
\[
J_{\mathbb P^1}=e^{(t_0+t_1H)/\hbar}
\left(1+\sum_{d=1}^{\infty}\left(
\frac{e^{dt_1}}{\hbar^{2d}}\langle\tau_{2d-1}H,1\rangle_{0,d}1
+\frac{e^{dt_1}}{\hbar^{2d+1}}\langle\tau_{2d},1\rangle_{0,d}H\right)\right).
\]
此外，式~\eqref{eq:givental-complete-intersection-i-function} 在 $n=1$、$\ell=0$ 时为
\[
I_{\mathbb P^1}=e^{(t_0+t_1H)/\hbar}\sum_{d=0}^{\infty}
\frac{e^{dt_1}}{\bigl((H+\hbar)(H+2\hbar)\cdots(H+d\hbar)\bigr)^2}.
\]
例题~\ref{ex:projective-line-gravitational} 计算了 $\langle\tau_{2d-1}H,1\rangle_{0,d}$ 与 $\langle\tau_{2d},1\rangle_{0,d}$；例题~\ref{ex:projective-line-j-function} 使用这些公式证明了 $I_{\mathbb P^1}=J_{\mathbb P^1}$。

从命题~\ref{prop:projective-space-i-equals-j} 的角度看，也可反过来进行这一过程。一旦知道 $I_{\mathbb P^1}=J_{\mathbb P^1}$，按 $H$ 的幂展开 $I_{\mathbb P^1}$（注意，$H^*(\mathbb P^1)$ 中 $H^2=0$），就容易得到 $\langle\tau_{2d-1}H,1\rangle_{0,d}$ 与 $\langle\tau_{2d},1\rangle_{0,d}$ 的公式。因此，可把命题~\ref{prop:projective-space-i-equals-j} 看作对 $\mathbb P^1$ 的\IndexLocation{indexloc:0463:380}\IndexLocation{indexloc:0469:380}引力关联数的计算。
\end{example}

接下来，定义上述 $\mathcal V=\bigoplus_{i=1}^{\ell}\mathcal O_{\mathbb P^n}(a_i)$ 对应的 $J_{\mathcal V}$。先引入一些记号。给定任意 $d>0$，当 $f:C\to\mathbb P^n$ 在 $\overline M_{0,k}(\mathbb P^n,d)$ 中变化时，截面空间 $H^0(C,f^*\mathcal V)$ 构成该模空间上的向量丛 $\mathcal V_{d,k}$。更准确地，令 $\pi_{k+1}:\overline M_{0,k+1}(\mathbb P^n,d)\to\overline M_{0,k}(\mathbb P^n,d)$ 为忘掉第 $k+1$ 个标记点的映射，令 $e_{k+1}:\overline M_{0,k+1}(\mathbb P^n,d)\to\mathbb P^n$ 为该标记点处的\IndexLocation{indexloc:0276:380}赋值映射。于是
\[
\mathcal V_{d,k}=\pi_{k+1*}e_{k+1}^*\mathcal V.
\]
\SourcePage{381}
特别地，$\mathcal V_{d,0}=\pi_{1*}e_1^*\mathcal V$，恰是本章先前记作 $\mathcal V_d$ 的丛。还可证明，当 $k>0$ 时，$\mathcal V_{d,k}=\pi_k^*(\mathcal V_{d,k-1})$。

现在固定 $1\le i\le k$。对 $\overline M_{0,k}(\mathbb P^n,d)$ 中每个 $f:C\to\mathbb P^n$，考虑 $H^0(C,f^*\mathcal V)$ 中在第 $i$ 个标记点处消失的截面。它们构成 $\overline M_{0,k}(\mathbb P^n,d)$ 上丛 $\mathcal V'_{d,k,i}$ 的纤维，且有正合列
\begin{equation}\label{eq:marked-bundle-evaluation-exact-sequence}
0\longrightarrow\mathcal V'_{d,k,i}\longrightarrow\mathcal V_{d,k}
\longrightarrow e_i^*\mathcal V\longrightarrow0.
\end{equation}
这给出
\begin{equation}\label{eq:marked-bundle-euler-factorization}
\operatorname{Euler}(\mathcal V_{d,k})
=\operatorname{Euler}(\mathcal V'_{d,k,i})\cup e_i^*\operatorname{Euler}(\mathcal V).
\end{equation}
特别地，得到 $\overline M_{0,2}(\mathbb P^n,d)$ 上的丛 $\mathcal V'_{d,2,1}$。利用此丛，以以下公式定义 $J_{\mathcal V}$：
\begin{equation}\label{eq:givental-projective-bundle-j-function}
\begin{aligned}
J_{\mathcal V}&=e^{(t_0+t_1H)/\hbar}\operatorname{Euler}(\mathcal V)\times\\
&\quad\left(1+\sum_{d=1}^{\infty}\sum_{b=0}^{\infty}
\frac{e^{dt_1}}{\hbar^{b+1}}
 e_{1!}\bigl(c_1(\mathcal L_1)^b\cup\operatorname{Euler}(\mathcal V'_{d,2,1})\bigr)\right).
\end{aligned}
\end{equation}
用上述式~\eqref{eq:complete-intersection-ordinary-j-function} 的\IndexLocation{indexloc:0411:381}符号记法，可写成
\begin{equation}\label{eq:givental-projective-bundle-j-symbolic}
J_{\mathcal V}=e^{(t_0+t_1H)/\hbar}\operatorname{Euler}(\mathcal V)
\left(1+\sum_{d=1}^{\infty}e^{dt_1}e_{1!}
\left(\frac{\operatorname{Euler}(\mathcal V'_{d,2,1})}{\hbar-c}\right)\right).
\end{equation}
注意 $J_{\mathcal V}$ 的两个重要性质。首先，当 $\mathcal V=0$ 时，对应完全交为 $X=\mathbb P^n$。因为 $\mathbb P^n$ 为凸，$\overline M_{0,2}(\mathbb P^n,d)$ 的\IndexLocation{indexloc:1406:381}虚拟基本类就是通常基本类，因而式~\eqref{eq:complete-intersection-ordinary-j-function} 的第二个公式简化为
\[
J_{\mathbb P^n}=e^{(t_0+t_1H)/\hbar}
\left(1+\sum_{d=1}^{\infty}e^{dt_1}e_{1!}\left(\frac1{\hbar-c}\right)\right).
\]
零丛的 Euler 类为 $1$，故这与式~\eqref{eq:givental-projective-bundle-j-symbolic} 在 $\mathcal V=0$ 时一致。因此，当 $\mathcal V=0$ 时，$J_{\mathcal V}=J_{\mathbb P^n}$。

再看第二个性质。考虑式~\eqref{eq:givental-projective-bundle-j-function} 的括号表达式；其 $\hbar^{-1}$ 的系数包含因子
\[
\operatorname{Euler}(\mathcal V)\cup e_{1!}\bigl(\operatorname{Euler}(\mathcal V'_{d,2,1})\bigr)
=e_{1!}\bigl(\operatorname{Euler}(\mathcal V_{d,2})\bigr)=0.
\]
第一个等号使用了式~\eqref{eq:marked-bundle-euler-factorization}。第二个等号源于 $e_1$ 可经由忘掉第二个标记点的映射 $\pi_2:\overline M_{0,2}(\mathbb P^n,d)\to\overline M_{0,1}(\mathbb P^n,d)$ 分解。上文已说明 $\mathcal V_{d,2}=\pi_2^*\mathcal V_{d,1}$，故
\[
\pi_{2!}\bigl(\operatorname{Euler}(\mathcal V_{d,2})\bigr)
=\pi_{2!}\pi_2^*\bigl(\operatorname{Euler}(\mathcal V_{d,1})\bigr)=0,
\]
因为 $\pi_2$ 的纤维具有正维数。这给出
\begin{equation}\label{eq:givental-bundle-j-normalization}
J_{\mathcal V}=e^{(t_0+t_1H)/\hbar}\operatorname{Euler}(\mathcal V)
\bigl(1+o(\hbar^{-1})\bigr),
\end{equation}
类似于 Givental $J$ 函数的性质~\eqref{eq:j-function-asymptotic-normalization}。

现在可以陈述 Givental 关于 $\mathbb P^n$ 中\IndexLocation{indexloc:0856:381}完全交的镜像定理。
\begin{theorem}\label{thm:givental-projective-mirror}\index{镜像定理!射影完全交|hyperpage}
设 $X\subset\mathbb P^n$ 为 $\ell$ 个次数为 $a_i$ 的超曲面的完全交，且 $\sum_{i=1}^{\ell}a_i\le n+1$。若 $\mathcal V=\bigoplus_{i=1}^{\ell}\mathcal O_{\mathbb P^n}(a_i)$，则取形式上同调值的函数 $I_{\mathcal V}$、$J_{\mathcal V}$ 在以下加权齐次的三角变量变换后相同：
\[
t_0\longmapsto t_0+f(e^{t_1})\hbar+h(e^{t_1}),\qquad
 t_1\longmapsto t_1+g(e^{t_1}),
\]
\SourcePage{382}
其中 $f,g,h$ 是加权齐次幂级数，满足 $\deg f=\deg g=0$、$\deg h=1$，且 $\deg e^{t_1}=n+1-\sum_i a_i$。
\end{theorem}

如先前所述，条件 $\sum_{i=1}^{\ell}a_i\le n+1$ 表示 $X$ 的反典范类为 nef。特别地，$\sum_{i=1}^{\ell}a_i$ 的最大可能值为 $n+1$；椭圆曲线、K3 曲面或维数至少为 $3$ 的 Calabi--Yau 流形都出现这种情形。还应注意，该定理包含命题~\ref{prop:projective-space-i-equals-j} 作为特例（$\mathcal V=0$）。将在\hyperref[subsec:givental-general-bundle-proof]{第~\ref*{subsec:givental-general-bundle-proof} 小节}概述定理~\ref{thm:givental-projective-mirror} 的证明。

另一个有用的说明是，定理~\ref{thm:givental-projective-mirror} 中的幂级数 $f,g,h$ 可由 $I_{\mathcal V}$ 的系数唯一确定。为此，对 $J_{\mathcal V}$ 作变量代换 $t_0\mapsto t_0+f(e^{t_1})\hbar+h(e^{t_1})$、$t_1\mapsto t_1+g(e^{t_1})$。略去 $I_{\mathcal V}$ 与 $J_{\mathcal V}$ 中的因子 $e^{t_0/\hbar}$，则由式~\eqref{eq:givental-bundle-j-normalization}，代换后的 $J_{\mathcal V}$ 为
\[
J_{\mathcal V}=e^f\operatorname{Euler}(\mathcal V)
\left(1+h\hbar^{-1}1+(t_1+g)\hbar^{-1}H+o(\hbar^{-1})\right).
\]
令它与 $I_{\mathcal V}$ 相等，即得
\begin{equation}\label{eq:givental-mirror-variable-coefficients}
\begin{aligned}
e^f&\longleftrightarrow\widetilde I_{\mathcal V}\text{ 中 }1\text{ 的系数},\\
e^fh&\longleftrightarrow\widetilde I_{\mathcal V}\text{ 中 }\hbar^{-1}1\text{ 的系数},\\
e^f(t_1+g)&\longleftrightarrow\widetilde I_{\mathcal V}\text{ 中 }\hbar^{-1}H\text{ 的系数},
\end{aligned}
\end{equation}
其中 $\widetilde I_{\mathcal V}$ 是式~\eqref{eq:givental-complete-intersection-i-function} 去掉因子 $\operatorname{Euler}(\mathcal V)$ 后的函数（故 $I_{\mathcal V}=\operatorname{Euler}(\mathcal V)\widetilde I_{\mathcal V}$）。如所断言，$I_{\mathcal V}$ 确定了镜像定理中的坐标变换。特别地，如果在原有变量 $t_0,t_1$ 下就已有
\[
I_{\mathcal V}=e^{(t_0+t_1H)/\hbar}\operatorname{Euler}(\mathcal V)
\bigl(1+o(\hbar^{-1})\bigr),
\]
则必有 $f=g=h=0$，从而 $I_{\mathcal V}=J_{\mathcal V}$。以下是一个这样的例子。
\SourceExampleNumber{11.2.1}{2}
\begin{example}\label{ex:fano-hypersurface-i-j-equality}
设 $X\subset\mathbb P^n$ 是次数 $\ell<n$ 的超曲面。关于 $\ell$ 的假设容易给出 $I_{\mathcal V}=e^{(t_0+t_1H)/\hbar}\operatorname{Euler}(\mathcal V)(1+o(\hbar^{-1}))$，故由上述讨论，在定理~\ref{thm:givental-projective-mirror} 中有 $f=g=h=0$。因此
\[
J_{\mathcal V}=I_{\mathcal V}=\ell H e^{(t_0+t_1H)/\hbar}
\left(1+\sum_{d=1}^{\infty}e^{dt_1}
\frac{(\ell H+\hbar)(\ell H+2\hbar)\cdots(\ell H+\ell d\hbar)}
{\bigl((H+\hbar)(H+2\hbar)\cdots(H+d\hbar)\bigr)^{n+1}}\right).
\]
很快将看到，这对 $X$ 的\IndexLocation{indexloc:1129:382}小量子上同调环有一些有趣的推论。
\end{example}

为充分理解定理~\ref{thm:givental-projective-mirror}，须先解释 $J_X$ 与 $J_{\mathcal V}$ 的关系。为此，使用包含映射 $i:X\hookrightarrow\mathbb P^n$ 诱导的 Gysin 映射 $i_!:H^*(X)\to H^{*+2\ell}(\mathbb P^n)$。有以下结果。
\begin{proposition}\label{prop:complete-intersection-j-pushforward}
设 $i:X\hookrightarrow\mathbb P^n$ 为 $\mathcal V=\bigoplus_{i=1}^{\ell}\mathcal O_{\mathbb P^n}(a_i)$ 所确定的完全交，且 $\dim X\ge3$。则 $J_{\mathcal V}=i_!(J_X)$。
\end{proposition}
\begin{proof}
$\dim X\ge3$ 的假设给出 $H^2(\mathbb P^n)\simeq H^2(X)$，因而 $J_{\mathcal V}$ 与 $J_X$ 都是 $t_0,t_1$ 的形式函数。还应注意，$i:X\hookrightarrow\mathbb P^n$ 诱导包含映射 $j:\overline M_{0,2}(X,d)\hookrightarrow\overline M_{0,2}(\mathbb P^n,d)$，使下图交换：
\[
\begin{tikzcd}[column sep=large,row sep=large]
\overline M_{0,2}(X,d)\arrow[r,"j"]\arrow[d,"e_1"']&\overline M_{0,2}(\mathbb P^n,d)\arrow[d,"e_1"]\\
X\arrow[r,"i"]&\mathbb P^n
\end{tikzcd}
\]
\SourcePage{383}
这里 $e_1$ 为赋值映射。

关键事实是，可以计算 $\overline M_{0,2}(X,d)$ 的虚拟基本类。按定义，$X$ 为 $\mathcal V$ 的一个截面 $s$ 的零点集，而 $s$ 诱导了上述丛 $\mathcal V_{d,k}$ 的一个截面 $\widetilde s$。与例题~\ref{ex:quintic-virtual-class} 一样，$\widetilde s$ 的零点概形（按代数叠的意义）自然等同于 $\overline M_{0,2}(X,d)$。于是可对向量丛 $\mathcal V_{d,2}$ 使用法锥构造，得到 $[\overline M_{0,2}(X,d)]^{\mathrm{virt}}$。推广式~\eqref{eq:quintic-virtual-class-pushforward}，可得
\begin{equation}\label{eq:complete-intersection-virtual-pushforward}
j_*\bigl([\overline M_{0,2}(X,d)]^{\mathrm{virt}}\bigr)
=\operatorname{Euler}(\mathcal V_{d,2})\cap[\overline M_{0,2}(\mathbb P^n,d)].
\end{equation}
因此，$\overline M_{0,2}(X,d)$ 的虚拟基本类是 $\mathcal V_{d,2}$ 的 Euler 类的一个精细化。

现在证明命题。从式~\eqref{eq:complete-intersection-ordinary-j-function} 给出的 $J_X$ 的第二个公式出发，利用投影公式及上述交换图，容易看到
\[
\begin{aligned}
i_!(J_X)&=e^{(t_0+t_1H)/\hbar}\left(\operatorname{Euler}(\mathcal V)+\sum_{d=1}^{\infty}e^{dt_1}
 e_{1!}PD^{-1}j_*\left(\frac1{\hbar-c}\cap[\overline M_{0,2}(X,d)]^{\mathrm{virt}}\right)\right).
\end{aligned}
\]
在符号记法 $1/(\hbar-c)$ 中，$c$ 表示 $\overline M_{0,2}(X,d)$ 上丛 $\mathcal L_1$ 的 Chern 类 $c_1(\mathcal L_1)$。它是 $\overline M_{0,2}(\mathbb P^n,d)$ 上同样记为 $\mathcal L_1$ 的丛经 $j$ 拉回所得。因此，由 $j$ 的投影公式，
\[
\begin{aligned}
j_*\left(\frac1{\hbar-c}\cap[\overline M_{0,2}(X,d)]^{\mathrm{virt}}\right)
&=\frac1{\hbar-c}\cap j_*\bigl([\overline M_{0,2}(X,d)]^{\mathrm{virt}}\bigr)\\
&=\frac1{\hbar-c}\cap\bigl(\operatorname{Euler}(\mathcal V_{d,2})\cap[\overline M_{0,2}(\mathbb P^n,d)]\bigr)\\
&=\frac{\operatorname{Euler}(\mathcal V_{d,2})}{\hbar-c}\cap[\overline M_{0,2}(\mathbb P^n,d)].
\end{aligned}
\]
于是 $i_!(J_X)$ 的公式简化为
\[
i_!(J_X)=e^{(t_0+t_1H)/\hbar}
\left(\operatorname{Euler}(\mathcal V)+\sum_{d=1}^{\infty}e^{dt_1}
 e_{1!}\left(\frac{\operatorname{Euler}(\mathcal V_{d,2})}{\hbar-c}\right)\right).
\]
不过，式~\eqref{eq:marked-bundle-euler-factorization} 给出
\[
\operatorname{Euler}(\mathcal V_{d,2})
=\operatorname{Euler}(\mathcal V'_{d,2,1})\cup e_1^*\operatorname{Euler}(\mathcal V).
\]
将它代入刚得到的 $i_!(J_X)$ 公式，并用 $e_1$ 的投影公式，得到
\[
i_!(J_X)=e^{(t_0+t_1H)/\hbar}\operatorname{Euler}(\mathcal V)
\left(1+\sum_{d=1}^{\infty}e^{dt_1}e_{1!}
\left(\frac{\operatorname{Euler}(\mathcal V'_{d,2,1})}{\hbar-c}\right)\right),
\]
这恰是 $J_{\mathcal V}$ 的定义。命题得证。
\end{proof}

既然已经知道 $i_!(J_X)=J_{\mathcal V}$，接下来研究从 $J_X$ 到 $J_{\mathcal V}$ 是否会丢失信息。使用以下两点观察：
\begin{itemize}
\item $X$ 是丰沛超曲面的完全交，公式 $i_!(i^*\alpha)=\alpha\cup\operatorname{Euler}(\mathcal V)$ 表明，$i_!$ 在 $i^*$ 的像上是单射。注意，$i^*H^*(\mathbb P^n)\subset H^{\mathrm{even}}(X)$。
\item \hyperref[subsec:gravitational-correlator-properties]{第~\ref*{subsec:gravitational-correlator-properties} 小节}的\IndexLocation{indexloc:0461:383}次数公理表明，$J_X$ 只涉及 $X$ 的偶次上同调。
\end{itemize}
\SourcePage{384}
当 $\dim X\ge3$ 为奇数时，由 Lefschetz 定理，$H^{\mathrm{even}}(\mathbb P^n)\to H^{\mathrm{even}}(X)$ 为满射。第一点说明，$i_!$ 在 $H^{\mathrm{even}}(X)$ 上为单射；结合第二点，可知能从 $J_{\mathcal V}=i_!(J_X)$ 恢复 $J_X$。

另一方面，当 $\dim X\ge3$ 为偶数时，记 $\dim X=2m$。此时，$J_X$ 仍取值于 $H^{\mathrm{even}}(X)$，但 $i^*:H^{\mathrm{even}}(\mathbb P^n)\to H^{\mathrm{even}}(X)$ 未必为满射。问题在于，$H^{2m}(X)$ 有 Lefschetz 分解
\[
H^{2m}(X)=\mathbb C\cdot i^*(H^m)\oplus H_0^{2m}(X),
\]
其中 $H_0^{2m}(X)$ 是本原上同调。$i^*$ 的像完全不包含 $H_0^{2m}(X)$ 中的非零元素；若 $J_X$ 有本原部分的分量，它便不在这个像中。对 $i^*:H^*(\mathbb P^n)\to H^*(X)$ 的像，有不同记号。按\hyperref[subsec:toric-mirror-conjecture]{第~\ref*{subsec:toric-mirror-conjecture} 小节}的术语，因为 $\mathbb P^n$ 为环面簇，把此像记作 $H_{\mathrm{toric}}^*(X)$。\SourceCite{Givental3} 则记作 $H^*(\mathcal V)$，这里采用后一记号。利用上述 $H^{2m}(X)$ 的分解，得到
\[
H^{\mathrm{even}}(X)=H^*(\mathcal V)\oplus H_0^{2m}(X),
\]
因为 $X$ 是完全交。这诱导了分解
\begin{equation}\label{eq:complete-intersection-j-primitive-decomposition}
J_X=J_X^{\mathcal V}+J_X^0,
\end{equation}
其中 $J_X^{\mathcal V}$ 取值于 $H^*(\mathcal V)$，$J_X^0$ 是 $J_X$ 的“本原部分”。为联系它与 $i_!(J_X)$，回顾 $H_0^{2m}(X)$ 正是 $i_!:H^{2m}(X)\to H^{2m+2\ell}(\mathbb P^n)$ 的核。对式~\eqref{eq:complete-intersection-j-primitive-decomposition} 应用 $i_!$，即得 $J_{\mathcal V}=i_!(J_X^{\mathcal V})$。因此，对于偶数维完全交，$J_{\mathcal V}$ 反映的是 $J_X$ 中来自 $\mathbb P^n$ 的部分。

可以把上述观察统一起来：式~\eqref{eq:complete-intersection-j-primitive-decomposition} 的分解 $J_X=J_X^{\mathcal V}+J_X^0$ 对任意维数都有意义（所以，当 $\dim X$ 为奇数时，$J_X^0=0$）。上述讨论表明，对所有 $X$，均可从 $J_{\mathcal V}=i_!(J_X^{\mathcal V})$ 恢复 $J_X^{\mathcal V}$。因此，镜像定理告诉我们如何计算 $X$ 的一部分亏格零引力关联数，即涉及来自环境空间的上同调类的关联数。

还可给出一个完全在 $H^*(\mathcal V)$ 内表述的镜像定理版本。若写成 $I_{\mathcal V}=\operatorname{Euler}(\mathcal V)\widetilde I_{\mathcal V}$，则 $i^*(\widetilde I_{\mathcal V})$ 与 $J_X^{\mathcal V}$ 都是取值于 $H^*(\mathcal V)$ 的形式函数。又因为 $i_!$ 在 $H^*(\mathcal V)$ 上为单射，分别把 $i^*(\widetilde I_{\mathcal V})$ 与 $J_X^{\mathcal V}$ 映到 $I_{\mathcal V}$ 与 $J_{\mathcal V}$，故定理~\ref{thm:givental-projective-mirror} 等价于断言：作适当类型的坐标变换后，$i^*(\widetilde I_{\mathcal V})$ 与 $J_X^{\mathcal V}$ 相等。这是 Givental 在 \SourceCite{Givental3} 中陈述镜像定理的方式，不过他允许环境空间为环面簇。\hyperref[subsec:givental-toric-mirror]{第~\ref*{subsec:givental-toric-mirror} 小节}将讨论这个更一般的情形。还应指出，这里使用 $H^*(\mathbb P^n)$ 而非 $H^*(\mathcal V)$ 的方法，基于 \SourceCite{Kim2}。

为看看镜像定理在具体情形下的样子，考虑五次三维簇。
\SourceExampleNumber{11.2.1}{3}
\begin{example}\leavevmode\phantomsection\label{ex:givental-quintic-mirror-potential}%
五次三维簇 $V\subset\mathbb P^4$ 对应于 $\mathcal V=\mathcal O_{\mathbb P^4}(5)$，故 $\operatorname{Euler}(\mathcal V)=5H$。这里 $H$ 是 $H^*(\mathbb P^4)$ 中的超平面类。于是
\[
I_{\mathcal V}=e^{(t_0+t_1H)/\hbar}5H\sum_{d=0}^{\infty}e^{dt_1}
\frac{(5H+\hbar)\cdots(5H+5d\hbar)}
{\bigl((H+\hbar)\cdots(H+d\hbar)\bigr)^5}.
\]
\SourcePage{385}
与式~\eqref{eq:quintic-hypergeometric-period-expansion} 比较，公式相同，只是 $\hbar$ 的符号不同，并多了一个因子 $e^{t_0/\hbar}$。因此
\[
I_{\mathcal V}=e^{t_0/\hbar}5H
\left(y_0+y_1\frac H\hbar+y_2\frac{H^2}{\hbar^2}+y_3\frac{H^3}{\hbar^3}\right),
\]
其中 $y_0,y_1,y_2,y_3$ 是\hyperref[subsec:quintic-mirror-theorem-statement]{第~\ref*{subsec:quintic-mirror-theorem-statement} 小节}讨论的五次镜像的 Picard--Fuchs 方程的一组解基。

镜像定理告诉我们，作适当变量变换后，$I_{\mathcal V}$ 等于 $J_{\mathcal V}$；由式~\eqref{eq:givental-mirror-variable-coefficients}，可从 $I_{\mathcal V}$ 读出 $f,g,h$。因此，上述 $I_{\mathcal V}$ 的公式给出
\[
\begin{aligned}
e^f=y_0&\ \Longrightarrow\ f=\log y_0,\\
e^fh=0&\ \Longrightarrow\ h=0,\\
e^f(t_1+g)=y_1&\ \Longrightarrow\ g=\frac{y_1}{y_0}-t_1.
\end{aligned}
\]
于是 $t_0\mapsto t_0+f\hbar=t_0+(\log y_0)\hbar$，$t_1\mapsto t_1+g=y_1/y_0$。注意，后者恰是\hyperref[subsec:quintic-mirror-theorem-statement]{第~\ref*{subsec:quintic-mirror-theorem-statement} 小节}使用的镜像映射。现在写 $J_{\mathcal V}=J_{\mathcal V}(t_0,t_1)$、$I_{\mathcal V}=I_{\mathcal V}(t_0,t_1)$。定理~\ref{thm:givental-projective-mirror} 给出
\[
J_{\mathcal V}\bigl(t_0+(\log y_0)\hbar,y_1/y_0\bigr)=I_{\mathcal V}(t_0,t_1).
\]
令 $s=y_1/y_0$，并注意 $e^{(t_0+\log(y_0)\hbar)/\hbar}=y_0e^{t_0/\hbar}$，上式就简化为
\[
J_{\mathcal V}(t_0,s)=\frac1{y_0}I_{\mathcal V}(t_0,t_1).
\]
结合上述 $I_{\mathcal V}$ 的公式，得到
\begin{equation}\label{eq:givental-quintic-j-periods}
J_{\mathcal V}(t_0,s)=e^{t_0/\hbar}5H
\left(1+s\frac H\hbar+\frac{y_2}{y_0}\frac{H^2}{\hbar^2}
+\frac{y_3}{y_0}\frac{H^3}{\hbar^3}\right).
\end{equation}
不过，\hyperref[subsec:cy-threefold-j-function]{第~\ref*{subsec:cy-threefold-j-function} 小节}还给出了 $J_V$ 的一些漂亮公式。使用 $H^{\mathrm{even}}(V)$ 的上同调基
$T_0=1$、$T_1=i^*H$、$T^1=i^*H^2/5$、$T^0=i^*H^3/5$。注意 $T_0,T^0$ 互为对偶，$T_1,T^1$ 亦然。若 $\Phi(s)$ 为五次三维簇的 Gromov--Witten 势，则命题~\ref{prop:cy-threefold-j-potential} 给出
\[
J_V(t_0,s)=e^{t_0/\hbar}
\left(1+\hbar^{-1}sT_1+\hbar^{-2}\Phi'(s)T^1
+\hbar^{-3}\bigl(s\Phi'(s)-2\Phi(s)\bigr)T^0\right).
\]
应用 $i_!$ 并使用命题~\ref{prop:complete-intersection-j-pushforward}，得到
\begin{equation}\label{eq:givental-quintic-j-potential}
J_{\mathcal V}(t_0,s)=e^{t_0/\hbar}5H
\left(1+s\frac H\hbar+\frac{\Phi'(s)}5\frac{H^2}{\hbar^2}
+\frac{s\Phi'(s)-2\Phi(s)}5\frac{H^3}{\hbar^3}\right).
\end{equation}
现在令 $s=\Psi(t_1)$。将式~\eqref{eq:givental-quintic-j-potential} 与式~\eqref{eq:givental-quintic-j-periods} 中的 $J_{\mathcal V}(t_0,s)$ 公式比较，容易得到
\[
\Phi(\Psi(t_1))=\frac52\left(\frac{y_1}{y_0}\frac{y_2}{y_0}-\frac{y_3}{y_0}\right).
\]
这恰是\hyperref[subsec:lly-quintic-mirror-proof]{第~\ref*{subsec:lly-quintic-mirror-proof} 小节}讨论的五次三维簇镜像定理。如定理~\ref{thm:quintic-mirror-potential-equivalence} 所指出，它蕴含\hyperref[sec:mirror-symmetry-revisited]{第~\ref*{sec:mirror-symmetry-revisited} 节}陈述的 Hodge 理论版本的镜像对称。

还须完成\hyperref[chap:quintic]{第二章}开始的讨论。回顾\hyperref[subsec:givental-ij]{第~\ref*{subsec:givental-ij} 小节}，那里给出了上述 $I_{\mathcal V}$、$J_{\mathcal V}$ 的公式，并讨论定理~\ref{thm:givental-projective-mirror} 如何
\SourcePage{386}
联系镜像对称。回看该小节，一个关键公式是式~\eqref{eq:iv-second-derivative}，它指出，当 $t_0=0$、$\hbar=1$ 时，有
\[
(I_{\mathcal V}/y_0)''=(5+2875q+\cdots)(H^3+sH^4),
\]
其中 $'$ 表示 $d/ds$。随后证明了括号中的级数 $5+2875q+\cdots$，正是五次镜像的归一化 Yukawa 耦合 $Y$。因此
\begin{equation}\label{eq:givental-quintic-yukawa-identity}
(I_{\mathcal V}/y_0)''=Y(H^3+sH^4).
\end{equation}
用它说明，Givental 的镜像定理如何给出五次三维簇的“经典”镜像对称表述。如上，定理~\ref{thm:givental-projective-mirror} 给出 $I_{\mathcal V}/y_0(t_0,t_1)=J_{\mathcal V}(t_0,s)$。因此
\[
(I_{\mathcal V}/y_0)''=J_{\mathcal V}'',
\]
其中 $'=d/ds$。不过，命题~\ref{prop:cy-threefold-j-potential} 表明
\[
J_V''=\hbar^{-2}e^{(t_0+sH)/\hbar}H*H,
\]
其中 $*$ 是 $H^*(V)$ 中的小量子乘积。严格地说，应将 $H$ 写为 $i^*H$，但这样会使记号过于繁杂。回顾式~\eqref{eq:one-parameter-quantum-yukawa}，$H*H=\langle H,H,H\rangle C$，其中 $C\in H^4(V)$ 是 $V$ 中一条直线的类。因此
\[
J_V''=\hbar^{-2}e^{(t_0+sH)/\hbar}\langle H,H,H\rangle C.
\]
令 $t_0=0$、$\hbar=1$。由于 $i_!(C)=H^3\in H^6(\mathbb P^4)$，有
\[
J_{\mathcal V}''=i_!(J_V)''=e^{sH}\langle H,H,H\rangle H^3
=\langle H,H,H\rangle(H^3+sH^4)
\]
在 $H^*(\mathbb P^4)$ 中成立。与式~\eqref{eq:givental-quintic-yukawa-identity} 比较，立即得到 $Y=\langle H,H,H\rangle$，从而恢复\hyperref[chap:quintic]{第二章}讨论的镜像对称版本。还应注意，五次三维簇的镜像对称可看作计算 $V$ 的瞬子数 $n_d$ 的一种方法。\SourceCite{Givental2} 解释了从镜像定理提取 $n_d$ 的另一种方式。
\end{example}

最后，解释 $J_{\mathcal V}$ 与 $J_X$ 的另一种联系，作为镜像定理讨论的收尾。其思想是，在 $H^*(\mathbb P^n)$ 上定义一个修改的“量子乘积” \SourceCite{Givental2,Kim2}，使得可用它构造 $J_{\mathcal V}$，就如 $J$ 是由通常的小量子乘积、式~\eqref{eq:giventals-flat-section-symbolic} 与式~\eqref{eq:giventals-j-pairing-definition} 得到的一样。\SourceCite{Pandharipande3} 详细定义并研究了这一乘积，记为 $*_X$。对我们而言，$*_X$ 最值得注意的性质如下。
\begin{proposition}\label{prop:modified-quantum-product-restriction}
设 $X\subset\mathbb P^n$ 是完全交。则
\[
i^*:H^*(\mathbb P^n)\otimes\mathbb C[[q]]\longrightarrow H^*(X)\otimes\mathbb C[[q]]
\]
是环同态，其中 $H^*(X)$ 采用小量子乘积 $*$，$H^*(\mathbb P^n)$ 采用修改的量子乘积 $*_X$。
\end{proposition}
\begin{proof}
证明见 \SourceCite[命题 4]{Pandharipande3}；该文将此结果归于 T. Graber。
\end{proof}
此外，函数 $J_{\mathcal V}$ 给出环 $(H^*(\mathbb P^n)\otimes\mathbb C[[q]],*_X)$ 中的关系，正如\hyperref[subsec:giventals-j-function]{第~\ref*{subsec:giventals-j-function} 小节}中 $J_X$ 给出 $X$ 的\IndexLocation{indexloc:1129:386}小\IndexLocation{indexloc:1127:386}\IndexLocation{indexloc:1128:386}量子上同调中的关系一样。
\begin{proposition}\label{prop:modified-quantum-j-relations}
设 $P(\hbar\partial/\partial t_1,e^{t_1},\hbar)J_{\mathcal V}=0$，其中 $P$ 是关于 $\hbar\partial/\partial t_1$ 的多项式，系数为 $e^{t_1}$、$\hbar$ 的幂级数。则关系 $P(H,q,0)=0$ 在 $(H^*(\mathbb P^n)\otimes\mathbb C[[q]],*_X)$ 中成立，即关系 $P(H,q,0)=0$ 在 $(H^*(\mathbb P^n)\otimes\mathbb C[[q]],*_X)/\ker i^*$ 中成立。
\end{proposition}
\SourcePage{387}
\begin{proof}
证明见 \SourceCite[引理 3]{Pandharipande3}。
\end{proof}
将命题~\ref{prop:modified-quantum-product-restriction} 与命题~\ref{prop:modified-quantum-j-relations} 合并，得到如下有用的推论。
\begin{corollary}\label{cor:complete-intersection-j-quantum-relations}
设 $P(\hbar\partial/\partial t_1,e^{t_1},\hbar)J_{\mathcal V}=0$，其中 $P$ 为上述形式的幂级数。则关系 $P(i^*H,q,0)=0$ 在 $X$ 的小量子上同调环中成立。
\end{corollary}
\begin{proof}
由命题~\ref{prop:modified-quantum-j-relations}，首先得到 $P(H,q,0)=0$，其中乘积采用 $*$。再由命题~\ref{prop:modified-quantum-product-restriction}，有 $i^*(P(H,q,0))=P(i^*H,q,0)$，右边按 $X$ 的小量子上同调环计算。因此，右边也为零。
\end{proof}
以下例子说明如何结合定理~\ref{thm:givental-projective-mirror} 与推论~\ref{cor:complete-intersection-j-quantum-relations}。
{\SourceExampleNumber{11.2.1}{2}
\renewcommand{\theHexam}{11.2.1.2.continued}
\begin{example}[续]\label{ex:fano-hypersurface-quantum-relation}
设 $X\subset\mathbb P^n$ 为次数 $\ell<n$ 的超曲面。在此例的第一部分，已经解释了 $\mathcal V=\mathcal O_{\mathbb P^n}(\ell)$ 的镜像定理如何给出等式
\[
J_{\mathcal V}=I_{\mathcal V}=\ell H e^{(t_0+t_1H)/\hbar}
\left(1+\sum_{d>0}e^{dt_1}
\frac{(\ell H+\hbar)(\ell H+2\hbar)\cdots(\ell H+\ell d\hbar)}
{\bigl((H+\hbar)(H+2\hbar)\cdots(H+d\hbar)\bigr)^{n+1}}\right).
\]
简单计算表明，$I_{\mathcal V}$ 被以下算子消去：
\[
\left(\hbar\frac d{dt_1}\right)^n
-e^{t_1}\ell\left(\ell\hbar\frac d{dt_1}+\hbar\right)
\cdots\left(\ell\hbar\frac d{dt_1}+(\ell-1)\hbar\right).
\]
因为 $J_{\mathcal V}=I_{\mathcal V}$，此算子也消去 $J_{\mathcal V}$。于是推论~\ref{cor:complete-intersection-j-quantum-relations} 给出超曲面 $X$ 的小量子上同调环中的关系
\[
H^n-q\ell(\ell H)^{\ell-1}=0,
\]
其中简记 $H=i^*H$。这与 \SourceCite{Beauville1} 所得到的关系一致。
\end{example}}

\subsection{量子超平面截面原理}\label{subsec:quantum-hyperplane-section-principle}\index{量子超平面截面原理|hyperpage}
迄今讨论的是完全交 $X\subset\mathbb P^n$ 的镜像定理。现在希望通过完全交的\emph{量子超平面截面原理}，把它放在更广泛的背景中。此原理由 \SourceCite{Givental3} 提出，并由 \SourceCite{Kim2} 给出正式表述；后者解释了广义旗流形中完全交的量子超平面截面原理。这一原理与把线性 $\sigma$ 模型的\IndexLocation{indexloc:0183:387}\IndexLocation{indexloc:0357:387}关联函数限制到反典范超曲面的量子限制公式 \SourceCite{MP1} 精神相近，但适用范围更广。

我们将在一个尚未获证的一般情形中陈述此原理。设 $Y$ 为光滑代数簇，$X\subset Y$ 为完全交，即向量丛 $\mathcal V=\bigoplus_{i=1}^{\ell}\mathcal L_i$ 的一个截面的零点集，其中 $\mathcal L_i$ 是 $Y$ 上的凸线丛。这里“凸”的含义与\hyperref[subsec:hypergeometric-mirror-transformations]{第~\ref*{subsec:hypergeometric-mirror-transformations} 小节}相同。因此，对所有无标记的亏格零\IndexLocation{indexloc:1280:387}稳定映射 $f:C\to Y$，均有 $H^1(C,f^*\mathcal L_i)=0$。

给定 $\beta\in H_2(Y,\mathbb Z)$，将\hyperref[subsec:givental-nef-projective-mirror]{第~\ref*{subsec:givental-nef-projective-mirror} 小节}的一些构造推广如下。特别地，令 $\mathcal V_{\beta,k}$ 为 $\overline M_{0,k}(Y,\beta)$ 上的向量丛，其在 $f$ 处的纤维为截面空间 $H^0(C,f^*\mathcal V)$。设 $\pi_k:\overline M_{0,k}(Y,\beta)\to\overline M_{0,k-1}(Y,\beta)$ 是忘掉第 $k$ 个标记点的映射，则 $\pi_k^*\mathcal V_{\beta,k-1}=\mathcal V_{\beta,k}$。对于 $1\le i\le k$，$H^0(C,f^*\mathcal V)$ 中在第 $i$ 个标记点消失的截面构成 $\overline M_{0,k}(Y,\beta)$ 上一个丛 $\mathcal V'_{\beta,k,i}$ 的纤维。这个丛有如下正合列
\[
0\longrightarrow\mathcal V'_{\beta,k,i}\longrightarrow\mathcal V_{\beta,k}\longrightarrow e_i^*\mathcal V\longrightarrow0.
\]
\SourcePage{388}
为简便起见，假设 $Y$ 为凸簇。现在定义取值于 $H^*(Y)$ 的形式函数 $J_{\mathcal V}$ 如下：
\begin{equation}\label{eq:givental-general-j-function}
J_{\mathcal V}=e^{(t_0+\delta)/\hbar}\operatorname{Euler}(\mathcal V)
\left(1+\sum_{\beta\ne0}q^\beta e_{1!}\left(\frac{\operatorname{Euler}(\mathcal V'_{\beta,2,1})}{\hbar-c}\right)\right).
\end{equation}
这里 $e_1$ 如常为赋值映射，$1/(\hbar-c)$ 是式~\eqref{eq:givental-projective-bundle-j-symbolic} 中所用的符号记法。采用约定 $\delta=\sum_{i=1}^r t_iT_i$，其中 $T_i$ 遍历 $H^2(Y)$ 的一组基，并令 $q^\beta=e^{\int_\beta\delta}$。当 $Y$ 为凸簇 $\mathbb P^n$ 时，这便化为式~\eqref{eq:givental-projective-bundle-j-symbolic}。此外，当 $Y$ 为一般的凸簇且 $\mathcal V=0$ 时，不难证明 $J_{\mathcal V}=J_Y$。

令 $\mathcal L_j(\beta)=\int_\beta c_1(\mathcal L_j)$，并定义另一个上同调值形式函数
\[
\begin{aligned}
\overline I_{\mathcal V}={}&e^{(t_0+\delta)/\hbar}\operatorname{Euler}(\mathcal V)\times\\
&\left(1+\sum_{0\ne\beta\in M(Y)}q^\beta
\prod_{j=1}^{\ell}\frac{\prod_{m=-\infty}^{\mathcal L_j(\beta)}(c_1(\mathcal L_j)+m\hbar)}{\prod_{m=-\infty}^0(c_1(\mathcal L_j)+m\hbar)}
 e_{1!}\left(\frac1{\hbar-c}\right)\right).
\end{aligned}
\]
如\hyperref[subsec:nef-gkz]{第~\ref*{subsec:nef-gkz} 小节}所述，乘积中除有限个因子外，其余全部约去。

现在陈述\IndexLocation{indexloc:1159:388}\IndexLocation{indexloc:1160:388}\IndexLocation{indexloc:1161:388}量子超平面截面原理。假设 $Y$ 的 Kähler 锥是由类 $T_i$ 生成的单纯锥。如常引入相应的变量 $t_i$ 与 $q_i=e^{t_i}$，用以下等式定义 $q_i$ 的次数：
\[
c_1(Y)-c_1(\mathcal V)=\sum_{i=1}^r(\deg q_i)T_i.
\]
注意，若 $Y=\mathbb P^n$、$\mathcal V=\bigoplus_{i=1}^{\ell}\mathcal O(a_i)$，则 $\deg q=n+1-\sum_{i=1}^{\ell}a_i$。
\begin{conjecture}\label{conj:quantum-hyperplane-section-principle}
如上，设 $X\subset Y$ 为与 $\mathcal V$ 相应的完全交。假设 $Y$ 为凸簇且所有 $i$ 均有 $\deg q_i\ge0$。则形式向量值函数 $\overline I_{\mathcal V}$ 与 $J_{\mathcal V}$ 在唯一的加权齐次三角变量变换后相等：
\[
t_0\longmapsto t_0+f_0(q)\hbar+h(q),\qquad t_i\longmapsto t_i+f_i(q),
\]
其中 $h,f_0,\ldots,f_r$ 是无常数项的加权齐次幂级数，且 $\deg f_0=\deg f_i=0$、$\deg h=1$。
\end{conjecture}
先考察此猜想在 $Y=\mathbb P^n$ 时的含义。从命题~\ref{prop:projective-space-i-equals-j} 开始，该命题断言 $I_{\mathbb P^n}=J_{\mathbb P^n}$。由式~\eqref{eq:givental-complete-intersection-i-function}，
\begin{equation}\label{eq:projective-space-i-hypergeometric}
I_{\mathbb P^n}=e^{(t_0+t_1H)/\hbar}\sum_{d=0}^{\infty}
\frac{q^d}{\prod_{m=1}^d(H+m\hbar)^{n+1}},
\end{equation}
其中 $q=e^{t_1}$。在式~\eqref{eq:givental-projective-bundle-j-symbolic} 后的说明中，已指出 $J_{\mathbb P^n}=J_{\mathcal V}$，这里 $\mathcal V=0$。于是式~\eqref{eq:givental-projective-bundle-j-symbolic} 给出
\begin{equation}\label{eq:projective-space-j-descendants}
J_{\mathbb P^n}=e^{(t_0+t_1H)/\hbar}\left(1+\sum_{d=1}^{\infty}e^{dt_1}e_{1!}\left(\frac1{\hbar-c}\right)\right).
\end{equation}
这是因为零秩平凡丛的 Euler 类为 $1$。比较两个表达式，得到
\begin{equation}\label{eq:projective-space-descendant-pushforward}
e_{1!}\left(\frac1{\hbar-c}\right)=\frac1{\prod_{m=1}^d(H+m\hbar)^{n+1}}.
\end{equation}
\SourcePage{389}
将它代入上述 $\overline I_{\mathcal V}$ 的公式，并注意此时 $\mathcal L_j(\beta)$ 就是 $a_jd$，便容易得出 $\overline I_{\mathcal V}$ 与式~\eqref{eq:givental-complete-intersection-i-function} 定义的 $I_{\mathcal V}$ 相同。

这表明：对于 $\mathbb P^n$，一旦证明命题~\ref{prop:projective-space-i-equals-j}，猜想~\ref{conj:quantum-hyperplane-section-principle} 陈述的量子超平面截面原理便等价于定理~\ref{thm:givental-projective-mirror} 的镜像定理。这是猜想~\ref{conj:quantum-hyperplane-section-principle} 发挥作用的典型方式：将式~\eqref{eq:projective-space-descendant-pushforward} 与量子超平面截面原理结合，就能得知 $\mathbb P^n$ 中完全交的镜像猜想应当是什么。更一般的情形也可作同样推理，因此，一旦知道凸簇 $Y$ 对应的式~\eqref{eq:projective-space-descendant-pushforward} 的右边，就立即知道 $Y$ 中完全交的相应镜像定理应当是什么。

作为猜想~\ref{conj:quantum-hyperplane-section-principle} 的证据，上面的讨论表明，证明定理~\ref{thm:givental-projective-mirror} 就能验证 $Y=\mathbb P^n$ 的情形。\SourceCite{Kim2} 给出了广义旗流形中完全交的证明；\SourceCite{BCFKvS1} 将其应用于 Grassmann 簇及完全旗流形中的 Calabi--Yau 完全交三维簇，计算有理曲线的虚数目。一般而言，量子超平面截面原理似乎对多种环境空间都给出合理结果。除 \SourceCite{Kim1,BCFKvS1} 研究的完全旗流形外，量子超平面截面原理还对某些二次曲面网簇产生了可以验证的结果（见 \SourceCite{Tjotta}）。

\subsection{命题~\ref{prop:projective-space-i-equals-j} 的证明}\label{subsec:givental-projective-space-proof}
现在概述命题~\ref{prop:projective-space-i-equals-j} 的证明。它断言
\[
I_{\mathbb P^n}=J_{\mathbb P^n},
\]
其中 $I_{\mathbb P^n}$ 与 $J_{\mathbb P^n}$ 分别由式~\eqref{eq:projective-space-i-hypergeometric} 和~\eqref{eq:projective-space-j-descendants} 给出。这里遵循 \SourceCite{Kim2} 的约定，其方法受到 \SourceCite{Givental4} 的启发。证明的关键步骤是引理~\ref{lem:projective-space-j-recursion} 给出的递推。\hyperref[subsec:givental-general-bundle-proof]{第~\ref*{subsec:givental-general-bundle-proof} 小节}将把这里的方法与一些重要思想结合，证明镜像定理（定理~\ref{thm:givental-projective-mirror}）。

策略是构造 $I_{\mathbb P^n}$ 和 $J_{\mathbb P^n}$ 的等变版本，再用类似于\hyperref[chap:localization]{第九章}和\hyperref[sec:quintic-mirror-theorem]{第~\ref*{sec:quintic-mirror-theorem} 节}的局部化技巧证明二者相等。采用例题~\ref{ex:projective-equivariant-localization} 所述 $T=(\mathbb C^*)^{n+1}$ 在 $\mathbb P^n$ 上的标准作用。它诱导 $\overline M_{0,2}(\mathbb P^n,d)$ 上的作用，使赋值映射 $e_1:\overline M_{0,2}(\mathbb P^n,d)\to\mathbb P^n$ 为等变映射。

先定义式~\eqref{eq:projective-space-j-descendants} 形式的 $J_{\mathbb P^n}$ 的等变版本：
\begin{equation}\label{eq:projective-space-equivariant-j}
J_T=e^{(t_0+t_1p)/\hbar}\left(1+\sum_{d=1}^{\infty}q^d e_{1!}\left(\frac1{\hbar-c_T}\right)\right),
\end{equation}
其中 $c_T$ 是等变 Chern 类 $c_1^T(\mathcal L_1)$ 的符号。如常设 $q=e^{t_1}$，$p$ 是 $\mathbb P^n$ 的\IndexLocation{indexloc:0250:389}等变超平面类。$J_T$ 取值于 $H_T^*(\mathbb P^n)$，取\IndexLocation{indexloc:0945:389}非等变极限 $\lambda_i\to0$、$p\to H$、$c_T\to c_1(\mathcal L_1)$，便能恢复 $J_{\mathbb P^n}$。更确切地，$J_{\mathbb P^n}=i_{\mathbb P^n}^*(J_T)$，其中 $i_{\mathbb P^n}$ 如\hyperref[sec:localization-theorem]{第~\ref*{sec:localization-theorem} 节}中定义。

去掉式~\eqref{eq:projective-space-equivariant-j} 的指数因子会更方便。令
\[
S=1+\sum_{d=1}^{\infty}q^d e_{1!}\left(\frac1{\hbar-c_T}\right),
\]
于是 $J_T=e^{(t_0+t_1p)/\hbar}S$。
\SourcePage{390}
接着把 $S$ 分解为较易研究的部分。令 $Z_i=\int_{(\mathbb P^n)_T}S\cup\phi_i$，其中 $\phi_i=\prod_{k\ne i}(p-\lambda_k)$ 是例题~\ref{ex:projective-equivariant-localization} 中 $H_T^*(\mathbb P^n)\otimes_{H^*(BT)}\mathcal R_T$ 的一组基。关键在于，$Z_i$ 足以确定 $S$，从而确定 $J_T$，最终也确定 $J_{\mathbb P^n}$：因为可以由 $S=\sum_{i=0}^nZ_i\cup\phi^i$ 恢复 $S$，其中 $\phi^i$ 是与 $\phi_i$ 对偶的基。

投影公式给出
\[
\int_{(\mathbb P^n)_T}e_{1!}\left(\frac1{\hbar-c_T}\right)\cup\phi_i
=\int_{\overline M_{0,2}(\mathbb P^n,d)_T}\frac{e_1^*(\phi_i)}{\hbar-c_T},
\]
所以可将 $Z_i$ 写为
\begin{equation}\label{eq:projective-space-j-fixed-components}
Z_i=1+\sum_{d=1}^{\infty}q^d\int_{\overline M_{0,2}(\mathbb P^n,d)_T}\frac{e_1^*(\phi_i)}{\hbar-c_T}.
\end{equation}
这里有一个简化。$\dim\overline M_{0,2}(\mathbb P^n,d)=(n+1)d+n-1$，因此其上按复次数计小于 $(n+1)d+n-1$ 的上同调类，积分均为零。另外，$\phi_i$ 的复次数为 $n$。将 $1/(\hbar-c_T)$ 按 $c_T/\hbar$ 的幂展开为几何级数，再代入式~\eqref{eq:projective-space-j-fixed-components}，可知含 $c_T^k$ 的项只有在 $k\ge(n+1)d-1$ 时才可能有贡献。因此，截去几何级数的低次项不会改变结果。出于稍后将明白的理由，我们截去的项略少于所能截去的项，从 $c_T^d$ 开始展开；这是可行的，因为 $d\le(n+1)d-1$。利用恒等式
\begin{equation}\label{eq:projective-space-descendant-truncation}
\sum_{k=d}^{\infty}\frac{c_T^k}{\hbar^{k+1}}=\left(\frac{c_T}{\hbar}\right)^d\frac1{\hbar-c_T},
\end{equation}
得到代替式~\eqref{eq:projective-space-j-fixed-components} 的另一形式：
\begin{equation}\label{eq:projective-space-j-truncated-components}
Z_i=Z_i(q,\hbar)=1+\sum_{d=1}^{\infty}\left(\frac q\hbar\right)^d
\int_{\overline M_{0,2}(\mathbb P^n,d)_T}\frac{e_1^*(\phi_i)c_T^d}{\hbar-c_T}.
\end{equation}
采用推论~\ref{cor:orbifold-localization-integration} 的形式，在 $\overline M_{0,2}(\mathbb P^n,d)$ 上局部化，计算 $Z_i$ 中出现的积分。每个积分
\[
\int_{\overline M_{0,2}(\mathbb P^n,d)_T}\frac{e_1^*(\phi_i)c_T^d}{\hbar-c_T}
\]
都会成为固定点集各分支上的积分之和；由\hyperref[subsec:kontsevich-localization]{第~\ref*{subsec:kontsevich-localization} 小节}，这些分支用某些带标记图 $\Gamma$ 描述。

设图 $\Gamma$ 对应于固定点分支 $\overline M_\Gamma\subset\overline M_{0,2}(\mathbb P^n,d)$。如\hyperref[subsec:kontsevich-localization]{第~\ref*{subsec:kontsevich-localization} 小节}，有包含映射 $i_\Gamma:\overline M_\Gamma\to\overline M_{0,2}(\mathbb P^n,d)$。首要任务是确定哪些分支在局部化后对上述积分有贡献。回顾 $T$ 在 $\mathbb P^n$ 上作用的固定点为 $q_j$，$0\le j\le n$。于是，除非 $e_1(\overline M_\Gamma)=q_i$，否则 $i_\Gamma^*(e_1^*(\phi_i))$ 为零。这是因为，若 $e_1(\overline M_\Gamma)=q_j$、$j\ne i$，则
\[
i_\Gamma^*(e_1^*(\phi_i))=i_{q_j}^*(\phi_i)=0,
\]
最后的等号由式~\eqref{eq:projective-dual-class-restriction} 得到。因此，局部化时只须对满足 $e_1(\overline M_\Gamma)=q_i$ 的 $\Gamma$ 求和。（注意，这与定理~\ref{thm:geometric-euler-data-gluing} 证明开头的做法相似。）
\SourcePage{391}
因为等变超平面类 $p$ 在 $q_i$ 上限制为 $\lambda_i$，式~\eqref{eq:projective-dual-class-restriction} 给出 $i_\Gamma^*(e_1^*\phi_i)=i_{q_i}^*(\phi_i)=\prod_{k\ne i}(\lambda_i-\lambda_k)$。也用符号 $c_\Gamma$ 表示 $c_T$ 在固定点分支上的限制。更准确地，定义
\[
\frac1{\hbar-c_\Gamma}=\sum_{k=0}^{\infty}\hbar^{-(k+1)}c_1^T(i_\Gamma^*\mathcal L_1)^k.
\]
最后，令 $a_\Gamma$ 为具有图 $\Gamma$ 的一般稳定映射的自同构群的阶。于是推论~\ref{cor:orbifold-localization-integration} 给出
\begin{equation}\label{eq:projective-space-j-localization}
Z_i(q,\hbar)=1+\sum_{d=1}^{\infty}\left(\frac q\hbar\right)^d\sum_\Gamma^{(i)}
\int_{(\overline M_\Gamma)_T}
\frac{c_\Gamma^d\prod_{k\ne i}(\lambda_i-\lambda_k)}{a_\Gamma(\hbar-c_\Gamma)\operatorname{Euler}_T(N_\Gamma)},
\end{equation}
其中 $\sum_\Gamma^{(i)}$ 表示仅对满足 $e_1(\overline M_\Gamma)=q_i$ 的 $\Gamma$ 求和。

还可以进一步减少需要考虑的分支。若 $(f,C,p_1,p_2)\in\overline M_\Gamma$，则 $e_1(\overline M_\Gamma)=q_i$ 意味着 $e_1(f)=f(p_1)=q_i$。令 $C'\subset C$ 为含 $p_1$ 的分支，区分以下两种情形：
\begin{itemize}
\item A 型：含 $p_1$ 的分支 $C'\subset C$ 被 $f$ 映到一条曲线上。
\item B 型：含 $p_1$ 的分支 $C'\subset C$ 被 $f$ 映到一个点上。
\end{itemize}
断言：B 型图 $\Gamma$ 对 $Z_i$ 没有贡献，因为此时 $c_\Gamma^d=0$。

为证明 B 型图的 $c_\Gamma^d=0$，须回顾\hyperref[subsec:kontsevich-localization]{第~\ref*{subsec:kontsevich-localization} 小节}对 $\overline M_\Gamma$ 的描述。图 $\Gamma$ 的顶点对应于 $f^{-1}(\{q_0,\ldots,q_n\})$ 的连通分支。因为 $\Gamma$ 为 B 型，$C'$ 包含在 $f^{-1}(q_i)$ 的某个连通分支 $C_v$ 内，$v$ 为相应顶点。如该小节，$n(v)$ 是 $C_v$ 中标记点的数目，加上将 $C_v$ 连到 $f$ 的次数 $d_e\ge1$ 的分支 $C_e$ 的结点数目。我们研究的是 $\overline M_{0,2}(\mathbb P^n,d)$，所以标记点至多有两个。另外，由图 $\Gamma$ 的定义，$\sum_ed_e=d$，所以这种结点至多有 $d$ 个。于是 $n(v)\le d+2$。

$n(v)$ 很重要，因为如该小节所解释，$\overline M_{0,n(v)}$ 是给出 $\overline M_\Gamma$ 的乘积中的一个因子。我们将研究 $c_\Gamma^d$ 在 $\overline M_{0,n(v)}$ 上的限制。因为 $n(v)\le d+2$，有 $\dim\overline M_{0,n(v)}=n(v)-3\le d-1$。此外，此时 $c_\Gamma$ 来自 $\overline M_{0,n(v)}$ 上相应丛 $\mathcal L_1$ 的等变 Chern 类 $c_1^T(\mathcal L_1)$。$T$ 在 $\overline M_{0,n(v)}$ 上的作用平凡，故 $H_T^*(\overline M_{0,n(v)})\simeq H^*(\overline M_{0,n(v)})\otimes H^*(BT)$。$T$ 在 $\mathcal L_1$ 上也平凡作用，因此在此同构下 $c_1^T(\mathcal L_1)=c_1(\mathcal L_1)\otimes1$。由于 $c_1(\mathcal L_1)$ 的实次数为 $2$，且 $\dim\overline M_{0,n(v)}\le d-1$，有 $c_1(\mathcal L_1)^d=0$。由此证明所断言的 $c_\Gamma^d=0$。

B 型图对积分没有贡献，所以用 $\sum_\Gamma^A$ 表示对 A 型图求和。仍假设 $e_1(\overline M_\Gamma)=q_i$，尽管记号中省略了下标 $i$。于是式~\eqref{eq:projective-space-j-localization} 成为
\begin{equation}\label{eq:projective-space-j-type-a-localization}
Z_i(q,\hbar)=1+\sum_{d=1}^{\infty}\sum_\Gamma^A\left(\frac q\hbar\right)^d
\int_{(\overline M_\Gamma)_T}\frac{c_\Gamma^d\prod_{k\ne i}(\lambda_i-\lambda_k)}{a_\Gamma(\hbar-c_\Gamma)\operatorname{Euler}_T(N_\Gamma)}.
\end{equation}
由此可证明以下引理。
\SourcePage{392}
\begin{lemma}[\IndexLocation{indexloc:0861:392}正则性引理]\label{lem:projective-space-j-regularity}
表达式 $Z_i$ 是环 $\mathbb Q(\lambda_0,\ldots,\lambda_n,\hbar)[[q]]$ 的元素。每个 $q^d$ 的系数都是 $\lambda_k$ 和 $\hbar$ 的有理函数，并且对于所有 $j\ne i$、$k\ge1$，在 $\hbar=(\lambda_i-\lambda_j)/k$ 处正则。
\end{lemma}
\begin{proof}
须考察式~\eqref{eq:projective-space-j-type-a-localization} 中 $q^d$ 的系数。为此，设 $\Gamma$ 为 A 型。分母中的 $\operatorname{Euler}_T(N_\Gamma)$ 属于 $H_T^*(\overline M_\Gamma)$，其局部化贡献关于 $\lambda_k$ 为有理函数，故只剩下分析 $c_\Gamma$。

回顾 $c_\Gamma=c_1^T(i_\Gamma^*\mathcal L_1)$。\hyperref[subsec:gravitational-correlator-definition]{第~\ref*{subsec:gravitational-correlator-definition} 小节}定义的 $\mathcal L_1$ 是 $\overline M_{0,2}(\mathbb P^n,d)$ 上的线丛，其在 $(f,C,p_1,p_2)$ 处的纤维为余切空间 $T_{p_1}^*C$。为计算 $c_1^T(i_\Gamma^*\mathcal L_1)$，须求 $(f,C,p_1,p_2)\in\overline M_\Gamma$ 时 $T$ 在 $T_{p_1}^*C$ 上作用的权。如前，令 $C'\subset C$ 为含 $p_1$ 的不可约分支。$\Gamma$ 为 A 型，因此 $C'$ 被映到直线 $\overline{q_iq_j}$ 上，次数为 $d'$。在 $q_i\in\overline{q_iq_j}\simeq\mathbb P^1$ 处，余切空间 $T_{q_i}^*\mathbb P^1$ 的权为 $\lambda_j-\lambda_i$；这可由例题~\ref{ex:projective-equivariant-localization} 的方法得出。因 $f|_{C'}$ 的次数为 $d'$，$T$ 在 $T_{p_1}^*C=T_{p_1}^*C'$ 上作用的权为 $(\lambda_j-\lambda_i)/d'$。所以 $c_\Gamma=(\lambda_j-\lambda_i)/d'$，从而
\[
\int_{(\overline M_\Gamma)_T}\frac{c_\Gamma^d\prod_{k\ne i}(\lambda_i-\lambda_k)}{a_\Gamma(\hbar-c_\Gamma)\operatorname{Euler}_T(N_\Gamma)}
=\int_{(\overline M_\Gamma)_T}\frac{((\lambda_j-\lambda_i)/d')^d\prod_{k\ne i}(\lambda_i-\lambda_k)}{a_\Gamma(\hbar-(\lambda_j-\lambda_i)/d')\operatorname{Euler}_T(N_\Gamma)}.
\]
由此容易验证所需结论。
\end{proof}
刚证明的引理表明，可以在 $Z_j$ 中代入 $\hbar=(\lambda_j-\lambda_i)/d$，得到有定义的表达式。这一点用于下面的递推引理。
\begin{lemma}[递推引理]\label{lem:projective-space-j-recursion}
对 $0\le i,j\le n$、$i\ne j$，令
\[
C_{i,j}(d)=\frac1{d!\displaystyle\prod_{\substack{1\le m\le d,\ k\ne i\\(k,m)\ne(j,d)}}(\lambda_i-\lambda_k+m(\lambda_j-\lambda_i)/d)}.
\]
则
\[
Z_i(q,\hbar)=1+\sum_{j\ne i}\sum_{d=1}^{\infty}\left(\frac q\hbar\right)^d
\frac{C_{i,j}(d)}{\lambda_i-\lambda_j+d\hbar}
Z_j\left(\frac q\hbar\frac{\lambda_j-\lambda_i}{d},\frac{\lambda_j-\lambda_i}{d}\right).
\]
\end{lemma}
\begin{proof}
从式~\eqref{eq:projective-space-j-type-a-localization} 推出此递推。如式~\eqref{eq:projective-space-j-fixed-components} 后的讨论，原本可把 $1/(\hbar-c_T)$ 的几何级数截到式~\eqref{eq:projective-space-descendant-truncation} 中从 $c_T^{(n+1)d-1}/\hbar^{(n+1)d}$ 开始。若这样做，会得到 \SourceCite{Givental2} 中的递推。这里的选择导出 \SourceCite{Kim2} 的递推关系，并避免了 \SourceCite{Givental2} 所用的 $q$ 的变量变换。

式~\eqref{eq:projective-space-j-type-a-localization} 中的图 $\Gamma$ 全为 A 型。对 $(f,C,p_1,p_2)\in\overline M_\Gamma$，有 $f(p_1)=q_i$。若 $C'\subset C$ 是含 $p_1$ 的分支，则存在 $j\ne i$，使 $f|_{C'}:C'\to\overline{q_iq_j}$ 为次数 $d'$ 的覆盖。若还有余下曲线，此时图 $\Gamma$ 的定义给出 $C=C'\cup C''$，其中 $C''$ 对应于图 $\Gamma''$，且 $p_2\in C''$。令 $p=C'\cap C''$，则 $q_j=f(p)$。把 $p$ 看作 $C''$ 上第一个标记点，则 $(f|_{C''},C'',p,p_2)$ 是
\SourcePage{393}
$\overline M_{0,2}(\mathbb P^n,d'')$ 中具有图 $\Gamma''$ 的二点稳定映射，其中 $d=d'+d''$。当 $C=C'$ 时，取 $d''=0$，相应贡献按常数项约定为 $1$。对于有余下曲线的图，断言
\begin{equation}\label{eq:projective-space-edge-removal-euler}
\begin{aligned}
a_\Gamma&=d'a_{\Gamma''},\\
c_\Gamma&=\frac{\lambda_j-\lambda_i}{d'},\\
\operatorname{Euler}_T(N_\Gamma)&=\operatorname{Euler}_T(N_{\Gamma''})
\left(\frac{\lambda_j-\lambda_i}{d'}-c_{\Gamma''}\right)
\prod_{m,k}^{(i)}\left(m\frac{\lambda_j-\lambda_i}{d'}+\lambda_i-\lambda_k\right),
\end{aligned}
\end{equation}
其中 $\prod_{m,k}^{(i)}$ 遍历 $0\le m\le d'-1$、$0\le k\le n$，但排除 $(m,k)=(0,i)$。

$a_\Gamma$ 的公式来自 $f|_{C'}$ 的自同构构成阶为 $d'$ 的循环群，正如式~\eqref{eq:fixed-graph-automorphisms} 及其后的讨论。$c_\Gamma$ 的公式已在引理~\ref{lem:projective-space-j-regularity} 的证明中得到。最后，对等变法丛 $N_\Gamma$ 作滤过，计算 $\operatorname{Euler}_T(N_\Gamma)$。因子 $\operatorname{Euler}_T(N_{\Gamma''})$ 来自 $f|_{C''}$ 的形变。因子
\[
\frac{\lambda_j-\lambda_i}{d'}-c_{\Gamma''}
\]
来自在结点 $p$ 处应用引理~\ref{lem:node-smoothing-tangent-tensor}，类似于定理~\ref{thm:geometric-euler-data-gluing} 的证明。最后的因子来自保持点 $p$ 不动的 $f|_{C'}$ 的形变。这些形变由向量场 $v_{m,k}=z_i^{d'-m}z_j^m\partial/\partial x_k$ 给出，其中 $0\le m<d'$。这里 $x_0,\ldots,x_n$ 是 $\mathbb P^n$ 上的坐标，$z_i,z_j$ 是 $C'$ 上的坐标，使 $f|_{C'}$ 由 $x_i=z_i^{d'}$、$x_j=z_j^{d'}$ 给出。$v_{m,k}$ 的权显然为
\[
(d'-m)\frac{\lambda_i}{d'}+m\frac{\lambda_j}{d'}-\lambda_k
=m\frac{\lambda_j-\lambda_i}{d'}+\lambda_i-\lambda_k.
\]
注意，向量场 $v_{d',k}$ 不保持 $p$ 不动，故没有贡献。另外，$v_{0,i}$ 除一个非零标量因子外，是 $C'$ 上保持 $p_1$ 和 $p$ 不动的向量场 $z_i\partial/\partial z_i$ 的像。它来自稳定映射的自同构，所以也须排除 $v_{0,i}$。把剩余各权相乘，便得到 $\operatorname{Euler}_T(N_\Gamma)$ 公式中的最后一个因子。

将式~\eqref{eq:projective-space-edge-removal-euler} 代入式~\eqref{eq:projective-space-j-type-a-localization}，计算得
\begin{equation}\label{eq:projective-space-j-recursion-factorization}
\begin{aligned}
Z_i=1+\sum_{j\ne i}\sum_{d'=1}^{\infty}
&\frac{(q/\hbar)^{d'}((\lambda_j-\lambda_i)/d')^{d'}\prod_{k\ne i}(\lambda_i-\lambda_k)}
{\left(\prod_{k\ne j}(\lambda_j-\lambda_k)\right)(\hbar-(\lambda_j-\lambda_i)/d')d'}\\
&\quad\times\frac{1}{\prod_{m,k}^{(i)}(m(\lambda_j-\lambda_i)/d'+\lambda_i-\lambda_k)}\\
&\quad\times\Biggl(1+\sum_{d''=1}^{\infty}\sum_{\Gamma''}^{(j)}\left(\frac q\hbar\right)^{d''}\\
&\qquad\quad\times\int_{(\overline M_{\Gamma''})_T}
\frac{((\lambda_j-\lambda_i)/d')^{d''}\prod_{k\ne j}(\lambda_j-\lambda_k)}
{\operatorname{Euler}_T(N_{\Gamma''})a_{\Gamma''}((\lambda_j-\lambda_i)/d'-c_{\Gamma''})}\Biggr).
\end{aligned}
\end{equation}
得到此式时，插入了相互约去的因子 $\prod_{k\ne j}(\lambda_j-\lambda_k)$，并使用 $d=d'+d''$。

现在从式~\eqref{eq:projective-space-j-recursion-factorization} 出发，注意以下各点，即可证明所需的 $Z_i$ 公式。此式最后一行正是
\[
Z_j\left(q(\lambda_j-\lambda_i)/(d'\hbar),(\lambda_j-\lambda_i)/d'\right),
\]
\SourcePage{394}
这可由式~\eqref{eq:projective-space-j-localization} 看出。至于第一行，其中的大分式含有因子
\[
\left(\frac q\hbar\right)^{d'}\frac1{\lambda_i-\lambda_j+d'\hbar}.
\]
所以，引理归结为证明剩余因子为 $C_{i,j}(d')$。也就是须证明
\begin{equation}\label{eq:projective-space-recursion-coefficient-cancellation}
C_{i,j}(d')=\frac{((\lambda_j-\lambda_i)/d')^{d'}\prod_{k\ne i}(\lambda_i-\lambda_k)}{\left(\prod_{k\ne j}(\lambda_j-\lambda_k)\right)\prod_{m,k}^{(i)}(m(\lambda_j-\lambda_i)/d'+\lambda_i-\lambda_k)}.
\end{equation}
为此，令 $A(m,k)=m(\lambda_j-\lambda_i)/d'+\lambda_i-\lambda_k$，注意三个涉及 $A(m,k)$ 的恒等式：
\[
\begin{aligned}
(d')!\left(\frac{\lambda_j-\lambda_i}{d'}\right)^{d'}&=\prod_{m=1}^{d'}A(m,i),\\
\lambda_i-\lambda_k&=A(0,k),\\
\lambda_j-\lambda_k&=A(d',k).
\end{aligned}
\]
于是可用 $A(m,k)$ 表示式~\eqref{eq:projective-space-recursion-coefficient-cancellation} 两边，并在分母中保留因子 $(d')!$。略去右边若干相互约去的 $A(m,k)$ 因子，就得到所需结果。
\end{proof}
初看时，引理~\ref{lem:projective-space-j-recursion} 的恒等式未必明显呈现为递推。将 $Z_i$ 写成
\[
Z_i(q,\hbar)=\sum_{d=0}^{\infty}Z_{i,d}(\hbar)q^d,
\]
其中 $Z_{i,0}(\hbar)=1$，递推便清楚了。对于 $d\ge1$，引理给出
\[
Z_{i,d}(\hbar)=\sum_{\substack{j\ne i,\ d=d'+d''\\d'\ge1}}
\frac{C_{i,j}(d')}{\lambda_i-\lambda_j+d'\hbar}
\left(\frac{\lambda_j-\lambda_i}{d'\hbar}\right)^{d''}\hbar^{-d'}
Z_{j,d''}\left(\frac{\lambda_j-\lambda_i}{d'}\right).
\]
因为 $d'\ge1$，$Z_{i,d}$ 由 $d''<d$ 的 $Z_{j,d''}$ 确定。再由引理给出 $Z_{i,0}=1$，可见 $Z_i$ 由引理中的恒等式唯一刻画。

有了\IndexLocation{indexloc:0860:394}递推引理，终于可以证明命题~\ref{prop:projective-space-i-equals-j}。
\begin{proof}[命题~\ref{prop:projective-space-i-equals-j} 的证明]
计划是构造 $I_{\mathbb P^n}$ 的等变版本 $I_T$，再证明它满足上述递推。回顾式~\eqref{eq:projective-space-i-hypergeometric}，
\[
I_{\mathbb P^n}=e^{(t_0+t_1H)/\hbar}\sum_d\frac{q^d}{((H+\hbar)\cdots(H+d\hbar))^{n+1}}.
\]
考虑等变表达式
\[
I_T=e^{(t_0+t_1p)/\hbar}\sum_{d=0}^{\infty}\frac{q^d}{\prod_{j=0}^n\prod_{m=1}^d(p-\lambda_j+m\hbar)},
\]
\SourcePage{395}
其非等变极限为 $I_{\mathbb P^n}$。目标是证明 $I_T=J_T$。先按对 $J_T$ 所作的同样方式分解 $I_T$。去掉指数因子，令
\[
S'=\sum_{d=0}^{\infty}\frac{q^d}{\prod_{j=0}^n\prod_{m=1}^d(p-\lambda_j+m\hbar)},
\]
再令 $Z'_i=\int_{(\mathbb P^n)_T}S'\cup\phi_i$。计算如下：把式~\eqref{eq:linear-model-restriction-integral} 的做法用于 $\mathbb P^n$，可知 $Z'_i=i_{q_i}^*(S')=S'(\lambda_i)$，从而
\[
Z'_i=1+\sum_{d=1}^{\infty}\frac{q^d}{\prod_{j=0}^n\prod_{m=1}^d(\lambda_i-\lambda_j+m\hbar)}.
\]
此式清楚表明 $Z'_i$ 属于 $\mathbb Q(\lambda_0,\ldots,\lambda_n,\hbar)[[q]]$，且满足引理~\ref{lem:projective-space-j-regularity} 的正则性。另外，直接计算可知 $Z'_i$ 满足引理~\ref{lem:projective-space-j-recursion} 的递推。由该引理证明后的讨论中指出的唯一性，所有 $i$ 均有 $Z_i=Z'_i$。为结束证明，回顾式~\eqref{eq:projective-space-equivariant-j} 后的讨论，
\[
J_T=e^{(t_0+t_1p)/\hbar}\sum_{i=0}^nZ_i\cup\phi^i.
\]
类似推理给出
\[
I_T=e^{(t_0+t_1p)/\hbar}\sum_{i=0}^nZ'_i\cup\phi^i,
\]
故 $I_T=J_T$。取非等变极限 $\lambda_i\to0$、$p\to H$，即得所需的 $I_{\mathbb P^n}=J_{\mathbb P^n}$。
\end{proof}

\subsection{定理~\ref{thm:givental-projective-mirror} 的证明}\label{subsec:givental-general-bundle-proof}
为证明\hyperref[subsec:givental-nef-projective-mirror]{第~\ref*{subsec:givental-nef-projective-mirror} 小节}陈述的镜像定理，将推广命题~\ref{prop:projective-space-i-equals-j} 证明中的各个步骤。不过，此时得到的递推（引理~\ref{lem:complete-intersection-j-recursion}）弱得多，所以还需要下面的引理~\ref{lem:complete-intersection-double-construction}、\ref{lem:complete-intersection-j-uniqueness} 和~\ref{lem:complete-intersection-p-transformations}。论述基于 \SourceCite{Givental2,Pandharipande3}，但也采用 \SourceCite{Kim2,Givental4} 的一些约定。这一讨论同时包含了射影空间中完全交的猜想~\ref{conj:quantum-hyperplane-section-principle} 的证明。

与命题~\ref{prop:projective-space-i-equals-j} 的证明相同，先给出 $J_{\mathcal V}$ 的等变版本。由式~\eqref{eq:givental-projective-bundle-j-symbolic} 中 $J_{\mathcal V}$ 的定义，令
\begin{equation}\label{eq:complete-intersection-equivariant-j}
J_T=e^{(t_0+t_1p)/\hbar}\operatorname{Euler}_T(\mathcal V)
\left(1+\sum_{d=1}^{\infty}q^d e_{1!}\left(\frac{\operatorname{Euler}_T(\mathcal V'_{d,2,1})}{\hbar-c_T}\right)\right).
\end{equation}
零秩平凡丛的 Euler 类为 $1$，所以当 $\mathcal V=0$ 时，此式化为式~\eqref{eq:projective-space-equivariant-j}。注意，$T$ 在 $\mathcal V=\bigoplus_{i=1}^{\ell}\mathcal O_{\mathbb P^n}(a_i)$ 上自然作用，故 $\operatorname{Euler}_T(\mathcal V'_{d,2,1})$ 有意义。如\hyperref[subsec:givental-projective-space-proof]{第~\ref*{subsec:givental-projective-space-proof} 小节}，去掉指数因子，令
\begin{equation}\label{eq:complete-intersection-equivariant-s}
S_{\mathcal V}=S_{\mathcal V}(q,\hbar)=\operatorname{Euler}_T(\mathcal V)
\left(1+\sum_{d=1}^{\infty}q^d e_{1!}\left(\frac{\operatorname{Euler}_T(\mathcal V'_{d,2,1})}{\hbar-c_T}\right)\right).
\end{equation}
也令
\[
\widetilde S_{\mathcal V}=1+\sum_{d>0}q^d e_{1!}\left(\frac{\operatorname{Euler}_T(\mathcal V'_{d,2,1})}{\hbar-c_T}\right),
\]
\SourcePage{396}
于是 $S_{\mathcal V}=\operatorname{Euler}_T(\mathcal V)\widetilde S_{\mathcal V}$。如例题~\ref{ex:projective-equivariant-localization} 所示，张量到 $\lambda_i$ 的有理函数域后，$p$ 可逆，故 $\operatorname{Euler}_T(\mathcal V)=(\prod_{i=1}^{\ell}a_i)p^\ell$ 也可逆。特别地，$S_{\mathcal V}$ 唯一确定 $\widetilde S_{\mathcal V}$，且
\[
\widetilde S_{\mathcal V}=(\operatorname{Euler}_T(\mathcal V))^{-1}S_{\mathcal V}.
\]
这里再次允许 $\lambda_i$ 的有理函数，使表达式有意义。本章余下部分将经常采用类似记号，用符号上的波浪号表示除以 $\operatorname{Euler}_T(\mathcal V)$。

最后，令
\[
Z_{i,\mathcal V}=\int_{(\mathbb P^n)_T}\widetilde S_{\mathcal V}\cup\phi_i=i_{q_i}^*(\widetilde S_{\mathcal V}).
\]
要研究 $Z_{i,\mathcal V}$ 是否与\hyperref[subsec:givental-projective-space-proof]{第~\ref*{subsec:givental-projective-space-proof} 小节}研究的函数 $Z_i$ 具有同样的性质。第一步很容易，因为可直接把式~\eqref{eq:projective-space-j-fixed-components} 推广为
\begin{equation}\label{eq:complete-intersection-j-fixed-components}
Z_{i,\mathcal V}=1+\sum_{d=1}^{\infty}q^d
\int_{\overline M_{0,2}(\mathbb P^n,d)_T}\frac{\operatorname{Euler}_T(\mathcal V'_{d,2,1})e_1^*(\phi_i)}{\hbar-c_T}.
\end{equation}
但由于 $\operatorname{Euler}_T(\mathcal V'_{d,2,1})$ 的出现，式~\eqref{eq:projective-space-descendant-truncation} 和~\eqref{eq:projective-space-j-truncated-components} 所作的截断一般不再适用。因此，不按式~\eqref{eq:projective-space-j-truncated-components} 到~\eqref{eq:projective-space-j-localization} 的方式局部化，而直接对式~\eqref{eq:complete-intersection-j-fixed-components} 中的积分应用局部化，容易得到
\begin{equation}\label{eq:complete-intersection-j-localization}
Z_{i,\mathcal V}=1+\sum_{d=1}^{\infty}q^d\sum_\Gamma^{(i)}
\int_{(\overline M_\Gamma)_T}\frac{i_\Gamma^*(\operatorname{Euler}_T(\mathcal V'_{d,2,1}))\prod_{k\ne i}(\lambda_i-\lambda_k)}{a_\Gamma(\hbar-c_\Gamma)\operatorname{Euler}_T(N_\Gamma)},
\end{equation}
其中 $\sum_\Gamma^{(i)}$ 与式~\eqref{eq:projective-space-j-localization} 中含义相同。

继续证明之前，暂停一下，对截断作一个有趣的观察。关键是，一般情形无法截断，这正是镜像定理需要变量变换的原因。例如，若 $\mathcal V=\bigoplus_{i=1}^{\ell}\mathcal O(a_i)$ 满足 $\sum_{i=1}^{\ell}a_i<n$，截断仍然可行。直接验证时用到 $\mathcal V'_{d,2,1}$ 的秩为 $d\sum_{i=1}^{\ell}a_i<dn$。所以 $d+\operatorname{rank}(\mathcal V'_{d,2,1})<d+nd$，从而 $d+\operatorname{rank}(\mathcal V'_{d,2,1})\le(n+1)d-1$，正好满足与式~\eqref{eq:projective-space-descendant-truncation} 前的论证类似的要求。此时，下面的证明显著简化，变得很像\hyperref[subsec:givental-projective-space-proof]{第~\ref*{subsec:givental-projective-space-proof} 小节}所作的证明。结果是，当 $\sum_{i=1}^{\ell}a_i<n$ 时，镜像定理断言 $I_{\mathcal V}=J_{\mathcal V}$。换言之，此时定理~\ref{thm:givental-projective-mirror} 中 $f_0=f_1=h=0$。例题~\ref{ex:fano-hypersurface-i-j-equality} 已给出这一情形的一个特例。一般地，当 $\sum_{i=1}^{\ell}a_i\le n+1$ 时，将看到无法截断如何使递推变复杂，并迫使我们改变变量。

回到镜像定理的证明，注意以下 \SourceCite{Pandharipande3} 的结果，它推广了引理~\ref{lem:projective-space-j-regularity}。
\begin{lemma}[\IndexLocation{indexloc:0861:396}正则性引理]\label{lem:complete-intersection-j-regularity}
表达式 $Z_{i,\mathcal V}$ 是环 $\mathbb Q(\lambda_0,\ldots,\lambda_n,\hbar)[[q]]$ 的元素。每个 $q^d$ 的系数都是 $\lambda_k$ 和 $\hbar$ 的有理函数，且对于所有 $j\ne i$、$k\ge1$，在 $\hbar=(\lambda_i-\lambda_j)/k$ 处正则。
\end{lemma}
\SourcePage{397}
\begin{proof}
对于 $\sum_\Gamma^{(i)}$ 中出现的图 $\Gamma$，须研究积分
\begin{equation}\label{eq:complete-intersection-regularity-local-integral}
\int_{(\overline M_\Gamma)_T}\frac{i_\Gamma^*(\operatorname{Euler}_T(\mathcal V'_{d,2,1}))\prod_{k\ne i}(\lambda_i-\lambda_k)}{a_\Gamma(\hbar-c_\Gamma)\operatorname{Euler}_T(N_\Gamma)}.
\end{equation}
\hyperref[subsec:givental-projective-space-proof]{第~\ref*{subsec:givental-projective-space-proof} 小节}把图 $\Gamma$ 分成 A 型和 B 型。对于 A 型图，引理~\ref{lem:projective-space-j-regularity} 的证明表明，存在 $d'\ge1$、$j\ne i$，使 $c_\Gamma=(\lambda_j-\lambda_i)/d'$。因此，这些分支的上述积分成为
\[
\int_{(\overline M_\Gamma)_T}\frac{i_\Gamma^*(\operatorname{Euler}_T(\mathcal V'_{d,2,1}))\prod_{k\ne i}(\lambda_i-\lambda_k)}{a_\Gamma(\hbar-(\lambda_j-\lambda_i)/d')\operatorname{Euler}_T(N_\Gamma)}.
\]
与引理~\ref{lem:projective-space-j-regularity} 相同，此表达式显然关于 $\lambda_k,\hbar$ 为有理函数，并对 $k\ge1$ 在 $\hbar=(\lambda_i-\lambda_j)/k$ 处正则。

再设 $\Gamma$ 为 B 型。在引理~\ref{lem:projective-space-j-regularity} 前的讨论中，已证明此时 $c_\Gamma^d=0$。于是，把式~\eqref{eq:complete-intersection-regularity-local-integral} 分母中的 $1/(\hbar-c_\Gamma)$ 展开为几何级数时，可忽略所有含 $c_\Gamma^k/\hbar^{k+1}$、$k\ge d$ 的项。因此，只须考虑 $k<d$ 的项。每项都是 $\lambda_i$ 的有理函数乘以 $\hbar^{-1}$ 的正次幂，所以关于 $\hbar$ 和 $\lambda_i$ 为有理函数，且满足所断言的正则性。
\end{proof}
如\hyperref[subsec:givental-projective-space-proof]{第~\ref*{subsec:givental-projective-space-proof} 小节}，此引理表明，可以在 $Z_{j,\mathcal V}$ 中代入 $\hbar=(\lambda_j-\lambda_i)/d$，得到有定义的表达式。
\begin{lemma}[递推引理]\label{lem:complete-intersection-j-recursion}
设 $\mathcal V=\bigoplus_{k=1}^{\ell}\mathcal O_{\mathbb P^n}(a_k)$。对于 $0\le i,j\le n$、$i\ne j$，令
\[
C_{i,j}^{\mathcal V}(d)=\frac{\prod_{k=1}^{\ell}\prod_{m=1}^{a_kd}(a_k\lambda_i+m(\lambda_j-\lambda_i)/d)}{d!\displaystyle\prod_{\substack{1\le m\le d,\ k\ne i\\(k,m)\ne(j,d)}}(\lambda_i-\lambda_k+m(\lambda_j-\lambda_i)/d)}.
\]
则差
\[
\begin{aligned}
R_i={}&Z_{i,\mathcal V}(q,\hbar)-\sum_{j\ne i}\sum_{d=1}^{\infty}\left(\frac q\hbar\right)^d
\frac{C_{i,j}^{\mathcal V}(d)}{\lambda_i-\lambda_j+d\hbar}\\
&\hspace{2em}\times Z_{j,\mathcal V}\left(\frac q\hbar\frac{\lambda_j-\lambda_i}{d},\frac{\lambda_j-\lambda_i}{d}\right)
\end{aligned}
\]
属于 $\mathcal R_T[\hbar^{-1}][[q]]$。
\end{lemma}
\begin{proof}
恒等式
\[
\frac1{\hbar-c_\Gamma}=\left(\frac{c_\Gamma}{\hbar}\right)^d\frac1{\hbar-c_\Gamma}
+\sum_{k=0}^{d-1}c_\Gamma^k\hbar^{-(k+1)}
\]
表明，式~\eqref{eq:complete-intersection-j-localization} 可改写为
\[
Z_{i,\mathcal V}=1+\sum_{d=1}^{\infty}\sum_\Gamma^{(i)}\left(\frac q\hbar\right)^d
\int_{(\overline M_\Gamma)_T}\frac{c_\Gamma^d i_\Gamma^*(\operatorname{Euler}_T(\mathcal V'_{d,2,1}))\prod_{k\ne i}(\lambda_i-\lambda_k)}{a_\Gamma(\hbar-c_\Gamma)\operatorname{Euler}_T(N_\Gamma)}
+\text{关于 $\hbar^{-1}$ 为多项式的项}.
\]
此外，引理~\ref{lem:complete-intersection-j-regularity} 的证明指出，B 型图有 $c_\Gamma^d=0$。故上式成为
\[
Z_{i,\mathcal V}=1+\sum_{d=1}^{\infty}\sum_\Gamma^A\left(\frac q\hbar\right)^d
\int_{(\overline M_\Gamma)_T}\frac{c_\Gamma^d i_\Gamma^*(\operatorname{Euler}_T(\mathcal V'_{d,2,1}))\prod_{k\ne i}(\lambda_i-\lambda_k)}{a_\Gamma(\hbar-c_\Gamma)\operatorname{Euler}_T(N_\Gamma)}
+\text{关于 $\hbar^{-1}$ 为多项式的项},
\]
\SourcePage{398}
其中如式~\eqref{eq:projective-space-j-type-a-localization}，$\sum_\Gamma^A$ 仅对满足 $e_1(\overline M_\Gamma)=q_i$ 的 A 型图 $\Gamma$ 求和。模去关于 $\hbar^{-1}$ 的多项式后，上式与式~\eqref{eq:projective-space-j-type-a-localization} 相同，只是分子多了因子 $i_\Gamma^*(\operatorname{Euler}_T(\mathcal V'_{d,2,1}))$。因此，按引理~\ref{lem:projective-space-j-recursion} 的论证，得到类似于式~\eqref{eq:projective-space-j-recursion-factorization} 的表达式，分子中额外出现若干因子。通过计算 $T$ 在丛 $\mathcal V'_{d,2,1}$ 上作用的权，可将这些因子与 $C_{i,j}^{\mathcal V}(d')$ 的分子等同；这里省略这一论证的细节。于是，$R_i$ 模去关于 $\hbar^{-1}$ 的多项式后为零，引理得证。
\end{proof}
将此引理与引理~\ref{lem:projective-space-j-recursion} 给出的 $\mathcal V=0$ 情形比较，可见该情形下 $R_i=1$。所以，对于零秩平凡丛，得到真正的递推：$Z_i$ 中 $q^d$ 的系数由各 $Z_j$ 中 $q^{d'}$、$d'<d$ 的系数决定。但当 $\mathcal V$ 任意时，只知道 $R_i$ 是 $\hbar^{-1}$ 的多项式，因而无法用 $\mathcal V=0$ 时的同样方式得到递推。为完成定理~\ref{thm:givental-projective-mirror} 的证明，还需要新内容。基本思想是，用下面引理~\ref{lem:complete-intersection-double-construction} 的“双重构造”补充引理~\ref{lem:complete-intersection-j-regularity} 和~\ref{lem:complete-intersection-j-recursion} 的正则性及递推条件，使这些性质足以保证唯一性。引理~\ref{lem:complete-intersection-j-uniqueness} 最终将实现这一点。
\begin{lemma}[\IndexLocation{indexloc:0850:398}双重构造引理]\label{lem:complete-intersection-double-construction}
设 $z$ 为变量，$\mathcal V=\bigoplus_{j=1}^{\ell}\mathcal O_{\mathbb P^n}(a_j)$。则表达式
\[
W(z,\hbar)=\sum_{i=0}^n\frac{E_i}{\prod_{j\ne i}(\lambda_i-\lambda_j)}e^{\lambda_i z}
Z_{i,\mathcal V}(qe^{z\hbar},\hbar)Z_{i,\mathcal V}(q,-\hbar),\qquad E_i=\prod_{j=1}^{\ell}(a_j\lambda_i),
\]
属于 $H^*(BT)[\hbar][[q,z]]$。
\end{lemma}
\begin{proof}
回顾\hyperref[subsec:stable-moduli-linear-models]{第~\ref*{subsec:stable-moduli-linear-models} 小节}讨论的映射 $\varphi:M_d\to N_d$ 及丛 $\mathcal V_d$，其中 $M_d=\overline M_{0,0}(\mathbb P^1\times\mathbb P^n,(1,d))$。这里采用此构造的一个变体：考虑 $\overline M_{0,2}(\mathbb P^1\times\mathbb P^n,(1,d))$，它具有\IndexLocation{indexloc:0276:398}赋值映射 $e_1,e_2$，取值于 $\mathbb P^1\times\mathbb P^n$。遵循 \SourceCite{Givental2}，定义
\[
L_d=e_1^{-1}(\{0\}\times\mathbb P^n)\cap e_2^{-1}(\{\infty\}\times\mathbb P^n).
\]
令 $i_{L_d}$ 为 $L_d$ 到 $\overline M_{0,2}(\mathbb P^1\times\mathbb P^n,(1,d))$ 的包含映射。映射 $\mu:L_d\to N_d$ 定义为 $\mu=\varphi\circ\pi_1\circ\pi_2\circ i_{L_d}$，其中 $\pi_k$ 是忘掉第 $k$ 个标记点的映射，$k=1,2$。\SourceCite{Givental2} 讨论了此映射，并将 $N_d$ 记为 $L'_d$。

还回顾由自然投影 $\mathbb P^1\times\mathbb P^n\to\mathbb P^n$ 诱导的映射 $\pi:M_d\to\overline M_{0,0}(\mathbb P^n,d)$。由它得到映射 $\widetilde\pi:L_d\to\overline M_{0,0}(\mathbb P^n,d)$，定义为 $\widetilde\pi=\pi\circ\pi_1\circ\pi_2\circ i_{L_d}$。对于 $d=0$，直接以 $\mathbb P^n$ 上的丛 $\mathcal V$ 定义相应零次项。令 $\widetilde\chi_d^{\mathcal V}=\operatorname{Euler}_G(\widetilde\pi^*\mathcal V_d)$。断言以下恒等式：
\begin{equation}\label{eq:complete-intersection-double-polynomiality-integral}
\sum_{d=0}^{\infty}q^d\int_{(L_d)_G}e^{\mu^*(\kappa)z}\widetilde\chi_d^{\mathcal V}
=\sum_{i=0}^n\frac{E_i}{\prod_{j\ne i}(\lambda_i-\lambda_j)}e^{\lambda_i z}
Z_{i,\mathcal V}(qe^{z\hbar},\hbar)Z_{i,\mathcal V}(q,-\hbar).
\end{equation}
一旦证明此式，就得到所需结论，因为 $\mu^*(\kappa)$ 与相应丛的 $G$ 等变 Euler 类属于等变上同调，因而关于 $\hbar$ 和 $\lambda_i$ 多项式地依赖。

将用推论~\ref{cor:orbifold-localization-integration} 的局部化公式证明此式。论证与定理~\ref{thm:geometric-euler-data-gluing} 的证明非常相似，因此仅概述计算，省略大部分细节。在 $\sum_ja_j=n+1$ 的情形，此式的完整证明见 \SourceCite[引理 7]{Pandharipande3}。
\SourcePage{399}
先考察式~\eqref{eq:complete-intersection-double-polynomiality-integral} 左边 $q^d$ 的系数
\begin{equation}\label{eq:complete-intersection-double-degree-integral}
\int_{(L_d)_G}e^{\mu^*(\kappa)z}\widetilde\chi_d^{\mathcal V}.
\end{equation}
第一步是确定 $L_d$ 的固定点分支。这些分支与定理~\ref{thm:geometric-euler-data-gluing} 证明中出现的固定点分支十分相似。

先固定 $0\le i\le n$，令 $q_i$ 为 $\mathbb P^n$ 中相应的 $T$ 固定点。对每个 $0<r\le d$，令 $F'_r$ 为 $\overline M_{0,2}(\mathbb P^n,r)^T$ 的一个分支，包含第一个标记点映到 $q_i$ 的稳定映射。若还有 $r<d$，选择 $\overline M_{0,2}(\mathbb P^n,d-r)^T$ 的一个分支 $F'_{d-r}$，其稳定映射的第一个标记点同样映到 $q_i$。将描述 $L_d^G$ 的一个同构于 $F'_r\times F'_{d-r}$ 的分支。略微滥用记号，仍用 $F'_r\times F'_{d-r}$ 表示它。一般映射 $f\in F'_r\times F'_{d-r}$ 的源曲线为 $C=C_0\cup C_1\cup C_2$，其中 $C_0=\mathbb P^1$，而 $f:C\to\mathbb P^1\times\mathbb P^n$ 给定为
\[
\begin{aligned}
f|_{C_0}(z)&=(z,q_i),\\
f|_{C_1}(z)&=(0,f_1(z)),\\
f|_{C_2}(z)&=(\infty,f_2(z)).
\end{aligned}
\]
曲线 $C_1$ 有两个标记点 $(p'_1,p_1)$，$f_1:(C_1,p'_1,p_1)\to\mathbb P^n$ 是 $F'_r$ 的元素。另外，$C_2$ 有两个标记点 $(p'_2,p_2)$，映射 $f_2:(C_2,p'_2,p_2)\to\mathbb P^n$ 属于 $F'_{d-r}$。最后，把 $0\in C_0$ 与 $p'_1\in C_1$ 粘合，把 $\infty\in C_0$ 与 $p'_2\in C_2$ 粘合。于是 $f:(C,p_1,p_2)\to\mathbb P^1\times\mathbb P^n$ 属于 $L_d$，所有这类映射构成的集合正是所要描述的分支 $F'_r\times F'_{d-r}$。$r=0$ 或 $r=d$ 时分支的描述更简单，留给读者。

接下来用局部化描述此分支对积分~\eqref{eq:complete-intersection-double-degree-integral} 的贡献。先限制被积式中的各项到 $F'_r\times F'_{d-r}$。注意，$\mu(F'_r\times F'_{d-r})$ 是 $N_d$ 的 $G$ 固定点 $p_{i,r}$，如定理~\ref{thm:geometric-euler-data-gluing} 的讨论。因此，由\hyperref[subsec:stable-moduli-linear-models]{第~\ref*{subsec:stable-moduli-linear-models} 小节}，$\mu^*(\kappa)$ 在 $F'_r\times F'_{d-r}$ 上限制为 $\lambda_i+r\hbar$。

再计算 $\widetilde\chi_d^{\mathcal V}$ 的限制。令 $\rho_1,\rho_2$ 为 $F'_r\times F'_{d-r}$ 到两个因子的投影，则有正合列
\[
\begin{aligned}
0\longrightarrow{}&\rho_1^*(\mathcal V'_{r,2,1}|_{F'_r})\oplus\rho_2^*(\mathcal V'_{d-r,2,1}|_{F'_{d-r}})\\
\longrightarrow{}&(\widetilde\pi^*\mathcal V_d)|_{F'_r\times F'_{d-r}}
\longrightarrow\mathcal V_{q_i}\otimes\mathcal O_{F'_r\times F'_{d-r}}\longrightarrow0.
\end{aligned}
\]
考察纤维就能说明此正合列的来源。在 $f\in F'_r\times F'_{d-r}$ 上，最后一个映射为限制映射
\[
H^0(C,f^*\mathcal V)\longrightarrow H^0(C_0,f|_{C_0}^*\mathcal V)=\mathcal V_{q_i}.
\]
其核由 $f_1^*\mathcal V$ 与 $f_2^*\mathcal V$ 的一对截面构成，二者分别在标记点 $p'_1$、$p'_2$ 消失。这些正是 $\rho_1^*(\mathcal V'_{r,2,1}|_{F'_r})$ 和 $\rho_2^*(\mathcal V'_{d-r,2,1}|_{F'_{d-r}})$ 的纤维。因此，正合列表明 $\widetilde\chi_d^{\mathcal V}$ 限制为
\[
E_i\operatorname{Euler}_G(\mathcal V'_{r,2,1}|_{F'_r})\operatorname{Euler}_G(\mathcal V'_{d-r,2,1}|_{F'_{d-r}}),
\]
因为 $\mathcal V_{q_i}=(\bigoplus_j\mathcal O_{\mathbb P^n}(a_j))_{q_i}$ 的各个权为 $a_j\lambda_i$，其 Euler 类为 $E_i=\prod_j(a_j\lambda_i)$。
\SourcePage{400}
最后，计算法丛的等变 Euler 类，结果为
\[
\begin{aligned}
\operatorname{Euler}_G(N(F'_r\times F'_{d-r}))={}&\prod_{j\ne i}(\lambda_i-\lambda_j)\\
&\times\frac{(\hbar-c_1^T(\mathcal L_{r,1}))\operatorname{Euler}_G(N(F'_r))}{\prod_{j\ne i}(\lambda_i-\lambda_j)}\\
&\times\frac{(-\hbar-c_1^T(\mathcal L_{d-r,1}))\operatorname{Euler}_G(N(F'_{d-r}))}{\prod_{j\ne i}(\lambda_i-\lambda_j)}.
\end{aligned}
\]
（暂不作显然的约分，其理由很快就会清楚。）此计算与定理~\ref{thm:geometric-euler-data-gluing} 中 $\operatorname{Euler}_G(N(F_r\times F_{d-r}))$ 的计算十分相似。唯一的区别是，后者有一个来自 $\mathbb P^1$ 上使 $0$ 与 $\infty$ 移动的自同构形变的因子 $-\hbar^2$。$\operatorname{Euler}_G(N(F'_r\times F'_{d-r}))$ 中没有此因子，因为按 $L_d$ 的定义，其中所有映射与投影到 $\mathbb P^1$ 因子的映射复合后，两个标记点的像始终为 $0$ 和 $\infty$。

汇总以上各项，可见 $F'_r\times F'_{d-r}$ 对积分~\eqref{eq:complete-intersection-double-degree-integral} 的贡献为
\[
\begin{aligned}
\frac{E_i}{\prod_{j\ne i}(\lambda_i-\lambda_j)}e^{(\lambda_i+r\hbar)z}
&\int_{(F'_r)_G}\frac{i_{F'_r}^*(\operatorname{Euler}_G(\mathcal V'_{r,2,1}))\prod_{j\ne i}(\lambda_i-\lambda_j)}{a_r(\hbar-c_{F'_r})\operatorname{Euler}_G(N(F'_r))}\\
\times{}&\int_{(F'_{d-r})_G}\frac{i_{F'_{d-r}}^*(\operatorname{Euler}_G(\mathcal V'_{d-r,2,1}))\prod_{j\ne i}(\lambda_i-\lambda_j)}{a_{d-r}(-\hbar-c_{F'_{d-r}})\operatorname{Euler}_G(N(F'_{d-r}))}.
\end{aligned}
\]
其中 $a_r,a_{d-r}$ 是相应有限自同构群的阶；$c_{F'_r},c_{F'_{d-r}}$ 分别为 $c_1^T(\mathcal L_{r,1})$ 和 $c_1^T(\mathcal L_{d-r,1})$ 在 $F'_r$、$F'_{d-r}$ 上的限制。

由推论~\ref{cor:orbifold-localization-integration}，积分~\eqref{eq:complete-intersection-double-degree-integral} 是上述贡献对所有 $F'_r$、$F'_{d-r}$ 求和，再对 $i=0,\ldots,n$ 求和的结果。随后乘以 $q^d$，再对所有 $d\ge0$ 求和，就得到式~\eqref{eq:complete-intersection-double-polynomiality-integral} 的左边。

为把它与该式右边联系起来，回顾式~\eqref{eq:complete-intersection-j-localization} 中 $Z_{i,\mathcal V}$ 的公式：
\[
Z_{i,\mathcal V}(q,\hbar)=1+\sum_{d=1}^{\infty}q^d\sum_\Gamma^{(i)}
\int_{(\overline M_\Gamma)_T}\frac{i_\Gamma^*(\operatorname{Euler}_T(\mathcal V'_{d,2,1}))\prod_{j\ne i}(\lambda_i-\lambda_j)}{a_\Gamma(\hbar-c_\Gamma)\operatorname{Euler}_T(N_\Gamma)}.
\]
如式~\eqref{eq:projective-space-j-localization} 所解释，$\sum_\Gamma^{(i)}$ 对 $\overline M_{0,2}(\mathbb P^n,d)$ 固定点集内、经 $e_1$ 映到 $q_i$ 的全部分支求和。将 $d$ 改记为 $r$，则上式可写成
\[
Z_{i,\mathcal V}(q,\hbar)=1+\sum_{r=1}^{\infty}q^r\sum_{F'_r}
\int_{(F'_r)_T}\frac{i_{F'_r}^*(\operatorname{Euler}_T(\mathcal V'_{r,2,1}))\prod_{j\ne i}(\lambda_i-\lambda_j)}{a_r(\hbar-c_{F'_r})\operatorname{Euler}_T(N(F'_r))}.
\]
此外，因 $\mathbb C^*$ 在 $\mathbb P^n$ 上平凡作用，可以将 $T$ 替换为 $G=\mathbb C^*\times T$。若用此 $Z_{i,\mathcal V}$ 公式展开式~\eqref{eq:complete-intersection-double-polynomiality-integral} 的右边，就容易看到结果恰是上一段所述的左边表达式。于是该式成立，引理也按前述论证得到证明。
\end{proof}
\SourcePage{401}
刚给出的 $Z_{i,\mathcal V}$ 公式表明，它自然可视为 $\hbar^{-1}$ 的形式幂级数。因此，引理~\ref{lem:complete-intersection-double-construction} 定义的 $W(z,\hbar)$ 竟是关于 $\hbar$ 的多项式，这颇为意外。将看到，这一点在唯一性引理（引理~\ref{lem:complete-intersection-j-uniqueness}）的证明中发挥重要作用。

接着将前三个引理所表述的 $Z_{i,\mathcal V}$ 的性质统一起来。考虑集合 $\{Z_i\}_{i=0}^n$，其中所有 $i$ 均有 $Z_i\in\mathbb Q(\lambda_0,\ldots,\lambda_n)[[\hbar^{-1},q]]$。遵循 \SourceCite{Pandharipande3}，作如下定义。
\begin{definition}\label{def:complete-intersection-class-p}\index{Givental 方法!$\mathcal P$ 类|hyperpage}
若集合 $\{Z_i\}_{i=0}^n$ 满足以下条件，就称它属于 $\mathcal P$ 类：
\begin{enumerate}[label=(\roman*)]
\item 所有 $i$ 均有 $Z_i(0,\hbar)=1$。
\item $Z_i$ 满足引理~\ref{lem:complete-intersection-j-regularity} 中 $Z_{i,\mathcal V}$ 所满足的正则性条件。
\item $Z_i$ 满足引理~\ref{lem:complete-intersection-j-recursion} 中 $Z_{i,\mathcal V}$ 所满足的递推条件。
\item 在引理~\ref{lem:complete-intersection-double-construction} 中用 $Z_i$ 代替 $Z_{i,\mathcal V}$ 定义 $W(z,\hbar)$，则 $W(z,\hbar)\in H^*(BT)[\hbar][[q,z]]$。
\end{enumerate}
\end{definition}
由此定义提及的三个引理，$\{Z_{i,\mathcal V}\}_{i=0}^n$ 属于 $\mathcal P$ 类。现在可陈述关键的唯一性引理。
\begin{lemma}[唯一性引理]\label{lem:complete-intersection-j-uniqueness}
设 $\{Z_i\}_{i=0}^n$ 与 $\{Z'_i\}_{i=0}^n$ 均属于 $\mathcal P$ 类。若所有 $0\le i\le n$ 的表达式按 $\hbar^{-1}$ 形式展开时，都有 $Z_i\equiv Z'_i\pmod{\hbar^{-2}}$，则所有 $i$ 均有 $Z_i=Z'_i$。
\end{lemma}
\begin{proof}
遵循 \SourceCite{Givental4,Kim2} 的证明。令 $Z_i=\sum_d Z_{i,d}q^d$，类似地写出 $Z'_i$。因两个集合均属于 $\mathcal P$ 类，有 $Z_{i,0}=Z'_{i,0}=1$。因此，可归纳假设：对于某个 $d_0>0$，所有 $0\le d<d_0$、$0\le i\le n$ 均有 $Z_{i,d}=Z'_{i,d}$。属于 $\mathcal P$ 类意味着引理~\ref{lem:complete-intersection-j-recursion} 的递推同时适用于两集合。结合归纳假设，对所有 $i$，$Z'_{i,d_0}-Z_{i,d_0}$ 可展开为 $\hbar^{-1}$ 的多项式。令
\[
D_i(\hbar)=Z'_{i,d_0}-Z_{i,d_0}=\hbar^{-2r_i}\bigl(A_i\hbar^{-1}+B_i+O(\hbar)\bigr),
\]
其中 $A_i,B_i\in\mathbb C(\lambda_0,\ldots,\lambda_n)$，$r_i\ge0$。由 $Z_i\equiv Z'_i\pmod{\hbar^{-2}}$，只要证明所有 $r_i$ 均可取为零，就可得 $Z'_{i,d_0}=Z_{i,d_0}$。

所以须证明：$r_i\ge1$ 时，最高阶的系数 $A_i,B_i$ 为零。为此，用引理~\ref{lem:complete-intersection-double-construction} 的公式，分别由 $\{Z_i\}$ 和 $\{Z'_i\}$ 构造 $W$、$W'$。归纳假设容易推出，$W'-W$ 中 $q^{d_0}$ 的系数为
\[
\sum_{i=0}^n\frac{E_i}{\prod_{j\ne i}(\lambda_i-\lambda_j)}e^{\lambda_i z}
\bigl(e^{d_0\hbar z}D_i(\hbar)+D_i(-\hbar)\bigr).
\]
将 $D_i$ 的上述公式代入并简化，可见它具有形式
\[
\sum_{i=0}^n\frac{E_i}{\prod_{j\ne i}(\lambda_i-\lambda_j)}e^{\lambda_i z}
\hbar^{-2r_i}\bigl(A_id_0z+2B_i+O(\hbar)\bigr).
\]
因为两个集合均属于 $\mathcal P$ 类，$W$、$W'$ 关于 $\hbar$ 都是多项式。取最大的 $r_i=r\ge1$，上述表达式中 $\hbar^{-2r}$ 的系数必须为零。对于互异的 $\lambda_i$，，而 $E_i/\prod_{j\ne i}(\lambda_i-\lambda_j)$ 在局部化后非零。因此，所有满足 $r_i=r$ 的 $A_i,B_i$ 均为零。逐次降低最高负幂，得到所有 $r_i\ge1$ 的系数消失，引理得证。
\end{proof}
证明\IndexLocation{indexloc:0847:401}\IndexLocation{indexloc:0850:401}\IndexLocation{indexloc:0860:401}\IndexLocation{indexloc:0861:401}\IndexLocation{indexloc:0863:401}镜像定理之前，还需两个准备步骤。首先，须给出现的变量赋予权重：赋予
\SourcePage{402}
$\lambda_i,p,\hbar$ 各权重 $1$，赋予 $q$ 权重 $n+1-\sum_i a_i$。下面“齐次”一词均相对于这些权重而言。以下引理表明，某些坐标变换不改变属于 $\mathcal P$ 类的性质。
\begin{lemma}\label{lem:complete-intersection-p-transformations}
考虑形式幂级数 $f(q)\in q\mathbb Q[[q]]$、$g(q),h(q)\in qH^*(BT)[[q]]$，均为齐次，且 $\deg f=\deg g=0$、$\deg h=1$。若 $\{Z_i\}_{i=0}^n$ 属于 $\mathcal P$ 类，则 $\{Z''_i\}_{i=0}^n$ 亦然，其中
\[
Z''_i(q,\hbar)=e^{f(q)+h(q)/\hbar+g(q)\lambda_i/\hbar}Z_i(qe^{g(q)},\hbar).
\]
\end{lemma}
\begin{proof}
当 $\sum_ja_j=n+1$ 时，详尽证明见 \SourceCite[引理 10]{Pandharipande3}。该论证容易推广到 $\sum_ja_j\le n+1$ 的情形，因此省略细节。
\end{proof}
终于可以概述 Givental 关于 $\mathbb P^n$ 中 nef 完全交的镜像定理的证明。
\begin{proof}[定理~\ref{thm:givental-projective-mirror} 的证明]
须证明 $I_{\mathcal V}$ 与 $J_{\mathcal V}$ 在适当的齐次变量变换后相等。与 $\mathcal V=0$ 的情形一样，策略是构造 $I_{\mathcal V}$ 与 $J_{\mathcal V}$ 的等变版本，随后证明定理的等变版本。

$J_{\mathcal V}$ 的等变版本是式~\eqref{eq:complete-intersection-equivariant-j} 中的 $J_T$。去掉指数因子得到式~\eqref{eq:complete-intersection-equivariant-s} 中的 $S_{\mathcal V}$，再去掉 $\operatorname{Euler}_T(\mathcal V)$ 得到 $\widetilde S_{\mathcal V}$。最后，由公式
\[
Z_{i,\mathcal V}=\int_{(\mathbb P^n)_T}\widetilde S_{\mathcal V}\cup\phi_i=i_{q_i}^*(\widetilde S_{\mathcal V})
\]
得到 $Z_{i,\mathcal V}$。知道 $J_T$ 由这些函数唯一确定，且由引理~\ref{lem:complete-intersection-j-regularity}、\ref{lem:complete-intersection-j-recursion} 和~\ref{lem:complete-intersection-double-construction}，集合 $\{Z_{i,\mathcal V}\}_{i=0}^n$ 属于 $\mathcal P$ 类。还注意，上述权重下 $\widetilde S_{\mathcal V}$ 是权重为零的齐次函数。这可由式~\eqref{eq:complete-intersection-equivariant-s} 及三个简单事实看出：丛 $\mathcal V'_{d,2,1}$ 的秩为 $\sum_{j=1}^{\ell}a_jd$；$e_1:\overline M_{0,2}(\mathbb P^n,d)\to\mathbb P^n$ 的纤维维数为 $(n+1)d-1$；$1/(\hbar-c_T)$ 展开中的一般项为 $c_T^k/\hbar^{k+1}$。细节留给读者。

转而考虑 $I_{\mathcal V}$ 的等变版本，定义
\begin{equation}\label{eq:complete-intersection-equivariant-i-series}
\widetilde S'_{\mathcal V}(q,\hbar)=1+\sum_{d=1}^{\infty}q^d
\frac{\prod_{j=1}^{\ell}\prod_{m=1}^{a_jd}(a_jp+m\hbar)}{\prod_{j=0}^n\prod_{m=1}^d(p-\lambda_j+m\hbar)},
\end{equation}
再令
\[
I_T=e^{(t_0+t_1p)/\hbar}\operatorname{Euler}_T(\mathcal V)\widetilde S'_{\mathcal V}(q,\hbar).
\]
当 $\lambda_i\to0$ 时，$I_T$ 的非等变极限显然为 $I_{\mathcal V}$。$\widetilde S'_{\mathcal V}$ 也是权重为零的齐次函数。最后定义
\[
Z'_{i,\mathcal V}=\int_{(\mathbb P^n)_T}\widetilde S'_{\mathcal V}\cup\phi_i=i_{q_i}^*(\widetilde S'_{\mathcal V}).
\]
可验证 $\{Z'_{i,\mathcal V}\}_{i=0}^n$ 属于 $\mathcal P$ 类。这一验证直接但略长，省略细节。

为比较 $I_T$ 与 $J_T$，先注意
\[
J_T=e^{(t_0+t_1p)/\hbar}\operatorname{Euler}_T(\mathcal V)\bigl(1+o(\hbar^{-1})\bigr).
\]
\SourcePage{403}
这与证明式~\eqref{eq:givental-bundle-j-normalization} 时的论证相同，该式给出渐近展开 $J_{\mathcal V}=e^{(t_0+t_1H)/\hbar}\operatorname{Euler}(\mathcal V)(1+o(\hbar^{-1}))$。

现在考虑变量变换 $t_0\mapsto t_0+f(q)\hbar+h(q)$、$t_1\mapsto t_1+g(q)$，令 $J''_T$ 为在 $J_T$ 中作此代换后的表达式。上述渐近展开给出
\[
J''_T=e^{(t_0+t_1p)/\hbar}\operatorname{Euler}_T(\mathcal V)
\bigl(e^{f(q)}+e^{f(q)}h(q)\hbar^{-1}+e^{f(q)}g(q)p\hbar^{-1}+o(\hbar^{-1})\bigr).
\]
以同样方式展开 $I_T$，可见存在唯一一组 $f,g,h$，使 $I_T\equiv J''_T\pmod{\hbar^{-2}}$；这类似于式~\eqref{eq:givental-mirror-variable-coefficients} 中所作的计算。注意，$f,g$ 的次数为零，$h$ 的次数为一，因为 $\widetilde S_{\mathcal V}$ 与 $\widetilde S'_{\mathcal V}$ 的次数均为零。

为证明定理，须证明 $I_T=J''_T$。首先注意，$t_1\mapsto t_1+g(q)$ 将 $q=e^{t_1}$ 变为 $qe^{g(q)}$。因此，在
\[
J_T=e^{(t_0+t_1p)/\hbar}\operatorname{Euler}_T(\mathcal V)\widetilde S_{\mathcal V}(q,\hbar)
\]
中作变量变换，得到
\[
\begin{aligned}
J''_T&=e^{(t_0+f(q)\hbar+h(q)+(t_1+g(q))p)/\hbar}\operatorname{Euler}_T(\mathcal V)\widetilde S_{\mathcal V}(qe^{g(q)},\hbar)\\
&=e^{(t_0+t_1p)/\hbar}\operatorname{Euler}_T(\mathcal V)e^{f(q)+h(q)/\hbar+g(q)p/\hbar}\widetilde S_{\mathcal V}(qe^{g(q)},\hbar)\\
&=e^{(t_0+t_1p)/\hbar}\operatorname{Euler}_T(\mathcal V)\widetilde S''_{\mathcal V}(q,\hbar),
\end{aligned}
\]
其中
\[
\widetilde S''_{\mathcal V}(q,\hbar)=e^{f(q)+h(q)/\hbar+g(q)p/\hbar}\widetilde S_{\mathcal V}(qe^{g(q)},\hbar).
\]
若令
\[
Z''_{i,\mathcal V}=\int_{(\mathbb P^n)_T}\widetilde S''_{\mathcal V}\cup\phi_i=i_{q_i}^*(\widetilde S''_{\mathcal V}),
\]
则 $I_T=J''_T$ 等价于所有 $i$ 均有 $Z'_{i,\mathcal V}=Z''_{i,\mathcal V}$。

再断言集合 $\{Z''_{i,\mathcal V}\}_{i=0}^n$ 属于 $\mathcal P$ 类。为此，注意上述 $\widetilde S''_{\mathcal V}$、$Z''_{i,\mathcal V}$ 的公式给出
\[
Z''_{i,\mathcal V}(q,\hbar)=e^{f(q)+h(q)/\hbar+g(q)\lambda_i/\hbar}Z_{i,\mathcal V}(qe^{g(q)},\hbar),
\]
因为 $i_{q_i}^*(p)=\lambda_i$。由引理~\ref{lem:complete-intersection-p-transformations}，$\{Z''_{i,\mathcal V}\}_{i=0}^n$ 属于 $\mathcal P$ 类，因为 $\{Z_{i,\mathcal V}\}_{i=0}^n$ 属于此类。

现在定理容易得证。已知 $\{Z'_{i,\mathcal V}\}$ 与 $\{Z''_{i,\mathcal V}\}$ 均属于 $\mathcal P$ 类，而 $I_T\equiv J''_T\pmod{\hbar^{-2}}$ 意味着所有 $i$ 均有 $Z'_{i,\mathcal V}\equiv Z''_{i,\mathcal V}\pmod{\hbar^{-2}}$。因此，唯一性引理（引理~\ref{lem:complete-intersection-j-uniqueness}）的假设成立，可得所有 $i$ 均有 $Z'_{i,\mathcal V}=Z''_{i,\mathcal V}$。如前所述，这意味着 $I_T=J''_T$。取非等变极限 $\lambda_i\to0$，表明所需的变量变换后 $I_{\mathcal V}$ 等于 $J_{\mathcal V}$。镜像定理的证明完成。
\end{proof}
对本节与\hyperref[sec:quintic-mirror-theorem]{第~\ref*{sec:quintic-mirror-theorem} 节}所讨论的镜像定理证明方法之间的关系，再作几点评论。引理~\ref{lem:complete-intersection-j-recursion} 的递推关系，与\hyperref[subsec:lly-euler-data]{第~\ref*{subsec:lly-euler-data} 小节}中引理~\ref{lem:euler-data-gluing} 的粘合恒等式起着类似作用。把递推关系与引理~\ref{lem:complete-intersection-double-construction} 的双重构造的多项式性质结合，对应于该小节讨论的 Euler 数据性质。此外，引理~\ref{lem:complete-intersection-j-uniqueness} 的唯一性对应于\hyperref[subsec:linked-euler-data]{第~\ref*{subsec:linked-euler-data} 小节}的定理~\ref{thm:linked-euler-data-uniqueness}。注意，该定理次数界 $5d-2$ 中的 $-2$，
\SourcePage{404}
对应于这里唯一性引理证明中两个系数（$1$、$\hbar^{-1}$ 的系数）的匹配。最后，两种方法都使用镜像变换，将两个不同的形式级数等同。\SourceCite{Givental5} 的一个长篇脚注更详细地比较了这两种方法，该文也给出了亏格一的镜像定理。

\subsection{环面完全交}\label{subsec:givental-toric-mirror}\index{镜像定理!环面完全交|hyperpage}
接下来讨论镜像定理到光滑环面簇中完全交的推广 \SourceCite{Givental4}。主要结果将在下面作为定理~\ref{thm:givental-toric-mirror} 重述。随后讨论环面镜像定理与本书一个主要主题的关系，即 Batyrev 镜像构造所给出的 Calabi--Yau 流形的镜像对称。最后给出几个有趣例子。

设 $X_\Sigma$ 为扇 $\Sigma$ 对应的光滑环面簇，$\mathcal L_1,\ldots,\mathcal L_\ell$ 是其上由全局截面生成的线丛。令 $\mathcal V=\bigoplus_{i=1}^{\ell}\mathcal L_i$，$X\subset X_\Sigma$ 为 $\mathcal V$ 的一般全局截面的零点集所定义的光滑完全交。在此情形，要定义上同调值函数 $I_{\mathcal V}$ 和 $J_{\mathcal V}$，推广\hyperref[subsec:givental-nef-projective-mirror]{第~\ref*{subsec:givental-nef-projective-mirror} 小节}的构造。

先约定一些记号。对每个 $\rho\in\Sigma(1)$，略微滥用记号，也用 $D_\rho$ 表示相应除子在 $H^2(X_\Sigma)$ 中的上同调类。遵循 \SourceCite{Givental4}，令 $\mathcal L_i(\beta)=\int_\beta c_1(\mathcal L_i)$、$D_\rho(\beta)=\int_\beta D_\rho$。还选取 $H^2(X_\Sigma,\mathbb Z)$ 的整基 $T_1,\ldots,T_r$，使各基元素位于 Kähler 锥的闭包内。如常令 $\delta=\sum_{i=1}^r t_iT_i$。

定义两个上同调值形式函数。先定义 $I_{\mathcal V}$：
\begin{equation}\label{eq:toric-complete-intersection-i-function}
\begin{aligned}
I_{\mathcal V}={}&e^{(t_0+\delta)/\hbar}\operatorname{Euler}(\mathcal V)\\
&\times\sum_{\beta\in M(X_\Sigma)}q^\beta
\frac{\prod_{i=1}^{\ell}\prod_{m=-\infty}^{\mathcal L_i(\beta)}(c_1(\mathcal L_i)+m\hbar)\prod_\rho\prod_{m=-\infty}^0(D_\rho+m\hbar)}{\prod_{i=1}^{\ell}\prod_{m=-\infty}^0(c_1(\mathcal L_i)+m\hbar)\prod_\rho\prod_{m=-\infty}^{D_\rho(\beta)}(D_\rho+m\hbar)}.
\end{aligned}
\end{equation}
其中 $q_i=e^{t_i}$、$q^\beta=\prod_{i=1}^r q_i^{\int_\beta T_i}$。注意，若 $\Sigma$ 是 $\mathbb P^n$ 的标准扇，就恢复式~\eqref{eq:givental-complete-intersection-i-function}。再定义 $J_{\mathcal V}$：
\begin{equation}\label{eq:toric-complete-intersection-j-function}
\begin{aligned}
J_{\mathcal V}={}&e^{(t_0+\delta)/\hbar}\operatorname{Euler}(\mathcal V)\\
&\times\left(1+\sum_{\beta\ne0}q^\beta PD^{-1}e_{1*}
\left(\frac{\operatorname{Euler}(\mathcal V'_{\beta,2,1})}{\hbar-c}\cap[\overline M_{0,2}(X_\Sigma,\beta)]^{\mathrm{virt}}\right)\right),
\end{aligned}
\end{equation}
其中 $[\overline M_{0,2}(X_\Sigma,\beta)]^{\mathrm{virt}}$ 是 $\overline M_{0,2}(X_\Sigma,\beta)$ 的虚拟基本类，$PD$ 是 Poincaré 对偶。注意，当 $X_\Sigma$ 为凸簇时，虚拟基本类就是通常的基本类，所以 $J_{\mathcal V}$ 的公式可简化。例如，当 $X_\Sigma$ 为凸簇 $\mathbb P^n$ 时，式~\eqref{eq:toric-complete-intersection-j-function} 化为式~\eqref{eq:givental-general-j-function}。

此时，变量 $q_i$ 有次数。如\hyperref[subsec:quantum-hyperplane-section-principle]{第~\ref*{subsec:quantum-hyperplane-section-principle} 小节}，用等式
\[
c_1(X_\Sigma)-c_1(\mathcal V)=\sum_{i=1}^r(\deg q_i)T_i
\]
定义 $\deg q_i$。假设 $X\subset X_\Sigma$ 是\hyperref[subsec:nef-gkz]{第~\ref*{subsec:nef-gkz} 小节}意义下的 nef 完全交，即 $-(K_{X_\Sigma}+\sum_{i=1}^{\ell}\mathcal L_i)$ 在 $X_\Sigma$ 上为 nef。此时假设基 $T_1,\ldots,T_r$ 已选取为使
\SourcePage{405}
$-(K_{X_\Sigma}+\sum_{i=1}^{\ell}\mathcal L_i)$ 位于 $T_i$ 所生成的锥中。在 nef 情形总可作这种选择。于是所有 $i$ 均有 $\deg q_i\ge0$。

现在可陈述 Givental 版本的环面镜像定理。
\begin{theorem}\label{thm:givental-toric-mirror}
设 $X\subset X_\Sigma$ 是 nef 完全交，$I_{\mathcal V}$、$J_{\mathcal V}$ 分别由式~\eqref{eq:toric-complete-intersection-i-function} 和~\eqref{eq:toric-complete-intersection-j-function} 定义。则二者在加权齐次三角变量变换后相等：
\[
t_0\longmapsto t_0+f_0(q)\hbar+h(q),\qquad t_i\longmapsto t_i+f_i(q),\quad1\le i\le r,
\]
其中 $f_0,\ldots,f_r,h$ 为加权齐次幂级数，且 $\deg f_0=\deg f_i=0$、$\deg h=1$。
\end{theorem}
\begin{proof}
如可预料，证明的策略是构造 $I_{\mathcal V}$ 和 $J_{\mathcal V}$ 的等变版本。这正是环面簇的重要性所在，因为 $T\subset X_\Sigma$ 提供了所需的群作用。不过，式~\eqref{eq:toric-complete-intersection-j-function} 表明，还需要虚拟基本类的等变版本。局部化在证明中起着如此核心的作用，因此等变虚拟基本类的局部化公式不可缺少。这便联系到\hyperref[subsec:virtual-localization]{第~\ref*{subsec:virtual-localization} 小节}讨论的一些问题。

这里不证明此定理。\SourceCite{Givental4} 的证明在写作时尚不完整。特别地，如上所述，证明以存在适当的等变虚拟基本类局部化公式为前提。\SourceCite{GPa} 已建立了这样的公式。考虑到此定理的重要性，我们希望其余证明细节能被仔细写出。我们对 Givental 建立的概念框架完全有信心，相信这些细节将会出现，就如 \SourceCite{BDPP,Pandharipande3} 对命题~\ref{prop:projective-space-i-equals-j} 和定理~\ref{thm:givental-projective-mirror} 所做的一样。
\end{proof}
还可问：式~\eqref{eq:toric-complete-intersection-j-function} 定义的形式函数 $J_{\mathcal V}$，与\hyperref[subsec:giventals-j-function]{第~\ref*{subsec:giventals-j-function} 小节}的 Givental 函数 $J_X$ 有什么关系？这里 $i:X\hookrightarrow X_\Sigma$ 是 $\mathcal V$ 的一般截面的零点集的包含映射。基本思想是，$J_{\mathcal V}$ 应计算来自环境空间的上同调类所对应的亏格零\IndexLocation{indexloc:0459:405}引力关联数。因此，预期 $J_{\mathcal V}$ 包含 $J_X$ 的部分信息。有一个情形下二者关系很简单：当 $\mathcal V=0$ 时，$X=X_\Sigma$，不难证明 $J_X=J_{\mathcal V}$。当 $\mathcal V$ 非零时，情况更复杂，因为 $J_{\mathcal V}$ 取值于 $H^*(X_\Sigma)$，而 $J_X$ 取值于 $H^*(X)$。对于 $X_\Sigma=\mathbb P^n$，命题~\ref{prop:complete-intersection-j-pushforward} 已证明 $J_{\mathcal V}=i_!(J_X)$。更一般地，当 $X_\Sigma$ 为凸簇且 $\mathcal V$ 丰沛时，也成立。不过，尚不知道在此定理的情形下是否仍有 $J_{\mathcal V}=i_!(J_X)$。一个问题是，$H^2(X_\Sigma)\to H^2(X)$ 未必为满射，故 $J_X$ 可能涉及比 $J_{\mathcal V}$ 更多的变量。若令 $H^*(\mathcal V)$ 为 $i^*:H^*(X_\Sigma)\to H^*(X)$ 的像，可能须先把 $J_X$ 适当地投影到 $H^*(\mathcal V)$ 上，把它称为 $J_X$ 的“环面”部分，再对该部分应用 $i_!$。一般而言，我们认为 $J_{\mathcal V}$ 与 $J_X$ 的确切关系值得进一步研究。

最后，还应评论此定理在 Calabi--Yau 完全交情形下的含义。如\hyperref[chap:mirror-constructions]{第四章}，设 $\Delta$ 为\IndexLocation{indexloc:1218:405}自反多胞形，$X_\Sigma$ 为对相应扇作\IndexLocation{indexloc:1096:405}极大射影细分得到的环面簇。进一步假设 $X_\Sigma$ 光滑。
\SourcePage{406}
为简便起见，只考虑 $V\subset X_\Sigma$ 为反典范除子的情形，并可假设 $V$ 光滑。于是定理~\ref{thm:givental-toric-mirror} 的 nef 假设当然成立，$J_{\mathcal V}$ 与 $I_{\mathcal V}$ 因而由坐标变换联系。如刚解释的，$J_{\mathcal V}$ 反映来自 $V$ 的“环面”部分的亏格零引力关联数。所以，形式函数 $I_{\mathcal V}$ 想必应告诉我们一些关于镜像 $V^\circ$ 的信息。

\hyperref[sec:hypergeometric]{第~\ref*{sec:hypergeometric} 节}研究过这一点，现在回顾当时的发现。首先，由\hyperref[subsec:nef-gkz]{第~\ref*{subsec:nef-gkz} 小节}，$I_{\mathcal V}$ 满足由式~\eqref{eq:nef-gkz-configuration} 中集合 $\mathcal A$ 确定的 GKZ 方程。因 $\Sigma(1)=\Delta^\circ\cap N-\{0\}$，可见 $\mathcal A=(\Delta^\circ\cap N)\times\{1\}$。令 $V^\circ$ 为 $V$ 的 Batyrev 镜像，它由自反多胞形 $\Delta^\circ$ 确定。再回顾\hyperref[subsec:toric-gkz]{第~\ref*{subsec:toric-gkz} 小节}，$\mathcal A=(\Delta^\circ\cap N)\times\{1\}$ 的 GKZ 方程是 $V^\circ$ 的 Picard--Fuchs 方程。看起来很好，但有两个困难。首先，GKZ 方程并不包含全部 Picard--Fuchs 方程，因此尚未证明 $I_{\mathcal V}$ 来自 $V^\circ$ 的周期，虽然 \SourceCite{Givental4} 未加证明地声称如此。其次，该小节并不处理 $V^\circ$ 的完整模空间，而只处理由环境环面簇所诱导的 $V^\circ$ 的形变给出的较小模空间。

因此，还未完全理解此定理与本书前几章所讨论镜像对称的关系。当 Calabi--Yau 超曲面 $V$ 为三维簇时，我们希望\hyperref[subsec:toric-mirror-conjecture]{第~\ref*{subsec:toric-mirror-conjecture} 小节}的猜想可由此定理推出。下面的例题~\ref{ex:octic-toric-mirror-theorem} 将说明这是怎样实现的。一般而言，镜像对称与此定理的关系，确实应进一步研究。

用两个例题结束本小节，展示此定理的效力。第一个是本书中已多次见过的例子，即 $\mathbb P(1,1,2,2,2)$ 的极小解奇异中的 Calabi--Yau 超曲面。
\SourceExampleNumber{11.2.5}{1}
\begin{example}\label{ex:octic-toric-mirror-theorem}
设 $\Sigma$ 是从 $\mathbb P(1,1,2,2,2)$ 的通常扇得到的扇：细分由 $(-1,-2,-2,-2)$ 与 $(1,0,0,0)$ 张成的锥，插入由 $(0,-1,-1,-1)$ 张成的射线。对于此扇，
\begin{equation}\label{eq:octic-toric-mirror-rays}
\begin{aligned}
\Sigma(1)=\{&(-1,-2,-2,-2),(1,0,0,0),(0,1,0,0),\\
&(0,0,1,0),(0,0,0,1),(0,-1,-1,-1)\},
\end{aligned}
\end{equation}
其中稍微滥用记号，把一维锥与其本原生成元等同。

此扇是 $\mathbb P(1,1,2,2,2)$ 标准扇的唯一极大射影细分。注意，式~\eqref{eq:octic-toric-mirror-rays} 中各点的凸包，正是 Batyrev 构造中使用的多胞形 $\Delta^\circ$。在此背景下，得到 Calabi--Yau 超曲面 $V\subset X_\Sigma$。例题~\ref{ex:weighted-extremal-transition} 证明，Kähler 锥由环面除子 $D_1,D_3$ 的类张成，这里按式~\eqref{eq:octic-toric-mirror-rays} 的次序为各环面除子编号。该例还指出线性等价关系 $D_1\sim D_2$、$D_3\sim D_4\sim D_5$、$D_6\sim-2D_1+D_3$，由此可用 $D_1,D_3$ 表示全部除子类。此外，$V\in|4D_3|$。按本小节的记号，令 $\mathcal V=\mathcal O_{X_\Sigma}(4D_3)$。

为与前几章的记号一致，将 Kähler 锥生成元排序为 $(D_3,D_1)$。令 $\{\beta_1,\beta_2\}$ 为 $M(X_\Sigma)$ 中与 $\{D_3,D_1\}$ 对偶的基，于是一般元素 $\beta\in M(X_\Sigma)$ 为 $\beta=d_1\beta_1+d_2\beta_2$。例题~\ref{ex:weighted-extremal-transition} 所计算的交数
\SourcePage{407}
给出 $\beta_1=D_3^2D_1$、$\beta_2=D_3^3-2D_3^2D_1$；这里通过 Poincaré 对偶，将 $\beta_i$ 看作 $H^6(X_\Sigma)$ 的元素。该例还指出，$|D_1|$ 定义映射 $X_\Sigma\to\mathbb P^1$，故 $D_1^2=0$。

为考察环面镜像定理在此例中的含义，先写出 $I_{\mathcal V}$。将式~\eqref{eq:toric-complete-intersection-i-function} 用于本例，得到
\begin{equation}\label{eq:octic-toric-i-function}
\begin{aligned}
I_{\mathcal V}={}&e^{(t_0+t_1D_3+t_2D_1)/\hbar}4D_3\sum_{d_1,d_2\ge0}z_1^{d_1}z_2^{d_2}\\
&\times\frac{\prod_{m=1}^{4d_1}(4D_3+m\hbar)\prod_{m=-\infty}^0(D_3-2D_1+m\hbar)}{\prod_{m=1}^{d_2}(D_1+m\hbar)^2\prod_{m=1}^{d_1}(D_3+m\hbar)^3\prod_{m=-\infty}^{d_1-2d_2}(D_3-2D_1+m\hbar)}.
\end{aligned}
\end{equation}
注意，对 $i=1,2$ 设 $z_i=e^{t_i}$。用 $z_i$ 而非 $q_i$，是为了与\hyperref[chap:hodge-yukawa]{第五章}、\hyperref[chap:moduli]{第六章}一致。这些坐标 $z_i$ 在例题~\ref{ex:weighted-complex-moduli} 中定义，并用于这两章中的多个例题。

环面镜像定理断言，适当变量变换后 $I_{\mathcal V}$ 等于 $J_{\mathcal V}$。此外，因
\[
J_{\mathcal V}=e^{(t_0+t_1D_3+t_2D_1)/\hbar}4D_3\bigl(1+o(\hbar^{-1})\bigr),
\]
可展开 $I_{\mathcal V}$，读出变量变换。由式~\eqref{eq:octic-toric-i-function}，
\begin{equation}\label{eq:octic-toric-i-period-expansion}
I_{\mathcal V}=e^{t_0/\hbar}4D_3\left(y_0+y_{1,1}\frac{D_3}{\hbar}+y_{1,2}\frac{D_1}{\hbar}
+y_{2,1}\frac{D_1D_3}{\hbar^2}+y_{2,2}\frac{D_3^2}{\hbar^2}+y_3\frac{D_1D_3^2}{\hbar^3}\right),
\end{equation}
其中
\[
y_0=\sum_{d_1\ge2d_2\ge0}z_1^{d_1}z_2^{d_2}\frac{(4d_1)!}{(d_2!)^2(d_1!)^3(d_1-2d_2)!},
\]
\[
\begin{aligned}
y_{1,1}={}&t_1y_0+\sum_{d_1\ge2d_2\ge0}z_1^{d_1}z_2^{d_2}\frac{(4d_1)!}{(d_2!)^2(d_1!)^3(d_1-2d_2)!}
\bigl(4\Psi(4d_1)-3\Psi(d_1)-\Psi(d_1-2d_2)\bigr)\\
&+\sum_{0\le d_1<2d_2}z_1^{d_1}z_2^{d_2}\frac{(-1)^{d_1+1}(2d_2-d_1-1)!(4d_1)!}{(d_2!)^2(d_1!)^3},
\end{aligned}
\]
\[
\begin{aligned}
y_{1,2}={}&t_2y_0+\sum_{d_1\ge2d_2\ge0}z_1^{d_1}z_2^{d_2}\frac{(4d_1)!}{(d_2!)^2(d_1!)^3(d_1-2d_2)!}
\bigl(2\Psi(d_1-2d_2)-2\Psi(d_2)\bigr)\\
&-2\sum_{0\le d_1<2d_2}z_1^{d_1}z_2^{d_2}\frac{(-1)^{d_1+1}(2d_2-d_1-1)!(4d_1)!}{(d_2!)^2(d_1!)^3}.
\end{aligned}
\]
这里使用了经典函数 $\Psi(n)=\sum_{i=1}^n1/i$，并约定 $\Psi(0)=0$。注意 $y_0,y_{1,1},y_{1,2}$ 是 $t_1,t_2$ 的函数。

由此容易看出，定理~\ref{thm:givental-toric-mirror} 中的 $f_0,f_1,f_2,h$ 为
\[
e^{f_0}=y_0,\qquad f_1=\frac{y_{1,1}}{y_0}-t_1,\qquad f_2=\frac{y_{1,2}}{y_0}-t_2,\qquad h=0,
\]
\SourcePage{408}
因此坐标变换为 $t_0\mapsto t_0+f_0\hbar$，以及 $t_i\mapsto t_i+f_i=y_{1,i}/y_0$，$i=1,2$。于是，对于
\[
(s_1,s_2)=\left(\frac{y_{1,1}(t_1,t_2)}{y_0(t_1,t_2)},\frac{y_{1,2}(t_1,t_2)}{y_0(t_1,t_2)}\right),
\]
环面镜像定理给出
\begin{equation}\label{eq:octic-toric-mirror-i-j-equality}
J_{\mathcal V}(s_1,s_2)=\frac{I_{\mathcal V}(t_1,t_2)}{y_0(t_1,t_2)}.
\end{equation}
令 $q_i=e^{s_i}$，上述变量变换成为
\[
(q_1,q_2)=\left(\exp\left(\frac{y_{1,1}}{y_0}\right),\exp\left(\frac{y_{1,2}}{y_0}\right)\right),
\]
这就是本例的镜像映射。例题~\ref{ex:weighted-hypergeometric-mirror-map} 用另一种方法计算了坐标 $q_2$，结果完全一致。

接着说明式~\eqref{eq:octic-toric-mirror-i-j-equality} 与用 A 模型、B 模型关联函数表述的“经典”镜像对称概念的关系。策略是证明 $I_{\mathcal V}/y_0$ 满足某些含 B 模型关联函数的微分方程。因此，$J_{\mathcal V}$ 满足相同的微分方程，进而得到所需的关联函数相等。

从 $I_{\mathcal V}$ 开始。由命题~\ref{prop:nef-i-gkz}，知道它满足来自 $V^\circ$ 的 GKZ 方程，但如前所述，这些并非 $V^\circ$ 的全部 Picard--Fuchs 方程。不过，例题~\ref{ex:weighted-gkz} 已计算了缺少的 Picard--Fuchs 方程，见式~\eqref{eq:weighted-gkz-second} 和~\eqref{eq:weighted-gkz-third}。直接计算表明 $I_{\mathcal V}$ 满足这些方程。（对此应当有一个一般定理，但我们尚不知道证明。）

例题~\ref{ex:weighted-gkz} 所构造的 Picard--Fuchs 方程采用了一个特定的全纯 $3$ 形式 $\Omega$。还知道，$y_0$ 是除常数倍外唯一在 $z_1=z_2=0$ 处全纯的周期。于是 $\widetilde\Omega=\Omega/y_0$ 是除常数倍外的归一化全纯 $3$ 形式，所以 $I_{\mathcal V}/y_0$ 满足 $\widetilde\Omega$ 的 Picard--Fuchs 方程。从变量 $z_i=e^{t_i}$ 变到 $q_i=e^{s_i}$ 后，仍然如此。

在坐标 $q_i$ 下，$V^\circ$ 的归一化 Yukawa 耦合导出一些漂亮的 Picard--Fuchs 方程。此时，归一化耦合为
\[
Y_{ijk}=-\frac1{(2\pi i)^3}\int_{V^\circ}\widetilde\Omega\wedge\nabla_{\delta_i}\nabla_{\delta_j}\nabla_{\delta_k}\widetilde\Omega,
\]
其中 $\delta_j=\partial/\partial s_j$，$\nabla$ 是镜像族的 Gauss--Manin 联络。这与\hyperref[subsec:toric-mirror-conjecture]{第~\ref*{subsec:toric-mirror-conjecture} 小节}的定义在适当的常数归一化下相对应；那里写 $q_j=e^{2\pi iu_j}$，而这里使用 $q_j=e^{s_j}$。\footnote{这在先验上可能不同于定义~\ref{def:normalized-yukawa} 中的归一化耦合。那里假设整性猜想，并使用\hyperref[subsec:mirror-map-definition]{第~\ref*{subsec:mirror-map-definition} 小节}定义的镜像映射。这里不假设整性猜想，而使用\hyperref[subsec:toric-mirror-map]{第~\ref*{subsec:toric-mirror-map} 小节}给出的环面版本镜像映射。利用 \SourceCite{CdFKM} 中对单值性的详细描述，可验证本例中二者相同。但我们更希望说明，即使不了解单值性的全部细节，仍能得到良好结果。}采用 $H^2(V)$ 的上同调基 $T_1=i^*(D_3)$、$T_2=i^*(D_1)$，目标是证明
\begin{equation}\label{eq:octic-toric-mirror-correlation-equality}
\langle T_i,T_j,T_k\rangle=Y_{ijk}\qquad\text{对所有 }i,j,k.
\end{equation}
\SourcePage{409}
回顾例题~\ref{ex:octic-toric-mirror-correlations}，可显式计算 $Y_{ijk}$。于是该式给出 $V$ 的全部 A 模型关联函数 $\langle T_i,T_j,T_k\rangle$ 的显式公式。

为证明此式，先注意可用 $Y_{ijk}$ 描述 Gauss--Manin 联络。更确切地，断言 $H^3(V^\circ,\mathbb C)$ 有一组基 $e_0,e_1,e_2,e^1,e^2,e^0$，使 $e_0=\widetilde\Omega$，且
\[
\nabla_{\delta_i}(e_0)=e_i,\qquad\nabla_{\delta_i}(e_j)=\sum_kY_{ijk}e^k,\qquad
\nabla_{\delta_i}(e^j)=\delta_{ij}e^0,\qquad\nabla_{\delta_i}(e^0)=0,
\]
与式~\eqref{eq:multimoduli-connection} 相同。对于\hyperref[subsec:normalized-yukawa-mirror]{第~\ref*{subsec:normalized-yukawa-mirror} 小节}定义的归一化 Yukawa 耦合，命题~\ref{prop:normalized-yukawa-pf} 及其证明后的讨论已给出这一结果。虽然证明要求整性猜想，但只要在 $\mathbb C$ 上而非 $\mathbb Z$ 上工作，所有内容都适用于此情形。

利用上述基，直接计算可得
\[
\nabla_{\delta_1}\left(\frac{Y_{112}\nabla_{\delta_1}^2-Y_{111}\nabla_{\delta_1}\nabla_{\delta_2}}{Y_{112}^2-Y_{111}Y_{122}}\right)\widetilde\Omega=0,
\]
因为 $\widetilde\Omega=e_0$。因此，这是一个 Picard--Fuchs 方程，而 $I_{\mathcal V}/y_0$ 满足 Picard--Fuchs 方程，故
\[
\frac\partial{\partial s_1}\left(\frac{Y_{112}\frac{\partial^2}{\partial s_1^2}-Y_{111}\frac{\partial^2}{\partial s_1\partial s_2}}{Y_{112}^2-Y_{111}Y_{122}}\right)\frac{I_{\mathcal V}}{y_0}=0.
\]
由环面镜像定理~\eqref{eq:octic-toric-mirror-i-j-equality}，得到
\begin{equation}\label{eq:octic-toric-j-yukawa-differential-equation}
\frac\partial{\partial s_1}\left(\frac{Y_{112}\frac{\partial^2}{\partial s_1^2}-Y_{111}\frac{\partial^2}{\partial s_1\partial s_2}}{Y_{112}^2-Y_{111}Y_{122}}\right)J_{\mathcal V}(s_1,s_2)=0,
\end{equation}
这里忽略了 $t_0$ 的依赖。还应指出，显式计算表明 $Y_{112}^2-Y_{111}Y_{122}$ 不恒为零。

下一步，将它与 $V$ 的 Givental 函数 $J_V$ 联系。包含映射 $i:V\subset X_\Sigma$ 诱导同构 $i^*:H^2(X_\Sigma)\to H^2(V)$，例题~\ref{ex:weighted-extremal-transition} 已说明，$V$ 的 Kähler 锥由 $T_1=i^*(D_3)$、$T_2=i^*(D_1)$ 张成。因此，$J_{\mathcal V}$ 与 $J_V$ 都可看作 $s_1,s_2$ 的函数，仍忽略 $t_0$。不过，尚不知道 $i_!(J_V)$ 是否等于 $J_{\mathcal V}$。环境空间为 $\mathbb P^n$ 时，命题~\ref{prop:complete-intersection-j-pushforward} 证明了二者相等。但证明用了 $\mathbb P^n$ 的凸性，此时的 $X_\Sigma$ 并不凸。考察该命题的证明，可验证：只要
\begin{equation}\label{eq:octic-virtual-class-pushforward-assumption}
j_*\bigl([\overline M_{0,2}(V,\beta)]^{\mathrm{virt}}\bigr)
=\operatorname{Euler}(\mathcal V_{\beta,2})\cap[\overline M_{0,2}(X_\Sigma,\beta)]^{\mathrm{virt}},
\end{equation}
仍有 $i_!(J_V)=J_{\mathcal V}$。我们推测此式为真，但尚无完整证明。因此，本例余下部分将\emph{假设}该式成立。特别地，这允许假设 $i_!(J_V)=J_{\mathcal V}$。

还可验证，本例中 $i_!$ 在偶次上同调上为单射。结合式~\eqref{eq:octic-toric-j-yukawa-differential-equation} 与 $i_!(J_V)=J_{\mathcal V}$，得到
\begin{equation}\label{eq:octic-ordinary-j-yukawa-quantum-equation}
\hbar\frac\partial{\partial s_1}\left(\frac{Y_{112}\hbar^2\frac{\partial^2}{\partial s_1^2}-Y_{111}\hbar^2\frac{\partial^2}{\partial s_1\partial s_2}}{Y_{112}^2-Y_{111}Y_{122}}\right)J_V(s_1,s_2)=0.
\end{equation}
加入各个 $\hbar$ 因子，是为了将此式写成\hyperref[sec:quantum-cohomology-relations]{第~\ref*{sec:quantum-cohomology-relations} 节}意义下的\IndexLocation{indexloc:1150:409}量子微分方程。为应用定理~\ref{thm:quantum-differential-relations}，展开此式，
\SourcePage{410}
得到
\[
\left(\frac{Y_{112}\hbar^3\frac{\partial^3}{\partial s_1^3}-Y_{111}\hbar^3\frac{\partial^3}{\partial s_1^2\partial s_2}}{Y_{112}^2-Y_{111}Y_{122}}
+\hbar P(\hbar\partial/\partial s,q,\hbar)\right)J_V(s_1,s_2)=0,
\]
其中 $P(\hbar\partial/\partial s,q,\hbar)$ 是某个微分算子。定理~\ref{thm:quantum-differential-relations} 表明，代入 $\hbar\partial/\partial s_i\mapsto T_i$、$\hbar\mapsto0$，得到 $V$ 的小量子上同调中的多项式关系。因此，$H^*(V)$ 中有
\[
\frac{Y_{112}T_1*T_1*T_1-Y_{111}T_1*T_1*T_2}{Y_{112}^2-Y_{111}Y_{122}}=0.
\]
由式~\eqref{eq:three-point-quantum-integral}，
\[
\int_V T_i*T_j*T_k=\langle T_i,T_j,T_k\rangle,
\]
所以上式给出
\[
Y_{112}\langle T_1,T_1,T_1\rangle-Y_{111}\langle T_1,T_1,T_2\rangle=0.
\]
刚才的论证适用于一般下标，故对任意 $1\le i,j,k,a,b,c\le2$，均有
\[
Y_{ijk}\langle T_a,T_b,T_c\rangle-Y_{abc}\langle T_i,T_j,T_k\rangle=0.
\]
由此可知，有一个共同因子 $c$，使
\[
\langle T_i,T_j,T_k\rangle=cY_{ijk}\qquad\text{对所有 }i,j,k.
\]
两侧都来自势函数，故各四阶混合导数的相容性给出 $(\partial_1c)Y_{12k}=(\partial_2c)Y_{11k}$，$k=1,2$。因为 $Y_{112}^2-Y_{111}Y_{122}$ 不恒为零，得到 $\partial_1c=\partial_2c=0$，所以 $c$ 为常数。

把 $\langle T_i,T_j,T_k\rangle$ 看作 $q_1,q_2$ 的形式函数，有 $\langle T_1,T_1,T_1\rangle(0,0)=\int_VT_1^3=8$。因此，为证明 $c=1$，须证明 $Y_{111}(0,0)=8$。令 $N_1$ 为 $V^\circ$ 的复模空间中边界除子 $z_1=0$ 周围单值性变换的对数，沿用例题~\ref{ex:weighted-yukawa} 的记号。\SourceCite{CdFKM} 的显式计算表明，有 $\mathcal H$ 的整平坦截面 $\beta^2,\alpha_2$，它们是一组 $\mathbb Z$ 基的部分元素，且满足 $\langle\alpha_2,\beta^2\rangle=1$、$N_1^3(\beta^2)=-8\alpha_2$。注意，\SourceCite{CdFKM} 中单值性的取向与这里相反。类似于例题~\ref{ex:quintic-normalized-yukawa} 的论证给出 $Y_{111}(0,0)=8$。故 $c=1$，式~\eqref{eq:octic-toric-mirror-correlation-equality} 得证。

例题~\ref{ex:octic-toric-mirror-correlations} 把 $\langle T_1,T_1,T_1\rangle$ 记为 $\langle H,H,H\rangle$，现在已证明该例中这个关联函数的公式。还可把式~\eqref{eq:octic-toric-mirror-correlation-equality} 与\hyperref[subsec:toric-mirror-conjecture]{第~\ref*{subsec:toric-mirror-conjecture} 小节}中镜像猜想的 Hodge 理论版本联系。如式~\eqref{eq:toric-mirror-correlations} 附近的讨论，A 模型与 B 模型关联函数在相差常数的意义下相等，意味着本例中猜想~\ref{conj:toric-hodge-mirror-surjective} 成立。更确切地，对于这里考虑的三维簇 $V,V^\circ$，已证明该猜想由定理~\ref{thm:givental-toric-mirror} 和式~\eqref{eq:octic-virtual-class-pushforward-assumption} 推出。此外，利用上一段，还可说明 $(V^\circ,V)$ 是定义~\ref{def:mathematical-mirror-pair} 意义下的数学镜像对。

不过，在充分发挥环面镜像定理的作用之前，还有更多内容可说。将看到，式~\eqref{eq:octic-toric-mirror-i-j-equality} 给出 $V$ 的 Gromov--Witten 势的一个漂亮公式。采用 $V$ 的偶次上同调的以下基：
\[
\begin{gathered}
T_0=1,\qquad T_1=i^*(D_3),\qquad T_2=i^*(D_1),\\
T^1=\frac{i^*(D_1D_3)}4,\qquad T^2=\frac{i^*(D_3^2-2D_1D_3)}4,\qquad T^0=\frac{i^*(D_1D_3^2)}4.
\end{gathered}
\]
\SourcePage{411}
可验证，杯积下 $T^1,T^2$ 与 $T_1,T_2$ 对偶，$T^0$ 与 $T_0$ 亦然。此时，变量为 $t_0,s_1,s_2$；由命题~\ref{prop:cy-threefold-j-potential}，$J_V$ 由公式
\begin{equation}\label{eq:octic-ordinary-j-potential}
\begin{aligned}
J_V=e^{t_0/\hbar}\Biggl(&1+\hbar^{-1}(s_1T_1+s_2T_2)
+\hbar^{-2}\left(\frac{\partial\Phi}{\partial s_1}T^1+\frac{\partial\Phi}{\partial s_2}T^2\right)\\
&+\hbar^{-3}\left(s_1\frac{\partial\Phi}{\partial s_1}+s_2\frac{\partial\Phi}{\partial s_2}-2\Phi\right)T^0\Biggr)
\end{aligned}
\end{equation}
给出，其中 $\Phi$ 是 $V$ 的 Gromov--Witten 势。

仍假设式~\eqref{eq:octic-virtual-class-pushforward-assumption}，故 $J_{\mathcal V}=i_!(J_V)$。因为 $i_!(i^*\alpha)=\alpha\cup4D_3$，对式~\eqref{eq:octic-ordinary-j-potential} 应用 $i_!$，得到
\[
\begin{aligned}
J_{\mathcal V}(s_1,s_2)=e^{t_0/\hbar}4D_3\Biggl(&1+s_1\frac{D_3}{\hbar}+s_2\frac{D_1}{\hbar}
+\frac{\partial\Phi}{\partial s_1}\frac{D_1D_3}{4\hbar^2}
+\frac{\partial\Phi}{\partial s_2}\frac{D_3^2-2D_1D_3}{4\hbar^2}\\
&+\left(s_1\frac{\partial\Phi}{\partial s_1}+s_2\frac{\partial\Phi}{\partial s_2}-2\Phi\right)\frac{D_1D_3^2}{4\hbar^3}\Biggr).
\end{aligned}
\]
为与 $I_{\mathcal V}/y_0$ 比较，注意式~\eqref{eq:octic-toric-i-period-expansion} 可写成
\[
\begin{aligned}
\frac{I_{\mathcal V}}{y_0}(t_1,t_2)=e^{t_0/\hbar}4D_3\Biggl(&1+\frac{y_{1,1}}{y_0}\frac{D_3}{\hbar}+\frac{y_{1,2}}{y_0}\frac{D_1}{\hbar}
+\frac{y_{2,1}+2y_{2,2}}{y_0}\frac{D_1D_3}{\hbar^2}\\
&+\frac{y_{2,2}}{y_0}\frac{D_3^2-2D_1D_3}{\hbar^2}+\frac{y_3}{y_0}\frac{D_1D_3^2}{\hbar^3}\Biggr).
\end{aligned}
\]
将镜像映射 $(s_1,s_2)=(y_{1,1}/y_0,y_{1,2}/y_0)$ 记为 $(s_1,s_2)=\Psi(t_1,t_2)$。如式~\eqref{eq:octic-toric-mirror-i-j-equality}，环面\IndexLocation{indexloc:0855:411}\IndexLocation{indexloc:1013:411}镜像定理表明，上述两个表达式在变量变换 $(s_1,s_2)=\Psi(t_1,t_2)$ 后相等。比较系数，得到 $V$ 的 Gromov--Witten 势的显式公式：
\begin{equation}\label{eq:octic-toric-mirror-potential}
\Phi(\Psi(t_1,t_2))=2\left(\frac{y_{1,1}}{y_0}\frac{y_{2,1}+2y_{2,2}}{y_0}
+\frac{y_{1,2}}{y_0}\frac{y_{2,2}}{y_0}-\frac{y_3}{y_0}\right),
\end{equation}
其中 $\Psi(t_1,t_2)=(y_{1,1}/y_0,y_{1,2}/y_0)$ 是镜像映射。这与\hyperref[sec:quintic-mirror-theorem]{第~\ref*{sec:quintic-mirror-theorem} 节}导出的\IndexLocation{indexloc:1190:411}五次三维簇的 Gromov--Witten 势公式十分相似。

最后一个观察是，比较上述 $J_{\mathcal V}$ 和 $I_{\mathcal V}/y_0$ 的公式，也给出 A \IndexLocation{indexloc:1003:411}模型关联函数的一些漂亮公式。例如，已知
\[
\langle T_1,T_1,T_1\rangle=\frac{\partial^3\Phi}{\partial s_1^3}.
\]
于是，把 $J_{\mathcal V}(s_1,s_2)=I_{\mathcal V}/y_0(t_1,t_2)$ 对 $s_1$ 求两次导数，再比较 $D_1D_3$ 的系数，得到
\begin{equation}\label{eq:octic-yukawa-period-second-derivative}
\langle T_1,T_1,T_1\rangle=4\frac{\partial^2}{\partial s_1^2}\left(\frac{y_{2,1}+2y_{2,2}}{y_0}\right).
\end{equation}
这个公式比例题~\ref{ex:octic-toric-mirror-correlations} 中考虑的公式\emph{容易推导得多}。也可将其看作归一化 Yukawa 耦合 $Y_{111}$ 的一个简洁公式。

此例很有启发性，因为其中许多计算应能推广到其他环面 Calabi--Yau 三维簇。特别地：
\SourcePage{412}
\begin{itemize}
\item 式~\eqref{eq:octic-toric-mirror-potential} 应能推广到任意作为环面簇 $X_\Sigma$ 中 nef 完全交的 Calabi--Yau 三维簇 $V$。由于 $H^2(X_\Sigma)\to H^2(V)$ 未必为满射，须定义 $V$ 的 Gromov--Witten 势的“环面”部分，再用类似于该式的公式，显式地以 $I_{\mathcal V}/y_0$ 表示它。
\item Gromov--Witten 势是 A 模型关联函数 $\langle T_i,T_j,T_k\rangle$ 的势函数。因此，该式表明，$I_{\mathcal V}/y_0$ 的系数给出 $V^\circ$ 的 B 模型关联函数 $Y_{ijk}$ 的一个显式势函数，含义如引理~\ref{lem:yukawa-potential}。应该可以直接证明这一点，而不使用环面镜像定理。
\item 式~\eqref{eq:octic-yukawa-period-second-derivative} 实质上给出了 $Y_{111}$ 的简单公式。一般而言，$V^\circ$ 的归一化 Yukawa 耦合应有类似的公式，以 $I_{\mathcal V}/y_0$ 表示。如上，也应该能够直接证明这一点。
\end{itemize}
我们认为，环面簇中的三维 nef 完全交还有许多有趣的未解问题。
\end{example}
作为定理~\ref{thm:givental-toric-mirror} 的第二个应用，研究 $\mathcal V=0$ 时的情形，此时 $X$ 就是环面簇 $X_\Sigma$。将看到，镜像定理与前面的结果很好地结合，描述 $X_\Sigma$ 的量子上同调。
\SourceExampleNumber{11.2.5}{2}
\begin{example}\label{ex:toric-fano-and-hirzebruch-quantum-ring}
本例将说明，若 $X=X_\Sigma$ 为光滑 Fano 环面簇，则例题~\ref{ex:batyrev-quantum-ring} 的 Batyrev 量子环 $H_\omega^*(X)$ 与小量子上同调环相同。将结合定理~\ref{thm:givental-toric-mirror} 与命题~\ref{prop:nef-i-gkz} 证明这一点。还将看到，当 $-K_X$ 为 nef 而非丰沛时，$H_\omega^*(X)$ 可能不同于量子上同调环，但该定理仍意味着两环之间存在关系。以 $X=\mathbb F_2$ 为例说明。

既然考虑环面簇 $X=X_\Sigma$ 本身，就在式~\eqref{eq:toric-complete-intersection-i-function} 中令 $\mathcal V$ 为零秩平凡丛，并分别将 $I_{\mathcal V}$、$J_{\mathcal V}$ 记为 $I_X$、$J_X$。注意，$X=\mathbb P^n$ 时，这与命题~\ref{prop:projective-space-i-equals-j} 的记号一致。有
\begin{equation}\label{eq:toric-zero-bundle-i-function}
I_X=e^{(t_0+\delta)/\hbar}\sum_{\beta\in M(X)}q^\beta
\frac{\prod_\rho\prod_{m=-\infty}^0(D_\rho+m\hbar)}{\prod_\rho\prod_{m=-\infty}^{D_\rho(\beta)}(D_\rho+m\hbar)},
\end{equation}
其中 $\delta=\sum_{i=1}^rt_iT_i$、$D_\rho(\beta)=\int_\beta D_\rho$。\hyperref[subsec:givental-nef-projective-mirror]{第~\ref*{subsec:givental-nef-projective-mirror} 小节}的式~\eqref{eq:complete-intersection-ordinary-j-function} 还给出了 $J_X$ 的公式。于是得到命题~\ref{prop:projective-space-i-equals-j} 的以下推广。
\begin{proposition}\label{prop:toric-fano-i-equals-j}
若 $X=X_\Sigma$ 为光滑 Fano 环面簇，则 $I_X=J_X$。
\end{proposition}
\begin{proof}
由式~\eqref{eq:j-function-asymptotic-normalization}，$J_X=e^{(t_0+\delta)/\hbar}(1+o(\hbar^{-1}))$。若知道 $X$ 为 Fano 时 $I_X$ 也如此，定理~\ref{thm:givental-toric-mirror} 的变量变换便是恒等变换，命题就由环面镜像定理立即推出。因此，只须证明
\begin{equation}\label{eq:toric-fano-i-normalization}
I_X=e^{(t_0+\delta)/\hbar}(1+o(\hbar^{-1})).
\end{equation}
为此，固定 Mori 锥 $M(X)$ 中的非零整类 $\beta$ 与 $\rho\in\Sigma(1)$，考察式~\eqref{eq:toric-zero-bundle-i-function} 中的分式
\[
\frac{\prod_{m=-\infty}^0(D_\rho+m\hbar)}{\prod_{m=-\infty}^{D_\rho(\beta)}(D_\rho+m\hbar)}.
\]
\SourcePage{413}
将此分式按 $\hbar^{-1}$ 的幂展开，研究首项。若 $D_\rho(\beta)>0$，首项容易看出为
\[
\frac1{D_\rho(\beta)!}\hbar^{-D_\rho(\beta)}.
\]
若 $D_\rho(\beta)=0$，分式就是 $1$；若 $D_\rho(\beta)=-1$，分式就是 $D_\rho$。最后，若 $D_\rho(\beta)<-1$，首项为
\[
(-1)^{-D_\rho(\beta)-1}D_\rho\,(-D_\rho(\beta)-1)!\,\hbar^{-D_\rho(\beta)-1}.
\]
所有情形下，首项中 $\hbar^{-1}$ 的幂次都至少为 $D_\rho(\beta)$，而且当且仅当 $D_\rho(\beta)<0$ 时严格大于它。

固定 $\beta$，对所有 $\rho$ 相乘，可见式~\eqref{eq:toric-zero-bundle-i-function} 展开中 $q^\beta$ 的系数的首项，其 $\hbar^{-1}$ 的幂次至少为 $\sum_\rho D_\rho(\beta)$。\hyperref[chap:toric]{第三章}已说明 $\sum_\rho D_\rho=-K_X$，因此该幂次至少为 $-\int_\beta K_X$。

$X$ 为 Fano，故 $-\int_\beta K_X$ 严格为正，所以首项中 $\hbar^{-1}$ 的幂次至少为一。但为证明式~\eqref{eq:toric-fano-i-normalization}，须证明至少为二。上述讨论表明，首项幂次为一的唯一可能是 $-\int_\beta K_X=1$，且所有 $\rho$ 均有 $D_\rho(\beta)\ge0$。因 $-K_X=\sum_\rho D_\rho$，可见存在 $\rho_0$，使 $D_{\rho_0}(\beta)=1$，而所有 $\rho\ne\rho_0$ 的 $D_\rho(\beta)=0$。须证明这不可能发生。

使用\hyperref[subsec:nef-gkz]{第~\ref*{subsec:nef-gkz} 小节}的引理~\ref{lem:nef-relation-lattice}，取 $k=0$。对于 $\beta\in H_2(X,\mathbb Z)$，该引理给出
\begin{equation}\label{eq:toric-curve-divisor-ray-relation}
\sum_\rho D_\rho(\beta)v_\rho=0,
\end{equation}
其中 $v_\rho$ 是 $\rho\cap N$ 的本原生成元，因为 $D_\rho(\beta)=D_\rho\cdot\beta$。上一段的情形会导致 $v_{\rho_0}=0$，矛盾。于是式~\eqref{eq:toric-fano-i-normalization} 得证，命题立即由定理~\ref{thm:givental-toric-mirror} 推出。
\end{proof}
接下来用 $I_X=J_X$ 研究光滑 Fano 环面簇 $X=X_\Sigma$ 的\IndexLocation{indexloc:1141:413}\IndexLocation{indexloc:1146:413}小量子上同调环。基本思想是：由\hyperref[subsec:nef-gkz]{第~\ref*{subsec:nef-gkz} 小节}，$I_X$ 满足一个 GKZ 方程组，故 $J_X$ 满足相同的方程；随后用\hyperref[sec:quantum-cohomology-relations]{第~\ref*{sec:quantum-cohomology-relations} 节}的\IndexLocation{indexloc:1150:413}量子微分方程理论，得到所需的量子上同调关系。

为实施此策略，须看看命题~\ref{prop:nef-i-gkz} 在这里的形式。这里 $\mathcal V=\bigoplus_{i=1}^k\mathcal L_i$ 为零丛，故 $k=0$。于是式~\eqref{eq:nef-gkz-configuration} 化为
\[
\mathcal A=\{v_\rho:\rho\in\Sigma(1)\},
\]
而引理~\ref{lem:nef-relation-lattice} 中的格 $\Lambda$，恰为 $v_\rho$ 之间整关系构成的格。此外，该引理通过将 $\beta$ 映到式~\eqref{eq:toric-curve-divisor-ray-relation} 的映射，把 $H_2(X,\mathbb Z)$ 与 $\Lambda$ 等同。

$\Lambda$ 同时也是相应 $\mathcal A$ 方程组的关系格，所以对于每个 $\beta\in\Lambda$，式~\eqref{eq:nef-gkz-box} 给出算子
\[
\Box_\beta=\prod_{D_\rho(\beta)>0}\partial_\rho^{D_\rho(\beta)}
-\prod_{D_\rho(\beta)<0}\partial_\rho^{-D_\rho(\beta)},
\]
\SourcePage{414}
因为 $\beta$ 对应于关系~\eqref{eq:toric-curve-divisor-ray-relation}。这里 $\partial_\rho=\partial/\partial\lambda_\rho$，由式~\eqref{eq:nef-gkz-q-coordinates}，用 $\lambda_\rho$ 定义 $q^\beta$ 如下：
\begin{equation}\label{eq:toric-quantum-monomial-lambda}
q^\beta=\prod_\rho\lambda_\rho^{D_\rho(\beta)}.
\end{equation}
若令
\[
\widetilde I=e^\delta\sum_{\beta\in M(X)}q^\beta
\frac{\prod_\rho\prod_{m=-\infty}^0(D_\rho+m)}{\prod_\rho\prod_{m=-\infty}^{D_\rho(\beta)}(D_\rho+m)},
\]
命题~\ref{prop:nef-i-gkz} 表明，所有 $\beta$ 均有 $\Box_\beta\widetilde I=0$。

注意，$\widetilde I$ 是在 $I_X$ 中置 $t_0=0$、$\hbar=1$ 所得。只须乘以 $e^{t_0/\hbar}$ 即可恢复 $t_0$，但恢复 $\hbar$ 须更小心：使用把 $\Box_\beta$ 中每个 $\partial_\rho$ 换成 $\hbar\partial_\rho$ 所得的算子 $\Box_{\beta,\hbar}$，即
\[
\Box_{\beta,\hbar}=\prod_{D_\rho(\beta)>0}(\hbar\partial_\rho)^{D_\rho(\beta)}
-\prod_{D_\rho(\beta)<0}(\hbar\partial_\rho)^{-D_\rho(\beta)}.
\]
容易证明 $\Box_\beta\widetilde I=0$ 蕴含 $\Box_{\beta,\hbar}I_X=0$，证明省略。

再定义
\[
\Box'_{\beta,\hbar}=\left(\prod_{D_\rho(\beta)>0}\lambda_\rho^{D_\rho(\beta)}\right)\Box_{\beta,\hbar},
\]
则 $\Box'_{\beta,\hbar}I_X=0$。令 $\delta_\rho=\lambda_\rho\partial_\rho$，将式~\eqref{eq:weighted-logarithmic-identities} 用于本情形，得到
\begin{equation}\label{eq:toric-quantum-gkz-operator}
\begin{aligned}
\Box'_{\beta,\hbar}={}&\prod_{D_\rho(\beta)>0}
\hbar\delta_\rho(\hbar\delta_\rho-\hbar)\cdots(\hbar\delta_\rho-(D_\rho(\beta)-1)\hbar)\\
&-q^\beta\prod_{D_\rho(\beta)<0}
\hbar\delta_\rho(\hbar\delta_\rho-\hbar)\cdots(\hbar\delta_\rho-(-D_\rho(\beta)-1)\hbar).
\end{aligned}
\end{equation}
$X$ 光滑且为 Fano，故命题~\ref{prop:toric-fano-i-equals-j} 给出 $I_X=J_X$，从而 $\Box'_{\beta,\hbar}J_X=0$。不过，在应用\hyperref[sec:quantum-cohomology-relations]{第~\ref*{sec:quantum-cohomology-relations} 节}的理论之前，须把 $\Box'_{\beta,\hbar}$ 写成关于 $\hbar\partial/\partial t_j$ 的微分算子。回顾 $t_j$ 是与 $H^2(X,\mathbb Z)$ 的基 $T_j$ 对应的变量。它们与 $\lambda_\rho$ 的联系如下：令 $\beta_j\in H_2(X,\mathbb Z)$ 为与 $T_j$ 对偶的基，再令
\[
q_j=\prod_\rho\lambda_\rho^{D_\rho(\beta_j)}.
\]
容易证明，$\prod_jq_j^{\int_\beta T_j}$ 等于式~\eqref{eq:toric-quantum-monomial-lambda} 所定义的 $q^\beta$。令 $q_j=e^{t_j}$，可见 $q^\beta$ 相对于 $t_j$ 具有通常的含义。

一般不能把式~\eqref{eq:toric-quantum-gkz-operator} 中的 $\hbar\delta_\rho$ 表示为 $\hbar\partial/\partial t_j$。关键观察是，对偶环面 $T^\circ=M\otimes\mathbb C^*$ 通过 $(m\otimes t)\cdot\lambda_\rho=t^{\langle m,v_\rho\rangle}\lambda_\rho$ 作用于这些变量。$q_j$ 在此作用下不变。如式~\eqref{eq:nef-logarithmic-identities} 所述，
\SourcePage{415}
对于 $T^\circ$ 不变的函数，有
\begin{equation}\label{eq:toric-quantum-gkz-coordinate-derivative}
\hbar\delta_\rho=\sum_j D_\rho(\beta_j)\hbar\frac\partial{\partial t_j}.
\end{equation}
此恒等式使式~\eqref{eq:toric-quantum-gkz-operator} 可用 $\hbar\partial/\partial t_j$ 写出，即 $\Box'_{\beta,\hbar}$ 是\hyperref[sec:quantum-cohomology-relations]{第~\ref*{sec:quantum-cohomology-relations} 节}意义下的量子微分算子。于是定理~\ref{thm:quantum-differential-relations} 表明，作代换
\[
\hbar\frac\partial{\partial t_j}\longmapsto T_j,\qquad\hbar\longmapsto0,
\]
就由该算子得到量子上同调关系。应用于式~\eqref{eq:toric-quantum-gkz-operator}，得到 $X$ 的小量子上同调中的关系
\begin{equation}\label{eq:toric-fano-quantum-stanley-reisner-relation}
\prod_{D_\rho(\beta)>0}D_\rho^{D_\rho(\beta)}
=q^\beta\prod_{D_\rho(\beta)<0}D_\rho^{-D_\rho(\beta)},
\end{equation}

可用它证明 Batyrev 量子环 $H_\omega^*(X)$ 同构于 $X$ 的小量子上同调环。令 $H^*(X)$ 采用小量子乘积 $*$，$\mathbb C[x_\rho]$ 为每个 $\rho\in\Sigma(1)$ 对应一个变量 $x_\rho$ 的多项式环。于是 $x_\rho\mapsto D_\rho$ 定义环同态 $\mathbb C[x_\rho]\to H^*(X)$。由上式，理想
\[
\left\langle\prod_{D_\rho(\beta)>0}x_\rho^{D_\rho(\beta)}
-q^\beta\prod_{D_\rho(\beta)<0}x_\rho^{-D_\rho(\beta)}:\beta\in H_2(X,\mathbb Z)\right\rangle
\]
包含于此映射的核。随着 $\beta$ 遍历 $H_2(X,\mathbb Z)$，关系~\eqref{eq:toric-curve-divisor-ray-relation}，即
\[
\sum_{D_\rho(\beta)>0}D_\rho(\beta)v_\rho
=\sum_{D_\rho(\beta)<0}-D_\rho(\beta)v_\rho,
\]
张成关系格 $\Lambda$。因此，上述理想就是式~\eqref{eq:quantum-sr-all-relations} 的理想，也就是 \SourceCite[定义 5.1]{Batyrev2} 给出的量子 Stanley--Reisner 理想 $SR_\omega(\Sigma)$ 的形式。于是 $SR_\omega(\Sigma)$ 包含于映射 $\mathbb C[x_\rho]\to H^*(X)$ 的核。此核还包含理想
\[
P(\Sigma)=\left\langle\sum_{i=1}^s\langle m,v_i\rangle x_i:m\in M\right\rangle,
\]
故有良定义的环同态
\[
\mathbb C[x_\rho]/(P(\Sigma)+SR_\omega(\Sigma))\longrightarrow H^*(X).
\]
例题~\ref{ex:batyrev-quantum-ring} 将左侧商环定义为 Batyrev 量子环 $H_\omega^*(X)$，所以得到环同态
\[
\varphi:H_\omega^*(X)\longrightarrow H^*(X),
\]
提醒读者，右侧采用小量子乘积。

现在断言 $\varphi$ 为同构。为此，回顾总可假设 $H^2(X)$ 的基 $T_j$ 位于 $X$ 的 Kähler 锥内部。这允许取极限 $q_j\to0$，此时 $\varphi$ 变为描述环面簇 $X$ 的通常上同调环的同构。因处理的是有限维向量空间，$\varphi$ 在 $q_j=0$ 的某个邻域中仍为同构。此外，若将 $q^\beta$ 看作形式变量，也得到同构。因此，$X$ 光滑且为 Fano 时，Batyrev 量子环与 $X$ 的通常小量子环相同，如所断言。
\SourcePage{416}
若削弱 $X$ 为 Fano 的假设，只假设 $-K_X$ 为 nef，则定理~\ref{thm:givental-toric-mirror} 仍适用，但结果更复杂，因为可能没有 $I_X=J_X$，而涉及非平凡变量变换。关键是，$-K_X$ 仅为 nef 时，可能有
\[
I_X\ne e^{(t_0+\delta)/\hbar}(1+o(\hbar^{-1})).
\]
于是该定理断言 $I_X$ 与 $J_X$ 由非平凡变量变换联系。这种情形下，仍可得到 $X$ 的小量子环中的有趣关系。命题~\ref{prop:nef-i-gkz} 仍适用，故 $\Box'_{\beta,\hbar}I_X=0$，其中算子由式~\eqref{eq:toric-quantum-gkz-operator} 给出。但不像 $\Box'_{\beta,\hbar}J_X=0$ 那样，直接用定理~\ref{thm:quantum-differential-relations} 得到式~\eqref{eq:toric-fano-quantum-stanley-reisner-relation}，这里须先改变变量，在 $H^*(X)$ 中得到该关系的一个“扭曲”版本。这正是 $H_\omega^*(X)$ 与 $X$ 的小量子上同调相联系的含义。

\begingroup\linespread{1.18}\selectfont
用 $X=\mathbb F_2$ 说明这一点。考虑 $\mathbb F_2$ 的标准扇 $\Sigma$，
\[
\Sigma(1)=\{(1,0),(0,1),(-1,2),(0,-1)\},
\]
如常把一维锥与其本原生成元等同。用 $D_1,D_2,D_3,D_4$ 表示相应环面除子，也用 $D_i$ 表示 $H^2(\mathbb F_2)$ 中的除子类。$\mathbb F_2$ 的 Kähler 锥由纤维类 $f$ 和满足 $H^2=2$ 的截面类 $H$ 张成。用环面除子表示，$D_1=D_3=f$、$D_2=H-2f$、$D_4=H$，因此 $-K_{\mathbb F_2}=\sum_iD_i=2H$，它为 nef。注意 $D_2$ 是 $\mathbb F_2$ 上的自交数为 $-2$ 的曲线。Mori 锥由这条曲线与纤维生成，即由 $H-2f$ 与 $f$ 生成。

须比较函数 $I_{\mathbb F_2}$ 与 $J_{\mathbb F_2}$。因为需要改变坐标，将沿用上一例题的约定，对 $I_{\mathbb F_2}$ 使用 $z_i=e^{t_i}$，对 $J_{\mathbb F_2}$ 使用 $q_i=e^{s_i}$。$H^2(\mathbb F_2)$ 的上同调基为 $T_1=f$、$T_2=H$。

先研究 $I_{\mathbb F_2}$，其公式为
\[
\begin{aligned}
I_{\mathbb F_2}(t_1,t_2)={}&e^{(t_0+t_1f+t_2H)/\hbar}\sum_{\beta\in M(\mathbb F_2)}z^\beta\\
&\times\frac{\prod_{m=-\infty}^0(f+m\hbar)^2\prod_{m=-\infty}^0(H-2f+m\hbar)\prod_{m=-\infty}^0(H+m\hbar)}{\prod_{m=-\infty}^{f\cdot\beta}(f+m\hbar)^2\prod_{m=-\infty}^{(H-2f)\cdot\beta}(H-2f+m\hbar)\prod_{m=-\infty}^{H\cdot\beta}(H+m\hbar)}.
\end{aligned}
\]
如命题~\ref{prop:toric-fano-i-equals-j} 的证明，特别关心其展开中 $\hbar^{-1}$ 的系数。该证明表明，此系数由满足 $-K_{\mathbb F_2}\cdot\beta=0$，且恰有一个 $i$ 满足 $D_i\cdot\beta<0$ 的有效类 $\beta$ 决定。容易计算，唯一的这些类为 $\beta=kD_2$，$k>0$。此时 $D_1\cdot\beta=D_3\cdot\beta=k$、$D_2\cdot\beta=-2k$、$D_4\cdot\beta=0$。

令 $\beta=kD_2=k(H-2f)$，$k\ge1$。计算得 $z^\beta=z_1^{f\cdot\beta}z_2^{H\cdot\beta}=z_1^k$，上述 $I_{\mathbb F_2}$ 公式中 $z_1^k$ 的系数为
\[
\frac{\prod_{m=-2k+1}^0(H-2f+m\hbar)}{\prod_{m=1}^k(f+m\hbar)^2}
=-(H-2f)\frac{(2k-1)!}{(k!)^2}\hbar^{-1}+o(\hbar^{-1}).
\]
因此，
\[
I_{\mathbb F_2}(t_1,t_2)=e^{(t_0+t_1f+t_2H)/\hbar}
\bigl(1-(H-2f)F(z_1)\hbar^{-1}+o(\hbar^{-1})\bigr),
\]
其中
\[
F(z_1)=\sum_{k=1}^{\infty}\frac{(2k-1)!}{(k!)^2}z_1^k.
\]
\SourcePage{417}
定理~\ref{thm:givental-toric-mirror} 要求
\[
f_0=0,\qquad f_1=2F(z_1),\qquad f_2=-F(z_1),\qquad h=0,
\]
随后环面镜像定理给出 $J_{\mathbb F_2}(s_1,s_2)=I_{\mathbb F_2}(t_1,t_2)$，其中
\[
(s_1,s_2)=(t_1+2F(e^{t_1}),t_2-F(e^{t_1})).
\]
置 $z_i=e^{t_i}$、$q_i=e^{s_i}$，即
\begin{equation}\label{eq:hirzebruch-two-mirror-coordinate-change}
(q_1,q_2)=(z_1e^{2F(z_1)},z_2e^{-F(z_1)}).
\end{equation}
令人意外的是，变量变换 $s_1=t_1+2F(z_1)$ 早已作为\hyperref[chap:moduli]{第六章}的式~\eqref{eq:weighted-boundary-log-series} 出现。与式~\eqref{eq:weighted-boundary-log-closed} 比较，得到
\[
s_1=\log\left(\frac{1-2z_1-\sqrt{1-4z_1}}{2z_1}\right),\qquad
q_1=\frac{1-2z_1-\sqrt{1-4z_1}}{2z_1}.
\]
于是 $z_1=q_1/(1+q_1)^2$，再容易得到上述坐标变换的逆为
\begin{equation}\label{eq:hirzebruch-two-inverse-mirror-coordinates}
(z_1,z_2)=\left(\frac{q_1}{(1+q_1)^2},q_2(1+q_1)\right),
\end{equation}
以及
\begin{equation}\label{eq:hirzebruch-two-inverse-log-coordinates}
(t_1,t_2)=\bigl(s_1-2\log(1+e^{s_1}),s_2+\log(1+e^{s_1})\bigr).
\end{equation}
这些坐标变换也出现在 \SourceCite{Givental4} 的一个例子中，只是变量名不同；那里考虑了 $\mathbb P^1$ 上的射影丛
$\mathbb P(\mathcal O_{\mathbb P^1}\oplus\mathcal O_{\mathbb P^1}(-1)\oplus\mathcal O_{\mathbb P^1}(-1))$ 与
$\mathbb P(\mathcal O_{\mathbb P^1}\oplus\mathcal O_{\mathbb P^1}(-2))$。

\looseness=-1 现在计算 $H_\omega^*(\mathbb F_2)$ 和 $\mathbb F_2$ 的小量子上同调中的一些关系。令 $\beta\in H_2(\mathbb F_2,\mathbb Z)$ 为 $f$ 的 Poincaré 对偶，计算算子 $\Box'_{\beta,\hbar}$。令 $\lambda_i$ 对应于 $D_i$，$\delta_i=\lambda_i\partial/\partial\lambda_i$。由式~\eqref{eq:toric-quantum-gkz-operator}，计算得
\[
\Box'_{\beta,\hbar}=\hbar\delta_2\hbar\delta_4-z_2,
\]
因为 $z^\beta=z_2$；回顾现在用 $z_i$ 而非 $q_i$。由式~\eqref{eq:toric-quantum-gkz-coordinate-derivative}，
\[
\hbar\delta_2=\hbar\frac\partial{\partial t_2}-2\hbar\frac\partial{\partial t_1},\qquad
\hbar\delta_4=\hbar\frac\partial{\partial t_2},
\]
随后命题~\ref{prop:nef-i-gkz} 给出
\begin{equation}\label{eq:hirzebruch-two-i-quantum-equation}
\left(\hbar\frac\partial{\partial t_2}\left(\hbar\frac\partial{\partial t_2}-2\hbar\frac\partial{\partial t_1}\right)-z_2\right)I_{\mathbb F_2}(t_1,t_2)=0.
\end{equation}
在 Batyrev 量子环 $H_\omega^*(\mathbb F_2)$ 中，这对应于关系
\begin{equation}\label{eq:hirzebruch-two-batyrev-quantum-relation}
H\mathbin{\cdot}(H-2f)=z_2.
\end{equation}
\looseness=-1 再应用环面镜像定理，考察它对 $J_{\mathbb F_2}$ 的含义。关键是，变量变换给出一个略有不同的微分方程。对式~\eqref{eq:hirzebruch-two-inverse-log-coordinates} 应用链式法则，得到
\[
\left(\hbar\frac\partial{\partial s_2}-2\hbar\frac\partial{\partial s_1}\right)
=\frac{1-q_1}{1+q_1}\left(\hbar\frac\partial{\partial t_2}-2\hbar\frac\partial{\partial t_1}\right),\qquad
\hbar\frac\partial{\partial s_2}=\hbar\frac\partial{\partial t_2}.
\]
\SourcePage{418}
结合这些等式，并利用式~\eqref{eq:hirzebruch-two-mirror-coordinate-change} 中 $q_1$ 与 $t_2$ 无关的事实，得到
\[
\begin{aligned}
&\frac{1-q_1}{1+q_1}\left(\hbar\frac\partial{\partial t_2}\left(\hbar\frac\partial{\partial t_2}-2\hbar\frac\partial{\partial t_1}\right)-z_2\right)\\
&\qquad=\frac{1-q_1}{1+q_1}\hbar\frac\partial{\partial t_2}\left(\hbar\frac\partial{\partial t_2}-2\hbar\frac\partial{\partial t_1}\right)-\frac{1-q_1}{1+q_1}z_2\\
&\qquad=\hbar\frac\partial{\partial s_2}\left(\hbar\frac\partial{\partial s_2}-2\hbar\frac\partial{\partial s_1}\right)-q_2(1-q_1),
\end{aligned}
\]
最后一行用了式~\eqref{eq:hirzebruch-two-inverse-mirror-coordinates} 中的 $z_2=q_2(1+q_1)$。

上式左侧算子消去 $I_{\mathbb F_2}(t_1,t_2)$，所以由定理~\ref{thm:givental-toric-mirror}，右侧算子消去 $J_{\mathbb F_2}(s_1,s_2)$。因此
\[
\hbar\frac\partial{\partial s_2}\left(\hbar\frac\partial{\partial s_2}-2\hbar\frac\partial{\partial s_1}\right)-q_2(1-q_1)
\]
是量子微分算子。由定理~\ref{thm:quantum-differential-relations}，代换
\[
\hbar\partial/\partial s_1\longmapsto T_1=f,\qquad\hbar\partial/\partial s_2\longmapsto T_2=H
\]
给出 $\mathbb F_2$ 的小量子上同调中的关系
\begin{equation}\label{eq:hirzebruch-two-small-quantum-relation}
H*(H-2f)=q_2(1-q_1).
\end{equation}
\looseness=-1 另一方面，式~\eqref{eq:hirzebruch-two-batyrev-quantum-relation} 表明，在 $\mathbb F_2$ 的 Batyrev 量子环中，$H\mathbin{\cdot}(H-2f)=z_2=q_2(1+q_1)$。因此，这两个环并不相同。

\looseness=-1 还应指出，计算 $\mathbb F_2$ 的小量子上同调有一种更初等的方法，因为 $\mathbb F_2$ 是 $\mathbb F_0=\mathbb P^1\times\mathbb P^1$ 的一个形变。再次用 $f$ 表示 $\mathbb F_0$ 的纤维，$s$ 表示由另一份 $\mathbb P^1$ 上的纤维化给出的截面。$\mathbb F_2$ 与 $\mathbb F_0$ 形变等价，所以由形变公理，它们的量子上同调相同。将它们的上同调等同，可见纤维类 $f$ 对应，且 $H=s+f$。$\mathbb P^1\times\mathbb P^1$ 为 Fano，可用 $\mathbb F_0$ 的 Batyrev 量子环计算其小量子上同调，得到关系
\[
s*s=q^f,\qquad f*f=q^s.
\]
再计算
\[
H*(H-2f)=(s+f)*(s-f)=s*s-f*f=q^f-q^s=q^f-q^{H-f}.
\]
因 $q^f=q_1^{f\cdot f}q_2^{H\cdot f}=q_2$，而 $q^{H-f}=q_1^{f\cdot(H-f)}q_2^{H\cdot(H-f)}=q_1q_2$，此计算与式~\eqref{eq:hirzebruch-two-small-quantum-relation} 一致。

\endgroup
\end{example}

% 第八批：原 PDF 物理页420–423中部。
\chapter{结语}\label{chap:conclusion}
\SourcePage{420}
现在已来到本书的结尾。一路上见识了一些精彩的数学，也希望已对镜像对称现象有所了解。在这最后一章，将回顾已完成的工作及尚待完成的工作，还会简要讨论前几章未涉及的镜像对称的众多方面中的若干主题。

\section{回顾与未解问题}\label{sec:reflections-open-problems}
研究镜像对称的过程中，用到了数量惊人的代数几何知识：讨论过环面簇、Hodge 结构变分、混合 Hodge 结构、具有正则奇点的微分方程、超几何方程、Frobenius 方法、Kähler 锥、种种模空间、辛几何、代数叠、计数几何、Clemens 猜想、等变上同调和局部化，以上还只是其中一部分。镜像对称把代数几何的如此多部分汇聚起来，表明我们共同从事的这项事业，其内在联系比我们所认识的更为深刻。

不过，镜像对称不仅是集中运用已有数学。在本书中，还见证了新对象进入代数几何研究，包括 Kähler 模、稳定映射、Gromov--Witten 不变量、量子上同调、A 模型 Hodge 结构变分、量子微分方程、Euler 数据以及上同调值形式函数。把它们与物理学中一些尚未有严密数学定义的思想结合起来，就能窥见二十一世纪代数几何可能带给我们的内容。

还要强调，镜像对称这一领域比本书所呈现的丰富得多。除了几乎未触及的庞大物理文献，还有许多数学方面完全没有涉及。第 \hyperref[sec:other-aspects-mirror-symmetry]{\ref*{sec:other-aspects-mirror-symmetry}} 节将列出其中一部分。不过，即使只限于前几章呈现的镜像对称观点，仍有很多工作尚待完成，下面就来说明。

\subsection{什么是镜像对称？}\label{subsec:what-is-mirror-symmetry}
基本问题在于，尚未提出一个完全令人满意的镜像定理版本，甚至尚未提出这样的镜像猜想。回顾迄今的工作，可辨认出三种不同的镜像对称形式：
\begin{itemize}
\item \hyperref[sec:mirror-symmetry-revisited]{第~\ref*{sec:mirror-symmetry-revisited} 节}探讨了 \SourceCite{Morrison7} 的建议：数学版本的镜像对称，应表述为一个光滑
\SourcePage{421}
Calabi--Yau 流形的 A 模型 Hodge 结构变分与其镜像的几何（即 B 模型）Hodge 结构变分之间的同构。
\item \hyperref[sec:quintic-mirror-theorem]{第~\ref*{sec:quintic-mirror-theorem} 节}研究了 \SourceCite{LLY} 的工作，它把镜像定理表述为 Euler 数据 $\widehat P$ 与 $\widehat Q$ 在镜像变换后的等价。这意味着某些超几何的上同调值函数之间存在关系，从中可计算五次三维簇的 Gromov--Witten 势。
\item \hyperref[sec:givental-mirror-approach]{第~\ref*{sec:givental-mirror-approach} 节}讨论了 \SourceCite{Givental2,Givental4} 的定理，其中镜像定理表述为上同调值函数 $I_{\mathcal V}$、$J_{\mathcal V}$ 在镜像变换后相等。所得等式给出了五次三维簇的 Gromov--Witten 势的一个等价描述。
\end{itemize}
\hyperref[sec:quintic-other-methods]{第~\ref*{sec:quintic-other-methods} 节}预览过这三种方法，它们在 $\mathbb P^4$ 中五次三维簇的情形下巧妙地汇合。不过，这些方法也有明显差别，例如：
\begin{itemize}
\item \hyperref[sec:mirror-symmetry-revisited]{第~\ref*{sec:mirror-symmetry-revisited} 节}的 Hodge 理论方法只对 Calabi--Yau 流形有意义。
\item \SourceCite{LLY} 的方法适用于多种情形，包括五次三维簇，以及\hyperref[subsec:multiple-cover-localization]{第~\ref*{subsec:multiple-cover-localization} 小节}研究的多重覆盖公式。但在本书写作时，这些结果的环境空间仍限于 $\mathbb P^n$。\footnote{我们刚获悉（1998 年 11 月），\SourceCite{LLY} 的作者即将宣布该文内容的一个环面版本。预印本发表后，应能得到环面镜像定理的详细证明。}
\item \SourceCite{Givental2,Givental4} 的定理适用于光滑环面簇 $X_\Sigma$ 中的 nef 完全交 $X$。$X=X_\Sigma$ 时也有非平凡结果，但这些方法例如不适用于多重覆盖公式。
\end{itemize}
后两种方法令人兴奋，因为终于能够实际证明镜像对称中的一些漂亮公式。回顾\hyperref[chap:mirror-theorem]{第十一章}镜像定理的证明中如何使用等变上同调及局部化，就能看出其局限。这些方法需要群作用，这解释了环面簇为何起着如此突出的作用。不过，结果不仅适用于 Calabi--Yau 情形。因此，\SourceCite{LLY} 与 \SourceCite{Givental2} 的方法可能属于一个更广泛的“环面”版本的镜像对称，Calabi--Yau 环面完全交则是其特例。

另一方面，\hyperref[sec:mirror-symmetry-revisited]{第~\ref*{sec:mirror-symmetry-revisited} 节}的 Hodge 理论版本与环面簇无关。任意射影 Calabi--Yau 流形 $V$ 都有 A 模型 Hodge 结构变分，可以问：它是否与某个 Calabi--Yau 流形的几何 Hodge 结构变分相同？若是，后者便是镜像 $V^\circ$。此外，如例题~\ref{ex:octic-toric-mirror-theorem} 后的评论所示，Hodge 理论版本的镜像对称与 Givental 版本的镜像定理之间的关系，仍有待弄清。最后，如第 \hyperref[subsec:homological-categorical-mirror]{\ref*{subsec:homological-categorical-mirror}} 小节将展示的，有些镜像对称版本涉及的内容远远超出前几章。所有这些都表明，距离完全理解镜像对称的数学，我们仍很遥远。

最后，提醒读者：按照物理学，镜像对称是两个物理理论的等价，而这些理论目前尚未得到严密的数学定义。
\SourcePage{422}
随着这些理论被置于更牢固的数学基础之上，将能更深入地理解镜像对称的数学。这可能填补刚才提到的一些空缺，也可能把我们引向完全出乎意料的方向。

\subsection{尚未解决的猜想}\label{subsec:open-conjectures}
如刚解释的，镜像对称中主要的未解问题是提出并证明一般的镜像定理。与此有关，现有镜像定理版本也有一些有趣问题。例如，应更好地理解\hyperref[subsec:givental-toric-mirror]{第~\ref*{subsec:givental-toric-mirror} 小节} Givental 的环面镜像定理版本，与\hyperref[subsec:toric-mirror-conjecture]{第~\ref*{subsec:toric-mirror-conjecture} 小节}的环面镜像猜想之间的联系。例题~\ref{ex:octic-toric-mirror-theorem} 结合此结果与式~\eqref{eq:octic-virtual-class-pushforward-assumption}，导出了式~\eqref{eq:octic-toric-mirror-potential} 和~\eqref{eq:octic-yukawa-period-second-derivative}。我们很希望知道，这些公式可以推广到何种程度。

不过，镜像定理提出的问题并非全部。读者还应记得：为了陈述本书所考虑的镜像对称版本，在建立所需工具的途中，也遇到了一些有趣的未解问题。现在回顾其中一些猜想，希望鼓励人们研究它们。

先列出\hyperref[chap:hodge-yukawa]{第五章}和\hyperref[chap:moduli]{第六章}讨论的、关于复模与 Kähler 模的四个猜想：
\begin{itemize}
\item \hyperref[subsec:higher-dimensional-mum]{第~\ref*{subsec:higher-dimensional-mum} 小节}的\emph{整性猜想}描述了最大幺幂边界点处，权滤过的整数结构应如何与单值性相互作用。
\item \hyperref[subsec:toric-complex-moduli]{第~\ref*{subsec:toric-complex-moduli} 小节}在 Calabi--Yau 环面超曲面的简化模空间上构造了某些特殊边界点。猜想~\ref{conj:toric-mum-boundary} 断言，它们是最大幺幂边界点。
\item \hyperref[subsec:large-radius-limits]{第~\ref*{subsec:large-radius-limits} 小节}的\emph{锥猜想}断言，维数大于二的 Calabi--Yau 流形 $V$ 的 Kähler 锥，模去 $\operatorname{Aut}(V)$ 的作用后应为有理多面锥。
\item \hyperref[subsec:enlarged-kahler-moduli]{第~\ref*{subsec:enlarged-kahler-moduli} 小节}的猜想~\ref{conj:toric-kahler-cone} 描述 Calabi--Yau 环面超曲面的 Kähler 锥的环面部分。
\end{itemize}
这些猜想自身都很有趣，同时也是镜像对称的重要组成部分。

除了这些模空间猜想，还有许多与镜像对称有关的未解问题值得代数几何学家思考。以下是我们觉得特别有趣的一些：
\begin{itemize}
\item \hyperref[subsec:toric-gkz]{第~\ref*{subsec:toric-gkz} 小节}研究 GKZ 方程时，看到这些方程给出了部分 Picard--Fuchs 方程，却并非全部。应当能够描述缺少的 Picard--Fuchs 方程。此外，上同调值形式函数 $I_{\mathcal V}$ 应满足镜像的 Picard--Fuchs 方程。
\item \hyperref[subsec:normalized-yukawa-connection]{第~\ref*{subsec:normalized-yukawa-connection} 小节}讨论了最大幺幂边界点处 Calabi--Yau 三维簇的归一化 Yukawa 耦合与势函数。当三维簇为环面完全交时，应能像例题~\ref{ex:octic-toric-mirror-theorem} 那样，用适当的 $I_{\mathcal V}$ 计算 Yukawa 耦合和势函数。
\item 关于虚拟基本类 $[\overline M_{g,n}(X,\beta)]^{\mathrm{virt}}$，有一些问题值得进一步研究。\hyperref[subsec:gw-equivalence]{第~\ref*{subsec:gw-equivalence} 小节}讨论了 Gromov--Witten 不变量与虚拟基本类的各种定义之间的等价性。另一个问题是，当 $X$ 为 $Y$ 中的完全交时，$[\overline M_{g,n}(X,\beta)]^{\mathrm{virt}}$ 与
\SourcePage{423}
$[\overline M_{g,n}(Y,\beta)]^{\mathrm{virt}}$ 有何联系。式~\eqref{eq:octic-virtual-class-pushforward-assumption} 表明，这对理解 Givental 函数 $J_X$ 与 $J_{\mathcal V}$ 的关系很重要。
\item \IndexLocation{indexloc:1131:423}量子上同调的定义涉及无穷和，\hyperref[subsec:small-quantum-definition]{第~\ref*{subsec:small-quantum-definition} 小节}猜测，此和在大半径极限点的某个邻域中收敛。注意，只有此猜想成立，A 模型 Hodge 结构变分才有意义。
\end{itemize}
如前，即使不考虑与镜像对称的联系，我们也认为这些问题很有趣。

本书中一些最重要的未解问题涉及计数几何。起初，\SourceCite{CdGP} 对五次三维簇的预言近乎神奇：物理学似乎毫不费力地解决了远超代数几何标准方法能力的计数问题。此后，镜像对称已成为计算 Gromov--Witten 不变量的有力工具，但理解一个 Gromov--Witten 不变量的计数含义，仍是非平凡的问题。即使对于五次三维簇 $V$，瞬子数 $n_{10}$ 与 $V$ 上十次有理曲线的关系，也可能很复杂，因为十次情形的 Clemens 猜想尚未获证。即使获证，$n_{10}$ 也不是这类曲线的数目，如\hyperref[subsec:instanton-numbers]{第~\ref*{subsec:instanton-numbers} 小节}所解释。另一个很好的问题是猜想~\ref{conj:symplectic-instantons}，它试图直接定义瞬子数，而不是像本书那样，间接地通过 Gromov--Witten 不变量 $\langle I_{0,0,\beta}\rangle$ 定义它们。不过，瞬子数的问题只是全貌的一部分。如\hyperref[subsec:instanton-numbers]{第~\ref*{subsec:instanton-numbers} 小节}和\hyperref[sec:computing-gw-ii]{第~\ref*{sec:computing-gw-ii} 节}的例题及参考文献所示，镜像对称使计数几何论文数量激增，其中许多与 Calabi--Yau 流形或 Gromov--Witten 不变量无关。

如前言所说，本书应视为对镜像对称这一仍在进行中的研究的报告。这个领域的某些部分已有牢固基础，其他部分仍在迅速变化。镜像对称的权威性著作尚有待写出。不过，希望本书所呈现的虽有限的镜像对称观点，已经传达了代数几何这一惊人领域的一些奇妙与激动人心之处。

% 第九批从原 PDF 423 中部接续；该页上部锚点已在第八批。
\hypertarget{source:423-batch9}{}
\section{镜像对称的其他方面}\label{sec:other-aspects-mirror-symmetry}
主要目标现已完成，但还有许多镜像对称的方面，本书未曾涉及。在本书写作时，镜像对称仍是活跃的研究领域，许多新思想正不断发展，并由此带来数学与物理学中的其他新进展。这里甚至不打算给出完整清单，只举几个例子。下面讨论的主题，有些性质纯属数学，有些则与物理学联系更密切，但它们都包含值得关注的数学内容。

\subsection{K3 曲面上的有理曲线}\label{subsec:rational-curves-k3}
固定一个一般的 K3 曲面，其中有除子类 $C$ 满足 $C^2=2g-2$。于是，$|C|$ 是一个 $g$ 维线性系，其中的曲线具有算术亏格 $g$。令 $n(g)$ 为这个线性系中有理曲线的数目。
\SourcePage{424}
Yau 与 Zaslow \SourceCite{YZ} 提出了如下惊人的公式猜想：
\[
1+\sum_{g=1}^{\infty}n(g)q^g=q/\Delta(q),
\]
其中 $\Delta(q)$ 是通常的权为 $12$ 的模形式。最早的数学阐述见 \SourceCite{Beauville2}，其中也证明了若干中间步骤。基本想法是：对这个线性系中出现的每条有理曲线，赋予一个重数，它等于该曲线的紧化 Jacobian 的 Euler 示性数（同一个 Euler 示性数也出现在 \SourceCite{YZ} 中）。这个重数可以与稳定映射联系起来 \SourceCite{FGvS}。此外，\SourceCite{FGvS} 还证明，此重数总为正数。因此，Yau--Zaslow 公式有明确的计数含义，即按适当重数计数有理曲线。

\SourceCite{Gottsche} 把 Yau--Zaslow 公式推广为一个关于代数曲面上结点曲线的猜想，提出一个生成函数，用以计数曲线的 $n$ 维族中具有 $n$ 个结点的曲线。当曲面是 K3 曲面或 Abel 曲面时，猜想中的生成函数可用拟模形式表示。对 K3 曲面，这个猜想已在 \SourceCite{BL} 中得到证明；下面稍加说明。考虑满足 $C^2=2g+2n-2$ 的曲线线性系 $|C|$。它的维数为 $g+n$，其中的曲线具有算术亏格 $g+n$。预期 $|C|$ 中具有 $n$ 个结点的曲线构成一个 $g$ 维子族。令 $N_g(n)$ 表示这个线性系中经过 $g$ 个一般点、几何亏格为 $g$ 且有 $n$ 个结点的曲线的数目。\SourceCite{BL} 证明了
\[
\sum_{n=0}^{\infty}N_g(n)q^n
=\left(\frac{d}{dq}G_2(q)\right)^g\frac{q}{\Delta(q)},
\]
其中 $G_2(q)$ 是 Eisenstein 级数
\[
G_2(q)=-\frac1{24}+\sum_{k=1}^{\infty}\sigma(k)q^k,
\]
$\sigma(k)$ 是算术函数 $\sum_{d\mid k}d$。注意，$G_2$ 及其导数都是拟模函数。
\index{K3曲面@K3 曲面!有理曲线|hyperpage}\index{Yau--Zaslow公式@Yau--Zaslow 公式|hyperpage}\index{拟模形式|hyperpage}

\subsection{K3 曲面的镜像对称与 Arnold 奇异对偶性}\label{subsec:k3-arnold-strange-duality}
二十世纪七十年代末，不同研究者指出，K3 曲面的一种对偶性可以解释例外单模临界点的 Arnold 奇异对偶性。近年人们认识到，K3 曲面的这种对偶性属于二维版本的镜像对称 \SourceCite{Dolgachev3,Kobayashi,Ebeling,GN1,GN3}。

K3 曲面的镜像对称有不同的特点。对维数大于二的 Calabi--Yau 流形 $V$，Kähler 锥在 $H^2(V,\mathbb R)$ 中是开集；对 K3 曲面 $S$，情形便不再如此，而且 Kähler 锥与格 $H^2(S,\mathbb Z)$ 之间的关系相当微妙。因此，格论在 K3 曲面的镜像对称中起重要作用，并不令人意外。\hyperref[sec:voisin-borcea]{第~\ref*{sec:voisin-borcea} 节}给出的 Voisin--Borcea 构造，已让我们见到这一点的一些迹象。

也可以在 K3 曲面 $S$ 上研究弦理论，\SourceCite{AM3} 对此作了说明。此时仍有 SCFT 模空间，但由于 $h^{2,0}(S)\ne0$，复结构与 Kähler 结构的形变不再解耦，因此不再得到
\SourcePage{425}
式~\eqref{eq:scft-foliations} 所示的局部乘积结构。这一主题内容广泛，无法在几个段落中概述。关于更多细节及相关文献，请参阅 \SourceCite{AM3,Aspinwall}。
\index{Arnold奇异对偶性@Arnold 奇异对偶性|hyperpage}\index{K3曲面@K3 曲面!镜像对称|hyperpage}

\subsection{McKay 对应}\label{subsec:mckay-correspondence}
若 $G\subset\operatorname{SL}(n,\mathbb C)$ 是有限群，且 $\mathbb C^n/G$ 有一个满足 $K_Y=0$ 的消解 $Y$，则 \emph{McKay 对应} \SourceCite{Reid6} 断言，$G$ 的不可约表示应给出 $H^*(Y,\mathbb Z)$ 的一组基，使 $G$ 的特征标表对应于 $Y$ 上的 Poincaré 对偶性。此猜想对物理学中的 Euler 数 \SourceCite{Vafa1,Roan2}，以及 Batyrev 与 Dais 的弦论 Hodge 数 \SourceCite{BD}，都有相应推论。
\index{McKay对应@McKay 对应|hyperpage}

\subsection{超 Kähler 流形}\label{subsec:hyperkahler-manifolds}
若紧致复流形 $M$ 有一个闭的、非退化的全纯二形式，就称 $M$ 是\emph{全纯辛}的。由 Calabi--Yau 定理，$M$ 有一个度量及一个关于 Riemann 联络平行的四元数作用，因此 $M$ 是超 Kähler 流形。\SourceCite{Verbitsky} 中，Verbitsky 证明，超 Kähler 流形是它自身的镜像。
\index{全纯辛流形|hyperpage}\index{超Kähler流形@超 Kähler 流形|hyperpage}

\subsection{D 膜与 T 对偶}\label{subsec:d-branes-t-duality}
本书将镜像对称理解为与镜像对 $(V,\omega)$ 和 $(V^\circ,\omega^\circ)$ 相联系的两个 SCFT 相等（见\hyperref[chap:introduction]{第一章}的讨论）。SCFT 中的态称为\emph{微扰态}，其中包括与 $V$ 和 $V^\circ$ 上的上同调类对应的态。近年人们认识到，这个理论在某种意义上可以扩充，以包含\emph{非微扰态}，包括与 \emph{D 膜}相联系的态。D 膜本质上是一个 $p$ 维物体在时空中传播时扫出的\emph{世界体}。为明确维数，有时将它们称为 D$p$ 膜；由于多了一个时间维，它们的世界体是 $p+1$ 维的。例如，D1 膜的样子与弦的世界面完全相同（如\hyperref[chap:introduction]{第一章}所绘）。态本身来自边界落在 D 膜上的开弦。D 膜的综述见 \SourceCite{Polchinski1}。人们猜测，与 $(V,\omega)$ 和 $(V^\circ,\omega^\circ)$ 相联系的物理理论，\emph{连同非微扰态在内}都相等，从而使镜像对称成为更强的断言。

\SourceCite{SYZ} 中，Strominger、Yau 与 Zaslow 提出，镜像对称就是物理学家所说的 \emph{T 对偶}。下面简述这个想法。注意，$V^\circ$ 上点的模空间与 $V^\circ$ 本身同构，而它又自然等同于 $V^\circ$ 上 D0 膜模空间的一部分。为理解这一点，回顾\hyperref[chap:introduction]{第一章}：弦理论是十维理论，其中 $V^\circ$ 贡献六个维度；时间维是余下四个维度之一。对任一点 $p\in V^\circ$，当 $p$ 随时间传播时，便扫出一条 D0 膜的世界线。镜像对称预言，它应对应于 $V$ 上 D3 膜模空间的一部分。进一步分析表明，Calabi--Yau 流形 $V$ 的镜像 $V^\circ$，应是 $V$ 上某些特殊 Lagrange 环面的模空间，经过适当复化与紧化得到的空间。如 \SourceCite{Morrison8} 所解释，T 对偶或许最终可以这样给出一个纯数学的镜像对称构造。在这个描述中，$V$ 与 $V^\circ$ 都映到同一个实三维底空间，一般纤维是实三维环面。人们猜测，$V$ 与 $V^\circ$ 的这两个纤维化中的对应纤维互为对偶环面，“T 对偶”之名也由此而来。这可能导出镜像流形的拓扑构造。上述断言的某些方面，已在 Voisin--Borcea 三维簇 \SourceCite{GW} 和 K3 曲面
\SourcePage{426}
\SourceCite{GW,Morrison8} 的情形下得到检验。对椭圆 K3 曲面，T 对偶与 Fourier--Mukai 变换有关 \SourceCite{BBRP}。\SourceCite{Gross3} 中还提出了若干值得关注的拓扑推论猜想。
\index{D膜@D 膜|hyperpage}\index{T对偶@T 对偶|hyperpage}\index{世界体|hyperpage}\index{特殊Lagrange环面@特殊 Lagrange 环面|hyperpage}

\subsection{同调镜像对称或范畴镜像对称}\label{subsec:homological-categorical-mirror}
1994 年，Kontsevich \SourceCite{Kontsevich3} 提出了镜像对称的“同调”版本，将镜像猜想表述为范畴等价。此时，$V^\circ$ 是 $V$ 的镜像，意味着 $V^\circ$ 上凝聚层的导出范畴，应等价于 $V$ 上带酉局部系统的 Lagrange 子流形的导出范畴。由 Lagrange 子流形构造一个范畴的基本想法来自 \SourceCite{Fukaya}。\SourceCite{PZ} 建议，或许应限制到 $V$ 的特殊 Lagrange 子流形。注意，Lagrange 子流形使用 $V$ 的 Kähler 结构（A 模型），凝聚层则使用 $V^\circ$ 的复结构（B 模型）。

\SourceCite{Kontsevich3} 解释，与 $V$ 相联系的 Lagrange 范畴的某些形变，应当使我们能够将 $V$ 的量子上同调构造为通常杯积的形变。这里的形变自然地在整个 $H^*(V)$ 上进行，而非像\hyperref[sec:a-model-vhs]{第~\ref*{sec:a-model-vhs} 节} A 模型联络的定义那样，只在 $H^{1,1}(V)$ 上进行。由于它应与 $V^\circ$ 上凝聚层的形变理论有关，这暗示着一种超出通常复结构形变的丰富结构。下面的公式可以让我们窥见其中可能涉及的内容：
\begin{equation}\label{eq:hochschild-polyvector-cohomology}
HH^n(V^\circ)=\bigoplus_{p+q=n}H^q(V^\circ,\wedge^p\Theta_{V^\circ}),
\end{equation}
其中 $HH$ 表示 Hochschild 上同调。这个上同调群应与 $V^\circ$ 的复模空间的某个加厚空间的切空间等同。注意，式~\eqref{eq:hochschild-polyvector-cohomology} 的各直和项，恰好是式~\eqref{eq:eigenspaces} 第一行所写 $N=2$ 超共形代数中算子 $(Q,\overline Q)$ 的 $(p,q)$ 特征子空间。以式~\eqref{eq:hochschild-polyvector-cohomology} 为切空间来加厚模空间的想法，最早出现在 \SourceCite{Witten4}。这些空间能够在数学中自然出现，是一件很美妙的事。如 \SourceCite{Kontsevich3} 所解释，$HH^2(V^\circ)$ 包含 $V^\circ$ 的复结构的通常一阶形变、代数层 $\mathcal O_{V^\circ}$ 的某些非交换形变，以及从形变观点看仍颇为神秘的一些项。由于维数至少为三的 Calabi--Yau 流形满足 $H^2(V^\circ,\mathcal O_{V^\circ})=0$ 和 $H^0(V^\circ,\wedge^2\Theta_{V^\circ})=0$，最后这两项在此时不成问题；不过，$n\ne2$ 时 $HH^n(V^\circ)$ 的含义仍有些神秘。此外，\SourceCite{Kontsevich3} 指出，猜想 $V$ 的 Lagrange 范畴具有 $A_\infty$ 结构。（不过，$A_\infty$ 范畴并非通常意义下的范畴，其复合不再严格满足结合律，因而须谨慎处理。）这再次暗示，在镜像一侧，$V^\circ$ 上凝聚层的范畴应有更丰富的结构。这些都表明，我们尚未发现镜像对称的全部内容。

\SourceCite{PZ} 最近给出了椭圆曲线情形下这种范畴等价的一个例子。该文还指出，T 对偶可能是同调镜像对称的推论。在这个背景下，给定一点 $p\in V^\circ$，子簇 $\{p\}\subset V^\circ$ 的凝聚层 $\mathcal O_{\{p\}}$，应对应于 $V$ 的某个特殊 Lagrange 子流形。这是
\SourcePage{427}
\hyperref[subsec:d-branes-t-duality]{第~\ref*{subsec:d-branes-t-duality} 小节}所述思想的一个版本：镜像 $V^\circ$ 是 $V$ 中某些实 Lagrange 环面的模空间。
\index{同调镜像对称|hyperpage}\index{范畴镜像对称|hyperpage}\index{Hochschild上同调@Hochschild 上同调|hyperpage}\index{A无穷范畴@$A_\infty$ 范畴|hyperpage}

\subsection{K3 纤维化与椭圆纤维化的 Calabi--Yau 流形}\label{subsec:k3-elliptic-fibered-cy}
具有 K3 纤维化或椭圆纤维化的 Calabi--Yau 流形，是大量数学与物理论文的研究对象。从数学上说，对这些纤维化感兴趣相当自然。在研究 Calabi--Yau 流形上的线性系时，K3 纤维化与椭圆纤维化起特殊作用，因为它们自然地出现在 Calabi--Yau 流形的对数极小模型纲领中。在通常的极小模型纲领中，人们试图把簇理解为纤维是极小模型的纤维化。对 Calabi--Yau 簇的对数极小模型，相应情形是纤维为较低维的 Calabi--Yau 簇。\SourceCite{Oguiso} 已开始对这些纤维化进行分类。

人们也研究了 K3 纤维化的镜像对称，猜测纤维的镜像应影响总空间的镜像。\SourceCite{GN2} 讨论了这个想法，\SourceCite{AKMS,HLY3} 中可找到这类空间的例子。

K3 纤维化与椭圆纤维化也出现在物理学中，但原因不同。推动研究 K3 纤维化的主要动力是\emph{对偶性}。更具体地说，人们相信，与 Calabi--Yau 三维簇相联系的 \emph{II 型}弦理论，对偶于与 K3 曲面和椭圆曲线的乘积相联系的杂化弦理论。这意味着，各理论的参数取某些值时，两个理论成为同一个理论。K3 曲面上的向量丛，是定义这类特定杂化弦理论所需数据的一部分。对这种对偶性，已有大量数学检验。

椭圆纤维化的 Calabi--Yau 三维簇，在物理学中是可以用来构造 \emph{F 理论}模型的空间 \SourceCite{Vafa3}。由于丛在 F 理论中也起关键作用，这促使人们进一步研究具有椭圆纤维化的 Calabi--Yau 流形上的主丛 \SourceCite{FMW1}。关于这些丛的数学研究，见 \SourceCite{Donagi,FMW2}。与此密切相关的是 F 理论和杂化弦理论之间的对偶性，它出现在由同时具有 K3 纤维化和椭圆纤维化的 Calabi--Yau 三维簇得到的理论中。这种对偶性可以视为低一维情形下类似结果的推广。与椭圆纤维化的 K3 曲面相联系的 F 理论，和椭圆曲线上的杂化弦理论之间的对偶性，有一个优美的数学支持：椭圆曲线上丛的模，可与某些类型的椭圆纤维化 K3 曲面的 Weierstrass 模型的模等同。
\index{K3纤维化@K3 纤维化|hyperpage}\index{椭圆纤维化|hyperpage}\index{杂化弦理论|hyperpage}\index{F理论@F 理论|hyperpage}

\subsection{锥形奇点跃迁}\label{subsec:conifold-transitions-conclusion}
当 Calabi--Yau 三维簇出现结点奇点，并有一个仍为 Calabi--Yau 的消解时，就发生了\emph{锥形奇点跃迁}。由杂化弦得到的 SCFT 会在这一跃迁过程中变得奇异。不过，\SourceCite{GMS} 表明，\emph{II 型}弦理论可以避免锥形奇点跃迁处的奇异性。粗略地说，与原流形中的一个消失圈（把它看作一个三维膜\footnote{\hyperref[subsec:d-branes-t-duality]{第~\ref*{subsec:d-branes-t-duality} 小节}所考虑的 D3 膜，是三维膜的一种特殊类型。直到 \SourceCite{GMS} 写成之后，D 膜才得到普遍接受。现在人们一致认为，\SourceCite{GMS} 中的三维膜实际上就是 D3 膜。}）相联系的带电黑洞，在 Calabi--Yau 簇变得奇异时质量变为零；这一奇异性补偿了前面所说理论中的奇异性。随后，这些黑洞在消解上成为基本粒子。若假定 Reid 幻想 \SourceCite{Reid5} 的一个 Kähler 版本成立，
\SourcePage{428}
便可以连接任意两个 Calabi--Yau 三维簇所产生的物理理论。这已在作为加权射影空间超曲面的 Calabi--Yau 三维簇情形下得到验证 \SourceCite{ACJM,CGGK}，尽管其奇点可能不是结点。\hyperref[subsec:kahler-moduli-boundary]{第~\ref*{subsec:kahler-moduli-boundary} 小节}讨论过锥形奇点跃迁与镜像对称有关的某些方面。

\SourceCite{Morrison9} 猜测，锥形奇点跃迁的镜像也是锥形奇点跃迁；更一般地，极端跃迁的镜像也是极端跃迁。\SourceCite{BCFKvS1,BCFKvS2} 已在一些例子中检验了这一点。

\subsection{旗流形中的 Calabi--Yau 完全交}\label{subsec:flag-cy-complete-intersections}
这里的关键观察来自 \SourceCite{Sturmfels2}：Grassmann 簇 $G(k,n)$ 可以平坦形变到一个射影环面簇 $P(k,n)$。此外，根据 \SourceCite{BCFKvS2}，$P(k,n)$ 是 Fano 簇，即它来自一个自反多胞形。利用这一点，\SourceCite{BCFKvS1} 的作者将\hyperref[chap:mirror-constructions]{第四章}的 Batyrev--Borisov 镜像构造，与\hyperref[subsec:kahler-moduli-boundary]{第~\ref*{subsec:kahler-moduli-boundary} 小节}及\hyperref[subsec:conifold-transitions-conclusion]{第~\ref*{subsec:conifold-transitions-conclusion} 小节}讨论的极端跃迁结合起来，给出 $G(k,n)$ 中 Calabi--Yau 完全交 $V$ 的镜像构造。

基本想法如下。设 $V\subset G(k,n)$ 是次数为 $d_1,\ldots,d_r$ 的超曲面的完全交。将 $G(k,n)$ 退化到 $P(k,n)$ 时，$V$ 退化到环面完全交 $\overline V\subset P(k,n)$。利用光滑的保典范消解 $\widehat P(k,n)\to P(k,n)$，可得到一个消解 $V'\to\overline V$。若将它理解为先从 $V'$ 经收缩到 $\overline V$，再从 $\overline V$ 经光滑化到 $V$，就得到了\hyperref[subsec:kahler-moduli-boundary]{第~\ref*{subsec:kahler-moduli-boundary} 小节}讨论的一种极端收缩（此时，由于 $\overline V$ 的奇点尤其温和，这一跃迁是锥形奇点跃迁）。如该小节所述，镜像对称暗示应当有一个“镜像”极端跃迁：$V'$ 的镜像 $V'^\circ$ 退化到 $\overline V^\circ$，然后消解为 $V$ 的镜像 $V^\circ$。

在此情形中，$V'\subset\widehat P(k,n)$ 的镜像 $V'^\circ$，可由\hyperref[sec:toric-complete-intersections]{第~\ref*{sec:toric-complete-intersections} 节}所述的 Batyrev--Borisov 镜像构造得到。通过求出相应的退化与光滑化，\SourceCite{BCFKvS1} 能够明确描述 Grassmann 簇中 Calabi--Yau \IndexLocation{indexloc:0802:428}完全交的镜像构造。

为利用这个构造计算 Gromov--Witten 不变量，\SourceCite{BCFKvS1} 引入了\emph{阶乘技巧}；该文将它描述为\hyperref[subsec:quantum-hyperplane-section-principle]{第~\ref*{subsec:quantum-hyperplane-section-principle} 小节}的\IndexLocation{indexloc:1159:428}量子超平面截面原理的一个朴素版本。\SourceCite{BCFKvS1} 使用的另一工具，是量子上同调与量子微分方程之间的关系，我们已在\hyperref[sec:quantum-cohomology-relations]{第~\ref*{sec:quantum-cohomology-relations} 节}讨论过。由于 Grassmann 簇的量子上同调已知，就能获得关于量子微分方程的大量信息（\SourceCite{BCFKvS1} 将其称为\emph{量子 $\mathcal D$ 模}）。再利用\hyperref[subsec:givental-toric-mirror]{第~\ref*{subsec:givental-toric-mirror} 小节}的\IndexLocation{indexloc:0862:428}环面镜像定理和 \SourceCite{Kim2} 中给出的推广，便可计算 Grassmann 簇中许多三维 Calabi--Yau 完全交的瞬子数。

上述工作都可在更一般的情形下进行。令 $F(n_1,\ldots,n_\ell,n)$ 表示由所有满足
\[
0\subset W_1\subset\cdots\subset W_\ell\subset\mathbb C^n,
\qquad \dim W_i=n_i
\]
的旗组成的部分旗流形。\SourceCite{GL} 证明，$F(n_1,\ldots,n_\ell,n)$ 退化到一个环面簇 $P(n_1,\ldots,n_\ell,n)$，推广了上面提到的 Sturmfels 的结果。由于 \SourceCite{BCFKvS2} 表明 $P(n_1,\ldots,n_\ell,n)$ 是 Fano 簇，以上构造便可应用。还应提及 \SourceCite{Batyrev5}，它回顾了刚才描述的构造，并考察一般情形下 Fano 簇退化到 Fano 环面簇时会发生什么。
\index{旗流形!Calabi--Yau完全交@Calabi--Yau 完全交|hyperpage}\index{阶乘技巧|hyperpage}\index{量子D模@量子 $\mathcal D$ 模|hyperpage}

\SourcePage{429}
\subsection{物理学中的其他主题}\label{subsec:other-physics-topics}
镜像对称的物理学近来有一长串新进展，是前几章没有讨论的。其中一些已在本节提到，但还有许多。例如，物理学论证使人们预期：当 Calabi--Yau 流形中的某些曲线被收缩时，可将超对称场论作为弦理论的一个极限得到。因此，Calabi--Yau 几何提供了分析这些超对称理论的工具，尤其是分析其规范理论方面的工具。这些场论有可利用镜像对称求解的“周期”，从而例如可以导出 Seiberg--Witten 曲线及其微分，而无需假定任何对偶猜想。这方面的镜像对称应用见 \SourceCite{KKV,KMV}；利用 Calabi--Yau 几何构造量子场论的其他情形见 \SourceCite{KV,BJPSV,BV}。这样的理论称为\emph{几何工程理论}，其入门介绍见 \SourceCite{Mayr}。在数学上，几何工程理论明确预言了\hyperref[subsec:kahler-moduli-boundary]{第~\ref*{subsec:kahler-moduli-boundary} 小节}所讨论的极端跃迁中光滑化分支的维数。最简单的例子是锥形奇点跃迁，此时 \SourceCite{GMS} 中来自物理学的预言与数学相符；这一情形的光滑化维数，可根据 \SourceCite{Clemens} 作数学计算。在本书写作时，物理学所给一般预言的正确性尚无数学证明。\SourceCite{KMP} 已在一些例子中检验了这个预言。
\index{几何工程理论|hyperpage}\index{Seiberg--Witten曲线@Seiberg--Witten 曲线|hyperpage}

\appendix
\chapter{奇异簇}\label{app:singular-varieties}
\SourcePage{430}
虽然复流形是镜像对称的重要组成部分，但在许多情形中，奇异簇也会自然出现。例如，\hyperref[chap:mirror-constructions]{第四章}的 Batyrev 构造，处理的是单纯环面簇中可能有奇点的超曲面。本附录回顾本书将遇到的奇点类型。

本书考虑的簇都是 Cohen--Macaulay 簇，即它们的所有局部环都是 Cohen--Macaulay 环。Cohen--Macaulay 簇 $X$ 的关键性质是，它有一个对偶化层，通常记为 $\omega_X$。当 $X$ 光滑时，$\omega_X$ 就是通常的层 $\Omega_X^d$，其中 $d=\dim(X)$；因此，我们常把一般 Cohen--Macaulay 簇的对偶化层写作 $\widehat\Omega_X^d$。Cohen--Macaulay 簇的良好性质在于，它们在对偶理论中表现良好；详细讨论见 \SourceCite{Oda}。$\widehat\Omega_X^d=\omega_X$ 的另一种记法是 $\mathcal O_X(K_X)$，其中 $K_X$ 是典范除子。注意，一般而言，$K_X$ 只是 Weil 除子。

由标准交换代数可知，对 Cohen--Macaulay 簇 $X$，下列条件等价：对偶化层 $\omega_X$ 是线丛；$K_X$ 是 Cartier 除子；所有局部环都是 Gorenstein 环。此时称 $X$ 是 \emph{Gorenstein 簇}。
\index{Cohen--Macaulay簇@Cohen--Macaulay 簇|hyperpage}\index{对偶化层|hyperpage}\index{Gorenstein簇@Gorenstein 簇|hyperpage}

\section{典范奇点与终端奇点}\label{sec:canonical-terminal-singularities}
若正规簇 $X$ 的某一点附近，存在正整数 $r$ 使 $rK_X$ 为 Cartier 除子，并存在一个局部奇点消解 $f:Y\to X$，使得
\begin{equation}\label{eq:canonical-terminal-discrepancies}
rK_Y=f^*(rK_X)+\sum_i a_iE_i,
\end{equation}
其中求和遍及 $f$ 的全部例外除子 $E_i$，且所有 $a_i\geq0$，就称该点是\emph{典范奇点}。若进一步所有 $a_i>0$，就称该奇点是\emph{终端奇点}。若 $X$ 是 Gorenstein 簇，则上式中可取 $r=1$。

这些奇点类型自然地出现在三维簇和更高维簇的典范模型与极小模型研究中。关于奇点本身的更多背景，见 \SourceCite{Reid1,Reid4}。对于 Calabi--Yau 簇，它们的重要性很明显，因为典范类在 Calabi--Yau 的定义中起突出作用。例如，定义~\ref{def:calabi-yau-variety} 中的极小 Calabi--Yau 簇 $V$，具有 Gorenstein、$\mathbb Q$ 析因的终端奇点。因此，若 $V$ 不光滑，则它的任何消解 $f:W\to V$ 都不再是 Calabi--Yau 簇。这是因为 $V$ 的奇点是终端奇点，所以
\[
K_W=f^*K_V+\sum_i a_iE_i
\]
不平凡：$K_V=0$，而 $a_i>0$。

这些奇点的另一个来源是环面簇。\SourceCite{Reid1} 证明，Gorenstein 环面簇至多具有典范奇点。下面将看到，Gorenstein 轨形也有同样的性质。
\index{典范奇点|hyperpage}\index{终端奇点|hyperpage}

\SourcePage{431}
\section{轨形}\label{sec:orbifolds}
先回顾一个簇是轨形的含义。
\begin{definition}\label{def:orbifold-local-chart}
若 $d$ 维簇 $X$ 的每一点 $p\in X$ 都有一个邻域，在解析意义下等价于 $U/G$ 中原点的一个邻域，其中 $G\subset\operatorname{GL}(d,\mathbb C)$ 是除恒等元外不含复反射的有限子群，$U\subset\mathbb C^d$ 是原点的一个 $G$ 不变邻域，则称 $X$ 为\emph{轨形}。
\end{definition}

在这个定义中，\emph{复反射}是 $\operatorname{GL}(d,\mathbb C)$ 的一个有限阶元素，其 $d-1$ 个特征值等于 $1$。定义~\ref{def:orbifold-local-chart} 中的群 $G$ 称为 $\operatorname{GL}(d,\mathbb C)$ 的一个\emph{小子群}，$(U/G,0)$ 称为 $X$ 在 $p$ 处的\emph{局部坐标片}。由 \SourceCite{Prill} 的一个定理，$G$ 在共轭意义下唯一。

“轨形”与“$V$ 流形”这两个名称可以互换，也就是拟光滑的含义。对于环面簇，这些条件等价于其扇是单纯的，此时群 $G$ 是 Abel 群。因此，\IndexLocation{indexloc:1375:431}单纯环面簇至多有有限 Abel 群商奇点。

一个重要直观是：“在 $\mathbb Q$ 上”，轨形的表现如同流形。例如，紧致轨形的奇异上同调满足 Poincaré 对偶性，但须使用 $\mathbb Q$ 系数。同样，对于完备环面簇，$X$ 光滑时 $H^*(X,\mathbb Z)$ 的组合描述，在 $X$ 单纯时仍适用于 $H^*(X,\mathbb Q)$。

接着研究 Gorenstein 轨形。众所周知，每个 Gorenstein 轨形至多具有典范奇点 \SourceCite{Reid1}。不过，还有下面这个广为流传的结果，表明在终端奇点情形下，Gorenstein 轨形的奇异部分相当小。
\begin{proposition}\label{prop:terminal-gorenstein-orbifold-codimension}
若 Gorenstein 轨形至多具有终端奇点，则其奇异轨迹的余维至少为 $4$。
\end{proposition}
\begin{proof}
可归结到 $\mathbb C^d/G$，其中 $G$ 为小子群。给定阶为 $r$ 的非恒等元 $g\in G$，令 $\zeta$ 为一个本原 $r$ 次单位根，将 $g$ 的特征值写成 $\zeta^{a_i}$，其中 $0\leq a_i<r$。由 \SourceCite{Reid1}，Gorenstein 条件意味着 $\sum_{i=1}^d a_i\equiv0\pmod r$；由 \SourceCite{Reid2}，终端条件意味着 $\sum_{i=1}^d a_i>r$。

根据 \SourceCite{Prill}，$\mathbb C^d/G$ 的奇异轨迹是
\[
\{x\in\mathbb C^d:\text{存在 }g\in G\setminus\{1\}\text{ 使 }g(x)=x\}/G.
\]
若奇异轨迹的余维小于 $4$，则能找到 $g\ne1$，使其特征值经适当重排后为 $\zeta^a,\zeta^b,\zeta^c,1,\ldots,1$。Gorenstein 条件与终端条件给出 $r\mid(a+b+c)$ 和 $a+b+c>r$，因而 $a+b+c=2r$。再对 $g^{-1}$ 应用终端判别条件，很容易得到矛盾。细节留给读者。
\end{proof}

轨形同时也是环面簇时，\SourceCite[定理 2.2.9]{Batyrev4} 已证明命题~\ref{prop:terminal-gorenstein-orbifold-codimension}。在三维情形，立即得到如下推论。
\begin{corollary}\label{cor:terminal-gorenstein-threefold-orbifold-smooth}
至多具有终端奇点的三维 Gorenstein 轨形是光滑的。
\end{corollary}

按照\hyperref[chap:introduction]{第一章}中 Calabi--Yau 的定义，$d$ 维轨形 $V$ 是 Calabi--Yau 簇，当且仅当其典范类平凡，且 $H^i(V,\mathcal O)=0$，$i=1,\ldots,d-1$。这是因为典范类平凡意味着 Gorenstein，而对轨形，Gorenstein 又意味着其奇点为典范奇点。此外，上述
\SourcePage{432}
推论表明，三维极小 Calabi--Yau 轨形是光滑的，因为“极小”意味着奇点至多为终端奇点。

接着定义轨形的子轨形。
\begin{definition}\label{def:suborbifold}
给定轨形 $X$，若子簇 $Y\subset X$ 满足：每一点 $p\in Y$ 处，存在 $X$ 的局部坐标片 $(U/G,0)$，使 $Y$ 在 $U$ 中的逆像在 $0$ 处光滑，则称 $Y$ 为\emph{子轨形}。
\end{definition}

容易看出，轨形的子轨形仍是轨形；反之则不成立。$X$ 的一个子簇，即便本身是轨形，也未必是子轨形。这是因为子轨形的奇点与环境空间的奇点有密切联系。
\index{轨形|hyperpage}\index{轨形!子轨形|hyperpage}\index{复反射|hyperpage}\index{小子群|hyperpage}

\section{轨形上的微分形式}\label{sec:orbifold-differential-forms}
轨形上有很好的微分形式理论。给定 $p\in X$，考虑局部坐标片 $(U/G,0)$。将 $U/G$ 上的 $C^\infty$ 形式定义为 $U$ 上的 $G$ 不变 $C^\infty$ 形式。不同坐标片上定义的形式，有自然的粘合方式，由此可定义 de Rham 上同调群 $H^*_{DR}(X,\mathbb R)$；它们与通常的上同调群同构（见 \SourceCite{Satake}）。

轨形上也可进行 Hodge 理论。轨形上的全纯 $p$ 形式用局部坐标片 $(U/G,0)$ 定义，称为 $X$ 上的 \emph{Zariski $p$ 形式}。它们确定一个层，记为 $\widehat\Omega_X^p$。Zariski 层有如下简单刻画。
\begin{proposition}\label{prop:zariski-forms-smooth-locus-pushforward}
若 $X$ 是轨形，$j:X_0\hookrightarrow X$ 为其光滑轨迹的包含映射，则 $\widehat\Omega_X^p=j_*(\Omega_{X_0}^p)$，其中 $\Omega_{X_0}^p$ 为复流形 $X_0$ 上通常的全纯 $p$ 形式层。
\end{proposition}

轨形 $X$ 是 Cohen--Macaulay 簇；可以证明，若 $d=\dim(X)$，则 $\widehat\Omega_X^d$ 为 $X$ 的\IndexLocation{indexloc:0225:432}对偶化层。因此，这一记号与前节一致。层 $\widehat\Omega_X^p$ 还有许多其他良好性质，包括：
\begin{itemize}
\item 存在微分 $d:\widehat\Omega_X^p\to\widehat\Omega_X^{p+1}$，使 $(\widehat\Omega_X^*,d)$ 成为常值层 $\mathbb C_X$ 的一个消解。
\item 存在自然乘积 $\widehat\Omega_X^p\otimes\widehat\Omega_X^q\to\widehat\Omega_X^{p+q}$，使自然映射
\[
\widehat\Omega_X^p\longrightarrow\mathcal Hom_{\mathcal O_X}
(\widehat\Omega_X^{d-p},\widehat\Omega_X^d)
\]
是同构。
\end{itemize}
也可定义 Hodge 群 $H^{p,q}(X)$，此时 Dolbeault 定理
\[
H^q(X,\widehat\Omega_X^p)\simeq H^{p,q}(X)
\]
仍然成立（见 \SourceCite{Baily1}）。对轨形，通往相交上同调的自然映射
\[
H^*(X,\mathbb Q)\longrightarrow IH^*(X,\mathbb Q)
\]
是同构。由于射影簇的相交同调有自然 Hodge 结构（见 \SourceCite{Saito}），可知当 $X$ 为射影轨形时，$H^k(X,\mathbb C)=\bigoplus_{p+q=k}H^{p,q}(X)$ 有纯 Hodge 结构（但须在 $\mathbb Q$ 上考虑），并满足强 Lefschetz 定理。这对\IndexLocation{indexloc:1081:432}单纯多胞形的组合性质有有趣的推论（见 \SourceCite[第 5.6 节]{Fulton3}）。轨形 Hodge 理论的其他参考资料有 \SourceCite{Steenbrink1,Steenbrink2}。

轨形上的 \emph{Kähler 形式}，是一个实的、光滑且闭的 $(1,1)$ 形式，并且在每一点都正定（在坐标片 $(U/G,0)$ 上，这意味着它到 $U$ 的拉回正定）。轨形上 Kähler 形式的更多细节，见 \SourceCite{AGM1,Baily2}。
\index{Zariski形式@Zariski 形式|hyperpage}\index{轨形!Hodge理论@Hodge 理论|hyperpage}\index{轨形!Kähler形式@Kähler 形式|hyperpage}

\SourcePage{433}
\section{轨形的切层}\label{sec:orbifold-tangent-sheaf}
任意代数簇 $X$ 都有切层 $\Theta_X$。当 $X$ 光滑时，还知道 $\Theta_X$ 与 $\widehat\Omega_X^1$ 对偶，且 $H^1(X,\Theta_X)$ 分类 $X$ 的无穷小形变。$X$ 为轨形时，第一个结论仍然成立。
\begin{proposition}\label{prop:orbifold-tangent-dual-zariski-forms}
若 $X$ 为轨形，则 $\Theta_X\simeq\mathcal Hom_{\mathcal O_X}(\widehat\Omega_X^1,\mathcal O_X)$。
\end{proposition}
\begin{proof}
令 $\Omega_X^1$ 为 $X$ 上 Kähler 一形式层。由 $\Omega_X^1$ 的泛性质，$\Theta_X=\mathcal Hom_{\mathcal O_X}(\Omega_X^1,\mathcal O_X)$。若 $j:X_0\to X$ 为 $X$ 的光滑部分的包含映射，则 \SourceCite[第 128 页]{Oda} 的论证表明，可先限制到 $X_0$ 来计算 $\mathcal Hom_{\mathcal O_X}(-,\mathcal O_X)$。命题由此容易推出。
\end{proof}

至于 $H^1(X,\Theta_X)$，回顾一般簇 $X$ 的无穷小形变由 $\operatorname{Ext}^1(\Omega_X^1,\mathcal O_X)$ 分类。对任意轨形，它可能不同于 $H^1(X,\Theta_X)$；但有一种情形，两者一致。
\begin{proposition}\label{prop:terminal-gorenstein-orbifold-deformations}
若 $X$ 是至多具有终端奇点的 Gorenstein 轨形，则 $H^1(X,\Theta_X)$ 分类 $X$ 的无穷小形变。
\end{proposition}
\begin{proof}
\SourceCite{Schlessinger} 的一个结果表明，若轨形的奇异轨迹的余维至少为三，则其层 $\mathcal Ext^1(\Omega_X^1,\mathcal O_X)$ 为零。由命题~\ref{prop:terminal-gorenstein-orbifold-codimension}，这里满足该条件；再由 $\operatorname{Ext}$ 的局部到整体谱序列，立即得到结论。
\end{proof}

对于至多具有终端\IndexLocation{indexloc:1263:433}环面型奇点且在余维三处为 $\mathbb Q$ 析因的簇（参见 \SourceCite[定理 2.2.9]{Batyrev4}），这一命题也适用，而不必假定它们是轨形。

对于任意轨形 $X$，还可证明 $\Theta_X=j_*\Theta_{X_0}$，并且 $\widehat\Omega_X^1$ 是 $\Omega_X^1$ 的双对偶。若 $X$ 为 Gorenstein 轨形，则同构 $\widehat\Omega_X^{d-1}\simeq\mathcal Hom_{\mathcal O_X}(\widehat\Omega_X^1,\widehat\Omega_X^d)$ 与命题~\ref{prop:orbifold-tangent-dual-zariski-forms} 给出
\begin{equation}\label{eq:orbifold-tangent-top-forms}
\widehat\Omega_X^{d-1}\simeq\Theta_X\otimes\widehat\Omega_X^d.
\end{equation}
可如下将它应用于 Calabi--Yau 簇。
\begin{proposition}\label{prop:minimal-cy-orbifold-deformations}
$d$ 维极小 Calabi--Yau 轨形 $V$ 的无穷小形变，由
\[
H^1(V,\Theta_V)\simeq H^1(V,\widehat\Omega_V^{d-1})
\]
分类。
\end{proposition}
\begin{proof}
极小 Calabi--Yau 轨形具有 Gorenstein 终端奇点，所以由命题~\ref{prop:terminal-gorenstein-orbifold-deformations}，$H^1(V,\Theta_V)$ 分类其无穷小形变。由于 $\widehat\Omega_V^d$ 平凡，再由式~\eqref{eq:orbifold-tangent-top-forms} 即得结论。
\end{proof}
\index{轨形!切层|hyperpage}\index{轨形!无穷小形变|hyperpage}

\section{辛轨形}\label{sec:symplectic-orbifolds}
最后，可在 $C^\infty$ 范畴中定义轨形，并有\emph{轨形\IndexLocation{indexloc:0993:433}微分同胚}的概念。\emph{辛轨形}是带有闭的、非退化的二形式的 $C^\infty$ 轨形（二形式用局部坐标片定义）。此外，辛约化过程在此情形下自然适用。细节见 \SourceCite[II-3.6]{Audin}。
\index{辛轨形|hyperpage}

\chapter{物理理论}\label{app:physical-theories}
\SourcePage{434}
本附录概述本书提到的物理理论中的若干要点。\hyperref[sec:general-field-theories]{第一节}介绍一些基本物理理论的背景，按照
\[
\text{经典力学}\longrightarrow\text{经典场论}\longrightarrow\text{量子场论}
\]
的顺序，逐步过渡到量子场论。每一种理论都有 Lagrange 表述和 Hamilton 表述。在量子场论中，将看到 Lagrange 表述更适合路径积分，Hamilton 表述则导出 Hilbert 空间上的自伴算子代数。由 Lagrange 表述转到 Hamilton 表述很容易，因此总是从 Lagrange 表述出发。

读者若对经典力学的基础知识有所了解（例如 \SourceCite{Arnold} 中的内容），会有所帮助。我们将讨论经典场论、量子场论以及一些规范理论。本附录余下各节介绍非线性 $\sigma$ 模型、共形场论、Landau--Ginzburg 理论、规范线性 $\sigma$ 模型，最后介绍拓扑量子场论。

对这里某些理论更详细、也更偏数学的论述，可参阅 1996--1997 年 IAS 量子场论项目的讲义 \SourceCite{DEFJKMMW}。

\section{一般场论}\label{sec:general-field-theories}
本节解释什么是场论。这里不求自成体系，而是通过几个简单例子，说明正文中提及的理论属于何种类型；一些必要的背景材料则不展开。更多细节见 \SourceCite{Rabin}；这里的介绍大体沿用其思路，但并非完全照搬。

先考虑时空 $M$ 上的经典场论。$M$ 的典型例子包括四维 Minkowski 空间，或弦的世界面（即 Riemann 曲面）。粗略地说，$M$ 上的\emph{经典场}可以是函数、微分形式、丛的截面或丛上的联络。此时不宜用更精确的定义来限制讨论范围；有时还会将场当作函数来谈论，使这种表述更加粗略。

物理动力学由 \emph{Lagrange 密度}决定，它是场及其导数的泛函。\footnote{这种说法不够精确，因为我们需要对 Lagrange 密度积分。读者会注意到，积分前总会把 Lagrange 密度乘以适当的体积形式。}Lagrange 量的选择受到对称性的强烈约束，并由场论所要描述的具体物理过程决定。动力学的基本原理是：场必须使作用量取驻值；作用量是
\SourcePage{435}
Lagrange 密度在区域 $V\subset M$ 上的积分。将不同点处的场值视为独立的动力学变量，便可由作用量积分的 Euler--Lagrange 方程确定经典场方程。Lagrange 量有对称性时，这个处理需要修改；本节稍后将讨论。

第一个例子取 $M=M^4$ 为四维 Minkowski 空间。约定在 $M$ 上使用坐标 $x=(x^0,x^1,x^2,x^3)$，其中 $x^0$ 表示时间。有时也写成 $x=(t,\vec x)$，其中 $\vec x$ 表示空间变量，$t$ 表示时间。再令 $\partial_i=\partial/\partial x^i$，并置 $\vec\partial=(\partial_0,\partial_1,\partial_2,\partial_3)$。还将使用 $M^4$ 上号型为 $(1,3)$ 的 Lorentz 内积；例如，波算子 $\partial_0^2-\sum_{i=1}^3\partial_i^2$ 可记为 $\Box=\vec\partial\cdot\vec\partial$。符号 $d^4x$ 表示 $M$ 上的 Euclid 体积密度。

现在考虑质量为 $m$ 的标量场。在数学上，它是 $M^4$ 上的实值函数 $\phi(x)$，满足 Klein--Gordon 方程
\begin{equation}\label{eq:klein-gordon-field}
(\Box+m^2)\phi(x)=0.
\end{equation}
此方程也可由区域 $V\subset M$ 上的作用量得到，作用量定义为
\begin{equation}\label{eq:scalar-field-action}
S=S(\phi)=\frac12\int_Vd^4x\,
(\vec\partial\phi\cdot\vec\partial\phi-m^2\phi^2)
=\int_Vd^4x\,\mathcal L(\phi,\vec\partial\phi),
\end{equation}
其中
\begin{equation}\label{eq:scalar-field-lagrangian-density}
\mathcal L(\phi,\vec\partial\phi)
=\frac12(\vec\partial\phi\cdot\vec\partial\phi-m^2\phi^2)
\end{equation}
是 \emph{Lagrange 密度}。照常，$\phi$ 使作用量 $S$ 取驻值的条件是满足 Euler--Lagrange 方程
\begin{equation}\label{eq:field-euler-lagrange}
\frac{\delta\mathcal L}{\delta\phi}
-\vec\partial\cdot\frac{\delta\mathcal L}{\delta(\vec\partial\phi)}=0,
\end{equation}
其中 $\delta\mathcal L/\delta\phi$ 和 $\delta\mathcal L/\delta(\vec\partial\phi)$ 表示变分导数。将它应用于式~\eqref{eq:scalar-field-lagrangian-density}，便得到式~\eqref{eq:klein-gordon-field}。

接着回顾经典力学的 Hamilton 表述。从 Lagrange 函数 $L=L(q,\dot q,t)$ 出发，其中 $q=q_i(t)$ 是坐标，$t$ 是时间，点号表示时间导数。内蕴地说，$L$ 是构型流形 $M$ 的切丛 $T_M$ 上的函数。$q_i$ 的\emph{共轭动量}定义为 $p_i=\partial L/\partial\dot q_i$。这个名称很自然，因为对于势场中的自由点粒子，坐标函数 $q_i$ 的共轭动量，就是动量的相应分量。注意，$(p_i,q_i)$ 构成余切丛 $T_M^*$ 上的一组局部坐标。流形 $T_M^*$ 有典范辛结构，其局部坐标表达式为
\[
\omega=\sum_i dp_i\wedge dq_i.
\]
方程 $p_i=\partial L/\partial\dot q_i$ 局部定义一个映射 $T_M\to T_M^*$，我们假设它是同构。

定义 Hamilton 量为
\[
H=\sum_i p_i\dot q_i-L.
\]
\index{经典场|hyperpage}\index{Lagrange密度@Lagrange 密度|hyperpage}\index{作用量|hyperpage}\index{共轭动量|hyperpage}\index{Klein--Gordon方程@Klein--Gordon 方程|hyperpage}

\SourcePage{436}
Hamilton 量是 $T_M^*$ 上的函数，这里利用上述同构，将 $L$ 也视为定义在 $T_M^*$ 上。Hamilton 方程为
\[
\dot q_i=\frac{\partial H}{\partial p_i},\qquad
\dot p_i=-\frac{\partial H}{\partial q_i}.
\]
不过，对我们的目的，等价形式
\begin{equation}\label{eq:classical-poisson-evolution}
\dot A=[A,H]=\sum_{i=1}^3
\left(\frac{\partial A}{\partial q_i}\frac{\partial H}{\partial p_i}
-\frac{\partial A}{\partial p_i}\frac{\partial H}{\partial q_i}\right),
\end{equation}
其中 $A=q_i$ 或 $p_i$，更容易推广。表达式 $[A,H]$ 是 \emph{Poisson 括号}的一个例子。内蕴地说，辛形式 $\omega$ 确定同构 $I:T_M^*\to T_M$，两个函数 $f,g$ 的 Poisson 括号可改写为
\[
[f,g]=\langle df,I(dg)\rangle.
\]
Hamilton 方程可由 Lagrange 函数 $L$ 所确定作用量的 Euler--Lagrange 方程推出。还要注意，按上述 Poisson 括号的定义，有交换关系
\begin{equation}\label{eq:classical-canonical-poisson-brackets}
\begin{aligned}
[q_i,q_j]&=[p_i,p_j]=0,\\
[q_i,p_j]&=\delta_{ij}.
\end{aligned}
\end{equation}

现在将这个经典系统量子化。得到一个态的 Hilbert 空间 $\mathcal H$，$T_M^*$ 上的函数 $f$ 被替换为 $\mathcal H$ 上的自伴算子 $\mathbf f$。$f$ 的经典取值体现为算子的特征值。例如，若某个态 $\psi\in\mathcal H$ 满足 $\mathbf q_i\cdot\psi=q_i^0\psi$，则将 $\psi$ 理解为一个粒子的第 $i$ 个坐标为 $q_i^0$ 时所对应的量子态。同样，算子 $\mathbf p_i$ 的特征值与动量第 $i$ 个分量的经典取值等同。Hilbert 空间 $\mathcal H$ 上的内积，可用来确定观测到某量子态具有某个确定经典取值的概率。

按照\emph{正则量子化}方法，将式~\eqref{eq:classical-canonical-poisson-brackets} 中的 Poisson 括号替换为对易子
\[
\begin{aligned}
[\mathbf q_i,\mathbf q_j]&=[\mathbf p_i,\mathbf p_j]=0,\\
[\mathbf q_i,\mathbf p_j]&=i\delta_{ij}.
\end{aligned}
\]
运动方程~\eqref{eq:classical-poisson-evolution} 则被下述陈述替代：假定 $H$ 不显含时间，一个态经过时间间隔 $t$ 的传播，由酉算子
\begin{equation}\label{eq:quantum-unitary-time-evolution}
U_t=e^{-i\mathbf Ht}
\end{equation}
给出，这里选取单位使 Planck 常数 $\hbar=1$。这是对量子力学 \emph{Schrödinger 绘景}的极简描述：算子不变，态随时间演化。量子力学也可等价地用 \emph{Heisenberg 绘景}表述：态不变，算子随时间演化。通过 $U_t$ 的共轭作用，可以比较 Schrödinger 绘景与 Heisenberg 绘景。特别地，Schrödinger 绘景中的算子 $\mathbf A$，在 Heisenberg 绘景中替换为随时间变化的算子 $\mathbf A(t)=U_t^{-1}\mathbf A U_t$。计算得
\begin{equation}\label{eq:heisenberg-operator-evolution}
\frac{d}{dt}\mathbf A(t)=-i[\mathbf A(t),\mathbf H],
\end{equation}
\SourcePage{437}
它是式~\eqref{eq:classical-poisson-evolution} 的量子力学版本。式~\eqref{eq:heisenberg-operator-evolution} 中使用通常的算子对易子定义 $[\mathbf A,\mathbf H]=\mathbf A\mathbf H-\mathbf H\mathbf A$。在量子场论中，我们将使用 Heisenberg 绘景，不再特别说明。

将上述过程用于场论，定义\emph{共轭动量} $\pi(x)$ 为
\[
\pi(x)=\frac{\delta\mathcal L}{\delta(\partial_0\phi)}.
\]
假设只有一个场 $\phi(x)$，则 \emph{Hamilton 量}（或能量）$H$ 定义为 $\pi(x)\partial_0\phi(x)-\mathcal L$ 的空间积分，即对不含时间变量的空间变量积分，其中 $\mathcal L$ 是 Lagrange 密度。Hamilton 量是 $\phi(x)$ 和 $\pi(x)$ 的泛函，$\phi(x)$ 的时间导数已被消去。Euler--Lagrange 方程~\eqref{eq:field-euler-lagrange} 给出任意可观测量 $O$ 随时间传播的 Hamilton 方程；在数学上，$O$ 就是场 $\phi(x)$ 和 $\pi(x)$ 的泛函。这些方程类似于式~\eqref{eq:classical-poisson-evolution}，即
\[
\dot O=[O,H]=\int
\left(\frac{\delta O}{\delta\phi(x)}\frac{\delta H}{\delta\pi(x)}
-\frac{\delta O}{\delta\pi(x)}\frac{\delta H}{\delta\phi(x)}\right)
dx^1dx^2dx^3.
\]
令 $O=\phi$ 或 $\pi$，便得到 Hamilton 场方程
\[
\dot\phi=[\phi,H],\qquad\dot\pi=[\pi,H].
\]
对标量场 $\phi(x)$，共轭动量为 $\pi(x)=\partial_0\phi(x)$，Hamilton 量为
\begin{equation}\label{eq:scalar-field-hamiltonian}
H=\frac12\int
\bigl(\pi(x)^2+|\vec\nabla\phi(x)|^2+m^2\phi^2\bigr)
dx^1dx^2dx^3.
\end{equation}
Poisson 括号为
\begin{equation}\label{eq:classical-equal-time-field-brackets}
\begin{aligned}
[\phi(\vec x,t),\phi(\vec y,t)]&=[\pi(\vec x,t),\pi(\vec y,t)]=0,\\
[\phi(\vec x,t),\pi(\vec y,t)]&=\delta(\vec x-\vec y).
\end{aligned}
\end{equation}
它们是式~\eqref{eq:classical-canonical-poisson-brackets} 的场论对应。

下一步转向量子场。在我们希望达到的一般性下，这一步尚无良好的数学定义。其想法是寻找态的 Hilbert 空间 $\mathcal H$，将经典场 $\phi(x)$ 替换为量子场，记为 $\Phi(x)$。它们是依赖于 $x$ 的 $\mathcal H$ 上的算子，并在变量 $x$ 中作为分布。场的经典取值体现为算子的特征值。例如，若存在态 $\psi\in\mathcal H$，使某个函数 $f(x)$ 满足 $\Phi(x)\cdot\psi=f(x)\psi$，就把这个量子态理解为对应于经典场 $\phi(x)=f(x)$。与通常的量子力学类似，内积用来确定相对概率。

采用正则量子化，令 $\Pi(\vec x,t)$ 为对应于 $\pi(\vec x,t)$ 的算子值分布，并将式~\eqref{eq:classical-equal-time-field-brackets} 中的 Poisson 括号替换为
\[
\begin{aligned}
[\Phi(\vec x,t),\Phi(\vec y,t)]&=[\Pi(\vec x,t),\Pi(\vec y,t)]=0,\\
[\Phi(\vec x,t),\Pi(\vec y,t)]&=i\delta(\vec x-\vec y).
\end{aligned}
\]
这些等时对易关系是算子值分布的恒等式。

下面说明，这一形式体系如何将量子场 $\Phi(x)$ 解释为由粒子（量子）构成。回到 Klein--Gordon 方程~\eqref{eq:klein-gordon-field}，使用求解偏微分方程的标准 Fourier 变换方法，可将经典场 $\phi$ 写成
\begin{equation}\label{eq:classical-scalar-field-fourier-expansion}
\phi(x)=\int\bigl(
a(\vec k)e^{i(\vec k\cdot\vec x-E_{\vec k}t)}
+a^\dagger(\vec k)e^{-i(\vec k\cdot\vec x-E_{\vec k}t)}
\bigr)\frac{dk^1dk^2dk^3}{\sqrt{(2\pi)^3\,2E_{\vec k}}}.
\end{equation}
\index{Poisson括号@Poisson 括号|hyperpage}\index{正则量子化|hyperpage}\index{Schrödinger绘景@Schrödinger 绘景|hyperpage}\index{Heisenberg绘景@Heisenberg 绘景|hyperpage}\index{等时对易关系|hyperpage}

\SourcePage{438}
这里 $a$ 和 $a^\dagger$ 是变量 $\vec k=(k^1,k^2,k^3)$ 的某两个互为复共轭的函数。式~\eqref{eq:classical-scalar-field-fourier-expansion} 中置 $E_{\vec k}=\sqrt{\vec k\cdot\vec k+m^2}$。由于 $E_{\vec k}^2=\vec k\cdot\vec k+m^2$，Klein--Gordon 方程~\eqref{eq:klein-gordon-field} 成立；在 $E_{\vec k}=\sqrt{\vec k\cdot\vec k+m^2}$ 定义的\emph{质壳}上，表达式 $dk^1dk^2dk^3/E_{\vec k}$ 在 $(E_{\vec k},k^1,k^2,k^3)$ 的 Lorentz 变换下不变。$E_{\vec k}$ 是质量为 $m$、动量为 $\vec k$ 的相对论粒子的总能量。对不熟悉这一点的数学读者，说明一下：当 $|\vec k|\ll m$ 时，有近似式 $E_{\vec k}\sim m+(\vec k\cdot\vec k)/(2m)$。第一项是粒子的静止能量（通常写成 $E=mc^2$，但这里取单位使 $c=1$），第二项则是质量为 $m$、动量为 $\vec k$ 的粒子的经典动能。

回到量子理论，我们断言，$a(\vec k)$ 和 $a^\dagger(\vec k)$ 有自然的粒子解释。首先，它们是算子值分布；可以证明，它们满足对易关系
\begin{equation}\label{eq:scalar-creation-annihilation-commutators}
\begin{aligned}
[a(\vec k),a(\vec k')]&=[a^\dagger(\vec k),a^\dagger(\vec k')]=0,\\
[a(\vec k),a^\dagger(\vec k')]&=\delta(\vec k-\vec k').
\end{aligned}
\end{equation}
此外，将展开式代回 Hamilton 量~\eqref{eq:scalar-field-hamiltonian}，可得
\[
H=\int\left(E_{\vec k}a^\dagger(\vec k)a(\vec k)
+\frac12E_{\vec k}\delta(\vec0)\right)dk^1dk^2dk^3,
\]
其中 $\vec0$ 为三维零向量。最后一项在数学上没有意义，因为它是一个无穷大的标量算子。处理办法是直接略去此项，因为从 Hamilton 量中减去一个常数，不影响动力学。最终得到
\begin{equation}\label{eq:normal-ordered-scalar-hamiltonian}
H=\int E_{\vec k}a^\dagger(\vec k)a(\vec k)\,dk^1dk^2dk^3.
\end{equation}

这有很好的直观描述。对易关系~\eqref{eq:scalar-creation-annihilation-commutators} 和 Hamilton 量~\eqref{eq:normal-ordered-scalar-hamiltonian} 表明，作用算子 $a(\vec k)$ 会降低能量。真空应当是最低能量态，因此可将真空态 $V$ 定义为被所有 $a(\vec k)$ 湮灭的态。利用式~\eqref{eq:scalar-creation-annihilation-commutators} 及 $a(\vec k)(V)=0$，可知态 $a^\dagger(\vec k')(V)$ 是能量（即特征值）为 $E_{\vec k'}$ 的特征态：
\[
\begin{aligned}
H\bigl(a^\dagger(\vec k')(V)\bigr)
&=\int E_{\vec k}a^\dagger(\vec k)a(\vec k)a^\dagger(\vec k')(V)\,d^3k\\
&=\int E_{\vec k}a^\dagger(\vec k)
\bigl(a^\dagger(\vec k')a(\vec k)+\delta(\vec k-\vec k')\bigr)(V)\,d^3k\\
&=E_{\vec k'}a^\dagger(\vec k')(V).
\end{aligned}
\]
于是，将 $a^\dagger(\vec k)$ 解释为\emph{\IndexLocation{indexloc:0203:438}产生算子}：它的作用产生一个动量为 $\vec k$ 的标量粒子。同样，$a(\vec k)$ 是\emph{\IndexLocation{indexloc:0049:438}湮灭算子}，其作用湮灭一个动量为 $\vec k$ 的粒子。

结论是：一般而言，式~\eqref{eq:classical-scalar-field-fourier-expansion} 将量子场 $\Phi(x)$ 表示为基本粒子态的叠加。

我们并不指望上述说明能使初次接触这些内容的读者完全信服，但希望明确指出：从 Lagrange 密度和形式上的量子化程序出发，已经得到
\SourcePage{439}
物质的量子结构，并以这些量子重新表述了动力学。换句话说，量子场的动力学完全由 Lagrange 密度决定。还应提到，上述内容中的一部分（但并非全部）可以建立在牢固的数学基础之上，见 \SourceCite[第 8 讲]{Rabin}。

迄今主要关注量子场论的 Hamilton 表述。如前所述，也有涉及\emph{路径积分}概念的 Lagrange 表述。设在时刻 $t<t'$，经典场的取值分别为 $\phi_0(\vec x,t)$ 和 $\phi_1(\vec x,t')$。以 $|\phi_0(\vec x,t)\rangle$ 和 $|\phi_1(\vec x,t')\rangle$ 表示在相应时刻的算子值分布 $\Phi$ 下，特征值分别为 $\phi_0(\vec x,t)$、$\phi_1(\vec x,t')$ 的量子态。也就是说，
\begin{equation}\label{eq:field-eigenstate-notation}
\Phi(\vec x,t)|\phi_0(\vec x,t)\rangle
=\phi_0(\vec x,t)|\phi_0(\vec x,t)\rangle,
\end{equation}
$\phi_1$ 也有类似表达式。式~\eqref{eq:field-eigenstate-notation} 左边涉及算子的作用，右边涉及乘以一个分布。

这两个态的 Hilbert 内积为 $\langle\phi_1(\vec x,t')\mid\phi_0(\vec x,t)\rangle$。其物理含义是：它的模平方给出态 $\phi_0(\vec x,t)$ 传播为 $\phi_1(\vec x,t')$ 的概率密度。形式计算表明，若 $S(\phi)$ 是推广式~\eqref{eq:scalar-field-action} 的作用量，则
\[
\langle\phi_1(\vec x,t')\mid\phi_0(\vec x,t)\rangle
=N\int[D\phi]e^{iS(\phi)},
\]
其中右侧\emph{路径积分}遍及场空间中所有初始点为 $\phi_0$、终止点为 $\phi_1$ 的路径。$N$ 是适当的归一化因子，$[D\phi]$ 是路径空间上的适当测度。目前，路径积分（或泛函积分）还没有良好的数学定义。

路径积分还可计算某些物理\emph{\IndexLocation{indexloc:0183:439}\IndexLocation{indexloc:0192:439}关联函数}，有时称为 \emph{$n$ 点函数}。选取 $n$ 个点 $x_1,\ldots,x_n\in M$，形式上将 $n$ 点函数定义为
\begin{equation}\label{eq:path-integral-n-point-function}
\langle\Phi(x_1)\cdots\Phi(x_n)\rangle
=N\int[D\phi]\,\phi(x_1)\cdots\phi(x_n)e^{iS(\phi)}.
\end{equation}
它们是点 $x_1,\ldots,x_n$ 及各点所对应场的类型的函数，而非场的具体取值的函数，因为这些取值已被积分。在上述标量场例子中，只有一种场；但例如 Yukawa 耦合是一类三点函数，其中不同点对应不同类型的场。Yukawa 耦合是\emph{拓扑量子场论}中的三点函数，这大大简化了问题。\hyperref[sec:topological-qft]{附录~\ref*{sec:topological-qft}}将讨论拓扑量子场论。

这些量本身具有物理意义。例如，二点函数就是理论中的 Green 函数。在很实在的意义上，可以说 $n$ 点函数包含这个理论的全部物理预言。

虽然式~\eqref{eq:path-integral-n-point-function} 这类路径积分尚无法严格表述，但有一些获得部分论证、被普遍接受的启发式计算方法。最重要的是\emph{驻相法}：把路径积分“局部化”为 $S$ 的驻点空间上的积分，该空间往往是有限维的。一个粗略解释是：若 $\phi_0$ 为 $S$ 的驻点，则在 $\phi_0$ 附近，$e^{iS(\phi)}$ 的相位变化缓慢，
\SourcePage{440}
因而 $\phi_0$ 附近的局部路径积分可以有非零贡献。若 $\phi_0$ 不是驻点，则相位振荡得更快，往往导致相互抵消。因此，$\phi_0$ 附近的态更可能出现，并对物理过程作出贡献。这为经典场论和量子场论建立了联系。驻相法对有限维参数空间是严格的，但路径积分所涉及的空间是无限维的。

这里有一个微妙之处。在路径积分~\eqref{eq:path-integral-n-point-function} 中，各项 $\phi(x_i)$ 是数值分布，彼此交换；但\emph{算子} $\Phi(x_i)$ 一般不交换，似乎有矛盾。解决办法是定义\emph{编时乘积} $T(\Phi(x_1)\cdots\Phi(x_n))$：按时间先后依次作用各 $\Phi(x_i)$，得到相应算子。于是
\[
\langle V\mid T(\Phi(x_1)\cdots\Phi(x_n))\mid V\rangle
=N\int[D\phi]\,\phi(x_1)\cdots\phi(x_n)e^{iS(\phi)},
\]
其中 $V$ 仍表示真空态。

还需要对\emph{规范理论}有所了解。此时，作用量在 $M$ 上一个连续的局部变换群下不变。这个群通常是无限维的；典型例子是 Yang--Mills 理论，规范群是从 $M$ 到某个有限维 Lie 群的映射空间。这里的目标很有限：在量子场论的讨论背景下，解释所谓\emph{规范固定}的作用，以及守恒量的概念。

最简单的例子是不含带电物质的经典电磁理论。此处只是让读者体会其风貌，并不力求自成体系。更多细节可参阅 \SourceCite{BjoD}。取 $M=M^4$ 为四维 Minkowski 空间，使用通常坐标 $(x^0,\ldots,x^3)=(t,\vec x)$。场为实一形式 $\phi=A_i\,dx^i$，\footnote{采用 Einstein 求和约定，对重复出现的上、下指标求和。通过与度量张量 $g_{ij}=\operatorname{diag}(1,-1,-1,-1)$ 缩并来升降指标，不再特别说明。}称为\emph{电磁势}。再令 $\vec A=(A_1,A_2,A_3)$。$\phi$ 的\emph{场强}定义为
\[
F=d\phi=\frac12F_{ij}\,dx^i\wedge dx^j,\qquad
F_{ij}=\partial_iA_j-\partial_jA_i.
\]
电场 $\vec E$ 和磁场 $\vec B$ 通过下式从场强中提取：
\[
F=-dt\wedge\left(\sum_{i=1}^3E_i\,dx^i\right)
+\iota_{\vec B}(dx^1\wedge dx^2\wedge dx^3),
\]
其中 $\iota_{\vec B}$ 是与视为向量场 $B^i\partial/\partial x^i$ 的 $\vec B$ 作内乘。按这些约定，电磁场强 $F$ 是闭二形式，等价于两个 Maxwell 方程
\[
\vec\nabla\times\vec E=-\dot{\vec B},\qquad
\vec\nabla\cdot\vec B=0.
\]
另两个方程来自 Lagrange 密度
\begin{equation}\label{eq:electromagnetic-lagrangian-density}
\mathcal L=-\frac14F_{ij}F^{ij}
=\frac12(|\vec E|^2-|\vec B|^2).
\end{equation}
此时，Euler--Lagrange 方程写成
\[
\vec\nabla\cdot\vec E=0,\qquad
\vec\nabla\times\vec B=\dot{\vec E},
\]
\SourcePage{441}
计算所得 Hamilton 量为
\[
H=\frac12\int(|\vec E|^2+|\vec B|^2)\,dx^1\wedge dx^2\wedge dx^3.
\]

给 $\phi=A_i\,dx^i$ 加上一个恰当形式，不改变场强 $F=d\phi$，因而也不改变相应的 Lagrange 量 $\mathcal L$。即它们在下列替换下不变：
\begin{equation}\label{eq:electromagnetic-gauge-transformation}
A_i\longmapsto A_i+\frac{\partial\Lambda(x)}{\partial x^i},
\end{equation}
其中 $\Lambda(x)$ 为 $M$ 上任意实函数。此替换对应于局部变换 $e^{i\Lambda(x)}$ 的\emph{规范群}。事实上，若将 $\sqrt{-1}A_i\,dx^i$ 视为 $M$ 上主 $U(1)$ 丛的联络形式，则该丛上的规范变换 $e^{i\Lambda(x)}$ 诱导式~\eqref{eq:electromagnetic-gauge-transformation}。因此，只需在 $\phi=A_i\,dx^i$ 的参数空间中选取一个截片，使每个可能的 $\phi$ 都通过规范变换~\eqref{eq:electromagnetic-gauge-transformation} 等价于截片内的某个 $\phi$。这样的截片称为\emph{规范选择}。\footnote{从路径积分的观点看，显然必须限制到这样的截片；否则，规范对称性会使同一条路径被重复计数，产生不希望出现的无穷大。}下面分两步作规范选择。首先，可设
\begin{equation}\label{eq:temporal-electromagnetic-gauge}
A_0=0.
\end{equation}
但不能只作这个规范选择就立即量子化，因为还需满足相应的 Euler--Lagrange 方程 $\delta S/\delta A_0=0$。它是 Gauss 定律 $\vec\nabla\cdot\vec E=0$；由于 $A_0=0$，它变为 $\vec\nabla\cdot\partial\vec A/\partial t=0$。因此，$\vec\nabla\cdot\vec A$ 与时间无关。现在解 Poisson 方程，找到与时间无关的函数 $f$，满足 $\vec\nabla\cdot\vec\nabla f=\vec\nabla\cdot\vec A$。将 $\phi=A_i\,dx^i$ 替换为 $\phi-df$ 不影响 $A_0$，所以还可使规范选择满足
\begin{equation}\label{eq:transverse-electromagnetic-gauge}
\vec\nabla\cdot\vec A=\vec\nabla\cdot(A_1,A_2,A_3)=0.
\end{equation}
一般而言，作规范选择之后，不能简单忽略已被“通过规范变换消去”的场。

固定规范选择~\eqref{eq:temporal-electromagnetic-gauge} 和~\eqref{eq:transverse-electromagnetic-gauge} 后，可采用与 Klein--Gordon 场类似的量子化过程。作 Fourier 展开
\begin{equation}\label{eq:electromagnetic-polarization-expansion}
\vec A(x)=\int\sum_{j=1}^2\vec\epsilon_j(\vec k)
\bigl(a_j(\vec k)e^{-ik\cdot x}+a_j^\dagger(\vec k)e^{ik\cdot x}\bigr)
\frac{dk^1\wedge dk^2\wedge dk^3}{\sqrt{2k_0(2\pi)^3}}.
\end{equation}
其中记号约定如下。$\vec\epsilon_j$ 是三维极化向量；由第二个规范条件~\eqref{eq:transverse-electromagnetic-gauge} 的 Fourier 变换可知，它们垂直于 $\vec k$，因而有两个独立的极化向量。最后，$k=(k_0,\vec k)$ 满足 $k_0=|\vec k|$，因为 $\vec A(x)$ 的每个分量都是四维波动方程的解（这由 Maxwell 方程推出）。

于是将 $a_j(\vec k)$ 和 $a_j^\dagger(\vec k)$ 分别解释为湮灭和产生一个动量为 $\vec k$、极化为 $\vec\epsilon_j(\vec k)$ 的电磁场量子的算子。这与前面对 $a(\vec k)$ 和 $a^\dagger(\vec k)$ 的讨论一致。
\index{产生算子|hyperpage}\index{湮灭算子|hyperpage}\index{质壳|hyperpage}\index{路径积分|hyperpage}\index{驻相法|hyperpage}\index{编时乘积|hyperpage}\index{规范固定|hyperpage}\index{电磁势|hyperpage}

\SourcePage{442}
一个关键事实是 Noether 定理：连续对称群给出守恒量。对使 Lagrange 密度不变的无穷小对称变换，设场变为 $\phi\mapsto\phi+\psi\varepsilon$，其中 $\varepsilon$ 为无穷小参数，则在场方程成立时，量
\begin{equation}\label{eq:noether-conserved-current}
\frac{\delta\mathcal L}{\delta(\partial_i\phi)}\psi
\end{equation}
的散度为零，称为\emph{守恒流}。因此，守恒流时间分量在固定时刻的空间积分与时间无关，称为\emph{守恒荷}。

在\hyperref[sec:conformal-field-theories]{附录~\ref*{sec:conformal-field-theories}}讨论共形场论时，规范群将是共形群。将看到，也需作类似的\IndexLocation{indexloc:0349:442}规范选择；规范群还产生类似的守恒量，它们在理论中起重要作用。

迄今讨论的\IndexLocation{indexloc:1155:442}量子场论最适合\emph{玻色子}，即自旋为整数的粒子。若干经典的（即非量子的）玻色场，可通过乘法组合，此乘法是交换的。但为描述\emph{费米子}（自旋为半整数的粒子）的场论，需要从乘法反交换的经典场出发。有几种表述方式，其中一种通过超空间来公理化反交换变量；\hyperref[sec:landau-ginzburg-models]{附录~\ref*{sec:landau-ginzburg-models}}将作简述。费米场的量子化过程，需要修改上面对玻色场概述的过程。下一节的非线性 $\sigma$ 模型，将给出\IndexLocation{indexloc:0313:442}费米变量的一些例子。关于玻色系统与费米系统的区别，很好的入门介绍见 \SourceCite[第 1、2 讲]{Alvarez}。
\index{守恒流|hyperpage}\index{守恒荷|hyperpage}\index{费米场|hyperpage}\index{玻色场|hyperpage}

\section{非线性 \texorpdfstring{$\sigma$}{σ} 模型}\label{sec:nonlinear-sigma-models}
为定义\emph{非线性 $\sigma$ 模型}（简称 \emph{$\sigma$ 模型}），从 Kähler 流形 $(X,g)$ 和闭二形式 $B$ 出发。$\sigma$ 模型建立在从 Riemann 曲面 $\Sigma$ 到 $X$ 的映射之上。把映射 $f:\Sigma\to X$ 理解为描述弦在 $X$ 中的传播，$\Sigma$ 就是弦传播时扫出的\emph{世界面}。理论中不固定 Riemann 曲面 $\Sigma$，而是允许任意曲面，因此允许一条弦分裂成两条或多条弦、或若干条弦结合等任意复杂的过程。换句话说，允许任意亏格的 Riemann 曲面，就自动将弦之间的相互作用纳入理论。这是所有弦理论的一般特征，并非 $\sigma$ 模型独有。

非线性 $\sigma$ 模型还包含反交换的费米场，它们是 $\Sigma$ 上某些丛的截面，如下表所示：
\begin{equation}\label{eq:sigma-model-fermionic-bundles}
\begin{array}{c|l}
\text{场}&\text{丛}\\\hline
\psi_+^i&K^{1/2}\otimes f^*(T_X^{1,0})\\
\psi_+^{\bar i}&K^{1/2}\otimes f^*(T_X^{0,1})\\
\psi_-^i&\overline K^{1/2}\otimes f^*(T_X^{1,0})\\
\psi_-^{\bar i}&\overline K^{1/2}\otimes f^*(T_X^{0,1})
\end{array}
\end{equation}
这里，$K$ 是 $\Sigma$ 的典范丛，$\overline K$ 是其复共轭（反全纯）丛，$K^{1/2}$ 是所选的 $K$ 的一个平方根（即 $\Sigma$ 上的旋量丛），$T_X$
\SourcePage{443}
是 $X$ 的复化切丛，并有通常的分解 $T_X=T_X^{1,0}\oplus T_X^{0,1}$。指标 $i,\bar i$ 取值从 $1$ 到 $n=\dim_{\mathbb C}X$，约定 $i$ 为全纯指标，$\bar i$ 为反全纯指标。选用这些记号，是为了可以将上表中的对象与它们的某些系数等同；具体含义由上下文辨明。做法如下：在 $\Sigma$ 上选局部全纯坐标 $z$，在 $X$ 上选局部复坐标 $x^i$。例如，$\psi_+^i$ 既可表示 $K^{1/2}\otimes f^*(T_X^{1,0})$ 的一个截面，也可表示该截面在 $(dz)^{1/2}\otimes\partial/\partial x^i$ 前的系数，其中 $(dz)^{1/2}$ 表示 $K^{1/2}$ 的局部生成元。记号上，$x^{\bar i}$ 与 $\overline{x^i}$ 含义相同。

现在可定义作用量 $S$。再令指标 $I,J$ 遍及 $i,\bar i$ 的所有可能取值，即同时遍及全纯与反全纯指标；于是 $x^I$ 遍及 $2n$ 个坐标 $x^i,x^{\bar i}$，这里 $n$ 仍是 $X$ 的复维数。置
\begin{equation}\label{eq:supersymmetric-sigma-action}
\begin{aligned}
S=\int_\Sigma\Bigl(&\frac12(g_{IJ}+iB_{IJ})\partial_zx^I\partial_{\bar z}x^J\\
&+\frac{\sqrt{-1}}2g_{i\bar i}\psi_-^{\bar i}D_z\psi_-^i
+\frac{\sqrt{-1}}2g_{i\bar i}\psi_+^{\bar i}D_{\bar z}\psi_+^i\\
&+\frac14R_{i\bar i j\bar j}\psi_+^i\psi_+^{\bar i}\psi_-^j\psi_-^{\bar j}\Bigr)
\sqrt{-1}\,dz\wedge d\bar z,
\end{aligned}
\end{equation}
其中 $R$ 是 $X$ 上度量 $g$ 的曲率张量，$D_z$ 来自将 $g$ 的 Levi--Civita 联络拉回到 $f^*T_X$ 所得的联络。注意，$S$ 同时依赖映射 $f:\Sigma\to X$ 和 $4n$ 个费米场 $\psi_\pm^i,\psi_\pm^{\bar i}$。

这个看似奇特的作用量其实很自然。对满足共形约束的经典解，玻色部分中的度量项化为 $\Sigma$ 的面积，$B$ 项则给出拓扑贡献。加入费米项的方式，既保持\emph{共形不变性}（本附录稍后讨论），又使理论具有\emph{超对称性}：作用量有一个无穷小对称变换，将费米子沿玻色方向扰动，反之亦然。事实上，这样的对称性有两个，因为 $\sigma$ 模型具有 $N=2$ 超对称。\IndexLocation{indexloc:1312:443}超对称变换相当复杂，可写成：
\begin{equation}\label{eq:sigma-model-supersymmetry-transformations}
\begin{aligned}
\delta x^i&=\sqrt{-1}\alpha_-\psi_+^i+\sqrt{-1}\alpha_+\psi_-^i,\\
\delta x^{\bar i}&=\sqrt{-1}\widetilde\alpha_-\psi_+^{\bar i}
+\sqrt{-1}\widetilde\alpha_+\psi_-^{\bar i},\\
\delta\psi_+^i&=-\widetilde\alpha_-\partial_zx^i
-\sqrt{-1}\alpha_+\psi_-^J\Gamma^i_{Jm}\psi_+^m,\\
\delta\psi_+^{\bar i}&=-\alpha_-\partial_zx^{\bar i}
-\sqrt{-1}\widetilde\alpha_+\psi_-^J\Gamma^{\bar i}_{J\bar m}\psi_+^{\bar m},\\
\delta\psi_-^i&=-\widetilde\alpha_+\partial_{\bar z}x^i
-\sqrt{-1}\alpha_-\psi_+^J\Gamma^i_{Jm}\psi_-^m,\\
\delta\psi_-^{\bar i}&=-\alpha_+\partial_{\bar z}x^{\bar i}
-\sqrt{-1}\widetilde\alpha_-\psi_+^J\Gamma^{\bar i}_{J\bar m}\psi_-^{\bar m}.
\end{aligned}
\end{equation}
这里，$\alpha_-$ 和 $\widetilde\alpha_-$ 是 $K^{-1/2}$ 的无穷小全纯截面，而 $\alpha_+$ 和 $\widetilde\alpha_+$ 是 $\overline K^{-1/2}$ 的无穷小反全纯截面。应用变换~\eqref{eq:sigma-model-supersymmetry-transformations}，并利用 Kähler 条件，可证明式~\eqref{eq:supersymmetric-sigma-action} 的被积二形式只改变一个恰当形式，因而作用量 $S$ 不变。

在杂化弦理论中，须适当修改 Lagrange 量，使某些费米场 $\psi_+^I$ 取值于包含 $E_8\times E_8$ 根格的十六维空间。于是所得理论可包含 $E_8\times E_8$ 规范
\SourcePage{444}
群；这在 Yukawa 耦合的定义中起作用。Yukawa 耦合就是我们理论中的关联函数（或 $n$ 点函数）。还应提到，在所考虑的杂化理论中，费米子的运动方程解耦为左行与右行方程，此处可将它们理解为全纯与反全纯部分。因此，这个 $\sigma$ 模型给出 $(2,2)$ 弦理论，即有两个全纯超对称和两个反全纯超对称。

借助\emph{扭曲理论}，即 A 模型与 B 模型，可以更容易地计算我们关心的关联函数。它们不同于 $\sigma$ 模型，但可以证明，其某些关联函数（包括\hyperref[chap:introduction]{第一章}提到的那些）与 $\sigma$ 模型中相应的关联函数一致。

这些理论的作用量在形式上与式~\eqref{eq:supersymmetric-sigma-action} 相同，但场 $\psi_\pm^I$ 的含义不同。A 模型将式~\eqref{eq:sigma-model-fermionic-bundles} 替换为
\[
\begin{array}{c|l}
\text{场}&\text{丛}\\\hline
\psi_+^i&f^*(T_X^{1,0})\\
\psi_+^{\bar i}&K\otimes f^*(T_X^{0,1})\\
\psi_-^i&\overline K\otimes f^*(T_X^{1,0})\\
\psi_-^{\bar i}&f^*(T_X^{0,1})
\end{array}
\]
而 B 模型使用
\[
\begin{array}{c|l}
\text{场}&\text{丛}\\\hline
\psi_+^i&K\otimes f^*(T_X^{1,0})\\
\psi_+^{\bar i}&f^*(T_X^{0,1})\\
\psi_-^i&\overline K\otimes f^*(T_X^{1,0})\\
\psi_-^{\bar i}&f^*(T_X^{0,1})
\end{array}
\]
换句话说，这些场局部上与 $\sigma$ 模型中的对应场一致，但在整体上用丛 $K^{1/2}$、$\overline K^{1/2}$ 或它们的逆作了扭曲。

这些模型仍然超对称，但此时只有 $N=1$ 超对称。此外，若 $V$ 和 $V^\circ$ 构成镜像对，则由 Calabi--Yau 流形 $V$ 得到的 A 模型，与由镜像流形 $V^\circ$ 得到的 B 模型镜像对称，其中须对应地选取 $V^\circ$ 的复结构和 $V$ 的 Kähler 结构。

先看 A 模型。此时，可将场 $\psi_+^i$ 和 $\psi_-^{\bar i}$ 视为 $f^*(T_X)$ 的一个截面的全纯与反全纯部分。此外，A 模型给 $X$ 上每个 $k$ 形式 $\eta$ 赋予一个\emph{局部算子} $\mathcal O_\eta(x)$，定义如下。可将 $k$ 形式写成 $\eta=h_{I_1\cdots I_k}(x)\,dx^{I_1}\wedge\cdots\wedge dx^{I_k}$，其中每个 $dx^I$ 或为 $dx^i$，或为 $dx^{\bar i}=d\overline{x^i}$。相应的局部算子定义为
\begin{equation}\label{eq:a-model-local-form-operator}
\mathcal O_\eta(x)=h_{I_1\cdots I_k}(x)\chi^{I_1}(x)\cdots\chi^{I_k}(x),
\end{equation}
其中 $\chi^I$ 或为 $\psi_+^i$，或为 $\psi_-^{\bar i}$；应将它们视为 Hilbert 空间 $\mathcal H$ 上的反交换算子。

式~\eqref{eq:supersymmetric-sigma-action} 的玻色部分的拓扑贡献由 $\int_\Sigma f^*(\omega)$ 给出，其中照常 $\omega=B+iJ$，$J$ 为度量 $g$ 的 Kähler 类。使用运动方程，即上面讨论的广义 Euler--Lagrange 方程后，作用量有所简化；其余项可以用一个费米算子，
\SourcePage{445}
即 \emph{BRST 算子} $Q$，紧凑地表示。这个算子是超对称从原来的 $N=2$ 理论减为扭曲后的 $N=1$ 理论时保留下来的超对称变换。事实上，模去运动方程后，有
\[
S=2\pi i\int_\Sigma\{Q,V\}\left(\frac i2dz\wedge d\bar z\right)
+2\pi\int_\Sigma f^*(\omega),
\]
其中 $V$ 是某个场，其细节无须关注；$\{Q,V\}$ 为反对易子。这里不展开细节，只指出 BRST 算子在量子场论中起重要作用。$Q^2=0$，算子 $d_Q$ 定义为 $d_Q(O)=\{Q,O\}$，这里花括号按分次对易子理解。它实际上确定一个复形，其上同调称为量子场论的 \emph{BRST 上同调}。有公式
\[
\{Q,\mathcal O_\eta\}=-\mathcal O_{d\eta},
\]
因此 $\eta\mapsto\mathcal O_\eta$ 给出从 $X$ 的 de Rham 上同调 $H^k(X)$ 到 BRST 上同调的映射。

令 $\omega_1,\ldots,\omega_n\in H^2(X)$，转向计算 $n$ 点关联函数 $\langle\mathcal O_{\omega_1}\cdots\mathcal O_{\omega_n}\rangle$。对式~\eqref{eq:supersymmetric-sigma-action} 作 BRST 局部化，相关固定点方程为 $\partial x^i/\partial\bar z=0$，故 $f$ 全纯。相应地，路径积分化为全纯映射空间上的积分。再置 $\gamma=f_*[\Sigma]$，将积分写成对 $\gamma$ 的求和，并把 $\gamma=0$ 项单独写出。得到
\[
\begin{aligned}
\langle\mathcal O_{\omega_1}\cdots\mathcal O_{\omega_n}\rangle
&=\int_X\omega_1\wedge\cdots\wedge\omega_n\\
&\quad+\sum_{\gamma\ne0}\int[D\phi][D\chi][D\psi]\,
e^{-2\pi\{Q,\int V\}}\prod_i\mathcal O_{\omega_i}
e^{2\pi i\int_\gamma\omega}\\
&=\int_X\omega_1\wedge\cdots\wedge\omega_n
+\sum_{\gamma\ne0}N_{\gamma,\omega_i}q^\gamma,
\end{aligned}
\]
其中 $q^\gamma$ 的记号如\hyperref[chap:gromov-witten]{第七章}所述，$N_{\gamma,\omega_i}$ 粗略定义如下：选取点 $p_1,\ldots,p_n\in\mathbb P^1$，并选取代表圈 $Z_i$，使其基本类与 $\omega_i$ Poincaré 对偶；$N_{\gamma,\omega_i}$ 表示满足 $f_*[\mathbb P^1]=\gamma$，且每个 $i$ 都有 $f(p_i)\in Z_i$ 的全纯映射 $f:\mathbb P^1\to X$ 的数目。这一定义还不严格，因为这样的映射可能有无穷多个，但可用\hyperref[chap:gromov-witten]{第七章}讨论的 Gromov--Witten 不变量概念精确化。如\hyperref[chap:quantum-cohomology]{第八章}所述，这些不变量是构造量子上同调的关键成分。\hyperref[sec:topological-qft]{附录~\ref*{sec:topological-qft}}将指出，A 模型是拓扑量子场论，因而三点关联函数定义结合环结构。上面第二个等式由在 $\chi$ 和 $\psi$ 的临界点空间上作适当积分得到，是\hyperref[sec:general-field-theories]{附录~\ref*{sec:general-field-theories}}所述驻相法的一个例子。非常值全纯映射 $f:\Sigma\to X$ 称为\emph{世界面瞬子}，上述 $n$ 点函数的形式称为\emph{瞬子和}。

对 B 模型，场 $\eta^{\bar i}=\psi_+^{\bar i}+\psi_-^{\bar i}$ 按 $f^*(T_X^{0,1})$ 的截面变换；类似地，$\theta_i=g_{i\bar i}(\psi_+^{\bar i}-\psi_-^{\bar i})$ 按 $(1,0)$ 形式变换。B 模型也有局部算子，来自取值于 $\wedge^pT_X^{1,0}$ 的 $(0,q)$ 形式。若将这样的形式 $\theta$ 局部写成
\SourcePage{446}
\[
\theta=h^{j_1\cdots j_p}_{\bar i_1\cdots\bar i_q}\,
dx^{\bar i_1}\wedge\cdots\wedge dx^{\bar i_q}
\otimes\frac{\partial}{\partial x^{j_1}}\wedge\cdots\wedge
\frac{\partial}{\partial x^{j_p}},
\]
则局部算子 $\mathcal O_\theta$ 定义为
\[
\mathcal O_\theta=h^{j_1\cdots j_p}_{\bar i_1\cdots\bar i_q}
\eta^{\bar i_1}\cdots\eta^{\bar i_q}\theta_{j_1}\cdots\theta_{j_p}.
\]
与 A 模型一样，考虑运动方程后，作用量有简化形式。此理论的 BRST 算子仍记为 $Q$，满足 $\{Q,\mathcal O_\theta\}=-\mathcal O_{\bar\partial\theta}$。因此，$\theta\mapsto\mathcal O_\theta$ 将 Dolbeault 上同调 $H^q(X,\wedge^p(T_X^{1,0}))$ 映到 B 模型的 BRST 上同调。

B 模型的关联函数，可用刚提及的作用量简化形式计算得到。原则上，此关联函数可能包含瞬子和，但某些\emph{\IndexLocation{indexloc:0944:446}非重整化定理}断言并非如此。另一种论证借助\IndexLocation{indexloc:0830:446}镜像对称：A 模型的关联函数不依赖复模，因此镜像流形的 B \IndexLocation{indexloc:0184:446}\IndexLocation{indexloc:0189:446}模型关联函数不依赖 Kähler 模；然而，瞬子和中每个 $\gamma\ne0$ 项，都因含 $q^\gamma$ 而显式依赖 Kähler 模。\hyperref[sec:three-point]{第~\ref*{sec:three-point} 节}给出了\IndexLocation{indexloc:0623:446}瞬子修正消失的另一种启发式论证。

最后，为使我们的 $\sigma$ 模型避免某些\IndexLocation{indexloc:1110:446}量子反常，$(X,g)$ 实际上须为 Calabi--Yau 三维簇。要求 $X$ 的实维数为六的原因，将在\hyperref[sec:conformal-field-theories]{附录~\ref*{sec:conformal-field-theories}}解释；Kähler 条件则是 $N=2$ 超对称的必要条件。\hyperref[sec:calabi-yau-physics]{第~\ref*{sec:calabi-yau-physics} 节}简述过的 Calabi--Yau 条件更为微妙，需要转入量子场论重整化的讨论，超出了这些附录的范围。更多细节见 \SourceCite{Huebsch}，超对称 $\sigma$ 模型的讨论也见 \SourceCite{GO}。
\index{非线性sigma模型@非线性 $\sigma$ 模型|hyperpage}\index{BRST算子@BRST 算子|hyperpage}\index{BRST上同调@BRST 上同调|hyperpage}\index{世界面瞬子|hyperpage}\index{瞬子和|hyperpage}

\section{共形场论}\label{sec:conformal-field-theories}
粗略地说，共形理论是定义在 Riemann 曲面上的场的理论，它尊重曲面的共形变换。下面尝试把这一说法讲得更精确一些。详细介绍见 \SourceCite{Ginsparg}。它与弦论的联系在于，这些 Riemann 曲面正是弦随时间传播时扫出的曲面。因此，我们将这种 Riemann 曲面称为\emph{世界面}。

Riemann 曲面的\emph{共形映射}是保持角度和定向的自映射；这样的映射是全纯的。我们希望在带复坐标 $z=x+iy$ 的二维圆盘 $\Delta$ 上，局部研究共形映射。遗憾的是，局部共形变换并不构成群，因为邻域未必保持不变，映射因而未必能够复合。不过，\IndexLocation{indexloc:0161:446}共形代数是良好定义的：无穷小共形变换在底层实曲面上诱导一个实向量场，在局部全纯坐标中可写成 $f(z)\partial/\partial z+\overline{f(z)}\partial/\partial\bar z$，其中 $f(z)$ 全纯。通常选取基 $\ell_n=-z^{n+1}\partial_z$ 与 $\bar\ell_n=-\bar z^{n+1}\partial_{\bar z}$，将全纯部分和反全纯部分分开；它们满足 Virasoro 代数关系
\[
[\ell_n,\ell_m]=(n-m)\ell_{n+m},\qquad
[\bar\ell_n,\bar\ell_m]=(n-m)\bar\ell_{n+m}.
\]
\SourcePage{447}
共形代数可由 Virasoro 代数恢复为由 $\ell_n+\bar\ell_n$ 和 $i(\ell_n-\bar\ell_n)$ 生成的子代数。共形场论包含 Virasoro 代数的一个表示；不过必须谨慎，因为某些表示会出现量子反常。

接下来定义共形场论中我们关心的场。
\begin{definition}\label{def:primary-field}
\emph{权为 $(h,\bar h)$ 的初级场}是一个场 $\Phi(z,\bar z)$，在共形变换 $z\mapsto f(z)$ 下按
\[
\Phi(z,\bar z)\longmapsto
\left(\frac{\partial f}{\partial z}\right)^h
\left(\frac{\partial\bar f}{\partial\bar z}\right)^{\bar h}
\Phi(f(z),\bar f(\bar z))
\]
变换。
\end{definition}
换言之，$\Phi(z,\bar z)$ 为初级场时，表达式 $\Phi(z,\bar z)\,dz^h\otimes d\bar z^{\bar h}$ 不变。

共形场论的规范群是局部共形群。如\hyperref[sec:general-field-theories]{附录~\ref*{sec:general-field-theories}}所述，这意味着必须作规范选择。在局部坐标 $z=x+iy$ 中，可通过把世界面上的度量 $h$ 固定为 $h=dx^2+dy^2$ 来确定规范。但与该节一样，仍须施加 Euler--Lagrange 方程给出的条件 $\delta S/\delta h^{ij}=0$。

遗憾的是，量子化并不像电磁场的情形那样直接。例如，为保持共形不变性，在量子化过程中须引入称为\emph{鬼场}的辅助场；超对称性还需要额外的鬼场。另外，虽然 Euler--Lagrange 方程 $\delta S/\delta h^{ij}=0$ 在经典层面成立，但在量子层面，其表现略有不同。为解释这一点，引入分量为
\[
T_{ij}=\frac1{\sqrt h}\frac{\delta S}{\delta h^{ij}}
\]
的\emph{能动张量}；这是一个对称张量。这里，照常用 $h$ 表示矩阵 $h_{ij}$ 的行列式，$h^{ij}$ 表示逆矩阵。在全纯坐标中，能动张量的分量为 $T(z)=T_{zz}$、$\bar T(\bar z)=T_{\bar z\bar z}$ 和 $T_{z\bar z}$。类似于式~\eqref{eq:noether-conserved-current} 中的处理，由任意正函数 $\phi$ 给出的共形缩放 $h\mapsto\phi h$ 的不变性，导出 $T_{z\bar z}=0$。再由能动量守恒，得到 $T(z)$ 全纯而 $\bar T(\bar z)$ 反全纯，这也解释了上述记号。因此，经典运动方程，即 Euler--Lagrange 方程，还包括约束 $T(z)=\bar T(\bar z)=0$。

在量子理论中，用\emph{\IndexLocation{indexloc:1108:447}径向量子化}来修改这些方程。这意味着，在局部坐标中，将等时曲线选为同心圆，而固定空间坐标的曲线则为径向线。弦的世界面局部包含一个空间坐标和一个时间坐标，但空间与时间的区分并非不变量。形式运算以及 $T(z)$ 为守恒量这一事实给出：若 $\Phi(z,\bar z)$ 是权为 $(h,\bar h)$ 的初级场，则有如下算子恒等式，称为\emph{\IndexLocation{indexloc:0989:447}算子乘积展开}：
\[
T(z)\Phi(w,\bar w)
=\frac h{(z-w)^2}\Phi(w,\bar w)
+\frac{\partial_w\Phi(w,\bar w)}{z-w}+R(z,w,\bar w),
\]
其中 $R(z,w,\bar w)$ 在 $z=w$ 时保持有限。$\bar T(\bar z)\Phi(w,\bar w)$ 有类似公式。$T$ 的这一变换性质似乎说明它应是权为 $(2,0)$ 的初级场。遗憾的是，若真如此，就会产生
\SourcePage{448}
量子反常。因此，实际公式在算子乘积展开中需要多一项：
\begin{equation}\label{eq:stress-tensor-ope}
T(z)T(w)\sim\frac{c/2}{(z-w)^4}
+\frac{2T(w)}{(z-w)^2}+\frac{\partial_wT(w)}{z-w},
\end{equation}
其中 $c$ 为称作\emph{中心荷}的常数，符号 $\sim$ 表示只写出奇异项。展开 $T(z)=\sum_nL_nz^{-(n+2)}$，代入式~\eqref{eq:stress-tensor-ope}，通过围道积分得到
\begin{equation}\label{eq:virasoro-central-extension}
[L_n,L_m]=(n-m)L_{n+m}+\frac c{12}(n^3-n)\delta_{n+m,0},
\end{equation}
其中 $n+m=0$ 时，$\delta_{n+m,0}$ 为恒等算子；否则为零。也就是说，在量子层面得到的是 Virasoro 代数的中心扩张的表示。还应提及，式~\eqref{eq:virasoro-central-extension} 也出现在\hyperref[subsec:virasoro-conjecture]{第~\ref*{subsec:virasoro-conjecture} 小节}关于 Virasoro 猜想的讨论中。

为使微扰弦论自洽，需要一个没有反常的 Virasoro 代数本身的表示。由式~\eqref{eq:virasoro-central-extension}，只要理论的总中心荷 $c$ 为零，就不会出现中心扩张。在实维数为 $2k$ 的流形 $X$ 上的 $\sigma$ 模型所给出的共形场论中，每个玻色子 $x^i$ 对 $c$ 的贡献为 $+1$，而 $2k$ 对费米子 $\psi_+^i,\psi_-^i$ 中，每一对的贡献为 $+1/2$，总贡献为 $2k\cdot1+2k\cdot\tfrac12=3k$。不过，已经看到，量子化需要鬼场，才能保持世界面重新参数化，即共形变换，以及超对称变换的不变性。补偿世界面重新参数化的鬼场对 $c$ 的贡献为 $-26$，补偿超对称性的鬼场贡献为 $+11$。方程 $3k-26+11=0$ 仅在 $2k=10$ 时满足。其中四个维度视为我们熟悉的世界的时空维度，例如常用四维 Minkowski 空间作为模型，剩下六个实维度。这就是将讨论限于复维数为三的 Calabi--Yau 流形的原因。

类似而更简单的论证表明，自洽的玻色弦必须在 $26$ 维空间中传播。

最后，对 $\sigma$ 模型 $(X,g)$ 给出的 $N=2$ 超共形代数表示作几点说明。这个表示包含上述算子 $T(z)$、$\bar T(\bar z)$。此外，全纯超对称性给出一个 $U(1)=SO(2)$ 对称性，自然作用于 $\sigma$ 模型的费米子
\[
\begin{pmatrix}\psi_+^i\\\psi_-^{\bar i}\end{pmatrix}.
\]
这个 $U(1)$ 对称性所产生的守恒流记为 $J(z)$。$J$ 的空间积分就是\hyperref[sec:physics]{第~\ref*{sec:physics} 节}中出现的算子 $Q$；它是理论中右行（即全纯）部分的生成元，对两种超对称性作无穷小旋转。同样，由反全纯超对称性及 $U(1)$ 在
\[
\begin{pmatrix}\psi_+^{\bar i}\\\psi_-^i\end{pmatrix}
\]
上的作用，得到 $\bar J$，进而得到左行，即反全纯部分的算子 $\bar Q$。上面已经说明，$H^1(X,T_X)$ 和 $H^1(X,\Omega_X^1)$ 如何成为理论中的算子；将它们作用于真空态，便得到对应于这两个空间的量子态。这些态是算子 $Q,\bar Q$ 的特征向量，
\SourcePage{449}
如\hyperref[sec:physics]{第~\ref*{sec:physics} 节}所述。镜像对称的一项主要动机，正是改变生成元 $\bar Q$ 的符号对相应特征值所产生的影响。
\index{共形场论|hyperpage}\index{初级场|hyperpage}\index{鬼场|hyperpage}\index{能动张量|hyperpage}\index{径向量子化|hyperpage}\index{算子乘积展开|hyperpage}\index{中心荷|hyperpage}\index{Virasoro代数@Virasoro 代数|hyperpage}

\section{Landau--Ginzburg 模型}\label{sec:landau-ginzburg-models}
这里描述的 Landau--Ginzburg（LG）理论也是超对称理论。为明确起见，将讨论限制于具有 $(2,2)$ 超对称性的理论。这些理论可用超空间和超场的形式体系描述。下面暂且离开主线，介绍这一体系。

\emph{超流形}是带有局部坐标
\[
(x^1,\ldots,x^n,y^1,\ldots,y^m)
\]
的拓扑空间 $M$，其中变量 $x^i$ 彼此对易，也与 $y^j$ 对易，而 $y^i$ 彼此反对易。虽然可以给出精确的数学定义，见 \SourceCite{Manin1}，但上述粗略定义已足够用于目前的讨论。可以将 $M$ 看成一个带局部坐标 $(x^1,\ldots,x^n)$ 的流形，再附加反对易变量 $y^i$。\emph{超场}就是超空间上的函数。超场可唯一写成和 $\sum_I f_I(x)y^I$，其中 $f_I(x)$ 是 $x^1,\ldots,x^n$ 的函数，$I=\{i_1,\ldots,i_k\}\subset\{1,\ldots,m\}$ 为指标集；若按 $i_1<\cdots<i_k$ 排列，则 $y^I=y^{i_1}\cdots y^{i_k}$。\hyperref[subsec:cohomology-supermanifold]{第~\ref*{subsec:cohomology-supermanifold} 小节}已经说明，$H^*(X,\mathbb C)$ 具有自然的复超流形结构。

这里的底层流形是世界面 $\Sigma$，故取 $n=2$。为了得到 $(2,2)$ 超对称性，还须取 $m=4$。以局部复坐标 $z$ 取代变量 $x^1,x^2$，并将 $y^i$ 记为 $\theta^i,\bar\theta^i$，其中 $i=1,2$。$\theta^i,\bar\theta^i$ 不只是局部表达式，它们还是\emph{旋量}，即按二维旋转群的双覆盖 $\operatorname{Spin}(2)$ 的某个一维旋量表示变换；$\theta^1,\theta^2$ 属于其中一个表示，$\bar\theta^1,\bar\theta^2$ 属于另一个表示。\IndexLocation{indexloc:1267:449}旋量几何的更多数学细节见 \SourceCite{LM}。

因此，超场 $\Phi$ 可在局部表示为
\[
\Phi=f(z)+\phi_i(z)\theta^i+\bar\phi_i(z)\bar\theta^i
+F_{ij}(z)\theta^i\theta^j+\cdots,
\]
其中省写了反对易变量的部分二次项和所有更高次项。超对称变换 $Q^i,\bar Q^i$ 可以显式写出，其反对易关系形式简单；四维超对称的类似公式可参见 \SourceCite{WB}。由此可以构造 $N=2$ 超对称代数。注意，不要将这些超对称变换与前面讨论过的 $U(1)$ 荷 $Q,\bar Q$ 混淆。

所有超场组成的空间给出 $N=2$ 超对称代数的一个表示。\emph{手征超场}构成其不变子空间。具体构造是，在超空间中显式写出与超对称算子分次对易的微分算子 $\bar D_i$，并将手征超场定义为满足 $\bar D_i\Phi=0$、$i=1,2$ 的超场 $\Phi$；手征超场的定义例如见 \SourceCite{WB}。算子 $\bar D_i$ 涉及对反对易变量的微分。事实上，$\bar D_i$ 包含一项 $\partial/\partial\bar\theta^i$，可按自然的方式定义。因此，手征超场的行为与普通空间中的全纯函数颇为相似。

Landau--Ginzburg 理论是某些超对称理论，可简单地用\emph{超势}描述。超势是二维手征超场的\IndexLocation{indexloc:1431:449}加权齐次多项式
\SourcePage{450}
$F(\Phi_1,\ldots,\Phi_k)$。这里的加权齐次，是指为 $\Phi_i$ 指定权 $d_i$。在我们关心的情形中，选择 $F$ 的加权次数为 $d=\sum_i d_i$。

Landau--Ginzburg 理论的作用量中的重要项为
\begin{equation}\label{eq:landau-ginzburg-superpotential-term}
\int d^2z\,d^2\theta\,F(\Phi_i),
\end{equation}
其中 $z$ 是世界面上的局部全纯坐标，$\Phi_i$ 为手征超场，$\theta^j$ 为费米超空间坐标。将 $\Phi_i$ 的标量玻色分量记为 $\phi_i$；经典真空的方程给出 $\partial F/\partial\phi_i=0$。显然有经典解 $\phi_i=0$，在这里考虑的 Landau--Ginzburg 理论中，这是唯一的解。

转到量子理论，与唯一的经典真空相联系的是唯一的真空态。将真空态记为 $V$，对任意多项式 $P$，都可构造新态 $P(\phi_i)(V)$。这些手征态组成的空间与 Jacobian 环 $\mathbb C[\phi_1,\ldots,\phi_k]/(\partial F/\partial\phi_i)$ 同构；此环也称为\emph{手征环}。

注意上一段中，我们很快从听起来相当陌生的对象，如费米超空间坐标等，转到了代数几何学家熟悉的对象，即加权齐次多项式的 Jacobian 环。方程 $\partial F/\partial\phi_i=0$ 只有原点这一解，等价于加权射影空间中的超曲面 $F=0$ 为拟光滑。这是标准数学对象突然出现在这些物理理论中的一个典型例子。

实际感兴趣的是相关的理论，即 \emph{Landau--Ginzburg 轨形}。在我们关心的情形中，令 $\eta$ 为本原 $d$ 次单位根，并对 $(\eta^{d_1},\ldots,\eta^{d_k})$ 的作用取轨形。这意味着，不仅考虑在世界面 $\Sigma$ 上单值定义的超场，还考虑弦两端的场值按这一作用的第 $i$ 次幂变换的超场。所得态称为属于\emph{第 $i$ 个扭曲部分}，并产生扭曲真空 $V_i$。再形成 $P(\phi_j)(V_i)$ 这样的表达式，就能构造更多的态。Landau--Ginzburg 轨形包含其中在 $\eta$ 的作用下不变的态。

下面将看到，Calabi--Yau 流形与 Landau--Ginzburg 轨形密切相关。现在先从另一个角度举例说明。
\SourceExampleNumber{B.4}{1}
\begin{example}\label{ex:lg-octic-twisted-sector}
设 $\Phi_1,\ldots,\Phi_5$ 的权依次为 $1,1,2,2,2$，令 $F$ 为这些变量中加权次数为 $8$ 的齐次多项式。在计算这个理论中的某些态之前，先注意，此例与 $\mathbb P^4(1,1,2,2,2)$ 中的八次超曲面有关；\SourceCite{CdFKM,HKTY1} 研究过这些超曲面。这个加权射影空间沿 $x_1=x_2=0$ 定义的 $\mathbb P^2$ 有 $A_1$ 奇点。将它爆破，例外除子是 $\mathbb P^2$ 上的一个直纹三维簇。限制到超曲面后，例外除子成为一条四次曲线（即亏格为三的曲线）上的直纹曲面。因此，纤维族的 Abel--Jacobi 映射的像是 $H^{2,1}$ 的一个三维子空间。可以证明，这些元素对应的 Calabi--Yau 三维簇形变，不能实现为 $\mathbb P^4(1,1,2,2,2)$ 中的超曲面。这个三维子空间，与由超曲面形变组成的余维为三的子空间 $H^{2,1}_{\mathrm{poly}}(V^\circ)\subset H^{2,1}(V^\circ)$ 互补。例题~\ref{ex:weighted-extremal-transition} 给出了这些形变的几何描述。

\SourcePage{451}
现在回到 LG 理论。我们有态 $P(\phi_i)(V)$。理论的模参数所对应的态，恰是 $P$ 的次数为 $8$ 的那些态。因此得到 Jacobian 环的八次部分，其维数为 $83$。这是代数几何中熟悉的构造。

那么，缺少的三个形变在哪里？它们出现在某个扭曲部分中，确切地说，是第四个扭曲部分。扭曲真空 $V_4$ 不是不变的；要使 $P(\phi_i)(V_4)$ 对应于形变，$P$ 必须具有权 $2$。变量 $\phi_3,\phi_4,\phi_5$ 在 $\eta^4$ 的作用下不变，而 $\phi_1,\phi_2$ 则不在该作用下保持不变。因此，这些态就是 $P(\phi_3,\phi_4,\phi_5)(V_4)$，组成三维空间，恰如预期。
\end{example}
上述讨论基于 \SourceCite{Vafa1} 中的计算；该文还从 Landau--Ginzburg 理论出发，给出了令人惊叹的 Calabi--Yau 流形上同调公式。
\index{Landau--Ginzburg模型@Landau--Ginzburg 模型|hyperpage}\index{超空间|hyperpage}\index{超场|hyperpage}\index{手征超场|hyperpage}\index{超势|hyperpage}\index{手征环|hyperpage}\index{Landau--Ginzburg轨形@Landau--Ginzburg 轨形|hyperpage}\index{扭曲部分|hyperpage}

\section{规范线性 \texorpdfstring{$\sigma$}{σ} 模型}\label{sec:gauged-linear-sigma-models}
Witten 的规范线性 $\sigma$ 模型（GLSM）十分有用；在适当意义下，它在 Calabi--Yau 模型与 Landau--Ginzburg 理论之间作“插值”。与 Landau--Ginzburg 模型一样，这一理论可用\hyperref[sec:landau-ginzburg-models]{附录~\ref*{sec:landau-ginzburg-models}}概述的手征超场来表述。更多细节见 \SourceCite{Witten5}。GLSM 与\hyperref[sec:nonlinear-sigma-models]{附录~\ref*{sec:nonlinear-sigma-models}}讨论的非线性 $\sigma$ 模型之间的数学比较，在\hyperref[chap:mirror-theorem]{第~\ref*{chap:mirror-theorem} 章}的镜像定理中起关键作用。

GLSM 是规范群为某个 $U(1)^s$ 的规范理论。与规范群相联系的是二维\emph{向量超场} $V_a$，$a=1,\ldots,s$。理论还包含 $k+r$ 个手征超场 $\Phi_i$、$i=1,\ldots,k$ 和 $P_j$、$j=1,\ldots,r$；其荷分别记为 $Q_{i,a}$ 和 $-q_{j,a}$，其中每个 $a$ 都有 $\sum_jq_{j,a}=\sum_iQ_{i,a}$。荷说明超场如何在对称群 $U(1)_a$，即 $U(1)^s$ 的第 $a$ 个因子的作用下变换。理论还由超势 $W$ 决定；我们将看到，超势把 Calabi--Yau 超曲面的方程与 Landau--Ginzburg 势联系起来。超势是超场的全纯函数，在 $U(1)^s$ 的作用下不变。这里限于形如 $W=\sum_{j=1}^rP_jG_j(\Phi_1,\ldots,\Phi_k)$ 的超势。

作用量本身是四项之和：
\begin{equation}\label{eq:glsm-action-decomposition}
S=S_{\mathrm{kin}}+S_W+S_{\mathrm{gauge}}+S_{D,\theta},
\end{equation}
具体形式无须关注。这里只强调，作用量是所有超场的函数。作用量中有几个常数，衡量理论中某些耦合的强度：$e_j$ 为 $U(1)_j$ 的规范耦合常数，$r_j$ 决定某种超对称相互作用的强度。这些常数将在下面的式~\eqref{eq:glsm-bosonic-potential} 和~\eqref{eq:glsm-d-term} 中显式出现。

暂且作简化假设 $r=1$，并以 $P$ 代替 $P_1$。因此超势为 $W=PG$，其中 $G$ 为某个多项式。本节稍后再回到 $r>1$ 的情形。将作用量中的项分成对易的\IndexLocation{indexloc:0084:451}玻色变量部分和反对易的\IndexLocation{indexloc:0313:451}费米变量部分，并考虑运动方程后，\IndexLocation{indexloc:0356:451}玻色部分的势为
\begin{equation}\label{eq:glsm-bosonic-potential}
U(\phi_i,p)=\sum_j\frac1{2e_j^2}D_j^2+\sum_\alpha|F_\alpha|^2,
\end{equation}
\SourcePage{452}
其中
\begin{equation}\label{eq:glsm-d-term}
D_j=-e_j^2\left(\sum_iQ_{i,j}|\phi_i|^2-q_j|p|^2-r_j\right),
\end{equation}
而
\begin{equation}\label{eq:glsm-f-term}
\sum_\alpha|F_\alpha|^2=|G|^2+|p|^2\sum_i\left|\frac{\partial G}{\partial\phi_i}\right|^2.
\end{equation}
在式~\eqref{eq:glsm-d-term} 和~\eqref{eq:glsm-f-term} 中，$\phi_i,p$ 是超场 $\Phi_i,P$ 的某些标量玻色分量（顶分量）。换言之，可以将 $\phi_i,p$ 看成普通函数。

我们将看到，$r_j$ 取值不同时，经典运动方程的解十分不同。Witten 将这些过渡称为理论中的\emph{相变}，见 \SourceCite{Witten5}。
\SourceExampleNumber{B.5}{1}
\begin{example}\label{ex:glsm-quintic-phases}
设 $k=5$、$s=1$，每个 $i$ 都有 $Q_{i,1}=1$。选择超势 $W=PG(\Phi_1,\ldots,\Phi_5)$，其中 $G$ 是 $\Phi_i$ 的一般五次多项式。为使势最小，式~\eqref{eq:glsm-d-term} 的每一项与式~\eqref{eq:glsm-f-term} 右端的每一项都必须为零。注意，式~\eqref{eq:glsm-d-term} 为零意味着 $\sum_{i=1}^5|\phi_i|^2-5|p|^2-r_1=0$。须考虑两种情形：
\begin{itemize}
\item $r_1>0$。此时，式~\eqref{eq:glsm-d-term} 为零表明 $(\phi_1,\ldots,\phi_5)\ne0$。再由 $|p|^2\sum_i|(\partial G/\partial\phi_i)|^2=0$，得到 $p=0$，因为 $(\phi_1,\ldots,\phi_5)\ne0$，且 $G$ 一般，即定义光滑超曲面。式~\eqref{eq:glsm-d-term} 为零现在意味着 $\sum_i|\phi_i|^2=r_1$。接着考虑规范对称性，对 $U(1)$ 取商；根据 $Q_i$ 的选择，$U(1)$ 通过标量乘法作用于 $(\phi_1,\ldots,\phi_5)$，所得空间为 $\mathbb P^4$。最后，由式~\eqref{eq:glsm-f-term} 的右端还须令 $G=0$。因此，经典解的空间就是五次三维簇。
\item $r_1<0$。此时，式~\eqref{eq:glsm-d-term} 为零表明 $p\ne0$。由式~\eqref{eq:glsm-f-term} 右端各项为零，首先得到每个 $\partial G/\partial\phi_i$ 为零。因 $G$ 一般，这意味着 $(\phi_1,\ldots,\phi_5)=(0,\ldots,0)$。再由式~\eqref{eq:glsm-d-term} 为零，得到 $|p|=\sqrt{-r_1/5}$。$p$ 本身只在规范变换 $p\mapsto u^{-5}p$、$u\in U(1)$ 的意义下确定。注意，$u$ 为五次单位根时，$p$ 保持不变。因此得到一个 Landau--Ginzburg 轨形。
\end{itemize}
所以，GLSM 在 Calabi--Yau 理论与 Landau--Ginzburg 理论之间作插值。
\end{example}

更一般地，可用环面几何描述 GLSM。$U(1)^s$ 在 $\Phi_i$ 上的作用给出在坐标为 $\phi_1,\ldots,\phi_k$ 的 $\mathbb C^k$ 上的作用。用\hyperref[subsec:toric-symplectic]{第~\ref*{subsec:toric-symplectic} 小节}的术语说，这个作用是 Hamilton 作用，其矩映射 $\mu:\mathbb C^k\to\mathbb R^s$ 为
\[
\mu(\phi_1,\ldots,\phi_k)
=\frac12\left(\sum_iQ_{i,1}|\phi_i|^2,\ldots,\sum_iQ_{i,s}|\phi_i|^2\right).
\]
可以证明，GLSM 的 Calabi--Yau 相，粗略地对应于 $\vec r=(r_1,\ldots,r_s)$ 位于 $\mathbb R^s$ 的某个锥中；该锥是环面簇
\[
X_{\vec r}=\mu^{-1}(\vec r/2)/U(1)^s
\]
的 Kähler 锥。
\SourcePage{453}
这里使用的是\hyperref[subsec:toric-reduction]{第~\ref*{subsec:toric-reduction} 小节}讨论的环面簇的辛构造。这种构造方法是 GLSM 的关键成分。此外，$G\in\mathbb C[\phi_1,\ldots,\phi_k]$ 是\hyperref[subsec:homogeneous-ring]{第~\ref*{subsec:homogeneous-ring} 小节}意义下的齐次多项式，因此 $G=0$ 给出 $X_{\vec r}$ 中的一个 Calabi--Yau 超曲面。所以，随着 $\vec r$ 改变，GLSM 的 Calabi--Yau 相对应于\hyperref[sec:gkz-decomposition]{第~\ref*{sec:gkz-decomposition} 节}中 GKZ 分解的某些锥。不过，理论还有其他相，如 Landau--Ginzburg 相，来自 GKZ 分解之外的锥。它们对应于 \SourceCite{MP1} 中描述的\emph{扩大 GKZ 分解}或\emph{扩大次级扇}中的锥。

到目前为止，讨论中一直假设 $r=1$。当 $r>1$ 时，超势为 $W=\sum_aP_aG_a$，其中 $G_a$ 仍为环面意义下的齐次多项式。这样的 GLSM 的一个 Calabi--Yau 相仍给出环面簇，但现在要考虑 Calabi--Yau 完全交 $G_1=\cdots=G_r=0$。

线性 $\sigma$ 模型与非线性 $\sigma$ 模型有一项重要差别。为明确起见，考虑上面的例子，假定 $r_1>0$，即处于 Calabi--Yau 相。在非线性 $\sigma$ 模型中，场是 $\mathbb P^4$ 上的局部坐标；在线性模型中，场是 $\mathbb C^5$ 上的坐标，再对标量作用取商。因此，对非线性 $\sigma$ 模型，须考虑全纯瞬子，即亏格为零的曲线 $C$ 到 $\mathbb P^4$ 的次数为 $d$、像落在 $G=0$ 上的全纯映射。对 GLSM，则使用五元组 $(f_1(s,t),\ldots,f_5(s,t))$，其中 $f_i$ 是固定次数 $d$ 的齐次多项式，满足 $G(f_1,\ldots,f_5)=0$；这里将 $(s,t)$ 视为 $\mathbb P^1$ 的齐次坐标。但这并非精确的一一对应：五元组 $(f_1(s,t),\ldots,f_5(s,t))$ 未必来自一个次数为 $d$ 的全纯映射，因为 $f_i$ 可能有公因子；若尝试将其解释为映射，在去掉基点后次数会降低。这些五元组，或称\emph{朴素瞬子}，更容易处理；在 GLSM 的框架中，它们已被用来为\IndexLocation{indexloc:0360:453}\IndexLocation{indexloc:0836:453}镜像对称提供有力支持，见 \SourceCite{MP1}。\hyperref[sec:quintic-mirror-theorem]{第~\ref*{sec:quintic-mirror-theorem} 节}定义的映射 $\varphi:M_d\to N_d$，可描述为全纯瞬子与朴素瞬子之间的比较。
\index{规范线性sigma模型@规范线性 $\sigma$ 模型|hyperpage}\index{相变|hyperpage}\index{扩大GKZ分解@扩大 GKZ 分解|hyperpage}\index{扩大次级扇|hyperpage}\index{朴素瞬子|hyperpage}

\section{拓扑量子场论}\label{sec:topological-qft}
\hyperref[sec:general-field-theories]{附录~\ref*{sec:general-field-theories}}简要介绍了一般量子场论。这种理论由态的 Hilbert 空间 $\mathcal H$ 及其上的某些算子描述，算子编码粒子的传播和相互作用。现在解释什么是\emph{拓扑量子场论}（TQFT）。我们将看到，Frobenius 代数和\IndexLocation{indexloc:1111:453}量子上同调自然地出现在这一背景中。

先讨论点粒子理论。令 $\mathcal H$ 为态空间，$\mathsf H$ 为通过量子化过程与 Hamilton 函数 $H$ 对应的算子。如式~\eqref{eq:quantum-unitary-time-evolution} 所述，态随时间的传播由酉算子 $\exp(-i\mathsf Ht)$ 给出。可以用长度为 $t$ 的区间示意这一传播过程：将 Hilbert 空间 $\mathcal H$ 与区间的每个端点相联系，将算子
\[
e^{-i\mathsf Ht}:\mathcal H\longrightarrow\mathcal H
\]
与区间相联系。这是一个 $(0+1)$ 维量子场论；$0$ 表示与 $\mathcal H$ 对应的端点的维数，$+1$ 表示连接这些端点所需要的额外时间维度。

\SourcePage{454}
为过渡到 TQFT，考虑哪些 $(0+1)$ 维理论的算子 $\exp(-i\mathsf Ht)$ 只依赖区间的拓扑。这个问题容易确定，因为我们要问的就是，$\exp(-i\mathsf Ht)$ 何时与 $t$ 无关。在 $t=0$ 处求导，得到必须有 $\mathsf H=0$。因此，任意 $t$ 都有 $\exp(-i\mathsf Ht)$ 为恒等算子。换言之，这个系统完全没有动力学，作为物理理论完全平凡，毫无趣味。

这些理论就是 $(0+1)$ 维拓扑量子场论。将会看到，当 $d>0$ 时，$(d+1)$ 维 TQFT 远非平凡。

现在给出拓扑量子场论的公理化定义。为明确起见，使用复系数。若要包含费米子，须使用超向量空间，即满足自然符号规则的 $\mathbb Z/2\mathbb Z$ 分次向量空间，见 \SourceCite{Manin1}。必要的修改留给感兴趣的读者。更广泛的论述见 \SourceCite{Quinn}。
\begin{definition}\label{def:topological-qft}
一个 $(d+1)$ 维\emph{拓扑量子场论}由下列资料组成：
\begin{itemize}
\item 对每个闭的定向 $d$ 维流形 $Y$，指定一个有限维复向量空间 $Z(Y)$。这一向量空间在 $Y$ 的同构下按函子方式变换。
\item 对每个定向 $(d+1)$ 维流形 $X$，其边界 $\partial X$ 是闭的定向 $d$ 维流形，指定一个元素 $Z_X\in Z(\partial X)$。这一元素在 $X$ 的同构下按函子方式变换。
\end{itemize}
这些资料还满足下列公理：
\begin{enumerate}[label=\textbf{公理 \arabic*：},leftmargin=4.7em]
\item $Z(Y_1\coprod Y_2)=Z(Y_1)\otimes Z(Y_2)$。
\item 将空集视为闭的定向 $d$ 维流形时，有 $Z(\varnothing)=\mathbb C$。
\item 将空集视为边界为空的定向 $(d+1)$ 维流形时，有 $Z_{\varnothing}=1\in Z(\varnothing)=\mathbb C$。
\item 以 $\bar Y$ 表示 $Y$ 配备相反定向所得的流形，则 $Z(\bar Y)\simeq Z(Y)^*$，此同构在 $Y$ 的同构下按函子方式变换。特别地，若 $\partial X=(\coprod_{i=1}^k\bar Y_i)\coprod(\coprod_{j=1}^{\ell}Y'_j)$，则可将 $Z_X$ 视为
\[
Z_X\in\operatorname{Hom}_{\mathbb C}
\left(\bigotimes_{i=1}^kZ(Y_i),\bigotimes_{j=1}^{\ell}Z(Y'_j)\right),
\]
这里使用了本条公理与公理 1。还假设，在自然的识别下，$Z_{\bar X}$ 是 $Z_X$ 的对偶映射。
\item 以 $I$ 表示定向区间 $[0,1]$。令 $Y$ 为任意闭的、无边界定向 $d$ 维流形，并置 $X=Y\times I$，则 $\partial X=\bar Y\coprod Y$。要求
\[
Z_X=1_{Z(Y)}\in\operatorname{Hom}_{\mathbb C}(Z(Y),Z(Y)),
\]
其中使用了公理 4。
\item 若 $\partial X=\bar Y_1\coprod Y_2$ 且 $\partial X'=\bar Y_2\coprod Y_3$，则可沿 $Y_2$ 将 $X$ 与 $X'$ 粘合，得到边界为 $\bar Y_1\coprod Y_3$ 的新流形 $X\cup_{Y_2}X'$。由公理 4，可写成
\[
\begin{aligned}
Z(\partial X)&=\operatorname{Hom}_{\mathbb C}(Z(Y_1),Z(Y_2)),\\
Z(\partial X')&=\operatorname{Hom}_{\mathbb C}(Z(Y_2),Z(Y_3)),\\
Z(\partial(X\cup_{Y_2}X'))&=\operatorname{Hom}_{\mathbb C}(Z(Y_1),Z(Y_3)).
\end{aligned}
\]
要求 $Z_{X\cup_{Y_2}X'}=Z_{X'}\circ Z_X$。
\end{enumerate}
\end{definition}
当 $d=0$ 时，这些公理很容易理解。可置 $\mathcal H=Z(\text{点})$，这是 $\mathbb C$ 上的有限维向量空间。将公理 5 用于
\begingroup\linespread{1.14}\selectfont
\SourcePage{455}
$Y=\text{点}$，得到 $Z_I=1_{\mathcal H}$；这正是本节开头讨论 $(0+1)$ 维 TQFT 所得到的结论。此外，$S^1$ 是边界为空的定向一维流形，所以 $Z(S^1)\in Z(\varnothing)=\mathbb C$。可以如下计算这个数：将 $S^1$ 实现为两个单位区间 $I$ 的并。
\begin{figure}[H]
\centering\includegraphics[width=0.125523\textwidth]{b6-circle.png}
\end{figure}
现在使用公理 6，取
\[
X=I,\quad Y_1=\varnothing,\quad Y_2=\overline{\text{点}}\coprod\text{点},
\quad X'=\bar I,\quad Y_3=\varnothing.
\]
结合公理 2、4、5，算得 $Z(S^1)=\dim\mathcal H$。由于任意紧致一维流形都是 $S^1,\bar S^1,I,\bar I$ 的不交并，上述公理说明，$(0+1)$ 维 TQFT 由 $\mathcal H=Z(\text{点})$ 唯一确定。还可证明，所有有限维复向量空间 $\mathcal H$ 都能以这种方式来自 $(0+1)$ 维 TQFT。

现在转向 $d=1$ 的情形。目标是刻画 $(1+1)$ 维 TQFT，为此需要下面的定义。
\begin{definition}\label{def:commutative-frobenius-algebra}
\emph{交换 Frobenius 代数}是一个有限维、有单位元 $1$ 的交换结合代数 $(A,*)$，配备非退化对称双线性配对 $\langle\ ,\ \rangle$，满足对所有 $a,b,c\in A$ 都有
\begin{equation}\label{eq:frobenius-invariant-pairing}
\langle a*b,c\rangle=\langle a,b*c\rangle.
\end{equation}
\end{definition}
文献中有时使用其他等价定义。例如，某些公理体系给出满足一定性质的\emph{迹映射} $\epsilon:A\to\mathbb C$。按这里的约定，只须置 $\epsilon(a)=\langle a,1\rangle$。

Frobenius 代数有十分丰富的结构。可定义三点关联函数 $\langle\ ,\ ,\ \rangle:A^{\otimes3}\to\mathbb C$，令所有 $a,b,c\in A$ 都满足
\[
\langle a,b,c\rangle=\langle a*b,c\rangle.
\]
由式~\eqref{eq:frobenius-invariant-pairing} 和交换性，\IndexLocation{indexloc:0341:455}三点函数对各自变量完全对称。类似地，对任意 $n$，可定义 $n$ 点关联函数为
\begin{equation}\label{eq:frobenius-n-point-function}
\langle a_1,\ldots,a_n\rangle=\langle a_1*\cdots*a_{n-1},a_n\rangle,
\end{equation}
也可证明，它对各自变量完全对称。

现在假设有一个 $(1+1)$ 维\IndexLocation{indexloc:1336:455}拓扑量子场论。固定一个标准的定向闭圆盘 $\Delta=\{z:|z|\leq1\}$，以 $S^1$ 表示其边界。置 $\mathcal H=Z(S^1)$，并将元素 $Z_\Delta\in Z(\partial\Delta)=Z(S^1)=\mathcal H$ 记为 $T_0$。我们将证明，$\mathcal H$ 有自然的 Frobenius 代数结构，以 $T_0$ 为单位元。

容易在 $\mathcal H$ 上构造配对。复共轭给出的同构 $\bar S^1\simeq S^1$，通过公理 4 诱导同构 $\mathcal H\simeq\mathcal H^*$。以 $\langle\ ,\ \rangle$ 表示相应的非退化双线性配对。

\SourcePage{456}
接着须构造乘法 $*:\mathcal H\otimes\mathcal H\to\mathcal H$。为此，考虑如下的“裤形曲面” $\Sigma$：
\begin{figure}[H]
\centering
\refstepcounter{equation}\label{eq:tqft-pair-of-pants}
\makebox[\textwidth]{\makebox[0pt][l]{(\theequation)}\hfill
\raisebox{-.5\height}{\includegraphics[width=0.278940\textwidth]{b6-pants.png}}\hfill\mbox{}}
\end{figure}
选择 $\Sigma$ 的边界分支与 $\bar S^1\coprod\bar S^1\coprod S^1$ 的定向同构，两个 $\bar S^1$ 对应于式~\eqref{eq:tqft-pair-of-pants} 左侧的边界圆，$S^1$ 对应于右侧的边界圆。由公理 4，得到元素 $Z_\Sigma\in\operatorname{Hom}_{\mathbb C}(\mathcal H\otimes\mathcal H,\mathcal H)$。以 $*$ 表示 $Z_\Sigma$ 在 $\mathcal H$ 上定义的乘法。
\begin{proposition}\label{prop:tqft-frobenius-algebra}
乘法 $*$ 交换且结合。元素 $T_0$ 是 $*$ 的单位元，对所有 $a,b,c\in\mathcal H$ 都有恒等式
\[
\langle a*b,c\rangle=\langle a,b*c\rangle.
\]
因此，$\mathcal H$ 在 $*$ 下是 Frobenius 代数。
\end{proposition}
\begin{proof}
将选定的同构 $\partial\Sigma\simeq\bar S^1\coprod\bar S^1\coprod S^1$，与交换前两个分支的 $\bar S^1\coprod\bar S^1\coprod S^1$ 的同构复合，便可得到另一个这样的同构。由函子性得到交换性。现在，将 $\Delta$ 粘到式~\eqref{eq:tqft-pair-of-pants} 中的一个边界 $\bar S^1$ 上，所得空间记为 $X$。由公理 6，可将 $Z_X$ 识别为 $\mathcal H$ 的自同态 $\phi\mapsto T_0*\phi$。另一方面，$X$ 与圆柱 $S^1\times I$ 定向保持同胚。由公理 5，有 $Z_X=1_{\mathcal H}$。比较两种计算，得到 $T_0$ 为 $*$ 的单位元。再选取 $\partial\Sigma$ 与 $\bar S^1\coprod\bar S^1\coprod\bar S^1$ 的同构，则 $Z_\Sigma\in\operatorname{Hom}_{\mathbb C}(\mathcal H\otimes\mathcal H\otimes\mathcal H,\mathbb C)$；这里仍用同一个符号 $Z_\Sigma$，稍微混用了记号。若在第三个边界分支上，将这一边界识别与同构 $\bar S^1\simeq S^1$ 复合，则算得 $Z_\Sigma(a,b,c)=\langle a*b,c\rangle$；若在第一个边界分支上作这一识别，则得到 $Z_\Sigma(a,b,c)=\langle a,b*c\rangle$。等式~\eqref{eq:frobenius-invariant-pairing} 随之成立。最后，为证明结合性，令 $\Sigma'$ 为挖去四个圆盘的二维球面，将 $\partial\Sigma'$ 与 $\bar S^1\coprod\bar S^1\coprod\bar S^1\coprod S^1$ 识别，得到 $Z_{\Sigma'}\in\operatorname{Hom}_{\mathbb C}(\mathcal H\otimes\mathcal H\otimes\mathcal H,\mathcal H)$。现在将 $\Sigma'$ 如下分解成两个裤形曲面：
\begin{figure}[H]
\centering\includegraphics[width=0.659693\textwidth]{b6-assoc-left.png}
\end{figure}
这表明 $Z_{\Sigma'}(a,b,c)=(a*b)*c$。不过，由于所处理的是\IndexLocation{indexloc:1336:456}拓扑量子场论，也可将 $\Sigma'$ 以
\SourcePage{457}
另一种方式分解成两个裤形曲面，即：
\begin{figure}[H]
\centering\includegraphics[width=0.659693\textwidth]{b6-assoc-right.png}
\end{figure}
第二幅图给出 $Z_{\Sigma'}(a,b,c)=a*(b*c)$。结合前面对 $Z_{\Sigma'}(a,b,c)$ 的计算，得到 $(a*b)*c=a*(b*c)$，这正是所需的结论。
\end{proof}

\looseness=-1 不难看出，$(1+1)$ 维 TQFT 由命题~\ref{prop:tqft-frobenius-algebra} 中描述的有限维 Frobenius 代数 $\mathcal H$ 唯一确定。这里不作证明，而是说明这些公理如何在实际中使用。例如，二维环面 $T^2$ 可由两个圆柱粘合得到。略微修改 $(0+1)$ 维情形中 $Z(S^1)$ 的计算，就得到 $Z(T^2)=\dim\mathcal H$。将两个圆盘粘合得到 $S^2$，则得到 $Z(S^2)=\langle T_0,T_0\rangle$。对其他闭的定向二维曲面 $\Sigma$，计算 $Z(\Sigma)$ 同样直接。

\looseness=-1 得到 $(1+1)$ 维 TQFT 的 Frobenius 代数 $\mathcal H$ 后，还可在 TQFT 中解释式~\eqref{eq:frobenius-n-point-function} 的 $n$ 点函数。例如，$n=4$ 时，令 $\Sigma'$ 为挖去四个圆盘的二维球面，并将其边界与 $\bar S^1\coprod\bar S^1\coprod\bar S^1\coprod\bar S^1$ 识别。仿照命题~\ref{prop:tqft-frobenius-algebra} 证明中的论证，可以说明，对所有 $a,b,c,d\in\mathcal H$，四点函数由
\[
\langle a,b,c,d\rangle=Z_{\Sigma'}(a,b,c,d)
\]
给出。更一般地，$n$ 点函数可用挖去 $n$ 个圆盘的二维球面来定义。

\looseness=-1 如\hyperref[sec:nonlinear-sigma-models]{附录~\ref*{sec:nonlinear-sigma-models}}所述，A 模型是依赖于复化 Kähler 类选择的 $(1+1)$ 维 TQFT。那么，它的底层 Frobenius 代数是什么？在该节中已经将 $\mathcal H$ 与 $H^*(X)$ 识别。配对由
\[
\langle T_1,T_2\rangle=\int_XT_1\cup T_2
\]
定义，而乘法 $*$ 就是量子乘法！\footnote{这里实际上必须使用超向量空间。这一点很明显，因为量子乘法是超交换的，见定理~\ref{thm:small-quantum-ring}。}这样便得到结合的量子上同调环。迹映射为 $\epsilon(T)=\int_XT$。若 $X$ 一般，即未必是 Calabi--Yau 流形，这个量子乘法来自称为\emph{拓扑 $\sigma$ 模型}的物理理论。已经看到，量子乘法由式~\eqref{eq:tqft-pair-of-pants} 的裤形曲面编码；将这些曲面粘合，可以得到挖去若干圆盘的更高亏格 Riemann 曲面 $\Sigma_g$。为边界圆适当选择定向，便得到亏格为 $g$ 的关联函数 $Z_{\Sigma_g}:\mathcal H^{\otimes n}\to\mathbb C$。它们推广了上述 $n$ 点关联函数，后者全部是亏格为零的关联函数。将 $\mathcal H$ 与 $H^*(X,\mathbb C)$ 识别后，这些更高亏格的关联函数，与 \SourceCite{RT1} 中的一些不变量相识别。注意，\SourceCite{RT1} 中固定了复结构，
\SourcePage{458}
因而这些不变量是纯拓扑的，不依赖复结构。

Frobenius 代数结构解释了，为什么在数学家能够证明量子上同调存在之前，物理学家就已知道它的存在。\hyperref[chap:quantum-cohomology]{第~\ref*{chap:quantum-cohomology} 章}已对量子上同调作了广泛的数学处理。量子上同调的物理解释见 \SourceCite{LVW}。

最后还应提及，$(2+1)$ 维拓扑量子场论与 Jones 多项式有关。\SourceCite{Atiyah} 给出了很好的介绍。
\index{拓扑量子场论|hyperpage}\index{Frobenius代数@Frobenius 代数|hyperpage}\index{迹映射|hyperpage}\index{裤形曲面|hyperpage}\index{拓扑sigma模型@拓扑 $\sigma$ 模型|hyperpage}
\par\endgroup

\chapter*{附录 C\quad 译注汇编}
\addcontentsline{toc}{chapter}{附录 C\quad 译注汇编}
\markboth{附录 C\quad 译注汇编}{附录 C\quad 译注汇编}

本汇编收录逐页对照原书英文扫描件时写下的译注，共 92 条，按原书 PDF 物理页排序（与原书印刷页相差 23）。内容分三类：原文笔误与 OCR 不可辨处的处理、术语与记号的歧义、作者未展开的跳步说明。正文未插入脚注，以免影响版面与原书页码对位；读者可据页号回查。


\subsection*{第 8 章}

\noindent\textbf{\footnotesize p282}\quad 第 5 行「$H^{\mathrm{middle}}(V^\circ)$ 在杯积下由 $H^{1,1}(V^\circ)$ 生成」疑为原文笔误：按 Lefschetz 分解，由 $H^{1,1}$ 在杯积下生成的是偶上同调 $H^{\mathrm{even}}$（CY 三维簇的中维 $H^3$ 含 $H^{2,1}$，不可能由 $H^{1,1}$ 生成，且后文 $\mathcal H^{V^\circ}$ 的基 $T_0,T_j,T^j,T^0$ 也对应 $H^{\mathrm{even}}$）。译文照录原文，读者可留意。


\subsection*{第 9 章}

\noindent\textbf{\footnotesize p299}\quad 原文 "the action of G on X lifts to an action of E which is linear on fibers" 中 "an action of E" 应为 "an action on E"（原文或 OCR 笔误）；译文「提升为 $E$ 上的作用」已按正确含义处理，可保留。

\noindent\textbf{\footnotesize p300}\quad 命题 9.1.2 中 $H_T^*(X)\otimes\mathcal R_T$、$H_T^*(Z_j)\otimes\mathcal R_T$ 的张量积省略了下标 $H^*(BT)$（$\mathcal R_T$ 即其分式域），原文亦如此省略，不是漏译；严谨读法应为 $\otimes_{H^*(BT)}\mathcal R_T$。

\noindent\textbf{\footnotesize p301}\quad 记号冲突：$T$ 在本章既指环面（$T$-作用、$H_T^*$、$\mathcal R_T$），又指 [GH] 记法下的 tautological 丛——本节中它即 $\mathcal O(-1)$、$c_1(T)=-p$（第 7 行）。译文保留原文说明，但同一字母两义，读公式时需区分。

\noindent\textbf{\footnotesize p303}\quad 原文措辞不严（疑笔误）：第 90 行“沿包含映射 $i_{\mathrm{point}}$ 拉回”与交换图 (9.11) 的类型不符——$i_{\mathrm{point}}$ 是 $\mathrm{point}\to BT$ 的包含，把 $c_k^T(E)$ 送回 $H^*(X)$ 的应是 $i_X$；可用的整体运算是第 98 行的 $i_{\mathrm{point}}^*\circ\int_{X_T}$。译文照原文直译。

\noindent\textbf{\footnotesize p304}\quad 记号冲突：$\lambda_1,\ldots,\lambda_5$ 在原文中既是 $(\mathbb C^*)^5$ 的群元素坐标（OCR: (入1,…,A5)·(x1,…)=(X-c1,…)），又是后文“把每个 $\rho_i$ 替换为 $\lambda_i$”所得的特征标权重；译文将群坐标改用 $t_i$、只保留 $\lambda_i$ 作权重，属对原文记号的整理，对照原文时需留意。

\noindent\textbf{\footnotesize p306}\quad 正合列 (9.13) 在英文 OCR 中两端作 $0$，译文作 $1$；群的短正合列以 $1$ 记平凡群，属记号层面的规范化，非误译。

\noindent\textbf{\footnotesize p308}\quad 记号不统一：同一对象（图 Γ 对应的模空间）在第 3、5 行写作 $\overline M_\Gamma$，在第 38、55、114 行写作 $\overline{\mathcal M}_\Gamma$；英文 OCR 均作 M\_Γ，建议全篇统一。

\noindent\textbf{\footnotesize p308}\quad 原文 OCR "We put i(F) = v" 显系笔误，应为 i(F)=i\_v（否则 ω\_F=(λ\_\{i(F)\}−λ\_\{j(F)\})/d\_e 无意义）；译文按 i(F)=i\_v 处理，正确。

\noindent\textbf{\footnotesize p310}\quad 英文（en 第 31 行）作 'at the point p'，与本节局部记号 q=(0,0) 不符，疑原文笔误；译文统一作 q，未按原样照译。

\noindent\textbf{\footnotesize p311}\quad 英文（en 第 48 行）'the denominator of e\^{}F (resp. the first factor of e\^{}F)' 两处均作 e\^{}F，但紧随的 q\_\{i(v)\} 表明第二处应为 e\^{}v；译文作 e\^{}v，属合理校订。

\noindent\textbf{\footnotesize p314}\quad 三式可相乘核对：e\^{}v 的分子、分母同为 (λ1−λ0)(λ0−λ1) 型，去掉两处 1/2 后 e\^{}v=4；于是 e\^{}F·e\^{}v·e\^{}e=(λ0−λ1)(λ1−λ0)（注意 (λ1−λ0)=−(λ0−λ1)），与文中 Euler\_T(N\_Γ3) 一致。作者未展示这步乘法，符号易错。

\noindent\textbf{\footnotesize p315}\quad 第 85 行「该分支中的一般稳定映射」与「将 C 分解为分支 C\_i」相邻出现，同一个「分支」兼指 fixed point component 与 curve component；后者宜改作「分量」，以免指代含混。

\noindent\textbf{\footnotesize p318}\quad 作者跳步：由法丛正合列 $0\to T_{\mathbb P^1}\to g^*T_V\to N_g\to0$ 直接断言 $h^0(N_g)=0$ 当且仅当 $h^0(g^*T_V)=3$，其中用到 $h^0(\mathbb P^1,T_{\mathbb P^1})=3$（即 $\operatorname{Aut}(\mathbb P^1)$ 的维数）与 $h^1(\mathbb P^1,T_{\mathbb P^1})=0$，原文与译文均未点明，读者易在此卡住。

\noindent\textbf{\footnotesize p319}\quad 译文用“分支”兼译 component 与 branch（如“两个光滑有理分支的并”与“结点附近的两个分支”），同一证明内两义并存，读者对照原文时易混；建议 component 作“连通分支”、branch 作“支”。

\noindent\textbf{\footnotesize p320}\quad 作者跳步：由 $\operatorname{Ext}^1(f_i^*\Omega_V^1\to\Omega_C^1,\mathcal O_C)$ 的维数为 0 直接断言 $f_i$ 是“预期维数的光滑点”，隐含用到五次三维簇 $c_1(T_V)=0$、$g=0$、次数 10 时虚拟维数 $(3-3)(1-0)+c_1\cdot\beta=0$ 这一事实，原文未点明。

\noindent\textbf{\footnotesize p322}\quad 英文中的交叉引用（例 7.1.6.1、例 9.2.1.3、推论 9.1.4、例 7.4.5.3、定理 11.2.16）在译文中均为 \textbackslash{}ref/\textbackslash{}hyperref，从本单元 tex 无法核对编号是否与原书一致，建议编译后抽查。


\subsection*{第 10 章}

\noindent\textbf{\footnotesize p325}\quad 原文给出虚基本类次数 $2(1-g)(d-3)-2\int_\beta\omega_X+2n$ 时未交代 $\omega_X$ 何指；应为典范类（$c_1(X)=-\omega_X$，于是次数即 $2(1-g)(d-3)+2\int_\beta c_1(X)+2n$）。译文补出“$\omega_X$ 为典范类”，与公式自洽，属作者跳步而非误译。

\noindent\textbf{\footnotesize p326}\quad 原文称 $1\in H^0(X)$ 为“the fundamental class of X”，严格说 $1$ 是基本类 $[X]$ 在 Poincaré 对偶下的上同调单位（原文笔误）；译文改述为“基本类所对应的上同调单位”，属合理修正。

\noindent\textbf{\footnotesize p327}\quad 记号冲突提醒：第 51 行维数公式里用 $\omega_X$，但本单元只在第 4 行定义了 $\omega$（复化 Kähler 类），$\omega_X$ 未加说明；若它指典范类 $K_X$，则 $-2\int_\beta\omega_X=+2\int_\beta c_1(X)$ 才与虚拟维数一致，读者极易误读成 Kähler 类。

\noindent\textbf{\footnotesize p328}\quad 原文（OCR）称 $\mathcal L_1$ 在 $\pi_{n+1}$ 的纤维上限制为“在 $p_i$ 处消失的对偶微分”，紧接着却说该层次数为 $2g-2+n$：消失给出的是 $2g-2-n$，二者矛盾。译文改作“允许在各 $p_i$ 处有至多一阶极点”（即 $\omega_C(\sum p_i)$），与次数一致，疑原文漏字或 OCR 失真，宜核对原书。

\noindent\textbf{\footnotesize p332}\quad 第 118 行译文作「c\_k(X)=c\_k(Θ\_X) 是切丛的 Chern 类」，英文 OCR 作「cotangent bundle」（余切丛）。按上下文（GW 障碍理论用 R\^{}1π\_*e\^{}*T\_X，且本页末结果 (1/24)∫\_X c\_d(X) 需 ∫\_X c\_d(T\_X)=χ(X)）Θ\_X 应为切丛，译文正确；疑原文或 OCR 有误，请勿据英文改回「余切丛」。

\noindent\textbf{\footnotesize p335}\quad 第 95 行“考虑从 $\overline M_{0,n}(X,\beta)$ 到一点的映射，记为 $\pi:\overline M_{0,n}(X,\beta)\to\overline M_{0,3}$”自相矛盾：既说“到一点”，又写映到 $\overline M_{0,3}$。按上下文 $\pi$ 应是删除标记点并收缩不稳定分支的遗忘映射，疑为原文笔误（或 OCR 脱字），译文照译未作说明。

\noindent\textbf{\footnotesize p338}\quad 记号冲突：δ 既表示除子类 δ=Σ(i=1..r) t\_iT\_i（命题 10.2.3 起），又在 (10.22) 中作 Kronecker 符号 δ\_\{k,n+1\}，同页并用，易混。

\noindent\textbf{\footnotesize p338}\quad 作者跳步：由 (10.24) 及 Σ(n≥0) ħ\^{}\{-(n+1)\}/(n+1)!·δ\^{}\{n+1\}=e\^{}\{δ/ħ\}-1，(10.23) 第一行括号内实为 T\_a∪(e\^{}\{δ/ħ\}-1)，再加开头的 T\_a 才得 T\_a∪e\^{}\{δ/ħ\}；原文以「连同 T\_a」一笔带过。

\noindent\textbf{\footnotesize p339}\quad 定理 10.2.4 陈述 $N_j(s(T))=-s(T_j\cup T)$（行 16），而证明中引作 $N_j(s(T))=s((-T_j)\cup T)$（行 56）；两者由 $s$ 的线性（$s(-X)=-s(X)$）等价，非原文笔误。

\noindent\textbf{\footnotesize p341}\quad 式 (10.27) 把推论 10.2.5 中按上同调全基 $\{T_0,T_j,T^j,T^0\}$ 的求和一次展开成四项（行 126–130）；其中 $\langle 1/(\hbar-c),T^0\rangle$ 与 $\langle 1/(\hbar-c),T^j\rangle$ 两类项分别由次数公理、基本类公理恒为零，作者未点明这一展开，属跳步。

\noindent\textbf{\footnotesize p344}\quad 作者跳步：$\widehat C=\mathbb C[[q_1,\dots,q_r]]$ 只含 $q_1,\dots,q_r$，而正文列出的 $P$ 的量表含 $e^{t_0}$，原文未交代其代换。译文第 108 行自行补出"若 $P$ 还含 $e^{t_0}$，则在系数环中相应增添标量参数 $q_0$"一句（英文 OCR 无对应句）；补足内容合理，建议改为脚注并标明系译者补足。

\noindent\textbf{\footnotesize p347}\quad 例题 10.3.1.1 中 $H^2$ 先后表示两种乘法：先写小量子上同调关系 $H^2=q$（量子上积 $H*H=q=e^{t_1}$），随后又写「由于 $H^2=0$」（经典上积 $H\cup H=0$，$H^4(\mathbb P^1)=0$）。同一记号 $H^2$ 在同一例中含义不同，作者未加区分，读者易误读。

\noindent\textbf{\footnotesize p348}\quad 记号提示：全节用上/下标区分对偶基——T\_0=1、T\_1,…,T\_r 张成 H\^{}2，T\^{}a∈H\^{}4、T\^{}0∈H\^{}6 为对偶类；随后作者让 β 既表示 H\_2(V,C) 中的曲线类，又表示它在 H\^{}4(V,C) 中的 Poincaré 对偶，原文自述为「略微滥用记号」，译文照录。

\noindent\textbf{\footnotesize p349}\quad 作者在命题 10.3.4 证明中把 (10.42) 三条恒等式的验证留给读者；其中第二条 (1/6)δ\^{}3=(1/3)Σt\_a∂\_a(1/6∫\_Vδ\^{}3)T\^{}0 依赖 δ 关于 t 的三次齐次性（欧拉定理），原文只以「第二行使用 Euler 公式」一句带过。

\noindent\textbf{\footnotesize p351}\quad 记号不一致：第 3 行作 (ħ d/dt₁)⁴J≠0，而第 9 行以 (ħ d/dt₁)⁴⟨s\_j,1⟩ 表述同一事实；J 与 ⟨s\_j,1⟩ 的关系原文未交代，疑原文用字有误（或 OCR 误识），译文照译。

\noindent\textbf{\footnotesize p352}\quad 作者跳步：由 (10.40) 直接写出 J=1+ħ⁻¹ΣA\_jT\_j+ħ⁻²ΣB\_jT\^{}j+ħ⁻³CT⁰，省略了先提出公因子 e\^{}\{t₀/ħ\}（对 t₁,…,t\_r 的微分不变）这一步，译文补出。另外英文 OCR 在 “and (10.46) follows.” 之后漏掉一行（By (10.46), we can specialize ħ to any nonzero constant），故译文中“由式(10.46)，可将 ħ 特化为任意非零常数”一句系据原文补足，不是增译。


\subsection*{第 11 章}

\noindent\textbf{\footnotesize p354}\quad 原文笔误：首段 "some versions of the the Mirror Theorem" 重复冠词 the（OCR 亦保留），译文按单冠词处理，无需改动。

\noindent\textbf{\footnotesize p355}\quad 记号冲突：Gromov–Witten 不变量在第七章记为 N\_d，本章改记为 K\_d，而把 N\_d 移用于线性 σ 模型的射影空间（N\_d 参数化 P\^{}1 上次数 d 的齐次式元组）；作者正文已作提示，跨章阅读时仍需留意同一符号两种含义。

\noindent\textbf{\footnotesize p357}\quad 记号冲突：本节 n 表示目标 P\^{}n 的维数，而引用第 9.2 节时 \textbackslash{}overline M\_\{0,n\}(P\^{}r,d) 中的 n 表示标记点数、r 表示射影空间维数。两种 n 在同一句（第 30 行）并列，正是 (C*)\^{}\{n+1\} 误写的根源；阅读该句时应按“第 9.2 节的 n 与本节无关”来理解。

\noindent\textbf{\footnotesize p359}\quad 原文（p.359，(11.6) 之后）称 α\_i 的公因子 L(w\_0,w\_1)\^{}\{d'\} 中 L 在 p 处消失；但按 (11.6) 的参数化 (w\_0,w\_1)↦((w\_1,w\_0),…)，源坐标与目标 P\^{}1 坐标互换，L 实际应在 p 的交换坐标处消失。译文已按交换后的点表述。此属原文笔误（或至少是记号含混），读者引用时宜注意。

\noindent\textbf{\footnotesize p360}\quad 原书 (11.10)、(11.11) 中 e\_1 一名三用：既标 M̄\_\{0,1\}(Pⁿ,r)→Pⁿ、M̄\_\{0,1\}(Pⁿ,d−r)→Pⁿ 两条边，又标 M̄(r,d−r)→Pⁿ；三者靠拉回图的交换性互相认同（r=0 时 p₁ 与 e₁ 也互换角色），原书未加区分，属记号冲突。

\noindent\textbf{\footnotesize p361}\quad 式 (11.12) 写作 κ̄=κ−dℏ、ℏ̄=−ℏ，但例 11.1.3.1 与 (11.13) 把上划线用于不动点限制 i*\_\{p\_\{i,0\}\}(Q\_r) 时，它实际只把 ℏ 换成 −ℏ（κ 已代成 λ\_i）；恒等式“P\_r(λ\_i) 的上划线 = ∏\_\{m=0\}\^{}\{ℓr\}(ℓλ\_i+mℏ)”依赖这一约定，原书未交代此跳步。

\noindent\textbf{\footnotesize p363}\quad 原文/OCR 记号冲突：式 (11.16) 的因子 OCR 作 (ℓλ\_i)\^{}\{-1\}，而末段比较时作 “lp”（即 ℓp）。只有乘以 ℓλ\_i 才能消去该因子并恰好得到下一恒等式 i\^{}*\_\{p\_\{i,r\}\}(φ)∏(λ\_i−λ\_j)=上划线 i\^{}*\_\{p\_\{i,0\}\}(φ)·i\^{}*\_\{p\_\{i,0\}\}(φ)，故译文统一作 ℓλ\_i；英文侧疑有笔误或 OCR 失真。

\noindent\textbf{\footnotesize p365}\quad 作者跳步：称最后恒等式“容易验证”，实需分 j=i 与 j≠i 两类、并按 s↔r−s 配对核验（右端两乘积的 s 分别取 0..d−r 与 0..r）；译文照录，未补说明。

\noindent\textbf{\footnotesize p368}\quad 唯一性证明（定理 11.1.7 的证明）跳过一步：「仅 r=0,d 的项可能贡献」依赖如下事实——当 0<r<d 时 i*\_\{p\_\{i,r\}\}(P\_d) 可由 P\_1,…,P\_\{d-1\} 表出，故归纳假设 P\_j=Q\_j（j<d）使这些不动点的贡献相消，只剩 r=0,d。作者未展开，译文照原文未补。

\noindent\textbf{\footnotesize p371}\quad 第 (ii) 部分的次数界：作者称由第 (i) 部分即得，但 (i) 只给出 $I_d^*(\widehat Q_d)=\phi^d e_{1!}(\alpha^d)$ 与 $\deg_\hbar\phi^d=(n+1)d$；要降到 $(n+1)d-2$ 还须 $e_{1!}(\alpha^d)$ 的 $\hbar$-次数估计，此步被略去。

\noindent\textbf{\footnotesize p373}\quad 定理 11.1.13 由 (11.29) 的推导原文称“留给读者”：需把 $\sum_d K_de^{dt}$、$\sum_d dK_de^{dt}$ 借 $\Phi=\frac{5t^3}{6}+\sum_d K_de^{dt}$ 换成 $\Phi,\Phi'$，再把 $e^{-Ht/\hbar}$ 展开并在 $H^5=0$（$H\in H^*(\mathbb P^4)$）处截断，才得定理中的四项式。

\noindent\textbf{\footnotesize p374}\quad 由“在 $\pi_1$ 的一条纤维上积分”得出 $\hbar^{-3}(-dt+2)$ 与 $\hbar^{-2}d$，实际用到 $\pi_{1!}(e_1^*H)=d$、$\pi_{1!}(c_1(\mathcal L_{d,1}))=2$（光滑有理纤维上 $\int c_1(T^*C)=2$），原文未交代，属作者跳步。

\noindent\textbf{\footnotesize p375}\quad 原文 P375 的 concavex 疑为 concave 与 convex 的合成词（指丛为凸型或凹型两类），译文作「凸凹」；同页 Example 11.1.6.1 又把 concave 译作「凹丛」，两处概念层次不同，读者宜对照原文。

\noindent\textbf{\footnotesize p376}\quad 式 (11.32) 及最后比较式用 H（非等变超平面类），而第 107、112、119、126 行起的等变展开一律改用 p；p 在非等变极限 λ\_i→0 下即为 H，作者未说明这一记号切换。

\noindent\textbf{\footnotesize p379}\quad 原文作 "where J is a variant of the function J discussed in Section 10.3"，两个 J 实为 $J_{\mathcal V}$ 与 $J$（原文脱落下标，或 OCR 漏字），应读作「$J_{\mathcal V}$ 是第 10.3 节所论函数 $J$ 的变体」。

\noindent\textbf{\footnotesize p383}\quad 第128–133行：(11.49) 的第二个等号除「j 的投影公式」外还用到 (11.47)（j\_* 把虚基本类推前为 Euler(V\_\{d,2\}) 与通常基本类的帽积），原文只提投影公式，未点明此处跳步。

\noindent\textbf{\footnotesize p386}\quad 英文 OCR 此页两处把 (H\^{}3+sH\^{}4) 作 (H\^{}2+sH\^{}4)（(I\_V/y\_0)'' 与 J\_V''=i\_!(J\_V)'' 两式），与紧随其后的 (11.51) 及 e\^{}\{sH\}H\^{}3=H\^{}3+sH\^{}4+⋯ 矛盾，应为 H\^{}3；译文统一作 H\^{}3，属正确重构，非误译（疑原文笔误或 OCR 误识）。

\noindent\textbf{\footnotesize p387}\quad 作者跳步：仅称“简单计算表明”该算子消去 $I_{\mathcal V}$，未给出超几何级数系数递推的验证；又由该算子作用结合推论得到 $H^n-q\ell(\ell H)^{\ell-1}=0$ 时，隐含取经典极限 $\hbar\to0$、$\hbar\partial_{t_1}\mapsto i^*H$（并记 $e^{t_1}=q$），原文均未交代。

\noindent\textbf{\footnotesize p387}\quad 记号冲突：同一字母 $H$ 既表示 $\mathbb P^n$ 的超平面类（如 $J_{\mathbb P^n}$ 的公式中），又表示其在 $X$ 上的限制 $i^*H$；作者是在给出关系 $H^n-q\ell(\ell H)^{\ell-1}=0$ 之后才补注“简记 $H=i^*H$”，阅读该式时需按后一含义理解。

\noindent\textbf{\footnotesize p390}\quad (11.57)→(11.59) 的截断依据是次数计数：被积类 e\_1\^{}*(φ\_i)c\_T\^{}k 的复次数为 n+k，故只有 k=(n+1)d-1 的项可能非零；原文「k≥(n+1)d-1 才可能有贡献」只是必要条件，作者未点明 k>(n+1)d-1 时次数超过维数自动为零，故从 c\_T\^{}d 起截断是安全的。

\noindent\textbf{\footnotesize p392}\quad 引理 11.2.8 称 Z\_i 在 ℏ=(λ\_i-λ\_j)/k 处正则，而正文 c\_Γ=(λ\_j-λ\_i)/d' 使分母 ℏ-c\_Γ 在 ℏ=(λ\_j-λ\_i)/d' 处为零；两处相差一负号并不矛盾——引理只断言 (λ\_i-λ\_j)/k 这些点正则，正是为随后把 ℏ=(λ\_j-λ\_i)/d 代入 Z\_j 做准备，作者未点破这层关系，读者易误判为符号笔误。

\noindent\textbf{\footnotesize p393}\quad 记号冲突：本页把边收缩处的结点记作 p（「在结点 p 处应用引理」「保持点 p 不动」），而 p.394–395 又用 p 作等变变量（I\_T、J\_T 中的 e\^{}\{(t\_0+t\_1p)/ħ\}，末尾取极限 p→H），同一单元内同名异指。

\noindent\textbf{\footnotesize p394}\quad 作者跳步：由 (11.63) 推到 (11.68) 需先把末行整体认作 Z\_j(q(λ\_j-λ\_i)/(d'ħ),(λ\_j-λ\_i)/d')，再把 (q/ħ)\^{}\{d'\} 与末行的 q\^{}\{d''\} 并成 q\^{}\{d'+d''\}、留下 ħ\^{}\{-d'\}；原文仅以「递推便清楚了」一句带过。

\noindent\textbf{\footnotesize p398}\quad 引理 11.2.11 证明末句承认省略了实质一步：要说明式 (11.63) 类比式分子中多出的因子恰是 C\^{}V\_\{i,j\}(d′) 的分子，需计算 T 在丛 V′\_\{d,2,1\} 上的权，原文即写 “we omit the details of this argument”，读者需自行补足。

\noindent\textbf{\footnotesize p401}\quad 作者跳步/原文笔误：原文称 “ħ\^{}\{-2r\_i\}(A\_i d\_0 z + 2B\_i) is a polynomial in ħ”，严格说成立的是整个和 ∑\_i (E\_i/∏\_\{j≠i\}(λ\_i-λ\_j))e\^{}\{λ\_i z\}ħ\^{}\{-2r\_i\}(A\_i d\_0 z+2B\_i+O(ħ)) 为 ħ 的多项式，单项并不必然；作者省略了“取最高负幂、再用 \{e\^{}\{λ\_i z\}, z e\^{}\{λ\_i z\}\} 线性无关排除相消”这一步。

\noindent\textbf{\footnotesize p403}\quad 原文 OCR 作“the Eulerity property”，疑为“Euler data property”的 OCR 误识或原文笔误；译文作“Euler 数据性质”，与 11.1.3 小节（subsec:lly-euler-data）的 Euler 数据（LLY）主题一致，建议对照纸书确认。

\noindent\textbf{\footnotesize p407}\quad 式(11.76) 中的 ∏\_\{m=−∞\}\^{}\{0\} 与 ∏\_\{m=−∞\}\^{}\{d\_1−2d\_2\} 是形式无穷乘积：两段相消后，d\_1≥2d\_2 时余 1/∏\_\{m=1\}\^{}\{d\_1−2d\_2\}(D\_3−2D\_1+mħ) 留在分母；d\_1<2d\_2 时余 ∏\_\{j=0\}\^{}\{2d\_2−d\_1−1\}(D\_3−2D\_1−jħ) 留在分子。原文未交代此截断约定，而这正是 y\_0 与 y\_\{1,1\}、y\_\{1,2\} 按 d\_1≥2d\_2、d\_1<2d\_2 分两段求和以及 (−1)\^{}\{d\_1+1\}(2d\_2−d\_1−1)! 的来源。

\noindent\textbf{\footnotesize p409}\quad 记号提醒：$J_{\mathcal V}$ 与 $J_V$ 指不同对象的 Givental 函数——(11.78)、(11.80) 用 $J_{\mathcal V}$（对应 $I_{\mathcal V}/y_0$），借助 $i_!(J_V)=J_{\mathcal V}$ 才有 (11.82) 的 $J_V$；OCR 原文两者均写作 $J_V$，阅读时勿混。

\noindent\textbf{\footnotesize p410}\quad 作者跳步：第 89–93 行由 $Y_{ijk}\langle T_a,T_b,T_c\rangle=Y_{abc}\langle T_i,T_j,T_k\rangle$ 直接断言存在比例 $c$，并称「两侧都来自势函数」即可由四阶混合导数相容性得 $(\partial_1c)Y_{12k}=(\partial_2c)Y_{11k}$，中间推导原文未展开。

\noindent\textbf{\footnotesize p411}\quad 由 (11.83) 推 (11.85) 时作者只说“比较 $D_1D_3$ 的系数”，但 (11.83)/(11.77) 中含 $D_3^2-2D_1D_3$ 的项也贡献 $D_1D_3$；须先用 $D_3^2$ 的系数定出 $\partial\Phi/\partial s_2$（两项的 $-2$ 部分相消），才能得到 (11.85)。属作者跳步。

\noindent\textbf{\footnotesize p412}\quad 原文 "when $-K_X$ is nef instead of Fano" 措辞不确（Fano 即 $-K_X$ 丰沛）；译文对应作“nef 而非丰沛”，语义等价，可视为原文笔误。

\noindent\textbf{\footnotesize p414}\quad δ 一字两用：第 9 行 e\^{}δ 的 δ 是 H²(X) 中的类（Kähler 类，置 t\_0=0、ħ=1 后），第 25 行的 δ\_ρ=λ\_ρ∂\_ρ 则是带下标的对数导数算子；原文如此，阅读时需按上下文区分。

\noindent\textbf{\footnotesize p416}\quad 英文 OCR 第 145 行作 D\_2·β = −k；由 β=kD\_2 且 D\_2²=−2 应得 −2k，译文取 −2k（数学上正确）。此处疑为原书笔误或 OCR 漏字。

\noindent\textbf{\footnotesize p418}\quad 记号冲突：本单元前半以 $s_1,s_2$ 记 $\mathbb F_2$ 的对数镜像坐标（$q_i=e^{s_i}$），末段又用同一个 $s$ 表示 $\mathbb F_0$ 中另一族纤维的截面类；$f$ 也同时指 $\mathbb F_2$ 与 $\mathbb F_0$ 的纤维类。作者默认二者对应，未加区分，读者需自行辨认。


\subsection*{第 12 章}

\noindent\textbf{\footnotesize p420}\quad 小节标题 "Reflections and Open Problems" 中 Reflections 兼有“回顾/反思”与“（镜面）反射”双关，译文只取“回顾”，双关无法保留，属原文修辞而非翻译缺陷。

\noindent\textbf{\footnotesize p421}\quad 原文例、式两套编号并存：Example 11.2.5.1 为“章.节.小节.例”，而式 (11.81)、(11.84)、(11.85) 为“章.式序号”。译文原样保留；读者须知 11.2.5.1 是第 11.2.5 小节的第 1 个例题，而非小节号。

\noindent\textbf{\footnotesize p423}\quad (11.81) 处的 Givental 函数在 OCR 中作 "Jx and Jv"，第二个下标是小写 v 还是 V 不明；译文用花体 $J_{\mathcal V}$，与上下文 $V$ 的记号不一致，建议核对原书后统一为 $J_V$ 或保留原记号。

\noindent\textbf{\footnotesize p425}\quad §12.2.4 原文从"全纯辛"直接引用 Calabi--Yau 定理推出度量与平行的四元数作用，省略了 $M$ 须为紧 Kähler 流形这一前提（作者跳步）；建议以译注说明，而不是在正文中补出条件。

\noindent\textbf{\footnotesize p426}\quad 式 (12.1) 之后讨论两个消失项的那句（译文 L13）应为 H\^{}2(V°,O)=0 与 H\^{}0(V°,∧\^{}2Θ)=0（Θ 为切层）；英文 OCR 把后者印成 ∧\^{}2O\_\{V°\}，但 ∧\^{}2O=0 恒成立会让该句变成废话，且只有按 HH\^{}2=H\^{}0(∧\^{}2Θ)⊕H\^{}1(Θ)⊕H\^{}2(O) 才能对应“非交换形变”与“神秘项”，故译文用 ∧\^{}2Θ 正确，校对时勿据 OCR 回改为 O。

\noindent\textbf{\footnotesize p428}\quad 英文 OCR 把 Sturmfels 的 “deformed flatly”（平坦形变）识作 “deformed fatly”；译文“平坦形变”正确（Grassmann 簇到 P(k,n) 是平坦形变），勿据 OCR 另译。

\noindent\textbf{\footnotesize p429}\quad 原文「discussed in Section 6.2.4」在译文中改写为对 subsec:kahler-moduli-boundary 标签的交叉引用；若该标签实际小节号不等于 6.2.4，读者看到的小节号会与原书不符（本单元内无从核对）。


\subsection*{附录 A}

\noindent\textbf{\footnotesize p431}\quad 证明末句“再对 g\^{}\{-1\} 应用终端判别条件”省略了计算：g\^{}\{-1\} 的特征指数为 r−a\_i（a\_i=0 时仍记 0），其和 = 3r−(a+b+c) = r，故终端条件 Σa\_i>r 不成立，即得矛盾。作者跳过了这一步。

\noindent\textbf{\footnotesize p433}\quad 原文 A.4.2 证明中的 Ext(Ω¹,𝒪) 兼指整体 Ext 与层 Ext（OCR 无法分辨字形）；译文分别用 \textbackslash{}operatorname\{Ext\}、\textbackslash{}mathcal Ext 并加『其层』限定，属对原文记号的合理消歧，不记为缺陷。

\noindent\textbf{\footnotesize p433}\quad 命题 A.4.2 的证明由 ℰxt¹=0 推出 Ext¹(Ω¹,𝒪)≅H¹(X,Θ\_X) 实际依赖局部—整体谱序列的五项正合列，原文作『follows immediately』，属作者跳步，译文照译无误。


\subsection*{附录 B}

\noindent\textbf{\footnotesize p434}\quad 第 18 行「场必须使作用量取驻值」对应原文 p.435 开头「must be minima of the action」，而同段末尾原文又作「Φ is stationary for the action S」；原文自身措辞不一致，译文统一作「取驻值」（数学上正确的说法），建议保留，无需回改。

\noindent\textbf{\footnotesize p437}\quad 原文同一对方括号 [ , ] 兼表经典 Poisson 括号（(B.10)）与量子对易子（等时对易关系），两者只差一个 i 因子、仅靠上下文辨别；译文照搬该记号，文字上用「Poisson 括号/对易子」区分，读者需留意。

\noindent\textbf{\footnotesize p438}\quad 作者在 (B.9)→(B.13) 处以「减去常数不影响动力学」为由直接删去 ½E\_kδ(0)，但该项实为无穷大标量算子，严格处理需正规序（normal ordering）；作者未点明这一步，译文如实照译。

\noindent\textbf{\footnotesize p441}\quad 原文（OCR）把 A₀ 的 Euler–Lagrange 方程写作 δL/δA₀=0，译文作 δS/δA₀=0；二者对 A₀ 的变分等价，只是记号与原文不同。

\noindent\textbf{\footnotesize p442}\quad 原文（OCR）称守恒流的空间积分不随时间改变；严格说守恒的是其时间分量的空间积分（∫J⃗d³x 并不守恒）。译文已补「时间分量」，与 Noether 定理一致。

\noindent\textbf{\footnotesize p443}\quad 原文说 D\_z 由 'the Levi-Civita connection of the metric f*(g) on Σ' 给出，译文改述为「把 g 的 Levi–Civita 联络拉回到 f*T\_X」；两者指向不同对象（世界面诱导度量的自旋联络 vs 目标空间联络的拉回）。D\_z 作用在 ψ 上必须含目标空间联络，疑原文此句有笔误或 OCR 脱字，建议核对原书。

\noindent\textbf{\footnotesize p445}\quad 记号 N\_\{γ,ω\_i\} 把自由指标 i 写进符号下标，严格应为 N\_γ(ω\_1,…,ω\_n)；原文（OCR 作 N\_\{γ,ω\}）亦如此，属原文记号不严谨，非译文新引入。

\noindent\textbf{\footnotesize p446}\quad 本单元开头承接上一单元的句子（OCR 的 "such a form θ locally as"），第 84 行也在句中断开，均属分单元切分所致，不是译文缺漏。

\noindent\textbf{\footnotesize p447}\quad 记号冲突（作者惯例）：h 既指世界面上的度量 h\_\{ij\}，又指其行列式（第 51 行「照常用 h 表示矩阵 h\_\{ij\} 的行列式」），h\^{}\{ij\} 为逆矩阵分量；读能动张量定义 T\_\{ij\} = (1/√h)·δS/δh\^{}\{ij\} 时需按此约定理解。

\noindent\textbf{\footnotesize p450}\quad 例 B.4.1 直接给出「八次部分维数 83」和「缺少三个形变」，未交待算法：权 (1,1,2,2,2) 的 8 次单项式共 105 个，Jacobian 理想的 8 次部分由 Φ\_a·∂\_iF（2×2）与 (Φ\_1²,Φ\_1Φ\_2,Φ\_2²,Φ\_3,Φ\_4,Φ\_5)·(∂\_3,∂\_4,∂\_5)F（6×3）张成，共 22 维，故 105−22=83；另 3 维来自第 4 扭曲部分，合计与文中「余维为三的 H\^{}\{2,1\}\_poly ⊂ H\^{}\{2,1\}」相符（h\^{}\{2,1\}=86）。

\noindent\textbf{\footnotesize p451}\quad 同节先用 r 记 P\_j 超场的个数（故设 r=1、以 P 代 P\_1），稍后又说常数 r\_j 决定某超对称相互作用的强度；r 与 r\_j 形近而含义无关，原文未加区分，读者易误认为同一参数。

\noindent\textbf{\footnotesize p452}\quad 第19行规范变换 p ↦ u\^{}\{-5\}p 中的指数 −5：G 关于 Φ\_1,…,Φ\_5 五次齐次，且 W = PG 须规范不变，故 P 的 U(1) 权为 −Σ\_i Q\_\{i,1\} = −5。原文未加说明（作者跳步），译文照译未补。

\noindent\textbf{\footnotesize p452}\quad 第16行「式 (B.31) 的每一项……都必须为零」是原文 each term (B.31) vanishes 的省略说法，实指势中每个 D\_j\^{}2 项为零（即每个 D\_j = 0）；译文照搬省略，易被读成 (B.31) 括号内各项分别取零，紧接的 D\_1 = 0 显式方程可澄清。

\noindent\textbf{\footnotesize p456}\quad 证明中的 $\Sigma'$（挖去四个圆盘的二维球面）与后文 n=4 段落再次引入的 $\Sigma'$ 是同一曲面，但两处对边界的定向识别不同：前者为 $\bar S^1\coprod\bar S^1\coprod\bar S^1\coprod S^1$，得 $Z_{\Sigma'}\in\operatorname{Hom}(\mathcal H^{\otimes3},\mathcal H)$；后者为四个 $\bar S^1$，得 $Z_{\Sigma'}\in\operatorname{Hom}(\mathcal H^{\otimes4},\mathbb C)$。同一符号 $Z_{\Sigma'}$ 指两个不同张量空间中的元素，属作者记号复用（正文只对 $Z_\Sigma$ 的混用作了说明）。

\noindent\textbf{\footnotesize p458}\quad 英文页首 was fixed. 为上一页跨页残句，本单元自 Hence 起译，该半句应属 u152，勿按漏译处理。


% 原书完整英文书目，物理页 460–474；全部条目保留。
\renewcommand{\bibname}{参考文献}
\begin{thebibliography}{Kontsevich2}
\addcontentsline{toc}{chapter}{参考文献}
\SourcePage{460}
\item[]
We use the following conventions for labeling references in the text and in the bibliography:
\begin{itemize}
\item References with one author are labeled by the author's last name.
\item References with multiple authors are labeled using the first letter of the authors' last names in alphabetical order. When different groups of authors would share the same label, we distinguish between these by using the first two letters of one of the authors' last names.
\item If there is more than one reference for the same author or group of authors, they are ordered by publication year.
\end{itemize}
In addition to the usual bibliographic information, we include references to the e-print archives \textbf{math.AG} (which includes \textbf{alg-geom}), \textbf{math.DG} (which includes \textbf{dg-ga}), and \textbf{hep-th} whenever possible. These archives are accessible on the WWW at the URL \url{http://xxx.lanl.gov/}.


\bibitem[AKa]{AKa}\hypertarget{bib:AKa}{}
A. Albano and S. Katz, Lines on the Fermat quintic threefold and the infinitesimal generalized Hodge conjecture, Trans. AMS 324 (1991), 353--368.

\bibitem[Alexeev]{Alexeev}\hypertarget{bib:Alexeev}{}
V. Alexeev, Moduli spaces $M_{g,n}(W)$ for surfaces, in Higher-Dimensional Complex Varieties (Trento, 1994), de Gruyter, Berlin, 1996, 1--22, \href{https://arxiv.org/abs/alg-geom/9410003}{alg-geom/9410003}.

\bibitem[AN]{AN}\hypertarget{bib:AN}{}
V. Alexeev and V. Nikulin, Classification of Del Pezzo surfaces with Log-terminal
singularities of index $\leq2$, and involutions of K3 surfaces,
Soviet Math. Dokl. 39 (1989), 507--511.

\bibitem[AKl]{AKl}\hypertarget{bib:AKl}{} A. Altman and S. Kleiman, \emph{Foundations of the Theory of Fano Schemes}, Compositio Math. \textbf{34} (1977), 3--47.

\bibitem[Aluffi]{Aluffi}\hypertarget{bib:Aluffi}{} \emph{Quantum Cohomology at the Mittag-Leffler Institute}, (P. Aluffi ed.), Appunti della Scuola Normale Superiore, Pisa, 1997.

\bibitem[Alvarez]{Alvarez}\hypertarget{bib:Alvarez}{} O. Alvarez, \emph{Lectures on quantum mechanics and the index theorem}, in \emph{Geometry and Quantum Field Theory} (D. Freed and K. Uhlenbeck, eds.), IAS/Park City Mathematics Series 1, AMS, Providence, RI, 1995, 273--322.

\bibitem[Arnold]{Arnold}\hypertarget{bib:Arnold}{} V. Arnold, \emph{Mathematical Methods of Classical Mechanics}, Springer-Verlag, New York, 1989.

\bibitem[AMRT]{AMRT}\hypertarget{bib:AMRT}{}
A. Ash, D. Mumford, M. Rapoport and Y.-S. Tai, Smooth Compactifications of Locally Symmetric Varieties, Math Sci Press, Brookline, MA, 1975.

\bibitem[Aspinwall]{Aspinwall}\hypertarget{bib:Aspinwall}{} P. Aspinwall, \emph{K3 Surfaces and String Duality}, in \emph{Fields, Strings and Duality (TASI 96, Boulder, CO)} (C. Efthimiou and B. Greene, eds.), World Scientific, River Edge, NJ, 1997, 421--540, \href{https://arxiv.org/abs/hep-th/9611137}{hep-th/9611137}.

\bibitem[AGM1]{AGM1}\hypertarget{bib:AGM1}{}
P. Aspinwall, B. Greene and D. R. Morrison,
The monomial-divisor mirror map,
Int. Math. Res. Notices (1993), 319--337, \href{https://arxiv.org/abs/hep-th/9309007}{hep-th/9309007}.

\bibitem[AGM2]{AGM2}\hypertarget{bib:AGM2}{}
P. Aspinwall, B. Greene and D. Morrison, Multiple mirror manifolds and topology
change in string theory, Phys. Lett. B303 (1993), 249--259, \href{https://arxiv.org/abs/hep-th/9301043}{hep-th/9301043}.

\bibitem[AGM3]{AGM3}\hypertarget{bib:AGM3}{}
P. Aspinwall, B. Greene and D. Morrison, Measuring small distances in $N=2$ sigma models, Nuclear Physics B420 (1994) 184, \href{https://arxiv.org/abs/hep-th/9311042}{hep-th/9311042}.

\SourcePage{461}

\bibitem[AM1]{AM1}\hypertarget{bib:AM1}{} P. Aspinwall and D. Morrison, \emph{Topological field theory and rational curves}, Comm. in Math. Phys. \textbf{151} (1993), 245--262, \href{https://arxiv.org/abs/hep-th/9110048}{hep-th/9110048}.

\bibitem[AM2]{AM2}\hypertarget{bib:AM2}{} P. Aspinwall and D. Morrison, \emph{Chiral rings do not suffice: $N=(2,2)$ theories with nonzero fundamental group}, Phys. Lett. \textbf{B344} (1994), 79--86, \href{https://arxiv.org/abs/hep-th/9406032}{hep-th/9406032}.

\bibitem[AM3]{AM3}\hypertarget{bib:AM3}{} P. Aspinwall and D. Morrison, \emph{String theory of K3 surfaces}, \emph{Mirror Symmetry II} (B. Greene and S.-T. Yau, eds.), AMS/IP Stud. Adv. Math. 1, AMS, Providence, RI, 1997, 703--716, \href{https://arxiv.org/abs/hep-th/9404151}{hep-th/9404151}.

\bibitem[AS]{AS}\hypertarget{bib:AS}{} A. Astashkevich and V. Sadov, \emph{Quantum cohomology of partial flag manifolds}, Comm. Math. Phys. \textbf{170} (1995), 503--528, \href{https://arxiv.org/abs/hep-th/9401103}{hep-th/9401103}.

\bibitem[Atiyah]{Atiyah}\hypertarget{bib:Atiyah}{} M. Atiyah, \emph{The Geometry and Physics of Knots}, Cambridge University Press, Cambridge, 1990.

\bibitem[AB]{AB}\hypertarget{bib:AB}{}
M. Atiyah and R. Bott, The moment map and equivariant cohomology, Topology 23 (1984), 1--28.

\bibitem[Audin]{Audin}\hypertarget{bib:Audin}{}
M. Audin, The Topology of Torus Actions on Symplectic Manifolds,
Progress in Math. 93, Birkhäuser, Boston--Basel--Berlin, 1991.

\bibitem[ACJM]{ACJM}\hypertarget{bib:ACJM}{}
A. Avram, P. Candelas, D. Jančić and M. Mandelberg, On the connectedness of the moduli space of Calabi--Yau manifolds, Nuclear Physics B465 (1996), 458--472, \href{https://arxiv.org/abs/hep-th/9511230}{hep-th/9511230}.

\bibitem[AKMS]{AKMS}\hypertarget{bib:AKMS}{} A. Avram, M. Kreuzer, M. Mandelberg and H. Skarke, \emph{Searching for K3 fibrations}, Nuclear Physics \textbf{B494} (1997), 567--589, \href{https://arxiv.org/abs/hep-th/9610154}{hep-th/9610154}.

\bibitem[Baily1]{Baily1}\hypertarget{bib:Baily1}{} W. Baily, Jr., \emph{The decomposition theorem for V-manifolds}, Amer. J. Math. \textbf{78} (1956), 862--888.

\bibitem[Baily2]{Baily2}\hypertarget{bib:Baily2}{}
W. Baily, Jr., On the embedding of V-manifolds in projective space,
Amer. J. Math. 79 (1957), 403--430.

\bibitem[BBRP]{BBRP}\hypertarget{bib:BBRP}{} C. Bartocci, U. Bruzzo, D. Hernández Ruipérez and J. Muñoz Porras, \emph{Mirror symmetry on elliptic K3 surfaces via Fourier-Mukai transform}, Comm. Math. Phys. \textbf{195} (1998), 79--93, \href{https://arxiv.org/abs/alg-geom/9704023}{alg-geom/9704023}.

\bibitem[Batyrev1]{Batyrev1}\hypertarget{bib:Batyrev1}{}
V. Batyrev, On the classification of smooth projective toric varieties,
Tohoku Math. J. 43 (1991), 569--585.

\bibitem[Batyrev2]{Batyrev2}\hypertarget{bib:Batyrev2}{}
V. Batyrev, Quantum cohomology rings of toric manifolds,
in Journées de Géométrie Algébrique d'Orsay (Orsay, 1992),
Astérisque 218, Soc. Math. France, Paris, 1993, 9--34,
\href{https://arxiv.org/abs/alg-geom/9310004}{alg-geom/9310004}.

\bibitem[Batyrev3]{Batyrev3}\hypertarget{bib:Batyrev3}{}
V. Batyrev, Variations of the mixed Hodge structure of affine hypersurfaces in algebraic tori, Duke Math. J. 69 (1993), 349--409.

\bibitem[Batyrev4]{Batyrev4}\hypertarget{bib:Batyrev4}{}
V. Batyrev, Dual polyhedra and mirror symmetry for Calabi--Yau
hypersurfaces in toric varieties,
J. Algebraic Geometry 3 (1994), 493--535, \href{https://arxiv.org/abs/alg-geom/9310003}{alg-geom/9310003}.

\bibitem[Batyrev5]{Batyrev5}\hypertarget{bib:Batyrev5}{} V. Batyrev, \emph{Toric degenerations of Fano varieties and constructing mirror manifolds}, preprint, 1997, \href{https://arxiv.org/abs/alg-geom/9712034}{math.AG/9712034}.

\bibitem[Batyrev6]{Batyrev6}\hypertarget{bib:Batyrev6}{}
V. Batyrev, On the classification of toric Fano 4-folds, preprint, 1998,
\href{https://arxiv.org/abs/math.AG/9801107}{math.AG/9801107}.

\bibitem[BB1]{BB1}\hypertarget{bib:BB1}{}
V. Batyrev and L. Borisov, On Calabi--Yau complete intersections in toric varieties,
in Higher-Dimensional Complex Varieties (Trento, 1994), de Gruyter, Berlin, 1996,
\href{https://arxiv.org/abs/alg-geom/9412017}{alg-geom/9412017}.

\bibitem[BB2]{BB2}\hypertarget{bib:BB2}{}
V. Batyrev and L. Borisov, Mirror duality and string-theoretic Hodge numbers,
Invent. Math. 126 (1996), 183--203, \href{https://arxiv.org/abs/alg-geom/9509009}{alg-geom/9509009}.

\bibitem[BB3]{BB3}\hypertarget{bib:BB3}{}
V. Batyrev and L. Borisov, Dual cones and mirror symmetry for generalized Calabi--Yau
manifolds, in Mirror Symmetry II (B. Greene and S.-T. Yau, eds.),
AMS/IP Stud. Adv. Math. 1, AMS, Providence, RI, 1997, 71--86, \href{https://arxiv.org/abs/alg-geom/9402002}{alg-geom/9402002}.

\bibitem[BCFKvS1]{BCFKvS1}\hypertarget{bib:BCFKvS1}{}
V. Batyrev, I. Ciocan-Fontanine, B. Kim, and D. van Straten, \emph{Conifold transition and mirror symmetry for complete intersections in Grassmannians}, Nuclear Physics \textbf{B514} (1998), 640--666, \href{https://arxiv.org/abs/alg-geom/9710022}{alg-geom/9710022}.

\bibitem[BCFKvS2]{BCFKvS2}\hypertarget{bib:BCFKvS2}{} V. Batyrev, I. Ciocan-Fontanine, B. Kim, and D. van Straten, \emph{Mirror symmetry and toric degenerations of partial flag manifolds}, preprint, 1998, \href{https://arxiv.org/abs/math.AG/9803108}{math.AG/9803108}.

\SourcePage{462}

\bibitem[BC]{BC}\hypertarget{bib:BC}{}
V. Batyrev and D. Cox, On the Hodge structure of projective hypersurfaces
in toric varieties, Duke Math. J. 75 (1994), 293--338,
\href{https://arxiv.org/abs/alg-geom/9306011}{alg-geom/9306011}.

\bibitem[BD]{BD}\hypertarget{bib:BD}{}
V. Batyrev and D. Dais, Strong McKay correspondence, string-theoretic Hodge numbers
and mirror symmetry, Topology 35 (1996), 901--929, \href{https://arxiv.org/abs/alg-geom/9410001}{alg-geom/9410001}.

\bibitem[BvS]{BvS}\hypertarget{bib:BvS}{}
V. Batyrev and D. van Straten, Generalized hypergeometric functions
and rational curves on Calabi--Yau complete intersections in toric varieties,
Comm. Math. Phys. 168 (1995), 493--533, \href{https://arxiv.org/abs/alg-geom/9307010}{alg-geom/9307010}.

\bibitem[BS]{BS}\hypertarget{bib:BS}{}
D. Bayer and M. Stillman, Macaulay: a computer algebra system for commutative algebra and algebraic geometry, Available by anonymous ftp from ftp.math.harvard.edu.

\bibitem[Beauville1]{Beauville1}\hypertarget{bib:Beauville1}{} A. Beauville, \emph{Quantum cohomology of complete intersections}, Mat. Fiz. Anal. Geom. \textbf{2} (1995), 384--398, \href{https://arxiv.org/abs/alg-geom/9501008}{alg-geom/9501008}.

\bibitem[Beauville2]{Beauville2}\hypertarget{bib:Beauville2}{} A. Beauville, \emph{Counting rational curves on K3 surfaces}, preprint, 1997, \href{https://arxiv.org/abs/alg-geom/9701019}{alg-geom/9701019}.

\bibitem[Behrend]{Behrend}\hypertarget{bib:Behrend}{}
K. Behrend, Gromov--Witten invariants in algebraic geometry, Invent. Math. 127 (1997), 601--617, \href{https://arxiv.org/abs/alg-geom/9601011}{alg-geom/9601011}.

\bibitem[BF]{BF}\hypertarget{bib:BF}{}
K. Behrend and B. Fantechi, The intrinsic normal cone, Invent. Math. 128 (1997), 45--88, \href{https://arxiv.org/abs/alg-geom/9601010}{alg-geom/9601010}.

\bibitem[BM]{BM}\hypertarget{bib:BM}{}
K. Behrend and Yu. Manin, Stacks of stable maps and Gromov--Witten invariants, Duke J. Math. 85 (1996), 1--60, \href{https://arxiv.org/abs/alg-geom/9506023}{alg-geom/9506023}.

\bibitem[BEFFGK]{BEFFGK}\hypertarget{bib:BEFFGK}{}
K. Behrend, D. Edidin, B. Fantechi, W. Fulton, L. Göttsche, and A. Kresch, Introduction to Stacks, in preparation.

\bibitem[BK1]{BK1}\hypertarget{bib:BK1}{}
P. Berglund and S. Katz,
Mirror symmetry for hypersurfaces in weighted projective space
and topological couplings,
Nuclear Physics B416 (1994), 481--538, \href{https://arxiv.org/abs/hep-th/9311014}{hep-th/9311014}.

\bibitem[BK2]{BK2}\hypertarget{bib:BK2}{}
P. Berglund and S. Katz, Mirror symmetry constructions: a review, in Mirror Symmetry II
(B. Greene and S.-T. Yau, eds.), AMS/IP Stud. Adv. Math. 1, AMS,
Providence, RI, 1997, 87--113, \href{https://arxiv.org/abs/hep-th/9406008}{hep-th/9406008}.

\bibitem[BKK]{BKK}\hypertarget{bib:BKK}{}
P. Berglund, S. Katz and A. Klemm,
Mirror symmetry and the moduli space for generic hypersurfaces in toric varieties,
Nuclear Physics B456 (1995), 153--204, \href{https://arxiv.org/abs/hep-th/9506091}{hep-th/9506091}.

\bibitem[BCOV1]{BCOV1}\hypertarget{bib:BCOV1}{} M. Bershadsky, S. Cecotti, H. Ooguri and C. Vafa, \emph{Holomorphic anomalies in topological field theories} (with an appendix by S. Katz), Nuclear Physics \textbf{B405} (1993), 279--304, reprinted in \emph{Mirror Symmetry II} (B. Greene and S.-T. Yau, eds.), AMS/IP Stud. Adv. Math. 1, AMS, Providence, RI, 1997, 655--682, \href{https://arxiv.org/abs/hep-th/9302103}{hep-th/9302103}.

\bibitem[BCOV2]{BCOV2}\hypertarget{bib:BCOV2}{}
M. Bershadsky, S. Cecotti, H. Ooguri and C. Vafa,
Kodaira--Spencer theory of gravity and exact results for quantum string amplitudes,
Commun. Math. Phys. 165 (1994), 311--428, \href{https://arxiv.org/abs/hep-th/9309140}{hep-th/9309140}.

\bibitem[BJPSV]{BJPSV}\hypertarget{bib:BJPSV}{} M. Bershadsky, A. Johansen, T. Pantev, V. Sadov and C. Vafa, \emph{F-theory, Geometric Engineering and $N=1$ Dualities}, Nuclear Physics \textbf{B505} (1997), 153--164, \href{https://arxiv.org/abs/hep-th/9612052}{hep-th/9612052}.

\bibitem[BV]{BV}\hypertarget{bib:BV}{} M. Bershadsky and C. Vafa, \emph{Global Anomalies and Geometric Engineering of Critical Theories in Six Dimensions}, preprint, 1997, \href{https://arxiv.org/abs/hep-th/9703167}{hep-th/9703167}.

\bibitem[Bertram]{Bertram}\hypertarget{bib:Bertram}{} A. Bertram, \emph{Quantum Schubert calculus}, Adv. in Math. \textbf{28} (1997), 289--305, \href{https://arxiv.org/abs/alg-geom/9410024}{alg-geom/9410024}.

\bibitem[BCF]{BCF}\hypertarget{bib:BCF}{} A. Bertram, I. Ciocan-Fontanine and W. Fulton, \emph{Quantum multiplication with Schur polynomials}, preprint, 1997, \href{https://arxiv.org/abs/alg-geom/9705024}{alg-geom/9705024}.

\bibitem[BDW]{BDW}\hypertarget{bib:BDW}{}
A. Bertram, G. Daskalopoulos, and R. Wentworth, Gromov invariants for holomorphic maps from Riemann surfaces to Grassmannians, J. AMS 9 (1996), 529--571, \href{https://arxiv.org/abs/alg-geom/9306005}{alg-geom/9306005}.

\bibitem[BDPP]{BDPP}\hypertarget{bib:BDPP}{}
G. Bini, C. De Concini, M. Polito, and C. Procesi, \emph{Givental\textquotesingle s work}, preprint, 1998, \href{https://arxiv.org/abs/math/9805097}{math.AG/9805097}.

\bibitem[BD]{BjoD}\hypertarget{bib:BjoD}{} J. Bjorken and S. Drell, \emph{Relativistic Quantum Fields}, International Series in Pure and Applied Physics, McGraw-Hill, New York 1965.

\SourcePage{463}

\bibitem[Bogomolov]{Bogomolov}\hypertarget{bib:Bogomolov}{}
F. Bogomolov, Hamiltonian Kähler manifolds, Dokl. Akad. Nauk. SSSR 243 (1978),
no. 5, 1101--1104.

\bibitem[Borcea1]{Borcea1}\hypertarget{bib:Borcea1}{}
C. Borcea, On desingularized Horrocks--Mumford quintics, J. Reine Angew. Math. 421 (1991), 23--41.

\bibitem[Borcea2]{Borcea2}\hypertarget{bib:Borcea2}{}
C. Borcea, K3 surfaces with involution and mirror pairs of Calabi--Yau manifolds,
in Mirror Symmetry II (B. Greene and S.-T. Yau, eds.), AMS/IP Stud. Adv. Math. 1,
AMS, Providence, RI, 1997, 717--743.

\bibitem[Borisov1]{Borisov1}\hypertarget{bib:Borisov1}{}
L. Borisov, Towards the mirror symmetry for Calabi--Yau complete intersections
in Gorenstein Fano toric varieties, preprint, 1993, \href{https://arxiv.org/abs/alg-geom/9310001}{math.AG/9310001}.

\bibitem[Borisov2]{Borisov2}\hypertarget{bib:Borisov2}{}
L. Borisov, \emph{On Betti numbers and Chern classes of varieties with trivial odd cohomology groups}, preprint, 1997, \href{https://arxiv.org/abs/alg-geom/9703023}{alg-geom/9703023}.

\bibitem[BL]{BL}\hypertarget{bib:BL}{}
J. Bryan and N. C. Leung, The enumerative geometry of K3 surfaces and modular forms, preprint, 1997, \href{https://arxiv.org/abs/alg-geom/9711031}{alg-geom/9711031}.

\bibitem[BG]{BG}\hypertarget{bib:BG}{}
R. Bryant and P. Griffiths, Some observations on the infinitesimal period relations
for regular threefolds with trivial canonical bundle, in Arithmetic and Geometry
Volume II (M. Artin and J. Tate, eds.), Progress in Math. 39, Birkhäuser,
Boston--Basel--Berlin, 1983.

\bibitem[Cd]{Cd}\hypertarget{bib:Cd}{}
P. Candelas and X. de la Ossa,
Moduli Space of Calabi--Yau manifolds,
Nuclear Physics B355 (1991), 455--481.

\bibitem[CdFKM]{CdFKM}\hypertarget{bib:CdFKM}{}
P. Candelas, X. de la Ossa, A. Font, S. Katz and D. Morrison,
Mirror symmetry for two parameter models---I,
Nuclear Physics B416 (1994), 481--538;
reprinted in Mirror Symmetry II (B. Greene and S.-T. Yau, eds.),
AMS/IP Stud. Adv. Math. 1, AMS, Providence, RI, 1997, 483--543,
\href{https://arxiv.org/abs/hep-th/9308083}{hep-th/9308083}.

\bibitem[CFKM]{CFKM}\hypertarget{bib:CFKM}{}
P. Candelas, A. Font, S. Katz and D. Morrison,
Mirror symmetry for two parameter models---II,
Nuclear Physics B429 (1994), 626--674, \href{https://arxiv.org/abs/hep-th/9403187}{hep-th/9403187}.

\bibitem[CdGP]{CdGP}\hypertarget{bib:CdGP}{}
P. Candelas, X. de la Ossa, P. Green and L. Parkes,
A pair of Calabi--Yau manifolds as an exactly soluble superconformal field theory,
Nuclear Physics B359 (1991), 21--74; and in Essays on Mirror Manifolds
(S.-T. Yau, ed.), International Press, Hong Kong, 1992, 31--95.

\bibitem[CdK]{CdK}\hypertarget{bib:CdK}{}
P. Candelas, X. de la Ossa and S. Katz,
Mirror symmetry for Calabi--Yau hypersurfaces in weighted
\(\mathbb P^4\) and extensions of Landau--Ginzburg theory,
Nuclear Physics B450 (1995), 267--292, \href{https://arxiv.org/abs/hep-th/9412117}{hep-th/9412117}.

\bibitem[CDP]{CDP}\hypertarget{bib:CDP}{}
P. Candelas, E. Derrick and L. Parkes, Generalized Calabi--Yau manifolds and the
mirror of a rigid manifold, Nuclear Physics B407 (1993), 115--154, \href{https://arxiv.org/abs/hep-th/9304045}{hep-th/9304045}.

\bibitem[Caporaso]{Caporaso}\hypertarget{bib:Caporaso}{} L. Caporaso, \emph{Counting curves on surfaces: A guide to new techniques and results}, preprint, 1996, \href{https://arxiv.org/abs/alg-geom/9611029}{alg-geom/9611029}.

\bibitem[CH1]{CH1}\hypertarget{bib:CH1}{} L. Caporaso and J. Harris, \emph{Parameter spaces for curves on surfaces and enumeration of rational curves}, Compositio Math. \textbf{113} (1998), 155--208, \href{https://arxiv.org/abs/alg-geom/9608023}{alg-geom/9608023}.

\bibitem[CH2]{CH2}\hypertarget{bib:CH2}{} L. Caporaso and J. Harris, \emph{Enumerating rational curves: the rational fibration method}, Compositio Math. \textbf{113} (1998), 209--236, \href{https://arxiv.org/abs/alg-geom/9608024}{alg-geom/9608024}.

\bibitem[CH3]{CH3}\hypertarget{bib:CH3}{} L. Caporaso and J. Harris, \emph{Counting plane curves of any genus}, Invent. Math. \textbf{131} (1998), 345--392, \href{https://arxiv.org/abs/alg-geom/9608025}{alg-geom/9608025}.

\bibitem[CCD]{CCD}\hypertarget{bib:CCD}{} E. Cattani, D. Cox, A. Dickenstein, \emph{Residues in toric varieties}, Compositio Math. \textbf{108} (1997), 35--76, \href{https://arxiv.org/abs/alg-geom/9506024}{alg-geom/9506024}.

\bibitem[CaK]{CaK}\hypertarget{bib:CaK}{}
E. Cattani and A. Kaplan, Degenerating variations of Hodge structures, in Actes du
Colloque de Théorie de Hodge, Luminy 1987, Astérisque 179--180,
Soc. Math. France, Paris, 1989, 67--96.

\bibitem[CKS1]{CKS1}\hypertarget{bib:CKS1}{}
E. Cattani, A. Kaplan, and W. Schmid, Degeneration of Hodge structures, Annals of Math. 123 (1986), 457--535.

\bibitem[CKS2]{CKS2}\hypertarget{bib:CKS2}{}
E. Cattani, A. Kaplan, and W. Schmid, $L^2$ and intersection cohomologies for a polarizable variation of Hodge structure, Inv. Math. 87 (1987), 217--252.

\bibitem[CGGK]{CGGK}\hypertarget{bib:CGGK}{}
T.-M. Chiang, B. Greene, M. Gross and Y. Kanter, Black hole condensation and the web of Calabi--Yau manifolds, Nuclear Physics Proc. Supp. 46 (1996), 82--95, \href{https://arxiv.org/abs/hep-th/9511204}{hep-th/9511204}.

\SourcePage{464}

\bibitem[CiocanF]{CiocanF}\hypertarget{bib:CiocanF}{} I. Ciocan-Fontanine, \emph{Quantum cohomology of flag varieties}, Internat. Math. Res. Notices (1995), 263--277, \href{https://arxiv.org/abs/alg-geom/9505002}{math.AG/9505002}.

\bibitem[Clemens]{Clemens}\hypertarget{bib:Clemens}{} C. H. Clemens, \emph{Double Solids}, Advances in Math. \textbf{47} (1983) 671--689.

\bibitem[CG]{CG}\hypertarget{bib:CG}{}
C. Clemens and P. Griffiths, The intermediate Jacobian of the cubic threefold,
Annals of Math. 95 (1972), 281--356.

\bibitem[ClK]{ClK}\hypertarget{bib:ClK}{}
H. Clemens and H. Kley, Counting Curves which Move with Threefolds, preprint, 1998, \href{https://arxiv.org/abs/math.AG/9807109}{math.AG/9807109}.

\bibitem[CKM]{CKM}\hypertarget{bib:CKM}{} H. Clemens, J. Kollár, and S. Mori, \emph{Higher dimensional complex geometry}, Astérisque \textbf{166} (1988), Soc. Math. France, Paris, 1989.

\bibitem[CL]{CL}\hypertarget{bib:CL}{}
E. Coddington and N. Levinson, Theory of Ordinary Differential Equations,
McGraw-Hill, New York--Toronto--London, 1955.

\bibitem[CJ]{CJ}\hypertarget{bib:CJ}{} A. Collino and M. Jinzenji, \emph{On the structure of small quantum cohomology rings for projective hypersurfaces}, preprint, 1996, \href{https://arxiv.org/abs/hep-th/9611053}{hep-th/9611053}.

\bibitem[Cox]{Cox}\hypertarget{bib:Cox}{}
D. Cox, The homogeneous coordinate ring of a toric variety,
J. Algebraic Geometry 4 (1995), 17--50, \href{https://arxiv.org/abs/alg-geom/9210008}{alg-geom/9210008}.

\bibitem[CM]{CM}\hypertarget{bib:CM}{} B. Crauder and R. Miranda, \emph{Quantum cohomology of rational curves}, in \emph{The Moduli Space of Curves (Texel Island, 1994)} (R. Dijkgraaf, C. Faber and G. van der Geer, eds.), Progress in Math. \textbf{129}, Birkhäuser, Boston-Basel-Berlin, 1995, 33--80.

\bibitem[Danilov]{Danilov}\hypertarget{bib:Danilov}{}
V. Danilov, The geometry of toric varieties,
Russian Math. Surveys 33 (1978), 97--154.

\bibitem[DK]{DK}\hypertarget{bib:DK}{}
V. Danilov and A. Khovanskii, Newton polyhedra and an algorithm for computing
Hodge-Deligne numbers, Math. USSR Izvestiya 29 (1987), 279--298.

\bibitem[Deligne1]{Deligne1}\hypertarget{bib:Deligne1}{}
P. Deligne, Équations Différentielles à Points Singuliers Réguliers,
Lecture Notes in Mathematics 163, Springer-Verlag, New York, 1970.

\bibitem[Deligne2]{Deligne2}\hypertarget{bib:Deligne2}{}
P. Deligne, Local behavior of Hodge structures at infinity, in Mirror Symmetry II
(B. Greene and S.-T. Yau, eds.), AMS/IP Stud. Adv. Math. 1, AMS,
Providence, RI, 1997, 683--699.

\bibitem[DEFJKMMW]{DEFJKMMW}\hypertarget{bib:DEFJKMMW}{} \emph{Quantum Fields and Strings: A Course for Mathematicians}, (P. Deligne, P. Etingof, D. S. Freed, L. Jeffrey, D. Kazhdan, J. Morgan, D. R. Morrison and E. Witten, eds.), two volumes, American Mathematical Society, to appear. Lecture notes and problem sets from the 1996--97 IAS Quantum Field Theory program are at \url{http://www.math.ias.edu/~drm/QFT/}.

\bibitem[DM]{DM}\hypertarget{bib:DM}{}
P. Deligne and D. Mumford, The irreducibility of the space of curves of a given genus, Publ. Math. IHES 36 (1969), 75--110.

\bibitem[DHSS]{DHSS}\hypertarget{bib:DHSS}{}
J. De Loera, S. Hosten, F. Santos and B. Sturmfels, The polytope of all triangulations
of a point configuration, Documenta Mathematica 1 (1996), 103--119,
\url{http://www.mathematik.uni-bielefeld.de/documenta}.

\bibitem[Demazure]{Demazure}\hypertarget{bib:Demazure}{}
M. Demazure, Sous-groupes algébriques de rang maximum du groupe de Cremona,
Ann. Sci. École Norm. Sup. 3 (1970), 507--588.

\bibitem[DI]{DI}\hypertarget{bib:DI}{} P. Di Francesco and C. Itzykson, \emph{Quantum intersection rings}, in \emph{The Moduli Space of Curves (Texel Island, 1994)} (R. Dijkgraaf, C. Faber and G. van der Geer, eds.), Progress in Math. \textbf{129}, Birkhäuser, Boston-Basel-Berlin, 1995, 81--148.

\bibitem[DVV1]{DVV1}\hypertarget{bib:DVV1}{} R. Dijkgraaf, E. Verlinde, and H. Verlinde, \emph{Topological Strings in $d<1$}, Nuclear Physics \textbf{B352} (1991), 59--86.

\bibitem[DVV2]{DVV2}\hypertarget{bib:DVV2}{} R. Dijkgraaf, E. Verlinde, and H. Verlinde, \emph{Notes on Topological String Theory and 2D Gravity}, in \emph{String Theory and Quantum Gravity}, Proceedings of the Trieste Spring School 1990, (M. Green et al., eds.), World Scientific, Singapore, 1991, 91--156.

\bibitem[Distler]{Distler}\hypertarget{bib:Distler}{} J. Distler, \emph{Notes on $N=2$ $\sigma$-models}, Lecture given at 1992 Trieste Spring School, 1992, \href{https://arxiv.org/abs/hep-th/9212062}{hep-th/9212062}.

\bibitem[Dolgachev1]{Dolgachev1}\hypertarget{bib:Dolgachev1}{}
I. Dolgachev, Weighted projective varieties, in Group Actions and Vector Fields
(J. B. Carrell ed.), Lecture Notes in Mathematics 956, Springer-Verlag,
New York, 1982, 34--71.

\SourcePage{465}

\bibitem[Dolgachev2]{Dolgachev2}\hypertarget{bib:Dolgachev2}{}
I. Dolgachev, Integral quadratic forms: Applications to algebraic geometry
(d'après Nikulin), Séminaire Bourbaki 1982/83, Astérisque 105--106,
Soc. Math. France, Paris, 1983, 251--278.

\bibitem[Dolgachev3]{Dolgachev3}\hypertarget{bib:Dolgachev3}{} I. Dolgachev, \emph{Mirror symmetry for lattice polarized K3 surfaces}, Algebraic geometry, 4, J. Math. Sci. \textbf{81} (1996), 2599--2630, \href{https://arxiv.org/abs/alg-geom/9502005}{alg-geom/9502005}.

\bibitem[Donagi]{Donagi}\hypertarget{bib:Donagi}{} R. Donagi, \emph{Principal bundles on elliptic fibrations}, Asian J. Math. \textbf{1} (1997), 214--223, \href{https://arxiv.org/abs/alg-geom/9702002}{alg-geom/9702002}.

\bibitem[Dubrovin1]{Dubrovin1}\hypertarget{bib:Dubrovin1}{} B. Dubrovin, \emph{Integrable systems in topological field theory}, Nuclear Physics \textbf{B379} (1992), \href{https://people.sissa.it/~dubrovin/Papers/1992/dubrovin_1992_npb.pdf}{627--689}.

\bibitem[Dubrovin2]{Dubrovin2}\hypertarget{bib:Dubrovin2}{} B. Dubrovin, \emph{The geometry of 2D topological field theories}, in \emph{Integrable systems and quantum groups}, Lecture Notes in Mathematics \textbf{1620}, Springer-Verlag, New York, 1996, 120--348.

\bibitem[Ebeling]{Ebeling}\hypertarget{bib:Ebeling}{} W. Ebeling, \emph{Strange duality, mirror symmetry, and the Leech lattice}, preprint, 1996, \href{https://arxiv.org/abs/alg-geom/9612010}{alg-geom/9612010}.

\bibitem[EG]{EG}\hypertarget{bib:EG}{}
D. Edidin and W. Graham, Equivariant Intersection Theory, Inv. Math. (to appear).

\bibitem[EHX]{EHX}\hypertarget{bib:EHX}{} T. Eguchi, K. Hori and C.-S. Xiong, \emph{Quantum cohomology and Virasoro algebra}, Phys. Lett. \textbf{B402} (1997), 71--80, \href{https://arxiv.org/abs/hep-th/9703086}{hep-th/9703086}.

\bibitem[EJX]{EJX}\hypertarget{bib:EJX}{}
T. Eguchi, M. Jinzenji, and C.-S. Xiong, \emph{Quantum Cohomology and Free Field Representation}, Nuclear Physics \textbf{B510} (1998), 608--622, \href{https://arxiv.org/abs/hep-th/9709152}{hep-th/9709152}.

\bibitem[EX]{EX}\hypertarget{bib:EX}{}
T. Eguchi and C.-S. Xiong, \emph{Quantum Cohomology at Higher Genus: Topological Recursion Relations and Virasoro Conditions}, Adv. Theor. Math. Phys. \textbf{2} (1998), 219--229, \href{https://arxiv.org/abs/hep-th/9801010}{hep-th/9801010}.

\bibitem[EJS]{EJS}\hypertarget{bib:EJS}{} T. Ekedahl, T. Johnsen and D. Sommervoll, \emph{Isolated rational curves on K3-fibered Calabi-Yau threefolds}, preprint, 1997, \href{https://arxiv.org/abs/alg-geom/9710010}{alg-geom/9710010}.

\bibitem[ES1]{ES1}\hypertarget{bib:ES1}{}
G. Ellingsrud and S. Strømme,
The number of twisted cubics on the general quintic threefold,
Math. Scand. 76 (1995).

\bibitem[ES2]{ES2}\hypertarget{bib:ES2}{} G. Ellingsrud and S. Strømme, \emph{Bott's formula and enumerative geometry}, J. AMS \textbf{9} (1996), 175--193, \href{https://arxiv.org/abs/alg-geom/9411005}{alg-geom/9411005}.

\bibitem[EK1]{EK1}\hypertarget{bib:EK1}{} L. Ernström and G. Kennedy, \emph{Recursive formulas for the characteristic numbers of rational plane curves}, J. Algebraic Geometry \textbf{7} (1998), 141--181, \href{https://arxiv.org/abs/alg-geom/9604019}{alg-geom/9604019}.

\bibitem[EK2]{EK2}\hypertarget{bib:EK2}{} L. Ernström and G. Kennedy, \emph{Contact cohomology of the projective plane}, preprint, 1997, \href{https://arxiv.org/abs/alg-geom/9703013}{alg-geom/9703013}.

\bibitem[FGvS]{FGvS}\hypertarget{bib:FGvS}{} B. Fantechi, L. Göttsche, and D. van Straten, \emph{Euler number of the compactified Jacobian and multiplicity of rational curves}, preprint, 1997, \href{https://arxiv.org/abs/alg-geom/9708012}{alg-geom/9708012}.

\bibitem[Font]{Font}\hypertarget{bib:Font}{}
A. Font, Periods and duality symmetries in Calabi--Yau compactifications,
Nuclear Physics B391 (1993), 358--388, \href{https://arxiv.org/abs/hep-th/9203084}{hep-th/9203084}.

\bibitem[FMW1]{FMW1}\hypertarget{bib:FMW1}{} R. Friedman, J. Morgan and E. Witten, \emph{Vector bundles and F theory}, Comm. Math. Phys. \textbf{187}, (1997), 679--743, \href{https://arxiv.org/abs/hep-th/9701162}{hep-th/9701162}.

\bibitem[FMW2]{FMW2}\hypertarget{bib:FMW2}{} R. Friedman, J. Morgan and E. Witten, \emph{Vector bundles over elliptic fibrations}, preprint, 1997, \href{https://arxiv.org/abs/alg-geom/9709029}{alg-geom/9709029}.

\bibitem[Fukaya]{Fukaya}\hypertarget{bib:Fukaya}{} K. Fukaya, \emph{Morse homotopy, $A^\infty$-category, and Floer homologies}, in \emph{Proceedings of GARC Workshop on Geometry and Topology '93 (Seoul, 1993)} (H. Kim, ed.), Lecture Notes Ser., 18, Seoul Nat. Univ., Seoul, 1993, 1--102.

\bibitem[FO]{FO}\hypertarget{bib:FO}{}
K. Fukaya and K. Ono, Arnold conjecture and Gromov--Witten invariants, preprint, 1996.

\bibitem[Fulton1]{Fulton1}\hypertarget{bib:Fulton1}{}
W. Fulton, Introduction to Intersection Theory in Algebraic Geometry, CBMS Series, Number 54, AMS, Providence, RI, 1984.

\bibitem[Fulton2]{Fulton2}\hypertarget{bib:Fulton2}{}
W. Fulton, Intersection Theory, Springer-Verlag, New York-Berlin-Heidelberg, 1984.

\bibitem[Fulton3]{Fulton3}\hypertarget{bib:Fulton3}{}
W. Fulton, Introduction to Toric Varieties,
Princeton University Press, Princeton, 1993.

\bibitem[FP]{FP}\hypertarget{bib:FP}{}
W. Fulton and R. Pandharipande, Notes on stable maps and quantum cohomology, in Algebraic Geometry---Santa Cruz 1995 (Kollár, Lazarsfeld, and Morrison, eds.), Proceedings of the 1995 Summer Research Institute in Algebraic Geometry, Proceedings of Symposia in Pure Mathematics, Volume 62 Part 2, 45--96.

\SourcePage{466}

\bibitem[Gathmann1]{Gathmann1}\hypertarget{bib:Gathmann1}{} A. Gathmann, \emph{Counting rational curves with multiple points and Gromov-Witten invariants of blow-ups}, preprint, 1996, \href{https://arxiv.org/abs/alg-geom/9609010}{alg-geom/9609010}.

\bibitem[Gathmann2]{Gathmann2}\hypertarget{bib:Gathmann2}{} A. Gathmann, \emph{Gromov-Witten invariants of blow-ups}, preprint, 1998, \href{https://arxiv.org/abs/math.AG/9804043}{math.AG/9804043}.

\bibitem[GKZ1]{GKZ1}\hypertarget{bib:GKZ1}{}
I. Gelfand, M. Kapranov and A. Zelevinsky, Generalized Euler integrals and
$\mathcal A$-hypergeometric functions, Adv. in Math. 84 (1990), 255--271.

\bibitem[GKZ2]{GKZ2}\hypertarget{bib:GKZ2}{}
I. Gelfand, M. Kapranov and A. Zelevinsky, Discriminants, Resultants and
Multidimensional Determinants, Birkhäuser, Boston--Basel--Berlin, 1994.

\bibitem[Getzler1]{Getzler1}\hypertarget{bib:Getzler1}{}
E. Getzler, Intersection theory on $\overline M_{1,4}$ and elliptic Gromov--Witten invariants, J. AMS 10 (1997), 973--998, \href{https://arxiv.org/abs/alg-geom/9612004}{alg-geom/9612004}.

\bibitem[Getzler2]{Getzler2}\hypertarget{bib:Getzler2}{}
E. Getzler, \emph{Topological Recursion Relations in Genus 2}, preprint, 1998, \href{https://arxiv.org/abs/math/9801003}{math.AG/9801003}.

\bibitem[GeP]{GeP}\hypertarget{bib:GeP}{} E. Getzler and R. Pandharipande, \emph{The number of elliptic space curves}, in preparation.

\bibitem[Ginsparg]{Ginsparg}\hypertarget{bib:Ginsparg}{} P. Ginsparg, \emph{Applied conformal field theory}, in \emph{Fields, Strings, and Critical Phenomena: Proceedings of the Les Houches Summer School 1988}, Elsevier Science Publishers, 1989.

\bibitem[Givental1]{Givental1}\hypertarget{bib:Givental1}{}
A. Givental, Homological geometry I, projective hypersurfaces, Selecta Math. 1 (1995), 325--345.

\bibitem[Givental2]{Givental2}\hypertarget{bib:Givental2}{}
A. Givental, Equivariant Gromov--Witten Invariants,
Internat. Math. Res. Notices (1996), 613--663, \href{https://arxiv.org/abs/alg-geom/9603021}{alg-geom/9603021}.

\bibitem[Givental3]{Givental3}\hypertarget{bib:Givental3}{} A. Givental, \emph{Stationary phase integrals, quantum Toda lattices, flag manifolds and the mirror conjecture}, preprint, 1996, \href{https://arxiv.org/abs/alg-geom/9612001}{alg-geom/9612001}.

\bibitem[Givental4]{Givental4}\hypertarget{bib:Givental4}{}
A. Givental, A mirror theorem for toric complete intersections,
preprint, 1997, \href{https://arxiv.org/abs/alg-geom/9701016}{alg-geom/9701016}.

\bibitem[Givental5]{Givental5}\hypertarget{bib:Givental5}{}
A. Givental, \emph{Elliptic Gromov--Witten invariants and the generalized mirror conjecture}, preprint, 1998, \href{https://arxiv.org/abs/math/9803053}{math.AG/9803053}.

\bibitem[Givental6]{Givental6}\hypertarget{bib:Givental6}{}
A. Givental, \emph{The mirror formula for quintic threefolds}, \href{https://arxiv.org/abs/math/9807070}{math.AG/9807070}.

\bibitem[GK]{GK}\hypertarget{bib:GK}{} A. Givental and B. Kim, \emph{Quantum cohomology of flag manifolds and Toda lattices}, Comm. Math. Phys. \textbf{168} (1995), 609--641.

\bibitem[GL]{GL}\hypertarget{bib:GL}{} N. Gonciulea and V. Lakshmibai, \emph{Degenerations of flag and Schubert varieties to toric varieties}, Transform. Groups \textbf{1} (1996), 215--248.

\bibitem[Göttsche]{Gottsche}\hypertarget{bib:Gottsche}{} L. Göttsche, \emph{A conjectural generating function for numbers of curves on surfaces}, preprint, 1997, \href{https://arxiv.org/abs/alg-geom/9711012}{math.AG/9711012}.

\bibitem[GöP]{GoP}\hypertarget{bib:GoP}{} L. Göttsche and R. Pandharipande, \emph{The quantum cohomology of blow-ups of $\mathbb P^2$ and enumerative geometry}, J. Diff. Geom. \textbf{48} (1998), 61--90, \href{https://arxiv.org/abs/alg-geom/9611012}{alg-geom/9611012}.

\bibitem[GPa]{GPa}\hypertarget{bib:GPa}{} T. Graber and R. Pandharipande, \emph{Localization of virtual classes}, preprint, 1997, \href{https://arxiv.org/abs/alg-geom/9708001}{alg-geom/9708001}.

\bibitem[GM]{GM}\hypertarget{bib:GM}{}
A. Grassi and D. Morrison, Automorphisms and the Kähler cone of certain Calabi--Yau manifolds, Duke Math. J. 71 (1993), 831--838.

\bibitem[GSW]{GSW}\hypertarget{bib:GSW}{}
M. Green, J. Schwarz and E. Witten, Superstring theory, vols. 1 and 2,
Cambridge University Press, Cambridge, 1987.

\bibitem[GMS]{GMS}\hypertarget{bib:GMS}{}
B. Greene, D. Morrison, and A. Strominger, Black hole condensation and the unification of string vacua, Nuclear Physics B451 (1995) 109, \href{https://arxiv.org/abs/hep-th/9504145}{hep-th/9504145}.

\bibitem[GO]{GO}\hypertarget{bib:GO}{} B. Greene and H. Ooguri, \emph{Geometry and quantum field theory: A brief introduction}, in \emph{Mirror Symmetry II} (B. Greene and S.-T. Yau, eds.), AMS/IP Stud. Adv. Math. 1, AMS, Providence, RI, 1997, 3--27.

\bibitem[GPl]{GPl}\hypertarget{bib:GPl}{}
B. Greene and M. Plesser, Duality in Calabi--Yau moduli space,
Nuclear Physics B338 (1990), 15--37.

\bibitem[Griffiths1]{Griffiths1}\hypertarget{bib:Griffiths1}{}
P. Griffiths, On the periods of certain rational integrals I, II,
Ann. of Math. 90 (1969), 460--495, 498--541.

\SourcePage{467}

\bibitem[Griffiths2]{Griffiths2}\hypertarget{bib:Griffiths2}{}
P. Griffiths (ed.), Topics in Transcendental Algebraic Geometry,
Ann. of Math. Stud. 106, Princeton University Press, Princeton, NJ, 1984.

\bibitem[GH]{GH}\hypertarget{bib:GH}{}
P. Griffiths and J. Harris, Principles of Algebraic Geometry, Wiley-Interscience, New York, 1978.

\bibitem[GN1]{GN1}\hypertarget{bib:GN1}{} V. Gritsenko and V. Nikulin, \emph{K3 surfaces, Lorentzian Kac-Moody algebras and mirror symmetry}, Math. Res. Lett. \textbf{3} (1996), 211--229, \href{https://arxiv.org/abs/alg-geom/9510008}{alg-geom/9510008}.

\bibitem[GN2]{GN2}\hypertarget{bib:GN2}{} V. Gritsenko and V. Nikulin, \emph{The arithmetic mirror symmetry and Calabi--Yau manifolds}, preprint, 1996, \href{https://arxiv.org/abs/alg-geom/9612002}{alg-geom/9612002}.

\bibitem[GN3]{GN3}\hypertarget{bib:GN3}{} V. Gritsenko and V. Nikulin, \emph{K3 surfaces with interesting groups of automorphisms}, preprint, 1997, \href{https://arxiv.org/abs/alg-geom/9701011}{alg-geom/9701011}.

\bibitem[Gromov]{Gromov}\hypertarget{bib:Gromov}{} M. Gromov, \emph{Pseudoholomorphic curves in symplectic manifolds}, Invent. Math. \textbf{82} (1985), 307--347.

\bibitem[Gross1]{Gross1}\hypertarget{bib:Gross1}{} M. Gross, \emph{Deforming Calabi--Yau threefolds}, Math. Ann. \textbf{308} (1997), 187--220, \href{https://arxiv.org/abs/alg-geom/9506022}{alg-geom/9506022}.

\bibitem[Gross2]{Gross2}\hypertarget{bib:Gross2}{} M. Gross, \emph{Primitive Calabi--Yau threefolds}, J. Differential Geom. \textbf{45} (1997), 288--318, \href{https://arxiv.org/abs/alg-geom/9512002}{alg-geom/9512002}.

\bibitem[Gross3]{Gross3}\hypertarget{bib:Gross3}{} M. Gross, \emph{Special Lagrangian fibrations I: topology}, preprint, 1997, \href{https://arxiv.org/abs/alg-geom/9710006}{alg-geom/9710006}.

\bibitem[GW]{GW}\hypertarget{bib:GW}{} M. Gross and P. Wilson, \emph{Mirror symmetry via 3-tori for a class of Calabi--Yau threefolds}, Math. Ann. \textbf{309} (1997), 505--531, \href{https://arxiv.org/abs/alg-geom/9608004}{alg-geom/9608004}.

\bibitem[GD]{GD}\hypertarget{bib:GD}{} A. Grothendieck and J. Dieudonné, \emph{Éléments de Géométrie Algébrique, IV, Quatrième Partie}, Publ. Math. I. H. E. S. \textbf{32} (1967).

\bibitem[Guillemin]{Guillemin}\hypertarget{bib:Guillemin}{}
V. Guillemin, Moment Maps and Combinatorial Invariants of Hamiltonian
$T^n$-spaces, Progress in Math. 122, Birkhäuser, Boston--Basel--Berlin, 1994.

\bibitem[HM]{HM}\hypertarget{bib:HM}{}
J. Harris and D. Mumford, On the Kodaira dimension of the moduli space of curves. With an appendix by W. Fulton, Invent. Math. 67 (1982), 23--88.

\bibitem[Hartshorne]{Hartshorne}\hypertarget{bib:Hartshorne}{} R. Hartshorne, \emph{Algebraic Geometry}, Springer-Verlag, 1977.

\bibitem[Hori]{Hori}\hypertarget{bib:Hori}{}
K. Hori, \emph{Constraints for Topological Strings in $D\geq1$}, Nuclear Physics \textbf{B439} (1995), 395, \href{https://arxiv.org/abs/hep-th/9411135}{hep-th/9411135}.

\bibitem[HKTY1]{HKTY1}\hypertarget{bib:HKTY1}{}
S. Hosono, A. Klemm, S. Theisen and S.-T. Yau,
Mirror symmetry, mirror map, and application to Calabi--Yau hypersurfaces,
Comm. Math. Phys. 167 (1995), 301--350, \href{https://arxiv.org/abs/hep-th/9308122}{hep-th/9308122}.

\bibitem[HKTY2]{HKTY2}\hypertarget{bib:HKTY2}{}
S. Hosono, A. Klemm, S. Theisen and S.-T. Yau,
Mirror symmetry, mirror map and applications to complete intersection
Calabi--Yau spaces, Nuclear Physics B433 (1995), 501--554;
reprinted in Mirror Symmetry II (B. Greene and S.-T. Yau, eds.),
AMS/IP Stud. Adv. Math. 1, AMS, Providence, RI, 1997, 545--606,
\href{https://arxiv.org/abs/hep-th/9406055}{hep-th/9406055}.

\bibitem[HLY1]{HLY1}\hypertarget{bib:HLY1}{}
S. Hosono, B. Lian and S.-T. Yau, GKZ-Generalized hypergeometric systems in mirror symmetry of Calabi--Yau hypersurfaces, Comm. Math. Phys. 182 (1996), 535--577, \href{https://arxiv.org/abs/alg-geom/9511001}{alg-geom/9511001}.

\bibitem[HLY2]{HLY2}\hypertarget{bib:HLY2}{}
S. Hosono, B. Lian and S.-T. Yau, Maximal degeneracy points of GKZ systems, J. Amer. Math. Soc. 10 (1997), 427--443, \href{https://arxiv.org/abs/alg-geom/9603014}{alg-geom/9603014}.

\bibitem[HLY3]{HLY3}\hypertarget{bib:HLY3}{} S. Hosono, B. Lian and S.-T. Yau, \emph{Calabi--Yau varieties and pencils of K3 surfaces}, preprint, 1996, \href{https://arxiv.org/abs/alg-geom/9603020}{alg-geom/9603020}.

\bibitem[HSS]{HSS}\hypertarget{bib:HSS}{} S. Hosono, M.-H. Saito and J. Stienstra, \emph{On mirror symmetry conjecture for Schoen's Calabi-Yau 3 folds}, preprint, 1997, \href{https://arxiv.org/abs/alg-geom/9709027}{alg-geom/9709027}.

\bibitem[Hübsch]{Huebsch}\hypertarget{bib:Huebsch}{}
T. Hübsch, Calabi--Yau Manifolds: A Bestiary for Physicists,
World Scientific, Singapore, 1992.

\bibitem[Illusie]{Illusie}\hypertarget{bib:Illusie}{}
L. Illusie, Complexe Cotangent et Déformations I, II, Lecture Notes in Mathematics 239, 283, Springer-Verlag, New York-Berlin-Heidelberg, 1971, 1972.

\bibitem[Ionel]{Ionel}\hypertarget{bib:Ionel}{} E. Ionel, \emph{Genus 1 enumerative invariants in $\mathbb P^n$ with fixed $j$ invariant}, Duke Math. J. \textbf{94} (1998), 279--324, \href{https://arxiv.org/abs/alg-geom/9608030}{alg-geom/9608030}.

\bibitem[IP]{IP}\hypertarget{bib:IP}{} E. Ionel and T. Parker, \emph{The Gromov Invariants of Ruan-Tian and Taubes}, Math. Res. Lett. \textbf{4} (1997), 521--532, \href{https://arxiv.org/abs/alg-geom/9702008}{alg-geom/9702008}.

\bibitem[Iversen]{Iversen}\hypertarget{bib:Iversen}{}
B. Iversen, A fixed point formula for action of tori on algebraic varieties, Inv. Math. 16 (1972) 229--236.

\SourcePage{468}

\bibitem[JK1]{JK1}\hypertarget{bib:JK1}{}
T. Johnsen and S. Kleiman,
Rational curves of degree at most 9 on a general quintic threefold,
Comm. Algebra 24 (1996), 2721--2753, \href{https://arxiv.org/abs/alg-geom/9510015}{alg-geom/9510015}.

\bibitem[JK2]{JK2}\hypertarget{bib:JK2}{} T. Johnsen and S. Kleiman, \emph{Toward Clemens' conjecture in degrees between 10 and 24}, preprint, 1996, \href{https://arxiv.org/abs/alg-geom/9601024}{alg-geom/9601024}.

\bibitem[KK]{KK}\hypertarget{bib:KK}{} A. Kabanov and T. Kimura, \emph{Intersection numbers and rank one cohomological field theories in genus one}, Comm. Math. Phys. \textbf{194} (1998), 651--674, \href{https://arxiv.org/abs/alg-geom/9706003}{alg-geom/9706003}.

\bibitem[KSZ1]{KSZ1}\hypertarget{bib:KSZ1}{}
M. Kapranov, B. Sturmfels, and A. Zelevinsky, Quotients of toric varieties, Math. Ann. 290 (1991), 643--655.

\bibitem[KSZ2]{KSZ2}\hypertarget{bib:KSZ2}{}
M. Kapranov, B. Sturmfels, and A. Zelevinsky, Chow polytopes and general resultants, Duke Math. J. 67 (1992), 189--218.

\bibitem[Katz1]{Katz1}\hypertarget{bib:Katz1}{}
S. Katz, Degenerations of quintic threefolds and their lines. Duke Math. J. 50 (1983) 1127--1135.

\bibitem[Katz2]{Katz2}\hypertarget{bib:Katz2}{}
S. Katz, On the finiteness of rational curves on quintic threefolds,
Compositio Math. 60 (1986), 151--162.

\bibitem[Katz3]{Katz3}\hypertarget{bib:Katz3}{}
S. Katz, Gromov-Witten Invariants via algebraic geometry, Nuclear Physics Proc. Supp. 46 (1996), 108--115, \href{https://arxiv.org/abs/hep-th/9510218}{hep-th/9510218}.

\bibitem[KKV]{KKV}\hypertarget{bib:KKV}{} S. Katz, A. Klemm, and C. Vafa, \emph{Geometric Engineering of Quantum Field Theories}, Nuc. Phys. \textbf{B497} (1997) 173--195, \href{https://arxiv.org/abs/hep-th/9609239}{hep-th/9609239}.

\bibitem[KMV]{KMV}\hypertarget{bib:KMV}{} S. Katz, P. Mayr, and C. Vafa, \emph{Mirror symmetry and Exact Solution of 4D $N=2$ Gauge Theories, I}, Adv. Theor. Math. Phys. \textbf{1} (1998), 53--114, \href{https://arxiv.org/abs/hep-th/9706110}{hep-th/9706110}.

\bibitem[KMP]{KMP}\hypertarget{bib:KMP}{}
S. Katz, D. Morrison and M. Plesser, Enhanced gauge symmetry in type II string theory, Nuclear Physics B477 (1996), 105--140, \href{https://arxiv.org/abs/hep-th/9601108}{hep-th/9601108}.

\bibitem[KQR]{KQR}\hypertarget{bib:KQR}{} S. Katz, Z. Qin and Y. Ruan, \emph{Enumeration of nodal genus-2 plane curves with fixed complex structure}, J. Alg. Geom. \textbf{7} (1998), 569--587, \href{https://arxiv.org/abs/alg-geom/9606014}{alg-geom/9606014}.

\bibitem[KS]{KaS}\hypertarget{bib:KaS}{} S. Katz and S. Strømme, \textsc{Schubert}, a Maple package for intersection theory, \url{http://www.math.okstate.edu/~katz/schubert.html}.

\bibitem[KV]{KV}\hypertarget{bib:KV}{} S. Katz and C. Vafa, \emph{Geometric Engineering of $N=1$ Quantum Field Theories}, Nucl. Phys. \textbf{B497} (1997) 196--204, \href{https://arxiv.org/abs/hep-th/9611090}{hep-th/9611090}.

\bibitem[Kaufmann]{Kaufmann}\hypertarget{bib:Kaufmann}{} R. Kaufmann, \emph{The intersection form in $H^*(\overline M_{0n})$ and the explicit Künneth formula in quantum cohomology}, Internat. Math. Res. Notices (1996), 929--952, \href{https://arxiv.org/abs/alg-geom/9608026}{alg-geom/9608026}.

\bibitem[Kawamata1]{Kawamata1}\hypertarget{bib:Kawamata1}{}
Y. Kawamata, Crepant blowing-up of 3-dimensional canonical singularities and its application to degenerations of surfaces, Ann. Math. 127 (1988), 93--163.

\bibitem[Kawamata2]{Kawamata2}\hypertarget{bib:Kawamata2}{}
Y. Kawamata, On the cone of divisors of Calabi--Yau fiber spaces, Internat. J. Math. 8 (1997), 665--687, \href{https://arxiv.org/abs/alg-geom/9701006}{alg-geom/9701006}.

\bibitem[Keel]{Keel}\hypertarget{bib:Keel}{} S. Keel, \emph{Intersection theory on moduli spaces of stable $n$-pointed curves of genus zero}, Trans. AMS \textbf{330} (1992), 545--574.

\bibitem[Kim1]{Kim1}\hypertarget{bib:Kim1}{} B. Kim, \emph{Quantum cohomology of flag manifolds $G/B$ and quantum Toda lattices}, preprint, 1996, \href{https://arxiv.org/abs/alg-geom/9607001}{alg-geom/9607001}.

\bibitem[Kim2]{Kim2}\hypertarget{bib:Kim2}{}
B. Kim, \emph{Quantum Hyperplane Section Theorem for Homogeneous Spaces}, preprint, 1998, \href{https://arxiv.org/abs/alg-geom/9712008}{alg-geom/9712008}.

\bibitem[KlM]{KlM}\hypertarget{bib:KlM}{} A. Klemm and P. Mayr, \emph{Strong coupling singularities and non-abelian gauge symmetries in $N=2$ string theory}, Nuclear Physics \textbf{B469} (1996) 37, \href{https://arxiv.org/abs/hep-th/9601014}{hep-th/9601014}.

\bibitem[KS]{KS}\hypertarget{bib:KS}{}
A. Klemm and R. Schimmrigk, Landau-Ginzburg string vacua,
Nuclear Physics B411 (1994), 559--583, \href{https://arxiv.org/abs/hep-th/9204060}{hep-th/9204060}.

\bibitem[KT1]{KT1}\hypertarget{bib:KT1}{}
A. Klemm and S. Theisen,
Considerations of one modulus Calabi--Yau compactifications:
Picard--Fuchs equations, Kähler potentials and mirror maps,
Nuclear Physics B389 (1993), 153--180.

\bibitem[KT2]{KT2}\hypertarget{bib:KT2}{}
A. Klemm and S. Theisen, Mirror maps and instanton sums for complete intersections in weighted projective space, Mod. Phys. Lett. A9 (1994), 1807--1818, \href{https://arxiv.org/abs/hep-th/9304034}{hep-th/9304034}.

\bibitem[Knudsen]{Knudsen}\hypertarget{bib:Knudsen}{}
F. Knudsen, Projectivity of the moduli space of stable curves, II, Math. Scand. 52 (1983), \href{https://www.mscand.dk/article/view/12001}{161--199}.

\SourcePage{469}

\bibitem[Kobayashi]{Kobayashi}\hypertarget{bib:Kobayashi}{} M. Kobayashi, \emph{Duality of weights, mirror symmetry and Arnold's strange duality}, preprint, 1995, \href{https://arxiv.org/abs/alg-geom/9502004}{alg-geom/9502004}.

\bibitem[Kollár]{Kollar}\hypertarget{bib:Kollar}{}
J. Kollár, Flops, Nagoya Math. J. 113 (1989), 14--36.

\bibitem[Kontsevich1]{Kontsevich1}\hypertarget{bib:Kontsevich1}{}
M. Kontsevich, \emph{Intersection theory on the moduli space of curves and the matrix Airy function}, Comm. Math. Phys. \textbf{147} (1992), 1--23.

\bibitem[Kontsevich2]{Kontsevich2}\hypertarget{bib:Kontsevich2}{}
M. Kontsevich, Enumeration of rational curves via torus actions,
in The Moduli Space of Curves (Texel Island, 1994)
(R. Dijkgraaf, C. Faber and G. van der Geer, eds.),
Progress in Math. 129, Birkhäuser, Boston--Basel--Berlin, 1995,
335--368, \href{https://arxiv.org/abs/hep-th/9405035}{hep-th/9405035}.

\bibitem[Kontsevich3]{Kontsevich3}\hypertarget{bib:Kontsevich3}{}
M. Kontsevich, Homological algebra of mirror symmetry,
in Proceedings of the International Congress of Mathematicians,
1994, Zürich, Birkhäuser, Boston--Basel--Berlin, 1995,
120--139, \href{https://arxiv.org/abs/alg-geom/9411018}{alg-geom/9411018}.

\bibitem[KoM1]{KoM1}\hypertarget{bib:KoM1}{}
M. Kontsevich and Yu. Manin, Gromov--Witten classes, quantum cohomology, and enumerative geometry, Comm. Math. Phys. 164 (1994) 525--562, reprinted in Mirror Symmetry II (B. Greene and S.-T. Yau, eds.), AMS/IP Stud. Adv. Math. 1, AMS, Providence, RI, 1997, 607--653, \href{https://arxiv.org/abs/hep-th/9402147}{hep-th/9402147}.

\bibitem[KoM2]{KoM2}\hypertarget{bib:KoM2}{} M. Kontsevich and Yu. Manin, \emph{Quantum cohomology of a product (with Appendix by R. Kaufmann)}, Invent. Math. \textbf{124} (1996), 313--339.

\bibitem[KoM3]{KoM3}\hypertarget{bib:KoM3}{}
M. Kontsevich and Yu. Manin, \emph{Relations between the correlators of the topological sigma-model coupled to gravity}, preprint, 1997, \href{https://arxiv.org/abs/alg-geom/9708024}{alg-geom/9708024}.

\bibitem[Kostant]{Kostant}\hypertarget{bib:Kostant}{} B. Kostant, \emph{Flag manifold quantum cohomology, the Toda lattice, and the representation of highest weight $\rho$}, Selecta Math. New Series \textbf{2} (1996), 43--91.

\bibitem[Kresch]{Kresch}\hypertarget{bib:Kresch}{} A. Kresch, \emph{Associativity relations in quantum cohomology}, preprint, 1997, \href{https://arxiv.org/abs/alg-geom/9703015}{alg-geom/9703015}.

\bibitem[KS2]{KS2}\hypertarget{bib:KS2}{}
M. Kreuzer and H. Skarke, No mirror symmetry in Landau-Ginzburg Spectra!,
Nuclear Physics B388 (1993), 113--130, \href{https://arxiv.org/abs/hep-th/9205004}{hep-th/9205004}.

\bibitem[KS3]{KS3}\hypertarget{bib:KS3}{}
M. Kreuzer and H. Skarke, On the classification of reflexive polyhedra,
Comm. Math. Phys. 185 (1997), 495--508, \href{https://arxiv.org/abs/hep-th/9512204}{hep-th/9512204}.

\bibitem[Kuranishi]{Kuranishi}\hypertarget{bib:Kuranishi}{}
M. Kuranishi, On the locally complete families of complex analytic structures, Ann. Math. 75 (1962), 536--577.

\bibitem[Landman]{Landman}\hypertarget{bib:Landman}{}
A. Landman, On the Picard-Lefschetz transformations,
Trans. Amer. Math. Soc. 181 (1973), 89--126.

\bibitem[LM]{LM}\hypertarget{bib:LM}{} H. Lawson, Jr. and M.-L. Michelsohn, \emph{Spin Geometry}, Princeton University Press, Princeton, NJ, 1989.

\bibitem[LVW]{LVW}\hypertarget{bib:LVW}{} W. Lerche, C. Vafa, and N. Warner, Nuclear Physics \textbf{B324} (1989) 427.

\bibitem[LTi1]{LTi1}\hypertarget{bib:LTi1}{}
J. Li and G. Tian, The quantum cohomology of homogeneous varieties, J. Alg. Geom. 6 (1997), 269--305, \href{https://arxiv.org/abs/alg-geom/9504009}{alg-geom/9504009}.

\bibitem[LTi2]{LTi2}\hypertarget{bib:LTi2}{}
J. Li and G. Tian, Virtual moduli cycles and Gromov--Witten invariants of algebraic varieties, J. AMS 11 (1998), 119--174, \href{https://arxiv.org/abs/alg-geom/9602007}{alg-geom/9602007}.

\bibitem[LTi3]{LTi3}\hypertarget{bib:LTi3}{}
J. Li and G. Tian, Virtual moduli cycles and Gromov--Witten invariants of general symplectic manifolds, preprint, 1996, \href{https://arxiv.org/abs/alg-geom/9608032}{alg-geom/9608032}.

\bibitem[LTi4]{LTi4}\hypertarget{bib:LTi4}{}
J. Li and G. Tian, Comparison of the algebraic and symplectic definitions of GW invariants, preprint, 1997, \href{https://arxiv.org/abs/alg-geom/9712035}{alg-geom/9712035}.

\bibitem[LLY]{LLY}\hypertarget{bib:LLY}{}
B. Lian, K. Liu and S.-T. Yau, Mirror principle I,
Asian J. Math. 1 (1997), 729--763, \href{https://arxiv.org/abs/alg-geom/9712011}{alg-geom/9712011}.

\bibitem[LY1]{LY1}\hypertarget{bib:LY1}{}
B. Lian and S.-T. Yau,
Arithmetic properties of mirror maps and quantum couplings,
Comm. Math. Phys. 176 (1996), 163--191, \href{https://arxiv.org/abs/hep-th/9411234}{hep-th/9411234}.

\bibitem[LY2]{LY2}\hypertarget{bib:LY2}{}
B. Lian and S.-T. Yau, Mirror maps, modular relations and
hypergeometric series I, preprint, 1995, \href{https://arxiv.org/abs/hep-th/9507151}{hep-th/9507151}.

\bibitem[LTe]{LTe}\hypertarget{bib:LTe}{}
A. Libgober and J. Teitelbaum, Lines on Calabi--Yau complete intersections,
mirror symmetry, and Picard-Fuchs equations, Int. Math. Res. Not. 1 (1993),
29--39, \href{https://arxiv.org/abs/alg-geom/9301001}{alg-geom/9301001}.

\SourcePage{470}

\bibitem[LiuT]{LiuT}\hypertarget{bib:LiuT}{}
X. Liu and G. Tian, \emph{Virasoro Constraints for Quantum Cohomology}, preprint, 1998, \href{https://arxiv.org/abs/math/9806028}{math.AG/9806028}.

\bibitem[Looijenga]{Looijenga}\hypertarget{bib:Looijenga}{}
E. Looijenga, New compactifications of locally symmetric varieties, in Proceedings of the 1984 Vancouver Conference in Algebraic Geometry (J. Carrell et al., eds.), CMS Conference Proceedings 6, AMS, Providence, RI, 1986, 341--364.

\bibitem[Manin1]{Manin1}\hypertarget{bib:Manin1}{} Yu. Manin, \emph{Gauge Field Theory and Complex Geometry}, Grund. Math. Wiss. 289, Springer-Verlag, Berlin 1988.

\bibitem[Manin2]{Manin2}\hypertarget{bib:Manin2}{} Yu. Manin, \emph{Generating functions in algebraic geometry and sums over trees}, in \emph{The Moduli Space of Curves (Texel Island, 1994)} (R. Dijkgraaf, C. Faber and G. van der Geer, eds.), Progr. Math. \textbf{129}, Birkhäuser, Boston-Basel-Berlin, 1995.

\bibitem[Manin3]{Manin3}\hypertarget{bib:Manin3}{} Yu. Manin, \emph{Frobenius manifolds, quantum cohomology, and moduli spaces (Chapters I, II, III)}, MPI preprint 96-113, Max-Planck-Institut für Mathematik, Bonn, 1996.

\bibitem[Mayr]{Mayr}\hypertarget{bib:Mayr}{} P. Mayr, \emph{Geometric construction of $N=2$ gauge theories}, preprint, 1998, \href{https://arxiv.org/abs/hep-th/9807096}{hep-th/9807096}.

\bibitem[MS]{MS}\hypertarget{bib:MS}{}
D. McDuff and D. Salamon, $J$-holomorphic Curves and Quantum Cohomology, University Lecture Series 6, AMS, Providence, RI, 1994.

\bibitem[Meurer]{Meurer}\hypertarget{bib:Meurer}{} P. Meurer, \emph{The number of rational quartics on Calabi-Yau hypersurfaces in weighted projective space $\mathbb P(2,1^4)$}, Math. Scand. \textbf{78} (1996), 63--83, \href{https://arxiv.org/abs/alg-geom/9409001}{alg-geom/9409001}.

\bibitem[MR]{MR}\hypertarget{bib:MR}{}
T. Mora and L. Robbiano, The Gröbner fan of an ideal, J. Symbolic Comput. 6 (1988), 183--208.

\bibitem[Morrison1]{Morrison1}\hypertarget{bib:Morrison1}{}
D. Morrison, Picard--Fuchs equations and mirror maps for hypersurfaces,
in Essays on Mirror Manifolds (S.-T. Yau, ed.),
International Press, Hong Kong, 1992, \href{https://arxiv.org/abs/alg-geom/9202026}{alg-geom/9202026}.

\bibitem[Morrison2]{Morrison2}\hypertarget{bib:Morrison2}{}
D. Morrison, Mirror symmetry and rational curves on quintic threefolds:
A guide for mathematicians, Jour. AMS 6 (1993), 223--247,
\href{https://arxiv.org/abs/alg-geom/9202004}{alg-geom/9202004}.

\bibitem[Morrison3]{Morrison3}\hypertarget{bib:Morrison3}{}
D. Morrison, Compactifications of Moduli Spaces Inspired by Mirror Symmetry,
in Journées de Géométrie Algébrique d'Orsay (Orsay, 1992), Astérisque 218,
Soc. Math. France, Paris, 1993, 243--271, \href{https://arxiv.org/abs/alg-geom/9304007}{alg-geom/9304007}.

\bibitem[Morrison4]{Morrison4}\hypertarget{bib:Morrison4}{} D. Morrison, \emph{Where is the large radius limit?}, preprint, 1993, \href{https://arxiv.org/abs/hep-th/9311049}{hep-th/9311049}.

\bibitem[Morrison5]{Morrison5}\hypertarget{bib:Morrison5}{}
D. Morrison, Beyond the Kähler cone, in Proceedings of the Hirzebruch 65 Conference on Algebraic Geometry (Ramat Gan, 1993), Israel Math. Conf. Proc. 9, Bar-Ilan Univ., Ramat Gan, 1996, 361--376, \href{https://arxiv.org/abs/alg-geom/9407007}{alg-geom/9407007}.

\bibitem[Morrison6]{Morrison6}\hypertarget{bib:Morrison6}{}
D. Morrison, Making enumerative predictions by means of mirror symmetry,
in Mirror Symmetry II (B. Greene and S.-T. Yau, eds.), AMS/IP Stud. Adv. Math. 1,
AMS, Providence, RI, 1997, 457--481, \href{https://arxiv.org/abs/alg-geom/9504013}{alg-geom/9504013}.

\bibitem[Morrison7]{Morrison7}\hypertarget{bib:Morrison7}{}
D. Morrison, Mathematical Aspects of Mirror Symmetry, in Complex algebraic geometry (Park City, UT, 1993), 265--327, IAS/Park City Math. Ser. 3, AMS, Providence, RI, 1997, \href{https://arxiv.org/abs/alg-geom/9609021}{alg-geom/9609021}.

\bibitem[Morrison8]{Morrison8}\hypertarget{bib:Morrison8}{} D. Morrison, \emph{The geometry underlying mirror symmetry}, preprint, 1996, \href{https://arxiv.org/abs/alg-geom/9608006}{alg-geom/9608006}.

\bibitem[Morrison9]{Morrison9}\hypertarget{bib:Morrison9}{}
D. Morrison, Through the Looking Glass, \href{https://arxiv.org/abs/alg-geom/9705028}{alg-geom/9705028}.

\bibitem[MP1]{MP1}\hypertarget{bib:MP1}{}
D. Morrison and M. Plesser,
Summing the instantons: quantum cohomology and mirror symmetry in toric varieties,
Nuclear Physics B440 (1995), 279--354, \href{https://arxiv.org/abs/hep-th/9412236}{hep-th/9412236}.

\bibitem[MP2]{MP2}\hypertarget{bib:MP2}{}
D. Morrison and M. Plesser, Towards mirror symmetry as duality for two-dimensional
abelian gauge theories, Nuclear Physics Proc. Suppl. 46 (1996),
177--186, \href{https://arxiv.org/abs/hep-th/9508107}{hep-th/9508107}.

\bibitem[Mumford1]{Mumford1}\hypertarget{bib:Mumford1}{}
D. Mumford, Geometric Invariant Theory, Springer-Verlag, Berlin--Heidelberg--New York, 1965.

\bibitem[Mumford2]{Mumford2}\hypertarget{bib:Mumford2}{}
D. Mumford, Towards an Enumerative Geometry of the Moduli Space of Curves, in Arithmetic and Geometry, Volume II (M. Artin and J. Tate, eds.), Progress in Math. 36, Birkhäuser, Boston--Basel--Berlin, 1983, 271--328.

\SourcePage{471}

\bibitem[Nikulin]{Nikulin}\hypertarget{bib:Nikulin}{}
V. Nikulin, Discrete Reflection groups in Lobachevsky spaces and algebraic surfaces,
in Proceedings of the International Congress of Mathematicians, Berkeley 1986,
AMS, Providence, RI, 1987, 654--671.

\bibitem[Oda]{Oda}\hypertarget{bib:Oda}{}
T. Oda, Convex Bodies and Algebraic Geometry,
Springer-Verlag, Berlin--Heidelberg--New York, 1988.

\bibitem[OdP]{OdP}\hypertarget{bib:OdP}{}
T. Oda and H. Park,
Linear Gale transforms and Gelfand--Kapranov--Zelevinskij decompositions,
Tohoku Math. J. 43 (1991), 375--399.

\bibitem[Oguiso]{Oguiso}\hypertarget{bib:Oguiso}{}
K. Oguiso, On algebraic fiber space structures on a Calabi--Yau threefold, Int. J. Math. 4 (1993), 439--465.

\bibitem[OgP]{OgP}\hypertarget{bib:OgP}{}
K. Oguiso and T. Peternell, Calabi--Yau threefolds with positive second Chern class, Comm. Anal. Geom. 6 (1998), 153--172.

\bibitem[Pandharipande1]{Pandharipande1}\hypertarget{bib:Pandharipande1}{} R. Pandharipande, \emph{Counting elliptic plane curves with fixed $j$-invariant}, Proc. AMS \textbf{125} (1997), 3471--3479, \href{https://arxiv.org/abs/alg-geom/9505023}{alg-geom/9505023}.

\bibitem[Pandharipande2]{Pandharipande2}\hypertarget{bib:Pandharipande2}{} R. Pandharipande, \emph{A geometric construction of Getzler's elliptic relation}, preprint, 1997, \href{https://arxiv.org/abs/alg-geom/9705016}{alg-geom/9705016}.

\bibitem[Pandharipande3]{Pandharipande3}\hypertarget{bib:Pandharipande3}{}
R. Pandharipande, Rational curves on hypersurfaces (after A. Givental), preprint, 1998, to appear in Séminaire Bourbaki, \href{https://arxiv.org/abs/math/9806133}{math.AG/9806133}.

\bibitem[PS]{PS}\hypertarget{bib:PS}{} C. Peters and J. Steenbrink, \emph{Infinitesimal variations of Hodge structure and the generic Torelli problem for projective hypersurfaces (after Carlson, Donagi, Green, Griffiths, Harris)}, in \emph{Classification of Algebraic and Analytic Manifolds} (K. Ueno, ed.), Progress in Math. 39, Birkhäuser, Boston-Basel-Berlin, 1983, 399--463.

\bibitem[Polchinski1]{Polchinski1}\hypertarget{bib:Polchinski1}{} J. Polchinski, \emph{Lectures on D-Branes}, in \emph{Fields, Strings and Duality (TASI 96, Boulder, CO)} (C. Efthimiou and B. Greene, eds.), World Scientific, River Edge, NJ, 1997, 293--356, \href{https://arxiv.org/abs/hep-th/9611050}{hep-th/9611050}.

\bibitem[Polchinski2]{Polchinski2}\hypertarget{bib:Polchinski2}{}
J. Polchinski, String theory, vols. 1 and 2,
Cambridge University Press, Cambridge, 1998.

\bibitem[PZ]{PZ}\hypertarget{bib:PZ}{} A. Polishchuk and E. Zaslow, \emph{Categorical mirror symmetry: the elliptic curve}, Adv. Theor. Math. Phys. \textbf{2} (1998), 443--470, \href{https://arxiv.org/abs/math.AG/9801119}{math.AG/9801119}.

\bibitem[Prill]{Prill}\hypertarget{bib:Prill}{} D. Prill, \emph{Local classification of quotients of complex manifolds by discontinuous groups}, Duke Math. J. \textbf{34} (1967), 375--386.

\bibitem[QR]{QR}\hypertarget{bib:QR}{} Z. Qin and Y. Ruan, \emph{Quantum cohomology of projective bundles over $\mathbb P^n$}, Trans. AMS \textbf{350} (1998), 3615--3638, \href{https://arxiv.org/abs/alg-geom/9510010}{alg-geom/9510010}.

\bibitem[Quinn]{Quinn}\hypertarget{bib:Quinn}{} F. Quinn, \emph{Lectures on axiomatic topological quantum field theory}, in \emph{Geometry and Quantum Field Theory} (D. Freed and K. Uhlenbeck, eds.), IAS/Park City Mathematics Series 1, AMS, 1995, 325--459.

\bibitem[Rabin]{Rabin}\hypertarget{bib:Rabin}{} J. Rabin, \emph{Introduction to quantum field theory for mathematicians}, in \emph{Geometry and Quantum Field Theory} (D. Freed and K. Uhlenbeck, eds.), IAS/Park City Mathematics Series 1, AMS, 1995, 185--269.

\bibitem[Ran1]{Ran1}\hypertarget{bib:Ran1}{} Z. Ran, \emph{Bend, break and count}, preprint, 1997, \href{https://arxiv.org/abs/alg-geom/9704004}{alg-geom/9704004}.

\bibitem[Ran2]{Ran2}\hypertarget{bib:Ran2}{}
Z. Ran, Unobstructedness of Calabi--Yau orbi-Kleinfolds, J. Math. Phys. 39 (1998), 625--629.

\bibitem[Ran3]{Ran3}\hypertarget{bib:Ran3}{} Z. Ran, \emph{Bend, break and count II: elliptics, cuspidals, linear genera}, preprint, 1997, \href{https://arxiv.org/abs/alg-geom/9708013}{alg-geom/9708013}.

\bibitem[Reid1]{Reid1}\hypertarget{bib:Reid1}{}
M. Reid, Canonical 3-folds,
in Journées de géométrie algébrique d'Angers (A. Beauville, ed.),
Sijthoff \& Noordhoff, Alphen aan den Rijn, 1980, 273--310.

\bibitem[Reid2]{Reid2}\hypertarget{bib:Reid2}{}
M. Reid, Minimal models of canonical threefolds, in Algebraic Varieties and Analytic Varieties, Adv. Stud. Pure Math. 1, North-Holland, Amsterdam--New York and Kinokuniya, Tokyo, 1983, 131--180.

\bibitem[Reid3]{Reid3}\hypertarget{bib:Reid3}{}
M. Reid, Decomposition of toric morphisms,
in Arithmetic and Geometry, Volume II (M. Artin and J. Tate, eds.),
Progress in Math. 36, Birkhäuser, Boston--Basel--Berlin, 1983, 395--418.

\bibitem[Reid4]{Reid4}\hypertarget{bib:Reid4}{}
M. Reid, Young person's guide to canonical singularities,
in Algebraic Geometry, Bowdoin 1985, Part 1 (S. Bloch, ed.),
Proc. Sym. Pure Math. 46, AMS, 1987, 345--414.

\SourcePage{472}

\bibitem[Reid5]{Reid5}\hypertarget{bib:Reid5}{}
M. Reid, The moduli space of 3-folds with $K=0$ may nevertheless be irreducible, Math. Ann. 278 (1987), 329--334.

\bibitem[Reid6]{Reid6}\hypertarget{bib:Reid6}{} M. Reid, \emph{McKay correspondence}, preprint, 1997, \href{https://arxiv.org/abs/alg-geom/9702016}{alg-geom/9702016}.

\bibitem[Roan1]{Roan1}\hypertarget{bib:Roan1}{}
S.-S. Roan, The mirror of Calabi--Yau orbifolds,
Int. J. Math. 2 (1991), 439--455.

\bibitem[Roan2]{Roan2}\hypertarget{bib:Roan2}{} S.-S. Roan, \emph{Orbifold Euler characteristic}, in \emph{Mirror Symmetry II} (B. Greene and S.-T. Yau, eds.), AMS/IP Stud. Adv. Math. 1, AMS, Providence, RI, 1997, 129--140.

\bibitem[Ruan1]{Ruan1}\hypertarget{bib:Ruan1}{} Y. Ruan, \emph{Symplectic topology and extremal rays}, Geom. Funct. Anal. \textbf{3} (1993), 395--430.

\bibitem[Ruan2]{Ruan2}\hypertarget{bib:Ruan2}{}
Y. Ruan, Topological sigma model and Donaldson type invariants in Gromov theory, Duke Math. J. 83 (1996), 461--500.

\bibitem[Ruan3]{Ruan3}\hypertarget{bib:Ruan3}{}
Y. Ruan, Virtual neighborhoods and pseudo-holomorphic curves, preprint, 1996, \href{https://arxiv.org/abs/alg-geom/9611021}{alg-geom/9611021}.

\bibitem[RT1]{RT1}\hypertarget{bib:RT1}{}
Y. Ruan and G. Tian, A mathematical theory of quantum cohomology, J. Diff. Geom. 42 (1995), 259--367.

\bibitem[RT2]{RT2}\hypertarget{bib:RT2}{}
Y. Ruan and G. Tian, Higher genus symplectic invariants and sigma model coupled with gravity, Inv. Math. 130 (1997), 455--516, \href{https://arxiv.org/abs/alg-geom/9601005}{alg-geom/9601005}.

\bibitem[Saito]{Saito}\hypertarget{bib:Saito}{} M. Saito, \emph{Mixed Hodge modules}, Publ. Res. Inst. Math. Sci. (Kyoto Univ.) \textbf{26} (1990), 221--333.

\bibitem[Satake]{Satake}\hypertarget{bib:Satake}{} I. Satake, \emph{On a generalization of the notion of a manifold}, Proc. Nat. Acad. Sci. U.S.A. \textbf{42} (1956), 359--363.

\bibitem[Schimmrigk]{Schimmrigk}\hypertarget{bib:Schimmrigk}{}
R. Schimmrigk, Critical superstring vacua from noncritical manifolds: a novel framework
for string compactifications and mirror symmetry,
Phys. Review Letters 70 (1993), 3688--3691, \href{https://arxiv.org/abs/hep-th/9210062}{hep-th/9210062}.

\bibitem[Schlessinger]{Schlessinger}\hypertarget{bib:Schlessinger}{} M. Schlessinger, \emph{Rigidity of quotient singularities}, Invent. Math. \textbf{14} (1971), 17--26.

\bibitem[Schmid]{Schmid}\hypertarget{bib:Schmid}{}
W. Schmid, Variation of Hodge structure: the singularities of the period mapping,
Inv. Math. 22 (1973), 211--319.

\bibitem[Siebert1]{Siebert1}\hypertarget{bib:Siebert1}{}
B. Siebert, Gromov-Witten invariants for general symplectic manifolds, preprint, 1996, \href{https://arxiv.org/abs/dg-ga/9608005}{dg-ga/9608005}.

\bibitem[Siebert2]{Siebert2}\hypertarget{bib:Siebert2}{}
B. Siebert, Algebraic and symplectic Gromov-Witten invariants coincide, preprint, 1998, \href{https://arxiv.org/abs/math.AG/9804108}{math.AG/9804108}.

\bibitem[ST]{ST}\hypertarget{bib:ST}{} B. Siebert and G. Tian, \emph{On quantum cohomology rings of Fano manifolds and a formula of Vafa and Intriligator}, Asian J. Math. \textbf{1} (1997), 679--695, \href{https://arxiv.org/abs/alg-geom/9403010}{alg-geom/9403010}.

\bibitem[Skarke]{Skarke}\hypertarget{bib:Skarke}{}
H. Skarke, Weight systems for toric Calabi--Yau varieties and reflexivity of Newton
polyhedra, Mod. Phys. Lett. A11 (1996), 1637--1652, \href{https://arxiv.org/abs/alg-geom/9603007}{alg-geom/9603007}.

\bibitem[Steenbrink1]{Steenbrink1}\hypertarget{bib:Steenbrink1}{} J. Steenbrink, \emph{Intersection form for quasi-homogeneous singularities}, Compositio Math. \textbf{34} (1977), 211--223.

\bibitem[Steenbrink2]{Steenbrink2}\hypertarget{bib:Steenbrink2}{} J. Steenbrink, \emph{Mixed Hodge structure on the vanishing cohomology}, in \emph{Real and Complex Singularities: Proceedings Nordic Summer School, Oslo, 1976}, Sijthoff \& Noordhoff, Alphen aan den Rijn, 1977, 525--563.

\bibitem[SYZ]{SYZ}\hypertarget{bib:SYZ}{}
A. Strominger, S.-T. Yau and E. Zaslow,
Mirror symmetry is T-duality,
Nuclear Physics B479 (1996), 243--259, \href{https://arxiv.org/abs/hep-th/9606040}{hep-th/9606040}.

\bibitem[Sturmfels1]{Sturmfels1}\hypertarget{bib:Sturmfels1}{}
B. Sturmfels, Gröbner bases of toric varieties, Tôhoku Math. J. 43 (1991), 249--261.

\bibitem[Sturmfels2]{Sturmfels2}\hypertarget{bib:Sturmfels2}{} B. Sturmfels, \emph{Gröbner Bases and Convex Polytopes}, AMS, Providence, 1996.

\bibitem[Szendroi]{Szendroi}\hypertarget{bib:Szendroi}{}
B. Szendrői, On the period map of Calabi--Yau Manifolds, preprint, 1997, \href{https://arxiv.org/abs/alg-geom/9708011}{alg-geom/9708011}.

\bibitem[Szendroi2]{Szendroi2}\hypertarget{bib:Szendroi2}{} B. Szendr\H{o}i, \emph{On a conjecture of Cox and Katz}, Math.\ Z.\ \textbf{240} (2002), 233--241, \href{https://arxiv.org/abs/math.AG/0110166}{math.AG/0110166}.

\bibitem[Taubes]{Taubes}\hypertarget{bib:Taubes}{} C. Taubes, \emph{Counting pseudo-holomorphic submanifolds in dimension 4}, J. Diff. Geom. \textbf{44} (1996), 818--893.

\bibitem[Tian]{Tian}\hypertarget{bib:Tian}{}
G. Tian, Smoothness of the universal deformation space of compact Calabi--Yau manifolds
and its Peterson-Weil metric, in Mathematical Aspects of String Theory
(S.-T. Yau, ed.), World Scientific, Singapore, 1987, 629--646.

\SourcePage{473}

\bibitem[TX]{TX}\hypertarget{bib:TX}{} G. Tian and G. Xu, \emph{On the semi-simplicity of the quantum cohomology algebras of complete intersections}, Math. Res. Lett. \textbf{4} (1997), 481--488, \href{https://arxiv.org/abs/alg-geom/9611035}{alg-geom/9611035}.

\bibitem[Tjøtta]{Tjotta}\hypertarget{bib:Tjotta}{}
E. Tjøtta, \emph{Rational curves on the space of determinantal nets of conics}, preprint, 1998, \href{https://arxiv.org/abs/math/9802037}{math.AG/9802037}.

\bibitem[Todorov]{Todorov}\hypertarget{bib:Todorov}{}
A. Todorov, The Weil-Petersson geometry of the moduli space of $SU(n\geq3)$
(Calabi--Yau) manifolds. I, Comm. Math. Phys. 126 (1989), 325--346.

\bibitem[Vafa1]{Vafa1}\hypertarget{bib:Vafa1}{} C. Vafa, \emph{String vacua and orbifoldized LG models}, Modern Phys. Lett. \textbf{A4} (1989) 1169--1185.

\bibitem[Vafa2]{Vafa2}\hypertarget{bib:Vafa2}{} C. Vafa, \emph{Topological Landau-Ginzburg models}, Modern Phys. Lett. \textbf{A6} (1991), 337--346.

\bibitem[Vafa3]{Vafa3}\hypertarget{bib:Vafa3}{} C. Vafa, \emph{Evidence for F-theory}, Nuclear Physics \textbf{B469} (1996), 403--418, \href{https://arxiv.org/abs/hep-th/9602022}{hep-th/9602022}.

\bibitem[Vainsencher]{Vainsencher}\hypertarget{bib:Vainsencher}{}
I. Vainsencher,
Enumeration of \(n\)-fold tangent hyperplanes to a surface,
J. Algebraic Geom. 4 (1995), 503--526, \href{https://arxiv.org/abs/alg-geom/9312012}{alg-geom/9312012}.

\bibitem[Vakil1]{Vakil1}\hypertarget{bib:Vakil1}{} R. Vakil, \emph{Counting curves of any genus on rational ruled surfaces}, preprint, 1997, \href{https://arxiv.org/abs/alg-geom/9709003}{alg-geom/9709003}.

\bibitem[Vakil2]{Vakil2}\hypertarget{bib:Vakil2}{} R. Vakil, \emph{Genus $g$ Gromov-Witten invariants of Del Pezzo surfaces: Counting plane curves with fixed multiple points}, preprint, 1997, \href{https://arxiv.org/abs/alg-geom/9709004}{math.AG/9709004}.

\bibitem[Vakil3]{Vakil3}\hypertarget{bib:Vakil3}{} R. Vakil, \emph{The enumerative geometry of rational and elliptic curves in projective space}, preprint, 1997, \href{https://arxiv.org/abs/alg-geom/9709007}{alg-geom/9709007}.

\bibitem[Verbitsky]{Verbitsky}\hypertarget{bib:Verbitsky}{} M. Verbitsky, \emph{Mirror symmetry for hyperkaehler manifolds}, preprint, 1995, \href{https://arxiv.org/abs/hep-th/9512195}{hep-th/9512195}.

\bibitem[Viehweg]{Viehweg}\hypertarget{bib:Viehweg}{}
E. Viehweg, Weak positivity and the stability of certain Hilbert points, III, Inv. Math. 101 (1990), 521--543.

\bibitem[Vistoli]{Vistoli}\hypertarget{bib:Vistoli}{}
A. Vistoli, Intersection theory on algebraic stacks and their moduli spaces, Invent. Math. 97, 613--670, 1989.

\bibitem[Voisin1]{Voisin1}\hypertarget{bib:Voisin1}{}
C. Voisin, Miroirs et involutions sur les surfaces K3,
in Journées de Géométrie Algébrique d'Orsay (Orsay, 1992), Astérisque 218,
Soc. Math. France, Paris, 1993, 273--323.

\bibitem[Voisin2]{Voisin2}\hypertarget{bib:Voisin2}{} C. Voisin, \emph{A mathematical proof of a formula of Aspinwall and Morrison}, Compositio Math. \textbf{104} (1996), 135--151.

\bibitem[Voisin3]{Voisin3}\hypertarget{bib:Voisin3}{}
C. Voisin, Symétrie Miroir, Panoramas et Synthèses 2,
Soc. Math. France, Paris, 1996.

\bibitem[WB]{WB}\hypertarget{bib:WB}{} J. Wess and J. Bagger, \emph{Supersymmetry and Supergravity}, Princeton U. Press, Princeton, NJ, 1992.

\bibitem[WW]{WW}\hypertarget{bib:WW}{}
E. Whittaker and G. Watson, A Course of Modern Analysis, Fourth Edition, Cambridge U. Press, Cambridge, 1963.

\bibitem[Wilson1]{Wilson1}\hypertarget{bib:Wilson1}{}
P. Wilson, Calabi--Yau manifolds with large Picard Number, Inv. Math. 98 (1989), 139--155.

\bibitem[Wilson2]{Wilson2}\hypertarget{bib:Wilson2}{}
P. Wilson, The Kähler cone on Calabi--Yau threefolds,
Inv. Math. 107 (1992), 561--583.

\bibitem[Wilson3]{Wilson3}\hypertarget{bib:Wilson3}{}
P. Wilson, Symplectic Deformations of Calabi--Yau threefolds, J. Diff. Geom. 45 (1997), 611--637, \href{https://arxiv.org/abs/alg-geom/9707007}{alg-geom/9707007}.

\bibitem[Wilson4]{Wilson4}\hypertarget{bib:Wilson4}{}
P. Wilson, Flops, Type III contractions and Gromov--Witten invariants on Calabi--Yau threefolds, preprint, 1997, \href{https://arxiv.org/abs/alg-geom/9707008}{alg-geom/9707008}.

\bibitem[Witten1]{Witten1}\hypertarget{bib:Witten1}{} E. Witten, \emph{Topological Sigma Models}, Comm. Math. Phys. \textbf{118} (1988), 411--449.

\bibitem[Witten2]{Witten2}\hypertarget{bib:Witten2}{} E. Witten, \emph{On the structure of the topological phase of two-dimensional gravity}, Nuclear Physics \textbf{B340} (1990), 281--332.

\bibitem[Witten3]{Witten3}\hypertarget{bib:Witten3}{}
E. Witten, Two dimensional gravity and intersection theory on moduli space, in Surveys in Differential Geometry, Proceedings of the Conference on Geometry and Topology held at Harvard University, Cambridge, MA, April 27--29, 1990, Supplement to the Journal of Differential Geometry, no. 1, Lehigh University, Bethlehem, PA; distributed by the AMS, Providence, RI, 1991, 243--310.

\SourcePage{474}

\bibitem[Witten4]{Witten4}\hypertarget{bib:Witten4}{} E. Witten, \emph{Mirror Manifolds and Topological Field Theory}, in \emph{Essays on Mirror Manifolds} (S.-T. Yau, ed.), International Press, Hong Kong 1992, 120--159.

\bibitem[Witten5]{Witten5}\hypertarget{bib:Witten5}{}
E. Witten, Phases of \(N=2\) theories in two dimensions,
Nuclear Physics B403 (1993), 159--222;
reprinted in Mirror Symmetry II (B. Greene and S.-T. Yau, eds.),
AMS/IP Stud. Adv. Math. 1, AMS, Providence, RI, 1997, 143--211,
\href{https://arxiv.org/abs/hep-th/9301042}{hep-th/9301042}.

\bibitem[Witten6]{Witten6}\hypertarget{bib:Witten6}{} E. Witten, \emph{The Verlinde algebra and the cohomology of the Grassmannian}, in \emph{Geometry, Topology, \& Physics}, Conf. Proc. Lecture Notes Geom. Topology, VI, Internat. Press, Cambridge, MA, 1995, 357--422, \href{https://arxiv.org/abs/hep-th/9312104}{hep-th/9312104}.

\bibitem[YZ]{YZ}\hypertarget{bib:YZ}{} S.-T. Yau and E. Zaslow, \emph{BPS states, string duality, and nodal curves on K3}, Nuclear Physics \textbf{B471} (1996), 503--512, \href{https://arxiv.org/abs/hep-th/9512121}{hep-th/9512121}.
\end{thebibliography}

% 由逐条人工译稿生成；保留原索引全部条目与层级，页码指向译本。
\clearpage
\begingroup
\makeatletter
\@nameuse{imki@set@main}
\def\indexname{索引}
\def\imki@maybeaddtotoc{\phantomsection\addcontentsline{toc}{chapter}{索引}}
\makeatother
\setlength{\multicolsep}{8pt}
\begin{theindex}
\raggedcolumns
\renewcommand{\indexspace}{\par\vskip9.5pt plus5pt minus3pt\relax}
\item \leavevmode\hypertarget{source:476}{}\hypertarget{idx:0038}{$\mathcal A$ 超几何方程}，\textit{参见} \hyperlink{idx:0055}{$\mathcal A$ 方程组}
\indexspace
\item \leavevmode\hypertarget{idx:0055}{$\mathcal A$ 方程组},\enspace\IndexPageList[0055]{113}，\textit{另见} \hyperlink{idx:0449}{GKZ 方程组}；\hyperlink{idx:0615}{超几何方程}
\subitem \leavevmode\hypertarget{idx:0056}{定义},\enspace\IndexPageList[0056]{114}
\subitem \leavevmode\hypertarget{idx:0057}{例题},\enspace\IndexPageList[0057]{116,117,118,119,148,185}
\subitem \leavevmode\hypertarget{idx:0058}{在最大幂幺边界点处的全纯解},\enspace\IndexPageList[0058]{183}
\subitem \leavevmode\hypertarget{idx:0059}{Calabi--Yau 环面超曲面的方程组}，\textit{参见} \hyperlink{idx:0099}{Calabi--Yau，Calabi--Yau 环面超曲面}
\subitem \leavevmode\hypertarget{idx:0060}{光滑 Fano 环面簇的方程组},\enspace\IndexPageList[0060]{413,414}
\subitem \leavevmode\hypertarget{idx:0061}{解},\enspace\IndexPageList[0061]{121,123,187}
\indexspace
\item \leavevmode\hypertarget{idx:0002}{A 模型},\enspace\IndexPageList[0002]{30,445,446}
\subitem \leavevmode\hypertarget{idx:0003}{联络},\enspace\IndexPageList[0003]{105,134,281,282,288,289,290,294}
\subsubitem \leavevmode\hypertarget{idx:0004}{渐近性态},\enspace\IndexPageList[0004]{268}
\subsubitem \leavevmode\hypertarget{idx:0005}{典范延拓},\enspace\IndexPageList[0005]{272,274,277,339}
\subsubitem \leavevmode\hypertarget{idx:0006}{收敛性假设},\enspace\IndexPageList[0006]{268,272,284,339}
\subsubitem \leavevmode\hypertarget{idx:0007}{定义},\enspace\IndexPageList[0007]{267}
\subsubitem \leavevmode\hypertarget{idx:0008}{平坦截面},\enspace\IndexPageList[0008]{271,272,338,339,340,341}
\subsubitem \leavevmode\hypertarget{idx:0009}{平坦性},\enspace\IndexPageList[0009]{267}
\subsubitem \leavevmode\hypertarget{idx:0010}{具有最大幂幺单值性},\enspace\IndexPageList[0010]{273}
\subsubitem \leavevmode\hypertarget{idx:0011}{单值性},\enspace\IndexPageList[0011]{269,271,272,339}
\subsubitem \leavevmode\hypertarget{idx:0012}{单值性权滤过},\enspace\IndexPageList[0012]{272,273,276}，\textit{另见} \hyperlink{idx:0907}{单值性权滤过}
\subsubitem \leavevmode\hypertarget{idx:0013}{Calabi--Yau 三维簇的联络},\enspace\IndexPageList[0013]{271,278,279,280,341,351,352}
\subsubitem \leavevmode\hypertarget{idx:0014}{对应的 Picard--Fuchs 算子}，\textit{参见} \hyperlink{idx:1062}{Picard--Fuchs 算子}
\subsubitem \leavevmode\hypertarget{idx:0015}{基 $T_0,\ldots,T_m$ 的作用},\enspace\IndexPageList[0015]{274}
\subsubitem \leavevmode\hypertarget{idx:0016}{满足整性猜想},\enspace\IndexPageList[0016]{273,281}
\subitem \leavevmode\hypertarget{idx:0017}{关联函数}，\textit{参见} \hyperlink{idx:0184}{关联函数，A 模型关联函数}
\subitem \leavevmode\hypertarget{idx:0018}{是 $(1+1)$ 维拓扑量子场论},\enspace\IndexPageList[0018]{457}
\indexspace
\item \leavevmode\hypertarget{idx:0037}{$\mathcal A$ 判别式},\enspace\IndexPageList[0037]{145,148}，\textit{另见} \hyperlink{idx:0214}{判别式轨迹}；\hyperlink{idx:1093}{主 $\mathcal A$ 行列式}
\indexspace
\item \leavevmode\hypertarget{idx:0019}{A 型 Hodge 结构变分},\enspace\IndexPageList[0019]{48,130,245,276,277}，\textit{另见} \hyperlink{idx:0003}{A 模型，联络}
\subitem \leavevmode\hypertarget{idx:0020}{定义},\enspace\IndexPageList[0020]{274,278}
\subitem \leavevmode\hypertarget{idx:0021}{Griffiths 横截性},\enspace\IndexPageList[0021]{275,277}
\subitem \leavevmode\hypertarget{idx:0022}{Hodge 滤过},\enspace\IndexPageList[0022]{274}
\subitem \leavevmode\hypertarget{idx:0023}{整数格},\enspace\IndexPageList[0023]{272,274,277,283}
\subitem \leavevmode\hypertarget{idx:0024}{中间 A 型 Hodge 结构变分},\enspace\IndexPageList[0024]{281}
\subsubitem \leavevmode\hypertarget{idx:0025}{定义},\enspace\IndexPageList[0025]{278}
\subsubitem \leavevmode\hypertarget{idx:0026}{具有最大幂幺单值性},\enspace\IndexPageList[0026]{278}
\subsubitem \leavevmode\hypertarget{idx:0027}{混合 Hodge 结构为 Hodge--Tate 型},\enspace\IndexPageList[0027]{278,287}
\subsubitem \leavevmode\hypertarget{idx:0028}{实结构},\enspace\IndexPageList[0028]{278,287}
\subsubitem \leavevmode\hypertarget{idx:0029}{环面情形},\enspace\IndexPageList[0029]{294}
\subitem \leavevmode\hypertarget{idx:0030}{混合 Hodge 结构},\enspace\IndexPageList[0030]{276,278}，\textit{另见} \hyperlink{idx:0012}{A 模型，联络，单值性权滤过}
\subitem \leavevmode\hypertarget{idx:0031}{极化},\enspace\IndexPageList[0031]{275,277,281}
\indexspace
\item \leavevmode\hypertarget{idx:0001}{$A_\infty$ 范畴},\enspace\IndexPageList[0001]{426}
\indexspace
\item \leavevmode\hypertarget{idx:0032}{Abel--Jacobi 映射},\enspace\IndexPageList[0032]{90,450}
\indexspace
\item \leavevmode\hypertarget{idx:0054}{Arnold 奇异对偶性},\enspace\IndexPageList[0054]{424}
\indexspace
\item \leavevmode\hypertarget{idx:0062}{B 模型},\enspace\IndexPageList[0062]{31,445,446}
\subitem \leavevmode\hypertarget{idx:0063}{关联函数}，\textit{参见} \hyperlink{idx:0189}{关联函数，B 模型关联函数}
\subitem \leavevmode\hypertarget{idx:0064}{势函数},\enspace\IndexPageList[0064]{131,132,284,286}
\subsubitem \leavevmode\hypertarget{idx:0065}{用周期表示},\enspace\IndexPageList[0065]{412,422}
\subsubitem \leavevmode\hypertarget{idx:0066}{五次镜像的势函数}，\textit{参见} \hyperlink{idx:1168}{五次镜像，B 模型势函数}
\subitem \leavevmode\hypertarget{idx:0067}{环面情形},\enspace\IndexPageList[0067]{292,411}
\indexspace
\item \leavevmode\hypertarget{source:477}{}\hypertarget{idx:0068}{Batyrev 镜像}，\textit{参见} \hyperlink{idx:0799}{镜像构造，Batyrev 构造}
\indexspace
\item \leavevmode\hypertarget{idx:0069}{Batyrev 量子环},\enspace\IndexPageList[0069]{246,248,412}
\subitem \leavevmode\hypertarget{idx:0070}{定义},\enspace\IndexPageList[0070]{247}
\subitem \leavevmode\hypertarget{idx:0071}{$-K_X$ 为 nef 的情形},\enspace\IndexPageList[0071]{416}
\subitem \leavevmode\hypertarget{idx:0072}{$\mathbb F_2$ 的情形},\enspace\IndexPageList[0072]{416,417,418}
\subitem \leavevmode\hypertarget{idx:0073}{$\mathbb P^n$ 的情形},\enspace\IndexPageList[0073]{247}
\subitem \leavevmode\hypertarget{idx:0074}{Fano 环面簇的情形},\enspace\IndexPageList[0074]{415}
\indexspace
\item \leavevmode\hypertarget{idx:0085}{BRST 上同调},\enspace\IndexPageList[0085]{325,445,446}
\indexspace
\item \leavevmode\hypertarget{idx:0086}{BRST 算子},\enspace\IndexPageList[0086]{445,446}
\indexspace
\item \leavevmode\hypertarget{idx:0087}{Calabi--Yau}
\subitem \leavevmode\hypertarget{idx:0088}{自同构群},\enspace\IndexPageList[0088]{35,150,152,153,154,161}
\subitem \leavevmode\hypertarget{idx:0089}{由 K3 曲面构造},\enspace\IndexPageList[0089]{89}
\subitem \leavevmode\hypertarget{idx:0090}{广义 Calabi--Yau},\enspace\IndexPageList[0090]{88}
\subitem \leavevmode\hypertarget{idx:0091}{Calabi--Yau 流形},\enspace\IndexPageList[0091]{24,33,208,241,245,266,269,287,288,289,290,338,340,382,404,444,450}
\subsubitem \leavevmode\hypertarget{idx:0092}{具有 K3 或椭圆纤维化},\enspace\IndexPageList[0092]{427}
\subitem \leavevmode\hypertarget{idx:0093}{极小 Calabi--Yau},\enspace\IndexPageList[0093]{290}
\subitem \leavevmode\hypertarget{idx:0094}{模}，\textit{参见} \hyperlink{idx:0870}{模}
\subitem \leavevmode\hypertarget{idx:0095}{Calabi--Yau 轨形},\enspace\IndexPageList[0095]{35,40,78,431,433}
\subitem \leavevmode\hypertarget{idx:0096}{刚性 Calabi--Yau},\enspace\IndexPageList[0096]{88,284,290}
\subitem \leavevmode\hypertarget{idx:0097}{Calabi--Yau 三维簇},\enspace\IndexPageList[0097]{26,27,29,30,32,33,76,97,102,125,126,127,128,129,130,131,132,133,134,135,153,155,156,157,166,174,224,225,226,227,228,229,230,231,232,233,248,249,250,261,270,271,274,278,279,280,281,282,283,284,285,290,305,311,312,313,314,315,316,331,341,348,349,350,351,352,375,389,411,412,428,446,448,450}，\textit{另见} \hyperlink{idx:0184}{关联函数，A 模型关联函数}；\hyperlink{idx:0870}{模}；\hyperlink{idx:1049}{Picard--Fuchs 方程}；\hyperlink{idx:1444}{Yukawa 耦合}
\subitem \leavevmode\hypertarget{idx:0098}{Calabi--Yau 环面完全交}，\textit{参见} \hyperlink{idx:1348}{环面完全交，Calabi--Yau 环面完全交}
\subitem \leavevmode\hypertarget{idx:0099}{Calabi--Yau 环面超曲面},\enspace\IndexPageList[0099]{76,79,83,139,140,141,142,143,144,145,146,147,148,149,150,157,158,159,160,161,162,163,164,165,166,178,179,180,181,182,183,184,185,186,187,188,189,290,293,406,453}
\subsubitem \leavevmode\hypertarget{idx:0100}{$\mathcal A$ 方程组},\enspace\IndexPageList[0100]{114,115,116}
\subsubitem \leavevmode\hypertarget{idx:0101}{复模}，\textit{参见} \hyperlink{idx:0871}{模，复模}
\subsubitem \leavevmode\hypertarget{idx:0102}{Kähler 模}，\textit{参见} \hyperlink{idx:0882}{模，Kähler 模}
\subsubitem \leavevmode\hypertarget{idx:0103}{非退化情形},\enspace\IndexPageList[0103]{145}
\subsubitem \leavevmode\hypertarget{idx:0104}{Picard--Fuchs 方程},\enspace\IndexPageList[0104]{114,115,406}
\subitem \leavevmode\hypertarget{idx:0105}{Calabi--Yau 簇},\enspace\IndexPageList[0105]{78,83,86,430}
\subsubitem \leavevmode\hypertarget{idx:0106}{定义},\enspace\IndexPageList[0106]{34}
\subsubitem \leavevmode\hypertarget{idx:0107}{极小 Calabi--Yau 簇},\enspace\IndexPageList[0107]{34,78,86,139,430,432,433}
\subitem \leavevmode\hypertarget{idx:0108}{加权射影 Calabi--Yau 超曲面},\enspace\IndexPageList[0108]{93,230,428}
\subsubitem \leavevmode\hypertarget{idx:0109}{7555 个超曲面的列表},\enspace\IndexPageList[0109]{83,95,168}
\indexspace
\item \leavevmode\hypertarget{idx:0110}{Calabi--Yau 锥},\enspace\IndexPageList[0110]{162}
\subitem \leavevmode\hypertarget{idx:0111}{例题},\enspace\IndexPageList[0111]{163}
\indexspace
\item \leavevmode\hypertarget{idx:0115}{Cauchy--Riemann 方程},\enspace\IndexPageList[0115]{208,210,214}
\subitem \leavevmode\hypertarget{idx:0116}{非齐次方程},\enspace\IndexPageList[0116]{210}
\subitem \leavevmode\hypertarget{idx:0117}{扰动方程},\enspace\IndexPageList[0117]{210}
\indexspace
\item \leavevmode\hypertarget{idx:0118}{Cayley 技巧},\enspace\IndexPageList[0118]{87}
\indexspace
\item \leavevmode\hypertarget{idx:0122}{Chow 商},\enspace\IndexPageList[0122]{140,142,143,171}
\indexspace
\item \leavevmode\hypertarget{idx:0124}{Clemens 猜想},\enspace\IndexPageList[0124]{38,39,225,229,237,423}
\subitem \leavevmode\hypertarget{idx:0125}{陈述},\enspace\IndexPageList[0125]{225}
\indexspace
\item \leavevmode\hypertarget{source:478}{}\hypertarget{idx:0199}{$\operatorname{cpl}(\Sigma)$},\enspace\IndexPageList[0199]{64,65,66,67,68,158,159,160,162,184}
\subitem \leavevmode\hypertarget{idx:0200}{在 $A^+(\Sigma)$ 中的定义},\enspace\IndexPageList[0200]{66}
\subitem \leavevmode\hypertarget{idx:0201}{在 $A^+_{n-1}(X_\Sigma)\otimes\mathbb R$ 中的定义},\enspace\IndexPageList[0201]{62}
\subitem \leavevmode\hypertarget{idx:0202}{$\Sigma$ 凸性},\enspace\IndexPageList[0202]{66}
\indexspace
\item \leavevmode\hypertarget{idx:0207}{D 膜},\enspace\IndexPageList[0207]{425,427}
\indexspace
\item \leavevmode\hypertarget{idx:0226}{Dubrovin 联络},\enspace\IndexPageList[0226]{262,334,335}
\subitem \leavevmode\hypertarget{idx:0227}{定义},\enspace\IndexPageList[0227]{262}
\subitem \leavevmode\hypertarget{idx:0228}{Euler 向量场},\enspace\IndexPageList[0228]{264}
\subitem \leavevmode\hypertarget{idx:0229}{平坦性},\enspace\IndexPageList[0229]{263,267}
\subitem \leavevmode\hypertarget{idx:0230}{形式联络},\enspace\IndexPageList[0230]{265}
\subitem \leavevmode\hypertarget{idx:0231}{齐次性},\enspace\IndexPageList[0231]{264}
\subitem \leavevmode\hypertarget{idx:0232}{势函数},\enspace\IndexPageList[0232]{262}
\subitem \leavevmode\hypertarget{idx:0233}{挠率},\enspace\IndexPageList[0233]{263}
\subitem \leavevmode\hypertarget{idx:0234}{WDVV 方程},\enspace\IndexPageList[0234]{263,264}
\indexspace
\item \leavevmode\hypertarget{idx:0235}{Dubrovin 形式体系},\enspace\IndexPageList[0235]{257,262,263,264,265,267,327}，\textit{另见} \hyperlink{idx:0126}{上同调场论}
\subitem \leavevmode\hypertarget{idx:0236}{二元乘积},\enspace\IndexPageList[0236]{262}
\subitem \leavevmode\hypertarget{idx:0237}{联络}，\textit{参见} \hyperlink{idx:0226}{Dubrovin 联络}
\indexspace
\item \leavevmode\hypertarget{source:479}{}\hypertarget{idx:0271}{Euler $\Gamma$ 函数},\enspace\IndexPageList[0271]{185,186,188}
\indexspace
\item \leavevmode\hypertarget{idx:0273}{Euler 常数},\enspace\IndexPageList[0273]{185}
\indexspace
\item \leavevmode\hypertarget{idx:0272}{Euler 积分},\enspace\IndexPageList[0272]{183}，\textit{另见} \hyperlink{idx:0449}{GKZ 方程组}
\indexspace
\item \leavevmode\hypertarget{idx:0257}{Euler 类},\enspace\IndexPageList[0257]{196,312,374,381,383}
\subitem \leavevmode\hypertarget{idx:0258}{等变 Euler 类},\enspace\IndexPageList[0258]{299,315,359,361,364,365,367,400}
\subsubitem \leavevmode\hypertarget{idx:0259}{$N_\Gamma$ 的 Euler 类},\enspace\IndexPageList[0259]{308,310,311,312,314,316,393}
\subitem \leavevmode\hypertarget{idx:0260}{零丛的 Euler 类},\enspace\IndexPageList[0260]{381,388,395}
\subitem \leavevmode\hypertarget{idx:0261}{代数叠与粗模空间之间的比较},\enspace\IndexPageList[0261]{204,207,325,355,360}
\indexspace
\item \leavevmode\hypertarget{idx:0256}{Euler 示性数},\enspace\IndexPageList[0256]{333}
\indexspace
\item \leavevmode\hypertarget{idx:0262}{Euler 数据},\enspace\IndexPageList[0262]{50,51,368,371,372,373,403}
\subitem \leavevmode\hypertarget{idx:0263}{定义},\enspace\IndexPageList[0263]{362}
\subitem \leavevmode\hypertarget{idx:0264}{例题}，\textit{参见} \hyperlink{idx:1043}{$\widehat P$}；\hyperlink{idx:1101}{$\widehat Q$}
\subitem \leavevmode\hypertarget{idx:0265}{超几何函数},\enspace\IndexPageList[0265]{369}，\textit{参见} \hyperlink{idx:0584}{$HG$，Euler 数据的 $HG$}
\subitem \leavevmode\hypertarget{idx:0266}{相关联的 Euler 数据},\enspace\IndexPageList[0266]{366,368,376,378}
\subsubitem \leavevmode\hypertarget{idx:0267}{定义},\enspace\IndexPageList[0267]{365}
\subsubitem \leavevmode\hypertarget{idx:0268}{镜像变换},\enspace\IndexPageList[0268]{371,375,376,404}，\textit{另见} \hyperlink{idx:0581}{$HG$，镜像变换}
\subitem \leavevmode\hypertarget{idx:0269}{多重覆盖解释},\enspace\IndexPageList[0269]{366}
\subitem \leavevmode\hypertarget{idx:0270}{唯一性定理},\enspace\IndexPageList[0270]{367,371,376,378,403}
\indexspace
\item \leavevmode\hypertarget{source:487}{}\hypertarget{idx:0966}{$\Omega$--Euler 数据}，\textit{参见} \hyperlink{idx:0262}{Euler 数据}
\indexspace
\item \leavevmode\hypertarget{idx:0274}{Euler--Lagrange 方程},\enspace\IndexPageList[0274]{435,436,437,441,444,445,447}
\subitem \leavevmode\hypertarget{idx:0275}{电磁场的方程},\enspace\IndexPageList[0275]{440}
\indexspace
\item \leavevmode\hypertarget{idx:0295}{F 理论},\enspace\IndexPageList[0295]{427}
\indexspace
\item \leavevmode\hypertarget{idx:0308}{Fano 簇},\enspace\IndexPageList[0308]{69,241,242,246,251,428}
\indexspace
\item \leavevmode\hypertarget{idx:0315}{Feynman 积分},\enspace\IndexPageList[0315]{25,30}，\textit{另见} \hyperlink{idx:1037}{路径积分}
\indexspace
\item \leavevmode\hypertarget{idx:0314}{Feynman 图},\enspace\IndexPageList[0314]{219}
\indexspace
\item \leavevmode\hypertarget{idx:0336}{Fourier 变换},\enspace\IndexPageList[0336]{437,441}
\indexspace
\item \leavevmode\hypertarget{idx:0337}{Frobenius 代数},\enspace\IndexPageList[0337]{244,257,263,453,458}
\subitem \leavevmode\hypertarget{idx:0338}{定义},\enspace\IndexPageList[0338]{455}
\subitem \leavevmode\hypertarget{idx:0339}{给出 $(1+1)$ 维拓扑量子场论},\enspace\IndexPageList[0339]{455,457}
\subitem \leavevmode\hypertarget{idx:0340}{$n$ 点函数},\enspace\IndexPageList[0340]{455,457}
\subitem \leavevmode\hypertarget{source:480}{}\hypertarget{idx:0341}{三点函数},\enspace\IndexPageList[0341]{455}
\subitem \leavevmode\hypertarget{idx:0342}{迹映射},\enspace\IndexPageList[0342]{455}
\indexspace
\item \leavevmode\hypertarget{idx:0344}{Frobenius 方法},\enspace\IndexPageList[0344]{49,100,184,185,186,356}
\subitem \leavevmode\hypertarget{idx:0345}{例题},\enspace\IndexPageList[0345]{185,188}
\subitem \leavevmode\hypertarget{idx:0346}{在 $\mathbb P^4$ 的上同调环中},\enspace\IndexPageList[0346]{187}
\indexspace
\item \leavevmode\hypertarget{idx:0343}{Frobenius 流形},\enspace\IndexPageList[0343]{264}
\indexspace
\item \leavevmode\hypertarget{idx:0348}{Gale 变换},\enspace\IndexPageList[0348]{67}
\indexspace
\item \leavevmode\hypertarget{idx:0365}{Gauss--Manin 联络},\enspace\IndexPageList[0365]{31,48,98,105,174,273,282,288,289,290}
\subitem \leavevmode\hypertarget{idx:0366}{联络矩阵}，\textit{参见} \hyperlink{idx:0174}{联络矩阵}
\subitem \leavevmode\hypertarget{idx:0367}{定义},\enspace\IndexPageList[0367]{97}
\subitem \leavevmode\hypertarget{idx:0368}{平坦整截面},\enspace\IndexPageList[0368]{105,129,134,173,174,410}
\subitem \leavevmode\hypertarget{idx:0369}{Calabi--Yau 三维簇的 Gauss--Manin 联络},\enspace\IndexPageList[0369]{129,131,132,274,280,281}
\subitem \leavevmode\hypertarget{idx:0370}{与归一化 Yukawa 耦合的关系},\enspace\IndexPageList[0370]{128}
\indexspace
\item \leavevmode\hypertarget{idx:0381}{GIT 商},\enspace\IndexPageList[0381]{140,143,145}
\indexspace
\item \leavevmode\hypertarget{idx:0382}{Givental $I$ 函数}
\subitem \leavevmode\hypertarget{idx:0383}{$I_{\mathcal V}$},\enspace\IndexPageList[0383]{49,50,123,381,402,403}
\subsubitem \leavevmode\hypertarget{idx:0384}{定义},\enspace\IndexPageList[0384]{379}
\subsubitem \leavevmode\hypertarget{idx:0385}{例题},\enspace\IndexPageList[0385]{387}
\subsubitem \leavevmode\hypertarget{idx:0386}{凸簇的 $I_{\mathcal V}$},\enspace\IndexPageList[0386]{404}
\subsubitem \leavevmode\hypertarget{idx:0387}{环面簇的 $I_{\mathcal V}$},\enspace\IndexPageList[0387]{404,406}
\subsubitem \leavevmode\hypertarget{idx:0388}{$\mathbb P(1,1,2,2,2)$ 中超曲面的 $I_{\mathcal V}$},\enspace\IndexPageList[0388]{407}
\subsubitem \leavevmode\hypertarget{idx:0389}{五次三维簇的 $I_{\mathcal V}$},\enspace\IndexPageList[0389]{49,121,124,187,384,385}
\subsubitem \leavevmode\hypertarget{idx:0390}{满足镜像的 Picard--Fuchs 方程},\enspace\IndexPageList[0390]{124,144,151,406,408,422}
\subitem \leavevmode\hypertarget{idx:0391}{$\overline I_{\mathcal V}$},\enspace\IndexPageList[0391]{388,389}
\subitem \leavevmode\hypertarget{idx:0392}{$\widetilde I$},\enspace\IndexPageList[0392]{121,122,123,124,151,187,382}
\subsubitem \leavevmode\hypertarget{idx:0393}{满足 $\mathcal A$ 方程组},\enspace\IndexPageList[0393]{121,414}
\subitem \leavevmode\hypertarget{idx:0394}{$I_{\mathbb P^n}$ 的等变版本},\enspace\IndexPageList[0394]{394}
\subitem \leavevmode\hypertarget{idx:0395}{$I_{\mathcal V}$ 的等变版本},\enspace\IndexPageList[0395]{402,405}
\subitem \leavevmode\hypertarget{idx:0396}{$\mathbb F_2$ 的 $I$ 函数},\enspace\IndexPageList[0396]{416,417}
\subitem \leavevmode\hypertarget{idx:0397}{$\mathbb P^1$ 的 $I$ 函数},\enspace\IndexPageList[0397]{380}
\subitem \leavevmode\hypertarget{idx:0398}{$\mathbb P^n$ 的 $I$ 函数},\enspace\IndexPageList[0398]{380,388,394}
\subitem \leavevmode\hypertarget{idx:0399}{环面簇的 $I$ 函数},\enspace\IndexPageList[0399]{412}
\subsubitem \leavevmode\hypertarget{idx:0400}{当 $-K_X$ 为 nef 时},\enspace\IndexPageList[0400]{416}
\indexspace
\item \leavevmode\hypertarget{idx:0401}{Givental $J$ 函数},\enspace\IndexPageList[0401]{48,343,344,345,346,347,348,381,382}
\subitem \leavevmode\hypertarget{idx:0402}{$J_{\mathcal V}$},\enspace\IndexPageList[0402]{49,381,382,402,403}
\subsubitem \leavevmode\hypertarget{idx:0403}{定义},\enspace\IndexPageList[0403]{381}
\subsubitem \leavevmode\hypertarget{idx:0404}{例题},\enspace\IndexPageList[0404]{387}
\subsubitem \leavevmode\hypertarget{idx:0405}{凸簇的 $J_{\mathcal V}$},\enspace\IndexPageList[0405]{388}
\subsubitem \leavevmode\hypertarget{idx:0406}{环面簇的 $J_{\mathcal V}$},\enspace\IndexPageList[0406]{404,405,406}
\subsubitem \leavevmode\hypertarget{idx:0407}{$\mathbb P(1,1,2,2,2)$ 中超曲面的 $J_{\mathcal V}$},\enspace\IndexPageList[0407]{407,411}
\subsubitem \leavevmode\hypertarget{idx:0408}{五次三维簇的 $J_{\mathcal V}$},\enspace\IndexPageList[0408]{49,385}
\subsubitem \leavevmode\hypertarget{idx:0409}{给出修正量子上同调中的关系},\enspace\IndexPageList[0409]{386}
\subsubitem \leavevmode\hypertarget{idx:0410}{给出量子上同调中的关系},\enspace\IndexPageList[0410]{387}
\subsubitem \leavevmode\hypertarget{idx:0411}{符号记法},\enspace\IndexPageList[0411]{381}
\subitem \leavevmode\hypertarget{idx:0412}{定义},\enspace\IndexPageList[0412]{343}
\subitem \leavevmode\hypertarget{idx:0413}{$J_{\mathbb P^n}$ 的等变版本},\enspace\IndexPageList[0413]{389}
\subsubitem \leavevmode\hypertarget{idx:0414}{$S$ 的定义},\enspace\IndexPageList[0414]{389}
\subsubitem \leavevmode\hypertarget{idx:0415}{$Z_i$ 的定义},\enspace\IndexPageList[0415]{390}
\subitem \leavevmode\hypertarget{idx:0416}{$J_{\mathcal V}$ 的等变版本},\enspace\IndexPageList[0416]{395,405}
\subsubitem \leavevmode\hypertarget{idx:0417}{$\widehat S_{\mathcal V}$ 的定义},\enspace\IndexPageList[0417]{395,396}
\subsubitem \leavevmode\hypertarget{idx:0418}{$S_{\mathcal V}$ 的定义},\enspace\IndexPageList[0418]{395,396}
\subsubitem \leavevmode\hypertarget{idx:0419}{$Z_{i,\mathcal V}$ 的定义},\enspace\IndexPageList[0419]{396}
\subitem \leavevmode\hypertarget{idx:0420}{$\mathbb F_2$ 的 $J$ 函数},\enspace\IndexPageList[0420]{416,417}
\subitem \leavevmode\hypertarget{idx:0421}{$\mathbb P^1$ 的 $J$ 函数},\enspace\IndexPageList[0421]{347,380}
\subitem \leavevmode\hypertarget{idx:0422}{$\mathbb P^n$ 的 $J$ 函数},\enspace\IndexPageList[0422]{347,348,369,380,388}
\subitem \leavevmode\hypertarget{idx:0423}{Calabi--Yau 三维簇的 $J$ 函数},\enspace\IndexPageList[0423]{348,349,350,351,352}
\subsubitem \leavevmode\hypertarget{idx:0424}{公式},\enspace\IndexPageList[0424]{348,349}
\subsubitem \leavevmode\hypertarget{idx:0425}{用 Gromov--Witten 势表示},\enspace\IndexPageList[0425]{349}
\subsubitem \leavevmode\hypertarget{idx:0426}{二阶偏导数},\enspace\IndexPageList[0426]{349,386}
\subitem \leavevmode\hypertarget{idx:0427}{环面簇的 $J$ 函数},\enspace\IndexPageList[0427]{412}
\subsubitem \leavevmode\hypertarget{idx:0428}{当 $-K_X$ 为 nef 时},\enspace\IndexPageList[0428]{416}
\subitem \leavevmode\hypertarget{idx:0429}{$\mathbb P(1,1,2,2,2)$ 中超曲面的 $J$ 函数},\enspace\IndexPageList[0429]{409,411}
\subitem \leavevmode\hypertarget{idx:0430}{五次三维簇的 $J$ 函数},\enspace\IndexPageList[0430]{48,385}
\subitem \leavevmode\hypertarget{idx:0431}{公式},\enspace\IndexPageList[0431]{345,379}
\subitem \leavevmode\hypertarget{idx:0432}{给出量子上同调中的关系},\enspace\IndexPageList[0432]{344}
\subitem \leavevmode\hypertarget{idx:0433}{本原部分},\enspace\IndexPageList[0433]{384}
\subitem \leavevmode\hypertarget{idx:0434}{$J_X$ 与 $J_{\mathcal V}$ 的关系},\enspace\IndexPageList[0434]{382,383,405,409,423}
\subitem \leavevmode\hypertarget{idx:0435}{符号记法},\enspace\IndexPageList[0435]{345,379,383}
\subitem \leavevmode\hypertarget{idx:0436}{环面部分},\enspace\IndexPageList[0436]{405}
\indexspace
\item \leavevmode\hypertarget{idx:0437}{Givental 联络},\enspace\IndexPageList[0437]{339}
\subitem \leavevmode\hypertarget{idx:0438}{定义},\enspace\IndexPageList[0438]{334}
\subitem \leavevmode\hypertarget{idx:0439}{对偶联络},\enspace\IndexPageList[0439]{344,352}
\subsubitem \leavevmode\hypertarget{source:481}{}\hypertarget{idx:0440}{平坦截面},\enspace\IndexPageList[0440]{334,335,336,342}
\IndexSubSubSubItem \leavevmode\hypertarget{idx:0441}{限制},\enspace\IndexPageList[0441]{337,345}
\IndexSubSubSubItem \leavevmode\hypertarget{idx:0442}{符号记法},\enspace\IndexPageList[0442]{335}
\subitem \leavevmode\hypertarget{idx:0443}{与 Dubrovin 联络的关系},\enspace\IndexPageList[0443]{268}
\subitem \leavevmode\hypertarget{idx:0444}{在 $H^0(X,\mathbb C)\oplus H^2(X,\mathbb C)$ 上的限制},\enspace\IndexPageList[0444]{337,343}
\subitem \leavevmode\hypertarget{idx:0445}{在 $H^2(V,\mathbb C)$ 上的限制},\enspace\IndexPageList[0445]{338}
\indexspace
\item \leavevmode\hypertarget{idx:0449}{GKZ 方程组},\enspace\IndexPageList[0449]{118,119,123,124,145,146,148,181,182,183,184,185,186,187,188,189,356,406}，\textit{另见} \hyperlink{idx:0055}{$\mathcal A$ 方程组}；\hyperlink{idx:0615}{超几何方程}
\indexspace
\item \leavevmode\hypertarget{idx:0446}{GKZ 分解},\enspace\IndexPageList[0446]{66,67,83,143,144,162,166,178,184,248,453}
\subitem \leavevmode\hypertarget{idx:0447}{扩大 GKZ 分解},\enspace\IndexPageList[0447]{66,453}
\subitem \leavevmode\hypertarget{idx:0448}{例题},\enspace\IndexPageList[0448]{73,74,75,147,150,163}
\indexspace
\item \leavevmode\hypertarget{idx:0450}{GLSM}，\textit{参见} \hyperlink{idx:0354}{规范线性 $\sigma$ 模型}
\indexspace
\item \leavevmode\hypertarget{idx:0452}{Gr\"obner 基方法},\enspace\IndexPageList[0452]{107,109,110,111,113,248}
\indexspace
\item \leavevmode\hypertarget{idx:0453}{Gr\"obner 扇},\enspace\IndexPageList[0453]{144,248}，\textit{另见} \hyperlink{idx:1239}{次级扇}
\indexspace
\item \leavevmode\hypertarget{idx:0454}{Grassmann 簇},\enspace\IndexPageList[0454]{190,196,242,250,251,304,316,428}
\subitem \leavevmode\hypertarget{idx:0455}{其中的 Calabi--Yau 完全交},\enspace\IndexPageList[0455]{428}
\subsubitem \leavevmode\hypertarget{idx:0456}{三维簇},\enspace\IndexPageList[0456]{389}
\indexspace
\item \leavevmode\hypertarget{idx:0483}{Griffiths 横截性},\enspace\IndexPageList[0483]{97,99,125,126,134,275,277}
\subitem \leavevmode\hypertarget{idx:0484}{射影超曲面的 Griffiths 横截性},\enspace\IndexPageList[0484]{106}
\subitem \leavevmode\hypertarget{idx:0485}{$A$ 型 Hodge 结构变分的 Griffiths 横截性}，\textit{参见} \hyperlink{idx:0021}{A 型 Hodge 结构变分，Griffiths 横截性}
\indexspace
\item \leavevmode\hypertarget{idx:0486}{Griffiths--Dwork 方法},\enspace\IndexPageList[0486]{106,107,108,109,110,117,119}，\textit{另见} \hyperlink{idx:1049}{Picard--Fuchs 方程}
\subitem \leavevmode\hypertarget{idx:0487}{方法说明},\enspace\IndexPageList[0487]{107}
\subitem \leavevmode\hypertarget{idx:0488}{例题},\enspace\IndexPageList[0488]{110,111,112,113}
\subitem \leavevmode\hypertarget{idx:0489}{用于射影超曲面},\enspace\IndexPageList[0489]{106,107,108}
\subitem \leavevmode\hypertarget{idx:0490}{用于环面超曲面},\enspace\IndexPageList[0490]{108,109,110}
\subitem \leavevmode\hypertarget{idx:0491}{局限},\enspace\IndexPageList[0491]{109}
\subitem \leavevmode\hypertarget{idx:0492}{用于加权射影超曲面},\enspace\IndexPageList[0492]{108}
\indexspace
\item \leavevmode\hypertarget{idx:0493}{Gromov--Uhlenbeck 紧化},\enspace\IndexPageList[0493]{210,232}
\indexspace
\item \leavevmode\hypertarget{idx:0515}{Gromov--Witten 不变量},\enspace\IndexPageList[0515]{30,32,38,39,47,48,51,288,289,296,304,321,322,323,324,330}
\subitem \leavevmode\hypertarget{idx:0516}{代数定义},\enspace\IndexPageList[0516]{191,192,193,194,195,196,197,198,199,200,201,202,203,204,205,206,207}，\textit{另见} \hyperlink{idx:1406}{虚拟基本类}
\subsubitem \leavevmode\hypertarget{source:482}{}\hypertarget{idx:0517}{一般情形},\enspace\IndexPageList[0517]{206}
\subsubitem \leavevmode\hypertarget{idx:0518}{最简单情形},\enspace\IndexPageList[0518]{194}
\subsubitem \leavevmode\hypertarget{idx:0519}{直观思想},\enspace\IndexPageList[0519]{191}
\subsubitem \leavevmode\hypertarget{idx:0520}{基本类的作用},\enspace\IndexPageList[0520]{195}
\subitem \leavevmode\hypertarget{idx:0521}{在物理中的出现},\enspace\IndexPageList[0521]{445}
\subitem \leavevmode\hypertarget{idx:0522}{公理},\enspace\IndexPageList[0522]{215,216}，\textit{另见} \hyperlink{idx:0495}{Gromov--Witten 类，公理}
\subitem \leavevmode\hypertarget{idx:0523}{计算},\enspace\IndexPageList[0523]{219,220,221,222,223,224,225,226,227,228,229,230,231,232,233,234,235,236,237,238}
\subsubitem \leavevmode\hypertarget{idx:0524}{$\mathbb P^2$ 的椭圆不变量},\enspace\IndexPageList[0524]{233,234,235,333}
\subsubitem \leavevmode\hypertarget{idx:0525}{$\mathbb P^3$ 的椭圆不变量},\enspace\IndexPageList[0525]{236}
\subsubitem \leavevmode\hypertarget{idx:0526}{五次三维簇的椭圆不变量},\enspace\IndexPageList[0526]{236,237}
\subsubitem \leavevmode\hypertarget{idx:0527}{$\mathbb P^2$ 的亏格 $0$ 不变量},\enspace\IndexPageList[0527]{220,221,222,223,224,259,260,333}
\subsubitem \leavevmode\hypertarget{idx:0528}{Calabi--Yau 三维簇的亏格 $0$ 不变量},\enspace\IndexPageList[0528]{224,225,226,227,228,229,230,231,232,233,331,348}，\textit{另见} \hyperlink{idx:0628}{瞬子数}；\hyperlink{idx:0915}{多重覆盖}
\subsubitem \leavevmode\hypertarget{idx:0529}{五次三维簇的亏格 $0$ 不变量},\enspace\IndexPageList[0529]{207,224,225,226}，\textit{另见} \hyperlink{idx:1186}{五次三维簇，Gromov--Witten 不变量}
\subsubitem \leavevmode\hypertarget{idx:0530}{$\mathbb P^2$ 的亏格 $2$ 不变量},\enspace\IndexPageList[0530]{334}
\subsubitem \leavevmode\hypertarget{idx:0531}{其他例题},\enspace\IndexPageList[0531]{224,237,428}
\subsubitem \leavevmode\hypertarget{idx:0532}{$\mathbb P^2$ 的不变量的递推关系},\enspace\IndexPageList[0532]{221,223,234,243,255,260}
\subsubitem \leavevmode\hypertarget{idx:0533}{$\mathbb P^3$ 的不变量的递推关系},\enspace\IndexPageList[0533]{236}
\subitem \leavevmode\hypertarget{idx:0534}{枚举意义},\enspace\IndexPageList[0534]{218,219,221,225,234,235,237,285,423}，\textit{另见} \hyperlink{idx:0632}{瞬子数，枚举意义}
\subitem \leavevmode\hypertarget{idx:0535}{代数定义与辛定义的等价性},\enspace\IndexPageList[0535]{218}
\subitem \leavevmode\hypertarget{idx:0536}{等变不变量},\enspace\IndexPageList[0536]{219,321,322,323,342}
\subsubitem \leavevmode\hypertarget{idx:0537}{定义},\enspace\IndexPageList[0537]{321}
\subsubitem \leavevmode\hypertarget{idx:0538}{非等变极限},\enspace\IndexPageList[0538]{321}
\subsubitem \leavevmode\hypertarget{idx:0539}{五次三维簇的等变不变量},\enspace\IndexPageList[0539]{322}
\subsubitem \leavevmode\hypertarget{idx:0540}{等变丛零截面的等变不变量},\enspace\IndexPageList[0540]{322}
\subitem \leavevmode\hypertarget{idx:0541}{辛定义},\enspace\IndexPageList[0541]{207,211,212,213,214,228,231,232}
\subsubitem \leavevmode\hypertarget{idx:0542}{定义的等价性},\enspace\IndexPageList[0542]{214,218}
\subsubitem \leavevmode\hypertarget{idx:0543}{直观思想},\enspace\IndexPageList[0543]{209}
\subsubitem \leavevmode\hypertarget{idx:0544}{Li--Tian 方法},\enspace\IndexPageList[0544]{213}
\subsubitem \leavevmode\hypertarget{idx:0545}{混合不变量与辛不变量的关系},\enspace\IndexPageList[0545]{213}
\subsubitem \leavevmode\hypertarget{idx:0546}{Ruan--Tian 混合不变量},\enspace\IndexPageList[0546]{211,212,233,235}
\subsubitem \leavevmode\hypertarget{idx:0547}{Ruan--Tian 辛不变量},\enspace\IndexPageList[0547]{212,233,235}
\subitem \leavevmode\hypertarget{idx:0548}{树级不变量},\enspace\IndexPageList[0548]{219,253}
\subitem \leavevmode\hypertarget{idx:0549}{用于定义 $A$ 模型的关联函数},\enspace\IndexPageList[0549]{244}
\subitem \leavevmode\hypertarget{idx:0550}{用于定义 Gromov--Witten 势},\enspace\IndexPageList[0550]{253}
\subitem \leavevmode\hypertarget{idx:0551}{用于定义量子上同调},\enspace\IndexPageList[0551]{241}
\indexspace
\item \leavevmode\hypertarget{idx:0494}{Gromov--Witten 类},\enspace\IndexPageList[0494]{223,235,243,255,328,332}
\subitem \leavevmode\hypertarget{idx:0495}{公理},\enspace\IndexPageList[0495]{214,215,216,217,218,327}
\subsubitem \leavevmode\hypertarget{idx:0496}{复合律},\enspace\IndexPageList[0496]{217}
\subsubitem \leavevmode\hypertarget{idx:0497}{形变公理},\enspace\IndexPageList[0497]{217,418}
\subsubitem \leavevmode\hypertarget{idx:0498}{次数公理},\enspace\IndexPageList[0498]{215,221,224,288}
\subsubitem \leavevmode\hypertarget{idx:0499}{除子公理},\enspace\IndexPageList[0499]{216,221,223,224,225,226,235,258,259,262,269,288,330}
\subsubitem \leavevmode\hypertarget{idx:0500}{有效性公理},\enspace\IndexPageList[0500]{215,251}
\subsubitem \leavevmode\hypertarget{idx:0501}{等变性公理},\enspace\IndexPageList[0501]{215,222,224,242,243,244,256,264}
\subsubitem \leavevmode\hypertarget{idx:0502}{基本类公理},\enspace\IndexPageList[0502]{216,221,222,224,235,242,257,258,288}
\subsubitem \leavevmode\hypertarget{idx:0503}{线性公理},\enspace\IndexPageList[0503]{215}
\subsubitem \leavevmode\hypertarget{idx:0504}{动机公理},\enspace\IndexPageList[0504]{217}
\subsubitem \leavevmode\hypertarget{idx:0505}{点映射公理},\enspace\IndexPageList[0505]{216,220,223,224,227,234,242,254,257,258,262,330,333,337}
\subsubitem \leavevmode\hypertarget{idx:0506}{约化公理},\enspace\IndexPageList[0506]{217,235,333}
\subsubitem \leavevmode\hypertarget{idx:0507}{分裂公理},\enspace\IndexPageList[0507]{217,219,221,222,242,243,256,264}
\subitem \leavevmode\hypertarget{idx:0508}{最简单情形下的定义},\enspace\IndexPageList[0508]{194}
\subitem \leavevmode\hypertarget{idx:0509}{定义},\enspace\IndexPageList[0509]{206}，\textit{另见} \hyperlink{idx:0516}{Gromov--Witten 不变量，代数定义}；\hyperlink{idx:1406}{虚拟基本类}
\subitem \leavevmode\hypertarget{idx:0510}{直观定义},\enspace\IndexPageList[0510]{191}
\subitem \leavevmode\hypertarget{idx:0511}{$\mathbb P^2$ 的 Gromov--Witten 类},\enspace\IndexPageList[0511]{221,222}
\subitem \leavevmode\hypertarget{idx:0512}{基本类的作用},\enspace\IndexPageList[0512]{195}
\subitem \leavevmode\hypertarget{idx:0513}{辛 Gromov--Witten 类},\enspace\IndexPageList[0513]{213}
\subitem \leavevmode\hypertarget{idx:0514}{树级 Gromov--Witten 类},\enspace\IndexPageList[0514]{219}
\indexspace
\item \leavevmode\hypertarget{idx:0552}{Gromov--Witten 势},\enspace\IndexPageList[0552]{51,132,253,254,255,256,257,267,270,326,327,412}
\subitem \leavevmode\hypertarget{idx:0553}{经典部分},\enspace\IndexPageList[0553]{258}
\subitem \leavevmode\hypertarget{idx:0554}{定义},\enspace\IndexPageList[0554]{253}
\subitem \leavevmode\hypertarget{idx:0555}{等变势},\enspace\IndexPageList[0555]{342,343}
\subitem \leavevmode\hypertarget{idx:0556}{亏格 $0$ 自由能},\enspace\IndexPageList[0556]{254}
\subitem \leavevmode\hypertarget{idx:0557}{亏格 $g$ 的势},\enspace\IndexPageList[0557]{254}
\subitem \leavevmode\hypertarget{idx:0558}{齐次性},\enspace\IndexPageList[0558]{258}
\subitem \leavevmode\hypertarget{idx:0559}{$\mathbb P^2$ 的势},\enspace\IndexPageList[0559]{259}
\subitem \leavevmode\hypertarget{idx:0560}{Calabi--Yau 三维簇的势},\enspace\IndexPageList[0560]{261,262,279,284,341,349}
\subsubitem \leavevmode\hypertarget{idx:0561}{可用镜像的周期表示},\enspace\IndexPageList[0561]{356,412}
\subsubitem \leavevmode\hypertarget{idx:0562}{环面部分},\enspace\IndexPageList[0562]{412}
\subitem \leavevmode\hypertarget{idx:0563}{椭圆曲线的势},\enspace\IndexPageList[0563]{254,258}
\subitem \leavevmode\hypertarget{idx:0564}{五次三维簇的势},\enspace\IndexPageList[0564]{51,260,261,356}
\subitem \leavevmode\hypertarget{idx:0565}{量子部分},\enspace\IndexPageList[0565]{258}
\subitem \leavevmode\hypertarget{idx:0566}{结构},\enspace\IndexPageList[0566]{257,258,259}
\subitem \leavevmode\hypertarget{idx:0567}{三阶偏导数},\enspace\IndexPageList[0567]{254,257}
\subitem \leavevmode\hypertarget{idx:0568}{WDVV 方程},\enspace\IndexPageList[0568]{255,256,257,259,260,262}
\indexspace
\item \leavevmode\hypertarget{idx:0570}{Gysin 谱序列},\enspace\IndexPageList[0570]{82}
\indexspace
\item \leavevmode\hypertarget{idx:0569}{Gysin 映射},\enspace\IndexPageList[0569]{194}
\indexspace
\item \leavevmode\hypertarget{idx:0571}{Hamilton 方程},\enspace\IndexPageList[0571]{436,437}
\indexspace
\item \leavevmode\hypertarget{idx:0572}{Hamilton 量},\enspace\IndexPageList[0572]{435,436,437,438,453}
\subitem \leavevmode\hypertarget{idx:0573}{电磁理论的 Hamilton 量},\enspace\IndexPageList[0573]{441}
\indexspace
\item \leavevmode\hypertarget{idx:0576}{Hamilton 向量场},\enspace\IndexPageList[0576]{63}
\indexspace
\item \leavevmode\hypertarget{idx:0574}{Hamilton 作用},\enspace\IndexPageList[0574]{63,452}
\subitem \leavevmode\hypertarget{idx:0575}{矩映射}，\textit{参见} \hyperlink{idx:0898}{矩映射}
\indexspace
\item \leavevmode\hypertarget{idx:0578}{Heisenberg 绘景},\enspace\IndexPageList[0578]{436}
\indexspace
\item \leavevmode\hypertarget{idx:0579}{$HG$}
\subitem \leavevmode\hypertarget{idx:0580}{定义},\enspace\IndexPageList[0580]{369}
\subitem \leavevmode\hypertarget{idx:0581}{镜像变换},\enspace\IndexPageList[0581]{371,372,376,377}，\textit{另见} \hyperlink{idx:0268}{Euler 数据，相关联的 Euler 数据，镜像变换}
\subitem \leavevmode\hypertarget{idx:0582}{$\widehat P$ 的 $HG$},\enspace\IndexPageList[0582]{51,124,187,370,375,376,377,380}
\subitem \leavevmode\hypertarget{idx:0583}{$\widehat Q$ 的 $HG$},\enspace\IndexPageList[0583]{51,370,373,375,376,377,380}
\subitem \leavevmode\hypertarget{idx:0584}{Euler 数据的 $HG$},\enspace\IndexPageList[0584]{369}
\indexspace
\item \leavevmode\hypertarget{idx:0585}{Hilbert 概形},\enspace\IndexPageList[0585]{231,305}
\indexspace
\item \leavevmode\hypertarget{idx:0586}{Hirzebruch--Riemann--Roch 定理},\enspace\IndexPageList[0586]{198}
\indexspace
\item \leavevmode\hypertarget{idx:0587}{Hochschild 上同调},\enspace\IndexPageList[0587]{426}
\indexspace
\item \leavevmode\hypertarget{idx:0593}{Hodge 结构}
\subitem \leavevmode\hypertarget{idx:0594}{退化},\enspace\IndexPageList[0594]{98}
\subitem \leavevmode\hypertarget{idx:0595}{Hodge 滤过},\enspace\IndexPageList[0595]{96,128,176}
\subsubitem \leavevmode\hypertarget{idx:0596}{极限 Hodge 滤过},\enspace\IndexPageList[0596]{99,134}
\subitem \leavevmode\hypertarget{idx:0597}{整数格},\enspace\IndexPageList[0597]{96,99,283,291}
\subitem \leavevmode\hypertarget{idx:0598}{混合 Hodge 结构},\enspace\IndexPageList[0598]{99,276}
\subsubitem \leavevmode\hypertarget{idx:0599}{Hodge--Tate 型},\enspace\IndexPageList[0599]{103,128,278,287,289,291}
\subsubitem \leavevmode\hypertarget{idx:0600}{极限混合 Hodge 结构},\enspace\IndexPageList[0600]{98,128,278,289}，\textit{另见} \hyperlink{idx:0907}{单值性权滤过}
\subitem \leavevmode\hypertarget{idx:0601}{轨形的 Hodge 结构},\enspace\IndexPageList[0601]{432}
\subitem \leavevmode\hypertarget{source:483}{}\hypertarget{idx:0602}{权为 $k$ 的 Hodge 结构},\enspace\IndexPageList[0602]{96}
\subitem \leavevmode\hypertarget{idx:0603}{极化 Hodge 结构},\enspace\IndexPageList[0603]{97,130,276,281,290}
\subitem \leavevmode\hypertarget{idx:0604}{纯 Hodge 结构},\enspace\IndexPageList[0604]{82}
\subitem \leavevmode\hypertarget{idx:0605}{Hodge 结构变分},\enspace\IndexPageList[0605]{48,97,130,132,175,181,276,277,281,290}
\subsubitem \leavevmode\hypertarget{idx:0606}{$\mathbb C$ 变分},\enspace\IndexPageList[0606]{291,294,295}
\subsubitem \leavevmode\hypertarget{idx:0607}{来自几何的变分},\enspace\IndexPageList[0607]{285}
\subsubitem \leavevmode\hypertarget{idx:0608}{形式退化的变分},\enspace\IndexPageList[0608]{284}
\subsubitem \leavevmode\hypertarget{idx:0609}{实变分},\enspace\IndexPageList[0609]{278,287}
\indexspace
\item \leavevmode\hypertarget{idx:0588}{Hodge 菱形图},\enspace\IndexPageList[0588]{27,28}
\indexspace
\item \leavevmode\hypertarget{idx:0589}{Hodge 数},\enspace\IndexPageList[0589]{27,29,89,250}
\subitem \leavevmode\hypertarget{idx:0590}{Calabi--Yau 环面完全交的 Hodge 数},\enspace\IndexPageList[0590]{86}
\subitem \leavevmode\hypertarget{idx:0591}{Calabi--Yau 环面超曲面的 Hodge 数},\enspace\IndexPageList[0591]{79,80,81,82,83,142,168,169}
\subitem \leavevmode\hypertarget{idx:0592}{弦论 Hodge 数},\enspace\IndexPageList[0592]{83,86,296,425}
\indexspace
\item \leavevmode\hypertarget{idx:0610}{Hodge--Riemann 双线性关系},\enspace\IndexPageList[0610]{97,277,291}
\indexspace
\item \leavevmode\hypertarget{idx:0611}{Hodge--Tate 型}，\textit{参见} \hyperlink{idx:0599}{Hodge 结构，混合 Hodge 结构，Hodge--Tate 型}
\indexspace
\item \leavevmode\hypertarget{idx:0614}{Hopf 纤维化},\enspace\IndexPageList[0614]{72}
\indexspace
\item \leavevmode\hypertarget{idx:0641}{$I_{\mathcal V}$}，\textit{参见} \hyperlink{idx:0383}{Givental $I$ 函数，$I_{\mathcal V}$}
\indexspace
\item \leavevmode\hypertarget{idx:0644}{$J$ 全纯曲线},\enspace\IndexPageList[0644]{232}
\subitem \leavevmode\hypertarget{idx:0645}{定义},\enspace\IndexPageList[0645]{231}
\indexspace
\item \leavevmode\hypertarget{idx:0646}{$J$ 全纯映射},\enspace\IndexPageList[0646]{208}
\subitem \leavevmode\hypertarget{idx:0647}{模空间}
\subsubitem \leavevmode\hypertarget{idx:0648}{维数},\enspace\IndexPageList[0648]{209,210}
\subsubitem \leavevmode\hypertarget{idx:0649}{基本类},\enspace\IndexPageList[0649]{211}
\subsubitem \leavevmode\hypertarget{idx:0650}{Li--Tian 方法},\enspace\IndexPageList[0650]{211,213}
\subsubitem \leavevmode\hypertarget{idx:0651}{Ruan 方法},\enspace\IndexPageList[0651]{209}
\subsubitem \leavevmode\hypertarget{idx:0652}{Ruan--Tian 方法},\enspace\IndexPageList[0652]{210}
\subitem \leavevmode\hypertarget{idx:0653}{多重覆盖},\enspace\IndexPageList[0653]{209}
\subitem \leavevmode\hypertarget{idx:0654}{受扰动的 $J$ 全纯映射},\enspace\IndexPageList[0654]{210,213,228}
\subitem \leavevmode\hypertarget{idx:0655}{简单 $J$ 全纯映射},\enspace\IndexPageList[0655]{209,231,233}
\subitem \leavevmode\hypertarget{idx:0656}{稳定 $J$ 全纯映射},\enspace\IndexPageList[0656]{211}
\indexspace
\item \leavevmode\hypertarget{idx:0658}{$J_{\mathcal V}$}，\textit{参见} \hyperlink{idx:0402}{Givental $J$ 函数，$J_{\mathcal V}$}
\indexspace
\item \leavevmode\hypertarget{idx:0643}{Jacobian 环},\enspace\IndexPageList[0643]{450,451}
\indexspace
\item \leavevmode\hypertarget{idx:0642}{Jacobian 理想},\enspace\IndexPageList[0642]{107,109,113}
\indexspace
\item \leavevmode\hypertarget{idx:0657}{Jones 多项式},\enspace\IndexPageList[0657]{458}
\indexspace
\item \leavevmode\hypertarget{idx:0677}{$K_d$},\enspace\IndexPageList[0677]{356,373}，\textit{另见} \hyperlink{idx:0927}{$N_d$，五次三维簇的 Gromov--Witten 不变量}
\subitem \leavevmode\hypertarget{idx:0678}{临界丛的 $K_d$},\enspace\IndexPageList[0678]{375}
\subitem \leavevmode\hypertarget{idx:0679}{多重覆盖情形的 $K_d$},\enspace\IndexPageList[0679]{375,378}
\indexspace
\item \leavevmode\hypertarget{idx:0659}{K3 曲面},\enspace\IndexPageList[0659]{382}
\subitem \leavevmode\hypertarget{idx:0660}{镜像对称},\enspace\IndexPageList[0660]{424}
\subitem \leavevmode\hypertarget{idx:0661}{其上的有理曲线},\enspace\IndexPageList[0661]{423}
\subitem \leavevmode\hypertarget{idx:0662}{其上的向量丛},\enspace\IndexPageList[0662]{427}
\subitem \leavevmode\hypertarget{idx:0663}{带有对合的 K3 曲面},\enspace\IndexPageList[0663]{88,89,90}
\indexspace
\item \leavevmode\hypertarget{idx:0664}{Kähler}
\subitem \leavevmode\hypertarget{idx:0665}{Kähler 锥},\enspace\IndexPageList[0665]{35,143,150,151,153,173,248,250,388}
\subsubitem \leavevmode\hypertarget{idx:0666}{复结构变动时的变化},\enspace\IndexPageList[0666]{157}
\subsubitem \leavevmode\hypertarget{idx:0667}{闭包},\enspace\IndexPageList[0667]{151,155,158,165,176,266,268,290,339,340,404}
\subsubitem \leavevmode\hypertarget{idx:0668}{面},\enspace\IndexPageList[0668]{155,156,157,167,169,181,250}
\subsubitem \leavevmode\hypertarget{idx:0669}{局部多面锥性质},\enspace\IndexPageList[0669]{155}
\subsubitem \leavevmode\hypertarget{idx:0670}{Calabi--Yau 三维簇的 Kähler 锥},\enspace\IndexPageList[0670]{155,167,169}，\textit{另见} \hyperlink{idx:0913}{可动锥}
\subsubitem \leavevmode\hypertarget{idx:0671}{多面锥性质},\enspace\IndexPageList[0671]{62,152}
\subsubitem \leavevmode\hypertarget{idx:0672}{单纯锥性质},\enspace\IndexPageList[0672]{120}
\subsubitem \leavevmode\hypertarget{idx:0673}{环面部分},\enspace\IndexPageList[0673]{158,159,161,162,422}
\subitem \leavevmode\hypertarget{idx:0674}{轨形上的 Kähler 形式},\enspace\IndexPageList[0674]{432}
\subitem \leavevmode\hypertarget{idx:0675}{Kähler 模}，\textit{参见} \hyperlink{idx:0882}{模，Kähler 模}
\subitem \leavevmode\hypertarget{idx:0676}{Kähler 空间}，\textit{参见} \hyperlink{idx:0143}{复化，复化 Kähler 空间}
\indexspace
\item \leavevmode\hypertarget{idx:0680}{Klein--Gordon 方程},\enspace\IndexPageList[0680]{435,437,438,441}
\indexspace
\item \leavevmode\hypertarget{idx:0681}{Kodaira--Spencer 类},\enspace\IndexPageList[0681]{156}
\indexspace
\item \leavevmode\hypertarget{idx:0682}{Kodaira--Spencer 映射},\enspace\IndexPageList[0682]{176}
\indexspace
\item \leavevmode\hypertarget{idx:0683}{Kuranishi 映射},\enspace\IndexPageList[0683]{198,200}
\indexspace
\item \leavevmode\hypertarget{source:484}{}\hypertarget{idx:0705}{$L_d$},\enspace\IndexPageList[0705]{359,398}
\subitem \leavevmode\hypertarget{idx:0706}{不动点轨迹},\enspace\IndexPageList[0706]{399}
\subsubitem \leavevmode\hypertarget{idx:0707}{法丛},\enspace\IndexPageList[0707]{400}
\subitem \leavevmode\hypertarget{idx:0708}{到 $N_d$ 的映射 $\mu$},\enspace\IndexPageList[0708]{398}
\indexspace
\item \leavevmode\hypertarget{idx:0715}{$\mathcal L_i$}，\textit{参见} \hyperlink{idx:0198}{余切线 $\mathcal L_i$}
\indexspace
\item \leavevmode\hypertarget{idx:0686}{Lagrange 范畴},\enspace\IndexPageList[0686]{426}
\indexspace
\item \leavevmode\hypertarget{idx:0685}{Lagrange 量},\enspace\IndexPageList[0685]{24,435,443}
\indexspace
\item \leavevmode\hypertarget{idx:0687}{Lagrange 密度},\enspace\IndexPageList[0687]{434,435,437,438,439}
\subitem \leavevmode\hypertarget{idx:0688}{电磁理论的 Lagrange 密度},\enspace\IndexPageList[0688]{440,441}
\indexspace
\item \leavevmode\hypertarget{idx:0684}{Lagrange 提升},\enspace\IndexPageList[0684]{372}，\textit{另见} \hyperlink{idx:0268}{Euler 数据，相关联的 Euler 数据，镜像变换}
\indexspace
\item \leavevmode\hypertarget{idx:0689}{Lagrange 子流形},\enspace\IndexPageList[0689]{426}
\subitem \leavevmode\hypertarget{idx:0690}{特殊 Lagrange 子流形},\enspace\IndexPageList[0690]{34,425,426}
\indexspace
\item \leavevmode\hypertarget{idx:0691}{Landau--Ginzburg}
\subitem \leavevmode\hypertarget{idx:0692}{模型},\enspace\IndexPageList[0692]{449}
\subitem \leavevmode\hypertarget{idx:0693}{轨形},\enspace\IndexPageList[0693]{29,166,450,452}
\subitem \leavevmode\hypertarget{idx:0694}{势},\enspace\IndexPageList[0694]{451}
\subitem \leavevmode\hypertarget{idx:0695}{理论},\enspace\IndexPageList[0695]{449,450,451,452}
\indexspace
\item \leavevmode\hypertarget{idx:0703}{Laurent 单项式},\enspace\IndexPageList[0703]{58,141,149}
\indexspace
\item \leavevmode\hypertarget{idx:0704}{Laurent 多项式},\enspace\IndexPageList[0704]{58,139,140,145,178}
\indexspace
\item \leavevmode\hypertarget{idx:0710}{Lefschetz 定理},\enspace\IndexPageList[0710]{384}
\indexspace
\item \leavevmode\hypertarget{idx:0709}{Lefschetz 分解},\enspace\IndexPageList[0709]{282,287,384}
\indexspace
\item \leavevmode\hypertarget{idx:0714}{Levi--Civita 联络},\enspace\IndexPageList[0714]{443}
\indexspace
\item \leavevmode\hypertarget{idx:0744}{Lorentz 变换},\enspace\IndexPageList[0744]{438}
\indexspace
\item \leavevmode\hypertarget{idx:0746}{$M(X)$}，\textit{参见} \hyperlink{idx:0911}{Mori 锥}
\indexspace
\item \leavevmode\hypertarget{idx:0725}{$\overline M_{0,n}(\mathbb P^n,d)$ 中的局部化},\enspace\IndexPageList[0725]{305,306,307,308,309,310,311,312,313,314,315,316}
\subitem \leavevmode\hypertarget{idx:0726}{不动点轨迹},\enspace\IndexPageList[0726]{305,363,371,390,391,399,400}
\subsubitem \leavevmode\hypertarget{idx:0727}{描述},\enspace\IndexPageList[0727]{306}
\subsubitem \leavevmode\hypertarget{idx:0728}{例题},\enspace\IndexPageList[0728]{307,308,312}
\subsubitem \leavevmode\hypertarget{idx:0729}{对应的图},\enspace\IndexPageList[0729]{306,312,313,390,391,393}
\subsubitem \leavevmode\hypertarget{idx:0730}{法丛},\enspace\IndexPageList[0730]{308,309,310,312,314,316,393}，\textit{另见} \hyperlink{idx:0259}{Euler 类，等变 Euler 类，$N_\Gamma$ 的 Euler 类}
\subsubitem \leavevmode\hypertarget{idx:0731}{A 型图},\enspace\IndexPageList[0731]{391,392,397}
\subsubitem \leavevmode\hypertarget{idx:0732}{B 型图},\enspace\IndexPageList[0732]{391,397}
\subitem \leavevmode\hypertarget{idx:0733}{其上的环面作用},\enspace\IndexPageList[0733]{305}
\indexspace
\item \leavevmode\hypertarget{idx:0777}{$M_d$},\enspace\IndexPageList[0777]{357}
\subitem \leavevmode\hypertarget{idx:0778}{不动点轨迹},\enspace\IndexPageList[0778]{357,362}
\subsubitem \leavevmode\hypertarget{idx:0779}{描述},\enspace\IndexPageList[0779]{357}
\subsubitem \leavevmode\hypertarget{idx:0780}{图 $\Gamma$},\enspace\IndexPageList[0780]{362}
\subsubitem \leavevmode\hypertarget{idx:0781}{法丛},\enspace\IndexPageList[0781]{364}
\subitem \leavevmode\hypertarget{idx:0782}{到 $N_d$ 的映射 $\varphi$},\enspace\IndexPageList[0782]{359,360,362,363,364,367,398}
\subitem \leavevmode\hypertarget{idx:0783}{其上的环面作用},\enspace\IndexPageList[0783]{357}
\indexspace
\item \leavevmode\hypertarget{idx:0747}{Macaulay},\enspace\IndexPageList[0747]{320}
\indexspace
\item \leavevmode\hypertarget{idx:0775}{Maxwell 方程},\enspace\IndexPageList[0775]{440,441}
\indexspace
\item \leavevmode\hypertarget{idx:0776}{McKay 对应},\enspace\IndexPageList[0776]{425}
\indexspace
\item \leavevmode\hypertarget{idx:0785}{Minkowski 和},\enspace\IndexPageList[0785]{56,86}
\indexspace
\item \leavevmode\hypertarget{source:486}{}\hypertarget{idx:0911}{Mori 锥},\enspace\IndexPageList[0911]{62,121,251}
\subitem \leavevmode\hypertarget{idx:0912}{环面簇的 Mori 锥}，\textit{参见} \hyperlink{idx:1372}{环面簇，Mori 锥}
\indexspace
\item \leavevmode\hypertarget{idx:0925}{$N_d$}
\subitem \leavevmode\hypertarget{idx:0926}{$\mathbb P^2$ 的 Gromov--Witten 不变量},\enspace\IndexPageList[0926]{221}
\subitem \leavevmode\hypertarget{idx:0927}{五次三维簇的 Gromov--Witten 不变量},\enspace\IndexPageList[0927]{224}
\subitem \leavevmode\hypertarget{idx:0928}{射影空间},\enspace\IndexPageList[0928]{357,398}
\subsubitem \leavevmode\hypertarget{idx:0929}{等变上同调},\enspace\IndexPageList[0929]{358}
\subsubitem \leavevmode\hypertarget{idx:0930}{等变超平面类},\enspace\IndexPageList[0930]{358}
\subsubitem \leavevmode\hypertarget{idx:0931}{不动点},\enspace\IndexPageList[0931]{358,363}
\subsubitem \leavevmode\hypertarget{idx:0932}{其上的环面作用},\enspace\IndexPageList[0932]{358}
\subsubitem \leavevmode\hypertarget{idx:0933}{权},\enspace\IndexPageList[0933]{358}
\indexspace
\item \leavevmode\hypertarget{idx:0935}{nef 分拆},\enspace\IndexPageList[0935]{85,86,120,123}
\subitem \leavevmode\hypertarget{idx:0936}{对偶 nef 分拆},\enspace\IndexPageList[0936]{86}
\indexspace
\item \leavevmode\hypertarget{idx:0934}{nef 完全交}，\textit{参见} \hyperlink{idx:1349}{环面完全交，nef 环面完全交}
\indexspace
\item \leavevmode\hypertarget{idx:0941}{Noether 定理},\enspace\IndexPageList[0941]{442}
\indexspace
\item \leavevmode\hypertarget{idx:0955}{Novikov 环},\enspace\IndexPageList[0955]{252}
\indexspace
\item \leavevmode\hypertarget{source:488}{}\hypertarget{idx:1043}{$\widehat P$},\enspace\IndexPageList[1043]{50,51,124,187,366,369,376}
\subitem \leavevmode\hypertarget{idx:1044}{定义},\enspace\IndexPageList[1044]{360}
\subitem \leavevmode\hypertarget{idx:1045}{凸凹丛的 $\widehat P$},\enspace\IndexPageList[1045]{373}
\subitem \leavevmode\hypertarget{idx:1046}{多重覆盖情形的 $\widehat P$},\enspace\IndexPageList[1046]{378}
\subitem \leavevmode\hypertarget{idx:1047}{是 Euler 数据},\enspace\IndexPageList[1047]{362}
\indexspace
\item \leavevmode\hypertarget{idx:0123}{$\mathcal P$ 类}，\textit{参见} \hyperlink{idx:0847}{镜像定理，Givental 镜像定理，$\mathcal P$ 类}
\indexspace
\item \leavevmode\hypertarget{idx:1001}{$\mathbb P(1,1,2,2,2)$ 及其环面爆破},\enspace\IndexPageList[1001]{72,112}
\subitem \leavevmode\hypertarget{idx:1002}{其中的 Calabi--Yau 超曲面},\enspace\IndexPageList[1002]{450}
\subsubitem \leavevmode\hypertarget{idx:1003}{A 模型关联函数},\enspace\IndexPageList[1003]{292,409,411}
\subsubitem \leavevmode\hypertarget{idx:1004}{Givental 函数 $I_{\mathcal V}$},\enspace\IndexPageList[1004]{407}
\subsubitem \leavevmode\hypertarget{idx:1005}{Givental 函数 $J_X$},\enspace\IndexPageList[1005]{409,411}
\subsubitem \leavevmode\hypertarget{idx:1006}{Givental 函数 $J_{\mathcal V}$},\enspace\IndexPageList[1006]{407,411}
\subsubitem \leavevmode\hypertarget{idx:1007}{Gromov--Witten 势},\enspace\IndexPageList[1007]{410,411}
\subsubitem \leavevmode\hypertarget{idx:1008}{Hodge 数},\enspace\IndexPageList[1008]{169}
\subsubitem \leavevmode\hypertarget{idx:1009}{瞬子数},\enspace\IndexPageList[1009]{292,293}
\subsubitem \leavevmode\hypertarget{idx:1010}{Kähler 锥},\enspace\IndexPageList[1010]{169,292,406}
\subsubitem \leavevmode\hypertarget{idx:1011}{Kähler 模},\enspace\IndexPageList[1011]{169,170}
\subsubitem \leavevmode\hypertarget{idx:1012}{镜像映射},\enspace\IndexPageList[1012]{181,182,187,188,189,293,408,411}
\subsubitem \leavevmode\hypertarget{idx:1013}{镜像定理},\enspace\IndexPageList[1013]{406,409,411}
\subsubitem \leavevmode\hypertarget{idx:1014}{模坐标},\enspace\IndexPageList[1014]{292}
\subsubitem \leavevmode\hypertarget{idx:1015}{超共形场论模},\enspace\IndexPageList[1015]{170}
\subsubitem \leavevmode\hypertarget{idx:1016}{环面 Kähler 锥},\enspace\IndexPageList[1016]{171}
\subsubitem \leavevmode\hypertarget{idx:1017}{环面 Kähler 模},\enspace\IndexPageList[1017]{171,292}
\subitem \leavevmode\hypertarget{idx:1018}{GKZ 分解},\enspace\IndexPageList[1018]{73,147,169}
\subitem \leavevmode\hypertarget{idx:1019}{镜像猜想},\enspace\IndexPageList[1019]{292}
\subitem \leavevmode\hypertarget{idx:1020}{其中 Calabi--Yau 超曲面的镜像},\enspace\IndexPageList[1020]{111,118,126,146,168,187}
\subsubitem \leavevmode\hypertarget{idx:1021}{$\mathcal A$ 方程组},\enspace\IndexPageList[1021]{119,408}
\subsubitem \leavevmode\hypertarget{idx:1022}{判别式轨迹},\enspace\IndexPageList[1022]{170}
\subsubitem \leavevmode\hypertarget{idx:1023}{Frobenius 方法},\enspace\IndexPageList[1023]{188}
\subsubitem \leavevmode\hypertarget{idx:1024}{Gauss--Manin 联络},\enspace\IndexPageList[1024]{292,408,409}
\subsubitem \leavevmode\hypertarget{idx:1025}{Hodge 数},\enspace\IndexPageList[1025]{169}
\subsubitem \leavevmode\hypertarget{idx:1026}{齐次坐标环},\enspace\IndexPageList[1026]{112}
\subsubitem \leavevmode\hypertarget{idx:1027}{模坐标},\enspace\IndexPageList[1027]{112,118,171,181,410}
\subsubitem \leavevmode\hypertarget{idx:1028}{单值性},\enspace\IndexPageList[1028]{293,408,410}
\subsubitem \leavevmode\hypertarget{idx:1029}{周期},\enspace\IndexPageList[1029]{187,188,408}
\subsubitem \leavevmode\hypertarget{idx:1030}{Picard--Fuchs 方程},\enspace\IndexPageList[1030]{113,119,126,181,188,408,409}
\subsubitem \leavevmode\hypertarget{idx:1031}{满足整性猜想},\enspace\IndexPageList[1031]{293,408}
\subsubitem \leavevmode\hypertarget{idx:1032}{简化模},\enspace\IndexPageList[1032]{170,171}
\subsubitem \leavevmode\hypertarget{idx:1033}{Yukawa 耦合},\enspace\IndexPageList[1033]{126,127,148,408,411}
\subitem \leavevmode\hypertarget{idx:1034}{多胞形},\enspace\IndexPageList[1034]{111}
\indexspace
\item \leavevmode\hypertarget{idx:1049}{Picard--Fuchs 方程},\enspace\IndexPageList[1049]{42,98,100,103,104,107,125,126}
\subitem \leavevmode\hypertarget{idx:1050}{通过 $\mathcal A$ 方程组计算}，\textit{参见} \hyperlink{idx:0055}{$\mathcal A$ 方程组}
\subitem \leavevmode\hypertarget{idx:1051}{通过 Griffiths--Dwork 方法计算}，\textit{参见} \hyperlink{idx:0486}{Griffiths--Dwork 方法}
\subitem \leavevmode\hypertarget{idx:1052}{例题},\enspace\IndexPageList[1052]{110,111,112,113}
\subitem \leavevmode\hypertarget{idx:1053}{$\mathcal A$ 方程组中缺失的方程},\enspace\IndexPageList[1053]{117,118,124,406,408,422}
\subitem \leavevmode\hypertarget{idx:1054}{归一化 Picard--Fuchs 方程},\enspace\IndexPageList[1054]{50,128,129,130}
\subitem \leavevmode\hypertarget{idx:1055}{Calabi--Yau 三维簇的 Picard--Fuchs 方程},\enspace\IndexPageList[1055]{102,128,350}
\subitem \leavevmode\hypertarget{idx:1056}{Calabi--Yau 环面超曲面的 Picard--Fuchs 方程}，\textit{参见} \hyperlink{idx:0099}{Calabi--Yau，Calabi--Yau 环面超曲面}
\subitem \leavevmode\hypertarget{idx:1057}{Calabi--Yau 环面完全交的镜像的 Picard--Fuchs 方程},\enspace\IndexPageList[1057]{123,124}
\subitem \leavevmode\hypertarget{idx:1058}{五次镜像的 Picard--Fuchs 方程}，\textit{参见} \hyperlink{idx:1175}{五次镜像，Picard--Fuchs 方程}
\subitem \leavevmode\hypertarget{idx:1059}{解},\enspace\IndexPageList[1059]{42,119,124,172,179,183,184,286,291}，\textit{另见} \hyperlink{idx:0344}{Frobenius 方法}
\indexspace
\item \leavevmode\hypertarget{idx:1060}{Picard--Fuchs 理想},\enspace\IndexPageList[1060]{97,100,117}
\subitem \leavevmode\hypertarget{idx:1061}{相应的 $\mathcal D$ 模},\enspace\IndexPageList[1061]{97,98,124}
\indexspace
\item \leavevmode\hypertarget{idx:1062}{Picard--Fuchs 算子},\enspace\IndexPageList[1062]{351,352}
\subitem \leavevmode\hypertarget{idx:1063}{$\hbar$ 齐次化},\enspace\IndexPageList[1063]{351,352}
\subitem \leavevmode\hypertarget{idx:1064}{给出量子上同调中的关系},\enspace\IndexPageList[1064]{352}
\subitem \leavevmode\hypertarget{idx:1065}{主部},\enspace\IndexPageList[1065]{351}
\indexspace
\item \leavevmode\hypertarget{idx:1066}{Picard--Lefschetz 变换}，\textit{参见} \hyperlink{idx:0901}{单值性变换}
\indexspace
\item \leavevmode\hypertarget{idx:1067}{Planck 常数 $\hbar$},\enspace\IndexPageList[1067]{436}
\indexspace
\item \leavevmode\hypertarget{idx:1068}{Poincaré 对偶},\enspace\IndexPageList[1068]{105,404}
\subitem \leavevmode\hypertarget{idx:1069}{轨形的 Poincaré 对偶},\enspace\IndexPageList[1069]{194,195,206,328,431}
\indexspace
\item \leavevmode\hypertarget{idx:1070}{Poisson 括号},\enspace\IndexPageList[1070]{436,437}
\indexspace
\item \leavevmode\hypertarget{idx:1101}{$\widehat Q$},\enspace\IndexPageList[1101]{50,51,363,366,367,370,371,373,374,376}
\subitem \leavevmode\hypertarget{idx:1102}{定义},\enspace\IndexPageList[1102]{360}
\subitem \leavevmode\hypertarget{idx:1103}{凸凹丛的 $\widehat Q$},\enspace\IndexPageList[1103]{372}
\subitem \leavevmode\hypertarget{idx:1104}{多重覆盖情形的 $\widehat Q$},\enspace\IndexPageList[1104]{375,378}
\subitem \leavevmode\hypertarget{idx:1105}{是 Euler 数据},\enspace\IndexPageList[1105]{362}
\indexspace
\item \leavevmode\hypertarget{source:490}{}\hypertarget{idx:1228}{Reid 幻想},\enspace\IndexPageList[1228]{166,428}
\indexspace
\item \leavevmode\hypertarget{idx:1232}{Ricci 曲率张量},\enspace\IndexPageList[1232]{33}
\indexspace
\item \leavevmode\hypertarget{idx:1233}{Riemann--Roch 定理},\enspace\IndexPageList[1233]{312,317,318}
\indexspace
\item \leavevmode\hypertarget{idx:1237}{Satake--Baily--Borel 紧化},\enspace\IndexPageList[1237]{155}
\indexspace
\item \leavevmode\hypertarget{idx:1238}{Schrödinger 绘景},\enspace\IndexPageList[1238]{436}
\indexspace
\item \leavevmode\hypertarget{idx:1243}{Seiberg--Witten 曲线},\enspace\IndexPageList[1243]{429}
\indexspace
\item \leavevmode\hypertarget{idx:1246}{Serre--Grothendieck 对偶},\enspace\IndexPageList[1246]{78}
\indexspace
\item \leavevmode\hypertarget{idx:1247}{Severi 簇},\enspace\IndexPageList[1247]{221}
\indexspace
\item \leavevmode\hypertarget{idx:1284}{Stanley--Reisner 理想},\enspace\IndexPageList[1284]{246,247}
\subitem \leavevmode\hypertarget{idx:1285}{量子 Stanley--Reisner 理想},\enspace\IndexPageList[1285]{247,248,415}
\indexspace
\item \leavevmode\hypertarget{source:491}{}\hypertarget{idx:1325}{T 对偶},\enspace\IndexPageList[1325]{425,426}
\indexspace
\item \leavevmode\hypertarget{idx:1335}{Toda 格},\enspace\IndexPageList[1335]{251}
\indexspace
\item \leavevmode\hypertarget{idx:1381}{TQFT}，\textit{参见} \hyperlink{idx:1336}{拓扑量子场论}
\indexspace
\item \leavevmode\hypertarget{idx:1382}{TRR}，\textit{参见} \hyperlink{idx:1344}{拓扑递推关系}
\indexspace
\item \leavevmode\hypertarget{idx:1384}{$V$ 流形},\enspace\IndexPageList[1384]{431}，\textit{另见} \hyperlink{idx:0990}{轨形}
\indexspace
\item \leavevmode\hypertarget{idx:1388}{$\mathcal V_{\beta,k}$},\enspace\IndexPageList[1388]{387}，\textit{另见} \hyperlink{idx:1395}{$\mathcal V_{d,k}$}
\indexspace
\item \leavevmode\hypertarget{idx:1389}{$\mathcal V'_{\beta,k,i}$},\enspace\IndexPageList[1389]{387}，\textit{另见} \hyperlink{idx:1396}{$\mathcal V'_{d,k,i}$}
\indexspace
\item \leavevmode\hypertarget{idx:1390}{$\mathcal V_d$}
\subitem \leavevmode\hypertarget{idx:1391}{$\mathcal O_{\mathbb P^n}(\ell)$ 的 $\mathcal V_d$},\enspace\IndexPageList[1391]{360,363,367}
\subitem \leavevmode\hypertarget{idx:1392}{凸凹丛的 $\mathcal V_d$},\enspace\IndexPageList[1392]{372,374,381,398}
\subitem \leavevmode\hypertarget{idx:1393}{五次三维簇的 $\mathcal V_d$},\enspace\IndexPageList[1393]{204,207,224,311,322,325,355}
\subsubitem \leavevmode\hypertarget{idx:1394}{多重覆盖情形},\enspace\IndexPageList[1394]{375}
\indexspace
\item \leavevmode\hypertarget{idx:1395}{$\mathcal V_{d,k}$},\enspace\IndexPageList[1395]{380,381,383}
\indexspace
\item \leavevmode\hypertarget{idx:1396}{$\mathcal V'_{d,k,i}$},\enspace\IndexPageList[1396]{381,395,396}
\indexspace
\item \leavevmode\hypertarget{idx:1386}{Vafa--Intriligator 公式},\enspace\IndexPageList[1386]{250}
\indexspace
\item \leavevmode\hypertarget{idx:1398}{Verlinde 代数},\enspace\IndexPageList[1398]{250}
\indexspace
\item \leavevmode\hypertarget{idx:1401}{Virasoro 猜想},\enspace\IndexPageList[1401]{333,448}
\subitem \leavevmode\hypertarget{idx:1402}{对下列对象的推论}
\subsubitem \leavevmode\hypertarget{idx:1403}{引力后裔},\enspace\IndexPageList[1403]{334}
\subsubitem \leavevmode\hypertarget{idx:1404}{Gromov--Witten 不变量},\enspace\IndexPageList[1404]{333,334}
\indexspace
\item \leavevmode\hypertarget{idx:1399}{Virasoro 代数},\enspace\IndexPageList[1399]{333,334,446,447,448}
\subitem \leavevmode\hypertarget{idx:1400}{中心扩张},\enspace\IndexPageList[1400]{333,448}
\indexspace
\item \leavevmode\hypertarget{idx:1405}{Virasoro 约束},\enspace\IndexPageList[1405]{234}
\indexspace
\item \leavevmode\hypertarget{source:492}{}\hypertarget{idx:1430}{WDVV 方程}，\textit{参见} \hyperlink{idx:0226}{Dubrovin 联络}；\hyperlink{idx:0552}{Gromov--Witten 势}
\indexspace
\item \leavevmode\hypertarget{idx:1439}{Witten 猜想},\enspace\IndexPageList[1439]{333}
\indexspace
\item \leavevmode\hypertarget{idx:1443}{Yang--Mills 理论},\enspace\IndexPageList[1443]{440}
\indexspace
\item \leavevmode\hypertarget{idx:1444}{Yukawa 耦合},\enspace\IndexPageList[1444]{30,31,32,33,439,444}，\textit{另见} \hyperlink{idx:0183}{关联函数}
\subitem \leavevmode\hypertarget{idx:1445}{归一化 Yukawa 耦合},\enspace\IndexPageList[1445]{127,128,129,130,131,132,133,134,135,280,282,291,409}
\subsubitem \leavevmode\hypertarget{idx:1446}{定义},\enspace\IndexPageList[1446]{133}
\subsubitem \leavevmode\hypertarget{idx:1447}{用周期表示},\enspace\IndexPageList[1447]{412,422}
\subsubitem \leavevmode\hypertarget{idx:1448}{负号},\enspace\IndexPageList[1448]{31,44,130}
\subsubitem \leavevmode\hypertarget{idx:1449}{五次镜像的归一化 Yukawa 耦合}，\textit{参见} \hyperlink{idx:1179}{五次镜像，Yukawa 耦合}
\subsubitem \leavevmode\hypertarget{idx:1450}{与镜像映射的关系},\enspace\IndexPageList[1450]{132,133,134,135,174}
\subsubitem \leavevmode\hypertarget{idx:1451}{与 Gauss--Manin 联络的关系},\enspace\IndexPageList[1451]{128,131}
\subitem \leavevmode\hypertarget{idx:1452}{Calabi--Yau 三维簇的 Yukawa 耦合},\enspace\IndexPageList[1452]{125,126,127,128,129,130,131,132,133,134,135,279,282,350}
\subsubitem \leavevmode\hypertarget{idx:1453}{一维模},\enspace\IndexPageList[1453]{125,128,129,130}
\subsubitem \leavevmode\hypertarget{idx:1454}{二维模},\enspace\IndexPageList[1454]{125,126,127}
\subsubitem \leavevmode\hypertarget{idx:1455}{$r$ 维模},\enspace\IndexPageList[1455]{130,131,132}
\subitem \leavevmode\hypertarget{idx:1456}{环面 Yukawa 耦合},\enspace\IndexPageList[1456]{291,292,408}，\textit{另见} \hyperlink{idx:0190}{关联函数，B 模型关联函数，环面情形}
\subsubitem \leavevmode\hypertarget{idx:1457}{与归一化耦合的区别},\enspace\IndexPageList[1457]{291}
\indexspace
\item \leavevmode\hypertarget{idx:1458}{Zariski $p$ 形式},\enspace\IndexPageList[1458]{432}
\indexspace
\item \leavevmode\hypertarget{idx:0144}{凹丛},\enspace\IndexPageList[0144]{372}
\subitem \leavevmode\hypertarget{idx:0145}{多重覆盖情形},\enspace\IndexPageList[0145]{375}
\indexspace
\item \leavevmode\hypertarget{idx:0965}{八面体},\enspace\IndexPageList[0965]{74}
\indexspace
\item \leavevmode\hypertarget{idx:1244}{半丰沛},\enspace\IndexPageList[1244]{77}
\indexspace
\item \leavevmode\hypertarget{idx:1245}{半正},\enspace\IndexPageList[1245]{208,210,212,233}
\indexspace
\item \leavevmode\hypertarget{idx:0204}{保典范},\enspace\IndexPageList[0204]{77}
\indexspace
\item \leavevmode\hypertarget{idx:1087}{本原集合},\enspace\IndexPageList[1087]{61,73,246,247,248}
\subitem \leavevmode\hypertarget{idx:1088}{丰沛性判据},\enspace\IndexPageList[1088]{63}
\indexspace
\item \leavevmode\hypertarget{idx:1085}{本原上同调},\enspace\IndexPageList[1085]{106,109,276,287,384}
\subitem \leavevmode\hypertarget{idx:1086}{定义},\enspace\IndexPageList[1086]{97}
\indexspace
\item \leavevmode\hypertarget{idx:1089}{本原收缩},\enspace\IndexPageList[1089]{155,166,171}
\subitem \leavevmode\hypertarget{idx:1090}{I 型},\enspace\IndexPageList[1090]{156,157,166,167,170,249,250}
\subitem \leavevmode\hypertarget{idx:1091}{II 型},\enspace\IndexPageList[1091]{156,157,166}
\subitem \leavevmode\hypertarget{idx:1092}{III 型},\enspace\IndexPageList[1092]{156,157,166,169,250}
\indexspace
\item \leavevmode\hypertarget{idx:1334}{编时乘积},\enspace\IndexPageList[1334]{440}
\indexspace
\item \leavevmode\hypertarget{idx:1438}{变量的权},\enspace\IndexPageList[1438]{381,388,401,402}
\indexspace
\item \leavevmode\hypertarget{idx:1434}{表示的权},\enspace\IndexPageList[1434]{299,301,303,308,315,367,392,393}
\subitem \leavevmode\hypertarget{idx:1435}{障碍权},\enspace\IndexPageList[1435]{312,313,314,315,320}
\subitem \leavevmode\hypertarget{idx:1436}{法丛的权},\enspace\IndexPageList[1436]{308,312}
\subitem \leavevmode\hypertarget{idx:1437}{切权},\enspace\IndexPageList[1437]{301,304,308,310,313,320}
\indexspace
\item \leavevmode\hypertarget{idx:0081}{玻色}
\subitem \leavevmode\hypertarget{idx:0082}{玻色分量},\enspace\IndexPageList[0082]{450,452}
\subitem \leavevmode\hypertarget{idx:0083}{玻色场},\enspace\IndexPageList[0083]{442}
\subitem \leavevmode\hypertarget{idx:0084}{玻色变量},\enspace\IndexPageList[0084]{450,451}
\indexspace
\item \leavevmode\hypertarget{idx:0080}{玻色子},\enspace\IndexPageList[0080]{25,442,443,448}
\indexspace
\item \leavevmode\hypertarget{idx:0203}{产生算子},\enspace\IndexPageList[0203]{438,441}
\indexspace
\item \leavevmode\hypertarget{idx:0316}{场}
\subitem \leavevmode\hypertarget{idx:0317}{玻色场}，\textit{参见} \hyperlink{idx:0083}{玻色，玻色场}
\subitem \leavevmode\hypertarget{idx:0318}{经典场},\enspace\IndexPageList[0318]{434,437,439}
\subitem \leavevmode\hypertarget{idx:0319}{费米场}，\textit{参见} \hyperlink{idx:0312}{费米，费米场}
\subitem \leavevmode\hypertarget{idx:0320}{初级场},\enspace\IndexPageList[0320]{447}
\subitem \leavevmode\hypertarget{idx:0321}{量子场},\enspace\IndexPageList[0321]{437,439}
\subitem \leavevmode\hypertarget{idx:0322}{场论},\enspace\IndexPageList[0322]{434,437}，\textit{另见} \hyperlink{idx:0162}{共形场论}；\hyperlink{idx:1155}{量子场论}
\indexspace
\item \leavevmode\hypertarget{idx:0323}{场强},\enspace\IndexPageList[0323]{440,441}
\indexspace
\item \leavevmode\hypertarget{idx:0620}{超 Kähler},\enspace\IndexPageList[0620]{425}
\indexspace
\item \leavevmode\hypertarget{idx:1296}{超场},\enspace\IndexPageList[1296]{449,450,452}
\subitem \leavevmode\hypertarget{idx:1297}{手征超场}，\textit{参见} \hyperlink{idx:0121}{手征超场}
\indexspace
\item \leavevmode\hypertarget{idx:1309}{超对称场论},\enspace\IndexPageList[1309]{429,443,449}
\indexspace
\item \leavevmode\hypertarget{idx:1310}{超对称性},\enspace\IndexPageList[1310]{24,25,26,30,33,447,448,449}，\textit{另见} \hyperlink{idx:1254}{$\sigma$ 模型，超对称 $\sigma$ 模型}
\subitem \leavevmode\hypertarget{idx:1311}{$N=2$ 超对称代数},\enspace\IndexPageList[1311]{449}
\subitem \leavevmode\hypertarget{idx:1312}{超对称变换},\enspace\IndexPageList[1312]{25,443,445,449}
\indexspace
\item \leavevmode\hypertarget{idx:0280}{超额法丛},\enspace\IndexPageList[0280]{196,197}
\indexspace
\item \leavevmode\hypertarget{idx:0277}{超额维数},\enspace\IndexPageList[0277]{331}
\subitem \leavevmode\hypertarget{idx:0278}{稳定映射模空间的超额维数},\enspace\IndexPageList[0278]{205}
\subitem \leavevmode\hypertarget{idx:0279}{相交理论中的超额维数},\enspace\IndexPageList[0279]{205}
\indexspace
\item \leavevmode\hypertarget{idx:1292}{超共形}
\subitem \leavevmode\hypertarget{idx:1293}{$N=2$ 超共形代数},\enspace\IndexPageList[1293]{25}
\subitem \leavevmode\hypertarget{idx:1294}{$N=2$ 超共形表示},\enspace\IndexPageList[1294]{79}
\subitem \leavevmode\hypertarget{idx:1295}{超共形场论（SCFT）},\enspace\IndexPageList[1295]{25,27,29,30,32,33,35,40,76,79}
\indexspace
\item \leavevmode\hypertarget{idx:0615}{超几何方程},\enspace\IndexPageList[0615]{113,114,115,116,117,118,119,120,121,122,123,124,184,356}，\textit{另见} \hyperlink{idx:0055}{$\mathcal A$ 方程组}；\hyperlink{idx:0449}{GKZ 方程组}
\indexspace
\item \leavevmode\hypertarget{idx:0616}{超几何函数},\enspace\IndexPageList[0616]{118,187}
\subitem \leavevmode\hypertarget{idx:0617}{$p$ 进超几何函数},\enspace\IndexPageList[0617]{43}
\subitem \leavevmode\hypertarget{idx:0618}{广义超几何函数},\enspace\IndexPageList[0618]{43}
\subitem \leavevmode\hypertarget{idx:0619}{$HG$}，\textit{参见} \hyperlink{idx:0579}{$HG$}
\indexspace
\item \leavevmode\hypertarget{idx:1291}{超交换},\enspace\IndexPageList[1291]{242,244,252,253,256,257,263,457}
\indexspace
\item \leavevmode\hypertarget{idx:1299}{超流形},\enspace\IndexPageList[1299]{252,253,262,449}
\subitem \leavevmode\hypertarget{idx:1300}{超流形上的偏导数},\enspace\IndexPageList[1300]{252,327}
\subitem \leavevmode\hypertarget{idx:1301}{切丛},\enspace\IndexPageList[1301]{263}
\indexspace
\item \leavevmode\hypertarget{idx:1302}{超势},\enspace\IndexPageList[1302]{449,451,453}
\subitem \leavevmode\hypertarget{idx:1303}{五次超势},\enspace\IndexPageList[1303]{452}
\indexspace
\item \leavevmode\hypertarget{idx:1304}{超弦理论}
\subitem \leavevmode\hypertarget{idx:1305}{对偶性},\enspace\IndexPageList[1305]{427}
\subitem \leavevmode\hypertarget{idx:1306}{杂化超弦理论},\enspace\IndexPageList[1306]{25,26,427,443,444}
\subitem \leavevmode\hypertarget{idx:1307}{I 型超弦理论},\enspace\IndexPageList[1307]{25}
\subitem \leavevmode\hypertarget{idx:1308}{II 型超弦理论},\enspace\IndexPageList[1308]{25,168,427}
\indexspace
\item \leavevmode\hypertarget{idx:1313}{超向量空间},\enspace\IndexPageList[1313]{454,457}
\indexspace
\item \leavevmode\hypertarget{idx:1298}{超引力},\enspace\IndexPageList[1298]{33}
\indexspace
\item \leavevmode\hypertarget{idx:1242}{次级多胞形},\enspace\IndexPageList[1242]{143}
\indexspace
\item \leavevmode\hypertarget{idx:1239}{次级扇},\enspace\IndexPageList[1239]{66,140,143,144,158,159,165,166,171,178,179}，\textit{另见} \hyperlink{idx:0446}{GKZ 分解}
\subitem \leavevmode\hypertarget{idx:1240}{扩大次级扇},\enspace\IndexPageList[1240]{453}
\subitem \leavevmode\hypertarget{idx:1241}{例题},\enspace\IndexPageList[1241]{73,74,147,149,163}
\indexspace
\item \leavevmode\hypertarget{idx:0696}{大半径极限点},\enspace\IndexPageList[0696]{138,151,152,153,154,159,161,172,173,175,176,178,179,268,274,280,287,290,292}
\subitem \leavevmode\hypertarget{idx:0697}{等价性},\enspace\IndexPageList[0697]{154,175}
\subitem \leavevmode\hypertarget{idx:0698}{例题},\enspace\IndexPageList[0698]{164}
\subitem \leavevmode\hypertarget{idx:0699}{该点处的切空间},\enspace\IndexPageList[0699]{181}
\indexspace
\item \leavevmode\hypertarget{idx:0040}{代数叠},\enspace\IndexPageList[0040]{192}
\subitem \leavevmode\hypertarget{idx:0041}{代数叠上的丛}，\textit{参见} \hyperlink{idx:0261}{Euler 类，代数叠与粗模空间之间的比较}
\subitem \leavevmode\hypertarget{idx:0042}{Deligne--Mumford 代数叠},\enspace\IndexPageList[0042]{201}
\subitem \leavevmode\hypertarget{idx:0043}{代数叠上的相交理论},\enspace\IndexPageList[0043]{201,302}
\subitem \leavevmode\hypertarget{idx:0044}{商叠},\enspace\IndexPageList[0044]{201}
\subitem \leavevmode\hypertarget{idx:0045}{光滑代数叠},\enspace\IndexPageList[0045]{193,302}
\indexspace
\item \leavevmode\hypertarget{idx:0039}{代数空间},\enspace\IndexPageList[0039]{193}
\indexspace
\item \leavevmode\hypertarget{idx:0910}{单项式—除子镜像映射},\enspace\IndexPageList[0910]{80,83,181,291,294,295}
\indexspace
\item \leavevmode\hypertarget{idx:0909}{单项式序},\enspace\IndexPageList[0909]{248}
\indexspace
\item \leavevmode\hypertarget{idx:0901}{单值性变换},\enspace\IndexPageList[0901]{98,100,104,137,339}
\subitem \leavevmode\hypertarget{idx:0902}{对数},\enspace\IndexPageList[0902]{98,104,134,137,165,269,270,272,339}
\subitem \leavevmode\hypertarget{idx:0903}{A 模型联络的单值性变换}，\textit{参见} \hyperlink{idx:0011}{A 模型，联络，单值性}
\subitem \leavevmode\hypertarget{idx:0904}{拟幂幺单值性变换},\enspace\IndexPageList[0904]{98}
\subitem \leavevmode\hypertarget{idx:0905}{幂幺单值性变换},\enspace\IndexPageList[0905]{98,100,101,137,269,270}
\subsubitem \leavevmode\hypertarget{idx:0906}{最大幂幺}，\textit{参见} \hyperlink{idx:0765}{最大幂幺}
\indexspace
\item \leavevmode\hypertarget{idx:0900}{单值性定理},\enspace\IndexPageList[0900]{98}
\indexspace
\item \leavevmode\hypertarget{idx:0907}{单值性权滤过},\enspace\IndexPageList[0907]{99,103,104,127,128,138,172,173,177,184,273,291}
\subitem \leavevmode\hypertarget{idx:0908}{A 模型联络的单值性权滤过}，\textit{参见} \hyperlink{idx:0012}{A 模型，联络，单值性权滤过}
\indexspace
\item \leavevmode\hypertarget{idx:0243}{等变 Chern 类},\enspace\IndexPageList[0243]{299,303,389}
\indexspace
\item \leavevmode\hypertarget{idx:0250}{等变超平面类},\enspace\IndexPageList[0250]{51,124,358,369,389,391}
\subitem \leavevmode\hypertarget{idx:0251}{定义},\enspace\IndexPageList[0251]{300}
\subitem \leavevmode\hypertarget{idx:0252}{$N_d$ 的等变超平面类}，\textit{参见} \hyperlink{idx:0930}{$N_d$，射影空间，等变超平面类}
\indexspace
\item \leavevmode\hypertarget{idx:0253}{等变积分},\enspace\IndexPageList[0253]{303,358,390,391,396,399,400}
\subitem \leavevmode\hypertarget{idx:0254}{定义},\enspace\IndexPageList[0254]{302}
\indexspace
\item \leavevmode\hypertarget{idx:0244}{等变上同调},\enspace\IndexPageList[0244]{50,124}
\subitem \leavevmode\hypertarget{idx:0245}{定义},\enspace\IndexPageList[0245]{299}
\subitem \leavevmode\hypertarget{idx:0246}{参数 $\hbar$ 的说明},\enspace\IndexPageList[0246]{342}
\subitem \leavevmode\hypertarget{idx:0247}{局部化},\enspace\IndexPageList[0247]{300,359}，\textit{另见} \hyperlink{idx:0739}{局部化定理，等变上同调的局部化定理}
\subitem \leavevmode\hypertarget{idx:0248}{$\mathbb P^n$ 的等变上同调},\enspace\IndexPageList[0248]{300,301,358}
\subitem \leavevmode\hypertarget{idx:0249}{限制映射},\enspace\IndexPageList[0249]{366,367}
\indexspace
\item \leavevmode\hypertarget{idx:0242}{等变特征类},\enspace\IndexPageList[0242]{361}
\indexspace
\item \leavevmode\hypertarget{idx:0255}{等变向量丛},\enspace\IndexPageList[0255]{299,321}
\indexspace
\item \leavevmode\hypertarget{idx:0745}{低能极限},\enspace\IndexPageList[0745]{33}
\indexspace
\item \leavevmode\hypertarget{idx:1210}{递推关系}
\subitem \leavevmode\hypertarget{idx:1211}{微分方程的递推关系},\enspace\IndexPageList[1211]{118,185,186}
\subitem \leavevmode\hypertarget{idx:1212}{Gromov--Witten 不变量的递推关系}，\textit{参见} \hyperlink{idx:0523}{Gromov--Witten 不变量，计算}
\indexspace
\item \leavevmode\hypertarget{source:489}{}\hypertarget{idx:1209}{递推引理}，\textit{参见} \hyperlink{idx:0860}{镜像定理，Givental 镜像定理，递推引理}
\indexspace
\item \leavevmode\hypertarget{idx:0112}{典范延拓},\enspace\IndexPageList[0112]{98,177,270,272,283}
\subitem \leavevmode\hypertarget{idx:0113}{A 模型联络的典范延拓}，\textit{参见} \hyperlink{idx:0005}{A 模型，联络，典范延拓}
\indexspace
\item \leavevmode\hypertarget{idx:0238}{电磁势},\enspace\IndexPageList[0238]{440}
\indexspace
\item \leavevmode\hypertarget{idx:1283}{叠}，\textit{参见} \hyperlink{idx:0040}{代数叠}
\indexspace
\item \leavevmode\hypertarget{idx:0225}{对偶化层},\enspace\IndexPageList[0225]{34,69,430,432}
\indexspace
\item \leavevmode\hypertarget{idx:0224}{对偶锥},\enspace\IndexPageList[0224]{54}
\indexspace
\item \leavevmode\hypertarget{idx:1071}{多胞形},\enspace\IndexPageList[1071]{56}
\subitem \leavevmode\hypertarget{idx:1072}{刻面},\enspace\IndexPageList[1072]{56}
\subitem \leavevmode\hypertarget{idx:1073}{整多胞形},\enspace\IndexPageList[1073]{56}
\subitem \leavevmode\hypertarget{idx:1074}{Newton 多胞形},\enspace\IndexPageList[1074]{58,83}
\subitem \leavevmode\hypertarget{idx:1075}{法扇},\enspace\IndexPageList[1075]{57,74,76,158}
\subitem \leavevmode\hypertarget{idx:1076}{归一化体积},\enspace\IndexPageList[1076]{82}
\subitem \leavevmode\hypertarget{idx:1077}{极多胞形},\enspace\IndexPageList[1077]{57,69,74,76,143}
\subsubitem \leavevmode\hypertarget{idx:1078}{极对偶},\enspace\IndexPageList[1078]{58}
\subitem \leavevmode\hypertarget{idx:1079}{自反多胞形},\enspace\IndexPageList[1079]{83}，\textit{参见} \hyperlink{idx:1218}{自反多胞形}
\subitem \leavevmode\hypertarget{idx:1080}{多胞形环},\enspace\IndexPageList[1080]{57,59,60}
\subitem \leavevmode\hypertarget{idx:1081}{单纯多胞形},\enspace\IndexPageList[1081]{432}
\subitem \leavevmode\hypertarget{idx:1082}{对应的环面簇},\enspace\IndexPageList[1082]{57}
\indexspace
\item \leavevmode\hypertarget{idx:0915}{多重覆盖},\enspace\IndexPageList[0915]{225,230,232,311,366}，\textit{另见} \hyperlink{idx:0628}{瞬子数}
\subitem \leavevmode\hypertarget{idx:0916}{对 Gromov--Witten 不变量的贡献},\enspace\IndexPageList[0916]{225,227,311,312,313,314,315,316,317,375}
\subsubitem \leavevmode\hypertarget{idx:0917}{Lian--Liu--Yau 的证明},\enspace\IndexPageList[0917]{378}
\subsubitem \leavevmode\hypertarget{idx:0918}{Pandharipande 的证明},\enspace\IndexPageList[0918]{314,315,316}
\subsubitem \leavevmode\hypertarget{idx:0919}{$d=2$ 时的证明},\enspace\IndexPageList[0919]{312,313,314}
\subitem \leavevmode\hypertarget{idx:0920}{对虚拟基本类的贡献},\enspace\IndexPageList[0920]{225,228,317}
\subitem \leavevmode\hypertarget{idx:0921}{椭圆多重覆盖},\enspace\IndexPageList[0921]{232,237,323}
\subitem \leavevmode\hypertarget{idx:0922}{六结点五次曲线的多重覆盖},\enspace\IndexPageList[0922]{316,317}，\textit{另见} \hyperlink{idx:0218}{六结点五次曲线的二重覆盖}
\subitem \leavevmode\hypertarget{idx:0923}{非孤立曲线的多重覆盖},\enspace\IndexPageList[0923]{231}
\indexspace
\item \leavevmode\hypertarget{idx:0924}{多重镜像},\enspace\IndexPageList[0924]{68,83,144}
\indexspace
\item \leavevmode\hypertarget{idx:0937}{二次曲面网},\enspace\IndexPageList[0937]{389}
\indexspace
\item \leavevmode\hypertarget{idx:0947}{法锥},\enspace\IndexPageList[0947]{195,197,200,201,204,383}
\subitem \leavevmode\hypertarget{idx:0948}{内蕴法锥},\enspace\IndexPageList[0948]{201,202}
\subsubitem \leavevmode\hypertarget{idx:0949}{例题},\enspace\IndexPageList[0949]{202}
\subitem \leavevmode\hypertarget{idx:0950}{细化 Euler 类},\enspace\IndexPageList[0950]{196}
\subitem \leavevmode\hypertarget{idx:0951}{虚法锥},\enspace\IndexPageList[0951]{200}
\indexspace
\item \leavevmode\hypertarget{idx:0332}{翻越},\enspace\IndexPageList[0332]{66,156,170,248,249}
\subitem \leavevmode\hypertarget{idx:0333}{例题},\enspace\IndexPageList[0333]{164}
\subitem \leavevmode\hypertarget{idx:0334}{广义翻越},\enspace\IndexPageList[0334]{160}
\subitem \leavevmode\hypertarget{idx:0335}{平凡翻越},\enspace\IndexPageList[0335]{160}
\indexspace
\item \leavevmode\hypertarget{idx:0328}{翻转},\enspace\IndexPageList[0328]{160,162}
\subitem \leavevmode\hypertarget{idx:0329}{例题},\enspace\IndexPageList[0329]{164}
\subitem \leavevmode\hypertarget{idx:0330}{几何意义},\enspace\IndexPageList[0330]{160}，\textit{另见} \hyperlink{idx:0334}{翻越，广义翻越}
\subitem \leavevmode\hypertarget{idx:0331}{平凡翻转},\enspace\IndexPageList[0331]{160,161,164,165,179}
\indexspace
\item \leavevmode\hypertarget{idx:0050}{反典范}
\subitem \leavevmode\hypertarget{idx:0051}{反典范除子},\enspace\IndexPageList[0051]{69,71,72,74,77}
\subitem \leavevmode\hypertarget{idx:0052}{反典范线性系},\enspace\IndexPageList[0052]{78,79}
\indexspace
\item \leavevmode\hypertarget{idx:0053}{反对易子},\enspace\IndexPageList[0053]{445}
\indexspace
\item \leavevmode\hypertarget{idx:0114}{范畴商},\enspace\IndexPageList[0114]{60}
\indexspace
\item \leavevmode\hypertarget{idx:0945}{非等变极限},\enspace\IndexPageList[0945]{51,369,373,374,375,376,377,380,389,395,402,403}
\subitem \leavevmode\hypertarget{idx:0946}{定义},\enspace\IndexPageList[0946]{299}
\indexspace
\item \leavevmode\hypertarget{idx:0943}{非微扰态},\enspace\IndexPageList[0943]{425}
\indexspace
\item \leavevmode\hypertarget{idx:0942}{非线性 $\sigma$ 模型}，\textit{参见} \hyperlink{idx:1249}{$\sigma$ 模型}
\indexspace
\item \leavevmode\hypertarget{idx:0944}{非重整化定理},\enspace\IndexPageList[0944]{30,31,446}
\indexspace
\item \leavevmode\hypertarget{idx:0311}{费米}
\subitem \leavevmode\hypertarget{idx:0312}{费米场},\enspace\IndexPageList[0312]{442,443}，\textit{另见} \hyperlink{idx:1253}{$\sigma$ 模型，费米场}
\subitem \leavevmode\hypertarget{idx:0313}{费米变量},\enspace\IndexPageList[0313]{442,450,451}
\indexspace
\item \leavevmode\hypertarget{idx:0309}{费米子},\enspace\IndexPageList[0309]{25,442,443,444,448}
\subitem \leavevmode\hypertarget{idx:0310}{与超交换变量的关系},\enspace\IndexPageList[0310]{253,442}
\indexspace
\item \leavevmode\hypertarget{idx:0139}{复反射},\enspace\IndexPageList[0139]{431}
\indexspace
\item \leavevmode\hypertarget{idx:0140}{复化}
\subitem \leavevmode\hypertarget{idx:0141}{复化 Kähler 类},\enspace\IndexPageList[0141]{26,29,30,31,32,34,35,38,41,45,227,240,241,248,251,253,326,327,457}
\subitem \leavevmode\hypertarget{idx:0142}{复化 Kähler 模}，\textit{参见} \hyperlink{idx:0882}{模，Kähler 模}
\subitem \leavevmode\hypertarget{idx:0143}{复化 Kähler 空间},\enspace\IndexPageList[0143]{35,41,150,151,175,266,268,274}
\indexspace
\item \leavevmode\hypertarget{idx:0276}{赋值映射},\enspace\IndexPageList[0276]{195,202,212,228,233,245,321,324,355,380,383,388,398}
\indexspace
\item \leavevmode\hypertarget{idx:0171}{共轭动量},\enspace\IndexPageList[0171]{435,437}
\indexspace
\item \leavevmode\hypertarget{idx:0167}{共形不变性},\enspace\IndexPageList[0167]{443,447}
\indexspace
\item \leavevmode\hypertarget{idx:0162}{共形场论},\enspace\IndexPageList[0162]{446,447,448,449}
\subitem \leavevmode\hypertarget{idx:0163}{规范的选择},\enspace\IndexPageList[0163]{447}
\subitem \leavevmode\hypertarget{idx:0164}{规范群},\enspace\IndexPageList[0164]{447}
\subitem \leavevmode\hypertarget{idx:0165}{$\sigma$ 模型的共形场论},\enspace\IndexPageList[0165]{448}
\indexspace
\item \leavevmode\hypertarget{idx:0161}{共形代数},\enspace\IndexPageList[0161]{25,446,447}
\indexspace
\item \leavevmode\hypertarget{idx:0166}{共形群},\enspace\IndexPageList[0166]{447,448}
\indexspace
\item \leavevmode\hypertarget{idx:0183}{关联函数},\enspace\IndexPageList[0183]{24,29,79,175,181,264,283,326,387,439,444}
\subitem \leavevmode\hypertarget{idx:0184}{A 模型关联函数},\enspace\IndexPageList[0184]{30,31,32,34,130,226,282,289,291,408,410,412,445,446}，\textit{另见} \hyperlink{idx:1183}{五次三维簇，A 模型关联函数}
\subsubitem \leavevmode\hypertarget{idx:0185}{定义},\enspace\IndexPageList[0185]{244}
\subsubitem \leavevmode\hypertarget{idx:0186}{翻越的影响},\enspace\IndexPageList[0186]{250}
\subsubitem \leavevmode\hypertarget{idx:0187}{用瞬子数表示},\enspace\IndexPageList[0187]{249}
\subsubitem \leavevmode\hypertarget{idx:0188}{Calabi--Yau 三维簇的关联函数},\enspace\IndexPageList[0188]{227,249,278,279,350}
\subitem \leavevmode\hypertarget{idx:0189}{B 模型关联函数},\enspace\IndexPageList[0189]{31,32,133,279,282,288,291,408,410,446}，\textit{另见} \hyperlink{idx:0064}{B 模型，势函数}；\hyperlink{idx:1444}{Yukawa 耦合}；\hyperlink{idx:1164}{五次镜像}
\subsubitem \leavevmode\hypertarget{idx:0190}{环面情形},\enspace\IndexPageList[0190]{291}
\subitem \leavevmode\hypertarget{idx:0191}{亏格 $g$ 的关联函数},\enspace\IndexPageList[0191]{457}
\subitem \leavevmode\hypertarget{idx:0192}{$n$ 点关联函数},\enspace\IndexPageList[0192]{439,457}
\subitem \leavevmode\hypertarget{idx:0193}{Frobenius 代数的关联函数},\enspace\IndexPageList[0193]{455}
\subitem \leavevmode\hypertarget{idx:0194}{五次情形}，\textit{参见} \hyperlink{idx:1164}{五次镜像}；\hyperlink{idx:1181}{五次三维簇}
\subitem \leavevmode\hypertarget{idx:0195}{三点关联函数},\enspace\IndexPageList[0195]{29,30,31,32,33,244,249,250,278,279,445}
\subitem \leavevmode\hypertarget{idx:0196}{Voisin--Borcea 镜像构造的关联函数},\enspace\IndexPageList[0196]{92}
\indexspace
\item \leavevmode\hypertarget{idx:1383}{管状域},\enspace\IndexPageList[1383]{153}
\indexspace
\item \leavevmode\hypertarget{idx:0371}{广义 Berglund--H\"ubsch 转置法则},\enspace\IndexPageList[0371]{93}
\indexspace
\item \leavevmode\hypertarget{idx:0373}{广义 Fredholm 轨形丛},\enspace\IndexPageList[0373]{214}
\subitem \leavevmode\hypertarget{idx:0374}{定向 Euler 类},\enspace\IndexPageList[0374]{214}
\indexspace
\item \leavevmode\hypertarget{idx:0372}{广义翻越}，\textit{参见} \hyperlink{idx:0334}{翻越，广义翻越}
\indexspace
\item \leavevmode\hypertarget{idx:0953}{归一化 $3$ 形式},\enspace\IndexPageList[0953]{44,127,128,129,133,135,174,281,282,291,351,408}，\textit{另见} \hyperlink{idx:1445}{Yukawa 耦合，归一化 Yukawa 耦合}
\indexspace
\item \leavevmode\hypertarget{idx:0954}{归一化 $d$ 形式},\enspace\IndexPageList[0954]{177}
\indexspace
\item \leavevmode\hypertarget{idx:0350}{规范固定},\enspace\IndexPageList[0350]{440}
\indexspace
\item \leavevmode\hypertarget{idx:0353}{规范理论},\enspace\IndexPageList[0353]{440,451}
\indexspace
\item \leavevmode\hypertarget{idx:0351}{规范群},\enspace\IndexPageList[0351]{440,442,443,447,452}
\subitem \leavevmode\hypertarget{idx:0352}{电磁理论的规范群},\enspace\IndexPageList[0352]{441}
\indexspace
\item \leavevmode\hypertarget{idx:0354}{规范线性 $\sigma$ 模型},\enspace\IndexPageList[0354]{35,65,68,357,451,452,453}
\subitem \leavevmode\hypertarget{idx:0355}{作用量},\enspace\IndexPageList[0355]{451}
\subitem \leavevmode\hypertarget{idx:0356}{玻色部分},\enspace\IndexPageList[0356]{451}
\subitem \leavevmode\hypertarget{idx:0357}{关联函数},\enspace\IndexPageList[0357]{387}
\subitem \leavevmode\hypertarget{idx:0358}{例题},\enspace\IndexPageList[0358]{452}
\subitem \leavevmode\hypertarget{idx:0359}{规范群},\enspace\IndexPageList[0359]{451}
\subitem \leavevmode\hypertarget{idx:0360}{镜像对称},\enspace\IndexPageList[0360]{453}
\subitem \leavevmode\hypertarget{idx:0361}{相},\enspace\IndexPageList[0361]{166}，\textit{另见} \hyperlink{idx:1042}{相}
\subsubitem \leavevmode\hypertarget{idx:0362}{Calabi--Yau 相},\enspace\IndexPageList[0362]{452}
\subsubitem \leavevmode\hypertarget{idx:0363}{Landau--Ginzburg 相},\enspace\IndexPageList[0363]{453}
\subitem \leavevmode\hypertarget{idx:0364}{与环面几何的关系},\enspace\IndexPageList[0364]{452}
\indexspace
\item \leavevmode\hypertarget{idx:0349}{规范选择},\enspace\IndexPageList[0349]{441,442,447}
\indexspace
\item \leavevmode\hypertarget{idx:0990}{轨形},\enspace\IndexPageList[0990]{35,55,61,77,194,430}
\subitem \leavevmode\hypertarget{idx:0991}{Calabi--Yau 轨形}，\textit{参见} \hyperlink{idx:0095}{Calabi--Yau，Calabi--Yau 轨形}
\subitem \leavevmode\hypertarget{idx:0992}{定义},\enspace\IndexPageList[0992]{431}
\subitem \leavevmode\hypertarget{idx:0993}{微分同胚},\enspace\IndexPageList[0993]{65,68,433}
\subitem \leavevmode\hypertarget{idx:0994}{Gorenstein 轨形},\enspace\IndexPageList[0994]{430,431,433}
\subitem \leavevmode\hypertarget{idx:0995}{轨形上的 Hodge 理论},\enspace\IndexPageList[0995]{432}
\subitem \leavevmode\hypertarget{idx:0996}{无穷小形变},\enspace\IndexPageList[0996]{433}
\subitem \leavevmode\hypertarget{idx:0997}{轨形上的积分},\enspace\IndexPageList[0997]{363}
\subitem \leavevmode\hypertarget{idx:0998}{轨形上的 Kähler 形式},\enspace\IndexPageList[0998]{432}
\subitem \leavevmode\hypertarget{idx:0999}{辛轨形},\enspace\IndexPageList[0999]{65,433}
\subitem \leavevmode\hypertarget{idx:1000}{切层},\enspace\IndexPageList[1000]{433}
\indexspace
\item \leavevmode\hypertarget{idx:0380}{鬼场},\enspace\IndexPageList[0380]{447}
\indexspace
\item \leavevmode\hypertarget{idx:0078}{黑洞},\enspace\IndexPageList[0078]{168,427}
\subitem \leavevmode\hypertarget{idx:0079}{无质量黑洞},\enspace\IndexPageList[0079]{168}
\indexspace
\item \leavevmode\hypertarget{idx:1207}{互易引理},\enspace\IndexPageList[1207]{367}
\indexspace
\item \leavevmode\hypertarget{idx:1353}{环面簇},\enspace\IndexPageList[1353]{55,405,430}
\subitem \leavevmode\hypertarget{idx:1354}{仿射环面簇},\enspace\IndexPageList[1354]{55,152}
\subitem \leavevmode\hypertarget{idx:1355}{自同构群},\enspace\IndexPageList[1355]{70,71,81,84,85,139,140,141,148}
\subsubitem \leavevmode\hypertarget{idx:1356}{维数},\enspace\IndexPageList[1356]{71}
\subsubitem \leavevmode\hypertarget{idx:1357}{根},\enspace\IndexPageList[1357]{71,147,149}
\subitem \leavevmode\hypertarget{idx:1358}{Batyrev 量子环}，\textit{参见} \hyperlink{idx:0069}{Batyrev 量子环}
\subitem \leavevmode\hypertarget{idx:1359}{Chow 群},\enspace\IndexPageList[1359]{55,59}
\subitem \leavevmode\hypertarget{idx:1360}{上同调环},\enspace\IndexPageList[1360]{246,415}
\subitem \leavevmode\hypertarget{idx:1361}{完备环面簇},\enspace\IndexPageList[1361]{55}
\subitem \leavevmode\hypertarget{idx:1362}{环面簇上的除子},\enspace\IndexPageList[1362]{56}
\subsubitem \leavevmode\hypertarget{idx:1363}{丰沛除子},\enspace\IndexPageList[1363]{56,58,63}
\subsubitem \leavevmode\hypertarget{idx:1364}{由整体截面生成的除子},\enspace\IndexPageList[1364]{56}
\subsubitem \leavevmode\hypertarget{idx:1365}{整体截面},\enspace\IndexPageList[1365]{57}
\subsubitem \leavevmode\hypertarget{idx:1366}{半丰沛除子}，\textit{参见} \hyperlink{idx:1244}{半丰沛}
\subitem \leavevmode\hypertarget{idx:1367}{例题},\enspace\IndexPageList[1367]{72,73,74,75}
\subitem \leavevmode\hypertarget{idx:1368}{Fano 环面簇},\enspace\IndexPageList[1368]{69,70,72,74,76,412,413,414,415,416,418}，\textit{另见} \hyperlink{idx:1218}{自反多胞形}
\subitem \leavevmode\hypertarget{idx:1369}{齐次坐标环},\enspace\IndexPageList[1369]{59,70,71,72,73,108,115,140,148,453}
\subitem \leavevmode\hypertarget{idx:1370}{理想与簇的对应},\enspace\IndexPageList[1370]{61}
\subitem \leavevmode\hypertarget{idx:1371}{Kähler 锥},\enspace\IndexPageList[1371]{61,62,64,415,416,452}
\subitem \leavevmode\hypertarget{idx:1372}{Mori 锥},\enspace\IndexPageList[1372]{62,63,406,412,416}
\subitem \leavevmode\hypertarget{idx:1373}{Picard 群},\enspace\IndexPageList[1373]{55,247}
\subitem \leavevmode\hypertarget{idx:1374}{量子上同调}，\textit{参见} \hyperlink{idx:1141}{量子上同调，小量子上同调，光滑环面簇的小量子上同调}
\subitem \leavevmode\hypertarget{idx:1375}{单纯环面簇},\enspace\IndexPageList[1375]{55,60,61,64,70,71,431}
\subitem \leavevmode\hypertarget{idx:1376}{光滑环面簇},\enspace\IndexPageList[1376]{55,404}
\subitem \leavevmode\hypertarget{idx:1377}{用扇描述}，\textit{参见} \hyperlink{idx:0307}{扇，扇的环面簇}
\subitem \leavevmode\hypertarget{idx:1378}{用齐次坐标描述},\enspace\IndexPageList[1378]{60}
\subitem \leavevmode\hypertarget{idx:1379}{用多胞形描述}，\textit{参见} \hyperlink{idx:1082}{多胞形，对应的环面簇}
\subitem \leavevmode\hypertarget{idx:1380}{用辛约化描述},\enspace\IndexPageList[1380]{64,65,68,453}
\indexspace
\item \leavevmode\hypertarget{idx:1350}{环面镜像定理}，\textit{参见} \hyperlink{idx:0862}{镜像定理，Givental 镜像定理，环面镜像定理}
\indexspace
\item \leavevmode\hypertarget{idx:1347}{环面完全交},\enspace\IndexPageList[1347]{29,85}
\subitem \leavevmode\hypertarget{idx:1348}{Calabi--Yau 环面完全交},\enspace\IndexPageList[1348]{29,85,405,453}
\subitem \leavevmode\hypertarget{idx:1349}{nef 环面完全交},\enspace\IndexPageList[1349]{119,120,151,404,405,412}
\indexspace
\item \leavevmode\hypertarget{idx:0347}{基本类},\enspace\IndexPageList[0347]{195}
\indexspace
\item \leavevmode\hypertarget{idx:0763}{极大射影保典范部分解消},\enspace\IndexPageList[0763]{78}
\indexspace
\item \leavevmode\hypertarget{idx:0764}{极大射影细分}，\textit{参见} \hyperlink{idx:1096}{射影细分，极大射影细分}
\indexspace
\item \leavevmode\hypertarget{idx:0285}{极端跃迁},\enspace\IndexPageList[0285]{166,167,169,428,429}
\subitem \leavevmode\hypertarget{idx:0286}{与镜像对称的相容性},\enspace\IndexPageList[0286]{167,171}
\subitem \leavevmode\hypertarget{idx:0287}{对偶跃迁},\enspace\IndexPageList[0287]{170,171,181,182,428}
\subitem \leavevmode\hypertarget{idx:0288}{例题},\enspace\IndexPageList[0288]{169,170,171,181}
\subitem \leavevmode\hypertarget{idx:0289}{镜像跃迁}，\textit{参见} \hyperlink{idx:0287}{极端跃迁，对偶跃迁}
\subitem \leavevmode\hypertarget{idx:0290}{与超共形场论模的关系},\enspace\IndexPageList[0290]{167}
\subitem \leavevmode\hypertarget{idx:0291}{环面情形},\enspace\IndexPageList[0291]{168}
\subitem \leavevmode\hypertarget{idx:0292}{I 型},\enspace\IndexPageList[0292]{166,168}
\subitem \leavevmode\hypertarget{idx:0293}{II 型},\enspace\IndexPageList[0293]{166}
\subitem \leavevmode\hypertarget{idx:0294}{III 型},\enspace\IndexPageList[0294]{166,168}
\indexspace
\item \leavevmode\hypertarget{idx:0716}{极限权滤过}，\textit{参见} \hyperlink{idx:0907}{单值性权滤过}
\indexspace
\item \leavevmode\hypertarget{idx:0784}{极小模型},\enspace\IndexPageList[0784]{29}
\indexspace
\item \leavevmode\hypertarget{idx:0379}{几何工程理论},\enspace\IndexPageList[0379]{429}
\indexspace
\item \leavevmode\hypertarget{idx:0378}{几何商},\enspace\IndexPageList[0378]{60,140}
\indexspace
\item \leavevmode\hypertarget{idx:0046}{几乎复结构},\enspace\IndexPageList[0046]{228,231}
\subitem \leavevmode\hypertarget{idx:0047}{相容的几乎复结构},\enspace\IndexPageList[0047]{208,211}
\subitem \leavevmode\hypertarget{idx:0048}{受驯服的几乎复结构},\enspace\IndexPageList[0048]{208}
\indexspace
\item \leavevmode\hypertarget{idx:1431}{加权齐次多项式},\enspace\IndexPageList[1431]{449,450}
\indexspace
\item \leavevmode\hypertarget{idx:1432}{加权射影空间},\enspace\IndexPageList[1432]{29,72,73,76,93,94,163}
\subitem \leavevmode\hypertarget{idx:1433}{成为 Fano 空间的判据},\enspace\IndexPageList[1433]{70,73}
\indexspace
\item \leavevmode\hypertarget{idx:0638}{交同调},\enspace\IndexPageList[0638]{432}
\indexspace
\item \leavevmode\hypertarget{idx:0940}{结点形变},\enspace\IndexPageList[0940]{364}
\indexspace
\item \leavevmode\hypertarget{idx:0868}{镜像变换}，\textit{参见} \hyperlink{idx:0268}{Euler 数据，相关联的 Euler 数据，镜像变换}；\hyperlink{idx:0581}{$HG$，镜像变换}
\indexspace
\item \leavevmode\hypertarget{source:485}{}\hypertarget{idx:0786}{镜像猜想},\enspace\IndexPageList[0786]{351}
\subitem \leavevmode\hypertarget{idx:0787}{Calabi--Yau 三维簇的镜像猜想},\enspace\IndexPageList[0787]{281,284}
\subitem \leavevmode\hypertarget{idx:0788}{五次三维簇的镜像猜想},\enspace\IndexPageList[0788]{285,356}
\subitem \leavevmode\hypertarget{idx:0789}{高维镜像猜想},\enspace\IndexPageList[0789]{290}
\subitem \leavevmode\hypertarget{idx:0790}{与 Givental 镜像定理的关系},\enspace\IndexPageList[0790]{406}
\subitem \leavevmode\hypertarget{idx:0791}{环面镜像猜想},\enspace\IndexPageList[0791]{135,290,291,292,293,294,295,296,406}，\textit{另见} \hyperlink{idx:0862}{镜像定理，Givental 镜像定理，环面镜像定理}
\subsubitem \leavevmode\hypertarget{idx:0792}{例题},\enspace\IndexPageList[0792]{292,293}
\subsubitem \leavevmode\hypertarget{idx:0793}{三维簇的环面镜像猜想},\enspace\IndexPageList[0793]{290}
\subsubitem \leavevmode\hypertarget{idx:0794}{高维环面镜像猜想},\enspace\IndexPageList[0794]{296}
\subsubitem \leavevmode\hypertarget{idx:0795}{Hodge 理论环面镜像猜想},\enspace\IndexPageList[0795]{79,135,291,295,296,410}
\subsubitem \leavevmode\hypertarget{idx:0796}{数学环面镜像猜想},\enspace\IndexPageList[0796]{294}
\subsubitem \leavevmode\hypertarget{idx:0797}{表述},\enspace\IndexPageList[0797]{291,295}
\indexspace
\item \leavevmode\hypertarget{idx:0844}{镜像定理},\enspace\IndexPageList[0844]{29,296,297}
\subitem \leavevmode\hypertarget{idx:0845}{更广泛的理解},\enspace\IndexPageList[0845]{378,421,426}
\subitem \leavevmode\hypertarget{idx:0846}{Givental 镜像定理},\enspace\IndexPageList[0846]{379,421}
\subsubitem \leavevmode\hypertarget{idx:0847}{$\mathcal P$ 类},\enspace\IndexPageList[0847]{401,402,403}
\subsubitem \leavevmode\hypertarget{idx:0848}{计算亏格 $0$ 引力关联数},\enspace\IndexPageList[0848]{384}
\subsubitem \leavevmode\hypertarget{idx:0849}{坐标变换引理},\enspace\IndexPageList[0849]{402,403,404}
\subsubitem \leavevmode\hypertarget{idx:0850}{双重构造引理},\enspace\IndexPageList[0850]{398,401,403}
\subsubitem \leavevmode\hypertarget{idx:0851}{例题},\enspace\IndexPageList[0851]{382,387}
\subsubitem \leavevmode\hypertarget{idx:0852}{$\mathbb F_2$ 的镜像定理},\enspace\IndexPageList[0852]{417}
\subsubitem \leavevmode\hypertarget{idx:0853}{$\mathbb P^1$ 的镜像定理},\enspace\IndexPageList[0853]{347,380}
\subsubitem \leavevmode\hypertarget{idx:0854}{$\mathbb P^n$ 的镜像定理},\enspace\IndexPageList[0854]{348,380,389,394,395}
\subsubitem \leavevmode\hypertarget{idx:0855}{$\mathbb P(1,1,2,2,2)$ 中超曲面的镜像定理},\enspace\IndexPageList[0855]{406,407,408,409,411}
\subsubitem \leavevmode\hypertarget{idx:0856}{$\mathbb P^n$ 中 nef 完全交的镜像定理},\enspace\IndexPageList[0856]{379,381,382,396,402,403}
\subsubitem \leavevmode\hypertarget{idx:0857}{亏格 $1$ 版本},\enspace\IndexPageList[0857]{404}
\subsubitem \leavevmode\hypertarget{idx:0858}{如何求变量变换},\enspace\IndexPageList[0858]{382}
\subsubitem \leavevmode\hypertarget{idx:0859}{五次三维簇的镜像定理},\enspace\IndexPageList[0859]{49,384,385,386}
\subsubitem \leavevmode\hypertarget{idx:0860}{递推引理},\enspace\IndexPageList[0860]{392,394,395,397,398,401,403}
\subsubitem \leavevmode\hypertarget{idx:0861}{正则性引理},\enspace\IndexPageList[0861]{392,396,398,401}
\subsubitem \leavevmode\hypertarget{idx:0862}{环面镜像定理},\enspace\IndexPageList[0862]{404,405,410,412,414,428}，\textit{另见} \hyperlink{idx:1141}{量子上同调，小量子上同调，光滑环面簇的小量子上同调}
\subsubitem \leavevmode\hypertarget{idx:0863}{唯一性引理},\enspace\IndexPageList[0863]{401,403}
\subitem \leavevmode\hypertarget{idx:0864}{Lian--Liu--Yau 镜像定理},\enspace\IndexPageList[0864]{51,286,355,357,372,377,403,421}
\subsubitem \leavevmode\hypertarget{idx:0865}{对多重覆盖的应用},\enspace\IndexPageList[0865]{378}
\subsubitem \leavevmode\hypertarget{idx:0866}{环面镜像定理},\enspace\IndexPageList[0866]{421}
\subitem \leavevmode\hypertarget{idx:0867}{五次三维簇的镜像定理},\enspace\IndexPageList[0867]{47,48,49,50,51,52,286,355,357,375,377,385}
\indexspace
\item \leavevmode\hypertarget{idx:0613}{镜像定理中的齐次坐标变换}，\textit{参见} \hyperlink{idx:1438}{变量的权}
\indexspace
\item \leavevmode\hypertarget{idx:0826}{镜像对},\enspace\IndexPageList[0826]{28,29,32,76,79,425,444}，\textit{另见} \hyperlink{idx:0798}{镜像构造}
\subitem \leavevmode\hypertarget{idx:0827}{Hodge 理论镜像对}，\textit{参见} \hyperlink{idx:0754}{数学镜像对，Hodge 理论镜像对}
\subitem \leavevmode\hypertarget{idx:0828}{数学镜像对}，\textit{参见} \hyperlink{idx:0748}{数学镜像对}
\subitem \leavevmode\hypertarget{idx:0829}{物理定义},\enspace\IndexPageList[0829]{27}
\indexspace
\item \leavevmode\hypertarget{idx:0830}{镜像对称},\enspace\IndexPageList[0830]{24,27,28,29,32,33,39,58,69,76,105,130,134,138,139,151,152,154,165,167,170,175,273,279,284,289,351,446,449}，\textit{另见} \hyperlink{idx:0786}{镜像猜想}；\hyperlink{idx:0844}{镜像定理}
\subitem \leavevmode\hypertarget{idx:0831}{范畴镜像对称}，\textit{参见} \hyperlink{idx:0840}{镜像对称，同调镜像对称}
\subitem \leavevmode\hypertarget{idx:0832}{经典镜像对称},\enspace\IndexPageList[0832]{32,36,47,50,356,386,408}
\subitem \leavevmode\hypertarget{idx:0833}{最早的证据},\enspace\IndexPageList[0833]{83}
\subitem \leavevmode\hypertarget{idx:0834}{枚举方面},\enspace\IndexPageList[0834]{237}
\subitem \leavevmode\hypertarget{idx:0835}{K3 曲面的镜像对称},\enspace\IndexPageList[0835]{424}
\subitem \leavevmode\hypertarget{idx:0836}{规范线性 $\sigma$ 模型的镜像对称},\enspace\IndexPageList[0836]{453}
\subitem \leavevmode\hypertarget{idx:0837}{五次三维簇的镜像对称}，\textit{参见} \hyperlink{idx:1197}{五次三维簇，镜像对称}
\subitem \leavevmode\hypertarget{idx:0838}{刚性 Calabi--Yau 簇的镜像对称},\enspace\IndexPageList[0838]{88}
\subitem \leavevmode\hypertarget{idx:0839}{Hodge 理论镜像对称},\enspace\IndexPageList[0839]{48,177,385}，\textit{另见} \hyperlink{idx:0754}{数学镜像对，Hodge 理论镜像对}
\subitem \leavevmode\hypertarget{idx:0840}{同调镜像对称},\enspace\IndexPageList[0840]{34,426}
\subitem \leavevmode\hypertarget{idx:0841}{数学镜像对称},\enspace\IndexPageList[0841]{34,135,240,265,280,297,420,425}，\textit{另见} \hyperlink{idx:0748}{数学镜像对}；\hyperlink{idx:0786}{镜像猜想}
\subitem \leavevmode\hypertarget{idx:0842}{与 Givental 镜像定理的关系},\enspace\IndexPageList[0842]{406}
\subitem \leavevmode\hypertarget{idx:0843}{环面镜像对称},\enspace\IndexPageList[0843]{79}
\indexspace
\item \leavevmode\hypertarget{idx:0798}{镜像构造},\enspace\IndexPageList[0798]{284}
\subitem \leavevmode\hypertarget{idx:0799}{Batyrev 构造},\enspace\IndexPageList[0799]{76,79,83,93,95,143,146,147,148,158,162,166,167,169,171,178,284,290,291,294,295,404,406}
\subsubitem \leavevmode\hypertarget{idx:0800}{五次三维簇},\enspace\IndexPageList[0800]{84}
\subitem \leavevmode\hypertarget{idx:0801}{Batyrev--Borisov 构造},\enspace\IndexPageList[0801]{85,86,123,170,428}，\textit{另见} \hyperlink{idx:0935}{nef 分拆}
\subitem \leavevmode\hypertarget{idx:0802}{旗流形中的 Calabi--Yau 完全交的镜像构造},\enspace\IndexPageList[0802]{428}
\subitem \leavevmode\hypertarget{idx:0803}{Greene--Plesser 构造},\enspace\IndexPageList[0803]{76,79,112}
\subitem \leavevmode\hypertarget{idx:0804}{Voisin--Borcea 构造},\enspace\IndexPageList[0804]{88,91,92,95,424,425}
\subsubitem \leavevmode\hypertarget{idx:0805}{复模与 Kähler 模},\enspace\IndexPageList[0805]{91,92}
\subsubitem \leavevmode\hypertarget{idx:0806}{与 Batyrev 构造的关系},\enspace\IndexPageList[0806]{93}
\subsubitem \leavevmode\hypertarget{idx:0807}{Yukawa 耦合},\enspace\IndexPageList[0807]{92}
\indexspace
\item \leavevmode\hypertarget{idx:0809}{镜像映射},\enspace\IndexPageList[0809]{28,32,79,83,105,132,151,166,170,171,172,173,174,175,176,177,178,179,180,181,182,183,184,185,186,187,188,189,280,408}
\subitem \leavevmode\hypertarget{idx:0810}{定义},\enspace\IndexPageList[0810]{173,180}
\subsubitem \leavevmode\hypertarget{idx:0811}{需要作出的选择},\enspace\IndexPageList[0811]{173,175,177}
\subsubitem \leavevmode\hypertarget{idx:0812}{定义所需的数据},\enspace\IndexPageList[0812]{173}
\subsubitem \leavevmode\hypertarget{idx:0813}{定义的不确定性},\enspace\IndexPageList[0813]{173}
\subsubitem \leavevmode\hypertarget{idx:0814}{物理依据},\enspace\IndexPageList[0814]{181}
\subsubitem \leavevmode\hypertarget{idx:0815}{整性猜想的作用},\enspace\IndexPageList[0815]{173,174}
\subitem \leavevmode\hypertarget{idx:0816}{导数},\enspace\IndexPageList[0816]{175,176,177,181,289,290}，\textit{另见} \hyperlink{idx:0811}{镜像映射，定义，需要作出的选择}
\subitem \leavevmode\hypertarget{idx:0817}{五次三维簇的镜像映射}，\textit{参见} \hyperlink{idx:1195}{五次三维簇，镜像映射}
\subitem \leavevmode\hypertarget{idx:0818}{镜像坐标},\enspace\IndexPageList[0818]{133,135,175,282,288}
\subsubitem \leavevmode\hypertarget{idx:0819}{定义},\enspace\IndexPageList[0819]{172}
\subitem \leavevmode\hypertarget{idx:0820}{环面镜像映射},\enspace\IndexPageList[0820]{171,174,178,179,180,181,182,290,291,293,295}
\subsubitem \leavevmode\hypertarget{idx:0821}{定义},\enspace\IndexPageList[0821]{178,179}
\subsubitem \leavevmode\hypertarget{idx:0822}{五次三维簇的环面镜像映射},\enspace\IndexPageList[0822]{179,180}
\subsubitem \leavevmode\hypertarget{idx:0823}{不为恒等映射},\enspace\IndexPageList[0823]{171}
\subsubitem \leavevmode\hypertarget{idx:0824}{消除非唯一性},\enspace\IndexPageList[0824]{179}
\subitem \leavevmode\hypertarget{idx:0825}{通过超几何函数构造},\enspace\IndexPageList[0825]{182,183,184,185,186,187,188,189}
\indexspace
\item \leavevmode\hypertarget{idx:0808}{镜像坐标}，\textit{参见} \hyperlink{idx:0809}{镜像映射}
\indexspace
\item \leavevmode\hypertarget{idx:0723}{局部化},\enspace\IndexPageList[0723]{302,303,389,405}，\textit{另见} \hyperlink{idx:0734}{局部化定理}
\subitem \leavevmode\hypertarget{idx:0724}{叠的局部化},\enspace\IndexPageList[0724]{312}
\indexspace
\item \leavevmode\hypertarget{idx:0734}{局部化定理}
\subitem \leavevmode\hypertarget{idx:0735}{在枚举几何中的应用},\enspace\IndexPageList[0735]{305}
\subitem \leavevmode\hypertarget{idx:0736}{应用于 $\mathbb P^n$},\enspace\IndexPageList[0736]{301,302}
\subitem \leavevmode\hypertarget{idx:0737}{计算瞬子数 $n_4$},\enspace\IndexPageList[0737]{311}
\subitem \leavevmode\hypertarget{idx:0738}{计算五次三维簇上的直线},\enspace\IndexPageList[0738]{305}
\subitem \leavevmode\hypertarget{idx:0739}{等变上同调的局部化定理},\enspace\IndexPageList[0739]{300,342,361,362,370,371,372}
\subitem \leavevmode\hypertarget{idx:0740}{Chern 类多项式的局部化定理},\enspace\IndexPageList[0740]{304,314,315,316}
\subitem \leavevmode\hypertarget{idx:0741}{光滑叠的局部化定理},\enspace\IndexPageList[0741]{303,304,312,391,396,398,400}
\subitem \leavevmode\hypertarget{idx:0742}{等变积分的局部化定理},\enspace\IndexPageList[0742]{302,358,362,363,364,365,368,371,399}
\subitem \leavevmode\hypertarget{idx:0743}{证明多重覆盖公式},\enspace\IndexPageList[0743]{312}
\indexspace
\item \leavevmode\hypertarget{idx:0722}{局部算子},\enspace\IndexPageList[0722]{444,446}
\indexspace
\item \leavevmode\hypertarget{idx:0898}{矩映射},\enspace\IndexPageList[0898]{64,67,68,72,166,452}
\subitem \leavevmode\hypertarget{idx:0899}{例题},\enspace\IndexPageList[0899]{73}
\indexspace
\item \leavevmode\hypertarget{idx:0760}{矩阵积分},\enspace\IndexPageList[0760]{333}
\indexspace
\item \leavevmode\hypertarget{idx:0761}{矩阵模型},\enspace\IndexPageList[0761]{333}
\indexspace
\item \leavevmode\hypertarget{idx:0913}{可动锥},\enspace\IndexPageList[0913]{156,165}
\subitem \leavevmode\hypertarget{idx:0914}{反射可动锥},\enspace\IndexPageList[0914]{157,165}
\indexspace
\item \leavevmode\hypertarget{idx:0956}{可观测量},\enspace\IndexPageList[0956]{24,25,59,437}
\indexspace
\item \leavevmode\hypertarget{idx:1035}{裤形曲面},\enspace\IndexPageList[1035]{456,457}
\indexspace
\item \leavevmode\hypertarget{idx:0376}{亏格 $g$ 的引力 Gromov--Witten 势},\enspace\IndexPageList[0376]{326}
\subitem \leavevmode\hypertarget{idx:0377}{定义},\enspace\IndexPageList[0377]{327}
\indexspace
\item \leavevmode\hypertarget{idx:0375}{亏格 $g$ 耦合},\enspace\IndexPageList[0375]{326,335}
\indexspace
\item \leavevmode\hypertarget{idx:0172}{联络 $\nabla^c$},\enspace\IndexPageList[0172]{270}，\textit{另见} \hyperlink{idx:0003}{A 模型，联络}
\subitem \leavevmode\hypertarget{idx:0173}{平坦截面},\enspace\IndexPageList[0173]{270,271,272,277,339,341}
\indexspace
\item \leavevmode\hypertarget{idx:0174}{联络矩阵},\enspace\IndexPageList[0174]{48,50,99,102,129,131,269}
\indexspace
\item \leavevmode\hypertarget{idx:1109}{量子 $\mathcal D$ 模},\enspace\IndexPageList[1109]{428}
\indexspace
\item \leavevmode\hypertarget{idx:1155}{量子场论},\enspace\IndexPageList[1155]{436,437,438,439,440,441,442,445,453}
\subitem \leavevmode\hypertarget{idx:1156}{Hamilton 表述},\enspace\IndexPageList[1156]{439}
\subitem \leavevmode\hypertarget{idx:1157}{Lagrange 表述},\enspace\IndexPageList[1157]{439}
\subitem \leavevmode\hypertarget{idx:1158}{二维量子场论},\enspace\IndexPageList[1158]{24}
\indexspace
\item \leavevmode\hypertarget{idx:1159}{量子超平面截面原理},\enspace\IndexPageList[1159]{49,387,388,389,428}
\subitem \leavevmode\hypertarget{idx:1160}{$\mathbb P^n$ 的量子超平面截面原理},\enspace\IndexPageList[1160]{388}
\subitem \leavevmode\hypertarget{idx:1161}{凸簇的量子超平面截面原理},\enspace\IndexPageList[1161]{388}
\subitem \leavevmode\hypertarget{idx:1162}{与 Givental 镜像定理的关系},\enspace\IndexPageList[1162]{389}
\indexspace
\item \leavevmode\hypertarget{idx:1110}{量子反常},\enspace\IndexPageList[1110]{446,447,448}
\indexspace
\item \leavevmode\hypertarget{idx:1106}{量子化},\enspace\IndexPageList[1106]{33,436,438,441,442,447,448,453}
\subitem \leavevmode\hypertarget{idx:1107}{正则量子化},\enspace\IndexPageList[1107]{436,437}
\subitem \leavevmode\hypertarget{idx:1108}{径向量子化},\enspace\IndexPageList[1108]{447}
\indexspace
\item \leavevmode\hypertarget{idx:1111}{量子上同调},\enspace\IndexPageList[1111]{30,48,105,124,223,445,453,458}
\subitem \leavevmode\hypertarget{idx:1112}{大量子上同调},\enspace\IndexPageList[1112]{252,253,327,335,337}，\textit{另见} \hyperlink{idx:0552}{Gromov--Witten 势}
\subsubitem \leavevmode\hypertarget{idx:1113}{结合性},\enspace\IndexPageList[1113]{255,257,260}，\textit{另见} \hyperlink{idx:0568}{Gromov--Witten 势，WDVV 方程}
\subsubitem \leavevmode\hypertarget{idx:1114}{收敛性},\enspace\IndexPageList[1114]{253}
\subsubitem \leavevmode\hypertarget{idx:1115}{定义},\enspace\IndexPageList[1115]{254}
\subsubitem \leavevmode\hypertarget{idx:1116}{次数条件},\enspace\IndexPageList[1116]{258}
\subsubitem \leavevmode\hypertarget{idx:1117}{K\"unneth 公式},\enspace\IndexPageList[1117]{257}
\subsubitem \leavevmode\hypertarget{idx:1118}{$\mathbb P^2$ 的大量子上同调},\enspace\IndexPageList[1118]{260}
\subsubitem \leavevmode\hypertarget{idx:1119}{Calabi--Yau 三维簇的大量子上同调},\enspace\IndexPageList[1119]{262}
\subsubitem \leavevmode\hypertarget{idx:1120}{五次三维簇的大量子上同调},\enspace\IndexPageList[1120]{261}
\subsubitem \leavevmode\hypertarget{idx:1121}{与小量子乘积的关系},\enspace\IndexPageList[1121]{265,266}
\subsubitem \leavevmode\hypertarget{idx:1122}{限制},\enspace\IndexPageList[1122]{265,266,267,337,343,350}
\subsubitem \leavevmode\hypertarget{idx:1123}{环结构},\enspace\IndexPageList[1123]{256}，\textit{另见} \hyperlink{idx:0337}{Frobenius 代数}
\subsubitem \leavevmode\hypertarget{idx:1124}{用于定义 Givental 联络},\enspace\IndexPageList[1124]{334}
\subitem \leavevmode\hypertarget{idx:1125}{系数},\enspace\IndexPageList[1125]{251,252,253}
\subitem \leavevmode\hypertarget{idx:1126}{引力量子上同调}，\textit{参见} \hyperlink{idx:0481}{引力量子乘积}
\subitem \leavevmode\hypertarget{idx:1127}{修正量子上同调},\enspace\IndexPageList[1127]{386}
\subsubitem \leavevmode\hypertarget{idx:1128}{其中的关系},\enspace\IndexPageList[1128]{386}
\subitem \leavevmode\hypertarget{idx:1129}{小量子上同调},\enspace\IndexPageList[1129]{240,241,242,243,244,245,246,247,248,249,250,251,253,337,338,343,349,382,386,387}
\subsubitem \leavevmode\hypertarget{idx:1130}{结合性},\enspace\IndexPageList[1130]{242,251}
\subsubitem \leavevmode\hypertarget{idx:1131}{收敛性},\enspace\IndexPageList[1131]{241,250,251,267,268,423}
\subsubitem \leavevmode\hypertarget{idx:1132}{定义},\enspace\IndexPageList[1132]{241}
\subsubitem \leavevmode\hypertarget{idx:1133}{次数条件},\enspace\IndexPageList[1133]{244,245}
\subsubitem \leavevmode\hypertarget{idx:1134}{翻越的影响},\enspace\IndexPageList[1134]{249,250}
\subsubitem \leavevmode\hypertarget{idx:1135}{$(1+1)$ 维拓扑量子场论的小量子上同调},\enspace\IndexPageList[1135]{457}
\subsubitem \leavevmode\hypertarget{idx:1136}{$\mathbb F_2$ 的小量子上同调},\enspace\IndexPageList[1136]{416,417,418}
\subsubitem \leavevmode\hypertarget{idx:1137}{$\mathbb P^n$ 的小量子上同调},\enspace\IndexPageList[1137]{245,246,260,348}
\subsubitem \leavevmode\hypertarget{idx:1138}{Calabi--Yau 流形的小量子上同调},\enspace\IndexPageList[1138]{241,245,275}
\subsubitem \leavevmode\hypertarget{idx:1139}{Calabi--Yau 三维簇的小量子上同调},\enspace\IndexPageList[1139]{248,249,250,262,350,352}
\subsubitem \leavevmode\hypertarget{idx:1140}{Grassmann 簇的小量子上同调},\enspace\IndexPageList[1140]{250,428}
\subsubitem \leavevmode\hypertarget{idx:1141}{光滑环面簇的小量子上同调},\enspace\IndexPageList[1141]{246,413,415,416}，\textit{另见} \hyperlink{idx:0069}{Batyrev 量子环}
\subsubitem \leavevmode\hypertarget{idx:1142}{五次三维簇的小量子上同调},\enspace\IndexPageList[1142]{248,249,261}
\subsubitem \leavevmode\hypertarget{idx:1143}{其他例题},\enspace\IndexPageList[1143]{250,251}
\subsubitem \leavevmode\hypertarget{idx:1144}{性质},\enspace\IndexPageList[1144]{244}
\subsubitem \leavevmode\hypertarget{idx:1145}{与大量子乘积的关系},\enspace\IndexPageList[1145]{265,266}
\subsubitem \leavevmode\hypertarget{idx:1146}{其中的关系},\enspace\IndexPageList[1146]{344,345,348,387,410,413,415}
\subsubitem \leavevmode\hypertarget{idx:1147}{环结构},\enspace\IndexPageList[1147]{242}，\textit{另见} \hyperlink{idx:0337}{Frobenius 代数}
\subsubitem \leavevmode\hypertarget{idx:1148}{用于定义 A 模型联络},\enspace\IndexPageList[1148]{266,267}
\indexspace
\item \leavevmode\hypertarget{idx:1150}{量子微分方程},\enspace\IndexPageList[1150]{344,347,350,352,379,409,413,428}，\textit{另见} \hyperlink{idx:1151}{量子微分算子}
\indexspace
\item \leavevmode\hypertarget{idx:1151}{量子微分算子},\enspace\IndexPageList[1151]{344,345,352,415,418}
\subitem \leavevmode\hypertarget{idx:1152}{由量子上同调中的关系导出},\enspace\IndexPageList[1152]{345,347}
\subitem \leavevmode\hypertarget{idx:1153}{例题},\enspace\IndexPageList[1153]{347,348,350}
\subitem \leavevmode\hypertarget{idx:1154}{给出量子上同调中的关系},\enspace\IndexPageList[1154]{344}，\textit{另见} \hyperlink{idx:1146}{量子上同调，小量子上同调，其中的关系}
\indexspace
\item \leavevmode\hypertarget{idx:1149}{量子修正},\enspace\IndexPageList[1149]{35,41,171}
\indexspace
\item \leavevmode\hypertarget{idx:0205}{临界丛}，\textit{参见} \hyperlink{idx:0147}{凸凹丛，临界丛}
\indexspace
\item \leavevmode\hypertarget{idx:1231}{留数矩阵 $\operatorname{Res}(\nabla)$},\enspace\IndexPageList[1231]{100,102,269}
\indexspace
\item \leavevmode\hypertarget{idx:1230}{留数映射},\enspace\IndexPageList[1230]{106,109}
\indexspace
\item \leavevmode\hypertarget{idx:0218}{六结点五次曲线的二重覆盖},\enspace\IndexPageList[0218]{316,317,318,319,320}
\subitem \leavevmode\hypertarget{idx:0219}{对 Gromov--Witten 不变量及虚拟基本类的贡献},\enspace\IndexPageList[0219]{320}
\subitem \leavevmode\hypertarget{idx:0220}{经正规化分解的覆盖},\enspace\IndexPageList[0220]{318,319,320}
\subitem \leavevmode\hypertarget{idx:0221}{不经正规化分解的覆盖},\enspace\IndexPageList[0221]{318,319}
\subitem \leavevmode\hypertarget{idx:0222}{模},\enspace\IndexPageList[0222]{318}
\subitem \leavevmode\hypertarget{idx:0223}{切空间},\enspace\IndexPageList[0223]{319}
\indexspace
\item \leavevmode\hypertarget{idx:1037}{路径积分},\enspace\IndexPageList[1037]{439,440,441,445}，\textit{另见} \hyperlink{idx:0315}{Feynman 积分}
\indexspace
\item \leavevmode\hypertarget{idx:0241}{枚举几何},\enspace\IndexPageList[0241]{39,238}，\textit{另见} \hyperlink{idx:0534}{Gromov--Witten 不变量，枚举意义}
\indexspace
\item \leavevmode\hypertarget{idx:0938}{幂零轨道},\enspace\IndexPageList[0938]{277}
\indexspace
\item \leavevmode\hypertarget{idx:0939}{幂零轨道定理},\enspace\IndexPageList[0939]{99}
\indexspace
\item \leavevmode\hypertarget{idx:0870}{模},\enspace\IndexPageList[0870]{24,28}
\subitem \leavevmode\hypertarget{idx:0871}{复模},\enspace\IndexPageList[0871]{26,27,28,32,79,104,105,136,137,138,139,140,141,142,143,144,145,146,147,148,149,150,151,154,167,171,172,287,446}
\subsubitem \leavevmode\hypertarget{idx:0872}{一维复模},\enspace\IndexPageList[0872]{101,102,105,107,110,125,128,130,132,280}
\subsubitem \leavevmode\hypertarget{idx:0873}{紧化},\enspace\IndexPageList[0873]{137,138,140,142,143,144}
\subsubitem \leavevmode\hypertarget{idx:0874}{坐标},\enspace\IndexPageList[0874]{137,146,147,178,182,290}
\subsubitem \leavevmode\hypertarget{idx:0875}{维数},\enspace\IndexPageList[0875]{141,142}
\subsubitem \leavevmode\hypertarget{idx:0876}{例题},\enspace\IndexPageList[0876]{146,147,148,149,150,162,170}
\subsubitem \leavevmode\hypertarget{idx:0877}{最大幂幺边界点}，\textit{参见} \hyperlink{idx:0766}{最大幂幺，最大幂幺边界点}
\subsubitem \leavevmode\hypertarget{idx:0878}{Calabi--Yau 三维簇的复模},\enspace\IndexPageList[0878]{101,102,105,125,128,130,132,144,280,281,283}
\subsubitem \leavevmode\hypertarget{idx:0879}{Calabi--Yau 环面超曲面的复模},\enspace\IndexPageList[0879]{115,116,139,140,141,142,143,144,145,146,147,148,149,150}
\subsubitem \leavevmode\hypertarget{idx:0880}{多项式复模},\enspace\IndexPageList[0880]{84,139,140,141,142,147,148,149,161,170,406,450}
\subsubitem \leavevmode\hypertarget{idx:0881}{简化复模},\enspace\IndexPageList[0881]{118,140,141,142,143,144,145,146,147,148,149,150,158,161,165,178,179,290,294}
\subitem \leavevmode\hypertarget{idx:0882}{Kähler 模},\enspace\IndexPageList[0882]{26,27,28,34,35,79,105,139,150,151,152,153,154,155,156,157,158,159,160,161,162,163,164,165,166,167,168,169,170,171,172,175,241,268,287,446}
\subsubitem \leavevmode\hypertarget{idx:0883}{一维 Kähler 模},\enspace\IndexPageList[0883]{172}
\subsubitem \leavevmode\hypertarget{idx:0884}{边界},\enspace\IndexPageList[0884]{166}
\subsubitem \leavevmode\hypertarget{idx:0885}{紧化},\enspace\IndexPageList[0885]{155,158,159,165,166,167,175}
\subsubitem \leavevmode\hypertarget{idx:0886}{坐标},\enspace\IndexPageList[0886]{151,152,290}
\subsubitem \leavevmode\hypertarget{idx:0887}{扩大 Kähler 模},\enspace\IndexPageList[0887]{155,156,157,159,165}
\subsubitem \leavevmode\hypertarget{idx:0888}{例题},\enspace\IndexPageList[0888]{162,163,164,165,169}
\subsubitem \leavevmode\hypertarget{idx:0889}{全局 Kähler 模},\enspace\IndexPageList[0889]{165}
\subsubitem \leavevmode\hypertarget{idx:0890}{大半径极限点}，\textit{参见} \hyperlink{idx:0696}{大半径极限点}
\subsubitem \leavevmode\hypertarget{idx:0891}{Calabi--Yau 三维簇的 Kähler 模},\enspace\IndexPageList[0891]{155,156,157,165,280,281}
\subsubitem \leavevmode\hypertarget{idx:0892}{Calabi--Yau 环面超曲面的 Kähler 模},\enspace\IndexPageList[0892]{157,158,159,160,161,162,163,164,165,166}
\subsubitem \leavevmode\hypertarget{idx:0893}{部分紧化},\enspace\IndexPageList[0893]{151,152,153,154,155,157,161,164,165,173}
\subsubitem \leavevmode\hypertarget{idx:0894}{环面 Kähler 模},\enspace\IndexPageList[0894]{83,143,157,158,161,164,178,179,290,294}
\subitem \leavevmode\hypertarget{idx:0895}{超共形场论模},\enspace\IndexPageList[0895]{26,27,28,151,167,168,287,424}
\subsubitem \leavevmode\hypertarget{idx:0896}{例题},\enspace\IndexPageList[0896]{170}
\subitem \leavevmode\hypertarget{idx:0897}{稳定映射的模}，\textit{参见} \hyperlink{idx:1276}{稳定映射，模空间}
\indexspace
\item \leavevmode\hypertarget{idx:1249}{$\sigma$ 模型},\enspace\IndexPageList[1249]{26,27,29,30,32,33,34,38,191,325,357,442,443,444,445,446,451,452,453}
\subitem \leavevmode\hypertarget{idx:1250}{A 模型},\enspace\IndexPageList[1250]{444}
\subitem \leavevmode\hypertarget{idx:1251}{作用量}，\textit{参见} \hyperlink{idx:0034}{作用量，$\sigma$ 模型的作用量}
\subitem \leavevmode\hypertarget{idx:1252}{B 模型},\enspace\IndexPageList[1252]{444}
\subitem \leavevmode\hypertarget{idx:1253}{费米场},\enspace\IndexPageList[1253]{442,444,448}
\subitem \leavevmode\hypertarget{idx:1254}{超对称 $\sigma$ 模型},\enspace\IndexPageList[1254]{444,446}
\subitem \leavevmode\hypertarget{idx:1255}{拓扑 $\sigma$ 模型},\enspace\IndexPageList[1255]{457}
\subitem \leavevmode\hypertarget{idx:1256}{扭曲 $\sigma$ 模型},\enspace\IndexPageList[1256]{444,445}
\indexspace
\item \leavevmode\hypertarget{idx:0640}{内蕴法层},\enspace\IndexPageList[0640]{201,202}，\textit{另见} \hyperlink{idx:0948}{法锥，内蕴法锥}
\indexspace
\item \leavevmode\hypertarget{idx:0639}{内蕴法锥}，\textit{参见} \hyperlink{idx:0948}{法锥，内蕴法锥}
\indexspace
\item \leavevmode\hypertarget{idx:0240}{能动张量},\enspace\IndexPageList[0240]{447}
\indexspace
\item \leavevmode\hypertarget{idx:0239}{能动量守恒},\enspace\IndexPageList[0239]{447}
\indexspace
\item \leavevmode\hypertarget{idx:1163}{拟光滑},\enspace\IndexPageList[1163]{83,109,139,431}，\textit{另见} \hyperlink{idx:0990}{轨形}
\indexspace
\item \leavevmode\hypertarget{idx:0214}{判别式轨迹},\enspace\IndexPageList[0214]{127,145}，\textit{另见} \hyperlink{idx:0103}{Calabi--Yau，Calabi--Yau 环面超曲面，非退化情形}
\subitem \leavevmode\hypertarget{idx:0215}{例题},\enspace\IndexPageList[0215]{148,170,181,182}
\subitem \leavevmode\hypertarget{idx:0216}{主分量},\enspace\IndexPageList[0216]{145,148}
\indexspace
\item \leavevmode\hypertarget{idx:1036}{配分函数},\enspace\IndexPageList[1036]{333}
\indexspace
\item \leavevmode\hypertarget{idx:1257}{奇点},\enspace\IndexPageList[1257]{430,431,432,433}
\subitem \leavevmode\hypertarget{idx:1258}{$\mathbb Q$ 析因奇点},\enspace\IndexPageList[1258]{34,78,430}
\subitem \leavevmode\hypertarget{idx:1259}{典范奇点},\enspace\IndexPageList[1259]{34,35,430,431}
\subitem \leavevmode\hypertarget{idx:1260}{Cohen--Macaulay 奇点},\enspace\IndexPageList[1260]{430}
\subitem \leavevmode\hypertarget{idx:1261}{Gorenstein 奇点},\enspace\IndexPageList[1261]{34,35,78,139,430}
\subitem \leavevmode\hypertarget{idx:1262}{终端奇点},\enspace\IndexPageList[1262]{34,77,78,139,430,431,433}
\subitem \leavevmode\hypertarget{idx:1263}{环面型奇点},\enspace\IndexPageList[1263]{433}
\indexspace
\item \leavevmode\hypertarget{idx:0324}{旗流形},\enspace\IndexPageList[0324]{242,251,277}
\subitem \leavevmode\hypertarget{idx:0325}{旗流形中的完全交},\enspace\IndexPageList[0325]{387,389}
\subsubitem \leavevmode\hypertarget{idx:0326}{Calabi--Yau 完全交},\enspace\IndexPageList[0326]{428}
\indexspace
\item \leavevmode\hypertarget{idx:0577}{强 Lefschetz 定理},\enspace\IndexPageList[0577]{273,432}
\indexspace
\item \leavevmode\hypertarget{idx:1326}{切函子},\enspace\IndexPageList[1326]{199}
\indexspace
\item \leavevmode\hypertarget{idx:1327}{切权}，\textit{参见} \hyperlink{idx:1437}{表示的权，切权}
\indexspace
\item \leavevmode\hypertarget{idx:1328}{切—障碍复形},\enspace\IndexPageList[1328]{199,201,202,203,308}
\subitem \leavevmode\hypertarget{idx:1329}{用于虚拟基本类},\enspace\IndexPageList[1329]{201}
\subitem \leavevmode\hypertarget{idx:1330}{完美切—障碍复形},\enspace\IndexPageList[1330]{197,204}
\subsubitem \leavevmode\hypertarget{idx:1331}{定义},\enspace\IndexPageList[1331]{199}
\subsubitem \leavevmode\hypertarget{idx:1332}{例题},\enspace\IndexPageList[1332]{200}
\indexspace
\item \leavevmode\hypertarget{idx:0206}{曲率张量},\enspace\IndexPageList[0206]{443}
\indexspace
\item \leavevmode\hypertarget{idx:0612}{全纯辛},\enspace\IndexPageList[0612]{425}
\indexspace
\item \leavevmode\hypertarget{idx:0869}{缺失的镜像},\enspace\IndexPageList[0869]{95}
\indexspace
\item \leavevmode\hypertarget{idx:1333}{三点函数}，\textit{参见} \hyperlink{idx:0195}{关联函数，三点关联函数}
\indexspace
\item \leavevmode\hypertarget{idx:0297}{扇},\enspace\IndexPageList[0297]{54}，\textit{另见} \hyperlink{idx:1094}{射影细分}
\subitem \leavevmode\hypertarget{idx:0298}{完备扇},\enspace\IndexPageList[0298]{55}
\subitem \leavevmode\hypertarget{idx:0299}{扇的翻转}，\textit{参见} \hyperlink{idx:0328}{翻转}
\subitem \leavevmode\hypertarget{idx:0300}{Gröbner 扇}，\textit{参见} \hyperlink{idx:0453}{Gr\"obner 扇}
\subitem \leavevmode\hypertarget{idx:0301}{次级扇}，\textit{参见} \hyperlink{idx:1239}{次级扇}
\subitem \leavevmode\hypertarget{idx:0302}{单纯扇},\enspace\IndexPageList[0302]{55}
\subitem \leavevmode\hypertarget{idx:0303}{光滑扇},\enspace\IndexPageList[0303]{55}
\subitem \leavevmode\hypertarget{idx:0304}{支集},\enspace\IndexPageList[0304]{54}
\subitem \leavevmode\hypertarget{idx:0305}{支撑于线性回路上的扇}，\textit{参见} \hyperlink{idx:0720}{线性回路，由扇支撑}
\subitem \leavevmode\hypertarget{idx:0306}{对称性},\enspace\IndexPageList[0306]{71}
\subitem \leavevmode\hypertarget{idx:0307}{扇的环面簇},\enspace\IndexPageList[0307]{55,74}
\indexspace
\item \leavevmode\hypertarget{idx:1317}{扇的支集}，\textit{参见} \hyperlink{idx:0304}{扇，支集}
\indexspace
\item \leavevmode\hypertarget{idx:0126}{上同调场论},\enspace\IndexPageList[0126]{264}，\textit{另见} \hyperlink{idx:0235}{Dubrovin 形式体系}
\subitem \leavevmode\hypertarget{idx:0127}{关联函数},\enspace\IndexPageList[0127]{264}
\subitem \leavevmode\hypertarget{idx:0128}{高亏格情形},\enspace\IndexPageList[0128]{264}
\subitem \leavevmode\hypertarget{idx:0129}{Künneth 公式},\enspace\IndexPageList[0129]{264}
\indexspace
\item \leavevmode\hypertarget{idx:1351}{上同调的环面部分},\enspace\IndexPageList[1351]{157,166,179,181,290,384,405}
\subitem \leavevmode\hypertarget{idx:1352}{定义},\enspace\IndexPageList[1352]{80,294}
\indexspace
\item \leavevmode\hypertarget{idx:0967}{尚未解决的猜想},\enspace\IndexPageList[0967]{422,423}，\textit{另见} \hyperlink{idx:0786}{镜像猜想}
\subitem \leavevmode\hypertarget{idx:0968}{Calabi--Yau 三维簇}
\subsubitem \leavevmode\hypertarget{idx:0969}{用 $I_{\mathcal V}$ 表示 B 模型势函数},\enspace\IndexPageList[0969]{412,422}
\subsubitem \leavevmode\hypertarget{idx:0970}{用 $I_{\mathcal V}$ 表示 Yukawa 耦合},\enspace\IndexPageList[0970]{412,422}
\subitem \leavevmode\hypertarget{idx:0971}{复几何}
\subsubitem \leavevmode\hypertarget{idx:0972}{整性猜想},\enspace\IndexPageList[0972]{105,422}
\subsubitem \leavevmode\hypertarget{idx:0973}{简化模空间中的最大幂幺边界点},\enspace\IndexPageList[0973]{143,422}
\subitem \leavevmode\hypertarget{idx:0974}{量子上同调的收敛性},\enspace\IndexPageList[0974]{241,423}
\subitem \leavevmode\hypertarget{idx:0975}{枚举几何}
\subsubitem \leavevmode\hypertarget{idx:0976}{Clemens 猜想},\enspace\IndexPageList[0976]{225,423}
\subsubitem \leavevmode\hypertarget{idx:0977}{Gromov--Witten 不变量的枚举意义},\enspace\IndexPageList[0977]{219,423}
\subsubitem \leavevmode\hypertarget{idx:0978}{瞬子数的辛定义},\enspace\IndexPageList[0978]{231,423}
\subitem \leavevmode\hypertarget{idx:0979}{Kähler 几何}
\subsubitem \leavevmode\hypertarget{idx:0980}{锥猜想},\enspace\IndexPageList[0980]{153,422}
\subsubitem \leavevmode\hypertarget{idx:0981}{Calabi--Yau 环面超曲面的 Kähler 锥的环面部分},\enspace\IndexPageList[0981]{161,422}
\subitem \leavevmode\hypertarget{idx:0982}{Picard--Fuchs 方程}
\subsubitem \leavevmode\hypertarget{idx:0983}{$I_{\mathcal V}$ 是解},\enspace\IndexPageList[0983]{124,422}
\subsubitem \leavevmode\hypertarget{idx:0984}{$\mathcal A$ 方程组中缺失的方程},\enspace\IndexPageList[0984]{117,422}
\subitem \leavevmode\hypertarget{idx:0985}{$J_X$ 与 $J_{\mathcal V}$ 的关系},\enspace\IndexPageList[0985]{409,423}
\subitem \leavevmode\hypertarget{idx:0986}{虚拟基本类}
\subsubitem \leavevmode\hypertarget{idx:0987}{定义的等价性},\enspace\IndexPageList[0987]{199,422}
\subsubitem \leavevmode\hypertarget{idx:0988}{完全交的虚拟基本类},\enspace\IndexPageList[0988]{409,423}
\indexspace
\item \leavevmode\hypertarget{idx:1094}{射影细分},\enspace\IndexPageList[1094]{77,78,160}
\subitem \leavevmode\hypertarget{idx:1095}{定义},\enspace\IndexPageList[1095]{76}
\subitem \leavevmode\hypertarget{idx:1096}{极大射影细分},\enspace\IndexPageList[1096]{77,78,79,83,85,86,139,158,290,405}
\subsubitem \leavevmode\hypertarget{idx:1097}{定义},\enspace\IndexPageList[1097]{76}
\subitem \leavevmode\hypertarget{idx:1098}{简化射影细分},\enspace\IndexPageList[1098]{159,162,163,164,166,169}
\subsubitem \leavevmode\hypertarget{idx:1099}{定义},\enspace\IndexPageList[1099]{158}
\indexspace
\item \leavevmode\hypertarget{idx:1204}{升指标},\enspace\IndexPageList[1204]{241,262,440}
\indexspace
\item \leavevmode\hypertarget{idx:1440}{世界面},\enspace\IndexPageList[1440]{24,25,33,442,445,446,447,448,449,450}
\subitem \leavevmode\hypertarget{idx:1441}{非微扰修正},\enspace\IndexPageList[1441]{30}
\indexspace
\item \leavevmode\hypertarget{idx:1442}{世界体},\enspace\IndexPageList[1442]{425}
\indexspace
\item \leavevmode\hypertarget{idx:1083}{势函数}，\textit{参见} \hyperlink{idx:0064}{B 模型，势函数}；\hyperlink{idx:0552}{Gromov--Witten 势}
\indexspace
\item \leavevmode\hypertarget{idx:0121}{手征超场},\enspace\IndexPageList[0121]{449,450,451}
\indexspace
\item \leavevmode\hypertarget{idx:0120}{手征环},\enspace\IndexPageList[0120]{59,240,450}
\indexspace
\item \leavevmode\hypertarget{idx:0175}{守恒流},\enspace\IndexPageList[0175]{442}
\indexspace
\item \leavevmode\hypertarget{idx:0748}{数学镜像对},\enspace\IndexPageList[0748]{27,34,48,79,265,410,420}
\subitem \leavevmode\hypertarget{idx:0749}{全局数学镜像对},\enspace\IndexPageList[0749]{281}
\subitem \leavevmode\hypertarget{idx:0750}{高维数学镜像对},\enspace\IndexPageList[0750]{287,288,289,290}
\subsubitem \leavevmode\hypertarget{idx:0751}{定义},\enspace\IndexPageList[0751]{287}
\subsubitem \leavevmode\hypertarget{idx:0752}{四维簇的数学镜像对},\enspace\IndexPageList[0752]{288,289}
\subsubitem \leavevmode\hypertarget{idx:0753}{局限},\enspace\IndexPageList[0753]{288}
\subitem \leavevmode\hypertarget{idx:0754}{Hodge 理论镜像对},\enspace\IndexPageList[0754]{281,287}
\subitem \leavevmode\hypertarget{idx:0755}{Calabi--Yau 三维簇的数学镜像对},\enspace\IndexPageList[0755]{280,281,282,283,284,285}
\subsubitem \leavevmode\hypertarget{idx:0756}{定义},\enspace\IndexPageList[0756]{281}
\subsubitem \leavevmode\hypertarget{idx:0757}{用关联函数表述},\enspace\IndexPageList[0757]{282,285}
\subsubitem \leavevmode\hypertarget{idx:0758}{用势函数表述},\enspace\IndexPageList[0758]{283,285}
\subitem \leavevmode\hypertarget{idx:0759}{五次镜像与五次三维簇},\enspace\IndexPageList[0759]{285,286,356}
\indexspace
\item \leavevmode\hypertarget{idx:0077}{双变 Chow 群},\enspace\IndexPageList[0077]{203,205}
\indexspace
\item \leavevmode\hypertarget{idx:0076}{双有理锥猜想}，\textit{参见} \hyperlink{idx:0158}{锥猜想，双有理锥猜想}
\indexspace
\item \leavevmode\hypertarget{idx:0075}{双棱锥构造},\enspace\IndexPageList[0075]{93,95}
\indexspace
\item \leavevmode\hypertarget{idx:0622}{瞬子}
\subitem \leavevmode\hypertarget{idx:0623}{瞬子修正},\enspace\IndexPageList[0623]{32,446}
\subitem \leavevmode\hypertarget{idx:0624}{全纯瞬子},\enspace\IndexPageList[0624]{31,453}
\subitem \leavevmode\hypertarget{idx:0625}{朴素瞬子},\enspace\IndexPageList[0625]{453}
\subitem \leavevmode\hypertarget{idx:0626}{瞬子求和},\enspace\IndexPageList[0626]{445,446}
\subitem \leavevmode\hypertarget{idx:0627}{世界面瞬子},\enspace\IndexPageList[0627]{445}
\indexspace
\item \leavevmode\hypertarget{idx:0628}{瞬子数},\enspace\IndexPageList[0628]{31,226,227,228,229,230,231,232,233,249,331}
\subitem \leavevmode\hypertarget{idx:0629}{定义},\enspace\IndexPageList[0629]{226}
\subitem \leavevmode\hypertarget{idx:0630}{整除性质},\enspace\IndexPageList[0630]{47}
\subitem \leavevmode\hypertarget{idx:0631}{椭圆情形的类似概念},\enspace\IndexPageList[0631]{236,237}
\subitem \leavevmode\hypertarget{idx:0632}{枚举意义},\enspace\IndexPageList[0632]{228,229,230,231}
\subitem \leavevmode\hypertarget{idx:0633}{例题},\enspace\IndexPageList[0633]{228,229,230,231}，\textit{另见} \hyperlink{idx:1191}{五次三维簇，瞬子数}
\subitem \leavevmode\hypertarget{idx:0634}{五次三维簇的瞬子数}，\textit{参见} \hyperlink{idx:1191}{五次三维簇，瞬子数}
\subitem \leavevmode\hypertarget{idx:0635}{辛定义},\enspace\IndexPageList[0635]{231,232,233,423}
\indexspace
\item \leavevmode\hypertarget{idx:0989}{算子乘积展开},\enspace\IndexPageList[0989]{447,448}
\indexspace
\item \leavevmode\hypertarget{idx:1345}{拓扑 $\sigma$ 模型},\enspace\IndexPageList[1345]{190}
\subitem \leavevmode\hypertarget{idx:1346}{与引力耦合},\enspace\IndexPageList[1346]{190,210}
\indexspace
\item \leavevmode\hypertarget{idx:1344}{拓扑递推关系},\enspace\IndexPageList[1344]{334,335}
\indexspace
\item \leavevmode\hypertarget{idx:1336}{拓扑量子场论},\enspace\IndexPageList[1336]{190,191,240,439,445,453,454,455,456,457,458}
\subitem \leavevmode\hypertarget{idx:1337}{$(0+1)$ 维拓扑量子场论},\enspace\IndexPageList[1337]{454,455}
\subitem \leavevmode\hypertarget{idx:1338}{$(1+1)$ 维拓扑量子场论},\enspace\IndexPageList[1338]{455}
\subsubitem \leavevmode\hypertarget{idx:1339}{给出 Frobenius 代数},\enspace\IndexPageList[1339]{455,457}
\subitem \leavevmode\hypertarget{idx:1340}{$(2+1)$ 维拓扑量子场论},\enspace\IndexPageList[1340]{458}
\subitem \leavevmode\hypertarget{idx:1341}{$(d+1)$ 维拓扑量子场论},\enspace\IndexPageList[1341]{454}
\subitem \leavevmode\hypertarget{idx:1342}{定义},\enspace\IndexPageList[1342]{454}
\subitem \leavevmode\hypertarget{idx:1343}{与 A 模型的关系},\enspace\IndexPageList[1343]{457}
\indexspace
\item \leavevmode\hypertarget{idx:0176}{凸}
\subitem \leavevmode\hypertarget{idx:0177}{凸丛},\enspace\IndexPageList[0177]{372}，\textit{另见} \hyperlink{idx:0146}{凸凹丛}
\subitem \leavevmode\hypertarget{idx:0178}{凸线丛},\enspace\IndexPageList[0178]{387}
\subitem \leavevmode\hypertarget{idx:0179}{凸簇},\enspace\IndexPageList[0179]{321,381,388,389,404,405}
\subsubitem \leavevmode\hypertarget{idx:0180}{定义},\enspace\IndexPageList[0180]{203}
\indexspace
\item \leavevmode\hypertarget{idx:1248}{$\Sigma$ 凸}，\textit{参见} \hyperlink{idx:0202}{$\operatorname{cpl}(\Sigma)$，$\Sigma$ 凸性}
\indexspace
\item \leavevmode\hypertarget{idx:0146}{凸凹丛},\enspace\IndexPageList[0146]{360,372,373,375}
\subitem \leavevmode\hypertarget{idx:0147}{临界丛},\enspace\IndexPageList[0147]{373,375}
\subsubitem \leavevmode\hypertarget{idx:0148}{定义},\enspace\IndexPageList[0148]{374}
\indexspace
\item \leavevmode\hypertarget{idx:0181}{凸函数},\enspace\IndexPageList[0181]{62}，\textit{另见} \hyperlink{idx:0199}{$\operatorname{cpl}(\Sigma)$}
\subitem \leavevmode\hypertarget{idx:0182}{严格凸函数},\enspace\IndexPageList[0182]{62,65,66,68}
\indexspace
\item \leavevmode\hypertarget{idx:0327}{图 $\Gamma$ 的旗},\enspace\IndexPageList[0327]{308}，\textit{另见} \hyperlink{idx:0726}{$\overline M_{0,n}(\mathbb P^n,d)$ 中的局部化，不动点轨迹}
\indexspace
\item \leavevmode\hypertarget{idx:0131}{完全交},\enspace\IndexPageList[0131]{387}
\subitem \leavevmode\hypertarget{idx:0132}{$\mathbb P^n$ 中的完全交},\enspace\IndexPageList[0132]{238,387}
\subitem \leavevmode\hypertarget{idx:0133}{凸簇中的完全交},\enspace\IndexPageList[0133]{389}
\subsubitem \leavevmode\hypertarget{idx:0134}{nef 完全交},\enspace\IndexPageList[0134]{388}
\subitem \leavevmode\hypertarget{idx:0135}{旗流形中的完全交}，\textit{参见} \hyperlink{idx:0324}{旗流形}
\subitem \leavevmode\hypertarget{idx:0136}{Grassmann 簇中的完全交}，\textit{参见} \hyperlink{idx:0454}{Grassmann 簇}
\subitem \leavevmode\hypertarget{idx:0137}{环面完全交}，\textit{参见} \hyperlink{idx:1347}{环面完全交}
\subitem \leavevmode\hypertarget{idx:0138}{虚拟基本类}，\textit{参见} \hyperlink{idx:1419}{虚拟基本类，完全交的虚拟基本类}
\indexspace
\item \leavevmode\hypertarget{idx:1040}{微扰态},\enspace\IndexPageList[1040]{425}
\indexspace
\item \leavevmode\hypertarget{idx:1100}{伪流形},\enspace\IndexPageList[1100]{212}
\indexspace
\item \leavevmode\hypertarget{idx:1269}{稳定 $C^\ell$ 映射},\enspace\IndexPageList[1269]{211,214}
\subitem \leavevmode\hypertarget{idx:1270}{等价性},\enspace\IndexPageList[1270]{211}
\indexspace
\item \leavevmode\hypertarget{idx:1271}{稳定映射}
\subitem \leavevmode\hypertarget{idx:1272}{自同构},\enspace\IndexPageList[1272]{304,306,307,313,314,315,363,364,391,393,400}
\subitem \leavevmode\hypertarget{idx:1273}{定义},\enspace\IndexPageList[1273]{192}
\subitem \leavevmode\hypertarget{idx:1274}{稳定映射族},\enspace\IndexPageList[1274]{192}
\subitem \leavevmode\hypertarget{idx:1275}{图构造},\enspace\IndexPageList[1275]{342,357}
\subitem \leavevmode\hypertarget{idx:1276}{模空间},\enspace\IndexPageList[1276]{39,50,196}，\textit{另见} \hyperlink{idx:0725}{$\overline M_{0,n}(\mathbb P^n,d)$ 中的局部化}；\hyperlink{idx:0777}{$M_d$}
\subsubitem \leavevmode\hypertarget{idx:1277}{粗模空间},\enspace\IndexPageList[1277]{192,193,201,204,324}
\subsubitem \leavevmode\hypertarget{idx:1278}{维数},\enspace\IndexPageList[1278]{193}
\subsubitem \leavevmode\hypertarget{idx:1279}{预期维数},\enspace\IndexPageList[1279]{194,198,203,210,215,234,329,331}
\subsubitem \leavevmode\hypertarget{idx:1280}{模叠},\enspace\IndexPageList[1280]{193,198,203,204,324,325,380,381,382,387}
\subsubitem \leavevmode\hypertarget{idx:1281}{切空间},\enspace\IndexPageList[1281]{198,308,313}
\subitem \leavevmode\hypertarget{idx:1282}{万有稳定映射},\enspace\IndexPageList[1282]{202,324,325,328}
\indexspace
\item \leavevmode\hypertarget{idx:1164}{五次镜像},\enspace\IndexPageList[1164]{40,145,356}
\subitem \leavevmode\hypertarget{idx:1165}{$\mathcal A$ 方程组},\enspace\IndexPageList[1165]{116,117}
\subsubitem \leavevmode\hypertarget{idx:1166}{解},\enspace\IndexPageList[1166]{124}
\subitem \leavevmode\hypertarget{idx:1167}{B 模型关联函数},\enspace\IndexPageList[1167]{40,44,45,46,50,52}，\textit{另见} \hyperlink{idx:0189}{关联函数，B 模型关联函数}
\subitem \leavevmode\hypertarget{idx:1168}{B 模型势函数},\enspace\IndexPageList[1168]{286,356}
\subitem \leavevmode\hypertarget{idx:1169}{复模},\enspace\IndexPageList[1169]{41,285}
\subitem \leavevmode\hypertarget{idx:1170}{构造},\enspace\IndexPageList[1170]{40,84}
\subitem \leavevmode\hypertarget{idx:1171}{Frobenius 方法},\enspace\IndexPageList[1171]{185}
\subitem \leavevmode\hypertarget{idx:1172}{Gauss--Manin 联络},\enspace\IndexPageList[1172]{48}
\subitem \leavevmode\hypertarget{idx:1173}{齐次坐标环},\enspace\IndexPageList[1173]{85}
\subitem \leavevmode\hypertarget{idx:1174}{单值性},\enspace\IndexPageList[1174]{134,135,285}
\subitem \leavevmode\hypertarget{idx:1175}{Picard--Fuchs 方程},\enspace\IndexPageList[1175]{43,44,49,104,110,111,118,186,356}
\subsubitem \leavevmode\hypertarget{idx:1176}{解},\enspace\IndexPageList[1176]{46,124,133,180,183,185,187,356,376,385}
\subitem \leavevmode\hypertarget{idx:1177}{自反多胞形},\enspace\IndexPageList[1177]{84}
\subitem \leavevmode\hypertarget{idx:1178}{满足整性猜想},\enspace\IndexPageList[1178]{105,133,173,180}
\subitem \leavevmode\hypertarget{idx:1179}{Yukawa 耦合},\enspace\IndexPageList[1179]{44,180}
\subsubitem \leavevmode\hypertarget{idx:1180}{归一化 Yukawa 耦合},\enspace\IndexPageList[1180]{44,45,46,49,50,133,134,285,386}，\textit{另见} \hyperlink{idx:1167}{五次镜像，B 模型关联函数}
\indexspace
\item \leavevmode\hypertarget{idx:1181}{五次三维簇},\enspace\IndexPageList[1181]{38,39,40,41,42,43,44,45,46,47,48,49,50,51,52,375,452}
\subitem \leavevmode\hypertarget{idx:1182}{其上的 $2875$ 条直线},\enspace\IndexPageList[1182]{46,47,180,196,304,305}
\subitem \leavevmode\hypertarget{idx:1183}{A 模型关联函数},\enspace\IndexPageList[1183]{38,39,45,46,50,51,226,249,285,356,386}，\textit{另见} \hyperlink{idx:0184}{关联函数，A 模型关联函数}
\subitem \leavevmode\hypertarget{idx:1184}{大量子上同调},\enspace\IndexPageList[1184]{261}
\subitem \leavevmode\hypertarget{idx:1185}{Fermat 五次三维簇},\enspace\IndexPageList[1185]{166,197}
\subitem \leavevmode\hypertarget{idx:1186}{Gromov--Witten 不变量},\enspace\IndexPageList[1186]{196,197,224,225,226,316,355}
\subsubitem \leavevmode\hypertarget{idx:1187}{等变不变量},\enspace\IndexPageList[1187]{322}
\subsubitem \leavevmode\hypertarget{idx:1188}{公式},\enspace\IndexPageList[1188]{207,224,304,355,373}
\subsubitem \leavevmode\hypertarget{idx:1189}{通过局部化计算},\enspace\IndexPageList[1189]{311}
\subitem \leavevmode\hypertarget{idx:1190}{Gromov--Witten 势},\enspace\IndexPageList[1190]{51,260,261,356,373,385,411}
\subitem \leavevmode\hypertarget{idx:1191}{瞬子数},\enspace\IndexPageList[1191]{39,46,47,48,228,230,285,316,356,386}，\textit{另见} \hyperlink{idx:0628}{瞬子数}
\subsubitem \leavevmode\hypertarget{idx:1192}{$n_4$ 的计算},\enspace\IndexPageList[1192]{311}
\subsubitem \leavevmode\hypertarget{idx:1193}{$n_{10}$ 的细微之处},\enspace\IndexPageList[1193]{39,229,320}
\subitem \leavevmode\hypertarget{idx:1194}{Kähler 模},\enspace\IndexPageList[1194]{38,41,45,180}
\subitem \leavevmode\hypertarget{idx:1195}{镜像映射},\enspace\IndexPageList[1195]{41,42,43,45,46,49,51,133,179,180,185,286,356,376,385}
\subitem \leavevmode\hypertarget{idx:1196}{镜像}，\textit{参见} \hyperlink{idx:1164}{五次镜像}
\subitem \leavevmode\hypertarget{idx:1197}{镜像对称},\enspace\IndexPageList[1197]{38,41,43,45,46,285,286,311,356}，\textit{另见} \hyperlink{idx:0867}{镜像定理，五次三维簇的镜像定理}
\subitem \leavevmode\hypertarget{idx:1198}{镜像定理}，\textit{参见} \hyperlink{idx:0867}{镜像定理，五次三维簇的镜像定理}
\subitem \leavevmode\hypertarget{idx:1199}{其上的有理曲线},\enspace\IndexPageList[1199]{24,38,39,47,225,229,311,316}，\textit{另见} \hyperlink{idx:0124}{Clemens 猜想}
\subsubitem \leavevmode\hypertarget{idx:1200}{其上的六结点五次曲线},\enspace\IndexPageList[1200]{39,229,316,318,320}
\subsubitem \leavevmode\hypertarget{idx:1201}{法丛},\enspace\IndexPageList[1201]{317}
\subitem \leavevmode\hypertarget{idx:1202}{小量子上同调},\enspace\IndexPageList[1202]{248,249}
\subitem \leavevmode\hypertarget{idx:1203}{虚拟基本类},\enspace\IndexPageList[1203]{196,203,204,207}
\indexspace
\item \leavevmode\hypertarget{idx:1048}{物理学家的 Euler 数},\enspace\IndexPageList[1048]{83,425}，\textit{另见} \hyperlink{idx:0592}{Hodge 数，弦论 Hodge 数}
\indexspace
\item \leavevmode\hypertarget{idx:1265}{稀疏结式},\enspace\IndexPageList[1265]{145}
\indexspace
\item \leavevmode\hypertarget{idx:1287}{弦理论},\enspace\IndexPageList[1287]{24,25,30,442}，\textit{另见} \hyperlink{idx:1304}{超弦理论}
\subitem \leavevmode\hypertarget{idx:1288}{极限},\enspace\IndexPageList[1288]{429}
\subitem \leavevmode\hypertarget{idx:1289}{周期},\enspace\IndexPageList[1289]{429}
\indexspace
\item \leavevmode\hypertarget{idx:0721}{线性 $\sigma$ 模型}，\textit{参见} \hyperlink{idx:0354}{规范线性 $\sigma$ 模型}
\indexspace
\item \leavevmode\hypertarget{idx:0700}{线性关系格},\enspace\IndexPageList[0700]{62,64,67,159,161,163,178,182,247,413}
\subitem \leavevmode\hypertarget{idx:0701}{例题},\enspace\IndexPageList[0701]{73,116,118,147,149,183,187}
\subitem \leavevmode\hypertarget{idx:0702}{与 $\mathcal A$ 方程组的关系},\enspace\IndexPageList[0702]{114,120,122}
\indexspace
\item \leavevmode\hypertarget{idx:0717}{线性回路},\enspace\IndexPageList[0717]{159,161}
\subitem \leavevmode\hypertarget{idx:0718}{例题},\enspace\IndexPageList[0718]{159,163,164}
\subitem \leavevmode\hypertarget{idx:0719}{收缩 GKZ 分解},\enspace\IndexPageList[0719]{162}
\subitem \leavevmode\hypertarget{idx:0720}{由扇支撑},\enspace\IndexPageList[0720]{160,162}
\indexspace
\item \leavevmode\hypertarget{idx:1042}{相},\enspace\IndexPageList[1042]{69,166}，\textit{另见} \hyperlink{idx:0354}{规范线性 $\sigma$ 模型}
\indexspace
\item \leavevmode\hypertarget{idx:1041}{相变},\enspace\IndexPageList[1041]{452}
\indexspace
\item \leavevmode\hypertarget{idx:1397}{向量超场},\enspace\IndexPageList[1397]{451}
\indexspace
\item \leavevmode\hypertarget{idx:1387}{消失圈},\enspace\IndexPageList[1387]{42}
\indexspace
\item \leavevmode\hypertarget{idx:1264}{小子群},\enspace\IndexPageList[1264]{431}
\indexspace
\item \leavevmode\hypertarget{idx:1318}{辛}
\subitem \leavevmode\hypertarget{idx:1319}{辛作用},\enspace\IndexPageList[1319]{63}
\subitem \leavevmode\hypertarget{idx:1320}{辛流形},\enspace\IndexPageList[1320]{63,207,208,232,251}
\subsubitem \leavevmode\hypertarget{idx:1321}{半正辛流形}，\textit{参见} \hyperlink{idx:1245}{半正}
\subitem \leavevmode\hypertarget{idx:1322}{辛轨形},\enspace\IndexPageList[1322]{65,433}
\subitem \leavevmode\hypertarget{idx:1323}{辛约化},\enspace\IndexPageList[1323]{64,65,67,68,166,433}
\indexspace
\item \leavevmode\hypertarget{idx:1324}{辛商}，\textit{参见} \hyperlink{idx:1323}{辛，辛约化}
\indexspace
\item \leavevmode\hypertarget{idx:0208}{形变}
\subitem \leavevmode\hypertarget{idx:0209}{无穷小形变},\enspace\IndexPageList[0209]{80,176,198}
\subitem \leavevmode\hypertarget{idx:0210}{极小 Calabi--Yau 轨形的形变},\enspace\IndexPageList[0210]{433}
\subitem \leavevmode\hypertarget{idx:0211}{轨形的形变},\enspace\IndexPageList[0211]{433}
\subitem \leavevmode\hypertarget{idx:0212}{非多项式形变},\enspace\IndexPageList[0212]{81,170}
\subitem \leavevmode\hypertarget{idx:0213}{多项式形变},\enspace\IndexPageList[0213]{80,84}
\indexspace
\item \leavevmode\hypertarget{idx:1427}{虚法丛},\enspace\IndexPageList[1427]{305,323}
\indexspace
\item \leavevmode\hypertarget{idx:1428}{虚法锥}，\textit{参见} \hyperlink{idx:0951}{法锥，虚法锥}
\indexspace
\item \leavevmode\hypertarget{idx:1406}{虚拟基本类},\enspace\IndexPageList[1406]{39,195,197,198,199,200,201,202,203,204,205,206,214,230,311,317,322,325,328,329,331,336,381,383,404}
\subitem \leavevmode\hypertarget{idx:1407}{与遗忘映射的相容性},\enspace\IndexPageList[1407]{205,328}
\subitem \leavevmode\hypertarget{idx:1408}{用 Euler 类计算},\enspace\IndexPageList[1408]{205,231,234,312,331,332}
\subitem \leavevmode\hypertarget{idx:1409}{定义}，\textit{另见} \hyperlink{idx:0959}{障碍理论}；\hyperlink{idx:1328}{切—障碍复形}
\subsubitem \leavevmode\hypertarget{idx:1410}{Behrend--Fantechi 方法},\enspace\IndexPageList[1410]{201,202}
\subsubitem \leavevmode\hypertarget{idx:1411}{Behrend--Fantechi 定义},\enspace\IndexPageList[1411]{202}
\subsubitem \leavevmode\hypertarget{idx:1412}{定义的等价性},\enspace\IndexPageList[1412]{199,218,422}
\subsubitem \leavevmode\hypertarget{idx:1413}{Li--Tian 方法},\enspace\IndexPageList[1413]{199,200,201,204}
\subsubitem \leavevmode\hypertarget{idx:1414}{Li--Tian 定义},\enspace\IndexPageList[1414]{200,201}
\subitem \leavevmode\hypertarget{idx:1415}{次数等于预期维数},\enspace\IndexPageList[1415]{201,215,245,327}
\subitem \leavevmode\hypertarget{idx:1416}{等变虚拟基本类},\enspace\IndexPageList[1416]{322,323,405}
\subsubitem \leavevmode\hypertarget{idx:1417}{局部化公式},\enspace\IndexPageList[1417]{308,323,405}
\subitem \leavevmode\hypertarget{idx:1418}{例题},\enspace\IndexPageList[1418]{200,202}
\subitem \leavevmode\hypertarget{idx:1419}{完全交的虚拟基本类},\enspace\IndexPageList[1419]{409,423}
\subsubitem \leavevmode\hypertarget{idx:1420}{$\mathbb P^n$ 中完全交的虚拟基本类},\enspace\IndexPageList[1420]{383}
\subitem \leavevmode\hypertarget{idx:1421}{Calabi--Yau 三维簇的虚拟基本类},\enspace\IndexPageList[1421]{311,312}
\subitem \leavevmode\hypertarget{idx:1422}{Fermat 五次三维簇的虚拟基本类},\enspace\IndexPageList[1422]{197}
\subitem \leavevmode\hypertarget{idx:1423}{五次三维簇的虚拟基本类},\enspace\IndexPageList[1423]{196,203,204,383}
\subitem \leavevmode\hypertarget{idx:1424}{性质},\enspace\IndexPageList[1424]{203,204,205,206,218,327}
\subitem \leavevmode\hypertarget{idx:1425}{相对虚拟基本类},\enspace\IndexPageList[1425]{203}
\indexspace
\item \leavevmode\hypertarget{idx:1426}{虚邻域},\enspace\IndexPageList[1426]{211}
\indexspace
\item \leavevmode\hypertarget{idx:1429}{虚切空间},\enspace\IndexPageList[1429]{323}
\indexspace
\item \leavevmode\hypertarget{idx:1268}{旋量},\enspace\IndexPageList[1268]{449}
\indexspace
\item \leavevmode\hypertarget{idx:1266}{旋量丛},\enspace\IndexPageList[1266]{442}
\indexspace
\item \leavevmode\hypertarget{idx:1267}{旋量几何},\enspace\IndexPageList[1267]{449}
\indexspace
\item \leavevmode\hypertarget{idx:0049}{湮灭算子},\enspace\IndexPageList[0049]{438,441}
\indexspace
\item \leavevmode\hypertarget{idx:0479}{引力 Gromov--Witten 势}，\textit{参见} \hyperlink{idx:0376}{亏格 $g$ 的引力 Gromov--Witten 势}
\indexspace
\item \leavevmode\hypertarget{idx:0459}{引力关联数},\enspace\IndexPageList[0459]{48,219,234,324,325,326,327,328,329,330,331,332,333,334,343,405,406}，\textit{另见} \hyperlink{idx:0376}{亏格 $g$ 的引力 Gromov--Witten 势}；\hyperlink{idx:0375}{亏格 $g$ 耦合}
\subitem \leavevmode\hypertarget{idx:0460}{公理}，\textit{另见} \hyperlink{idx:0457}{引力类}
\subsubitem \leavevmode\hypertarget{idx:0461}{次数公理},\enspace\IndexPageList[0461]{327,333,334,340,341,348,383}
\subsubitem \leavevmode\hypertarget{idx:0462}{伸缩公理},\enspace\IndexPageList[0462]{329,330,331,333,341}
\subsubitem \leavevmode\hypertarget{idx:0463}{除子公理},\enspace\IndexPageList[0463]{328,329,330,331,333,337,341,380}
\subsubitem \leavevmode\hypertarget{idx:0464}{基本类公理},\enspace\IndexPageList[0464]{328,329,330,331,341,347,348}
\subsubitem \leavevmode\hypertarget{idx:0465}{分裂公理},\enspace\IndexPageList[0465]{328,329,330,336}
\subitem \leavevmode\hypertarget{idx:0466}{系数},\enspace\IndexPageList[0466]{326}
\subitem \leavevmode\hypertarget{idx:0467}{定义},\enspace\IndexPageList[0467]{325}
\subitem \leavevmode\hypertarget{idx:0468}{用 Gromov--Witten 不变量表示亏格 $0$ 引力关联数},\enspace\IndexPageList[0468]{329}
\subitem \leavevmode\hypertarget{idx:0469}{$\mathbb P^1$ 的亏格 $0$ 引力关联数},\enspace\IndexPageList[0469]{329,380}
\subsubitem \leavevmode\hypertarget{idx:0470}{递推关系},\enspace\IndexPageList[0470]{330,331}
\subitem \leavevmode\hypertarget{idx:0471}{Calabi--Yau 三维簇的亏格 $0$ 引力关联数},\enspace\IndexPageList[0471]{331}
\subitem \leavevmode\hypertarget{idx:0472}{亏格 $1$ 引力关联数 $\langle\tau_1\rangle_{1,0}$},\enspace\IndexPageList[0472]{331,332}
\subsubitem \leavevmode\hypertarget{idx:0473}{与 Euler 示性数的关系},\enspace\IndexPageList[0473]{333}
\subitem \leavevmode\hypertarget{idx:0474}{由拓扑递推关系得到的递推式},\enspace\IndexPageList[0474]{334}，\textit{另见} \hyperlink{idx:1344}{拓扑递推关系}
\subitem \leavevmode\hypertarget{idx:0475}{与 Gromov--Witten 不变量的关系},\enspace\IndexPageList[0475]{326}
\subitem \leavevmode\hypertarget{idx:0476}{符号记法},\enspace\IndexPageList[0476]{338,379}
\indexspace
\item \leavevmode\hypertarget{idx:0477}{引力后裔},\enspace\IndexPageList[0477]{325,334}
\subitem \leavevmode\hypertarget{idx:0478}{与引力关联数的关系},\enspace\IndexPageList[0478]{326}
\indexspace
\item \leavevmode\hypertarget{idx:0457}{引力类},\enspace\IndexPageList[0457]{329,332}，\textit{另见} \hyperlink{idx:0459}{引力关联数}
\subitem \leavevmode\hypertarget{idx:0458}{定义},\enspace\IndexPageList[0458]{328}
\indexspace
\item \leavevmode\hypertarget{idx:0481}{引力量子乘积},\enspace\IndexPageList[0481]{257,327}
\subitem \leavevmode\hypertarget{idx:0482}{系数},\enspace\IndexPageList[0482]{327}
\indexspace
\item \leavevmode\hypertarget{idx:0480}{引力势},\enspace\IndexPageList[0480]{333}
\indexspace
\item \leavevmode\hypertarget{idx:1205}{有理多面锥},\enspace\IndexPageList[1205]{54}
\subitem \leavevmode\hypertarget{idx:1206}{强凸有理多面锥},\enspace\IndexPageList[1206]{54}
\indexspace
\item \leavevmode\hypertarget{idx:1234}{右行}
\subitem \leavevmode\hypertarget{idx:1235}{右行费米子},\enspace\IndexPageList[1235]{25,444}
\subitem \leavevmode\hypertarget{idx:1236}{右行扇区},\enspace\IndexPageList[1236]{448}
\indexspace
\item \leavevmode\hypertarget{idx:0197}{余切复形},\enspace\IndexPageList[0197]{201}
\indexspace
\item \leavevmode\hypertarget{idx:0198}{余切线 $\mathcal L_i$},\enspace\IndexPageList[0198]{325,392}
\indexspace
\item \leavevmode\hypertarget{idx:0281}{预期维数}
\subitem \leavevmode\hypertarget{idx:0282}{Kuranishi 理论中的预期维数},\enspace\IndexPageList[0282]{197}
\subitem \leavevmode\hypertarget{idx:0283}{稳定映射模空间的预期维数}，\textit{参见} \hyperlink{idx:1276}{稳定映射，模空间}
\subitem \leavevmode\hypertarget{idx:0284}{零点轨迹的预期维数},\enspace\IndexPageList[0284]{196}
\indexspace
\item \leavevmode\hypertarget{idx:1084}{预稳定映射},\enspace\IndexPageList[1084]{202}
\indexspace
\item \leavevmode\hypertarget{idx:0451}{粘合引理},\enspace\IndexPageList[0451]{361,365,403}
\indexspace
\item \leavevmode\hypertarget{idx:0217}{占优猜想},\enspace\IndexPageList[0217]{140,178}
\indexspace
\item \leavevmode\hypertarget{idx:0957}{障碍丛},\enspace\IndexPageList[0957]{231,312,313,315}，\textit{另见} \hyperlink{idx:1435}{表示的权，障碍权}
\indexspace
\item \leavevmode\hypertarget{idx:0958}{障碍空间},\enspace\IndexPageList[0958]{198}
\indexspace
\item \leavevmode\hypertarget{idx:0959}{障碍理论},\enspace\IndexPageList[0959]{201,323}
\subitem \leavevmode\hypertarget{idx:0960}{完美障碍理论},\enspace\IndexPageList[0960]{197,203,322}
\subsubitem \leavevmode\hypertarget{idx:0961}{定义},\enspace\IndexPageList[0961]{202}
\subsubitem \leavevmode\hypertarget{idx:0962}{等变完美障碍理论},\enspace\IndexPageList[0962]{322}
\subsubitem \leavevmode\hypertarget{idx:0963}{例题},\enspace\IndexPageList[0963]{202}
\subsubitem \leavevmode\hypertarget{idx:0964}{用于虚拟基本类},\enspace\IndexPageList[0964]{203}
\indexspace
\item \leavevmode\hypertarget{idx:1385}{真空态},\enspace\IndexPageList[1385]{438,440,448,450}
\indexspace
\item \leavevmode\hypertarget{idx:0636}{整性猜想},\enspace\IndexPageList[0636]{128,129,130,132,133,135,174,178,180,273,281,282,285,289,291,293,408,409,422}
\subitem \leavevmode\hypertarget{idx:0637}{表述},\enspace\IndexPageList[0637]{105,173}
\indexspace
\item \leavevmode\hypertarget{idx:0952}{正规化正合列},\enspace\IndexPageList[0952]{313,315}
\indexspace
\item \leavevmode\hypertarget{idx:1225}{正则奇点},\enspace\IndexPageList[1225]{43,98,100,269,270}
\indexspace
\item \leavevmode\hypertarget{idx:1226}{正则三角剖分},\enspace\IndexPageList[1226]{66,67}
\indexspace
\item \leavevmode\hypertarget{idx:1227}{正则性引理}，\textit{参见} \hyperlink{idx:0861}{镜像定理，Givental 镜像定理，正则性引理}
\indexspace
\item \leavevmode\hypertarget{idx:1314}{支撑函数},\enspace\IndexPageList[1314]{55}
\subitem \leavevmode\hypertarget{idx:1315}{上凸支撑函数},\enspace\IndexPageList[1315]{56}
\subsubitem \leavevmode\hypertarget{idx:1316}{严格上凸支撑函数},\enspace\IndexPageList[1316]{56,58}
\indexspace
\item \leavevmode\hypertarget{idx:0621}{指标方程},\enspace\IndexPageList[0621]{100,102,103,186}
\indexspace
\item \leavevmode\hypertarget{idx:0119}{中心荷},\enspace\IndexPageList[0119]{333,448}
\indexspace
\item \leavevmode\hypertarget{idx:1208}{重构定理},\enspace\IndexPageList[1208]{219,220,223}
\indexspace
\item \leavevmode\hypertarget{idx:1229}{重整化},\enspace\IndexPageList[1229]{25,30,31,446}，\textit{另见} \hyperlink{idx:0944}{非重整化定理}
\indexspace
\item \leavevmode\hypertarget{idx:1038}{周期},\enspace\IndexPageList[1038]{42,43,50,98,115,117,118,182,356}，\textit{另见} \hyperlink{idx:1059}{Picard--Fuchs 方程，解}
\indexspace
\item \leavevmode\hypertarget{idx:1039}{周期域},\enspace\IndexPageList[1039]{277}
\indexspace
\item \leavevmode\hypertarget{idx:1093}{主 $\mathcal A$ 行列式},\enspace\IndexPageList[1093]{145}
\indexspace
\item \leavevmode\hypertarget{idx:1286}{驻相法},\enspace\IndexPageList[1286]{440,445}
\indexspace
\item \leavevmode\hypertarget{idx:0149}{锥}
\subitem \leavevmode\hypertarget{idx:0150}{Calabi--Yau 锥}，\textit{参见} \hyperlink{idx:0110}{Calabi--Yau 锥}
\subitem \leavevmode\hypertarget{idx:0151}{对偶锥}，\textit{参见} \hyperlink{idx:0224}{对偶锥}
\subitem \leavevmode\hypertarget{idx:0152}{锥的面}，\textit{参见} \hyperlink{idx:0296}{锥的面}
\subitem \leavevmode\hypertarget{idx:0153}{Gorenstein 锥},\enspace\IndexPageList[0153]{87}
\subsubitem \leavevmode\hypertarget{idx:0154}{自反 Gorenstein 锥}，\textit{参见} \hyperlink{idx:1213}{自反 Gorenstein 锥}
\subitem \leavevmode\hypertarget{idx:0155}{齐次自伴锥},\enspace\IndexPageList[0155]{154}
\subitem \leavevmode\hypertarget{idx:0156}{有理多面锥}，\textit{参见} \hyperlink{idx:1205}{有理多面锥}
\indexspace
\item \leavevmode\hypertarget{idx:0157}{锥猜想},\enspace\IndexPageList[0157]{153,157,173,175,268,422}
\subitem \leavevmode\hypertarget{idx:0158}{双有理锥猜想},\enspace\IndexPageList[0158]{156}
\subitem \leavevmode\hypertarget{idx:0159}{推论},\enspace\IndexPageList[0159]{153}
\subitem \leavevmode\hypertarget{idx:0160}{已获证明的情形},\enspace\IndexPageList[0160]{153}
\indexspace
\item \leavevmode\hypertarget{idx:0296}{锥的面},\enspace\IndexPageList[0296]{54}
\indexspace
\item \leavevmode\hypertarget{idx:0168}{锥形奇点},\enspace\IndexPageList[0168]{145,166}
\subitem \leavevmode\hypertarget{idx:0169}{锥形奇点跃迁},\enspace\IndexPageList[0169]{166,427,428,429}，\textit{另见} \hyperlink{idx:0292}{极端跃迁，I 型}
\subsubitem \leavevmode\hypertarget{idx:0170}{例题},\enspace\IndexPageList[0170]{167}
\indexspace
\item \leavevmode\hypertarget{idx:1290}{子轨形},\enspace\IndexPageList[1290]{78,432}
\indexspace
\item \leavevmode\hypertarget{idx:1213}{自反 Gorenstein 锥},\enspace\IndexPageList[1213]{85,87}
\subitem \leavevmode\hypertarget{idx:1214}{对偶性},\enspace\IndexPageList[1214]{87}
\subitem \leavevmode\hypertarget{idx:1215}{指数},\enspace\IndexPageList[1215]{87}
\subitem \leavevmode\hypertarget{idx:1216}{与 nef 分拆的关系},\enspace\IndexPageList[1216]{87}
\subitem \leavevmode\hypertarget{idx:1217}{与自反多胞形的关系},\enspace\IndexPageList[1217]{87}
\indexspace
\item \leavevmode\hypertarget{idx:1224}{自反层},\enspace\IndexPageList[1224]{60}
\indexspace
\item \leavevmode\hypertarget{idx:1218}{自反多胞形},\enspace\IndexPageList[1218]{29,58,70,76,77,78,86,119,139,141,157,158,167,168,178,290,405,406,428}
\subitem \leavevmode\hypertarget{idx:1219}{Batyrev 对偶},\enspace\IndexPageList[1219]{69,79}
\subitem \leavevmode\hypertarget{idx:1220}{分类},\enspace\IndexPageList[1220]{70}
\subitem \leavevmode\hypertarget{idx:1221}{定义},\enspace\IndexPageList[1221]{69}
\subitem \leavevmode\hypertarget{idx:1222}{例题},\enspace\IndexPageList[1222]{72,74,94,111,147,148,162,169}
\subitem \leavevmode\hypertarget{idx:1223}{五次镜像的自反多胞形},\enspace\IndexPageList[1223]{84}
\indexspace
\item \leavevmode\hypertarget{idx:0130}{组合对偶性},\enspace\IndexPageList[0130]{65,69}
\indexspace
\item \leavevmode\hypertarget{idx:0762}{最大退化},\enspace\IndexPageList[0762]{144}，\textit{另见} \hyperlink{idx:0766}{最大幂幺，最大幂幺边界点}
\indexspace
\item \leavevmode\hypertarget{idx:0765}{最大幂幺}
\subitem \leavevmode\hypertarget{idx:0766}{最大幂幺边界点},\enspace\IndexPageList[0766]{105,116,125,127,128,132,133,136,137,138,142,144,151,154,161,172,174,175,176,177,178,179,182,183,184,274,280,289,290,291,292}
\subsubitem \leavevmode\hypertarget{idx:0767}{定义},\enspace\IndexPageList[0767]{104}
\subsubitem \leavevmode\hypertarget{idx:0768}{等价性},\enspace\IndexPageList[0768]{138,143}
\subsubitem \leavevmode\hypertarget{idx:0769}{例题},\enspace\IndexPageList[0769]{112,113,148,150,165}
\subsubitem \leavevmode\hypertarget{idx:0770}{存在性},\enspace\IndexPageList[0770]{139,143,422}
\subsubitem \leavevmode\hypertarget{idx:0771}{该点处的切空间},\enspace\IndexPageList[0771]{181}
\subitem \leavevmode\hypertarget{idx:0772}{最大幂幺单值性},\enspace\IndexPageList[0772]{41,42,101,102,103,104,105,132,136,150,176,186,287}
\subsubitem \leavevmode\hypertarget{idx:0773}{一维模空间},\enspace\IndexPageList[0773]{101,102,103,104}
\subsubitem \leavevmode\hypertarget{idx:0774}{$r$ 维模空间},\enspace\IndexPageList[0774]{104,105}
\indexspace
\item \leavevmode\hypertarget{idx:0711}{左行}
\subitem \leavevmode\hypertarget{idx:0712}{左行费米子},\enspace\IndexPageList[0712]{25,444}
\subitem \leavevmode\hypertarget{idx:0713}{左行扇区},\enspace\IndexPageList[0713]{448}
\indexspace
\item \leavevmode\hypertarget{idx:0033}{作用量},\enspace\IndexPageList[0033]{435,439}
\subitem \leavevmode\hypertarget{idx:0034}{$\sigma$ 模型的作用量},\enspace\IndexPageList[0034]{443,444}
\subsubitem \leavevmode\hypertarget{idx:0035}{玻色部分},\enspace\IndexPageList[0035]{443}
\subsubitem \leavevmode\hypertarget{idx:0036}{费米部分},\enspace\IndexPageList[0036]{443}
\end{theindex}
\endgroup

\end{document}




